\documentclass[11pt]{article}

\usepackage[margin=1in]{geometry}
\usepackage{amsmath,amssymb,amsthm,mathtools,bm}
\usepackage{enumitem}
\usepackage{booktabs,tabularx}
\usepackage{xcolor}
\usepackage{microtype}
\usepackage[colorlinks=true,linkcolor=blue,citecolor=blue,urlcolor=blue]{hyperref}
\usepackage[nameinlink,capitalise,noabbrev]{cleveref}
\allowdisplaybreaks[3]
\setlength{\emergencystretch}{4em}
\raggedbottom
\makeatletter
\patchcmd{\l@section}{1.5em}{2.3em}{}{} 
\makeatother

\newtheorem{theorem}{Theorem}[section]
\newtheorem{proposition}[theorem]{Proposition}
\newtheorem{lemma}[theorem]{Lemma}
\newtheorem{corollary}[theorem]{Corollary}
\newtheorem{criterion}[theorem]{Criterion}
\theoremstyle{definition}
\newtheorem{definition}[theorem]{Definition}
\newtheorem{problem}[theorem]{Open problem}
\theoremstyle{remark}
\newtheorem{remark}[theorem]{Remark}
\newtheorem{assessment}[theorem]{Assessment}
\newtheorem{firewall}[theorem]{Firewall}

\newcommand{\R}{\mathbb R}
\newcommand{\C}{\mathbb C}
\newcommand{\N}{\mathbb N}
\newcommand{\cH}{\mathcal H}
\newcommand{\cZ}{\mathcal Z}
\newcommand{\cV}{\mathcal V}
\newcommand{\dd}{\,\mathrm d}
\newcommand{\norm}[1]{\left\lVert #1\right\rVert}
\newcommand{\abs}[1]{\left\lvert #1\right\rvert}
\newcommand{\ip}[2]{\left\langle #1,#2\right\rangle}
\newcommand{\supp}{\operatorname{supp}}
\newcommand{\diag}{\operatorname{diag}}
\newcommand{\Ran}{\operatorname{Ran}}
\newcommand{\Ker}{\operatorname{Ker}}
\newcommand{\spec}{\operatorname{spec}}

\title{Renewal GRH Cofinal Visibility Attack Notes\\
\large Fixed-bank obstruction, a normalized modulation repair, and the honest remaining ledger\\
\normalsize Version V1}
\author{}
\date{10 September 2026}

\begin{document}
\maketitle

\begin{abstract}
These notes begin a new, independently named GRH attack series.  They audit the
GRH programme in the Grand Tensor monograph V074, the attached arithmetic
renewal manuscript, and the current notes in \texttt{latex\_notes}.  The first
target is the finite-height Mellin--Prony tomography route.  We prove that no
fixed finite bank of smooth compactly supported profiles, translations, or
fixed-order scale jets can have the cutoff-uniform Mellin visibility demanded
by a cofinal argument.  We then give a predeclared, zero-independent modulation
bank whose source-normalized profile norm stays constant and whose visibility
loss is only the explicit factor \(J_T^{-1/2}\asymp T^{-1/2}\).  The construction
is exactly compatible with the translation recurrence used by finite zero
tomography.  A quantitative Hankel perturbation lemma isolates what is still
missing: the explicit-formula multiplier, Vandermonde conditioning, row
subtraction at the resulting scale, physical admission of the growing bank,
the independent restoring lower bridge, and cofinal descent.  No assertion of
GRH is made.
\end{abstract}

\tableofcontents

\section{Context, scope, and nonduplication boundary}
\label{sec:context}

The user requested a renewed attack on GRH using the added Grand Tensor
structure, with rigorous proofs and without repeating earlier notes.  This V1
is the first file in the series
\path{Renewal_GRH_Cofinal_Visibility_Attack_V1.tex}; successors use the same
stem with V2, V3, and so on.
At each later iteration the preceding file is to be copied exactly, new
material is to be inserted before its terminal document-closing marker, and the
preceding version is to be removed after validation.  At every tenth version
the current companion notes are to be swept again.  After V100 the series is
to roll to a suffix \texttt{f2}, with numbering restarted at V1.

The mathematical sources inspected before beginning this note were:
\begin{enumerate}[label=\textup{(S\arabic*)},leftmargin=3em]
\item the attached
\texttt{Renewal\_Grand\_Tensor\_Monograph\_V074(1).tex}, especially its
canonical arithmetic loading, completed--Euler short, prime-scale restoring
short, boundary-labelled bank, finite zero tomography, integrated hard-edge
router, and explicit open-problem ledger;
\item the attached arithmetic renewal manuscript, whose conditional GRH
routes include the Mertens action, rectangular determinant, theta--Jensen,
and Weil--Schur formulations;
\item every filename in \texttt{latex\_notes}, with the current endpoints of
the Navier--Stokes, Standard Model--Einstein, main Yang--Mills, and parallel
Yang--Mills series inspected for transferable results and collision of scope;
\item the attached strategic memorandum recommending that the programme now
concentrate on the quantitative completed--Euler upper bridge, the independent
restoring lower bridge, and the cofinal passage.
\end{enumerate}

No existing note in the folder develops the Mellin-visibility obstruction or
the modulation repair below.  The parallel programmes do, however, impose two
useful disciplines that are retained here: a finite operator is interpreted
only in a complete normalized source frame, and a finite positive certificate
is not promoted to a cofinal theorem without a uniform tail or compactness
argument.

V074 has already removed a substantial ambiguity.  Arithmetic count,
multiplication, the incidence zeta history, M\"obius and von Mangoldt sources,
characters, and the relevant heat--Mellin coordinates are canonically derived
on the protected chronology.  Its finite-height tomography theorem gives, for
zero nodes \(\rho\) with \(\abs{\Im\rho}\le T\),
\[
 Y_k=\sum_{\rho\in\cZ_T}
 e^{k\tau(\rho-1/2)}v_\rho,
 \qquad 0\le k<2M,
 \qquad 0<\tau<\frac{\pi}{T},
 \tag{1.1}
\]
where \(M=\#\cZ_T\) and
\[
 v_\rho=m(\rho)\Xi_G(\rho),
 \qquad
 \Xi_G(s)=\int_\R e^{(s-1/2)t}G(t)\dd t.
 \tag{1.2}
\]
At a fixed height this is an exact Prony system.  Its cofinal use requires,
among other things, a lower bound for
\(\min_{\rho\in\cZ_T}\norm{v_\rho}\).  The present note isolates the part of
that demand contributed by \(\Xi_G\).

\begin{firewall}[Logical status]
\label{fire:status}
Theorems below concern the analytic visibility and conditioning of the
finite-height measurement system.  They neither prove that all reconstructed
nodes lie on the unit circle nor supply the completed--Euler upper bridge,
the prime-scale restoring inequality, or the source-anchored cofinal passage.
Every unproved implication remains listed in \Cref{sec:ledger}.
\end{firewall}

\section{A fixed smooth bank cannot be cofinally visible}
\label{sec:no-go}

For a finite-dimensional Hilbert space \(\cH\) and
\(G\in C_c^\infty(\R;\cH)\), define its critical-line Mellin trace by
\[
 \mathcal M_G(\gamma)
 :=\Xi_G\left(\frac12+i\gamma\right)
 =\int_\R e^{i\gamma t}G(t)\dd t.
 \tag{2.1}
\]
Thus the visibility problem is a vector-valued Fourier problem.

\begin{lemma}[Quantitative integration by parts]
\label{lem:ibp}
Let \(G\in W^{r,1}(\R;\cH)\) have compact support, where \(r\ge1\).  Then for
every \(\gamma\ne0\),
\[
 \boxed{
 \norm{\mathcal M_G(\gamma)}_\cH
 \le \abs{\gamma}^{-r}\norm{G^{(r)}}_{L^1(\R;\cH)}.}
 \tag{2.2}
\]
\end{lemma}

\begin{proof}
For each \(u\in\cH\), scalar integration by parts gives
\[
 \ip{\mathcal M_G(\gamma)}{u}
 =(-i\gamma)^{-r}
 \int_\R e^{i\gamma t}\ip{G^{(r)}(t)}{u}\dd t.
\]
Compact support removes all boundary terms.  Taking absolute values, then the
supremum over unit vectors \(u\), proves \textup{(2.2)}.
\end{proof}

\begin{theorem}[Fixed finite-bank visibility obstruction]
\label{thm:fixed-bank-no-go}
Let \(G_1,\ldots,G_q\) be a fixed finite collection in
\(C_c^\infty(\R;\cH_j)\), and form the direct-sum bank
\[
 \mathcal B(\gamma)
 =\bigl(\mathcal M_{G_1}(\gamma),\ldots,
 \mathcal M_{G_q}(\gamma)\bigr).
\]
Then, for every \(r\ge1\),
\[
 \boxed{
 \norm{\mathcal B(\gamma)}
 \le C_r(1+\abs{\gamma})^{-r}}
 \tag{2.3}
\]
for a constant \(C_r\) depending only on the fixed bank.  In particular,
\[
 \inf_{\abs{\gamma}\le T}\norm{\mathcal B(\gamma)}
 \le \norm{\mathcal B(T)}\longrightarrow0,
 \tag{2.4}
\]
so the bank has no positive cutoff-uniform Mellin-visibility floor.
\end{theorem}

\begin{proof}
Apply \Cref{lem:ibp} to each component and sum the squared estimates.  On
\(\abs{\gamma}\ge1\) one may take
\[
 C_r^2=\sum_{j=1}^q\norm{G_j^{(r)}}_{L^1}^2.
\]
On the compact interval \(\abs{\gamma}\le1\), enlarge \(C_r\) using
\(\norm{\mathcal M_{G_j}(\gamma)}\le\norm{G_j}_{L^1}\).  This proves
\textup{(2.3)}; substituting \(\gamma=T\) proves \textup{(2.4)}.
\end{proof}

The point is not merely the Riemann--Lebesgue lemma.  Smooth compact support
forces decay faster than every fixed inverse power, so a finite list of
profiles cannot meet a cofinal lower-frame requirement.

\begin{corollary}[Fixed translations and fixed jets do not repair visibility]
\label{cor:fixed-operations}
Begin with finitely many profiles in \(C_c^\infty\), and apply finitely many
of the following operations with cutoff-independent parameters:
\begin{enumerate}[label=\textup{(\roman*)},leftmargin=2.5em]
\item translation \(G(t)\mapsto G(t-a)\);
\item multiplication by a polynomial in \(t\);
\item a finite difference in the translation parameter;
\item a derivative of fixed order;
\item a finite-dimensional linear combination or direct sum.
\end{enumerate}
The resulting finite bank still satisfies \textup{(2.3)} for every \(r\), with
new constants.  Hence finitely many centre translations, taper moments, or
fixed-order scale jets cannot by themselves give cofinal uniform visibility.
\end{corollary}

\begin{proof}
Every listed operation sends \(C_c^\infty\) to a finite collection of
\(C_c^\infty\) profiles.  Alternatively, on the Fourier side a translation
multiplies by \(e^{i\gamma a}\), multiplication by \(t\) differentiates in
\(\gamma\), and a fixed difference contributes a bounded trigonometric
multiplier.  Theorem~\ref{thm:fixed-bank-no-go} applies to the resulting fixed
collection.
\end{proof}

\begin{theorem}[Frequency-budget law]
\label{thm:frequency-budget}
Let \(G_T\in W^{r,1}(\R;\cH_T)\) be a cutoff-dependent family with compact
support, where the finite-dimensional target space may depend on \(T\).  If
for some \(a_T>0\)
\[
 \norm{\mathcal M_{G_T}(T)}\ge a_T,
\]
then necessarily
\[
 \boxed{\norm{G_T^{(r)}}_{L^1(\R;\cH_T)}\ge a_TT^r.}
 \tag{2.5}
\]
Consequently, a family uniformly bounded in one fixed \(W^{r,1}\) norm has
visibility at ordinate \(T\) at most \(O(T^{-r})\).
\end{theorem}

\begin{proof}
This is \Cref{lem:ibp} with \(\gamma=T\), rearranged.
\end{proof}

\begin{assessment}[Upstream consequence]
\label{ass:no-go}
The fixed-profile interpretation of V074's cofinal Mellin-visibility request
is impossible.  Any viable physical centre--taper bank must grow in number,
frequency, derivative size, support scale, or some combination of these.
That growth is a source cost and must be retained in the completed Gram; it
cannot be hidden by declaring each finite-height profile separately.
\end{assessment}

\section{A predeclared normalized modulation bank}
\label{sec:repair}

The obstruction has a simple analytic repair.  It does not use the locations
of any zeros.

Choose a nonzero scalar \(\phi\in C_c^\infty(\R)\) with \(\phi\ge0\), and set
\[
 F_\phi(\beta,\xi)
 =\int_\R e^{(\beta-1/2)t}e^{i\xi t}\phi(t)\dd t,
 \qquad \widehat\phi(\xi)=F_\phi(1/2,\xi).
\]
The function \(F_\phi\) is continuous and
\(F_\phi(\beta,0)>0\) for \(0\le\beta\le1\).  Compactness and uniform
continuity therefore give numbers \(r_0>0\) and \(c_0>0\) such that
\[
 \abs{F_\phi(\beta,\xi)}\ge c_0
 \qquad(0\le\beta\le1,\ \abs{\xi}\le r_0/2).
 \tag{3.1}
\]
For \(T\ge1\), put
\[
 L_T=\left\lceil\frac{T}{r_0}\right\rceil+1,
 \qquad
 \Lambda_T=\{jr_0:-L_T\le j\le L_T\},
 \qquad J_T=2L_T+1,
 \tag{3.2}
\]
and let \(\cH_T=\ell^2(\Lambda_T)\).  Define
\[
 \boxed{
 G_T(t)=\frac{\phi(t)}{\sqrt{J_T}}
 \bigl(e^{-i\lambda t}\bigr)_{\lambda\in\Lambda_T}.}
 \tag{3.3}
\]

\begin{theorem}[Zero-independent source-normalized visibility]
\label{thm:modulation-bank}
The bank \textup{(3.3)} is predeclared from \(T\) and \(\phi\), independently of
the zeros of any \(L\)-function.  It satisfies
\[
 \boxed{
 \norm{G_T(t)}_{\cH_T}=\abs{\phi(t)},
 \qquad
 \norm{G_T}_{L^p(\R;\cH_T)}=\norm{\phi}_{L^p(\R)}
 }
 \tag{3.4}
\]
for \(1\le p\le\infty\), and for every \(\gamma\in[-T,T]\),
\[
 \boxed{
 \norm{\mathcal M_{G_T}(\gamma)}_{\cH_T}
 \ge\frac{c_0}{\sqrt{J_T}}.}
 \tag{3.5}
\]
Since \(J_T\le 2T/r_0+5\), the loss in \textup{(3.5)} is at worst an explicit
constant times \(T^{-1/2}\).
\end{theorem}

\begin{proof}
At each \(t\), all \(J_T\) components in \textup{(3.3)} have modulus
\(\abs{\phi(t)}/\sqrt{J_T}\), proving \textup{(3.4)}.  For any
\(\gamma\in[-T,T]\), choose a nearest lattice point
\(\lambda(\gamma)\in\Lambda_T\).  The extra endpoint in the definition of
\(L_T\) ensures
\(\abs{\gamma-\lambda(\gamma)}\le r_0/2\).  The corresponding coordinate of
the Mellin trace is
\[
 \bigl[\mathcal M_{G_T}(\gamma)\bigr]_{\lambda(\gamma)}
 =\frac{1}{\sqrt{J_T}}
 \widehat\phi\bigl(\gamma-\lambda(\gamma)\bigr).
\]
Equation \textup{(3.1)} now gives \textup{(3.5)}.
\end{proof}

\begin{remark}[Normalization tradeoff]
\label{rem:normalization}
Removing the factor \(J_T^{-1/2}\) gives a constant visibility floor but
makes the pointwise source norm and every homogeneous Gram scale like
\(J_T^{1/2}\) and \(J_T\), respectively.  Source normalization therefore
converts the growing number of frequency channels into the unavoidable
visibility factor \(J_T^{-1/2}\).  This is an accounting identity, not a loss
caused by the proof.
\end{remark}

\begin{proposition}[The modulation bank pays the required frequency budget]
\label{prop:bank-derivatives}
For every integer \(r\ge0\), there is a constant \(C_{r,\phi,r_0}\) such that
\[
 \norm{G_T^{(r)}}_{L^1(\R;\cH_T)}
 \le C_{r,\phi,r_0}(1+T)^r.
 \tag{3.6}
\]
Thus the construction has precisely the polynomial high-frequency growth
permitted, and qualitatively required, by \Cref{thm:frequency-budget}.
\end{proposition}

\begin{proof}
Leibniz' rule writes the \(\lambda\)-coordinate of \(G_T^{(r)}\) as
\[
 \frac{e^{-i\lambda t}}{\sqrt{J_T}}
 \sum_{a=0}^r\binom ra(-i\lambda)^a\phi^{(r-a)}(t).
\]
Taking the \(\ell^2(\Lambda_T)\) norm, using
\(\max_{\lambda\in\Lambda_T}\abs{\lambda}\le T+2r_0\), and integrating in
\(t\) gives \textup{(3.6)}.
\end{proof}

\section{Compatibility with translation tomography}
\label{sec:translation}

The repair would be irrelevant if it destroyed the exponential recurrence.
It does not.

\begin{proposition}[Exact translation covariance]
\label{prop:translation}
For \(k\in\N\) and \(\tau>0\), define
\(G_{T,k}(t)=G_T(t-k\tau)\).  Then for every \(s\in\C\),
\[
 \boxed{
 \Xi_{G_{T,k}}(s)
 =e^{k\tau(s-1/2)}\Xi_{G_T}(s).}
 \tag{4.1}
\]
In particular, after the standard prime, pole, gamma/conductor, contour, and
edge rows have been subtracted, the visible zero output retains exactly the
Prony form \textup{(1.1)} with a vector-valued weight in \(\cH_T\).
\end{proposition}

\begin{proof}
With \(u=t-k\tau\),
\[
 \int_\R e^{(s-1/2)t}G_T(t-k\tau)\dd t
 =e^{k\tau(s-1/2)}
 \int_\R e^{(s-1/2)u}G_T(u)\dd u.
\]
This is \textup{(4.1)}.
\end{proof}

For completeness, the conditioning statement used in the tomography route is
proved next in the exact norm needed for perturbations.

\begin{theorem}[Vector Prony injectivity and singular-value floor]
\label{thm:prony-floor}
Let \(z_1,\ldots,z_M\) be distinct nonzero complex numbers and let
\(v_1,\ldots,v_M\) be nonzero vectors in a finite-dimensional Hilbert space
\(\cH\).  Put
\[
 Y_k=\sum_{j=1}^M z_j^kv_j,
 \qquad 0\le k<2M,
\]
and define the block-Hankel map
\[
 \mathsf H_Y:\C^M\longrightarrow\cH^M,
 \qquad
 (\mathsf H_Ya)_k=\sum_{\ell=0}^{M-1}a_\ell Y_{k+\ell},
 \quad 0\le k<M.
 \tag{4.2}
\]
If \(V=(z_j^k)_{0\le k<M,1\le j\le M}\), then
\[
 \boxed{
 \sigma_{\min}(\mathsf H_Y)
 \ge
 \left(\min_{1\le j\le M}\norm{v_j}\right)
 \sigma_{\min}(V)^2.}
 \tag{4.3}
\]
Consequently \(\mathsf H_Y\) is injective, and the monic annihilating
polynomial has degree exactly \(M\) and roots \(z_1,\ldots,z_M\).
\end{theorem}

\begin{proof}
For \(a=(a_0,\ldots,a_{M-1})\), let
\(p_a(z)=\sum_{\ell=0}^{M-1}a_\ell z^\ell\).  Then
\[
 (\mathsf H_Ya)_k
 =\sum_{j=1}^Mz_j^kp_a(z_j)v_j.
\]
Thus \(\mathsf H_Y\) factors as
\[
 \C^M\xrightarrow{V^{\mathsf T}}\C^M
 \xrightarrow{D_v}\cH^M
 \xrightarrow{V\otimes I_\cH}\cH^M,
\]
where \(D_v(c_1,\ldots,c_M)=(c_1v_1,\ldots,c_Mv_M)\).  The three least
singular values are respectively
\(\sigma_{\min}(V)\), at least \(\min_j\norm{v_j}\), and
\(\sigma_{\min}(V)\).  Their product proves \textup{(4.3)}.  Injectivity rules
out an annihilator of degree below \(M\), while
\(\prod_{j=1}^M(z-z_j)\) plainly annihilates the sequence.
\end{proof}

\begin{corollary}[What the modulation bank actually closes]
\label{cor:combined-floor}
Let \(\cZ_T=\{\rho_1,\ldots,\rho_M\}\) be a finite set with
\(\abs{\Im\rho_j}\le T\), let
\(z_j=e^{\tau(\rho_j-1/2)}\), and suppose \(0<\tau<\pi/T\) so that the
retained nodes are not aliased.  With
\(v_j=m(\rho_j)\Xi_{G_T}(\rho_j)\), define
\[
 \mu_T=\min_{1\le j\le M}\abs{m(\rho_j)}.
\]
Then
\[
 \boxed{
 \min_j\norm{v_j}\ge\frac{c_0\mu_T}{\sqrt{J_T}},
 \qquad
 \sigma_{\min}(\mathsf H_Y)
 \ge\frac{c_0\mu_T}{\sqrt{J_T}}\sigma_{\min}(V)^2.}
 \tag{4.4}
\]
The bare Mellin-visibility factor has therefore been made explicit and
polynomial.  No lower bound for \(\mu_T\) or \(\sigma_{\min}(V)\) is asserted.
\end{corollary}

\begin{proof}
For \(\rho_j=\beta_j+i\gamma_j\), the critical-line trace in
\Cref{thm:modulation-bank} is modified by the bounded positive weight
\(e^{(\beta_j-1/2)t}\) on the fixed support of \(\phi\).  On the GRH-testing
line \(\beta_j=1/2\), \textup{(3.5)} applies directly.  More generally, on any
fixed critical strip one may choose \(\phi\ge0\) and shrink \(r_0\), uniformly
over the compact horizontal parameter, to retain a positive constant; relabel
that constant as \(c_0\).  Multiplication by \(m(\rho_j)\) gives the first
bound in \textup{(4.4)}, and \Cref{thm:prony-floor} gives the second.
\end{proof}

\begin{remark}[Why the strip-uniform sentence is legitimate]
\label{rem:strip}
If \(\phi\ge0\) is supported in a sufficiently short interval about zero,
then
\[
 F(\beta,\xi)=\int e^{(\beta-1/2)t}e^{i\xi t}\phi(t)\dd t
\]
is continuous and nonzero at every \((\beta,0)\) in the compact strip
\(0\le\beta\le1\).  Uniform continuity supplies one common neighborhood
\(\abs{\xi}\le r_0/2\) and one positive lower bound.  This choice may be made
before any zero is known.
\end{remark}

\begin{proposition}[Stable recovery of recurrence coefficients]
\label{prop:perturbation}
Let the exact block-Hankel map \(\mathsf H\) in \textup{(4.2)} have least
singular value \(\sigma>0\), and let \(q\in\C^M\) solve the monic recurrence
\[
 \mathsf Hq=-b,
 \qquad b=(Y_M,\ldots,Y_{2M-1})\in\cH^M.
\]
Suppose measured data give \(\widetilde{\mathsf H}=\mathsf H+E\) and
\(\widetilde b=b+e\), with \(\norm{E}<\sigma\), and let
\(\widetilde q\) be the least-squares solution.  Then
\[
 \boxed{
 \norm{\widetilde q-q}
 \le
 \frac{\norm e+\norm E\,\norm q}{\sigma-\norm E}.}
 \tag{4.5}
\]
\end{proposition}

\begin{proof}
Weyl's singular-value inequality gives
\(\sigma_{\min}(\widetilde{\mathsf H})\ge\sigma-\norm E>0\).  Subtracting
\(\widetilde{\mathsf H}\widetilde q=-\widetilde b\) from
\(\widetilde{\mathsf H}q=-b+Eq\) yields
\[
 \widetilde{\mathsf H}(\widetilde q-q)=-(e+Eq).
\]
Apply the norm of the left inverse, bounded by
\((\sigma-\norm E)^{-1}\), to obtain \textup{(4.5)}.
\end{proof}

\begin{assessment}[Tomography after V1]
\label{ass:tomography}
The finite system is algebraically exact, and a zero-independent profile bank
now prevents superpolynomial loss from the Mellin transform itself.  Stable
cofinal recovery still requires the total row-subtraction error to be little-o
of the right-hand side of \textup{(4.4)}.  This exposes, rather than hides, the
condition-number problem.  A numerical reconstruction at each fixed height
would not establish that uniform estimate.
\end{assessment}

\section{Consequences for the direction of the GRH attack}
\label{sec:direction}

The strategic memorandum proposed finite zero tomography as one of several
bridges.  V1 changes its role.  Tomography is useful as an exact finite
diagnostic, but a fixed physical profile cannot carry it cofinally.  A growing
modulation bank can repair visibility, at the explicit cost of increasing
frequency complexity.  Therefore any attempt to use tomography as the main
GRH proof must establish all of the following on one canonical cutoff family:
\begin{enumerate}[label=\textup{(T\arabic*)},leftmargin=3em]
\item admission of the predeclared modulation channels as actual arithmetic
centre--taper or scale-jet rows, rather than a post hoc fit to zero locations;
\item a source-normalized bound for the completed explicit-formula multiplier
\(\mu_T\);
\item a lower bound for the Vandermonde/eigenframe factor
\(\sigma_{\min}(V)^2\), or a different stable recurrence representation;
\item subtraction of prime, pole, gamma/conductor, completed-contour, and
cutoff-edge rows with error smaller than the full floor in \textup{(4.4)};
\item a positive Stein or equivalent unit-circle certificate which does not
assume the critical-line conclusion;
\item the independent prime-scale restoring inequality and the final
source-anchored cofinal descent.
\end{enumerate}

Items (T2)--(T4) may erase the gain from \Cref{thm:modulation-bank}.  They are
not bookkeeping details.  In particular, if the multiplier \(m(\rho)\) is a
fixed Gaussian smoother, its high-ordinate decay remains in \(\mu_T\) and
must be compensated or replaced by a cutoff-matched family.  The modulation
bank corrects the profile transform, not an independently inserted smoothing
factor.

This suggests a division of labor.  The completed--Euler relative-metric
route and the prime-scale restoring short should remain the main proof
battlefield, because they formulate GRH as operator inequalities before the
ill-conditioned task of locating individual zeros.  The modulation--Prony
route should be retained as a diagnostic and as a possible compact-height
closure only if its source cost can be routed into those same canonical
metrics.

\begin{problem}[Source-native admission of the modulation bank]
\label{prob:admission}
Construct the vector profile \textup{(3.3)} from a predeclared family of actual
chronology-derived arithmetic experiments.  Prove that its complete mixed
score Gram, after the common nuisance short, is the normalized direct-sum Gram
used above, and that the derivative cost \textup{(3.6)} is paid by an explicit
source row rather than discarded under scalar normalization.
\end{problem}

\begin{problem}[Conditioning without individual zero localization]
\label{prob:conditioning}
Replace the monomial Vandermonde basis, if possible, by a predeclared
orthogonal-polynomial, unitary-pencil, or localized-frame realization whose
least singular value is controlled from prime-side data.  Any basis change
must include its own condition number; algebraic similarity alone cannot
improve the physical recovery floor.
\end{problem}

\begin{problem}[Router to the restoring low sector]
\label{prob:router}
On the same source metric, construct a target-independent map from the
frequency-resolved or modulated completed--Euler current to the translated,
taper-boundary, or stable-return rows in the low-\(A_+\) restoring short.  Prove
the map before applying positivity.  A factorization chosen after observing a
negative restoring witness is inadmissible.
\end{problem}

\section{Rigorous status ledger after V1}
\label{sec:ledger}

\begin{center}
\begin{tabularx}{\textwidth}{@{}p{0.24\textwidth}p{0.18\textwidth}X@{}}
\toprule
Component & Status & Reason \\
\midrule
Canonical arithmetic source & inherited finite theorem & V074 derives the
normalized arithmetic loading; no new source is postulated here. \\
Fixed finite-profile cofinal visibility & disproved &
\Cref{thm:fixed-bank-no-go,cor:fixed-operations}. \\
Growing predeclared bare Mellin visibility & proved &
\Cref{thm:modulation-bank} gives a zero-independent normalized bank with
floor \(c_0J_T^{-1/2}\). \\
Translation/Prony compatibility & proved &
\Cref{prop:translation,thm:prony-floor}. \\
Perturbative recurrence recovery & proved conditionally &
\Cref{prop:perturbation} requires error below the complete singular-value
floor. \\
Explicit-formula multiplier floor & open & The modulation bank does not
remove decay in \(m(\rho)\). \\
Vandermonde / eigenframe conditioning & open & It is isolated in
\(\sigma_{\min}(V)^2\). \\
Physical score-row admission & open & \Cref{prob:admission}. \\
Completed--Euler upper bridge & open & No proof that
\(\eta_X^+\to0\) is supplied. \\
Prime-scale restoring lower bridge & open & No proof of the low-\(A_+\)
Schur inequality is supplied. \\
Cofinal GRH passage & open & Finite tomography and finite positivity do not
exclude remote escape. \\
GRH & open & The required upper, lower, and cofinal estimates have not all
been proved. \\
\bottomrule
\end{tabularx}
\end{center}

\section{Next iteration target}
\label{sec:next}

V2 should go upstream of the monomial Vandermonde system.  The first task is
to determine whether the translation nodes can be represented by a unitary
pencil or by localized Fourier frames with a conditioning estimate stated
directly in terms of separation and local zero density.  The estimate must be
proved uniformly and must keep the change-of-basis condition number.  In
parallel, the explicit multiplier \(m(\rho)\) used by V074 must be extracted
from the completed formula so that its height dependence is no longer hidden
inside the symbol \(\mu_T\).  If either cost is necessarily exponential, the
tomography route should be demoted to a diagnostic and the attack should move
upstream to the completed--Euler relative metric and the restoring low-sector
router.

\begin{assessment}[V1 acquisition]
The first iteration yields one negative and one positive result.  Negatively,
fixed smooth finite tomography banks cannot meet the cofinal visibility
requirement.  Positively, a predeclared modulation bank repairs the bare
Mellin visibility with constant total source norm and only an explicit
square-root channel-count loss.  The result narrows the GRH attack: the real
tomographic obstacles are now the completed multiplier, eigenframe
conditioning, physical admission and subtraction accuracy, not the existence
of finite-height moments.
\end{assessment}


% =============================================================================
\section{V2: the Gaussian multiplier and the arc-conditioning obstruction}
\label{sec:grh-v2}
% =============================================================================

\begin{center}
\large\bfseries Version V2 continuation
\end{center}

\subsection{Acquisition target and nonduplication check}

V1 left two analytic symbols unresolved in the finite-height tomography
route: the completed explicit-formula multiplier \(m(\rho)\) and the
Vandermonde factor.  V2 returns to the literal cutoff-matched Gaussian prime
port in V074.  It proves the following.
\begin{enumerate}[label=\textup{(V2.\arabic*)},leftmargin=3.5em]
\item The Gaussian zero multiplier can be evaluated exactly on the completed
full-line zero density.  Its apparent high-ordinate decay is removed by a
predeclared grid of Mellin centres, at the same explicit channel-count cost as
the V1 modulation bank.
\item The monomial translation Vandermonde is exponentially ill-conditioned
whenever the translation step has a fixed relative safety margin below the
Nyquist value \(\pi/T\), even in the best-case situation in which every zero
is already on the critical line.
\item A change of polynomial basis cannot remove that exponential loss without
paying an exponentially large synthesis norm.  Thus it does not by itself
produce a source-native conditioning estimate.
\end{enumerate}

The active companion files added after the V1 audit were also searched for
Gaussian-multiplier, Mellin-visibility, Prony, and Vandermonde results.  Their
new Nyquist discussion concerns a Yang--Mills cardinal grid, not the arithmetic
translation recurrence, and none contains the results below.  The compulsory
full companion sweep remains scheduled for V10.

\subsection{Exact completed-bulk Gaussian zero coefficient}

V074 fixes \(A>0\),
\[
 G_A(v)=\frac{1}{\sqrt{4\pi A}}e^{-v^2/(4A)},
\]
and uses the travelling port
\[
 \mathcal G_{A,\omega}(r)
 =e^{i\omega r}\int_\R e^{-(1/2+i\omega)u}
 G_A(u-r)\dd\mathfrak e(u).
 \tag{V2.1}
\]
The letter \(\omega\) is used for the Mellin centre in V2, reserving \(h\) for
the later translation step.  The completed logarithmic-derivative explicit
formula has nontrivial-zero density
\[
 \dd\mathfrak e_{\rm zero}(u)
 =-\sum_\rho m_\rho e^{\rho u}\dd u,
 \tag{V2.2}
\]
where \(m_\rho\) is the multiplicity.  The prime half-line, pole, trivial-zero,
gamma/conductor, and endpoint differences are separate rows in V074.  The
calculation below concerns the completed full-line zero row; the difference
between that row and the finite cutoff is part of the required subtraction
error.

\begin{theorem}[Completed-bulk Gaussian multiplier]
\label{thm:v2-gaussian-multiplier}
The contribution of one distinct nontrivial zero
\(\rho=\beta+i\gamma\), with multiplicity \(m_\rho\), to
\textup{(V2.1)}--\textup{(V2.2)} is
\[
 \boxed{
 -m_\rho K_{A,\omega}(\rho)e^{(\rho-1/2)r},
 \qquad
 K_{A,\omega}(\rho)
 =\exp\!\left(A(\rho-\tfrac12-i\omega)^2\right).}
 \tag{V2.3}
\]
Consequently
\[
 \boxed{
 \abs{K_{A,\omega}(\rho)}
 =\exp\!\left(A\left[(\beta-\tfrac12)^2
                    -(\gamma-\omega)^2\right]\right).}
 \tag{V2.4}
\]
On the critical line this is exactly \(e^{-A(\gamma-\omega)^2}\).
\end{theorem}

\begin{proof}
Put \(a=\rho-1/2-i\omega\).  The zero's contribution is
\[
 \begin{aligned}
 &-m_\rho e^{i\omega r}
 \int_\R e^{-(1/2+i\omega)u}G_A(u-r)e^{\rho u}\dd u\\
 &\qquad=-m_\rho e^{(\rho-1/2)r}
 \int_\R G_A(v)e^{av}\dd v,
 \end{aligned}
\]
after \(u=r+v\).  Completing the square in the Gaussian gives
\[
 \int_\R G_A(v)e^{av}\dd v=e^{Aa^2},
\]
which proves \textup{(V2.3)}.  Since
\[
 \Re(a^2)=(\beta-1/2)^2-(\gamma-\omega)^2,
\]
equation \textup{(V2.4)} follows.
\end{proof}

\begin{firewall}[Full-line coefficient versus finite cutoff]
\label{fire:v2-bulk-cutoff}
Theorem~\ref{thm:v2-gaussian-multiplier} does not declare the finite prime
port equal to its completed zero row.  The equality is exact for the completed
full-line zero density.  V074's finite theorem additionally requires uniform
subtraction of the prime, pole, gamma/conductor, completed-contour, and
cutoff-edge rows.  Their error must be compared with the final conditioning
floor, not merely shown to tend to zero.
\end{firewall}

\subsection{A normalized grid of Gaussian Mellin centres}

Fix a centre spacing \(d>0\).  For \(T\ge1\), let
\[
 \Omega_T=\{jd:|j|\le L_T\},
 \qquad L_T=\lceil T/d\rceil+1,
 \qquad J_T=2L_T+1.
 \tag{V2.5}
\]
Take the normalized direct sum of the ports with
\(\omega\in\Omega_T\).  The grid is fixed by \(T,A,d\), independently of
the locations of the zeros.

\begin{theorem}[Uniform Gaussian-multiplier visibility]
\label{thm:v2-gaussian-bank}
For every \(\rho=\beta+i\gamma\) with
\(0\le\beta\le1\) and \(\abs{\gamma}\le T\), the normalized vector of
completed-bulk zero coefficients satisfies
\[
 \boxed{
 \left\|
 \frac{m_\rho}{\sqrt{J_T}}
 \bigl(K_{A,\omega}(\rho)\bigr)_{\omega\in\Omega_T}
 \right\|_{\ell^2(\Omega_T)}
 \ge
 \frac{m_\rho}{\sqrt{J_T}}e^{-Ad^2/4}.}
 \tag{V2.6}
\]
Since \(m_\rho\ge1\) and \(J_T\le2T/d+5\), the Gaussian multiplier itself
causes at most a constant times \(T^{-1/2}\) loss in the normalized bank.
Moreover,
\[
 \boxed{
 \left\|
 \frac{m_\rho}{\sqrt{J_T}}
 \bigl(K_{A,\omega}(\rho)\bigr)_{\omega\in\Omega_T}
 \right\|_{\ell^2}
 \le \frac{m_\rho}{\sqrt{J_T}}e^{A/4}C_{A,d},}
 \tag{V2.7}
\]
where \(C_{A,d}<\infty\) is independent of \(T\) and \(\rho\).
\end{theorem}

\begin{proof}
Choose a nearest grid centre \(\omega(\gamma)\).  Then
\(\abs{\gamma-\omega(\gamma)}\le d/2\), and
\Cref{thm:v2-gaussian-multiplier} gives
\[
 \abs{K_{A,\omega(\gamma)}(\rho)}
 \ge e^{-Ad^2/4},
\]
because \((\beta-1/2)^2\ge0\).  One coordinate yields
\textup{(V2.6)}.

For the upper bound, group the centres by
\(n\le\abs{\gamma-\omega}/d<n+1\).  At most two centres occur in each shell,
and \((\beta-1/2)^2\le1/4\).  Hence
\[
 \sum_{\omega\in\Omega_T}\abs{K_{A,\omega}(\rho)}^2
 \le e^{A/2}
 \left(1+2\sum_{n\ge0}e^{-2Ad^2n^2}\right)
 =e^{A/2}C_{A,d}^2.
\]
Taking square roots and restoring the normalization factor \(J_T^{-1/2}\) proves \textup{(V2.7)}.
\end{proof}

\begin{assessment}[Multiplier status]
\label{ass:v2-multiplier-status}
For the literal Gaussian logarithmic-derivative port, the symbolic multiplier
floor \(\mu_T\) from V1 is no longer open: a predeclared grid of Mellin centres
gives the explicit lower bound \(e^{-Ad^2/4}J_T^{-1/2}\).  What remains open is
whether the full direct-sum port is physically admitted with its complete
mixed score Gram and whether every nonzero row can be subtracted below the
much smaller recurrence-conditioning scale derived next.
\end{assessment}

\subsection{A proper-arc Vandermonde is exponentially ill-conditioned}

Let
\[
 V_M(z_1,\ldots,z_M)=(z_j^k)_{0\le k<M,\,1\le j\le M}.
\]
The next estimate is deterministic; it uses no zero-spacing conjecture.

\begin{theorem}[Proper-arc Vandermonde upper bound]
\label{thm:v2-arc-vandermonde}
Suppose \(M\ge2\), \(\abs{z_j}=1\), and, after a common rotation,
\[
 \abs{\arg z_j}\le\alpha<\pi
 \qquad(1\le j\le M).
\]
Then
\[
 \boxed{
 \sigma_{\min}(V_M)
 \le
 \sqrt{M(2M-1)}
 \sin(\alpha/2)^{M-1}.}
 \tag{V2.8}
\]
Thus a square monomial Vandermonde supported on any fixed proper arc has an
exponentially small least singular value as \(M\to\infty\).
\end{theorem}

\begin{proof}
Put \(n=M-1\) and consider
\[
 p(z)=(z-1)^n=\sum_{k=0}^na_kz^k.
\]
The binomial identity gives
\[
 \norm a_2^2=\sum_{k=0}^n\binom nk^2=\binom{2n}{n}.
\]
Since the largest of the \(2n+1\) binomial coefficients in
\((1+1)^{2n}\) is \(\binom{2n}{n}\),
\[
 \binom{2n}{n}\ge\frac{4^n}{2n+1},
 \qquad
 \norm a_2\ge\frac{2^n}{\sqrt{2n+1}}.
 \tag{V2.9}
\]
For \(\abs{\arg z_j}\le\alpha\),
\[
 \abs{p(z_j)}=\abs{z_j-1}^n
 \le\bigl(2\sin(\alpha/2)\bigr)^n.
\]
Therefore
\[
 \frac{\norm{V_M^{\mathsf T}a}_2}{\norm a_2}
 \le\sqrt{M(2n+1)}\sin(\alpha/2)^n.
\]
The transpose of a complex matrix has the same singular values as the matrix.
Substituting \(n=M-1\) proves \textup{(V2.8)}.
\end{proof}

\begin{corollary}[Fixed sub-Nyquist margin is exponentially unstable]
\label{cor:v2-subnyquist}
Assume, for this conditioning test, that all retained zeros satisfy
\(\Re\rho=1/2\) and \(\abs{\Im\rho}\le T\).  Let the tomography translation
step be \(h\), and suppose
\[
 0<h\le(1-\delta)\frac{\pi}{T}
 \qquad(0<\delta<1).
\]
For the nodes \(z_\rho=e^{h(\rho-1/2)}\),
\[
 \boxed{
 \sigma_{\min}(V_M)
 \le\sqrt{M(2M-1)}
 \cos(\pi\delta/2)^{M-1}.}
 \tag{V2.10}
\]
In particular the instability occurs even on the desired GRH locus; it is not
evidence for an off-line zero.
\end{corollary}

\begin{proof}
The nodes lie on the unit circle with arguments in
\([-hT,hT]\).  Apply \Cref{thm:v2-arc-vandermonde} with
\(\alpha=hT\le(1-\delta)\pi\) and use
\(\sin((1-\delta)\pi/2)=\cos(\pi\delta/2)\).
\end{proof}

The preceding statement controls the Vandermonde factor appearing in V074.
The actual vector Hankel system is also unstable under a mild upper bound on
its zero weights.

\begin{proposition}[Vector Hankel upper-conditioning obstruction]
\label{prop:v2-hankel-upper}
With the notation of \Cref{thm:prony-floor}, put
\(v_{\max}=\max_j\norm{v_j}\).  Then
\[
 \boxed{
 \sigma_{\min}(\mathsf H_Y)
 \le \norm{V_M}\,v_{\max}\,\sigma_{\min}(V_M)
 \le Mv_{\max}\sigma_{\min}(V_M).}
 \tag{V2.11}
\]
Consequently, under the hypotheses of \Cref{cor:v2-subnyquist}, any
subexponential upper bound on \(v_{\max}\) leaves the block-Hankel system
exponentially ill-conditioned for fixed \(\delta>0\).
\end{proposition}

\begin{proof}
The factorization in \Cref{thm:prony-floor} is
\[
 \mathsf H_Y=(V_M\otimes I)D_vV_M^{\mathsf T}.
\]
Choose a unit right singular vector \(a\) of \(V_M^{\mathsf T}\) associated
with its least singular value.  Then
\[
 \norm{\mathsf H_Ya}
 \le\norm{V_M}\,\norm{D_v}\,
       \norm{V_M^{\mathsf T}a}
 =\norm{V_M}v_{\max}\sigma_{\min}(V_M).
\]
Finally \(\norm{V_M}\le\norm{V_M}_{\rm F}=M\), because every entry has
modulus one.
\end{proof}

\begin{corollary}[Near-Nyquist necessity]
\label{cor:v2-near-nyquist}
Write
\[
 \delta_M=1-\frac{hT}{\pi}.
\]
If \(M\delta_M^2/\log M\to\infty\), then the upper bound in
\textup{(V2.10)} is smaller than every inverse power of \(M\).  Therefore a
monomial-Prony scheme seeking only polynomial conditioning must take
\[
 \boxed{
 \delta_M=O\!\left(\sqrt{\frac{\log M}{M}}\right),}
 \tag{V2.12}
\]
up to the limitations of this obstruction.  The translation step must approach
the aliasing boundary \(\pi/T\) with a vanishing relative margin.
\end{corollary}

\begin{proof}
For \(0\le x\le\pi/2\), \(\cos x\le e^{-x^2/2}\).  Hence
\[
 \cos(\pi\delta_M/2)^{M-1}
 \le\exp\!\left(-\frac{\pi^2}{8}(M-1)\delta_M^2\right).
\]
The polynomial prefactor in \textup{(V2.10)} is absorbed by any fixed power
of \(M\).  The stated growth condition makes the exponential smaller than
\(M^{-C}\) for every fixed \(C\).  Its contrapositive gives the necessary
scale \textup{(V2.12)} for avoiding this particular superpolynomial loss.
\end{proof}

\subsection{A basis change must pay for the missing arc}

One might replace monomials by an orthogonal-polynomial or cardinal basis.
The next elementary statement records the cost which must not be omitted.

\begin{proposition}[Basis-change conditioning firewall]
\label{prop:v2-basis-firewall}
Let \(C_M\) be an invertible coefficient-change matrix and let
\(W_M=C_MV_M\).  Then
\[
 \boxed{
 \sigma_{\min}(W_M)
 \le\norm{C_M}\,\sigma_{\min}(V_M),
 \qquad
 \sigma_{\min}(V_M)
 \le\norm{C_M^{-1}}\,\sigma_{\min}(W_M).}
 \tag{V2.13}
\]
Thus if \(V_M\) has the exponential loss in \textup{(V2.10)} and
\(W_M\) is uniformly or polynomially conditioned, then at least one of the
physical synthesis/analysis norms \(\norm{C_M}\), \(\norm{C_M^{-1}}\) must
carry the compensating loss.  A coordinate rewrite alone is not a uniform
source-frame theorem.
\end{proposition}

\begin{proof}
The general inequality \(\sigma_{\min}(AB)\le\norm A\,
\sigma_{\min}(B)\) gives the first statement.  Apply it to
\(V_M=C_M^{-1}W_M\) for the second.
\end{proof}

\begin{firewall}[What the proper-arc theorem does not say]
\label{fire:v2-arc-scope}
Theorem~\ref{thm:v2-arc-vandermonde} does not prove that every conceivable
finite-zero representation is ill-conditioned.  It proves that the literal
square monomial translation recurrence in V074 is exponentially unstable at
any fixed sub-Nyquist safety margin.  A unitary pencil or redundant localized
frame could evade the monomial geometry, but only after its physical source
maps, redundancy, and analysis norms are proved and charged explicitly.
\end{firewall}

\subsection{Updated GRH direction and exact remaining ledger}

Combining V1 and V2 yields a sharper decision.  The bare Mellin transform and
the Gaussian zero multiplier admit polynomial-cost, zero-independent repairs.
The square monomial recurrence does not admit a robust fixed-margin cofinal
limit.  Pushing \(h\) to the Nyquist boundary trades the proper-arc loss for a
vanishing aliasing margin and still leaves row-subtraction stability to be
proved.  This makes finite zero tomography valuable as a bounded-height
diagnostic, but presently a poor primary route to GRH.

The main attack should therefore move upstream to objects that test positivity
without resolving every zero individually:
\begin{enumerate}[label=\textup{(D\arabic*)},leftmargin=3em]
\item the completed--Euler relative metric
\(H_{E/M,X}\), seeking a source-native bound
\((\norm{H_{E/M,X}}-1)_+\to0\);
\item the low-\(A_+\) restoring Schur inequality, using the translated,
boundary, and stable-return rows before high-sector relaxation;
\item the localized Weil--Schur form, but only as a finite low-mode scanner
whose negative witness is routed into one of those actual restoring rows;
\item the source-anchored cofinal descent after both finite bridges close.
\end{enumerate}

The Gaussian grid remains useful in this upstream programme: it is a
predeclared frequency-localized test bank for the completed--Euler metric and
can expose the first frequency block on which domination fails, without
requiring polynomial reconstruction of all zeros.

\begin{center}
\begin{tabularx}{\textwidth}{@{}p{0.27\textwidth}p{0.18\textwidth}X@{}}
\toprule
V2 component & Status & Exact conclusion \\
\midrule
Gaussian bulk multiplier & proved & Formula \textup{(V2.3)} and modulus
\textup{(V2.4)}. \\
Normalized Gaussian-centre bank & proved & Uniform strip visibility
\(e^{-Ad^2/4}J_T^{-1/2}\), independent of zero locations. \\
Finite-cutoff row error & open & It must be below the final recurrence or
operator-inequality margin. \\
Fixed-margin monomial Prony conditioning & disproved & Exponential loss
\textup{(V2.10)} occurs even on the critical line. \\
Polynomially stable monomial regime & boundary constrained & It requires the
near-Nyquist scale \textup{(V2.12)} at minimum. \\
Alternative unitary/redundant representation & open & Its complete source
frame and basis-change cost must be constructed. \\
Completed--Euler upper bridge & open & No cutoff-uniform metric domination is
proved here. \\
Restoring lower bridge & open & No regulator-uniform low-sector row floor is
proved here. \\
Cofinal passage and GRH & open & Neither finite visibility nor finite
conditioning supplies the two bridges and source-anchored descent. \\
\bottomrule
\end{tabularx}
\end{center}

\subsection{V3 target}

V3 should stop trying to recover individual zeros.  It should place the
Gaussian centre bank directly into the completed--Euler relative metric and
derive an exact block decomposition by frequency cells.  The target is a
finite theorem of the form
\[
 \norm{H_{E/M,X}}
 \le
 \sup_j a_{j,X}+\text{explicit inter-cell leakage},
\]
with every term computed before target-dependent shorting.  If the off-diagonal
Gaussian overlap has a summable Toeplitz envelope, Schur's test may reduce the
upper bridge to uniform diagonal cell estimates plus a controllable bandwidth
tail.  The diagonal estimate itself must come from the canonical completed and
Euler sources; Gaussian almost-orthogonality alone cannot determine its sign.

\begin{assessment}[V2 acquisition]
V2 closes the multiplier question for the literal Gaussian completed-bulk
port and proves a new obstruction to its proposed Prony continuation.  A
predeclared Mellin-centre grid prevents high zeros from disappearing under the
Gaussian, but a fixed sub-Nyquist translation margin compresses all unit-circle
nodes into a proper arc and forces exponential monomial conditioning.  The
mathematically justified response is to use the Gaussian bank as an upstream
frequency localizer for the completed--Euler and restoring inequalities, not
as the primary cofinal zero-reconstruction engine.
\end{assessment}


% =============================================================================
\section{V3: a frequency-localized completed--Euler bridge}
\label{sec:grh-v3}
% =============================================================================

\begin{center}
\large\bfseries Version V3 continuation
\end{center}

\subsection{Why the attack moves to the relative metric}

V2 proved that the Gaussian multiplier can be kept visible by a predeclared
Mellin-centre bank, but that the square monomial Prony recurrence is
exponentially ill-conditioned at every fixed sub-Nyquist margin.  V3 therefore
stops trying to reconstruct individual zeros.  It places a Gaussian frequency
partition directly on the completed--Euler relative operator
\[
 H=M^{-1/2}EM^{-1/2},
 \tag{V3.1}
\]
where \(M\succ0\) is the nuisance-shorted completed metric and \(E\succeq0\)
is the nuisance-shorted Euler metric on their common supported source space.
The desired finite upper bridge is \(H\preceq I\).

The new result is an exact operator-valued localization identity.  It divides
the bridge into a frequency-localized completely positive part and one
double-commutator transport residual.  The latter has a block formula and a
Schur bound in terms of the frequency commutator \([D,H]\).  This is not the
generic head--tail estimate already present in V074: it identifies the precise
off-frequency response that a Gaussian arithmetic bank must control.

\subsection{Quadratic partitions and an exact IMS identity}

Let \(\mathcal K\) be finite dimensional.  A \emph{quadratic partition of the
identity} is a finite or norm-summable family of self-adjoint operators
\((F_j)_j\) such that
\[
 \sum_jF_j^2=I.
 \tag{V3.2}
\]

\begin{theorem}[Exact frequency-localization identity]
\label{thm:v3-ims}
For every self-adjoint \(H\) and every quadratic partition
\textup{(V3.2)},
\[
 \boxed{
 H=\mathcal A_F(H)+\mathcal R_F(H),
 \qquad
 \mathcal A_F(H)=\sum_jF_jHF_j,}
 \tag{V3.3}
\]
where
\[
 \boxed{
 \mathcal R_F(H)
 =\frac12\sum_j[F_j,[F_j,H]].}
 \tag{V3.4}
\]
If \(H\succeq0\), then \(\mathcal A_F(H)\succeq0\).  If
\[
 \norm{\mathcal A_F(H)}\le a,
 \qquad
 \norm{\mathcal R_F(H)}\le\varepsilon,
 \tag{V3.5}
\]
then
\[
 \boxed{
 H\preceq(a+\varepsilon)I,
 \qquad
 (\norm H-1)_+\le(a+\varepsilon-1)_+.}
 \tag{V3.6}
\]
In particular, \(a\le1-\delta\) and \(\varepsilon\le\delta\) prove the exact
completed--Euler upper bridge.
\end{theorem}

\begin{proof}
Expanding one double commutator gives
\[
 [F_j,[F_j,H]]=F_j^2H-2F_jHF_j+HF_j^2.
\]
Summation and \textup{(V3.2)} yield
\[
 \frac12\sum_j[F_j,[F_j,H]]
 =H-\sum_jF_jHF_j,
\]
which proves \textup{(V3.3)}--\textup{(V3.4)}.  Positivity of each
\(F_jHF_j\) follows by congruence.  Finally
\[
 H\preceq\mathcal A_F(H)+\norm{\mathcal R_F(H)}I
 \preceq(a+\varepsilon)I,
\]
and \textup{(V3.6)} follows.
\end{proof}

\begin{remark}[The residual has a sign only after it is estimated]
The map \(H\mapsto\mathcal A_F(H)\) is completely positive and unital, but
\(\mathcal R_F(H)\) need not be positive.  Replacing it by a positive source
without proving an order comparison would change the bridge.  The safe
quantity is its operator norm or a proved one-sided upper bound.
\end{remark}

\subsection{Exact spectral blocks of the transport residual}

Let \(D=D^*\) be the physically loaded Mellin-frequency observable on the
completed-whitened source space.  Write its distinct-eigenvalue resolution as
\[
 D=\sum_{p=1}^N\xi_pP_p.
 \tag{V3.7}
\]
Suppose \(F_j=f_j(D)\), where the real functions \(f_j\) obey
\(\sum_jf_j(x)^2=1\) on \(\spec D\).

\begin{theorem}[Transport-kernel formula]
\label{thm:v3-transport-kernel}
Define
\[
 w_F(x,y)=\frac12\sum_j\bigl(f_j(x)-f_j(y)\bigr)^2.
 \tag{V3.8}
\]
Then
\[
 \boxed{
 P_p\mathcal R_F(H)P_q
 =w_F(\xi_p,\xi_q)P_pHP_q.}
 \tag{V3.9}
\]
In particular, the residual has zero frequency-diagonal blocks.  If all
\(f_j\ge0\), then
\[
 0\le w_F(x,y)\le1.
 \tag{V3.10}
\]
Moreover, the finite block Schur test gives
\[
 \boxed{
 \norm{\mathcal R_F(H)}
 \le
 \max_p\sum_q
 w_F(\xi_p,\xi_q)\norm{P_pHP_q}.}
 \tag{V3.11}
\]
\end{theorem}

\begin{proof}
Functional calculus gives
\[
 P_p[F_j,H]P_q
 =\bigl(f_j(\xi_p)-f_j(\xi_q)\bigr)P_pHP_q.
\]
Applying the commutator once more and summing proves
\textup{(V3.9)}.  If \(f(x)=(f_j(x))_j\), then
\(\norm{f(x)}_2=1\) and
\[
 w_F(x,y)=1-\ip{f(x)}{f(y)}.
\]
Nonnegativity of the entries puts the inner product in \([0,1]\), proving
\textup{(V3.10)}.  For a block vector \(u=(u_q)_q\),
\[
 \norm{(\mathcal R_F(H)u)_p}
 \le\sum_qw_F(\xi_p,\xi_q)
       \norm{P_pHP_q}\norm{u_q}.
\]
The scalar coefficient matrix is symmetric.  Its maximum row sum bounds its
\(\ell^2\) norm, proving \textup{(V3.11)}.
\end{proof}

\subsection{The canonical normalized Gaussian partition}

Fix \(A,d>0\) and, for \(j\in\mathbb Z\), put
\[
 q_j(x)=e^{-A(x-jd)^2},
 \qquad
 Z(x)=\sum_{k\in\mathbb Z}q_k(x)^2,
 \qquad
 f_j(x)=\frac{q_j(x)}{\sqrt{Z(x)}}.
 \tag{V3.12}
\]
The series and all of its derivatives converge locally uniformly.  The
function \(Z\) is positive, continuous, and \(d\)-periodic.  Hence
\[
 0<Z_-:=\min_{[0,d]}Z
 \le Z(x)\le
 Z_+:=\max_{[0,d]}Z<\infty.
 \tag{V3.13}
\]

\begin{proposition}[Gaussian overlap kernel and localization length]
\label{prop:v3-gaussian-partition}
The functions in \textup{(V3.12)} form a nonnegative quadratic partition.
Their overlap is
\[
 \boxed{
 s_F(x,y):=\sum_jf_j(x)f_j(y)
 =e^{-A(x-y)^2/2}
 \frac{Z((x+y)/2)}{\sqrt{Z(x)Z(y)}}.}
 \tag{V3.14}
\]
Consequently
\[
 0\le s_F(x,y)
 \le\frac{Z_+}{Z_-}e^{-A(x-y)^2/2},
 \qquad
 w_F(x,y)=1-s_F(x,y).
 \tag{V3.15}
\]
The quantity
\[
 L_{A,d}:=\sup_{x\in\R}
 \left(\sum_j\abs{f_j'(x)}^2\right)^{1/2}
 \tag{V3.16}
\]
is finite.  More explicitly, if
\(\omega_j(x)=f_j(x)^2\) and
\(\overline c(x)=\sum_j jd\,\omega_j(x)\), then
\[
 L_{A,d}^2
 =4A^2\sup_x\sum_j\omega_j(x)
       \bigl(jd-\overline c(x)\bigr)^2.
 \tag{V3.16a}
\]
Moreover,
\[
 \boxed{
 w_F(x,y)
 \le\min\left\{1,\frac{L_{A,d}^2}{2}\abs{x-y}^2\right\}.}
 \tag{V3.17}
\]
\end{proposition}

\begin{proof}
The partition identity is the definition of \(Z\).  Completing the square
gives
\[
 \begin{aligned}
 \sum_jq_j(x)q_j(y)
 &=\sum_j\exp\!\left(-A[(x-jd)^2+(y-jd)^2]\right)\\
 &=e^{-A(x-y)^2/2}
   \sum_j e^{-2A((x+y)/2-jd)^2},
 \end{aligned}
\]
which is \textup{(V3.14)}.  Bounds \textup{(V3.13)} give
\textup{(V3.15)}.

Differentiating the normalized Gaussian gives
\[
 f_j'(x)=2A\bigl(jd-\overline c(x)\bigr)f_j(x),
\]
which proves \textup{(V3.16a)}.  The map
\(x\mapsto f(x)=(f_j(x))_j\) is continuously differentiable into
\(\ell^2(\mathbb Z)\), and shifting \(x\) by \(d\) merely shifts its
coordinates.  Thus \(\norm{f'(x)}_2\) is continuous and \(d\)-periodic, so
the supremum in \textup{(V3.16)} is finite.  The fundamental theorem of
calculus yields
\[
 \norm{f(x)-f(y)}_2\le L_{A,d}\abs{x-y}.
\]
Since \(w_F=\tfrac12\norm{f(x)-f(y)}_2^2\), this and
\textup{(V3.10)} prove \textup{(V3.17)}.
\end{proof}

Equation \textup{(V3.14)} makes the split transparent.  The localized map
\(\mathcal A_F\) multiplies the \((p,q)\) frequency block of \(H\) by the
rapidly decreasing factor \(s_F(\xi_p,\xi_q)\); the residual retains the
complementary genuine frequency transport.

\subsection{A source-native commutator criterion}

For the spectral resolution \textup{(V3.7)}, define the block row-sum norm
\[
 \mathfrak C_D(H)
 :=\max_p\sum_{q\ne p}
 \norm{P_p[D,H]P_q}.
 \tag{V3.18}
\]
This is a stronger quantity than \(\norm{[D,H]}\); the distinction retains the
possible accumulation of many small frequency transfers.

\begin{theorem}[Gaussian frequency-transport bridge]
\label{thm:v3-commutator-bridge}
For the normalized Gaussian partition \textup{(V3.12)},
\[
 \boxed{
 \norm{\mathcal R_F(H)}
 \le\frac{L_{A,d}}{\sqrt2}\,\mathfrak C_D(H).}
 \tag{V3.19}
\]
Therefore, if
\[
 a_{A,d}(H):=\norm{\mathcal A_F(H)}
 \tag{V3.20}
\]
satisfies
\[
 \boxed{
 a_{A,d}(H)
 +\frac{L_{A,d}}{\sqrt2}\mathfrak C_D(H)\le1,}
 \tag{V3.21}
\]
then \(E\preceq M\).  Along a cutoff family, the weaker asymptotic condition
\[
 \limsup_X a_{A_X,d_X}(H_X)\le1,
 \qquad
 L_{A_X,d_X}\mathfrak C_{D_X}(H_X)\longrightarrow0
 \tag{V3.22}
\]
implies the required defect
\((\norm{H_X}-1)_+\to0\).
\end{theorem}

\begin{proof}
For \(p\ne q\),
\[
 P_p[D,H]P_q=(\xi_p-\xi_q)P_pHP_q.
\]
Put \(r=\abs{\xi_p-\xi_q}\).  By \textup{(V3.17)},
\[
 \begin{aligned}
 w_F(\xi_p,\xi_q)\norm{P_pHP_q}
 &\le
 \min\left\{\frac1r,\frac{L_{A,d}^2r}{2}\right\}
 \norm{P_p[D,H]P_q}\\
 &\le\frac{L_{A,d}}{\sqrt2}
 \norm{P_p[D,H]P_q},
 \end{aligned}
\]
because \(\min\{r^{-1},L^2r/2\}\le L/\sqrt2\).  Sum over \(q\) and apply
\textup{(V3.11)} to prove \textup{(V3.19)}.  Theorem~\ref{thm:v3-ims}
then proves \textup{(V3.21)}--\textup{(V3.22)}.
\end{proof}

\begin{corollary}[Margin form of the completed--Euler bridge]
\label{cor:v3-margin}
If, on one exhaustive supported source family,
\[
 a_{A_X,d_X}(H_X)\le1-\delta_X,
 \qquad
 \mathfrak C_{D_X}(H_X)
 \le\frac{\sqrt2\,\delta_X}{L_{A_X,d_X}},
 \tag{V3.23}
\]
then \(\widehat{\mathbb E}_X\preceq\widehat M_X\) at every cutoff.  If the
right inequality holds with \(\delta_X+o(1)\) while
\(a_{A_X,d_X}(H_X)\le1-\delta_X+o(1)\), then the multiplicative upper-bridge
defect tends to zero.
\end{corollary}

\begin{proof}
Apply \Cref{thm:v3-commutator-bridge} and conjugate
\(H_X\preceq I\) by \(\widehat M_X^{1/2}\) on its support.
\end{proof}

\subsection{The bare Gaussian source frame is explicitly conditionable}

The V2 Mellin-centre vectors have a simple Toeplitz Gram before the arithmetic
response is applied.  This separates source-frame conditioning from the
relative-response transport in \textup{(V3.19)}.

\begin{proposition}[Gaussian Toeplitz source Gram]
\label{prop:v3-toeplitz}
Let
\[
 g_j(x)=\left(\frac{2A}{\pi}\right)^{1/4}e^{-A(x-jd)^2}
 \qquad(j\in\mathbb Z).
\]
Then \(\norm{g_j}_2=1\) and
\[
 \boxed{
 \ip{g_i}{g_j}_{L^2(\R)}
 =e^{-Ad^2(i-j)^2/2}.}
 \tag{V3.24}
\]
For every finite subbank, put
\[
 r_{A,d}=2\sum_{n\ge1}e^{-Ad^2n^2/2}.
 \tag{V3.25}
\]
Its Gram matrix \(G\) obeys
\[
 \boxed{
 (1-r_{A,d})I\preceq G\preceq(1+r_{A,d})I.}
 \tag{V3.26}
\]
When \(r_{A,d}<1\), the canonical whitening \(G^{-1/2}\) has condition number
at most
\(\sqrt{(1+r_{A,d})/(1-r_{A,d})}\), uniformly in the number of centres.
\end{proposition}

\begin{proof}
Gaussian integration gives the norm and, after completing the square,
\[
 \int_\R g_i(x)g_j(x)\dd x=e^{-Ad^2(i-j)^2/2}.
\]
The off-diagonal absolute row sum is at most \(r_{A,d}\).  The Schur test, or
Gershgorin's theorem for the Hermitian Gram, gives
\textup{(V3.26)}.  Functional calculus gives the whitening bound.
\end{proof}

\begin{firewall}[Gaussian overlap is not Euler domination]
\label{fire:v3-overlap}
The Toeplitz estimate \textup{(V3.26)} conditions the \emph{bare profile
source}.  It does not bound \(\mathfrak C_D(H)\) and does not imply
\(E\preceq M\).  Indeed, take two orthogonal frequency cells and
\[
 H=\begin{pmatrix}1&1\\1&1\end{pmatrix}\succeq0.
\]
Each diagonal cell equals one and the source overlap is zero, while
\(\norm H=2\).  The missing unit is precisely the off-frequency transport.
Any GRH argument that replaces the mixed completed--Euler response by the
Gaussian profile Gram loses the quantity it needs to estimate.
\end{firewall}

\subsection{Interpretation on the V074 arithmetic source}

To apply \Cref{thm:v3-commutator-bridge}, three constructions must occur in
the stated order.
\begin{enumerate}[label=\textup{(A\arabic*)},leftmargin=3em]
\item Form the complete completed, Euler, and nuisance mixed score Gram on one
canonical arithmetic carrier, and short the common nuisance family.
\item Whiten by the supported completed metric to obtain \(H_X\).  The
Mellin-frequency observable must be transported through the same whitening;
declaring a diagonal frequency coordinate before this step is insufficient.
\item Apply the predeclared Gaussian quadratic partition to that transported
observable and evaluate both \(a_{A,d}(H_X)\) and the full block row-sum
commutator \(\mathfrak C_{D_X}(H_X)\).
\end{enumerate}

The frequency-localized term \(a_{A,d}(H_X)\) asks whether the completed
gamma/conductor metric dominates the Euler response within one Gaussian
Mellin cell.  The commutator term asks whether the relative response transports
source mass between separated cells.  Neither is determined by finite gamma
totality, local Euler-cell Schur positivity, or profile overlap.

There is nevertheless a concrete analytic opening.  In the logarithmic
variable, multiplication by \(\log n\) differentiates Mellin frequency.
Consequently \([D_X,H_X]\) is the correct place to seek a first-moment or
boundary-Ward estimate for the mixed completed--Euler kernel.  The estimate
must be proved after completed whitening and must control the block row sum,
not only the operator norm of the commutator.

\subsection{V3 status ledger and next target}

\begin{center}
\begin{tabularx}{\textwidth}{@{}p{0.27\textwidth}p{0.18\textwidth}X@{}}
\toprule
V3 component & Status & Exact conclusion \\
\midrule
Relative-metric frequency split & proved & Exact IMS identity
\textup{(V3.3)}--\textup{(V3.4)}. \\
Gaussian localization kernel & proved & Exact overlap
\textup{(V3.14)} and transport weight \textup{(V3.15)}. \\
Transport residual bound & proved & The concrete commutator criterion
\textup{(V3.19)}--\textup{(V3.23)}. \\
Bare Gaussian source conditioning & proved & Uniform Toeplitz frame bound
when \(r_{A,d}<1\). \\
Localized completed--Euler norm & open & The arithmetic estimate
\(a_{A,d}(H_X)\le1-\delta_X\) is not supplied by localization. \\
Frequency commutator row sum & open & Must be obtained from the actual
completed--Euler mixed response after whitening. \\
Restoring lower bridge & open & Independent of the upper localization theorem. \\
Cofinal passage and GRH & open & Require both bridges and the same-history
source-anchored descent. \\
\bottomrule
\end{tabularx}
\end{center}

V4 should compute the commutator in an explicit kernel model.  Write the
completed-whitened relative response in the Mellin-frequency representation as
\(H_X(\xi,\eta)\).  The exact identity
\[
 [D_X,H_X](\xi,\eta)=(\xi-\eta)H_X(\xi,\eta)
\]
turns \(\mathfrak C_{D_X}(H_X)\) into a first off-diagonal moment.  The next
attack is to derive this kernel from the Gaussian Plancherel formula and the
joint completed--Euler score panel, then determine whether the prime-side
kernel has a summable first moment after the gamma/conductor normalization.
If it does not, the first non-summable frequency block should be returned as a
specific upper-bridge witness and compared with the translated or
taper-boundary row of the restoring short.

\begin{assessment}[V3 acquisition]
V3 converts the qualitative request for ``Gaussian frequency localization''
into one exact operator inequality.  The completed--Euler defect is bounded by
a localized completely positive norm plus a source-native frequency-transport
commutator.  Gaussian profile overlap is now proved to condition the bare
source but is also proved insufficient for Euler domination.  The next genuine
GRH estimate is therefore a first off-diagonal moment bound for the actual
completed-whitened mixed response, together with a local cell margin; no
individual-zero reconstruction is required.
\end{assessment}


\section{V4: the first-moment kernel and the whitening tax}
\label{sec:grh-v4}
% =============================================================================

\begin{center}
\large\bfseries Version V4 continuation
\end{center}

\subsection{Why the first-moment question must be asked in a row algebra}

V3 reduced the Gaussian transport residual to
\(\mathfrak C_D(H)\), a block row sum of the commutator of the
completed-whitened relative metric.  An ordinary operator-norm estimate of a
commutator does not control that quantity.  The appropriate finite algebra can
be written explicitly.

Let
\[
 D=\sum_{p=1}^N\xi_pP_p
 \tag{V4.1}
\]
be self-adjoint on a finite-dimensional Hilbert space.  For an operator \(T\)
put
\[
 \norm{T}_{1,D}:=\max_p\sum_q\norm{P_pTP_q}.
 \tag{V4.2}
\]

\begin{lemma}[Block-row Banach algebra]
\label{lem:v4-row-algebra}
For all \(A,B\),
\[
 \boxed{\norm{AB}_{1,D}\le\norm A_{1,D}\norm B_{1,D}.}
 \tag{V4.3}
\]
If \(T^*=T\) or \(T^*=-T\), then
\[
 \boxed{\norm T\le\norm T_{1,D}.}
 \tag{V4.4}
\]
For self-adjoint \(H\), the V3 transport quantity is exactly
\[
 \boxed{
 \mathfrak C_D(H)=\norm{[D,H]}_{1,D}.}
 \tag{V4.5}
\]
\end{lemma}

\begin{proof}
Insert \(I=\sum_rP_r\) between \(A\) and \(B\).  For each \(p\),
\[
 \begin{aligned}
 \sum_q\norm{P_pABP_q}
 &\le\sum_r\norm{P_pAP_r}
              \sum_q\norm{P_rBP_q}\\
 &\le\norm B_{1,D}\sum_r\norm{P_pAP_r}.
 \end{aligned}
\]
Taking the maximum proves \textup{(V4.3)}.  When \(T\) is self-adjoint or
skew-adjoint, the scalar matrix
\(a_{pq}=\norm{P_pTP_q}\) is symmetric.  The block Schur test therefore gives
\(\norm T\le\max_p\sum_qa_{pq}\), proving \textup{(V4.4)}.  Finally,
\[
 P_p[D,H]P_q=(\xi_p-\xi_q)P_pHP_q,
\]
and the diagonal blocks vanish.  Comparison with \textup{(V3.18)} proves
\textup{(V4.5)}.
\end{proof}

\subsection{Exact propagation of a commutator through whitening}

The square root in \(H=M^{-1/2}EM^{-1/2}\) cannot be treated as a scalar
normalization.  The next theorem gives the precise extra term and retains the
row norm required by V3.

\begin{theorem}[Row-resolvent whitening formula]
\label{thm:v4-whitening-row}
Let \(M\succ0\) and \(E=E^*\) on the space of \textup{(V4.1)}, put
\[
 R=M^{-1/2},\qquad H=RER,
 \tag{V4.6}
\]
and define
\[
 r_D(t):=\norm{(M+tI)^{-1}}_{1,D},
 \qquad
 \kappa_D(M):=\frac1\pi\int_0^\infty
 t^{-1/2}r_D(t)^2\dd t.
 \tag{V4.7}
\]
Then the integral is finite and
\[
 \boxed{
 [D,R]=-\frac1\pi\int_0^\infty t^{-1/2}
 (M+tI)^{-1}[D,M](M+tI)^{-1}\dd t,}
 \tag{V4.8}
\]
where the integral converges in every finite-dimensional norm.  Consequently
\[
 \boxed{
 \begin{aligned}
 \mathfrak C_D(H)
 \le{}&\norm R_{1,D}^{,2}\norm{[D,E]}_{1,D}\\
 &+2\norm R_{1,D}\norm E_{1,D}\,
 \kappa_D(M)\norm{[D,M]}_{1,D}.
 \end{aligned}}
 \tag{V4.9}
\]
If only operator norm is retained and \(M\succeq mI\), then
\[
 \boxed{
 \norm{[D,H]}
 \le m^{-1}\norm{[D,E]}
 +m^{-2}\norm E\norm{[D,M]}.}
 \tag{V4.10}
\]
The weaker estimate \textup{(V4.10)} does not by itself imply
\textup{(V4.9)}.
\end{theorem}

\begin{proof}
The Stieltjes representation
\[
 M^{-1/2}=\frac1\pi\int_0^\infty
 t^{-1/2}(M+tI)^{-1}\dd t
\]
and
\[
 [D,(M+tI)^{-1}]
 =-(M+tI)^{-1}[D,M](M+tI)^{-1}
\]
give \textup{(V4.8)}.  In finite dimension \(r_D(t)\) is bounded near zero
and is \(O(t^{-1})\) at infinity, so \(\kappa_D(M)<\infty\).  Lemma
\ref{lem:v4-row-algebra} gives
\[
 \norm{[D,R]}_{1,D}
 \le\kappa_D(M)\norm{[D,M]}_{1,D}.
\]
Expanding
\[
 [D,H]=[D,R]ER+R[D,E]R+RE[D,R]
 \tag{V4.11}
\]
and using the row-algebra inequality proves \textup{(V4.9)} together with
\textup{(V4.5)}.

For the operator-norm statement,
\[
 \norm{[D,R]}
 \le\frac{\norm{[D,M]}}\pi
 \int_0^\infty\frac{t^{-1/2}}{(m+t)^2}\dd t
 =\frac{\norm{[D,M]}}{2m^{3/2}}.
\]
Insert this and \(\norm R\le m^{-1/2}\) into \textup{(V4.11)}.
\end{proof}

The two terms in \textup{(V4.9)} have different meanings.  The first is the
prime-side first moment in the declared frequency coordinate.  The second is
the \emph{whitening tax}: frequency transport already present in the completed
metric and propagated through its inverse square root.  Neither term may be
deleted merely because \(M\) is positive.

\subsection{The canonical joint-score version of the kernel}

The preceding theorem begins with a self-adjoint \(D\) on the whitened source
space.  V074's simultaneous nuisance short canonically supplies the metrics,
but it does not automatically supply a frequency observable.  That observable
must itself be a jointly visible score moment.

Let \(\widehat A\) and \(\widehat S\) be the nuisance-shorted completed and
Euler syntheses, let
\[
 M=\widehat A^*\widehat A,
 \qquad E=\widehat S^*\widehat S,
 \tag{V4.12}
\]
and work on \(\supp M\), where \(M\succ0\).  If \(Q=Q^*\) is the declared
Mellin-frequency observable on the completed score carrier, its compressed
first-moment form is
\[
 K_Q:=\widehat A^*Q\widehat A.
 \tag{V4.13}
\]
Thus \(K_Q\) is operational only when the frequency-weighted completed score
row is retained in the same joint panel.

\begin{theorem}[Completed-score commutator kernel]
\label{thm:v4-joint-score-kernel}
With \(R=M^{-1/2}\), define
\[
 D_Q=RK_QR,
 \qquad H=RER.
 \tag{V4.14}
\]
Then both operators are self-adjoint and
\[
 \boxed{
 [D_Q,H]
 =R\bigl(K_QM^{-1}E-EM^{-1}K_Q\bigr)R.}
 \tag{V4.15}
\]
If \(D_Q=\sum_p\xi_pP_p\), its exact first off-diagonal moment is therefore
\[
 \boxed{
 \mathfrak C_{D_Q}(H)
 =\max_p\sum_q
 \norm{P_pR
 (K_QM^{-1}E-EM^{-1}K_Q)RP_q}.}
 \tag{V4.16}
\]
In particular, the V3 commutator is a statistic of the complete triple
\((M,E,K_Q)\), not of the two marginal metrics alone.
\end{theorem}

\begin{proof}
Self-adjointness follows by congruence.  Since \(R^2=M^{-1}\), direct
multiplication gives
\[
 D_QH=RK_QM^{-1}ER,
 \qquad
 HD_Q=REM^{-1}K_QR,
\]
which proves \textup{(V4.15)}.  Apply \textup{(V4.5)} in the spectral
resolution of \(D_Q\) to obtain \textup{(V4.16)}.
\end{proof}

\begin{firewall}[Frequency spectrum without joint orientation is insufficient]
\label{fire:v4-frequency-orientation}
Even the eigenvalues of the missing moment form do not determine the
commutator.  Take
\[
 M=I,qquad
 E=\begin{pmatrix}1&0\\0&0\end{pmatrix},qquad
 K_0=\begin{pmatrix}0&0\\0&1\end{pmatrix},qquad
 K_1=\frac12\begin{pmatrix}1&1\\1&1\end{pmatrix}.
\]
Both \(K_0\) and \(K_1\) are positive with spectrum \(\{0,1\}\).  But
\[
 [K_0,E]=0,
 \qquad
 [K_1,E]=\frac12\begin{pmatrix}0&-1\\1&0\end{pmatrix},
 \qquad
 \norm{[K_1,E]}=\frac12.
\]
Thus a list of Mellin frequencies attached separately to the completed source
does not fix their orientation relative to the Euler metric.  The row
\textup{(V4.13)} must be retained before whitening, exactly as the mixed
completed--Euler row must be retained before nuisance shorting.
\end{firewall}

\subsection{The raw localized Euler kernel has a divergent absolute first moment}

The localized Weil chart in the arithmetic manuscript makes the prime-side
question completely explicit.  This chart is used here only to audit the
transport mechanism; its signed Euler response is not identified with the
positive Euler score metric \(E\) in \textup{(V4.12)}.

Let \(I_A=(-A,A)\), let \(Xg(x)=xg(x)\), and write
\[
 (T_{A,u}g)(x)
 =\mathbf 1_{I_A}(x)\mathbf 1_{I_A}(x+u)g(x+u).
 \tag{V4.17}
\]
For \(u_n=\log n\) and \(\lambda_n=\Lambda(n)/\sqrt n\), define the finite
signed Euler response
\[
 \mathsf E_{\chi,A}
 :=\sum_{\substack{2\le n<e^{2A}\\ \gcd(n,f)=1}}
 \lambda_n\bigl(
 \chi(n)T_{A,u_n}+\overline{\chi(n)}T_{A,-u_n}
 \bigr).
 \tag{V4.18}
\]

\begin{theorem}[Exact shift commutator and central-row obstruction]
\label{thm:v4-raw-prime-moment}
On \(L^2(I_A)\),
\[
 \boxed{[X,T_{A,u}]=-uT_{A,u},}
 \tag{V4.19}
\]
and hence
\[
 \boxed{
 [X,\mathsf E_{\chi,A}]
 =\sum_{\substack{2\le n<e^{2A}\\ \gcd(n,f)=1}}
 \lambda_nu_n
 \bigl(-\chi(n)T_{A,u_n}
 +\overline{\chi(n)}T_{A,-u_n}\bigr).}
 \tag{V4.20}
\]
Regard the finite shift sum as an operator-valued atomic kernel.  The total
variation of its row at \(x\in I_A\) is exactly
\[
 \boxed{
 \begin{aligned}
 c_{\chi,A}(x)
 =\sum_{\substack{2\le n<e^{2A}\\ \gcd(n,f)=1}}
 \frac{\Lambda(n)\log n}{\sqrt n}
 \bigl(&\mathbf 1_{I_A}(x+\log n)\\
      &+\mathbf 1_{I_A}(x-\log n)\bigr).
 \end{aligned}}
 \tag{V4.21}
\]
In particular,
\[
 \boxed{
 c_{\chi,A}(0)
 =2\sum_{\substack{2\le n<e^A\\ \gcd(n,f)=1}}
 \frac{\Lambda(n)\log n}{\sqrt n}
 \longrightarrow\infty
 \qquad(A\to\infty).}
 \tag{V4.22}
\]
Therefore the raw prime kernel has no cofinally summable absolute first-row
moment.  Any proof of the V3 bridge that takes absolute values prime by prime
before completed whitening must fail.
\end{theorem}

\begin{proof}
Inside the surviving overlap interval,
\[
 XT_{A,u}g(x)=xg(x+u),
 \qquad
 T_{A,u}Xg(x)=(x+u)g(x+u),
\]
which proves \textup{(V4.19)}; the cutoff multipliers commute with \(X\).
Summing gives \textup{(V4.20)}.  Distinct positive integers give distinct
positive shifts, and the positive and negative shift atoms are disjoint.
Since \(\abs{\chi(n)}=1\) for \((n,f)=1\), taking row total variation gives
\textup{(V4.21)}.  At (x=0), both shift atoms survive precisely when
\(\log n<A\), proving the equality in \textup{(V4.22)}.

For divergence, retain only unramified primes.  Apart from finitely many
exceptions,
\[
 \frac{(\log p)^2}{\sqrt p}\ge\frac1p.
\]
Euler's divergence of \(\sum_p1/p\), unchanged after deleting the finitely
many primes dividing \(f\), completes the proof.
\end{proof}

\begin{remark}[What the obstruction does and does not say]
The theorem rules out an \(\ell^1\)-in-prime triangle proof.  It does not show
that the signed operator norm of \textup{(V4.20)} diverges, nor that the
completed-whitened commutator \textup{(V4.15)} is large.  Character phases,
overlap geometry, and the nonlocal inverse square root may cancel.  Such
cancellation must be kept until a square-function, Cotlar, or joint-score
estimate is applied; total variation deliberately erases it.
\end{remark}

\subsection{An explicit archimedean Ward kernel}

The arithmetic manuscript's parity-renormalized form also permits the
completed-side commutator to be computed before whitening.  Let
\(\mathsf A_{\chi,A}\) be the operator associated with its form
\textup{(W.21)}.  On a smooth compactly supported core it is the sum of the
singular difference operator
\[
 (\mathsf D_Ag)(x)
 =\frac12\int_{I_A}\frac{g(x)-g(y)}{\abs{x-y}}\dd y,
 \tag{V4.23}
\]
a multiplication operator, and the smooth integral operator with kernel
\(s_\alpha(\abs{x-y})\).

\begin{proposition}[Completed-side frequency commutator]
\label{prop:v4-archimedean-ward}
On the smooth core,
\[
 \boxed{
 ([X,\mathsf A_{\chi,A}]g)(x)
 =\int_{I_A}b_{\alpha,A}(x,y)g(y)\dd y,}
 \tag{V4.24}
\]
where
\[
 \boxed{
 b_{\alpha,A}(x,y)
 =-\frac12\operatorname{sgn}(x-y)
 +(x-y)s_\alpha(\abs{x-y}).}
 \tag{V4.25}
\]
This commutator extends to a Hilbert--Schmidt operator and obeys
\[
 \boxed{
 \norm{[X,\mathsf A_{\chi,A}]}
 \le A+\sqrt{\frac83},A^2
 \sup_{0\le u\le2A}\abs{s_\alpha(u)}.}
 \tag{V4.26}
\]
Thus the completed-side whitening tax is a concrete fixed-support kernel, not
an unspecified error.  The bound is not uniform in \(A\).
\end{proposition}

\begin{proof}
The conductor, boundary-logarithm, and other multiplication terms commute with
\(X\).  From \textup{(V4.23)},
\[
 ([X,\mathsf D_A]g)(x)
 =-\frac12\int_{I_A}\operatorname{sgn}(x-y)g(y)\dd y.
\]
For the smooth kernel the commutator inserts the factor \(x-y\).  This proves
\textup{(V4.24)}--\textup{(V4.25)}.  The first kernel has Hilbert--Schmidt norm
\(A\).  Since \(s_\alpha\) is continuous after removal of its singularity,
the second has Hilbert--Schmidt norm at most
\[
 \sup_{[0,2A]}\abs{s_\alpha}
 \left(\int_{I_A^2}(x-y)^2\dd x\dd y\right)^{1/2}
 =\sqrt{\frac83}\,A^2\sup_{[0,2A]}\abs{s_\alpha}.
\]
The operator norm is at most the Hilbert--Schmidt norm.
\end{proof}

\subsection{V4 decision and upstream target}

The first-moment route now has a precise logical shape.
\begin{enumerate}[label=\textup{(B\arabic*)},leftmargin=3em]
\item The quantity required by the Gaussian IMS estimate is exactly the
block-row norm of the \emph{post-whitening} commutator, by
\textup{(V4.5)}.
\item Whitening contributes the resolvent-row term in \textup{(V4.9)}.  A
spectral floor for \(M\) controls only operator norm and cannot replace the
missing row-resolvent estimate.
\item In the canonical score realization, the exact numerator is
\(K_QM^{-1}E-EM^{-1}K_Q\).  The frequency moment \(K_Q\) must be exported from
the same operational panel; separate frequency labels do not determine it.
\item The raw Weil prime shifts have the divergent absolute row moment
\textup{(V4.22)}.  Hence no proof may take absolute values before grouping the
signed prime cells and applying completed whitening.
\item The completed-side commutator has the explicit Ward kernel
\textup{(V4.25)} at every fixed support wall.  Cofinal control still requires
its propagation through the inverse square root in the stronger row algebra.
\end{enumerate}

\begin{center}
\begin{tabularx}{\textwidth}{@{}p{0.28\textwidth}p{0.17\textwidth}X@{}}
\toprule
V4 component & Status & Exact conclusion \\
\midrule
Block-row commutator algebra & proved &
\(\mathfrak C_D(H)=\norm{[D,H]}_{1,D}\) and the norm is submultiplicative. \\
Whitening propagation & proved & Exact resolvent formula
\textup{(V4.8)} and sufficient row estimate \textup{(V4.9)}. \\
Joint-score frequency kernel & proved & Exact numerator
\(K_QM^{-1}E-EM^{-1}K_Q\) in \textup{(V4.15)}. \\
Raw prime first moment & disproved as an absolute-summability route & The
central row diverges by \textup{(V4.22)}. \\
Completed Ward kernel & proved at fixed support & Formula
\textup{(V4.25)} and Hilbert--Schmidt bound \textup{(V4.26)}. \\
Signed post-whitening row bound & open & Requires cancellation before absolute
values and a uniform row-resolvent estimate. \\
Localized completed--Euler margin & open & Still requires
\(a_{A,d}(H_X)\le1-\delta_X\). \\
Restoring bridge and cofinal descent & open & Unchanged and logically
independent. \\
\bottomrule
\end{tabularx}
\end{center}

The literal prime-by-prime first-moment attack has reached a rigorous
obstruction rather than a closure estimate.  V5 should therefore move one
level upstream.  The first option is a signed square-function or Cotlar bound
for the grouped joint-score numerator in \textup{(V4.15)}, retaining character
phases through whitening.  If the required row-resolvent localization cannot
be proved from the score carrier, the cleaner route is V074's Petz--Euler
certificate: formulate a quantitative modular-invariance defect for the Euler
subalgebra inside the completed algebra and determine whether it yields an
approximate metric contraction.  Either route must be tested on the actual
joint source panel, not on the marginal prime coefficients.

\begin{assessment}[V4 acquisition]
V4 computes the proposed first-moment mechanism far enough to decide its
naive form.  The unwhitened Euler shift kernel has a divergent absolute central
row, while noncommuting completed whitening contributes a second, explicitly
resolvent-mediated transport term.  The correct finite observable is the
joint-score commutator \textup{(V4.15)}.  This narrows the upper bridge to a
signed post-whitening cancellation theorem or an upstream modular
conditional-expectation estimate; it does not prove GRH.
\end{assessment}


% =============================================================================
\section{V5: a quantitative Petz predictor for the upper bridge}
\label{sec:grh-v5}
% =============================================================================

\begin{center}
\large\bfseries Version V5 continuation
\end{center}

\subsection{Moving upstream from the divergent absolute prime row}

V4 proved that the raw localized Euler shift kernel has a divergent absolute
first-row moment.  It also proved that completed whitening inserts a second
resolvent-row commutator.  Thus an estimate obtained by taking absolute values
prime by prime cannot close the V3 IMS criterion.  V074 contains a genuinely
upstream alternative: a state-preserving map from the completed algebra to the
Euler algebra is automatically contractive in the physical score norm.

Exact modular invariance produces a conditional expectation, as recorded in
V074.  The useful quantitative observation is that a canonical unital
state-preserving completely positive \emph{predictor} exists even without
modular invariance.  It need not fix the Euler algebra, but its mismatch with
the actual Euler score can be measured in the same joint source panel.  This
section turns that mismatch into an explicit bound for the relative metric
defect.

\subsection{The always-defined finite Petz predictor}

Let \(\mathcal M\) be a finite-dimensional \(C^*\)-algebra with faithful
tracial state \(\tau\), let \(\mathcal N\subseteq\mathcal M\) be a unital
subalgebra, and let
\[
 \mathcal E_{\mathcal N}:\mathcal M\longrightarrow\mathcal N
 \tag{V5.1}
\]
be the \(\tau\)-preserving conditional expectation.  Let
\(\rho\in\mathcal M\) be positive and invertible with \(\tau(\rho)=1\), put
\[
 \varphi(x)=\tau(\rho x),
 \qquad
 \rho_{\mathcal N}=\mathcal E_{\mathcal N}(\rho),
 \tag{V5.2}
\]
and equip \(\mathcal M\) with
\[
 \ip{x}{y}_{\rho}:=\varphi(x^*y).
 \tag{V5.3}
\]

\begin{theorem}[Canonical Petz score predictor]
\label{thm:v5-petz-predictor}
The map
\[
 \boxed{
 \Pi_{\rho,\mathcal N}(x)
 :=\rho_{\mathcal N}^{-1/2}
 \mathcal E_{\mathcal N}
   \bigl(\rho^{1/2}x\rho^{1/2}\bigr)
 \rho_{\mathcal N}^{-1/2}}
 \tag{V5.4}
\]
is unital, completely positive, has range in \(\mathcal N\), and preserves
the physical state:
\[
 \boxed{\varphi(\Pi_{\rho,\mathcal N}(x))=\varphi(x).}
 \tag{V5.5}
\]
It is a contraction on the physical score Hilbert space,
\[
 \boxed{
 \norm{\Pi_{\rho,\mathcal N}(x)}_{2,\rho}
 \le\norm{x}_{2,\rho},}
 \tag{V5.6}
\]
and the same assertion holds at every matrix amplification.

If \(\mathcal N\) is invariant under the modular group
\(\rho^{it}(\cdot)\rho^{-it}\), then \(\Pi_{\rho,\mathcal N}\) is the
\(\varphi\)-preserving conditional expectation and fixes \(\mathcal N\).
Without that invariance it is generally not idempotent and must be treated as
a predictor, not as an exact source identification.
\end{theorem}

\begin{proof}
Every factor in \textup{(V5.4)} is completely positive.  Moreover
\[
 \Pi_{\rho,\mathcal N}(I)
 =\rho_{\mathcal N}^{-1/2}
  \mathcal E_{\mathcal N}(\rho)
  \rho_{\mathcal N}^{-1/2}=I.
\]
For \(y\in\mathcal N\), tracial orthogonality of
\(\mathcal E_{\mathcal N}\) gives
\(\tau(\rho y)=\tau(\rho_{\mathcal N}y)\).  Since the output of
\(\Pi_{\rho,\mathcal N}\) lies in \(\mathcal N\), cyclicity of \(\tau\)
yields
\[
 \begin{aligned}
 \varphi(\Pi_{\rho,\mathcal N}(x))
 &=\tau\!\left(
 \rho_{\mathcal N}^{1/2}
 \mathcal E_{\mathcal N}(\rho^{1/2}x\rho^{1/2})
 \rho_{\mathcal N}^{-1/2}\right)\\
 &=\tau(\rho^{1/2}x\rho^{1/2})
 =\varphi(x).
 \end{aligned}
\]
Kadison--Schwarz and \textup{(V5.5)} now give
\[
 \begin{aligned}
 \norm{\Pi_{\rho,\mathcal N}(x)}_{2,\rho}^2
 &\le
 \varphi\!\left(\Pi_{\rho,\mathcal N}(x^*x)\right)\\
 &=\varphi(x^*x).
 \end{aligned}
\]
Applying the same proof to
\(\operatorname{id}_{M_k}\otimes\Pi_{\rho,\mathcal N}\) proves complete
score contraction.  The last assertion is the finite-dimensional Takesaki
criterion; under modular invariance the state-preserving expectation is
unique.  Its failure is not used in the contraction proof.
\end{proof}

\begin{remark}[Centering is preserved]
If \(x\) is a centered score, \(\varphi(x)=0\), then
\(\varphi(\Pi_{\rho,\mathcal N}(x))=0\).  Thus the predictor can be inserted
after score centering without adding a constant row.
\end{remark}

\subsection{Petz residual control of the completed--Euler metric}

Let \(\mathcal K\) be a finite coefficient space.  Regard the completed and
Euler score syntheses as maps into the physical \(L^2(\rho)\) carrier.  Let
\(\mathsf Q\) be the common orthogonal projection away from the declared
nuisance score span and put
\[
 \widehat A=\mathsf Q A,
 \qquad
 \widehat S=\mathsf Q S,
 \qquad
 M=\widehat A^*\widehat A,
 \qquad
 E=\widehat S^*\widehat S.
 \tag{V5.7}
\]
Work on \(\supp M\), so \(M\succ0\).  Define the nuisance-shorted Petz
prediction and its physical residual by
\[
 T=\mathsf Q\Pi_{\rho,\mathcal N}\widehat A,
 \qquad
 R_E=\widehat S-T.
 \tag{V5.8}
\]
Here the map in \textup{(V5.4)} is applied componentwise to the score
synthesis.

\begin{theorem}[Quantitative Petz--Euler bridge]
\label{thm:v5-petz-bridge}
Put
\[
 P=T^*T,\qquad
 \Delta_E=R_E^*R_E,
 \tag{V5.9}
\]
and define the two source-normalized quantities
\[
 q:=\norm{TM^{-1/2}},
 \qquad
 \varepsilon:=\norm{R_EM^{-1/2}}
 =\norm{M^{-1/2}\Delta_EM^{-1/2}}^{1/2}.
 \tag{V5.10}
\]
Then
\[
 \boxed{P\preceq M,\qquad 0\le q\le1,}
 \tag{V5.11}
\]
and the completed--Euler relative metric
\[
 H=M^{-1/2}EM^{-1/2}
 \tag{V5.12}
\]
obeys
\[
 \boxed{
 \norm H\le(q+\varepsilon)^2\le(1+\varepsilon)^2.}
 \tag{V5.13}
\]
Consequently its sharp upper-bridge defect satisfies
\[
 \boxed{
 (\norm H-1)_+\le2\varepsilon+\varepsilon^2.}
 \tag{V5.14}
\]
If \(q+\varepsilon\le1\), then \(E\preceq M\) exactly.

For every \(\theta>0\), there is also the order estimate
\[
 \boxed{
 E\preceq(1+\theta)M+
 (1+\theta^{-1})\Delta_E.}
 \tag{V5.15}
\]
All conclusions persist at matrix amplifications when the same nuisance
projection and source labels are amplified.
\end{theorem}

\begin{proof}
Theorem~\ref{thm:v5-petz-predictor} and orthogonality of \(\mathsf Q\) give
\[
 \norm{Tc}_{2,\rho}
 \le\norm{\widehat Ac}_{2,\rho}
 \qquad(c\in\mathcal K).
\]
This is exactly \(P\preceq M\), which proves \textup{(V5.11)} after
whitening.  Since \(\widehat S=T+R_E\),
\[
 \norm H^{1/2}
 =\norm{\widehat SM^{-1/2}}
 \le\norm{TM^{-1/2}}+\norm{R_EM^{-1/2}}
 =q+\varepsilon.
\]
This proves \textup{(V5.13)} and hence \textup{(V5.14)}.  If
\(q+\varepsilon\le1\), then \(\norm H\le1\), which is equivalent to the
metric order on \(\supp M\).

Finally, for every coefficient vector \(c\), Young's inequality gives
\[
 \norm{Tc+R_Ec}^2
 \le(1+\theta)\norm{Tc}^2
 +(1+\theta^{-1})\norm{R_Ec}^2.
\]
Use \(T^*T\preceq M\) to obtain \textup{(V5.15)}.  Complete score
contraction from \Cref{thm:v5-petz-predictor} gives the amplified statement.
\end{proof}

\begin{corollary}[Cofinal residual criterion]
\label{cor:v5-cofinal-petz}
Along an exhaustive cutoff family, suppose
\[
 \Ker M_X\subseteq\Ker E_X
 \quad\hbox{and}\quad
 \varepsilon_X
 :=\norm{R_{E,X}M_X^{\dagger/2}}\longrightarrow0.
 \tag{V5.16}
\]
Then
\[
 \boxed{
 \bigl(\norm{H_{E/M,X}}-1\bigr)_+\longrightarrow0.}
 \tag{V5.17}
\]
Thus the completed--Euler upper bridge reduces to a single normalized
joint-score residual.  The restoring lower bridge and the final cofinal
descent remain independent requirements.
\end{corollary}

\begin{proof}
Restrict to \(\supp M_X\) and apply \textup{(V5.14)}.  The kernel condition
handles the orthogonal complement.
\end{proof}

\subsection{The residual is itself a reconstructible joint Gram}

The criterion above is useful only if \(\Delta_E\) is available before the
target is inspected.  It is: one additional mixed score row suffices.

\begin{proposition}[Exact Petz residual Gram]
\label{prop:v5-residual-gram}
Let
\[
 C_{TS}:=T^*\widehat S.
 \tag{V5.18}
\]
Then
\[
 \boxed{
 \Delta_E
 =E+P-C_{TS}-C_{TS}^*.}
 \tag{V5.19}
\]
The block
\[
 \boxed{
 \begin{pmatrix}
 P&C_{TS}\\
 C_{TS}^*&E
 \end{pmatrix}\succeq0}
 \tag{V5.20}
\]
is the joint Gram of the predicted and actual Euler syntheses.  Hence the
number required in \textup{(V5.16)} is
\[
 \boxed{
 \varepsilon^2
 =\norm{
 M^{-1/2}
 (E+P-C_{TS}-C_{TS}^*)
 M^{-1/2}}.}
 \tag{V5.21}
\]
It is invariant under a common unitary change of physical score carrier and
under source-coordinate changes transported through all four blocks.
\end{proposition}

\begin{proof}
Expand
\[
 R_E^*R_E
 =(\widehat S-T)^*(\widehat S-T).
\]
This gives \textup{(V5.19)}.  Formula \textup{(V5.20)} is the Gram of the
column synthesis \((T,\widehat S)\).  Congruence by \(M^{-1/2}\) gives
\textup{(V5.21)}, and Gram covariance proves the invariance statements.
\end{proof}

\begin{firewall}[Marginal score norms do not determine the Petz residual]
\label{fire:v5-marginals}
The cross row \(C_{TS}\) cannot be reconstructed from \(P\) and \(E\).
On a one-dimensional coefficient space, take a two-dimensional physical
score carrier and \(T c=ce_1\).  The two choices
\[
 \widehat S_0c=ce_1,
 \qquad
 \widehat S_1c=ce_2
\]
have the same marginal Grams \(P=E=1\).  But
\[
 \Delta_{E,0}=0,
 \qquad
 \Delta_{E,1}=2.
\]
Choosing the orientation after seeing the target would therefore manufacture
the desired residual.  The predicted--actual mixed pairing must be exported
from the original operational record.
\end{firewall}

\subsection{Explicit atomic formula and the modular defect}

The predictor becomes directly computable when the Euler algebra is an atomic
record algebra.  Let
\[
 \mathcal N=\operatorname{span}\{e_1,\ldots,e_J\},
 \qquad
 e_ae_b=\delta_{ab}e_a,
 \qquad
 \sum_ae_a=I,
 \tag{V5.22}
\]
where the \(e_a\) are the minimal projections of \(\mathcal N\).

\begin{proposition}[Atomic Petz score average]
\label{prop:v5-atomic-petz}
For the algebra in \textup{(V5.22)},
\[
 \boxed{
 \Pi_{\rho,\mathcal N}(x)
 =\sum_{a=1}^J
 \frac{\tau\!\left(
 e_a\rho^{1/2}x\rho^{1/2}\right)}
 {\tau(e_a\rho)}\,e_a.}
 \tag{V5.23}
\]
Moreover,
\[
 \boxed{
 \rho^{it}\mathcal N\rho^{-it}=\mathcal N
 \quad(t\in\R)
 \iff [\rho,e_a]=0\quad(1\le a\le J).}
 \tag{V5.24}
\]
In the exact branch, \textup{(V5.23)} fixes every \(e_a\) and is the
\(\varphi\)-preserving conditional expectation.  Outside that branch,
\textup{(V5.23)} is still a canonical finite weighted score average, so its
residual Gram can be measured without asserting modular invariance.
\end{proposition}

\begin{proof}
The trace expectation is
\[
 \mathcal E_{\mathcal N}(x)
 =\sum_a\frac{\tau(e_ax)}{\tau(e_a)}e_a,
\]
and
\[
 \rho_{\mathcal N}
 =\sum_a\frac{\tau(e_a\rho)}{\tau(e_a)}e_a.
\]
Substitution in \textup{(V5.4)} proves \textup{(V5.23)}.

If the modular group leaves \(\mathcal N\) invariant, then
\(t\mapsto\rho^{it}e_a\rho^{-it}\) is a continuous path among the finitely
many projections of the atomic abelian algebra having the same rank as
\(e_a\).  It is therefore constant, so \([\rho,e_a]=0\).  The converse is
immediate.  When the commutators vanish, direct substitution in
\textup{(V5.23)} shows that each \(e_a\) is fixed.
\end{proof}

There is also a strong but explicit quantitative comparison.  Put
\[
 \widetilde\rho:=\sum_ae_a\rho e_a,
 \qquad
 r_{\min}:=\min_a\frac{\tau(e_a\rho)}{\tau(e_a)}>0,
 \qquad
 \delta_{\rm mod}
 :=\norm{\rho^{1/2}-\widetilde\rho^{1/2}}.
 \tag{V5.25}
\]

\begin{proposition}[Atomic modular-predictor stability]
\label{prop:v5-modular-stability}
For every \(x\in\mathcal M\),
\[
 \boxed{
 \norm{
 \Pi_{\rho,\mathcal N}(x)
 -\Pi_{\widetilde\rho,\mathcal N}(x)}
 \le
 \frac{\delta_{\rm mod}}{r_{\min}}
 \bigl(\norm{\rho^{1/2}}
      +\norm{\widetilde\rho^{1/2}}\bigr)\norm x.}
 \tag{V5.26}
\]
The pinched density \(\widetilde\rho\) commutes with every \(e_a\), so
\(\Pi_{\widetilde\rho,\mathcal N}\) is an exact
\(\widetilde\varphi\)-preserving conditional expectation.  Estimate
\textup{(V5.26)} is a finite diagnostic, not a claim that
\(\delta_{\rm mod}\to0\) in the arithmetic cutoff.
\end{proposition}

\begin{proof}
The denominators in \textup{(V5.23)} are unchanged by pinching.  Write
\(b=\rho^{1/2}\) and \(\widetilde b=\widetilde\rho^{1/2}\).  Then
\[
 bxb-\widetilde b x\widetilde b
 =(b-\widetilde b)xb
 +\widetilde b x(b-\widetilde b).
\]
For every \(a\),
\[
 \abs{\tau(e_ay)}
 \le\tau(e_a)\norm y,
 \qquad
 \frac{\tau(e_a)}{\tau(e_a\rho)}
 \le r_{\min}^{-1}.
\]
Apply these inequalities coefficientwise in \textup{(V5.23)}.  The
commutation statement follows from the block-diagonal definition of
\(\widetilde\rho\) and \textup{(V5.24)}.
\end{proof}

\begin{remark}[Why density pinching is only a diagnostic]
Small \(\delta_{\rm mod}\) is stronger than approximate invariance of a few
selected Euler words, and the physical Euler synthesis need not equal the
pinched-state conditional expectation.  The primary bridge quantity remains
the actual residual \(\varepsilon\) in \textup{(V5.10)}.  Formula
\textup{(V5.26)} is useful only when the arithmetic loading independently
identifies the actual Euler row with, or close to, the atomic conditional
expectation row.
\end{remark}

\subsection{V5 status and next arithmetic computation}

\begin{center}
\begin{tabularx}{\textwidth}{@{}p{0.28\textwidth}p{0.17\textwidth}X@{}}
\toprule
V5 component & Status & Exact conclusion \\
\midrule
Always-defined Petz predictor & proved & Unital complete positivity,
state preservation, and complete \(L^2(\rho)\) contraction. \\
Residual upper bridge & proved & Defect bound
\((\norm H-1)_+\le2\varepsilon+\varepsilon^2\). \\
Operational residual Gram & proved & One predicted--actual mixed row gives
\(\Delta_E\) by \textup{(V5.19)}. \\
Atomic Euler record formula & proved & Explicit weighted score average
\textup{(V5.23)} and exact modular criterion \textup{(V5.24)}. \\
Quantitative atomic pinching & proved at finite cutoff & Predictor stability
\textup{(V5.26)}; no cofinal decay is asserted. \\
Arithmetic Petz residual & open & The actual V074 completed and Euler score
rows must be inserted into \textup{(V5.21)}. \\
Localized cell margin & open & The V3 Gaussian diagonal estimate remains
an alternative sufficient route. \\
Restoring lower bridge and descent & open & Independent of the Petz residual
criterion. \\
\bottomrule
\end{tabularx}
\end{center}

V5 supplies a quantitative version of the upstream Petz route without
assuming the missing modular theorem.  V6 should now inspect the actual finite
Euler algebra generated by V074's arithmetic record at one cutoff.  If it is
atomic, compute the projections \(e_a\), the weights
\(\tau(e_a\rho)\), the predicted score coefficients in \textup{(V5.23)}, and
the mixed Gram \(C_{TS}\).  The decisive number is the source-normalized
residual \textup{(V5.21)}, not the raw prime first moment.  If the Euler
algebra is nonabelian, the same computation should be performed blockwise
with \textup{(V5.4)}.  Failure of residual decay should be returned to the
translated or boundary restoring row rather than hidden by a target-chosen
orientation.

\begin{assessment}[V5 acquisition]
The exact Petz--Euler theorem in V074 required modular invariance and exact
mapping of completed words to Euler words.  V5 separates those two demands.
The canonical Petz predictor is always a complete physical-score contraction;
all failure of the upper metric bridge is confined to one measurable
predicted--actual residual.  Vanishing of its completed-normalized norm is a
new sufficient cofinal upper-bridge criterion.  Establishing that vanishing,
the restoring lower bridge, and the final descent remain open, so no assertion
of GRH is made.
\end{assessment}


% =============================================================================
\section{V6: the Euler-algebra closure audit and a finite acquisition card}
\label{sec:grh-v6}
% =============================================================================

\begin{center}
\large\bfseries Version V6 continuation
\end{center}

\subsection{What the arithmetic construction fixes, and what it does not}

V5 reduced the completed--Euler upper bridge to a Petz-predicted versus actual
Euler residual.  To evaluate that residual one needs a specific inclusion
\(\mathcal N_X^{\rm Eul}\subseteq\mathcal M_X^{\rm comp}\), a faithful
physical density, and the two score syntheses on the same nuisance-shorted
carrier.  The V074 arithmetic record section rigorously constructs instead the
following coefficient-free objects:
\[
 \mathcal R_X=\C^X,\quad
 C^*(S_X^{\rm rec})=M_X(\C),\quad
 \mathcal D_X:=\operatorname{span}\{P_{n,X}:1\le n\le X\},
 \tag{V6.1}
\]
together with the Peano histories \(L_{a,X}^{\rm P}\), the finite zeta
history, and the von Mangoldt writer
\[
 \mathcal V_X
 :=[H_X,\log\mathsf Z_X]
 =\sum_{2\le n\le X}\Lambda(n)L_{n,X}^{\rm P}.
 \tag{V6.2}
\]
It later declares completed, Euler, and nuisance score syntheses abstractly.
The symbols \(\mathcal N_X^{\rm Eul}\), \(\mathcal M_X^{\rm comp}\), and the
faithful density in the Petz theorem occur as hypotheses, not as outputs of
the finite Peano construction.

This is not a cosmetic missing definition.  Two natural meanings of the term
Euler algebra have opposite closures.

\subsection{The diagonal/full closure dichotomy}

\begin{theorem}[Prime-transition \(C^*\)-closure is the full record algebra]
\label{thm:v6-prime-closure}
Let
\[
 \mathcal U_X
 :=C^*\bigl(L_{p,X}^{\rm P}:p\le X,\ p\ {\rm prime}\bigr)
 \subseteq M_X(\C).
 \tag{V6.3}
\]
Then
\[
 \boxed{\mathcal U_X=M_X(\C).}
 \tag{V6.4}
\]
By contrast, \(\mathcal D_X\) in \textup{(V6.1)} is the atomic abelian
endpoint-read algebra.  Thus retaining only Euler coefficient Reads gives
\(\mathcal D_X\), whereas taking the \(C^*\)-closure of the actual Peano prime
transitions gives the entire matrix algebra.
\end{theorem}

\begin{proof}
For every prime \(p\le X\), the range projection
\[
 R_{p,X}
 :=L_{p,X}^{\rm P}(L_{p,X}^{\rm P})^*
 =\sum_{\substack{m\le X\\p\mid m}}P_{m,X}
 \tag{V6.5}
\]
belongs to \(\mathcal U_X\).  These diagonal projections commute.  Every
integer \(m\in\{2,\ldots,X\}\) has a prime divisor not exceeding \(X\), while
\(1\) has none.  Therefore
\[
 \boxed{
 P_{1,X}=\prod_{\substack{p\le X\\p\ {\rm prime}}}(I-R_{p,X})
 \in\mathcal U_X.}
 \tag{V6.6}
\]
If \(n=\prod_pp^{v_p(n)}\le X\), multiplicativity of the Peano histories gives
\[
 L_{n,X}^{\rm P}
 =\prod_p(L_{p,X}^{\rm P})^{v_p(n)}
 \in\mathcal U_X.
 \tag{V6.7}
\]
Since \(L_{n,X}^{\rm P}P_{1,X}=E_{n1}\), one obtains every matrix unit:
\[
 \boxed{
 E_{nm}
 =L_{n,X}^{\rm P}P_{1,X}(L_{m,X}^{\rm P})^*
 \in\mathcal U_X.}
 \tag{V6.8}
\]
The matrix units span \(M_X(\C)\), proving \textup{(V6.4)}.  The assertion
about \(\mathcal D_X\) follows directly from its mutually orthogonal minimal
projections.
\end{proof}

\begin{corollary}[Neither extreme gives the missing nontrivial Petz bridge]
\label{cor:v6-extreme-algebras}
Let \(\tau_X=X^{-1}\operatorname{Tr}\).
\begin{enumerate}[label=\textup{(\roman*)},leftmargin=2.8em]
\item For the tracial density \(\rho=I\) and
      \(\mathcal N=\mathcal D_X\), the Petz predictor is diagonal pinching and
\[
 \boxed{
 \Pi_{I,\mathcal D_X}(\mathcal V_X)=0.}
 \tag{V6.9}
\]
Nevertheless
\[
 \boxed{
 \mathcal V_X\eta_X^{\rm rec}
 =\sum_{2\le n\le X}\Lambda(n)e_n,\qquad
 \norm{\mathcal V_X\eta_X^{\rm rec}}^2
 =\sum_{2\le n\le X}\Lambda(n)^2.}
 \tag{V6.10}
\]
Thus endpoint pinching deletes every nontrivial prime transition.
\item For \(\mathcal N=\mathcal U_X=M_X(\C)\), the Petz predictor is the
      identity for every faithful density and modular invariance is automatic.
      It gives no contraction comparing two distinct completed and Euler score
      syntheses; the required mapping condition becomes their entire
      difference.
\end{enumerate}
Consequently V074's Petz theorem needs a separately constructed physical
Euler response algebra, not either closure inferred solely from the Peano
record.
\end{corollary}

\begin{proof}
For the tracial density, \(\Pi_{I,\mathcal D_X}\) is the trace-orthogonal
projection onto the diagonal.  Each \(L_{n,X}^{\rm P}\) with \(n\ge2\) is
strictly lower triangular in the endpoint ordering, so its diagonal is zero.
This proves \textup{(V6.9)}.  Applying the Peano writer to the anchor gives
\textup{(V6.10)}; orthonormality of the endpoint vectors gives its norm.
For the full algebra, the trace expectation is the identity and direct
substitution in the Petz formula gives \(\Pi_{\rho,M_X}=I\).
\end{proof}

\begin{remark}[The growing bare norm is not yet the completed metric]
The last sum in \textup{(V6.10)} diverges with \(X\), already on its prime
terms.  This is a statement in the canonical endpoint frame.  It cannot be
substituted for the completed-normalized residual of V5, because V074 has not
identified the completed score synthesis or its physical density with this
bare frame.
\end{remark}

\subsection{Endpoint probabilities do not determine the Petz predictor}

Even after selecting the atomic algebra, its classical record law does not
determine the ambient physical coherence.  A two-endpoint corner gives an
exact obstruction.

\begin{theorem}[Coherent-state ambiguity with identical endpoint law]
\label{thm:v6-state-ambiguity}
Let \(\mathcal M=M_2(\C)\), let
\(\mathcal N=\operatorname{span}\{E_{11},E_{22}\}\), and use
\(\tau_2=\frac12\operatorname{Tr}\).  Put
\[
 \sigma_x=
 \begin{pmatrix}0&1\\1&0\end{pmatrix},
 \qquad
 \sigma_z=
 \begin{pmatrix}1&0\\0&-1\end{pmatrix},
 \qquad
 \rho_c=I+c\sigma_x,
 \quad 0\le c<1.
 \tag{V6.11}
\]
Every \(\rho_c\) is a faithful density and all states
\(\varphi_c(x)=\tau_2(\rho_cx)\) have the same restriction to the endpoint
algebra:
\[
 \boxed{\varphi_c|_{\mathcal N}=\tau_2|_{\mathcal N}.}
 \tag{V6.12}
\]
Nevertheless
\[
 \boxed{
 \Pi_{\rho_c,\mathcal N}(\sigma_z)
 =\sqrt{1-c^2}\,\sigma_z.}
 \tag{V6.13}
\]
If both the completed and actual Euler one-column score syntheses equal
\(\sigma_z\), then their marginal metrics are
\[
 M_c=E_c=1
 \tag{V6.14}
\]
for every \(c\), while the Petz prediction parameters are
\[
 \boxed{
 q_c=\sqrt{1-c^2},\qquad
 \varepsilon_c=1-\sqrt{1-c^2}.}
 \tag{V6.15}
\]
Thus the Petz residual can range from \(0\) to arbitrarily near \(1\) while
the endpoint law and both marginal score metrics remain fixed.
\end{theorem}

\begin{proof}
The eigenvalues of \(\rho_c\) are \(1\pm c\), so it is faithful and
\(\tau_2(\rho_c)=1\).  Its diagonal equals \(I\), proving
\textup{(V6.12)} and
\[
 (\rho_c)_{\mathcal N}=I.
 \tag{V6.16}
\]
Write
\[
 \rho_c^{1/2}=a_cI+b_c\sigma_x,
 \quad
 a_c=\frac{\sqrt{1+c}+\sqrt{1-c}}2,
 \quad
 b_c=\frac{\sqrt{1+c}-\sqrt{1-c}}2.
 \tag{V6.17}
\]
Since \(\sigma_x\sigma_z+\sigma_z\sigma_x=0\) and
\(\sigma_x\sigma_z\sigma_x=-\sigma_z\),
\[
 \rho_c^{1/2}\sigma_z\rho_c^{1/2}
 =(a_c^2-b_c^2)\sigma_z
 =\sqrt{1-c^2}\,\sigma_z.
\]
This matrix is diagonal, so the trace expectation leaves it unchanged and
\textup{(V6.13)} follows from \textup{(V5.4)}.  Finally
\[
 \norm{\sigma_z}_{2,\rho_c}^2
 =\varphi_c(I)=1.
\]
The predicted score is \(\sqrt{1-c^2}\sigma_z\) and the actual score is
\(\sigma_z\), giving \textup{(V6.14)}--\textup{(V6.15)}.
\end{proof}

\begin{remark}[The V5 estimate sees the exact compensation]
In this example
\[
 q_c+\varepsilon_c=1,
\]
so the sharp estimate \(\norm H\le(q+\varepsilon)^2\) is an equality even
though the residual varies.  Residual decay is a sufficient cofinal
criterion, not a necessary condition for a bridge already supplied by another
mechanism.  The example proves underdetermination of the Petz decomposition,
not failure of the metric order.
\end{remark}

\begin{corollary}[The V074 record tuple does not fix the Petz datum]
\label{cor:v6-no-canonical-petz}
The pointed successor, endpoint projections, count operator, Peano histories,
and endpoint probability law do not determine a unique Petz predictor.  The
two-dimensional construction embeds as a corner of every \(M_X(\C)\) with
\(X\ge2\).  Hence a physical density or an equivalent complete state row must
be retained in addition to the algebraic arithmetic loading.
\end{corollary}

\begin{proof}
All operators in the two-endpoint record tuple are held fixed in
\Cref{thm:v6-state-ambiguity}, and the restrictions of the states to the
endpoint Reads agree.  Yet their Petz images of the same score differ.
Taking the direct sum with any fixed faithful density on the complementary
corner gives the assertion for larger \(X\).
\end{proof}

\subsection{A basis-independent finite acquisition algorithm}

The missing datum is finite and testable; it should not be replaced by a
qualitative assertion that an Euler algebra exists.  The next theorem gives
an exact acquisition card.

\begin{theorem}[Finite Petz residual algorithm]
\label{thm:v6-petz-algorithm}
At one cutoff, suppose the following matrices are supplied on a common
finite carrier:
\begin{enumerate}[label=\textup{(\alph*)},leftmargin=2.8em]
\item a faithful density \(\rho\) and faithful trace \(\tau\);
\item a linearly independent basis
      \(n_1,\ldots,n_d\) of a declared unital \(C^*\)-subalgebra
      \(\mathcal N^{\rm Eul}\);
\item completed, actual Euler, and common nuisance score syntheses
      \(A,S,N\).
\end{enumerate}
Let
\[
 G_{ab}:=\tau(n_a^*n_b),
 \qquad
 b_a(x):=\tau(n_a^*x).
 \tag{V6.18}
\]
Then the trace expectation is exactly
\[
 \boxed{
 \mathcal E_{\mathcal N}(x)
 =\sum_{a,b=1}^d
 n_b(G^{-1})_{ba}b_a(x).}
 \tag{V6.19}
\]
Compute successively
\[
 \rho_{\mathcal N}=\mathcal E_{\mathcal N}(\rho),
 \qquad
 \Pi_{\rho,\mathcal N}\ \hbox{from \textup{(V5.4)}},
 \tag{V6.20}
\]
the common nuisance projection
\[
 \mathsf Q=I-N(N^*N)^\dagger N^*,
 \tag{V6.21}
\]
and
\[
 \widehat A=\mathsf QA,\quad
 \widehat S=\mathsf QS,\quad
 T=\mathsf Q\Pi_{\rho,\mathcal N}\widehat A.
 \tag{V6.22}
\]
With
\[
 M=\widehat A^*\widehat A,\quad
 E=\widehat S^*\widehat S,\quad
 P=T^*T,\quad
 C_{TS}=T^*\widehat S,
 \tag{V6.23}
\]
the exact normalized residual is
\[
 \boxed{
 \varepsilon^2
 =\lambda_{\max}\!\left(
 M^{\dagger/2}
 (E+P-C_{TS}-C_{TS}^*)
 M^{\dagger/2}\right)}
 \tag{V6.24}
\]
on \(\supp M\), together with the independent kernel test
\(\Ker M\subseteq\Ker E\).

Every quantity in \textup{(V6.18)}--\textup{(V6.24)} is invariant under a
change of basis of \(\mathcal N^{\rm Eul}\) and under a common unitary change
of physical carrier.  Therefore a supplied acquisition card produces one
unambiguous finite upper-bridge bound via V5.
\end{theorem}

\begin{proof}
The right side of \textup{(V6.19)} is the orthogonal projection in the
Hilbert space \((\mathcal M,\ip{\cdot}{\cdot}_{\tau})\) onto the span of the
\(n_a\).  Indeed, its coefficient vector solves the Gram normal equations
\(Gc=b(x)\).  A unital \(C^*\)-subalgebra has a unique trace-preserving
conditional expectation, so this orthogonal projection is
\(\mathcal E_{\mathcal N}\).  The remaining formulas are the finite Petz map,
the Moore--Penrose nuisance short, and the residual Gram identity
\textup{(V5.19)}.  The largest eigenvalue in \textup{(V6.24)} is exactly the
square of \(\norm{R_EM^{\dagger/2}}\).  Orthogonal projections and Gram
congruences do not depend on the chosen algebra basis or physical coordinates.
\end{proof}

\begin{firewall}[The acquisition card cannot be filled by notation]
\label{fire:v6-acquisition}
Writing symbols \(A_X,S_X,N_X,\rho_X\) does not populate
\textup{(V6.18)}--\textup{(V6.24)}.  In particular:
\begin{enumerate}[label=\textup{(\roman*)},leftmargin=2.8em]
\item the canonical endpoint Gram \(I_X\) is not a faithful physical density
      on the completed score algebra;
\item the diagonal coefficient algebra omits the Peano prime transitions;
\item the \(C^*\)-closure of all those transitions is already \(M_X(\C)\);
\item separate completed and Euler marginal Grams omit \(C_{TS}\);
\item choosing an algebra, density, or mixed row after observing the GRH
      target violates the source-fixed requirement.
\end{enumerate}
Accordingly the V5 residual is a rigorous conditional certificate, but V074
does not presently supply its numerical or asymptotic input.
\end{firewall}

\subsection{V6 status and return to an explicitly populated chart}

\begin{center}
\begin{tabularx}{\textwidth}{@{}p{0.28\textwidth}p{0.17\textwidth}X@{}}
\toprule
V6 component & Status & Exact conclusion \\
\midrule
Euler algebra audit & proved & Diagonal Reads give \(\mathcal D_X\), while
prime-transition \(C^*\)-closure gives \(M_X(\C)\). \\
Diagonal Petz port & disproved as a transition-preserving route &
\(\Pi_{I,\mathcal D_X}(\mathcal V_X)=0\). \\
Full-algebra Petz port & vacuous & The predictor is the identity and the
mapping condition retains the whole completed--Euler difference. \\
Physical state determination & disproved from endpoint data & The family
\(\rho_c\) has one endpoint law but different Petz predictors. \\
Finite acquisition card & proved & Once algebra, density, scores, nuisance,
and mixed row are supplied, \(\varepsilon\) is given by
\textup{(V6.24)}. \\
Arithmetic Petz residual & unpopulated & V074 does not specify the input card
needed for a cofinal estimate. \\
Restoring bridge and GRH descent & open & No change in logical status. \\
\bottomrule
\end{tabularx}
\end{center}

The Petz route has therefore reached a source-data boundary, not an analytic
estimate.  Continuing to manipulate an unspecified \(\rho_X\) would be
circular.  V7 should return to the arithmetic manuscript's localized
Weil--Schur chart, where every operator is explicit.  The next concrete
calculation is the complete Legendre matrix of one truncated prime shift:
for \(0\le\delta=\log n/A<2\),
\[
 \ip{e_m}{T_{\delta}e_k}
 =\frac{\sqrt{(2m+1)(2k+1)}}2
 \int_{-1}^{1-\delta}P_m(t)P_k(t+\delta)\dd t.
 \tag{V6.25}
\]
Computing these entries and the archimedean low-mode block gives a genuinely
populated finite Schur obstruction.  A negative low-mode witness can then be
routed back to the translated or boundary restoring row; no invented state or
Euler algebra is required.

\begin{assessment}[V6 acquisition]
V6 determines why the promising V5 certificate cannot yet be evaluated from
V074.  The unconditional Peano record does not select a nontrivial Euler
\(C^*\)-subalgebra, and its endpoint law does not select the faithful ambient
state.  These are independent mathematical choices, as shown by exact closure
and two-state counterexamples.  The note supplies a complete finite
acquisition algorithm for any future physical panel and redirects the active
attack to the fully explicit localized Weil matrix.  GRH is not proved.
\end{assessment}


% =============================================================================
\section{V7: exact Legendre acquisition and a certified Schur bracket}
\label{sec:grh-v7}
% =============================================================================

V6 returned the attack to the localized Weil chart because that chart is
populated entirely by the source: the support wall, conductor, parity, and
active prime powers determine every operator.  This section carries out the
announced first calculation.  It gives the complete matrix of a truncated
translation in the Legendre basis, exposes an exact parity splitting of every
prime cell, reduces the smooth archimedean block to moments of the same
translation polynomials, and replaces the formal tail resolvent by a
genuinely finite certified bracket.

A targeted nonduplication search found the harmonic diagonal and the abstract
Schur complement in the arithmetic manuscript, and the ambient shift audit in
V4, but no formula for the Legendre translation matrix, no parity law for the
complex-character cell, and no finite-data enclosure for the tail resolvent.
The results below therefore address the live V6 boundary rather than restating
the manuscript's finite-history reduction.

\subsection{The exact truncated-translation matrix}

On \(L^2(-1,1)\), extend functions by zero and define, for
\(0\le\delta\le2\),
\[
 (T_\delta h)(t)
 := \mathbf 1_{(-1,\,1-\delta)}(t)h(t+\delta).
 \tag{V7.1}
\]
This is the scaled version of \(T_{A,u}\) with
\(\delta=u/A\).  Recall
\[
 e_j(t)=\sqrt{\frac{2j+1}{2}}P_j(t),
 \qquad
 b_{mk}(\delta):=\ip{e_m}{T_\delta e_k},
 \qquad
 B_N(\delta):=(b_{mk}(\delta))_{0\le m,k<N}.
 \tag{V7.2}
\]

Write the monomial coefficients of the Legendre polynomial as
\[
 P_j(t)=\sum_{a=0}^{\lfloor j/2\rfloor}c_{j,a}t^{j-2a},
 \qquad
 c_{j,a}
 =\frac{(-1)^a(2j-2a)!}
 {2^j a!(j-a)!(j-2a)!}.
 \tag{V7.3}
\]

\begin{theorem}[Complete polynomial acquisition of one prime shift]
\label{thm:v7-translation-matrix}
For \(0\le\delta\le2\), put \(r=m-2a\) and \(s=k-2b\).  Then
\[
 \boxed{
 \begin{aligned}
 b_{mk}(\delta)
 ={}&\frac{\sqrt{(2m+1)(2k+1)}}2
 \sum_{a=0}^{\lfloor m/2\rfloor}
 \sum_{b=0}^{\lfloor k/2\rfloor}c_{m,a}c_{k,b}\\
 &\times\sum_{\ell=0}^{s}\binom{s}{\ell}\delta^{s-\ell}
 \frac{(1-\delta)^{r+\ell+1}
       -(-1)^{r+\ell+1}}
      {r+\ell+1}.
 \end{aligned}}
 \tag{V7.4}
\]
In particular, every entry is an explicit polynomial with rational
coefficients times \(\sqrt{(2m+1)(2k+1)}\).  It satisfies
\[
 \boxed{
 B_N(0)=I_N,
 \qquad B_N(2)=0,
 \qquad
 b_{km}(\delta)=(-1)^{m+k}b_{mk}(\delta).}
 \tag{V7.5}
\]
For each fixed \(N\),
\[
 \boxed{
 b_{mk}(\delta)
 =(-1)^m\frac{\sqrt{(2m+1)(2k+1)}}2(2-\delta)
 +O_{m,k}((2-\delta)^2)}
 \tag{V7.6}
\]
as \(\delta\uparrow2\), and hence
\(
 \norm{B_N(\delta)}=O_N(2-\delta)
\).
\end{theorem}

\begin{proof}
Expanding \(P_m(t)\) and \(P_k(t+\delta)\) by \textup{(V7.3)} gives
\[
 P_k(t+\delta)
 =\sum_b c_{k,b}\sum_{\ell=0}^{s}
   \binom{s}{\ell}t^\ell\delta^{s-\ell}.
\]
Termwise integration over \((-1,1-\delta)\) proves
\textup{(V7.4)}.  At \(\delta=0\), Legendre orthonormality gives the first
identity in \textup{(V7.5)}, while at \(\delta=2\) the overlap has zero
measure.

Let \(Rh(t)=h(-t)\) be reflection.  Then
\[
 T_\delta^*=T_{-\delta}=RT_\delta R,
 \qquad Re_j=(-1)^je_j.
 \tag{V7.7}
\]
Taking matrix entries gives the parity identity in \textup{(V7.5)}.  Finally,
write \(h=2-\delta\) and \(t=-1+y\) on the overlap \(0<y<h\).  Uniformly there,
\[
 e_m(-1+y)e_k(1-h+y)
 =(-1)^m\frac{\sqrt{(2m+1)(2k+1)}}2+O_{m,k}(h).
\]
Integration over an interval of length \(h\) proves \textup{(V7.6)}.
\end{proof}

\begin{corollary}[The first protected block]
\label{cor:v7-first-B-block}
The first three Legendre modes give
\[
 \boxed{
 B_3(\delta)=
 \begin{pmatrix}
  1-\frac\delta2
  &\frac{\sqrt3}{4}\delta(2-\delta)
  &-\frac{\sqrt5}{4}\delta(\delta-1)(\delta-2)\\[2mm]
  -\frac{\sqrt3}{4}\delta(2-\delta)
  &1-\frac32\delta+\frac14\delta^3
  &\frac{\sqrt{15}}2
    \left(\delta-\delta^2+\frac18\delta^4\right)\\[2mm]
  -\frac{\sqrt5}{4}\delta(\delta-1)(\delta-2)
  &-\frac{\sqrt{15}}2
    \left(\delta-\delta^2+\frac18\delta^4\right)
  &1-\frac52\delta+\frac54\delta^3-\frac3{16}\delta^5
 \end{pmatrix}.}
 \tag{V7.8}
\]
\end{corollary}

\begin{proof}
Insert \(P_0=1\), \(P_1=t\), and
\(P_2=(3t^2-1)/2\) into \textup{(V7.2)}, or specialize
\textup{(V7.4)}.  The skew signs in the odd-parity entries provide an
independent check through \textup{(V7.5)}.
\end{proof}

\begin{remark}[Low modes remove the ambient threshold jump]
V4 recalled that \(\norm{T_\delta}=1\) for every \(0\le\delta<2\), followed
by a jump to zero at \(\delta=2\).  Formula \textup{(V7.6)} shows that this
jump lives entirely in modes escaping every fixed Legendre window.  Each
source-fixed protected block instead activates continuously and linearly in
the overlap length.  This distinction is essential for any incremental
support renewal.
\end{remark}

\subsection{Prime cells and the exact parity split}

Let
\[
 J_N:=\diag(1,-1,1,-1,\ldots)
 \tag{V7.9}
\]
on the first \(N\) Legendre modes.  For an active prime power \(n\), put
\[
 \delta_n=\frac{\log n}{A}\in(0,2),
 \qquad
 \mathcal P_{\chi,N}(n)
 :=P_N\mathsf P_{\chi,A}^{\rm sc}(n)P_N.
 \tag{V7.10}
\]

\begin{theorem}[Exact complex-character prime block]
\label{thm:v7-prime-parity}
One has
\[
 \boxed{
 \mathcal P_{\chi,N}(n)
 =\chi(n)B_N(\delta_n)
  +\overline{\chi(n)}B_N(\delta_n)^{\mathsf T}
 =\chi(n)B_N(\delta_n)
  +\overline{\chi(n)}J_NB_N(\delta_n)J_N.}
 \tag{V7.11}
\]
Entrywise,
\[
 \boxed{
 (\mathcal P_{\chi,N}(n))_{mk}
 =\begin{cases}
  2\operatorname{Re}\chi(n)\,b_{mk}(\delta_n),&m+k\text{ even},\\
  2\mathrm i\operatorname{Im}\chi(n)\,b_{mk}(\delta_n),&m+k\text{ odd}.
 \end{cases}}
 \tag{V7.12}
\]
Thus a real character preserves the even and odd Legendre sectors.  For a
complex character the diagonal unitary which multiplies every odd mode by
\(\mathrm i\) converts the whole prime block to a real symmetric matrix.
Moreover, for fixed \(N\),
\[
 \norm{\mathcal P_{\chi,N}(n)}
 =O_N(2-\delta_n)
 \tag{V7.13}
\]
at the activation wall.
\end{theorem}

\begin{proof}
The negative shift is the adjoint of the positive shift, so its real
Legendre matrix is \(B_N(\delta_n)^{\mathsf T}\).  Apply the parity identity
\textup{(V7.5)} to obtain \textup{(V7.11)}--\textup{(V7.12)}.  The diagonal
unitary changes a purely imaginary even--odd entry into a real entry and
preserves the Hermitian symmetry.  The final estimate follows from
\textup{(V7.6)}.
\end{proof}

\subsection{The archimedean low block from the same polynomials}

The universal wall multiplier in the scaled archimedean operator is
\[
 (\mathsf Wh)(t)=-\frac12\log(1-t^2)h(t).
\]
Let \(W_{mk}=\ip{e_m}{\mathsf We_k}\).  For \(q\ge0\), define
\[
 \mathfrak j_q
 :=\int_{-1}^1t^{2q}\log(1-t^2)\dd t
 =\frac2{2q+1}
 \left(2\log2-2H_{2q+2}+H_{q+1}\right).
 \tag{V7.14}
\]

\begin{theorem}[Exact logarithmic-wall matrix]
\label{thm:v7-wall-matrix}
The wall matrix respects parity.  If \(m+k\) is odd, then \(W_{mk}=0\).
If \(m<k\) and \(m+k\) is even, then
\[
 \boxed{
 W_{mk}=W_{km}
 =\frac{\sqrt{(2m+1)(2k+1)}}
 {(k-m)(k+m+1)}.}
 \tag{V7.15}
\]
The diagonal is the finite exact sum
\[
 \boxed{
 W_{mm}
 =-\frac{2m+1}{4}
 \sum_{a,b=0}^{\lfloor m/2\rfloor}
 c_{m,a}c_{m,b}\mathfrak j_{m-a-b}.}
 \tag{V7.16}
\]
In particular,
\[
 \boxed{
 (W_{mk})_{0\le m,k<3}
 =\begin{pmatrix}
  1-\log2&0&\frac{\sqrt5}{6}\\
  0&\frac43-\log2&0\\
  \frac{\sqrt5}{6}&0&\frac{41}{30}-\log2
 \end{pmatrix}.}
 \tag{V7.17}
\]
\end{theorem}

\begin{proof}
The parity assertion is immediate.  Formula \textup{(V7.14)} follows from
the substitution \(y=t^2\) and the beta-function derivative, or directly by
integrating its elementary power series.  Expanding \(P_m^2\) by
\textup{(V7.3)} then proves \textup{(V7.16)}.

For the off-diagonal formula, put
\(I_{mk}=\int_{-1}^1P_mP_k\log(1-t^2)\dd t\).  Subtract the two Legendre
Sturm--Liouville equations after multiplying by the opposite polynomial and
by \(\log(1-t^2)\).  The endpoint term vanishes and, for \(m<k\) of the same
parity, the identities \(P_n'=\sum_{j<n,\ n-j\ \text{odd}}(2j+1)P_j\) and
\(tP_j=((j+1)P_{j+1}+jP_{j-1})/(2j+1)\), together with
orthogonality, give
\[
 \int_{-1}^1t(P_mP_k'-P_kP_m')\dd t=2.
\]
Hence
\[
 [k(k+1)-m(m+1)]I_{mk}=-4,
\]
so
\[
 I_{mk}=-\frac4{(k-m)(k+m+1)}.
\]
Multiplication by the normalization in \textup{(V7.2)} and by \(-1/2\)
gives \textup{(V7.15)}.  The first three diagonal values from
\textup{(V7.16)} yield \textup{(V7.17)}.
\end{proof}

The parity-dependent smooth kernel is
\[
 s_\alpha(y)=\frac1{2y}
 -\frac{e^{-2\alpha y}}{1-e^{-2y}},
 \qquad
 s_\alpha(0):=\alpha-\frac12,
 \qquad
 \alpha=\frac{1+2\kappa}{4}.
 \tag{V7.18}
\]
Write
\[
 S_{mk}^{(\alpha,A)}
 :=\ip{e_m}{\mathsf S_{\alpha,A}^{\rm sc}e_k}.
\]

\begin{theorem}[One-dimensional acquisition of the smooth gamma block]
\label{thm:v7-smooth-block}
The smooth block is determined by the same translation polynomials:
\[
 \boxed{
 S_{mk}^{(\alpha,A)}
 =A\int_0^2s_\alpha(Ar)
   \bigl(b_{mk}(r)+b_{km}(r)\bigr)\dd r
 =A\bigl(1+(-1)^{m+k}\bigr)
   \int_0^2s_\alpha(Ar)b_{mk}(r)\dd r.}
 \tag{V7.19}
\]
Consequently \(S_{mk}^{(\alpha,A)}=0\) for \(m+k\) odd.  If
\[
 b_{mk}(r)=\sum_{q=0}^{m+k+1}\beta_{mkq}r^q,
 \qquad
 \mu_{\alpha,A,q}:=A\int_0^2r^qs_\alpha(Ar)\dd r,
 \tag{V7.20}
\]
then
\[
 \boxed{
 S_{mk}^{(\alpha,A)}
 =\bigl(1+(-1)^{m+k}\bigr)
  \sum_q\beta_{mkq}\mu_{\alpha,A,q}.}
 \tag{V7.21}
\]
Thus the complete \(N\)-mode smooth block requires only the \(2N\) scalar
moments with \(0\le q\le2N-1\) of an explicit continuous function.
\end{theorem}

\begin{proof}
Split the square \((-1,1)^2\) into \(v=t+r\) and \(t=v+r\), with
\(0<r<2\).  The two inner integrals are exactly \(b_{mk}(r)\) and
\(b_{km}(r)\).  The scaled kernel contributes \(A s_\alpha(Ar)\), proving
\textup{(V7.19)}.  Apply \textup{(V7.5)} and then expand the polynomial from
\textup{(V7.4)} to obtain \textup{(V7.21)}.
\end{proof}

\begin{corollary}[The populated protected matrix]
\label{cor:v7-populated-low-matrix}
Let \(n<e^{2A}\) range over the active unramified prime powers and put
\(
 \lambda_n=\Lambda(n)/\sqrt n
\).  The entire low block of the completed localized Weil operator is the
explicit Hermitian matrix
\[
 \boxed{
 \begin{aligned}
 L_{\chi,A,N}
 :={}&P_N\mathsf H_{\chi,A}^{\rm sc}P_N\\
 ={}&\diag(H_0,\ldots,H_{N-1})+c_{f,A}I_N
      +(W_{mk})_{m,k<N}+(S_{mk}^{(\alpha,A)})_{m,k<N}\\
 &-\sum_{\substack{2\le n<e^{2A}\\ n\ {\rm unramified}}}
   \lambda_n\mathcal P_{\chi,N}(n),
 \qquad
 c_{f,A}=\log\frac f{2\pi A}-\gamma.
 \end{aligned}}
 \tag{V7.22}
\]
No unspecified state, conditional expectation, or mixed Gram row occurs in
\textup{(V7.22)}.
\end{corollary}

\subsection{A genuinely finite Schur certificate}

There is an important distinction between a Schur complement whose
\emph{range} is finite dimensional and one whose entries are obtainable from
finite source data.  The matrix in the manuscript's \textup{(W.28)} contains
\[
 (Q_N\mathsf H Q_N)^{-1},
\]
an inverse on an infinite-dimensional tail.  Thus the low block
\textup{(V7.22)} alone does not populate that Schur matrix.  The next result
replaces the inverse by a rigorous finite enclosure.

Put
\[
 \mathsf V_{\chi,A}
 :=\mathsf W+\mathsf S_{\alpha,A}^{\rm sc}
 -\sum_n\lambda_n\mathsf P_{\chi,A}^{\rm sc}(n),
 \qquad
 \mathsf H=\mathsf D+c_{f,A}I+\mathsf V_{\chi,A}.
 \tag{V7.23}
\]

\begin{theorem}[Finite Schur bracket]
\label{thm:v7-finite-schur-bracket}
Assume the manuscript's tail estimate
\[
 Q_N\mathsf H Q_N\succeq d_NQ_N,
 \qquad d_N>0.
 \tag{V7.24}
\]
Define the finite positive matrix
\[
 \boxed{
 R_{\chi,A,N}
 :=P_N\mathsf V_{\chi,A}Q_N
       \mathsf V_{\chi,A}P_N.}
 \tag{V7.25}
\]
Then the exact Schur complement \(F_{\chi,A,N}\) obeys
\[
 \boxed{
 L_{\chi,A,N}-d_N^{-1}R_{\chi,A,N}
 \preceq F_{\chi,A,N}
 \preceq L_{\chi,A,N}.}
 \tag{V7.26}
\]
Accordingly:
\begin{enumerate}[label=\textup{(\roman*)},leftmargin=2.8em]
\item if
\(
 L_{\chi,A,N}-d_N^{-1}R_{\chi,A,N}\succeq0
\), then the full localized Weil operator is nonnegative;
\item if \(L_{\chi,A,N}\not\succeq0\), then the full operator is not
nonnegative and a negative eigenvector gives an explicit polynomial Weil
witness;
\item if the two finite endpoints straddle zero, this certificate is
inconclusive and the tail resolvent cannot be omitted.
\end{enumerate}
\end{theorem}

\begin{proof}
Because \(\mathsf D+c_{f,A}I\) is diagonal in the Legendre basis,
\[
 P_N\mathsf H Q_N=P_N\mathsf V_{\chi,A}Q_N.
 \tag{V7.27}
\]
The tail lower bound gives
\[
 0\preceq(Q_N\mathsf H Q_N)^{-1}\preceq d_N^{-1}Q_N.
\]
Insert this order interval into the Feshbach formula
\[
 F_{\chi,A,N}
 =L_{\chi,A,N}
 -P_N\mathsf H Q_N(Q_N\mathsf H Q_N)^{-1}
  Q_N\mathsf H P_N
\]
and use \textup{(V7.27)}.  This proves \textup{(V7.26)}.  The first claim
follows from the Schur criterion, and the second follows by testing
\(\mathsf H\) directly on the negative low-mode vector.
\end{proof}

\begin{proposition}[The correction matrix is source computable]
\label{prop:v7-R-acquisition}
For \(0\le k<N\), let
\[
 v_k:=\mathsf V_{\chi,A}e_k.
 \tag{V7.28}
\]
Then
\[
 \boxed{
 (R_{\chi,A,N})_{mk}
 :=\ip{R_{\chi,A,N}e_k}{e_m}
 =\ip{Q_Nv_k}{Q_Nv_m}
 =\ip{v_k}{v_m}
  -\sum_{\ell=0}^{N-1}
    \ip{v_k}{e_\ell}\ip{e_\ell}{v_m}.}
 \tag{V7.29}
\]
Every function in this formula is explicit:
\[
 \boxed{
 \begin{aligned}
 v_k(t)={}&-\frac12\log(1-t^2)e_k(t)
 +A\int_{-1}^1s_\alpha(A|t-u|)e_k(u)\dd u\\
 &-\sum_n\lambda_n\left[
 \chi(n)\mathbf 1_{t<1-\delta_n}e_k(t+\delta_n)
 +\overline{\chi(n)}\mathbf 1_{t>-1+\delta_n}e_k(t-\delta_n)
 \right].
 \end{aligned}}
 \tag{V7.30}
\]
Thus \(R_{\chi,A,N}\) is obtained from finitely many integrals of elementary
polynomials, logarithms, and the explicit smooth function
\(s_\alpha\).  In particular, \textup{(V7.26)} is a finite-data certificate,
unlike the unevaluated tail inverse.
\end{proposition}

\begin{proof}
Equation \textup{(V7.29)} is Pythagoras after subtracting the first \(N\)
Legendre coefficients.  The logarithmic factor is square integrable against
each polynomial, the smooth kernel is bounded on the compact square, and
there are finitely many active shifts.  Hence \(v_k\in L^2(-1,1)\), and
substitution of \textup{(V7.23)} gives \textup{(V7.30)}.
\end{proof}

\begin{firewall}[What a low-mode computation can certify]
\label{fire:v7-low-mode}
A negative eigenvalue of \(L_{\chi,A,N}\) is a genuine negative Weil witness,
not a defect that may be repaired by choosing a state, basis, or nuisance
row.  A nonnegative \(L_{\chi,A,N}\), however, is not by itself a positivity
proof because tail feedback subtracts a positive matrix.  The valid positive
certificate is the lower endpoint in \textup{(V7.26)}.  Any numerical scan
must therefore report both \(L_{\chi,A,N}\) and
\(R_{\chi,A,N}/d_N\), with rigorous integration and eigenvalue error bounds.
\end{firewall}

\subsection{The cost of the absolute-prime-mass tail bound}

Let
\[
 \mathcal L_f(A)
 :=\sum_{\substack{2\le n<e^{2A}\\ n\ {\rm unramified}}}
 \frac{\Lambda(n)}{\sqrt n}.
 \tag{V7.31}
\]
The manuscript's available coercivity constant is
\[
 d_N=H_N+c_{f,A}
 -\norm{\mathsf S_{\alpha,A}^{\rm sc}}
 -2\mathcal L_f(A).
 \tag{V7.32}
\]

\begin{theorem}[Double-exponential dimension forced by the present bound]
\label{thm:v7-dimension-cost}
Any \(N\) certified by \textup{(V7.32)} with \(d_N>0\) necessarily satisfies
\[
 \boxed{
 \log N>2\mathcal L_f(A)-c_{f,A}-1.}
 \tag{V7.33}
\]
By the prime number theorem and removal of the finitely many ramified prime
powers,
\[
 \mathcal L_f(A)=2e^A+o(e^A),
 \tag{V7.34}
\]
so along \(A\to\infty\),
\[
 \boxed{
 N\ge\exp\!\left((4+o(1))e^A\right).}
 \tag{V7.35}
\]
This is a limitation of the absolute-prime-mass certificate
\textup{(V7.32)}, not a lower bound for every possible high-mode coercivity
argument.
\end{theorem}

\begin{proof}
If \(d_N>0\), then
\[
 H_N>2\mathcal L_f(A)-c_{f,A}
      +\norm{\mathsf S_{\alpha,A}^{\rm sc}}
 \ge2\mathcal L_f(A)-c_{f,A}.
\]
Since \(H_N\le1+\log N\), this proves \textup{(V7.33)}.  Partial summation
of the prime number theorem gives
\[
 \sum_{n\le x}\frac{\Lambda(n)}{\sqrt n}
 =2\sqrt x+o(\sqrt x).
\]
Take \(x=e^{2A}\); deletion of prime powers over the finitely many primes
dividing \(f\) contributes only \(O_f(A)\).  Since
\(-c_{f,A}=\log A+O_f(1)\), equations
\textup{(V7.34)}--\textup{(V7.35)} follow.
\end{proof}

\begin{remark}[Where the next estimate must improve]
The exact low block sees a newly activated shift with size
\(O_N(2-\delta_n)\), whereas \textup{(V7.32)} charges it the full ambient
norm \(2\lambda_n\).  The double-exponential cost is precisely the price of
discarding this support geometry and all character phase information.  A
cofinal argument must replace the absolute mass by a multiscale tail estimate
that distinguishes boundary activations from common-interior oscillation.
\end{remark}

\subsection{V7 status and the next upstream target}

\begin{center}
\begin{tabularx}{\textwidth}{@{}p{0.28\textwidth}p{0.17\textwidth}X@{}}
\toprule
V7 component & Status & Exact conclusion \\
\midrule
One-shift Legendre matrix & proved & Every entry is the polynomial
\textup{(V7.4)}, with reflection parity and linear activation law. \\
Complex-character cell & proved & Formula \textup{(V7.12)} separates
same-parity real parts from cross-parity imaginary parts. \\
Archimedean low block & populated & The logarithmic wall is given by
\textup{(V7.15)}--\textup{(V7.17)} and the smooth gamma block by finitely
many moments in \textup{(V7.21)}. \\
Finite Schur acquisition & repaired & The exact tail inverse is enclosed by
the source-computable bracket \textup{(V7.26)}. \\
Absolute-mass high tail & quantitatively inadequate cofinally & Its certified
dimension is at least double exponential in the support wall. \\
Fixed-wall positivity & two-sided finite test & The lower bracket proves
positivity when nonnegative; a negative low block gives an explicit
counter-witness; intermediate cases retain the tail problem. \\
Cofinal positivity and GRH & open & No uniform multiscale tail estimate has
yet replaced the absolute prime mass. \\
\bottomrule
\end{tabularx}
\end{center}

V8 should go upstream from the double-exponential absolute-mass bound.  The
next source-native object is the high-frequency symbol of a dyadic band of
compressed shifts, separated into:
\begin{enumerate}[label=\textup{(\alph*)},leftmargin=2.8em]
\item a boundary layer, controlled by the overlap length and the exact
activation polynomial \textup{(V7.6)}; and
\item a common-interior layer, where modulated localized packets convert the
prime cells into the Dirichlet polynomial
\[
 \sum_{n\in\mathcal B}\frac{\Lambda(n)\chi(n)}{\sqrt n}
 e^{\mathrm i\xi\log n/A}.
 \tag{V7.36}
\]
\end{enumerate}
The required theorem is a block estimate strong enough to be absorbed by the
harmonic diagonal without assuming the desired zero-free information.  If
such an estimate fails, its maximizing packet is the actual translated or
boundary restoring row demanded by the source-fixed firewall.

\begin{assessment}[V7 acquisition]
V7 turns the localized Weil reduction into explicit finite arithmetic.  The
prime, logarithmic-wall, and smooth gamma low blocks are all acquired, and
the infinite tail resolvent is replaced by a rigorous finite Schur bracket.
The calculation also isolates the cofinal obstruction: the existing high-mode
bound throws away activation geometry and forces a protected dimension of
order at least \(\exp((4+o(1))e^A)\).  The next attack must therefore be
multiscale and phase sensitive.  GRH is not proved.
\end{assessment}


% =============================================================================
\section{V8: prime-log resonance, the archimedean clock, and a uniform-tail no-go}
\label{sec:grh-v8}
% =============================================================================

V7 proposed replacing absolute prime mass by a phase-sensitive estimate on
high-frequency packets.  That proposal must first survive a simultaneous
resonance test.  It does not survive in uniform operator norm: at every fixed
support wall, arbitrarily high modulations align all unramified prime phases
with the character.  This section proves that no fixed Legendre tail gains
uniform character cancellation.  It also proves the complementary fact that
the same resonant peaks are extremely atypical in frequency mean square.

The resulting picture is more precise than either an absolute-mass estimate
or an average cancellation slogan.  The completed archimedean terms produce
an exact logarithmic clock, the typical prime polynomial has size \(O(A)\),
and a sparse recurrent set can attain size \(2e^A/A+o(e^A/A)\).  A valid
renewal must retain that exceptional set as an actual restoring row rather
than suppress it by a supremum estimate or ignore it by averaging.

A targeted search of the arithmetic manuscript and the current note corpus
found no use of rational independence of prime logarithms, no simultaneous
character-phase alignment theorem, and no frequency mean-square audit of the
localized Weil prime polynomial.  Thus the argument below changes the V7
direction rather than duplicating an existing route.

\subsection{A source-native modulated row}

Fix a primitive character \(\chi\) of conductor \(f\), a support wall
\(A>0\), and its parity parameter
\(\alpha=(1+2\kappa)/4\).  Let
\[
 \mathcal A_f(A)
 :=\{n:\ 2\le n<e^{2A},\ n\ {\rm is\ an\ unramified\ prime\ power}\}.
 \tag{V8.1}
\]
For \(n\in\mathcal A_f(A)\), put
\[
 \lambda_n=\frac{\Lambda(n)}{\sqrt n},
 \qquad
 a_A(n)=1-\frac{\log n}{2A}>0.
 \tag{V8.2}
\]
Define the finite Dirichlet polynomial and its positive triangular mass by
\[
 \boxed{
 \begin{aligned}
 \mathcal P_{\chi,A}(\tau)
 &:=\sum_{n\in\mathcal A_f(A)}
  \lambda_na_A(n)\chi(n)e^{\mathrm i\tau\log n},\\
 \mathcal M_f(A)
 &:=\sum_{n\in\mathcal A_f(A)}\lambda_na_A(n).
 \end{aligned}}
 \tag{V8.3}
\]
Let the full active prime operator be
\[
 \mathsf E_{\chi,A}
 :=\sum_{n\in\mathcal A_f(A)}
   \lambda_n\mathsf P_{\chi,A}^{\rm sc}(n).
 \tag{V8.4}
\]

For \(\xi\in\mathbb R\), introduce the normalized constant-modulus window
\[
 g_\xi(t):=2^{-1/2}e^{\mathrm i\xi t},
 \qquad -1<t<1.
 \tag{V8.5}
\]

\begin{lemma}[Exact modulated prime row and Legendre leakage]
\label{lem:v8-modulated-prime-row}
The window \(g_\xi\) belongs to the closed logarithmic form domain, has unit
norm, and for \(0\le\delta\le2\),
\[
 \boxed{
 \ip{T_\delta g_\xi}{g_\xi}
 =\left(1-\frac\delta2\right)e^{\mathrm i\xi\delta}.}
 \tag{V8.6}
\]
Consequently
\[
 \boxed{
 \ip{\mathsf E_{\chi,A}g_{A\tau}}{g_{A\tau}}
 =2\operatorname{Re}\mathcal P_{\chi,A}(\tau).}
 \tag{V8.7}
\]
If \(P_N\) is the first-\(N\) Legendre projection and \(j_m\) is the
spherical Bessel function, then
\[
 \boxed{
 \norm{P_Ng_\xi}^2
 =\sum_{m=0}^{N-1}(2m+1)j_m(\xi)^2
 =O_N(|\xi|^{-2})}
 \tag{V8.8}
\]
as \(|\xi|\to\infty\).
\end{lemma}

\begin{proof}
The correlation in \textup{(V8.6)} is the exponential phase times the
overlap length \(2-\delta\), divided by the normalization \(2\).  Formula
\textup{(V8.7)} follows from the positive/negative-shift pairing in the prime
cell and \(\delta_n=\log n/A\).

The classical finite Fourier transform of a Legendre polynomial is
\[
 \int_{-1}^1e^{\mathrm i\xi t}P_m(t)\dd t
 =2\mathrm i^m j_m(\xi).
 \tag{V8.9}
\]
After the two normalizations in \textup{(V8.5)} and the definition of
\(e_m\), the squared coefficient is
\((2m+1)j_m(\xi)^2\).  Every fixed \(j_m(\xi)=O_m(|\xi|^{-1})\), proving
\textup{(V8.8)}.  Finally, \(g_\xi\) has finite difference energy and
integrable logarithmic-wall energy.  Smooth cutoff approximation at the two
endpoints places it in the closed form domain.
\end{proof}

The modulated row also permits an exact resummation of the completion.  Write
\(\operatorname{Ci}\) for the cosine integral normalized by
\[
 \operatorname{Ci}(x)
 =\gamma+\log x+\int_0^x\frac{\cos u-1}{u}\dd u,
 \qquad x>0,
 \tag{V8.10}
\]
and define the smooth gamma remainder
\[
 \Sigma_{\alpha,A}(\tau)
 :=\int_0^{2A}
 \left(2-\frac yA\right)s_\alpha(y)\cos(\tau y)\dd y.
 \tag{V8.11}
\]

\begin{theorem}[Exact archimedean clock identity]
\label{thm:v8-clock}
For \(\tau\ne0\), the complete localized Weil operator of V7 satisfies
\[
 \boxed{
 \begin{aligned}
 \ip{\mathsf H_{\chi,A}^{\rm sc}g_{A\tau}}{g_{A\tau}}
 ={}&\log\frac{f|\tau|}{2\pi}
 -\operatorname{Ci}(2A|\tau|)
 +\frac{\sin(2A\tau)}{2A\tau}\\
 &+\Sigma_{\alpha,A}(\tau)
 -2\operatorname{Re}\mathcal P_{\chi,A}(\tau).
 \end{aligned}}
 \tag{V8.12}
\]
For fixed \(A\),
\[
 \Sigma_{\alpha,A}(\tau)\longrightarrow0
 \tag{V8.13}
\]
as \(|\tau|\to\infty\).  Thus the leading completed cost of a modulation is
the source-native clock \(\log(f|\tau|/(2\pi))\), with no residual support-wall
term.
\end{theorem}

\begin{proof}
Using the difference form in the manuscript's \textup{(W.23)},
\[
 \begin{aligned}
 \ip{\mathsf Dg_\xi}{g_\xi}
 &=\frac12\int_0^2
   \frac{(2-r)(1-\cos(\xi r))}{r}\dd r\\
 &=\gamma+\log(2|\xi|)-1-\operatorname{Ci}(2|\xi|)
   +\frac{\sin(2\xi)}{2\xi}.
 \end{aligned}
 \tag{V8.14}
\]
The logarithmic wall contributes
\[
 \ip{\mathsf Wg_\xi}{g_\xi}=1-\log2.
 \tag{V8.15}
\]
The constant row of the scaled smooth kernel is
\[
 \ip{\mathsf S_{\alpha,A}^{\rm sc}g_\xi}{g_\xi}
 =A\int_0^2(2-r)s_\alpha(Ar)\cos(\xi r)\dd r.
 \tag{V8.16}
\]
Put \(\xi=A\tau\), change variables \(y=Ar\), add
\(c_{f,A}=\log(f/(2\pi A))-\gamma\), and use
\textup{(V8.7)}.  All constants and all appearances of \(A\) in the main
logarithm cancel, yielding \textup{(V8.12)}.  The integrand in
\textup{(V8.11)} is continuous on a compact interval, so
\textup{(V8.13)} is the Riemann--Lebesgue lemma.
\end{proof}

\begin{corollary}[A finite dangerous clock window for this row]
\label{cor:v8-clock-window}
Put
\[
 C_{\alpha,A}
 :=1+\sup_{x\ge2}|\operatorname{Ci}(x)|
 +\int_0^{2A}\left|2-\frac yA\right||s_\alpha(y)|\dd y.
 \tag{V8.17}
\]
If
\[
 |\tau|\ge A^{-1}
 \quad\text{and}\quad
 \log\frac{f|\tau|}{2\pi}
 \ge2\mathcal M_f(A)+C_{\alpha,A},
 \tag{V8.18}
\]
then
\[
 \ip{\mathsf H_{\chi,A}^{\rm sc}g_{A\tau}}{g_{A\tau}}\ge0.
 \tag{V8.19}
\]
Hence this particular modulated row can be negative only in an explicit
finite clock interval.
\end{corollary}

\begin{proof}
Use \(|\mathcal P_{\chi,A}(\tau)|\le\mathcal M_f(A)\), the bound
\(|\sin(2A\tau)/(2A\tau)|\le1\), and \textup{(V8.17)} in
\textup{(V8.12)}.
\end{proof}

\subsection{Simultaneous prime-log resonance}

\begin{theorem}[Arbitrarily high exact-supremum resonance]
\label{thm:v8-Kronecker}
For every \(\varepsilon>0\) and \(T_0>0\), there is \(\tau>T_0\) such that
\[
 |\chi(p)e^{\mathrm i\tau\log p}-1|<\varepsilon
 \tag{V8.20}
\]
simultaneously for every unramified prime \(p<e^{2A}\).  Consequently
\[
 \boxed{
 \limsup_{\tau\to+\infty}
 \operatorname{Re}\mathcal P_{\chi,A}(\tau)
 =\sup_{\tau\in\mathbb R}|\mathcal P_{\chi,A}(\tau)|
 =\mathcal M_f(A).}
 \tag{V8.21}
\]
\end{theorem}

\begin{proof}
The numbers \(\{\log p\}\), over any finite collection of distinct primes,
are rationally independent.  Indeed, an integer relation would exponentiate
to a relation among prime factorizations.  Kronecker's theorem therefore
makes the continuous flow
\[
 \tau\longmapsto(\tau\log p\bmod 2\pi)_p
\]
dense in the corresponding torus, and recurrence supplies arbitrarily large
positive hitting times.  Choose the target angle
\(-\arg\chi(p)\) in every coordinate.

If \(n=p^k\), then
\(\chi(n)e^{\mathrm i\tau\log n}
 =(\chi(p)e^{\mathrm i\tau\log p})^k\).
There are only finitely many relevant exponents, so the base-prime
approximation makes every summand in \(\mathcal P_{\chi,A}\) arbitrarily
close to its positive modulus.  This proves the lower equality in
\textup{(V8.21)}; the triangle inequality gives the reverse bound.
\end{proof}

\begin{theorem}[No uniform character cancellation in any Legendre tail]
\label{thm:v8-tail-resonance}
For every fixed \(A>0\) and every finite \(N\),
\[
 \boxed{
 \norm{Q_N\mathsf E_{\chi,A}Q_N}
 \ge2\mathcal M_f(A).}
 \tag{V8.22}
\]
In particular, projecting away any finite collection of low Legendre modes
does not create a uniform phase gain for the active prime operator.
\end{theorem}

\begin{proof}
Choose a resonant sequence \(\tau_j\to\infty\) from
\Cref{thm:v8-Kronecker}.  By \textup{(V8.8)},
\[
 h_j:=\frac{Q_Ng_{A\tau_j}}{\norm{Q_Ng_{A\tau_j}}}
 \quad\text{satisfies}\quad
 \norm{h_j-g_{A\tau_j}}\longrightarrow0.
 \tag{V8.23}
\]
At fixed \(A\), the active prime operator is a finite bounded sum.  Hence
\textup{(V8.7)}, \textup{(V8.21)}, and \textup{(V8.23)} give
\[
 \ip{Q_N\mathsf E_{\chi,A}Q_Nh_j}{h_j}
 \longrightarrow2\mathcal M_f(A).
\]
Taking the supremum of the Rayleigh quotient proves \textup{(V8.22)}.
\end{proof}

\begin{corollary}[The resonance mass is exponentially large]
\label{cor:v8-resonance-mass}
As \(A\to\infty\),
\[
 \boxed{
 \mathcal M_f(A)=\frac{2e^A}{A}+o\!\left(\frac{e^A}{A}\right),}
 \tag{V8.24}
\]
and therefore, for every finite \(N\),
\[
 \boxed{
 \norm{Q_N\mathsf E_{\chi,A}Q_N}
 \ge\frac{4e^A}{A}
   +o\!\left(\frac{e^A}{A}\right).}
 \tag{V8.25}
\]
\end{corollary}

\begin{proof}
Put \(x=e^{2A}\).  Partial summation of the classical
de la Vall\'ee Poussin error term in the prime number theorem gives
\[
 \begin{aligned}
 \mathcal M_f(A)
 &=\sum_{\substack{n<x\\n\ {\rm unramified}}}
   \frac{\Lambda(n)}{\sqrt n}
   \left(1-\frac{\log n}{\log x}\right)\\
 &=\int_1^x t^{-1/2}
   \left(1-\frac{\log t}{\log x}\right)\dd t
   +o(\sqrt x/\log x)\\
 &=\frac{4\sqrt x}{\log x}
   +o(\sqrt x/\log x).
 \end{aligned}
 \tag{V8.26}
\]
The finitely many ramified prime towers contribute \(O_f(A)\).  Substitution
of \(\log x=2A\) proves \textup{(V8.24)}, and
\Cref{thm:v8-tail-resonance} proves \textup{(V8.25)}.
\end{proof}

\begin{corollary}[Uniform operator-norm coercivity remains double exponential]
\label{cor:v8-norm-cost}
Suppose one attempts to certify the high tail only through
\[
 H_N+c_{f,A}
 -\norm{\mathsf S_{\alpha,A}^{\rm sc}}
 -\norm{Q_N\mathsf E_{\chi,A}Q_N}>0.
 \tag{V8.27}
\]
Then necessarily
\[
 \boxed{
 \log N\ge\frac{4e^A}{A}
 +o\!\left(\frac{e^A}{A}\right),\qquad
 N\ge\exp\!\left(
 \frac{(4+o(1))e^A}{A}\right).}
 \tag{V8.28}
\]
Thus exact character-sensitive operator norm can improve the V7 absolute
bound by at most the displayed factor \(A\) at this stage; it cannot change
the double-exponential character of the protected dimension.
\end{corollary}

\begin{proof}
Condition \textup{(V8.27)} implies
\[
 H_N>\norm{Q_N\mathsf E_{\chi,A}Q_N}-c_{f,A}.
\]
Apply \textup{(V8.25)}, \(H_N\le1+\log N\), and
\(-c_{f,A}=\log A+O_f(1)\).
\end{proof}

\subsection{Mean-square cancellation and the exceptional set}

The uniform-tail no-go does not mean that a typical modulation sees the
resonant mass.  The exact frequency mean square is dramatically smaller.
Define
\[
 \mathcal V_f(A)
 :=\sum_{n\in\mathcal A_f(A)}
   \lambda_n^2a_A(n)^2.
 \tag{V8.29}
\]

\begin{theorem}[Exact frequency mean square]
\label{thm:v8-mean-square}
For every fixed \(A\),
\[
 \boxed{
 \lim_{T\to\infty}\frac1T\int_0^T
 |\mathcal P_{\chi,A}(\tau)|^2\dd\tau
 =\mathcal V_f(A).}
 \tag{V8.30}
\]
As \(A\to\infty\),
\[
 \boxed{
 \mathcal V_f(A)=\frac{A^2}{3}+o(A^2).}
 \tag{V8.31}
\]
Consequently, for every \(Y>0\),
\[
 \boxed{
 \limsup_{T\to\infty}\frac1T
 \operatorname{meas}\left\{\tau\in[0,T]:
 |\mathcal P_{\chi,A}(\tau)|\ge Y\right\}
 \le\frac{\mathcal V_f(A)}{Y^2}.}
 \tag{V8.32}
\]
In particular, after \(T\to\infty\), the upper density of the set where
\(|\mathcal P_{\chi,A}|\ge A^{1+\eta}\) is \(O(A^{-2\eta})\) for every
fixed \(\eta>0\).
\end{theorem}

\begin{proof}
On expanding the square, every off-diagonal term contains
\[
 \frac1T\int_0^T e^{\mathrm i\tau(\log m-\log n)}\dd\tau,
\]
which tends to zero for \(m\ne n\).  The diagonal terms give
\textup{(V8.30)}.

Prime powers of exponent at least two contribute \(O(1)\) to
\(\mathcal V_f(A)\).  For the primes, partial summation of the prime number
theorem gives, with \(L=2A\),
\[
 \begin{aligned}
 \sum_{p<e^L}\frac{(\log p)^2}{p}
 \left(1-\frac{\log p}{L}\right)^2
 &=(1+o(1))\int_0^L y\left(1-\frac yL\right)^2\dd y\\
 &=\frac{L^2}{12}+o(L^2)
 =\frac{A^2}{3}+o(A^2).
 \end{aligned}
 \tag{V8.33}
\]
This proves \textup{(V8.31)}.  Markov's inequality applied to
\textup{(V8.30)} proves \textup{(V8.32)}.
\end{proof}

\begin{corollary}[Quantified typical--resonant separation]
\label{cor:v8-separation}
The root-mean-square size and recurrent supremum satisfy
\[
 \boxed{
 \mathcal V_f(A)^{1/2}
 =\frac A{\sqrt3}+o(A),
 \qquad
 \frac{\mathcal M_f(A)}{\mathcal V_f(A)^{1/2}}
 =\frac{2\sqrt3\,e^A}{A^2}(1+o(1)).}
 \tag{V8.34}
\]
Thus the obstruction to a small uniform norm is an exponentially separated
recurrent resonance set, not the typical frequency scale.
\end{corollary}

\subsection{V8 firewall and the revised upstream target}

\begin{firewall}[Average cancellation is not positivity]
\label{fire:v8-average}
The mean-square identity \textup{(V8.30)} cannot be substituted into a
pointwise lower bound for the Weil form.  Conversely, the recurrence theorem
\textup{(V8.21)} does not itself produce a negative Weil witness: along an
unspecified resonant sequence the archimedean clock
\(\log(f|\tau|/(2\pi))\) may already exceed the aligned prime mass.  A valid
argument must compare the location and width of exceptional resonances with
that clock, or retain those frequencies in an explicit protected block.
\end{firewall}

\begin{firewall}[Uniform phase cancellation is unavailable]
\label{fire:v8-uniform}
No upper estimate of the form
\[
 \norm{Q_N\mathsf E_{\chi,A}Q_N}\le e^{o(A)}
 \tag{V8.35}
\]
can hold for any finite choice \(N=N(A)\) as \(A\to\infty\); the lower bound
\textup{(V8.25)} rules it out.  Increasing \(N\) does not remove the
recurrences, because every finite Legendre tail contains arbitrarily high
modulations.  Therefore V7's proposed character-sensitive
\emph{uniform-operator-norm} replacement must be abandoned.
\end{firewall}

\begin{center}
\begin{tabularx}{\textwidth}{@{}p{0.28\textwidth}p{0.17\textwidth}X@{}}
\toprule
V8 component & Status & Exact conclusion \\
\midrule
Modulated completed row & proved & Formula \textup{(V8.12)} couples the
prime polynomial to the conductor clock with all support constants canceled. \\
Prime-log phase alignment & proved & Arbitrarily high Kronecker recurrences
attain the full triangular prime mass. \\
Finite-tail character cancellation & disproved in uniform norm & Every
\(Q_N\) compression retains norm at least \(2\mathcal M_f(A)\). \\
Resonant mass & evaluated & It is \(2e^A/A+o(e^A/A)\), still exponential in
the support wall. \\
Typical frequency size & evaluated & The exact mean square is
\(A^2/3+o(A^2)\). \\
Exceptional set & sparse but unavoidable & Its density has the bound
\textup{(V8.32)}, but positivity cannot discard it. \\
Cofinal positivity and GRH & open & A restoring mechanism for exceptional
frequency rows is not yet proved. \\
\bottomrule
\end{tabularx}
\end{center}

V9 should no longer seek a uniform phase gain.  The next source-native
problem is an exceptional-frequency renewal:
\begin{enumerate}[label=\textup{(\alph*)},leftmargin=2.8em]
\item partition modulation space by a predeclared frequency tiling;
\item use the diagonal mean square \textup{(V8.30)} on typical tiles;
\item retain every tile exceeding its clock budget as a protected modulation
row rather than averaging it away;
\item use the exact Legendre leakage \textup{(V8.8)} and Schur bracket
\textup{(V7.26)} to quantify the coupling of those rows back to the polynomial
low block.
\end{enumerate}
The key missing estimate is not cancellation everywhere, but a
source-independent capacity bound for the clock-deficient exceptional rows.
If their number or Gram capacity is not absorbable, that failure itself
identifies the next genuine restoring sector.

\begin{assessment}[V8 acquisition]
V8 closes the naive phase-sensitive tail route.  Rational independence of
prime logarithms forces full character alignment arbitrarily far into every
finite Legendre tail, so a uniform operator norm cannot improve the problem
below an exponential prime mass.  At the same time, exact frequency
orthogonality proves an \(A^2/3\) mean square and isolates the obstruction in
a sparse recurrent set.  The completed scalar row \textup{(V8.12)} supplies
the correct logarithmic clock against which those recurrences must be judged.
GRH is not proved.
\end{assessment}


% =============================================================================
\section{V9: Fourier-lattice acquisition and the resolution--clock mismatch}
\label{sec:grh-v9}
% =============================================================================

V8 proved that uniform phase cancellation is impossible but that the
frequency mean square is small.  The proposed repair was to retain
clock-deficient exceptional frequencies as protected modulation rows.  This
section tests whether that proposal supplies a finite-capacity bridge.

The natural predeclared rows are the Fourier basis of the scaled support
interval.  In that basis every truncated translation has an exact sinc
matrix, the diagonal is precisely the V8 Dirichlet polynomial, and discrete
frequency averaging has an explicit finite-horizon error.  The calculation
reveals a resolution--clock mismatch: the averaging horizon needed to
separate adjacent prime logarithms is of order \(Ae^{2A}\), while the
archimedean clock at that horizon is already \(2A+O(\log A)\), larger than
the root-mean-square prime row.  Thus the mean square becomes quantitatively
diagonal only after it has mixed the unresolved early region with a much
larger high-clock region.

A targeted nonduplication search found general large-sieve references in the
arithmetic manuscript but no Fourier matrix for the compressed Weil shifts,
no alias-free sampling result at the support-adapted lattice, and no
comparison between prime-log resolution time and the completed clock.  The
results below therefore populate the V8 exceptional-row proposal rather than
restate an existing sieve step.

\subsection{Exact Fourier matrix of a compressed shift}

Let
\[
 \phi_k(t):=2^{-1/2}e^{\mathrm i\pi kt},
 \qquad k\in\mathbb Z,
 \tag{V9.1}
\]
so that \((\phi_k)_{k\in\mathbb Z}\) is a source-fixed orthonormal basis of
\(L^2(-1,1)\).  Put
\[
 \operatorname{sinc}z:=
 \begin{cases}
  \sin z/z,&z\ne0,\\
  1,&z=0.
 \end{cases}
 \tag{V9.2}
\]

\begin{theorem}[Complete Fourier matrix of one truncated translation]
\label{thm:v9-Fourier-shift}
For \(0\le\delta\le2\), let \(a=1-\delta/2\).  Then
\[
 \boxed{
 \ip{T_\delta\phi_\ell}{\phi_k}
 =a\,e^{\mathrm i\pi(k+\ell)\delta/2}
 \operatorname{sinc}\!\bigl(\pi(\ell-k)a\bigr).}
 \tag{V9.3}
\]
Consequently the self-adjoint prime-power cell has the real symmetric matrix
\[
 \boxed{
 \begin{aligned}
 \ip{\mathsf P_{\chi,A}^{\rm sc}(n)\phi_\ell}{\phi_k}
 ={}&2a_A(n)
 \operatorname{sinc}\!\bigl(\pi(\ell-k)a_A(n)\bigr)\\
 &\times\operatorname{Re}\!\left[
 \chi(n)e^{\mathrm i\pi(k+\ell)\delta_n/2}\right],
 \end{aligned}}
 \tag{V9.4}
\]
where \(\delta_n=\log n/A\).  In particular, its diagonal is
\[
 \boxed{
 \ip{\mathsf P_{\chi,A}^{\rm sc}(n)\phi_k}{\phi_k}
 =2a_A(n)\operatorname{Re}\!\left[
 \chi(n)e^{\mathrm i\pi k\log n/A}\right].}
 \tag{V9.5}
\]
For \(k\ne\ell\),
\[
 \left|\ip{T_\delta\phi_\ell}{\phi_k}\right|
 \le\min\left\{a,\frac1{\pi|k-\ell|}\right\}.
 \tag{V9.6}
\]
\end{theorem}

\begin{proof}
The overlap interval has length \(2-\delta=2a\) and midpoint
\(-\delta/2\).  Therefore
\[
 \begin{aligned}
 \ip{T_\delta\phi_\ell}{\phi_k}
 &=\frac12e^{\mathrm i\pi\ell\delta}
   \int_{-1}^{1-\delta}
   e^{\mathrm i\pi(\ell-k)t}\dd t\\
 &=a\,e^{\mathrm i\pi(k+\ell)\delta/2}
   \operatorname{sinc}\!\bigl(\pi(\ell-k)a\bigr),
 \end{aligned}
 \]
which proves \textup{(V9.3)}.  The negative shift is the adjoint, so adding
the two character-weighted matrices gives \textup{(V9.4)}.  Its
\(k=\ell\) specialization is \textup{(V9.5)}, and
\(|\sin z|\le\min\{|z|,1\}\) gives \textup{(V9.6)}.
\end{proof}

\begin{corollary}[The V8 scalar row is the exact Fourier diagonal]
\label{cor:v9-Fourier-diagonal}
For the active prime operator \(\mathsf E_{\chi,A}\) of
\textup{(V8.4)}, set
\[
 \tau_k:=\frac{\pi k}{A}.
 \tag{V9.7}
\]
Then
\[
 \boxed{
 \ip{\mathsf E_{\chi,A}\phi_k}{\phi_k}
 =2\operatorname{Re}\mathcal P_{\chi,A}(\tau_k).}
 \tag{V9.8}
\]
Thus the support-adapted modulation lattice is not an auxiliary sampling
device: it is the diagonal of the actual compressed prime operator in an
orthonormal physical basis.
\end{corollary}

\subsection{Alias-free prime-log sampling}

For \(n\in\mathcal A_f(A)\), define
\[
 \theta_n:=\frac{\pi\log n}{A}\in(0,2\pi),
 \qquad
 c_n:=\lambda_na_A(n)\chi(n).
 \tag{V9.9}
\]
Then
\[
 \mathcal P_{\chi,A}(\tau_k)=\sum_nc_ne^{\mathrm ik\theta_n}.
 \tag{V9.10}
\]
For angles on the circle, let
\(\|x\|_{2\pi}:=\min_{j\in\mathbb Z}|x-2\pi j|\).

\begin{lemma}[Explicit separation of the sampled prime logs]
\label{lem:v9-angle-separation}
Assume \(A\ge1\).  If \(m\ne n\) belong to \(\mathcal A_f(A)\), then
\[
 \boxed{
 \|\theta_m-\theta_n\|_{2\pi}
 \ge\frac{\pi e^{-2A}}{A}.}
 \tag{V9.11}
\]
In particular the sampled frequencies are pairwise distinct modulo
\(2\pi\), so there is no diagonal aliasing.
\end{lemma}

\begin{proof}
Suppose \(m<n<e^{2A}\).  The ordinary angular difference satisfies
\[
 \frac{\pi}{A}\log\frac nm
 \ge\frac{\pi}{A}\log\left(1+\frac1m\right)
 \ge\frac{\pi}{A(m+1)}
 >\frac{\pi e^{-2A}}A.
\]
If the shorter circular distance goes across the endpoint, it is
\[
 \frac{\pi}{A}\left(2A-\log\frac nm\right)
 =\frac{\pi}{A}\log\frac{e^{2A}m}{n}
 >\frac{\pi\log2}{A},
\]
which is larger than the claimed bound.
\end{proof}

Let
\[
 D_K(x):=\sum_{k=-K}^Ke^{\mathrm ikx}
 \tag{V9.12}
\]
be the Dirichlet kernel.

\begin{theorem}[Quantitative discrete mean square]
\label{thm:v9-discrete-mean-square}
There is an absolute constant \(C_0\) such that, for \(A\ge1\) and every
integer \(K\ge1\),
\[
 \boxed{
 \left|
 \sum_{k=-K}^K|\mathcal P_{\chi,A}(\tau_k)|^2
 -(2K+1)\mathcal V_f(A)
 \right|
 \le\Delta_A\mathcal V_f(A),}
 \tag{V9.13}
\]
where one may take
\[
 \boxed{\Delta_A=C_0A^2e^{2A}.}
 \tag{V9.14}
\]
Consequently
\[
 \boxed{
 \sum_{k=-K}^K|\mathcal P_{\chi,A}(\tau_k)|^2
 \le(2K+1+\Delta_A)\mathcal V_f(A).}
 \tag{V9.15}
\]
\end{theorem}

\begin{proof}
Expanding the square in \textup{(V9.10)} gives the Gram matrix
\[
 G_{mn}=D_K(\theta_m-\theta_n).
 \]
Its diagonal is \(2K+1\).  For \(x\) represented in \([-\pi,\pi]\),
\[
 |D_K(x)|
 \le\min\left\{2K+1,\frac1{|\sin(x/2)|}\right\}
 \le\frac{\pi}{|x|}.
 \tag{V9.16}
\]
Order the at most \(e^{2A}\) angles around the circle.  By
\Cref{lem:v9-angle-separation}, the \(j\)-th neighbor on either side is at
circular distance at least a constant times
\(j\pi e^{-2A}/A\).  Therefore every off-diagonal absolute row sum is at
most
\[
 C\frac{Ae^{2A}}{\pi}
 \sum_{j\le e^{2A}}\frac1j
 \le C_0A^2e^{2A}.
 \tag{V9.17}
\]
The Schur test bounds the operator norm of the off-diagonal Gram matrix by
the same quantity.  Since
\(\sum_n|c_n|^2=\mathcal V_f(A)\), equation
\textup{(V9.13)} follows.
\end{proof}

\begin{corollary}[Finite exceptional-row upper ledger]
\label{cor:v9-exception-upper}
For \(Y>0\), put
\[
 \mathcal E_{A,K}(Y)
 :=\{k\in\mathbb Z:\ |k|\le K,\
 |\mathcal P_{\chi,A}(\tau_k)|\ge Y\}.
 \tag{V9.18}
\]
Then
\[
 \boxed{
 \#\mathcal E_{A,K}(Y)
 \le
 \frac{(2K+1+\Delta_A)\mathcal V_f(A)}{Y^2}.}
 \tag{V9.19}
\]
If \(K\gg A^2e^{2A}\) and \(Y=A^{1+\eta}\), this gives the relative bound
\[
 \frac{\#\mathcal E_{A,K}(A^{1+\eta})}{2K+1}
 =O(A^{-2\eta}).
 \tag{V9.20}
\]
\end{corollary}

\begin{proof}
Every exceptional term contributes at least \(Y^2\) to the left side of
\textup{(V9.15)}.  Apply \textup{(V8.31)} for the final asymptotic.
\end{proof}

\begin{corollary}[Many moderately large absolute rows after resolution]
\label{cor:v9-exception-lower}
If \(2K+1\ge2\Delta_A\), then
\[
 \boxed{
 \#\left\{k:\ |k|\le K,\
 |\mathcal P_{\chi,A}(\tau_k)|
 \ge\frac12\mathcal V_f(A)^{1/2}\right\}
 \ge
 \frac{(2K+1)\mathcal V_f(A)}
 {4\mathcal M_f(A)^2}.}
 \tag{V9.21}
\]
The lower relative density in \textup{(V9.21)} is
\[
 \boxed{
 \frac{\mathcal V_f(A)}{4\mathcal M_f(A)^2}
 =\frac{A^4}{48e^{2A}}(1+o(1)).}
 \tag{V9.22}
\]
\end{corollary}

\begin{proof}
Under the hypothesis, \textup{(V9.13)} gives a total square sum at least
\((2K+1)\mathcal V_f(A)/2\).  If \(r\) samples exceed
\(\mathcal V_f(A)^{1/2}/2\), the triangle bound
\(|\mathcal P_{\chi,A}|\le\mathcal M_f(A)\) gives
\[
 \sum_{k=-K}^K|\mathcal P_{\chi,A}(\tau_k)|^2
 \le r\mathcal M_f(A)^2
 +(2K+1)\frac{\mathcal V_f(A)}4.
\]
Rearrangement proves \textup{(V9.21)}.  Insert
\textup{(V8.24)} and \textup{(V8.31)} to obtain
\textup{(V9.22)}.
\end{proof}

\subsection{The resolution--clock mismatch}

The quantitative error \(\Delta_A\) is the price of resolving the closest
distinct logarithms \(\log n\) below \(2A\).  Set
\[
 K_{\rm res}(A):=\left\lceil2\Delta_A\right\rceil,
 \qquad
 \tau_{\rm res}(A):=\frac{\pi K_{\rm res}(A)}A.
 \tag{V9.23}
\]

\begin{theorem}[Prime-log resolution occurs after the typical clock crossing]
\label{thm:v9-resolution-clock}
As \(A\to\infty\),
\[
 \boxed{
 \log\frac{f\tau_{\rm res}(A)}{2\pi}
 =2A+\log A+O_f(1),}
 \tag{V9.24}
\]
whereas the root-mean-square upper scale of the adverse diagonal prime term is
\[
 \boxed{
 2\mathcal V_f(A)^{1/2}
 =\frac{2A}{\sqrt3}+o(A).}
 \tag{V9.25}
\]
In particular, the finite-horizon orthogonality estimate
\textup{(V9.13)} becomes uniformly relative only after the endpoint
archimedean clock has already passed the root-mean-square prime scale by a
positive linear margin.
\end{theorem}

\begin{proof}
Equations \textup{(V9.14)} and \textup{(V9.23)} give
\[
 \tau_{\rm res}(A)
 =2\pi C_0Ae^{2A}(1+o(1)),
\]
which proves \textup{(V9.24)}.  Equation \textup{(V9.25)} is
\textup{(V8.31)}.
\end{proof}

\begin{firewall}[Why the mean square does not control the early window]
\label{fire:v9-resolution}
The interval average in \textup{(V9.13)} includes all rows from frequency
zero up to \(\tau_{\rm res}(A)\).  Its small mean is dominated numerically by
the enormous population of late rows whose completed clock is already of
size \(2A\).  It therefore gives no pointwise or capacity estimate for the
earlier clock-deficient band.  Applying the limiting mean square from V8
inside that band without the finite-horizon error \(\Delta_A\) would discard
the very prime-log near-collisions that delay orthogonality.
\end{firewall}

\subsection{Off-diagonal leakage and the capacity firewall}

Even a perfect diagonal exceptional-set count would not yet control the
prime operator.  Formula \textup{(V9.4)} gives, for the full active operator,
\[
 \boxed{
 \begin{aligned}
 \ip{\mathsf E_{\chi,A}\phi_\ell}{\phi_k}
 =2\sum_{n\in\mathcal A_f(A)}
 &\lambda_na_A(n)
 \operatorname{sinc}\!\bigl(\pi(\ell-k)a_A(n)\bigr)\\
 &\times\operatorname{Re}\!\left[
 \chi(n)e^{\mathrm i\pi(k+\ell)\delta_n/2}\right].
 \end{aligned}}
 \tag{V9.26}
\]
The matrix depends simultaneously on the sum and difference of the Fourier
indices and is not diagonal, Toeplitz, or a scalar multiplier.

\begin{proposition}[Diagonal retention is not a Schur certificate]
\label{prop:v9-diagonal-not-enough}
Let \(\Pi_{\mathcal E}\) project onto any selected set of Fourier rows.  The
diagonal data
\[
 \left\{
 \ip{\mathsf E_{\chi,A}\phi_k}{\phi_k}:k\in\mathbb Z
 \right\}
 \tag{V9.27}
\]
do not determine any of the three quantities
\[
 \Pi_{\mathcal E}\mathsf E_{\chi,A}(I-\Pi_{\mathcal E}),
 \qquad
 (I-\Pi_{\mathcal E})\mathsf E_{\chi,A}(I-\Pi_{\mathcal E}),
 \qquad
 \norm{\mathsf E_{\chi,A}}.
 \tag{V9.28}
\]
For the actual Weil operator these missing entries are the nonzero sinc
couplings in \textup{(V9.26)}.  Therefore an exceptional-row renewal based
only on \(\mathcal E_{A,K}(Y)\) cannot be inserted into the V7 Schur bracket.
\end{proposition}

\begin{proof}
Even among \(2\times2\) self-adjoint matrices, the family
\[
 \begin{pmatrix}d_1&z\\\overline z&d_2\end{pmatrix}
 \]
has fixed diagonal and arbitrary off-diagonal coupling \(z\), hence arbitrary
Schur feedback and varying norm.  Equation \textup{(V9.26)} shows that this is
not merely an abstract ambiguity: the compressed translations supply
off-diagonal entries unless their individual sinc factors happen to vanish.
\end{proof}

\begin{firewall}[Sparse density is not small protected dimension]
\label{fire:v9-capacity}
Because the \(\phi_k\) are orthonormal, retaining \(r\) lattice rows costs
exactly \(r\) protected dimensions before any off-diagonal completion.  The
upper ledger \textup{(V9.19)} is proportional to the frequency horizon
\(K\); small relative density does not make the absolute protected rank
cofinally small.  Moreover, \textup{(V9.21)} shows that after logarithmic
resolution there are linearly many rows with moderately large absolute
prime diagonal, albeit with exponentially small density.  Neither fact
decides their sign against the clock, and neither supplies the missing sinc
cross block.
\end{firewall}

\subsection{V9 status and the mandated V10 sweep}

\begin{center}
\begin{tabularx}{\textwidth}{@{}p{0.28\textwidth}p{0.17\textwidth}X@{}}
\toprule
V9 component & Status & Exact conclusion \\
\midrule
Fourier shift matrix & proved & Formula \textup{(V9.3)} acquires every
diagonal and off-diagonal truncated-shift entry. \\
Prime Fourier diagonal & identified & It is exactly the V8 Dirichlet
polynomial on the support-adapted lattice. \\
Finite-horizon mean square & proved & Equation \textup{(V9.13)} has explicit
resolution error \(O(A^2e^{2A})\). \\
Exceptional-row count & bounded above and below in absolute magnitude &
Equations \textup{(V9.19)} and \textup{(V9.21)} quantify the lattice ledger. \\
Resolution versus completion & obstructed & Reliable diagonal averaging
begins only after the typical archimedean clock crossing. \\
Diagonal exceptional renewal & insufficient & The sinc cross block
\textup{(V9.26)} is indispensable for a Schur certificate. \\
Cofinal positivity and GRH & open & No sign-sensitive capacity estimate for
the early Fourier block has been proved. \\
\bottomrule
\end{tabularx}
\end{center}

V10 is the next version at which the user requested a full sweep of parallel
notes.  That sweep should be performed before choosing another analytic
route.  Subject to anything genuinely new found there, the active
source-native target is the \emph{early Fourier block}
\[
 |\tau|\lesssim Ae^{2A},
 \tag{V9.29}
\]
with all entries acquired by \textup{(V9.26)} and with the completed
diagonal supplied by \textup{(V8.12)}.  The required theorem is a
sign-sensitive Schur or block-capacity bound; a diagonal large-sieve count is
now proved insufficient.

\begin{assessment}[V9 acquisition]
V9 makes the exceptional-frequency proposal finite and exact, then identifies
its missing bridge.  The Fourier lattice gives a closed sinc formula for the
entire prime operator and a quantitative discrete mean square.  However,
prime-log orthogonality is not resolved until the completed clock has already
crossed the typical prime scale, and diagonal exceptional counts omit the
actual sinc feedback needed by the Schur complement.  The next version must
first sweep the parallel note corpus, then attack the sign-sensitive early
Fourier block rather than repeat a density-only argument.  GRH is not proved.
\end{assessment}


% =============================================================================
% BEGIN APPENDED V10 MATERIAL
% =============================================================================

\section{V10: compulsory corpus sweep and the support-wall compactness trap}
\label{sec:grh-v10}

V9 made the complete prime operator visible in the Fourier basis, but it
also proved that diagonal frequency density is not a Schur certificate.  At
this tenth-version boundary we first carry out the requested sweep of the
parallel note corpus.  The sweep changes the order of attack: a relative
finite-matrix atlas is useful only after the tail object has been put in the
correct clock-relative geometry.  We therefore do not begin another grid of
low-mode positivity tests.  Instead we prove what happens to a prime cell at
the instant it enters through the support wall.

The result has two parts.  First, a newly active shift converges strongly to
zero and disappears from every fixed finite compression, while its norm on
the complementary tail converges to one.  This is a genuine compactness
trap, not merely a loose constant in the V7 estimate.  Second, the norm-one
vectors live in shrinking endpoint layers.  The harmonic-number operator
charges such vectors a logarithmically divergent amount, and an explicit
relative-form inequality records that compensation.  This identifies the
next source-native object: an aggregated, character-sensitive boundary form
bound relative to the harmonic clock.

\subsection{The mandated parallel-note sweep}

The checkpoint was taken on 10 September 2026 after the rolling Standard
Model note advanced during the sweep.  Excluding the present GRH file, the
folder contained twenty-nine \texttt{.tex} files.  In lexicographic order,
the canonical manifest rows were
\[
 \texttt{filename|byte-count|SHA--256},
 \tag{V10.1}
\]
terminated by a newline.  The SHA--256 of that complete UTF--8 manifest was
\[
 \boxed{
 \texttt{C083EEA4A927780D578EE60938FEE10ABB35EDDEA093FAF54CE980236AA54EA8}.}
 \tag{V10.2}
\]
The archive groups represented in the manifest were
\[
\begin{array}{p{0.20\textwidth}|p{0.68\textwidth}}
\text{programme}&\text{files inspected}\cr\hline
\text{Navier--Stokes}
 &\texttt{V100\_f1},\ \texttt{V100\_f2},\ \texttt{V100\_f3},
   \texttt{V26}\cr
\text{SM--Einstein}
 &\texttt{V157\_f1},\ \texttt{V100\_f2}--\texttt{V100\_f13},
   \texttt{V58}\cr
\text{Yang--Mills}
 &\texttt{V135\_f1},\ \texttt{V100\_f2}--\texttt{V100\_f7},
   \texttt{V65\_f8},\ \texttt{V100\_f9},\ \texttt{V19}\cr
\text{parallel Yang--Mills}
 &\text{parallel YM V5}\cr
\end{array}
\tag{V10.3}
\]
The four live terminal records at the checkpoint were
\[
\begin{array}{p{0.49\textwidth}|r|p{0.31\textwidth}}
\text{checkpoint}&\text{bytes}&\text{SHA--256}\cr\hline
\text{NS V26}
 &278283&{\scriptsize \texttt{82C887F8C2C69E4F28FAD7C3DC8DE5B}\allowbreak\texttt{9099CB5A4BF431EBA1FFF3E91935978AD}}\cr
\text{SM--Einstein V58}
 &739567&{\scriptsize \texttt{1AD41098D4E68CDA311125084968836D}\allowbreak\texttt{1F8210DCC5396D0753ACB7BA3A015394}}\cr
\text{Yang--Mills V19}
 &152282&{\scriptsize \texttt{8CF5FEE84D9E0E8AB13D9DDD906EC9D}\allowbreak\texttt{FE8AC500EBBE7B4FBDD73DB49ABE8A593}}\cr
\text{parallel YM V5}
 &54174&{\scriptsize \texttt{306F1BB84CFA8A8D47EA569EC17E9545}\allowbreak\texttt{F9867C8D32B548EC9D9C444B4F3C0512}}\cr
\end{array}
\tag{V10.4}
\]

The archived cycles were read through their result ledgers, terminal audits,
acquisition tables, decisions, and claim boundaries; the live continuations
were checked at their new mathematical sections.  This prevents a version
label from being mistaken for a new theorem and prevents an embedded
programme instruction from being treated as an instruction for this GRH
lineage.

\begin{assessment}[Transfer ledger from the sweep]
\label{ass:v10-sweep}
The following are the only transferable proof disciplines found in the
parallel corpus.
\begin{enumerate}[label=\textup{(\roman*)},leftmargin=3em]
\item Yang--Mills V18--V19 and compact V100 show that a Hermitian pencil
should be normalized by an anchor congruence before its coefficientwise
variation is estimated.  They also show by exact witnesses that locally
valid boxes are not a theorem until their union is proved to cover the
declared parameter space.
\item Compact Yang--Mills V98--V100 shows that a norm must be taken on the
actual constrained range and that a direct or reduced Schur object is
preferable to an auxiliary split whose pieces create a gluing obligation.
Its improved scalar radii still fail to cover the torus.
\item SM--Einstein thirteenth-cycle V99--V100 distinguishes a uniform
inverse margin on a cofinal cell from pointwise invertibility.  Its current
fourteenth-cycle V55--V58 goes further upstream when a required positive
Hessian is absent: it proves the inertia obstruction rather than continuing
to optimize a downstream cross term.
\item Navier--Stokes V26 computes norms on the realized rank-one input
class.  One derivative improves and one does not; restriction must be
calculated, not assumed to help.
\item Parallel Yang--Mills V5 proves that exact local energy control does
not imply a volume-uniform tail for an extensive global sector variable.
This is a warning against promoting local or density information to a
cofinal global statement.
\end{enumerate}
No companion clock, Hessian, kernel constant, torus radius, sector estimate,
or positivity conclusion is arithmetic Weil data.  None is imported.
\end{assessment}

The direct relative-box lemma itself is already proved in the companion
notes and is not duplicated here.  Its GRH use would be legitimate only for
the populated lower Schur endpoint in \textup{(V7.26)}, with a proved cover
and a uniform positive floor.  The next calculation shows why applying that
finite-dimensional architecture before repairing the tail would miss the
active obstruction.

\subsection{A prime cell at the support wall}

For \(0<\delta<2\), let
\[
 (T_\delta u)(t)
 :=\mathbf 1_{\{-1<t<1-\delta\}}u(t+\delta),
 \qquad u\in L^2(-1,1),
 \tag{V10.5}
\]
and put
\[
 U_{\zeta,\delta}:=\zeta T_\delta+\overline\zeta T_\delta^*,
 \qquad |\zeta|=1.
 \tag{V10.6}
\]
For an unramified prime power \(n\), this is the unweighted cell in
\textup{(V9.4)} with
\[
 \delta=\delta_n(A):=\frac{\log n}{A},\qquad
 a=a_n(A):=1-\frac{\delta_n(A)}2,\qquad
 A_n:=\frac12\log n.
 \tag{V10.7}
\]
Thus \(a=(A-A_n)/A\) for \(A>A_n\).  Extend the cell by zero for
\(A\le A_n\).

Let
\[
 I_\delta^-:=(-1,1-\delta),\qquad
 I_\delta^+:=(-1+\delta,1),\qquad
 E_\delta:=I_\delta^-\cup I_\delta^+.
 \tag{V10.8}
\]
When \(1\le\delta<2\), these intervals are disjoint and
\[
 |I_\delta^-|=|I_\delta^+|=2a,qquad |E_\delta|=4a.
 \tag{V10.9}
\]

\begin{theorem}[Strong disappearance but persistent tail norm]
\label{thm:v10-strong-not-norm}
For \(1\le\delta<2\),
\[
 \boxed{\norm{U_{\zeta,\delta}}=1.}
 \tag{V10.10}
\]
As \(\delta\uparrow2\),
\[
 U_{\zeta,\delta}\longrightarrow0
 \quad\hbox{strongly but not in operator norm}.             \tag{V10.11}
\]
If \(P\) is any fixed finite-rank orthogonal projection and \(Q=I-P\),
then
\[
 \boxed{
 \norm{PU_{\zeta,\delta}P}\longrightarrow0,qquad
 \norm{QU_{\zeta,\delta}P}\longrightarrow0,qquad
 \norm{QU_{\zeta,\delta}Q}\longrightarrow1.}
 \tag{V10.12}
\]
Hence a new prime cell can disappear from both the protected block and its
cross block while retaining full norm in the complementary tail.
\end{theorem}

\begin{proof}
Translation identifies \(L^2(I_\delta^+)\) unitarily with
\(L^2(I_\delta^-)\).  On
\(L^2(I_\delta^-)\oplus L^2(I_\delta^+)\), the operator
\(U_{\zeta,\delta}\) is an off-diagonal unitary involution, up to the phases
\(\zeta\) and \(\overline\zeta\), and it vanishes on
\(L^2(E_\delta^c)\).  This proves \textup{(V10.10)}.

For fixed \(u\in L^2(-1,1)\),
\[
 \norm{T_\delta u}^2
 =\int_{I_\delta^+}|u(s)|^2\dd s,
 \qquad
 \norm{T_\delta^*u}^2
 =\int_{I_\delta^-}|u(s)|^2\dd s.                         \tag{V10.13}
\]
Absolute continuity of the integral makes both expressions tend to zero,
so the convergence in \textup{(V10.11)} is strong.  It is not norm
convergence by \textup{(V10.10)}.

Strong convergence is uniform on the unit sphere of a fixed
finite-dimensional range.  Thus
\(
 \norm{U_{\zeta,\delta}P}\to0
\), and self-adjointness gives
\(
 \norm{PU_{\zeta,\delta}}\to0
\).  The first two limits in \textup{(V10.12)} follow.  Finally,
\[
 \norm{U_{\zeta,\delta}-QU_{\zeta,\delta}Q}
 \le
 \norm{PU_{\zeta,\delta}}
 +\norm{QU_{\zeta,\delta}P}\longrightarrow0.             \tag{V10.14}
\]
Combine this with \textup{(V10.10)} to obtain the last limit.
\end{proof}

The theorem explains a feature that looked merely wasteful in V7.  A
finite matrix sees the support length \(a\), but the ambient operator norm
does not.  The absolute tail debit is therefore responding to a real
non-norm-continuity.  Its factor and its failure to use the clock may still
be very wasteful.

\subsection{Quantitative Fourier invisibility}

Let
\[
 \mathcal F_K:=\operatorname{span}\{\phi_k:-K\le k\le K\},
 \qquad \Pi_K:L^2(-1,1)\longrightarrow\mathcal F_K,
 \qquad m:=2K+1.
 \tag{V10.15}
\]

\begin{proposition}[Explicit low/cross loss and tail retention]
\label{prop:v10-Fourier-wall}
For \(1\le\delta<2\),
\[
 \boxed{
 \norm{\Pi_KU_{\zeta,\delta}\Pi_K}_{\rm F}\le2ma,
 \qquad
 \norm{U_{\zeta,\delta}\Pi_K}\le2\sqrt{ma}.}
 \tag{V10.16}
\]
Writing \(Q_K=I-\Pi_K\), one also has
\[
 \boxed{
 \norm{Q_KU_{\zeta,\delta}Q_K}
 \ge1-4\sqrt{ma}.}
 \tag{V10.17}
\]
For the weighted prime cell, multiply the right sides of
\textup{(V10.16)} by
\(
 \lambda_n=\Lambda(n)/\sqrt n
\), while \textup{(V10.17)} becomes a lower bound multiplied by
\(\lambda_n\).
\end{proposition}

\begin{proof}
Formula \textup{(V9.4)} and \(|\operatorname{sinc}x|\le1\) bound every
entry of the \(m\)-dimensional compression by \(2a\).  Its Frobenius norm is
therefore at most \(2ma\).

Each normalized Fourier vector has constant squared modulus \(1/2\).
Equation \textup{(V10.13)} consequently gives
\[
 \norm{T_\delta\phi_k}=\norm{T_\delta^*\phi_k}=\sqrt a,
 \qquad
 \norm{U_{\zeta,\delta}\phi_k}\le2\sqrt a.               \tag{V10.18}
\]
Summing the column squares gives
\(
 \norm{U_{\zeta,\delta}\Pi_K}
 \le\norm{U_{\zeta,\delta}\Pi_K}_{\rm HS}
 \le2\sqrt{ma}
\).
As in \textup{(V10.14)}, self-adjointness yields
\[
 \norm{U_{\zeta,\delta}-Q_KU_{\zeta,\delta}Q_K}
 \le2\norm{U_{\zeta,\delta}\Pi_K}
 \le4\sqrt{ma}.                                           \tag{V10.19}
\]
Reverse triangle inequality and \textup{(V10.10)} prove
\textup{(V10.17)}.
\end{proof}

\begin{corollary}[Piecewise analytic low cells with a linear wall law]
\label{cor:v10-low-cell-analytic}
For fixed \(K\), the matrix
\(
 \Pi_KU_{\chi(n),\delta_n(A)}\Pi_K
\), extended by zero for \(A\le A_n\), is continuous at \(A_n\), analytic
for \(A>A_n\), and obeys
\[
 \norm{\Pi_KU_{\chi(n),\delta_n(A)}\Pi_K}_{\rm F}
 \le2m\frac{A-A_n}{A}.                                     \tag{V10.20}
\]
Inside the active cell its individual matrix entries \(u_{k\ell}^{(n)}\)
satisfy
\[
 \boxed{
 \left|\frac{\dd}{\dd A}u_{k\ell}^{(n)}(A)\right|
 \le \frac{\log n}{A^2}\bigl(1+2\pi K a_n(A)\bigr).}
 \tag{V10.21}
\]
Thus a congruence-normalized interval certificate can control any fixed low
matrix across finitely many support walls.  It still says nothing about the
last quantity in \textup{(V10.12)}.
\end{corollary}

\begin{proof}
The first claims and \textup{(V10.20)} follow from
Proposition~\ref{prop:v10-Fourier-wall}.  For the derivative, put
\[
 b_d(a):=a\operatorname{sinc}(\pi da)
 =\begin{cases}
 a,&d=0,\\
 \sin(\pi da)/(\pi d),&d\ne0.
 \end{cases}                                               \tag{V10.22}
\]
Then \(|b_d(a)|\le a\), \(|b_d'(a)|\le1\), and
\[
 a_n'(A)=\frac{\log n}{2A^2}.
 \tag{V10.23}
\]
The phase in \textup{(V9.4)} has derivative of absolute value at most
\(
 \pi|k+\ell|a_n'(A)
\).  Differentiating
\(
 2b_{\ell-k}(a_n(A))
 \operatorname{Re}[\chi(n)e^{\mathrm i\pi(k+\ell)\delta_n(A)/2}]
\)
and using \(|k+\ell|\le2K\) proves \textup{(V10.21)}.
\end{proof}

There is an analogous precise statement for the V7 correction matrix.  If
\(
 X=Q_K\mathsf V\Pi_K
\) and the \(n\)-cell is omitted from \(X^{(-n)}\), then
\[
 \epsilon_n(A)
 :=\norm{X-X^{(-n)}}
 \le2\lambda_n\sqrt{ma_n(A)}.                             \tag{V10.24}
\]
Consequently
\[
 \boxed{
 \norm{X^*X-(X^{(-n)})^*X^{(-n)}}
 \le2\norm{X^{(-n)}}\epsilon_n(A)+\epsilon_n(A)^2.}
 \tag{V10.25}
\]
This proves continuity of the finite correction at activation, but only
with a square-root wall modulus in the available bound.  Meanwhile the new
cell remains norm one on the tail.  A finite direct-pencil atlas therefore
cannot substitute for a tail theorem.

\subsection{The harmonic clock sees the boundary concentration}

The preceding obstruction is an operator-norm obstruction.  The completed
Weil operator contains the unbounded harmonic clock
\[
 \mathsf De_j=H_je_j,
 \tag{V10.26}
\]
in the normalized Legendre basis.  We now measure the wall cell relative to
that clock.

\begin{lemma}[Endpoint mass versus harmonic energy]
\label{lem:v10-endpoint-clock}
Let \(1\le\delta<2\), let \(a=1-\delta/2\), and let
\(U=U_{\zeta,\delta}\).  For every \(u\) in the quadratic-form domain of
\(\mathsf D\) and every integer \(N\ge1\),
\[
 \boxed{
 |\ip{Uu}{u}|
 \le4aN^2\norm{u}^2+\frac2{H_N}\ip{\mathsf Du}{u}.}
 \tag{V10.27}
\]
If \(\norm u=1\) and \(u\) is supported in \(E_\delta\), then
\[
 \boxed{
 \ip{\mathsf Du}{u}
 \ge H_N(1-2aN^2).}
 \tag{V10.28}
\]
In particular, for \(0<a\le1/64\) and
\(
 N=\lfloor(2\sqrt a)^{-1}\rfloor
\),
\[
 \boxed{
 \ip{\mathsf Du}{u}
 \ge \frac12\log\frac1{4\sqrt a}.}
 \tag{V10.29}
\]
Thus every normalized family supported on the wall pair has harmonic
energy tending to infinity as \(a\downarrow0\).
\end{lemma}

\begin{proof}
For the normalized Legendre functions,
\[
 |e_j(t)|^2\le\frac{2j+1}{2},
 \qquad
 \sum_{j=0}^{N-1}|e_j(t)|^2\le\frac{N^2}{2}.                \tag{V10.30}
\]
Writing \(P_N\) for their first \(N\) modes and using
\(|E_\delta|=4a\), we obtain
\[
 \norm{\mathbf1_{E_\delta}P_Nu}^2
 \le2aN^2\norm u^2.                                      \tag{V10.31}
\]
Also
\[
 \norm{(I-P_N)u}^2
 \le \frac1{H_N}\ip{\mathsf Du}{u}.                     \tag{V10.32}
\]
The operator \(U\) is supported on \(E_\delta\) and has norm one, so
\[
 |\ip{Uu}{u}|
 \le\norm{\mathbf1_{E_\delta}u}^2.
 \tag{V10.33}
\]
Apply \((x+y)^2\le2x^2+2y^2\) to the decomposition
\(
 \mathbf1_Eu=\mathbf1_EP_Nu+\mathbf1_E(I-P_N)u
\), then use \textup{(V10.31)}--\textup{(V10.32)}.  This proves
\textup{(V10.27)}.

If (u) is supported in \(E_\delta\), Bessel's inequality and
\textup{(V10.30)} give
\(
 \norm{P_Nu}^2\le2aN^2
\).  Hence
\(
 \ip{\mathsf Du}{u}\ge H_N\norm{(I-P_N)u}^2
\), which is \textup{(V10.28)}.
For the stated choice of \(N\), \(2aN^2\le1/2\).  Since
\(a\le1/64\), one has
\(
 N\ge(4\sqrt a)^{-1}
\), and \(H_N\ge\log(N+1)\).  Substitution proves
\textup{(V10.29)}.
\end{proof}

\begin{corollary}[A support-aware one-cell form debit]
\label{cor:v10-one-cell-form}
For every unramified \(n\) with
\(
 1\le\delta_n(A)<2
\), every \(N\ge1\), and every form-domain vector \(u\),
\[
 \boxed{
 -\lambda_n\ip{U_{\chi(n),\delta_n(A)}u}{u}
 \ge
 -\frac{2\lambda_n}{H_N}\ip{\mathsf Du}{u}
 -4\lambda_na_n(A)N^2\norm u^2.}
 \tag{V10.34}
\]
Unlike the ambient debit \(2\lambda_n\norm u^2\) used in
\textup{(V7.32)}, the scalar part of \textup{(V10.34)} vanishes linearly at
activation.  The price is a displayed relative expenditure of harmonic
clock.
\end{corollary}

\begin{proof}
Multiply \textup{(V10.27)} by \(\lambda_n\) and take the adverse sign.
\end{proof}

The corollary is not summed over \(n\).  Summing it term by term with a
common \(N\) spends
\(
 2H_N^{-1}\sum_n\lambda_n
\)
of the clock and therefore returns to the exponential absolute-mass loss.
Choosing a separate \(N_n\) creates a genuine multiscale optimization, but
an absolute sum still discards the character phases.  The remaining theorem
must aggregate the wall cells before taking absolute values.

\subsection{What the relative-pencil import can and cannot now do}

On a closed interval containing only finitely many support walls, the low
Fourier matrix is continuous and piecewise analytic by
Corollary~\ref{cor:v10-low-cell-analytic}.  If a populated Hermitian lower
Schur endpoint \(B_K(A)\) and an invertible anchor \(R\) satisfy
\[
 R^*B_K(A_0)R\succeq\beta I,
 \qquad R^*R\preceq\gamma I,
 \tag{V10.35}
\]
then the companion relative method would certify
\[
 B_K(A)\succeq\frac{\beta-\eta}{\gamma}I                    \tag{V10.36}
\]
whenever a rigorous coefficientwise or interval calculation proves
\[
 \norm{R^*(B_K(A)-B_K(A_0))R}_{\rm op}\le\eta<\beta.       \tag{V10.37}
\]
This is simply the already-audited congruence argument, specialized to the
GRH matrix; no new generic lemma is claimed.

The V7 endpoint cannot yet be used cofinally in this way.  Its tail scalar
\(d_N\) is obtained by an unweighted operator-norm sum, and positivity of
that scalar already forces the double-exponential dimension in
\textup{(V7.35)}.  Theorem~\ref{thm:v10-strong-not-norm} proves why a fixed
finite compression cannot make the tail cells disappear in operator norm.
Lemma~\ref{lem:v10-endpoint-clock} proves, equally importantly, that the
unweighted norm is the wrong final geometry: the corresponding edge states
pay divergent harmonic energy.

\begin{firewall}[Finite atlases do not see the cofinal tail]
\label{fire:v10-atlas}
A positive relative box for a finite Fourier or Legendre compression proves
only that finite matrix statement.  Even vanishing low and cross blocks at
a support wall coexist with tail norm tending to one by
\textup{(V10.12)}.  Conversely, the norm-one wall cell does not prove a
negative direction for the completed operator, because its localized
vectors pay the harmonic cost \textup{(V10.29)}.  Neither side may be
discarded.
\end{firewall}

\subsection{V10 acquisition and upstream direction}

\begin{center}
\begin{tabularx}{\textwidth}{@{}p{0.29\textwidth}p{0.17\textwidth}X@{}}
\toprule
V10 component & Status & Exact conclusion \\
\midrule
Parallel corpus sweep & completed & Twenty-nine non-GRH notes were
inventoried; their terminal results and live continuations were classified,
with exact checkpoint provenance. \\
Relative-pencil import & delimited & Anchor congruence, constrained-range
norms, and cover audits transfer as proof methods; no companion constant or
physical conclusion transfers. \\
Support-wall operator & proved & A new prime cell tends strongly to zero but
retains norm one, and every fixed finite tail retains that norm. \\
Fourier visibility & quantified & The low block is \(O_K(a)\), the cross
block is \(O_K(\sqrt a)\), and the tail norm is
\(1-O_K(\sqrt a)\). \\
Clock compensation & proved & Endpoint-supported unit vectors have harmonic
energy at least logarithmic in \(a^{-1}\), and \textup{(V10.34)} is an
explicit one-cell relative-form debit. \\
Cofinal aggregation & open & The one-cell debit has not been summed with
character-sensitive cancellation and does not prove tail coercivity. \\
GRH & open & No cofinal nonnegativity theorem for the completed Weil operator
has been proved. \\
\bottomrule
\end{tabularx}
\end{center}

The next target is no longer a finite low-mode atlas.  For the boundary
family
\[
 \mathsf B_{\chi,A}
 :=\sum_{e^A\le n<e^{2A}}
 \lambda_nU_{\chi(n),\delta_n(A)},                         \tag{V10.38}
\]
one needs a character-sensitive estimate of the form
\[
 \boxed{
 |\ip{\mathsf B_{\chi,A}u}{u}|
 \le\varepsilon_A\ip{\mathsf Du}{u}
 +C_A\norm u^2,}                                          \tag{V10.39}
\]
with a clock expenditure \(\varepsilon_A<1\) and a scalar remainder that
can be paid by the completed archimedean floor.  The cells must be summed
before absolute values are taken.  If no such character-sensitive
aggregation can be proved, the correct upstream move is to replace the
sharp support wall by a positive smooth localization whose Mellin and
Fourier geometry is compatible with a relative-form large sieve; it is not
to enumerate more finite boxes.

\begin{assessment}[V10 acquisition]
The compulsory sweep has produced a genuine change of direction rather than
an imported analogy.  The support-wall theorem isolates a compactness trap:
fixed finite matrices lose a newly activated prime cell while the tail keeps
its full norm.  The endpoint-clock lemma then shows why the ambient norm is
too pessimistic: the escaping norm-one states concentrate, and the harmonic
clock detects that concentration quantitatively.  Equation
\textup{(V10.34)} is the first support-aware relative-form replacement for
one term of the absolute-mass tail debit.  Its character-sensitive cofinal
aggregation remains open, so GRH is not proved.
\end{assessment}

\section*{Version V10 append-only record}
\addcontentsline{toc}{section}{Version V10 append-only record}
The complete V9 source was retained before this continuation.  It had
\(155710\) bytes, \(4712\) lines, and SHA--256
\[
 \texttt{43B01F624FAA9CD388ABB18630358966C488678B99CAD6F180EEE9A69CBE70F9}.
\]
V10 completed the mandated corpus sweep, proved the support-wall strong/not-
norm dichotomy, quantified finite Fourier invisibility, and derived the
one-cell harmonic-clock relative-form bound.

% =============================================================================
% END APPENDED V10 MATERIAL
% =============================================================================

% =============================================================================
% BEGIN APPENDED V11 MATERIAL
% =============================================================================

\section{V11: magnetic renewal energy and a clock-restricted gap}
\label{sec:grh-v11}

V10 proved that a prime cell entering through the support wall is invisible
to every fixed finite compression while retaining full norm in the tail.  It
also proved that the corresponding norm-one vectors pay divergent harmonic
energy.  The next question is therefore not whether the character phases
reduce the ambient operator norm---V8 proved that they do not---but whether
the prime cells possess a positive form component that can be coupled to the
harmonic clock before absolute values are taken.

They do.  The prime operator is exactly a character-twisted graph adjacency.
Completing its squares produces a nonnegative magnetic renewal energy and an
explicit multiplication degree.  For every nonprincipal primitive
character, two fixed unramified prime shifts make the magnetic energy
injective once the support wall is large enough.  It has no unconstrained
spectral gap, by the already-proved Kronecker resonance, but it has a
strictly positive minimum on every bounded harmonic-energy ball.  We also
give a finite Legendre lower certificate for that minimum.  This is the
first sign-sensitive aggregate statement in the programme that survives the
V8 norm no-go.

The manuscript already proves the harmonic Dirichlet-form identity
\[
 \ip{\mathsf Du}{u}
 =\frac14\int_{-1}^1\int_{-1}^1
   \frac{|u(t)-u(v)|^2}{|t-v|}\dd t\dd v
 \tag{V11.1}
\]
and its Legendre diagonalization
\(\mathsf De_j=H_je_j\).  Those results are used, not reproved.  The new
content below is the magnetic square completion, its zero-mode rigidity,
and the clock-constrained finite acquisition.

\subsection{Exact magnetic square completion}

Keep the active history
\[
 \mathcal A_f(A)
 =\{n:\ 2\le n<e^{2A},\
       n\text{ an unramified prime power}\}
 \tag{V11.2}
\]
and write
\[
 \lambda_n=\frac{\Lambda(n)}{\sqrt n},\qquad
 \delta_n=\frac{\log n}{A},\qquad
 a_n=1-\frac{\delta_n}{2}.
 \tag{V11.3}
\]
For \(u\in L^2(-1,1)\), define the \(n\)-cell magnetic energy
\[
 \boxed{
 \mathcal E_{\chi,n,A}[u]
 :=\int_{-1}^{1-\delta_n}
 \left|u(t+\delta_n)-\overline{\chi(n)}u(t)\right|^2\dd t.}
 \tag{V11.4}
\]
Its aggregate is the bounded nonnegative quadratic form
\[
 \boxed{
 \mathcal G_{\chi,A}[u]
 :=\sum_{n\in\mathcal A_f(A)}
   \lambda_n\mathcal E_{\chi,n,A}[u].}
 \tag{V11.5}
\]
Let
\[
 d_{n,A}(t)
 :=\mathbf1_{\{-1<t<1-\delta_n\}}
   +\mathbf1_{\{-1+\delta_n<t<1\}},
 \qquad
 V_A(t):=\sum_{n\in\mathcal A_f(A)}\lambda_nd_{n,A}(t).
 \tag{V11.6}
\]
The function \(V_A\) is the weighted incidence degree of the renewal graph;
it is independent of the character.

\begin{theorem}[Magnetic renewal decomposition]
\label{thm:v11-magnetic-decomposition}
For every \(u\in L^2(-1,1)\),
\[
 \boxed{
 -\ip{\mathsf E_{\chi,A}u}{u}
 =\mathcal G_{\chi,A}[u]
  -\int_{-1}^1V_A(t)|u(t)|^2\dd t.}
 \tag{V11.7}
\]
Consequently the completed localized Weil form is exactly
\[
 \boxed{
 \begin{aligned}
 \ip{\mathsf H_{\chi,A}^{\rm sc}u}{u}
 ={}&\ip{\mathsf Du}{u}+\mathcal G_{\chi,A}[u]\\
 &+\int_{-1}^1
 \bigl(c_{f,A}-\tfrac12\log(1-t^2)-V_A(t)\bigr)
 |u(t)|^2\dd t\\
 &+\ip{\mathsf S_{\alpha,A}^{\rm sc}u}{u}.
 \end{aligned}}
 \tag{V11.8}
\]
No prime phase or off-diagonal shift has been discarded in this identity.
\end{theorem}

\begin{proof}
For one active \(n\), expansion of the square gives
\[
 \begin{aligned}
 \mathcal E_{\chi,n,A}[u]
 ={}&
 \int_{-1}^{1-\delta_n}
 \bigl(|u(t+\delta_n)|^2+|u(t)|^2\bigr)\dd t\\
 &-2\operatorname{Re}\left[
 \chi(n)\int_{-1}^{1-\delta_n}
 u(t+\delta_n)\overline{u(t)}\dd t\right].
 \end{aligned}
 \tag{V11.9}
\]
The first line is
\(\int d_{n,A}(t)|u(t)|^2\dd t\).  The second line is the negative of
\(\ip{\mathsf P_{\chi,A}^{\rm sc}(n)u}{u}\), by the prime-cell convention
in \textup{(V8.4)} and \textup{(V10.5)}--\textup{(V10.6)}.  Multiply by
\(\lambda_n\), sum the finite history, and rearrange to prove
\textup{(V11.7)}.  Substitution into \textup{(V7.23)} gives
\textup{(V11.8)}.
\end{proof}

For the boundary family
\[
 \mathcal B_f(A)
 :=\{n\in\mathcal A_f(A):e^A\le n<e^{2A}\},
 \tag{V11.10}
\]
the two incidence intervals are disjoint.  Put
\[
 V_A^\partial(t)
 :=\sum_{n\in\mathcal B_f(A)}\lambda_nd_{n,A}(t).
 \tag{V11.11}
\]

\begin{proposition}[Exact boundary degree]
\label{prop:v11-boundary-degree}
For almost every \(t\in(-1,1)\),
\[
 \boxed{
 V_A^\partial(t)
 =
 \sum_{\substack{e^A\le n<e^{A(1+|t|)}\\
                  n\ {\rm unramified\ prime\ power}}}
 \frac{\Lambda(n)}{\sqrt n}.}
 \tag{V11.12}
\]
For every fixed \(t\) with \(0<|t|\le1\), the prime number theorem gives
\[
 \boxed{
 V_A^\partial(t)
 =
 2e^{A(1+|t|)/2}-2e^{A/2}
 +o_f\!\left(e^{A(1+|t|)/2}\right)}
 \tag{V11.13}
\]
as \(A\to\infty\).
\end{proposition}

\begin{proof}
If \(t<0\), membership in the left incidence interval is
\(\delta_n<1-t=1+|t|\); if \(t>0\), membership in the right interval is
\(\delta_n<1+t=1+|t|\).  For \(\delta_n\ge1\), the other interval does not
contain \(t\).  This proves \textup{(V11.12)} away from endpoints of
intervals, a null set.  Partial summation of
\(\sum_{n\le x}\Lambda(n)=x+o(x)\) gives
\[
 \sum_{n\le x}\frac{\Lambda(n)}{\sqrt n}
 =2\sqrt x+o(\sqrt x).
 \tag{V11.14}
\]
Apply this at \(x=e^{A(1+|t|)}\) and \(x=e^A\).  Removing the finitely many
ramified prime towers contributes only \(O_f(A)\), proving
\textup{(V11.13)}.
\end{proof}

The exponentially large degree in \textup{(V11.13)} shows why dropping
\(\mathcal G_{\chi,A}\) cannot improve the V7 tail estimate.  The positive
magnetic squares carry the same scale and must remain paired with the
degree.

\subsection{Two-shift zero-mode rigidity}

We first isolate the finite-interval translation fact needed to prove that
the magnetic squares have no common nonzero zero mode.

\begin{lemma}[Local quasiperiods force a common frequency]
\label{lem:v11-local-quasiperiods}
Let \(0<\rho,\sigma<2\) with \(\rho+\sigma\le2\), and let
\(c_\rho,c_\sigma\in\mathbb T\).  Suppose there is no
\(\xi\in\mathbb R\) satisfying
\[
 e^{\mathrm i\xi\rho}=c_\rho,\qquad
 e^{\mathrm i\xi\sigma}=c_\sigma.
 \tag{V11.15}
\]
If \(u\in L^2(-1,1)\) obeys
\[
 \begin{aligned}
 u(t+\rho)&=c_\rho u(t)
 &&\text{for a.e. }t\in(-1,1-\rho),\\
 u(t+\sigma)&=c_\sigma u(t)
 &&\text{for a.e. }t\in(-1,1-\sigma),
 \end{aligned}
 \tag{V11.16}
\]
then \(u=0\).
\end{lemma}

\begin{proof}
Use the first relation to extend \(u\) to an
\(\rho\)-quasiperiodic function \(\widetilde u\in L^2_{\rm loc}(\mathbb R)\):
\[
 \widetilde u(x+\rho)=c_\rho\widetilde u(x).
 \tag{V11.17}
\]
The extension agrees with \(u\) on \((-1,1)\).  Define
\[
 h(x):=\widetilde u(x+\sigma)-c_\sigma\widetilde u(x).
 \tag{V11.18}
\]
Then \(h(x+\rho)=c_\rho h(x)\), and the second relation in
\textup{(V11.16)} says that \(h\) vanishes on an interval of length
\(2-\sigma\ge\rho\).  Every class modulo \(\rho\) has a representative in
that interval, so quasiperiodicity makes \(h=0\) almost everywhere on
\(\mathbb R\).

Choose \(\theta\) with \(e^{\mathrm i\theta\rho}=c_\rho\).  The function
\[
 v(x):=e^{-\mathrm i\theta x}\widetilde u(x)
 \tag{V11.19}
\]
is \(\rho\)-periodic and has an \(L^2\) Fourier series.  Its frequencies in
\(\widetilde u\) are
\[
 \xi_k=\theta+\frac{2\pi k}{\rho},\qquad k\in\mathbb Z.
 \tag{V11.20}
\]
The global \(\sigma\)-relation forces every nonzero Fourier coefficient to
satisfy \(e^{\mathrm i\xi_k\sigma}=c_\sigma\).  Together with the definition
of \(\xi_k\), this is precisely \textup{(V11.15)}, which has no solution.
All Fourier coefficients vanish, hence \(u=0\).
\end{proof}

\begin{theorem}[Injectivity of the nonprincipal magnetic energy]
\label{thm:v11-magnetic-injective}
Let \(\chi\) be a nonprincipal primitive Dirichlet character.  There exist
two distinct primes \(p,q\nmid f\), with \(\chi(p)\ne1\), such that for every
\[
 A>\frac12\log(pq)
 \tag{V11.21}
\]
the bounded positive operator associated with
\(\mathcal G_{\chi,A}\) has trivial kernel:
\[
 \boxed{\mathcal G_{\chi,A}[u]=0\quad\Longrightarrow\quad u=0.}
 \tag{V11.22}
\]
\end{theorem}

\begin{proof}
Because \(\chi\) is nonprincipal, choose an integer \(b\), coprime to \(f\),
with \(\chi(b)\ne1\).  In the prime factorization of \(b\), at least one
prime \(p\nmid f\) satisfies \(\chi(p)\ne1\).  Choose any distinct prime
\(q\nmid f\).  Both are active when \textup{(V11.21)} holds, and
\[
 \rho:=\frac{\log p}{A},\qquad
 \sigma:=\frac{\log q}{A}
 \tag{V11.23}
\]
satisfy \(\rho+\sigma<2\).

We claim that no \(\xi\) satisfies
\[
 e^{\mathrm i\xi\rho}=\overline{\chi(p)},\qquad
 e^{\mathrm i\xi\sigma}=\overline{\chi(q)}.
 \tag{V11.24}
\]
Let \(r\) be the finite order of \(\chi\).  If \textup{(V11.24)} held,
raising both equations to the \(r\)-th power would give integers \(k,\ell\)
such that
\[
 \frac{r\xi}{A}\log p=2\pi k,\qquad
 \frac{r\xi}{A}\log q=2\pi\ell.
 \tag{V11.25}
\]
If \(\xi\ne0\), this makes \(\log p/\log q=k/\ell\) rational, contradicting
unique prime factorization after exponentiation.  If \(\xi=0\), both
character values in \textup{(V11.24)} are one, contradicting
\(\chi(p)\ne1\).

If \(\mathcal G_{\chi,A}[u]=0\), every nonnegative cell in the finite sum
\textup{(V11.5)} vanishes.  The \(p\)- and \(q\)-cells give
\textup{(V11.16)} with the phases in \textup{(V11.24)}.  Apply
Lemma~\ref{lem:v11-local-quasiperiods}.
\end{proof}

\begin{firewall}[The principal character is a different block]
\label{fire:v11-principal}
For the primitive principal character of conductor one, constants are
common zero modes of all magnetic differences.  Theorem
\ref{thm:v11-magnetic-injective} is therefore deliberately restricted to
nonprincipal characters.  The zeta block must retain its separate
archimedean and pole terms; no nonprincipal magnetic gap is transferred to
it.
\end{firewall}

\subsection{Injective but not bounded below}

For the normalized modulations \(g_{A\tau}\) of \textup{(V8.5)}, direct
substitution in \textup{(V11.4)} gives
\[
 \boxed{
 \mathcal G_{\chi,A}[g_{A\tau}]
 =
 \sum_{n\in\mathcal A_f(A)}
 \lambda_na_n
 \left|\chi(n)e^{\mathrm i\tau\log n}-1\right|^2.}
 \tag{V11.26}
\]

\begin{proposition}[The magnetic operator has zero global gap]
\label{prop:v11-zero-global-gap}
For every fixed \(A\) and every finite Legendre projection \(P_N\), with
\(Q_N=I-P_N\),
\[
 \boxed{
 \inf_{\norm u=1}\mathcal G_{\chi,A}[u]=0,\qquad
 \inf_{\substack{\norm u=1\\u\in\operatorname{Ran}Q_N}}
 \mathcal G_{\chi,A}[u]=0.}
 \tag{V11.27}
\]
Thus the injectivity in \textup{(V11.22)} does not furnish an
unconstrained coercive constant.
\end{proposition}

\begin{proof}
Use the resonant sequence \(\tau_j\to\infty\) from
Theorem~\ref{thm:v8-Kronecker}.  Every summand in \textup{(V11.26)} tends to
zero, so the first infimum is zero.  By \textup{(V8.8)},
\[
 h_j:=\frac{Q_Ng_{A\tau_j}}{\norm{Q_Ng_{A\tau_j}}}
 \longrightarrow g_{A\tau_j}
 \quad\text{in norm}.
 \tag{V11.28}
\]
The magnetic operator is a finite bounded sum, so its quadratic form on
\(h_j\) also tends to zero.  This proves the second equality.
\end{proof}

This proposition is the positive-square version of the V8 tail-resonance
no-go.  Character phases do create rigidity, but only after high-clock
approximate gauges have been excluded.

\subsection{A strictly positive clock-restricted magnetic gap}

For \(E\ge0\), define the clock sublevel minimum
\[
 \boxed{
 \Gamma_{\chi,A}(E)
 :=
 \inf\left\{
 \mathcal G_{\chi,A}[u]:
 \norm u=1,\ 
 u\in\mathcal Q(\mathsf D),\
 \ip{\mathsf Du}{u}\le E
 \right\}.}
 \tag{V11.29}
\]

\begin{theorem}[Clock-restricted magnetic gap]
\label{thm:v11-clock-gap}
Under the hypotheses and wall condition of
Theorem~\ref{thm:v11-magnetic-injective}, one has
\[
 \boxed{\Gamma_{\chi,A}(E)>0\qquad\text{for every finite }E.}
 \tag{V11.30}
\]
At the same time,
\[
 \boxed{\Gamma_{\chi,A}(E)\downarrow0\quad\text{as }E\to\infty.}
 \tag{V11.31}
\]
\end{theorem}

\begin{proof}
Since \(H_j\to\infty\), the unit ball of the form domain with bounded
\(\mathsf D\)-energy is relatively compact in \(L^2(-1,1)\).  The set in
\textup{(V11.29)} is closed under \(L^2\) limits with uniformly bounded
energy, by lower semicontinuity of the closed form, and is therefore
compact.  The bounded quadratic form \(\mathcal G_{\chi,A}\) attains its
minimum there.  A zero minimum would give a unit vector in its kernel,
contradicting \textup{(V11.22)}.  This proves \textup{(V11.30)}.

The function \(\Gamma_{\chi,A}\) is nonincreasing.  Equation
\textup{(V11.26)} and the resonant sequence in the proof of
Proposition~\ref{prop:v11-zero-global-gap} give unit form-domain vectors
whose magnetic energy tends to zero, each with some finite harmonic energy.
Letting the allowed energy pass through those values proves
\textup{(V11.31)}.
\end{proof}

The theorem is qualitative but not nonconstructive in the sense of the
programme: the gap has a convergent finite-dimensional lower-certificate
scheme.

Let \(P_N\) be the first-\(N\) Legendre projection.  For \(E'\ge0\), define
the finite trust-region value
\[
 \Gamma_{\chi,A,N}(E')
 :=
 \min\left\{
 \mathcal G_{\chi,A}[v]:
 v\in\operatorname{Ran}P_N,\ 
 \norm v=1,\
 \ip{\mathsf Dv}{v}\le E'
 \right\}.
 \tag{V11.32}
\]
It is the minimum of an explicit Hermitian matrix on the intersection of a
sphere and an ellipsoid.  All of its entries are acquired by the polynomial
translation formula \textup{(V7.4)}.

\begin{theorem}[Finite Legendre lower certificate]
\label{thm:v11-finite-gap-certificate}
Put
\[
 \mathcal L_f(A)
 :=\sum_{n\in\mathcal A_f(A)}\lambda_n.
 \tag{V11.33}
\]
The magnetic operator obeys
\[
 \norm{\mathcal G_{\chi,A}}\le4\mathcal L_f(A).
 \tag{V11.34}
\]
Fix \(E<\infty\) and choose \(N\) with \(H_N>E\).  Set
\[
 r_N:=\sqrt{\frac{E}{H_N}}<1.
 \tag{V11.35}
\]
Then
\[
 \boxed{
 \Gamma_{\chi,A}(E)
 \ge
 (1-r_N^2)
 \Gamma_{\chi,A,N}\!\left(\frac{E}{1-r_N^2}\right)
 -8\mathcal L_f(A)r_N.}
 \tag{V11.36}
\]
For every fixed \(A,E\) satisfying
\textup{(V11.21)}, the right side of \textup{(V11.36)} is positive for all
sufficiently large \(N\).
\end{theorem}

\begin{proof}
For one cell, the map
\[
 u\longmapsto
 \mathbf1_{(-1,1-\delta_n)}
 \bigl(u(\,\cdot+\delta_n)-\overline{\chi(n)}u\bigr)
 \tag{V11.37}
\]
has norm at most two.  Its squared operator therefore has norm at most four.
Summation proves \textup{(V11.34)}.

Take an admissible unit vector \(u\), and write
\[
 p=P_Nu,\qquad q=Q_Nu.
 \tag{V11.38}
\]
The harmonic diagonal gives
\[
 \norm q^2\le\frac{E}{H_N}=r_N^2.
 \tag{V11.39}
\]
Thus \(\norm p^2\ge1-r_N^2\).  If \(v=p/\norm p\), then
\[
 \ip{\mathsf Dv}{v}
 \le\frac{E}{1-r_N^2},
 \tag{V11.40}
\]
and hence
\[
 \mathcal G_{\chi,A}[p]
 \ge
 (1-r_N^2)
 \Gamma_{\chi,A,N}\!\left(\frac{E}{1-r_N^2}\right).
 \tag{V11.41}
\]
Positivity of the magnetic operator and the cross-term estimate give
\[
 \mathcal G_{\chi,A}[u]
 \ge\mathcal G_{\chi,A}[p]
 -2\norm{\mathcal G_{\chi,A}}\norm p\norm q.
 \tag{V11.42}
\]
Use \(\norm p\le1\), \(\norm q\le r_N\), and
\textup{(V11.34)} to prove \textup{(V11.36)}.

Finally, the compactness argument in
Theorem~\ref{thm:v11-clock-gap} implies convergence of the truncated
trust-region problems to the positive value
\(\Gamma_{\chi,A}(E)\).  The tail norm in \textup{(V11.39)} tends to zero.
Thus a finite interval computation using the exact low and cross matrices
eventually produces a positive lower enclosure.  The displayed
\textup{(V11.36)} is one sufficient enclosure, not a claim that its crude
factor \(8\mathcal L_f(A)\) is cofinally efficient.
\end{proof}

\begin{firewall}[Fixed-wall acquisition is not cofinal control]
\label{fire:v11-fixed-wall}
Theorems~\ref{thm:v11-clock-gap} and
\ref{thm:v11-finite-gap-certificate} solve a fixed-\(A\), bounded-clock
compactness problem.  They do not give a useful lower asymptotic for
\(\Gamma_{\chi,A}(E)\) as \(A\) and \(E\) grow.  The crude error in
\textup{(V11.36)} contains the same absolute mass that caused the V7
double-exponential dimension.  A numerical certificate at isolated walls,
or a finite list of relative boxes, would not prove the required cofinal
uniform statement.
\end{firewall}

\subsection{The new cofinal scalar and the V12 target}

Equation \textup{(V11.8)} shows that the relevant competition is not
\(\mathcal G_{\chi,A}\) alone but its cancellation of the incidence degree.
Define the clock-restricted prime-adjacency profile
\[
 \boxed{
 \mathfrak A_{\chi,A}(E)
 :=
 \sup\left\{
 \int_{-1}^1V_A(t)|u(t)|^2\dd t
 -\mathcal G_{\chi,A}[u]:
 \norm u=1,\ \ip{\mathsf Du}{u}\le E
 \right\}.}
 \tag{V11.43}
\]
By \textup{(V11.7)}, this is exactly the largest prime-operator Rayleigh
quotient on the harmonic-energy sublevel; it retains every character phase
and every off-diagonal shift.  The completed form on that sublevel obeys
\[
 \ip{\mathsf H_{\chi,A}^{\rm sc}u}{u}
 \ge
 \ip{\mathsf Du}{u}
 +c_{f,A}
 +\inf_t\!\left[-\frac12\log(1-t^2)\right]
 -\mathfrak A_{\chi,A}(E)
 -\norm{\mathsf S_{\alpha,A}^{\rm sc}},
 \tag{V11.44}
\]
where the displayed logarithmic infimum is zero.  This bound is deliberately
coarse only in the archimedean multiplication row; unlike V7, its prime
quantity \(\mathfrak A_{\chi,A}(E)\) is sign sensitive and clock
restricted.

\begin{assessment}[V12 attacks the quantitative rigidity profile]
\label{dec:v11-next}
The next version should make the two-shift rigidity proof quantitative.  Use
the \(\rho\)-quasiperiodic extension in
Lemma~\ref{lem:v11-local-quasiperiods} to reduce approximate common zero
modes to the inhomogeneous Diophantine profile
\[
 \min_{k\in\mathbb Z}
 \left|
 e^{\mathrm i(\theta+2\pi k/\rho)\sigma}
 -c_\sigma
 \right|.
 \tag{V11.45}
\]
Then couple the allowed Fourier indices to harmonic energy, rather than
taking an unrestricted infimum that V8 proves is zero.  The first acceptable
output is an explicit lower bound for
\(\Gamma_{\chi,A}(E)\) or a proved asymptotic obstruction showing that two
prime shifts are too weak.  Only after that decision should the full degree
profile \(\mathfrak A_{\chi,A}(E)\) be attacked.  Do not return to an
unconstrained prime-operator norm or a diagonal density count.
\end{assessment}

\begin{center}
\begin{tabularx}{\textwidth}{@{}p{0.29\textwidth}p{0.17\textwidth}X@{}}
\toprule
V11 component & Status & Exact conclusion \\
\midrule
Prime square completion & proved & The prime adjacency equals incidence
degree minus nonnegative magnetic renewal energy. \\
Boundary degree & acquired & Formula \textup{(V11.12)} and its PNT scale
\textup{(V11.13)} are explicit. \\
Magnetic zero modes & excluded for nonprincipal \(\chi\) & Two fixed prime
shifts and local quasiperiod rigidity give a trivial kernel for large walls. \\
Unrestricted magnetic gap & disproved & Kronecker resonance makes the
infimum zero even on every fixed Legendre tail. \\
Clock-restricted gap & proved positive & Every finite harmonic-energy
sublevel has a positive magnetic minimum. \\
Finite acquisition & proved & Equation \textup{(V11.36)} is a rigorous
Legendre lower certificate, though not cofinally efficient. \\
Cofinal quantitative profile and GRH & open & No uniform asymptotic for
\(\mathfrak A_{\chi,A}(E)\) or completed positivity has been proved. \\
\bottomrule
\end{tabularx}
\end{center}

\begin{assessment}[V11 acquisition]
V11 replaces the proposed character-sensitive bound by the exact object it
must control.  The magnetic renewal form is positive and injective for every
nonprincipal primitive character at sufficiently large support walls, but
it has no global gap because arbitrarily costly harmonic modulations
approximately gauge away every finite character history.  Bounded harmonic
energy restores a strictly positive and finitely certifiable gap.  The
remaining obstruction is quantitative: the gap must be compared with the
explicit incidence degree on a cofinal clock scale.  GRH is not proved.
\end{assessment}

\section*{Version V11 append-only record}
\addcontentsline{toc}{section}{Version V11 append-only record}
The complete V10 source was retained before this continuation.  It had
\(177229\) bytes, \(5317\) lines, and SHA--256
\[
 \texttt{6D8761CE74890D6B0766E5D3324B040CE3F641D90E0D0A63478F69288CBAF06A}.
\]
V11 derived the magnetic renewal decomposition, proved two-prime zero-mode
rigidity for nonprincipal characters, separated injectivity from the V8
zero-gap resonance, and established a positive clock-restricted gap with a
finite Legendre lower certificate.

% =============================================================================
% END APPENDED V11 MATERIAL
% =============================================================================

% =============================================================================
% BEGIN APPENDED V12 MATERIAL
% =============================================================================

\section{V12: quantitative phase rigidity and a square-root resonance clock}
\label{sec:grh-v12}

V11 proved that two fixed prime shifts make the nonprincipal magnetic renewal
energy injective, while V8 resonance makes its unrestricted gap zero.  The
present version quantifies that tension and then tests whether the two-prime
mechanism has the correct cofinal scale.

It does not.  The two-prime phase mismatch has an explicit, but extremely
weak, exponential lower bound on a finite modulation band.  More decisively,
the total energy of two fixed cells is bounded independently of the support
wall, whereas the full renewal degree on the constant zero-clock vector is
asymptotic to \(4e^A/A\).  Thus the V11 injectivity proof is a valid
fixed-wall compactness argument but cannot be the cofinal payment mechanism.

We then restore the entire character reader before estimating it.  The
classical fixed-character zero-free region stated among the manuscript's
analytic inputs yields an explicit bound for the triangular prime
polynomial.  It proves that an approximate simultaneous resonance of the
full active history cannot occur before frequency
\(\exp(c_\chi\sqrt A)\), and therefore cannot occur below a harmonic clock
of order \(\sqrt A\).  This is a genuine quantitative improvement over bare
Kronecker recurrence, but its remaining exponential prime bound is still far
larger than the archimedean clock.

\subsection{The exact two-log inhomogeneous profile}

Let \(\chi\) have finite order \(r\), and choose the fixed primes \(p,q\)
from Theorem~\ref{thm:v11-magnetic-injective}, with \(\chi(p)\ne1\).  Write
\[
 \chi(p)=e^{2\pi\mathrm i a_p/r},\qquad
 \chi(q)=e^{2\pi\mathrm i a_q/r},
 \qquad 0\le a_p,a_q<r,
 \tag{V12.1}
\]
and put
\[
 \omega:=\frac{\log q}{\log p}\notin\mathbb Q.
 \tag{V12.2}
\]
The frequencies which align the \(p\)-cell exactly are
\[
 \boxed{
 \tau_k:=\frac{2\pi}{\log p}
 \left(k-\frac{a_p}{r}\right),\qquad k\in\mathbb Z.}
 \tag{V12.3}
\]
At such a frequency, the residual \(q\)-phase is
\[
 \boxed{
 \chi(q)e^{\mathrm i\tau_k\log q}
 =e^{2\pi\mathrm i\eta_k},\qquad
 \eta_k:=\frac{a_q}{r}
 +\left(k-\frac{a_p}{r}\right)\omega.}
 \tag{V12.4}
\]
Thus the profile proposed in \textup{(V11.45)} is precisely the
inhomogeneous rotation distance \(\|\eta_k\|_{\mathbb R/\mathbb Z}\).

\begin{lemma}[Elementary quantitative separation of the two logarithms]
\label{lem:v12-log-separation}
There are explicit constants \(c_{p,q,r}>0\) and \(C_{p,q,r}>0\) such that
\[
 \boxed{
 \|\eta_k\|_{\mathbb R/\mathbb Z}
 \ge c_{p,q,r}e^{-C_{p,q,r}|k|}
 \qquad(k\in\mathbb Z).}
 \tag{V12.5}
\]
For example one may take
\[
 C_{p,q,r}=2r\log q
 \tag{V12.6}
\]
after decreasing the positive prefactor to cover the finitely many small
\(k\).
\end{lemma}

\begin{proof}
Let \(\ell\) be a nearest integer to \(\eta_k\).  Multiplication by
\(r\log p\) gives the exact nonzero logarithmic form
\[
 \begin{aligned}
 r\log p\,(\eta_k-\ell)
 &=(rk-a_p)\log q-(r\ell-a_q)\log p\\
 &=\log\left(\frac{q^{\,rk-a_p}}{p^{\,r\ell-a_q}}\right).
 \end{aligned}
 \tag{V12.7}
\]
It is nonzero.  Indeed, unique factorization would force
\(rk-a_p=r\ell-a_q=0\), while \(0<a_p<r\) makes the first equality
impossible for an integer \(k\).

If \(X/Y\ne1\) is a positive rational in lowest terms, then
\[
 \left|\log\frac XY\right|
 \ge\frac1{\max\{X,Y\}}.
 \tag{V12.8}
\]
Indeed, if \(X>Y\), then
\(\log(X/Y)\ge(X-Y)/X\ge1/X\), and the other case is symmetric.
Applying this after moving negative exponents in \textup{(V12.7)} between
numerator and denominator gives
\[
 |r\log p\,(\eta_k-\ell)|
 \ge
 \exp\left(
 -|rk-a_p|\log q-|r\ell-a_q|\log p
 \right).
 \tag{V12.9}
\]
Since \(\ell\) is nearest to \(\eta_k\),
\[
 |\ell|
 \le |k|\omega+\omega+\frac32.
 \tag{V12.10}
\]
Substitution in \textup{(V12.9)} bounds the exponent by
\[
 2r|k|\log q+
 r\log q+r\left(\omega+\frac52\right)\log p.
 \tag{V12.11}
\]
Division by \(r\log p\) proves \textup{(V12.5)} with the displayed
\textup{(V12.6)} and an explicit positive prefactor.
\end{proof}

The lemma can be made uniform over all frequencies in a compact band, not
only over the exact \(p\)-alignment lattice.

\begin{theorem}[Finite-band two-prime phase floor]
\label{thm:v12-two-prime-band}
There are explicit \(c_1,c_2>0\), depending only on \(p,q,\chi\), such that
for every \(T\ge1\),
\[
 \boxed{
 \inf_{|\tau|\le T}
 \max\left\{
 |\chi(p)e^{\mathrm i\tau\log p}-1|,
 |\chi(q)e^{\mathrm i\tau\log q}-1|
 \right\}
 \ge c_1e^{-c_2T}.}
 \tag{V12.12}
\]
Consequently, whenever \(A>\frac12\log(pq)\),
\[
 \boxed{
 \inf_{|\tau|\le T}
 \mathcal G_{\{p,q\},A}[g_{A\tau}]
 \ge
 c_1^2e^{-2c_2T}
 \min\{\lambda_pa_p(A),\lambda_qa_q(A)\},}
 \tag{V12.13}
\]
after decreasing \(c_1\) by an absolute factor.  Here
\(\mathcal G_{\{p,q\},A}\) is the sum of only the \(p\)- and \(q\)-cells.
\end{theorem}

\begin{proof}
Put
\[
 x_p(\tau):=\frac{a_p}{r}+\frac{\tau\log p}{2\pi},
 \qquad
 x_q(\tau):=\frac{a_q}{r}+\frac{\tau\log q}{2\pi}.
 \tag{V12.14}
\]
Choose \(k\) nearest to \(x_p(\tau)\).  Then
\[
 |k|\le\frac{T\log p}{2\pi}+2
 \tag{V12.15}
\]
and the \(q\)-phase at \(\tau_k\) satisfies
\[
 \|\eta_k\|_{\mathbb R/\mathbb Z}
 \le
 \|x_q(\tau)\|_{\mathbb R/\mathbb Z}
 +\omega\|x_p(\tau)\|_{\mathbb R/\mathbb Z}.
 \tag{V12.16}
\]
Lemma~\ref{lem:v12-log-separation} and \textup{(V12.15)} imply that the
maximum of the two circle distances on the right is at least
\(c_1'e^{-c_2T}\).  For \(\|x\|_{\mathbb R/\mathbb Z}\le1/2\),
\[
 |e^{2\pi\mathrm ix}-1|
 =2|\sin\pi x|
 \ge4\|x\|_{\mathbb R/\mathbb Z}.
 \tag{V12.17}
\]
This proves \textup{(V12.12)}.  Formula \textup{(V11.26)}, restricted to
the two cells, proves \textup{(V12.13)}.
\end{proof}

\begin{assessment}[What the elementary profile actually gives]
\label{ass:v12-elementary-profile}
The lower bound in \textup{(V12.12)} is rigorous but exponentially small in
the physical frequency horizon \(T\).  Since the harmonic energy of
\(g_{A\tau}\) is logarithmic in \(A|\tau|\), substituting
\(T\asymp e^E/A\) turns this elementary two-log estimate into a
double-exponentially small function of the allowed clock energy \(E\).
Sharper lower bounds for linear forms in logarithms could improve this
profile, but no improvement of a fixed two-cell profile can repair the
cofinal scale failure proved next.
\end{assessment}

\subsection{A fixed two-prime certificate is cofinally too small}

For the constant normalized vector \(g_0=2^{-1/2}\), one has
\[
 \ip{\mathsf Dg_0}{g_0}=0.
 \tag{V12.18}
\]
The degree identity and \(|E_{\delta_n}|=4a_n\) give
\[
 \int_{-1}^1V_A(t)|g_0(t)|^2\dd t
 =2\sum_{n\in\mathcal A_f(A)}\lambda_na_n
 =2\mathcal M_f(A).
 \tag{V12.19}
\]

\begin{theorem}[Two fixed cells cannot pay the renewal degree]
\label{cth:v12-two-cell-scale}
For every \(u\in L^2(-1,1)\),
\[
 \mathcal G_{\{p,q\},A}[u]
 \le4(\lambda_p+\lambda_q)\norm u^2.
 \tag{V12.20}
\]
On the zero-clock vector \(g_0\),
\[
 \boxed{
 \int_{-1}^1V_A(t)|g_0(t)|^2\dd t
 -\mathcal G_{\{p,q\},A}[g_0]
 \ge
 \frac{4e^A}{A}
o_f\!\left(\frac{e^A}{A}\right)-O_{p,q}(1).}
 \tag{V12.21}
\]
Therefore no cofinal comparison of the full incidence degree with the
magnetic energy can be obtained from any fixed finite set of prime cells,
even on the harmonic ground-state sector.
\end{theorem}

\begin{proof}
Each difference map in \textup{(V11.37)} has norm at most two, proving
\textup{(V12.20)}.  Equations \textup{(V12.19)} and
\textup{(V8.24)} give
\[
 2\mathcal M_f(A)
 =\frac{4e^A}{A}+o_f(e^A/A).
 \tag{V12.22}
\]
Subtract the bound \textup{(V12.20)} at \(\norm{g_0}=1\).
\end{proof}

This is the requested asymptotic decision about the V11 two-shift route.
The rigidity theorem remains useful for proving compactness and excluding
exact zero modes, but quantitative cofinal work must restore a growing
prime family before extracting phase dispersion.

\subsection{The restored character reader before resonance}

For a fixed nonprincipal primitive character, define
\[
 \Psi_\chi(x):=\sum_{n\le x}\Lambda(n)\chi(n).
 \tag{V12.23}
\]
The classical fixed-character zero-free region gives constants
\(C_\chi,c_\chi>0\) such that
\[
 \boxed{
 |\Psi_\chi(x)|
 \le C_\chi x e^{-c_\chi\sqrt{\log x}}
 \qquad(x\ge2).}
 \tag{V12.24}
\]
This is the standard de la Vall\'ee Poussin input already declared in the
manuscript.  Its constants may depend on the fixed conductor; it assumes no
zero-free region near the critical line.

\begin{theorem}[Triangular character-polynomial bound]
\label{thm:v12-triangular-PNT}
After decreasing \(c_\chi\) and increasing \(C_\chi\), for every sufficiently
large \(A\) and every real \(\tau\),
\[
 \boxed{
 |\mathcal P_{\chi,A}(\tau)|
 \le
 C_\chi(1+|\tau|)e^A e^{-c_\chi\sqrt A}.}
 \tag{V12.25}
\]
\end{theorem}

\begin{proof}
Put \(X=e^{2A}\) and
\[
 w_{A,\tau}(x)
 :=x^{-1/2+\mathrm i\tau}
 \left(1-\frac{\log x}{2A}\right).
 \tag{V12.26}
\]
The triangular weight vanishes at \(X\).  Stieltjes partial summation gives
\[
 \mathcal P_{\chi,A}(\tau)
 =-\int_{1^-}^{X}\Psi_\chi(x)w_{A,\tau}'(x)\dd x,
 \tag{V12.27}
\]
where the finitely many ramified terms are already zero under the Dirichlet
character.  Direct differentiation yields
\[
 |w_{A,\tau}'(x)|
 \le x^{-3/2}
 \left[
 \left(\frac12+|\tau|\right)
 \frac1{2A}
 \right].
 \tag{V12.28}
\]
On \(2\le x\le e^A\), the resulting integral is
\(O_\chi((1+|\tau|)e^{A/2})\).  On \(e^A\le x\le e^{2A}\), use
\textup{(V12.24)} and
\(e^{-c_\chi\sqrt{\log x}}\le e^{-c_\chi\sqrt A}\) to obtain
\[
 O_\chi\!\left(
 (1+|\tau|)e^{-c_\chi\sqrt A}
 \int_{e^A}^{e^{2A}}x^{-1/2}\dd x
 \right)
 =
 O_\chi\!\left(
 (1+|\tau|)e^Ae^{-c_\chi\sqrt A}
 \right).
 \tag{V12.29}
\]
The first range is absorbed by the same bound for large \(A\), after a
decrease of \(c_\chi\).  This proves \textup{(V12.25)}.
\end{proof}

\begin{corollary}[Full-history square-root resonance barrier]
\label{cor:v12-resonance-barrier}
There is \(b_\chi>0\) such that, for all sufficiently large \(A\),
\[
 |\tau|\le e^{b_\chi\sqrt A}
 \quad\Longrightarrow\quad
 |\mathcal P_{\chi,A}(\tau)|
 \le\frac12\mathcal M_f(A)
 \tag{V12.30}
\]
and hence
\[
 \boxed{
 \mathcal G_{\chi,A}[g_{A\tau}]
 \ge\mathcal M_f(A).}
 \tag{V12.31}
\]
In particular, any modulation satisfying
\[
 \mathcal G_{\chi,A}[g_{A\tau}]
 <\mathcal M_f(A)
 \tag{V12.32}
\]
must obey
\[
 \boxed{\log|\tau|>b_\chi\sqrt A.}
 \tag{V12.33}
\]
\end{corollary}

\begin{proof}
Choose \(0<b_\chi<c_\chi\).  Equations \textup{(V12.25)} and
\textup{(V8.24)} give, uniformly in the stated band,
\[
 \frac{|\mathcal P_{\chi,A}(\tau)|}{\mathcal M_f(A)}
 \ll_\chi
 A e^{-(c_\chi-b_\chi)\sqrt A}
 \longrightarrow0.
 \tag{V12.34}
\]
This proves \textup{(V12.30)}.  Expanding \textup{(V11.26)} gives the exact
identity
\[
 \mathcal G_{\chi,A}[g_{A\tau}]
 =2\mathcal M_f(A)
 -2\operatorname{Re}\mathcal P_{\chi,A}(\tau).
 \tag{V12.35}
\]
The remaining claims follow.
\end{proof}

\begin{corollary}[The resonance pays a square-root harmonic clock]
\label{cor:v12-clock-barrier}
Under \textup{(V12.32)}, the exact harmonic-clock formula
\textup{(V8.14)} gives
\[
 \boxed{
 \ip{\mathsf Dg_{A\tau}}{g_{A\tau}}
 \ge b_\chi\sqrt A+\log A-O_\chi(1).}
 \tag{V12.36}
\]
\end{corollary}

\begin{proof}
For \(A|\tau|\ge2\), the cosine integral and the sine quotient in
\textup{(V8.14)} are bounded.  Thus
\[
 \ip{\mathsf Dg_{A\tau}}{g_{A\tau}}
 =\log(A|\tau|)+O(1).
 \tag{V12.37}
\]
Insert \textup{(V12.33)}.
\end{proof}

\begin{firewall}[Why the classical zero-free region does not prove the
needed positivity]
\label{fire:v12-zero-free}
The bound \textup{(V12.25)} is exponentially smaller than the absolute
triangular mass, but it is still of size
\[
 e^{A-c_\chi\sqrt A}
 \tag{V12.38}
\]
on a fixed low-frequency band.  This is exponentially larger than the
\(O(\sqrt A)\) clock in \textup{(V12.36)} and much larger than the
root-mean-square \(O(A)\) scale from V8.  Moreover,
\textup{(V12.30)}--\textup{(V12.36)} concern constant-modulus modulations,
not arbitrary vectors in a harmonic-energy ball.  Promoting this scalar
bound to the operator profile \(\mathfrak A_{\chi,A}(E)\) without controlling
off-diagonal superpositions would repeat the diagonal-only error ruled out
in V9.
\end{firewall}

\subsection{V12 status and the next operator target}

\begin{center}
\begin{tabularx}{\textwidth}{@{}p{0.29\textwidth}p{0.17\textwidth}X@{}}
\toprule
V12 component & Status & Exact conclusion \\
\midrule
Two-log profile & quantified & The exact inhomogeneous rotation has the
elementary separation \textup{(V12.5)} and the finite-band phase floor
\textup{(V12.12)}. \\
Fixed-cell cofinal mechanism & disproved & Any fixed finite prime family has
bounded magnetic energy, while the full degree is \(4e^A/A+o(e^A/A)\) on
the zero-clock vector. \\
Restored character reader & bounded & Classical fixed-character PNT gives
\textup{(V12.25)} for the full triangular polynomial. \\
Resonance clock & acquired for modulations & Near-resonance requires
\(\log|\tau|\gg_\chi\sqrt A\) and harmonic energy
\(\gg_\chi\sqrt A\). \\
Arbitrary clock-band vectors & open & No operator bound for coherent
superpositions has been deduced from the scalar modulation estimate. \\
Cofinal completed positivity and GRH & open & The remaining prime bound is
exponential and does not close the Weil form. \\
\bottomrule
\end{tabularx}
\end{center}

The next object must retain a growing family and the entire clock band.
Let
\[
 \Pi_E:=\mathbf1_{[0,E]}(\mathsf D).
 \tag{V12.39}
\]
Since \(H_j\sim\log j\), its dimension is of order \(e^E\).  V13 should
compute or bound the full compressed adjacency
\[
 \boxed{
 \Pi_E\mathsf E_{\chi,A}\Pi_E}
 \tag{V12.40}
\]
using the exact translation matrices from V7, with character cancellation
taken only after the matrix is assembled.  The acceptable alternatives are:
\begin{enumerate}[label=\textup{(\roman*)},leftmargin=3em]
\item a sign-sensitive large-sieve or Gram bound whose error is subexponential
in \(A\) on the dangerous clock scale;
\item a rigorous counterexample showing that the clock-band compression
still contains coherent packets of exponential adverse size; or
\item an upstream change of localization if the sharp support wall prevents
any such operator estimate.
\end{enumerate}
Another scalar frequency-density count is not an acceptable substitute,
because V9 already proved that it omits the sinc cross block.

\begin{assessment}[V12 acquisition]
V12 completes the quantitative audit of the V11 two-prime route.  Its
inhomogeneous phase profile is explicit, but the route is cofinally
under-scaled even on the harmonic ground state.  Restoring all prime cells
before applying the classical zero-free region produces the first explicit
lower bound on the time required for simultaneous character resonance:
the resonant modulation must pay a square-root harmonic clock.  The
classical estimate is nevertheless exponentially too large and remains
scalar.  The live problem is now the sign-sensitive compression of the full
prime adjacency to a harmonic clock band.  GRH is not proved.
\end{assessment}

\section*{Version V12 append-only record}
\addcontentsline{toc}{section}{Version V12 append-only record}
The complete V11 source was retained before this continuation.  It had
\(197939\) bytes, \(5956\) lines, and SHA--256
\[
 \texttt{9E5BACD5C000128AF713CE8DB056CC328EB5F5C53152E972FA4872E96D552E75}.
\]
V12 quantified the two-log phase profile, proved the fixed-cell cofinal
scale obstruction, derived the full triangular character-polynomial bound,
and established a square-root harmonic-clock barrier for simultaneous
resonance.

% =============================================================================
% END APPENDED V12 MATERIAL
% =============================================================================
% =============================================================================
% BEGIN APPENDED V13 MATERIAL
% =============================================================================

\section{Full Legendre-Band Compression: Matrix Abel Summation and Its Capacity Limit}
\label{sec:grh-v13}

V12 left one precise obligation: character cancellation had to be applied
\emph{after} assembling every off-diagonal matrix element of the prime
adjacency on a harmonic clock band.  This chapter performs that assembly.
The exact polynomial translation matrices of V7 make the reduction finite,
but expanding them coefficient by coefficient conceals a serious growth
cost.  A matrix-valued Abel transform avoids that coefficient explosion and
isolates the cost in one exact derivative estimate.

The result is both positive and diagnostic.  It is the first bound in this
record for the complete compressed adjacency
\(\Pi_E\mathsf E_{\chi,A}\Pi_E\), including all coherent off-diagonal
superpositions.  It also proves that the classical fixed-character
zero-free region, used through total variation, retains an unavoidable
\(e^A\) scale even on the one-dimensional clock band.  Thus the missing
ingredient is no longer ``control the cross terms'' in the abstract: it is a
sign-sensitive bilinear estimate which saves the critical half exponent
before total variation is taken.

\section{The assembled one-sided Legendre block}

Retain the normalized Legendre basis and truncated translation matrix from
V7:
\[
 e_j(t)=\sqrt{\frac{2j+1}{2}}P_j(t),\qquad
 B_N(\delta)=\bigl(b_{mk}(\delta)\bigr)_{0\le m,k<N},
 \tag{V13.1}
\]
where
\[
 b_{mk}(\delta)
 =\int_{-1}^{1-\delta}e_m(t)e_k(t+\delta)\,dt,
 \qquad 0\le\delta\le2.
 \tag{V13.2}
\]
Equivalently, \(B_N(\delta)\) is the matrix of
\(P_NT_\delta P_N\) on \(\operatorname{span}\{e_0,\ldots,e_{N-1}\}\).
In particular,
\[
 \|B_N(\delta)\|_{\mathrm{op}}\le1,
 \qquad B_N(2)=0.
 \tag{V13.3}
\]
Here and below \(P_N\) denotes the first-\(N\) Legendre projection, not a
prime sum.

Define the one-sided character block
\[
 \mathsf K_{\chi,A,N}
 :=\sum_{n<e^{2A}}
 \frac{\Lambda(n)\chi(n)}{\sqrt n}\,
 B_N\!\left(\frac{\log n}{A}\right).
 \tag{V13.4}
\]
The convention \(\chi(n)=0\) at ramified integers permits the unrestricted
sum.  The exact V7 prime-cell formula gives
\[
 P_N\mathsf E_{\chi,A}P_N
 =\mathsf K_{\chi,A,N}+\mathsf K_{\chi,A,N}^{*}.
 \tag{V13.5}
\]
Thus no diagonalization, scalar modulation, or omission of sinc cross terms
has occurred in \textup{(V13.5)}.

There is also a finite moment description.  Since every entry in
\textup{(V13.2)} is a polynomial of degree at most \(m+k+1\), the entire
matrix \(\mathsf K_{\chi,A,N}\) is determined by the moment bank
\[
 \mathcal M_{r}(A)
 :=\sum_{n<e^{2A}}
 \frac{\Lambda(n)\chi(n)}{\sqrt n}
 \left(\frac{\log n}{A}\right)^r,
 \qquad 0\le r\le2N-1.
 \tag{V13.6}
\]
Direct use of \textup{(V13.6)} would require tracking the monomial
coefficients of all Legendre correlations.  The next estimate keeps the
translation operator intact instead.

\section{An exact derivative budget for the translation matrix}

Let
\[
 v_N(t):=(e_0(t),\ldots,e_{N-1}(t))^{\mathsf T}.
 \tag{V13.7}
\]
The endpoint Christoffel identity and its elementary maximum bound give
\[
 \|v_N(t)\|_2^2
 =\sum_{j=0}^{N-1}|e_j(t)|^2
 \le \sum_{j=0}^{N-1}\frac{2j+1}{2}
 =\frac{N^2}{2},
 \qquad -1\le t\le1.
 \tag{V13.8}
\]
The inequality uses \(|P_j(t)|\le1\) on \([-1,1]\); equality holds at an
endpoint.

\begin{lemma}[Legendre differentiation budget]
Let \(\mathsf R_N\) be differentiation restricted to the first \(N\)
normalized Legendre modes, regarded as an operator from coefficient
\(\ell^2\) to \(L^2[-1,1]\).  Then
\[
 \|\mathsf R_N\|_{\mathrm{op}}
 \le\|\mathsf R_N\|_{\mathrm{HS}}
 =\frac{N\sqrt{N^2-1}}{2}
 <\frac{N^2}{2}.
 \tag{V13.9}
\]
\end{lemma}

\begin{proof}
The classical Legendre identity
\(\int_{-1}^{1}|P_j'(t)|^2dt=j(j+1)\) yields
\[
 \|e_j'\|_2^2=\frac{2j+1}{2}j(j+1).
 \tag{V13.10}
\]
Consequently
\[
 \|\mathsf R_N\|_{\mathrm{HS}}^2
 =\frac12\sum_{j=0}^{N-1}(2j+1)j(j+1)
 =\frac{N^2(N^2-1)}{4},
 \tag{V13.11}
\]
which proves the assertion.
\end{proof}

\begin{proposition}[Uniform derivative bound for the full matrix]
For every \(N\ge1\) and \(0<\delta<2\),
\[
 \boxed{\ \|B_N'(\delta)\|_{\mathrm{op}}\le N^2.\ }
 \tag{V13.12}
\]
The derivative exists entrywise up to the one-sided endpoint values.
\end{proposition}

\begin{proof}
Differentiating \textup{(V13.2)} gives the exact identity
\[
 b_{mk}'(\delta)
 =-e_m(1-\delta)e_k(1)
 +\int_{-1}^{1-\delta}e_m(t)e_k'(t+\delta)\,dt.
 \tag{V13.13}
\]
The first matrix on the right is the rank-one operator
\(-v_N(1-\delta)v_N(1)^{\mathsf T}\).  By \textup{(V13.8)}, its norm is at
most \(N^2/2\).  The second matrix is the compression of
\(T_\delta\mathsf R_N\).  Translation with zero extension is a contraction,
so \textup{(V13.9)} bounds its norm by \(N^2/2\).  The triangle inequality
proves \textup{(V13.12)}.
\end{proof}

The two summands in \textup{(V13.13)} have distinct geometric meanings.
The rank-one term is the moving support wall, while the second is internal
polynomial differentiation.  The support wall therefore appears explicitly
inside the clock-band regularity cost rather than being hidden in an
entrywise coefficient estimate.

\section{Matrix-valued Abel summation}

Write
\[
 \Psi_\chi(x):=\sum_{n\le x}\Lambda(n)\chi(n).
 \tag{V13.14}
\]
For the fixed primitive nonprincipal character under consideration, the
standard zero-free region supplies constants \(C_\chi,c_\chi>0\) such that
\[
 |\Psi_\chi(x)|
 \le C_\chi x\exp\!\bigl(-c_\chi\sqrt{\log x}\bigr)
 \tag{V13.15}
\]
for all sufficiently large \(x\), after changing \(C_\chi\) on a bounded
range.

\begin{theorem}[Full Legendre-band prime bound]
There are constants \(A_\chi,C_\chi,c_\chi>0\), independent of \(N\), such
that for \(A\ge A_\chi\),
\[
 \boxed{
 \|\mathsf K_{\chi,A,N}\|_{\mathrm{op}}
 \le C_\chi\left(1+\frac{N^2}{A}\right)
 e^A e^{-c_\chi\sqrt A}.}
 \tag{V13.16}
\]
Consequently the complete Hermitian prime adjacency satisfies
\[
 \boxed{
 \|P_N\mathsf E_{\chi,A}P_N\|_{\mathrm{op}}
 \le 2C_\chi\left(1+\frac{N^2}{A}\right)
 e^A e^{-c_\chi\sqrt A}.}
 \tag{V13.17}
\]
\end{theorem}

\begin{proof}
Set
\[
 F_{A,N}(x):=x^{-1/2}B_N\!\left(\frac{\log x}{A}\right),
 \qquad 1\le x\le e^{2A}.
 \tag{V13.18}
\]
Because \(B_N(2)=0\) and \(\Psi_\chi(1^-)=0\), Stieltjes partial summation
has no endpoint contribution and gives the exact matrix identity
\[
 \mathsf K_{\chi,A,N}
 =-\int_{1^-}^{e^{2A}}\Psi_\chi(x)\,dF_{A,N}(x).
 \tag{V13.19}
\]
Away from the harmless lower endpoint,
\[
 F_{A,N}'(x)
 =x^{-3/2}\left[
 -\frac12B_N\!\left(\frac{\log x}{A}\right)
 +\frac1A B_N'\!\left(\frac{\log x}{A}\right)
 \right].
 \tag{V13.20}
\]
Equations \textup{(V13.3)} and \textup{(V13.12)} imply
\[
 \|F_{A,N}'(x)\|_{\mathrm{op}}
 \le x^{-3/2}\left(\frac12+\frac{N^2}{A}\right).
 \tag{V13.21}
\]
On \(e^A\le x\le e^{2A}\), \textup{(V13.15)} and
\(\sqrt{\log x}\ge\sqrt A\) therefore give
\[
 \int_{e^A}^{e^{2A}}
 |\Psi_\chi(x)|\,\|F_{A,N}'(x)\|_{\mathrm{op}}\,dx
 \le C_\chi\left(1+\frac{N^2}{A}\right)
 e^A e^{-c_\chi\sqrt A}.
 \tag{V13.22}
\]
On \(1\le x\le e^A\), the trivial estimate
\(|\Psi_\chi(x)|\ll_\chi x\) contributes
\[
 O_\chi\!\left(\left(1+\frac{N^2}{A}\right)e^{A/2}\right),
 \tag{V13.23}
\]
which is absorbed by the right side of \textup{(V13.16)} for large \(A\)
after decreasing \(c_\chi\) and enlarging \(C_\chi\).  This proves
\textup{(V13.16)}.  Finally \textup{(V13.5)} and
\(\|K+K^*\|\le2\|K\|\) prove \textup{(V13.17)}.
\end{proof}

This argument applies Abel summation once in the operator norm after the
matrix has been assembled.  It is therefore immune to the diagonal-only
defect identified in V9.  For \(N=1\), one has
\(B_1(\delta)=1-\delta/2\), so \textup{(V13.16)} recovers the triangular
scalar estimate of V12.  The higher-dimensional theorem is a genuine
extension of that estimate, not a collection of scalar frequency tests.

\section{Conversion to the harmonic clock band}

Let
\[
 N(E):=\#\{j\ge0:H_j\le E\}.
 \tag{V13.24}
\]
Since the clock is diagonal in the Legendre basis,
\[
 \Pi_E=P_{N(E)}.
 \tag{V13.25}
\]
The elementary inequality \(H_j\ge\log(j+1)\) gives
\[
 N(E)\le e^E,
 \qquad E\ge0.
 \tag{V13.26}
\]

\begin{corollary}[Complete harmonic-band compression]
For \(A\ge A_\chi\) and \(E\ge0\),
\[
 \boxed{
 \|\Pi_E\mathsf E_{\chi,A}\Pi_E\|_{\mathrm{op}}
 \le C_\chi
 \left(1+\frac{e^{2E}}{A}\right)
 e^A e^{-c_\chi\sqrt A}.}
 \tag{V13.27}
\]
Here the constant has absorbed the factor two in \textup{(V13.17)}.  In
particular, if \(E\le\eta\sqrt A\) with \(2\eta<c_\chi\), then
\[
 \|\Pi_E\mathsf E_{\chi,A}\Pi_E\|_{\mathrm{op}}
 \le C_{\chi,\eta}
 e^{A-(c_\chi-2\eta)\sqrt A}.
 \tag{V13.28}
\]
\end{corollary}

\begin{proof}
Insert \(N=N(E)\) into \textup{(V13.17)} and use
\textup{(V13.26)}.  For the last assertion, absorb the factor
\(1+A^{-1}e^{2\eta\sqrt A}\) into a constant multiple of
\(e^{2\eta\sqrt A}\).
\end{proof}

Formula \textup{(V13.27)} settles the operator-theoretic gap left at the end
of V12: every superposition in the clock band is controlled at once.  The
dimension cost is explicit.  Clock height \(E\) corresponds to roughly
\(e^E\) polynomial modes, and differentiating the moving-wall translation
matrix costs at most their square.

\section{What the completed compression does and does not certify}

The estimate is still too large for the desired form domination.  Even on
the constant band \(E=0\), its right side has size
\[
 e^{A-c_\chi\sqrt A},
 \tag{V13.29}
\]
whereas the clock available there is zero and the admissible archimedean
debit is only scalar.  For every \(E=O(\sqrt A)\), the bound remains
\(e^{A-O_\chi(\sqrt A)}\), exponentially above an \(O(\sqrt A)\) clock.
Thus the completed matrix assembly removes the possibility that the V12
failure was merely caused by testing modulations one at a time.

\begin{firewall}[A large upper bound is not an adverse eigenvalue]
Equation \textup{(V13.27)} does not prove that the compressed adjacency has
an exponentially large positive eigenvalue.  It proves that the classical
zero-free-region plus total-variation method cannot certify the much
smaller clock-scale bound required for positivity.  Reversing this logical
direction would turn an insufficient upper estimate into a nonexistent
lower estimate.
\end{firewall}

\begin{proposition}[Capacity limit of total-variation Abel control]
Within the route \textup{(V13.19)}--\textup{(V13.23)}, the two independent
costs are
\[
 \underbrace{e^A}_{\text{critical }n^{-1/2}\text{ mass after Abel}}
 \quad\text{and}\quad
 \underbrace{1+N^2/A}_{\text{moving-wall clock-band variation}}.
 \tag{V13.30}
\]
The zero-free-region saving \(e^{-c_\chi\sqrt A}\) can absorb the
dimension factor only while \(\log N<c_\chi\sqrt A/2\), but it cannot
absorb the leading critical mass at any nonnegative clock height.
\end{proposition}

\begin{proof}
The first and second factors are exactly those in
\textup{(V13.22)}.  If \(N=e^{\eta\sqrt A}\), the derivative factor
contributes at most \(e^{2\eta\sqrt A}\) up to powers of \(A\), so it is
absorbed precisely for \(2\eta<c_\chi\).  The remaining exponent is
\(A-(c_\chi-2\eta)\sqrt A\), which tends to \(+\infty\) linearly in \(A\).
\end{proof}

This capacity calculation is useful because it specifies the strength of
the missing estimate.  More smoothness of the Legendre entries can improve
the polynomial factor in \(N\), but no total-variation application of
\textup{(V13.15)} removes the critical \(e^A\) term.  That term must be
cancelled before absolute values are inserted.

\section{The next sign-sensitive object}

For a unit vector \(u\in\operatorname{ran}P_N\), define its translation
correlation
\[
 C_u(\delta):=\langle u,T_\delta u\rangle.
 \tag{V13.31}
\]
The exact variational form of the compressed adjacency is
\[
 \|P_N\mathsf E_{\chi,A}P_N\|_{\mathrm{op}}
 =\sup_{\|u\|=1}
 \left|2\Re\sum_{n<e^{2A}}
 \frac{\Lambda(n)\chi(n)}{\sqrt n}
 C_u\!\left(\frac{\log n}{A}\right)\right|.
 \tag{V13.32}
\]
The family \(\{C_u:u\in\operatorname{ran}P_N\}\) is finite-dimensional but
quadratic in \(u\).  V13 controls it by the total variation of the whole
matrix curve \(B_N\).  A successful continuation must instead exploit
sign cancellation uniformly over this correlation class.

One exact route is the full Gram expansion
\[
 \|\mathsf K_{\chi,A,N}\|_{\mathrm{HS}}^2
 =\sum_{m,k<N}
 \left|\sum_{n<e^{2A}}
 \frac{\Lambda(n)\chi(n)}{\sqrt n}
 b_{mk}\!\left(\frac{\log n}{A}\right)\right|^2.
 \tag{V13.33}
\]
Expanding the square keeps the character phase
\(\chi(n)\overline{\chi(r)}\) and produces the exact translation Gram
kernel
\[
 \mathcal G_N(\delta,\varepsilon)
 :=\operatorname{Tr}\bigl(B_N(\delta)B_N(\varepsilon)^*\bigr).
 \tag{V13.34}
\]
V14 should compute or sharply bound \(\mathcal G_N\) before estimating the
double prime sum.  This is the first remaining route in which cancellation
between distinct prime cells survives long enough to attack the leading
\(e^A\) scale.  If its off-diagonal mass is still exponential on the
dangerous clock band, that will be a rigorous reason to move upstream to a
different localization.

\begin{assessment}[V13 acquisition]
V13 closes the operator-compression obligation posed in V12.  The exact
Legendre translation matrices admit a uniform derivative bound
\(\|B_N'\|\le N^2\), and matrix-valued Abel summation gives the first
complete bound for \(\Pi_E\mathsf E_{\chi,A}\Pi_E\), including every
off-diagonal coherent superposition.  The calculation also exposes a sharp
methodological obstruction: classical zero-free-region cancellation taken
through total variation retains the critical \(e^A\) mass, while the
moving support wall adds a controlled \(e^{2E}/A\) capacity cost.  The next
live object is the full translation Gram kernel, where phases between
distinct prime cells have not yet been discarded.  GRH is not proved.
\end{assessment}

\section*{Version V13 append-only record}
\addcontentsline{toc}{section}{Version V13 append-only record}
The complete V12 source was retained before this continuation.  It had
\(213848\) bytes and SHA--256
\[
 \texttt{7E628D29081A8EAE48E031DC014058B65E221B7F46A90C73AF5EDC7260CF82B7}.
\]
V13 assembled the full finite Legendre block, proved the exact \(N^2\)
translation-derivative budget, established the matrix-valued
fixed-character prime bound, converted it to a harmonic clock-band bound,
and isolated the critical-mass obstruction created by total variation.

% =============================================================================
% END APPENDED V13 MATERIAL
% =============================================================================
% =============================================================================
% BEGIN APPENDED V14 MATERIAL
% =============================================================================

\section{The Constant Legendre Reader Is Already a GRH Detector}
\label{sec:grh-v14}

V13 completed the full clock-band compression of the prime adjacency and
identified the critical \(e^A\) mass left by total-variation Abel summation.
Before investing in a sharper Gram estimate for that adjacency, one must ask
whether a subexponential absolute bound is genuinely an intermediate target.
It is not.  The constant Legendre mode alone is an exact Laplace-transform
detector for zeros to the right of the critical line.

This observation forces an upstream correction in strategy.  Proving the
desired saving for the isolated adjacency would already prove GRH.  A route
which is structurally weaker than GRH cannot supply it.  The remaining
geometric opportunity is instead to keep the magnetic defect, clock, and
archimedean debit together, so that the exponentially large local pieces
cancel before an absolute value is taken.

\section{The triangular reader and its Laplace transform}

Put \(y=2A\) and define the one-sided triangular reader
\[
 \mathcal R_\chi(y)
 :=\sum_{n<e^y}
 \frac{\Lambda(n)\chi(n)}{\sqrt n}
 \left(1-\frac{\log n}{y}\right),
 \qquad y>0.
 \tag{V14.1}
\]
At a boundary point \(n=e^y\) the weight vanishes, so the choice of strict or
nonstrict cutoff is immaterial.  Since
\(B_1(\delta)=1-\delta/2\), equations \textup{(V13.4)} and
\textup{(V13.5)} give
\[
 \mathcal R_\chi(2A)=\mathsf K_{\chi,A,1},
 \qquad
 \mathcal Q_\chi(2A)
 :=2\Re\mathcal R_\chi(2A)
 =\langle e_0,\mathsf E_{\chi,A}e_0\rangle.
 \tag{V14.2}
\]
Thus \(\mathcal Q_\chi\) is not an auxiliary scalar model: it is the exact
Hermitian prime adjacency on the lowest clock mode.

The identity
\[
 y\mathcal R_\chi(y)
 =\sum_{n\ge1}
 \frac{\Lambda(n)\chi(n)}{\sqrt n}
 (y-\log n)_+
 \tag{V14.3}
\]
turns the triangular cutoff into a second-order Laplace kernel.

\begin{theorem}[Exact Laplace identity]
For \(\Re z>1/2\),
\[
 \boxed{
 \int_0^\infty y\mathcal R_\chi(y)e^{-zy}\,dy
 =-\frac1{z^2}\frac{L'}{L}\!\left(\frac12+z,\chi\right).}
 \tag{V14.4}
\]
Consequently
\[
 \boxed{
 \int_0^\infty y\mathcal Q_\chi(y)e^{-zy}\,dy
 =-\frac1{z^2}\left[
 \frac{L'}{L}\!\left(\frac12+z,\chi\right)
 +\frac{L'}{L}\!\left(\frac12+z,\overline\chi\right)
 \right].}
 \tag{V14.5}
\]
\end{theorem}

\begin{proof}
For \(\Re z>1/2\), absolute convergence permits interchange of the sum in
\textup{(V14.3)} and the Laplace integral.  For \(a\ge0\),
\[
 \int_a^\infty(y-a)e^{-zy}\,dy=\frac{e^{-az}}{z^2}.
 \tag{V14.6}
\]
Taking \(a=\log n\) therefore gives
\[
 \int_0^\infty y\mathcal R_\chi(y)e^{-zy}\,dy
 =\frac1{z^2}\sum_{n\ge1}
 \frac{\Lambda(n)\chi(n)}{n^{1/2+z}}.
 \tag{V14.7}
\]
The Dirichlet series on the right is
\(-L'/L(1/2+z,\chi)\), proving \textup{(V14.4)}.  Since the weights in
\textup{(V14.1)} are real,
\(\overline{\mathcal R_\chi(y)}=\mathcal R_{\overline\chi}(y)\).
Adding the two versions of \textup{(V14.4)} proves \textup{(V14.5)}.
\end{proof}

\section{A precise zero-free consequence}

The use of the Hermitian reader in \textup{(V14.5)} loses no zero
information.  Logarithmic derivatives of two entire Dirichlet
\(L\)-functions cannot cancel a zero pole: the residues at a common zero
have the same, not opposite, sign.

\begin{theorem}[Exponential type controls the rightmost zero]
Let \(\chi\) be primitive and nonprincipal.  Suppose that for some
\(\alpha\ge0\),
\[
 |\mathcal Q_\chi(y)|\le C e^{\alpha y},
 \qquad y\ge1.
 \tag{V14.8}
\]
Then neither \(L(s,\chi)\) nor \(L(s,\overline\chi)\) has a zero in
\[
 \Re s>\frac12+\alpha.
 \tag{V14.9}
\]
More generally, if for every \(\varepsilon>0\),
\[
 \mathcal Q_\chi(y)=O_{\chi,\varepsilon}(e^{\varepsilon y}),
 \tag{V14.10}
\]
then GRH holds for \(L(s,\chi)\) and \(L(s,\overline\chi)\).
\end{theorem}

\begin{proof}
Under \textup{(V14.8)}, the left side of \textup{(V14.5)} converges locally
uniformly and defines a holomorphic function on \(\Re z>\alpha\).  Hence
the sum of logarithmic derivatives on the right of \textup{(V14.5)}, first
identified on \(\Re z>1/2\), has a holomorphic continuation to
\(\Re z>\alpha\).

If \(L(s,\chi)\) has a zero of multiplicity \(m\) at
\(s=\rho\), its logarithmic derivative has residue \(m\) there.  The
logarithmic derivative for \(\overline\chi\) is either holomorphic at
\(\rho\) or has another zero pole with a nonnegative integral residue.
Thus the two poles cannot cancel.  The same argument with the characters
interchanged applies to zeros of \(L(s,\overline\chi)\).  Taking
\(z=\rho-1/2\) proves \textup{(V14.9)}.

Under \textup{(V14.10)}, choose \(\varepsilon<\Re z\) on each compact
subset of \(\Re z>0\).  The same argument makes the Laplace transform
holomorphic throughout \(\Re z>0\), so there are no zeros with
\(\Re s>1/2\).  The functional equation reflects every nontrivial zero
with \(\Re s<1/2\) to one with \(\Re s>1/2\).  Therefore every nontrivial
zero lies on the critical line.
\end{proof}

The theorem is quantitative.  An estimate of exponential type \(\alpha\)
only excludes zeros to the right of \(1/2+\alpha\).  In particular, the V13
bound at \(N=1\), written with \(y=2A\), is
\[
 |\mathcal Q_\chi(y)|
 \ll_\chi
 \exp\!\left(\frac y2-c_\chi\sqrt y\right).
 \tag{V14.11}
\]
Its exponential type is still \(1/2\), so the Laplace detector extracts no
more than zero-freeness in \(\Re s>1\).  The square-root saving from the
classical zero-free region is invisible at the level of exponential type.

\section{The converse benchmark under GRH}

The subexponential condition is not merely necessary by an abstract
analytic-continuation argument.  GRH returns a polynomial bound for the same
reader.

\begin{proposition}[GRH gives polynomial triangular cancellation]
Assume GRH for \(L(s,\chi)\) and \(L(s,\overline\chi)\).  Then, for fixed
\(\chi\) and \(y\ge1\),
\[
 \mathcal R_\chi(y)=O_\chi(y^3),
 \qquad
 \mathcal Q_\chi(y)=O_\chi(y^3).
 \tag{V14.12}
\]
Consequently the subexponential condition \textup{(V14.10)} is equivalent
to GRH for this conjugate pair.
\end{proposition}

\begin{proof}
The standard GRH prime-number-theorem estimate for a fixed primitive
character is
\[
 \Psi_\chi(x)
 :=\sum_{n\le x}\Lambda(n)\chi(n)
 =O_\chi\!\left(x^{1/2}\log^2(2x)\right).
 \tag{V14.13}
\]
Set \(X=e^y\) and
\[
 f_y(x):=x^{-1/2}\left(1-\frac{\log x}{y}\right),
 \qquad 1\le x\le X.
 \tag{V14.14}
\]
Since \(f_y(X)=0\), partial summation has no upper endpoint term.  Moreover,
\[
 |f_y'(x)|
 \le x^{-3/2}\left(\frac12+\frac1y\right).
 \tag{V14.15}
\]
Equations \textup{(V14.13)}--\textup{(V14.15)} give
\[
 |\mathcal R_\chi(y)|
 \ll_\chi
 \int_1^{e^y}\frac{\log^2(2x)}{x}\,dx
 \ll_\chi y^3.
 \tag{V14.16}
\]
The same estimate for \(\overline\chi\) proves the assertion for
\(\mathcal Q_\chi\).
\end{proof}

\begin{corollary}[The absolute compression target contains GRH]
Any estimate of the form
\[
 \|P_N\mathsf E_{\chi,A}P_N\|_{\mathrm{op}}
 \le \exp(o(A))
 \tag{V14.17}
\]
which holds cofinally in \(A\) for every \(N\ge1\), with enough uniformity
to give \textup{(V14.10)} at \(N=1\), proves GRH for the conjugate pair.
Thus a proof of \textup{(V14.17)} cannot be treated as a preliminary
unconditional large-sieve lemma.
\end{corollary}

\begin{proof}
Take \(N=1\), use \textup{(V14.2)}, and apply the preceding theorem.
\end{proof}

\begin{firewall}[Cofinal subsequences need a growth bridge]
A bound only along a sparse sequence \(A_k\to\infty\) does not by itself
imply \textup{(V14.10)} for every \(y\).  To invoke the Laplace argument one
must control the gaps, for example by a uniform derivative or bounded-ratio
interpolation estimate.  The word ``cofinally'' must not silently replace
this missing bridge.
\end{firewall}

\section{Upstream pivot: estimate the defect before the adjacency}

The V11 magnetic identity was
\[
 -\langle\mathsf E_{\chi,A}u,u\rangle
 =\mathsf G_{\chi,A}[u]
 -\int_{-1}^{1}V_A(t)|u(t)|^2\,dt,
 \tag{V14.18}
\]
where
\[
 \mathsf G_{\chi,A}[u]
 =\sum_{n<e^{2A}}\frac{\Lambda(n)}{\sqrt n}
 \int_{-1}^{1-\log n/A}
 |u(t+\log n/A)-\overline{\chi(n)}u(t)|^2dt
 \ge0.
 \tag{V14.19}
\]
Separately estimating the adjacency discards the local cancellation between
the degree potential and this magnetic square.  V14 shows that recovering a
subexponential bound after that discard is already GRH-equivalent on the
constant mode.

The next viable target is therefore the relative form inequality
\[
 \boxed{
 \int_{-1}^{1}V_A(t)|u(t)|^2dt
 \le \mathsf G_{\chi,A}[u]
 +\langle\mathsf D u,u\rangle
 +C_\chi\|u\|_2^2,}
 \tag{V14.20}
\]
or the exact variant dictated by the archimedean term in the Weil form.
Unlike \textup{(V14.17)}, inequality \textup{(V14.20)} keeps the
exponentially large pieces paired before taking norms.  V10 already proved
that a single short cell has an \(O(a)\) relative debit against the harmonic
clock; V11 proved that every finite clock band has a positive magnetic gap.
The unresolved step is a cofinal packing theorem which sums those local
debits without replacing \(\mathsf G_{\chi,A}\) by its vanishing global
spectral gap.

An appropriate V15 attack is a dyadic support-wall decomposition.  Prime
cells with overlap length
\[
 a_n:=2-\frac{\log n}{A}
 \tag{V14.21}
\]
should be grouped by \(2^{-r-1}<a_n\le2^{-r}\).  On each layer, the exact
one-cell estimate from V10 should be combined with the corresponding terms
of \(\mathsf G_{\chi,A}\) before summing in \(r\).  The required output is a
Carleson-type packing bound in which each unit of endpoint trace is charged
only logarithmically many times.  If such a bound fails, the failure should
produce an explicit multiscale packet rather than another absolute
adjacency estimate.

\begin{assessment}[V14 acquisition]
V14 identifies the exact analytic strength of the V13 compression target.
The Hermitian constant-mode triangular reader has Laplace transform equal
to the sum of the logarithmic derivatives of
\(L(s,\chi)\) and \(L(s,\overline\chi)\).  Subexponential growth of this one
scalar observable is equivalent to GRH for the conjugate pair, while the
classical zero-free-region estimate has exponential type \(1/2\) and hence
recovers only the half-plane \(\Re s>1\).  Absolute compression is therefore
not a cheaper intermediate lemma.  The main line moves upstream to a
relative magnetic-defect estimate which preserves degree--adjacency
cancellation cell by cell.  GRH is not proved.
\end{assessment}

\section*{Version V14 append-only record}
\addcontentsline{toc}{section}{Version V14 append-only record}
The complete V13 source was retained before this continuation.  It had
\(228576\) bytes, \(6905\) lines, and SHA--256
\[
 \texttt{1229DD2ACFB826E602EC882ECBA18C937354BFD7BA3007D50842898AD89EE50A}.
\]
V14 proved the exact Laplace identity for the constant Legendre reader,
established its quantitative zero-free consequence and GRH equivalence,
and redirected the geometric attack from absolute adjacency compression to
a cofinal relative magnetic-defect packing inequality.

% =============================================================================
% END APPENDED V14 MATERIAL
% =============================================================================
% =============================================================================
% BEGIN APPENDED V15 MATERIAL
% =============================================================================

\section{Phase-Variance Coercivity and the Failure of Clock-Paid Shift Freezing}
\label{sec:grh-v15}

V14 moved the main line from an absolute adjacency norm to the relative
magnetic form.  The first plausible cofinal mechanism is to group nearby
prime cells whose character phases are balanced, replace their nearby
translations by one common translation, and use phase variance to control
the degree potential.  This chapter proves the exact phase-variance lemma
and the needed harmonic-clock translation modulus.  Their combination is
rigorous: phase balance is indeed dimension-free, but paying for the shift
freezing with the harmonic clock introduces the factor
\(W/\log(1/h)\), where \(W\) is the total prime weight of the cluster and
\(h\) its shift diameter.  On every scale where classical prime-number
theory supplies balanced clusters, this factor is exponentially large.

Thus the proposed dyadic freezing is not the cofinal packing theorem sought
in V14.  A successful magnetic argument must retain the distinct shifts and
extract frustration from cycles of the resulting translation graph; it
cannot first collapse an exponentially heavy cluster to one edge.

\section{Exact coercivity at a common shift}

Let \(\mathcal C\) be a finite weighted phase family
\[
 \{(w_r,\zeta_r):1\le r\le R\},
 \qquad w_r>0,\quad |\zeta_r|=1,
 \tag{V15.1}
\]
and put
\[
 W_{\mathcal C}:=\sum_{r=1}^R w_r,
 \qquad
 S_{\mathcal C}:=\sum_{r=1}^R w_r\zeta_r.
 \tag{V15.2}
\]

\begin{lemma}[Weighted phase variance]
For every \(a,b\in\mathbb C\),
\[
 \boxed{
 \sum_{r=1}^R w_r|b-\overline{\zeta_r}a|^2
 \ge (W_{\mathcal C}-|S_{\mathcal C}|)
 (|a|^2+|b|^2).}
 \tag{V15.3}
\]
The constant is sharp.
\end{lemma}

\begin{proof}
Direct expansion gives
\[
 \sum_r w_r|b-\overline{\zeta_r}a|^2
 =W_{\mathcal C}(|a|^2+|b|^2)
 -2\Re(S_{\mathcal C}b\overline a).
 \tag{V15.4}
\]
The Hermitian matrix of this quadratic form is
\[
 \begin{pmatrix}
 W_{\mathcal C}&-\overline{S_{\mathcal C}}\\
 -S_{\mathcal C}&W_{\mathcal C}
 \end{pmatrix},
 \tag{V15.5}
\]
whose eigenvalues are \(W_{\mathcal C}\pm|S_{\mathcal C}|\).  This proves
the estimate and its sharpness.
\end{proof}

For a common shift \(0\le\delta<2\), define
\[
 \begin{aligned}
 \mathsf G_{\mathcal C,\delta}[u]
 &:=\sum_{r=1}^R w_r
 \int_{-1}^{1-\delta}
 |u(t+\delta)-\overline{\zeta_r}u(t)|^2dt,\\
 \mathsf V_\delta[u]
 &:=\int_{-1}^{1-\delta}
 \bigl(|u(t)|^2+|u(t+\delta)|^2\bigr)dt.
 \end{aligned}
 \tag{V15.6}
\]

\begin{corollary}[Dimension-free common-shift gap]
One has
\[
 \boxed{
 \mathsf G_{\mathcal C,\delta}[u]
 \ge(W_{\mathcal C}-|S_{\mathcal C}|)\mathsf V_\delta[u].}
 \tag{V15.7}
\]
In particular, if \(|S_{\mathcal C}|\le\vartheta W_{\mathcal C}\) with
\(\vartheta<1\), then the magnetic squares control a fraction
\(1-\vartheta\) of the full cluster degree on every vector, with no clock
cutoff and no dependence on dimension.
\end{corollary}

\begin{proof}
Apply \textup{(V15.3)} pointwise with
\(a=u(t)\), \(b=u(t+\delta)\), and integrate.
\end{proof}

This is the desired algebraic cancellation.  For a nonprincipal character,
equal distribution of prime weight among reduced residue classes makes
\(S_{\mathcal C}\) small because
\(\sum_{a\bmod q}\chi(a)=0\).  The remaining issue is that distinct prime
cells never have exactly the same shift.

\section{A harmonic-clock modulus for zero-extended translation}

Extend \(u\in L^2[-1,1]\) by zero to the real line and let
\[
 (\mathsf S_hu)(t):=u(t+h),\qquad t\in\mathbb R.
 \tag{V15.8}
\]
The harmonic clock is
\[
 \mathsf De_j=H_je_j,
 \qquad H_j=\sum_{k=1}^j\frac1k.
 \tag{V15.9}
\]

\begin{proposition}[Logarithmic translation modulus]
For \(0<h\le2^{-8}\) and every \(u\) in the form domain of \(\mathsf D\),
\[
 \boxed{
 \|\mathsf S_hu-u\|_{L^2(\mathbb R)}^2
 \le3\sqrt h\,\|u\|_2^2
 +\frac{32}{\log(1/h)}\langle\mathsf Du,u\rangle.}
 \tag{V15.10}
\]
The same estimate holds with \(h\) replaced by \(-h\).
\end{proposition}

\begin{proof}
Choose
\[
 N:=\lfloor h^{-1/4}\rfloor,
 \qquad p=P_Nu,quad r=(I-P_N)u.
 \tag{V15.11}
\]
On the interior interval \([-1,1-h]\), the fundamental theorem of calculus
and Cauchy--Schwarz give
\[
 \int_{-1}^{1-h}|p(t+h)-p(t)|^2dt
 \le h^2\|p'\|_2^2.
 \tag{V15.12}
\]
The two boundary slivers created by zero extension have total mass at most
\[
 hN^2\|p\|_2^2,
 \tag{V15.13}
\]
because the Christoffel bound \textup{(V13.8)} gives
\(|p(t)|^2\le N^2\|p\|_2^2/2\).  The differentiation budget
\textup{(V13.9)} gives
\[
 \|p'\|_2^2\le\frac{N^4}{4}\|p\|_2^2.
 \tag{V15.14}
\]
Consequently
\[
 \|\mathsf S_hp-p\|_2^2
 \le\left(\frac{h^2N^4}{4}+hN^2\right)\|u\|_2^2.
 \tag{V15.15}
\]
Translation is unitary on \(L^2(\mathbb R)\), while the spectral tail obeys
\[
 \|r\|_2^2\le H_N^{-1}\langle\mathsf Du,u\rangle.
 \tag{V15.16}
\]
Using \(\|x+y\|^2\le2\|x\|^2+2\|y\|^2\) and
\(\|\mathsf S_hr-r\|\le2\|r\|\), we obtain
\[
 \|\mathsf S_hu-u\|_2^2
 \le\left(\frac{h^2N^4}{2}+2hN^2\right)\|u\|_2^2
 +\frac8{H_N}\langle\mathsf Du,u\rangle.
 \tag{V15.17}
\]
Now \(N\le h^{-1/4}\), so the scalar coefficient is at most
\(h/2+2\sqrt h\le3\sqrt h\).  Also
\(N+1>h^{-1/4}\) and \(H_N\ge\log(N+1)\), hence
\(H_N>\frac14\log(1/h)\).  This proves \textup{(V15.10)}.  Reflection gives
the assertion for negative \(h\).
\end{proof}

The logarithm in \textup{(V15.10)} is not an artifact of a derivative
clock: it is exactly the scale predicted by \(H_j\sim\log j\).  This makes
the cost of freezing a heavily weighted translation cluster explicit.

\section{Nearby shifts: the exact freezing remainder}

Let a cluster now contain triples
\[
 \mathcal C=\{(w_r,\zeta_r,\delta_r):1\le r\le R\},
 \qquad 0\le\delta_r<2,
 \tag{V15.18}
\]
with
\[
 |\delta_r-\delta_0|\le h\le2^{-8}.
 \tag{V15.19}
\]
Put \(\delta_*:=\max_r\delta_r\),
\(I_*=[-1,1-\delta_*]\), and define the common-core potential
\[
 \mathsf V_{*,0}[u]
 :=\int_{I_*}\bigl(|u(t)|^2+|u(t+\delta_0)|^2\bigr)dt.
 \tag{V15.20}
\]
Let \(\mathsf G_{\mathcal C}[u]\) be the sum of the actual magnetic cells,
each integrated on \([-1,1-\delta_r]\).

\begin{theorem}[Phase coercivity with a quantified shift remainder]
For every \(0<\eta<1\),
\[
 \boxed{
 \begin{aligned}
 \mathsf G_{\mathcal C}[u]
 \ge{}&(1-\eta)(W_{\mathcal C}-|S_{\mathcal C}|)
 \mathsf V_{*,0}[u]\\
 &-(\eta^{-1}-1)W_{\mathcal C}
 \left(
 3\sqrt h\,\|u\|_2^2
 +\frac{32}{\log(1/h)}\langle\mathsf Du,u\rangle
 \right).
 \end{aligned}}
 \tag{V15.21}
\]
\end{theorem}

\begin{proof}
Restrict every nonnegative cell integral to \(I_*\).  On that interval set
\[
 a(t)=u(t),\qquad b_0(t)=u(t+\delta_0),qquad
 e_r(t)=u(t+\delta_r)-u(t+\delta_0).
 \tag{V15.22}
\]
The elementary inequality
\[
 |x+e|^2\ge(1-\eta)|x|^2-(\eta^{-1}-1)|e|^2
 \tag{V15.23}
\]
and the phase-variance lemma yield
\[
 \begin{aligned}
 \mathsf G_{\mathcal C}[u]
 \ge{}&(1-\eta)(W_{\mathcal C}-|S_{\mathcal C}|)
 \mathsf V_{*,0}[u]\\
 &-(\eta^{-1}-1)
 \sum_r w_r\int_{I_*}|e_r(t)|^2dt.
 \end{aligned}
 \tag{V15.24}
\]
After a change of variable, each last integral is bounded by
\(\|\mathsf S_{\delta_r-\delta_0}u-u\|_{L^2(\mathbb R)}^2\).
Apply \textup{(V15.10)} and use
\(|\delta_r-\delta_0|\le h\).  Summing the weights proves
\textup{(V15.21)}.
\end{proof}

The favorable term in \textup{(V15.21)} is proportional to the cluster
weight, as it must be to control the degree.  The unfavorable clock term is
also proportional to that weight.  Phase balance does not reduce it.

\section{Capacity audit on prime-number-theorem bins}

For a shift interval \(J=[d,d+h]\subset(0,2)\), define
\[
 \begin{aligned}
 W_A(J)&:=\sum_{\substack{e^{Ad}\le n<e^{A(d+h)}\\ (n,q)=1}}
 \frac{\Lambda(n)}{\sqrt n},\\
 S_{\chi,A}(J)&:=\sum_{e^{Ad}\le n<e^{A(d+h)}}
 \frac{\Lambda(n)\chi(n)}{\sqrt n}.
 \end{aligned}
 \tag{V15.25}
\]
Ordinary and character prime-number theorems, followed by partial
summation, give for fixed \(0<d<2\)
\[
 W_A(J)\asymp Ah\,e^{Ad/2}
 \tag{V15.26}
\]
when \(Ah\le1\) and the main term exceeds the classical error, while
\[
 |S_{\chi,A}(J)|
 \ll_\chi e^{A(d+h)/2-c_\chi\sqrt A}.
 \tag{V15.27}
\]
For example, if
\[
 h=e^{-\kappa\sqrt A},
 \qquad 0<\kappa<c_\chi,
 \tag{V15.28}
\]
then
\[
 \frac{|S_{\chi,A}(J)|}{W_A(J)}
 \ll_\chi A^{-1}e^{-(c_\chi-\kappa)\sqrt A}=o(1),
 \tag{V15.29}
\]
so the bin is strongly phase-balanced.  Its weight, however, is
\[
 W_A(J)\asymp A\exp\!\left(\frac{Ad}{2}-\kappa\sqrt A\right).
 \tag{V15.30}
\]

\begin{proposition}[Clock-capacity failure of PNT-scale freezing]
Apply \textup{(V15.21)} to a bin satisfying
\textup{(V15.28)}.  The coefficient multiplying the harmonic clock in the
freezing remainder is at least on the scale
\[
 \frac{W_A(J)}{\log(1/h)}
 \asymp_{\kappa}
 \sqrt A\exp\!\left(\frac{Ad}{2}-\kappa\sqrt A\right),
 \tag{V15.31}
\]
which is exponentially large for every fixed \(d>0\).  Hence classical
PNT phase balance plus common-shift freezing cannot yield a relative form
bound with clock coefficient less than one.
\end{proposition}

\begin{proof}
Insert \(\log(1/h)=\kappa\sqrt A\) and
\textup{(V15.30)} into the clock remainder in
\textup{(V15.21)}.  Constants depending on \(\eta\) do not change the
exponential scale.
\end{proof}

More abstractly, a cluster of weight \(W\) can be frozen by
\textup{(V15.21)} with \(O(1)\) clock expenditure only if
\[
 \log(1/h)\gtrsim W.
 \tag{V15.32}
\]
For an exponentially heavy balanced cluster this requires a
double-exponentially small shift diameter.  Classical fixed-character PNT
balances shifts only on the square-root-exponential scale
\textup{(V15.28)}, far above this requirement.

\begin{firewall}[Failure of freezing is not failure of magnetic coercivity]
Proposition \textup{(V15.31)} disproves only the route which first replaces
all shifts in a balanced cluster by a common shift and then pays the
replacement through \textup{(V15.10)}.  It does not show that the original
distinct-shift magnetic form lacks coercivity.  Distinct shifts create
additional cycles and may increase frustration.  Discarding those cycles is
exactly what produced the exponential clock debit.
\end{firewall}

\section{Next object: translation-graph cycle frustration}

For two cells \((\delta_1,\zeta_1)\) and
\((\delta_2,\zeta_2)\), their magnetic residuals satisfy on the common
domain the exact relation
\[
 \begin{aligned}
 &u(t+\delta_2)
 -\overline{\zeta_2}\zeta_1u(t+\delta_1)\\
 &\qquad=
 \bigl(u(t+\delta_2)-\overline{\zeta_2}u(t)\bigr)
 -\overline{\zeta_2}\zeta_1
 \bigl(u(t+\delta_1)-\overline{\zeta_1}u(t)\bigr).
 \end{aligned}
 \tag{V15.33}
\]
Unlike common-shift freezing, \textup{(V15.33)} is paid directly by the two
magnetic squares.  It transports phase frustration to the difference shift
\(\delta_2-\delta_1\) without multiplying the harmonic clock by the total
cluster weight.

V16 should iterate \textup{(V15.33)} around short cycles of prime shifts and
derive a weighted graph Poincar\'e inequality.  The target is a lower bound
for the degree potential in terms of the original magnetic squares plus a
\emph{single} harmonic-clock remainder.  The arithmetic input should enter
through the holonomy
\[
 \prod_{e\in\gamma}\chi(n_e)^{\sigma_e}
 \tag{V15.34}
\]
of a translation cycle, not through an absolute sum of cell weights.  A
cycle family whose holonomies stay away from one on every dangerous scale
would provide the cofinal packing mechanism that common-shift bins cannot.

\begin{assessment}[V15 acquisition]
V15 proves the exact dimension-free phase-variance gap for a balanced
common-shift cluster and a quantitative logarithmic modulus of translation
for the harmonic clock.  Combining them gives a complete audit of the
natural dyadic-freezing proposal from V14.  On every bin where classical
fixed-character prime-number theory ensures phase balance, the required
clock expenditure is exponentially large.  The obstruction comes from
discarding distinct-shift cycles, not from a lack of phase variance.  The
next live object is therefore the weighted translation graph and its cycle
holonomy.  GRH is not proved.
\end{assessment}

\section*{Version V15 append-only record}
\addcontentsline{toc}{section}{Version V15 append-only record}
The complete V14 source was retained before this continuation.  It had
\(239879\) bytes, \(7227\) lines, and SHA--256
\[
 \texttt{692964B4A5125A78CB13B5575373114DAF6C141CB6ED86AA670BFB17B60EDDFA}.
\]
V15 proved common-shift phase-variance coercivity, established the harmonic
clock translation modulus, quantified the exact nearby-shift freezing
remainder, and ruled out PNT-scale dyadic freezing as a cofinal relative
form argument.

% =============================================================================
% END APPENDED V15 MATERIAL
% =============================================================================
% =============================================================================
% BEGIN APPENDED V16 MATERIAL
% =============================================================================

\section{Arithmetic Flatness of Exact Translation Cycles and the Effective-Conductance Target}
\label{sec:grh-v16}

V15 showed that collapsing a phase-balanced prime cluster to one common
translation spends an exponentially large amount of harmonic clock.  It
therefore proposed retaining the distinct shifts and extracting frustration
from translation-graph cycles.  The first audit of that proposal reveals a
rigid structural constraint: the shift and the magnetic phase are two images
of the same multiplicative label.  Every exact zero-displacement cycle is
therefore flat and has trivial character holonomy.

Nontrivial holonomy occurs only on an \emph{approximate} cycle, and its
unclosed displacement must be paid by spatial translation.  This chapter
proves the exact path telescoping inequality, quantifies the arithmetic
displacement of a nontrivial near-cycle, and derives a path Poincar\'e bound
whose coefficient is the reciprocal electrical resistance of the path.
The resulting network criterion is substantially sharper than the failed
cluster-freezing estimate: clock is weighted by effective conductance, not
by the full prime mass.  It also gives a concrete next problem with explicit
congestion and clock budgets.

\section{Multiplicative paths, displacement, and holonomy}

An oriented prime-cell path is a finite list
\[
 \gamma=((n_1,\sigma_1),\ldots,(n_L,\sigma_L)),
 \qquad \sigma_j\in\{-1,+1\},
 \tag{V16.1}
\]
where \(\Lambda(n_j)>0\) and \((n_j,q)=1\).  Put
\[
 \delta_n:=\frac{\log n}{A},
 \qquad
 s_j:=\sum_{k=1}^j\sigma_k\delta_{n_k},
 \qquad s_0=0,
 \tag{V16.2}
\]
and define
\[
 R_\gamma:=\prod_{j=1}^L n_j^{\sigma_j}\in\mathbb Q_{>0}.
 \tag{V16.3}
\]
The net displacement and magnetic holonomy are
\[
 \Delta_A(\gamma):=s_L=\frac1A\log R_\gamma,
 \qquad
 \operatorname{Hol}_\chi(\gamma)
 :=\prod_{j=1}^L\overline{\chi(n_j)}^{\sigma_j}
 =\overline{\chi(R_\gamma)}.
 \tag{V16.4}
\]
For a rational \(M/N\) prime to \(q\), the notation
\(\chi(M/N):=\chi(M)\overline{\chi(N)}\) is well defined after reducing the
fraction.

\begin{theorem}[Arithmetic flatness of exact cycles]
Every exact zero-displacement multiplicative path has trivial holonomy:
\[
 \boxed{
 \Delta_A(\gamma)=0
 \quad\Longrightarrow\quad
 R_\gamma=1
 \quad\Longrightarrow\quad
 \operatorname{Hol}_\chi(\gamma)=1.}
 \tag{V16.5}
\]
Thus no exact translation cycle can supply a character-phase gap.
\end{theorem}

\begin{proof}
The real logarithm is injective on \(\mathbb Q_{>0}\), so
\(\log R_\gamma=0\) implies \(R_\gamma=1\).  Complete multiplicativity then
gives \(\chi(R_\gamma)=1\).
\end{proof}

This is not merely the absence of a convenient short cycle.  It is a flat
connection identity for the entire multiplicative translation graph.  Any
cycle argument which assumes a nontrivial phase around an exactly closed
logarithmic path contradicts \textup{(V16.5)}.

\section{Arithmetic size of a nontrivial near-cycle}

Write \(R_\gamma=M_\gamma/N_\gamma\) in lowest terms and set the reduced
multiplicative height
\[
 \mathcal H_\times(\gamma)
 :=\log\max(M_\gamma,N_\gamma).
 \tag{V16.6}
\]

\begin{proposition}[Holonomy forces a resolvable displacement]
If \(\operatorname{Hol}_\chi(\gamma)\ne1\), then
\[
 \boxed{
 |\Delta_A(\gamma)|
 \ge\frac{e^{-\mathcal H_\times(\gamma)}}{A}.}
 \tag{V16.7}
\]
Moreover, nontrivial holonomy with exponentially small displacement really
occurs: there are paths \(\gamma_k\) and a constant
\(c_\chi>0\) such that
\[
 |1-\operatorname{Hol}_\chi(\gamma_k)|\ge c_\chi,
 \qquad
 |\Delta_A(\gamma_k)|\ll_q A^{-1}e^{-\mathcal H_\times(\gamma_k)}.
 \tag{V16.8}
\]
\end{proposition}

\begin{proof}
Nontrivial holonomy implies \(M_\gamma\ne N_\gamma\).  If \(M>N\), then
\[
 \log\frac MN
 \ge\log\left(1+\frac1N\right)
 \ge\frac1{N+1}
 \ge\frac1M.
 \tag{V16.9}
\]
The case \(N>M\) is symmetric.  Division by \(A\) proves
\textup{(V16.7)}.

Because \(\chi\) is nonprincipal, choose reduced residue classes
\(r,s\pmod q\) with \(\chi(r)\ne\chi(s)\).  Choose integer representatives
with fixed nonzero difference and put
\[
 M_k=kq+r,qquad N_k=kq+s
 \tag{V16.10}
\]
for all sufficiently large \(k\).  Both integers are prime to \(q\), and
\[
 \chi(M_k)\overline{\chi(N_k)}
 =\chi(r)\overline{\chi(s)}\ne1.
 \tag{V16.11}
\]
Factor \(M_k/N_k\) into primes and their inverses to obtain an oriented
prime-cell path.  Since \(|M_k-N_k|=|r-s|\),
\[
 \left|\log\frac{M_k}{N_k}\right|
 \ll_q k^{-1}
 \asymp_q e^{-\mathcal H_\times(\gamma_k)}.
 \tag{V16.12}
\]
The fixed phase in \textup{(V16.11)} supplies \(c_\chi>0\), proving
\textup{(V16.8)}.
\end{proof}

Thus approximate arithmetic cycles can carry a fixed phase defect while
missing closure by only the reciprocal multiplicative height.  The
harmonic clock is precisely logarithmic enough to see this displacement,
but only if the path energy is assembled with the correct conductance.

\section{Exact weighted path telescoping}

For a prime cell define its conductance
\[
 \lambda_n:=\frac{\Lambda(n)}{\sqrt n}.
 \tag{V16.13}
\]
For the oriented path \(\gamma\), let
\[
 \alpha_j:=\overline{\chi(n_j)}^{\sigma_j}.
 \tag{V16.14}
\]
An interval \(I_\gamma\subset[-1,1]\) is called admissible if, for every
\(t\in I_\gamma\), all vertices \(t+s_j\) lie in \([-1,1]\) and every
oriented edge is contained in the overlap domain of its underlying prime
cell.  Explicitly, for \(\sigma_j=+1\) require
\[
 t+s_{j-1}\in[-1,1-\delta_{n_j}],
 \tag{V16.15}
\]
and for \(\sigma_j=-1\) require
\[
 t+s_j\in[-1,1-\delta_{n_j}].
 \tag{V16.16}
\]
Define the path resistance
\[
 \mathcal R(\gamma):=\sum_{j=1}^L\lambda_{n_j}^{-1}
 \tag{V16.17}
\]
and the magnetic energy of its cells, counted with path multiplicity,
\[
 \mathsf G_\gamma[u]
 :=\sum_{j=1}^L\lambda_{n_j}
 \int_{-1}^{1-\delta_{n_j}}
 |u(x+\delta_{n_j})-\overline{\chi(n_j)}u(x)|^2dx.
 \tag{V16.18}
\]

\begin{theorem}[Weighted path inequality]
For every admissible \(I_\gamma\),
\[
 \boxed{
 \int_{I_\gamma}
 |u(t+\Delta_A(\gamma))
 -\operatorname{Hol}_\chi(\gamma)u(t)|^2dt
 \le\mathcal R(\gamma)\mathsf G_\gamma[u].}
 \tag{V16.19}
\]
\end{theorem}

\begin{proof}
Set
\[
 e_j(t):=u(t+s_j)-\alpha_ju(t+s_{j-1}).
 \tag{V16.20}
\]
For \(\sigma_j=+1\), this is the prime-cell residual evaluated at
\(t+s_{j-1}\).  For \(\sigma_j=-1\), it is, up to multiplication by the
unit scalar \(-\chi(n_j)\), the forward residual evaluated at \(t+s_j\).
The admissibility conditions therefore place every \(e_j\) inside the
corresponding integral in \textup{(V16.18)}.

Telescoping with the unit phases gives
\[
 u(t+s_L)-\left(\prod_{j=1}^L\alpha_j\right)u(t)
 =\sum_{j=1}^L
 \left(\prod_{k=j+1}^L\alpha_k\right)e_j(t).
 \tag{V16.21}
\]
Weighted Cauchy--Schwarz yields
\[
 \left|\sum_j c_je_j(t)\right|^2
 \le\left(\sum_j\lambda_{n_j}^{-1}\right)
 \left(\sum_j\lambda_{n_j}|e_j(t)|^2\right),
 \tag{V16.22}
\]
since \(|c_j|=1\).  Integrate over \(I_\gamma\), enlarge each edge integral
to its full overlap domain, and use \textup{(V16.4)}.
\end{proof}

The reciprocal resistance \(\mathcal R(\gamma)^{-1}\), rather than the sum
of cell weights, is the conductance carried by one path.  This reverses the
fatal weighting in the common-shift remainder of V15.

\section{Near-cycle Poincar\'e inequality with one clock debit}

Put
\[
 \kappa_\gamma
 :=|1-\operatorname{Hol}_\chi(\gamma)|.
 \tag{V16.23}
\]

\begin{corollary}[One-path magnetic coercivity]
Suppose \(0<|\Delta_A(\gamma)|\le2^{-8}\).  Then
\[
 \boxed{
 \begin{aligned}
 \mathsf G_\gamma[u]
 \ge{}&\frac{\kappa_\gamma^2}{2\mathcal R(\gamma)}
 \int_{I_\gamma}|u(t)|^2dt\\
 &-\frac{3\sqrt{|\Delta_A(\gamma)|}}{\mathcal R(\gamma)}\|u\|_2^2
 -\frac{32}{\mathcal R(\gamma)
 \log(1/|\Delta_A(\gamma)|)}
 \langle\mathsf Du,u\rangle.
 \end{aligned}}
 \tag{V16.24}
\]
\end{corollary}

\begin{proof}
For \(|H|=1\),
\[
 |u(t+\Delta)-Hu(t)|^2
 \ge\frac{|1-H|^2}{2}|u(t)|^2
 -|u(t+\Delta)-u(t)|^2.
 \tag{V16.25}
\]
Insert this in \textup{(V16.19)}, restrict the first term to
\(I_\gamma\), and bound the translation error by the full-line estimate
\textup{(V15.10)}.
\end{proof}

For an exact cycle, \(\Delta_A(\gamma)=0\) and arithmetic flatness forces
\(\kappa_\gamma=0\), so \textup{(V16.24)} correctly supplies no mass gap.
For a nontrivial near-cycle, the mass coefficient is its effective
conductance times its squared phase defect, while its clock coefficient is
smaller by the additional logarithmic resolution factor.

\section{The explicit network certificate}

Let \(\Gamma_A\) be a finite family of admissible nontrivial near-cycles.
For each prime cell \(n\), let \(m_\gamma(n)\) be its multiplicity in
\(\gamma\), and define the edge congestion
\[
 \operatorname{Cong}(\Gamma_A)
 :=\sup_n\sum_{\gamma\in\Gamma_A}m_\gamma(n).
 \tag{V16.26}
\]
Summing \textup{(V16.24)} and using
\[
 \sum_{\gamma\in\Gamma_A}\mathsf G_\gamma[u]
 \le\operatorname{Cong}(\Gamma_A)\mathsf G_{\chi,A}[u]
 \tag{V16.27}
\]
gives an exact sufficient-condition framework.  Define the pointwise cycle
conductance density
\[
 \mathcal C_{\Gamma_A}(t)
 :=\sum_{\substack{\gamma\in\Gamma_A\\t\in I_\gamma}}
 \frac{\kappa_\gamma^2}{2\mathcal R(\gamma)},
 \tag{V16.28}
\]
the scalar leakage
\[
 \mathcal S_{\Gamma_A}
 :=\sum_{\gamma\in\Gamma_A}
 \frac{3\sqrt{|\Delta_A(\gamma)|}}{\mathcal R(\gamma)},
 \tag{V16.29}
\]
and the clock load
\[
 \mathcal T_{\Gamma_A}
 :=\sum_{\gamma\in\Gamma_A}
 \frac{32}{\mathcal R(\gamma)
 \log(1/|\Delta_A(\gamma)|)}.
 \tag{V16.30}
\]

\begin{proposition}[Near-cycle network bound]
Every such family satisfies
\[
 \boxed{
 \begin{aligned}
 \operatorname{Cong}(\Gamma_A)\mathsf G_{\chi,A}[u]
 \ge{}&\int_{-1}^{1}\mathcal C_{\Gamma_A}(t)|u(t)|^2dt\\
 &-\mathcal T_{\Gamma_A}\langle\mathsf Du,u\rangle
 -\mathcal S_{\Gamma_A}\|u\|_2^2.
 \end{aligned}}
 \tag{V16.31}
\]
\end{proposition}

\begin{proof}
Sum \textup{(V16.24)} over \(\Gamma_A\), use Tonelli's theorem on the
nonnegative mass terms, and apply \textup{(V16.27)}.
\end{proof}

This converts the qualitative V14 target into three quantitative network
requirements.  A successful family must have bounded congestion, must make
\(\mathcal C_{\Gamma_A}(t)\) dominate the dangerous part of the degree
potential, and must satisfy
\[
 \mathcal T_{\Gamma_A}<\operatorname{Cong}(\Gamma_A)
 \tag{V16.32}
\]
with scalar leakage payable by the archimedean floor.  All three quantities
are explicit arithmetic sums over factorizations and can be audited without
forming an absolute adjacency norm.

\begin{firewall}[Formal paths must have spatial domain]
A multiplicative relation alone does not produce an admissible interval.
The partial logarithmic sums \(s_j\) must be ordered so that every translated
vertex remains in \([-1,1]\) and every edge uses its actual overlap domain.
Ignoring this ordering constraint would turn the full multiplicative group
into a fictitious translation graph with more cycles than the renewal
operator possesses.
\end{firewall}

\section{V17 target: parallel conductance under bounded congestion}

The single-path construction in \textup{(V16.10)} can have large resistance
because a large prime factor has small conductance
\(\lambda_p=(\log p)/\sqrt p\).  No one path should be expected to dominate
the exponential degree.  The possible gain is parallelism: exponentially
many arithmetically distinct near-cycles may contribute conductances while
using each prime edge only boundedly often.

V17 should begin with paths built from two factorizations of nearby reduced
residue classes and solve the following finite packing problem on a fixed
shift layer:
\[
 \boxed{
 \begin{gathered}
 \operatorname{Cong}(\Gamma_A)=O(1),\qquad
 \mathcal T_{\Gamma_A}=o(1),\\
 \mathcal C_{\Gamma_A}(t)
 \ge \operatorname{Cong}(\Gamma_A)V_A^{\mathrm{layer}}(t)-O_\chi(1)
 \quad\text{on the layer's common spatial core.}
 \end{gathered}}
 \tag{V16.33}
\]
If the conductance sum is too small, the failure will be an explicit
network-capacity obstruction, stronger than the common-shift freezing
failure of V15.  If it is large enough, \textup{(V16.31)} supplies the first
cofinal relative magnetic certificate with a single clock debit.

\begin{assessment}[V16 acquisition]
V16 corrects the cycle-frustration proposal at its arithmetic foundation.
Exact logarithmic cycles have trivial character holonomy, so only near-cycles
can be coercive.  Their nonclosure is bounded below by reciprocal
multiplicative height and this scale is essentially attained by nearby
residue classes.  Weighted path telescoping produces a rigorous magnetic
Poincar\'e inequality governed by electrical resistance, and summing paths
gives an explicit network certificate with separate conductance, congestion,
clock, and scalar budgets.  The live question is now whether enough
low-resistance near-cycles can be packed in parallel inside the actual
support geometry.  GRH is not proved.
\end{assessment}

\section*{Version V16 append-only record}
\addcontentsline{toc}{section}{Version V16 append-only record}
The complete V15 source was retained before this continuation.  It had
\(252788\) bytes, \(7640\) lines, and SHA--256
\[
 \texttt{0C029A3C0A3A2565C5DBACA6050602AC0E66F4C24A80E5AF5EE77A3EF099D8FC}.
\]
V16 proved arithmetic flatness of exact translation cycles, quantified
nontrivial near-cycle displacement, established the weighted path and
near-cycle Poincar\'e inequalities, and formulated the effective-conductance
network certificate for the next cofinal attack.

% =============================================================================
% END APPENDED V16 MATERIAL
% =============================================================================
% =============================================================================
% BEGIN APPENDED V17 MATERIAL
% =============================================================================

\section{Two-Edge Near-Cycles: Full Conductance but an Unavoidable Exponential Clock Load}
\label{sec:grh-v17}

V16 reduced the relative magnetic problem to an effective-conductance
network with bounded edge congestion and a subunit clock load.  The first
candidate network pairs two prime cells with different character phases and
nearly equal logarithmic shifts.  Such pairs are abundant in every
prime-number-theorem-scale bin, have bounded congestion, and carry total
conductance comparable to the full weight of the bin.

This chapter proves that they nevertheless cannot satisfy the V16 network
certificate.  The obstruction is arithmetic and independent of the
matching algorithm: two distinct active integers are separated by at least
one, so their scaled logarithmic displacement has resolution at most
\(2A+\log A\).  Any two-edge family carrying enough conductance to dominate
an exponentially heavy prime layer must therefore spend an exponentially
large amount of harmonic clock.  Longer multiplicative paths are not an
optional refinement; they are necessary to create the much finer
displacements required by the relative form.

\section{Exact geometry and resistance of a paired path}

For distinct active prime powers \(n,m<e^{2A}\), consider
\[
 \gamma_{n,m}=((n,+1),(m,-1)).
 \tag{V17.1}
\]
Its displacement, holonomy, resistance, and conductance are
\[
 \begin{aligned}
 \Delta_{n,m}&=\frac1A\log\frac nm,\\
 H_{n,m}&=\overline{\chi(n)}\chi(m),\\
 \mathcal R_{n,m}&=\lambda_n^{-1}+\lambda_m^{-1},\\
 c_{n,m}&:=\mathcal R_{n,m}^{-1}
 =\frac{\lambda_n\lambda_m}{\lambda_n+\lambda_m}.
 \end{aligned}
 \tag{V17.2}
\]
The partial displacements are
\[
 0,\qquad \delta_n,\qquad \delta_n-\delta_m,
 \tag{V17.3}
\]
so their range is exactly \(\max(\delta_n,\delta_m)\).  Hence the maximal
admissible interval has length
\[
 |I_{n,m}|=2-\max(\delta_n,\delta_m).
 \tag{V17.4}
\]
This is the same support scale as the larger of the two original prime
cells; pairing does not sacrifice an additional macroscopic boundary
layer.

If
\[
 d\le\delta_n,\delta_m\le d+h,
 \qquad 0<h<\frac{2-d}{2},
 \tag{V17.5}
\]
then every such paired path admits the common interval
\[
 I_{d,h}:=[-1+h,\,1-(d+h)],
 \qquad |I_{d,h}|=2-d-2h.
 \tag{V17.6}
\]
Indeed, if \(\delta_n\ge\delta_m\), its admissible interval is
\([-1,1-\delta_n]\); if \(\delta_n<\delta_m\), it is
\([-1+\delta_m-\delta_n,1-\delta_n]\).  Both contain
\textup{(V17.6)}.

\section{Bounded-congestion phase matching is abundant}

Let \(G_\chi\) be the finite image of \(\chi\) on
\((\mathbb Z/q\mathbb Z)^\times\), and put \(r_\chi:=|G_\chi|\ge2\).  Every
value in \(G_\chi\) has the same number of reduced residue-class preimages.
Define the prime layer
\[
 \mathcal P_A(d,h)
 :=\{p\text{ prime}:e^{Ad}\le p<e^{A(d+h)},\ p\nmid q\}
 \tag{V17.7}
\]
and its weight
\[
 W_A^{\mathrm p}(d,h)
 :=\sum_{p\in\mathcal P_A(d,h)}\frac{\log p}{\sqrt p}.
 \tag{V17.8}
\]

\begin{proposition}[A matching with full layer-scale conductance]
Fix \(0<d<2\).  There is \(c_\chi>0\) such that, whenever
\[
 h=e^{-\kappa\sqrt A},
 \qquad 0<\kappa<c_\chi,
 \tag{V17.9}
\]
all sufficiently large \(A\) admit a family \(\mathcal M_A(d,h)\) of
disjoint pairs \(\{p,m\}\subset\mathcal P_A(d,h)\) satisfying
\[
 \chi(p)\ne\chi(m),
 \qquad
 \sum_{\{p,m\}\in\mathcal M_A(d,h)}c_{p,m}
 \gg_{\chi,d}W_A^{\mathrm p}(d,h).
 \tag{V17.10}
\]
Every prime edge occurs at most once, and the phase defects obey
\[
 |1-H_{p,m}|\ge
 \kappa_\chi:=min_{\substack{z,w\in G_\chi\\z\ne w}}|z-w|>0.
 \tag{V17.11}
\]
\end{proposition}

\begin{proof}
Choose two distinct values \(z,w\in G_\chi\).  The fixed-modulus prime
number theorem in arithmetic progressions, subtracted at the two endpoints
of \textup{(V17.7)}, shows that the numbers of primes with character value
\(z\) and \(w\) are asymptotic and each is a fixed positive proportion of
the layer population.  Condition \textup{(V17.9)} makes the main term
larger than the classical
\(e^{-c_\chi\sqrt A}\) error.  Match the smaller class injectively into the
larger.

Within the layer, the weights
\(\lambda_p=(\log p)/\sqrt p\) differ by at most a constant factor because
\(Ah=o(1)\).  Thus
\[
 c_{p,m}=\frac{\lambda_p\lambda_m}{\lambda_p+\lambda_m}
 \asymp_d\lambda_p\asymp_d\lambda_m.
 \tag{V17.12}
\]
The matched classes contain a fixed positive fraction of the prime weight,
which proves \textup{(V17.10)}.  The matching is edge-disjoint, and
\textup{(V17.11)} follows from finiteness of \(G_\chi\).
\end{proof}

By partial summation in the ordinary prime number theorem,
\[
 W_A^{\mathrm p}(d,h)
 \asymp_d Ah\,e^{Ad/2}
 =A\exp\!\left(\frac{Ad}{2}-\kappa\sqrt A\right)
 \tag{V17.13}
\]
under \textup{(V17.9)}.  Equations \textup{(V17.6)},
\textup{(V17.10)}, and \textup{(V17.11)} show that paired paths can supply
the correct exponential order of mass conductance throughout a common
spatial core.  Conductance and congestion are not the obstruction.

\section{The integer-resolution ceiling}

\begin{lemma}[Every two-edge displacement has only linear logarithmic
resolution]
If \(n,m<e^{2A}\) are distinct positive integers and
\(0<|\Delta_{n,m}|<1\), then
\[
 \boxed{
 \log\frac1{|\Delta_{n,m}|}
 \le2A+\log A.}
 \tag{V17.14}
\]
\end{lemma}

\begin{proof}
Assume \(n>m\).  As in \textup{(V16.9)},
\[
 \log\frac nm\ge\frac1n>e^{-2A}.
 \tag{V17.15}
\]
Hence \(|\Delta_{n,m}|>A^{-1}e^{-2A}\), which is equivalent to
\textup{(V17.14)}.  The other ordering is symmetric.
\end{proof}

The bound uses only integrality and the activation cutoff.  Choosing
consecutive integers, exceptionally close primes, or an optimized matching
cannot improve its exponential order.

\section{Universal clock lower bound for paired networks}

Let \(\Gamma_A^{(2)}\) be any finite family of two-edge near-cycles with
nontrivial holonomy and
\(0<|\Delta_{n,m}|\le2^{-8}\).  Its V16 clock load is
\[
 \mathcal T_{\Gamma_A^{(2)}}
 =\sum_{(n,m)\in\Gamma_A^{(2)}}
 \frac{32c_{n,m}}{\log(1/|\Delta_{n,m}|)}.
 \tag{V17.16}
\]

\begin{theorem}[Two-edge clock-capacity obstruction]
Every such family satisfies
\[
 \boxed{
 \mathcal T_{\Gamma_A^{(2)}}
 \ge\frac{32}{2A+\log A}
 \sum_{(n,m)\in\Gamma_A^{(2)}}c_{n,m}.}
 \tag{V17.17}
\]
Moreover, at every point \(t\), its conductance density obeys
\[
 \mathcal C_{\Gamma_A^{(2)}}(t)
 \le2\sum_{(n,m)\in\Gamma_A^{(2)}}c_{n,m}.
 \tag{V17.18}
\]
Consequently, if at one point of a prime layer the network is required to
dominate its source-degree weight in the normalized form
\[
 \mathcal C_{\Gamma_A^{(2)}}(t)
 \ge\operatorname{Cong}(\Gamma_A^{(2)})W,
 \tag{V17.19}
\]
then
\[
 \boxed{
 \frac{\mathcal T_{\Gamma_A^{(2)}}}
 {\operatorname{Cong}(\Gamma_A^{(2)})}
 \ge\frac{16W}{2A+\log A}.}
 \tag{V17.20}
\]
\end{theorem}

\begin{proof}
Apply \textup{(V17.14)} term by term in \textup{(V17.16)} to obtain
\textup{(V17.17)}.  Since
\(|1-H_{n,m}|\le2\), the definition \textup{(V16.28)} gives
\[
 \frac{|1-H_{n,m}|^2}{2\mathcal R_{n,m}}
 \le2c_{n,m},
 \tag{V17.21}
\]
and summation proves \textup{(V17.18)}.  Combine
\textup{(V17.18)} with \textup{(V17.19)}, then substitute the resulting
lower bound for \(\sum c_{n,m}\) into \textup{(V17.17)}.
\end{proof}

Take \(W=W_A^{\mathrm p}(d,h)\) with \textup{(V17.9)}.  Equations
\textup{(V17.13)} and \textup{(V17.20)} give
\[
 \frac{\mathcal T_{\Gamma_A^{(2)}}}
 {\operatorname{Cong}(\Gamma_A^{(2)})}
 \gg_{\chi,d}
 \exp\!\left(\frac{Ad}{2}-\kappa\sqrt A\right),
 \tag{V17.22}
\]
which is exponentially larger than the subunit requirement
\textup{(V16.32)}.  This applies to every two-edge pairing, not just the
particular matching constructed above.

\begin{firewall}[Large clock load belongs to this certificate]
Theorem \textup{(V17.17)} proves that two-edge paths cannot close the V16
near-cycle network certificate, whose translation remainder is paid by the
harmonic modulus \textup{(V15.10)}.  It does not prove that the original
two-cell magnetic form is adverse.  A different inequality could recycle
the translation remainder into other magnetic edges instead of charging it
directly to the clock.
\end{firewall}

\section{Why longer paths are the next necessary test}

For a two-edge path, the numerator and denominator of the rational label
are themselves active integers below \(e^{2A}\), forcing
\(\log(1/|\Delta|)=O(A)\).  A longer path can have products of many active
integers and hence reduced multiplicative height far beyond \(e^{2A}\).
By V16, it may achieve
\[
 \log\frac1{|\Delta_A(\gamma)|}
 \asymp\mathcal H_\times(\gamma)+\log A,
 \tag{V17.23}
\]
large enough to reduce the clock-to-mass ratio.  The price is increased
resistance and edge congestion.

V18 should therefore derive the optimal tradeoff among
\[
 \boxed{
 \mathcal H_\times(\gamma),\qquad
 \mathcal R(\gamma),\qquad
 \max_j s_j-\min_j s_j,qquad
 \text{and edge congestion}.}
 \tag{V17.24}
\]
The first two-edge audit shows the required scale: to dominate a layer of
weight \(W_A\) with subunit clock, the conductance-weighted average of
\(\log(1/|\Delta|)\) must be at least of order \(W_A\), hence exponential
in \(A\).  V18 must decide whether paths of that multiplicative height can
retain enough effective conductance and an admissible interval under a
bounded-congestion packing.

\begin{assessment}[V17 acquisition]
V17 completely resolves the two-edge branch of the V16 network program.
Fixed-modulus prime distribution supplies edge-disjoint, phase-separated
pairings whose conductance is comparable to the full prime-layer weight on
a common spatial core.  Nevertheless, the integrality of two active labels
limits logarithmic displacement resolution to \(O(A)\).  Any paired network
with enough mass conductance therefore has exponentially excessive harmonic
clock load, independently of how the matching is chosen.  Longer
multiplicative near-cycles, or a mechanism that recycles translation errors
through magnetic edges, are now mathematically necessary.  GRH is not
proved.
\end{assessment}

\section*{Version V17 append-only record}
\addcontentsline{toc}{section}{Version V17 append-only record}
The complete V16 source was retained before this continuation.  It had
\(266555\) bytes, \(8061\) lines, and SHA--256
\[
 \texttt{B43004CAC26E5E550B1BA9E0C9CD0563668C945BA17F22BB4AFCB109620E86E4}.
\]
V17 constructed full-conductance bounded-congestion prime pairings, proved
the universal integer-resolution ceiling for two-edge near-cycles, and
established their unavoidable exponential clock-capacity obstruction.

% =============================================================================
% END APPENDED V17 MATERIAL
% =============================================================================
% =============================================================================
% BEGIN APPENDED V18 MATERIAL
% =============================================================================

\section{An All-Path Resolution--Conductance Capacity Theorem}
\label{sec:grh-v18}

V17 proved that two-edge near-cycles cannot carry enough layer conductance
while keeping the V16 clock load subunit.  Longer paths appear at first to
evade that result: their numerator and denominator can have enormous
multiplicative height, producing exponentially finer residual displacement.
This chapter proves that the apparent gain is exactly limited by electrical
resistance and edge congestion.

The key quantity is not path length by itself but resolution-weighted
conductance.  For every path, its logarithmic resolution is at most its
multiplicative path height, while its conductance is no larger than the
conductance of each edge it uses.  After summing under bounded congestion,
the entire resolution budget is bounded by
\[
 \operatorname{Cong}(\Gamma_A)
 \sum_{n<e^{2A}}\lambda_n\log n
 \ll \operatorname{Cong}(\Gamma_A)Ae^A.
 \tag{V18.1}
\]
A weighted Cauchy--Schwarz inequality then gives a universal lower bound on
the clock load.  On every fixed upper shift layer \(1<d<2\), that lower
bound is exponential.  Thus arbitrarily long paths do not rescue the V16
certificate if each residual displacement is ultimately paid directly by
the harmonic translation modulus.

\section{Path height and resolution}

For a nontrivial near-cycle \(\gamma\), retain the V16 notation and put
\[
 c_\gamma:=\mathcal R(\gamma)^{-1},
 \qquad
 L_\gamma:=\log\frac1{|\Delta_A(\gamma)|},
 \tag{V18.2}
\]
and define its unreduced path height
\[
 \mathcal H_{\mathrm{path}}(\gamma)
 :=\sum_{j=1}^{|\gamma|}\log n_j.
 \tag{V18.3}
\]
Since cancellation can only decrease numerator and denominator,
\[
 \mathcal H_\times(\gamma)
 \le\mathcal H_{\mathrm{path}}(\gamma).
 \tag{V18.4}
\]
The arithmetic displacement estimate \textup{(V16.7)} therefore gives
\[
 \boxed{
 L_\gamma
 \le\mathcal H_{\mathrm{path}}(\gamma)+\log A.}
 \tag{V18.5}
\]

For a finite path family \(\Gamma_A\), write
\[
 K_A:=\operatorname{Cong}(\Gamma_A),
 \qquad
 C_A:=\sum_{\gamma\in\Gamma_A}c_\gamma.
 \tag{V18.6}
\]
The global prime resolution stock is
\[
 \mathcal B_A
 :=\sum_{\substack{n<e^{2A}\\ (n,q)=1}}
 \lambda_n\log n
 =\sum_{\substack{n<e^{2A}\\ (n,q)=1}}
 \frac{\Lambda(n)\log n}{\sqrt n}.
 \tag{V18.7}
\]

\begin{lemma}[Congestion bounds resolution-weighted conductance]
Every finite near-cycle family satisfies
\[
 \boxed{
 \sum_{\gamma\in\Gamma_A}c_\gamma L_\gamma
 \le K_A\mathcal B_A+(\log A)C_A.}
 \tag{V18.8}
\]
\end{lemma}

\begin{proof}
If an edge \(n\) occurs in \(\gamma\), then
\[
 \mathcal R(\gamma)\ge\lambda_n^{-1},
 \qquad c_\gamma\le\lambda_n.
 \tag{V18.9}
\]
Let \(m_\gamma(n)\) be its path multiplicity.  Using
\textup{(V18.5)} and interchanging two finite sums,
\[
 \begin{aligned}
 \sum_\gamma c_\gamma L_\gamma
 &\le\sum_n(\log n)
 \sum_\gamma m_\gamma(n)c_\gamma
 +(\log A)C_A\\
 &\le\sum_n(\log n)\lambda_n
 \sum_\gamma m_\gamma(n)+(\log A)C_A\\
 &\le K_A\mathcal B_A+(\log A)C_A.
 \end{aligned}
 \tag{V18.10}
\]
\end{proof}

The estimate is insensitive to path length, repeated factors, factorization
choices, and cancellation inside the rational label.  All such choices
consume the same congested edge stock.

\section{A universal lower bound on clock load}

The V16 clock load is
\[
 \mathcal T_{\Gamma_A}
 =32\sum_{\gamma\in\Gamma_A}\frac{c_\gamma}{L_\gamma}.
 \tag{V18.11}
\]

\begin{theorem}[All-path clock-capacity inequality]
For every finite family of admissible nontrivial near-cycles,
\[
 \boxed{
 \mathcal T_{\Gamma_A}
 \ge\frac{32C_A^2}
 {K_A\mathcal B_A+(\log A)C_A}.}
 \tag{V18.12}
\]
If at some point \(t_0\) the cycle conductance density is required to
dominate a source-degree weight \(W>0\), in the normalized form
\[
 \mathcal C_{\Gamma_A}(t_0)\ge K_AW,
 \tag{V18.13}
\]
then
\[
 \boxed{
 \frac{\mathcal T_{\Gamma_A}}{K_A}
 \ge\frac{8W^2}{\mathcal B_A+\frac12W\log A}.}
 \tag{V18.14}
\]
\end{theorem}

\begin{proof}
Weighted Cauchy--Schwarz gives
\[
 C_A^2
 =\left(\sum_\gamma
 \sqrt{\frac{c_\gamma}{L_\gamma}}
 \sqrt{c_\gamma L_\gamma}\right)^2
 \le
 \left(\sum_\gamma\frac{c_\gamma}{L_\gamma}\right)
 \left(\sum_\gamma c_\gamma L_\gamma\right).
 \tag{V18.15}
\]
Apply \textup{(V18.8)} and multiply by \(32\) to prove
\textup{(V18.12)}.

Since every phase defect is at most two, definition \textup{(V16.28)}
implies
\[
 \mathcal C_{\Gamma_A}(t_0)
 \le2C_A.
 \tag{V18.16}
\]
Thus \textup{(V18.13)} forces \(C_A\ge K_AW/2\).  The function
\(x\mapsto x^2/(K_A\mathcal B_A+x\log A)\) is increasing for
\(x>0\).  Substitute \(x=K_AW/2\) in \textup{(V18.12)} and divide by
\(K_A\) to obtain \textup{(V18.14)}.
\end{proof}

Unlike the V17 result, \textup{(V18.14)} applies simultaneously to paths of
all lengths and multiplicative heights.  Increasing height can improve one
path's resolution, but under congestion the improvement is paid from the
finite stock \(\mathcal B_A\).

\section{Evaluation of the global resolution stock}

\begin{proposition}[Prime resolution stock]
For fixed \(q\),
\[
 \boxed{\mathcal B_A\ll_q Ae^A.}
 \tag{V18.17}
\]
In fact the prime number theorem gives
\(\mathcal B_A\asymp Ae^A\).
\end{proposition}

\begin{proof}
Since \(\log n<2A\),
\[
 \mathcal B_A
 \le2A\sum_{n<e^{2A}}\frac{\Lambda(n)}{\sqrt n}.
 \tag{V18.18}
\]
Partial summation with the Chebyshev bound
\(\sum_{n\le x}\Lambda(n)\ll x\) shows that the last sum is
\(O(e^A)\), proving the upper bound.  Using the prime number theorem on any
fixed terminal subinterval supplies the matching lower order.
\end{proof}

Now fix \(1<d<2\) and take the PNT-scale layer from V17 with
\[
 h=e^{-\kappa\sqrt A},
 \qquad
 W=W_A^{\mathrm p}(d,h)asymp_d Ah e^{Ad/2}.
 \tag{V18.19}
\]
Because \(d<2\), the term \(W\log A=o(\mathcal B_A)\).  Therefore
\textup{(V18.14)} and \textup{(V18.17)} imply
\[
 \boxed{
 \frac{\mathcal T_{\Gamma_A}}{K_A}
 \gg_{q,d}
 A\exp\!\left(A(d-1)-2\kappa\sqrt A\right).}
 \tag{V18.20}
\]
This tends to infinity exponentially for every fixed \(d>1\).

\begin{corollary}[No V16 path family closes an upper shift layer]
No family of admissible near-cycles, of arbitrary lengths and with arbitrary
factorizations, can simultaneously satisfy
\[
 \mathcal C_{\Gamma_A}(t_0)\ge K_AW_A^{\mathrm p}(d,h)
 \quad\text{and}\quad
 \mathcal T_{\Gamma_A}<K_A
 \tag{V18.21}
\]
on a fixed layer \(1<d<2\) for all sufficiently large \(A\).
\end{corollary}

\begin{proof}
The first condition and \textup{(V18.20)} force
\(\mathcal T_{\Gamma_A}/K_A\to\infty\), contradicting the second.
\end{proof}

\section{What has actually been ruled out}

The result closes the longer-path question posed in V17 for the precise V16
certificate.  It is not merely a failure of a particular factorization or
matching.  The proof uses only four unavoidable facts:
\begin{enumerate}[label=\textup{(\roman*)},leftmargin=3em]
\item nontrivial rational holonomy has resolution bounded by multiplicative
height;
\item path conductance does not exceed any used edge conductance;
\item each edge has bounded congestion; and
\item the terminal translation error is paid through the harmonic modulus
\textup{(V15.10)}.
\end{enumerate}

\begin{firewall}[The magnetic form has more structure than path disposal]
Corollary \textup{(V18.21)} does not disprove relative magnetic coercivity.
It disproves certificates which telescope each path to one residual
translation and then dispose of every such residual independently through
the clock.  The original magnetic graph can route a residual translation
into further edges, share intermediate vertices, and cancel flows before a
terminal clock charge is made.  Those operations are absent from the sum of
one-path inequalities.
\end{firewall}

The remaining alternatives are now sharply separated.  V19 should either:
\begin{enumerate}[label=\textup{(\alph*)},leftmargin=3em]
\item construct a multicommodity magnetic flow in which near-cycle residuals
are recycled into other prime edges and only a small final defect reaches
the clock; or
\item move upstream from the sharp support-wall localization to a positive
smooth localization whose relative form admits a global Fourier or Mellin
square before individual path residuals are created.
\end{enumerate}
Repeating the V16 path certificate with longer paths cannot succeed by
\textup{(V18.20)}.

\begin{assessment}[V18 acquisition]
V18 closes the arbitrary-length path branch of the effective-conductance
program.  Resolution-weighted path conductance is globally bounded by the
congested prime stock \(\mathcal B_A\ll Ae^A\).  Weighted
Cauchy--Schwarz converts this stock bound into a universal clock lower bound,
which is exponential on every upper shift layer \(1<d<2\) once the network
has enough conductance to dominate the degree.  The obstruction is uniform
in path length, multiplicative height, factorization, and congestion.
Further progress must recycle residual translations inside the magnetic
graph or alter the localization upstream.  GRH is not proved.
\end{assessment}

\section*{Version V18 append-only record}
\addcontentsline{toc}{section}{Version V18 append-only record}
The complete V17 source was retained before this continuation.  It had
\(277388\) bytes, \(8378\) lines, and SHA--256
\[
 \texttt{18E4A825438A6BA5DB6C8A8D2565608CB40E6413412C67CC35B3E6B1DAA86033}.
\]
V18 bounded the global resolution-weighted conductance of every congested
path family and proved the resulting all-path exponential clock-capacity
obstruction on upper prime-shift layers.

% =============================================================================
% END APPENDED V18 MATERIAL
% =============================================================================
% =============================================================================
% BEGIN APPENDED V19 MATERIAL
% =============================================================================

\section{A Gaussian-Averaged Weil Criterion Without a Localization-Scale Support Jump}
\label{sec:grh-v19}

V18 ruled out every certificate which telescopes magnetic paths and pays
their terminal translation residuals independently through the harmonic
clock.  The remaining upstream option was to replace the sharp physical
localization by a positive smooth one before splitting the prime form.  The
localized Guinand--Weil formula in the attached arithmetic manuscript is
already written in terms of the autocorrelation
\[
 C_g(u)=\int_{\mathbb R}g(x+u)\overline{g(x)}\,dx,
 \tag{V19.1}
\]
so a continuous Gaussian quadratic partition acts on every prime cell by an
exact scalar overlap factor.

This chapter constructs that partition and proves an exact averaged Weil
identity.  The prime side acquires the smooth multiplier
\(e^{-u^2/(8\sigma^2)}\), while the archimedean Fourier symbol is convolved
with the matching Gaussian probability density.  Positivity of these
averaged forms along any cofinal sequence \(\sigma\to\infty\) is equivalent
to the original Weil criterion and hence to GRH.  The construction does not
prove positivity, but it removes the localization-scale support jump without
weakening the target and supplies a new exact arena for a relative Fourier
square.

\section{Continuous Gaussian quadratic partition}

For \(\sigma>0\) and \(y\in\mathbb R\), put
\[
 \eta_{\sigma,y}(x)
 :=(2\pi\sigma^2)^{-1/4}
 \exp\!\left(-\frac{(x-y)^2}{4\sigma^2}\right).
 \tag{V19.2}
\]

\begin{lemma}[Exact quadratic partition and overlap]
For every \(x,x'\in\mathbb R\),
\[
 \int_{\mathbb R}\eta_{\sigma,y}(x)^2dy=1
 \tag{V19.3}
\]
and
\[
 \boxed{
 \int_{\mathbb R}\eta_{\sigma,y}(x)
 \eta_{\sigma,y}(x')dy
 =s_\sigma(x-x'),
 \qquad
 s_\sigma(u):=e^{-u^2/(8\sigma^2)}.}
 \tag{V19.4}
\]
\end{lemma}

\begin{proof}
The square in \textup{(V19.3)} is the centered Gaussian probability density
in the variable \(y\).  For \textup{(V19.4)}, complete the square:
\[
 (x-y)^2+(x'-y)^2
 =2\left(y-\frac{x+x'}2\right)^2+\frac{(x-x')^2}{2}.
 \tag{V19.5}
\]
The remaining normalized Gaussian integrates to one.
\end{proof}

If \(g\in C_c^\infty(\mathbb R)\), then every
\(\eta_{\sigma,y}g\) is again a smooth compactly supported Weil test.  The
partition preserves its norm:
\[
 \int_{\mathbb R}\|\eta_{\sigma,y}g\|_2^2dy=\|g\|_2^2.
 \tag{V19.6}
\]

\section{Exact damping of every prime autocorrelation}

\begin{proposition}[Gaussian autocorrelation localization]
For every \(u\in\mathbb R\),
\[
 \boxed{
 \int_{\mathbb R}C_{\eta_{\sigma,y}g}(u)dy
 =s_\sigma(u)C_g(u).}
 \tag{V19.7}
\]
Consequently the averaged prime form is
\[
 \boxed{
 \begin{aligned}
 &2\Re\sum_{n\ge2}\frac{\Lambda(n)\chi(n)}{\sqrt n}
 \int_{\mathbb R}C_{\eta_{\sigma,y}g}(\log n)dy\\
 &\qquad=
 2\Re\sum_{n\ge2}\frac{\Lambda(n)\chi(n)}{\sqrt n}
 e^{-(\log n)^2/(8\sigma^2)}C_g(\log n).
 \end{aligned}}
 \tag{V19.8}
\]
\end{proposition}

\begin{proof}
Insert the definition of autocorrelation, apply Fubini, and use
\textup{(V19.4)} with \(x'=x+u\).  The prime sum is finite because compact
support of \(g\) makes \(C_g(u)=0\) for all sufficiently large \(|u|\).
\end{proof}

Unlike the hard activation factor \(\mathbf1_{\{\log n<2A\}}\), the
Gaussian multiplier in \textup{(V19.8)} varies smoothly with both the prime
label and the localization scale.  The compact support of the individual
test still makes the formula finite, but changing \(\sigma\) creates no new
prime-cell norm jump.

\section{Exact Gaussian convolution of the archimedean symbol}

Write the fixed-character archimedean symbol as
\[
 a_\chi(t)
 :=\log\frac f\pi
 +\Re\frac{\Gamma'}{\Gamma}
 \!\left(\frac{1/2+\kappa+it}{2}\right),
 \tag{V19.9}
\]
so the attached manuscript's archimedean form is
\[
 \mathcal A_\chi[g]
 =\frac1{2\pi}\int_{\mathbb R}
 |\widehat g(t)|^2a_\chi(t)dt.
 \tag{V19.10}
\]
Define the Gaussian probability density
\[
 \rho_\sigma(v)
 :=\sqrt{\frac2\pi}\,\sigma e^{-2\sigma^2v^2},
 \qquad
 \int_{\mathbb R}\rho_\sigma(v)dv=1.
 \tag{V19.11}
\]

\begin{theorem}[Averaged archimedean localization]
For every compactly supported smooth \(g\),
\[
 \boxed{
 \int_{\mathbb R}\mathcal A_\chi[\eta_{\sigma,y}g]dy
 =\frac1{2\pi}\int_{\mathbb R}|\widehat g(t)|^2
 (a_\chi*\rho_\sigma)(t)dt.}
 \tag{V19.12}
\]
\end{theorem}

\begin{proof}
Let \(K_\chi(x-x')\) be the distribution kernel of the Fourier multiplier
\(a_\chi(D)\).  Averaging the localized quadratic form multiplies this
kernel by the overlap \(s_\sigma(x-x')\) from \textup{(V19.4)}.  With the
Fourier convention of the manuscript,
\[
 \widehat{s_\sigma}(v)
 =\sqrt{8\pi}\,\sigma e^{-2\sigma^2v^2},
 \qquad
 \frac1{2\pi}\widehat{s_\sigma}(v)=\rho_\sigma(v).
 \tag{V19.13}
\]
Multiplication of the kernel by \(s_\sigma\) therefore convolves its symbol
with \(\rho_\sigma\), proving \textup{(V19.12)}.  The logarithmic growth of
\(a_\chi\) and Gaussian decay justify the convolution; compact support and
smoothness of \(g\) justify the quadratic-form limit.
\end{proof}

\section{The exact Gaussian-averaged Weil form}

Define
\[
 \begin{aligned}
 \mathcal W_{\chi,\sigma}^{\mathrm{av}}[g]
 :={}&\frac1{2\pi}\int_{\mathbb R}|\widehat g(t)|^2
 (a_\chi*\rho_\sigma)(t)dt\\
 &-2\Re\sum_{n\ge2}
 \frac{\Lambda(n)\chi(n)}{\sqrt n}
 e^{-(\log n)^2/(8\sigma^2)}C_g(\log n).
 \end{aligned}
 \tag{V19.14}
\]

\begin{theorem}[Exact averaging identity]
For every \(\sigma>0\) and \(g\in C_c^\infty(\mathbb R)\),
\[
 \boxed{
 \mathcal W_{\chi,\sigma}^{\mathrm{av}}[g]
 =\int_{\mathbb R}
 \mathcal W_\chi[\eta_{\sigma,y}g]dy.}
 \tag{V19.15}
\]
\end{theorem}

\begin{proof}
Apply \textup{(V19.12)} to the archimedean term and
\textup{(V19.8)} to the prime term in the localized Guinand--Weil formula.
\end{proof}

This identity retains the completed conductor and gamma factor exactly; the
archimedean side is not replaced by a scalar floor.  It also retains all
off-diagonal prime correlations through \(C_g\).

\section{A smooth criterion equivalent to GRH}

\begin{theorem}[Gaussian-averaged Weil criterion]
For a fixed primitive character \(\chi\), the following are equivalent:
\begin{enumerate}[label=\textup{(\roman*)},leftmargin=3em]
\item GRH holds for the paired \(\chi,\overline\chi\) channel;
\item \(\mathcal W_{\chi,\sigma}^{\mathrm{av}}[g]\ge0\) for every
\(\sigma>0\) and every \(g\in C_c^\infty(\mathbb R)\);
\item there exists a sequence \(\sigma_j\to\infty\) such that
\(\mathcal W_{\chi,\sigma_j}^{\mathrm{av}}[g]\ge0\) for every \(j\) and
every \(g\in C_c^\infty(\mathbb R)\).
\end{enumerate}
\end{theorem}

\begin{proof}
Under GRH, the Weil criterion makes
\(\mathcal W_\chi[h]\ge0\) for every compactly supported smooth \(h\).
The averaging identity \textup{(V19.15)} then proves (ii), and (ii) implies
(iii).

Conversely, fix \(g\).  Since \(\rho_\sigma\) is an approximate identity as
\(\sigma\to\infty\),
\[
 (a_\chi*\rho_\sigma)(t)\longrightarrow a_\chi(t).
 \tag{V19.17}
\]
The convergence is dominated against the rapidly decaying
\(|\widehat g(t)|^2\).  Also
\(e^{-(\log n)^2/(8\sigma^2)}\to1\), and the prime sum is finite for fixed
\(g\).  Hence
\[
 \mathcal W_{\chi,\sigma}^{\mathrm{av}}[g]
 \longrightarrow\mathcal W_\chi[g].
 \tag{V19.18}
\]
Condition (iii) gives \(\mathcal W_\chi[g]\ge0\).  Since this holds for
every test \(g\), the Weil criterion proves GRH.
\end{proof}

The word ``cofinal'' is safe here because \textup{(V19.18)} supplies the
missing bridge explicitly; this is precisely the bridge absent from a sparse
sequence of sharp support walls in V14.

\section{Smooth scale resolution of the prime side}

For every \(u\ne0\),
\[
 \frac{\partial}{\partial\log\sigma}s_\sigma(u)
 =k_\sigma(u)
 :=\frac{u^2}{4\sigma^2}e^{-u^2/(8\sigma^2)}\ge0,
 \tag{V19.19}
\]
and
\[
 \boxed{
 \int_0^\infty k_\sigma(u)\frac{d\sigma}{\sigma}=1.}
 \tag{V19.20}
\]
Thus the unsmoothed prime form has the exact positive scale resolution
\[
 \begin{aligned}
 &2\Re\sum_{n\ge2}\frac{\Lambda(n)\chi(n)}{\sqrt n}C_g(\log n)\\
 &\quad=\int_0^\infty
 2\Re\sum_{n\ge2}\frac{\Lambda(n)\chi(n)}{\sqrt n}
 k_\sigma(\log n)C_g(\log n)\frac{d\sigma}{\sigma}.
 \end{aligned}
 \tag{V19.21}
\]
For compactly supported \(g\), all sums are finite, so no interchange issue
arises.  Each prime cell now enters every scale smoothly, concentrated near
\(\log n\asymp\sigma\), instead of appearing at one hard wall.

The absolute scale mass is still large.  The prime number theorem and the
change of variable \(t=\log x\) give the saddle estimate
\[
 \sum_{n\ge2}\frac{\Lambda(n)}{\sqrt n}k_\sigma(\log n)
 \asymp \sigma^3e^{\sigma^2/2}
 \tag{V19.22}
\]
up to fixed polynomial factors as \(\sigma\to\infty\); the exponent follows
from
\[
 \frac t2-\frac{t^2}{8\sigma^2}
 =\frac{\sigma^2}{2}
 -\frac{(t-2\sigma^2)^2}{8\sigma^2}.
 \tag{V19.23}
\]
Thus Gaussian localization removes the jump, not the need for signed
character cancellation or relative completed-form geometry.

\begin{firewall}[Smoothing is an equivalent reformulation, not positivity]
The Gaussian-averaged criterion does not establish that
\(\mathcal W_{\chi,\sigma}^{\mathrm{av}}\) is nonnegative.  It proves that
such nonnegativity is a GRH-equivalent target with smooth prime weights and
an exact convolved archimedean symbol.  Bounding the prime term absolutely
would reproduce an exponential mass at the saddle \textup{(V19.23)} and
would not close the form.
\end{firewall}

V20 is the next mandatory corpus sweep.  It should compare the exact
Gaussian form \textup{(V19.14)} against every parallel note created since
the V10 sweep, with special attention to heat kernels, Gaussian arithmetic
banks, completed gamma-factor estimates, and signed square-function
identities.  The post-sweep target is a relative Fourier-square or Schur
factorization for \(\mathcal W_{\chi,\sigma}^{\mathrm{av}}\), not another
absolute prime sum.

\begin{assessment}[V19 acquisition]
V19 executes the upstream localization change forced by V18.  A continuous
Gaussian quadratic partition damps each prime autocorrelation by
\(e^{-u^2/(8\sigma^2)}\) and convolves the exact archimedean symbol with the
matching Gaussian probability density.  The resulting averaged Weil forms
are exactly integrals of genuine Weil tests, converge back to the original
form, and give a criterion equivalent to GRH along any cofinal smooth scale
sequence.  A positive scale derivative resolves the prime side without
activation jumps.  The exponential absolute saddle remains, so the next
required ingredient is a relative signed factorization.  GRH is not proved.
\end{assessment}

\section*{Version V19 append-only record}
\addcontentsline{toc}{section}{Version V19 append-only record}
The complete V18 source was retained before this continuation.  It had
\(287214\) bytes, \(8680\) lines, and SHA--256
\[
 \texttt{D5D96681E851D5C678D20D782509ECDD607ECA3009DEA965A1A066EC8AFF3820}.
\]
V19 constructed the exact Gaussian-averaged Weil form, proved its
archimedean and prime localization identities, established its equivalence
to the original Weil criterion, and derived a positive smooth scale
resolution of the prime side.

% =============================================================================
% END APPENDED V19 MATERIAL
% =============================================================================
% =============================================================================
% BEGIN APPENDED V20 MATERIAL
% =============================================================================

\section{V20: the heat-time positivity threshold and a relative Loewner target}
\label{sec:grh-v20}

V19 replaced sharp support activation by an exact Gaussian quadratic
partition.  This mandatory tenth-version sweep finds that the resulting
family has more structure than was used there.  After the inverse change of
scale
\[
 t=\frac{1}{8\sigma^2},
 \tag{V20.1}
\]
the averaged Weil forms form a heat semigroup under Gaussian averages of
frequency modulation.  Positivity therefore propagates exactly toward
larger heat time.  Moreover, elementary lower bounds for the completed
gamma symbol show that every fixed-character form is coercive at all
sufficiently large heat times.  GRH is consequently equivalent to saying
that the first heat time at which positivity appears is zero.

This is still a reformulation, not a proof of GRH.  Its gain is that it
turns an unstructured cofinal limit into a one-parameter positivity-front
problem and isolates a precise relative differential inequality which would
force that front back to zero.

\subsection{The mandated parallel-note sweep}

The snapshot was taken on 10 September 2026.  Excluding the present GRH
lineage, the folder contained twenty-nine \texttt{.tex} files.  As in V10,
the canonical manifest consists of the lexicographically ordered rows
\[
 \texttt{filename|byte-count|SHA--256}
 \tag{V20.2}
\]
with a terminal newline.  Its UTF--8 SHA--256 is
\[
 \boxed{
 \texttt{3C744E6E540E721C4D6E9719133D53B733BBCBFC73644190FE3CE40509B63222}.}
 \tag{V20.3}
\]
The complete groups inspected were
\[
\begin{array}{p{0.20\textwidth}|p{0.68\textwidth}}
\text{programme}&\text{files inspected}\cr\hline
\text{Navier--Stokes}
 &\texttt{V100\_f1},\ \texttt{V100\_f2},\ \texttt{V100\_f3},
   \texttt{V26}\cr
\text{SM--Einstein}
 &\texttt{V157\_f1},\ \texttt{V100\_f2}--\texttt{V100\_f13},
   \texttt{V93}\cr
\text{Yang--Mills}
 &\texttt{V135\_f1},\ \texttt{V100\_f2}--\texttt{V100\_f7},
   \texttt{V65\_f8},\ \texttt{V100\_f9},\ \texttt{V29}\cr
\text{parallel Yang--Mills}
 &\text{parallel YM V5}\cr
\end{array}
\tag{V20.4}
\]
The live terminal rows at finalization were
\[
\begin{array}{p{0.46\textwidth}|r|p{0.34\textwidth}}
\text{checkpoint}&\text{bytes}&\text{SHA--256}\cr\hline
\text{NS V26}
 &278283&{\scriptsize \texttt{82C887F8C2C69E4F28FAD7C3DC8DE5B}
 \allowbreak\texttt{9099CB5A4BF431EBA1FFF3E91935978AD}}\cr
\text{SM--Einstein V93}
 &1103350&{\scriptsize \texttt{4F5EA5CA6367BD7E2ADDFBD71DB8B60A}
 \allowbreak\texttt{85B9A4C696CE2F6D66A41EB3DAA784F2}}\cr
\text{Yang--Mills V29}
 &232987&{\scriptsize \texttt{85AEAC9219611F19F1C2F98A2F92C425}
 \allowbreak\texttt{3B7F32F7120F17183099CCA237B0F0B8}}\cr
\text{parallel YM V5}
 &54174&{\scriptsize \texttt{306F1BB84CFA8A8D47EA569EC17E9545}
 \allowbreak\texttt{F9867C8D32B548EC9D9C444B4F3C0512}}\cr
\end{array}
\tag{V20.5}
\]

Twenty-seven archive rows are byte-for-byte unchanged from the V10
checkpoint, so their result ledgers, terminal audits, decisions, and claim
boundaries were hash-gated against the V10 reading rather than reinterpreted
as new results.  The two rolling lineages were read through their intervening
sections: SM--Einstein V59--V93 and Yang--Mills V20--V29.  The SM file
advanced from V91 through V93 during this sweep, so the manifest and the terminal
reading were refreshed before finalization.

\begin{assessment}[V20 transfer ledger]
\label{ass:v20-sweep}
Only the following proof mechanisms transfer to the present problem.
\begin{enumerate}[label=\textup{(\roman*)},leftmargin=3em]
\item SM--Einstein V58--V59 distinguishes positivity of the full reflected
      Gaussian cross-kernel from positivity of a local Hessian and performs
      the null quotient before taking a logarithmic generator.  This supports
      treating the complete Weil form, rather than its archimedean and prime
      pieces separately.
\item SM--Einstein V76 proves positivity by a heat-kernel factorization and
      diagonal congruence; V84 requires one common feature factorization
      before applying a Schur product.  Separate diagonal or frozen-background
      positivity is explicitly insufficient.  The relevant import here is
      the full-kernel heat semigroup, not any SM numerical bound.
\item SM--Einstein V91--V92 shows that a parameter derivative must retain its
      contact and covariance pieces and that convergence of an integrated
      quantity does not control concentration.  Applied here, the heat
      derivative must keep both the convolved gamma term and the signed prime
      response; neither may be discarded by an absolute estimate.  V93 then
      distinguishes a genuine local gap from one manufactured by a global
      refresh.  Analogously, large-heat positivity below is a control
      regularization, not evidence that the unsmoothed arithmetic form is
      positive.
\item Yang--Mills V20--V29 separates failure of an outer hull or a fixed
      splitting from failure of the physical pencil.  Its exact V29 witness
      closes one constant splitting but not the direct pencil.  Thus failure
      of one finite Fourier square for the heat-deformed Weil form should lead
      to an adaptive local factorization or the direct relative form, not to a
      negative conclusion about that form.
\end{enumerate}
The companion Hessians, graph gaps, covariance matrices, torus radii, and
certificate witnesses are not arithmetic data and are not imported.
\end{assessment}

\subsection{Heat-time form and the exact dephasing semigroup}

Let
\[
 H_t(v):=\frac{1}{\sqrt{4\pi t}}e^{-v^2/(4t)},
 \qquad t>0,
 \tag{V20.6}
\]
so that \(H_t*H_s=H_{t+s}\) and
\[
 \int_{\mathbb R}H_t(v)e^{ivu}\,dv=e^{-tu^2}.
 \tag{V20.7}
\]
Write
\[
 a_{\chi,t}:=a_\chi*H_t
 \tag{V20.8}
\]
and define, for \(g\in C_c^\infty(\mathbb R)\),
\[
 \begin{aligned}
 \mathcal W_{\chi,t}[g]
 :={}&\frac1{2\pi}\int_{\mathbb R}
       |\widehat g(\xi)|^2a_{\chi,t}(\xi)d\xi\\
 &-2\Re\sum_{n\ge2}\frac{\Lambda(n)\chi(n)}{\sqrt n}
       e^{-t(\log n)^2}C_g(\log n).
 \end{aligned}
 \tag{V20.9}
\]
Equations \textup{(V19.11)} and \textup{(V20.1)} give
\[
 H_t=\rho_\sigma,
 \qquad
 \mathcal W_{\chi,t}=\mathcal W_{\chi,\sigma}^{\mathrm{av}}.
 \tag{V20.10}
\]

For \(v\in\mathbb R\), let
\[
 (M_vg)(x):=e^{ivx}g(x).
 \tag{V20.11}
\]

\begin{theorem}[Exact Gaussian dephasing identity]
\label{thm:v20-dephasing}
For all \(s,t>0\) and \(g\in C_c^\infty(\mathbb R)\),
\[
 \boxed{
 \mathcal W_{\chi,t+s}[g]
 =\int_{\mathbb R}H_s(v)\mathcal W_{\chi,t}[M_vg],dv.}
 \tag{V20.12}
\]
Consequently, if \(\mathcal W_{\chi,t}\) is nonnegative on all Weil tests
at one time \(t\), then it is nonnegative at every time \(t+s\).
\end{theorem}

\begin{proof}
Modulation preserves compact support and satisfies
\[
 \widehat{M_vg}(\xi)=\widehat g(\xi-v),
 \qquad
 C_{M_vg}(u)=e^{ivu}C_g(u).
 \tag{V20.13}
\]
Averaging the first identity against \(H_s\) convolves
\(a_{\chi,t}\) once more and gives \(a_{\chi,t+s}\).  Averaging the second
and using \textup{(V20.7)} changes
\(e^{-tu^2}\) into \(e^{-(t+s)u^2}\).  This proves
\textup{(V20.12)}.  Its integrand is nonnegative when the time-\(t\) form is
nonnegative, proving the last assertion.
\end{proof}

At the kernel level, \textup{(V20.12)} is the unital completely positive map
\[
 \Phi_s(T)=\int_{\mathbb R}H_s(v)M_v^*TM_v,dv,
 \qquad
 \Phi_s\Phi_t=\Phi_{s+t}.
 \tag{V20.14}
\]
It multiplies a kernel \(T(x,x')\) by
\(e^{-s(x-x')^2}\).  Thus the V19 window average is simultaneously a
continuous Schur multiplier and a Gaussian mixture of unitary congruences.
This is the common full-kernel factorization required by the sweep; it does
not split the signed prime form from its completed diagonal.

\subsection{Eventual coercivity from the completed gamma symbol}

\begin{lemma}[Uniform heat lower bound]
\label{lem:v20-gamma-lower}
There is a character-dependent constant \(C_\chi<\infty\) such that
\[
 a_\chi(\xi)\ge \log(1+|\xi|)-C_\chi
 \tag{V20.15}
\]
for every \(\xi\in\mathbb R\).  If
\[
 p_0:=1-\frac{1}{2\sqrt\pi}>0,
 \tag{V20.16}
\]
then
\[
 \boxed{
 \inf_{\xi\in\mathbb R}a_{\chi,t}(\xi)
 \ge p_0\log\left(1+\frac{\sqrt t}{2}\right)-C_\chi.}
 \tag{V20.17}
\]
\end{lemma}

\begin{proof}
In the fixed half-plane containing
\((1/2+\kappa+i\xi)/2\), the standard digamma asymptotic is
\(\Re\psi(z)=\log|z|+O(|z|^{-1})\).  It gives
\textup{(V20.15)} outside a compact interval; continuity supplies the same
bound on the remaining interval after enlarging \(C_\chi\).

For fixed \(\xi\), an interval of radius \(\sqrt t/2\) has
\(H_t\)-mass at most its length times the maximum of \(H_t\), namely
\(1/(2\sqrt\pi)\).  Hence
\[
 \int_{|v-\xi|\ge\sqrt t/2}H_t(v)dv\ge p_0.
 \tag{V20.18}
\]
Convolve \textup{(V20.15)} with \(H_t\), use nonnegativity of
\(\log(1+|\cdot|)\), and restrict to the set in
\textup{(V20.18)}.
\end{proof}

Put
\[
 R(t):=2\sum_{n\ge2}\frac{\Lambda(n)}{\sqrt n}
                 e^{-t(\log n)^2}.
 \tag{V20.19}
\]
For every \(t>0\) this series converges, and \(R(t)\to0\) as
\(t\to\infty\).  Since translations are unitary on \(L^2(\mathbb R)\),
\[
 \left|2\Re\sum_{n\ge2}\frac{\Lambda(n)\chi(n)}{\sqrt n}
 e^{-t(\log n)^2}C_g(\log n)\right|
 \le R(t)\|g\|_2^2.
 \tag{V20.20}
\]
Plancherel and Lemma~\ref{lem:v20-gamma-lower} therefore give
\[
 \boxed{
 \mathcal W_{\chi,t}[g]
 \ge
 \left[
 p_0\log\left(1+\frac{\sqrt t}{2}\right)-C_\chi-R(t)
 \right]\|g\|_2^2.}
 \tag{V20.21}
\]
In particular, there is a finite \(T_\chi\) for which the form is coercive
for every \(t\ge T_\chi\).  This is only a large-heat statement: the
Gaussian damping has then erased almost all arithmetic off-diagonal data.

\subsection{The critical heat time}

Define the positivity set and its threshold by
\[
 \mathfrak P_\chi
 :=\{t>0:\mathcal W_{\chi,t}[g]\ge0
              \text{ for every }g\in C_c^\infty(\mathbb R)\},
 \qquad
 \vartheta_\chi:=\inf\mathfrak P_\chi.
 \tag{V20.22}
\]

\begin{theorem}[Heat-threshold criterion for GRH]
\label{thm:v20-threshold}
The set \(\mathfrak P_\chi\) is a nonempty upper interval.  If
\(\vartheta_\chi>0\), it contains its left endpoint and the threshold form
has no strictly positive uniform \(L^2\) floor.  Most importantly,
\[
 \boxed{
 \text{GRH for the paired }\chi,\overline\chi\text{ channel}
 \quad\Longleftrightarrow\quad
 \vartheta_\chi=0.}
 \tag{V20.23}
\]
\end{theorem}

\begin{proof}
Nonemptiness follows from \textup{(V20.21)}, and the upper-interval property
follows from Theorem~\ref{thm:v20-dephasing}.  On every compact subinterval
of \((0,\infty)\), the difference
\(a_{\chi,t}-a_{\chi,s}\) tends uniformly to zero as \(t\to s\): one may
differentiate the heat convolution and use boundedness of
\(a_\chi''\).  The prime-form difference tends to zero in operator norm by
the absolutely convergent majorant with an extra factor
\((\log n)^2\).  Thus the forms are locally norm-continuous in heat time.
It follows that a positive left endpoint is included.  A strictly positive
floor there would persist for slightly smaller time, contradicting the
definition of \(\vartheta_\chi\).

If GRH holds, the original Weil form is nonnegative and
\textup{(V19.15)}, or equivalently \textup{(V20.12)} with initial time zero,
makes every positive-time form nonnegative.  Hence
\(\vartheta_\chi=0\).  Conversely, if the threshold is zero, the
upper-interval property gives \(\mathfrak P_\chi=(0,\infty)\).  Letting
\(t\downarrow0\) in \textup{(V20.9)} gives
\(\mathcal W_\chi[g]\ge0\) for every compactly supported smooth \(g\), by
the same dominated convergence and finite-prime-sum argument as
\textup{(V19.17)}--\textup{(V19.18)}.  Weil's criterion proves GRH.
\end{proof}

The theorem turns a cofinal family into a front: generic dephasing flows can
cross the positive cone at a positive time.  Indeed, for
\[
 Q_0=\begin{pmatrix}1&b\\ b&1\end{pmatrix},
 \qquad b>1,
 \tag{V20.24}
\]
and two position labels separated by one, Gaussian dephasing gives
\[
 Q_t=\begin{pmatrix}1&be^{-t}\\ be^{-t}&1\end{pmatrix}
 \succeq0
 \quad\Longleftrightarrow\quad t\ge\log b.
 \tag{V20.25}
\]
Thus complete positivity and eventual coercivity alone do not force
\(\vartheta_\chi=0\); the missing input must use the arithmetic relation
between the completed diagonal and the signed prime shifts.

\subsection{A relative heat-flow inequality sufficient for GRH}

The heat equation and absolute convergence for \(t>0\) give the exact
derivative
\[
 \begin{aligned}
 \dot{\mathcal W}_{\chi,t}[g]
 ={}&\frac1{2\pi}\int_{\mathbb R}|\widehat g(\xi)|^2
       a_{\chi,t}''(\xi)d\xi\\
 &+2\Re\sum_{n\ge2}\frac{\Lambda(n)\chi(n)}{\sqrt n}
       (\log n)^2e^{-t(\log n)^2}C_g(\log n).
 \end{aligned}
 \tag{V20.26}
\]
The first line is the heat contact term and the second is the signed prime
response.  Their sum, rather than either absolute value, is the object to
control.

\begin{proposition}[Backward relative-Loewner criterion]
\label{prop:v20-relative-flow}
Suppose that for some \(T\ge T_\chi\) there is a locally integrable real
function \(c\) on \((0,T]\) such that
\[
 \boxed{
 \dot{\mathcal W}_{\chi,t}[g]
 \le c(t)\mathcal W_{\chi,t}[g]}
 \tag{V20.27}
\]
for every \(g\in C_c^\infty(\mathbb R)\) and every \(0<t\le T\).
Then \(\vartheta_\chi=0\), and GRH holds for the paired character channel.
\end{proposition}

\begin{proof}
Fix \(g\) and \(0<\varepsilon<T\).  With
\(C(t)=\int_\varepsilon^tc(r)dr\), inequality \textup{(V20.27)} says
\[
 \frac d{dt}\left(e^{-C(t)}\mathcal W_{\chi,t}[g]\right)\le0.
 \tag{V20.28}
\]
Since \(T\ge T_\chi\), the value at \(T\) is nonnegative.  Therefore the
value at \(\varepsilon\) is nonnegative.  This holds for every test and
every \(\varepsilon>0\), so \(\mathfrak P_\chi=(0,\infty)\).  Apply
Theorem~\ref{thm:v20-threshold}.
\end{proof}

Inequality \textup{(V20.27)} is the precise post-sweep target.  It is a
relative form inequality on the complete heat-deformed Weil operator, not an
absolute prime estimate.  Equivalently, it asks that the heat velocity lie
in the backward tangent cone of the positive Loewner cone.  At a hypothetical
positive threshold, normalized vectors with
\(\mathcal W_{\chi,\vartheta_\chi}[g_j]\to0\) must exist; a uniform version
of \textup{(V20.27)} would forbid the positive crossing exhibited by
\textup{(V20.25)}.

\begin{firewall}[The heat threshold has not been shown to vanish]
Neither the parallel notes nor the semigroup identity supplies
\textup{(V20.27)} for the arithmetic Weil form.  Large heat time is positive
because dephasing suppresses the prime shifts and the convolved gamma symbol
grows; reversing that smoothing is ill-conditioned in general, as
\textup{(V20.24)}--\textup{(V20.25)} prove.  The new result is an exact
critical-time formulation and a sufficient relative-Loewner target.  GRH
remains open.
\end{firewall}

\begin{assessment}[V20 acquisition]
V20 completes the mandatory corpus sweep, including live SM V93 and
Yang--Mills V29.  It identifies the V19 Gaussian family as a completely
positive dephasing semigroup, proves exact forward propagation of
positivity, proves eventual coercivity from the completed gamma symbol and
Gaussian prime damping, and defines a finite critical heat time whose
vanishing is equivalent to GRH.  It also reduces a valid backward attack to
the relative differential inequality \textup{(V20.27)} and gives a finite
counterexample showing why semigroup positivity alone cannot supply it.  The
next version should analyze the heat velocity on near-null directions of
the complete form, using a direct relative Schur object or an adaptive local
factorization rather than a fixed global square.
\end{assessment}

\section*{Version V20 append-only record}
\addcontentsline{toc}{section}{Version V20 append-only record}
The complete V19 source was retained before this continuation.  It had
\(298662\) bytes, \(9017\) lines, and SHA--256
\[
 \texttt{488EB91241A4A159AFC1CF51BC081FAB071115764FC1092205E1FB8BBED42DE5}.
\]
V20 performed the required full parallel-note sweep and constructed the
heat-time semigroup, the eventual-coercivity bound, the GRH-equivalent
critical threshold, and the backward relative-Loewner target.

% =============================================================================
% END APPENDED V20 MATERIAL
% =============================================================================
% =============================================================================
% BEGIN APPENDED V21 MATERIAL
% =============================================================================

\section{V21: finite-window critical packets and heat-flow carr\'e du champ}
\label{sec:grh-v21}

V20 reduced the paired-character GRH problem to the vanishing of a heat-time
positivity threshold, but its sufficient differential inequality was stated
on every test vector.  That is much stronger than the geometry of a first
spectral crossing requires.  This continuation restricts the complete form
to one fixed physical support window.  The logarithmic growth of the gamma
symbol then gives compact resolvent, so a positive critical time has an
actual finite-dimensional null space.  The Gaussian dephasing identity
computes the heat velocity exactly on that null space.

The outcome is a finite-window obstruction theorem: GRH fails if and only if
one support window possesses a positive critical time; at such a time the
kernel compression of the heat velocity is a nonzero positive Gram matrix.
For a simple critical kernel the lowest eigenvalue crosses zero
transversely, with an explicit positive derivative.  This replaces the
abstract near-null sequence in V20 by a compact critical packet while
retaining the full completed arithmetic form.

\subsection{Closed forms on a fixed support window}

For \(A>0\), let
\[
 \mathcal H_A
 :=\{g\in L^2(\mathbb R):\operatorname{supp}g\subset[-A,A]\},
 \tag{V21.1}
\]
and define the logarithmic form domain
\[
 \mathcal D_A
 :=\left\{g\in\mathcal H_A:
 \int_{\mathbb R}\log(2+|\xi|)|\widehat g(\xi)|^2d\xi<\infty
 \right\}.
 \tag{V21.2}
\]
The core \(C_c^\infty((-A,A))\) is dense in this domain.  Let
\(\mathcal B_{\chi,t}^{A}\) be the Hermitian polarization of
\(\mathcal W_{\chi,t}\) on \(\mathcal D_A\), including \(t=0\), and write
\[
 q_{A,t}[g]:=\mathcal B_{\chi,t}^{A}(g,g).
 \tag{V21.3}
\]

\begin{proposition}[Compact finite-window realization]
\label{prop:v21-compact}
For every \(A>0\) and \(t\ge0\), the form \(q_{A,t}\) is closed and bounded
below on the common domain \(\mathcal D_A\).  Its associated self-adjoint
operator \(T_{A,t}\) on \(\mathcal H_A\) has compact resolvent.  On every
compact heat-time interval the forms are norm-continuous, and on compact
subintervals of \((0,\infty)\) they form a real-analytic family of forms.
\end{proposition}

\begin{proof}
The digamma asymptotic used in Lemma~\ref{lem:v20-gamma-lower}, together with
the matching upper estimate, gives
\[
 a_{\chi,t}(\xi)=\log(2+|\xi|)+O_{\chi,T}(1)
 \tag{V21.4}
\]
uniformly in \(\xi\) and \(0\le t\le T\).  More directly, convolution with
a Gaussian of bounded variance
changes \(\log(2+|\xi|)\) by a uniformly bounded amount.  On
\(\mathcal H_A\), one has
\[
 C_g(u)=0\qquad(|u|>2A),
 \tag{V21.5}
\]
so even at \(t=0\) only finitely many prime powers occur.  Each translation
term is bounded on \(L^2\).  Thus \(q_{A,t}\) is a bounded perturbation of
the positive logarithmic Fourier form, proving closedness, a common domain,
and semiboundedness.

The unit ball of \(\mathcal D_A\) has Fourier tail
\[
 \int_{|\xi|>R}|\widehat g(\xi)|^2d\xi
 \le \frac{C}{\log(2+R)}.
 \tag{V21.6}
\]
The low-frequency truncation of functions supported in \([-A,A]\) is a
Hilbert--Schmidt time--frequency map, hence compact, while
\textup{(V21.6)} makes those truncations uniformly close in \(L^2\).
Therefore \(\mathcal D_A\hookrightarrow\mathcal H_A\) is compact and the
represented operator has compact resolvent.

Finally, \(a_\chi''\) is bounded on the fixed vertical line, so the heat
equation gives
\[
 \|a_{\chi,t}-a_{\chi,s}\|_\infty
 \le \|a_\chi''\|_\infty|t-s|.
 \tag{V21.7}
\]
The finite prime sum on the window is norm-continuous as well.  Gaussian
convolution and the finite exponential coefficients are real analytic for
positive heat time.  This proves the remaining claims.
\end{proof}

In particular the form domain itself is independent of heat time.

\subsection{Finite-window thresholds exhaust the global threshold}

Define
\[
 \mathfrak P_{\chi,A}
 :=\{t\ge0:q_{A,t}[g]\ge0\text{ for every }g\in\mathcal D_A\},
 \qquad
 \vartheta_\chi(A):=\inf\mathfrak P_{\chi,A}.
 \tag{V21.8}
\]

\begin{theorem}[Cofinal finite-window exhaustion]
\label{thm:v21-window-exhaustion}
For every \(A>0\), the set \(\mathfrak P_{\chi,A}\) is a nonempty closed
upper interval
\[
 \mathfrak P_{\chi,A}=[\vartheta_\chi(A),\infty).
 \tag{V21.9}
\]
The function \(A\mapsto\vartheta_\chi(A)\) is nondecreasing and
\[
 \boxed{
 \vartheta_\chi
 =\sup_{A>0}\vartheta_\chi(A)
 =\lim_{A\to\infty}\vartheta_\chi(A).}
 \tag{V21.10}
\]
Consequently,
\[
 \boxed{
 \text{GRH for the paired }\chi,\overline\chi\text{ channel fails}
 \quad\Longleftrightarrow\quad
 \vartheta_\chi(A)>0\text{ for some finite }A.}
 \tag{V21.11}
\]
\end{theorem}

\begin{proof}
The eventual coercivity bound \textup{(V20.21)} makes every window-positive
set nonempty.  The dephasing identity \textup{(V20.12)} remains inside
\(\mathcal H_A\), because modulation does not change support; hence each set
is upper.  Proposition~\ref{prop:v21-compact} gives closure.  Domain
inclusion \(\mathcal H_A\subset\mathcal H_B\) for \(A<B\) makes the
thresholds nondecreasing.

Positivity on all compactly supported smooth tests is equivalent to
positivity on every \(\mathcal D_A\), by core density.  Therefore the global
positive set is the intersection of the upper intervals in
\textup{(V21.9)}.  The left endpoint of that intersection is their supremum,
which proves \textup{(V21.10)}.

If GRH fails, the Weil criterion supplies a compactly supported smooth
\(g\) with \(q_{A,0}[g]<0\) for some finite support radius \(A\).  Thus zero
is outside \(\mathfrak P_{\chi,A}\), so
\(\vartheta_\chi(A)>0\).  Conversely, a positive finite-window threshold
means \(q_{A,0}\) is not nonnegative.  Core density supplies a smooth
compactly supported negative test, and the Weil criterion disproves GRH.
\end{proof}

This theorem removes a possible compactness ambiguity from V20.  A
hypothetical off-line zero cannot force the heat threshold solely through
test functions whose supports escape to infinity: one finite support window
already contains a complete spectral crossing.

\subsection{The exact null-space heat velocity}

Let \(X\) denote multiplication by the physical variable,
\[
 (Xg)(x)=xg(x).
 \tag{V21.12}
\]
On a fixed window, multiplication by every power of \(X\) preserves
\(\mathcal D_A\).  Polarizing Theorem~\ref{thm:v20-dephasing} and
differentiating the heat convolution at zero gives
\[
 \dot{\mathcal B}_{\chi,t}^{A}(g,h)
 =\left.\frac{d^2}{dv^2}
 \mathcal B_{\chi,t}^{A}(M_vg,M_vh)\right|_{v=0}.
 \tag{V21.13}
\]

\begin{theorem}[Critical-kernel carr\'e du champ]
\label{thm:v21-carre}
Assume \(q_{A,t}\ge0\), and let
\[
 \mathcal K_{A,t}:=\ker T_{A,t}.
 \tag{V21.14}
\]
Then \(\mathcal K_{A,t}\) is finite dimensional and, for every
\(g,h\in\mathcal K_{A,t}\),
\[
 \boxed{
 \dot{\mathcal B}_{\chi,t}^{A}(g,h)
 =2\mathcal B_{\chi,t}^{A}(Xg,Xh).}
 \tag{V21.15}
\]
Equivalently, if \(P_{A,t}\) is the orthogonal projection onto the kernel,
then in quadratic-form notation
\[
 \boxed{
 P_{A,t}\dot T_{A,t}P_{A,t}
 =2\bigl(T_{A,t}^{1/2}XP_{A,t}\bigr)^*
      \bigl(T_{A,t}^{1/2}XP_{A,t}\bigr)\succeq0.}
 \tag{V21.16}
\]
If the kernel is nonzero, the positive matrix in \textup{(V21.16)} is not
the zero matrix.
\end{theorem}

\begin{proof}
Compact resolvent makes the kernel finite dimensional.  Positivity makes it
the radical of the form: Cauchy--Schwarz gives
\[
 \mathcal B_{\chi,t}^{A}(g,w)=0
 \quad(g\in\mathcal K_{A,t},\ w\in\mathcal D_A).
 \tag{V21.17}
\]
Since
\[
 \left.\frac d{dv}M_vg\right|_{v=0}=iXg,
 \qquad
 \left.\frac {d^2}{dv^2}M_vg\right|_{v=0}=-X^2g,
 \tag{V21.18}
\]
the second derivative in \textup{(V21.13)} is
\[
 2\mathcal B_{\chi,t}^{A}(Xg,Xh)
 -\mathcal B_{\chi,t}^{A}(X^2g,h)
 -\mathcal B_{\chi,t}^{A}(g,X^2h).
 \tag{V21.19}
\]
The last two terms vanish by \textup{(V21.17)}, proving
\textup{(V21.15)} and its Gram factorization \textup{(V21.16)}.

If the matrix were zero, then
\(q_{A,t}[Xg]=0\) for every \(g\in\mathcal K_{A,t}\), so
\(X\mathcal K_{A,t}\subset\mathcal K_{A,t}\).  A nonzero finite-dimensional
complex space invariant under \(X\) contains an eigenvector of the
restriction of \(X\).  Such a vector would obey
\((x-\lambda)g(x)=0\) almost everywhere and hence be supported at one point,
which is impossible for a nonzero \(L^2\) function.  Therefore the Gram
matrix is nonzero.
\end{proof}

At a positive finite-window threshold, norm continuity and compact resolvent
give
\[
 q_{A,\vartheta_\chi(A)}\ge0,
 \qquad
 \ker T_{A,\vartheta_\chi(A)}\ne\{0\}.
 \tag{V21.20}
\]
Indeed, a positive spectral floor at the endpoint would persist to a smaller
heat time and contradict minimality.  Hence every failure of GRH supplies
the finite critical packet \textup{(V21.14)}--\textup{(V21.16)}.

\begin{corollary}[Simple critical crossings are transverse]
\label{cor:v21-transverse}
Suppose \(\vartheta_\chi(A)>0\) and the kernel in
\textup{(V21.20)} is one dimensional, spanned by a normalized \(g_A\).
Then the analytic lowest-eigenvalue branch satisfies
\[
 \boxed{
 \lambda_A'(\vartheta_\chi(A))
 =2q_{A,\vartheta_\chi(A)}[Xg_A]>0.}
 \tag{V21.21}
\]
In particular,
\[
 \lambda_A(\vartheta_\chi(A)-s)
 =-2s\,q_{A,\vartheta_\chi(A)}[Xg_A]+O(s^2)<0
 \tag{V21.22}
\]
for all sufficiently small \(s>0\).
\end{corollary}

\begin{proof}
Analytic perturbation of the simple eigenvalue and the form
Hellmann--Feynman identity give the equality in \textup{(V21.21)}.
If its right side vanished, positivity would put \(Xg_A\) in the
one-dimensional kernel, making \(g_A\) an \(L^2\) eigenvector of
multiplication by \(x\), which is impossible.  Taylor expansion gives
\textup{(V21.22)}.
\end{proof}

\subsection{An exact near-null identity}

The same calculation without imposing the null condition gives a useful
identity on the whole positive window:
\[
 \boxed{
 \dot q_{A,t}[g]
 =2q_{A,t}[Xg]
 -2\Re\mathcal B_{\chi,t}^{A}(g,X^2g).}
 \tag{V21.23}
\]
Therefore, whenever \(q_{A,t}\ge0\),
\[
 \boxed{
 \left|\dot q_{A,t}[g]-2q_{A,t}[Xg]\right|
 \le2\sqrt{q_{A,t}[g]q_{A,t}[X^2g]}.}
 \tag{V21.24}
\]
This is the precise near-null replacement for estimating the two lines of
\textup{(V20.26)} separately.  It is sign-sensitive and contains the
completed gamma and prime pieces automatically.

For example, on any positive heat interval, the pair of relative graph
bounds
\[
 q_{A,t}[Xg]\le\alpha_A(t)q_{A,t}[g],
 \qquad
 q_{A,t}[X^2g]\le\beta_A(t)q_{A,t}[g]
 \tag{V21.25}
\]
would imply
\[
 \dot q_{A,t}[g]
 \le2\bigl(\alpha_A(t)+\sqrt{\beta_A(t)}\bigr)q_{A,t}[g].
 \tag{V21.26}
\]
But Theorem~\ref{thm:v21-carre} shows exactly where such a bound must fail at
a positive critical packet: it would force the finite-dimensional kernel to
be invariant under \(X\).  Thus the obstruction is no longer an unspecified
large constant.  It is the nonzero Gram matrix
\[
 \mathfrak C_{A,t}
 :=\left[
 2\mathcal B_{\chi,t}^{A}(Xe_j,Xe_k)
 \right]_{e_j,e_k\in\mathcal K_{A,t}}.
 \tag{V21.27}
\]

\begin{firewall}[Positive null curvature is not a GRH proof]
The carr\'e-du-champ matrix \textup{(V21.27)} is automatically positive at a
positive-form boundary; it does not show that the boundary cannot occur.
Indeed it describes how Gaussian dephasing can remove a negative direction.
To prove GRH by this route one must use the explicit arithmetic structure to
exclude every finite critical packet, for example by proving an independent
upper tangency relation incompatible with the nonzero Gram matrix.  No such
arithmetic tangency theorem is proved here.
\end{firewall}

\begin{assessment}[V21 acquisition]
V21 replaces the global near-null ambiguity by compact finite-window
spectral data.  The global heat threshold is the increasing supremum of
finite-window thresholds, and failure of GRH is equivalent to one positive
finite-window threshold.  At every such critical endpoint the complete Weil
form has a finite-dimensional kernel, and its heat-velocity compression is
the nonzero positive Gram \textup{(V21.16)}.  A simple critical eigenvalue
crosses transversely with the exact slope \textup{(V21.21)}.  The near-null
identity \textup{(V21.23)} retains all signed arithmetic cancellation and
identifies the next load-bearing object as an arithmetic upper-tangency
estimate on the finite critical kernel, not a global absolute prime bound.
GRH remains open.
\end{assessment}

\section*{Version V21 append-only record}
\addcontentsline{toc}{section}{Version V21 append-only record}
The complete V20 source was retained before this continuation.  It had
\(315095\) bytes, \(9454\) lines, and SHA--256
\[
 \texttt{806546E1770CC7882418A9752D58808A2619D23B273F1A97C4BBADCA66D3512C}.
\]
V21 constructed the compact finite-window heat operators, exhausted the
global threshold by their critical times, and proved the exact null-space
carr\'e-du-champ and simple-crossing formulas.

% =============================================================================
% END APPENDED V21 MATERIAL
% =============================================================================
% =============================================================================
% BEGIN APPENDED V22 MATERIAL
% =============================================================================

\section{V22: the scalar Weil heat symbol and complex-zero cancellation}
\label{sec:grh-v22}

V21 converted a hypothetical failure of GRH into compact spectral crossings
on fixed support windows.  There is a further upstream simplification at
positive heat time.  Gaussian damping makes the full prime series absolutely
convergent, and both the archimedean term and every prime autocorrelation are
translation invariant.  The complete heat-deformed Weil operator is
therefore a scalar Fourier multiplier.  Its symbol solves the ordinary
one-dimensional heat equation.

The standard paired Weil explicit formula identifies that symbol with a
locally uniformly convergent sum of Gaussians whose centers are the reflected
complex coordinates of the nontrivial zeros.  Critical positivity is thus a
pointwise cancellation problem.  Critical-line zeros contribute positive
real Gaussians; an off-line reflected pair contributes an oscillatory cosine
Gaussian.  Any positive critical heat time must contain an exact cancellation
between these two types.

\subsection{The complete scalar heat symbol}

For \(t>0\), define
\[
 \boxed{
 \begin{aligned}
 m_{\chi,t}(\xi)
 :={}&a_{\chi,t}(\xi)\\
 &-2\Re\sum_{n\ge2}
 \frac{\Lambda(n)\chi(n)}{\sqrt n}
 e^{-t(\log n)^2}e^{i\xi\log n}.
 \end{aligned}}
 \tag{V22.1}
\]
The series and all of its \(\xi\)-derivatives converge uniformly on compact
subsets of \((0,\infty)\times\mathbb R\); in fact the series extends entire
in \(\xi\) for each fixed positive \(t\).

\begin{theorem}[Scalarization of the heat-deformed Weil form]
\label{thm:v22-scalarization}
For every \(t>0\) and \(g\in C_c^\infty(\mathbb R)\),
\[
 \boxed{
 \mathcal W_{\chi,t}[g]
 =\frac1{2\pi}\int_{\mathbb R}
 m_{\chi,t}(\xi)|\widehat g(\xi)|^2d\xi.}
 \tag{V22.2}
\]
Moreover,
\[
 \boxed{
 \partial_t m_{\chi,t}(\xi)
 =\partial_\xi^2m_{\chi,t}(\xi),}
 \tag{V22.3}
\]
and
\[
 \boxed{
 \mathcal W_{\chi,t}\ge0\text{ on }C_c^\infty(\mathbb R)
 \quad\Longleftrightarrow\quad
 m_{\chi,t}(\xi)\ge0\text{ for every }\xi\in\mathbb R.}
 \tag{V22.4}
\]
\end{theorem}

\begin{proof}
Plancherel gives
\[
 C_g(u)=\frac1{2\pi}\int_{\mathbb R}
 e^{i\xi u}|\widehat g(\xi)|^2d\xi
 \tag{V22.5}
\]
with the convention already used in V19.  The bound
\[
 \sum_{n\ge2}\frac{\Lambda(n)}{\sqrt n}
 (\log n)^k e^{-t(\log n)^2}<\infty
 \tag{V22.6}
\]
for every \(k\ge0\) and \(t>0\) justifies inserting
\textup{(V22.5)} into \textup{(V20.9)} and interchanging sum and integral.
This proves \textup{(V22.2)}.

The archimedean symbol satisfies the heat equation by
\(a_{\chi,t}=H_t*a_\chi\).  Twice differentiating
\(e^{i\xi u}\) gives \(-u^2e^{i\xi u}\), exactly matching the time
derivative of \(e^{-tu^2}\) in the prime term, proving
\textup{(V22.3)}.

Pointwise nonnegativity immediately implies form nonnegativity.  Conversely,
if \(m_{\chi,t}(\xi_0)<0\), choose a fixed nonzero
\(\phi\in C_c^\infty(\mathbb R)\) and put
\[
 g_R(x)=R^{-1/2}\phi(x/R)e^{i\xi_0x}.
 \tag{V22.7}
\]
After normalization, \((2\pi)^{-1}|\widehat g_R|^2d\xi\) is an approximate
identity at \(\xi_0\).  Equation \textup{(V22.2)} is then negative for all
sufficiently large \(R\), proving the converse.
\end{proof}

Thus the completely positive operator flow of V20 is the ordinary scalar
heat flow after Fourier diagonalization:
\[
 m_{\chi,t+s}=H_s*m_{\chi,t}.
 \tag{V22.8}
\]
This identity uses the complete symbol.  Neither the gamma component nor the
prime component separately has the required positivity meaning.

\subsection{The paired zero-side Gaussian expansion}

Let \(\mathfrak Z_\chi\) be the symmetrized multiset of nontrivial zeros in
the paired \(\chi,\overline\chi\) channel, with the conventional factor
one-half when the two character multisets are combined.  For
\(\rho=\beta+i\gamma\), write
\[
 z_\rho:=\frac{\rho-1/2}{i}
 =\gamma-i(\beta-1/2).
 \tag{V22.9}
\]
The multiset is invariant under complex conjugation.  In the principal
channel, the following statement is read after the standard pole subtraction
used in the Weil criterion.

For a Weil test \(g\), put
\[
 h_{t,g}(u):=e^{-tu^2}C_g(u).
 \tag{V22.10}
\]
If \(F_g(\xi)=|\widehat g(\xi)|^2\) on the real axis, then the entire Fourier
transform of \(h_{t,g}\) is
\[
 \widehat h_{t,g}(z)
 =\int_{\mathbb R}H_t(z-\xi)F_g(\xi)d\xi.
 \tag{V22.11}
\]

\begin{theorem}[Complex-zero heat-kernel representation]
\label{thm:v22-zero-heat}
The standard paired Weil explicit formula implies, for every \(t>0\),
\[
 \boxed{
 m_{\chi,t}(\xi)
 =2\pi\sum_{\rho\in\mathfrak Z_\chi}
 m_\rho H_t(\xi-z_\rho),}
 \tag{V22.12}
\]
where \(m_\rho\) is the zero multiplicity with the preceding pairing
normalization.  The series converges absolutely and locally uniformly for
real \(\xi\), together with all derivatives.  It is real valued because
conjugate centers occur with equal multiplicity.
\end{theorem}

\begin{proof}
The completed paired explicit formula, in the normalization of
\textup{(V19.9)}--\textup{(V19.14)}, is
\[
 \mathcal W_\chi[h]
 =\sum_{\rho\in\mathfrak Z_\chi}m_\rho\widehat h(z_\rho)
 \tag{V22.13}
\]
for the positive-type test \(h=C_g\), and by the same formula for
\(h=h_{t,g}\).  Its prime side is precisely \textup{(V20.9)}.  Substitute
\textup{(V22.11)} into \textup{(V22.13)} and interchange the zero sum with
the real integral.  The zero-counting estimate
\(N_\chi(Y)=O_\chi(Y\log(2+Y))\), the critical-strip bound
\(|\Re\rho-1/2|\le1/2\), and Gaussian decay in \(|\xi-\Im\rho|\) justify the
interchange locally uniformly.  Comparing with \textup{(V22.2)} gives
\textup{(V22.12)}.
\end{proof}

When \(\rho=1/2+i\gamma\), its summand in \textup{(V22.12)} is the positive
real Gaussian
\[
 H_t(\xi-\gamma)>0.
 \tag{V22.14}
\]
For a reflected off-line pair with horizontal displacement
\(0<\delta=|\beta-1/2|\le1/2\) and common ordinate \(\gamma\), the paired
contribution is
\[
 \boxed{
 \frac{2}{\sqrt{4\pi t}}
 \exp\!\left(-\frac{(\xi-\gamma)^2-\delta^2}{4t}\right)
 \cos\!\left(\frac{\delta(\xi-\gamma)}{2t}\right).}
 \tag{V22.15}
\]
Thus horizontal zero displacement is exactly an oscillatory phase and an
amplification factor \(e^{\delta^2/(4t)}\) in the scalar heat symbol.

\subsection{Critical contact equations and the off-line phase gate}

For fixed \(t>0\), the prime series in \textup{(V22.1)} is bounded uniformly
in \(\xi\), while \(a_{\chi,t}(\xi)=\log(2+|\xi|)+O_{\chi,t}(1)\).  Hence
\[
 m_{\chi,t}(\xi)\longrightarrow+\infty
 \qquad(|\xi|\to\infty).
 \tag{V22.16}
\]
Set
\[
 \mu_\chi(t):=\min_{\xi\in\mathbb R}m_{\chi,t}(\xi).
 \tag{V22.17}
\]
Theorem~\ref{thm:v22-scalarization} gives
\[
 \mathfrak P_\chi=\{t>0:\mu_\chi(t)\ge0\}.
 \tag{V22.18}
\]

\begin{theorem}[Scalar structure of a positive critical time]
\label{thm:v22-critical-contact}
Assume \(\vartheta_\chi>0\).  Then the real critical set
\[
 Z_{\chi,*}
 :=\{\xi:m_{\chi,\vartheta_\chi}(\xi)=0\}
 \tag{V22.19}
\]
is nonempty and finite.  Every \(\xi_*\in Z_{\chi,*}\) is a zero of even
order \(2r\), for some \(r\ge1\), and
\[
 \partial_\xi^j m_{\chi,\vartheta_\chi}(\xi_*)=0
 \quad(0\le j<2r),
 \qquad
 \partial_\xi^{2r}m_{\chi,\vartheta_\chi}(\xi_*)>0.
 \tag{V22.20}
\]
The heat equation gives the exact opening law
\[
 \boxed{
 \partial_t^r m_{\chi,t}(\xi_*)\big|_{t=\vartheta_\chi}
 =\partial_\xi^{2r}m_{\chi,\vartheta_\chi}(\xi_*)>0.}
 \tag{V22.21}
\]
In the nondegenerate case \(r=1\),
\[
 \partial_tm_{\chi,t}(\xi_*)\big|_{t=\vartheta_\chi}
 =m_{\chi,\vartheta_\chi}''(\xi_*)>0.
 \tag{V22.22}
\]
\end{theorem}

\begin{proof}
At the left endpoint of the positive upper interval,
\(m_{\chi,\vartheta_\chi}\ge0\).  If its minimum were positive, local
uniform continuity in heat time and the coercive growth at infinity would
preserve positivity for a smaller time, contradicting minimality.  Thus the
critical set is nonempty.  The symbol is real analytic and is not identically
zero by \textup{(V22.16)}.  Its zeros in the compact zero sublevel region are
therefore finite.  Nonnegativity forces each zero to have even order and
positive first nonzero derivative.  Iterating
\(\partial_t=\partial_\xi^2\) proves \textup{(V22.21)}.
\end{proof}

At every critical point, Theorem~\ref{thm:v22-zero-heat} gives the two exact
contact equations
\[
 \boxed{
 \sum_{\rho}m_\rho H_{\vartheta_\chi}(\xi_*-z_\rho)=0,
 \qquad
 \sum_{\rho}m_\rho(\xi_*-z_\rho)
 H_{\vartheta_\chi}(\xi_*-z_\rho)=0.}
 \tag{V22.23}
\]
The second is \(m'_{\chi,\vartheta_\chi}(\xi_*)=0\), with its harmless
factor \(-1/(2\vartheta_\chi)\) removed.

\begin{corollary}[Off-line negative-phase gate]
\label{cor:v22-phase-gate}
At every positive critical contact \((\vartheta_\chi,\xi_*)\), at least one
off-line reflected zero pair \((\delta,\gamma)\) satisfies
\[
 \cos\!\left(
 \frac{\delta(\xi_*-\gamma)}{2\vartheta_\chi}
 \right)<0.
 \tag{V22.24}
\]
Consequently that pair obeys
\[
 \boxed{
 |\xi_*-\gamma|>
 \frac{\pi\vartheta_\chi}{\delta}
 \ge2\pi\vartheta_\chi.}
 \tag{V22.25}
\]
The total negative off-line lobes in \textup{(V22.15)} exactly balance all
critical-line Gaussians and all positive off-line lobes in the first equation
of \textup{(V22.23)}.
\end{corollary}

\begin{proof}
Every critical-line contribution \textup{(V22.14)} is strictly positive.  If
all off-line cosines were nonnegative, the first sum in
\textup{(V22.23)} would be strictly positive.  Hence one cosine is negative.
Any real number with negative cosine has absolute value greater than
\(\pi/2\) modulo its first passage from zero, which gives the first inequality
in \textup{(V22.25)}.  The critical-strip bound \(\delta\le1/2\) gives the
second.
\end{proof}

If GRH holds, every \(z_\rho\) is real and \textup{(V22.12)} is a sum of
strictly positive Gaussians.  Hence
\[
 m_{\chi,t}(\xi)>0
 \quad(t>0,\ \xi\in\mathbb R),
 \tag{V22.26}
\]
which recovers the zero-threshold implication with strict positive-time
pointwise positivity.

\subsection{Strictly cofinal finite-window crossings}

The scalar symbol sharpens the relation between V21 and the global threshold.

\begin{theorem}[No finite window attains a positive global threshold]
\label{thm:v22-strict-cofinal}
If \(\vartheta_\chi>0\), then for every finite \(A>0\),
\[
 \boxed{
 \vartheta_\chi(A)<\vartheta_\chi,}
 \tag{V22.27}
\]
although V21 gives
\(\vartheta_\chi(A)\uparrow\vartheta_\chi\).  At the global critical time,
the form is strictly positive on every nonzero \(L^2\) vector in its form
domain, but it has no uniform global \(L^2\) floor.
\end{theorem}

\begin{proof}
The nonnegative real-analytic function
\(m_{\chi,\vartheta_\chi}\) has only finitely many real zeros, hence its zero
set has Lebesgue measure zero.  Equation \textup{(V22.2)} then gives
\[
 q_{A,\vartheta_\chi}[g]>0
 \tag{V22.28}
\]
for every nonzero \(g\in\mathcal D_A\).  The compact-resolvent realization
of Proposition~\ref{prop:v21-compact} upgrades this strict inequality to a
positive lowest eigenvalue on each fixed window.  Norm continuity in heat
time preserves that window floor for a short interval to the left of
\(\vartheta_\chi\), proving \textup{(V22.27)}.  Globally, the approximate
frequency packets \textup{(V22.7)} centered at any
\(\xi_*\in Z_{\chi,*}\) make the Rayleigh quotient tend to zero, so no
uniform global floor exists.
\end{proof}

More precisely, if \(\|g_j\|_2=1\) and
\(\mathcal W_{\chi,\vartheta_\chi}[g_j]\to0\), then for every open
neighbourhood \(U\) of \(Z_{\chi,*}\),
\[
 \boxed{
 \frac1{2\pi}\int_{\mathbb R\setminus U}
 |\widehat g_j(\xi)|^2d\xi\longrightarrow0.}
 \tag{V22.29}
\]
Indeed, the critical symbol has a strictly positive minimum on the complement
of \(U\), including its coercive tails.  Thus critical near-null sequences
concentrate in frequency at the finite contact set and cannot remain in any
fixed support window.  This is the exact mechanism by which the compact
critical packets of V21 approach, but never attain, a positive global
threshold.

\begin{firewall}[The complex Gaussian sum has not been proved positive]
Formula \textup{(V22.12)} is the heat-regularized zero side of the standard
Weil explicit formula.  Off-line centers produce the amplified oscillatory
terms \textup{(V22.15)}, so the representation is not positive without GRH.
The phase gate \textup{(V22.24)} is a necessary condition for a positive
critical time, not an exclusion of one.  Proving that the negative lobes
cannot satisfy both contact equations \textup{(V22.23)} requires new
arithmetic zero-density, separation, or signed-cancellation information.
GRH remains open.
\end{firewall}

\begin{assessment}[V22 acquisition]
V22 diagonalizes the complete positive-time Weil form into the scalar heat
symbol \textup{(V22.1)} and proves that form positivity is exactly pointwise
symbol positivity.  The paired explicit formula becomes the complex-center
Gaussian sum \textup{(V22.12)}.  A positive critical time has finitely many
even-order real contacts, obeys the exact zero-side value and derivative
equations \textup{(V22.23)}, and requires a negative off-line phase at the
distance scale \textup{(V22.25)}.  Every fixed support threshold lies
strictly below the global threshold, while global near-null sequences
concentrate at the finite critical frequency set.  The next proof target is
now the signed zero-side cancellation budget in \textup{(V22.23)}, rather
than an operator-wide Schur inequality.
\end{assessment}

\section*{Version V22 append-only record}
\addcontentsline{toc}{section}{Version V22 append-only record}
The complete V21 source was retained before this continuation.  It had
\(328273\) bytes, \(9828\) lines, and SHA--256
\[
 \texttt{3A3C2FBC1341A1DE8F88F0754B22D09A1A5B18A23D79F6BE57C344489400BD71}.
\]
V22 constructed the scalar Weil heat symbol, its complex-zero Gaussian
expansion, the critical contact and phase-gate equations, and the strict
cofinality of finite-window crossings.

% =============================================================================
% END APPENDED V22 MATERIAL
% =============================================================================
% =============================================================================
% BEGIN APPENDED V23 MATERIAL
% =============================================================================

\section{V23: exposed off-line zeros and the local negative-part criterion}
\label{sec:grh-v23}

V22 turned the positive-time Weil problem into pointwise positivity of a
scalar heat symbol and displayed each off-line reflected zero pair as an
amplified oscillatory Gaussian.  The proposed next step was to combine the
two critical contact equations with zero density and separation.  Before
attempting that estimate, one must determine what unsigned density
information can possibly prove.

This continuation gives a sharp answer.  If any off-line zero exists, a
finite local family of zero centers contains a uniquely exposed pair.  By
tuning the heat time to a negative cosine lobe, that one pair exponentially
dominates the entire remaining zero set and makes the scalar Weil symbol
negative.  This gives a constructive small-heat witness and an equivalent
criterion in terms of the exponential rate of the local negative part of the
prime-side symbol.  The absolute prime mass has exactly the full critical
strip exponent and therefore cannot establish that criterion.

\subsection{The Lorentzian score envelope of the zero set}

Group identical reflected pairs in the symmetrized zero multiset of V22.  A
pair is described by its ordinate \(\gamma_p\), horizontal displacement
\[
 0\le\delta_p:=|\Re\rho_p-1/2|\le1/2,
 \tag{V23.1}
\]
and positive multiplicity \(m_p\).  For a real observation point \(x\), put
\[
 S_p(x):=\delta_p^2-(x-\gamma_p)^2,
 \qquad
 E_\chi(x):=\max_p S_p(x).
 \tag{V23.2}
\]
The maximum is attained.  Indeed, the zero set is locally finite and
\(S_p(x)\to-\infty\) as \(|\gamma_p|\to\infty\), uniformly in
\(0\le\delta_p\le1/2\).

\begin{lemma}[Existence of a uniquely exposed off-line pair]
\label{lem:v23-exposed}
If the paired zero set contains an off-line zero, then there are
\(x_0\in\mathbb R\), a unique off-line pair \(p_*\), and \(\eta>0\) such
that
\[
 E_\chi(x_0)=S_{p_*}(x_0)>0
 \tag{V23.3}
\]
and
\[
 S_p(x_0)\le S_{p_*}(x_0)-\eta
 \qquad(p\ne p_*).
 \tag{V23.4}
\]
The same strict exposure, with gap at least \(\eta/2\), holds throughout a
sufficiently small neighbourhood of \(x_0\).
\end{lemma}

\begin{proof}
Take an off-line pair \(p_0\) and begin at \(x=\gamma_{p_0}\).  Then
\(S_{p_0}(x)=\delta_{p_0}^2>0\), whereas every critical-line pair has score
at most zero.  Hence every maximizer near this point is off-line and the
maximum remains positive.

Only finitely many pairs can compete with a fixed positive score on a compact
neighbourhood: pairs with distant ordinates have uniformly negative score.
For two distinct grouped pairs,
\[
 S_p(x)-S_q(x)
 =\delta_p^2-delta_q^2-(x-\gamma_p)^2+(x-\gamma_q)^2
 \tag{V23.5}
\]
is a nonconstant affine function of \(x\), unless the ordinates agree; in
the latter case one displacement is strictly larger and no tie occurs.
Thus the finite set of possible tie points has empty interior.  Choose
\(x_0\) nearby but outside it.  The maximizer is unique and off-line.
Finiteness of the nearby competitors and uniform decay of the distant scores
give the positive gap \(\eta\).  Continuity preserves half that gap on a
smaller neighbourhood.
\end{proof}

The exponent \(S_p(x)/(4t)\) is precisely the logarithmic amplitude of the
pair contribution \textup{(V22.15)}, apart from the common
\((4\pi t)^{-1/2}\) factor.  Thus \(E_\chi\) is the small-heat Laplace
envelope of the symmetrized zero set.  In particular,
\[
 \sup_{x\in\mathbb R}E_\chi(x)
 =\sup_\rho|\Re\rho-1/2|^2.
 \tag{V23.6}
\]

For normalization clarity, \textup{(V22.15)} displays the contribution of a
pair to the inner zero sum in \textup{(V22.12)}.  Its contribution to
\(m_{\chi,t}\) carries the additional common positive factor \(2\pi\).
This factor affects none of the sign or exponential-rate conclusions.

\subsection{Gaussian dominance after phase tuning}

The following elementary estimate is the only place where zero counting is
needed.

\begin{lemma}[Exponential suppression below an exposed score]
\label{lem:v23-remainder}
Under Lemma~\ref{lem:v23-exposed}, uniformly for \(x\) in a sufficiently
small fixed neighbourhood of \(x_0\),
\[
 \sum_{p\ne p_*}m_p
 \exp\!\left(\frac{S_p(x)}{4t}\right)
 =o\!\left(
 \exp\!\left(\frac{S_{p_*}(x)}{4t}\right)
 \right)
 \qquad(t\downarrow0).
 \tag{V23.7}
\]
\end{lemma}

\begin{proof}
For ordinates in a fixed bounded interval, only finitely many zeros occur,
and the exposure gap supplies a factor \(e^{-\eta/(8t)}\) after shrinking the
neighbourhood.  For the remaining ordinates, the strip bound
\(\delta_p^2\le1/4\) and the quadratic term
\(-(x-\gamma_p)^2\) give Gaussian decay.  Summation by unit ordinate shells
and the standard estimate
\[
 N_\chi(Y+1)-N_\chi(Y)=O_\chi(\log(2+Y))
 \tag{V23.8}
\]
bound the distant sum by a polynomial in \(t^{-1}\) times a strictly smaller
exponential.  Every polynomial is absorbed by the exposure gap as
\(t\downarrow0\), proving \textup{(V23.7)}.
\end{proof}

\begin{theorem}[Every off-line zero produces a local negative heat sequence]
\label{thm:v23-negative-sequence}
If the paired channel has an off-line zero, then there are a compact interval
\(I\), sequences \(t_k\downarrow0\) and \(\xi_k\in I\), and a constant
\(E_*>0\) such that
\[
 \boxed{
 m_{\chi,t_k}(\xi_k)<0}
 \tag{V23.9}
\]
for all sufficiently large \(k\), and quantitatively
\[
 \boxed{
 \liminf_{k\to\infty}
 4t_k\log\bigl(-m_{\chi,t_k}(\xi_k)\bigr)
 \ge E_* .}
 \tag{V23.10}
\]
One may take \(E_*=E_\chi(x_0)\) for an exposed point from
Lemma~\ref{lem:v23-exposed}.
\end{theorem}

\begin{proof}
Write \(p_*=(\gamma_*,\delta_*)\) and
\(d_*=x_0-\gamma_*\).  If \(d_*\ne0\), choose
\[
 t_k:=\frac{\delta_*|d_*|}{2(2k+1)\pi},
 \qquad \xi_k:=x_0.
 \tag{V23.11}
\]
Then
\[
 \cos\!\left(
 \frac{\delta_*(\xi_k-\gamma_*)}{2t_k}
 \right)=-1.
 \tag{V23.12}
\]
If \(d_*=0\), take any \(t_k\downarrow0\) and put
\[
 \xi_k:=x_0+\frac{2\pi t_k}{\delta_*};
 \tag{V23.13}
\]
again \textup{(V23.12)} holds and \(\xi_k\to x_0\).

The contribution of \(p_*\) in \textup{(V22.15)} is therefore negative with
magnitude
\[
 \frac{4\pi m_{p_*}}{\sqrt{4\pi t_k}}
 \exp\!\left(\frac{S_{p_*}(\xi_k)}{4t_k}\right).
 \tag{V23.14}
\]
The absolute value of the sum of every other pair is little-oh of this
quantity by Lemma~\ref{lem:v23-remainder}; critical-line zeros are included
among those other pairs.  Hence the full zero expansion
\textup{(V22.12)} is negative.  Since
\(S_{p_*}(\xi_k)\to E_\chi(x_0)>0\), taking logarithms proves
\textup{(V23.10)}.
\end{proof}

This is a constructive version of the hard direction of the Weil criterion.
It does not interpolate arbitrary values at zeros.  Instead, a real
frequency and a heat scale are selected before the scalar form test, and the
off-line pair is isolated by its Gaussian exponential rate.

\subsection{An equivalent local negative-part criterion}

For a compact interval \(I\subset\mathbb R\), define
\[
 \mathcal N_{\chi,I}(t)
 :=\sup_{\xi\in I}[-m_{\chi,t}(\xi)]_+.
 \tag{V23.15}
\]

\begin{theorem}[Local subexponential negative-part criterion]
\label{thm:v23-local-criterion}
The following are equivalent:
\begin{enumerate}[label=\textup{(\roman*)},leftmargin=3em]
\item GRH holds for the paired \(\chi,\overline\chi\) channel;
\item for every compact interval \(I\),
\[
 \boxed{
 \limsup_{t\downarrow0}
 4t\log\bigl(1+\mathcal N_{\chi,I}(t)\bigr)=0.}
 \tag{V23.16}
\]
\end{enumerate}
If GRH fails, some compact \(I\) instead satisfies
\[
 \limsup_{t\downarrow0}
 4t\log\bigl(1+\mathcal N_{\chi,I}(t)\bigr)>0.
 \tag{V23.17}
\]
\end{theorem}

\begin{proof}
Under GRH, \textup{(V22.26)} gives
\(\mathcal N_{\chi,I}(t)=0\), proving \textup{(V23.16)}.  If an off-line
zero exists, take the compact interval containing the sequence in
Theorem~\ref{thm:v23-negative-sequence}.  Equations
\textup{(V23.9)}--\textup{(V23.10)} give \textup{(V23.17)}, contradicting
\textup{(V23.16)}.
\end{proof}

By the prime-side definition \textup{(V22.1)}, the quantity in
\textup{(V23.15)} is exactly
\[
 \sup_{\xi\in I}
 \left[
 2\Re\sum_{n\ge2}\frac{\Lambda(n)\chi(n)}{\sqrt n}
 e^{-t(\log n)^2}e^{i\xi\log n}
 -a_{\chi,t}(\xi)
 \right]_+.
 \tag{V23.18}
\]
Thus \textup{(V23.16)} is a local, completed, prime-side criterion.  It asks
only that the negative excess have subexponential \(1/t\) growth on each
fixed frequency window; it does not ask for pointwise positivity at a
preassigned finite cutoff.

\subsection{Why absolute prime density saturates the strip ceiling}

The translation triangle inequality gives the absolute mass already used in
V20,
\[
 R(t)=2\sum_{n\ge2}\frac{\Lambda(n)}{\sqrt n}
 e^{-t(\log n)^2}.
 \tag{V23.19}
\]
The prime number theorem and the change of variables \(u=\log x\) give, at
the exponential scale,
\[
 R(t)=\exp\!\left(\frac{1+o(1)}{16t}\right)
 \quad(t\downarrow0),
 \tag{V23.20}
\]
because
\[
 \frac u2-tu^2
 =\frac{1}{16t}-t\left(u-\frac{1}{4t}\right)^2.
 \tag{V23.21}
\]
Equivalently,
\[
 \boxed{
 \lim_{t\downarrow0}16t\log R(t)=1.}
 \tag{V23.22}
\]
The saddle \(u=1/(4t)\) is exactly the V19 saddle after
\(t=(8\sigma^2)^{-1}\); V19 differentiated in scale and therefore carried
additional polynomial powers.

An off-line displacement \(\delta\) has zero-side exponential rate
\[
 \exp\!\left(\frac{\delta^2}{4t}\right),
 \qquad 0<\delta^2\le\frac14.
 \tag{V23.23}
\]
The absolute prime estimate \textup{(V23.20)} therefore has the rate
\(e^{(1/4)/(4t)}\), precisely the largest rate allowed anywhere in the
critical strip.  It cannot distinguish \(\delta=0\) from any positive
horizontal displacement.  Polynomial improvements to \(R(t)\), or unsigned
zero-counting estimates multiplying \textup{(V22.15)}, leave this exponential
obstruction unchanged.

\begin{theorem}[Unsigned density cannot prove the local criterion]
\label{ctr:v23-absolute-no-go}
Any proof of \textup{(V23.16)} which first replaces the character phases or
the complex-zero cosines by absolute values and then uses only polynomial
zero-density factors is quantitatively insufficient.  Its best universal
exponent is the strip ceiling \(1/(16t)\), whereas
\textup{(V23.16)} requires exponent zero for the negative part on every
fixed frequency interval.
\end{theorem}

\begin{proof}
On the prime side this is \textup{(V23.20)}--\textup{(V23.22)}.  On the zero
side, absolute values replace every cosine in \textup{(V22.15)} by at most
one and retain \(e^{\delta^2/(4t)}\).  The zero-count in Gaussian ordinate
shells is polynomial and hence does not change the \(1/t\) exponential rate.
Taking the strip supremum \(\delta\le1/2\) again gives \(1/(16t)\).
\end{proof}

The no-go is directional.  Zero density is sufficient to control the
remainder after one pair has been exposed, as in
Lemma~\ref{lem:v23-remainder}; it is not sufficient to prove positivity of
the whole sum.  The missing input must preserve the signed character phases
in \textup{(V23.18)} or produce a genuinely subexponential bound for its
negative excess.

\begin{firewall}[The local subexponential estimate is not established]
Theorem~\ref{thm:v23-local-criterion} is equivalent to GRH, and the attached
sources do not prove \textup{(V23.16)}.  Classical PNT cancellation for a
fixed nonprincipal character may improve lower-order factors or introduce
subleading exponent savings, but an absolute-value argument remains at the
full-strip rate \textup{(V23.22)}.  The exposed-pair theorem explains how an
off-line zero would manifest itself; it does not exclude such a zero.
\end{firewall}

\begin{assessment}[V23 acquisition]
V23 proves that every off-line zero has a uniquely exposed Gaussian-score
location and, after exact phase tuning, produces a negative scalar Weil
symbol with exponential rate \(E_\chi(x_0)/(4t)\).  This yields the local
negative-part criterion \textup{(V23.16)}, expressed directly by the
completed prime sum \textup{(V23.18)}.  It also proves that absolute prime
mass and unsigned zero density saturate the critical-strip ceiling
\(e^{1/(16t)}\) and therefore cannot establish the criterion.  The next
attack must remain on the signed prime side and seek subexponential control
of the negative excess on fixed frequency windows, rather than sum absolute
Gaussian zero tails.
\end{assessment}

\section*{Version V23 append-only record}
\addcontentsline{toc}{section}{Version V23 append-only record}
The complete V22 source was retained before this continuation.  It had
\(342423\) bytes, \(10232\) lines, and SHA--256
\[
 \texttt{B2D568CCC23C6E854E9B8AE84BDE6DC874967E6705F2FFF580BD4D2C4C650771}.
\]
V23 constructed the exposed-zero Laplace envelope, the phase-tuned negative
heat sequence, the local subexponential negative-part criterion, and the
absolute-density exponent no-go.

% =============================================================================
% END APPENDED V23 MATERIAL
% =============================================================================
% =============================================================================
% BEGIN APPENDED V24 MATERIAL
% =============================================================================

\section{V24: exact heat-rate tomography of the off-line zero set}
\label{sec:grh-v24}

V23 proved that one exposed off-line reflected pair forces exponentially
large negative values of the scalar heat symbol along a tuned sequence.  It
also showed that an absolute prime estimate has the full-strip rate
\(e^{1/(16t)}\).  The remaining quantitative question is whether the V23
lower bound is merely a witness or is the exact small-heat invariant.

This continuation proves exactness.  On every nondegenerate compact frequency
interval, the logarithmic rate of the negative part is precisely the largest
positive Gaussian score of an off-line zero pair meeting that interval.  The
result converts the signed, completed prime expression into a max-plus image
of the off-line zero geometry.  It also reconstructs the ordinate and
horizontal displacement of every locally exposed pair from the first
derivative of the rate field.  This is a sharper acquisition than the V14
triangular reader: V14 required a two-sided absolute growth estimate, whereas
the present invariant uses only the one-sided excess which can make the
completed Weil symbol negative.

\begin{remark}[Typographical correction in V23]
In \textup{(V23.5)}, the second displacement term is
\(\delta_q^2\).  Thus the displayed affine difference is
\[
 S_p(x)-S_q(x)
 =\delta_p^2-\delta_q^2-(x-\gamma_p)^2+(x-\gamma_q)^2.
 \tag{V24.1}
\]
The proof and every later use in V23 employ this corrected identity.
\end{remark}

\subsection{The interval score and a uniform Laplace upper bound}

Let \(\mathfrak P_\chi^{\mathrm{off}}\) be the collection of grouped
reflected pairs from V23 with \(\delta_p>0\).  For a nonempty compact interval
\(I\subset\mathbb R\), define
\[
 \begin{aligned}
 D_\chi(I)
 &:={}
 \max\left\{0,
 \sup_{\substack{p\in\mathfrak P_\chi^{\mathrm{off}}\\x\in I}}
 S_p(x)\right\}\\
 &=\max\left\{0,
 \sup_{p\in\mathfrak P_\chi^{\mathrm{off}}}
 \left(\delta_p^2-\operatorname{dist}(\gamma_p,I)^2\right)
 \right\}.
 \end{aligned}
 \tag{V24.2}
\]
Here and below the supremum over an empty off-line pair set is \(-\infty\),
so \(D_\chi(I)=0\) when GRH holds.
The supremum is a maximum whenever it is positive.  Indeed, pairs with
\(\lvert\gamma_p\rvert\to\infty\) have scores tending uniformly to
\(-\infty\) on \(I\), while the zero set is locally finite.

The zero-side formula \textup{(V22.12)} separates into a positive critical-line
part and the off-line reflected-pair part
\[
 \frac{4\pi}{\sqrt{4\pi t}}
 \sum_{p\in\mathfrak P_\chi^{\mathrm{off}}}m_p
 \exp\!\left(\frac{S_p(\xi)}{4t}\right)
 \cos\!\left(\frac{\delta_p(\xi-\gamma_p)}{2t}\right).
 \tag{V24.3}
\]

\begin{lemma}[Uniform interval upper rate]
\label{lem:v24-upper}
For every compact interval \(I\), there are constants
\(C_{\chi,I},B_{\chi,I}>0\) such that, for \(0<t\le1\),
\[
 \boxed{
 \mathcal N_{\chi,I}(t)
 \le C_{\chi,I}t^{-B_{\chi,I}}
 \exp\!\left(\frac{D_\chi(I)}{4t}\right).}
 \tag{V24.4}
\]
Consequently,
\[
 \limsup_{t\downarrow0}4t
 \log\bigl(1+\mathcal N_{\chi,I}(t)\bigr)
 \le D_\chi(I).
 \tag{V24.5}
\]
\end{lemma}

\begin{proof}
Every critical-line term in \textup{(V22.12)} is a positive real Gaussian.
Discarding that part can only increase the negative part.  Taking absolute
values only after this discard and using \textup{(V24.3)} gives
\[
 \mathcal N_{\chi,I}(t)
 \le\frac{4\pi}{\sqrt{4\pi t}}
 \sup_{\xi\in I}
 \sum_{p\in\mathfrak P_\chi^{\mathrm{off}}}m_p
 \exp\!\left(\frac{S_p(\xi)}{4t}\right).
 \tag{V24.6}
\]

Put \(M=1+\sup_{\xi\in I}|\xi|\) and split the pairs at
\(|\gamma_p|\le4M\).  The bounded-ordinate set is finite, and each of its
scores is at most \(D_\chi(I)\).  Its contribution to \textup{(V24.6)} is
therefore at most a constant times
\(t^{-1/2}e^{D_\chi(I)/(4t)}\).

For \(|\gamma_p|>4M\) and \(\xi\in I\), the critical-strip bound
\(\delta_p^2\le1/4\) gives
\[
 S_p(\xi)
 \le\frac14-(|\gamma_p|-M)^2
 \le-\frac18\gamma_p^2.
 \tag{V24.7}
\]
After enlarging the fixed splitting point if necessary, the last inequality
is uniform.  Summation over unit ordinate shells, using
\(N_\chi(Y+1)-N_\chi(Y)=O_\chi(\log(2+Y))\), shows that the remaining sum is
bounded by a polynomial in \(t^{-1}\) times \(e^{-c_I/t}\).  Since
\(D_\chi(I)\ge0\), this is absorbed into the right side of
\textup{(V24.4)}.  Equation \textup{(V24.5)} follows because
\(t\log(t^{-1})\to0\).
\end{proof}

The absolute value in this proof is harmless because it is applied only to
the zero pairs after their individual exponential scores have been retained.
It is not the prime-side absolute-value step ruled out by
Theorem~\ref{ctr:v23-absolute-no-go}.

\subsection{Near-maximal exposure inside an arbitrary interval}

The existence lemma in V23 selected one exposed pair somewhere on the real
axis.  Exact interval rates require a localized strengthening.

\begin{lemma}[Interior exposure arbitrarily close to the interval score]
\label{lem:v24-local-exposure}
Let \(I=[a,b]\) with \(a<b\), and suppose \(D_\chi(I)>0\).  For every
\(0<\varepsilon<D_\chi(I)\), there are \(x_0\in(a,b)\), an off-line pair
\(p_*\), and \(\eta>0\) such that
\[
 S_{p_*}(x_0)>D_\chi(I)-\varepsilon
 \tag{V24.8}
\]
and
\[
 S_p(x_0)\le S_{p_*}(x_0)-\eta
 \qquad(p\ne p_*).
 \tag{V24.9}
\]
The same strict inequality with gap \(\eta/2\) holds on a neighbourhood of
\(x_0\) contained in \((a,b)\).
\end{lemma}

\begin{proof}
By continuity in \(x\), an interior point can be chosen at which the
off-line envelope exceeds \(D_\chi(I)-\varepsilon/4\), even when the maximum
in \textup{(V24.2)} occurs at an endpoint.  Only finitely many pairs have
interval supremum at least \(D_\chi(I)-\varepsilon/2\): distant ordinates
have uniformly negative score.  Among these finitely many pairs, equality of
two distinct score functions occurs at at most one point, by
\textup{(V24.1)}, unless their ordinates agree.  For equal ordinates, the pair
with smaller displacement is everywhere strictly lower; identical data were
already grouped.

Perturb the chosen interior point by an arbitrarily small amount so that it
avoids all of the finitely many tie points and its envelope remains larger
than \(D_\chi(I)-\varepsilon/3\).  It then has a unique maximizing pair
\(p_*\).  Pairs outside the finite competitor set have score gap at least
\(\varepsilon/6\),
and the finitely many remaining differences have a positive minimum at the
chosen point.  This gives \(\eta>0\).  Continuity and the interior location
give the final neighbourhood assertion.  Since the winning score is
positive, \(p_*\) is off the critical line.
\end{proof}

\subsection{Exact local negative-part exponent}

\begin{theorem}[Exact interval heat-rate formula]
\label{thm:v24-exact-rate}
For every nondegenerate compact interval \(I\subset\mathbb R\),
\[
 \boxed{
 \limsup_{t\downarrow0}4t
 \log\bigl(1+\mathcal N_{\chi,I}(t)\bigr)
 =D_\chi(I).}
 \tag{V24.10}
\]
Thus the local negative-part rate is not merely bounded below by one exposed
zero: it is exactly the largest positive off-line Gaussian score visible from
the interval.
\end{theorem}

\begin{proof}
Lemma~\ref{lem:v24-upper} proves the upper bound.  If
\(D_\chi(I)=0\), the left side of \textup{(V24.10)} is nonnegative by
definition, so equality follows.

Suppose \(D_\chi(I)>0\), fix \(\varepsilon\in(0,D_\chi(I))\), and take
\(x_0,p_*,\eta\) from Lemma~\ref{lem:v24-local-exposure}.  Put
\(d_*=x_0-\gamma_{p_*}\).  If \(d_*\ne0\), use the V23 phase sequence
\[
 t_k=\frac{\delta_{p_*}|d_*|}{2(2k+1)\pi},
 \qquad \xi_k=x_0.
 \tag{V24.11}
\]
If \(d_*=0\), take any \(t_k\downarrow0\) and set
\[
 \xi_k=x_0+\frac{2\pi t_k}{\delta_{p_*}}.
 \tag{V24.12}
\]
For large \(k\), the second sequence remains in the exposure neighbourhood
and in \(I\).  In either case, the cosine of the \(p_*\)-term in
\textup{(V24.3)} is \(-1\).

The proof of Lemma~\ref{lem:v23-remainder}, applied with the local gap
\textup{(V24.9)}, shows that every other off-line term is little-oh of the
\(p_*\)-amplitude.  Every critical-line term has score at most zero and is
also exponentially smaller because \(S_{p_*}(x_0)>0\).  Therefore
\[
 m_{\chi,t_k}(\xi_k)<0
 \quad\hbox{and}\quad
 \liminf_{k\to\infty}4t_k
 \log\bigl(-m_{\chi,t_k}(\xi_k)\bigr)
 \ge S_{p_*}(x_0)>D_\chi(I)-\varepsilon.
 \tag{V24.13}
\]
Since \(\xi_k\in I\), the same lower bound holds for the limsup in
\textup{(V24.10)}.  Letting \(\varepsilon\downarrow0\) completes the proof.
\end{proof}

The use of a limsup is essential.  A fixed off-line pair oscillates between
positive and negative cosine lobes as \(t\downarrow0\), and no pointwise
limit of its signed exponential amplitude should be expected.

\subsection{The max-plus rate field and reconstruction of exposed pairs}

For \(x\in\mathbb R\), set
\[
 D_\chi(x)
 :=\max\left\{0,
 \sup_{p\in\mathfrak P_\chi^{\mathrm{off}}}
 \bigl(\delta_p^2-(x-\gamma_p)^2\bigr)
 \right\}.
 \tag{V24.14}
\]
Local finiteness and quadratic decay in the ordinate make this function
continuous.  Equivalently, by Theorem~\ref{thm:v24-exact-rate},
\[
 \boxed{
 D_\chi(x)
 =\lim_{r\downarrow0}
 \limsup_{t\downarrow0}4t
 \log\bigl(1+\mathcal N_{\chi,[x-r,x+r]}(t)\bigr).}
 \tag{V24.15}
\]
Indeed, \(D_\chi([x-r,x+r])\downarrow D_\chi(x)\) as \(r\downarrow0\).

\begin{corollary}[Influence intervals of off-line zeros]
\label{cor:v24-influence}
The positive set of the rate field is exactly
\[
 \boxed{
 \{x:D_\chi(x)>0\}
 =\bigcup_{p\in\mathfrak P_\chi^{\mathrm{off}}}
 (\gamma_p-\delta_p,\gamma_p+\delta_p).}
 \tag{V24.16}
\]
Consequently an off-line pair of displacement \(\delta_p\) forces positive
negative-part exponent throughout its open frequency influence interval, not
only at one specially selected observation point.
\end{corollary}

\begin{proof}
By \textup{(V24.14)}, \(D_\chi(x)>0\) exactly when
\(\delta_p^2-(x-\gamma_p)^2>0\) for at least one off-line pair \(p\).  This is
equivalent to \(|x-\gamma_p|<\delta_p\).
\end{proof}

The rate field is a tropical, or max-plus, projection of the zero set.  At a
point where a single pair is exposed, it contains enough differential data to
recover that pair.

\begin{corollary}[Differential reconstruction at a smooth exposed point]
\label{cor:v24-reconstruction}
Suppose \(D_\chi(x)>0\) and a unique pair \(p\) realizes the maximum in
\textup{(V24.14)} throughout a neighbourhood of \(x\).  Then
\[
 \boxed{
 \gamma_p=x+\frac12D_\chi'(x),
 \qquad
 \delta_p^2=D_\chi(x)+\frac14D_\chi'(x)^2.}
 \tag{V24.17}
\]
Moreover,
\[
 x^2+D_\chi(x)
 =\sup\left\{x^2,
 \sup_{p\in\mathfrak P_\chi^{\mathrm{off}}}
 \bigl(2\gamma_px+\delta_p^2-\gamma_p^2\bigr)
 \right\}
 \tag{V24.18}
\]
is convex.  Its affine exposed faces encode the locally visible reflected
zero pairs.
\end{corollary}

\begin{proof}
Under the uniqueness hypothesis,
\(D_\chi(x)=\delta_p^2-(x-\gamma_p)^2\) locally.  Differentiation gives
\(D_\chi'(x)=2(\gamma_p-x)\), and substitution proves
\textup{(V24.17)}.  Adding \(x^2\) to \textup{(V24.14)} gives
\textup{(V24.18)}, a supremum of the convex function \(x^2\) and affine
functions; hence it is convex.
\end{proof}

Put
\[
 \Delta_\chi
 :=\sup_{\rho\in\mathfrak Z_\chi}
 |\Re\rho-1/2|.
 \tag{V24.19}
\]
Taking \(x=\gamma_p\) and then the supremum over pairs gives the global
identity
\[
 \boxed{
 \sup_{x\in\mathbb R}D_\chi(x)=\Delta_\chi^2.}
 \tag{V24.20}
\]
Thus the largest horizontal violation of GRH is exactly the square root of
the largest local negative heat exponent.

\subsection{The exact signed-prime formulation}

Define the completed signed Gaussian prime excess
\[
 \mathcal E_{\chi,t}(\xi)
 :=2\Re\sum_{n\ge2}
 \frac{\Lambda(n)\chi(n)}{\sqrt n}
 e^{-t(\log n)^2}e^{i\xi\log n}
 -a_{\chi,t}(\xi).
 \tag{V24.21}
\]
By \textup{(V22.1)}, \(\mathcal E_{\chi,t}=-m_{\chi,t}\).  Therefore
Theorem~\ref{thm:v24-exact-rate} becomes the entirely prime-side identity
\[
 \boxed{
 \limsup_{t\downarrow0}4t\log\left(
 1+\sup_{\xi\in I}[\mathcal E_{\chi,t}(\xi)]_+
 \right)
 =D_\chi(I).}
 \tag{V24.22}
\]
In particular, a bound on the positive part of the completed prime excess is
not merely sufficient for GRH.  Its exact exponential type records the
off-line zero envelope.

\begin{corollary}[Quantitative zero exclusion from a local prime bound]
\label{cor:v24-zero-exclusion}
Suppose a nondegenerate compact interval \(I\) and a number \(\alpha\ge0\)
satisfy
\[
 \limsup_{t\downarrow0}4t\log\left(
 1+\sup_{\xi\in I}[\mathcal E_{\chi,t}(\xi)]_+
 \right)\le\alpha^2.
 \tag{V24.23}
\]
Then every off-line pair obeys
\[
 \boxed{
 \delta_p^2
 \le\alpha^2+\operatorname{dist}(\gamma_p,I)^2.}
 \tag{V24.24}
\]
In particular, if \(\gamma_p\in I\), then \(\delta_p\le\alpha\).
\end{corollary}

\begin{proof}
Equations \textup{(V24.2)}, \textup{(V24.22)}, and
\textup{(V24.23)} give
\(\delta_p^2-\operatorname{dist}(\gamma_p,I)^2\le\alpha^2\) for every
off-line pair.
\end{proof}

\subsection{What the classical fixed-character PNT yields at this scale}

The earlier V12 and V13 Abel estimates used triangular support weights.  The
same classical input can be transferred directly to the Gaussian reader, and
the resulting exponent locates the precise residual gap.

\begin{lemma}[Gaussian Abel bound from the classical zero-free region]
\label{lem:v24-gaussian-pnt}
Let \(I\subset\mathbb R\) be compact.  The fixed-character estimate
\[
 |\Psi_\chi(x)|
 \le C_\chi x e^{-c_\chi\sqrt{\log x}}
 \qquad(x\ge2)
 \tag{V24.25}
\]
implies that there are \(c'_{\chi}>0\), \(C_{\chi,I}<\infty\), and
\(t_{\chi,I}>0\) such that
\[
 \boxed{
 \sup_{\xi\in I}\left|
 \sum_{n\ge2}\frac{\Lambda(n)\chi(n)}{\sqrt n}
 e^{-t(\log n)^2}e^{i\xi\log n}
 \right|
 \le C_{\chi,I}
 \exp\!\left(\frac1{16t}-\frac{c'_{\chi}}{\sqrt t}\right)}
 \tag{V24.26}
\]
for \(0<t<t_{\chi,I}\).
\end{lemma}

\begin{proof}
Apply Stieltjes partial summation with
\[
 w_{t,\xi}(x)
 :=x^{-1/2+i\xi}e^{-t(\log x)^2}.
 \tag{V24.27}
\]
The boundary terms vanish, and, uniformly for \(\xi\in I\),
\[
 |w_{t,\xi}'(x)|
 \le C_Ix^{-3/2}
 \bigl(1+t\log x\bigr)e^{-t(\log x)^2}.
 \tag{V24.28}
\]
With \(u=\log x\), equations \textup{(V24.25)}--\textup{(V24.28)} reduce
the absolute value to a constant multiple of
\[
 \int_0^\infty(1+tu)
 \exp\!\left(\frac u2-tu^2-c_\chi\sqrt u\right)du.
 \tag{V24.29}
\]
Complete the square as in \textup{(V23.21)}.  On
\(u\ge1/(8t)\), the final term saves at least
\(c_\chi/\sqrt{8t}\).  On \(u<1/(8t)\), the quadratic part is at most
\(3/(64t)\), a saving of \(1/(64t)\) from the full saddle value
\(1/(16t)\).  Gaussian integration contributes only a power of \(t^{-1}\),
which is absorbed after decreasing the constant multiplying
\(t^{-1/2}\).  This proves \textup{(V24.26)}.
\end{proof}

The saving in \textup{(V24.26)} is substantial but subleading at the
\(1/t\) scale:
\[
 \limsup_{t\downarrow0}4t
 \log\left(1+\sup_{\xi\in I}\left|
 \sum_{n\ge2}\frac{\Lambda(n)\chi(n)}{\sqrt n}
 e^{-t(\log n)^2}e^{i\xi\log n}
 \right|\right)
 \le\frac14.
 \tag{V24.30}
\]
The completed gamma term is bounded on fixed \(I\) as \(t\downarrow0\), so
it does not alter this rate.  Theorem~\ref{thm:v24-exact-rate} shows exactly
what is missing: GRH requires rate zero for the positive part of the
completed excess, while the classical zero-free region leaves the maximal
strip rate \(1/4\).

\begin{firewall}[Heat tomography identifies but does not remove the zeros]
Equations \textup{(V24.10)} and \textup{(V24.22)} are exact identities, not an
unconditional proof that \(D_\chi(I)=0\).  The Gaussian Abel estimate
\textup{(V24.26)} improves the full absolute mass by
\(e^{-c/\sqrt t}\), but that improvement disappears after multiplication by
\(4t\) and passage to the logarithmic rate.  Proving that the positive part
of \(\mathcal E_{\chi,t}\) has rate zero remains equivalent to excluding all
off-line zeros.  No estimate established here performs that exclusion, so
GRH remains open.
\end{firewall}

\begin{assessment}[V24 acquisition]
V24 determines the exact small-heat invariant left unresolved in V23.  On
every nondegenerate compact interval, the exponential rate of the negative
part of the scalar Weil symbol equals the largest positive off-line zero
score \(D_\chi(I)\).  The induced rate field has support equal to the union of
the zero influence intervals, reconstructs every locally exposed pair by
\textup{(V24.17)}, and has global maximum \(\Delta_\chi^2\).  The same result
is the signed-prime identity \textup{(V24.22)}.  Finally, direct Gaussian Abel
summation proves that the classical fixed-character zero-free region supplies
only a subleading \(e^{-c/\sqrt t}\) saving and therefore cannot force the
rate to vanish.  The next attack should target a genuinely one-sided
mechanism for the completed excess---for example a positive-part moment or a
frequency-averaged inequality that retains cancellation before applying
convexity---rather than another absolute Gaussian estimate.
\end{assessment}

\section*{Version V24 append-only record}
\addcontentsline{toc}{section}{Version V24 append-only record}
The complete V23 source was retained before this continuation.  It had
\(355445\) bytes, \(10592\) lines, and SHA--256
\[
 \texttt{A78FDB2D15A1C163AE9334BC7755F0C8D358330FF68343A693B68839A051E13B}.
\]
V24 established the exact interval heat-rate formula, the max-plus rate
field and exposed-pair reconstruction, the quantitative local zero-exclusion
rule, and the fixed-character Gaussian Abel bound.

% =============================================================================
% END APPENDED V24 MATERIAL
% =============================================================================
% =============================================================================
% BEGIN APPENDED V25 MATERIAL
% =============================================================================

\section{V25: negative-lobe thickness and the averaging trichotomy}
\label{sec:grh-v25}

V24 identified the exact pointwise exponential rate of the negative part of
the scalar Weil heat symbol.  Its proposed next step was to ask whether a
positive-part moment or a frequency average could retain the arithmetic sign
while weakening the pointwise target.  This continuation decides the generic
forms of both proposals.

An exposed off-line pair does not create a negative value at an isolated
frequency.  At each phase-tuned heat time it creates an entire interval of
negative values of length comparable to \(t\), with the same exponential
height.  Consequently every fixed finite \(L^q\) moment of the negative part
has exactly the V24 max-plus rate.  Taking a moment after clipping therefore
does not make the GRH criterion easier.  On the other hand, Gaussian frequency
averaging is exactly additional heat evolution.  Averaging on the \(t\)-scale
only divides the exponent by a fixed factor, whereas averaging on a broader
scale makes the exponent disappear only because the wrong time variable is
being used.  A one-pair total-variation calculation exhibits the exponential
cancellation which an unclipped signed average conceals.

\subsection{A phase-tuned negative lobe has polynomial width}

Recall the interval score \(D_\chi(I)\) from \textup{(V24.2)}.  The following
strengthens the pointwise lower bound in Theorem~\ref{thm:v24-exact-rate}.

\begin{lemma}[Thick negative lobe below every visible rate]
\label{lem:v25-thick-lobe}
Let \(I=[a,b]\) with \(a<b\), and let
\[
 0<\alpha<D_\chi(I).
 \tag{V25.1}
\]
Then there are constants \(c_0,c_1>0\), a sequence \(t_k\downarrow0\), and
closed intervals \(J_k\subset I\) such that
\[
 |J_k|\ge c_0t_k
 \tag{V25.2}
\]
and, for all sufficiently large \(k\),
\[
 \boxed{
 -m_{\chi,t_k}(\xi)
 \ge c_1t_k^{-1/2}
 \exp\!\left(\frac{\alpha}{4t_k}
 +\frac{D_\chi(I)-\alpha}{16t_k}\right)
 \quad(\xi\in J_k).}
 \tag{V25.3}
\]
In particular,
\[
 \left|
 \left\{\xi\in I:
 -m_{\chi,t_k}(\xi)\ge e^{\alpha/(4t_k)}
 \right\}
 \right|
 \ge c_0t_k
 \tag{V25.4}
\]
for all sufficiently large \(k\).
\end{lemma}

\begin{proof}
Apply Lemma~\ref{lem:v24-local-exposure} with
\(\varepsilon=(D_\chi(I)-\alpha)/2\).  It gives an interior point \(x_0\),
a uniquely exposed off-line pair \(p_*\), and a fixed exposure neighbourhood
such that
\[
 S_{p_*}(x_0)
 >\frac{D_\chi(I)+\alpha}{2}.
 \tag{V25.5}
\]
Let \(\delta_*=\delta_{p_*}\), \(\gamma_*=\gamma_{p_*}\), and
\(d_*=x_0-\gamma_*\).  If \(d_*\ne0\), take
\[
 t_k=\frac{\delta_*|d_*|}{2(2k+1)\pi},
 \qquad \xi_k=x_0;
 \tag{V25.6}
\]
if \(d_*=0\), take any \(t_k\downarrow0\) and put
\[
 \xi_k=x_0+\frac{2\pi t_k}{\delta_*}.
 \tag{V25.7}
\]
In both cases the \(p_*\)-phase at \(\xi_k\) is an odd multiple of \(\pi\).

Choose a fixed \(c>0\) small enough that
\[
 \left|
 \frac{\delta_*(\xi-\xi_k)}{2t_k}
 \right|\le\frac14
 \qquad\text{whenever }|\xi-\xi_k|\le ct_k.
 \tag{V25.8}
\]
For large \(k\), the intervals
\(J_k=[\xi_k-ct_k,\xi_k+ct_k]\) lie in \(I\) and in the fixed exposure
neighbourhood.  On each \(J_k\), the dominant cosine is at most
\(-\cos(1/4)\).  Since \(S_{p_*}'\) is bounded on that neighbourhood,
\[
 S_{p_*}(\xi)=S_{p_*}(x_0)+O_{p_*,I}(t_k)
 \qquad(\xi\in J_k).
 \tag{V25.9}
\]

The uniform remainder estimate in Lemma~\ref{lem:v23-remainder}, with the
localized gap supplied by Lemma~\ref{lem:v24-local-exposure}, shows that all
other zero terms are little-oh of the dominant amplitude, uniformly on
\(J_k\).  Hence there is \(c_1>0\) such that
\[
 -m_{\chi,t_k}(\xi)
 \ge c_1t_k^{-1/2}
 \exp\!\left(\frac{S_{p_*}(x_0)}{4t_k}\right)
 \qquad(\xi\in J_k)
 \tag{V25.10}
\]
after absorbing the bounded exponential change in \textup{(V25.9)}.  From
\textup{(V25.5)},
\[
 \frac{D_\chi(I)+\alpha}{2}
 \ge \alpha+\frac{D_\chi(I)-\alpha}{4},
 \tag{V25.11}
\]
so \textup{(V25.10)} implies \textup{(V25.3)}.  Equation
\textup{(V25.4)} follows because the extra exponential factor in
\textup{(V25.3)} eventually dominates \(c_1^{-1}t_k^{1/2}\).
\end{proof}

The interval width in \textup{(V25.2)} is only polynomially small in heat
time.  It therefore has no cost at the \(1/t\) logarithmic scale.  This is the
mechanism behind the moment theorem below.

\subsection{Every fixed positive-part moment has the same exponent}

For \(0<q<\infty\), define
\[
 \mathcal M_{\chi,I}^{(q)}(t)
 :=\left(
 \int_I[-m_{\chi,t}(\xi)]_+^q\,d\xi
 \right)^{1/q}.
 \tag{V25.12}
\]
For \(q=\infty\), set
\(\mathcal M_{\chi,I}^{(\infty)}=\mathcal N_{\chi,I}\).

\begin{theorem}[Exact fixed-moment heat rate]
\label{thm:v25-moment-rate}
For every nondegenerate compact interval \(I\) and every
\(0<q\le\infty\),
\[
 \boxed{
 \limsup_{t\downarrow0}4t
 \log\bigl(1+\mathcal M_{\chi,I}^{(q)}(t)\bigr)
 =D_\chi(I).}
 \tag{V25.13}
\]
Equivalently, for \(0<q<\infty\),
\[
 \boxed{
 \limsup_{t\downarrow0}\frac{4t}{q}
 \log\left(1+
 \int_I[-m_{\chi,t}(\xi)]_+^q\,d\xi
 \right)
 =D_\chi(I).}
 \tag{V25.14}
\]
\end{theorem}

\begin{proof}
The case \(q=\infty\) is Theorem~\ref{thm:v24-exact-rate}.  For finite
\(q\),
\[
 \mathcal M_{\chi,I}^{(q)}(t)
 \le |I|^{1/q}\mathcal N_{\chi,I}(t),
 \tag{V25.15}
\]
so Lemma~\ref{lem:v24-upper} gives the upper rate \(D_\chi(I)\).

If \(D_\chi(I)=0\), nonnegativity of the left side proves equality.  If
\(D_\chi(I)>0\), fix \(0<\alpha<D_\chi(I)\) and use the sets in
\textup{(V25.4)}.  They give
\[
 \mathcal M_{\chi,I}^{(q)}(t_k)
 \ge(c_0t_k)^{1/q}e^{\alpha/(4t_k)}.
 \tag{V25.16}
\]
Thus the limsup in \textup{(V25.13)} is at least \(\alpha\).  Let
\(\alpha\uparrow D_\chi(I)\).  Formula \textup{(V25.14)} follows from
\[
 \log(1+X^q)=q\log(1+X)+O_q(1)
 \qquad(X\ge0),
 \tag{V25.17}
\]
with \(X=\mathcal M_{\chi,I}^{(q)}(t)\).
\end{proof}

\begin{corollary}[Fixed positive-part moments are GRH-equivalent]
\label{cor:v25-moment-grh}
Fix one \(q\in(0,\infty]\).  GRH holds for the paired channel if and only if
for every nondegenerate compact interval \(I\),
\[
 \limsup_{t\downarrow0}4t
 \log\bigl(1+\mathcal M_{\chi,I}^{(q)}(t)\bigr)=0.
 \tag{V25.18}
\]
No fixed choice of \(q\), including \(q<1\), reduces the exponential
difficulty of the local criterion.
\end{corollary}

\begin{proof}
Under GRH the scalar symbol is positive for every \(t>0\), so every moment
vanishes.  Conversely, if an off-line zero exists, its influence interval in
\textup{(V24.16)} contains a nondegenerate compact subinterval with positive
score.  Theorem~\ref{thm:v25-moment-rate} then contradicts
\textup{(V25.18)}.
\end{proof}

This rules out a generic post-processing shortcut: Hölder, Markov, or convex
moment inequalities applied after formation of \([-m_{\chi,t}]_+\) still face
the full off-line exponent.

\subsection{Gaussian averaging is exactly a change of heat time}

For \(s\ge0\), let \(H_s\) be the Gaussian kernel from \textup{(V20.6)}.
The semigroup identity \textup{(V22.8)} says
\[
 H_s*m_{\chi,t}=m_{\chi,t+s}.
 \tag{V25.19}
\]
Thus a Gaussian frequency average cannot be analyzed independently of its
averaging scale.

\begin{theorem}[Averaging-scale trichotomy]
\label{thm:v25-averaging-trichotomy}
Let \(I\) be a nondegenerate compact interval.

\begin{enumerate}[label=\textup{(\roman*)},leftmargin=3em]
\item For every fixed \(\lambda\ge0\), define
\[
 \widetilde m_{\chi,t}^{(\lambda)}
 :=H_{\lambda t}*m_{\chi,t}=m_{\chi,(1+\lambda)t}.
 \tag{V25.20}
\]
Then
\[
 \boxed{
 \limsup_{t\downarrow0}4t\log\left(
 1+\sup_{\xi\in I}[-\widetilde m_{\chi,t}^{(\lambda)}(\xi)]_+
 \right)
 =\frac{D_\chi(I)}{1+\lambda}.}
 \tag{V25.21}
\]

\item If \(0<\theta<1\) and
\[
 \widehat m_{\chi,t}^{(\theta)}
 :=H_{t^\theta}*m_{\chi,t}=m_{\chi,t+t^\theta},
 \tag{V25.22}
\]
then, whether or not GRH holds,
\[
 \boxed{
 \limsup_{t\downarrow0}4t\log\left(
 1+\sup_{\xi\in I}[-\widehat m_{\chi,t}^{(\theta)}(\xi)]_+
 \right)=0.}
 \tag{V25.23}
\]
With the correct effective time \(\tau=t+t^\theta\), however,
\[
 \boxed{
 \limsup_{t\downarrow0}4\tau\log\left(
 1+\sup_{\xi\in I}[-\widehat m_{\chi,t}^{(\theta)}(\xi)]_+
 \right)=D_\chi(I).}
 \tag{V25.24}
\]

\item For every fixed \(s_0>0\), the family
\(H_{s_0}*m_{\chi,t}=m_{\chi,s_0+t}\) is uniformly bounded below on
\(I\) as \(t\downarrow0\).  Hence its negative part has zero
\(1/t\)-exponential rate even when off-line zeros exist.
\end{enumerate}
\end{theorem}

\begin{proof}
For \textup{(i)}, substitute \(\tau=(1+\lambda)t\) in
Theorem~\ref{thm:v24-exact-rate}.  Multiplication by
\(t/\tau=(1+\lambda)^{-1}\) gives \textup{(V25.21)}.

For \textup{(ii)}, Lemma~\ref{lem:v24-upper} at
\(\tau=t+t^\theta\) gives
\[
 4t\log\left(1+\mathcal N_{\chi,I}(\tau)\right)
 \le\frac{t}{\tau}D_\chi(I)
 +O_{\chi,I}\!\left(t\log\frac1\tau\right).
 \tag{V25.25}
\]
Here \(t/\tau\to0\), so the right side tends to zero; the left side is
nonnegative.  This proves \textup{(V25.23)}.  The function
\(t\mapsto t+t^\theta\) is continuous, strictly increasing, and maps a
punctured neighbourhood of zero onto another such neighbourhood.  Applying
Theorem~\ref{thm:v24-exact-rate} after this change of variable proves
\textup{(V25.24)}.

For \textup{(iii)}, positive-time smoothness makes
\(m_{\chi,s_0+t}\) continuous jointly in \((t,\xi)\) on a compact
neighbourhood of \(\{0\}\times I\).  Its negative part is therefore bounded,
and multiplication of its logarithm by \(t\) gives zero.
\end{proof}

The vanishing in \textup{(V25.23)} is a normalization artifact, not progress
toward GRH.  Broader smoothing replaces the diagnostic time \(t\) by the
larger time \(t+t^\theta\); using the latter restores the complete V24 rate.

\subsection{One reflected pair: bounded signed mass, exponential variation}

The averaging trichotomy can be seen directly on one off-line pair.  Write
its full contribution to \(m_{\chi,t}\) as
\[
 K_{p,t}(x)
 :=\frac{4\pi m_p}{\sqrt{4\pi t}}
 \exp\!\left(
 \frac{\delta_p^2-(x-\gamma_p)^2}{4t}
 \right)
 \cos\!\left(
 \frac{\delta_p(x-\gamma_p)}{2t}
 \right),
 \qquad\delta_p>0.
 \tag{V25.26}
\]

\begin{proposition}[Exponential cancellation inside one pair]
\label{prop:v25-pair-variation}
For every off-line reflected pair,
\[
 \boxed{
 \int_{\mathbb R}K_{p,t}(x)\,dx=4\pi m_p}
 \tag{V25.27}
\]
for all \(t>0\), whereas
\[
 \boxed{
 \lim_{t\downarrow0}
 e^{-\delta_p^2/(4t)}
 \int_{\mathbb R}|K_{p,t}(x)|\,dx
 =8m_p.}
 \tag{V25.28}
\]
Consequently,
\[
 \boxed{
 \begin{aligned}
 \lim_{t\downarrow0}e^{-\delta_p^2/(4t)}
 \int_{\mathbb R}[K_{p,t}(x)]_+\,dx&=4m_p,\\
 \lim_{t\downarrow0}e^{-\delta_p^2/(4t)}
 \int_{\mathbb R}[-K_{p,t}(x)]_+\,dx&=4m_p.
 \end{aligned}}
 \tag{V25.29}
\]
Thus the bounded signed mass is the difference of positive and negative
masses which are each exponentially large.
\end{proposition}

\begin{proof}
The two complex-conjugate Gaussian terms comprising \(K_{p,t}\) each have
real-line integral one before multiplication by \(2\pi m_p\).  Equivalently,
the elementary Fourier transform of a real Gaussian cancels the factor
\(e^{\delta_p^2/(4t)}\).  This proves \textup{(V25.27)}.

For the variation, set \(y=(x-\gamma_p)/(2\sqrt t)\).  Then
\[
 e^{-\delta_p^2/(4t)}
 \int_{\mathbb R}|K_{p,t}(x)|\,dx
 =4\sqrt\pi m_p
 \int_{\mathbb R}e^{-y^2}
 \left|\cos\!\left(\frac{\delta_p y}{\sqrt t}\right)\right|dy.
 \tag{V25.30}
\]
The periodic function \(|\cos u|\) has mean \(2/\pi\).  Its Fourier series,
followed by the Riemann--Lebesgue lemma against \(e^{-y^2}\), gives
\[
 \int_{\mathbb R}e^{-y^2}
 \left|\cos\!\left(\frac{\delta_p y}{\sqrt t}\right)\right|dy
 \longrightarrow\frac2\pi\int_{\mathbb R}e^{-y^2}dy
 =\frac2{\sqrt\pi}.
 \tag{V25.31}
\]
This proves \textup{(V25.28)}.  Finally,
\([K]_+=(|K|+K)/2\) and \([-K]_+=(|K|-K)/2\).  The fixed mass
\textup{(V25.27)} vanishes after multiplication by
\(e^{-\delta_p^2/(4t)}\), yielding \textup{(V25.29)}.
\end{proof}

This calculation explains both sides of the averaging problem.  A signed
linear average can hide the entire off-line exponent by cancellation between
rapid lobes.  Taking absolute values or positive parts restores the exponent,
but Theorem~\ref{thm:v25-moment-rate} then shows that the resulting estimate is
again GRH-equivalent.

\begin{theorem}[Generic post-processing cannot supply an intermediate bound]
\label{thm:v25-postprocessing-no-go}
The following three generic modifications of the V24 local criterion do not
produce a strictly weaker route to GRH:
\begin{enumerate}[label=\textup{(\alph*)},leftmargin=3em]
\item replacing the local supremum of the negative part by any fixed
\(L^q\) moment, because the exponent remains \(D_\chi(I)\);
\item applying Gaussian smoothing of width comparable to \(t\), because the
exponent is only rescaled by \((1+\lambda)^{-1}\);
\item applying Gaussian smoothing broader than \(t\) and retaining the old
\(1/t\) normalization, because the apparent zero exponent is restored to
\(D_\chi(I)\) at the effective heat time.
\end{enumerate}
Any useful averaging mechanism must therefore introduce new arithmetic
structure before clipping or norm formation; it cannot be a generic scalar
post-processing inequality for the already assembled heat symbol.
\end{theorem}

\begin{proof}
Part \textup{(a)} is Theorem~\ref{thm:v25-moment-rate}, part \textup{(b)} is
\textup{(V25.21)}, and part \textup{(c)} is
\textup{(V25.23)}--\textup{(V25.24)}.  Proposition~\ref{prop:v25-pair-variation}
shows why passing instead to an unclipped signed average is also insufficient:
it may cancel positive and negative masses of the same exponential size.
\end{proof}

\begin{firewall}[The averaging audit is not a positivity estimate]
V25 proves that several natural weakenings retain or erase the V24
obstruction in exactly quantifiable ways.  It does not establish a new upper
bound for \([-m_{\chi,t}]_+\), and hence does not prove that
\(D_\chi(I)=0\).  The thick-lobe theorem is conditional on the existence of a
positive visible score and describes the contradiction that an off-line zero
would create.  GRH remains open.
\end{firewall}

\begin{assessment}[V25 acquisition]
V25 proves that an exposed off-line zero creates a negative lobe of width
\(\asymp t\), not an isolated spike.  Every fixed positive-part \(L^q\)
moment consequently has the exact max-plus exponent \(D_\chi(I)\).  The heat
semigroup gives a complete Gaussian averaging trichotomy: \(t\)-scale
smoothing attenuates but preserves the exponent, broader smoothing only
changes the correct time normalization, and fixed smoothing loses the
small-time detector.  A one-pair calculation shows that bounded signed mass
can conceal exponentially large positive and negative variation.  These
results close the generic moment/averaging branch.  The next attack must go
upstream and seek an arithmetic operation before scalar clipping---for
example a prime-cell bilinear identity, a one-sided entropy production law,
or a comparison coupling the completed excess to its heat derivative at the
same effective time.
\end{assessment}

\section*{Version V25 append-only record}
\addcontentsline{toc}{section}{Version V25 append-only record}
The complete V24 source was retained before this continuation.  It had
\(373132\) bytes, \(11096\) lines, and SHA--256
\[
 \texttt{37AD3070E899BBA5D4EC2E18B1B9999125C869670726CED8F30E832A3EC58725}.
\]
V25 established negative-lobe thickness, exact fixed-moment rates, the
Gaussian averaging-scale trichotomy, and the reflected-pair
signed-mass/variation cancellation law.

% =============================================================================
% END APPENDED V25 MATERIAL
% =============================================================================
% =============================================================================
% BEGIN APPENDED V26 MATERIAL
% =============================================================================

\section{V26: negative entropy, adaptive prime Grams, and the free boundary}
\label{sec:grh-v26}

V25 showed that fixed positive-part moments and generic frequency smoothing
do not provide an intermediate route to GRH.  It therefore proposed moving
the arithmetic operation before scalar clipping.  The natural candidate is
the entropy production of the complete heat symbol itself.  Because
\(m_{\chi,t}\) satisfies the heat equation, its negative part has an exact
dissipation law.  Differentiating the completed prime formula then turns the
dissipation into a positive-semidefinite Gram matrix whose integration region
is the negative set selected by the full arithmetic symbol.

This continuation derives that identity without separating the gamma and
prime pieces by an inequality.  It also determines the exact local
exponential rate of the gradient dissipation.  The result is both an
acquisition and a warning: the adaptive Gram is genuinely upstream of the
post-processing operations ruled out in V25, but a subexponential bound for
its total contraction is itself GRH-equivalent.  Its possible extra content
lies only in the free-boundary constraints at the endpoints of the negative
set.

\subsection{Compact negative sets and the entropy identity}

For \(t>0\), define
\[
 \Omega_{\chi,t}
 :=\{\xi\in\mathbb R:m_{\chi,t}(\xi)<0\},
 \qquad
 w_{\chi,t}:=[-m_{\chi,t}]_+.
 \tag{V26.1}
\]
By \textup{(V22.16)}, \(m_{\chi,t}(\xi)\to+\infty\) as
\(|\xi|\to\infty\), so \(\Omega_{\chi,t}\) is bounded.  Positive heat time
also makes the symbol real analytic in \(\xi\).  Since it is not identically
zero, its zero set in a compact interval is finite.  Thus
\(\Omega_{\chi,t}\) is a finite union of bounded open intervals.

Set
\[
 \mathcal S_\chi(t)
 :=\int_{\mathbb R}w_{\chi,t}(\xi)^2\,d\xi,
 \qquad
 \mathcal D_\chi(t)
 :=\int_{\Omega_{\chi,t}}
 |\partial_\xi m_{\chi,t}(\xi)|^2\,d\xi.
 \tag{V26.2}
\]

\begin{theorem}[Exact negative-entropy dissipation]
\label{thm:v26-entropy}
The function \(\mathcal S_\chi\) is locally absolutely continuous on
\((0,\infty)\), and for almost every \(t>0\),
\[
 \boxed{
 \mathcal S_\chi'(t)=-2\mathcal D_\chi(t).}
 \tag{V26.3}
\]
There is \(T_\chi<\infty\) such that \(w_{\chi,t}=0\) for
\(t\ge T_\chi\), and consequently
\[
 \boxed{
 \mathcal S_\chi(t)
 =2\int_t^{T_\chi}\mathcal D_\chi(s)\,ds
 =2\int_t^\infty\mathcal D_\chi(s)\,ds.}
 \tag{V26.4}
\]
\end{theorem}

\begin{proof}
On every compact positive heat-time interval, the prime series and all of its
\(t\)- and \(\xi\)-derivatives converge locally uniformly, while the
archimedean convolution has the same regularity.  The negative set remains
bounded locally in time by the logarithmic growth in \textup{(V22.16)}.
Standard Sobolev chain rules for the Lipschitz map \(r\mapsto[-r]_+\), or a
smooth convex approximation followed by dominated convergence, therefore
give local absolute continuity of \(\mathcal S_\chi\).

On \(\Omega_{\chi,t}\), one has \(w_{\chi,t}=-m_{\chi,t}\).  Using the heat
equation \textup{(V22.3)} and integrating by parts on every component of the
negative set gives
\[
 \begin{aligned}
 \mathcal S_\chi'(t)
 &=2\int_{\Omega_{\chi,t}}
 m_{\chi,t}\,\partial_t m_{\chi,t}\,d\xi\\
 &=2\int_{\Omega_{\chi,t}}
 m_{\chi,t}\,\partial_\xi^2m_{\chi,t}\,d\xi\\
 &=-2\int_{\Omega_{\chi,t}}
 |\partial_\xi m_{\chi,t}|^2\,d\xi.
 \end{aligned}
 \tag{V26.5}
\]
The boundary terms vanish because \(m_{\chi,t}=0\) at every finite endpoint.
The approximation argument covers degenerate contact times, which form no
obstruction to the almost-everywhere identity.

Eventual coercivity from Theorem~\ref{thm:v20-threshold}, combined with the
pointwise equivalence in Theorem~\ref{thm:v22-scalarization}, supplies a
finite \(T_\chi\) for which \(m_{\chi,t}\ge0\) whenever
\(t\ge T_\chi\).  Integrating \textup{(V26.3)} from \(t\) to
\(T_\chi\) proves \textup{(V26.4)}.
\end{proof}

More generally, for every \(q\ge2\), the same chain rule gives
\[
 \frac d{dt}\int_{\mathbb R}w_{\chi,t}^q\,d\xi
 =-q(q-1)\int_{\Omega_{\chi,t}}
 w_{\chi,t}^{q-2}|\partial_\xi m_{\chi,t}|^2\,d\xi
 \tag{V26.6}
\]
for almost every positive time.  The quadratic case is distinguished because
its production is an unweighted Gram square.

\begin{corollary}[Entropy-rate criterion for GRH]
\label{cor:v26-entropy-grh}
The following are equivalent:
\begin{enumerate}[label=\textup{(\roman*)},leftmargin=3em]
\item GRH holds for the paired channel;
\item \(\mathcal S_\chi(t)=0\) for every \(t>0\);
\item
\[
 \boxed{
 \limsup_{t\downarrow0}2t
 \log\bigl(1+\mathcal S_\chi(t)\bigr)=0.}
 \tag{V26.7}
\]
\end{enumerate}
If GRH fails, the limsup in \textup{(V26.7)} is strictly positive.
\end{corollary}

\begin{proof}
Under GRH, \textup{(V22.26)} gives \(w_{\chi,t}=0\).  Conversely,
\(\mathcal S_\chi(t)=0\) and continuity of the symbol imply
\(m_{\chi,t}\ge0\) pointwise, so the V22 criterion gives GRH.

If GRH fails, Corollary~\ref{cor:v24-influence} supplies a nondegenerate
compact interval \(I\) with \(D_\chi(I)>0\).  Since
\[
 \mathcal S_\chi(t)
 \ge\int_I[-m_{\chi,t}(\xi)]_+^2\,d\xi,
 \tag{V26.8}
\]
Theorem~\ref{thm:v25-moment-rate}, in the form \textup{(V25.14)} with
\(q=2\), makes the limsup in \textup{(V26.7)} at least
\(D_\chi(I)>0\).  This proves all implications.
\end{proof}

\subsection{The exact local rate of gradient production}

For a compact interval \(I\), put
\[
 \mathcal D_{\chi,I}(t)
 :=\int_{I\cap\Omega_{\chi,t}}
 |\partial_\xi m_{\chi,t}(\xi)|^2\,d\xi.
 \tag{V26.9}
\]

\begin{lemma}[Differentiated interval upper bound]
\label{lem:v26-gradient-upper}
For every compact interval \(I\), there are constants
\(C_{\chi,I},B_{\chi,I}>0\) such that
\[
 \sup_{\xi\in I}|\partial_\xi m_{\chi,t}(\xi)|
 \le C_{\chi,I}t^{-B_{\chi,I}}
 \exp\!\left(\frac{D_\chi(I)}{4t}\right)
 \qquad(0<t\le1).
 \tag{V26.10}
\]
\end{lemma}

\begin{proof}
Differentiate the absolutely convergent zero expansion
\textup{(V22.12)} term by term.  For an off-line pair, differentiating
\textup{(V24.3)} introduces only the factor
\[
 O\!\left(
 \frac{1+|\xi-\gamma_p|}{t}
 \right)
 \tag{V26.11}
\]
in addition to its Gaussian amplitude.  The bounded-ordinate pairs are
finite and have score at most \(D_\chi(I)\).  For distant ordinates, the
quadratic bound \textup{(V24.7)} absorbs both the factor in
\textup{(V26.11)} and the logarithmic zero count into a power of \(t^{-1}\).
Critical-line terms have nonpositive score on the real axis and obey the
same polynomial bound.  Summing the differentiated shells proves
\textup{(V26.10)}.
\end{proof}

\begin{lemma}[A sloping sublobe of an exposed pair]
\label{lem:v26-sloping-lobe}
Let \(I\) be nondegenerate and \(0<\alpha<D_\chi(I)\).  There are
\(c_0,c_1>0\), times \(t_k\downarrow0\), and intervals
\(L_k\subset I\cap\Omega_{\chi,t_k}\) such that
\[
 |L_k|\ge c_0t_k
 \tag{V26.12}
\]
and
\[
 \boxed{
 |\partial_\xi m_{\chi,t_k}(\xi)|
 \ge c_1t_k^{-3/2}e^{\alpha/(4t_k)}
 \qquad(\xi\in L_k)}
 \tag{V26.13}
\]
for all sufficiently large \(k\).
\end{lemma}

\begin{proof}
Use the exposed pair and phase-tuned sequence in the proof of
Lemma~\ref{lem:v25-thick-lobe}, choosing its score strictly larger than
\(\alpha\).  Its derivative is
\[
 \begin{aligned}
 \partial_\xi K_{p_*,t}(\xi)
 ={}&\frac{4\pi m_{p_*}}{\sqrt{4\pi t}}
 e^{S_{p_*}(\xi)/(4t)}\frac1{2t}\\
 &\times\left[
 -(\xi-\gamma_*)\cos\theta
 -\delta_*\sin\theta
 \right],
 \qquad
 \theta=\frac{\delta_*(\xi-\gamma_*)}{2t}.
 \end{aligned}
 \tag{V26.14}
\]

If \(x_0-\gamma_*\ne0\), a sufficiently small fixed phase interval about
the tuned odd multiple of \(\pi\) makes the bracket in
\textup{(V26.14)} bounded away from zero.  If
\(x_0=\gamma_*\), move instead to a fixed one-sided phase subinterval
\(\theta=(2j+1)\pi+h\) with, for example,
\(1/8\le h\le1/4\).  There the sine term is bounded away from zero, the
cosine remains negative, and the corresponding \(\xi\)-interval has length
comparable to \(t\).  In both cases this produces an interval \(L_k\) of
length at least \(c_0t_k\) on which the dominant pair is negative and its
derivative has the lower bound in \textup{(V26.13)}, before reduction of the
constant.

Differentiating the remainder estimate costs only polynomial powers of
\(t_k^{-1}\) and ordinate factors.  The fixed exposure gap absorbs all such
factors, exactly as in Lemma~\ref{lem:v26-gradient-upper}.  Hence the full
symbol is negative on \(L_k\), while its derivative remains at least one half
of the dominant derivative for large \(k\).  This proves the assertion.
\end{proof}

\begin{theorem}[Exact local entropy-production rate]
\label{thm:v26-production-rate}
For every nondegenerate compact interval \(I\),
\[
 \boxed{
 \limsup_{t\downarrow0}2t
 \log\bigl(1+\mathcal D_{\chi,I}(t)\bigr)
 =D_\chi(I).}
 \tag{V26.15}
\]
\end{theorem}

\begin{proof}
Equations \textup{(V26.9)} and \textup{(V26.10)} give
\[
 \mathcal D_{\chi,I}(t)
 \le |I|C_{\chi,I}^2t^{-2B_{\chi,I}}
 e^{D_\chi(I)/(2t)},
 \tag{V26.16}
\]
which proves the upper rate.  If \(D_\chi(I)=0\), nonnegativity gives
equality.  If \(D_\chi(I)>0\), choose
\(0<\alpha<D_\chi(I)\) in Lemma~\ref{lem:v26-sloping-lobe}.  Then
\[
 \mathcal D_{\chi,I}(t_k)
 \ge c_0c_1^2t_k^{-2}e^{\alpha/(2t_k)}.
 \tag{V26.17}
\]
The limsup is at least \(\alpha\); let
\(\alpha\uparrow D_\chi(I)\).
\end{proof}

The entropy production is therefore no cheaper at exponential scale than the
negative entropy it dissipates.  The gain in moving upstream is not a smaller
target; it is the bilinear algebra exposed next.

\subsection{The adaptive completed prime-cell Gram matrix}

For \(s>0\), write \(u_n=\log n\) and
\[
 c_n(s)
 :=\frac{\Lambda(n)\chi(n)u_n}{\sqrt n}e^{-su_n^2}.
 \tag{V26.18}
\]
The prime-side derivative of \textup{(V22.1)} is
\[
 \partial_\xi m_{\chi,s}(\xi)
 =b_0(s,\xi)+\sum_{n\ge2}b_n(s,\xi),
 \tag{V26.19}
\]
where
\[
 b_0(s,\xi):=\partial_\xi a_{\chi,s}(\xi),
 \qquad
 b_n(s,\xi):=2\Im\bigl(c_n(s)e^{iu_n\xi}\bigr).
 \tag{V26.20}
\]
The series is absolutely and locally uniformly convergent.

Index the gamma row by \(0\), and define the real symmetric matrix
\[
 \mathbb G_\chi(s)_{r\ell}
 :=\int_{\Omega_{\chi,s}}
 b_r(s,\xi)b_\ell(s,\xi)\,d\xi,
 \qquad r,\ell\in\{0,2,3,\ldots\}.
 \tag{V26.21}
\]

\begin{theorem}[Adaptive prime-cell Gram representation]
\label{thm:v26-adaptive-gram}
For every \(s>0\), \(\mathbb G_\chi(s)\) is positive semidefinite on
finitely supported coefficient vectors, and its absolutely convergent total
contraction satisfies
\[
 \boxed{
 \mathcal D_\chi(s)
 =\sum_{r,\ell\in\{0,2,3,\ldots\}}
 \mathbb G_\chi(s)_{r\ell}.}
 \tag{V26.22}
\]
Consequently, for every \(t>0\),
\[
 \boxed{
 \mathcal S_\chi(t)
 =2\int_t^{T_\chi}
 \sum_{r,\ell}
 \mathbb G_\chi(s)_{r\ell}\,ds.}
 \tag{V26.23}
\]
This identity retains the completed gamma--prime and prime--prime cross terms
before any absolute value is taken.
\end{theorem}

\begin{proof}
For every finitely supported complex vector \(z=(z_r)\),
\[
 \sum_{r,\ell}\overline{z_r}
 \mathbb G_\chi(s)_{r\ell}z_\ell
 =\int_{\Omega_{\chi,s}}
 \left|\sum_rz_rb_r(s,\xi)\right|^2d\xi\ge0.
 \tag{V26.24}
\]
Thus the matrix is positive semidefinite.  Absolute convergence of the
Gaussian-damped derivative series permits expansion of the square in
\textup{(V26.19)} under the bounded-domain integral.  Taking every
coefficient equal to one gives \textup{(V26.22)}.  Substitution into the
entropy identity \textup{(V26.4)} proves \textup{(V26.23)}.
\end{proof}

The matrix is explicit once the adaptive negative set is known.  Define its
Fourier transform
\[
 F_{\chi,s}(v)
 :=\int_{\Omega_{\chi,s}}e^{iv\xi}\,d\xi.
 \tag{V26.25}
\]
For prime indices \(n,m\ge2\), the identity
\(4\Im A\Im B=2\Re(A\overline B-AB)\) gives
\[
 \boxed{
 \begin{aligned}
 \mathbb G_\chi(s)_{nm}
 =2\Re\bigl[{}
 &c_n(s)\overline{c_m(s)}
 F_{\chi,s}(u_n-u_m)\\
 &-c_n(s)c_m(s)
 F_{\chi,s}(u_n+u_m)
 \bigr].
 \end{aligned}}
 \tag{V26.26}
\]
The first term is a Toeplitz difference-frequency block; the second is a
Hankel sum-frequency block.  Both occur with their arithmetic phases intact.

If
\[
 \Omega_{\chi,s}=\bigcup_{j=1}^{J(s)}(\alpha_j(s),\beta_j(s)),
 \tag{V26.27}
\]
then, for \(v\ne0\),
\[
 F_{\chi,s}(v)
 =\sum_{j=1}^{J(s)}
 \frac{e^{iv\beta_j(s)}-e^{iv\alpha_j(s)}}{iv},
 \tag{V26.28}
\]
and hence
\[
 \boxed{
 |F_{\chi,s}(v)|
 \le\min\left\{|\Omega_{\chi,s}|,
 \frac{2J(s)}{|v|}\right\}.}
 \tag{V26.29}
\]
This endpoint formula is the precise free-boundary information absent from a
fixed frequency average.

\subsection{Capacity and the remaining free-boundary target}

The positivity of \(\mathbb G_\chi(s)\) does not make its all-ones
contraction small.  Nor does \textup{(V26.29)} alone close the estimate.  Near
the Gaussian prime saddle, adjacent logarithms have spacing
\[
 \log(n+1)-\log n\asymp\frac1n,
 \tag{V26.30}
\]
so the difference-frequency denominator in \textup{(V26.29)} can be
exponentially large in \(1/s\).  Replacing the adaptive endpoints by their
count or total variation would therefore return the absolute-density barrier
of V23.

There is, however, additional structure not present in an arbitrary Gram
window.  Every endpoint in \textup{(V26.27)} obeys the completed arithmetic
contact equation
\[
 \boxed{
 a_{\chi,s}(x)
 =2\Re\sum_{n\ge2}
 \frac{\Lambda(n)\chi(n)}{\sqrt n}
 e^{-s(\log n)^2}e^{ix\log n},
 \qquad
 x\in\{\alpha_j(s),\beta_j(s)\}.}
 \tag{V26.31}
\]
At a simple endpoint \(x=x(s)\), implicit differentiation and the heat
equation also give the velocity law
\[
 \boxed{
 x'(s)
 =-\frac{\partial_sm_{\chi,s}(x(s))}
          {\partial_\xi m_{\chi,s}(x(s))}
 =-\frac{\partial_\xi^2m_{\chi,s}(x(s))}
          {\partial_\xi m_{\chi,s}(x(s))}.}
 \tag{V26.32}
\]
Thus the Fourier coefficients in \textup{(V26.26)} are not freely chosen:
their endpoints lie on, and move with, the zero level of the same completed
prime sum whose derivative generates the Gram.

\begin{theorem}[Entropy production is an exact but not weaker GRH target]
\label{thm:v26-entropy-no-shortcut}
The following condition is equivalent to GRH for the paired channel:
for every nondegenerate compact interval \(I\),
\[
 \limsup_{t\downarrow0}2t
 \log\bigl(1+\mathcal D_{\chi,I}(t)\bigr)=0.
 \tag{V26.33}
\]
Consequently neither the entropy identity nor positive semidefiniteness of the
adaptive Gram alone supplies an intermediate estimate.  Any advance through
\textup{(V26.23)} must use the contact and velocity constraints
\textup{(V26.31)}--\textup{(V26.32)} to control its all-ones contraction.
\end{theorem}

\begin{proof}
Under GRH the negative set is empty, so every local dissipation vanishes.
Conversely, Theorem~\ref{thm:v26-production-rate} identifies the limsup in
\textup{(V26.33)} with \(D_\chi(I)\).  If it vanishes for every interval,
Theorem~\ref{thm:v24-exact-rate} gives GRH.
\end{proof}

\begin{firewall}[The adaptive Gram has not been bounded]
Equations \textup{(V26.22)}--\textup{(V26.26)} reorganize the one-sided
problem into a positive bilinear object without discarding arithmetic phases.
They do not prove that its all-ones contraction has subexponential rate.
The endpoint identities \textup{(V26.31)}--\textup{(V26.32)} are exact, but no
inequality extracted here makes them incompatible with an off-line zero.
Accordingly GRH remains open.
\end{firewall}

\begin{assessment}[V26 acquisition]
V26 establishes the exact negative-entropy law
\(\mathcal S_\chi(t)=2\int_t^\infty\mathcal D_\chi(s)ds\), proves that local
gradient production has the same max-plus exponent as the negative part, and
expresses the production as an adaptive completed prime-cell Gram matrix.
Its prime block splits exactly into Toeplitz difference-frequency and Hankel
sum-frequency terms evaluated at the Fourier transform of the negative set.
Generic endpoint-count bounds fail at the exponentially fine prime-log
spacing, so the only new leverage is the self-consistent free boundary:
endpoints satisfy the completed contact equation and move by
\(-m_{\xi\xi}/m_\xi\).  V27 should analyze these endpoint constraints as a
coupled interpolation problem, rather than bounding the Gram entries
individually.
\end{assessment}

\section*{Version V26 append-only record}
\addcontentsline{toc}{section}{Version V26 append-only record}
The complete V25 source was retained before this continuation.  It had
\(388873\) bytes, \(11562\) lines, and SHA--256
\[
 \texttt{9EBC99E4A51C7D32F2FD7766AD15B17628307542FBFBFF1AD759EB47D244B23E}.
\]
V26 derived the exact negative-entropy dissipation law, the local production
rate, the adaptive completed prime-cell Gram representation, and the
free-boundary contact and velocity constraints.

% =============================================================================
% END APPENDED V26 MATERIAL
% =============================================================================
% =============================================================================
% BEGIN APPENDED V27 MATERIAL
% =============================================================================

\section{V27: boundary flux and zero-level reconstruction}
\label{sec:grh-v27}

V26 converted quadratic negative entropy into an adaptive prime-cell Gram and
isolated the zero level of the complete heat symbol as the only source of new
structure.  This continuation studies that free boundary directly.  The
linear negative mass has an exact dissipation law supported only at boundary
slopes.  More sharply, the largest slope at a zero-level contact on a fixed
frequency interval has exactly the V24 max-plus exponent \(D_\chi(I)\).

An exposed off-line pair produces two simple moving contacts around each
phase-tuned negative lobe.  Their spacing recovers the horizontal displacement
of the pair, their velocities recover its ordinate, and their slopes recover
its Gaussian score.  Thus the coupled endpoint interpolation problem is
GRH-complete in a very economical sense: it samples only the value, first
derivative, and heat velocity of the completed prime reader at its own
zero-level contacts.

\subsection{Linear negative mass dissipates through boundary slopes}

Retain the negative set \(\Omega_{\chi,t}\) and negative part
\(w_{\chi,t}\) from \textup{(V26.1)}.  Define
\[
 \mathcal A_\chi(t)
 :=\int_{\mathbb R}w_{\chi,t}(\xi)\,d\xi
 \tag{V27.1}
\]
and the total boundary flux
\[
 \mathcal B_\chi(t)
 :=\sum_{x\in\partial\Omega_{\chi,t}}
 |\partial_\xi m_{\chi,t}(x)|.
 \tag{V27.2}
\]
The sum is finite.  A degenerate endpoint contributes zero.

\begin{theorem}[Exact negative-mass boundary flux]
\label{thm:v27-boundary-flux}
The function \(\mathcal A_\chi\) is locally absolutely continuous on
\((0,\infty)\), and for almost every \(t>0\),
\[
 \boxed{
 \mathcal A_\chi'(t)=-\mathcal B_\chi(t).}
 \tag{V27.3}
\]
With the eventual positivity time \(T_\chi\) from V26,
\[
 \boxed{
 \mathcal A_\chi(t)
 =\int_t^{T_\chi}\mathcal B_\chi(s)\,ds
 =\int_t^\infty\mathcal B_\chi(s)\,ds.}
 \tag{V27.4}
\]
\end{theorem}

\begin{proof}
The same compactness and chain-rule argument used in
Theorem~\ref{thm:v26-entropy} applies to the Lipschitz convex function
\(r\mapsto[-r]_+\).  At a time for which the derivative exists,
\[
 \mathcal A_\chi'(t)
 =-\int_{\Omega_{\chi,t}}\partial_t m_{\chi,t}(\xi)\,d\xi
 =-\int_{\Omega_{\chi,t}}\partial_\xi^2m_{\chi,t}(\xi)\,d\xi.
 \tag{V27.5}
\]
Write the negative set as
\(\bigcup_{j=1}^{J(t)}(\alpha_j(t),\beta_j(t))\).  Since the symbol is
nonnegative immediately outside each component and negative inside,
\[
 \partial_\xi m_{\chi,t}(\alpha_j(t))\le0,
 \qquad
 \partial_\xi m_{\chi,t}(\beta_j(t))\ge0.
 \tag{V27.6}
\]
Integration of \textup{(V27.5)} on every component therefore gives
\[
 \mathcal A_\chi'(t)
 =-\sum_{j=1}^{J(t)}
 \left(
 |\partial_\xi m_{\chi,t}(\alpha_j(t))|
 +|\partial_\xi m_{\chi,t}(\beta_j(t))|
 \right),
 \tag{V27.7}
\]
which is \textup{(V27.3)}.  Degenerate endpoints have zero first derivative
and cause no boundary contribution.  Eventual positivity gives
\(\mathcal A_\chi(T_\chi)=0\); integration proves \textup{(V27.4)}.
\end{proof}

\begin{corollary}[Boundary-flux rate criterion]
\label{cor:v27-flux-grh}
GRH holds for the paired channel if and only if
\[
 \boxed{
 \limsup_{t\downarrow0}4t
 \log\left(
 1+\int_t^\infty\mathcal B_\chi(s)\,ds
 \right)=0.}
 \tag{V27.8}
\]
If GRH fails, the limsup is strictly positive.
\end{corollary}

\begin{proof}
By \textup{(V27.4)}, the expression is
\(4t\log(1+\mathcal A_\chi(t))\).  Under GRH the negative part vanishes.
If GRH fails, choose a compact influence interval with \(D_\chi(I)>0\).
Theorem~\ref{thm:v25-moment-rate} with \(q=1\) gives
\[
 \limsup_{t\downarrow0}4t\log\left(
 1+\int_I[-m_{\chi,t}(\xi)]_+\,d\xi
 \right)=D_\chi(I)>0,
 \tag{V27.9}
\]
and the global negative mass is at least this local mass.
\end{proof}

\subsection{The completed prime reader at a boundary contact}

For \(r=0,1,2\), put
\[
 P_{\chi,t}^{[r]}(x)
 :=\sum_{n\ge2}
 \frac{\Lambda(n)\chi(n)(\log n)^r}{\sqrt n}
 e^{-t(\log n)^2}e^{ix\log n}.
 \tag{V27.10}
\]
At every positive time these series converge absolutely.  The prime-side
definition of the symbol and its first two spatial derivatives are
\[
 \begin{aligned}
 m_{\chi,t}(x)
 &=a_{\chi,t}(x)-2\Re P_{\chi,t}^{[0]}(x),\\
 \partial_xm_{\chi,t}(x)
 &=\partial_xa_{\chi,t}(x)+2\Im P_{\chi,t}^{[1]}(x),\\
 \partial_x^2m_{\chi,t}(x)
 &=\partial_x^2a_{\chi,t}(x)+2\Re P_{\chi,t}^{[2]}(x).
 \end{aligned}
 \tag{V27.11}
\]
Consequently every simple boundary branch \(x=x(t)\) obeys the coupled
interpolation system
\[
 \boxed{
 \begin{aligned}
 2\Re P_{\chi,t}^{[0]}(x(t))
 &=a_{\chi,t}(x(t)),\\
 \partial_xm_{\chi,t}(x(t))
 &=\partial_xa_{\chi,t}(x(t))
   +2\Im P_{\chi,t}^{[1]}(x(t)),\\
 x'(t)
 &=-\frac{\partial_x^2a_{\chi,t}(x(t))
          +2\Re P_{\chi,t}^{[2]}(x(t))}
         {\partial_xa_{\chi,t}(x(t))
          +2\Im P_{\chi,t}^{[1]}(x(t))}.
 \end{aligned}}
 \tag{V27.12}
\]
This is the contact equation and velocity law of V26 written as three
successive Gaussian prime readers.  In particular,
\[
 \boxed{
 \mathcal B_\chi(t)
 =\sum_{x\in\partial\Omega_{\chi,t}}
 \left|
 \partial_xa_{\chi,t}(x)+2\Im P_{\chi,t}^{[1]}(x)
 \right|.}
 \tag{V27.13}
\]
No gamma--prime cancellation has been discarded in this formula.

\subsection{Actual endpoints generated by an exposed off-line pair}

Let \(p_*=(\gamma_*,\delta_*,m_{p_*})\) be uniquely exposed at \(x_0\), with
\[
 E_*:=S_{p_*}(x_0)>0
 \tag{V27.14}
\]
and exposure gap \(\eta>0\).  Use the phase-tuned sequences
\(t_k,\xi_k\) from \textup{(V23.11)}--\textup{(V23.13)}, so that the
dominant cosine at \(\xi_k\) is \(-1\).  Its two nominal adjacent zero
locations are
\[
 y_{k,\pm}:=\xi_k\pm\frac{\pi t_k}{\delta_*}.
 \tag{V27.15}
\]

\begin{theorem}[Exposed-pair free-boundary asymptotics]
\label{thm:v27-endpoint-asymptotics}
There are \(c>0\) and, for all sufficiently large \(k\), two simple zeros
\[
 x_{k,-}<\xi_k<x_{k,+}
 \tag{V27.16}
\]
of \(m_{\chi,t_k}\) such that
\[
 \boxed{
 x_{k,\pm}
 =y_{k,\pm}+O\!\left(t_ke^{-c/t_k}\right).}
 \tag{V27.17}
\]
The interval \((x_{k,-},x_{k,+})\) is a component of
\(\Omega_{\chi,t_k}\), and
\[
 \boxed{
 x_{k,+}-x_{k,-}
 =\frac{2\pi t_k}{\delta_*}
 +O\!\left(t_ke^{-c/t_k}\right).}
 \tag{V27.18}
\]
Its endpoint slopes satisfy
\[
 \boxed{
 \lim_{k\to\infty}4t_k
 \log|\partial_\xi m_{\chi,t_k}(x_{k,\pm})|
 =E_*.}
 \tag{V27.19}
\]
Each simple zero extends to a local analytic branch \(x_{k,\pm}(t)\).
At \(t=t_k\),
\[
 \boxed{
 x_{k,\pm}'(t_k)
 =\frac{x_{k,\pm}-\gamma_*}{t_k}
 +O\!\left(t_k^{-M}e^{-c/t_k}\right)}
 \tag{V27.20}
\]
for some fixed \(M<\infty\).
\end{theorem}

\begin{proof}
Write the full symbol near \(x_0\) as
\[
 m_{\chi,t}=K_{p_*,t}+R_{p_*,t},
 \tag{V27.21}
\]
where \(K_{p_*,t}\) is the exposed-pair contribution
\textup{(V25.26)}.  The exposure gap, Gaussian shell count, and termwise
differentiation give constants \(c>0\) and \(M_j<\infty\) such that, for
\(j=0,1,2\),
\[
 |\partial_\xi^jR_{p_*,t}(\xi)|
 \le t^{-M_j}e^{-c/t}
 \frac{4\pi m_{p_*}}{\sqrt{4\pi t}}
 e^{S_{p_*}(\xi)/(4t)}
 \tag{V27.22}
\]
uniformly on a fixed exposure neighbourhood.  Polynomial factors have been
absorbed by decreasing \(c\).

At \(y_{k,\pm}\), the dominant cosine vanishes and
\[
 |\partial_\xi K_{p_*,t_k}(y_{k,\pm})|
 =\frac{\delta_*}{2t_k}
 \frac{4\pi m_{p_*}}{\sqrt{4\pi t_k}}
 e^{S_{p_*}(y_{k,\pm})/(4t_k)}.
 \tag{V27.23}
\]
The quantitative implicit-function theorem applied to
\textup{(V27.21)}--\textup{(V27.23)} produces one unique simple root within
\(O(t_ke^{-c/t_k})\) of each nominal root.  At fixed phase distance from the
nominal roots, the exposed pair dominates with positive sign outside and
negative sign between them.  Hence these roots bound one negative component,
proving \textup{(V27.17)}--\textup{(V27.18)}.

Because \(y_{k,\pm}=x_0+O(t_k)\) and the root corrections are exponentially
smaller,
\[
 S_{p_*}(x_{k,\pm})=E_*+O(t_k).
 \tag{V27.24}
\]
Equations \textup{(V27.22)}--\textup{(V27.23)} show that the full endpoint
slope equals the dominant endpoint slope up to relative
\(O(t_k^{-M}e^{-c/t_k})\).  Taking logarithms proves
\textup{(V27.19)}.

At a zero of the pure pair, direct differentiation of
\textup{(V25.26)} gives
\[
 -\frac{\partial_\xi^2K_{p_*,t}(x)}
         {\partial_\xi K_{p_*,t}(x)}
 =\frac{x-\gamma_*}{t}.
 \tag{V27.25}
\]
The actual zeros are simple, so the implicit-function theorem and the heat
equation give
\(x'=-m_{\xi\xi}/m_\xi\).  Equations
\textup{(V27.17)}, \textup{(V27.22)}, and \textup{(V27.25)} then yield
\textup{(V27.20)}, after increasing \(M\) if necessary.
\end{proof}

\subsection{Tomography from the moving zero level}

The preceding theorem makes the free-boundary data invertible.

\begin{corollary}[Endpoint reconstruction of the exposed zero pair]
\label{cor:v27-endpoint-reconstruction}
Under Theorem~\ref{thm:v27-endpoint-asymptotics},
\[
 \boxed{
 \delta_*
 =\lim_{k\to\infty}
 \frac{2\pi t_k}{x_{k,+}-x_{k,-}},}
 \tag{V27.26}
\]
\[
 \boxed{
 \gamma_*
 =\lim_{k\to\infty}
 \left(
 x_{k,\pm}-t_kx_{k,\pm}'(t_k)
 \right),}
 \tag{V27.27}
\]
and
\[
 \boxed{
 E_*
 =\lim_{k\to\infty}4t_k
 \log|\partial_\xi m_{\chi,t_k}(x_{k,\pm})|.}
 \tag{V27.28}
\]
Thus adjacent endpoint spacing, one endpoint velocity, and one endpoint slope
recover respectively the horizontal displacement, ordinate, and visible
Gaussian score of the reflected zero pair.
\end{corollary}

\begin{proof}
Equations \textup{(V27.26)} and \textup{(V27.28)} are immediate from
\textup{(V27.18)} and \textup{(V27.19)}.  Multiplying the error term in
\textup{(V27.20)} by \(t_k\) makes it tend to zero, and
\textup{(V27.27)} follows.
\end{proof}

For a nondegenerate compact interval \(I\), define the maximal contact slope
\[
 \mathcal C_{\chi,I}(t)
 :=\sup\left\{
 |\partial_\xi m_{\chi,t}(x)|:
 x\in I,\ m_{\chi,t}(x)=0
 \right\},
 \tag{V27.29}
\]
with value zero when the contact set is empty.

\begin{theorem}[Exact zero-level slope rate]
\label{thm:v27-contact-rate}
For every nondegenerate compact interval \(I\),
\[
 \boxed{
 \limsup_{t\downarrow0}4t
 \log\bigl(1+\mathcal C_{\chi,I}(t)\bigr)
 =D_\chi(I).}
 \tag{V27.30}
\]
Consequently GRH holds for the paired channel if and only if the left side of
\textup{(V27.30)} vanishes for every such interval.
\end{theorem}

\begin{proof}
The differentiated upper bound \textup{(V26.10)} gives the rate at most
\(D_\chi(I)\).  If \(D_\chi(I)=0\), nonnegativity gives equality.

Suppose \(D_\chi(I)>0\) and fix \(\varepsilon>0\).  By
Lemma~\ref{lem:v24-local-exposure}, choose an exposed point in the interior of
\(I\) with score larger than \(D_\chi(I)-\varepsilon\).  The two actual
endpoints in Theorem~\ref{thm:v27-endpoint-asymptotics} then lie in \(I\) for
large \(k\), and \textup{(V27.19)} makes the limsup in
\textup{(V27.30)} at least \(D_\chi(I)-\varepsilon\).  Let
\(\varepsilon\downarrow0\).  The final equivalence follows from
Theorem~\ref{thm:v24-exact-rate}.
\end{proof}

Equation \textup{(V27.30)} is a strict reduction in sampling dimension: the
negative interval need not be integrated or maximized over.  It is enough to
control the derivative of the completed prime reader at points where its
value equals the gamma completion.  The exponential difficulty, however,
does not decrease.

\begin{criterion}[A discrete boundary target equivalent to GRH]
\label{crit:v27-boundary-target}
It suffices, and is necessary, to prove for every compact interval \(I\) and
every \(\varepsilon>0\) that
\[
 \left|
 \partial_xa_{\chi,t}(x)+2\Im P_{\chi,t}^{[1]}(x)
 \right|
 \le C_{\chi,I,\varepsilon}e^{\varepsilon/t}
 \tag{V27.31}
\]
whenever \(x\in I\), \(0<t\le1\), and
\[
 2\Re P_{\chi,t}^{[0]}(x)=a_{\chi,t}(x).
 \tag{V27.32}
\]
\end{criterion}

\begin{proof}
Equations \textup{(V27.11)} and \textup{(V27.32)} say precisely that \(x\)
is a zero-level contact, while the left side of \textup{(V27.31)} is its
slope.  Condition \textup{(V27.31)} gives zero rate in
\textup{(V27.30)}; conversely GRH makes the positive-time symbol strictly
positive by \textup{(V22.26)}, so there are no contacts and the condition is
vacuous.
\end{proof}

\subsection{What the endpoint equations do and do not exclude}

The endpoint system is more rigid than the adaptive set estimate
\textup{(V26.29)}, but it is not algebraically contradictory under failure of
GRH.  Theorem~\ref{thm:v27-endpoint-asymptotics} constructs contacts which
satisfy \textup{(V27.12)} and whose data converge to an off-line zero pair.
Thus the universal heat equation, contact equation, and implicit velocity law
are compatible with the forbidden geometry.  New input must be arithmetic:
it must bound the derivative reader conditionally on the value constraint
\textup{(V27.32)}, not merely relate successive derivatives of an arbitrary
caloric function.

Taking absolute values in \(P_{\chi,t}^{[1]}\) gives an even larger version
of the V23 full-strip saddle and cannot prove \textup{(V27.31)}.  Likewise,
Bernstein inequalities based only on the effective frequency
\(\log n\asymp1/t\) cost polynomial powers of \(t^{-1}\) relative to the
amplitude but do not reduce its exponential type.  The required saving must
come from the character phases under the contact constraint.

\begin{firewall}[Boundary reconstruction is not boundary exclusion]
V27 proves that a forbidden zero pair would be completely visible in the
moving zero-level geometry.  It does not prove the conditional derivative
bound \textup{(V27.31)}.  The constructed endpoint branches demonstrate that
contact and velocity identities alone do not yield a contradiction.  No
off-line zero has been excluded, and GRH remains open.
\end{firewall}

\begin{assessment}[V27 acquisition]
V27 derives the exact linear negative-mass flux
\(\mathcal A_\chi(t)=\int_t^\infty\mathcal B_\chi(s)ds\), supported only on
zero-level slopes.  It proves that an exposed off-line pair creates an actual
negative component whose endpoint spacing is \(2\pi t/\delta_*\), whose
endpoint slopes have exponent \(E_*/(4t)\), and whose velocities recover
\(\gamma_*\).  The maximal contact slope on any fixed interval has exact rate
\(D_\chi(I)\), yielding the discrete conditional prime-reader criterion
\textup{(V27.31)}--\textup{(V27.32)}.  The next attack should seek a genuinely
arithmetic value--derivative inequality for these two correlated Gaussian
prime readers, rather than another universal heat-flow identity.
\end{assessment}

\section*{Version V27 append-only record}
\addcontentsline{toc}{section}{Version V27 append-only record}
The complete V26 source was retained before this continuation.  It had
\(405830\) bytes, \(12072\) lines, and SHA--256
\[
 \texttt{01A08B710578E05328DF9D0CB2A6BE9A908B844FE2D06A7FADC3F6316288D3BA}.
\]
V27 established the negative-mass boundary flux, exposed-pair endpoint
asymptotics and reconstruction, the exact zero-level slope rate, and the
conditional value--derivative prime-reader criterion.

% =============================================================================
% END APPENDED V27 MATERIAL
% =============================================================================
% =============================================================================
% BEGIN APPENDED V28 MATERIAL
% =============================================================================

\section{V28: saddle-normalized contact covariance and its sharp capacity}
\label{sec:grh-v28}

V27 reduced the local obstruction to a conditional value--derivative problem:
bound the first Gaussian prime reader at points where the zeroth reader equals
the completed gamma term.  This continuation tests the strongest estimate
available from saddle centering, unit-modulus phases, and Cauchy--Schwarz
alone.

The absolute prime weights define a probability distribution on logarithmic
prime powers.  Its mean is the full-strip saddle \(1/(4t)\), and its standard
deviation is of order \(t^{-1/2}\).  The derivative reader splits exactly
into a carrier term and a centered phase covariance.  The contact condition
fixes only the real part of the mean phasor.  Optimizing over its imaginary
part gives a sharp covariance envelope whose exponential rate is still the
full strip value \(1/4\).  Even imposing a real mean phasor merely replaces
the carrier scale \(t^{-1}\) by the fluctuation scale \(t^{-1/2}\), a
polynomial gain with no effect on exponential type.

\subsection{The absolute saddle probability space}

For a fixed primitive nonprincipal character \(\chi\), define
\[
 Z_\chi(t)
 :=\sum_{n\ge2}
 \frac{\Lambda(n)|\chi(n)|}{\sqrt n}
 e^{-t(\log n)^2}.
 \tag{V28.1}
\]
Terms with \(\chi(n)=0\) are omitted.  On the remaining prime powers, put
\[
 \nu_{\chi,t}(n)
 :=\frac1{Z_\chi(t)}
 \frac{\Lambda(n)}{\sqrt n}e^{-t(\log n)^2},
 \qquad
 u_n:=\log n.
 \tag{V28.2}
\]
Thus \(\nu_{\chi,t}\) is a probability measure.  Write
\[
 \mu_\chi(t):=\mathbb E_{\nu_{\chi,t}}u,
 \qquad
 \sigma_\chi(t)^2
 :=\mathbb E_{\nu_{\chi,t}}(u-\mu_\chi(t))^2.
 \tag{V28.3}
\]

\begin{proposition}[Full-strip saddle moments]
\label{prop:v28-saddle-moments}
As \(t\downarrow0\),
\[
 \boxed{
 Z_\chi(t)
 =(1+o_\chi(1))
 \sqrt{\frac{\pi}{t}}
 \exp\!\left(\frac1{16t}\right),}
 \tag{V28.4}
\]
\[
 \boxed{
 \mu_\chi(t)=\frac1{4t}+o_\chi(t^{-1/2}),\qquad
 \sigma_\chi(t)^2=\frac1{2t}(1+o_\chi(1)).}
 \tag{V28.5}
\]
More generally, for every fixed integer \(r\ge1\),
\[
 \mathbb E_{\nu_{\chi,t}}
 |u-\mu_\chi(t)|^r
 \asymp_{\chi,r}t^{-r/2}.
 \tag{V28.6}
\]
\end{proposition}

\begin{proof}
The condition \(|\chi(n)|=0\) removes only powers of the finitely many primes
dividing the conductor.  Their total contribution to \textup{(V28.1)} and
every fixed logarithmic moment is bounded by a power of \(t^{-1}\), hence is
negligible compared with the claimed exponential scale.

Partial summation with the prime number theorem reduces the \(r\)-th raw
moment to
\[
 (1+o_\chi(1))
 \int_0^\infty u^r e^{u/2-tu^2}\,du.
 \tag{V28.7}
\]
The ordinary prime-number-theorem error is smaller at the saddle by
\(e^{-c_\chi/\sqrt t}\), as in Lemma~\ref{lem:v24-gaussian-pnt}.  Completing
the square gives
\[
 \frac u2-tu^2
 =\frac1{16t}-t\left(u-\frac1{4t}\right)^2.
 \tag{V28.8}
\]
The lower endpoint is exponentially far from the Gaussian center in standard
deviation units.  The zeroth integral is therefore
\((1+o(1))\sqrt{\pi/t}\,e^{1/(16t)}\), proving
\textup{(V28.4)}.  The centered Gaussian has mean \(1/(4t)\), variance
\(1/(2t)\), and fixed absolute centered moments comparable to \(t^{-r/2}\).
The PNT error is exponentially smaller than each such polynomial moment.
This proves \textup{(V28.5)}--\textup{(V28.6)}.
\end{proof}

\subsection{Mean phasor and the exact carrier--covariance split}

For \(\chi(n)\ne0\), define the unit phasor
\[
 \zeta_{n,x}
 :=\frac{\chi(n)}{|\chi(n)|}e^{ixu_n},
 \qquad |\zeta_{n,x}|=1,
 \tag{V28.9}
\]
and its mean
\[
 z_{\chi,t}(x)
 :=\mathbb E_{\nu_{\chi,t}}\zeta_{n,x}.
 \tag{V28.10}
\]
The zeroth and first readers from \textup{(V27.10)} become
\[
 P_{\chi,t}^{[0]}(x)
 =Z_\chi(t)z_{\chi,t}(x),
 \tag{V28.11}
\]
\[
 P_{\chi,t}^{[1]}(x)
 =Z_\chi(t)\left(
 \mu_\chi(t)z_{\chi,t}(x)+C_{\chi,t}(x)
 \right),
 \tag{V28.12}
\]
where the centered phase covariance is
\[
 C_{\chi,t}(x)
 :=\mathbb E_{\nu_{\chi,t}}
 \left[(u-\mu_\chi(t))\zeta_{n,x}\right].
 \tag{V28.13}
\]

\begin{lemma}[Sharp unit-phasor covariance bound]
\label{lem:v28-phasor-covariance}
For every \(t>0\) and \(x\in\mathbb R\),
\[
 \boxed{
 |C_{\chi,t}(x)|^2
 \le\sigma_\chi(t)^2
 \left(1-|z_{\chi,t}(x)|^2\right).}
 \tag{V28.14}
\]
\end{lemma}

\begin{proof}
Since \(\mathbb E(u-\mu_\chi)=0\),
\[
 C_{\chi,t}
 =\mathbb E\left[
 (u-\mu_\chi)(\zeta-z_{\chi,t})
 \right].
 \tag{V28.15}
\]
Cauchy--Schwarz gives the product of
\(\mathbb E(u-\mu_\chi)^2=\sigma_\chi^2\) and
\[
 \mathbb E|\zeta-z_{\chi,t}|^2
 =\mathbb E|\zeta|^2-|z_{\chi,t}|^2
 =1-|z_{\chi,t}|^2,
 \tag{V28.16}
\]
which proves the result.
\end{proof}

At a zero-level contact, \textup{(V27.12)} and \textup{(V28.11)} impose
\[
 \Re z_{\chi,t}(x)
 =r_{\chi,t}(x)
 :=\frac{a_{\chi,t}(x)}{2Z_\chi(t)}.
 \tag{V28.17}
\]
On every fixed compact interval, \(a_{\chi,t}\) remains bounded as
\(t\downarrow0\), while \(Z_\chi(t)\) grows as in \textup{(V28.4)}.  Hence
\[
 \boxed{
 |r_{\chi,t}(x)|
 \le\exp\!\left(
 -\frac1{16t}+O_{\chi,I}(\log(1/t))
 \right)\longrightarrow0}
 \tag{V28.18}
\]
uniformly at contacts \(x\in I\).

\begin{theorem}[Optimal contact covariance envelope]
\label{thm:v28-contact-envelope}
At every zero-level contact \(m_{\chi,t}(x)=0\),
\[
 \boxed{
 \left|
 \partial_xm_{\chi,t}(x)-\partial_xa_{\chi,t}(x)
 \right|
 \le
 2Z_\chi(t)
 \sqrt{\mu_\chi(t)^2+\sigma_\chi(t)^2}\,
 \sqrt{1-r_{\chi,t}(x)^2}.}
 \tag{V28.19}
\]
Along every sequence of existing contacts in a fixed compact interval, the
right side has exact absolute
saddle rate
\[
 \boxed{
 \lim_{t\downarrow0}4t\log\left[
 1+2Z_\chi(t)
 \sqrt{\mu_\chi(t)^2+\sigma_\chi(t)^2}\,
 \sqrt{1-r_{\chi,t}(x)^2}
 \right]=\frac14,}
 \tag{V28.20}
\]
uniformly over existing contacts.
\end{theorem}

\begin{proof}
Write
\[
 z_{\chi,t}(x)=r+iy.
 \tag{V28.21}
\]
Equations \textup{(V27.11)} and \textup{(V28.12)} give
\[
 \frac{
 \partial_xm_{\chi,t}(x)-\partial_xa_{\chi,t}(x)}
 {2Z_\chi(t)}
 =\Im\left(\mu_\chi(t)(r+iy)+C_{\chi,t}(x)\right).
 \tag{V28.22}
\]
The covariance lemma bounds its absolute value by
\[
 \mu_\chi|y|
 +\sigma_\chi\sqrt{1-r^2-y^2}.
 \tag{V28.23}
\]
For \(|y|\le\sqrt{1-r^2}\), Cauchy--Schwarz in \(\mathbb R^2\) gives
\[
 \mu_\chi|y|
 +\sigma_\chi\sqrt{1-r^2-y^2}
 \le\sqrt{\mu_\chi^2+\sigma_\chi^2}\sqrt{1-r^2}.
 \tag{V28.24}
\]
This proves \textup{(V28.19)}.  Equations
\textup{(V28.4)}, \textup{(V28.5)}, and \textup{(V28.18)} show that every
factor after \(Z_\chi(t)\) is only polynomial in \(t^{-1}\), with the final
square root tending to one.  Multiplication of its logarithm by \(4t\)
therefore gives \(1/4\).
\end{proof}

\subsection{Why the optimized envelope is genuinely sharp}

The preceding optimization is not an artifact of two successive loose
inequalities.  The carrier term alone asymptotically reaches the envelope
under the information retained by the contact constraint.

\begin{proposition}[Asymptotic saturation under contact data alone]
\label{prop:v28-capacity}
Let \(\nu\) be any probability measure with mean \(\mu>0\) and variance
\(\sigma^2\), and fix \(r\in(-1,1)\).  The constant unit phasor
\[
 \zeta\equiv r+i\sqrt{1-r^2}
 \tag{V28.25}
\]
satisfies \(\Re\mathbb E\zeta=r\), has centered covariance zero, and obeys
\[
 \Im\mathbb E(u\zeta)=\mu\sqrt{1-r^2}.
 \tag{V28.26}
\]
For the saddle measures \(\nu_{\chi,t}\),
\[
 \frac{\mu_\chi(t)}
 {\sqrt{\mu_\chi(t)^2+\sigma_\chi(t)^2}}
 \longrightarrow1.
 \tag{V28.27}
\]
Thus the order and exponential type in \textup{(V28.19)} cannot be improved
using only unit modulus, the real contact value, and the first two absolute
saddle moments.
\end{proposition}

\begin{proof}
The first two statements follow immediately from constancy of \(\zeta\).
Equation \textup{(V28.26)} is
\(\Im(\mu\zeta)=\mu\sqrt{1-r^2}\).  Finally,
\textup{(V28.5)} gives
\[
 \frac{\sigma_\chi(t)^2}{\mu_\chi(t)^2}
 =8t(1+o_\chi(1)),
 \tag{V28.28}
\]
which proves \textup{(V28.27)}.
\end{proof}

This is a capacity statement, not a model of the actual character phases.
Its role is to identify precisely what an arithmetic proof must use beyond
the generic probability data: the multiplicative dependence of
\(\zeta_{n,x}\) on \(n\).

Even an additional reality condition on the mean phasor does not change the
exponential obstruction.

\begin{corollary}[Centering without an imaginary carrier is still critical]
\label{cor:v28-centered-barrier}
Suppose, in addition to contact, that
\(\Im z_{\chi,t}(x)=0\).  Then
\[
 \left|
 \partial_xm_{\chi,t}(x)-\partial_xa_{\chi,t}(x)
 \right|
 \le2Z_\chi(t)\sigma_\chi(t)
 \sqrt{1-r_{\chi,t}(x)^2}.
 \tag{V28.29}
\]
The right side still satisfies
\[
 \lim_{t\downarrow0}4t\log\left(
 1+2Z_\chi(t)\sigma_\chi(t)
 \sqrt{1-r_{\chi,t}(x)^2}
 \right)=\frac14.
 \tag{V28.30}
\]
Thus removal of the carrier gains a factor of order \(t^{1/2}\), but no
exponential saving.
\end{corollary}

\begin{proof}
With \(y=0\), equations \textup{(V28.14)} and \textup{(V28.22)} give
\textup{(V28.29)}.  Proposition~\ref{prop:v28-saddle-moments} proves
\textup{(V28.30)}.
\end{proof}

\subsection{Every fixed saddle jet has the same absolute exponent}

The failure is not confined to first-order centering.  For \(r\ge0\), define
\[
 Z_{\chi,r}^{\mathrm{cen}}(t)
 :=Z_\chi(t)
 \left(
 \mathbb E_{\nu_{\chi,t}}
 |u-\mu_\chi(t)|^{2r}
 \right)^{1/2}.
 \tag{V28.31}
\]

\begin{corollary}[Finite centered jets do not lower exponential type]
\label{cor:v28-finite-jets}
For every fixed integer \(r\ge0\),
\[
 \boxed{
 \lim_{t\downarrow0}4t
 \log\bigl(1+Z_{\chi,r}^{\mathrm{cen}}(t)\bigr)
 =\frac14.}
 \tag{V28.32}
\]
Consequently any fixed number of saddle centerings followed by an absolute
Cauchy--Schwarz estimate retains the full-strip exponent.
\end{corollary}

\begin{proof}
By \textup{(V28.4)} and \textup{(V28.6)},
\[
 Z_{\chi,r}^{\mathrm{cen}}(t)
 =\exp\!\left(\frac1{16t}+O_{\chi,r}(\log(1/t))\right).
 \tag{V28.33}
\]
Equation \textup{(V28.32)} follows.
\end{proof}

The signed fixed-character PNT estimate does not repair this contact bound.
Repeating the Gaussian Abel argument of Lemma~\ref{lem:v24-gaussian-pnt} with
a fixed factor \((\log n)^r\) gives, for compact \(I\),
\[
 \sup_{x\in I}|P_{\chi,t}^{[r]}(x)|
 \le C_{\chi,I,r}t^{-M_r}
 \exp\!\left(\frac1{16t}-\frac{c_\chi}{\sqrt t}\right).
 \tag{V28.34}
\]
The logarithmic rate after multiplication by \(4t\) is still \(1/4\).

\subsection{The exact missing phase-regression input}

At a contact, define
\[
 \mathcal R_{\chi,t}(x)
 :=\Im\left(
 \mu_\chi(t)z_{\chi,t}(x)+C_{\chi,t}(x)
 \right).
 \tag{V28.35}
\]
Then \textup{(V28.22)} is the exact regression identity
\[
 \boxed{
 \partial_xm_{\chi,t}(x)
 =\partial_xa_{\chi,t}(x)
 +2Z_\chi(t)\mathcal R_{\chi,t}(x).}
 \tag{V28.36}
\]
Combining this with Theorem~\ref{thm:v27-contact-rate} yields the following
minimal formulation.

\begin{criterion}[Conditional phase-regression criterion]
\label{crit:v28-phase-regression}
GRH holds for the paired channel if and only if, on every nondegenerate
compact interval \(I\),
\[
 \limsup_{t\downarrow0}4t\log\left(
 1+\sup_{\substack{x\in I\\
 \Re z_{\chi,t}(x)=a_{\chi,t}(x)/(2Z_\chi(t))}}
 Z_\chi(t)|\mathcal R_{\chi,t}(x)|
 \right)=0.
 \tag{V28.37}
\]
The supremum is zero if the contact set is empty.
\end{criterion}

\begin{proof}
On fixed \(I\), \(\partial_xa_{\chi,t}\) is bounded as \(t\downarrow0\).
Equations \textup{(V28.17)} and \textup{(V28.36)} therefore identify
\textup{(V28.37)} with the zero-level slope condition
\textup{(V27.30)}, up to bounded terms.  Apply
Theorem~\ref{thm:v27-contact-rate}.
\end{proof}

The generic covariance envelope gives only
\[
 |\mathcal R_{\chi,t}(x)|
 \le\sqrt{\mu_\chi(t)^2+\sigma_\chi(t)^2}
 \sqrt{1-r_{\chi,t}(x)^2},
 \tag{V28.38}
\]
whose multiplication by \(Z_\chi(t)\) has rate \(1/4\).  Criterion
\textup{(V28.37)} instead requires cancellation of the normalized regression
by the exponential factor supplied by \(Z_\chi(t)^{-1}\).  Such cancellation
cannot come from a finite absolute-moment hierarchy; it must use the signed
multiplicative phases themselves.

\begin{theorem}[Saddle covariance cannot prove the contact criterion]
\label{thm:v28-covariance-no-go}
Any proof of the V27 conditional value--derivative bound which uses the
character phases only through
\[
 |\zeta_{n,x}|=1,\qquad
 \Re\mathbb E_{\nu_{\chi,t}}\zeta_{n,x}
 =r_{\chi,t}(x),
 \tag{V28.39}
\]
and finitely many absolute centered moments of \(u\), cannot improve the
universal logarithmic rate \(1/4\).  In particular, saddle centering and
finite covariance iteration alone cannot establish
\textup{(V28.37)}.
\end{theorem}

\begin{proof}
The optimal first covariance envelope has rate \(1/4\) by
Theorem~\ref{thm:v28-contact-envelope} and is asymptotically saturated under
the retained data by Proposition~\ref{prop:v28-capacity}.  Every additional
fixed centered-moment Cauchy budget has the same rate by
Corollary~\ref{cor:v28-finite-jets}.  Finite sums and polynomial
coefficients do not change logarithmic \(1/t\) type.  Hence the listed data
cannot yield the zero rate required in \textup{(V28.37)}.
\end{proof}

\begin{firewall}[The multiplicative phase regression is not bounded]
V28 derives the exact normalized regression statistic and proves that its
generic covariance capacity saturates the full critical strip.  It does not
establish cancellation of \(\mathcal R_{\chi,t}\) at arithmetic contacts.
The constant-phasor saturation example is a capacity witness, not a claim
about Dirichlet-character values.  The needed multiplicative estimate remains
unproved, and GRH remains open.
\end{firewall}

\begin{assessment}[V28 acquisition]
V28 normalizes the contact problem by the absolute Gaussian saddle measure.
It proves precise mass, mean, variance, and centered-moment asymptotics;
splits the derivative reader into an imaginary carrier and a centered phase
covariance; and optimizes the strongest contact bound obtainable from unit
phasors and Cauchy--Schwarz.  That bound is asymptotically sharp under the
retained data and has full-strip rate \(1/4\), even after removal of the
carrier or any fixed number of centerings.  The residual target is the
conditional multiplicative phase regression \textup{(V28.37)}.  V29 should
go upstream into the periodic character structure---for example residue-class
or conductor-Fourier decomposition of the saddle measure---and determine
whether it supplies a nonabsolute regression identity not already equivalent
to the logarithmic derivative of \(L(s,\chi)\).
\end{assessment}

\section*{Version V28 append-only record}
\addcontentsline{toc}{section}{Version V28 append-only record}
The complete V27 source was retained before this continuation.  It had
\(421024\) bytes, \(12536\) lines, and SHA--256
\[
 \texttt{A179CC5928910CFE7FF0C71B4885CEF4B7EB79A6334D5C8A4524785DEFC6B2AA}.
\]
V28 established the saddle probability asymptotics, the optimal
carrier--covariance contact envelope, its sharp capacity barrier, the
finite-centered-jet no-go, and the conditional phase-regression criterion.

% =============================================================================
% END APPENDED V28 MATERIAL
% =============================================================================
% =============================================================================
% BEGIN APPENDED V29 MATERIAL
% =============================================================================

\section{V29: conductor Fourier decomposition of the contact regression}
\label{sec:grh-v29}

V28 proved that saddle centering and finite absolute covariance information
cannot control the phase regression at a zero-level contact.  It proposed
using the actual periodic and multiplicative structure of the Dirichlet
character.  This continuation carries out the complete fixed-conductor
decomposition.

The common prime-density profile cancels exactly because a nonprincipal
character sums to zero over the reduced residue classes.  What remains is the
character projection of the arithmetic-progression error vector.  A Gauss
sum gives the equivalent additive Fourier description.  These are related by
a fixed finite-dimensional transform, so conductor orthogonality changes
constants but cannot change logarithmic \(1/t\) type.  Classical PNT in
progressions leaves the same \(e^{1/(16t)-c/\sqrt t}\) error found in V24 and
V28.  The needed subexponential contact regression is therefore not supplied
by periodicity alone; it is precisely a signed estimate for the projected
progression errors.

This analysis does not duplicate V15--V17.  Those versions grouped residue
classes to build magnetic translation energy on a physical support interval.
V29 instead decomposes the scalar Gaussian value and derivative readers at
the free-boundary contacts of V27.

\subsection{Residue-class readers and exact main-term cancellation}

Let \(q\) be the conductor of \(\chi\), and let
\[
 G_q:=(\mathbb Z/q\mathbb Z)^\times.
 \tag{V29.1}
\]
For \(a\in G_q\) and an integer \(r\ge0\), define
\[
 P_{a,t}^{[r]}(x)
 :=\sum_{\substack{n\ge2\\n\equiv a\;(\mathrm{mod}\ q)}}
 \frac{\Lambda(n)(\log n)^r}{\sqrt n}
 e^{-t(\log n)^2}e^{ix\log n}.
 \tag{V29.2}
\]
The common continuous main reader is
\[
 M_{q,t}^{[r]}(x)
 :=\frac1{\varphi(q)}
 \int_1^\infty
 \frac{(\log y)^r}{\sqrt y}
 e^{-t(\log y)^2}e^{ix\log y}\,dy,
 \tag{V29.3}
\]
and the progression error reader is
\[
 E_{a,t}^{[r]}(x)
 :=P_{a,t}^{[r]}(x)-M_{q,t}^{[r]}(x).
 \tag{V29.4}
\]

\begin{theorem}[Exact residue-error projection]
\label{thm:v29-residue-projection}
For every \(r\ge0\), \(t>0\), and \(x\in\mathbb R\),
\[
 \boxed{
 P_{\chi,t}^{[r]}(x)
 =\sum_{a\in G_q}\chi(a)E_{a,t}^{[r]}(x).}
 \tag{V29.5}
\]
Thus the common saddle profile has no component in the nonprincipal
character direction.
\end{theorem}

\begin{proof}
Grouping the absolutely convergent series by residue class gives
\[
 P_{\chi,t}^{[r]}(x)
 =\sum_{a\in G_q}\chi(a)P_{a,t}^{[r]}(x).
 \tag{V29.6}
\]
Because \(\chi\) is nonprincipal,
\[
 \sum_{a\in G_q}\chi(a)=0.
 \tag{V29.7}
\]
Subtracting the same quantity \(M_{q,t}^{[r]}(x)\) from every residue
component therefore leaves the projection unchanged, proving
\textup{(V29.5)}.
\end{proof}

Define the error vectors
\[
 \mathbf E_t^{[r]}(x)
 :=\bigl(E_{a,t}^{[r]}(x)\bigr)_{a\in G_q},
 \qquad
 \boldsymbol\chi
 :=\bigl(\chi(a)\bigr)_{a\in G_q}.
 \tag{V29.8}
\]
With the bilinear pairing
\(\boldsymbol\chi\cdot\mathbf E=\sum_a\chi(a)E_a\),
\textup{(V29.5)} reads
\[
 P_{\chi,t}^{[r]}=\boldsymbol\chi\cdot\mathbf E_t^{[r]}.
 \tag{V29.9}
\]
The V27 contact and slope equations become
\[
 \boxed{
 2\Re\bigl(
 \boldsymbol\chi\cdot\mathbf E_t^{[0]}(x)
 \bigr)=a_{\chi,t}(x),}
 \tag{V29.10}
\]
\[
 \boxed{
 \partial_xm_{\chi,t}(x)
 =\partial_xa_{\chi,t}(x)
 +2\Im\bigl(
 \boldsymbol\chi\cdot\mathbf E_t^{[1]}(x)
 \bigr).}
 \tag{V29.11}
\]
These identities hold at every \(x\), with \textup{(V29.10)} selecting the
zero-level contacts.

\subsection{The Gaussian PNT-in-progressions error}

For \(a\in G_q\), write
\[
 \Psi(y;q,a)
 :=\sum_{\substack{n\le y\\n\equiv a\;(\mathrm{mod}\ q)}}\Lambda(n),
 \qquad
 \Delta_a(y)
 :=\Psi(y;q,a)-\frac {y-1}{\varphi(q)}.
 \tag{V29.12}
\]
For fixed \(q\), the classical PNT in arithmetic progressions gives constants
\(C_q,c_q>0\) such that
\[
 |\Delta_a(y)|
 \le C_qy e^{-c_q\sqrt{\log y}}
 \qquad(y\ge2,\ a\in G_q).
 \tag{V29.13}
\]
The constants absorb any fixed exceptional real zero: its power \(y^\beta\),
with \(\beta<1\), is eventually smaller than the right side.

\begin{proposition}[Gaussian progression-error bound]
\label{prop:v29-ap-gaussian}
For every fixed integer \(r\ge0\) and compact interval \(I\), there are
\(C_{q,I,r},c_q'>0\), \(M_r<\infty\), and \(t_0>0\) such that
\[
 \boxed{
 \max_{a\in G_q}\sup_{x\in I}
 |E_{a,t}^{[r]}(x)|
 \le C_{q,I,r}t^{-M_r}
 \exp\!\left(
 \frac1{16t}-\frac{c_q'}{\sqrt t}
 \right)}
 \tag{V29.14}
\]
for \(0<t<t_0\).
\end{proposition}

\begin{proof}
Let
\[
 f_{t,x,r}(y)
 :=y^{-1/2+ix}(\log y)^r e^{-t(\log y)^2}.
 \tag{V29.15}
\]
Stieltjes partial summation and the definition of the common main reader give
\[
 E_{a,t}^{[r]}(x)
 =-\int_{1^-}^\infty\Delta_a(y)\,df_{t,x,r}(y),
 \tag{V29.16}
\]
with vanishing endpoint terms.  After \(u=\log y\), differentiation of
\(f_{t,x,r}\) bounds the absolute value by a fixed power of \(t^{-1}\)
times an integral of the form
\[
 \int_0^\infty(1+u)^{r+1}
 \exp\!\left(
 \frac u2-tu^2-c_q\sqrt u
 \right)du.
 \tag{V29.17}
\]
The split at \(u=1/(8t)\) used in the proof of
Lemma~\ref{lem:v24-gaussian-pnt} gives the saving
\(e^{-c_q'/\sqrt t}\); all fixed powers of \(u\) contribute only powers of
\(t^{-1}\).  This proves \textup{(V29.14)} uniformly over the finite set
\(G_q\).
\end{proof}

Projecting the vector estimate in \textup{(V29.14)} gives
\[
 \sup_{x\in I}|P_{\chi,t}^{[r]}(x)|
 \le\varphi(q)C_{q,I,r}t^{-M_r}
 \exp\!\left(
 \frac1{16t}-\frac{c_q'}{\sqrt t}
 \right),
 \tag{V29.18}
\]
which recovers \textup{(V28.34)} in residue coordinates.  The common main
profile has canceled exactly, but the remaining absolute error still obeys
\[
 \limsup_{t\downarrow0}4t\log\left(
 1+\max_{a\in G_q}\sup_{x\in I}
 |E_{a,t}^{[r]}(x)|
 \right)\le\frac14.
 \tag{V29.19}
\]
The square-root saving in the exponent is invisible at logarithmic \(1/t\)
scale.

\subsection{The equivalent additive Fourier basis}

Let
\[
 \tau(\overline\chi)
 :=\sum_{h\;(\mathrm{mod}\ q)}
 \overline{\chi(h)}e^{2\pi ih/q}
 \tag{V29.20}
\]
be the Gauss sum.  Primitivity gives
\[
 \chi(n)
 =\frac1{\tau(\overline\chi)}
 \sum_{h\;(\mathrm{mod}\ q)}
 \overline{\chi(h)}e^{2\pi ihn/q}.
 \tag{V29.21}
\]
Define the coprime additive readers
\[
 A_{h,t}^{[r]}(x)
 :=\sum_{\substack{n\ge2\\ \gcd(n,q)=1}}
 \frac{\Lambda(n)(\log n)^r}{\sqrt n}
 e^{-t(\log n)^2}e^{ix\log n}e^{2\pi ihn/q}.
 \tag{V29.22}
\]

\begin{proposition}[Conductor Fourier equivalence]
\label{prop:v29-fourier-equivalence}
For every \(r,t,x\),
\[
 \boxed{
 P_{\chi,t}^{[r]}(x)
 =\frac1{\tau(\overline\chi)}
 \sum_{h\;(\mathrm{mod}\ q)}
 \overline{\chi(h)}A_{h,t}^{[r]}(x).}
 \tag{V29.23}
\]
Moreover,
\[
 A_{h,t}^{[r]}(x)
 =\sum_{a\in G_q}
 e^{2\pi iha/q}P_{a,t}^{[r]}(x),
 \tag{V29.24}
\]
so the residue and additive reader vectors are related by the fixed discrete
Fourier transform modulo \(q\).
\end{proposition}

\begin{proof}
Insert \textup{(V29.21)} into the absolutely convergent definition of
\(P_{\chi,t}^{[r]}\) and interchange the finite and infinite sums to obtain
\textup{(V29.23)}.  Grouping the coprime integers in
\textup{(V29.22)} by their reduced residue classes gives
\textup{(V29.24)}.
\end{proof}

Because \(q\) is fixed, the Fourier transform and its inverse have operator
norms bounded independently of \(t\).  Passing between the two bases can
therefore change constants and powers of \(q\), but not a logarithmic
\(1/t\) exponent.

\subsection{The projected error is exactly the character prime error}

At the unsmoothed counting level,
\[
 \sum_{a\in G_q}\chi(a)\Delta_a(y)
 =\sum_{n\le y}\Lambda(n)\chi(n)
 =\Psi_\chi(y),
 \tag{V29.25}
\]
because the main terms cancel by \textup{(V29.7)}.  Thus the progression
error vector does not contain a new independent source of cancellation: its
character projection is precisely the classical character-weighted Chebyshev
error.

\begin{theorem}[Square-root progression projection is GRH-complete]
\label{thm:v29-projected-error-grh}
For the fixed conjugate pair \(\chi,\overline\chi\), the family of estimates
\[
 \boxed{
 \sum_{a\in G_q}\chi(a)\Delta_a(y)
 =O_{\chi,\varepsilon}(y^{1/2+\varepsilon})
 \quad\text{for every }\varepsilon>0}
 \tag{V29.26}
\]
is equivalent to GRH for the pair, after imposing the conjugate estimate as
well.
\end{theorem}

\begin{proof}
By \textup{(V29.25)}, condition \textup{(V29.26)} is the standard
square-root prime-number-theorem estimate for \(\Psi_\chi(y)\).  Partial
summation gives holomorphic continuation of
\(-L'/L(s,\chi)\) to every half-plane
\(\Re s>1/2+\varepsilon\), so no zero can lie to the right of the critical
line.  Applying the same statement to \(\overline\chi\) and using the
functional equation excludes zeros to the left.  Conversely, GRH gives
\(\Psi_\chi(y)=O_\chi(y^{1/2}\log^2(2y))\), and likewise for the conjugate,
which implies \textup{(V29.26)}.
\end{proof}

This theorem is the unsmoothed counterpart of the V28 phase-regression
criterion.  Neither residue coordinates nor the additive Gauss basis weaken
the needed projected estimate.

\subsection{Finite-conductor capacity at a contact}

The contact equation \textup{(V29.10)} fixes one real scalar projection of
the error vector.  Fixed-dimensional geometry leaves its imaginary
projection unrestricted.  Indeed, for any real \(c\) and \(L>0\), the
abstract vector
\[
 V_a
 :=\frac{\overline{\chi(a)}}{\varphi(q)}(c+iL),
 \qquad a\in G_q,
 \tag{V29.27}
\]
satisfies
\[
 \boldsymbol\chi\cdot\mathbf V=c+iL,
 \qquad
 \max_a|V_a|
 =\frac{\sqrt{c^2+L^2}}{\varphi(q)}.
 \tag{V29.28}
\]
Thus the condition
\(\Re(\boldsymbol\chi\cdot\mathbf V)=c\) is compatible with an imaginary
projection of the full componentwise size.

\begin{theorem}[Fixed-conductor Fourier inequalities cannot close the contact]
\label{thm:v29-conductor-no-go}
Suppose an argument uses the periodic decomposition only through:
\begin{enumerate}[label=\textup{(\roman*)},leftmargin=3em]
\item the componentwise progression bounds \textup{(V29.14)};
\item norm or orthogonality inequalities for the fixed residue or additive
Fourier vectors;
\item the real contact constraint \textup{(V29.10)}.
\end{enumerate}
Then its universal upper bound for the derivative projection in
\textup{(V29.11)} retains logarithmic rate \(1/4\).  These data alone cannot
prove the zero-rate contact criterion \textup{(V27.31)}.
\end{theorem}

\begin{proof}
Every norm comparison or discrete Fourier transform in the fixed dimension
\(\varphi(q)\) changes \textup{(V29.14)} only by a \(t\)-independent
constant.  Hence the resulting derivative bound has rate at best \(1/4\).
The capacity vector \textup{(V29.27)}--\textup{(V29.28)} shows that the real
contact projection cannot force an additional saving in its imaginary
counterpart.  Therefore no universal inequality using only the listed data
can yield a subexponential derivative projection.
\end{proof}

The capacity vector is not asserted to be an actual progression error vector.
It proves that the missing input is not finite conductor geometry.  One must
use correlations between the value and derivative errors inherited from the
same prime sequence, beyond componentwise PNT in progressions.

\begin{criterion}[Residue-error contact regression]
\label{crit:v29-residue-regression}
GRH is equivalent to the following statement: for every nondegenerate compact
interval \(I\),
\[
 \limsup_{t\downarrow0}4t\log\left(
 1+\sup_{\substack{x\in I\\
 2\Re(\boldsymbol\chi\cdot\mathbf E_t^{[0]}(x))
 =a_{\chi,t}(x)}}
 \left|
 \Im\bigl(\boldsymbol\chi\cdot\mathbf E_t^{[1]}(x)\bigr)
 \right|
 \right)=0.
 \tag{V29.29}
\]
The supremum is defined to be zero when the contact set is empty.
\end{criterion}

\begin{proof}
At contacts, \textup{(V29.11)} differs from twice the quantity in
\textup{(V29.29)} by the bounded archimedean derivative.  Conversely, GRH
leaves no positive-time contacts, so the displayed supremum is zero.
Necessity and sufficiency follow from
Theorem~\ref{thm:v27-contact-rate}.
\end{proof}


\begin{firewall}[Periodicity removes the main term, not the zero signal]
V29 proves exact cancellation of the common residue-class saddle and gives
both conductor Fourier coordinate systems.  It does not improve the
character projection of the progression errors beyond
\(e^{1/(16t)-c/\sqrt t}\).  The fixed-dimensional capacity theorem shows why
orthogonality and contact alone cannot do so.  No off-line zero is excluded,
and GRH remains open.
\end{firewall}

\begin{assessment}[V29 acquisition]
V29 resolves the periodic-structure branch proposed in V28.  The common
prime-density profile cancels exactly in the nonprincipal character
direction, leaving the scalar value and derivative readers as projections of
the progression-error vector.  Residue and additive Gauss coordinates are
fixed finite Fourier transforms and therefore preserve logarithmic
\(1/t\) type.  Classical PNT in progressions gives only the subleading
\(e^{-c/\sqrt t}\) saving, while the projected square-root estimate is
already GRH-complete.  The next version is the mandatory V30 corpus sweep.
It should compare this contact-regression obstruction with all parallel notes
created since the V20 sweep, looking specifically for a growing-dimensional
or multiscale correlation identity that is not reducible to fixed-conductor
orthogonality or a classical logarithmic-derivative criterion.
\end{assessment}

\section*{Version V29 append-only record}
\addcontentsline{toc}{section}{Version V29 append-only record}
The complete V28 source was retained before this continuation.  It had
\(436338\) bytes, \(13044\) lines, and SHA--256
\[
 \texttt{ED67657F0E61FFAB4D99FD7A848E7040D921354C09EEB322E7CFFBA08B92CB84}.
\]
V29 established the residue-error projection, Gaussian PNT-in-progressions
bound, additive Gauss-sum equivalence, fixed-conductor capacity no-go, and the
residue-error contact-regression target.

% =============================================================================
% END APPENDED V29 MATERIAL
% =============================================================================
% =============================================================================
% BEGIN APPENDED V30 MATERIAL
% =============================================================================

\section{V30: compulsory corpus sweep and the growing-modulus martingale}
\label{sec:grh-v30}

V29 showed that fixed-conductor Fourier coordinates remove the common prime
density but do not lower the logarithmic type of the progression-error
projection.  The mandated ten-version sweep has now found a useful upstream
mechanism in the parallel corpus: exact conditional projections and their
Pythagorean martingale decomposition.  This section imports that mechanism,
but none of the parallel programmes' physical constants or conclusions.

Applied to the absolute Gaussian prime saddle, conditional projection along
the nested residue moduli
\[
 q,q^2,q^3,\ldots
\]
gives an exact infinite decomposition of the V28 phase regression.  For each
fixed depth, however, the prime number theorem in arithmetic progressions
shows that the projection captures asymptotically none of the logarithmic
saddle variance.  Away from height zero it also captures asymptotically none
of the phase energy.  Thus the unresolved conditional tail retains the full
critical Cauchy capacity.  A successful use of this decomposition must control
a signed tail at a depth growing with the saddle; fixed-depth martingale
energy is no stronger than the fixed-conductor geometry already closed in
V29.

\subsection{Timestamped V30 corpus ledger}

The V20 sweep froze \(29\) non-GRH \texttt{.tex} files.  Its canonical
UTF--8 manifest consisted of the rows
\[
 \texttt{filename|byte-count|SHA--256}\backslash\texttt{n}
\]
in ordinal filename order and had SHA--256
\[
 \texttt{3C744E6E540E721C4D6E9719133D53B733BBCBFC73644190FE3CE40509B63222}.
 \tag{V30.1}
\]
The live V30 cutoff was taken at
\[
 \texttt{2026-09-10T23:07:52.4684869Z}.
 \tag{V30.2}
\]
It contained \(31\) non-GRH \texttt{.tex} files and its manifest hash was
\[
 \boxed{
 \texttt{34EE166478D15E763BD319E568AB2F2156EF42037622526D7D716D38CD727BE6}.}
 \tag{V30.3}
\]

Exactly \(27\) V20 rows were byte-identical and were hash-gated rather than
reread.  The two former rolling endpoints had disappeared:
\begin{center}
\tiny
\begin{tabular}{p{0.56\textwidth}|r|p{0.28\textwidth}}
\textbf{former endpoint}&\textbf{bytes}&\textbf{SHA--256}\\ \hline
\path{Renewal_SM_Einstein_Continuum_Research_Notes_V93.tex}
 &1103350&\texttt{4F5EA5CA6367BD7E2ADDFBD71DB8B60A}\newline \texttt{85B9A4C696CE2F6D66A41EB3DAA784F2}\\
\path{Renewal_Yang_Mills_Mass_Gap_Research_Notes_V29.tex}
 &232987&\texttt{85AEAC9219611F19F1C2F98A2F92C42}\newline \texttt{53B7F32F7120F17183099CCA237B0F0B8}
\end{tabular}
\end{center}
The four new rows, including the simultaneously visible YM handoff pair,
were:
\begin{center}
\tiny
\begin{tabular}{p{0.54\textwidth}|r|p{0.29\textwidth}}
\textbf{new row}&\textbf{bytes}&\textbf{SHA--256}\\ \hline
\path{Renewal_SM_Einstein_Continuum_Research_Notes_V100_f14.tex}
 &1166228&\texttt{C6608FD46EE95038605B57FE09CC5849}\newline \texttt{46F88DF6D480C7A53443385B94E64BD9}\\
\path{Renewal_SM_Einstein_Continuum_Research_Notes_V21.tex}
 &203941&\texttt{7A4873C14ECDDBA469074EA0C3217859}\newline \texttt{D0FFD77E23C2DF667DC425548FCDF5B0}\\
\path{Renewal_Yang_Mills_Mass_Gap_Research_Notes_V35.tex}
 &289778&\texttt{4EA177BF18DD0ED14474AE54E22DE26F}\newline \texttt{F178C904813F84F53D135BE7DA989C04}\\
\path{Renewal_Yang_Mills_Mass_Gap_Research_Notes_V36.tex}
 &304790&\texttt{D89406E870AB91AE25485633381382D4}\newline \texttt{55F5D624D6F861039089DB470C5AF9D3}
\end{tabular}
\end{center}
The rolling SM file advanced from V19 through V20 to V21 during the sweep.
V21 is append-only and contains both intervening sections, so the cutoff loses
no V19--V20 result.  Likewise V36 contains V35; both files happened to be
present at the timestamp in \textup{(V30.2)}.  The explicit timestamp and
manifest make this concurrent observation reproducible without pretending
that a moving directory is static.

\begin{assessment}[V30 transfer verdict]
\label{aud:v30-transfer}
The changed corpus contributes the following typed lessons.
\begin{enumerate}[label=\textup{(\roman*)},leftmargin=3em]
\item SM archive f14 supplies exact conditional-expectation projections,
      conditional-variance Pythagoras, and the warning that exact local
      projection gaps do not imply a global gap.  V30 imports only this
      Hilbert-space mechanism.
\item Rolling SM V20--V21 proves that a natural source coordinate can expose
      an obstruction hidden by an endpoint metric, and that even the corrected
      untruncated metric can fail under a partial update.  The GRH analogue is
      to measure the actual conditional phase regression, not merely residue
      vector norms.
\item YM V35--V36 proves by exact invariants that visually similar coordinate
      branches need not be conjugate.  It also exhibits the connection terms
      created by differentiating a moving frame.  No residue levels are
      identified by an unproved symmetry here, and the filtration below is
      not differentiated in \(t\).
\item The \(27\) unchanged rows, including NS V26 and the parallel YM V5 note,
      add no new operator or estimate to the V29 contact problem.
\end{enumerate}
No physical gap, matrix floor, transport radius, or continuum conclusion from
another programme is imported.
\end{assessment}

\subsection{The conductor-power filtration}

Retain the V28 probability space
\[
 \Omega_\chi:=\{n\ge2:\chi(n)\ne0\},\qquad
 \nu_{\chi,t}(n)
 =\frac{\Lambda(n)n^{-1/2}e^{-t(\log n)^2}}{Z_\chi(t)}.
 \tag{V30.4}
\]
Put
\[
 U(n):=\log n,qquad
 \Xi_x(n):=\frac{\chi(n)}{|\chi(n)|}e^{ix\log n}.
 \tag{V30.5}
\]
Thus
\[
 \mathbb E U=\mu_\chi(t),\qquad
 \mathbb E\Xi_x=z_{\chi,t}(x),\qquad |\Xi_x|=1.
 \tag{V30.6}
\]
For \(j\ge1\), let
\[
 Q_j:=q^j,qquad
 \mathcal F_j:=\sigma(n\bmod Q_j),qquad
 \mathcal F_0:=\{\varnothing,\Omega_\chi\}.
 \tag{V30.7}
\]
Because \(Q_j\mid Q_{j+1}\), these sigma algebras are increasing.  Define
\[
 U_j:=\mathbb E[U\mid\mathcal F_j],\qquad
 X_j(x):=\mathbb E[\Xi_x\mid\mathcal F_j],
 \tag{V30.8}
\]
and their martingale differences
\[
 d_jU:=U_j-U_{j-1},\qquad
 d_jX(x):=X_j(x)-X_{j-1}(x).
 \tag{V30.9}
\]

\begin{theorem}[Exact growing-modulus martingale regression]
\label{thm:v30-modulus-martingale}
For every \(t>0\) and \(x\in\mathbb R\), the filtrations in
\textup{(V30.7)} generate the full sigma algebra on the countable support
\(\Omega_\chi\), and
\[
 U_j\longrightarrow U,qquad X_j(x)\longrightarrow\Xi_x
 \quad\hbox{in }L^2(\nu_{\chi,t}).
 \tag{V30.10}
\]
Moreover,
\[
 \boxed{
 \mathbb E[U\Xi_x]
 =\mu_\chi(t)z_{\chi,t}(x)
 +\sum_{j\ge1}\mathbb E[d_jU\,d_jX(x)].}
 \tag{V30.11}
\]
The series is absolutely convergent, and its two square-function budgets are
\[
 \sum_{j\ge1}\|d_jU\|_2^2=\sigma_\chi(t)^2,
 \qquad
 \sum_{j\ge1}\|d_jX(x)\|_2^2=1-|z_{\chi,t}(x)|^2.
 \tag{V30.12}
\]
At every finite depth \(J\ge1\), one also has the exact head--tail identity
\[
 \boxed{
 \mathbb E[U\Xi_x]
 =H_{J,t}(x)+T_{J,t}(x),}
 \tag{V30.13}
\]
where
\[
 H_{J,t}(x):=\mathbb E[U_JX_J(x)],qquad
 T_{J,t}(x):=
 \mathbb E[(U-U_J)(\Xi_x-X_J(x))],
 \tag{V30.14}
\]
and
\[
 |T_{J,t}(x)|
 \le
 \|U-U_J\|_2\,\|\Xi_x-X_J(x)\|_2.
 \tag{V30.15}
\]
\end{theorem}

\begin{proof}
If \(m\ne n\), then \(q^j>|m-n|\) for all sufficiently large \(j\), so
\(m\not\equiv n\pmod{q^j}\).  The residue maps therefore separate the points
of the countable set \(\Omega_\chi\).  Their generated sigma algebra contains
every singleton and hence is the full sigma algebra.  Since \(U\in L^2\) by
the Gaussian weight and \(|\Xi_x|=1\), the \(L^2\) martingale convergence
theorem proves \textup{(V30.10)}.

For \(j<k\), the variable \(d_jU\) is
\(\mathcal F_{k-1}\)-measurable while
\(\mathbb E[d_kX\mid\mathcal F_{k-1}]=0\).  Hence
\[
 \mathbb E[d_jU\,d_kX]=0.
 \tag{V30.16}
\]
The same argument with \(U\) and \(X\) interchanged handles \(j>k\).
Expanding the two martingales and passing to the \(L^2\) limits leaves only
the diagonal terms, proving \textup{(V30.11)}.  Orthogonality of martingale
differences gives \textup{(V30.12)}.  In particular,
\[
 \sum_{j\ge1}|\mathbb E[d_jU\,d_jX]|
 \le
 \left(\sum_j\|d_jU\|_2^2\right)^{1/2}
 \left(\sum_j\|d_jX\|_2^2\right)^{1/2}<\infty.
 \tag{V30.17}
\]
Finally, conditional expectation is an orthogonal projection.  The cross
terms between the projected and residual parts vanish, giving
\textup{(V30.13)}; Cauchy--Schwarz gives \textup{(V30.15)}.
\end{proof}

The first level is exactly the conductor residue sigma algebra used in V29.
All later levels are genuine arithmetic refinements, not additional moments
of the same fixed residue vector.

\subsection{Fixed depth is blind to the Gaussian saddle}

The exact decomposition alone does not create cancellation.  The following
estimate identifies how much of the two square functions a fixed modulus can
see.

\begin{theorem}[Fixed-depth modulus blindness]
\label{thm:v30-fixed-depth-blindness}
Fix \(J\ge1\), and let \(K\Subset\mathbb R\setminus\{0\}\).  There are
constants \(A_J<\infty\), \(c_J>0\), and \(C_{J,K}<\infty\) such that, as
\(t\downarrow0\),
\[
 \boxed{
 \|U_J-\mu_\chi(t)\|_2
 \le C_{J,K}t^{-A_J}e^{-c_J/\sqrt t},}
 \tag{V30.18}
\]
and
\[
 \boxed{
 \sup_{x\in K}\|X_J(x)\|_2
 \le C_{J,K}t^{-A_J}e^{-c_J/\sqrt t}.}
 \tag{V30.19}
\]
Consequently, uniformly for \(x\in K\),
\[
 \|U-U_J\|_2^2
 =\frac1{2t}(1+o_{\chi,J}(1)),
 \qquad
 \|\Xi_x-X_J(x)\|_2^2=1+o_{\chi,J,K}(1).
 \tag{V30.20}
\]
Thus the absolute tail budget in \textup{(V30.15)} has the exact critical
rate
\[
 \boxed{
 \lim_{t\downarrow0}4t\log\left(
 1+Z_\chi(t)\|U-U_J\|_2\|\Xi_x-X_J(x)\|_2
 \right)=\frac14,}
 \tag{V30.21}
\]
uniformly on \(K\).
\end{theorem}

\begin{proof}
Put \(Q=q^J\).  Every residue occurring in \(\Omega_\chi\) is reduced
modulo \(Q\).  For \(a\in(\mathbb Z/Q\mathbb Z)^\times\) and \(r=0,1\), set
\[
 A_{a,t}^{[r]}(x)
 :=\sum_{\substack{n\ge2\\n\equiv a\;(\mathrm{mod}\ Q)}}
 \frac{\Lambda(n)(\log n)^r}{\sqrt n}
 e^{-t(\log n)^2}e^{ix\log n}.
 \tag{V30.22}
\]
The fixed-modulus PNT in progressions and the Gaussian Abel argument of
Proposition~\ref{prop:v29-ap-gaussian} give, uniformly for \(a\) and for \(x\)
in a fixed compact set,
\[
 A_{a,t}^{[r]}(x)
 =\frac1{\varphi(Q)}I_t^{[r]}(x)
 +O_{Q,K,r}\left(
 t^{-B_r}e^{1/(16t)-c_Q/\sqrt t}
 \right),
 \tag{V30.23}
\]
where
\[
 I_t^{[r]}(x)
 :=\int_0^\infty u^r e^{u/2-tu^2}e^{ixu}\,du.
 \tag{V30.24}
\]
In particular,
\[
 I_t^{[0]}(0)
 =(1+o(1))\sqrt{\frac\pi t}e^{1/(16t)},
 \qquad
 \frac{I_t^{[1]}(0)}{I_t^{[0]}(0)}
 =\frac1{4t}+o(t^{-1/2}).
 \tag{V30.25}
\]
Dividing the \(r=1\) estimate by the positive \(r=0\) estimate shows that
each residue-conditional mean differs from their common main ratio by at
most a polynomial times \(e^{-c_Q/\sqrt t}\).  The global mean is a convex
combination of these conditional means.  This proves \textup{(V30.18)}.

On the residue atom \(a\pmod Q\), the character is constant and
\[
 X_J(x)
 =\chi(a)\frac{A_{a,t}^{[0]}(x)}{A_{a,t}^{[0]}(0)}.
 \tag{V30.26}
\]
Completing the square on the full real line and bounding the omitted negative
half-line gives, uniformly on compact \(x\)-sets,
\[
 \left|\frac{I_t^{[0]}(x)}{I_t^{[0]}(0)}\right|
 \le C_K\left(e^{-x^2/(4t)}+t^{1/2}e^{-1/(16t)}\right).
 \tag{V30.27}
\]
Since \(K\) stays a positive distance from zero, insertion of
\textup{(V30.23)} into \textup{(V30.26)} proves
\textup{(V30.19)}.

Pythagoras for conditional expectation now yields
\[
 \|U-U_J\|_2^2
 =\sigma_\chi(t)^2-\|U_J-\mu_\chi(t)\|_2^2,
 \tag{V30.28}
\]
and
\[
 \|\Xi_x-X_J(x)\|_2^2
 =\|\Xi_x\|_2^2-\|X_J(x)\|_2^2.
 \tag{V30.29}
\]
Use Proposition~\ref{prop:v28-saddle-moments},
\textup{(V30.18)}, and \textup{(V30.19)} to obtain
\textup{(V30.20)}.  Combining this with \textup{(V28.4)} proves
\textup{(V30.21)}.
\end{proof}

The exclusion of \(x=0\) in Theorem~\ref{thm:v30-fixed-depth-blindness} is
essential: there \(\Xi_0=\chi(n)\) is already
\(\mathcal F_1\)-measurable, so its conditional tail is identically zero.
Nonzero-height contacts have the opposite behavior: every fixed conductor
power leaves asymptotically all phase energy in the unresolved tail.

\begin{corollary}[A fixed-depth Pythagorean estimate is still critical]
\label{cor:v30-fixed-depth-no-go}
On any compact \(K\Subset\mathbb R\setminus\{0\}\), an argument which bounds
the tail in \textup{(V30.13)} only by conditional Pythagoras and
Cauchy--Schwarz retains logarithmic rate \(1/4\).  The fixed head satisfies
only the classical estimate
\[
 |H_{J,t}(x)|
 \le C_{J,K}t^{-A'_J}e^{-c'_J/\sqrt t},
 \tag{V30.30}
\]
before multiplication by \(Z_\chi(t)\), so fixed-modulus PNT supplies no
zero-rate conclusion there either.  Therefore fixed-depth conditional energy,
even though exact, does not prove the V28 contact criterion.
\end{corollary}

\begin{proof}
Equation \textup{(V30.21)} proves the tail statement.  For the head, use
\[
 |H_{J,t}(x)|\le\|U_J\|_2\|X_J(x)\|_2
\]
together with \(\|U_J\|_2=O_\chi(t^{-1})\) and
\textup{(V30.19)}.  After multiplication by \(Z_\chi(t)\), the available
upper bound still contains \(e^{1/(16t)-c/\sqrt t}\), whose \(4t\)-logarithmic
rate is \(1/4\).  No assertion is made that the actual signed head or tail
saturates these bounds.
\end{proof}

\subsection{The exact signed-tail contact criterion}

By \textup{(V28.35)}, the normalized regression is simply
\[
 \mathcal R_{\chi,t}(x)=\Im\mathbb E[U\Xi_x].
 \tag{V30.31}
\]
For a compact interval \(I\), write
\[
 \mathcal C_{\chi,t}(I)
 :=\left\{x\in I:
 \Re z_{\chi,t}(x)=\frac{a_{\chi,t}(x)}{2Z_\chi(t)}
 \right\}.
 \tag{V30.32}
\]

\begin{criterion}[Growing-modulus signed-tail criterion]
\label{crit:v30-growing-modulus}
Let \(J(t)\ge1\) be any integer-valued function.  GRH for the paired channel
is equivalent to the statement that, for every nondegenerate compact interval
\(I\),
\[
 \boxed{
 \limsup_{t\downarrow0}4t\log\left(
 1+\sup_{x\in\mathcal C_{\chi,t}(I)}
 Z_\chi(t)
 \left|\Im\left(
 H_{J(t),t}(x)+T_{J(t),t}(x)
 \right)\right|
 \right)=0.}
 \tag{V30.33}
\]
The supremum is zero when the contact set is empty.
\end{criterion}

\begin{proof}
The head--tail identity \textup{(V30.13)} is exact at every depth, including
a depth depending on \(t\).  Equations \textup{(V30.31)} and
\textup{(V30.33)} are therefore exactly the V28 phase-regression criterion
\textup{(V28.37)}.  Apply Criterion~\ref{crit:v28-phase-regression}.
\end{proof}

Criterion~\ref{crit:v30-growing-modulus} is not advertised as a proof of
GRH.  Its gain over V29 is localization: the missing scalar is separated into
a residue head at modulus \(q^{J(t)}\) and an orthogonal within-residue tail,
with an exact multiscale expansion and exact energy ledger.  Theorem
\ref{thm:v30-fixed-depth-blindness} proves that \(J(t)=O(1)\) cannot improve
the old absolute method away from height zero.  A new arithmetic estimate
must do at least one of the following:
\begin{enumerate}[label=\textup{(\roman*)},leftmargin=3em]
\item obtain signed cancellation between \(H_{J(t),t}\) and \(T_{J(t),t}\);
\item let \(J(t)\to\infty\) and control both terms uniformly in the growing
      modulus \(q^{J(t)}\); or
\item bypass conditional energy and estimate the multiplicative regression
      directly.
\end{enumerate}
Merely summing absolute martingale increments returns
\textup{(V30.17)} and hence the V28 covariance envelope.

\begin{firewall}[Conditional Pythagoras is not arithmetic cancellation]
V30 completes the required corpus sweep and imports an exact conditional
projection mechanism.  It proves the conductor-power martingale identity and
the fixed-depth blindness theorem.  It does not bound the signed tail in
\textup{(V30.33)}, does not provide estimates uniform for a modulus growing
with \(t^{-1}\), and does not exclude any off-line zero.  The parallel SM and
YM conclusions are not GRH inputs.  GRH remains open.
\end{firewall}

\begin{assessment}[V30 acquisition and next upstream move]
V30 turns the qualitative request for a growing-dimensional correlation into
an exact object.  The V29 conductor vector is the first conditional level;
successive conductor powers form a complete arithmetic filtration; and the
phase regression is the absolutely convergent diagonal pairing of its two
martingale difference sequences.  Fixed depth sees exponentially little of
the saddle log-coordinate and, away from zero height, exponentially little of
the phase, leaving a tail whose absolute capacity has the same critical rate
\(1/4\).

V31 should derive the \(t\)-variation formula for the conditional projections
before choosing a growing depth.  Because the measure
\(\nu_{\chi,t}\) itself moves with \(t\), differentiating \(U_J\), \(X_J\), or
the head--tail split creates score-covariance connection terms, just as the YM
V36 audit warns for a moving frame.  The concrete next target is an exact
evolution identity for \(H_{J,t}+T_{J,t}\), followed by a test of whether any
choice \(q^{J(t)}\) below the Gaussian saddle scale yields a non-GRH-complete
signed estimate.  If it only repackages the logarithmic derivative, that route
must be closed explicitly.
\end{assessment}

\section*{Version V30 append-only record}
\addcontentsline{toc}{section}{Version V30 append-only record}
The complete V29 source was retained before this continuation.  It had
\(450550\) bytes, \(13483\) lines, and SHA--256
\[
 \texttt{7A7EE9D4A179E69CC5DF491C5D71C0DFCA0FBC452B05C518087AD61389F62922}.
\]
V30 completed the mandatory corpus sweep, established the exact
conductor-power martingale and head--tail identities, proved fixed-depth
modulus blindness at the Gaussian saddle, and isolated the growing-modulus
signed-tail criterion.

% =============================================================================
% END APPENDED V30 MATERIAL
% =============================================================================
% =============================================================================
% BEGIN APPENDED V31 MATERIAL
% =============================================================================

\section{V31: moving-measure connection and polynomial-modulus blindness}
\label{sec:grh-v31}

V30 decomposed the normalized derivative regression by conditional
expectations on the conductor-power filtration
\(\mathcal F_J=\sigma(n\bmod q^J)\).  It proved that every fixed depth leaves
the full critical absolute capacity in the orthogonal tail.  Two questions
were left:
\begin{enumerate}[label=\textup{(\roman*)},leftmargin=3em]
\item what connection terms appear because the saddle probability
      \(\nu_{\chi,t}\) moves with \(t\); and
\item how far the blindness theorem persists when \(J=J(t)\) grows.
\end{enumerate}

Both questions have exact answers.  Conditional differentiation is governed
by the conditional covariance with \(U^2\).  In the residual tail, the two
connection terms vanish after expectation by conditional orthogonality,
leaving a particularly simple evolution identity.  Uniform
Siegel--Walfisz, applied only on the actual saddle region, then extends V30
from fixed \(J\) to every modulus \(q^{J(t)}\) bounded by a fixed power of
\(t^{-1}\).  Consequently logarithmic filtration depth is still
asymptotically blind at every nonzero height.

This is not a new proof of GRH.  It closes two tempting shortcuts: a finite
score-connection hierarchy is just the old heat hierarchy, and
polynomial-size residue moduli do not resolve the phase tail.

\subsection{The score connection for a conditional expectation}

Write \(\mathbb E_t\) for expectation under \(\nu_{\chi,t}\), retain
\[
 U(n)=\log n,\qquad
 \Xi_x(n)=\frac{\chi(n)}{|\chi(n)|}e^{ixU(n)},
 \tag{V31.1}
\]
and put
\[
 M_2(t):=\mathbb E_tU^2.
 \tag{V31.2}
\]
The global logarithmic score of the saddle law is
\[
 S_t(n):=\partial_t\log\nu_{\chi,t}(n)
 =-\bigl(U(n)^2-M_2(t)\bigr).
 \tag{V31.3}
\]
For a fixed depth \(J\), denote conditional expectation by
\[
 \Pi_{J,t}F:=\mathbb E_t[F\mid\mathcal F_J]
 \tag{V31.4}
\]
and define the bilinear conditional covariance
\[
 \operatorname{Cov}_t(F,G\mid\mathcal F_J)
 :=\Pi_{J,t}(FG)-(\Pi_{J,t}F)(\Pi_{J,t}G).
 \tag{V31.5}
\]
No complex conjugation is placed in this covariance; this is the bilinear
pairing occurring in the character reader.

\begin{theorem}[Exact conditional score connection]
\label{thm:v31-score-connection}
Let \(F_t\) be a complex-valued family for which \(F_t\),
\(\partial_tF_t\), and their products with \(U^2\) are integrable locally in
\(t>0\).  Then, on every atom of \(\mathcal F_J\),
\[
 \boxed{
 \partial_t(\Pi_{J,t}F_t)
 =
 \Pi_{J,t}(\partial_tF_t)
 -\operatorname{Cov}_t(F_t,U^2\mid\mathcal F_J).}
 \tag{V31.6}
\]
Globally,
\[
 \boxed{
 \partial_t\mathbb E_tF_t
 =
 \mathbb E_t(\partial_tF_t)
 -\operatorname{Cov}_t(F_t,U^2).}
 \tag{V31.7}
\]
In particular, for the V30 projected variables,
\[
 \partial_tU_J
 =-\operatorname{Cov}_t(U,U^2\mid\mathcal F_J),
 \qquad
 \partial_tX_J(x)
 =-\operatorname{Cov}_t(\Xi_x,U^2\mid\mathcal F_J).
 \tag{V31.8}
\]
\end{theorem}

\begin{proof}
On an atom \(A\) of \(\mathcal F_J\), write
\[
 D_A(t):=\sum_{n\in A}w(n)e^{-tU(n)^2},
 \qquad
 N_A(t):=\sum_{n\in A}F_t(n)w(n)e^{-tU(n)^2},
 \tag{V31.9}
\]
where \(w(n)=\Lambda(n)n^{-1/2}\).  The conditional expectation is
\(N_A/D_A\).  Differentiation of the quotient gives
\[
 \frac{d}{dt}\frac{N_A}{D_A}
 =
 \Pi_{J,t}(\partial_tF_t)
 -\Pi_{J,t}(F_tU^2)
 (\Pi_{J,t}F_t)(\Pi_{J,t}U^2),
 \tag{V31.10}
\]
which is \textup{(V31.6)}.  Gaussian damping justifies termwise
differentiation.  Taking \(A=\Omega_\chi\) proves
\textup{(V31.7)}, and \(F=U,\Xi_x\) gives \textup{(V31.8)}.
\end{proof}

Equation \textup{(V31.6)} is the missing moving-frame term.  Treating
\(\Pi_{J,t}\) as a \(t\)-independent orthogonal projector would delete the
conditional covariance and is therefore false.

\subsection{Connection-completed head and tail evolution}

Suppress \(t,x\) temporarily and put
\[
 \widetilde U_J:=U-U_J,\qquad
 \widetilde X_J:=\Xi_x-X_J(x).
 \tag{V31.11}
\]
The V30 head and tail are
\[
 H_{J,t}(x)=\mathbb E_t[U_JX_J(x)],
 \qquad
 T_{J,t}(x)=\mathbb E_t[\widetilde U_J\widetilde X_J].
 \tag{V31.12}
\]

\begin{theorem}[Exact head--tail connection cancellation]
\label{thm:v31-head-tail-evolution}
For every fixed \(J\ge1\), \(t>0\), and \(x\in\mathbb R\),
\[
 \boxed{
 \partial_tT_{J,t}(x)
 =
 -\operatorname{Cov}_t(
 \widetilde U_J\widetilde X_J,U^2).}
 \tag{V31.13}
\]
Moreover,
\[
 \boxed{
 \partial_tH_{J,t}(x)
 =
 -\operatorname{Cov}_t(U\Xi_x,U^2)
 +\operatorname{Cov}_t(
 \widetilde U_J\widetilde X_J,U^2),}
 \tag{V31.14}
\]
and hence
\[
 \boxed{
 \partial_t(H_{J,t}+T_{J,t})
 =
 -\operatorname{Cov}_t(U\Xi_x,U^2).}
 \tag{V31.15}
\]
\end{theorem}

\begin{proof}
Theorem~\ref{thm:v31-score-connection} gives
\[
 \partial_t\widetilde U_J
 =\operatorname{Cov}_t(U,U^2\mid\mathcal F_J),
 \qquad
 \partial_t\widetilde X_J
 =\operatorname{Cov}_t(\Xi_x,U^2\mid\mathcal F_J).
 \tag{V31.16}
\]
Both right sides are \(\mathcal F_J\)-measurable, whereas
\[
 \mathbb E_t[\widetilde U_J\mid\mathcal F_J]
 =
 \mathbb E_t[\widetilde X_J\mid\mathcal F_J]=0.
 \tag{V31.17}
\]
Consequently the two differentiated-residual terms vanish:
\[
 \mathbb E_t[
 (\partial_t\widetilde U_J)\widetilde X_J]
 =
 \mathbb E_t[
 \widetilde U_J(\partial_t\widetilde X_J)]=0.
 \tag{V31.18}
\]
Applying the global score formula \textup{(V31.7)} to
\(\widetilde U_J\widetilde X_J\) proves \textup{(V31.13)}.
The exact identity
\[
 H_{J,t}+T_{J,t}=\mathbb E_t[U\Xi_x]
 \tag{V31.19}
\]
and \textup{(V31.7)} applied to \(U\Xi_x\) then prove
\textup{(V31.14)}--\textup{(V31.15)}.
\end{proof}

Normalization removes the global score centering.  Since
\[
 Z_\chi'(t)=-Z_\chi(t)M_2(t),
 \tag{V31.20}
\]
Theorem~\ref{thm:v31-head-tail-evolution} has the unnormalized form
\[
 \boxed{
 \partial_t\!\left(Z_\chi T_{J,t}\right)
 =
 -Z_\chi\,
 \mathbb E_t[
 \widetilde U_J\widetilde X_JU^2],}
 \tag{V31.21}
\]
\[
 \boxed{
 \partial_t\!\left(Z_\chi H_{J,t}\right)
 =
 -Z_\chi\,
 \mathbb E_t[
 (U\Xi_x-\widetilde U_J\widetilde X_J)U^2].}
 \tag{V31.22}
\]
Both left sides tend to zero as \(t\to\infty\), so integration gives the
exact conditional Duhamel identities
\[
 \boxed{
 Z_\chi(t)T_{J,t}(x)
 =
 \int_t^\infty
 Z_\chi(s)\,
 \mathbb E_s[
 \widetilde U_{J,s}\widetilde X_{J,s}U^2]\,ds,}
 \tag{V31.23}
\]
\[
 \boxed{
 Z_\chi(t)H_{J,t}(x)
 =
 \int_t^\infty
 Z_\chi(s)\,
 \mathbb E_s[
 (U\Xi_x-\widetilde U_{J,s}\widetilde X_{J,s})U^2]\,ds.}
 \tag{V31.24}
\]
Here the residuals inside the integrals are formed with
\(\nu_{\chi,s}\), not frozen at the lower endpoint.

\begin{proposition}[The completed connection is the old heat hierarchy]
\label{prop:v31-connection-collapse}
Adding \textup{(V31.21)} and \textup{(V31.22)} yields
\[
 \partial_tP_{\chi,t}^{[1]}(x)
 =-P_{\chi,t}^{[3]}(x)
 =\partial_x^2P_{\chi,t}^{[1]}(x).
 \tag{V31.25}
\]
More generally, for every fixed \(k\ge0\),
\[
 \boxed{
 \partial_t^kP_{\chi,t}^{[1]}(x)
 =(-1)^kP_{\chi,t}^{[2k+1]}(x).}
 \tag{V31.26}
\]
Thus finitely many score-connection derivatives add no closed equation
beyond the scalar heat hierarchy of V22.  An absolute moment estimate for
either Duhamel integrand retains logarithmic rate \(1/4\).
\end{proposition}

\begin{proof}
The sum of the two integrands in \textup{(V31.21)}--\textup{(V31.22)}
is \(U^3\Xi_x\).  Since
\[
 P_{\chi,t}^{[r]}(x)
 =Z_\chi(t)\mathbb E_t[U^r\Xi_x],
 \tag{V31.27}
\]
this proves the first equality in \textup{(V31.25)}.  Twice
differentiating \(e^{ixU}\) proves the second.  Iteration proves
\textup{(V31.26)}.

For the tail integrand, \(|\widetilde X_J|\le2\) and
Cauchy--Schwarz give
\[
 Z_\chi\left|
 \mathbb E_t[\widetilde U_J\widetilde X_JU^2]\right|
 \le
 2Z_\chi
 \|\widetilde U_J\|_2
 (\mathbb E_tU^4)^{1/2}.
 \tag{V31.28}
\]
The head integrand is bounded by the same expression plus
\(Z_\chi\mathbb E_tU^3\).  Proposition~\ref{prop:v28-saddle-moments}
shows that every factor after \(Z_\chi\) is polynomial in \(t^{-1}\);
\textup{(V28.4)} therefore gives rate \(1/4\).  Absolute integration
from \(t\) to a fixed positive time changes only a polynomial factor.
\end{proof}

The proposition is a nonduplication boundary.  V22 proved the complete
symbol heat equation.  V31 determines exactly how a moving arithmetic
conditional split realizes that equation and proves that the apparent new
connection does not close the regression hierarchy.

\subsection{Uniform progressions at polynomial conductor powers}

V30 used the fixed-modulus PNT in progressions.  We now allow the modulus to
grow.  For \(Q\) divisible by \(q\), define
\[
 A_{a,t,Q}^{[r]}(x)
 :=
 \sum_{\substack{n\ge2\\n\equiv a\;(\mathrm{mod}\ Q)}}
 \frac{\Lambda(n)(\log n)^r}{\sqrt n}
 e^{-t(\log n)^2}e^{ix\log n},
 \qquad (a,Q)=1.
 \tag{V31.29}
\]

\begin{lemma}[Gaussian Siegel--Walfisz on the saddle window]
\label{lem:v31-gaussian-sw}
Fix \(B>0\), an integer \(r\ge0\), and a compact interval
\(I\subset\mathbb R\).  Uniformly for
\[
 q\mid Q,\qquad Q\le t^{-B},\qquad (a,Q)=1,\qquad x\in I,
 \tag{V31.30}
\]
there are constants \(C,c>0\) and \(M<\infty\), depending on
\(\chi,B,r,I\), such that
\[
 \boxed{
 A_{a,t,Q}^{[r]}(x)
 =
 \frac1{\varphi(Q)}I_t^{[r]}(x)
 +O\!\left(
 t^{-M}e^{1/(16t)-c/\sqrt t}
 \right),}
 \tag{V31.31}
\]
where \(I_t^{[r]}\) is the continuous reader in
\textup{(V30.24)}.  The constants may inherit the usual
ineffectivity of Siegel--Walfisz.
\end{lemma}

\begin{proof}
For every fixed \(D>0\), Siegel--Walfisz gives
\[
 \Psi(y;Q,a)
 =
 \frac{y}{\varphi(Q)}
 +O_D\!\left(ye^{-c_D\sqrt{\log y}}\right)
 \tag{V31.32}
\]
uniformly for \(Q\le(\log y)^D\) and \((a,Q)=1\).
Choose \(D>B+1\), put \(u=\log y\), and split the Abel integral at
\[
 u_0=\frac1{8t}.
 \tag{V31.33}
\]
On \(u\ge u_0\), the condition \(Q\le t^{-B}\) implies
\(Q\le u^D\) for all sufficiently small \(t\), so
\textup{(V31.32)} applies uniformly.  The Gaussian Abel estimate of
V29 then gives the error on the right of \textup{(V31.31)}.

On \(0\le u\le u_0\), use the trivial bound by the total
Chebyshev mass.  The exponent obeys
\[
 \max_{0\le u\le1/(8t)}
 \left(\frac u2-tu^2\right)
 =\frac3{64t},
 \tag{V31.34}
\]
whereas the saddle exponent is \(1/(16t)=4/(64t)\).
This lower region is therefore smaller than the main saddle by
\(e^{-1/(64t)}\), up to polynomial factors and the polynomial loss
\(\varphi(Q)\le Q\le t^{-B}\).  It is absorbed by the stated error.
\end{proof}

\begin{theorem}[Polynomial-modulus blindness]
\label{thm:v31-polynomial-blindness}
Fix \(B>0\), and let \(J(t)\ge1\) be any integer-valued function satisfying
\[
 Q(t):=q^{J(t)}\le t^{-B}.
 \tag{V31.35}
\]
Let \(U_{J(t)}\) and \(X_{J(t)}(x)\) be the conditional projections from
V30 at the indicated depth.  For every
\(K\Subset\mathbb R\setminus\{0\}\), there are \(A<\infty\) and \(c>0\)
such that
\[
 \boxed{
 \|U_{J(t)}-\mu_\chi(t)\|_2
 \le t^{-A}e^{-c/\sqrt t},}
 \tag{V31.36}
\]
\[
 \boxed{
 \sup_{x\in K}\|X_{J(t)}(x)\|_2
 \le t^{-A}e^{-c/\sqrt t}.}
 \tag{V31.37}
\]
Consequently, uniformly on \(K\),
\[
 \|U-U_{J(t)}\|_2^2
 =\frac1{2t}(1+o(1)),
 \qquad
 \|\Xi_x-X_{J(t)}(x)\|_2^2=1+o(1),
 \tag{V31.38}
\]
and the absolute conditional-tail budget still has exact rate
\[
 \boxed{
 \lim_{t\downarrow0}4t\log\left(
 1+Z_\chi(t)
 \|U-U_{J(t)}\|_2
 \|\Xi_x-X_{J(t)}(x)\|_2
 \right)=\frac14.}
 \tag{V31.39}
\]
\end{theorem}

\begin{proof}
On a residue atom \(a\bmod Q(t)\),
\[
 U_{J(t)}
 =
 \frac{A_{a,t,Q(t)}^{[1]}(0)}
 {A_{a,t,Q(t)}^{[0]}(0)},
 \qquad
 X_{J(t)}(x)
 =
 \chi(a)
 \frac{A_{a,t,Q(t)}^{[0]}(x)}
 {A_{a,t,Q(t)}^{[0]}(0)}.
 \tag{V31.40}
\]
Lemma~\ref{lem:v31-gaussian-sw} applies uniformly.  The main denominator is
\[
 \frac1{\varphi(Q(t))}I_t^{[0]}(0)
 \asymp
 \frac1{\varphi(Q(t))}
 t^{-1/2}e^{1/(16t)}.
 \tag{V31.41}
\]
Division costs at most the polynomial factor
\(\varphi(Q(t))\le t^{-B}\).  The quotient for \(U_{J(t)}\) therefore
differs from the common saddle mean by a polynomial times
\(e^{-c/\sqrt t}\), uniformly in the residue.  Since the global mean is
their convex combination, this proves \textup{(V31.36)}.

For \(X_{J(t)}\), the continuous oscillatory ratio satisfies
\[
 \left|
 \frac{I_t^{[0]}(x)}{I_t^{[0]}(0)}
 \right|
 \le C_K\left(
 e^{-x^2/(4t)}+t^{1/2}e^{-1/(16t)}
 \right)
 \tag{V31.42}
\]
as in V30.  On \(K\), this and the uniform error in
\textup{(V31.31)} prove \textup{(V31.37)}.

Conditional Pythagoras gives
\[
 \|U-U_{J(t)}\|_2^2
 =
 \sigma_\chi(t)^2
 -\|U_{J(t)}-\mu_\chi(t)\|_2^2
 \tag{V31.43}
\]
and
\[
 \|\Xi_x-X_{J(t)}(x)\|_2^2
 =1-\|X_{J(t)}(x)\|_2^2.
 \tag{V31.44}
\]
Proposition~\ref{prop:v28-saddle-moments} and
\textup{(V31.36)}--\textup{(V31.37)} prove
\textup{(V31.38)}.  Equation \textup{(V28.4)} then proves
\textup{(V31.39)}.
\end{proof}

\begin{corollary}[Superlogarithmic depth is necessary for energy resolution]
\label{cor:v31-superlog-depth}
Suppose \(K\Subset\mathbb R\setminus\{0\}\).  If
\[
 J(t)\le C\log(1/t)
 \tag{V31.45}
\]
along a sequence \(t\downarrow0\), then the phase residual satisfies
\[
 \inf_{x\in K}
 \|\Xi_x-X_{J(t)}(x)\|_2\longrightarrow1
 \tag{V31.46}
\]
along that sequence.  Therefore any conductor-power strategy which seeks to
make the phase tail small by conditional measurability must have
\[
 \frac{J(t)}{\log(1/t)}\longrightarrow\infty
 \tag{V31.47}
\]
along every successful nonzero-height sequence.
\end{corollary}

\begin{proof}
Equation \textup{(V31.45)} gives
\[
 q^{J(t)}
 \le t^{-C\log q}.
 \tag{V31.48}
\]
Apply Theorem~\ref{thm:v31-polynomial-blindness} with
\(B=C\log q\).  This proves \textup{(V31.46)}.  Its contrapositive proves
\textup{(V31.47)}.
\end{proof}

This corollary concerns resolution by conditional energy.  It does not rule
out a signed cancellation between a small projected head and a large
orthogonal tail.  Such a cancellation is exactly the arithmetic information
still missing from the contact criterion.

\subsection{Growing-depth jumps and the sharpened stopping line}

The canonical polynomial-depth choice
\[
 J_B(t):=\max\{J\ge1:q^J\le t^{-B}\}
 \tag{V31.49}
\]
is a step function.  The connection formulas apply on each interval on
which \(J_B\) is constant.  At a jump from \(J\) to \(J+1\), the head and
tail may jump separately, but the exact V30 decomposition gives
\[
 \boxed{
 H_{J+1,t}-H_{J,t}
 =-\bigl(T_{J+1,t}-T_{J,t}\bigr).}
 \tag{V31.50}
\]
Thus no artificial delta term occurs in their sum.  Any future
\(t\)-differentiation of a growing filtration must retain both the continuous
score connection \textup{(V31.6)} and the cancelling discrete jumps
\textup{(V31.50)}.

\begin{criterion}[Connection-completed polynomial-depth contact criterion]
\label{crit:v31-polynomial-contact}
For any fixed \(B>0\), GRH for the paired channel is equivalent to
\[
 \limsup_{t\downarrow0}4t\log\left(
 1+\sup_{x\in\mathcal C_{\chi,t}(I)}
 Z_\chi(t)
 \left|
 \Im\bigl(
 H_{J_B(t),t}(x)+T_{J_B(t),t}(x)
 \bigr)\right|
 \right)=0
 \tag{V31.51}
\]
for every nondegenerate compact interval \(I\), with the empty-contact
convention of V30.
\end{criterion}

\begin{proof}
This is Criterion~\ref{crit:v30-growing-modulus} with the exact choice
\textup{(V31.49)}.  The jump identity \textup{(V31.50)} confirms that the
sum is independent of the selected depth.
\end{proof}

The criterion remains equivalent to GRH, but Theorem
\ref{thm:v31-polynomial-blindness} adds a genuine negative result: on compact
sets away from zero, no polynomial-size conductor power can make the tail
small in \(L^2\).  The fixed-conductor obstruction of V29 therefore persists
through every logarithmic number of conditional refinements.

\begin{firewall}[The connection and Siegel--Walfisz do not bound the signed tail]
V31 proves the exact moving-measure connection, the connection-completed
head--tail evolution, its collapse to the known heat hierarchy, and uniform
polynomial-modulus blindness.  It does not estimate the signed imaginary
part in \textup{(V31.51)}, does not control superpolynomial conductor powers,
and does not handle the exceptional height \(x=0\) by the blindness theorem.
No off-line zero is excluded, and GRH remains open.
\end{firewall}

\begin{assessment}[V31 acquisition and the next upstream scale]
The conditional projection route now has a sharp scale boundary.
Fixed depth and even \(J(t)=O(\log(1/t))\) are unconditionally too shallow
to resolve the nonzero-height character phase under the Gaussian saddle.
The moving-measure terms are fully accounted for and supply no hidden
finite hierarchy: after normalization they sum to
\(\partial_tP^{[1]}=-P^{[3]}\).

The next question is not another fixed conditional moment.  V32 should test
the intermediate superpolynomial regime by deriving a Gaussian
Bombieri--Vinogradov or large-sieve estimate for moduli divisible by \(q\),
then translating it into an averaged conditional-projection norm.  The
direction of the estimate must be checked carefully: equidistribution makes
\(X_J\) small and therefore leaves the tail large.  If all moduli accessible
below the square-root distribution level remain blind, the residue-filtration
strategy can resolve phase only near the individual-integer scale
\(\exp(1/(4t))\), and that obstruction should be stated before returning to
signed cross-scale cancellation.
\end{assessment}

\section*{Version V31 append-only record}
\addcontentsline{toc}{section}{Version V31 append-only record}
The complete V30 source was retained before this continuation.  It had
\(468439\) bytes, \(13978\) lines, and SHA--256
\[
 \texttt{474D296C475DC8B3D8485AE7C951D70F8094355072947F5BFD84201FBEA44753}.
\]
V31 established the conditional score connection, exact head--tail Duhamel
evolution, finite-connection collapse to the heat hierarchy, Gaussian
Siegel--Walfisz uniformity for polynomial conductor powers, and the
superlogarithmic depth obstruction.

% =============================================================================
% END APPENDED V31 MATERIAL
% =============================================================================
% =============================================================================
% BEGIN APPENDED V32 MATERIAL
% =============================================================================

\section{V32: the large-sieve wrong-way theorem below the square-root scale}
\label{sec:grh-v32}

V31 proved that every conductor-power modulus of polynomial size in
\(t^{-1}\) is blind to the nonzero-height phase energy of the Gaussian
saddle.  It proposed testing the superpolynomial range with a Gaussian
large sieve.  This section carries out that test for auxiliary prime
refinements of the conductor.

Let
\[
 Y_t:=\exp\!\left(\frac1{4t}\right)
 \tag{V32.1}
\]
be the integer scale at the center of the absolute prime saddle.  For every
fixed \(0<\theta<1/2\), take auxiliary primes
\[
 \ell\asymp Y_t^\theta
 =\exp\!\left(\frac{\theta}{4t}\right)
 \tag{V32.2}
\]
and condition on residues modulo \(q\ell\).  A smooth multiplicative
large-sieve estimate shows that, for a density-one subset of these primes,
the new conditional phase component beyond modulus \(q\) tends to zero
exponentially on the \(1/t\) scale.  The old modulus-\(q\) component tends
to zero by the classical fixed-character PNT at the slower
\(e^{-c/\sqrt t}\) rate.  Hence almost every such refinement still leaves
all phase energy in the orthogonal tail.

This conclusion points in the opposite direction from a GRH proof.  Average
equidistribution below the square-root scale produces blind conditional
coordinates; it does not resolve the multiplicative phase.  Any
residue-energy strategy must therefore use a sparse exceptional modulus,
reach the square-root transition, or abandon absolute projection energy in
favor of signed cross-scale cancellation.

\subsection{Auxiliary residue projections and their Fourier energy}

For every modulus \(Q\) divisible by \(q\), let
\[
 \mathcal G_Q:=\sigma(n\bmod Q),\qquad
 \Pi_{Q,t}F:=\mathbb E_t[F\mid\mathcal G_Q].
 \tag{V32.3}
\]
Since \(\mathcal G_q\subseteq\mathcal G_Q\), conditional Pythagoras gives
\[
 \|\Pi_{Q,t}F\|_2^2
 =
 \|\Pi_{q,t}F\|_2^2
 \|(\Pi_{Q,t}-\Pi_{q,t})F\|_2^2.
 \tag{V32.4}
\]

For \((a,Q)=1\), define
\[
 A_{a,t,Q}^{[r]}(x)
 :=
 \sum_{\substack{n\ge2\\n\equiv a\;(\mathrm{mod}\ Q)}}
 \frac{\Lambda(n)(\log n)^r}{\sqrt n}
 e^{-t(\log n)^2}e^{ix\log n}.
 \tag{V32.5}
\]
On the atom \(a\bmod Q\),
\[
 \Pi_{Q,t}\Xi_x
 =
 \chi(a)\frac{A_{a,t,Q}^{[0]}(x)}
 {A_{a,t,Q}^{[0]}(0)}.
 \tag{V32.6}
\]

Let \(\widehat G_Q\) be the Dirichlet-character group modulo \(Q\), and put
\[
 S_{\psi,t}^{[r]}(x)
 :=
 \sum_{n\ge2}
 \frac{\Lambda(n)\psi(n)(\log n)^r}{\sqrt n}
 e^{-t(\log n)^2}e^{ix\log n}.
 \tag{V32.7}
\]
Character orthogonality gives
\[
 A_{a,t,Q}^{[r]}(x)
 =
 \frac1{\varphi(Q)}
 \sum_{\psi\in\widehat G_Q}
 \overline{\psi(a)}S_{\psi,t}^{[r]}(x).
 \tag{V32.8}
\]
Writing \(\psi_{0,Q}\) for the principal character and retaining the
continuous main reader \(I_t^{[r]}(x)\) from V30, Parseval gives the exact
deviation identity
\[
 \boxed{
 \begin{aligned}
 &\varphi(Q)\sum_{a\in G_Q}
 \left|
 A_{a,t,Q}^{[r]}(x)
 -\frac{I_t^{[r]}(x)}{\varphi(Q)}
 \right|^2\\
 &\quad=
 \left|
 S_{\psi_{0,Q},t}^{[r]}(x)-I_t^{[r]}(x)
 \right|^2
 +\sum_{\substack{\psi\in\widehat G_Q\\\psi\ne\psi_{0,Q}}}
 |S_{\psi,t}^{[r]}(x)|^2.
 \end{aligned}}
 \tag{V32.9}
\]

\begin{lemma}[Fourier energy controls conditional phase energy]
\label{lem:v32-fourier-to-conditional}
Let
\[
 \mathcal D_{Q,t}^{[r]}(x)
 :=
 \left|
 S_{\psi_{0,Q},t}^{[r]}(x)-I_t^{[r]}(x)
 \right|^2
 +\sum_{\psi\ne\psi_{0,Q}}|S_{\psi,t}^{[r]}(x)|^2.
 \tag{V32.10}
\]
For every fixed compact \(K\subset\mathbb R\), there is \(M<\infty\) such
that, uniformly for the moduli considered below,
\[
 \boxed{
 \sup_{x\in K}\|\Pi_{Q,t}\Xi_x\|_2^2
 \le
 t^{-M}\left[
 \sup_{x\in K}
 \frac{|I_t^{[0]}(x)|^2+\mathcal D_{Q,t}^{[0]}(x)}
 {|I_t^{[0]}(0)|^2}
 +
 \frac{\mathcal D_{Q,t}^{[0]}(0)}
 {|I_t^{[0]}(0)|^2}
 \right].}
 \tag{V32.11}
\]
Likewise,
\[
 \boxed{
 \|\Pi_{Q,t}U-\mu_\chi(t)\|_2^2
 \le
 t^{-M}
 \left(
 \frac{
 \mathcal D_{Q,t}^{[0]}(0)
 +\mathcal D_{Q,t}^{[1]}(0)}
 {|I_t^{[0]}(0)|^2}
 \right)^{1/2}.}
 \tag{V32.12}
\]
\end{lemma}

\begin{proof}
Put \(m=I_t^{[0]}(0)/\varphi(Q)\) and write
\[
 A_{a,t,Q}^{[0]}(0)=m+e_a.
 \tag{V32.13}
\]
Call a residue good when \(|e_a|\le m/2\).  On good residues the
denominator in \textup{(V32.6)} is at least \(m/2\).  On bad residues,
positivity and \(|e_a|>m/2\) give
\[
 \sum_{\mathrm{bad}\ a}A_{a,t,Q}^{[0]}(0)
 \le \frac6m\sum_a|e_a|^2.
 \tag{V32.14}
\]
The conditional phasor has modulus at most one, so the bad-residue
contribution is bounded by their total probability.  On the good residues,
insert the decomposition
\[
 A_{a,t,Q}^{[0]}(x)
 =\frac{I_t^{[0]}(x)}{\varphi(Q)}+e_a(x)
 \tag{V32.15}
\]
and use \(|v+w|^2\le2|v|^2+2|w|^2\).
Equation \textup{(V32.9)} and the saddle asymptotic
\(Z_\chi(t)\asymp I_t^{[0]}(0)\) yield
\textup{(V32.11)}.  The factor \(\varphi(Q)\) cancels exactly between
\(m^{-1}\) and Parseval; no exponential modulus loss is hidden in
\(t^{-M}\).

For \(\Pi_{Q,t}U\), use the quotient
\[
 \frac{A_{a,t,Q}^{[1]}(0)}
 {A_{a,t,Q}^{[0]}(0)}
 \tag{V32.16}
\]
on good residues, subtract the common ratio
\(I_t^{[1]}(0)/I_t^{[0]}(0)\), and repeat the same argument with the
\(r=0,1\) deviations.  The global mean differs from the common ratio by
the fixed-\(q\) PNT error, which is already included after enlarging the
polynomial factor.  If \(B\) is the union of the bad residue atoms, Jensen
and Cauchy--Schwarz give
\[
 \sum_{\mathrm{bad}\ a}\nu_{\chi,t}(a)
 |\mathbb E_t[U\mid a]|^2
 \le
 \mathbb E_t[U^2\mathbf1_B]
 \le
 (\mathbb E_tU^4)^{1/2}\nu_{\chi,t}(B)^{1/2}.
\]
The bad mass is controlled by \textup{(V32.14)}, while
\(\mathbb E_tU^4\) is polynomial in \(t^{-1}\).  Enlarging the linear
good-residue bound to the displayed square-root bound proves
\textup{(V32.12)}.
\end{proof}

\subsection{A smooth multiplicative large sieve at the prime saddle}

Fix \(0<\alpha<1/8\), put
\[
 R_t:=\exp(\alpha/t),
 \qquad
 \mathcal L_t:=
 \{\ell\text{ prime}:R_t\le\ell\le2R_t,\ \ell\nmid q\}.
 \tag{V32.17}
\]
Thus \(R_t=Y_t^{4\alpha}\), and \(4\alpha<1/2\).
For a character \(\eta\bmod q\) and a nonprincipal character
\(\vartheta\bmod\ell\), define
\[
 \mathcal S_{\eta,\vartheta,t}^{[r]}(x)
 :=
 \sum_{n\ge2}
 \frac{\Lambda(n)\eta(n)\vartheta(n)(\log n)^r}{\sqrt n}
 e^{-t(\log n)^2}e^{ix\log n}.
 \tag{V32.18}
\]

\begin{theorem}[Gaussian large-sieve energy]
\label{thm:v32-gaussian-large-sieve}
For every fixed integer \(r\ge0\) and compact interval \(K\), there is
\(M<\infty\) such that
\[
 \boxed{
 \sum_{\ell\in\mathcal L_t}
 \sum_{\eta\bmod q}
 \sum_{\substack{\vartheta\bmod\ell\\\vartheta\ne\vartheta_0}}
 \sup_{x\in K}
 |\mathcal S_{\eta,\vartheta,t}^{[r]}(x)|^2
 \le
 t^{-M}\left(
 e^{1/(8t)}+R_t^2
 \right).}
 \tag{V32.19}
\]
\end{theorem}

\begin{proof}
For fixed \(\eta\) and \(x\), apply the multiplicative large sieve to a
dyadic block \(N<n\le2N\) with coefficients
\[
 b_n=
 \frac{\Lambda(n)\eta(n)(\log n)^r}{\sqrt n}
 e^{-t(\log n)^2}e^{ix\log n}.
 \tag{V32.20}
\]
Every nonprincipal character modulo a prime \(\ell\) is primitive, so the
standard inequality gives
\[
 \sum_{\ell\le2R_t}
 \frac{\ell}{\varphi(\ell)}
 \sum_{\vartheta\bmod\ell}^{*}
 \left|\sum_{N<n\le2N}b_n\vartheta(n)\right|^2
 \ll
 (N+R_t^2)\sum_{N<n\le2N}|b_n|^2.
 \tag{V32.21}
\]
The elementary bound
\(\sum_{n\le X}\Lambda(n)^2\ll X\log X\) yields, with \(u=\log N\),
\[
 \sum_{N<n\le2N}|b_n|^2
 \ll_r
 (1+u)^{2r+1}e^{-2tu^2}.
 \tag{V32.22}
\]
Minkowski over the dyadic blocks, followed by Cauchy with a polynomial
block weight, reduces the two parts of \textup{(V32.21)} to
\[
 \left(
 \sum_u(1+u)^{r+1}
 e^{u/2-tu^2}
 \right)^2
 R_t^2
 \left(
 \sum_u(1+u)^{r+1}e^{-tu^2}
 \right)^2.
 \tag{V32.23}
\]
The first Gaussian has maximum exponent \(1/(16t)\), so its square is
\(t^{-M}e^{1/(8t)}\); the second sum is polynomial in \(t^{-1}\).
There are only \(\varphi(q)\) choices of \(\eta\).

To obtain the supremum in \(x\), apply the one-dimensional Sobolev
inequality on a slightly larger compact interval.  The \(x\)-derivative
replaces \(r\) by \(r+1\), so it changes only the power \(M\).
This proves \textup{(V32.19)}.
\end{proof}

The two terms in \textup{(V32.19)} have distinct meanings.  The
\(e^{1/(8t)}\) term is the square of the prime-saddle mass, while \(R_t^2\)
is the dual large-sieve cost.  After averaging over
\[
 \#\mathcal L_t\asymp\frac{R_t}{\log R_t},
 \tag{V32.24}
\]
both give a saving precisely when \(0<\alpha<1/8\).

\subsection{Density-one blindness below the square-root scale}

For \(\ell\in\mathcal L_t\), put \(Q_\ell=q\ell\).  By the Chinese
remainder theorem, every character modulo \(Q_\ell\) is a product
\(\eta\vartheta\) with \(\eta\bmod q\) and
\(\vartheta\bmod\ell\).  The sector \(\vartheta\ne\vartheta_0\) is
controlled by Theorem~\ref{thm:v32-gaussian-large-sieve}.  The finitely many
sectors with \(\vartheta=\vartheta_0\) are fixed-\(q\) character readers,
apart from deletion of the powers of \(\ell\).  The fixed-character PNT and
the Gaussian Abel argument give, uniformly for \(\ell\in\mathcal L_t\),
\[
 |S_{\eta\vartheta_0,t}^{[r]}(x)
 -\mathbf1_{\eta=\eta_0}I_t^{[r]}(x)|
 \le
 t^{-M}e^{1/(16t)-c/\sqrt t}
 \tag{V32.25}
\]
on compact \(x\)-sets.  The deleted \(\ell\)-power terms are exponentially
smaller than this bound.

Define
\[
 \delta_\alpha:=
 \min\left\{\alpha,\frac18-\alpha\right\}>0.
 \tag{V32.26}
\]

\begin{theorem}[Large-sieve wrong-way theorem]
\label{thm:v32-wrong-way}
Fix \(0<\alpha<1/8\) and
\(K\Subset\mathbb R\setminus\{0\}\).  There are subsets
\(\mathcal L_t^{\mathrm{good}}\subseteq\mathcal L_t\) such that
\[
 \frac{\#(\mathcal L_t\setminus\mathcal L_t^{\mathrm{good}})}
 {\#\mathcal L_t}
 \le
 t^{-M}e^{-\delta_\alpha/(2t)}
 \longrightarrow0,
 \tag{V32.27}
\]
and, uniformly for \(\ell\in\mathcal L_t^{\mathrm{good}}\),
\[
 \boxed{
 \sup_{x\in K}
 \|\Pi_{Q_\ell,t}\Xi_x\|_2^2
 \le
 t^{-M}\left(
 e^{-\delta_\alpha/(2t)}
 +e^{-c/\sqrt t}
 \right)\longrightarrow0,}
 \tag{V32.28}
\]
\[
 \boxed{
 \|\Pi_{Q_\ell,t}U-\mu_\chi(t)\|_2^2
 \le
 t^{-M}\left(
 e^{-\delta_\alpha/(4t)}
 +e^{-c/(2\sqrt t)}
 \right)\longrightarrow0.}
 \tag{V32.29}
\]
Equivalently, almost every auxiliary prime refinement below the
square-root scale satisfies
\[
 \sup_{x\in K}
 \|\Xi_x-\Pi_{Q_\ell,t}\Xi_x\|_2^2=1-o(1),
 \qquad
 \|U-\Pi_{Q_\ell,t}U\|_2^2
 =\frac1{2t}(1+o(1)).
 \tag{V32.30}
\]
\end{theorem}

\begin{proof}
Divide \textup{(V32.19)} by
\[
 \#\mathcal L_t\,|I_t^{[0]}(0)|^2.
 \tag{V32.31}
\]
Using \textup{(V32.24)} and
\[
 |I_t^{[0]}(0)|^2
 \asymp t^{-1}e^{1/(8t)}
 \tag{V32.32}
\]
gives the average normalized auxiliary energy bound
\[
 t^{-M}\left(
 e^{-\alpha/t}
 +e^{-(1/8-\alpha)/t}
 \right)
 \le t^{-M}e^{-\delta_\alpha/t}.
 \tag{V32.33}
\]
Apply Markov's inequality with threshold
\(e^{-\delta_\alpha/(2t)}\), simultaneously for \(r=0,1\) and the
compact Sobolev supremum.  This produces
\(\mathcal L_t^{\mathrm{good}}\) and proves
\textup{(V32.27)}.

For a good \(\ell\), combine its auxiliary-character estimate with the
fixed sectors \textup{(V32.25)} in the exact Fourier identity
\textup{(V32.9)}.  Lemma~\ref{lem:v32-fourier-to-conditional} proves
\textup{(V32.28)}--\textup{(V32.29)}.  Finally conditional Pythagoras,
\(|\Xi_x|=1\), and the saddle variance
\(\sigma_\chi(t)^2=(2t)^{-1}(1+o(1))\) prove
\textup{(V32.30)}.
\end{proof}

\begin{corollary}[The new auxiliary martingale increment is negligible]
\label{cor:v32-increment}
For the same density-one family,
\[
 \sup_{x\in K}
 \|(\Pi_{Q_\ell,t}-\Pi_{q,t})\Xi_x\|_2^2
 \le
 t^{-M}\left(
 e^{-\delta_\alpha/(2t)}
 +e^{-c/\sqrt t}
 \right).
 \tag{V32.34}
\]
Thus adding a typical prime modulus
\(\ell=Y_t^\theta\), \(0<\theta<1/2\), captures asymptotically no new
phase energy beyond the conductor.
\end{corollary}

\begin{proof}
Apply \textup{(V32.4)} with \(F=\Xi_x\).  The increment norm is bounded by
\(\|\Pi_{Q_\ell,t}\Xi_x\|_2^2\), and
Theorem~\ref{thm:v32-wrong-way} supplies the result.
\end{proof}

The conclusion is uniform in \(x\in K\), so the auxiliary prime can be
chosen before locating a contact inside \(K\).  It is not a
contact-dependent selection.

\subsection{Implication for the contact regression}

At any depth or auxiliary modulus, the exact regression remains
\[
 \mathcal R_{\chi,t}(x)
 =
 \Im\left\{
 \mathbb E_t[
 (\Pi_{Q,t}U)(\Pi_{Q,t}\Xi_x)]
 +
 \mathbb E_t[
 (U-\Pi_{Q,t}U)(\Xi_x-\Pi_{Q,t}\Xi_x)]
 \right\}.
 \tag{V32.35}
\]
For the good moduli of Theorem~\ref{thm:v32-wrong-way}, the projected head
obeys
\[
 \sup_{x\in K}
 \left|
 \mathbb E_t[
 (\Pi_{Q_\ell,t}U)(\Pi_{Q_\ell,t}\Xi_x)]
 \right|
 \le
 t^{-M}\left(
 e^{-\delta_\alpha/(4t)}
 +e^{-c/(2\sqrt t)}
 \right).
 \tag{V32.36}
\]
After multiplication by \(Z_\chi(t)\), the second term on the right still
permits logarithmic rate \(1/4\); it is not the zero-rate contact estimate.
Meanwhile the residual phase norm tends to one and the residual
log-coordinate variance tends to \(1/(2t)\).  Its Cauchy budget therefore
has the exact critical rate \(1/4\), just as in V30--V31.

\begin{theorem}[Averaged distribution does not close the contact]
\label{thm:v32-average-no-contact}
Fix \(K\Subset\mathbb R\setminus\{0\}\) and
\(0<\theta<1/2\).  Conditioning on a density-one set of auxiliary prime
moduli \(q\ell\) with \(\ell\asymp Y_t^\theta\), and then using only:
\begin{enumerate}[label=\textup{(\roman*)},leftmargin=3em]
\item the averaged large-sieve energy;
\item conditional Pythagoras;
\item Cauchy--Schwarz on the residual tail; and
\item the real contact constraint,
\end{enumerate}
does not yield the zero-rate criterion of V28.  The projected head tends to
zero only in normalized probability scale, while the absolute tail budget
retains rate \(1/4\).
\end{theorem}

\begin{proof}
The head statement is \textup{(V32.36)}.  Equations
\textup{(V32.28)}--\textup{(V32.30)} give
\[
 Z_\chi(t)
 \|U-\Pi_{Q_\ell,t}U\|_2
 \|\Xi_x-\Pi_{Q_\ell,t}\Xi_x\|_2
 =
 \exp\!\left(
 \frac1{16t}+O(\log(1/t))
 \right),
 \tag{V32.37}
\]
uniformly on \(K\).  Its \(4t\)-logarithmic rate is \(1/4\).
The real contact constraint fixes no imaginary tail projection, by the same
capacity geometry as V28--V29.  Hence the listed absolute inputs do not
prove the contact criterion.
\end{proof}

This theorem does not say that every modulus below the square-root scale is
blind.  It says that blindness has relative density one in the auxiliary
prime family, and that the natural averaged theorem supplies no mechanism
for selecting a rare resolving modulus.  Nor does it address \(x=0\), where
\(\Xi_0=\chi(n)\) is already measurable modulo \(q\).

\begin{firewall}[The large sieve controls energy, not signed contact cancellation]
V32 proves a Gaussian multiplicative large-sieve estimate and translates it
into a density-one conditional-projection theorem for auxiliary prime
moduli below \(Y_t^{1/2}\).  The theorem makes the conditional head smaller
and the tail larger.  It does not bound the signed tail, does not control
rare exceptional moduli, and does not reach or cross the square-root
transition.  No off-line zero is excluded, and GRH remains open.
\end{firewall}

\begin{assessment}[V32 acquisition and the square-root transition]
The residue-filtration route now has two proved blind regimes.  Uniform
Siegel--Walfisz handles every polynomial conductor power, while the
multiplicative large sieve handles a density-one family of auxiliary prime
refinements \(q\ell\) for every
\(\ell=Y_t^\theta\) with \(0<\theta<1/2\).  In the second regime the new
martingale increment is exponentially small on the \(1/t\) scale.  Average
equidistribution therefore cannot be reinterpreted as phase resolution.

V33 should analyze the boundary \(\ell\asymp Y_t^{1/2}\), where the
prime-length and dual-modulus terms in the large sieve balance after
averaging and the saving in \textup{(V32.33)} disappears.  The first task is
an exact capacity example showing whether this loss is intrinsic to generic
large-sieve data.  If it is, the residue-energy route has reached its
natural uncertainty barrier and the programme should return upstream to a
signed bilinear relation between the value and derivative errors at a
contact, rather than seek another absolute modulus norm.
\end{assessment}

\section*{Version V32 append-only record}
\addcontentsline{toc}{section}{Version V32 append-only record}
The complete V31 source was retained before this continuation.  It had
\(486464\) bytes, \(14603\) lines, and SHA--256
\[
 \texttt{0183C1D078A4FDC33880D7A54DD35E0BFE6D7ECBF763BEEEEBCB636389132407}.
\]
V32 established the Gaussian multiplicative large-sieve energy bound,
Fourier-to-conditional conversion, density-one blindness for auxiliary
prime refinements below the square-root saddle scale, and the resulting
averaged-contact no-go theorem.

% =============================================================================
% END APPENDED V32 MATERIAL
% =============================================================================
% =============================================================================
% BEGIN APPENDED V33 MATERIAL
% =============================================================================

\section{V33: collision entropy and the sharp square-root capacity}
\label{sec:grh-v33}

V32 found that the Gaussian multiplicative large sieve loses its
exponential saving when the auxiliary modulus reaches
\[
 Q\asymp Y_t^{1/2},
 \qquad
 Y_t=e^{1/(4t)}.
 \tag{V33.1}
\]
This section determines whether that transition is intrinsic to the
absolute saddle data.

It is intrinsic to \(L^2\) energy.  The collision probability of the V28
saddle law is
\[
 \sum_n\nu_{\chi,t}(n)^2
 \sim\frac1{4\pi}e^{-1/(8t)},
 \tag{V33.2}
\]
so its inverse collision probability is
\[
 4\pi e^{1/(8t)}
 =4\pi Y_t^{1/2}.
 \tag{V33.3}
\]
This is precisely the large-sieve transition, including its exponential
scale.  A generic random-phase partition and a one-atom dual large-sieve
example both saturate it.

There is nevertheless a crucial saddle-specific distinction.  The
collision law lives at \(\log n\asymp t^{-1/2}\), far below the probability
mass saddle \(\log n\sim(4t)^{-1}\).  At the square-root modulus scale, the
atoms large enough to dominate their own residue cells carry exponentially
vanishing total probability.  Thus the large-sieve obstruction is a sharp
barrier for generic \(L^2\) information, but it is not proof that the actual
character phase becomes residue-measurable there.

\subsection{The squared von Mangoldt saddle}

Define the unnormalized collision mass
\[
 W_{\chi,2}(t)
 :=
 \sum_{\substack{n\ge2\\\chi(n)\ne0}}
 \frac{\Lambda(n)^2}{n}e^{-2t(\log n)^2}.
 \tag{V33.4}
\]

\begin{lemma}[Squared von Mangoldt Gaussian asymptotic]
\label{lem:v33-lambda-square}
As \(t\downarrow0\),
\[
 \boxed{
 W_{\chi,2}(t)=\frac1{4t}(1+o_\chi(1)).}
 \tag{V33.5}
\]
\end{lemma}

\begin{proof}
The prime number theorem and partial summation give
\[
 \sum_{n\le y}\Lambda(n)^2
 =y\log y\,(1+o(1)).
 \tag{V33.6}
\]
Indeed the prime contribution has this asymptotic, while prime powers
\(p^k\), \(k\ge2\), contribute \(O(y^{1/2}\log^2y)\).
Removing the powers of the finitely many primes dividing \(q\) changes
\textup{(V33.4)} by \(O_\chi(1)\).

Stieltjes partial summation and \(u=\log y\) therefore reduce the leading
term to
\[
 \int_0^\infty (u+1)e^{-2tu^2}\,du
 =
 \frac1{4t}+O(t^{-1/2}).
 \tag{V33.7}
\]
The \(o(y\log y)\) remainder is \(o(t^{-1})\) after splitting into a fixed
compact \(u\)-range, a scaled central range, and a Gaussian tail.  This
proves \textup{(V33.5)}.
\end{proof}

\begin{theorem}[Exact collision scale of the prime saddle]
\label{thm:v33-collision-scale}
Let
\[
 \mathfrak c_\chi(t)
 :=\sum_{n\in\Omega_\chi}\nu_{\chi,t}(n)^2,
 \qquad
 N_{\chi,2}(t):=\mathfrak c_\chi(t)^{-1}.
 \tag{V33.8}
\]
Then
\[
 \boxed{
 \mathfrak c_\chi(t)
 =\frac1{4\pi}e^{-1/(8t)}(1+o_\chi(1)),}
 \tag{V33.9}
\]
\[
 \boxed{
 N_{\chi,2}(t)
 =4\pi e^{1/(8t)}(1+o_\chi(1))
 =4\pi Y_t^{1/2}(1+o_\chi(1)).}
 \tag{V33.10}
\]
\end{theorem}

\begin{proof}
By definition,
\[
 \mathfrak c_\chi(t)=\frac{W_{\chi,2}(t)}{Z_\chi(t)^2}.
 \tag{V33.11}
\]
Insert Lemma~\ref{lem:v33-lambda-square} and the full-strip saddle mass
\[
 Z_\chi(t)^2
 =\frac\pi t e^{1/(8t)}(1+o_\chi(1))
 \tag{V33.12}
\]
from Proposition~\ref{prop:v28-saddle-moments}.  This proves both formulas.
\end{proof}

The normalized collision law is
\[
 \kappa_{\chi,t}(n)
 :=
 \frac{\nu_{\chi,t}(n)^2}{\mathfrak c_\chi(t)}
 =
 \frac{
 \Lambda(n)^2n^{-1}e^{-2t(\log n)^2}}
 {W_{\chi,2}(t)}.
 \tag{V33.13}
\]

\begin{theorem}[Rayleigh collision profile]
\label{thm:v33-rayleigh}
For every bounded continuous \(F:[0,\infty)\to\mathbb C\),
\[
 \boxed{
 \lim_{t\downarrow0}
 \mathbb E_{\kappa_{\chi,t}}
 F\!\left(\sqrt{2t}\,\log n\right)
 =
 \int_0^\infty F(v)\,2ve^{-v^2}\,dv.}
 \tag{V33.14}
\]
In particular, the collision mass is concentrated at
\(\log n\asymp t^{-1/2}\), whereas the original probability law is
concentrated at
\[
 \log n=\frac1{4t}+O_{\nu_{\chi,t}}(t^{-1/2}).
 \tag{V33.15}
\]
\end{theorem}

\begin{proof}
The proof of Lemma~\ref{lem:v33-lambda-square} reduces the collision
expectation to the normalized measure
\[
 \frac{u e^{-2tu^2}\,du}
 {\int_0^\infty u e^{-2tu^2}\,du}
 \tag{V33.16}
\]
up to an error tending to zero.  Put \(v=\sqrt{2t}\,u\).
Since the denominator is \(1/(4t)\), the transformed density is
\(2ve^{-v^2}\,dv\).  Equation \textup{(V33.15)} is the V28 saddle-moment
theorem.
\end{proof}

\subsection{The full Renyi support spectrum}

For \(p>1\), set
\[
 W_{\chi,p}(t)
 :=
 \sum_{\chi(n)\ne0}
 \frac{\Lambda(n)^p}{n^{p/2}}
 e^{-pt(\log n)^2},
 \qquad
 H_p(t):=
 \frac1{1-p}\log\sum_n\nu_{\chi,t}(n)^p.
 \tag{V33.17}
\]
Also let \(H_1(t)=-\sum_n\nu_{\chi,t}(n)\log\nu_{\chi,t}(n)\) and
\[
 H_\infty(t):=-\log\max_n\nu_{\chi,t}(n).
 \tag{V33.18}
\]

\begin{theorem}[Renyi phase diagram of the saddle]
\label{thm:v33-renyi}
For \(1<p<2\),
\[
 \log W_{\chi,p}(t)
 =
 \frac{(2-p)^2}{16pt}
 +O_{\chi,p}(\log(1/t)).
 \tag{V33.19}
\]
At \(p=2\), \(W_{\chi,2}(t)\sim(4t)^{-1}\), while for \(p>2\),
\[
 W_{\chi,p}(t)\longrightarrow
 \sum_{\chi(n)\ne0}\frac{\Lambda(n)^p}{n^{p/2}}
 \in(0,\infty).
 \tag{V33.20}
\]
Consequently,
\[
 \boxed{
 \lim_{t\downarrow0}4tH_p(t)
 =
 \begin{cases}
 1, &p=1,\\[2mm]
 p^{-1}, &1<p\le2,\\[2mm]
 \dfrac{p}{4(p-1)}, &2\le p<\infty,\\[3mm]
 \dfrac14, &p=\infty.
 \end{cases}}
 \tag{V33.21}
\]
The two formulas agree at \(p=2\).
\end{theorem}

\begin{proof}
For \(1<p<2\), the prime contribution to \(W_{\chi,p}\) has the Laplace
model
\[
 \int_0^\infty
 u^{p-1}
 \exp\!\left(
 \left(1-\frac p2\right)u-ptu^2
 \right)du.
 \tag{V33.22}
\]
Its center and exponential height are
\[
 u_p=\frac{2-p}{4pt},
 \qquad
 \frac{(2-p)^2}{16pt}.
 \tag{V33.23}
\]
The PNT, partial summation, and Gaussian localization show that prime
powers and all prefactors change only the logarithmic error, proving
\textup{(V33.19)}.  The case \(p=2\) is Lemma
\ref{lem:v33-lambda-square}.  For \(p>2\), the undamped series converges
absolutely and monotone convergence proves \textup{(V33.20)}.

Since
\[
 \sum_n\nu_{\chi,t}(n)^p
 =\frac{W_{\chi,p}(t)}{Z_\chi(t)^p},
 \tag{V33.24}
\]
and
\(\log Z_\chi(t)=1/(16t)+O(\log(1/t))\), equations
\textup{(V33.19)}--\textup{(V33.20)} give the middle two branches of
\textup{(V33.21)}.

For \(p=1\), use
\[
 \begin{aligned}
 H_1(t)
 ={}&\log Z_\chi(t)
 -\mathbb E_t\log\Lambda(n)
 +\frac12\mathbb E_tU+t\mathbb E_tU^2.
 \end{aligned}
 \tag{V33.25}
\]
The logarithmic von Mangoldt term is \(O_\chi(\log(1/t))\), while V28 gives
\[
 \mathbb E_tU=\frac1{4t}+o(t^{-1/2}),\qquad
 \mathbb E_tU^2=\frac1{16t^2}+O(t^{-1}).
 \tag{V33.26}
\]
Thus \(H_1(t)=1/(4t)+o(t^{-1})\).

Finally, the unnormalized atom weight
\(\Lambda(n)n^{-1/2}e^{-t(\log n)^2}\) is bounded above uniformly in small
\(t\), and a fixed prime gives a positive lower bound for its maximum.
Hence
\[
 H_\infty(t)=\log Z_\chi(t)+O_\chi(1)
 =\frac1{16t}+O_\chi(\log(1/t)),
 \tag{V33.27}
\]
which proves the endpoint.
\end{proof}

The support scale depends on the norm.  Shannon entropy sees \(Y_t\),
collision entropy sees \(Y_t^{1/2}\), and min-entropy sees
\(Y_t^{1/4}\).  Using one of these scales as though it controlled the
others loses the nonuniform atom geometry.

\subsection{Exact generic conditional-energy capacity}

The next lemma is independent of arithmetic.  It identifies exactly what
collision probability can and cannot control.

\begin{lemma}[Random-phase projection formula]
\label{lem:v33-random-phase}
Let \((p_n)\) be a probability mass function on a countable set, let
\(\mathcal P\) be a countable partition with cell masses
\[
 p_C:=\sum_{n\in C}p_n>0,
 \tag{V33.28}
\]
and let \(\varepsilon_n\) be independent uniform unit complex phases.
If \(\Pi_{\mathcal P}\) is conditional expectation onto the partition, then
\[
 \boxed{
 \mathbb E_\varepsilon
 \|\Pi_{\mathcal P}\varepsilon\|_{L^2(p)}^2
 =
 \sum_{C\in\mathcal P}
 \frac{\sum_{n\in C}p_n^2}{p_C}.}
 \tag{V33.29}
\]
Consequently,
\[
 \mathbb E_\varepsilon
 \|\varepsilon-\Pi_{\mathcal P}\varepsilon\|_{L^2(p)}^2
 =
 1-
 \sum_{C\in\mathcal P}
 \frac{\sum_{n\in C}p_n^2}{p_C}.
 \tag{V33.30}
\]
\end{lemma}

\begin{proof}
On a cell \(C\),
\[
 \Pi_{\mathcal P}\varepsilon
 =
 \frac1{p_C}\sum_{n\in C}p_n\varepsilon_n.
 \tag{V33.31}
\]
Independence and zero mean kill every off-diagonal term in the expected
squared modulus, leaving \(p_C^{-2}\sum_{n\in C}p_n^2\).  Multiplication by
the cell mass and summation prove \textup{(V33.29)}.  Conditional
Pythagoras and \(\|\varepsilon\|_2=1\) prove \textup{(V33.30)}.
\end{proof}

If a partition has \(Q\) cells and
\[
 \frac{c}{Q}\le p_C\le\frac{C}{Q}
 \tag{V33.32}
\]
for fixed \(0<c\le C<\infty\), then
\[
 \frac{Q}{C}\sum_np_n^2
 \le
 \mathbb E_\varepsilon
 \|\Pi_{\mathcal P}\varepsilon\|_2^2
 \le
 \frac{Q}{c}\sum_np_n^2.
 \tag{V33.33}
\]
Thus collision information predicts a transition at
\[
 Q\asymp\left(\sum_np_n^2\right)^{-1}.
 \tag{V33.34}
\]

\begin{proposition}[Sharp uniform capacity model]
\label{prop:v33-uniform-capacity}
Let \(p_n=1/N\) on \(N\) atoms, and partition them into \(Q\) equal cells,
where \(Q\mid N\).  Then
\[
 \boxed{
 \mathbb E_\varepsilon
 \|\Pi_{\mathcal P}\varepsilon\|_2^2
 =\frac QN.}
 \tag{V33.35}
\]
At \(Q=N\), the phase is completely resolved and the projection energy is
one.  Taking \(N=N_{\chi,2}(t)\) on the exponential scale shows that no
universal \(L^2\)-energy theorem based only on collision mass can force
blindness at \(Q\asymp Y_t^{1/2}\).
\end{proposition}

\begin{proof}
Every cell has mass \(1/Q\) and squared-atom mass \(1/(NQ)\).
Each summand in \textup{(V33.29)} equals \(1/N\), and there are \(Q\)
summands.
\end{proof}

This is a capacity model, not a model of the deterministic character phase.
Its purpose is to prove that the collision scale is the sharp stopping line
for generic conditional \(L^2\) data.

\subsection{Sharpness of the dual large-sieve term}

\begin{proposition}[One-atom saturation of the dual term]
\label{prop:v33-dual-sharpness}
Let \(R\to\infty\), choose an integer \(n_0\) coprime to every prime
\(\ell\in[R,2R]\), and take a coefficient sequence supported at \(n_0\)
with value \(A\).  Then
\[
 \sum_{\substack{R\le\ell\le2R\\\ell\ {\rm prime}}}
 \sum_{\substack{\vartheta\bmod\ell\\
 \vartheta\ne\vartheta_0}}
 \left|\sum_na_n\vartheta(n)\right|^2
 \asymp
 \frac{R^2}{\log R}|A|^2.
 \tag{V33.36}
\]
Hence the \(R^2\sum|a_n|^2\) term in the multiplicative large sieve cannot
be replaced by \(o(R^2)\sum|a_n|^2\) for arbitrary coefficients.
\end{proposition}

\begin{proof}
Every nonprincipal character in the sum has
\[
 \left|\sum_na_n\vartheta(n)\right|=|A|.
 \tag{V33.37}
\]
There are \(\ell-2\) such characters for each prime \(\ell\), and the prime
number theorem gives
\[
 \sum_{\substack{R\le\ell\le2R\\\ell\ {\rm prime}}}(\ell-2)
 \asymp\frac{R^2}{\log R}.
 \tag{V33.38}
\]
\end{proof}

Scaling \(|A|^2\) to the collision numerator
\(W_{\chi,2}(t)\asymp t^{-1}\) and taking
\(R\asymp N_{\chi,2}(t)\) makes the average dual energy per modulus
comparable to \(Z_\chi(t)^2\) on the exponential scale.  This is the exact
capacity mechanism behind the disappearance of the V32 saving at
\(Y_t^{1/2}\).  The one-atom vector is not asserted to equal the Gaussian
von Mangoldt coefficient sequence.

\subsection{Heavy atoms do not carry the saddle probability at the transition}

For a resolution scale \(Q(t)\), define the mass of atoms which can dominate
a uniform \(Q(t)\)-cell by
\[
 \mathfrak h_{\chi,t}(Q)
 :=
 \sum_{\substack{n\in\Omega_\chi\\
 \nu_{\chi,t}(n)\ge Q^{-1}}}
 \nu_{\chi,t}(n).
 \tag{V33.39}
\]

\begin{theorem}[Heavy-atom probability profile]
\label{thm:v33-heavy-profile}
Suppose
\[
 Q(t)=\exp\!\left(\frac{\alpha+o(1)}{t}\right).
 \tag{V33.40}
\]
If \(0<\alpha<1/16\), then the set in \textup{(V33.39)} is empty for all
sufficiently small \(t\).  If
\[
 \frac1{16}<\alpha<\frac14,
 \tag{V33.41}
\]
then
\[
 \boxed{
 \lim_{t\downarrow0}
 t\log\mathfrak h_{\chi,t}(Q)
 =
 -\left(\frac12-\sqrt\alpha\right)^2.}
 \tag{V33.42}
\]
In particular, at the collision and large-sieve transition
\(\alpha=1/8\),
\[
 \mathfrak h_{\chi,t}(Q)
 =
 \exp\!\left(
 -\frac{(1/2-1/(2\sqrt2))^2+o(1)}{t}
 \right)
 \longrightarrow0.
 \tag{V33.43}
\]
\end{theorem}

\begin{proof}
For \(u=\log n=c/t\), the logarithm of one prime atom has the exponential
scale
\[
 t\log\nu_{\chi,t}(n)
 =
 -\left(c^2+\frac c2+\frac1{16}\right)+o(1).
 \tag{V33.44}
\]
The factors \(\Lambda(n)\) and \(t^{-1/2}\) contribute only
\(o(t^{-1})\) to the logarithm.  The maximum unnormalized atom weight is
bounded above and below by positive constants, so if
\(\alpha<1/16\), even the largest normalized atom is smaller than \(Q^{-1}\)
for small \(t\).

For \(\alpha>1/16\), the upper logarithmic edge of the heavy set is
\[
 c_\alpha
 =
 \sqrt\alpha-\frac14,
 \tag{V33.45}
\]
the positive solution of
\[
 c^2+\frac c2+\frac1{16}=\alpha.
 \tag{V33.46}
\]
Since \(\alpha<1/4\), one has \(c_\alpha<1/4\), below the probability
saddle.  The PNT and endpoint Laplace asymptotics give
\[
 \sum_{\log n\le(c_\alpha+o(1))/t}
 \frac{\Lambda(n)}{\sqrt n}e^{-t(\log n)^2}
 =
 \exp\!\left(
 \frac{c_\alpha/2-c_\alpha^2+o(1)}{t}
 \right).
 \tag{V33.47}
\]
Division by
\(Z_\chi(t)=\exp((1/16+o(1))/t)\) yields
\[
 \begin{aligned}
 c_\alpha/2-c_\alpha^2-\frac1{16}
 &=
 \sqrt\alpha-\alpha-\frac14\\
 &=-\left(\frac12-\sqrt\alpha\right)^2,
 \end{aligned}
 \tag{V33.48}
\]
proving \textup{(V33.42)}.  Substitution of \(\alpha=1/8\) gives
\textup{(V33.43)}.
\end{proof}

The collision scale is therefore driven by atoms which are important after
squaring the probability but collectively negligible in ordinary
probability.  This is why the generic \(L^2\) capacity becomes critical at
\(Y_t^{1/2}\) without implying that the actual residue projection has
captured an order-one portion of the character phase.

\subsection{Consequences for the GRH contact route}

\begin{theorem}[The square-root transition is sharp but not sufficient]
\label{thm:v33-transition-verdict}
The following statements hold simultaneously.
\begin{enumerate}[label=\textup{(\roman*)},leftmargin=3em]
\item The exponential transition \(Q\asymp Y_t^{1/2}\) in the V32
      large-sieve bound equals the inverse collision probability of the
      absolute prime saddle.
\item Generic conditional \(L^2\) data cannot force blindness at that scale,
      by Proposition~\ref{prop:v33-uniform-capacity}.
\item Generic multiplicative large-sieve data cannot improve the dual
      \(Q^2\) term, by Proposition~\ref{prop:v33-dual-sharpness}.
\item For the actual saddle law, the atoms individually heavy at that scale
      carry exponentially vanishing probability, by
      Theorem~\ref{thm:v33-heavy-profile}.
\item None of these energy statements relates the real value constraint at
      a contact to the imaginary derivative regression.
\end{enumerate}
Therefore the square-root transition is an intrinsic stopping line for
generic \(L^2\) arguments, but neither a positive nor a negative resolution
of the signed GRH contact criterion.
\end{theorem}

\begin{proof}
Items (i)--(iv) are exactly Theorem~\ref{thm:v33-collision-scale},
Propositions~\ref{prop:v33-uniform-capacity} and
\ref{prop:v33-dual-sharpness}, and
Theorem~\ref{thm:v33-heavy-profile}.  Item (v) is the V28--V29 capacity
obstruction: the contact fixes one real scalar, while the derivative is a
signed imaginary projection.  The conclusions concern norms and entropy,
not that signed relation.
\end{proof}

\begin{firewall}[Collision entropy is not character cancellation]
V33 computes the exact collision constant, the collision-profile limit, the
full Renyi exponential spectrum, and the heavy-atom probability at the
large-sieve transition.  Its random-phase and one-atom examples are capacity
witnesses, not models of Dirichlet-character values.  No estimate for the
signed contact tail is obtained, no off-line zero is excluded, and GRH
remains open.
\end{firewall}

\begin{assessment}[V33 acquisition and the next saddle-specific test]
V33 proves that the V32 square-root boundary is the natural \(L^2\)
effective-support scale \(N_{\chi,2}\), not an arbitrary artifact of the
proof.  It also proves that the saddle is strongly nonuniform: its collision
law lies at logarithmic height \(t^{-1/2}\), its probability law lies at
height \((4t)^{-1}\), and its Renyi effective supports range from
\(Y_t^{1/4}\) to \(Y_t\).  At \(Q\asymp Y_t^{1/2}\), individually resolved
heavy atoms still have exponentially vanishing ordinary mass.

V34 should use this nonuniformity rather than another global \(L^2\) norm.
Split the saddle at the threshold \(\nu_{\chi,t}(n)=Q^{-1}\), derive an
exact heavy--light conditional-projection inequality, and insert
\textup{(V33.42)}.  The test is whether the light sector admits a
large-sieve or dispersion estimate beyond the collision transition after
the heavy sector is removed.  If the dual \(Q^2\) term is entirely carried
by the probability-light atoms, their removal may extend blindness; if not,
the residue-energy route is closed and the programme must return to the
signed value--derivative relation.
\end{assessment}

\section*{Version V33 append-only record}
\addcontentsline{toc}{section}{Version V33 append-only record}
The complete V32 source was retained before this continuation.  It had
\(503449\) bytes, \(15175\) lines, and SHA--256
\[
 \texttt{3201AA714AA330A2B3D0A8358EA5E51FB0C08D158E17A2E9E05D5C3175FFA8DA}.
\]
V33 established the exact collision scale and constant, Rayleigh collision
profile, Renyi support spectrum, sharp generic conditional and dual
large-sieve capacity examples, and the heavy-atom probability profile at
the square-root transition.

% =============================================================================
% END APPENDED V33 MATERIAL
% =============================================================================

% =============================================================================
% BEGIN APPENDED V34 MATERIAL
% =============================================================================

\section{V34: heavy--light dispersion beyond the collision scale}
\label{sec:grh-v34}

V33 isolated the obstruction in the global large-sieve norm.  At
\(R_t=\exp(\alpha/t)\), the dual term sees
\(R_t\sum_n\nu_{\chi,t}(n)^2\), which ceases to be small at
\(\alpha=1/8\).  But the collision mass is concentrated on atoms which
carry exponentially vanishing ordinary probability.  This section removes
those atoms before applying the large sieve.

The removal changes the stopping line.  If
\[
 \frac18\le\alpha<\frac14,
 \qquad R_t=e^{\alpha/t},
 \tag{V34.1}
\]
and
\[
 H_t:=\{n:\nu_{\chi,t}(n)\ge R_t^{-1}\},
 \qquad L_t:=\Omega_\chi\setminus H_t,
 \tag{V34.2}
\]
then both the heavy probability and the rescaled light collision mass have
the same exponential rate:
\[
 \nu_{\chi,t}(H_t)
 =\exp\!\left(-\frac{d_\alpha+o(1)}t\right),
 \qquad
 R_t\sum_{n\in L_t}\nu_{\chi,t}(n)^2
 =\exp\!\left(-\frac{d_\alpha+o(1)}t\right),
 \tag{V34.3}
\]
where
\[
 d_\alpha:=\left(\frac12-\sqrt\alpha\right)^2>0.
 \tag{V34.4}
\]
The first term pays for discarding the heavy atoms; the second is exactly
the averaged dual large-sieve cost of the light atoms.  Consequently, a
density-one family of prime moduli remains blind throughout the entire
range \(R_t=Y_t^{4\alpha}=o_{\exp}(Y_t)\), not merely below
\(Y_t^{1/2}\).

There is one typographical correction to the inherited V32 text.  Equation
\textup{(V32.4)} is the conditional Pythagoras identity and its right-hand
side must contain a plus sign:
\[
 \|\Pi_{Q,t}F\|_2^2
 =\|\Pi_{q,t}F\|_2^2
 +\|(\Pi_{Q,t}-\Pi_{q,t})F\|_2^2.
 \tag{V34.5}
\]
All uses of that identity in V32 have this meaning.

\subsection{An exact heavy--light projection inequality}

We first separate the probability geometry from the arithmetic.  Let
\((p_n)\) be a probability mass function, let \(\mathcal P\) be a finite
partition, and let \(\Pi_{\mathcal P}\) denote conditional expectation in
\(L^2(p)\).  Given a set \(H\), put \(L=H^c\),
\[
 h:=\sum_{n\in H}p_n,
 \qquad \lambda:=1-h,
 \qquad
 \ell_C:=\sum_{n\in C\cap L}p_n.
 \tag{V34.6}
\]

\begin{lemma}[Heavy--light conditional projection]
\label{lem:v34-heavy-light-projection}
For every \(f\in L^2(p)\),
\[
 \boxed{
 \|\Pi_{\mathcal P}f\|_2
 \le
 \|f\mathbf1_H\|_2
 \left(
 \sum_{C:\ell_C>0}
 \frac{\left|\sum_{n\in C\cap L}p_nf(n)\right|^2}{\ell_C}
 \right)^{1/2}.}
 \tag{V34.7}
\]
Suppose in addition that the cells form a finite abelian group \(G\) of
order \(J\), set \(m=\lambda/J\), and define
\[
 T_\psi^L(f):=\sum_{n\in L}p_nf(n)\psi(n),
 \qquad \psi\in\widehat G.
 \tag{V34.8}
\]
Then, for \(|f|\le1\),
\[
 \boxed{
 \sum_{C:\ell_C>0}
 \frac{\left|\sum_{n\in C\cap L}p_nf(n)\right|^2}{\ell_C}
 \le
 \frac2\lambda\sum_{\psi\in\widehat G}|T_\psi^L(f)|^2
 +\frac6\lambda
 \sum_{\substack{\psi\in\widehat G\\\psi\ne\psi_0}}
 |T_\psi^L(1)|^2.}
 \tag{V34.9}
\]
For general \(f\in L^4(p)\), the last term on the right of
\textup{(V34.9)} may be replaced by
\[
 \left(\mathbb E_p|f|^4\right)^{1/2}
 \left(
 \frac6\lambda
 \sum_{\psi\ne\psi_0}|T_\psi^L(1)|^2
 \right)^{1/2}.
 \tag{V34.10}
\]
\end{lemma}

\begin{proof}
Write \(p_C=\sum_{n\in C}p_n\).  Conditional expectation is a contraction,
so the heavy part contributes at most \(\|f\mathbf1_H\|_2\).  For the
light part, \(p_C\ge\ell_C\), and hence
\[
 \|\Pi_{\mathcal P}(f\mathbf1_L)\|_2^2
 =\sum_C
 \frac{\left|\sum_{n\in C\cap L}p_nf(n)\right|^2}{p_C}
 \le
 \sum_{C:\ell_C>0}
 \frac{\left|\sum_{n\in C\cap L}p_nf(n)\right|^2}{\ell_C}.
 \tag{V34.11}
\]
Minkowski proves \textup{(V34.7)}.

For the Fourier bound, write \(\ell_C=m+e_C\), and call a cell good when
\(|e_C|\le m/2\).  On good cells, \(\ell_C\ge m/2\).  On bad cells,
positivity and \(|e_C|>m/2\) give
\[
 \sum_{C\ {
m bad}}\ell_C
 \le\frac6m\sum_C|e_C|^2.
 \tag{V34.12}
\]
If \(|f|\le1\), the bad-cell quotient is at most \(\ell_C\).  Therefore
\[
 \sum_{C:\ell_C>0}
 \frac{\left|\sum_{n\in C\cap L}p_nf(n)\right|^2}{\ell_C}
 \le
 \frac2m\sum_C
 \left|\sum_{n\in C\cap L}p_nf(n)\right|^2
 +\frac6m\sum_C|e_C|^2.
 \tag{V34.13}
\]
Parseval on \(G\), together with \(m=\lambda/J\), is exactly
\textup{(V34.9)}.

For general \(f\), the bad-cell contribution is bounded by
\[
 \mathbb E_p\!\left[|f|^2\mathbf1_{L\cap B}\right]
 \le
 \left(\mathbb E_p|f|^4\right)^{1/2}
 \left(\sum_{C\ {
m bad}}\ell_C\right)^{1/2},
 \tag{V34.14}
\]
where \(B\) is the union of the bad cells.  Equations
\textup{(V34.12)} and Parseval prove \textup{(V34.10)}.
\end{proof}

The denominator in \textup{(V34.7)} is the light mass of a cell, even
though the original conditional expectation has the larger full-cell
denominator.  Thus heavy atoms can be discarded at their ordinary
\(L^2(p)\) cost, while only the light collision norm enters the Fourier
estimate.  This is the separation missing from the global V32 bound.

\subsection{The light collision tail}

Return to \(p_n=\nu_{\chi,t}(n)\) and the sets in
\textup{(V34.2)}.  Put
\[
 h_{\chi,t}(R_t):=\sum_{n\in H_t}\nu_{\chi,t}(n),
 \qquad
 \mathfrak c_{\chi,t}^{L}(R_t)
 :=\sum_{n\in L_t}\nu_{\chi,t}(n)^2.
 \tag{V34.15}
\]

\begin{theorem}[Matched heavy mass and light dual cost]
\label{thm:v34-matched-cost}
Fix \(1/8\le\alpha<1/4\), and let \(R_t=e^{\alpha/t}\).  Then
\[
 \boxed{
 \lim_{t\downarrow0}t\log h_{\chi,t}(R_t)
 =-d_\alpha,}
 \tag{V34.16}
\]
\[
 \boxed{
 \lim_{t\downarrow0}t\log
 \mathfrak c_{\chi,t}^{L}(R_t)
 =-\frac18-2\left(\sqrt\alpha-\frac14\right)^2,}
 \tag{V34.17}
\]
and consequently
\[
 \boxed{
 \lim_{t\downarrow0}t\log
 \left(R_t\mathfrak c_{\chi,t}^{L}(R_t)\right)
 =-d_\alpha.}
 \tag{V34.18}
\]
The same exponential rates hold after multiplying the summands in
\textup{(V34.15)} by any fixed power of \(1+\log n\).
\end{theorem}

\begin{proof}
Equation \textup{(V34.16)} is Theorem~\ref{thm:v33-heavy-profile}.  We
prove the new collision-tail formula.  At logarithmic height
\(u=\log n=c/t\), the atom threshold
\(\nu_{\chi,t}(n)<e^{-\alpha/t}\) has the exponential boundary
\[
 c>c_\alpha+o(1),
 \qquad
 c_\alpha:=\sqrt\alpha-\frac14,
 \tag{V34.19}
\]
because
\[
 t\log\nu_{\chi,t}(n)
 =-\left(c^2+\frac c2+\frac1{16}\right)+o(1).
 \tag{V34.20}
\]
The prime contribution to the unnormalized squared tail is, on the
exponential scale,
\[
 \int_{(c_\alpha+o(1))/t}^{\infty}
 u e^{-2tu^2}\,du
 =\exp\!\left(
 -\frac{2c_\alpha^2+o(1)}t
 \right).
 \tag{V34.21}
\]
The PNT gives matching upper and lower exponential bounds.  Prime powers
of exponent at least two have an additional negative linear exponential
factor and do not change the rate.  Division by
\[
 Z_\chi(t)^2
 =\exp\!\left(\frac{1/8+o(1)}t\right)
 \tag{V34.22}
\]
proves \textup{(V34.17)}.

Finally,
\[
 \begin{aligned}
 \alpha-\frac18-2c_\alpha^2
 &=\sqrt\alpha-\alpha-\frac14\\
 &=-\left(\frac12-\sqrt\alpha\right)^2
 =-d_\alpha,
 \end{aligned}
 \tag{V34.23}
\]
which proves \textup{(V34.18)}.  Polynomial logarithmic weights contribute
only \(o(t^{-1})\) to the logarithm.
\end{proof}

This identity is the decisive compensation.  The global collision norm
has rate \(e^{-1/(8t)}\), but after the atoms above \(R_t^{-1}\) are
removed, its dual cost still decays for every \(\alpha<1/4\).  Exactly the
same exponent pays for the removed probability.

\subsection{A truncated Gaussian large sieve}

Let \(\mathcal L_t\) be the auxiliary primes in \([R_t,2R_t]\) not
dividing \(q\).  For a fixed Dirichlet character \(\rho\bmod q\), an
integer \(r\ge0\), and \(x\in\mathbb R\), define
\[
 \mathcal T_{\rho,\eta,\vartheta,t}^{L,r}(x)
 :=
 \sum_{n\in L_t}
 \nu_{\chi,t}(n)\rho(n)\eta(n)\vartheta(n)
 (\log n-\mu_\chi(t))^r e^{ix\log n},
 \tag{V34.24}
\]
where \(\eta\bmod q\) and \(\vartheta\bmod\ell\).

\begin{theorem}[Light-sector Gaussian large sieve]
\label{thm:v34-light-large-sieve}
For fixed \(r\), fixed \(\rho\bmod q\), and compact \(K\subset\mathbb R\),
there is \(M<\infty\) such that
\[
 \boxed{
 \sum_{\ell\in\mathcal L_t}
 \sum_{\eta\bmod q}
 \sum_{\substack{\vartheta\bmod\ell\\
                  \vartheta\ne\vartheta_0}}
 \sup_{x\in K}
 \left|\mathcal T_{\rho,\eta,\vartheta,t}^{L,r}(x)\right|^2
 \le
 t^{-M}\left(
 1+R_t^2\mathfrak c_{\chi,t}^{L}(R_t)
 \right).}
 \tag{V34.25}
\]
Consequently, after division by
\(\#\mathcal L_t\asymp R_t/\log R_t\), the mean auxiliary-character
energy is
\[
 \boxed{
 \frac1{\#\mathcal L_t}
 \sum_{\ell\in\mathcal L_t}
 \sum_{\eta\bmod q}
 \sum_{\vartheta\ne\vartheta_0}
 \sup_{x\in K}
 \left|\mathcal T_{\rho,\eta,\vartheta,t}^{L,r}(x)\right|^2
 \le
 t^{-M}e^{-(d_\alpha+o(1))/t}.}
 \tag{V34.26}
\]
\end{theorem}

\begin{proof}
Apply the multiplicative large sieve on a dyadic block
\(N<n\le2N\) to the normalized coefficients in
\textup{(V34.24)}.  If \(E_N\) denotes their squared \(\ell^2\) mass on
the block, Minkowski over blocks gives two contributions bounded by
\[
 \left(\sum_N\sqrt{N E_N}\right)^2,
 \qquad
 R_t^2\left(\sum_N\sqrt{E_N}\right)^2.
 \tag{V34.27}
\]
The first is \(t^{-M}\): before the normalization by \(Z_\chi(t)^2\), its
Gaussian maximum is \(e^{1/(8t)}\), exactly canceled by the same
exponential factor in \(Z_\chi(t)^2\).  For the second, Cauchy--Schwarz
with a polynomial dyadic weight gives
\[
 \left(\sum_N\sqrt{E_N}\right)^2
 \le t^{-M}\mathfrak c_{\chi,t}^{L}(R_t),
 \tag{V34.28}
\]
where Theorem~\ref{thm:v34-matched-cost} absorbs the fixed powers of
\(\log n\).  There are only finitely many \(\eta\) and \(\rho\).

The compact supremum follows from the one-dimensional Sobolev inequality;
an \(x\)-derivative inserts one more logarithmic factor.  This proves
\textup{(V34.25)}.  Averaging, using Theorem
\ref{thm:v34-matched-cost}, gives
\[
 t^{-M}\left(R_t^{-1}
 +R_t\mathfrak c_{\chi,t}^{L}(R_t)\right)
 =t^{-M}e^{-(d_\alpha+o(1))/t},
 \tag{V34.29}
\]
because \(d_\alpha<\alpha\) in the range \textup{(V34.1)}.
\end{proof}

\subsection{Density-one blindness below the full saddle scale}

For \(\ell\in\mathcal L_t\), set \(Q_\ell=q\ell\).  The characters modulo
\(Q_\ell\) split as \(\eta\vartheta\) by the Chinese remainder theorem.
Theorem~\ref{thm:v34-light-large-sieve} controls all sectors with
\(\vartheta\ne\vartheta_0\).  In the finitely many sectors with
\(\vartheta=\vartheta_0\), the full fixed-\(q\) sums have the V32 PNT
bound.  Removing \(H_t\) changes an unweighted sum by at most
\(h_{\chi,t}(R_t)\), and changes a fixed logarithmic moment by at most a
polynomial in \(t^{-1}\) times that mass.  The nonunit atoms for
\(q\ell\) are exactly the powers \(n=\ell^k\).  Uniformly for
\(\ell\in\mathcal L_t\), their total probability, and each fixed
logarithmic moment of that probability, is at most
\[
 t^{-M}\exp\!\left(-\frac{(\alpha+1/4)^2+o(1)}t\right)
 =o\!\left(e^{-d_\alpha/t}\right).
 \tag{V34.30a}
\]
They may therefore be charged as an additional exceptional component in
Lemma~\ref{lem:v34-heavy-light-projection}.  Thus, uniformly for
\(\ell\in\mathcal L_t\),
\[
 \begin{aligned}
 &\sum_{\eta\bmod q}
 \sup_{x\in K}
 \left|\mathcal T_{\chi,\eta,\vartheta_0,t}^{L,0}(x)\right|^2
 \\
 &\qquad\le
 t^{-M}\left(e^{-c_K/\sqrt t}
 +h_{\chi,t}(R_t)^2\right),
 \qquad K\Subset\mathbb R\setminus\{0\},
 \end{aligned}
 \tag{V34.30}
\]
and the analogous estimates hold for the nonprincipal fixed sectors of
the light mass and of \(U-\mu_\chi(t)\).

\begin{theorem}[Heavy--light wrong-way theorem]
\label{thm:v34-wrong-way}
Fix \(1/8\le\alpha<1/4\) and
\(K\Subset\mathbb R\setminus\{0\}\).  There are subsets
\(\mathcal L_t^{\mathrm{HL}}\subseteq\mathcal L_t\) such that
\[
 \boxed{
 \frac{\#(\mathcal L_t\setminus\mathcal L_t^{\mathrm{HL}})}
 {\#\mathcal L_t}
 \le
 t^{-M}e^{-(d_\alpha+o(1))/(2t)}
 \longrightarrow0,}
 \tag{V34.31}
\]
and, uniformly for \(\ell\in\mathcal L_t^{\mathrm{HL}}\),
\[
 \boxed{
 \sup_{x\in K}
 \|\Pi_{Q_\ell,t}\Xi_x\|_2^2
 \le
 t^{-M}\left(
 e^{-(d_\alpha+o(1))/(2t)}+e^{-c_K/\sqrt t}
 \right)
 \longrightarrow0,}
 \tag{V34.32}
\]
\[
 \boxed{
 \|\Pi_{Q_\ell,t}U-\mu_\chi(t)\|_2^2
 \le
 t^{-M}\left(
 e^{-(d_\alpha+o(1))/(4t)}+e^{-c/\sqrt t}
 \right)
 \longrightarrow0.}
 \tag{V34.33}
\]
Consequently,
\[
 \sup_{x\in K}
 \|\Xi_x-\Pi_{Q_\ell,t}\Xi_x\|_2^2=1-o(1),
 \qquad
 \|U-\Pi_{Q_\ell,t}U\|_2^2
 =\frac1{2t}(1+o(1)).
 \tag{V34.34}
\]
Thus density-one residue blindness persists for every exponential modulus
strictly below the full saddle scale \(Y_t=e^{1/(4t)}\).
\end{theorem}

\begin{proof}
Apply Markov's inequality to \textup{(V34.26)}, simultaneously to the
finite collection of phase, mass, and first-moment readers required in
Lemma~\ref{lem:v34-heavy-light-projection}.  With threshold
\(e^{-(d_\alpha+o(1))/(2t)}\), this gives
\textup{(V34.31)} and the corresponding pointwise auxiliary Fourier
energies.

For \(f=\Xi_x\), use \(|f|=1\) in
Lemma~\ref{lem:v34-heavy-light-projection}, charging the nonunit
\(\ell\)-power atoms by \textup{(V34.30a)}.  The heavy term has squared
norm exactly \(h_{\chi,t}(R_t)\).  Equations
\textup{(V34.9)}, \textup{(V34.16)}, and \textup{(V34.30)} then give
\textup{(V34.32)}.

For \(f=U-\mu_\chi(t)\), use \textup{(V34.10)}.  The full fourth moment is
polynomial in \(t^{-1}\).  Moreover, \(H_t\) lies below
\((c_\alpha+o(1))/t\) on the exponential scale, so
\[
 \mathbb E_t\!\left[
 |U-\mu_\chi(t)|^2\mathbf1_{H_t}
 \right]
 \le t^{-M}h_{\chi,t}(R_t).
 \tag{V34.35}
\]
The square root in the bad-cell term \textup{(V34.10)} halves the Markov
exponent once more, proving \textup{(V34.33)}.  Conditional Pythagoras,
\(|\Xi_x|=1\), and
\(\operatorname{Var}_t(U)=(2t)^{-1}(1+o(1))\) prove
\textup{(V34.34)}.
\end{proof}

\begin{corollary}[The generic residue-energy route stops at Shannon scale]
\label{cor:v34-shannon-stopping-line}
The square-root scale \(Y_t^{1/2}\) is the stopping line only for an
unsplit global collision-norm argument.  After the exact heavy--light
decomposition, a density-one family of auxiliary prime refinements is
still asymptotically blind at every scale
\[
 R_t=Y_t^\theta,
 \qquad \frac12\le\theta<1.
 \tag{V34.36}
\]
The saving degenerates precisely as \(\theta\uparrow1\).  Therefore a
generic residue-energy strategy can acquire order-one phase information
only from exceptional moduli or in the critical Shannon window
\(\log R_t=1/(4t)+O(t^{-1/2})\).
\end{corollary}

\begin{proof}
Put \(\theta=4\alpha\) in
Theorem~\ref{thm:v34-wrong-way}.  Then \(1/8\le\alpha<1/4\) is equivalent to
\(1/2\le\theta<1\), and
\(d_\alpha\downarrow0\) as \(\alpha\uparrow1/4\).  The final statement is
the exhaustive alternative left by the density-one theorem; it does not
assert that either exceptional or critical moduli actually resolve the
phase.
\end{proof}

\subsection{Consequences for the signed contact programme}

V34 improves the negative structural information.  The collision barrier
was not the genuine endpoint of residue blindness; it was caused by a
small-probability population which is large only after squaring atom
weights.  Once that population is paid for in ordinary probability, the
light-sector dual cost stays small until the entropy scale itself.

This does not make the conditional head useful for GRH.  On the contrary,
for density-one refinements below \(Y_t\), both the phase reader and the
centered logarithmic reader remain negligible.  The unresolved orthogonal
tail retains asymptotically all phase energy and all variance.  The real
contact equation still supplies no relation forcing its signed imaginary
component to vanish.

\begin{firewall}[Extended blindness is not signed cancellation]
V34 proves a heavy--light projection inequality, a sharp exponential
formula for the light collision tail, a truncated Gaussian large sieve,
and density-one residue blindness up to the full saddle scale.  These are
absolute \(L^2\) and density-one statements.  They do not control a sparse
exceptional modulus, do not identify the phase in the Shannon critical
window, and do not relate the real contact value to the imaginary
derivative.  No off-line zero is excluded and GRH remains open.
\end{firewall}

\begin{assessment}[V34 acquisition and the critical-window test]
V34 moves the intrinsic generic stopping line from collision scale
\(Y_t^{1/2}\) to Shannon scale \(Y_t\).  The gain is exact on the
exponential scale because the discarded heavy probability and the
remaining light dual cost share the rate
\(d_\alpha=(1/2-\sqrt\alpha)^2\).  This closes the proposed V33 test in the
positive direction but makes the residue-energy route even less favorable:
almost every sub-Shannon refinement is blind.

V35 should resolve the critical window rather than repeat a fixed-exponent
large-sieve estimate.  Set
\[
 \log R_t=\frac1{4t}+\frac{s}{\sqrt t}+O(\log(1/t)),
 \tag{V34.37}
\]
retain the polynomial prefactor in the atom threshold, and compute the
Gaussian transition profiles of both the heavy ordinary mass and the light
dual collision cost.  The resulting profile will decide whether critical
residue partitions have a nontrivial limiting capacity or jump directly
from blindness to resolution.  If only a generic capacity law results,
the residue-energy branch should be closed and the next work should return
upstream to a signed value--derivative identity.
\end{assessment}

\section*{Version V34 append-only record}
\addcontentsline{toc}{section}{Version V34 append-only record}
The complete V33 source was retained before this continuation.  It had
\(521349\) bytes, \(15801\) lines, and SHA--256
\[
 \texttt{8C2B14F99CC087BCD375E20148D274953EC52E58E40BE82699562136AF55B19A}.
\]
V34 established the exact heavy--light conditional-projection inequality,
the matched heavy-probability and light-collision exponents, the truncated
Gaussian large sieve, and density-one blindness throughout the range
\(Y_t^{1/2}\le R_t<Y_t\).

% =============================================================================
% END APPENDED V34 MATERIAL
% =============================================================================


% =============================================================================
% BEGIN APPENDED V35 MATERIAL
% =============================================================================

\section{V35: the Shannon critical window and its Gaussian capacity envelope}
\label{sec:grh-v35}

V34 proved density-one blindness at every fixed exponential scale below
\(Y_t=e^{1/(4t)}\).  At the endpoint, the exponential rate
\(d_\alpha=(1/2-\sqrt\alpha)^2\) vanishes, so the polynomial normalization
of an individual saddle atom becomes decisive.  This section retains that
normalization and resolves the critical window.

The transition is not a jump from blindness to complete resolution.  A
modulus calibrated a bounded number of saddle standard deviations from the
center leaves a Gaussian fraction \(\Phi(s)\) of the probability in atoms
which must be paid for directly.  After those atoms are removed with a
small polynomial slack, the light-sector large-sieve energy still tends to
zero.  Thus the surviving generic obstruction has the continuous profile
\(\Phi(s)\).  If the window center drifts to the left,
\(s_t\to-\infty\), density-one blindness persists even though the modulus
is no longer separated from \(Y_t\) on the \(1/t\) exponential scale.

\subsection{Exact calibration of a critical atom}

Write
\[
 u_0(t):=\frac1{4t},
 \qquad
 u_t(s):=u_0(t)+\frac{s}{\sqrt{2t}},
 \tag{V35.1}
\]
and let
\[
 \phi(s):=\frac1{\sqrt{2\pi}}e^{-s^2/2},
 \qquad
 \Phi(s):=\int_{-\infty}^s\phi(z)\,dz.
 \tag{V35.2}
\]
The continuous prime atom at height \(u\) has weight
\[
 \mathfrak a_{\chi,t}(u)
 :=\frac{u e^{-u/2-tu^2}}{Z_\chi(t)}.
 \tag{V35.3}
\]
For fixed \(s\in\mathbb R\), calibrate the modulus scale by
\[
 R_{\chi,t}(s)
 :=\mathfrak a_{\chi,t}(u_t(s))^{-1}
 =\frac{Z_\chi(t)}{u_t(s)}
 e^{u_t(s)/2+t u_t(s)^2}.
 \tag{V35.4}
\]

\begin{lemma}[Critical modulus expansion]
\label{lem:v35-critical-expansion}
For fixed \(s\),
\[
 \boxed{
 \log R_{\chi,t}(s)
 =\frac1{4t}+\frac{s}{\sqrt{2t}}
 +\frac{s^2}{2}+\log\!\bigl(4\sqrt{\pi t}\bigr)+o_\chi(1).}
 \tag{V35.5}
\]
In particular,
\[
 \log R_{\chi,t}(s)
 =\frac1{4t}+\frac{s}{\sqrt{2t}}+O_s(\log(1/t)).
 \tag{V35.6}
\]
\end{lemma}

\begin{proof}
The V28 saddle asymptotic is
\[
 Z_\chi(t)
 =\sqrt{\frac\pi t}\,e^{1/(16t)}(1+o_\chi(1)).
 \tag{V35.7}
\]
Substitute \(u_t(s)=1/(4t)+s/\sqrt{2t}\) into
\textup{(V35.4)}.  The exponential part is
\[
 \frac1{16t}+\frac{u_t(s)}2+t u_t(s)^2
 =\frac1{4t}+\frac{s}{\sqrt{2t}}+\frac{s^2}{2},
 \tag{V35.8}
\]
while
\[
 \frac12\log\frac\pi t-\log u_t(s)
 =\log\!\bigl(4\sqrt{\pi t}\bigr)+o_s(1).
 \tag{V35.9}
\]
This proves the claim.
\end{proof}

Fix \(A>1/2\).  The polynomially relaxed heavy--light split is
\[
 H_{t}^{A}(s)
 :=\left\{n\in\Omega_\chi:
 \nu_{\chi,t}(n)\ge\frac{t^A}{R_{\chi,t}(s)}\right\},
 \qquad
 L_t^{A}(s):=\Omega_\chi\setminus H_t^{A}(s).
 \tag{V35.10}
\]
Let
\[
 h_{\chi,t}^{A}(s):=\nu_{\chi,t}(H_t^{A}(s)),
 \qquad
 \mathfrak c_{\chi,t}^{A,L}(s)
 :=\sum_{n\in L_t^{A}(s)}\nu_{\chi,t}(n)^2.
 \tag{V35.11}
\]
The slack \(t^A\) enlarges the heavy set only by
\(O(\log(1/t))\) in logarithmic height, which is
\(o(t^{-1/2})\) on the probability-saddle scale.

\begin{theorem}[Gaussian heavy profile and critical light cost]
\label{thm:v35-critical-profile}
For fixed \(s\in\mathbb R\) and fixed \(A>1/2\),
\[
 \boxed{
 h_{\chi,t}^{A}(s)\longrightarrow\Phi(s).}
 \tag{V35.12}
\]
If
\[
 G_t(n):=\sqrt{2t}\,\bigl(\log n-\mu_\chi(t)\bigr),
 \tag{V35.13}
\]
then
\[
 \boxed{
 \mathbb E_t\!\left[G_t^2\mathbf1_{H_t^{A}(s)}\right]
 \longrightarrow
 M_2(s):=\int_{-\infty}^s z^2\phi(z)\,dz
 =\Phi(s)-s\phi(s).}
 \tag{V35.14}
\]
Moreover,
\[
 \boxed{
 R_{\chi,t}(s)\,
 \mathfrak c_{\chi,t}^{A,L}(s)
 \sim_\chi
 \frac{t^{A+1/2}}{\sqrt\pi}e^{-s^2/2},}
 \tag{V35.15}
\]
and therefore
\[
 \boxed{
 R_{\chi,t}(s)\log R_{\chi,t}(s)\,
 \mathfrak c_{\chi,t}^{A,L}(s)
 \sim_\chi
 \frac{t^{A-1/2}}{4\sqrt\pi}e^{-s^2/2}
 \longrightarrow0.}
 \tag{V35.16}
\]
The analog of \textup{(V35.16)} with an extra factor \(G_t(n)^2\) in
the collision sum is
\(O_{\chi,s}(t^{A-1/2}e^{-s^2/2})\).
\end{theorem}

\begin{proof}
For primes, the function
\(u\mapsto u e^{-u/2-tu^2}\) is strictly decreasing throughout the
critical region.  Let \(v_t=v_t(s,A)\) be the unique critical-region
solution of
\[
 \mathfrak a_{\chi,t}(v_t)
 =t^A\mathfrak a_{\chi,t}(u_t(s)).
 \tag{V35.17}
\]
Logarithmic differentiation gives
\[
 v_t-u_t(s)=A\log(1/t)+o(\log(1/t)),
 \qquad
 \sqrt{2t}\,(v_t-u_0(t))=s+o(1).
 \tag{V35.18}
\]

The PNT and the change of variables
\(z=\sqrt{2t}(u-u_0(t))\) give the local saddle law
\[
 \frac1{Z_\chi(t)}e^{u/2-tu^2}\,du
 =\bigl(1+o_\chi(1)\bigr)\phi(z)\,dz
 \tag{V35.19}
\]
on bounded \(z\)-sets.  Prime powers of exponent at least two have
exponentially smaller total mass there.  Equations
\textup{(V35.18)}--\textup{(V35.19)} prove
\textup{(V35.12)}.  The same argument with the standardized centered
square gives \textup{(V35.14)}; replacing \(u_0(t)\) by
\(\mu_\chi(t)=u_0(t)+o(t^{-1/2})\) does not change the limit.

For the light collision tail, the prime contribution has the sharp
endpoint asymptotic
\[
 \mathfrak c_{\chi,t}^{A,L}(s)
 \sim_\chi
 \frac1{4tZ_\chi(t)^2}e^{-2t v_t^2}.
 \tag{V35.20}
\]
Again prime powers are exponentially smaller.  By
\textup{(V35.17)},
\[
 R_{\chi,t}(s)
 =\frac{t^A}{\mathfrak a_{\chi,t}(v_t)}.
 \tag{V35.21}
\]
Multiplying \textup{(V35.20)} by \textup{(V35.21)} yields
\[
 R_{\chi,t}(s)\mathfrak c_{\chi,t}^{A,L}(s)
 \sim_\chi
 \frac{t^A}{4t v_t Z_\chi(t)}e^{v_t/2-tv_t^2}.
 \tag{V35.22}
\]
Here \(4tv_t\to1\), and completing the square in the last quotient gives
\[
 \frac1{Z_\chi(t)}e^{v_t/2-tv_t^2}
 \sim_\chi
 \sqrt{\frac t\pi}e^{-s^2/2}.
 \tag{V35.23}
\]
This proves \textup{(V35.15)}.  Lemma
\ref{lem:v35-critical-expansion} proves \textup{(V35.16)}.  On the
collision endpoint layer, \(G_t=s+o_s(1)\), up to a geometrically
decaying tail, which proves the final assertion.
\end{proof}

The logarithm in \textup{(V35.16)} is important.  V34 could absorb it
into an unspecified power of \(t^{-1}\) because it had an exponential
saving.  In the critical window it is exactly the loss from averaging only
over prime moduli.  The condition \(A>1/2\) is the sharp requirement of
this polynomially relaxed prime-modulus argument.

\subsection{A polynomially sharp light-sector large sieve}

Let
\[
 \mathcal L_t(s)
 :=\{\ell\text{ prime}:R_{\chi,t}(s)\le\ell\le
 2R_{\chi,t}(s),\ \ell\nmid q\}.
 \tag{V35.24}
\]
For a fixed character \(\rho\bmod q\), take either of the readers
\[
 B_{\rho,x,t}^{(0)}(n):=\rho(n)e^{ix\log n},
 \qquad
 B_{\rho,t}^{(1)}(n):=\rho(n)G_t(n),
 \tag{V35.25}
\]
where \(x\) is fixed.  Define
\[
 \mathcal U_{\rho,\eta,\vartheta,t}^{A,j}
 :=\sum_{n\in L_t^A(s)}
 \nu_{\chi,t}(n)B_{\rho,t}^{(j)}(n)
 \eta(n)\vartheta(n),
 \tag{V35.26}
\]
with the understood \(x\)-dependence when \(j=0\).

\begin{theorem}[Critical light-sector large sieve]
\label{thm:v35-critical-large-sieve}
Fix \(s\in\mathbb R\), \(A>1/2\), \(j\in\{0,1\}\), a fixed reader in
\textup{(V35.25)}, and a fixed \(\rho\bmod q\).  Then
\[
 \boxed{
 \begin{aligned}
 &\frac1{\#\mathcal L_t(s)}
 \sum_{\ell\in\mathcal L_t(s)}
 \sum_{\eta\bmod q}
 \sum_{\substack{\vartheta\bmod\ell\\
                  \vartheta\ne\vartheta_0}}
 \left|\mathcal U_{\rho,\eta,\vartheta,t}^{A,j}\right|^2\\
 &\qquad\ll_{\chi,s,A,\rho,x}
 e^{-c/t}
 +(1+s^2)t^{A-1/2}e^{-s^2/2}.
 \end{aligned}}
 \tag{V35.27}
\]
In particular, the mean auxiliary Fourier energy tends to zero.
\end{theorem}

\begin{proof}
Repeat the dyadic large-sieve proof of Theorem
\ref{thm:v34-light-large-sieve}, but retain the polynomial factors.  If
\(E_N\) is the squared coefficient mass on \(N<n\le2N\), Minkowski gives
\[
 \left(\sum_N\sqrt{NE_N}\right)^2
 +R_{\chi,t}(s)^2
 \left(\sum_N\sqrt{E_N}\right)^2.
 \tag{V35.28}
\]
The first term is at most a fixed power of \(t^{-1}\).  After division by
\[
 \#\mathcal L_t(s)
 \sim\frac{R_{\chi,t}(s)}{\log R_{\chi,t}(s)},
 \tag{V35.29}
\]
it is \(O(e^{-c/t})\).

The light collision blocks begin at \(v_t\asymp t^{-1}\).  Successive
dyadic blocks then decay geometrically, since the logarithmic derivative
of \(e^{-2tu^2}\) is bounded above by a negative constant.  Hence
\[
 \left(\sum_N\sqrt{E_N}\right)^2
 \ll_{s,A}(1+s^2)
 \mathfrak c_{\chi,t}^{A,L}(s),
 \tag{V35.30}
\]
with the factor \(1+s^2\) only needed for the standardized-score reader.
The averaged dual term is therefore bounded by
\[
 (1+s^2)R_{\chi,t}(s)log R_{\chi,t}(s)
 \mathfrak c_{\chi,t}^{A,L}(s).
 \tag{V35.31}
\]
Theorem~\ref{thm:v35-critical-profile} proves
\textup{(V35.27)}.
\end{proof}

Unlike V34, this theorem is pointwise in \(x\).  Taking a compact
supremum by a raw Sobolev estimate inserts unnecessary powers of
\(t^{-1}\) and destroys the critical polynomial saving.  Pointwise control
is the relevant statement for the fixed ordinate of any proposed contact.

\subsection{The critical capacity envelope}

The light auxiliary Fourier sectors now vanish in density, but the heavy
mass no longer does.  The next theorem states exactly what the
heavy--light method retains.

\begin{theorem}[Gaussian envelope for critical residue projections]
\label{thm:v35-critical-envelope}
Fix \(s\in\mathbb R\), \(A>1/2\), and \(x\ne0\).  Put
\[
 \kappa_A:=\frac12\left(A-\frac12\right)>0.
 \tag{V35.32}
\]
There are subsets
\(\mathcal L_t^{\mathrm{crit}}(s)\subseteq\mathcal L_t(s)\) such that
\[
 \frac{\#(\mathcal L_t(s)\setminus
 \mathcal L_t^{\mathrm{crit}}(s))}
 {\#\mathcal L_t(s)}
 =O_{\chi,s,A,x}(t^{\kappa_A}),
 \tag{V35.33}
\]
and, for \(Q_\ell=q\ell\),
\[
 \boxed{
 \limsup_{t\downarrow0}
 \sup_{\ell\in\mathcal L_t^{\mathrm{crit}}(s)}
 \|\Pi_{Q_\ell,t}\Xi_x\|_2^2
 \le
 \min\!\left\{1,
 C_q\frac{\Phi(s)}{1-\Phi(s)}\right\}.}
 \tag{V35.34}
\]
Likewise,
\[
 \boxed{
 \limsup_{t\downarrow0}
 \sup_{\ell\in\mathcal L_t^{\mathrm{crit}}(s)}
 \|\Pi_{Q_\ell,t}G_t\|_2^2
 \le
 \min\!\left\{1,
 C_q\frac{\Phi(s)+M_2(s)}{1-\Phi(s)}\right\}.}
 \tag{V35.35}
\]
Here \(C_q<\infty\) depends only on the fixed base modulus.  The bounds
are capacity envelopes, not asserted limits for the deterministic
character phase.
\end{theorem}

\begin{proof}
Apply Markov's inequality to Theorem
\ref{thm:v35-critical-large-sieve} with threshold \(t^{\kappa_A}\), for
the finite set of phase, mass, and standardized-score readers needed by
Lemma~\ref{lem:v34-heavy-light-projection}.  This gives
\textup{(V35.33)} and makes every auxiliary light Fourier sector
\(o(1)\) on the retained family.

The finitely many sectors with
\(\vartheta=\vartheta_0\) are fixed-\(q\) readers.  For the full saddle
they tend to zero by the fixed-character PNT, and for \(x\ne0\) by the
Gaussian Fourier decay in the principal sector.  Restriction to the light
set changes an unweighted reader by at most
\(h_{\chi,t}^{A}(s)\).  For the standardized score, Cauchy--Schwarz bounds
the change by the square root of the product of
\(h_{\chi,t}^{A}(s)\) and its heavy second moment.  The nonunit
\(\ell\)-power atoms have exponentially vanishing mass and may be charged
separately as in V34.

Insert these estimates into Lemma
\ref{lem:v34-heavy-light-projection}.  Since the light mass tends to
\(1-\Phi(s)>0\), the phase heavy norm, the fixed-sector correction, and
the bad-cell mass are bounded by a constant multiple of
\(\Phi(s)/(1-\Phi(s))\).  Equations \textup{(V35.12)} and
\textup{(V35.14)} give the corresponding standardized-score bound.
Finally use that conditional expectation is a contraction to cap both
bounds by one.
\end{proof}

The envelope is continuous in \(s\).  It tends to zero as
\(s\to-\infty\), while the heavy population becomes order one for bounded
or positive \(s\).  Thus the generic method does not exhibit a sudden
resolution threshold.  Its light Fourier information remains negligible;
the only surviving obstruction is the Gaussian population of individually
large atoms.

\subsection{Blindness in the moderate precritical window}

The fixed-\(s\) theorem has a uniform drifting version.  The restriction
below ensures that the classical PNT error remains smaller than the
Gaussian moderate-deviation tail.

\begin{theorem}[Moderate-window density-one blindness]
\label{thm:v35-moderate-blindness}
Let \(s_t\to-\infty\) with
\[
 |s_t|=o(t^{-1/4}),
 \tag{V35.36}
\]
fix \(A>1/2\), and define \(R_{\chi,t}(s_t)\) by
\textup{(V35.4)}.  For each fixed \(x\ne0\), there are subsets
\[
 \mathcal L_t^{\mathrm{mod}}
 \subseteq\mathcal L_t(s_t),
 \qquad
 \frac{\#\mathcal L_t^{\mathrm{mod}}}
 {\#\mathcal L_t(s_t)}\longrightarrow1,
 \tag{V35.37}
\]
such that, uniformly for \(\ell\in\mathcal L_t^{\mathrm{mod}}\),
\[
 \boxed{
 \|\Pi_{q\ell,t}\Xi_x\|_2^2\longrightarrow0,
 \qquad
 \|\Pi_{q\ell,t}G_t\|_2^2\longrightarrow0.}
 \tag{V35.38}
\]
Equivalently,
\[
 \|\Xi_x-\Pi_{q\ell,t}\Xi_x\|_2^2=1-o(1),
 \qquad
 \|U-\Pi_{q\ell,t}U\|_2^2
 =\frac1{2t}(1+o(1)).
 \tag{V35.39}
\]
The modulus satisfies
\[
 \log R_{\chi,t}(s_t)
 =\frac1{4t}-\frac{|s_t|}{\sqrt{2t}}
 +O(s_t^2+\log(1/t)),
 \tag{V35.40}
\]
so this is a shrinking, moderate-deviation gap below the Shannon scale,
not a fixed exponential gap.
\end{theorem}

\begin{proof}
The Gaussian saddle calculation is uniform under
\textup{(V35.36)}.  Indeed
\(e^{-s_t^2/2}=e^{-o(t^{-1/2})}\), while the classical PNT remainder at
logarithmic height \(u\asymp t^{-1}\) has an
\(e^{-c/\sqrt t}\) saving.  Hence
\[
 h_{\chi,t}^{A}(s_t)
 \sim\Phi(s_t)\longrightarrow0,
 \qquad
 \mathbb E_t[G_t^2\mathbf1_{H_t^A(s_t)}]
 \sim M_2(s_t)\longrightarrow0.
 \tag{V35.41}
\]
The averaged light Fourier energy in \textup{(V35.27)} also tends to zero:
every fixed polynomial in \(|s_t|\), multiplied by
\(e^{-s_t^2/2}\), tends to zero along the negative tail, and the remaining
factor \(t^{A-1/2}\) tends to zero.  Choose a Markov threshold tending to zero
more slowly than this mean energy.  The proof of Theorem
\ref{thm:v35-critical-envelope}, now with the quantities in
\textup{(V35.41)} tending to zero, proves \textup{(V35.37)}--
\textup{(V35.38)}.  Conditional Pythagoras and
\(\operatorname{Var}_t(U)=(2t)^{-1}(1+o(1))\) give
\textup{(V35.39)}.  Lemma~\ref{lem:v35-critical-expansion} gives
\textup{(V35.40)}.
\end{proof}

\subsection{Consequences for the contact problem}

V35 exhausts what the generic residue-energy mechanism can say at the
Shannon boundary.  The light sector remains invisible even in the
critical window.  The obstruction is now localized: it is the signed
contribution of the prime atoms in the Gaussian prefix
\[
 \sqrt{2t}\left(\log n-\frac1{4t}\right)\le s+o(1).
 \tag{V35.42}
\]
Absolute energy cannot decide whether that prefix reinforces or cancels
the light tail at a contact.

\begin{firewall}[A Gaussian capacity envelope is not a phase limit]
V35 proves the critical normalization, Gaussian heavy-mass and score
profiles, a polynomially sharp pointwise large sieve, and
moderate-window density-one blindness.  The bounds
\textup{(V35.34)}--\textup{(V35.35)} are upper capacity envelopes only.
They do not prove that residue projections converge to those envelopes,
do not control exceptional moduli, and do not impose the signed
value--derivative relation required at a contact.  No off-line zero is
excluded and GRH remains open.
\end{firewall}

\begin{assessment}[V35 acquisition and the required upstream return]
The V34 stopping line has now been resolved through its central-limit
window.  Prime-modulus density costs require the polynomial slack
\(A>1/2\), after which the light dual energy vanishes like
\(t^{A-1/2}e^{-s^2/2}\).  The heavy ordinary mass converges to
\(\Phi(s)\), and its standardized derivative energy converges to
\(\Phi(s)-s\phi(s)\).  Thus the generic transition is a Gaussian envelope,
not a discontinuous acquisition of the character phase.

The residue-energy branch should now go upstream.  V36 should define the
signed critical-prefix readers obtained by truncating the prime saddle at
\textup{(V35.42)}, derive their exact value and logarithmic-derivative
identities by Gaussian Abel summation, and test whether the real contact
condition constrains the imaginary first moment of that prefix.  A purely
norm-based refinement would only reproduce the envelope already computed
here.
\end{assessment}

\section*{Version V35 append-only record}
\addcontentsline{toc}{section}{Version V35 append-only record}
The complete V34 source was retained before this continuation.  It had
\(539308\) bytes, \(16367\) lines, and SHA--256
\[
 \texttt{C300DC47A290E402E6BF1CE2D85C3B1856383EC35B21E2C4440BDFC553C723EE}.
\]
V35 established the exact Shannon-window calibration, Gaussian heavy-mass
and standardized-score profiles, the polynomially sharp critical
light-sector large sieve, the continuous capacity envelope, and
density-one blindness in the moderate precritical window.

% =============================================================================
% END APPENDED V35 MATERIAL
% =============================================================================


% =============================================================================
% BEGIN APPENDED V36 MATERIAL
% =============================================================================

\section{V36: signed critical prefixes and the zero-saddle Stokes split}
\label{sec:grh-v36}

V35 exhausted the generic residue-energy mechanism at the absolute
Shannon saddle \(\log n=1/(4t)\).  The next question is whether the signed
prime reader itself splits favorably at that location.  The answer has two
parts.

First, a zero with horizontal displacement \(\delta>0\) creates a signed
Gaussian response whose real saddle is
\[
 \log n=\frac{\delta}{2t},
 \tag{V36.1}
\]
strictly to the left of the absolute saddle.  A critical prefix therefore
captures the zero response in a central frequency cone.  Outside that cone,
however, the prefix and tail are exponentially larger than their sum and
cancel across a complex-error-function Stokes boundary.  Second, even in
the captured regime, the real contact value does not determine the
imaginary first moment.  A sharp median-prefix phase witness retains a
critical derivative while making the entire mean phasor real.

Thus the upstream return localizes the signed obstruction but does not
remove it.  Any successful continuation must use the actual multiplicative
phase or the geometry of exposed contacts, not another absolute prefix
norm.

\subsection{Exact truncated prime readers}

For \(v>0\), \(r\ge0\), and \(X=e^v\), define
\[
 P_{\chi,t,\le v}^{[r]}(x)
 :=\sum_{2\le n\le X}
 \frac{\Lambda(n)\chi(n)(\log n)^r}{\sqrt n}
 e^{-t(\log n)^2}e^{ix\log n},
 \tag{V36.2}
\]
\[
 P_{\chi,t,>v}^{[r]}(x)
 :=P_{\chi,t}^{[r]}(x)-P_{\chi,t,\le v}^{[r]}(x).
 \tag{V36.3}
\]
Let
\[
 \Psi_\chi(y):=\sum_{n\le y}\Lambda(n)\chi(n),
 \qquad
 F_{t,x,r}(y):=y^{-1/2+ix}(\log y)^r
 e^{-t(\log y)^2}.
 \tag{V36.4}
\]

\begin{lemma}[Signed prefix Abel identities]
\label{lem:v36-abel-prefix}
For every \(t>0\), \(v>0\), and \(r\ge0\),
\[
 \boxed{
 P_{\chi,t,\le v}^{[r]}(x)
 =F_{t,x,r}(X)\Psi_\chi(X)
 -\int_{1^-}^{X}\Psi_\chi(y)\,dF_{t,x,r}(y),}
 \tag{V36.5}
\]
\[
 \boxed{
 P_{\chi,t,>v}^{[r]}(x)
 =-F_{t,x,r}(X)\Psi_\chi(X)
 -\int_X^\infty\Psi_\chi(y)\,dF_{t,x,r}(y).}
 \tag{V36.6}
\]
The boundary terms cancel exactly in the sum.  With the cutoff \(v\) held
fixed,
\[
 \boxed{
 \partial_xP_{\chi,t,\le v}^{[r]}(x)
 =iP_{\chi,t,\le v}^{[r+1]}(x),
 \qquad
 \partial_tP_{\chi,t,\le v}^{[r]}(x)
 =-P_{\chi,t,\le v}^{[r+2]}(x),}
 \tag{V36.7}
\]
and the same identities hold for the tail.
\end{lemma}

\begin{proof}
The sums are Stieltjes integrals of \(F_{t,x,r}\) against
\(d\Psi_\chi\).  Integration by parts on \([1^-,X]\) gives
\textup{(V36.5)}.  On \([X,\infty)\), the Gaussian makes the upper
boundary vanish, giving \textup{(V36.6)}.  Absolute convergence at positive
time permits termwise differentiation, proving \textup{(V36.7)}.
\end{proof}

The warning ``with the cutoff held fixed'' is essential.  If \(v=v(t)\),
its motion produces the boundary flux visible with opposite signs in
\textup{(V36.5)}--\textup{(V36.6)}.  Treating a moving hard cutoff as though
it commuted with \(\partial_t\) would lose precisely the cancellation at
issue.

\subsection{The Gaussian response of one zero}

For \(\lambda\in\mathbb C\), define the model kernels
\[
 \mathcal K_t^{[r]}(\lambda;v)
 :=\int_0^v u^r e^{-tu^2+\lambda u}\,du,
 \qquad
 \mathcal T_t^{[r]}(\lambda;v)
 :=\int_v^\infty u^r e^{-tu^2+\lambda u}\,du.
 \tag{V36.8}
\]
Their sum is denoted \(\mathcal K_t^{[r]}(\lambda;\infty)\).  The zeroth
kernels have the exact error-function formulas
\[
 \mathcal K_t^{[0]}(\lambda;\infty)
 =\frac{\sqrt\pi}{2\sqrt t}e^{\lambda^2/(4t)}
 \operatorname{erfc}\!\left(-\frac{\lambda}{2\sqrt t}\right),
 \tag{V36.9}
\]
\[
 \mathcal T_t^{[0]}(\lambda;v)
 =\frac{\sqrt\pi}{2\sqrt t}e^{\lambda^2/(4t)}
 \operatorname{erfc}\!\left(
 \sqrt t\,v-\frac{\lambda}{2\sqrt t}
 \right).
 \tag{V36.10}
\]
Higher kernels are obtained by differentiating in \(\lambda\).

To connect this with the explicit formula, suppose a zero term in
\(\Psi_\chi\) is \(-y^\rho/\rho\).  Its Stieltjes differential is
\(-y^{\rho-1}dy\), so its contribution to
\textup{(V36.2)} is
\[
 -\mathcal K_t^{[r]}(\lambda;v),
 \qquad
 \lambda:=\rho-\frac12+ix.
 \tag{V36.11}
\]
Choosing the member of the conjugate zero pair with ordinate convention
matching the reader, write
\[
 \lambda=\delta+i\tau,
 \qquad
 \delta=\Re\rho-\frac12>0,
 \qquad
 \tau=x-\gamma.
 \tag{V36.12}
\]
The visible zero region is \(|\tau|<\delta\), where
\(\Re(\lambda^2)=\delta^2-\tau^2>0\).

Take a V35 critical cutoff
\[
 v_t(s)=\frac1{4t}+\frac{s}{\sqrt{2t}}+O(\log(1/t)),
 \qquad s\in\mathbb R\ \text{fixed},
 \tag{V36.13}
\]
where the logarithmic term may include the polynomial slack used there.
Put
\[
 \Delta_\delta:=\frac12-\delta>0.
 \tag{V36.14}
\]

\begin{theorem}[Zero-saddle Stokes split at the Shannon cutoff]
\label{thm:v36-stokes-split}
Fix \(0<\delta<1/2\), an integer \(r\ge0\), and a compact subset of the
visible region \(|\tau|<\delta\) avoiding
\(|\tau|=\Delta_\delta\).  Then
\[
 \boxed{
 \lim_{t\downarrow0}4t\log
 \left|
 \frac{\mathcal T_t^{[r]}(\delta+i\tau;v_t(s))}
 {\mathcal K_t^{[r]}(\delta+i\tau;\infty)}
 \right|
 =-\left(\Delta_\delta^2-\tau^2\right).}
 \tag{V36.15}
\]
Consequently:
\begin{enumerate}[label=\textup{(\roman*)},leftmargin=3em]
\item if \(|\tau|<\Delta_\delta\), then for every fixed \(r\),
\[
 \mathcal K_t^{[r]}(\delta+i\tau;v_t(s))
 =\mathcal K_t^{[r]}(\delta+i\tau;\infty)
 \left(1+O(t^{-M}e^{-c/t})\right);
 \tag{V36.16}
\]
the critical prefix captures both the zero value and all fixed
logarithmic derivatives;
\item if \(\Delta_\delta<|\tau|<\delta\), then
\[
 \left|
 \frac{\mathcal T_t^{[r]}(\delta+i\tau;v_t(s))}
 {\mathcal K_t^{[r]}(\delta+i\tau;\infty)}
 \right|\longrightarrow\infty,
 \qquad
 \frac{\mathcal K_t^{[r]}(\delta+i\tau;v_t(s))}
 {\mathcal T_t^{[r]}(\delta+i\tau;v_t(s))}
 \longrightarrow-1.
 \tag{V36.17}
\]
The prefix and tail are then exponentially larger than the full zero
response and cancel to produce it.
\end{enumerate}
\end{theorem}

\begin{proof}
In the visible region, the standard large-argument asymptotic for
\(\operatorname{erfc}\) in \textup{(V36.9)} gives
\[
 \mathcal K_t^{[r]}(\lambda;\infty)
 =t^{-r}\left(\frac\lambda2+O(t)\right)^r
 \sqrt{\frac\pi t}\,e^{\lambda^2/(4t)}(1+O(e^{-c/t})).
 \tag{V36.18}
\]
For the tail, the error-function argument is
\[
 z_t
 :=\sqrt t\,v_t(s)-\frac{\lambda}{2\sqrt t}
 =\frac{\Delta_\delta-i\tau}{2\sqrt t}
 +\frac{s}{\sqrt2}+O(\sqrt t\log(1/t)).
 \tag{V36.19}
\]
Its argument remains in a closed subsector of
\(|\arg z|<3\pi/4\).  Hence
\[
 \operatorname{erfc}(z_t)
 =\frac{e^{-z_t^2}}{\sqrt\pi z_t}
 \left(1+O(|z_t|^{-2})\right).
 \tag{V36.20}
\]
Polynomial factors from \(r\) and the bounded or logarithmic displacement
in \(v_t\) do not affect a \(4t\)-scaled logarithm.  Since
\[
 4t\Re(z_t^2)
 \longrightarrow\Delta_\delta^2-\tau^2,
 \tag{V36.21}
\]
equations \textup{(V36.9)}--\textup{(V36.10)} prove
\textup{(V36.15)}.  The first conclusion follows when the right side is
negative.  When it is positive, the identity
\(\mathcal K(\infty)=\mathcal K(v)+\mathcal T(v)\) shows that the prefix
and tail have ratio tending to \(-1\), proving the second conclusion.
\end{proof}

The two real saddles are now explicit.  The absolute prime mass is centered
at \(1/(4t)\), while the zero response is centered at \(\delta/(2t)\).
Their separation is
\[
 \frac{1-2\delta}{4t}=\frac{\Delta_\delta}{2t}.
 \tag{V36.22}
\]
The Stokes boundary compares this real separation with the Fourier offset
\(|\tau|\).  It is not visible in an absolute \(L^2\) analysis.

\begin{corollary}[The three-quarter threshold]
\label{cor:v36-three-quarter}
If \(0<\delta\le1/4\), equivalently
\(1/2<\Re\rho\le3/4\), then the entire visible interval
\(|\tau|<\delta\) lies in the prefix-capture cone
\(|\tau|<\Delta_\delta\), apart from the common boundary when
\(\delta=1/4\).  If \(1/4<\delta<1/2\), the visible interval splits into
\[
 |\tau|<\frac12-\delta
 \quad\text{and}\quad
 \frac12-\delta<|\tau|<\delta,
 \tag{V36.23}
\]
corresponding respectively to prefix capture and prefix--tail Stokes
cancellation.
\end{corollary}

\begin{proof}
Compare \(\delta\) with \(\Delta_\delta=1/2-\delta\).
\end{proof}

This is a localization result, not a zero-free region.  In the captured
regime the forbidden zero signal is already inside the critical prefix; in
the outer wing it survives only after a large signed cancellation across
the cutoff.  Neither alternative follows from V35's light-sector energy.

\subsection{Exact prefix contact and derivative coordinates}

For any cutoff \(v\), normalize the signed value pieces by
\[
 z_{\le v}(x):=\frac{P_{\chi,t,\le v}^{[0]}(x)}{Z_\chi(t)},
 \qquad
 z_{>v}(x):=\frac{P_{\chi,t,>v}^{[0]}(x)}{Z_\chi(t)},
 \tag{V36.24}
\]
and the standardized first-moment pieces by
\[
 g_{\le v}(x)
 :=\sqrt{2t}\left(
 \frac{P_{\chi,t,\le v}^{[1]}(x)}{Z_\chi(t)}
 -\mu_\chi(t)z_{\le v}(x)
 \right),
 \tag{V36.25}
\]
\[
 g_{>v}(x)
 :=\sqrt{2t}\left(
 \frac{P_{\chi,t,>v}^{[1]}(x)}{Z_\chi(t)}
 -\mu_\chi(t)z_{>v}(x)
 \right).
 \tag{V36.26}
\]

\begin{proposition}[Signed two-piece contact ledger]
\label{prop:v36-contact-ledger}
At every zero-level contact,
\[
 \boxed{
 \Re\bigl(z_{\le v}(x)+z_{>v}(x)\bigr)
 =\frac{a_{\chi,t}(x)}{2Z_\chi(t)},}
 \tag{V36.27}
\]
and the normalized derivative regression is exactly
\[
 \boxed{
 \mathcal R_{\chi,t}(x)
 =\mu_\chi(t)\Im\bigl(z_{\le v}(x)+z_{>v}(x)\bigr)
 +\frac1{\sqrt{2t}}
 \Im\bigl(g_{\le v}(x)+g_{>v}(x)\bigr).}
 \tag{V36.28}
\]
These identities hold for every cutoff, including the critical cutoff
\textup{(V36.13)}.
\end{proposition}

\begin{proof}
Equation \textup{(V36.27)} is the contact equation
\textup{(V28.17)} after the exact prefix--tail split.  By definition,
\[
 \frac{P_{\chi,t}^{[1]}(x)}{Z_\chi(t)}
 =\mu_\chi(t)(z_{\le v}+z_{>v})
 +\frac1{\sqrt{2t}}(g_{\le v}+g_{>v}).
 \tag{V36.29}
\]
Take imaginary parts and use \(\mathcal R_{\chi,t}(x)
=\Im(P_{\chi,t}^{[1]}(x)/Z_\chi(t))\).
\end{proof}

The contact ledger exposes the remaining logical gap.  It fixes one real
component of the sum in \textup{(V36.27)}.  It fixes neither the imaginary
carrier in the first term of \textup{(V36.28)} nor the imaginary
standardized moment in the second.  The Stokes theorem tells where a zero
signal sits, but supplies no algebraic relation between those components.

\subsection{A median-prefix witness with no imaginary carrier}

The carrier can be removed completely without solving the problem.  The
following limiting Gaussian construction strengthens the constant-phase
capacity witness of V28 by forcing the whole mean phasor to be real.

\begin{proposition}[Sharp signed median-prefix capacity]
\label{prop:v36-median-capacity}
Let \(Z\) be standard normal and fix \(-1<r<1\).  Define
\[
 \zeta_r(Z)
 :=r+i\sqrt{1-r^2}\,\operatorname{sgn}(Z).
 \tag{V36.30}
\]
Then \(|\zeta_r|=1\) almost surely,
\[
 \boxed{
 \mathbb E\zeta_r=r\in\mathbb R,}
 \tag{V36.31}
\]
but
\[
 \boxed{
 \Im\mathbb E[Z\zeta_r]
 =\sqrt{1-r^2}\,\sqrt{\frac2\pi}.}
 \tag{V36.32}
\]
Consequently, for the Gaussian saddle coordinate
\(U_t=\mu_\chi(t)+Z/\sqrt{2t}\),
\[
 \boxed{
 \Im\mathbb E[U_t\zeta_r]
 =\frac{\sqrt{1-r^2}}{\sqrt{\pi t}}.}
 \tag{V36.33}
\]
After multiplication by \(Z_\chi(t)\), this witness has logarithmic rate
\[
 \lim_{t\downarrow0}4t\log\left(
 1+Z_\chi(t)\left|\Im\mathbb E[U_t\zeta_r]\right|
 \right)=\frac14.
 \tag{V36.34}
\]
\end{proposition}

\begin{proof}
Symmetry gives \(\mathbb E\operatorname{sgn}(Z)=0\), proving
\textup{(V36.31)}.  Also
\[
 \mathbb E[Z\operatorname{sgn}(Z)]
 =\mathbb E|Z|=\sqrt{\frac2\pi},
 \tag{V36.35}
\]
which proves \textup{(V36.32)}.  Since the mean phasor is real, the
\(\mu_\chi(t)\) carrier contributes no imaginary part.  Scaling proves
\textup{(V36.33)}, and the V28 asymptotic
\(Z_\chi(t)=\exp(1/(16t)+O(\log(1/t)))\) proves
\textup{(V36.34)}.
\end{proof}

The two values of \(\zeta_r\) lie on opposite sides of the median critical
prefix \(Z=0\).  Thus even the stronger datum
\(\mathbb E\zeta_r=r\)---not merely its real part---leaves a first
standardized moment of order one and an unnormalized derivative of order
\(t^{-1/2}\).  This is a capacity witness, not an assertion about
Dirichlet-character phases, but it proves that critical-prefix mass,
Gaussian moments, and the contact value cannot by themselves yield the
needed signed saving.

\subsection{Consequences for the GRH route}

V36 changes the interpretation of the V35 endpoint.  The absolute saddle
at \(1/(4t)\) is not the saddle of a hypothetical zero contribution.  A
right-half zero at \(1/2+\delta+i\gamma\) lives at
\(\delta/(2t)\), deep inside every fixed V35 critical prefix.  For
\(\Re\rho\le3/4\), the prefix captures its entire visible response.  For
\(\Re\rho>3/4\), only the outer visible wings require an exponentially
large prefix--tail cancellation.

Neither case creates a contradiction with a boundary contact.  In the
capture cone, both the value and derivative of the zero response are
retained and can form the steep contacts already constructed in V27.  In
the Stokes wing, estimating the pieces separately destroys their essential
signed cancellation.  Proposition~\ref{prop:v36-median-capacity} shows that
even eliminating the imaginary carrier does not force the centered
derivative to vanish.

\begin{firewall}[Signed-saddle localization is not zero exclusion]
V36 derives exact Abel prefix identities, the complex-error-function
transfer law for one zero response, the three-quarter Stokes threshold,
the exact two-piece contact ledger, and a median-prefix capacity witness.
The one-zero kernel is a component of the explicit formula, not the full
prime reader; different zeros and archimedean terms still interact.  The
capacity witness is not a character sequence.  No bound for the actual
contact regression is proved, no off-line zero is excluded, and GRH
remains open.
\end{firewall}

\begin{assessment}[V36 acquisition and the exposure-cone test]
The signed obstruction has now been localized more sharply than by the
absolute entropy analysis.  Its decisive geometry is the comparison
\[
 |x-\gamma|
 \quad\text{versus}\quad
 \frac12-\delta,
 \tag{V36.36}
\]
not the collision or Shannon support size.  The critical prefix captures a
zero response inside this cone and participates in Stokes cancellation
outside it.  The real contact condition supplies no missing imaginary
moment in either regime.

V37 should compare the exposed-contact locations constructed in V23--V27
with the cone \textup{(V36.36)}.  If every exposed off-line zero admits a
phase-tuned contact inside its capture cone, the critical prefix gives a
complete local model but no exclusion and this branch can be closed.  If
some exposed contacts are forced into the Stokes wing, the exact
prefix--tail cancellation there becomes a new signed constraint worth
testing.  This geometric comparison is upstream of any further norm
estimate.
\end{assessment}

\section*{Version V36 append-only record}
\addcontentsline{toc}{section}{Version V36 append-only record}
The complete V35 source was retained before this continuation.  It had
\(556263\) bytes, \(16914\) lines, and SHA--256
\[
 \texttt{D583732F1DF7B8CA36281F0B889FC50650A31E29399B5496154D920AB18E6E18}.
\]
V36 established the exact signed prefix Abel ledger, the zero-saddle
Stokes split and three-quarter threshold, the critical prefix contact
coordinates, and the sharp median-prefix derivative-capacity witness.

% =============================================================================
% END APPENDED V36 MATERIAL
% =============================================================================


% =============================================================================
% BEGIN APPENDED V37 MATERIAL
% =============================================================================

\section{V37: centered exposure, captured contacts, and hard-prefix aliasing}
\label{sec:grh-v37}

V36 left a geometric alternative.  An exposed contact in the cone
\[
 |x-\gamma|<\frac12-\delta
 \tag{V37.1}
\]
is captured by the Shannon critical prefix, whereas a contact outside the
cone is produced through prefix--tail Stokes cancellation.  This section
determines which alternative is forced by the exposed-pair construction of
V23--V27.

There is always a centered exposed pair.  Starting from any off-line pair,
move to a competitor which dominates it at its own ordinate.  Each move
strictly increases \(\delta^2\), and the squared ordinate displacement is
at most that increase.  An infinite chain cannot remain bounded by local
finiteness.  It also cannot escape to infinity: crossing each unit ordinate
shell costs at least the reciprocal of the shell's \(O(\log T)\) zero
count in squared displacement, and those costs diverge.  The chain must
terminate at a pair which is uniquely exposed at \(x=\gamma\).

Its V27 phase-tuned contacts satisfy \(x-\gamma=O(t)\), so they are strictly
inside the V36 capture cone.  Thus the Stokes wing is never forced for the
dominant pair witnessing a failure of GRH.

This does not validate a hard prime-prefix proof.  A different zero with
\(\delta>1/4\), observed outside its capture cone, contributes to the hard
prefix and tail at logarithmic rate \(\delta-1/4\), independent of its
ordinate.  Those two large boundary responses cancel in the full reader.
Consequently full-symbol exposure need not pass to either hard piece.  The
critical hard cutoff captures the dominant zero but can simultaneously
alias remote high-displacement zeros.

There is also a typographical correction to the inherited V23 text.  In
equation \textup{(V23.5)}, the first displacement difference is
\(\delta_p^2-\delta_q^2\); the missing backslash in front of the second
\(\delta\) is purely typographical.

\subsection{A centered exposure theorem}

Retain the grouped reflected pairs
\(p=(\gamma_p,\delta_p,m_p)\) and their scores
\[
 S_p(x)=\delta_p^2-(x-\gamma_p)^2,
 \qquad 0\le\delta_p<\frac12.
 \tag{V37.2}
\]
Pairs with the same ordinate and smaller displacement may be discarded,
because their score is everywhere smaller.  The standard local zero count
gives
\[
 \#\{p:T\le|\gamma_p|<T+1\}=O_\chi(\log(2+T)).
 \tag{V37.3}
\]

\begin{theorem}[Existence of a center-exposed off-line pair]
\label{thm:v37-centered-exposure}
If the paired zero set contains an off-line pair, then there are an
off-line pair \(p_*=(\gamma_*,\delta_*,m_*)\) and \(\eta>0\) such that
\[
 \boxed{
 S_{p_*}(\gamma_*)=\delta_*^2>0,}
 \tag{V37.4}
\]
and
\[
 \boxed{
 S_p(\gamma_*)\le\delta_*^2-\eta
 \qquad(p\ne p_*).}
 \tag{V37.5}
\]
The same exposure, with gap at least \(\eta/2\), holds on a fixed
neighbourhood of \(\gamma_*\).
\end{theorem}

\begin{proof}
Start from any off-line pair \(p_0\).  If \(p_j\) is not the unique
maximizer of the score at its own center \(\gamma_j\), choose a distinct
pair \(p_{j+1}\) satisfying
\[
 S_{p_{j+1}}(\gamma_j)
 \ge S_{p_j}(\gamma_j)=\delta_j^2.
 \tag{V37.6}
\]
Such a competitor is off line, since its score is positive.  After the
same-ordinate dominated pairs have been removed, its ordinate differs from
\(\gamma_j\), and \textup{(V37.6)} gives
\[
 0<|\gamma_{j+1}-\gamma_j|<\frac12,
 \qquad
 |\gamma_{j+1}-\gamma_j|^2
 \le\delta_{j+1}^2-\delta_j^2.
 \tag{V37.7}
\]
Thus \(\delta_j^2\) is strictly increasing and
\[
 \sum_{j\ge0}|\gamma_{j+1}-\gamma_j|^2
 \le\frac14-\delta_0^2<\infty.
 \tag{V37.8}
\]

Suppose the construction does not terminate.  The ordinates cannot remain
in a bounded interval, because the zero set is locally finite and no pair
can recur while \(\delta_j^2\) strictly increases.  Hence a subsequence
escapes to one of the two ends of the real line; reverse signs if necessary
and consider escape to \(+\infty\).

For every sufficiently large integer \(k\) crossed by the polygonal path
\((\gamma_j)\), collect the steps whose line segments meet \([k,k+1]\).
Since every step has length less than \(1/2\), both endpoints of each
collected step lie in \([k-1,k+2]\).  By \textup{(V37.3)} and the fact that
the chain visits each pair at most once, there are at most
\(C_\chi\log(2+k)\) such steps.  Their total variation inside
\([k,k+1]\) is at least one.  Cauchy--Schwarz therefore gives
\[
 \sum_{\substack{j:\,[\gamma_j,\gamma_{j+1}]\\
                         \cap[k,k+1]\ne\varnothing}}
 |\gamma_{j+1}-\gamma_j|^2
 \ge\frac{c_\chi}{\log(2+k)}.
 \tag{V37.9}
\]
Each step is counted for at most two adjacent unit intervals.  Summing
\textup{(V37.9)} over the crossed shells contradicts
\textup{(V37.8)}, because
\(\sum_k1/\log(2+k)=\infty\).  Therefore the ascent terminates.

At the terminal pair \(p_*\), every other pair has strictly smaller score
at \(\gamma_*\).  Only finitely many pairs can have nonnegative score
there, and the remaining pairs are uniformly below zero.  Taking the
minimum of the finitely many positive gaps gives \(\eta>0\).
Continuity gives the neighbourhood statement.
\end{proof}

The theorem strengthens Lemma~\ref{lem:v23-exposed}: the exposed point need
not merely be chosen near some zero ordinate; it can be taken to be the
ordinate of the winning pair itself.

\subsection{The V27 contacts are always in the capture cone}

\begin{corollary}[Centered phase-tuned contacts are prefix-captured]
\label{cor:v37-captured-contacts}
Let \(p_*\) be the pair from Theorem
\ref{thm:v37-centered-exposure}.  There are phase-tuned times
\(t_k\downarrow0\), negative lobes centered at
\[
 \xi_k=\gamma_*+\frac{2\pi t_k}{\delta_*},
 \tag{V37.10}
\]
and adjacent simple contacts \(x_{k,-}<\xi_k<x_{k,+}\) satisfying
\[
 \boxed{
 x_{k,\pm}
 =\gamma_*+\frac{(2\pm1)\pi t_k}{\delta_*}
 +O\!\left(t_ke^{-c/t_k}\right).}
 \tag{V37.11}
\]
In particular,
\[
 \boxed{
 |x_{k,\pm}-\gamma_*|
 <\frac12-\delta_*}
 \tag{V37.12}
\]
for all sufficiently large \(k\).  At both contacts, every fixed
zero-response reader of \(p_*\) is captured by a V35 critical prefix with
relative error \(O(t_k^{-M}e^{-c'/t_k})\).
\end{corollary}

\begin{proof}
At the centered exposure point, the displacement \(d_*\) in V23 is zero.
Equation \textup{(V23.13)} gives \textup{(V37.10)}, and the nominal V27
contacts are \(\xi_k\pm\pi t_k/\delta_*\).  Theorem
\ref{thm:v27-endpoint-asymptotics} gives the exponentially smaller root
correction in \textup{(V37.11)}.  Since \(0<\delta_*<1/2\), the capture
radius \(1/2-\delta_*\) is fixed and positive, whereas the contact
displacements in \textup{(V37.11)} tend to zero.  This proves
\textup{(V37.12)}.  Apply Theorem~\ref{thm:v36-stokes-split} on a compact
subcone around \(\tau=0\).
\end{proof}

Thus the forbidden contacts required to witness GRH failure never need the
Stokes wing of their dominant pair.  The signed zero contribution creating
their value and slope is already present in the critical prefix.

\subsection{The ordinate-independent boundary alias}

The preceding statement concerns the dominant pair.  A hard prefix is a
sum over all prime powers, and its other zero components need not inherit
the full-symbol exposure gap.

\begin{theorem}[Hard-cutoff response outside the capture cone]
\label{thm:v37-boundary-alias}
Let \(\lambda=\delta+i\tau\), with \(0<\delta<1/2\), and let
\(v_t(s)\) be the critical cutoff \textup{(V36.13)}.  For every fixed
integer \(r\ge0\) and fixed \(\tau\),
\[
 \boxed{
 \lim_{t\downarrow0}4t\log
 \left|\mathcal T_t^{[r]}(\lambda;v_t(s))\right|
 =\delta-\frac14.}
 \tag{V37.13}
\]
Furthermore,
\[
 \limsup_{t\downarrow0}4t\log
 \left(1+|\mathcal K_t^{[r]}(\lambda;\infty)|\right)
 \le\max\{\delta^2-\tau^2,0\}.
 \tag{V37.14}
\]
Hence, if
\[
 \delta>\frac14,
 \qquad
 |\tau|>\frac12-\delta,
 \tag{V37.15}
\]
then
\[
 \boxed{
 \lim_{t\downarrow0}4t\log
 \left|\mathcal K_t^{[r]}(\lambda;v_t(s))\right|
 =\delta-\frac14,}
 \tag{V37.16}
\]
and
\[
 \boxed{
 \frac{\mathcal K_t^{[r]}(\lambda;v_t(s))}
 {\mathcal T_t^{[r]}(\lambda;v_t(s))}
 \longrightarrow-1.}
 \tag{V37.17}
\]
The positive exponential rate in \textup{(V37.13)} and
\textup{(V37.16)} is independent of the ordinate offset \(\tau\).
\end{theorem}

\begin{proof}
For \(r=0\), insert the large-argument expansion
\textup{(V36.20)} into the exact tail formula \textup{(V36.10)}.  The
exponential factor simplifies before taking real parts:
\[
 \frac{\lambda^2}{4t}
 -\left(\sqrt t\,v-\frac\lambda{2\sqrt t}\right)^2
 =-tv^2+\lambda v.
 \tag{V37.18}
\]
Since \(v=1/(4t)+O(t^{-1/2}+\log(1/t))\),
\[
 4t\Re(-tv^2+\lambda v)
 \longrightarrow\delta-\frac14.
 \tag{V37.19}
\]
The error-function denominator and the kernels with \(r>0\) contribute
only powers of \(t^{-1}\), proving \textup{(V37.13)}.

Equation \textup{(V36.9)} gives \textup{(V37.14)}.  In the saddle region
\(|\tau|<\delta\), its exponential rate is
\(\delta^2-\tau^2\); outside that region, the complementary-error-function
endpoint term has rate zero.  Under \textup{(V37.15)},
\[
 \delta-\frac14
 -\max\{\delta^2-\tau^2,0\}>0.
 \tag{V37.20}
\]
For \(|\tau|<\delta\), this difference is
\(\tau^2-(1/2-\delta)^2\); for \(|\tau|\ge\delta\), it is simply
\(\delta-1/4\).  Thus the tail is exponentially larger than the full
kernel.  Since
\(\mathcal K(v)=\mathcal K(\infty)-\mathcal T(v)\), equations
\textup{(V37.16)}--\textup{(V37.17)} follow.
\end{proof}

The rate \(\delta-1/4\) is the cutoff boundary value.  It is the signed
counterpart of the boundary terms in the Abel identities
\textup{(V36.5)}--\textup{(V36.6)}.  Its independence of \(\tau\) means
that the hard cutoff has destroyed Gaussian ordinate localization before
the prefix and tail are recombined.

\begin{proposition}[Full exposure does not transfer to a hard prefix]
\label{prop:v37-exposure-transfer-fails}
There are finite two-pair zero-kernel configurations for which one pair is
uniquely center-exposed in the full Gaussian reader, while the other pair
has exponentially larger critical-prefix and critical-tail components.
Specifically, choose
\[
 0<a<\frac12,
 \qquad
 \frac14<b<\frac12,
 \qquad
 a^2<b-\frac14,
 \tag{V37.21}
\]
put \((\gamma_p,\delta_p)=(0,a)\), and put
\((\gamma_q,\delta_q)=(L,b)\) with \(L>b\).  At \(x=0\), the full
\(p\)-response has positive rate \(a^2\), the full \(q\)-response has no
positive rate, but the hard prefix and tail of \(q\) both have rate
\[
 b-\frac14>a^2.
 \tag{V37.22}
\]
\end{proposition}

\begin{proof}
The full scores at \(x=0\) are \(a^2\) and \(b^2-L^2<0\), so \(p\) is
uniquely center-exposed.  Pair \(q\) satisfies
\(|x-\gamma_q|=L>1/2-b\), and Theorem
\ref{thm:v37-boundary-alias} gives \textup{(V37.22)}.
\end{proof}

This finite model obeys the strip and local-finiteness inputs used by the
exposure argument.  It does not assert that a particular Dirichlet
\(L\)-function has these two zeros.  It proves that the V23 exposure gap
alone cannot control a hard prime prefix or hard prime tail.

\subsection{Consequences for the contact route}

The exposure-cone test therefore has a definite but two-sided answer.
The dominant off-line pair can always be selected so that its actual V27
contacts lie in the prefix-capture cone.  For that pair, no Stokes
cancellation is needed.  But the full prime reader cannot be replaced by
the hard prefix: remote zeros with \(\Re\rho>3/4\) can enter through the
cutoff boundary at an ordinate-independent positive exponential rate and
cancel only against the hard tail.

This explains why the exact contact ledger in V36 did not become a local
prime criterion after truncation.  The obstruction is no longer the
dominant zero; it is boundary aliasing from the rest of the zero set.  A
componentwise exposure argument and a hard partition of the prime sum do
not commute.

\begin{firewall}[Centered capture is not hard-prefix domination]
V37 proves that some off-line pair is center-exposed and that its
phase-tuned contacts are inside the V36 capture cone.  It also proves the
single-kernel boundary-alias rate and gives a finite capacity configuration
showing failure of exposure transfer.  The zero-kernel decomposition of a
globally truncated explicit formula requires the usual symmetric
regularization; V37 does not assert termwise absolute summability over all
zeros.  No estimate for the actual signed prime prefix is obtained, no
off-line zero is excluded, and GRH remains open.
\end{firewall}

\begin{assessment}[V37 acquisition and the soft-split requirement]
The Stokes wing is optional for the dominant witness but unavoidable as a
potential contamination channel.  The new center-exposure theorem removes
the geometric uncertainty left by V36.  The boundary-alias theorem then
shows why this does not close the contact regression: a hard cutoff trades
Gaussian ordinate localization for a boundary response of rate
\((\delta-1/4)_+\).

V38 should replace the hard prefix by an analytic soft split whose Mellin
transform retains decay in the ordinate offset.  The first task is to
derive the exact transfer function for a Gaussian or error-function window
of width \(t^{-1/2}\), and to determine whether one can simultaneously
capture the centered exposed pair and suppress every remote
\(\delta>1/4\) response.  If an uncertainty tradeoff forbids both, it
should be stated as a sharp no-go theorem; if not, the resulting soft
reader is the first prefix localization not already defeated by
\textup{(V37.13)}.
\end{assessment}

\section*{Version V37 append-only record}
\addcontentsline{toc}{section}{Version V37 append-only record}
The complete V36 source was retained before this continuation.  It had
\(571982\) bytes, \(17402\) lines, and SHA--256
\[
 \texttt{54C6FE6E5DB8E4CB5B9B98273DB535EE50D90F756EC1F52F45D50B3B775BCB96}.
\]
V37 established center exposure under zero counting, capture of the
associated phase-tuned contacts, the ordinate-independent hard-boundary
alias rate, and the failure of full exposure to transfer to a hard prime
prefix.

% =============================================================================
% END APPENDED V37 MATERIAL
% =============================================================================


% =============================================================================
% BEGIN APPENDED V38 MATERIAL
% =============================================================================

\section{V38: analytic soft splitting restores ordinate localization}
\label{sec:grh-v38}

V37 proved that a hard cutoff at the absolute Shannon saddle creates an
ordinate-independent boundary alias of rate \((\delta-1/4)_+\).  It asked
whether this is an unavoidable localization tradeoff or merely a defect of
the discontinuous split.  This section gives a definitive answer: it is a
hard-boundary defect.

An error-function partition of unity has an exactly computable Mellin
transfer function.  With a fixed transition width on the saddle standard
deviation scale, the remote response regains Gaussian decay in the ordinate
offset.  With a mesoscopic width strictly between \(t^{-1/2}\) and
\(t^{-1}\), the soft-prefix exponent converges all the way back to the full
localized zero score
\[
 \delta^2-(x-\gamma)^2.
 \tag{V38.1}
\]
At the same time, the centered exposed pair of V37 is captured with
relative error tending to zero faster than every power of \(t\).

Thus there is no uncertainty no-go at this stage.  A soft split can both
capture the dominant centered zero and suppress remote high-displacement
aliases.  What remains is to find arithmetic control of the resulting soft
prime reader; the analytic localization problem itself is solved.

\subsection{The Gaussian-CDF partition}

Let
\[
 \phi(z):=\frac1{\sqrt{2\pi}}e^{-z^2/2},
 \qquad
 \Phi(z):=\frac12\operatorname{erfc}\!\left(-\frac z{\sqrt2}\right),
 \tag{V38.2}
\]
where \(\Phi\) is understood by its entire continuation.  For
\(\omega>0\), set
\[
 a_{t,\omega}:=\frac{\omega}{\sqrt{2t}},
 \tag{V38.3}
\]
and define, for real \(u,v\),
\[
 W^-_{t,\omega,v}(u)
 :=\Phi\!\left(\frac{v-u}{a_{t,\omega}}\right),
 \qquad
 W^+_{t,\omega,v}(u):=1-W^-_{t,\omega,v}(u).
 \tag{V38.4}
\]
Both windows take values in \([0,1]\) on the real axis, are entire in
\(u\), and satisfy \(W^-+W^+=1\) exactly.

For \(\lambda\in\mathbb C\), introduce the bilateral kernels
\[
 \mathcal G_{t,\omega}^{\pm}(\lambda;v)
 :=\int_{-\infty}^{\infty}
 e^{-tu^2+\lambda u}W^\pm_{t,\omega,v}(u)\,du,
 \tag{V38.5}
\]
and
\[
 \mathcal G_t(\lambda)
 :=\int_{-\infty}^{\infty}e^{-tu^2+\lambda u}\,du
 =\sqrt{\frac\pi t}\,e^{\lambda^2/(4t)}.
 \tag{V38.6}
\]

\begin{theorem}[Exact soft-window transfer function]
\label{thm:v38-exact-transfer}
For every \(t,\omega>0\), \(v\in\mathbb R\), and
\(\lambda\in\mathbb C\),
\[
 \boxed{
 \frac{\mathcal G_{t,\omega}^{-}(\lambda;v)}
 {\mathcal G_t(\lambda)}
 =\Phi(\zeta_{t,\omega}(\lambda;v)),}
 \tag{V38.7}
\]
\[
 \boxed{
 \frac{\mathcal G_{t,\omega}^{+}(\lambda;v)}
 {\mathcal G_t(\lambda)}
 =\Phi(-\zeta_{t,\omega}(\lambda;v)),}
 \tag{V38.8}
\]
where
\[
 \boxed{
 \zeta_{t,\omega}(\lambda;v)
 :=\frac{\sqrt{2t}\,v-\lambda/\sqrt{2t}}
 {\sqrt{1+\omega^2}}.}
 \tag{V38.9}
\]
The same formulas for a fixed logarithmic derivative of the kernel are
obtained by differentiating in \(\lambda\).
\end{theorem}

\begin{proof}
Complete the square in \textup{(V38.5)}.  Relative to the normalized
Gaussian with complex mean \(m=\lambda/(2t)\) and variance
\(\sigma^2=1/(2t)\), the left side of \textup{(V38.7)} is
\[
 \mathbb E\Phi\!\left(\frac{v-U}{a_{t,\omega}}\right).
 \tag{V38.10}
\]
The elementary Gaussian convolution identity
\[
 \mathbb E\Phi\!\left(\frac{v-U}{a}\right)
 =\Phi\!\left(
 \frac{v-m}{\sqrt{a^2+\sigma^2}}
 \right)
 \tag{V38.11}
\]
holds first for real \(m\) and then for complex \(m\) by analytic
continuation.  Substituting \(a=\omega/\sqrt{2t}\) gives
\textup{(V38.7)}.  The identity \(\Phi(z)+\Phi(-z)=1\) gives
\textup{(V38.8)}.  Differentiation under the Gaussian integral is
justified absolutely on compact \(\lambda\)-sets.
\end{proof}

For the critical cutoff
\[
 v_t(s)=\frac1{4t}+\frac{s}{\sqrt{2t}}+O(\log(1/t)),
 \tag{V38.12}
\]
and a zero parameter \(\lambda=\delta+i\tau\), put as before
\(\Delta_\delta=1/2-\delta\).  Then
\[
 \zeta_{t,\omega}(\delta+i\tau;v_t(s))
 =\frac{\Delta_\delta-i\tau}
 {\sqrt{2t(1+\omega^2)}}
 +\frac{s}{\sqrt{1+\omega^2}}
 +O\!\left(
 \frac{\sqrt t\log(1/t)}{\sqrt{1+\omega^2}}
 \right).
 \tag{V38.13}
\]

\subsection{Fixed-width restoration of ordinate decay}

\begin{theorem}[Soft Stokes law]
\label{thm:v38-soft-stokes}
Fix \(0<\delta<1/2\), \(\omega>0\), and \(s\in\mathbb R\).  Uniformly
on compact \(\tau\)-sets avoiding
\(|\tau|=\Delta_\delta\),
\[
 \boxed{
 \lim_{t\downarrow0}4t\log
 \left|
 \frac{\mathcal G_{t,\omega}^{+}
       (\delta+i\tau;v_t(s))}
      {\mathcal G_t(\delta+i\tau)}
 \right|
 =-\frac{\Delta_\delta^2-\tau^2}{1+\omega^2}.}
 \tag{V38.14}
\]
For the soft prefix,
\[
 \boxed{
 \lim_{t\downarrow0}4t\log
 \left|
 \frac{\mathcal G_{t,\omega}^{-}
       (\delta+i\tau;v_t(s))}
      {\mathcal G_t(\delta+i\tau)}
 \right|
 =\frac{(\tau^2-\Delta_\delta^2)_+}{1+\omega^2}.}
 \tag{V38.15}
\]
Consequently the bilateral soft-prefix response has rate
\[
 \boxed{
 \delta^2-\tau^2
 +\frac{(\tau^2-\Delta_\delta^2)_+}{1+\omega^2}.}
 \tag{V38.16}
\]
In the outer Stokes region this equals
\[
 \delta^2-\frac{\Delta_\delta^2}{1+\omega^2}
 -\frac{\omega^2}{1+\omega^2}\tau^2.
 \tag{V38.17}
\]
Thus every \(\omega>0\) restores a strictly negative quadratic ordinate
coefficient.  The hard-cutoff alias \(\delta-1/4\) is recovered only in
the singular limit \(\omega\downarrow0\).
\end{theorem}

\begin{proof}
Equations \textup{(V38.7)}--\textup{(V38.8)} reduce the problem to the
large-complex-argument asymptotics of \(\Phi(\pm\zeta)\).  When
\(|\tau|<\Delta_\delta\), the right-tail factor satisfies
\[
 \Phi(-\zeta)
 =\frac{e^{-\zeta^2/2}}{\sqrt{2\pi}\,\zeta}
 \left(1+O(|\zeta|^{-2})\right),
 \tag{V38.18}
\]
while \(\Phi(\zeta)=1+o(1)\).  When
\(|\tau|>\Delta_\delta\), both CDF pieces have exponentially large
opposite components, and
\[
 4t\log|\Phi(\pm\zeta)|
 \longrightarrow
 \frac{\tau^2-\Delta_\delta^2}{1+\omega^2}.
 \tag{V38.19}
\]
Equation \textup{(V38.13)} shows that the bounded critical displacement
does not affect these limits.  This proves
\textup{(V38.14)}--\textup{(V38.15)}.  Adding the full bilateral Gaussian
rate \(\delta^2-\tau^2\) gives \textup{(V38.16)}, and elementary algebra
gives \textup{(V38.17)}.
\end{proof}

The Stokes boundary itself does not move: it is still
\(|\tau|=1/2-\delta\).  What changes is the behavior beyond it.  A hard
cutoff erases the entire \(\tau^2\) penalty, whereas a soft cutoff retains
the fraction \(\omega^2/(1+\omega^2)\).

\subsection{The mesoscopic window}

Choose a positive function \(\omega_t\) satisfying
\[
 \boxed{
 \omega_t\longrightarrow\infty,
 \qquad
 t\omega_t^2\longrightarrow0.}
 \tag{V38.20}
\]
Equivalently, the physical transition width obeys
\[
 t^{-1/2}\ll a_{t,\omega_t}\ll t^{-1}.
 \tag{V38.21}
\]
For example, \(\omega_t=t^{-1/8}\) is admissible.

For the actual half-line prime transform, define
\[
 \mathcal S_{t}^{\pm,[r]}(\lambda;s)
 :=\int_0^\infty u^r e^{-tu^2+\lambda u}
 W^\pm_{t,\omega_t,v_t(s)}(u)\,du.
 \tag{V38.22}
\]
The difference between \(\mathcal S_t^{\pm,[r]}\) and the corresponding
bilateral kernel is bounded by a power of \(t^{-1}\) for fixed
\(\Re\lambda>0\), and hence has zero positive \(1/t\)-rate.

\begin{theorem}[Mesoscopic soft prefixes have no positive alias]
\label{thm:v38-no-alias}
Fix \(0<\delta<1/2\), \(s\in\mathbb R\), \(r\ge0\), and
\(\tau\in\mathbb R\).  Under \textup{(V38.20)},
\[
 \boxed{
 \limsup_{t\downarrow0}4t\log
 \left(1+\left|
 \mathcal S_t^{-,[r]}(\delta+i\tau;s)
 \right|\right)
 \le(\delta^2-\tau^2)_+.}
 \tag{V38.23}
\]
The same bound holds for the soft tail.  If \(|\tau|<\delta\), the soft
prefix has the exact positive rate
\[
 \boxed{
 \lim_{t\downarrow0}4t\log
 \left|
 \mathcal S_t^{-,[r]}(\delta+i\tau;s)
 \right|
 =\delta^2-\tau^2.}
 \tag{V38.24}
\]
At the center \(\tau=0\),
\[
 \boxed{
 \frac{\mathcal S_t^{+,[r]}(\delta;s)}
 {\mathcal S_t^{-,[r]}(\delta;s)}
 =O\!\left(
 t^{-M}\exp\!\left[
 -\frac{\Delta_\delta^2+o(1)}
 {4t(1+\omega_t^2)}
 \right]
 \right)\longrightarrow0.}
 \tag{V38.25}
\]
Thus the centered zero is fully captured, while every remote positive
exponential response retains its full Gaussian ordinate localization.
\end{theorem}

\begin{proof}
Because \(t(1+\omega_t^2)\to0\), the error-function arguments in
\textup{(V38.13)} still tend to infinity.  The proof of Theorem
\ref{thm:v38-soft-stokes} applies with \(\omega=\omega_t\).  The window
correction to the \(4t\)-scaled logarithm is bounded by
\[
 \frac{(\tau^2-\Delta_\delta^2)_+}{1+\omega_t^2}
 \longrightarrow0.
 \tag{V38.26}
\]
The bilateral rate therefore converges to
\(\delta^2-\tau^2\).  The omitted negative half-line has no positive
rate, which proves \textup{(V38.23)}--\textup{(V38.24)}.  At \(\tau=0\),
the Gaussian tail asymptotic gives the exponential factor in
\textup{(V38.25)}.  Its exponent tends to \(-\infty\) by
\textup{(V38.20)}, and polynomial factors are absorbed.
\end{proof}

The two conditions in \textup{(V38.20)} have distinct roles.
The divergence \(\omega_t\to\infty\) restores the full ordinate penalty.
The upper restriction \(t\omega_t^2\to0\) keeps the transition narrow
relative to the \(t^{-1}\) separation between the absolute and signed
saddles, so the centered zero remains captured.

\subsection{Soft prime readers and exposure transfer}

Define the exact soft prime split
\[
 P_{\chi,t}^{\pm,[r]}(x;s)
 :=\sum_{n\ge2}
 \frac{\Lambda(n)\chi(n)(\log n)^r}{\sqrt n}
 e^{-t(\log n)^2}e^{ix\log n}
 W^\pm_{t,\omega_t,v_t(s)}(\log n).
 \tag{V38.27}
\]
Then
\[
 \boxed{
 P_{\chi,t}^{-,[r]}(x;s)
 +P_{\chi,t}^{+,[r]}(x;s)
 =P_{\chi,t}^{[r]}(x).}
 \tag{V38.28}
\]
For fixed \(t\) and window parameters,
\[
 \partial_xP_{\chi,t}^{\pm,[r]}(x;s)
 =iP_{\chi,t}^{\pm,[r+1]}(x;s).
 \tag{V38.29}
\]
No \(t\)-derivative identity is asserted without differentiating the
moving window.

\begin{theorem}[Mesoscopic soft splitting preserves centered exposure]
\label{thm:v38-exposure-transfer}
Let \(p_*=(\gamma_*,\delta_*,m_*)\) be the center-exposed pair from
Theorem~\ref{thm:v37-centered-exposure}, with exposure gap \(\eta>0\).
For any admissible \(\omega_t\), the following hold on a sufficiently
small fixed neighbourhood of \(\gamma_*\).
\begin{enumerate}[label=\textup{(\roman*)},leftmargin=3em]
\item The soft-prefix response of \(p_*\) has the same value and every
fixed \(x\)-derivative as its full response, with relative error
\[
 O\!\left(
 t^{-M}\exp\!\left[
 -\frac{c_*+o(1)}{t(1+\omega_t^2)}
 \right]
 \right)=o(1).
 \tag{V38.30}
\]
\item The sum of the soft-prefix responses of all other zero pairs is
\(o(1)\) relative to the \(p_*\) response along the V37 phase-tuned
sequence.
\item The soft-prefix zero-side symbol therefore has the same negative
lobe, contact-scale derivative, and exponential rate
\(\delta_*^2/(4t)\) as the centered exposed full response.
\end{enumerate}
\end{theorem}

\begin{proof}
At \(x=\gamma_*+o(1)\), the offset of the winning pair tends to zero.
Equation \textup{(V38.25)} proves item (i), and differentiation in
\(\lambda\) changes only polynomial factors.

For every other fixed pair \(p\), Theorem~\ref{thm:v38-no-alias} bounds
the positive soft-prefix exponent by
\[
 \bigl(\delta_p^2-(x-\gamma_p)^2\bigr)_+.
 \tag{V38.31}
\]
The centered exposure gap makes this at most
\(\delta_*^2-\eta/2\) on a smaller neighbourhood whenever it is positive.
For distant ordinates, the fixed-width bound
\textup{(V38.17)}, with \(\omega_t\ge1\), gives a Gaussian envelope
\[
 \frac14-c(x-\gamma_p)^2
 \tag{V38.32}
\]
with an absolute \(c>0\).  The shell count
\textup{(V37.3)} therefore sums the distant responses with only a
polynomial loss.  The exposure gap absorbs that loss.  This proves item
(ii).  The phase tuning and the differentiated dominance argument of
V23--V27 then prove item (iii).
\end{proof}

This is the first prefix construction in the note which is compatible with
the full zero-side exposure geometry.  The hard split failed because its
outer response had no ordinate penalty.  The mesoscopic analytic split
retains that penalty and still captures the winning signed saddle.

\subsection{What the successful soft split does not yet prove}

The new reader does not create a contradiction.  Under failure of GRH, its
zero-side expansion contains the same dominant negative lobe and steep
derivative as the full symbol.  On the prime side it is still a signed
multiplicative sum of exponential size.  The real contact equation for the
full reader does not separately fix either soft component.

The gain is nevertheless substantive.  Remote high-displacement zeros can
no longer contaminate the prefix at an ordinate-independent rate.  Hence a
future signed estimate for \(P_{\chi,t}^{-,[r]}\) would test the centered
exposed pair itself, rather than a hard-boundary mixture of the entire zero
set.

\begin{firewall}[Soft localization is not a prime-side bound]
V38 derives the exact Gaussian-CDF transfer function, proves the soft
Stokes law, constructs a mesoscopic no-alias regime, and transfers centered
zero exposure to the analytic soft prefix.  These are transform and
zero-side localization results.  They do not bound the actual soft
character sum at contacts, do not prove positivity of its associated Weil
form, and do not exclude an off-line zero.  GRH remains open.
\end{firewall}

\begin{assessment}[V38 acquisition and the soft contact target]
The uncertainty question posed in V37 has a positive answer.  Transition
widths
\[
 t^{-1/2}\ll a_t\ll t^{-1}
 \tag{V38.33}
\]
simultaneously capture a centered off-line zero and restore the complete
Gaussian ordinate exponent.  Hard-prefix aliasing is therefore not an
intrinsic obstruction.

V39 should write the V36 contact ledger in the exact soft coordinates
\textup{(V38.27)}, isolate the soft-prefix value and derivative at the
centered phase-tuned contacts, and determine the weakest prime-side bound
on that localized reader which would contradict the V38 negative lobe.
That criterion should be compared with the GRH-complete unsmoothed bound of
V29 to determine whether soft localization has genuinely weakened the
arithmetic target.  V40 is the next scheduled parallel-note sweep.
\end{assessment}

\section*{Version V38 append-only record}
\addcontentsline{toc}{section}{Version V38 append-only record}
The complete V37 source was retained before this continuation.  It had
\(586289\) bytes, \(17784\) lines, and SHA--256
\[
 \texttt{9D4AFAC623AAAD4DCD1DA315360922FD47E49EC54882FB648A0E3B84030BE1D6}.
\]
V38 established the exact analytic soft-window transfer, the restored
ordinate-decay law, a mesoscopic no-alias regime, and preservation of the
centered exposed zero under soft prime splitting.

% =============================================================================
% END APPENDED V38 MATERIAL
% =============================================================================


% =============================================================================
% BEGIN APPENDED V39 MATERIAL
% =============================================================================

\section{V39: the soft-contact criterion and the half-mass tax}
\label{sec:grh-v39}

V38 constructed a mesoscopic analytic split which captures a centered
exposed zero and restores Gaussian ordinate localization.  This section
asks whether that localization weakens the prime-side estimate needed for
GRH.

It weakens the sampling requirement but not the presently available
cancellation exponent.  If GRH fails, the V37 phase-tuned contacts have a
soft-prefix derivative with the full positive rate \(\delta_*^2\), while
the soft-tail derivative is little-oh of it.  Hence a zero-rate bound for
the localized derivative only at actual contacts is equivalent to GRH.

On the other hand, the mesoscopic transition is much wider than the
\(t^{-1/2}\) absolute saddle fluctuations.  Across those fluctuations the
window is asymptotically constant with value \(1/2\).  Each soft piece
therefore retains one half of the absolute mass and one half of every fixed
absolute saddle moment.  Classical prime-number-theorem cancellation still
has logarithmic rate \(1/4\).  The soft split removes zero-side aliasing; it
does not by itself create a stronger character-sum bound.

\subsection{The exact soft contact ledger}

Fix an admissible mesoscopic parameter
\[
 \omega_t\to\infty,
 \qquad t\omega_t^2\to0,
 \tag{V39.1}
\]
a critical cutoff \(v_t(s)\) as in \textup{(V38.12)}, and the soft prime
readers \(P_{\chi,t}^{\pm,[r]}(x;s)\) from
\textup{(V38.27)}.  Write
\[
 z_t^\pm(x;s):=\frac{P_{\chi,t}^{\pm,[0]}(x;s)}{Z_\chi(t)},
 \tag{V39.2}
\]
and
\[
 g_t^\pm(x;s)
 :=\sqrt{2t}\left(
 \frac{P_{\chi,t}^{\pm,[1]}(x;s)}{Z_\chi(t)}
 -\mu_\chi(t)z_t^\pm(x;s)
 \right).
 \tag{V39.3}
\]

\begin{proposition}[Exact soft value--derivative coordinates]
\label{prop:v39-soft-ledger}
For every \(t>0\) and real \(x\),
\[
 \boxed{
 z_t^-(x;s)+z_t^+(x;s)=z_{\chi,t}(x),}
 \tag{V39.4}
\]
\[
 \boxed{
 \mathcal R_{\chi,t}(x)
 =\mu_\chi(t)\Im\bigl(z_t^-(x;s)+z_t^+(x;s)\bigr)
 +\frac1{\sqrt{2t}}
 \Im\bigl(g_t^-(x;s)+g_t^+(x;s)\bigr).}
 \tag{V39.5}
\]
At a zero-level contact these are supplemented by the single real
constraint
\[
 \boxed{
 \Re\bigl(z_t^-(x;s)+z_t^+(x;s)\bigr)
 =\frac{a_{\chi,t}(x)}{2Z_\chi(t)}.}
 \tag{V39.6}
\]
Equivalently, the contact slope is
\[
 \boxed{
 \partial_xm_{\chi,t}(x)
 =\partial_xa_{\chi,t}(x)
 +2\Im P_{\chi,t}^{-,[1]}(x;s)
 +2\Im P_{\chi,t}^{+,[1]}(x;s).}
 \tag{V39.7}
\]
\end{proposition}

\begin{proof}
The partition identity \(W^-+W^+=1\) gives
\(P^-+P^+=P\) for every reader order.  Divide the zeroth identity by
\(Z_\chi(t)\) to obtain \textup{(V39.4)}.  Center the first identity at
\(\mu_\chi(t)\), take imaginary parts, and use
\(\mathcal R_{\chi,t}=\Im(P_{\chi,t}^{[1]}/Z_\chi(t))\) to obtain
\textup{(V39.5)}.  Equations \textup{(V28.17)} and
\textup{(V27.11)} give \textup{(V39.6)}--\textup{(V39.7)}.
\end{proof}

The ledger is exact but underdetermined.  The contact fixes the real part
of the sum of the soft values, not either soft value separately and not an
imaginary derivative.  The V38 exposure theorem supplies additional
information only along the contact sequence forced by an off-line zero.

\subsection{Tail removal at centered exposed contacts}

\begin{theorem}[The localized derivative carries the forbidden slope]
\label{thm:v39-prefix-carries-slope}
Assume the paired channel has an off-line zero.  Let
\(p_*=(\gamma_*,\delta_*,m_*)\) be the center-exposed pair of V37 and let
\(t_k\downarrow0\), \(x_{k,\pm}\) be its phase-tuned simple contacts.
For either sign and every fixed critical parameter \(s\),
\[
 \boxed{
 \lim_{k\to\infty}4t_k\log
 \left|Im P_{\chi,t_k}^{-,[1]}(x_{k,\pm};s)\right|
 =\delta_*^2,}
 \tag{V39.8}
\]
and
\[
 \boxed{
 \frac{
 \Im P_{\chi,t_k}^{+,[1]}(x_{k,\pm};s)}{
 \Im P_{\chi,t_k}^{-,[1]}(x_{k,\pm};s)}
 \longrightarrow0.}
 \tag{V39.9}
\]
In particular, the soft-prefix derivative alone has the exact exponential
rate of the forbidden contact slope.
\end{theorem}

\begin{proof}
Corollary~\ref{cor:v37-captured-contacts} places the contacts at
\(x_{k,\pm}=\gamma_*+O(t_k)\), strictly inside the capture cone.
Theorem~\ref{thm:v38-exposure-transfer}, applied after one
\(x\)-derivative, shows that the soft-prefix response of \(p_*\) equals
its full response up to relative \(o(1)\), while all other soft-prefix zero
responses are relatively negligible.  Its soft tail is relatively
negligible by \textup{(V38.25)}.  The same exposure-gap and shell argument
applies to the other tail responses, because the full response is also
dominated by \(p_*\) and the prefix has already captured that leading
term.

At the simple contacts, Theorem~\ref{thm:v27-endpoint-asymptotics} gives
\[
 \lim_{k\to\infty}4t_k
 \log|\partial_xm_{\chi,t_k}(x_{k,\pm})|
 =\delta_*^2.
 \tag{V39.10}
\]
The archimedean derivative is bounded on the fixed contact neighbourhood.
Insert the soft split into \textup{(V39.7)}.  The captured dominant
derivative has a fixed nonzero phase at either endpoint, so the imaginary
prime-reader component has the same magnitude up to relative \(o(1)\).
This proves \textup{(V39.8)}--\textup{(V39.9)}.
\end{proof}

The theorem is stronger than a magnitude statement at arbitrary points.
It identifies the precise imaginary component that enters the contact
slope and proves that the localized reader, rather than its tail, carries
the forbidden growth.

\subsection{A contact-restricted criterion equivalent to GRH}

For a nondegenerate compact interval \(I\), let
\[
 \mathcal C_{\chi,t}(I)
 :=\{x\in I:m_{\chi,t}(x)=0\}.
 \tag{V39.11}
\]
As before, a supremum over the empty set is zero.

\begin{criterion}[Mesoscopic soft-prefix contact criterion]
\label{crit:v39-soft-contact}
GRH for the paired channel is equivalent to the following statement: for
every nondegenerate compact interval \(I\) and one, equivalently every,
fixed \(s\in\mathbb R\),
\[
 \boxed{
 \limsup_{t\downarrow0}4t\log\left(
 1+sup_{x\in\mathcal C_{\chi,t}(I)}
 \left|Im P_{\chi,t}^{-,[1]}(x;s)\right|
 \right)=0.}
 \tag{V39.12}
\]
Equivalently, it is necessary and sufficient to prove that for every
\(\varepsilon>0\),
\[
 \left|Im P_{\chi,t}^{-,[1]}(x;s)\right|
 \le C_{\chi,I,s,\varepsilon}e^{\varepsilon/t}
 \tag{V39.13}
\]
whenever \(x\in I\), \(m_{\chi,t}(x)=0\), and \(0<t\le1\).
\end{criterion}

\begin{proof}
Under GRH the positive-time symbol is strictly positive, so the contact
sets are empty and the condition is vacuous.  Conversely, if GRH fails,
Theorem~\ref{thm:v37-centered-exposure} and Corollary
\ref{cor:v37-captured-contacts} produce a compact interval containing the
contacts in Theorem~\ref{thm:v39-prefix-carries-slope}.  Equation
\textup{(V39.8)} makes the left side of \textup{(V39.12)} equal to at
least \(\delta_*^2>0\) along that sequence, a contradiction.  The
\(\varepsilon\)-form is the standard reformulation of zero
\(1/t\)-logarithmic rate.
\end{proof}

This is the weakest sampling criterion obtained in the note so far.  It
asks for one smoothed, localized, first logarithmic reader only on the
moving zero set of the full completed symbol.  It does not ask for a
uniform estimate in \(x\), in \(t\), or in the unsmoothed cutoff variable.

For comparison, there is also a uniform soft formulation.

\begin{theorem}[Uniform soft-reader criterion]
\label{thm:v39-uniform-soft}
GRH for the paired channel is equivalent to the family of estimates
\[
 \boxed{
 \sup_{x\in I}
 \left|P_{\chi,t}^{-,[r]}(x;s)\right|
 \le C_{\chi,I,r,s,\varepsilon}e^{\varepsilon/t}}
 \tag{V39.14}
\]
for every compact interval \(I\), every fixed \(r\ge0\), every
\(\varepsilon>0\), and \(0<t\le1\).  It is enough to impose
\textup{(V39.14)} for \(r=1\).
\end{theorem}

\begin{proof}
GRH gives the square-root Chebyshev estimate of V29.  Partial summation
against the smooth window and Gaussian weight then gives a fixed power of
\(t^{-1}\), which is bounded by \(C_\varepsilon e^{\varepsilon/t}\).
Conversely, if GRH fails, the contacts in Theorem
\ref{thm:v39-prefix-carries-slope} violate the \(r=1\) estimate on a fixed
compact interval.  Higher fixed readers follow from the same partial
summation under GRH.
\end{proof}

The uniform theorem is logically equivalent to the GRH-complete
unsmoothed estimate \textup{(V29.26)}.  The contact criterion is a much
sparser sampling statement, but no unconditional method currently exploits
that sparsity.

\subsection{The mesoscopic half-mass law}

The price of removing hard-boundary aliasing is visible under the absolute
saddle measure.  Let
\[
 W_t^\pm(n):=W^\pm_{t,\omega_t,v_t(s)}(\log n).
 \tag{V39.15}
\]

\begin{theorem}[Each no-alias window retains half the absolute saddle]
\label{thm:v39-half-mass}
Under \textup{(V39.1)}, for every fixed \(s\),
\[
 \boxed{
 \mathbb E_{\nu_{\chi,t}}W_t^-
 \longrightarrow\frac12,
 \qquad
 \mathbb E_{\nu_{\chi,t}}W_t^+
 \longrightarrow\frac12.}
 \tag{V39.16}
\]
More generally, for every fixed integer \(j\ge0\),
\[
 \boxed{
 \mathbb E_{\nu_{\chi,t}}
 \left[G_t^{,j}W_t^\pm\right]
 \longrightarrow
 \frac12\mathbb E[Z^j],}
 \tag{V39.17}
\]
where \(G_t=\sqrt{2t}(U-\mu_\chi(t))\) and \(Z\) is standard normal.
For fixed raw moment order \(r\ge0\),
\[
 \boxed{
 \sum_n
 \frac{\Lambda(n)|\chi(n)|}{\sqrt n}
 e^{-t(\log n)^2}(\log n)^rW_t^\pm(n)
 =\left(\frac12+o(1)\right)
 Z_\chi(t)\mu_\chi(t)^r.}
 \tag{V39.18}
\]
Every soft absolute reader therefore has logarithmic rate \(1/4\).
\end{theorem}

\begin{proof}
Under the absolute saddle law,
\(G_t\) converges in distribution to \(Z\) and has uniformly bounded
moments of every fixed order.  Also
\[
 \frac{v_t(s)-U}{a_{t,\omega_t}}
 =\frac{s-G_t+o_{\nu}(1)}{\omega_t}
 \longrightarrow0
 \quad\hbox{in probability}.
 \tag{V39.19}
\]
Continuity and boundedness of \(\Phi\) on the real axis give
\(W_t^\pm\to1/2\) in probability.  Bounded convergence proves
\textup{(V39.16)}.  Uniform integrability of fixed powers of \(G_t\)
proves \textup{(V39.17)}.

Finally, \(U/\mu_\chi(t)\to1\) in every fixed \(L^p\), because the
relative saddle fluctuations are \(O(\sqrt t)\).  Combining this with
\textup{(V39.16)} proves \textup{(V39.18)}.  Equations
\textup{(V28.4)}--\textup{(V28.5)} give the final rate.
\end{proof}

This is the half-mass tax.  The transition is narrow compared with the
distance between the absolute saddle and every fixed off-line signed
saddle, but broad compared with the fluctuations of the absolute saddle
itself.  It can therefore localize zero responses while looking almost
constant to absolute prime mass.

\subsection{What classical PNT cancellation gives}

\begin{proposition}[The unconditional soft-reader bound keeps full type]
\label{prop:v39-soft-pnt}
For every fixed \(r\ge0\), compact interval \(I\), and fixed critical
parameter \(s\), there are \(M<\infty\) and \(c_\chi>0\) such that
\[
 \boxed{
 \sup_{x\in I}
 \left|P_{\chi,t}^{\pm,[r]}(x;s)\right|
 \le t^{-M}
 \exp\!\left(
 \frac1{16t}-\frac{c_\chi}{\sqrt t}
 \right).}
 \tag{V39.20}
\]
This estimate has only the logarithmic consequence
\[
 \limsup_{t\downarrow0}4t\log\left(
 1+\sup_{x\in I}|P_{\chi,t}^{\pm,[r]}(x;s)|
 \right)\le\frac14,
 \tag{V39.21}
\]
not the zero rate required by \textup{(V39.12)} or
\textup{(V39.14)}.
\end{proposition}

\begin{proof}
Apply Stieltjes partial summation to \(\Psi_\chi(y)\) with the product of
the Gaussian reader and \(W_t^\pm(\log y)\).  The classical fixed-character
PNT gives
\(|\Psi_\chi(y)|\ll y e^{-c_\chi\sqrt{\log y}}\).  The window is bounded
by one on the real axis, and its \(u\)-derivative is
\[
 O\!\left(a_{t,\omega_t}^{-1}\right)
 =O\!\left(\frac{\sqrt t}{\omega_t}\right).
 \tag{V39.22}
\]
All remaining derivatives insert only fixed powers of \(u\) and
\(t^{-1}\).  The same Gaussian Abel estimate as in V24 and V29 therefore
gives \textup{(V39.20)}.  Equation \textup{(V39.21)} follows immediately.
\end{proof}

The half-mass theorem shows that this failure is not caused by a careless
absolute majorant: each soft reader genuinely retains the full absolute
exponential scale.  Improvement from rate \(1/4\) to rate zero must use the
signed multiplicative phase and the contact condition together.

\subsection{Comparison with the unsmoothed GRH target}

The V29 square-root estimate
\[
 \Psi_\chi(y)=O_{\chi,\varepsilon}(y^{1/2+\varepsilon})
 \tag{V39.23}
\]
implies the uniform soft bound by partial summation.  Conversely, the
uniform soft bound excludes every off-line zero through V37--V38 and hence
implies GRH, after which \textup{(V39.23)} follows by the standard explicit
formula.  They are logically equivalent.

The contact-restricted criterion \textup{(V39.12)} is formally weaker as a
sampling problem: it asks for neither all cutoffs nor all spatial points.
But the best unconditional input, \textup{(V39.20)}, does not detect that
restriction.  Soft localization has therefore reduced the locus of the
missing estimate without reducing its known exponential difficulty.

\begin{theorem}[Soft localization changes sampling, not absolute type]
\label{thm:v39-verdict}
The following statements hold simultaneously.
\begin{enumerate}[label=\textup{(\roman*)},leftmargin=3em]
\item The contact-restricted soft-prefix derivative bound
\textup{(V39.12)} is equivalent to GRH.
\item Under a forbidden centered exposure, the soft prefix carries the
entire positive contact-slope exponent and the soft tail is negligible.
\item Each soft piece retains asymptotically one half of the absolute
saddle mass and every fixed absolute moment.
\item Classical PNT cancellation leaves logarithmic rate \(1/4\), exactly
as for the unsplit reader.
\end{enumerate}
Thus the mesoscopic analytic split is a genuine zero-localization gain but
not, by itself, a prime-cancellation gain.
\end{theorem}

\begin{proof}
These are respectively Criterion~\ref{crit:v39-soft-contact}, Theorem
\ref{thm:v39-prefix-carries-slope}, Theorem~\ref{thm:v39-half-mass}, and
Proposition~\ref{prop:v39-soft-pnt}.
\end{proof}

\begin{firewall}[A weaker sampling criterion is not a proof]
V39 derives the exact soft contact ledger, proves tail removal at the
centered forbidden contacts, gives a soft-prefix criterion equivalent to
GRH, and quantifies the half-mass and PNT obstructions.  Equivalence does
not supply the required contact-conditioned character cancellation.  The
unconditional bound still has full type, no off-line zero is excluded, and
GRH remains open.
\end{firewall}

\begin{assessment}[V39 acquisition and the scheduled sweep]
The analytic localization branch now has a precise endpoint.  It removes
hard-boundary aliasing and reduces the missing estimate to one imaginary
soft-prefix derivative sampled only at contacts.  Yet its absolute mass
and classical PNT bound remain at rate \(1/4\).  Further window engineering
without new signed arithmetic would repeat this obstruction.

V40 is the scheduled parallel-note sweep.  It should inventory all notes
created since the V30 sweep, compare their strongest current results with
the soft-contact criterion \textup{(V39.12)}, and import only arguments
that act on the actual character phase or contact-conditioned derivative.
If no such input exists, V40 should close the window-localization branch
and select the most promising genuinely signed upstream route.
\end{assessment}

\section*{Version V39 append-only record}
\addcontentsline{toc}{section}{Version V39 append-only record}
The complete V38 source was retained before this continuation.  It had
\(601199\) bytes, \(18243\) lines, and SHA--256
\[
 \texttt{55299B2FB77AB94C3BD9A230636673E071D98D3BF3C94BF4343D43173CD48582}.
\]
V39 established the exact soft contact ledger, localization of the
forbidden slope in the soft prefix, the contact-restricted and uniform soft
criteria, the mesoscopic half-mass law, and the persistence of the
unconditional \(1/4\) type.

% =============================================================================
% END APPENDED V39 MATERIAL
% =============================================================================

% =============================================================================
% BEGIN APPENDED V40 MATERIAL
% =============================================================================

\section{V40: scheduled corpus sweep and the contact two-prime burst}
\label{sec:grh-v40}

V39 reduced the remaining prime-side problem to the imaginary first
logarithmic soft reader at actual zero-level contacts.  The scheduled sweep
below finds no new companion estimate which acts on that signed quantity.
It does, however, suggest a useful change of type.  The small-time lesson in
the SM--Einstein programme is that values do not control derivatives without
a uniform second-moment row, while the Yang--Mills atlas repeatedly separates
pointwise coverage from a global sign.  Here those warnings can be made exact:
the contact derivative is embedded into a local Sobolev energy, that energy
has an explicit two-prime expansion, and its diagonal is only polynomial.
Thus an off-line zero forces an exponentially large \emph{off-diagonal signed
prime correlation} at its phase-tuned contacts.

\subsection{Timestamped V40 corpus ledger}

The preceding scheduled sweep, V30, froze \(31\) non-GRH
\texttt{.tex} files at
\[
 \texttt{2026-09-10T23:07:52.4684869Z}
\]
with canonical manifest SHA--256
\[
 \texttt{34EE166478D15E763BD319E568AB2F2156EF42037622526D7D716D38CD727BE6}.
 \tag{V40.1}
\]
The V40 cutoff was taken at
\[
 \boxed{\texttt{2026-09-11T00:42:04.3520709Z}.}
 \tag{V40.2}
\]
There were \(30\) non-GRH \texttt{.tex} files.  For the canonical UTF--8
manifest with rows
\[
 \texttt{filename|byte-count|SHA--256}\backslash\texttt{n}
\]
in ordinal filename order, the hash was
\[
 \boxed{
 \texttt{4F42F003B6739E2AABC063B0AD8E818FC189DD06E02FA86F58DCE8B0D151583E}.}
 \tag{V40.3}
\]

Exactly \(27\) rows already hash-gated as unchanged at V30 remain
byte-identical.  The SM archive f14 is also unchanged:
\begin{center}
\tiny
\begin{tabular}{p{0.53\textwidth}|r|p{0.30\textwidth}}
\textbf{unchanged row}&\textbf{bytes}&\textbf{SHA--256}\\ \hline
\path{Renewal_SM_Einstein_Continuum_Research_Notes_V100_f14.tex}
 &1166228&\texttt{C6608FD46EE95038605B57FE09CC5849}\newline
 \texttt{46F88DF6D480C7A53443385B94E64BD9}
\end{tabular}
\end{center}
The rolling V30 endpoints
\path{Renewal_SM_Einstein_Continuum_Research_Notes_V21.tex} and
\path{Renewal_Yang_Mills_Mass_Gap_Research_Notes_V35.tex}--
\path{V36.tex} have disappeared.  Their append-only successors at the V40
cutoff were:
\begin{center}
\tiny
\begin{tabular}{p{0.53\textwidth}|r|p{0.30\textwidth}}
\textbf{V40 endpoint}&\textbf{bytes}&\textbf{SHA--256}\\ \hline
\path{Renewal_SM_Einstein_Continuum_Research_Notes_V61.tex}
 &595097&\texttt{CD34919A1695F4E2D4EDEE26E0EEECE}\newline
 \texttt{27A52357415C89E3511964BBF15575CA0}\\
\path{Renewal_Yang_Mills_Mass_Gap_Research_Notes_V39.tex}
 &329898&\texttt{9B44B0FBCA6037F3319E86036753B0F3}\newline
 \texttt{B76BDD06F6045F00747E05DC01E23F9D}
\end{tabular}
\end{center}
The SM successor contains V22--V61, including the V60 audit and V61
locally covariant descent; the YM successor contains V37--V39.  Hence the
rolling replacement loses no intervening append-only section.  The cutoff
and manifest deliberately record one instant of a concurrently moving
directory.

\begin{assessment}[V40 typed transfer verdict]
\label{ass:v40-transfer}
The new corpus contributes three relevant lessons, but no GRH estimate.
\begin{enumerate}[label=\textup{(\roman*)},leftmargin=3em]
\item SM V22--V61 develops smooth spectral locality, conditional shell
      entropy, reflected transfer semigroups, second-spectral-moment recovery,
      analytic local cores, joint spectral visibility, and locally covariant
      quotient algebras.  Its sharp V56 warning---finite-time values cannot
      be differentiated at zero without a uniform second moment---motivates
      the two-reader Sobolev packet below.  None of these abstract theorems
      estimates a Dirichlet-character phase at a moving contact.
\item YM V37--V39 constructs independent exact Wilson-pencil charts and
      proves pointwise positivity on
      \([0,\pi/4]^2\times[0,t]\).  It explicitly stops before a global
      covariant sign or continuum bridge.  Its finite atlas cannot be
      substituted for cancellation in an infinite signed prime sum.
\item V30 had already imported the only directly reusable structural idea,
      namely conditional-expectation Pythagoras together with the
      moving-measure connection.  Reimporting spectral locality, finite
      charts, or positive Gram geometry would remain invariant under phase
      changes which alter
      \(\Im P_{\chi,t}^{-,[1]}\).  Such phase-blind information cannot close
      Criterion~\ref{crit:v39-soft-contact}; the sharp phase-capacity model of
      V28 already certifies this obstruction.
\end{enumerate}
Accordingly V40 imports no companion conclusion.  It uses only the
methodological second-moment warning to expose a new, genuinely signed
two-prime target.
\end{assessment}

\subsection{A contact-centred Sobolev packet}

Fix the V38--V39 mesoscopic soft prefix and a critical parameter \(s\).
Put
\[
 u_n:=\log n,
 \qquad
 b_{n,t}^{[r]}(s):=
 \frac{\Lambda(n)u_n^r}{\sqrt n}e^{-tu_n^2}W_t^-(n),
 \tag{V40.4}
\]
so that
\[
 A_{t,r}(y;s):=P_{\chi,t}^{-,[r]}(y;s)
 =\sum_{n\ge2}b_{n,t}^{[r]}(s)\chi(n)e^{iyu_n}.
 \tag{V40.5}
\]
Take the polynomial observation radius
\[
 h_t:=t^{-2}
 \tag{V40.6}
\]
and, for \(\epsilon\in\{+1,-1\}\), define
\[
 \mathcal E_{t,r}^{\epsilon}(x;s)
 :=\frac1{2h_t}\int_{x-h_t}^{x+h_t}
 \begin{cases}
  (\Re A_{t,r}(y;s))^2,&\epsilon=+1,\\
  (\Im A_{t,r}(y;s))^2,&\epsilon=-1
 \end{cases}dy.
 \tag{V40.7}
\]
These are nonnegative deterministic energies; no probabilistic model of the
primes is being introduced.

Let
\[
 \operatorname{sinc}z:=
 \begin{cases}\sin z/z,&z\ne0,\\1,&z=0.\end{cases}
 \tag{V40.8}
\]

\begin{proposition}[Exact signed two-prime expansion]
\label{prop:v40-two-prime}
For every \(t>0\), \(x\in\mathbb R\), fixed \(r\ge0\), and
\(\epsilon\in\{+1,-1\}\),
\[
\boxed{
\begin{aligned}
 \mathcal E_{t,r}^{\epsilon}(x;s)
  ={}&\frac12\Re\sum_{m,n\ge2}
  b_{m,t}^{[r]}b_{n,t}^{[r]}
  \chi(m)\overline{\chi(n)}
  e^{ix(u_m-u_n)}
  \operatorname{sinc}\!\bigl(h_t(u_m-u_n)\bigr)\\
 &+\frac{\epsilon}{2}\Re\sum_{m,n\ge2}
  b_{m,t}^{[r]}b_{n,t}^{[r]}
  \chi(m)\chi(n)
  e^{ix(u_m+u_n)}
  \operatorname{sinc}\!\bigl(h_t(u_m+u_n)\bigr).
\end{aligned}}
 \tag{V40.9}
\]
Here and below the common argument \(s\) on the \(b\)'s is suppressed.
The difference-frequency diagonal in the first line is
\[
 \mathcal D_{t,r}(s)
 :=\frac12\sum_{n\ge2}
 |b_{n,t}^{[r]}(s)|^2|\chi(n)|^2,
 \tag{V40.10}
\]
and, for fixed \(r\),
\[
 \boxed{\mathcal D_{t,r}(s)\ll_{r}t^{-(2r+3)/2}.}
 \tag{V40.11}
\]
In particular the diagonal has zero \(1/t\)-logarithmic rate.
\end{proposition}

\begin{proof}
Write \(A=A_{t,r}\).  The elementary identities
\[
 (\Re A)^2=\frac12|A|^2+\frac12\Re A^2,
 \qquad
 (\Im A)^2=\frac12|A|^2-\frac12\Re A^2
 \tag{V40.12}
\]
reduce the claim to termwise integration.  Absolute convergence follows from
the Gaussian weight, and
\[
 \frac1{2h}\int_{x-h}^{x+h}e^{iyv}\,dy
 =e^{ixv}\operatorname{sinc}(hv).
 \tag{V40.13}
\]
This proves \textup{(V40.9)} and identifies \textup{(V40.10)}.

Since \(0\le W_t^-\le1\), \(|\chi(n)|\le1\), and
\(\Lambda(n)\le\log n\), split the integers into the blocks
\(k\le\log n<k+1\).  The elementary bound
\(\sum_{e^k\le n<e^{k+1}}n^{-1}\ll1\) gives
\[
 \mathcal D_{t,r}(s)
 \ll_r\sum_{k\ge0}(k+1)^{2r+2}e^{-2tk^2}
 \ll_r t^{-(2r+3)/2},
 \tag{V40.14}
\]
uniformly in the window.  No prime number theorem is used.
\end{proof}

The two lines of \textup{(V40.9)} are genuinely arithmetic.  They contain
the multiplicative phases
\(\chi(m)\overline{\chi(n)}\) and \(\chi(m)\chi(n)\), sampled at the
contact-dependent ratio and product phases.  This is precisely the phase
information absent from an absolute Gram or finite-atlas estimate.

\subsection{Off-line zeros force an off-diagonal correlation burst}

Combine the two energies needed to recover the missing first reader:
\[
 \mathcal B_t(x;s)
 :=\mathcal E_{t,1}^{-1}(x;s)
   +h_t^2\mathcal E_{t,2}^{+1}(x;s).
 \tag{V40.15}
\]

\begin{lemma}[Point evaluation from the two-reader packet]
\label{lem:v40-trace}
For every \(t>0\) and real \(x\),
\[
 \boxed{
 |\Im P_{\chi,t}^{-,[1]}(x;s)|^2
 \le4\mathcal B_t(x;s).}
 \tag{V40.16}
\]
\end{lemma}

\begin{proof}
Let \(q(y)=\Im A_{t,1}(y;s)\).  Then
\[
 q'(y)=\Re A_{t,2}(y;s).
 \tag{V40.17}
\]
For \(y\in[x-h_t,x+h_t]\), the fundamental theorem of calculus and
Cauchy--Schwarz give
\[
 |q(x)|^2\le2|q(y)|^2
 +2h_t\int_{x-h_t}^{x+h_t}|q'(v)|^2\,dv.
 \tag{V40.18}
\]
Averaging in \(y\) yields
\[
 |q(x)|^2
 \le2\mathcal E_{t,1}^{-1}(x;s)
 +4h_t^2\mathcal E_{t,2}^{+1}(x;s)
 \le4\mathcal B_t(x;s).
 \tag{V40.19}
\]
\end{proof}

\begin{theorem}[Contact two-prime burst under failure of GRH]
\label{thm:v40-burst}
Assume the paired channel has an off-line zero, and let
\(p_*=(\gamma_*,\delta_*,m_*)\), \(t_k\downarrow0\), and
\(x_{k,\pm}\) be the centered exposed pair and phase-tuned contacts of
V37--V39.  For either sign and every fixed \(s\),
\[
 \boxed{
 \limsup_{k\to\infty}4t_k\log
 \bigl(1+\mathcal B_{t_k}(x_{k,\pm};s)\bigr)
 \ge2\delta_*^2.}
 \tag{V40.20}
\]
Moreover, subtracting from \textup{(V40.15)} the two
difference-frequency diagonals
\[
 \mathcal D_{t,1}(s)+h_t^2\mathcal D_{t,2}(s)
 \tag{V40.21}
\]
leaves a signed off-diagonal correlation remainder with the same positive
lower logarithmic rate.  Thus the burst cannot be carried by the diagonal.
\end{theorem}

\begin{proof}
Theorem~\ref{thm:v39-prefix-carries-slope} gives
\[
 \lim_{k\to\infty}4t_k\log
 |\Im P_{\chi,t_k}^{-,[1]}(x_{k,\pm};s)|
 =\delta_*^2.
 \tag{V40.22}
\]
Square this identity and apply Lemma~\ref{lem:v40-trace}.  The constant
The constant \(4\) has zero logarithmic rate, proving \textup{(V40.20)}.
By Proposition~\ref{prop:v40-two-prime}, \textup{(V40.21)} is bounded by a
fixed power of \(t^{-1}\), because \(h_t=t^{-2}\).  Subtracting a
zero-rate quantity cannot remove the positive rate in
\textup{(V40.20)}.  The remaining terms are exactly the \(m\ne n\)
part of the difference-frequency sums together with the sum-frequency sums
in \textup{(V40.9)}.
\end{proof}

This converts a one-prime pointwise obstruction into a two-prime energy
obstruction without using an absolute majorant.  The long polynomial window
in \(x\) resolves logarithmic frequency differences through the sinc kernel;
it costs no exponential factor.  The second reader is not cosmetic: it is
the exact derivative row required by the local trace inequality.

\begin{criterion}[Contact-correlation criterion equivalent to GRH]
\label{crit:v40-correlation}
GRH for the paired channel is equivalent to the following statement: for
every nondegenerate compact interval \(I\) and one, equivalently every,
fixed \(s\in\mathbb R\),
\[
 \boxed{
 \limsup_{t\downarrow0}4t\log\left(
 1+\sup_{x\in\mathcal C_{\chi,t}(I)}\mathcal B_t(x;s)
 \right)=0.}
 \tag{V40.23}
\]
As usual the supremum over an empty contact set is zero.
\end{criterion}

\begin{proof}
Under GRH the positive-time symbol has no zero-level contacts, so the
condition holds.  If GRH fails, choose \(I\) containing the V37 phase-tuned
contacts.  Theorem~\ref{thm:v40-burst} contradicts \textup{(V40.23)}.
\end{proof}

Criterion~\ref{crit:v40-correlation} is not claimed to be easier than GRH.
Its value is diagnostic: any proof through local energy must cancel the
explicit signed off-diagonal sums in \textup{(V40.9)} at points constrained
by \(m_{\chi,t}(x)=0\).  Large-sieve energy averaged over characters, absolute
spectral locality, or a finite positive atlas does not estimate this
contact-conditioned two-prime remainder.

\begin{firewall}[A correlation target is not correlation cancellation]
V40 performs the scheduled corpus sweep, imports no unsupported companion
claim, derives an exact two-prime expansion of the contact Sobolev packet,
and proves that an off-line zero forces an exponential signed off-diagonal
burst while the diagonal remains polynomial.  It does not bound that
off-diagonal remainder, exclude an off-line zero, or prove GRH.
\end{firewall}

\begin{assessment}[V40 acquisition and the next upstream route]
The window-localization branch is now closed as an engineering programme.
V38--V39 found a no-alias reader and its weakest contact sampling criterion;
the V40 sweep found no external theorem acting on its character phase.
Further changes of the one-prime window would leave the half-mass tax and the
same phase obstruction.

The new viable object is the contact-conditioned two-prime remainder in
\textup{(V40.9)}.  V41 should separate its ratio-frequency and
product-frequency pieces, exploit the congruence identities of the fixed
Dirichlet character before taking absolute values, and determine whether the
contact equation supplies a genuine linear cancellation of either piece.
The first required test is exact: compute the contribution of the contact
constraint to the sum-frequency term and exhibit a phase-saturating model if
no improvement follows.  This moves upstream to multiplicative correlations
rather than repeating localization or generic energy bounds.
\end{assessment}

\section*{Version V40 append-only record}
\addcontentsline{toc}{section}{Version V40 append-only record}
The complete V39 source was retained before this continuation.  It had
\(617302\) bytes, \(18687\) lines, and SHA--256
\[
 \texttt{1C93432AF78201CD7852DD91E1AC70DEBDEB0D1201092AE474584B28977F4E63}.
\]
V40 performed the scheduled corpus sweep, proved the exact contact
two-prime energy expansion and polynomial diagonal bound, and showed that
failure of GRH forces an exponentially large signed off-diagonal correlation
burst at centered phase-tuned contacts.

% =============================================================================
% END APPENDED V40 MATERIAL
% =============================================================================
% =============================================================================
% BEGIN APPENDED V41 MATERIAL
% =============================================================================

\section{V41: ratio-frequency reduction and the fixed-character residue collapse}
\label{sec:grh-v41}

V40 converted the localized contact derivative into a two-reader energy
whose exact expansion contains a Toeplitz difference-frequency block and a
Hankel sum-frequency block.  This split resembles the adaptive negative-cell
Gram of V26, but the present contact geometry permits a further reduction
which was not available there: the Hankel block can be discarded altogether.
Failure of GRH already forces an exponential burst in the positive Hermitian
mean square of the soft reader.

The remaining ratio block admits an exact decomposition by multiplicative
residue ratio modulo the fixed conductor.  That decomposition does not
produce an automatic character cancellation.  At fixed \(\chi\), its residue
matrix has rank one, and the contact equation fixes only the real part of the
zeroth reader in that same one-dimensional channel.  An explicit
phase-saturating moment model shows that this scalar constraint is compatible
with arbitrarily coherent derivative correlations even when every atom has
one residue label and the exact fixed-character phase.

\subsection{The positive ratio packet}

Retain \(A_{t,r}(y;s)\), \(h_t=t^{-2}\), and the component energies from
\textup{(V40.5)}--\textup{(V40.7)}.  Define
\[
 \mathcal R_{t,r}(x;s)
 :=\frac1{2h_t}\int_{x-h_t}^{x+h_t}
 |A_{t,r}(y;s)|^2\,dy
 \tag{V41.1}
\]
and
\[
 \mathcal H_t(x;s)
 :=\mathcal R_{t,1}(x;s)+h_t^2\mathcal R_{t,2}(x;s).
 \tag{V41.2}
\]

\begin{proposition}[Exact ratio-frequency representation]
\label{prop:v41-ratio}
For every \(t>0\), \(x\in\mathbb R\), and fixed \(r\ge0\),
\[
 \boxed{
 \mathcal R_{t,r}(x;s)
 =\Re\sum_{m,n\ge2}
 b_{m,t}^{[r]}(s)b_{n,t}^{[r]}(s)
 \chi(m)\overline{\chi(n)}
 e^{ix(u_m-u_n)}
 \operatorname{sinc}\!\bigl(h_t(u_m-u_n)\bigr).}
 \tag{V41.3}
\]
Moreover,
\[
 0\le \mathcal E_{t,r}^{\epsilon}(x;s)
 \le\mathcal R_{t,r}(x;s),
 \qquad \epsilon\in\{+1,-1\},
 \tag{V41.4}
\]
and hence
\[
 \boxed{0\le\mathcal B_t(x;s)\le\mathcal H_t(x;s).}
 \tag{V41.5}
\]
The diagonal of \(\mathcal H_t\) is
\[
 2\mathcal D_{t,1}(s)+2h_t^2\mathcal D_{t,2}(s)
 \ll t^{-15/2},
 \tag{V41.6}
\]
so it has zero \(1/t\)-logarithmic rate.
\end{proposition}

\begin{proof}
Expand \(|A_{t,r}|^2\), use absolute convergence, and apply
\textup{(V40.13)}.  This gives \textup{(V41.3)}.  Each real or imaginary
component has square at most the squared modulus, proving
\textup{(V41.4)}--\textup{(V41.5)}.  The diagonal in
\textup{(V41.3)} is
\(\sum_n|b_{n,t}^{[r]}|^2|\chi(n)|^2=2\mathcal D_{t,r}\).
Proposition~\ref{prop:v40-two-prime} gives
\[
 \mathcal D_{t,1}\ll t^{-5/2},
 \qquad
 h_t^2\mathcal D_{t,2}
 \ll t^{-4}t^{-7/2}=t^{-15/2},
 \]
which proves \textup{(V41.6)}.
\end{proof}

The representation contains no product frequency \(u_m+u_n\).  Positivity
has removed it before any entrywise absolute value is taken.

\begin{theorem}[A forbidden contact forces a ratio-only burst]
\label{thm:v41-ratio-burst}
Assume the paired channel has an off-line zero.  Let
\(p_*=(\gamma_*,\delta_*,m_*)\), \(t_k\downarrow0\), and
\(x_{k,\pm}\) be the centered exposed pair and phase-tuned contacts of
V37--V40.  For either sign and every fixed \(s\),
\[
 \boxed{
 \limsup_{k\to\infty}4t_k\log
 \bigl(1+\mathcal H_{t_k}(x_{k,\pm};s)\bigr)
 \ge2\delta_*^2.}
 \tag{V41.7}
\]
After subtracting the diagonal \textup{(V41.6)}, the off-diagonal part of
the Hermitian ratio sum \textup{(V41.3)} retains the same positive lower
rate.
\end{theorem}

\begin{proof}
Theorem~\ref{thm:v40-burst} gives the lower rate \(2\delta_*^2\) for
\(\mathcal B_{t_k}\).  Inequality \textup{(V41.5)} gives the same lower
bound for \(\mathcal H_{t_k}\).  The diagonal is polynomial by
\textup{(V41.6)}, so subtracting it cannot remove a positive exponential
rate.
\end{proof}

\begin{criterion}[Ratio-correlation criterion equivalent to GRH]
\label{crit:v41-ratio}
GRH for the paired channel is equivalent to the following statement: for
every nondegenerate compact interval \(I\) and one, equivalently every,
fixed \(s\in\mathbb R\),
\[
 \boxed{
 \limsup_{t\downarrow0}4t\log\left(
 1+\sup_{x\in\mathcal C_{\chi,t}(I)}
 \mathcal H_t(x;s)\right)=0.}
 \tag{V41.8}
\]
\end{criterion}

\begin{proof}
Under GRH the contact set is empty.  If GRH fails, choose \(I\) containing
the phase-tuned contact sequence and apply
Theorem~\ref{thm:v41-ratio-burst}.
\end{proof}

Thus V40's product-frequency branch is not a necessary GRH target.  The
live bilinear problem can be stated entirely in terms of ratios \(m/n\).
This is stronger than merely bounding the Hankel block by the Toeplitz
block: the positive packet eliminates the former at the level of the
criterion.

\subsection{Exact multiplicative residue-ratio decomposition}

Let
\[
 G_q:=(\mathbb Z/q\mathbb Z)^\times.
 \tag{V41.9}
\]
For \(c\in G_q\), define
\[
 \begin{aligned}
 \mathcal K_{t,r,c}(x;s)
 :=\sum_{\substack{m,n\ge2\\ (mn,q)=1\\ m\equiv cn\;(\bmod q)}}
 &b_{m,t}^{[r]}(s)b_{n,t}^{[r]}(s)
 e^{ix\log(m/n)}\\
 &{}\times
 \operatorname{sinc}\!\bigl(h_t\log(m/n)\bigr).
 \end{aligned}
 \tag{V41.10}
\]

\begin{theorem}[Fixed-character residue collapse]
\label{thm:v41-residue-collapse}
For every \(t>0\), \(x\in\mathbb R\), and fixed \(r\),
\[
 \boxed{
 \mathcal R_{t,r}(x;s)
 =\sum_{c\in G_q}\chi(c)\mathcal K_{t,r,c}(x;s).}
 \tag{V41.11}
\]
The right side is real because
\[
 \mathcal K_{t,r,c^{-1}}(x;s)
 =\overline{\mathcal K_{t,r,c}(x;s)},
 \qquad
 \chi(c^{-1})=\overline{\chi(c)}.
 \tag{V41.12}
\]
Equivalently, if
\[
 F_{t,r,a}(y;s)
 :=\sum_{\substack{n\ge2\\n\equiv a\;(\bmod q)}}
 b_{n,t}^{[r]}(s)e^{iyu_n},
 \qquad a\in G_q,
 \tag{V41.13}
\]
then
\[
 A_{t,r}(y;s)=\sum_{a\in G_q}\chi(a)F_{t,r,a}(y;s),
 \tag{V41.14}
\]
and the residue kernel in its squared modulus is
\[
 \mathsf M_\chi(a,b)=\chi(a)\overline{\chi(b)}.
 \tag{V41.15}
\]
The matrix \(\mathsf M_\chi\) has rank one and its sole nonzero eigenvalue
is \(\varphi(q)\).
\end{theorem}

\begin{proof}
When \((mn,q)=1\),
\[
 \chi(m)\overline{\chi(n)}
 =\chi(mn^{-1}\bmod q).
 \tag{V41.16}
\]
Group the pairs in \textup{(V41.3)} by
\(c=mn^{-1}\bmod q\) to obtain \textup{(V41.11)}.  Swapping \(m\) and
\(n\) proves \textup{(V41.12)}.  Grouping the one-prime sum by its residue
gives \textup{(V41.14)}.  Finally
\(\mathsf M_\chi=vv^*\), where \(v_a=\chi(a)\); hence its rank is one and
its nonzero eigenvalue is
\(\|v\|^2=\sum_{a\in G_q}|\chi(a)|^2=\varphi(q)\).
\end{proof}

This calculation identifies what fixed-conductor periodicity actually
does.  It does not average over all character modes.  It selects one
rank-one Fourier channel of the residue vector.  Orthogonality over a
complete character family would diagonalize an average, but V32 already
showed why such an average does not control the exceptional fixed channel
at its own moving contacts.

\subsection{What the contact equation contributes}

Let \(F_{t,0,a}^{\pm}\) denote the residue sums
\textup{(V41.13)} formed with \(W_t^\pm\).  At an actual contact,
\textup{(V39.6)} becomes
\[
 \boxed{
 \Re\sum_{a\in G_q}\chi(a)
 \bigl(F_{t,0,a}^-(x;s)+F_{t,0,a}^+(x;s)\bigr)
 =\frac{a_{\chi,t}(x)}2.}
 \tag{V41.17}
\]
This is one real constraint on the zeroth-reader rank-one coordinate at
one point.  The packet \(\mathcal H_t\) instead uses the first and second
soft-prefix coordinates throughout an interval of length \(2h_t\).

Even for the full zeroth reader
\(P_{\chi,t}^{[0]}(x)=c+id\), the contact condition fixes only
\(c=a_{\chi,t}(x)/2\).  Its pointwise ratio and product contractions obey
\[
 |P_{\chi,t}^{[0]}(x)|^2=c^2+d^2,
 \qquad
 \Re\bigl(P_{\chi,t}^{[0]}(x)^2\bigr)=c^2-d^2.
 \tag{V41.18}
\]
Thus
\[
 |P_{\chi,t}^{[0]}(x)|^2+
 \Re\bigl(P_{\chi,t}^{[0]}(x)^2\bigr)=2c^2
 \tag{V41.19}
\]
The complementary difference of these contractions is \(2d^2\), and
neither identity constrains the first reader.  For the soft
prefix alone, \textup{(V41.17)} supplies still less: its real value can
trade against the soft tail.

The following finite model verifies that no omitted residue identity turns
\textup{(V41.17)} into derivative cancellation.

\begin{proposition}[Same-residue phase saturation]
\label{prop:v41-phase-model}
Fix \(a\in G_q\) and write \(\chi(a)=e^{i\theta}\).  There are finite
positive-weight moment readers, all assigned the single residue label
\(a\), for which the zeroth contact value is zero while the imaginary first
reader has coherent off-diagonal growth.

More precisely, choose a fixed integer \(M\) large enough that all
frequencies below are positive and, for \(1\le j\le N\), set
\[
 u_j=\frac\pi2-\theta+2\pi(M+N+j),
 \qquad
 v_j=-\frac\pi2-\theta+2\pi(M+j).
 \tag{V41.20}
\]
Define
\[
 Q_{N,r}(x)
 :=\sum_{j=1}^N\chi(a)
 \bigl(u_j^r e^{ixu_j}+v_j^r e^{ixv_j}\bigr).
 \tag{V41.21}
\]
Then
\[
 \boxed{Q_{N,0}(1)=0,}
 \qquad
 \boxed{\Im Q_{N,1}(1)=\pi N(2N+1).}
 \tag{V41.22}
\]
Meanwhile the first-reader diagonal satisfies
\[
 \sum_{j=1}^N(u_j^2+v_j^2)=O(N^3),
 \tag{V41.23}
\]
whereas
\[
 |\Im Q_{N,1}(1)|^2=4\pi^2N^4+O(N^3).
 \tag{V41.24}
\]
Hence the coherent excess over the diagonal is of order \(N^4\).  Any
prescribed real zeroth value can be obtained by adding one fixed positive
atom whose phase at \(x=1\) is \(0\) or \(\pi\).
\end{proposition}

\begin{proof}
By construction,
\[
 \chi(a)e^{iu_j}=i,
 \qquad
 \chi(a)e^{iv_j}=-i.
 \tag{V41.25}
\]
The zeroth terms therefore cancel pairwise.  The imaginary first reader is
\[
 \sum_{j=1}^N(u_j-v_j)
 =\sum_{j=1}^N(\pi+2\pi N)
 =\pi N(2N+1),
 \]
which proves \textup{(V41.22)}.  All frequencies are \(O(N)\), so the sum
of their squares is \(O(N^3)\), while squaring the displayed first reader
gives \textup{(V41.24)}.  For a real target \(c\), add an atom of weight
\(|c|\) and choose its positive frequency modulo \(2\pi\) so that its
phase after multiplication by \(\chi(a)\) is \(\operatorname{sgn}c\).
\end{proof}

The proposition is an algebraic capacity witness, not a model of the prime
distribution: its frequencies are synthetic and its weights are not the
Gaussian von Mangoldt weights.  Its precise conclusion is narrower and
useful.  Fixed-character congruence, exact moment differentiation, and the
contact hyperplane alone do not cancel the ratio off-diagonal.  Any such
cancellation must use the arithmetic placement and weights of the actual
prime powers.

\begin{theorem}[V41 ratio verdict]
\label{thm:v41-verdict}
The contact two-prime problem has the following exact reduction.
\begin{enumerate}[label=\textup{(\roman*)},leftmargin=3em]
\item The sum-frequency/Hankel block is unnecessary for a GRH-equivalent
      contact criterion.
\item Failure of GRH forces rate at least \(2\delta_*^2\) in the
      off-diagonal Hermitian ratio packet.
\item Modulo the conductor, that packet is one rank-one character
      projection of multiplicative residue-ratio correlations.
\item The contact equation fixes one real zeroth-reader coordinate and
      gives no algebraic cancellation of the derivative ratio packet.
\end{enumerate}
\end{theorem}

\begin{proof}
These are respectively Proposition~\ref{prop:v41-ratio} and
Criterion~\ref{crit:v41-ratio}, Theorem~\ref{thm:v41-ratio-burst},
Theorem~\ref{thm:v41-residue-collapse}, and
Proposition~\ref{prop:v41-phase-model}.
\end{proof}

\begin{firewall}[Residue grouping is not a fixed-channel bound]
V41 removes the product-frequency branch, derives the exact
multiplicative-ratio decomposition, and proves that neither fixed-character
rank-one algebra nor the scalar contact equation forces cancellation.
It does not bound the actual off-diagonal ratio sums, exclude an off-line
zero, or prove GRH.
\end{firewall}

\begin{assessment}[V41 acquisition and the near-ratio target]
The V40 target is now strictly narrower.  Only the Hermitian correlations
\[
 \chi(m)\overline{\chi(n)}
 e^{ix\log(m/n)}
 \operatorname{sinc}\!\bigl(t^{-2}\log(m/n)\bigr)
 \tag{V41.26}
\]
need be studied.  Their diagonal is harmless, their fixed-conductor
decomposition is exact, and failure of GRH forces their off-diagonal
contraction to have positive exponential rate at contacts.

V42 should replace the sharp observation interval by a positive
Fourier-localizing window, so that the ratio kernel is nonnegative and has
rapid decay.  It should then split genuinely near ratios from the far tail,
group the near pairs by \(m\equiv cn\pmod q\), and determine the smallest
observation radius which makes the diagonal dominant without paying a
positive \(1/t\)-rate in the point-evaluation inequality.  This is an
uncertainty calculation on actual prime-log spacings, not another generic
large-sieve average.
\end{assessment}

\section*{Version V41 append-only record}
\addcontentsline{toc}{section}{Version V41 append-only record}
The complete V40 source was retained before this continuation.  It had
\(631315\) bytes, \(19074\) lines, and SHA--256
\[
 \texttt{EC976BAAD476A3EFB939BB6F79F70C741864DEA596D5194CDF80C4B36B2CF375}.
\]
V41 eliminated the Hankel/product-frequency branch from the contact
criterion, isolated the Hermitian ratio-only burst, decomposed it by
multiplicative residue ratio, and proved that fixed-character rank-one
algebra plus the contact equation supplies no automatic derivative
cancellation.

% =============================================================================
% END APPENDED V41 MATERIAL
% =============================================================================
% =============================================================================
% BEGIN APPENDED V42 MATERIAL
% =============================================================================

\section{V42: Gaussian near-ratio localization and the no-free-resolution theorem}
\label{sec:grh-v42}

V41 reduced the contact obstruction to the positive Hermitian ratio packet
and decomposed it by multiplicative residue ratio.  The sharp observation
interval left a sinc kernel with a slowly decaying tail.  This section
replaces that interval by a Gaussian observation window.  The resulting
ratio kernel is positive and rapidly decreasing.

Two exact facts emerge.  First, with the zero-cost polynomial radius
\(h_t=t^{-2}\), every correlation outside
\[
 |\log(m/n)|\le t^{3/2}
\]
is exponentially negligible even against the full absolute prime mass.
Thus failure of GRH forces its V41 ratio burst into genuinely near ratios.
Second, this band is nowhere near diagonal: a still narrower core of
logarithmic width \(t^2\) contains exponentially many primes in each fixed
invertible residue class, and its same-residue contraction has the complete
squared absolute rate \(1/2\).  Resolving all such prime logarithms requires
an observation radius of rate at least \(1\), whereas point evaluation
requires a radius of rate zero.  Diagonalization and zero-cost recovery are
therefore incompatible.

\subsection{The Gaussian observation packet}

For \(h>0\), put
\[
 g_h(z):=\frac1{\sqrt\pi h}e^{-z^2/h^2}.
 \tag{V42.1}
\]
Its Fourier transform is
\[
 \int_{\mathbb R}g_h(y-x)e^{iyv}\,dy
 =e^{ixv}e^{-h^2v^2/4}.
 \tag{V42.2}
\]
Retain the soft reader \(A_{t,r}\) of \textup{(V40.5)} and set
\[
 h_t:=t^{-2},
 \tag{V42.3}
\]
\[
 \mathcal R^{\rm G}_{t,r}(x;s)
 :=\int_{\mathbb R}g_{h_t}(y-x)
 |A_{t,r}(y;s)|^2\,dy,
 \tag{V42.4}
\]
\[
 \mathcal H^{\rm G}_t(x;s)
 :=\mathcal R^{\rm G}_{t,1}(x;s)
 +h_t^2\mathcal R^{\rm G}_{t,2}(x;s).
 \tag{V42.5}
\]

\begin{proposition}[Positive Gaussian ratio representation]
\label{prop:v42-gaussian-ratio}
For every \(t>0\), \(x\in\mathbb R\), and fixed \(r\ge0\),
\[
 \boxed{
 \mathcal R^{\rm G}_{t,r}(x;s)
 =\Re\sum_{m,n\ge2}
 b_{m,t}^{[r]}(s)b_{n,t}^{[r]}(s)
 \chi(m)\overline{\chi(n)}
 e^{ix\log(m/n)}
 e^{-h_t^2\log^2(m/n)/4}.}
 \tag{V42.6}
\]
Its diagonal is \(2\mathcal D_{t,r}(s)\), hence polynomial by
Proposition~\ref{prop:v40-two-prime}.
\end{proposition}

\begin{proof}
Expand the squared modulus in \textup{(V42.4)} and use
\textup{(V42.2)}.  Gaussian damping gives absolute convergence of the
one-prime series, and hence of the expanded integral.  The terms with
\(m=n\) give \(\sum_n|b_{n,t}^{[r]}|^2|\chi(n)|^2
=2\mathcal D_{t,r}\).
\end{proof}

\begin{lemma}[Gaussian trace recovery]
\label{lem:v42-gaussian-trace}
There is an absolute constant \(C_{\rm tr}=2e\sqrt\pi\) such that
\[
 \boxed{
 |\Im P_{\chi,t}^{-,[1]}(x;s)|^2
 \le C_{\rm tr}\mathcal H^{\rm G}_t(x;s)}
 \tag{V42.7}
\]
for every \(t>0\) and real \(x\).
\end{lemma}

\begin{proof}
Let \(q(y)=\Im A_{t,1}(y;s)\), so that
\(q'(y)=\Re A_{t,2}(y;s)\).  On
\(I=[x-h_t,x+h_t]\), the argument of Lemma~\ref{lem:v40-trace} gives
\[
 |q(x)|^2
 \le\frac1{h_t}\int_I|q(y)|^2\,dy
 +2h_t\int_I|q'(y)|^2\,dy.
 \tag{V42.8}
\]
Since
\[
 g_{h_t}(y-x)\ge\frac{e^{-1}}{\sqrt\pi h_t}
 \qquad(y\in I),
 \tag{V42.9}
\]
the first integral in \textup{(V42.8)} is at most
\(e\sqrt\pi h_t\mathcal R^{\rm G}_{t,1}\), and the second is at most
\(e\sqrt\pi h_t\mathcal R^{\rm G}_{t,2}\).  Substitution proves
\textup{(V42.7)}.
\end{proof}

\begin{theorem}[Gaussian ratio burst at a forbidden contact]
\label{thm:v42-gaussian-burst}
Assume the paired channel has an off-line zero, and retain the centered
exposed pair \(p_*=(\gamma_*,\delta_*,m_*)\) and its phase-tuned contacts
\(x_{k,\pm}\) at \(t_k\downarrow0\).  For either sign and every fixed
\(s\),
\[
 \boxed{
 \limsup_{k\to\infty}4t_k\log
 \bigl(1+\mathcal H^{\rm G}_{t_k}(x_{k,\pm};s)\bigr)
 \ge2\delta_*^2.}
 \tag{V42.10}
\]
Consequently GRH is equivalent to the zero-rate bound for
\(\mathcal H^{\rm G}_t\) on every compact-interval contact set.
\end{theorem}

\begin{proof}
Combine Lemma~\ref{lem:v42-gaussian-trace} with the exact contact growth
\textup{(V39.8)}.  The multiplicative constant has zero logarithmic rate.
Under GRH the contact sets are empty; under failure the displayed contact
sequence violates a zero-rate bound.
\end{proof}

\subsection{The far ratio tail is exponentially negligible}

Define the absolute soft masses
\[
 \mathcal S_{t,r}(s)
 :=\sum_{n\ge2}|b_{n,t}^{[r]}(s)|.
 \tag{V42.11}
\]
The half-mass law \textup{(V39.18)} gives, for every fixed \(r\),
\[
 \lim_{t\downarrow0}4t\log\mathcal S_{t,r}(s)=\frac14.
 \tag{V42.12}
\]
Let
\[
 \rho_t:=t^{3/2}.
 \tag{V42.13}
\]
Split the double sum \textup{(V42.6)} into
\(\mathcal R^{{\rm G},{\rm near}}_{t,r}\) and
\(\mathcal R^{{\rm G},{\rm far}}_{t,r}\) according as
\[
 |\log(m/n)|\le\rho_t
 \quad\hbox{or}\quad
 |\log(m/n)|>\rho_t.
 \tag{V42.14}
\]
Define \(\mathcal H^{{\rm G},{\rm near}}_t\) and
\(\mathcal H^{{\rm G},{\rm far}}_t\) by the same first-plus-second-reader
combination as in \textup{(V42.5)}.

\begin{theorem}[Exact near-ratio localization]
\label{thm:v42-near-localization}
For every fixed \(s\), uniformly in \(x\in\mathbb R\),
\[
 \boxed{
 |\mathcal H^{{\rm G},{\rm far}}_t(x;s)|
 \le e^{-1/(4t)}
 \bigl(\mathcal S_{t,1}(s)^2+
       h_t^2\mathcal S_{t,2}(s)^2\bigr).}
 \tag{V42.15}
\]
In particular,
\[
 \boxed{
 \limsup_{t\downarrow0}4t\log
 |\mathcal H^{{\rm G},{\rm far}}_t(x;s)|
 \le-\frac12}
 \tag{V42.16}
\]
whenever the far part is nonzero.

If GRH fails, the contacts in Theorem~\ref{thm:v42-gaussian-burst} satisfy
\[
 \boxed{
 \limsup_{k\to\infty}4t_k\log
 \bigl(1+
 \mathcal H^{{\rm G},{\rm near}}_{t_k}(x_{k,\pm};s)\bigr)
 \ge2\delta_*^2.}
 \tag{V42.17}
\]
The diagonal is polynomial, so the same positive rate occurs in the
off-diagonal near-ratio remainder.
\end{theorem}

\begin{proof}
On the far set,
\[
 e^{-h_t^2\log^2(m/n)/4}
 \le e^{-h_t^2\rho_t^2/4}
 =e^{-1/(4t)}.
 \tag{V42.18}
\]
Take absolute values after this factor has been extracted to obtain
\textup{(V42.15)}.  Equation \textup{(V42.12)} says that either squared
absolute mass has \(4t\)-logarithmic rate \(1/2\), while \(h_t^2\) is only
polynomial.  The Gaussian factor contributes rate \(-1\), proving
\textup{(V42.16)}.  Subtract the far part from
\textup{(V42.10)} to obtain \textup{(V42.17)}.  Finally use the polynomial
diagonal bound from Proposition~\ref{prop:v42-gaussian-ratio}.
\end{proof}

Thus the contact obstruction is not merely a generic two-prime
correlation.  It is supported, at exponential scale, on prime powers whose
ratio tends to one at the explicit polynomial rate
\(m/n=\exp(O(t^{3/2}))\).

\subsection{Residue ratios inside the near band}

For \(c\in G_q\), let
\[
 \begin{aligned}
 \mathcal K^{{\rm G},{\rm near}}_{t,r,c}(x;s)
 :=\sum_{\substack{m,n\ge2\\ (mn,q)=1\\
                   m\equiv cn\;(\bmod q)\\
                   |\log(m/n)|\le\rho_t}}
 &b_{m,t}^{[r]}(s)b_{n,t}^{[r]}(s)
 e^{ix\log(m/n)}\\
 &{}\times e^{-h_t^2\log^2(m/n)/4}.
 \end{aligned}
 \tag{V42.19}
\]
Then the V41 grouping gives the exact identity
\[
 \boxed{
 \mathcal R^{{\rm G},{\rm near}}_{t,r}(x;s)
 =\sum_{c\in G_q}
 \chi(c)\mathcal K^{{\rm G},{\rm near}}_{t,r,c}(x;s).}
 \tag{V42.20}
\]

The identity class \(c=1\) is already exponentially populated.  Put
\[
 U_t:=\frac1{4t},
 \qquad
 \ell_t:=\frac12t^2,
 \tag{V42.21}
\]
and, for \(a\in G_q\), let
\[
 \mathcal P_{t,a}
 :=\{p\ {\rm prime}:e^{U_t}\le p\le e^{U_t+\ell_t},
 \ p\equiv a\pmod q\}.
 \tag{V42.22}
\]

\begin{lemma}[Exponentially many same-residue primes in one kernel core]
\label{lem:v42-prime-core}
For every fixed \(a\in G_q\),
\[
 \boxed{
 \#\mathcal P_{t,a}
 \sim\frac{e^{U_t}\ell_t}{\varphi(q)U_t},}
 \qquad
 \boxed{
 \lim_{t\downarrow0}4t\log\#\mathcal P_{t,a}=1.}
 \tag{V42.23}
\]
Uniformly for \(p\in\mathcal P_{t,a}\) and fixed \(r,s\),
\[
 \boxed{
 \lim_{t\downarrow0}4t\log b_{p,t}^{[r]}(s)
 =-\frac34.}
 \tag{V42.24}
\]
\end{lemma}

\begin{proof}
The fixed-modulus prime number theorem in progressions, with its classical
zero-free-region error, is uniform on the interval in
\textup{(V42.22)} because
\[
 \ell_t=t^2/2\gg e^{-c/\sqrt t}.
 \tag{V42.25}
\]
Subtracting the two endpoint asymptotics gives the first formula in
\textup{(V42.23)}.  Since \(U_t=1/(4t)\), its logarithm is
\(1/(4t)+O(\log(1/t))\), proving the second.

For \(u=\log p=U_t+O(t^2)\), the mesoscopic window satisfies
\[
 W_t^-(p)\longrightarrow\frac12
 \tag{V42.26}
\]
uniformly: the displacement from the critical cutoff divided by the
transition width is \(s/\omega_t+o(1)\).  Polynomial logarithmic factors
have zero rate, while
\[
 -\frac u2-tu^2
 =-\frac1{8t}-\frac1{16t}+o(1/t)
 =-\frac3{16t}+o(1/t).
 \tag{V42.27}
\]
Multiplication by \(4t\) proves \textup{(V42.24)}.
\end{proof}

\begin{theorem}[The identity residue-ratio block has full squared type]
\label{thm:v42-identity-block}
Fix a compact interval \(I\), \(r\ge0\), and \(s\).  Uniformly for
\(x\in I\), the prime-core subblock
\[
 \begin{aligned}
 \mathcal Q_{t,r,a}(x;s)
 :=\sum_{p,p'\in\mathcal P_{t,a}}
 &b_{p,t}^{[r]}(s)b_{p',t}^{[r]}(s)
 e^{ix\log(p/p')}\\
 &{}\times e^{-h_t^2\log^2(p/p')/4}
 \end{aligned}
 \tag{V42.28}
\]
is real and positive for all sufficiently small \(t\), and
\[
 \boxed{
 \lim_{t\downarrow0}4t\log\mathcal Q_{t,r,a}(x;s)
 =\frac12.}
 \tag{V42.29}
\]
Moreover, the complete \(c=1\) near block is positive on \(I\) for small
\(t\) and has the same exact rate:
\[
 \boxed{
 \lim_{t\downarrow0}4t\log
 \mathcal K^{{\rm G},{\rm near}}_{t,r,1}(x;s)
 =\frac12.}
 \tag{V42.30}
\]
\end{theorem}

\begin{proof}
Swapping \(p,p'\) shows that \(\mathcal Q_{t,r,a}\) is real.  In its
range,
\[
 |\log(p/p')|\le\ell_t=\frac1{2h_t},
 \]
so the Gaussian kernel is at least \(e^{-1/16}\).  Also
\(\cos(x\log(p/p'))\to1\) uniformly on \(I\).  Hence, for small \(t\),
\[
 \mathcal Q_{t,r,a}(x;s)
 \asymp_I
 \left(\sum_{p\in\mathcal P_{t,a}}
 b_{p,t}^{[r]}(s)\right)^2.
 \tag{V42.31}
\]
Equations \textup{(V42.23)}--\textup{(V42.24)} give rate
\[
 2\left(1-\frac34\right)=\frac12.
 \]

For every pair in the \(c=1\) near block,
\(|\log(m/n)|\le\rho_t\to0\).  Thus
\(\cos(x\log(m/n))>0\) uniformly on \(I\) for small \(t\); the block is
positive and dominates the prime core.  Its upper bound by
\(\mathcal S_{t,r}^2\) has rate \(1/2\) by \textup{(V42.12)}.  This proves
\textup{(V42.30)}.
\end{proof}

The theorem gives the exact cancellation burden in
\textup{(V42.20)}.  The identity ratio class is positive and already has
the full absolute squared rate.  Any zero-rate bound for the fixed
character must cancel it against the nonidentity residue-ratio classes
before taking absolute values.  Estimating the classes separately cannot
reach the GRH criterion.

\subsection{Resolving to the diagonal has an unavoidable rate cost}

The abundance above can be sharpened to a spacing obstruction.  Let
\(N_t=e^{1/(4t)}\).  The fixed-modulus PNT gives
\[
 \#\{p\in[N_t,2N_t]:p\equiv a\pmod q\}
 \asymp_q\frac{N_t}{\log N_t}.
 \tag{V42.32}
\]
Pigeonholing their logarithms in an interval of length \(\log2\) produces
two distinct such primes with
\[
 0<|\log(p/p')|
 \ll_q\frac{\log N_t}{N_t}.
 \tag{V42.33}
\]

\begin{theorem}[No-free-resolution theorem]
\label{thm:v42-no-free-resolution}
Let \(H_t\ge1\) be any Gaussian observation radius.  If its Fourier kernel
resolves every distinct same-residue prime pair in
\([N_t,2N_t]\), in the sense that
\[
 e^{-H_t^2\log^2(p/p')/4}\longrightarrow0
 \tag{V42.34}
\]
uniformly over those pairs, then
\[
 \boxed{
 \liminf_{t\downarrow0}4t\log H_t\ge1.}
 \tag{V42.35}
\]
By contrast, an observation radius has zero cost in the trace packet only
if
\[
 4t\log H_t\longrightarrow0.
 \tag{V42.36}
\]
The two requirements are incompatible.  Indeed, once diagonal dominance
has been obtained at a radius satisfying \textup{(V42.35)}, the derivative
factor \(H_t^2\) alone contributes logarithmic rate at least \(2\).
\end{theorem}

\begin{proof}
Apply the uniform resolution assumption to the close pair in
\textup{(V42.33)}.  Necessarily
\[
 H_t\frac{\log N_t}{N_t}\longrightarrow\infty.
 \tag{V42.37}
\]
Taking logarithms and using
\(\log N_t=1/(4t)\) gives \textup{(V42.35)}.  The trace proof of
Lemma~\ref{lem:v42-gaussian-trace}, with \(H_t\) in place of \(h_t\),
multiplies the second-reader energy by \(H_t^2\).  Avoiding a positive
exponential debit therefore requires \textup{(V42.36)}.  If the mean square
has been diagonalized, its second-reader diagonal is bounded below by one
fixed unramified prime term for all small \(t\), while
\[
 \liminf 4t\log H_t^2\ge2.
 \]
This proves the final assertion.
\end{proof}

\begin{firewall}[Near-ratio localization is not near-ratio cancellation]
V42 replaces the sinc packet by a positive Gaussian ratio kernel, removes
the far ratio tail at exponential scale, identifies an exponentially
populated same-residue core of exact rate \(1/2\), and proves that resolving
that core to the diagonal is incompatible with zero-cost point recovery.
It does not cancel the nonidentity residue-ratio classes, exclude an
off-line zero, or prove GRH.
\end{firewall}

\begin{assessment}[V42 acquisition and the cross-residue target]
The ratio obstruction is now localized both analytically and
arithmetically.  Only pairs with
\(|\log(m/n)|\le t^{3/2}\) can carry a forbidden contact burst.  Inside
that band, the identity residue class has full squared type \(1/2\);
therefore the desired zero-rate estimate is a cancellation statement
between multiplicative residue-ratio classes, not a diagonal estimate.

V43 should write the vector of near kernels
\[
 \bigl(\mathcal K^{{\rm G},{\rm near}}_{t,r,c}\bigr)_{c\in G_q}
 \tag{V42.38}
\]
in the complete character Fourier basis on \(G_q\).  It should determine
exactly which character modes retain the positive identity-class saddle,
compare the principal and target-character modes before absolute values,
and test whether the contact equation couples those modes or merely fixes
one zeroth-reader coordinate.  Averaging away the target mode is forbidden;
the objective is an exact mode-to-mode identity or a rigorous capacity
obstruction.
\end{assessment}

\section*{Version V42 append-only record}
\addcontentsline{toc}{section}{Version V42 append-only record}
The complete V41 source was retained before this continuation.  It had
\(644891\) bytes, \(19483\) lines, and SHA--256
\[
 \texttt{C0E520668D0B20F757924B475560308946A73EA96809439305CF2C7B609A56D7}.
\]
V42 introduced the positive Gaussian contact packet, localized every
forbidden ratio burst to \(|\log(m/n)|\le t^{3/2}\), proved that the
same-residue near block has exact rate \(1/2\), and established the
no-free-resolution theorem for actual prime logarithms.

% =============================================================================
% END APPENDED V42 MATERIAL
% =============================================================================
% =============================================================================
% BEGIN APPENDED V43 MATERIAL
% =============================================================================

\section{V43: complete character modes and closure of the ratio-correlation branch}
\label{sec:grh-v43}

V42 localized the Hermitian contact packet to near ratios and proved that
the identity residue-ratio class has exact rate \(1/2\).  The proposed next
step was to place the entire residue-ratio vector in the complete character
Fourier basis.  That calculation can be done exactly.

It closes, rather than advances, the correlation branch.  Every character
Fourier coefficient of the ratio vector is precisely the Gaussian mean
square of the corresponding one-prime character reader.  The principal
mode retains the full squared absolute type, while the target
nonprincipal mode is exactly the square of the progression-error projection
already isolated in V29.  Parseval controls the sum of all modes but gives
no control of one prescribed mode.  The moving contact equation also stays
inside the target mode and fixes only its real zeroth coordinate.  It
creates no coupling to the principal mode or to the first-reader energy.

\subsection{Full ratio kernels and their character transform}

Write
\[
 d:=|G_q|=\varphi(q),
 \qquad
 \widehat G_q:=\{\eta:G_q\to\mathbb C^\times
                 \text{ a group character}\}.
 \tag{V43.1}
\]
Extend every \(\eta\in\widehat G_q\) by zero on integers not coprime to
\(q\).  With the V42 Gaussian radius \(h_t=t^{-2}\), define the full
residue-ratio kernels
\[
 \begin{aligned}
 \mathcal K^{\rm G}_{t,r,c}(x;s)
 :=\sum_{\substack{m,n\ge2\\ (mn,q)=1\\
                   m\equiv cn\;(\bmod q)}}
 &b_{m,t}^{[r]}(s)b_{n,t}^{[r]}(s)
 e^{ix\log(m/n)}\\
 &{}\times e^{-h_t^2\log^2(m/n)/4}.
 \end{aligned}
 \tag{V43.2}
\]
Their finite group Fourier transform is
\[
 \widehat{\mathcal K}_{t,r}(\eta;x,s)
 :=\sum_{c\in G_q}\eta(c)
 \mathcal K^{\rm G}_{t,r,c}(x;s).
 \tag{V43.3}
\]
Define the soft reader in the \(\eta\)-channel by
\[
 A_{\eta,t,r}(y;s)
 :=\sum_{n\ge2}b_{n,t}^{[r]}(s)\eta(n)e^{iy\log n}
 \tag{V43.4}
\]
and its Gaussian mean square by
\[
 \mathcal R^{\rm G}_{\eta,t,r}(x;s)
 :=\int_{\mathbb R}g_{h_t}(y-x)
 |A_{\eta,t,r}(y;s)|^2\,dy.
 \tag{V43.5}
\]

\begin{theorem}[Character modes are one-prime mean squares]
\label{thm:v43-mode-square}
For every \(\eta\in\widehat G_q\),
\[
 \boxed{
 \widehat{\mathcal K}_{t,r}(\eta;x,s)
 =\mathcal R^{\rm G}_{\eta,t,r}(x;s)\ge0.}
 \tag{V43.6}
\]
Fourier inversion gives
\[
 \boxed{
 \mathcal K^{\rm G}_{t,r,c}(x;s)
 =\frac1d\sum_{\eta\in\widehat G_q}
 \overline{\eta(c)}
 \mathcal R^{\rm G}_{\eta,t,r}(x;s).}
 \tag{V43.7}
\]
In particular,
\[
 \boxed{
 \mathcal K^{\rm G}_{t,r,1}(x;s)
 =\frac1d\sum_{\eta\in\widehat G_q}
 \mathcal R^{\rm G}_{\eta,t,r}(x;s).}
 \tag{V43.8}
\]
\end{theorem}

\begin{proof}
For coprime \(m,n\),
\[
 \eta(m)\overline{\eta(n)}
 =\eta(mn^{-1}\bmod q).
 \]
Expanding \textup{(V43.5)}, applying the Gaussian Fourier identity
\textup{(V42.2)}, and grouping by \(c=mn^{-1}\bmod q\) gives
\textup{(V43.6)}.  The remaining formulas are finite Fourier inversion on
\(G_q\).  Positivity follows from the integral in \textup{(V43.5)}.
\end{proof}

Thus the positive identity block is the average of positive character-mode
energies.  It is not itself an estimate for the target mode.

\subsection{Parseval and exact modal freedom}

For \(a\in G_q\), let
\[
 F_{t,r,a}(y;s)
 :=\sum_{\substack{n\ge2\\n\equiv a\;(\bmod q)}}
 b_{n,t}^{[r]}(s)e^{iy\log n}.
 \tag{V43.9}
\]
Then
\[
 A_{\eta,t,r}(y;s)
 =\sum_{a\in G_q}\eta(a)F_{t,r,a}(y;s).
 \tag{V43.10}
\]

\begin{proposition}[Character Parseval has no exceptional-mode lower bound]
\label{prop:v43-parseval}
For every \(y,t,r,s\),
\[
 \boxed{
 \sum_{\eta\in\widehat G_q}
 |A_{\eta,t,r}(y;s)|^2
 =d\sum_{a\in G_q}|F_{t,r,a}(y;s)|^2.}
 \tag{V43.11}
\]
After Gaussian averaging,
\[
 \boxed{
 \sum_{\eta\in\widehat G_q}
 \mathcal R^{\rm G}_{\eta,t,r}(x;s)
 =d\sum_{a\in G_q}
 \int_{\mathbb R}g_{h_t}(y-x)
 |F_{t,r,a}(y;s)|^2\,dy.}
 \tag{V43.12}
\]
No inequality depending only on this identity can bound one prescribed
nonprincipal mode by the principal mode.  Indeed, for arbitrary complex
modal data \(z_\eta\), the residue vector
\[
 F_a=\frac1d\sum_{\eta\in\widehat G_q}
 \overline{\eta(a)}z_\eta
 \tag{V43.13}
\]
satisfies
\[
 \sum_{a\in G_q}\eta(a)F_a=z_\eta.
 \tag{V43.14}
\]
\end{proposition}

\begin{proof}
Character orthogonality gives
\[
 \sum_{\eta\in\widehat G_q}
 \eta(a)\overline{\eta(b)}
 =d\,\mathbf1_{a=b}.
 \]
Insert \textup{(V43.10)} to prove \textup{(V43.11)}, then integrate to
obtain \textup{(V43.12)}.  Fourier inversion proves
\textup{(V43.13)}--\textup{(V43.14)}.
\end{proof}

For example, a residue vector constant in \(a\) populates only the
principal mode.  A vector proportional to \(\overline{\chi(a)}\) populates
only the \(\chi\)-mode.  Positivity of all mode energies therefore does not
couple them.

\subsection{Principal type versus the target mode}

Let \(\mathbf1\) denote the principal character of \(G_q\), and set
\[
 \mathcal H^{\rm G}_{\eta,t}(x;s)
 :=\mathcal R^{\rm G}_{\eta,t,1}(x;s)
 +h_t^2\mathcal R^{\rm G}_{\eta,t,2}(x;s).
 \tag{V43.15}
\]

\begin{theorem}[Principal mode and target-mode alternatives]
\label{thm:v43-modal-rates}
For every compact interval \(I\) and fixed \(s\),
\[
 \boxed{
 \lim_{t\downarrow0}4t\log
 \mathcal H^{\rm G}_{\mathbf1,t}(x;s)=\frac12}
 \tag{V43.16}
\]
uniformly for \(x\in I\).

For a nonprincipal \(\eta\), GRH in that channel implies
\[
 \boxed{
 \sup_{x\in I}\mathcal H^{\rm G}_{\eta,t}(x;s)
 \ll_{\eta,I,s}t^{-M}}
 \tag{V43.17}
\]
for some \(M<\infty\).  On the other hand, if the target \(\chi\)-channel
has an off-line zero, its centered phase-tuned contacts satisfy
\[
 \boxed{
 \limsup_{k\to\infty}4t_k\log
 \mathcal H^{\rm G}_{\chi,t_k}(x_{k,\pm};s)
 \ge2\delta_*^2.}
 \tag{V43.18}
\]
\end{theorem}

\begin{proof}
For the principal mode, restrict the near-ratio expansion to the prime core
of Lemma~\ref{lem:v42-prime-core}.  Every coefficient is positive, the
Gaussian kernel is bounded below, and
\(\cos(x\log(p/p'))\to1\) uniformly on \(I\).  The lower rate is \(1/2\)
by Theorem~\ref{thm:v42-identity-block}; the upper bound by the squared
absolute masses \(\mathcal S_{t,r}^2\) has the same rate.  The far part is
negligible by Theorem~\ref{thm:v42-near-localization}.  Polynomial factors
\(h_t^2\) do not alter the rate.

Under GRH, the square-root Chebyshev estimate followed by partial summation
against the Gaussian and soft window gives
\[
 |A_{\eta,t,r}(y;s)|
 \ll t^{-M_r}(1+|y|)^{M_r}
 \tag{V43.19}
\]
for fixed \(r\).  Gaussian moments at scale \(h_t=t^{-2}\) are polynomial
in \(t^{-1}\), proving \textup{(V43.17)} after increasing \(M\).
Finally \textup{(V43.18)} is Theorem~\ref{thm:v42-gaussian-burst}.
\end{proof}

The principal saddle is therefore always present in character space, but it
does not force a nonprincipal mode to inherit its type.  Under GRH the
target mode is polynomial even though the principal mode has rate \(1/2\).

\subsection{The target mode is exactly the V29 progression error}

Define the empirical common component
\[
 \overline F_{t,r}(y;s)
 :=\frac1d\sum_{a\in G_q}F_{t,r,a}(y;s)
 \tag{V43.20}
\]
and the residue errors
\[
 E^-_{t,r,a}(y;s)
 :=F_{t,r,a}(y;s)-\overline F_{t,r}(y;s).
 \tag{V43.21}
\]
For every nonprincipal \(\eta\),
\[
 \boxed{
 A_{\eta,t,r}(y;s)
 =\sum_{a\in G_q}\eta(a)E^-_{t,r,a}(y;s).}
 \tag{V43.22}
\]
The empirical mean may be replaced by any common main reader, including
the continuous progression main term of V29, because
\(\sum_a\eta(a)=0\).

\begin{theorem}[Mode-square loop closure]
\label{thm:v43-loop-closure}
For the target nonprincipal character,
\[
 \boxed{
 \mathcal R^{\rm G}_{\chi,t,r}(x;s)
 =\int_{\mathbb R}g_{h_t}(y-x)
 \left|
 \sum_{a\in G_q}\chi(a)E^-_{t,r,a}(y;s)
 \right|^2dy.}
 \tag{V43.23}
\]
Thus the V42 cross-residue cancellation problem is exactly the Gaussian
square of the one-prime progression-error projection
\textup{(V29.5)}.  It is not an independent bilinear cancellation
mechanism.

At the unsmoothed counting level the same projection is
\[
 \sum_{a\in G_q}\chi(a)\Delta_a(Y)=\Psi_\chi(Y),
 \tag{V43.24}
\]
whose square-root bound is GRH-complete by
Theorem~\ref{thm:v29-projected-error-grh}.  Passing to the ratio expansion,
the residue-ratio Fourier basis, or its inverse changes the representation
but introduces no new arithmetic input.
\end{theorem}

\begin{proof}
Equation \textup{(V43.22)} inserted into the definition
\textup{(V43.5)} gives \textup{(V43.23)}.  Equation
\textup{(V43.24)} is \textup{(V29.25)}.  The remaining assertion follows
because Theorem~\ref{thm:v43-mode-square} identifies the entire two-prime
mode with this squared one-prime projection, without an inequality or loss.
\end{proof}

\subsection{The contact remains internal to one mode}

Let \(A^\pm_{\eta,t,0}\) be the soft prefix and tail in the
\(\eta\)-channel.  At a target-channel contact,
\[
 \boxed{
 2\Re\bigl(
 A^-_{\chi,t,0}(x;s)+A^+_{\chi,t,0}(x;s)
 \bigr)=a_{\chi,t}(x).}
 \tag{V43.25}
\]
The slope identity is
\[
 \boxed{
 \partial_xm_{\chi,t}(x)
 =\partial_xa_{\chi,t}(x)
 +2\Im\bigl(
 A^-_{\chi,t,1}(x;s)+A^+_{\chi,t,1}(x;s)
 \bigr).}
 \tag{V43.26}
\]
Neither equation contains
\(A_{\mathbf1,t,r}\) or any other character mode.

\begin{proposition}[No contact-induced mode coupling]
\label{prop:v43-no-contact-coupling}
Finite character orthogonality, positivity of the mode squares, Parseval
\textup{(V43.11)}, and the contact constraint \textup{(V43.25)} imply no
upper bound for
\(\mathcal H^{\rm G}_{\chi,t}\) in terms of the principal mode.
They also imply no bound for the imaginary first target coordinate.
\end{proposition}

\begin{proof}
At a fixed point and order, Proposition~\ref{prop:v43-parseval} permits
arbitrary modal coordinates \(z_\eta\).  Condition
\textup{(V43.25)} fixes one real component of \(z_\chi\) at order zero.
The imaginary and principal coordinates remain independent.  Across reader
orders, Proposition~\ref{prop:v41-phase-model} supplies an exact moment
reader with fixed zeroth contact value and arbitrarily coherent imaginary
first coordinate.  Squared-modulus positivity adds no relation between
these independent nonnegative mode energies.
\end{proof}

The actual prime residue vectors are of course not arbitrary.  The
proposition says precisely that any further relation must come from their
arithmetic construction, not from the finite conductor transform or the
contact hyperplane.

\begin{theorem}[V43 branch-closure verdict]
\label{thm:v43-verdict}
The complete character-mode audit yields:
\begin{enumerate}[label=\textup{(\roman*)},leftmargin=3em]
\item the identity residue-ratio block is the average of positive
      one-prime character-mode energies;
\item the principal mode has exact rate \(1/2\);
\item the target mode is exactly the Gaussian square of V29's projected
      progression error;
\item Parseval and the contact equation do not couple the target derivative
      to the principal or other modes.
\end{enumerate}
Consequently the cross-residue ratio programme supplies no arithmetic
estimate beyond the one-prime character error already known to be
GRH-complete.
\end{theorem}

\begin{proof}
Apply Theorems~\ref{thm:v43-mode-square},
\ref{thm:v43-modal-rates}, and \ref{thm:v43-loop-closure}, followed by
Proposition~\ref{prop:v43-no-contact-coupling}.
\end{proof}

\begin{firewall}[An exact Fourier diagonalization is not a target-mode bound]
V43 completely diagonalizes the V42 residue-ratio vector and proves that
the apparent two-prime target is the square of the V29 one-prime
progression-error projection.  It supplies no new bound for that projection,
does not exclude an off-line zero, and does not prove GRH.
\end{firewall}

\begin{assessment}[V43 acquisition and the multiplicative upstream move]
The ratio-correlation branch is closed to prevent circular work.  Gaussian
localization and residue grouping were useful diagnostics, but complete
character Fourier analysis returns exactly to \(\Psi_\chi\) and its
GRH-complete square-root estimate.  Further character-mode inequalities
would repackage V29 or average away the exceptional target mode.

V44 should move upstream of the residue transform and use the actual
multiplicative identity
\[
 \Lambda=\mu*\log.
 \tag{V43.27}
\]
It should insert this identity into the soft first reader before taking
absolute values, derive the exact two-variable Gaussian phase and saddle
geometry, and separate Type~I from genuinely bilinear ranges.  The required
decision is whether the contact equation constrains either multiplicative
factor, or whether Möbius cancellation at the resulting saddle is itself
equivalent to the original character error.  No further ratio-window
engineering should be attempted unless that multiplicative test produces a
new estimate.
\end{assessment}

\section*{Version V43 append-only record}
\addcontentsline{toc}{section}{Version V43 append-only record}
The complete V42 source was retained before this continuation.  It had
\(659875\) bytes, \(19977\) lines, and SHA--256
\[
 \texttt{2E61AD329D16E161C056236809D1D9CE27D8707A2A5AAA4FA54CD3F4D5520A12}.
\]
V43 diagonalized the near-ratio residue vector in the complete character
basis, identified the principal and target modes, and proved that the
two-prime correlation branch collapses exactly to the V29 one-prime
progression-error projection.

% =============================================================================
% END APPENDED V43 MATERIAL
% =============================================================================
% =============================================================================
% BEGIN APPENDED V44 MATERIAL
% =============================================================================

\section{V44: Möbius convolution and localization to large divisor factors}
\label{sec:grh-v44}

V43 closed the residue-ratio branch because its complete character
diagonalization returned exactly to the V29 one-prime progression error.
This section moves upstream of that projection and inserts the genuinely
multiplicative identity
\[
 \Lambda=\mu*\log
 \tag{V44.1}
\]
into the soft readers before taking absolute values.

The factorized Gaussian has a flat allocation saddle: its exponential
weight depends only on \(\log a+\log b\), so both the \(a=1\) slice and
every fixed balanced factor sector retain absolute rate \(1/4\).  There is
nevertheless one unconditional signed gain.  If the Möbius factor \(a\) is
restricted to any subexponential range, the inner \(b\)-sum is a smooth
nonprincipal character sum and bounded character partial sums reduce it to
zero \(1/t\)-rate.  Hence a forbidden contact slope cannot be carried by a
small Möbius factor; it must move into the complementary large-\(a\)
tail.

\subsection{Exact soft-reader convolution}

For \(\sigma\in\{+,-\}\), write
\[
 W_{t}^{\sigma}(ab;s)
 :=W^\sigma_{t,\omega_t,v_t(s)}(\log(ab)).
 \tag{V44.2}
\]

\begin{theorem}[Exact multiplicative soft-reader identity]
\label{thm:v44-exact-convolution}
For every fixed integer \(r\ge0\), \(t>0\), and real \(x\),
\[
 \boxed{
 \begin{aligned}
 P_{\chi,t}^{\sigma,[r]}(x;s)
 =\sum_{\substack{a,b\ge1\\ab\ge2}}
 &\frac{\mu(a)\chi(a)\chi(b)\log b\,(\log(ab))^r}
       {\sqrt{ab}}\\
 &{}\times e^{-t\log^2(ab)}e^{ix\log(ab)}
 W_t^\sigma(ab;s).
 \end{aligned}}
 \tag{V44.3}
\]
The double series converges absolutely for each \(t>0\).
\end{theorem}

\begin{proof}
The divisor identity
\[
 \sum_{ab=n}\mu(a)\log b=\Lambda(n)
 \tag{V44.4}
\]
and complete multiplicativity of the Dirichlet character give
\[
 \sum_{ab=n}\mu(a)\chi(a)\chi(b)\log b
 =\chi(n)\Lambda(n).
 \tag{V44.5}
\]
Every remaining factor in \textup{(V44.3)} depends only on \(ab=n\), so
grouping by the product proves the identity.  Absolute convergence follows
from Gaussian decay and the elementary divisor bound
\[
 \sum_{ab\le X}\frac{|\mu(a)|\log b}{\sqrt{ab}}
 \ll X^{1/2}(\log X)^2.
 \tag{V44.6}
\]
\end{proof}

The identity is fiberwise: for a fixed product \(n\), every analytic phase
and window factor is common to all divisors, and the divisor sum
\textup{(V44.4)} is the only cancellation inside that fiber.

\subsection{The flat factor-allocation saddle}

Define the absolute convolved mass
\[
 \begin{aligned}
 \mathcal A_{t,r}^{\sigma}(s)
 :=\sum_{\substack{a,b\ge1\\ab\ge2}}
 &\frac{|\mu(a)\chi(a)\chi(b)|\log b\,(\log(ab))^r}
       {\sqrt{ab}}\\
 &{}\times e^{-t\log^2(ab)}W_t^\sigma(ab;s).
 \end{aligned}
 \tag{V44.7}
\]
For a closed interval \(J\subset(0,1)\) of positive length, let
\(\mathcal A_{t,r}^{-,J}\) denote the same sum restricted by
\[
 \frac{\log a}{\log(ab)}\in J.
 \tag{V44.8}
\]
Let \(\mathcal A_{t,r}^{-,0}\) be the \(a=1\) slice.

\begin{theorem}[Every factor-allocation sector has full absolute type]
\label{thm:v44-flat-saddle}
For every fixed \(r,s\),
\[
 \boxed{
 \lim_{t\downarrow0}4t\log\mathcal A_{t,r}^{-}(s)
 =\frac14.}
 \tag{V44.9}
\]
The endpoint Type-I slice and every fixed balanced sector have the same
rate:
\[
 \boxed{
 \lim_{t\downarrow0}4t\log\mathcal A_{t,r}^{-,0}(s)
 =\frac14,}
 \tag{V44.10}
\]
\[
 \boxed{
 \lim_{t\downarrow0}4t\log\mathcal A_{t,r}^{-,J}(s)
 =\frac14.}
 \tag{V44.11}
\]
\end{theorem}

\begin{proof}
Put \(w=\log(ab)\).  Splitting products into
\(k\le w<k+1\), estimate
\[
 \sum_{\substack{a,b\ge1\\ab<e^{k+1}}}
 \frac{|\mu(a)|\log b\,(1+\log(ab))^r}{\sqrt{ab}}
 \ll_r e^{k/2}(1+k)^{r+3}.
 \tag{V44.12}
\]
Indeed, sum first over \(b\le e^{k+1}/a\), use
\(\sum_{b\le Y}b^{-1/2}\log b\ll Y^{1/2}\log(2Y)\), and then
\(\sum_{a\le e^{k+1}}a^{-1}\ll k+1\).  Therefore
\[
 \mathcal A_{t,r}^{-}(s)
 \ll_r\sum_{k\ge0}(1+k)^{r+3}e^{k/2-tk^2}.
 \tag{V44.13}
\]
The exponent \(k/2-tk^2\) has its maximum \(1/(16t)\) at
\[
 k=\frac1{4t}.
 \tag{V44.14}
\]
This proves the upper rate \(1/4\).

For the \(a=1\) lower bound, restrict \(b\) to one fixed invertible residue
class and to a logarithmic interval of fixed width about
\(1/(4t)\).  The PNT in that progression, or merely integer density for
the absolute character support, gives exponential count \(e^{1/(4t)}\);
the factor \(b^{-1/2}e^{-t\log^2b}\) contributes
\(e^{-3/(16t)+o(1/t)}\).  The mesoscopic prefix window tends to \(1/2\)
there.  The resulting rate is \(1-3/4=1/4\), proving
\textup{(V44.10)} and the matching lower bound in \textup{(V44.9)}.

For \(J\subset(0,1)\), choose an interior point \(\theta\in J\), restrict
\(\log a\) and \(\log b\) to fixed-width intervals about
\(\theta/(4t)\) and \((1-\theta)/(4t)\), and take \(a,b\) prime in fixed
invertible residue classes.  Prime number theory in fixed progressions
supplies the two exponential counts; all prime-density and logarithmic
factors are polynomial.  Their combined exponent is again
\[
 \frac{w}{2}-tw^2
 \quad\hbox{at}\quad w=\frac1{4t},
 \tag{V44.15}
\]
independently of \(\theta\).  This proves \textup{(V44.11)}.
\end{proof}

Thus ordinary Type-I/Type-II allocation does not lower the absolute
exponent.  Any gain must use a signed factor before the two-dimensional
mass is estimated.

\subsection{Small Möbius factors are unconditionally harmless}

For \(A\ge1\), define the truncated Type-I reader
\[
 \begin{aligned}
 \mathcal I_{A,t}^{\sigma,[r]}(x;s)
 :=\sum_{\substack{1\le a\le A\\b\ge1\\ab\ge2}}
 &\frac{\mu(a)\chi(a)\chi(b)\log b\,(\log(ab))^r}
       {\sqrt{ab}}\\
 &{}\times e^{-t\log^2(ab)}e^{ix\log(ab)}
 W_t^\sigma(ab;s),
 \end{aligned}
 \tag{V44.16}
\]
and let
\[
 \mathcal J_{A,t}^{\sigma,[r]}(x;s)
 :=P_{\chi,t}^{\sigma,[r]}(x;s)
 -\mathcal I_{A,t}^{\sigma,[r]}(x;s).
 \tag{V44.17}
\]

\begin{lemma}[Bounded character partial sums]
\label{lem:v44-character-partial}
Because \(\chi\) is nonprincipal and periodic modulo \(q\),
\[
 \boxed{
 \sup_{Y\ge1}\left|\sum_{b\le Y}\chi(b)\right|
 \le q.}
 \tag{V44.18}
\]
\end{lemma}

\begin{proof}
Every complete period has sum zero.  Only the final incomplete period
remains.
\end{proof}

\begin{theorem}[Subexponential Type-I cancellation]
\label{thm:v44-type-I}
Let \(A_t\ge1\) satisfy
\[
 \log A_t=o(1/t).
 \tag{V44.19}
\]
For every fixed \(r,s\) and compact interval \(I\),
\[
 \boxed{
 \limsup_{t\downarrow0}4t\log\left(
 1+\sup_{x\in I}
 |\mathcal I_{A_t,t}^{-,[r]}(x;s)|
 \right)=0.}
 \tag{V44.20}
\]
The same conclusion holds for the soft tail.
\end{theorem}

\begin{proof}
Fix \(a\) and put \(u=\log a\).  Apart from the constant factor
\(\mu(a)\chi(a)a^{-1/2}e^{ixu}\), the inner sum in \(b\) is
\[
 \sum_{b\ge1}\chi(b)F_{a,t,x,r}(\log b),
 \tag{V44.21}
\]
where
\[
 F_{a,t,x,r}(v)
 :=v(u+v)^r e^{-v/2-t(u+v)^2}
 e^{ixv}W^-_{t,\omega_t,v_t(s)}(u+v).
 \tag{V44.22}
\]
Stieltjes summation with Lemma~\ref{lem:v44-character-partial} gives
\[
 \left|\sum_{b\ge1}\chi(b)F_{a,t,x,r}(\log b)\right|
 \le q\int_0^\infty
 |\partial_vF_{a,t,x,r}(v)|\,dv.
 \tag{V44.23}
\]
The boundary terms vanish because of the factor \(v\) and the Gaussian
tail.  On a fixed \(x\)-interval, direct differentiation yields
\[
 \int_0^\infty|\partial_vF_{a,t,x,r}(v)|\,dv
 \ll_{I,r,s}t^{-M_r}(1+u)^{M_r}e^{-tu^2}.
 \tag{V44.24}
\]
Here \(0\le W^-\le1\), and its derivative is
\[
 O\!\left(a_{t,\omega_t}^{-1}\right)
 =O(\sqrt t/\omega_t),
 \tag{V44.25}
\]
so it creates no exponential cost.

Summing \textup{(V44.24)} over \(a\le A_t\) gives
\[
 \sup_{x\in I}|\mathcal I_{A_t,t}^{-,[r]}(x;s)|
 \ll t^{-M}
 \sum_{a\le A_t}a^{-1/2}(1+\log a)^M
 \ll t^{-M}A_t^{1/2}(1+\log A_t)^M.
 \tag{V44.26}
\]
Condition \textup{(V44.19)} proves \textup{(V44.20)}.  The proof for
\(W^+\) is identical.
\end{proof}

This cancellation is much stronger than applying the PNT to the original
prime reader.  It uses only periodicity of the bare character after the
Möbius factor has been frozen.

\subsection{A forbidden slope moves to the large Möbius tail}

\begin{theorem}[Large-factor localization of a forbidden contact]
\label{thm:v44-large-factor}
Assume the paired channel has an off-line zero.  Let
\(p_*=(\gamma_*,\delta_*,m_*)\), \(t_k\downarrow0\), and
\(x_{k,\pm}\) be the centered exposed pair and phase-tuned contacts from
V37--V39.  For every cutoff \(A_t\) satisfying
\textup{(V44.19)}, either sign of contact, and every fixed \(s\),
\[
 \boxed{
 \lim_{k\to\infty}4t_k\log
 \left|
 \Im\mathcal J_{A_{t_k},t_k}^{-,[1]}
 (x_{k,\pm};s)
 \right|
 =\delta_*^2.}
 \tag{V44.27}
\]
Thus the complete forbidden soft-prefix slope is carried by convolution
terms with
\[
 a>A_t=e^{o(1/t)}.
 \tag{V44.28}
\]
\end{theorem}

\begin{proof}
Theorem~\ref{thm:v39-prefix-carries-slope} gives rate
\(\delta_*^2\) for
\(\Im P_{\chi,t_k}^{-,[1]}(x_{k,\pm};s)\).
Theorem~\ref{thm:v44-type-I} gives zero rate for the subtracted Type-I
piece on a fixed compact interval containing the contacts.  Since
\(\delta_*^2>0\), the Type-I piece is relatively exponentially negligible.
Equation \textup{(V44.17)} then proves \textup{(V44.27)}.
\end{proof}

This is a genuine factor-space localization.  It does not yet distinguish
between a large Möbius factor paired with a small \(b\) and a genuinely
balanced bilinear pair.

\subsection{The contact equation does not separate convolution factors}

Apply \textup{(V44.16)}--\textup{(V44.17)} to both soft pieces and to
reader orders zero and one.  At every zero-level contact,
\[
 \boxed{
 2\Re\sum_{\sigma\in\{+,-\}}
 \left(
 \mathcal I_{A,t}^{\sigma,[0]}(x;s)
 +\mathcal J_{A,t}^{\sigma,[0]}(x;s)
 \right)
 =a_{\chi,t}(x).}
 \tag{V44.29}
\]
The slope ledger is
\[
 \boxed{
 \partial_xm_{\chi,t}(x)
 =\partial_xa_{\chi,t}(x)
 +2\Im\sum_{\sigma\in\{+,-\}}
 \left(
 \mathcal I_{A,t}^{\sigma,[1]}(x;s)
 +\mathcal J_{A,t}^{\sigma,[1]}(x;s)
 \right).}
 \tag{V44.30}
\]
These equations constrain only the sums of all divisor ranges.  They give
no independent value for a fixed \(a\)-range or for either multiplicative
factor.  The localization in Theorem~\ref{thm:v44-large-factor} comes from
the periodic Type-I estimate, not from an unproved separation of the
contact equation.

\begin{theorem}[V44 multiplicative verdict]
\label{thm:v44-verdict}
The Möbius convolution has the following exact profile.
\begin{enumerate}[label=\textup{(\roman*)},leftmargin=3em]
\item Its full and factor-balanced absolute masses all have rate \(1/4\).
\item Every subexponential small-\(a\) Type-I range has zero rate by bounded
      nonprincipal character partial sums.
\item Under failure of GRH, the large-\(a\) remainder alone retains the
      complete forbidden contact-slope rate \(\delta_*^2\).
\item The contact equation supplies no separate constraint on the two
      multiplicative factors.
\end{enumerate}
\end{theorem}

\begin{proof}
Combine Theorems~\ref{thm:v44-flat-saddle},
\ref{thm:v44-type-I}, and \ref{thm:v44-large-factor} with the exact ledger
\textup{(V44.29)}--\textup{(V44.30)}.
\end{proof}

\begin{firewall}[Large-factor localization is not Möbius cancellation]
V44 proves an exact multiplicative decomposition and removes every
subexponential small-Möbius-factor range from the forbidden contact slope.
It does not bound the complementary large-\(a\) tail, separate its
small-\(b\) and balanced components, exclude an off-line zero, or prove
GRH.
\end{firewall}

\begin{assessment}[V44 acquisition and the Type-I/II fork]
Unlike the finite residue transform, multiplicative convolution produces a
strict new localization: a forbidden slope cannot arise with
\(a\le e^{o(1/t)}\).  The remaining large-\(a\) tail has two qualitatively
different sectors.  If \(b\) is small, it is a smooth
Möbius-character sum in \(a\); if both factors are large, it is a genuine
bilinear sum.  The flat saddle theorem shows that absolute estimates for
either sector can still retain full type.

V45 should introduce a second subexponential cutoff \(B_t\), prove the exact
three-way decomposition
\[
 \{\text{small }a\}\ \cup\
 \{\text{large }a,\text{ small }b\}\ \cup\
 \{\text{large }a,\text{ large }b\},
 \tag{V44.31}
\]
and determine which of the last two sectors must carry the V44 forbidden
slope.  The small-\(b\) sector should be expressed in terms of
\[
 M_\chi(X):=\sum_{n\le X}\mu(n)\chi(n),
 \tag{V44.32}
\]
with its exact zero-free/GRH content stated before any estimate.  Only if
that sector can be removed without assuming GRH should the balanced Type-II
branch be pursued.
\end{assessment}

\section*{Version V44 append-only record}
\addcontentsline{toc}{section}{Version V44 append-only record}
The complete V43 source was retained before this continuation.  It had
\(673557\) bytes, \(20398\) lines, and SHA--256
\[
 \texttt{D9E09B043E27EAE4607BED952D086339D6E4C4A300EA5E571DB71E81A4755833}.
\]
V44 inserted \(\Lambda=\mu*\log\) into the soft readers, proved the flat
factor-allocation saddle, established zero-rate cancellation for every
subexponential small-Möbius-factor range, and localized every forbidden
contact slope to the complementary large-factor tail.

% =============================================================================
% END APPENDED V44 MATERIAL
% =============================================================================
% =============================================================================
% BEGIN APPENDED V45 MATERIAL
% =============================================================================

\section{V45: the small logarithmic factor is already a reciprocal-
\texorpdfstring{$L$}{L} problem}
\label{sec:grh-v45}

V44 localized every forbidden contact slope to terms with a large
Möbius factor.  The proposed next split is exact, but it does not force a
balanced Type-II problem.  Freezing the logarithmic factor at the first
integer visible to the character leaves the reciprocal Dirichlet series
\(1/L(z,\chi)\).  Thus the large-\(a\), small-\(b\) sector already contains
every zero of \(L(z,\chi)\), without any possible cancellation among
different \(b\)'s.

This section makes that obstruction quantitative.  A power exponent
\(\theta\) for the twisted Mertens sum transfers to the Gaussian reader
with normalized heat rate at most
\[
 \left(\theta-\frac12\right)^2.
 \tag{V45.1}
\]
Consequently the rate becomes zero precisely at the GRH-strength
Mertens exponent \(1/2\).  The three-way decomposition therefore identifies
an exact fork, but no unconditional passage from the small-factor branch to
the balanced branch.

\subsection{The exact three-way factor split}

For \(A,B\ge1\), retain \(\mathcal I_{A,t}^{\sigma,[r]}\) from
\textup{(V44.16)} and define
\[
 \begin{aligned}
 \mathcal K_{A,B,t}^{\sigma,[r]}(x;s)
 :=\sum_{\substack{a>A,\ 1\le b\le B\\ab\ge2}}
 &\frac{\mu(a)\chi(a)\chi(b)\log b\,(\log(ab))^r}
        {\sqrt{ab}}\\
 &{}\times e^{-t\log^2(ab)}e^{ix\log(ab)}W_t^\sigma(ab;s),
 \end{aligned}
 \tag{V45.2}
\]
and
\[
 \begin{aligned}
 \mathcal T_{A,B,t}^{\sigma,[r]}(x;s)
 :=\sum_{\substack{a>A,\ b>B}}
 &\frac{\mu(a)\chi(a)\chi(b)\log b\,(\log(ab))^r}
        {\sqrt{ab}}\\
 &{}\times e^{-t\log^2(ab)}e^{ix\log(ab)}W_t^\sigma(ab;s).
 \end{aligned}
 \tag{V45.3}
\]
The letter \(\mathcal K\) denotes the large-Möbius/small-logarithmic
factor sector, while \(\mathcal T\) denotes the sector in which both
factors exceed their cutoffs.

\begin{theorem}[Exact Type-I/small-factor/Type-II decomposition]
\label{thm:v45-three-way}
For all \(A,B,t,x,r,s\) and \(\sigma\in\{+,-\}\),
\[
 \boxed{
 P_{\chi,t}^{\sigma,[r]}(x;s)
 =\mathcal I_{A,t}^{\sigma,[r]}(x;s)
  +\mathcal K_{A,B,t}^{\sigma,[r]}(x;s)
  +\mathcal T_{A,B,t}^{\sigma,[r]}(x;s).}
 \tag{V45.4}
\]
If \(\log A_t=o(1/t)\), the first term has zero normalized rate by
Theorem~\ref{thm:v44-type-I}.
\end{theorem}

\begin{proof}
The three mutually disjoint ranges are
\[
 \{a\le A\},\qquad
 \{a>A,\ b\le B\},\qquad
 \{a>A,\ b>B\}.
 \tag{V45.5}
\]
They partition every pair in the absolutely convergent series
\textup{(V44.3)}.  The last assertion is exactly
Theorem~\ref{thm:v44-type-I}.
\end{proof}

At the phase-tuned contacts of Theorem~\ref{thm:v44-large-factor}, put
\[
 \kappa_k:=\Im\mathcal K_{A_{t_k},B_{t_k},t_k}^{-,[1]}
 (x_{k,\pm};s),
 \qquad
 \tau_k:=\Im\mathcal T_{A_{t_k},B_{t_k},t_k}^{-,[1]}
 (x_{k,\pm};s).
 \tag{V45.6}
\]

\begin{corollary}[Exact routing dichotomy]
\label{cor:v45-routing}
Under the off-line-zero hypothesis of
Theorem~\ref{thm:v44-large-factor}, for every positive cutoff \(B_t\),
\[
 \boxed{
 \max\left\{
 \limsup_{k\to\infty}4t_k\log(1+|\kappa_k|),
 \limsup_{k\to\infty}4t_k\log(1+|\tau_k|)
 \right\}\ge\delta_*^2.}
 \tag{V45.7}
\]
If either displayed limsup is strictly below \(\delta_*^2\), the other
term has the exact limit \(\delta_*^2\) after the harmless \(1+|\cdot|\)
regularization.
\end{corollary}

\begin{proof}
Equations \textup{(V44.17)} and \textup{(V45.4)} give
\(Im\mathcal J=\kappa_k+\tau_k\).  Theorem
\ref{thm:v44-large-factor} and the triangle inequality prove
\textup{(V45.7)}.  If one summand has strictly smaller exponential rate,
subtracting it from a quantity of rate \(\delta_*^2\) preserves the rate
of the other.
\end{proof}

The corollary is only a dichotomy.  To route the forbidden slope into
\(\mathcal T\), one would first need a rate strictly below
\(\delta_*^2\) for \(\mathcal K\).

\subsection{The exact arithmetic object in the small-factor branch}

Define the twisted Mertens sum
\[
 M_\chi(X):=\sum_{n\le X}\mu(n)\chi(n)
 \tag{V45.8}
\]
and, for a finite cutoff \(B\), the Dirichlet polynomial
\[
 Q_{\chi,B}(z):=
 \sum_{2\le b\le B}\frac{\chi(b)\log b}{b^z}.
 \tag{V45.9}
\]
Let
\[
 c_{\chi,B}(n)
 :=\sum_{\substack{ab=n\\b\le B}}
 \mu(a)\chi(a)\chi(b)\log b.
 \tag{V45.10}
\]

\begin{theorem}[Finite small-factor series]
\label{thm:v45-finite-series}
For \(\Re z>1\),
\[
 \boxed{
 \sum_{n\ge1}\frac{c_{\chi,B}(n)}{n^z}
 =\frac{Q_{\chi,B}(z)}{L(z,\chi)}.}
 \tag{V45.11}
\]
For a fixed cutoff \(A\), imposing \(a>A\) changes the right side only by
an entire Dirichlet polynomial.  Hence its poles at zeros of \(L\) remain
unless the numerator \(Q_{\chi,B}\) vanishes there to sufficient order.
\end{theorem}

\begin{proof}
Absolute convergence in \(\Re z>1\) permits factorization:
\[
 \begin{aligned}
 \sum_{a\ge1}\sum_{b\le B}
 \frac{\mu(a)\chi(a)\chi(b)\log b}{(ab)^z}
 &=\left(\sum_{a\ge1}\frac{\mu(a)\chi(a)}{a^z}\right)
   \left(\sum_{2\le b\le B}
   \frac{\chi(b)\log b}{b^z}\right)\\
 &=\frac{Q_{\chi,B}(z)}{L(z,\chi)}.
 \end{aligned}
 \tag{V45.12}
\]
Deleting the finitely many values \(a\le A\) subtracts
\[
 Q_{\chi,B}(z)
 \sum_{a\le A}\frac{\mu(a)\chi(a)}{a^z},
 \tag{V45.13}
\]
which is entire.
\end{proof}

There is a cutoff for which numerator cancellation is impossible.  Let
\[
 b_0:=\min\{b\ge2:(b,q)=1\}.
 \tag{V45.14}
\]
By minimality, \(\chi(b)=0\) for \(2\le b<b_0\), while
\(\chi(b_0)\ne0\).

\begin{theorem}[The first visible logarithmic factor sees every zero]
\label{thm:v45-first-visible}
For \(B=b_0\),
\[
 \boxed{
 Q_{\chi,b_0}(z)
 =\chi(b_0)\log b_0\,b_0^{-z},}
 \tag{V45.15}
\]
which is nowhere zero.  Consequently
\[
 \boxed{
 \sum_{n\ge1}\frac{c_{\chi,b_0}(n)}{n^z}
 =\frac{\chi(b_0)\log b_0\,b_0^{-z}}{L(z,\chi)}}
 \tag{V45.16}
\]
has a pole, with the same order, at every zero of \(L(z,\chi)\).  This
remains true after deleting any fixed finite range \(a\le A\).
\end{theorem}

\begin{proof}
All terms preceding \(b_0\) in \textup{(V45.9)} vanish, so
\textup{(V45.15)} is immediate.  Its exponential factor never vanishes.
The assertions now follow from Theorem~\ref{thm:v45-finite-series}.
\end{proof}

Thus even the fixed, hence subexponential, choice \(B_t=b_0\) makes
\(\mathcal K\) a Gaussian smoothing of a coefficient sequence whose
meromorphic transform contains every reciprocal-\(L\) pole.  The
small-factor branch is not an elementary periodic character sum like the
small-\(a\) branch in V44.

\subsection{The Mertens exponent and its heat-rate transfer}

For clarity, call
\[
 \mathsf{RZ}_\chi:qquad
 L(z,\chi)\ne0\quad(\Re z>1/2)
 \tag{V45.17}
\]
the right-half zero-free statement for \(\chi\).  For a complex
character, \(\mathsf{RZ}_\chi\) alone is one-sided; by the functional
equation, \(\mathsf{RZ}_\chi\) and
\(\mathsf{RZ}_{\overline\chi}\) together are the paired GRH statement.

\begin{theorem}[Twisted Littlewood criterion]
\label{thm:v45-littlewood}
For a fixed nonprincipal Dirichlet character \(\chi\), the following are
equivalent:
\begin{enumerate}[label=\textup{(\roman*)},leftmargin=3em]
\item \(\mathsf{RZ}_\chi\) holds.
\item For every \(\varepsilon>0\),
\[
 M_\chi(X)=O_{\chi,\varepsilon}(X^{1/2+\varepsilon}).
 \tag{V45.18}
\]
\end{enumerate}
Moreover, \textup{(V45.18)} implies, uniformly for \(x\) in every fixed
compact interval \(I\),
\[
 \sum_{n\le X}\mu(n)\chi(n)n^{ix}
 =O_{\chi,\varepsilon,I}(X^{1/2+\varepsilon}).
 \tag{V45.19}
\]
The corresponding two estimates for \(\chi\) and \(\overline\chi\) are
equivalent to paired GRH.
\end{theorem}

\begin{proof}
If \textup{(V45.18)} holds, Abel summation makes
\[
 \sum_{n\ge1}\frac{\mu(n)\chi(n)}{n^z}
 =z\int_1^\infty M_\chi(u)u^{-z-1}\,du
 \tag{V45.20}
\]
analytic in every half-plane \(\Re z>1/2+\varepsilon\).  It agrees with
\(1/L(z,\chi)\) for \(\Re z>1\), so analytic continuation forbids a zero
of \(L\) to the right of \(1/2\).

Conversely, under \(\mathsf{RZ}_\chi\), the reciprocal is analytic in
\(\Re z>1/2\).  The standard Littlewood contour argument for a fixed
Dirichlet \(L\)-function applies Perron's formula on
\(\Re z=1/2+\varepsilon\); the functional equation and finite order give
the required polynomial vertical growth.  Truncation and partial
summation yield \textup{(V45.18)}.  No conductor-uniform assertion is
being made.  Finally,
\[
 \sum_{n\le X}\mu(n)\chi(n)n^{ix}
 =X^{ix}M_\chi(X)
  -ix\int_1^X M_\chi(u)u^{ix-1}\,du,
 \tag{V45.21}
\]
which proves \textup{(V45.19)}.  The paired assertion follows from the
functional equation.
\end{proof}

The next theorem records exactly how any available power exponent passes
through the Gaussian saddle.

\begin{theorem}[Quadratic Mertens-to-heat transfer]
\label{thm:v45-rate-transfer}
Assume that some \(\theta\in[1/2,1]\) satisfies, for every
\(\varepsilon>0\),
\[
 M_\chi(X)=O_{\chi,\varepsilon}(X^{\theta+\varepsilon}).
 \tag{V45.22}
\]
Let \(B_t\ge1\) obey \(\log B_t=o(1/t)\); the cutoff \(A_t\) may be
arbitrary.  For every fixed \(r,s\) and compact interval \(I\),
\[
 \boxed{
 \limsup_{t\downarrow0}4t\log\left(
 1+\sup_{x\in I}
 |\mathcal K_{A_t,B_t,t}^{\sigma,[r]}(x;s)|
 \right)
 \le\left(\theta-\frac12\right)^2}
 \tag{V45.23}
\]
for both \(\sigma=+\) and \(\sigma=-\).
In particular, \(\mathsf{RZ}_\chi\) implies zero rate for every such
small-factor sector.
\end{theorem}

\begin{proof}
Fix \(b\le B_t\), write \(v=\log b\), and sum first in \(a>A_t\).
Stieltjes summation against \(M_\chi(u)\), including the lower endpoint,
bounds the inner sum by a polynomial in
\(t^{-1},1+v,1+|x|\) times
\[
 \sup_{y\ge0}
 \exp\left\{
 (\theta+\varepsilon-\tfrac12)y-t(y+v)^2
 \right\}.
 \tag{V45.24}
\]
The derivatives of either soft window have only subexponential cost, as
in \textup{(V44.25)}.  Since \(v\ge0\), the exponent in
\textup{(V45.24)} is at most
\[
 \frac{(\theta+\varepsilon-1/2)^2}{4t}.
 \tag{V45.25}
\]
The remaining sum
\(\sum_{b\le B_t}b^{-1/2}\log b\) is
\(O(B_t^{1/2}\log(2B_t))\), whose normalized logarithm tends to zero.
Thus the left side of \textup{(V45.23)} is at most
\((\theta+\varepsilon-1/2)^2\).  Let
\(\varepsilon\downarrow0\).  The last assertion follows from
Theorem~\ref{thm:v45-littlewood}.
\end{proof}

The trivial estimate \(|M_\chi(X)|\le X\) corresponds to
\(\theta=1\) and returns the full upper rate \(1/4\).  Any estimate of
the form
\[
 |M_\chi(X)|\le X\exp\{-o(\log X)\}
 \tag{V45.26}
\]
still has power exponent \(\theta=1\) and does not change this normalized
\(1/t\)-rate.  A fixed power saving lowers the rate, while zero rate by
this transfer requires every exponent \(1/2+\varepsilon\), exactly the
criterion in Theorem~\ref{thm:v45-littlewood}.

\subsection{The small-factor branch cannot be discarded upstream of GRH}

\begin{theorem}[V45 reciprocal-factor verdict]
\label{thm:v45-verdict}
The large-\(a\) remainder of V44 has the following exact structure.
\begin{enumerate}[label=\textup{(\roman*)},leftmargin=3em]
\item It splits exactly into the small-\(b\) sector \(\mathcal K\) and
      the two-large-factor sector \(\mathcal T\).
\item At a forbidden contact, at least one of these sectors has limsup
      normalized rate at least \(\delta_*^2\).
\item For the fixed subexponential cutoff \(B=b_0\), the whole
      small-\(b\) coefficient series is a nowhere-zero entire multiple of
      \(1/L(z,\chi)\), so it retains every zero as a pole.
\item A Mertens power exponent \(\theta\) gives heat rate at most
      \((\theta-1/2)^2\), and the zero-rate endpoint is the one-sided
      GRH criterion itself.
\end{enumerate}
Consequently no unconditional result obtained here forces the forbidden
slope into the balanced Type-II sector.
\end{theorem}

\begin{proof}
Combine Theorem~\ref{thm:v45-three-way},
Corollary~\ref{cor:v45-routing}, Theorem~\ref{thm:v45-first-visible}, and
Theorems~\ref{thm:v45-littlewood}--\ref{thm:v45-rate-transfer}.
\end{proof}

\begin{firewall}[A decomposition is not a Type-II reduction]
V45 does not prove a new bound for \(M_\chi\), show that
\(\mathcal K\) is negligible at a forbidden contact, force
\(\mathcal T\) to carry the slope, exclude an off-line zero, or prove
GRH.  Theorem~\ref{thm:v45-rate-transfer} is a conditional transfer
principle.  Feeding it the exponent \(1/2+\varepsilon\) is precisely
feeding it the right-half GRH statement one is trying to prove.
\end{firewall}

\begin{assessment}[V45 closes the naive balanced-factor fork]
The first logarithmic factor visible to \(\chi\) already produces
\(Q_{\chi,b_0}/L\) with a nowhere-vanishing numerator.  Thus increasing
the second cutoff and attacking a balanced bilinear sum is not compelled:
the complementary fixed small-factor slice contains the complete zero
obstruction.

The next upstream object is not a larger Type-II box but the cancellation
of reciprocal principal parts as the \(b\)-sum is completed.  In the
half-plane of absolute convergence,
\[
 \lim_{B\to\infty}Q_{\chi,B}(z)=-L'(z,\chi),
 \qquad
 \frac{Q_{\chi,B}(z)}{L(z,\chi)}longrightarrow
 -\frac{L'}{L}(z,\chi).
 \tag{V45.27}
\]
For a zero of multiplicity \(m\), \(1/L\) has a pole of order \(m\),
whereas \(-L'/L\) has only a simple pole.  V46 should track, with an
exact finite-\(B\) remainder, how the growing logarithmic-factor sum
cancels all higher reciprocal principal parts and changes the surviving
residue.  That pole-order descent is a structural constraint absent from
absolute Type-I/II estimates and is the appropriate next upstream test.
\end{assessment}

\section*{Version V45 append-only record}
\addcontentsline{toc}{section}{Version V45 append-only record}
The complete V44 source was retained before this continuation.  It had
\(686956\) bytes, \(20828\) lines, and SHA--256
\[
 \texttt{F733573FB45F6149D335BEEF27241721509CF7A7448DC2A0369F3884236ACA71}.
\]
V45 made the proposed three-way factor split exact, identified the first
visible small logarithmic factor as a nowhere-zero multiple of
\(1/L(z,\chi)\), proved the quadratic Mertens-to-heat rate transfer, and
showed that discarding the small-factor branch at zero rate would import
the one-sided GRH criterion itself.

% =============================================================================
% END APPENDED V45 MATERIAL
% =============================================================================
% =============================================================================
% BEGIN APPENDED V46 MATERIAL
% =============================================================================

\section{V46: finite-factor completion and reciprocal pole-order descent}
\label{sec:grh-v46}

V45 showed that the first logarithmic factor visible to \(\chi\) produces
a nowhere-zero multiple of \(1/L(z,\chi)\).  At a zero of multiplicity
\(m\), that quotient has a pole of order \(m\), whereas the completed
logarithmic derivative \(-L'/L\) has residue \(-m\) at a simple pole.
This section computes exactly how the finite logarithmic-factor sum
descends from the former local structure to the latter.

There are two conclusions.  First, bounded character partial sums give
local uniform convergence of every fixed derivative of
\(Q_{\chi,B}\) to the corresponding derivative of \(-L'\) throughout
\(\Re z>0\).  This forces every higher principal-part coefficient to
tend to zero and the simple coefficient to tend to \(-m\), with an
explicit \(B^{-\Re\rho}\) bound.  Second, a pole of any fixed order gives
a derivative of the same complex heat kernel.  Its order changes only a
polynomial factor in \(t^{-1}\), so all nonzero pole orders at the same
off-line zero have the same normalized heat rate.  Pole-order descent is
therefore exact but cannot itself supply a GRH contradiction.

\begin{remark}[Correction to the displayed arrow in V45]
In \textup{(V45.27)}, the control slash before the printed word
``longrightarrow'' was inadvertently omitted.  The intended statement,
proved quantitatively below, is
\[
 Q_{\chi,B}(z)\longrightarrow-L'(z,\chi),
 \qquad
 \frac{Q_{\chi,B}(z)}{L(z,\chi)}\longrightarrow
 -\frac{L'}{L}(z,\chi).
 \tag{V46.1}
\]
\end{remark}

\subsection{Uniform completion of the logarithmic factor}

Recall
\[
 Q_{\chi,B}(z)=
 \sum_{2\le b\le B}\chi(b)\log b\,b^{-z}.
 \tag{V46.2}
\]
For a compact set \(K\Subset\{\Re z>0\}\), put
\[
 \sigma_K:=\min_{z\in K}\Re z>0.
 \tag{V46.3}
\]

\begin{theorem}[Quantitative local completion]
\label{thm:v46-local-completion}
For every nonnegative integer \(J\),
\[
 \boxed{
 \max_{0\le j\le J}\sup_{z\in K}
 \left|
 Q_{\chi,B}^{(j)}(z)-(-L'(z,\chi))^{(j)}
 \right|
 \ll_{K,J,\chi}
 B^{-\sigma_K}(\log(2B))^{J+1}.}
 \tag{V46.4}
\]
In particular, \(Q_{\chi,B}\to-L'\) locally uniformly with all fixed
derivatives in \(\Re z>0\).
\end{theorem}

\begin{proof}
Let
\[
 A_\chi(X):=\sum_{n\le X}\chi(n).
 \tag{V46.5}
\]
Lemma~\ref{lem:v44-character-partial} gives
\(|A_\chi(X)|\le q\).  Termwise differentiation shows that the relevant
tail, up to a sign of modulus one, is
\[
 \sum_{n>B}\chi(n)(\log n)^{j+1}n^{-z}.
 \tag{V46.6}
\]
Apply Stieltjes summation to
\(f_{j,z}(u)=(\log u)^{j+1}u^{-z}\).  Uniformly for \(z\in K\),
\[
 |f_{j,z}'(u)|
 \ll_{K,J}u^{-\sigma_K-1}
 \bigl((\log(2u))^{j}+(\log(2u))^{j+1}\bigr).
 \tag{V46.7}
\]
The endpoint at infinity vanishes, and the endpoint at \(B\) plus the
integral of \textup{(V46.7)} is
\(O_{K,J,\chi}(B^{-\sigma_K}(\log(2B))^{j+1})\).
This proves convergence of the differentiated character series in
\(\Re z>0\).  It equals \(-L'(z,\chi)\) first in \(\Re z>1\), hence
throughout the larger half-plane by analytic continuation.
\end{proof}

This estimate uses only periodicity.  Its power \(B^{-\sigma_K}\) is
strong on an ordinary \(B\)-scale but must still be compared with the
singular heat scale \(e^{c/t}\).

\subsection{Exact descent of the local principal part}

Let \(\rho\) be a nontrivial zero of order \(m\ge1\), and write
\[
 L(z,\chi)=(z-\rho)^m h_\rho(z),
 \qquad h_\rho(\rho)\ne0.
 \tag{V46.8}
\]
Expand
\[
 Q_{\chi,B}(z)=\sum_{j\ge0}q_{B,j}(z-\rho)^j,
 \qquad
 \frac1{h_\rho(z)}=\sum_{\ell\ge0}d_{\rho,\ell}(z-\rho)^\ell,
 \tag{V46.9}
\]
where \(q_{B,j}=Q_{\chi,B}^{(j)}(\rho)/j!\).  For
\(1\le k\le m\), define
\[
 C_{B,-k}(\rho)
 :=\sum_{j+\ell=m-k}q_{B,j}d_{\rho,\ell}.
 \tag{V46.10}
\]

\begin{theorem}[Finite-\(B\) principal-part ledger]
\label{thm:v46-principal-ledger}
In a punctured neighbourhood of \(\rho\),
\[
 \boxed{
 \frac{Q_{\chi,B}(z)}{L(z,\chi)}
 =\sum_{k=1}^{m}\frac{C_{B,-k}(\rho)}{(z-\rho)^k}
  +H_{B,\rho}(z),}
 \tag{V46.11}
\]
where \(H_{B,\rho}\) is holomorphic.  As \(B\to\infty\),
\[
 \boxed{
 C_{B,-1}(\rho)\longrightarrow-m,
 \qquad
 C_{B,-k}(\rho)\longrightarrow0\quad(2\le k\le m).}
 \tag{V46.12}
\]
More quantitatively,
\[
 \boxed{
 |C_{B,-1}(\rho)+m|
 +\sum_{k=2}^{m}|C_{B,-k}(\rho)|
 \ll_{\rho,m,\chi}
 B^{-\Re\rho}(\log(2B))^{m}.}
 \tag{V46.13}
\]
For the first visible cutoff \(B=b_0\), the quotient has pole order
exactly \(m\).  For the completed value \(B=\infty\), only the simple
pole with residue \(-m\) remains.
\end{theorem}

\begin{proof}
Equations \textup{(V46.8)}--\textup{(V46.10)} give
\[
 \frac{Q_{\chi,B}(z)}{L(z,\chi)}
 =(z-\rho)^{-m}
 \left(\sum_{j\ge0}q_{B,j}(z-\rho)^j\right)
 \left(\sum_{\ell\ge0}d_{\rho,\ell}(z-\rho)^\ell\right),
 \tag{V46.14}
\]
which proves \textup{(V46.11)}.

Logarithmic differentiation of \textup{(V46.8)} gives
\[
 -\frac{L'}{L}(z,\chi)
 =-\frac{m}{z-\rho}-\frac{h_\rho'(z)}{h_\rho(z)}.
 \tag{V46.15}
\]
Theorem~\ref{thm:v46-local-completion}, evaluated through derivative order
\(m-1\) at \(\rho\), shows that the Taylor jet of \(Q_{\chi,B}\)
converges to that of \(-L'\), with error
\(O(B^{-\Re\rho}(\log(2B))^m)\).  Substitution in the finite convolution
\textup{(V46.10)} proves \textup{(V46.12)}--\textup{(V46.13)}.

For \(B=b_0\), Theorem~\ref{thm:v45-first-visible} says that
\(Q_{\chi,b_0}(\rho)\ne0\), so
\(C_{b_0,-m}=Q_{\chi,b_0}(\rho)/h_\rho(\rho)\ne0\).  The completed
claim is \textup{(V46.15)}.
\end{proof}

The result also gives the conditionally convergent jet identities
\[
 \sum_{n\ge1}\chi(n)(\log n)^{j+1}n^{-\rho}=0
 \qquad(0\le j\le m-2),
 \tag{V46.16}
\]
with the order-\(m-1\) sum equal to
\((-1)^{m-1}(-L')^{(m-1)}(\rho)\).  These are identities forced by the
multiplicity of the zero; they are not additional zero-free information.

\subsection{Pole order is invisible to normalized heat type}

The inverse Laplace transform of \((z-\rho)^{-k}\) is supported on the
positive logarithmic half-line and contains
\(u^{k-1}e^{\rho u}/(k-1)!\).  After removing the critical-line factor and
tuning the Fourier phase to the ordinate, the relevant scalar model is
therefore
\[
 \mathcal G_{k,t}(\delta)
 :=\frac1{(k-1)!}\int_0^\infty
 u^{k-1}e^{-tu^2+\delta u}\,du,
 \qquad \delta>0.
 \tag{V46.17}
\]

\begin{lemma}[Every fixed half-line pole order has the same heat type]
\label{lem:v46-pole-packet}
For every fixed integer \(k\ge1\) and \(\delta>0\),
\[
 \boxed{
 \mathcal G_{k,t}(\delta)
 =\frac{\sqrt\pi}{(k-1)!\sqrt t}
 \left(\frac{\delta}{2t}\right)^{k-1}
 e^{\delta^2/(4t)}(1+O_{k,\delta}(\sqrt t)),}
 \tag{V46.18}
\]
and hence
\[
 \boxed{
 \lim_{t\downarrow0}4t\log
 \mathcal G_{k,t}(\delta)=\delta^2.}
 \tag{V46.19}
\]
\end{lemma}

\begin{proof}
Complete the square:
\[
 -tu^2+\delta u
 =-t\left(u-\frac{\delta}{2t}\right)^2
  +\frac{\delta^2}{4t}.
 \tag{V46.20}
\]
The saddle \(u_0=\delta/(2t)\) is a distance
\(\delta/(2\sqrt t)\to\infty\) from the lower endpoint in Gaussian
units.  Put \(u=u_0+v/\sqrt t\).  On every fixed \(v\)-interval,
\[
 u^{k-1}=u_0^{k-1}(1+O_{k,\delta}(\sqrt t\,|v|)),
 \tag{V46.21}
\]
and the Gaussian tails permit dominated truncation.  Extending the lower
limit to \(-\infty\) costs \(O(e^{-\delta^2/(8t)})\) relative to the
main Gaussian integral.  Since
\(\int_{\mathbb R}e^{-v^2}\,dv=\sqrt\pi\), this proves
\textup{(V46.18)} and \textup{(V46.19)}.
\end{proof}

Thus changing the pole order changes the power of \(t^{-1}\), but not the
\(1/t\) exponent generated by the horizontal displacement
\(\delta=|\Re\rho-1/2|\).
\begin{theorem}[Pole-order descent preserves the forbidden heat rate]
\label{thm:v46-rate-invariance}
Fix an off-line zero \(\rho\) with displacement
\(\delta=|\Re\rho-1/2|>0\).
\begin{enumerate}[label=\textup{(\roman*)},leftmargin=3em]
\item For \(B=b_0\), the order-\(m\) reciprocal pole in
      \(Q_{\chi,b_0}/L\) produces a nonzero half-line Gaussian packet of
      normalized rate \(\delta^2\) at its phase-tuned ordinate.
\item For \(B=\infty\), the surviving simple pole of \(-L'/L\), with
      residue \(-m\), produces the same normalized rate \(\delta^2\).
\item The cancellation of pole orders \(m,m-1,\ldots,2\) changes only
      subexponential, in particular polynomial, prefactors at fixed
      multiplicity.
\end{enumerate}
\end{theorem}

\begin{proof}
The leading coefficient at \(B=b_0\) is nonzero by
Theorem~\ref{thm:v46-principal-ledger}; at completion the residue \(-m\)
is nonzero.  Inverse Laplace transformation turns the order-\(k\)
principal part into its nonzero coefficient times
\(u^{k-1}e^{\rho u}/(k-1)!\).  At the phase-tuned ordinate,
Lemma~\ref{lem:v46-pole-packet} applies.  In a finite principal part the
largest \(k\) with nonzero coefficient has the largest power of
\(t^{-1}\), so its leading term cannot be cancelled by the lower orders
for all sufficiently small \(t\).  These powers contribute only
\(O(\log(1/t))\) to the logarithm, which vanishes after multiplication by
\(4t\).\end{proof}

This is consistent with the completed-bulk multiplier in
Theorem~\ref{thm:v2-gaussian-multiplier}: the completed logarithmic
derivative already records multiplicity as the coefficient \(-m\), while
the off-line displacement determines the exponential type.

\subsection{The cutoff scale needed to see principal-part completion}

Let \(B=B_t\to\infty\).  From \textup{(V46.13)}, the deterministic
completion error in any principal-part coefficient is bounded by
\[
 \exp\{-\Re\rho\log B_t+O(\log\log(2B_t))\}.
 \tag{V46.22}
\]
If \(\log B_t=o(1/t)\), this is only
\(e^{-o(1/t)}\) on the heat scale.  In particular, the estimate supplies
no strictly negative normalized exponent.  To guarantee from
\textup{(V46.13)} a coefficient saving \(e^{-c/t}\), one must take
\[
 \log B_t\ge \frac{c+o(1)}{\Re\rho\,t}.
 \tag{V46.23}
\]
Such a cutoff lies on the same exponential factor scale as the Gaussian
saddle and is no longer a subexponential small-factor truncation.

\begin{theorem}[V46 pole-completion verdict]
\label{thm:v46-verdict}
Finite logarithmic-factor completion has the following exact profile.
\begin{enumerate}[label=\textup{(\roman*)},leftmargin=3em]
\item Every fixed derivative of \(Q_{\chi,B}\) converges locally
      uniformly to the corresponding derivative of \(-L'\) in
      \(\Re z>0\), with error \(B^{-\sigma}\) up to logarithms.
\item At a zero of multiplicity \(m\), all pole coefficients of orders
      at least two tend to zero, and the simple residue tends to \(-m\).
\item A subexponential cutoff does not turn the quantitative completion
      estimate into a strict gain on the normalized heat scale.
\item Both the initial order-\(m\) pole and the completed simple pole have
      the same off-line rate \(\delta^2\).
\end{enumerate}
Therefore pole-order descent supplies exact amplitude identities but no
new exponential separation between critical-line and off-line zeros.
\end{theorem}

\begin{proof}
Combine Theorems~\ref{thm:v46-local-completion},
\ref{thm:v46-principal-ledger}, and
\ref{thm:v46-rate-invariance} with
\textup{(V46.22)}--\textup{(V46.23)}.
\end{proof}

\begin{firewall}[Vanishing higher principal parts do not prove GRH]
V46 proves exact local convergence and pole-order cancellation.  It does
not show a forbidden heat packet vanishes: the completed simple residue
\(-m\) remains nonzero and carries the same rate \(\delta^2\).  Nor does
the upper bound \textup{(V46.13)} exclude accidental faster convergence
along selected cutoffs.  No off-line zero is excluded and GRH remains
open.
\end{firewall}

\begin{assessment}[V46 closes pole order as an exponential lever]
The multiplicative completion from \(1/L\) to \(-L'/L\) is now fully
accounted for.  Higher reciprocal poles cancel because the logarithmic
factor sum develops exactly the vanishing Taylor jet of \(-L'\) at a
multiple zero.  Those identities are equivalent to the assumed
multiplicity and do not constrain the zero's real part.  The simple
residue \(-m\) survives, and normalized heat type sees location rather
than pole order.

The next useful question must therefore act on the surviving simple
residue, not on the discarded higher orders.  V47 should test whether the
finite completion error
\[
 E_{\chi,B}(z):=-L'(z,\chi)-Q_{\chi,B}(z)
 \tag{V46.24}
\]
admits a source-side sign, phase, or value--derivative constraint at the
phase-tuned contact when \(\log B\asymp1/t\).  The test must be made on
the exact exponential cutoff frontier from \textup{(V46.23)} and must be
abandoned if it merely reconstructs the already-known prime reader or
requires a bound equivalent to GRH.
\end{assessment}

\section*{Version V46 append-only record}
\addcontentsline{toc}{section}{Version V46 append-only record}
The complete V45 source was retained before this continuation.  It had
\(701328\) bytes, \(21258\) lines, and SHA--256
\[
 \texttt{204FC833A97E6EB034CB6DD163483127BDB4EED4C754FA6ABFFDDC622EDD993E}.
\]
V46 proved quantitative local completion of the finite logarithmic-factor
series, computed the full principal-part descent to the residue \(-m\),
and showed that fixed pole order is invisible to the normalized off-line
heat exponent.

% =============================================================================
% END APPENDED V46 MATERIAL
% =============================================================================
% =============================================================================
% BEGIN APPENDED V47 MATERIAL
% =============================================================================

\section{V47: smooth factor completion on the exponential frontier}
\label{sec:grh-v47}

V46 found a \(B^{-\Re\rho}\) convergence rate for the sharp partial sum
\(Q_{\chi,B}\).  At \(B=e^{\lambda/t}\), this becomes visible on the
same \(1/t\)-scale as an off-line heat packet.  The sharp boundary is not,
however, intrinsic.  Because a nonprincipal periodic character has no
zero-frequency Fourier mode, it admits bounded periodic discrete
antiderivatives of every order.  A smooth factor cutoff can therefore be
summed by parts arbitrarily many times.

This gives a genuine improvement: at any fixed exponential cutoff
\(B_t=e^{\lambda/t}\), every unwanted reciprocal principal-part
coefficient can be made exponentially negligible compared with the
off-line packet.  The improvement still does not prove GRH.  What remains
after completion is exactly the simple pole of \(-L'/L\), with residue
\(-m\), and hence exactly the prime reader already isolated in V22--V39.
On the source side the smoothly truncated convolution retains the first
visible fixed-\(b\) slice and full absolute Gaussian type.

\subsection{Periodic discrete antiderivatives of all orders}

Write \(\Delta F(n):=F(n+1)-F(n)\).

\begin{lemma}[Arbitrarily deep periodic summation by parts]
\label{lem:v47-periodic-primitives}
For every integer \(N\ge1\), there is a bounded \(q\)-periodic sequence
\(A_{\chi,N}(n)\) such that
\[
 \boxed{\Delta^N A_{\chi,N}(n)=\chi(n).}
 \tag{V47.1}
\]
Consequently, if a sequence \(f(n)\) and its first \(N-1\) differences
vanish at both ends of its support, then
\[
 \boxed{
 \left|\sum_n\chi(n)f(n)\right|
 \le C_{\chi,N}\sum_n|\Delta^Nf(n)|.}
 \tag{V47.2}
\]
The same conclusion holds for a decaying infinite tail whenever the right
side converges.
\end{lemma}

\begin{proof}
Expand the periodic sequence in its discrete Fourier basis:
\[
 \chi(n)=\sum_{r=0}^{q-1}\widehat\chi(r)e^{2\pi irn/q}.
 \tag{V47.3}
\]
Nonprincipality gives
\(\widehat\chi(0)=q^{-1}\sum_{a=1}^q\chi(a)=0\).  Hence
\[
 A_{\chi,N}(n):=
 \sum_{r=1}^{q-1}
 \frac{\widehat\chi(r)}{(e^{2\pi ir/q}-1)^N}
 e^{2\pi irn/q}
 \tag{V47.4}
\]
is well defined, bounded, periodic, and satisfies \textup{(V47.1)}.
Apply ordinary discrete summation by parts \(N\) times.  The assumed
endpoint conditions remove the boundary terms; taking absolute values and
using boundedness of \(A_{\chi,N}\) proves \textup{(V47.2)}.  Truncation
and passage to the limit give the tail version.
\end{proof}

The existence of every periodic primitive, not merely the bounded first
partial sum in Lemma~\ref{lem:v44-character-partial}, is the source of the
arbitrary-order gain below.

\subsection{A smooth finite logarithmic-factor completion}

Fix once and for all \(\phi\in C^\infty([0,\infty))\) satisfying
\[
 0\le\phi\le1,
 \qquad
 \phi(u)=1\ (0\le u\le1),
 \qquad
 \phi(u)=0\ (u\ge2).
 \tag{V47.5}
\]
Define the finite smooth sum and its completed tail by
\[
 Q_{\chi,B}^{\phi}(z)
 :=\sum_{n\ge2}\chi(n)\log n\,n^{-z}\phi(n/B),
 \tag{V47.6}
\]
\[
 R_{\chi,B}^{\phi}(z)
 :=-L'(z,\chi)-Q_{\chi,B}^{\phi}(z)
 =\sum_{n\ge1}\chi(n)\log n\,n^{-z}
   (1-\phi(n/B)).
 \tag{V47.7}
\]
The first series is supported on \(n<2B\); the second identity holds by
the locally uniform character-series continuation in
Theorem~\ref{thm:v46-local-completion}.

\begin{theorem}[Superalgebraic smooth completion]
\label{thm:v47-smooth-completion}
Let \(K\Subset\{\Re z>0\}\), let
\(\sigma_K=\min_{z\in K}\Re z\), and fix integers \(J\ge0\) and
\(N\ge1\) with \(N+\sigma_K>1\).  Then
\[
 \boxed{
 \max_{0\le j\le J}\sup_{z\in K}
 \left|(R_{\chi,B}^{\phi})^{(j)}(z)\right|
 \ll_{K,J,N,\chi,\phi}
 B^{1-\sigma_K-N}(\log(2B))^{J+N+1}.}
 \tag{V47.8}
\]
Since \(N\) is arbitrary, the smooth completion gains every prescribed
fixed power of \(B^{-1}\).
\end{theorem}

\begin{proof}
For \(0\le j\le J\), termwise differentiation of
\textup{(V47.7)} reduces the assertion, up to a sign, to
\[
 \sum_{n\ge1}\chi(n)f_{B,z,j}(n),
 \qquad
 f_{B,z,j}(u):=(\log u)^{j+1}u^{-z}(1-\phi(u/B)).
 \tag{V47.9}
\]
This smooth function vanishes on \([1,B]\), tends to zero at infinity
together with its first \(N-1\) derivatives, and satisfies
\[
 \int_1^\infty|f_{B,z,j}^{(N)}(u)|\,du
 \ll_{K,J,N,\phi}
 B^{1-\sigma_K-N}(\log(2B))^{J+N+1}.
 \tag{V47.10}
\]
Indeed, the terms in which derivatives hit \(\phi(u/B)\) are supported
on \([B,2B]\) and each cutoff derivative contributes \(B^{-1}\); beyond
\(2B\), ordinary differentiation of
\((\log u)^{j+1}u^{-z}\) gives the same bound.  The condition
\(N+\sigma_K>1\) makes the final tail integral converge.

The finite-difference identity
\[
 \Delta^Nf(n)=
 \int_{[0,1]^N}f^{(N)}(n+u_1+\cdots+u_N)
 \,du_1\cdots du_N
 \tag{V47.11}
\]
implies
\(\sum_n|\Delta^Nf(n)|\ll_N\int_1^\infty|f^{(N)}(u)|du\), with an
inessential compact-endpoint adjustment.  Apply
Lemma~\ref{lem:v47-periodic-primitives} and \textup{(V47.10)}.
\end{proof}

Thus V46's first-power estimate is the correct sharp-cutoff bound, but it
is not a regularity barrier for factor completion.

\subsection{Principal parts on the exponential cutoff frontier}

Let \(\rho\) be a zero of multiplicity \(m\), write
\(\beta=\Re\rho>0\), and define the principal-part coefficients
\(C_{B,-k}^{\phi}(\rho)\) for
\(Q_{\chi,B}^{\phi}/L\) exactly as in
\textup{(V46.9)}--\textup{(V46.11)}.

\begin{theorem}[Arbitrary-order smooth pole descent]
\label{thm:v47-smooth-poles}
For every fixed \(N\ge1\),
\[
 \boxed{
 |C_{B,-1}^{\phi}(\rho)+m|
 +\sum_{k=2}^{m}|C_{B,-k}^{\phi}(\rho)|
 \ll_{\rho,m,N,\chi,\phi}
 B^{1-\beta-N}(\log(2B))^{m+N}.}
 \tag{V47.12}
\]
Let \(B_t=e^{\lambda/t}\) with \(\lambda>0\).  Then
\[
 \boxed{
 \limsup_{t\downarrow0}4t\log\left(
 |C_{B_t,-1}^{\phi}(\rho)+m|
 +\sum_{k=2}^{m}|C_{B_t,-k}^{\phi}(\rho)|
 \right)
 \le-4\lambda(N-1+\beta).}
 \tag{V47.13}
\]
\end{theorem}

\begin{proof}
The difference between the finite quotient and the completed logarithmic
derivative is
\[
 \frac{Q_{\chi,B}^{\phi}(z)}{L(z,\chi)}
 +\frac{L'}{L}(z,\chi)
 =-\frac{R_{\chi,B}^{\phi}(z)}{L(z,\chi)}.
 \tag{V47.14}
\]
As in \textup{(V46.10)}, every principal-part coefficient on the right is
a fixed finite linear combination of the first \(m-1\) derivatives of
\(R_{\chi,B}^{\phi}\) at \(\rho\).  Apply
Theorem~\ref{thm:v47-smooth-completion} with the compact set
\(K=\{\rho\}\) and \(J=m-1\).  This proves
\textup{(V47.12)}.  Substitution of \(B_t=e^{\lambda/t}\) gives
\[
 4t(1-\beta-N)\log B_t
 =-4\lambda(N-1+\beta),
 \tag{V47.15}
\]
while all logarithmic factors have zero normalized rate.
\end{proof}

\begin{corollary}[The completion error can be removed below every
off-line heat scale]
\label{cor:v47-error-removal}
Let \(\delta=|\beta-1/2|>0\).  The pole packet arising from the difference
in \textup{(V47.14)} has normalized rate at most
\[
 \boxed{
 \delta^2-4\lambda(N-1+\beta).}
 \tag{V47.16}
\]
For every fixed \(\lambda>0\), choosing \(N\) sufficiently large makes
this rate negative.  The conclusion uses only periodicity of \(\chi\),
not GRH.
\end{corollary}

\begin{proof}
Each order-\(k\) coefficient has the bound in
Theorem~\ref{thm:v47-smooth-poles}.  Its half-line Gaussian pole packet
has rate \(\delta^2\) and only a polynomial order-dependent prefactor by
Lemma~\ref{lem:v46-pole-packet}.  Adding the coefficient rate in
\textup{(V47.13)} proves \textup{(V47.16)}.  There are only \(m\)
principal-part terms.
\end{proof}

The corollary is local at a fixed zero.  It says that the smooth finite
factor sum reproduces the completed simple residue \(-m\) to arbitrary
exponential accuracy on the heat scale.  It does not say that the
surviving simple-pole packet is small.

\subsection{The exact source split and its unchanged absolute saddle}

Insert the same smooth partition into the multiplicative reader.  Define
\[
 \begin{aligned}
 \mathcal P_{B,t}^{\phi,\sigma,[r]}(x;s)
 :=\sum_{\substack{a,b\ge1\\ab\ge2}}
 &\frac{\mu(a)\chi(a)\chi(b)\log b\,(\log(ab))^r}{\sqrt{ab}}
 \phi(b/B)\\
 &{}\times e^{-t\log^2(ab)}e^{ix\log(ab)}W_t^\sigma(ab;s),
 \end{aligned}
 \tag{V47.17}
\]
and define \(\mathcal E_{B,t}^{\phi,\sigma,[r]}\) by replacing
\(\phi(b/B)\) with \(1-\phi(b/B)\).

\begin{theorem}[Exact smooth source partition]
\label{thm:v47-source-partition}
For all parameters,
\[
 \boxed{
 P_{\chi,t}^{\sigma,[r]}(x;s)
 =\mathcal P_{B,t}^{\phi,\sigma,[r]}(x;s)
  +\mathcal E_{B,t}^{\phi,\sigma,[r]}(x;s).}
 \tag{V47.18}
\]
For every \(B>b_0\), the partial term contains the complete fixed slice
\(b=b_0\).  The absolute mass of that slice has normalized rate \(1/4\):
\[
 \boxed{
 \lim_{t\downarrow0}4t\log
 \sum_{a\ge1}
 \frac{|\mu(a)\chi(a)\chi(b_0)|\log b_0
       (1+\log(ab_0))^r}{\sqrt{ab_0}}
 e^{-t\log^2(ab_0)}W_t^-(ab_0;s)
 =\frac14.}
 \tag{V47.19}
\]
Thus smooth completion does not create an absolute source-side gain.
\end{theorem}

\begin{proof}
Equation \textup{(V47.18)} is the identity
\(\phi+(1-\phi)=1\) inside the absolutely convergent convolution.  Since
\(b_0/B<1\), \textup{(V47.5)} gives \(\phi(b_0/B)=1\).

For \textup{(V47.19)}, squarefree integers coprime to \(q\) have positive
natural density.  Restrict \(\log a\) to a fixed-width interval about
\(1/(4t)-\log b_0\).  Their count has exponential order
\(e^{1/(4t)}\), while the factor
\(a^{-1/2}e^{-t\log^2(ab_0)}\) contributes
\(e^{-3/(16t)+o(1/t)}\); the soft prefix tends to a positive constant at
the saddle.  This gives the lower rate \(1/4\).  The matching upper bound
follows from the one-dimensional version of \textup{(V44.12)}--
\textup{(V44.14)}.
\end{proof}

The signed fixed slice is the reciprocal-\(L\) object of
Theorem~\ref{thm:v45-first-visible}.  Hence neither its signed control nor
its absolute control has become easier merely because the complementary
completion error is smooth.

\subsection{The error test reconstructs the original prime reader}

\begin{theorem}[V47 smooth-completion verdict]
\label{thm:v47-verdict}
At every fixed off-line zero and every exponential factor scale
\(B_t=e^{\lambda/t}\), the following hold.
\begin{enumerate}[label=\textup{(\roman*)},leftmargin=3em]
\item Smooth periodic summation can make all higher reciprocal
      principal parts and the simple-residue error exponentially
      negligible at any prescribed fixed rate.
\item The surviving local object is the simple pole of \(-L'/L\) with
      residue \(-m\), whose heat rate remains \(\delta^2\).
\item The source-side finite completion retains the fixed \(b_0\) slice
      and full absolute rate \(1/4\).
\item The zero-contact equation still constrains only the sum in
      \textup{(V47.18)}, not either smooth factor sector separately.
\end{enumerate}
Therefore the finite completion error admits arbitrary smooth suppression,
but suppressing it merely reconstructs the prime reader already present
before the Möbius convolution.
\end{theorem}

\begin{proof}
Items (i)--(iii) are Corollary~\ref{cor:v47-error-removal},
Theorem~\ref{thm:v46-rate-invariance}, and
Theorem~\ref{thm:v47-source-partition}.  Item (iv) follows by inserting
\textup{(V47.18)} into the exact value--slope ledger
\textup{(V44.29)}--\textup{(V44.30)}; no new scalar equation is created
by an algebraic partition of unity.
\end{proof}

\begin{firewall}[Arbitrary smooth completion is not cancellation of the
simple residue]
V47 removes the finite-factor approximation error near a fixed zero.  It
does not remove the completed residue \(-m\), bound the signed reciprocal
slice, separate either smooth source term at a contact, exclude an
off-line zero, or prove GRH.  The arbitrary summation depth \(N\) is fixed
before \(t\downarrow0\); no claim with \(N=N_t\to\infty\) or uniform
control of its constants is made.
\end{firewall}

\begin{assessment}[V47 closes finite factor completion as a new source of
coercivity]
The sharp-frontier obstruction in V46 was a boundary artifact.  Smooth
completion can approximate the local logarithmic derivative to arbitrary
exponential accuracy using only the periodicity of \(\chi\).  This is a
useful exact acquisition, but after the error is removed the same
quantized simple residue and the same prime-side contact remain.  The
multiplicative factorization has therefore been followed upstream and
back to its starting object without generating a new sign inequality.

V48 should leave the factor-completion branch.  The next nonduplicative
test is at the completed simple-pole level: derive the full parabolic jet
of the scalar Weil symbol at a zero-level contact and ask whether the
renewal geometry supplies a source-side inequality involving
\(m_t,m_{xx},m_{xxx},m_{xxxx}\) that is stronger than the already-used
value and first-derivative equations.  If the heat equation makes every
such relation an identity or if the required sign is equivalent to GRH,
that contact-jet route should also be closed explicitly.
\end{assessment}

\section*{Version V47 append-only record}
\addcontentsline{toc}{section}{Version V47 append-only record}
The complete V46 source was retained before this continuation.  It had
\(714792\) bytes, \(21639\) lines, and SHA--256
\[
 \texttt{716C6DE4784D2CA822E9BBF2DB52D681BBF7AD4829403EBC386BCD339A257BC8}.
\]
V47 introduced a smooth logarithmic-factor partition, proved
arbitrary-order periodic completion and exponential principal-part
suppression at \(B=e^{\lambda/t}\), and showed that the surviving object
is exactly the original simple-residue prime reader with unchanged
absolute source type.

% =============================================================================
% END APPENDED V47 MATERIAL
% =============================================================================
% =============================================================================
% BEGIN APPENDED V48 MATERIAL
% =============================================================================

\section{V48: the full parabolic contact jet and its caloric freedom}
\label{sec:grh-v48}

V47 returned the multiplicative completion branch to the completed
simple-pole heat symbol.  The next proposed source of rigidity was the
higher parabolic jet at a zero-level contact.  V22 already proved the
even-order opening law, V26--V27 derived the simple-boundary velocity, and
V31 showed that the prime-reader time connection collapses to the ordinary
heat hierarchy.  What remained was to assemble the full mixed jet and ask
whether its higher entries obey a new sign condition.

They do not at the level of heat kinematics.  Every mixed derivative is a
spatial derivative with the same parabolic weight.  At a nondegenerate
critical minimum, the first motion of the minimum is positive, but its
second derivative has no sign.  More strongly, explicit globally
nonnegative polynomial data produce heat flows with a genuine positive
threshold and arbitrary prescribed third and fourth spatial derivatives
at the contact.  The higher reader equations are therefore exact moment
ledgers, not additional scalar constraints.

\subsection{Collapse of the mixed parabolic jet}

Let \((t_*,x_*)\) be any positive-time point and put
\[
 c_j:=\partial_x^jm_{\chi,t_*}(x_*),
 \qquad j\ge0.
 \tag{V48.1}
\]

\begin{theorem}[Complete mixed-jet collapse]
\label{thm:v48-mixed-jet}
For every pair of nonnegative integers \(a,b\),
\[
 \boxed{
 \partial_t^a\partial_x^b m_{\chi,t}(x)
 \big|_{(t_*,x_*)}=c_{2a+b}.}
 \tag{V48.2}
\]
If \(x_*\) is a zero of order \(2r\) of the nonnegative critical symbol
\(m_{\chi,t_*}\), then
\[
 \boxed{
 \partial_t^a\partial_x^b m_{\chi,t}(x)
 \big|_{(t_*,x_*)}=0
 \quad\text{when }2a+b<2r,}
 \tag{V48.3}
\]
and every entry on the first nonzero parabolic diagonal is the same
positive number:
\[
 \boxed{
 \partial_t^a\partial_x^b m_{\chi,t}(x)
 \big|_{(t_*,x_*)}=c_{2r}>0
 \quad\text{when }2a+b=2r.}
 \tag{V48.4}
\]
There are no further identities among higher entries beyond equality of
entries with the same weight \(2a+b\).
\end{theorem}

\begin{proof}
The heat equation \textup{(V22.3)} and commutation of positive-time
derivatives give
\[
 \partial_t^a\partial_x^b m
 =\partial_x^{2a+b}m,
 \tag{V48.5}
\]
which proves \textup{(V48.2)}.  The critical-contact structure
\textup{(V22.20)} says
\(c_0=\cdots=c_{2r-1}=0\) and \(c_{2r}>0\), proving
\textup{(V48.3)}--\textup{(V48.4)}.  The last assertion will also follow
constructively from Theorem~\ref{thm:v48-caloric-freedom} below.
\end{proof}

For \(r=1\), the complete jet through parabolic weight four begins
\[
 m=m_x=0,
 \qquad
 m_t=m_{xx}=c_2>0,
 \qquad
 m_{tx}=m_{xxx}=c_3,
 \qquad
 m_{tt}=m_{txx}=m_{xxxx}=c_4.
 \tag{V48.6}
\]
Thus the third and fourth spatial derivatives are the only new scalars at
these weights, and the heat equation assigns them no sign.

\subsection{Motion and curvature of a nondegenerate critical minimum}

Assume now that \(c_2=\kappa>0\).  By the implicit-function theorem
there is a unique analytic critical-point curve \(x_c(t)\) near \(t_*\)
with
\[
 x_c(t_*)=x_*,
 \qquad
 \partial_xm_{\chi,t}(x_c(t))=0.
 \tag{V48.7}
\]
Put
\[
 \nu_c(t):=m_{\chi,t}(x_c(t)).
 \tag{V48.8}
\]

\begin{theorem}[Exact minimum-curve jet]
\label{thm:v48-minimum-curve}
At a nondegenerate critical contact,
\[
 \boxed{
 x_c'(t_*)=-\frac{c_3}{\kappa},}
 \tag{V48.9}
\]
\[
 \boxed{
 \nu_c(t_*)=0,
 \qquad
 \nu_c'(t_*)=\kappa>0,}
 \tag{V48.10}
\]
and
\[
 \boxed{
 \nu_c''(t_*)
 =c_4-\frac{c_3^2}{\kappa}.}
 \tag{V48.11}
\]
In particular, the heat equation and critical nonnegativity determine the
first opening velocity but give no sign for the curvature in
\textup{(V48.11)}.
\end{theorem}

\begin{proof}
Differentiate \(m_x(t,x_c(t))=0\):
\[
 0=m_{xt}+m_{xx}x_c'
 =m_{xxx}+m_{xx}x_c',
 \tag{V48.12}
\]
which proves \textup{(V48.9)}.  Along the curve,
\[
 \nu_c'=m_t+m_x\,x_c'=m_t=m_{xx},
 \tag{V48.13}
\]
giving \textup{(V48.10)}.  A second differentiation yields
\[
 \nu_c''
 =m_{tt}+2m_{tx}x_c'+m_{xx}(x_c')^2+m_x\,x_c''.
 \tag{V48.14}
\]
The final term vanishes, while
\(m_{tt}=m_{xxxx}=c_4\),
\(m_{tx}=m_{xxx}=c_3\), and
\(x_c'=-c_3/\kappa\).  Substitution proves
\textup{(V48.11)}.
\end{proof}

The upper-interval positivity of the arithmetic symbol forces
\(\nu_c(t)\ge0\) for \(t\ge t_*\) near the threshold.  Because the
linear term in \textup{(V48.10)} is already positive, this one-sided fact
places no restriction on the quadratic coefficient
\textup{(V48.11)}.

\subsection{A realization theorem showing higher-jet freedom}

The preceding lack of sign is not an artifact of formal differentiation.

\begin{theorem}[Caloric realization of arbitrary third and fourth jets]
\label{thm:v48-caloric-freedom}
Fix arbitrary real numbers \(c_3,c_4\), a number \(\kappa>0\), and a
positive time \(t_*>0\).  There is a real polynomial heat solution
\(u_t=u_{xx}\) on \(\mathbb R^2\) such that
\begin{enumerate}[label=\textup{(\roman*)},leftmargin=3em]
\item \(u(t,x)\ge0\) for every \(t\ge t_*\) and every real \(x\);
\item \(u(t_*,0)=u_x(t_*,0)=0\) and
\[
 u_{xx}(t_*,0)=\kappa,
 \quad
 u_{xxx}(t_*,0)=c_3,
 \quad
 u_{xxxx}(t_*,0)=c_4;
 \tag{V48.15}
\]
\item for every sufficiently small \(\varepsilon>0\),
\(u(t_*-\varepsilon,0)<0\).
\end{enumerate}
Hence \(t_*\) is a genuine local left endpoint of the positivity interval,
while \(c_3\), \(c_4\), the drift \textup{(V48.9)}, and the curvature
\textup{(V48.11)} are freely prescribable.
\end{theorem}

\begin{proof}
Choose \(K>0\) sufficiently large that
\[
 f(x):=\frac{\kappa}{2}x^2+\frac{c_3}{6}x^3
       +\frac{c_4}{24}x^4+Kx^6
 \tag{V48.16}
\]
is nonnegative on \(\mathbb R\).  Such a \(K\) exists: after factoring
\(x^2\), the remaining quartic is positive near zero because its constant
term is \(\kappa/2\), and \(Kx^4\) dominates the fixed lower-degree
terms away from a sufficiently small neighbourhood of zero.

For \(\tau=t-t_*\), set
\[
 u(t,x):=e^{\tau\partial_x^2}f(x)
 =\sum_{j=0}^{3}\frac{\tau^j}{j!}\partial_x^{2j}f(x).
 \tag{V48.17}
\]
The sum is finite, so it defines a polynomial for all real \(t,x\) and
satisfies \(u_t=u_{xx}\).  When \(\tau\ge0\), it is equivalently the
Gaussian expectation
\[
 u(t,x)=\mathbb E\,f(x+\sqrt{2\tau}\,Z)\ge0,
 \tag{V48.18}
\]
where \(Z\) is standard normal.  At \(\tau=0\), the derivatives through
order four are those of \(f\), giving \textup{(V48.15)}.  Finally,
\[
 u(t_*+\tau,0)
 =\kappa\tau+\frac{c_4}{2}\tau^2+120K\tau^3.
 \tag{V48.19}
\]
For sufficiently small negative \(\tau\), the first term dominates and
the value is negative.  This proves all assertions.
\end{proof}

By choosing \(c_4\) on either side of \(c_3^2/\kappa\), the model makes
\(\nu_c''(t_*)\) positive, zero, or negative without changing any lower
critical-contact condition.

\subsection{Simple zero-level boundaries have the same freedom}

The small-time phase-tuned contacts in V37--V39 are simple endpoints of
negative lobes rather than minima.  Let \(x_b(t)\) be any simple zero
curve, so
\[
 m_{\chi,t}(x_b(t))=0,
 \qquad
 s_*:=m_x(t_*,x_b(t_*))\ne0.
 \tag{V48.20}
\]

\begin{proposition}[Second boundary-velocity identity]
\label{prop:v48-boundary-acceleration}
At a simple zero-level boundary,
\[
 \boxed{
 x_b'=-\frac{m_{xx}}{m_x},}
 \tag{V48.21}
\]
and
\[
 \boxed{
 x_b''=-\frac{
 m_{xxxx}+2x_b'm_{xxx}+(x_b')^2m_{xx}}
 {m_x}.}
 \tag{V48.22}
\]
All derivatives are evaluated on the moving boundary.  Apart from the
orientation sign of \(m_x\), neither expression has a universal sign.
\end{proposition}

\begin{proof}
The first formula is \textup{(V26.32)}.  Differentiate
\(m(t,x_b(t))=0\) twice and use
\(m_t=m_{xx}\), \(m_{tt}=m_{xxxx}\), and
\(m_{tx}=m_{xxx}\):
\[
 0=m_{xxxx}+2x_b'm_{xxx}+x_b''m_x+(x_b')^2m_{xx}.
 \tag{V48.23}
\]
Solving for \(x_b''\) gives \textup{(V48.22)}.  Arbitrary spatial jets
at a simple crossing are realized by applying
\(e^{(t-t_*)\partial_x^2}\) to a polynomial beginning with
\(s_*x+c_2x^2/2+c_3x^3/6+c_4x^4/24\).  Thus caloricity supplies no
additional sign.
\end{proof}

\subsection{The higher prime-reader moment ledger}

For the unsplit reader of V28, positive-time differentiation gives
\[
 \partial_x^jP_{\chi,t}^{[0]}(x)=i^jP_{\chi,t}^{[j]}(x),
 \qquad
 \partial_t^aP_{\chi,t}^{[j]}(x)=(-1)^a
 P_{\chi,t}^{[j+2a]}(x).
 \tag{V48.24}
\]

\begin{theorem}[All higher contact equations are one spatial moment
ledger]
\label{thm:v48-reader-jet}
At a critical contact of order \(2r\),
\[
 \boxed{
 2\Re\left(i^jP_{\chi,t_*}^{[j]}(x_*)\right)
 =\partial_x^ja_{\chi,t_*}(x_*)
 \quad(0\le j<2r),}
 \tag{V48.25}
\]
while
\[
 \boxed{
 \partial_x^{2r}a_{\chi,t_*}(x_*)
 -2(-1)^r\Re P_{\chi,t_*}^{[2r]}(x_*)>0.}
 \tag{V48.26}
\]
Every equation obtained by applying any mixed derivative
\(\partial_t^a\partial_x^b\) is exactly the member of this same hierarchy
with index \(j=2a+b\).  Thus time differentiation creates no independent
reader constraint.
\end{theorem}

\begin{proof}
Differentiate
\[
 m_{\chi,t}(x)=a_{\chi,t}(x)
 -2\Re P_{\chi,t}^{[0]}(x)
 \tag{V48.27}
\]
and use the first identity in \textup{(V48.24)} together with the
even-order contact conditions \textup{(V22.20)}.  This gives
\textup{(V48.25)}--\textup{(V48.26)}.  The second identity in
\textup{(V48.24)} and the archimedean heat equation identify every mixed
derivative with the spatial index \(2a+b\), as in
Theorem~\ref{thm:v48-mixed-jet}.
\end{proof}

For a nondegenerate contact, \textup{(V48.25)} contains only the already
used value and first-derivative equations.  The second reader moment enters
through the opening inequality \textup{(V48.26)}, while the third and
fourth moments determine the freely signed drift and curvature data above.
Absolute bounds for these fixed reader orders retain the Gaussian saddle
rate \(1/4\), as already recorded in V31.

\subsection{The parabolic jet supplies no new kinematic coercivity}

\begin{theorem}[V48 contact-jet verdict]
\label{thm:v48-verdict}
The complete higher contact jet has the following exact status.
\begin{enumerate}[label=\textup{(\roman*)},leftmargin=3em]
\item Its mixed derivatives collapse to one spatial sequence indexed by
      parabolic weight \(2a+b\).
\item At a nondegenerate critical minimum, only
      \(m_t=m_{xx}>0\) has a forced sign; the minimum drift and second
      time curvature are unrestricted.
\item At a simple moving boundary, velocity and acceleration are exact
      rational functions of the spatial jet but have no new sign.
\item Globally nonnegative polynomial heat data realize arbitrary third
      and fourth jets at a genuine positive threshold.
\item On the source side, all higher equations are the existing reader
      moment hierarchy, with no additional equation from time evolution.
\end{enumerate}
Therefore neither the heat equation nor renewal contact kinematics gives
the missing inequality needed to force the heat threshold to zero.
\end{theorem}

\begin{proof}
Combine Theorems~\ref{thm:v48-mixed-jet},
\ref{thm:v48-minimum-curve}, and
\ref{thm:v48-caloric-freedom}, Proposition
\ref{prop:v48-boundary-acceleration}, and
Theorem~\ref{thm:v48-reader-jet}.
\end{proof}

\begin{firewall}[Caloric jet freedom does not model the arithmetic
coefficients]
The polynomial realization proves that no desired higher-jet sign follows
from the heat equation, positivity upper interval, and contact geometry
alone.  It does not claim that an arbitrary caloric polynomial is a Weil
symbol.  An additional inequality could still arise from the arithmetic
prime coefficients, but it must be proved from those coefficients and
would be new input beyond heat kinematics.  No off-line zero is excluded
and GRH remains open.
\end{firewall}

\begin{assessment}[V48 closes the generic higher-contact-jet route]
The full parabolic jet contains no hidden extra evolution equation.
Nondegenerate critical opening is universally linear, and all curvature
and drift information begins in freely prescribable higher spatial
derivatives.  On the prime side those derivatives are merely higher
logarithmic moments, whose absolute scale remains critical.

The next possible arithmetic test is finite-dimensional and sharper than
another heat derivative.  At a nondegenerate critical contact, normalize
the zeroth, first, and second prime readers by the Gaussian saddle measure
of V28 and solve their exact unit-phasor moment feasibility problem under
the value equation, the zero-slope equation, and the positive-opening
inequality.  V49 should determine whether this three-moment contact body
has a strict gap for arithmetic phases.  If the sharp moment body contains
the off-line contact data at full type, the contact-moment route should be
closed rather than extended to still higher derivatives.
\end{assessment}

\section*{Version V48 append-only record}
\addcontentsline{toc}{section}{Version V48 append-only record}
The complete V47 source was retained before this continuation.  It had
\(728563\) bytes, \(22019\) lines, and SHA--256
\[
 \texttt{5C3CD6510C782AE892EDED3D590A13CA1FA42A18AF7321670A4A520ED3594306}.
\]
V48 derived the full mixed parabolic contact jet, computed minimum and
boundary curvature formulas, constructed positive-threshold caloric
realizations with arbitrary higher jets, and proved that all source-side
time derivatives collapse to the existing reader moment hierarchy.

% =============================================================================
% END APPENDED V48 MATERIAL
% =============================================================================
% =============================================================================
% BEGIN APPENDED V49 MATERIAL
% =============================================================================

\section{V49: the sharp three-moment contact body}
\label{sec:grh-v49}

V48 showed that the second, third, and fourth contact derivatives carry no
universal caloric sign beyond the positive quadratic opening.  This section
tests whether the zeroth, first, and second prime readers nevertheless lie
in a unit-phasor moment body that excludes such an opening.

The body admits an exact Hilbert-space description through three
orthogonal polynomial coordinates.  In the Gaussian saddle limit it has a
simple interior witness: the real threshold phasor
\(\zeta=\operatorname{sgn}X\) has zero mean phasor, zero imaginary first
moment, and a strictly positive raw second-reader opening.  A two-parameter
threshold and constant phase rotation absorbs small archimedean value and
slope targets.  Hence the three contact conditions are robustly feasible
at full Gaussian type.  Any exclusion must use the actual multiplicative
curve \(\chi(n)e^{ix\log n}\), not merely unit modulus and the first four
absolute saddle moments.

\subsection{Standardized second-reader coordinates}

On the V28 probability space, put
\[
 X_t(n):=\frac{u_n-\mu_\chi(t)}{\sigma_\chi(t)},
 \qquad
 \zeta_{n,x}:=\frac{\chi(n)}{|\chi(n)|}e^{ixu_n}.
 \tag{V49.1}
\]
Define
\[
 z_t(x):=\mathbb E_t\zeta_{n,x},
 \qquad
 c_t(x):=\mathbb E_t[X_t\zeta_{n,x}],
 \qquad
 d_t(x):=\mathbb E_t[X_t^2\zeta_{n,x}],
 \tag{V49.2}
\]
where \(\mathbb E_t=\mathbb E_{\nu_{\chi,t}}\).  Then
\[
 \boxed{
 \frac{P_{\chi,t}^{[0]}}{Z_\chi}=z_t,}
 \tag{V49.3}
\]
\[
 \boxed{
 \frac{P_{\chi,t}^{[1]}}{Z_\chi}
 =\mu_\chi z_t+\sigma_\chi c_t,}
 \tag{V49.4}
\]
\[
 \boxed{
 \frac{P_{\chi,t}^{[2]}}{Z_\chi}
 =\mu_\chi^2z_t+2\mu_\chi\sigma_\chi c_t
  +\sigma_\chi^2d_t.}
 \tag{V49.5}
\]
All functions on the right are evaluated at \(x\).

Let
\[
 \kappa_{3,t}:=\mathbb E_tX_t^3,
 \qquad
 \kappa_{4,t}:=\mathbb E_tX_t^4,
 \qquad
 \tau_t^2:=\kappa_{4,t}-1-\kappa_{3,t}^2.
 \tag{V49.6}
\]

\begin{proposition}[Gaussian standardized saddle]
\label{prop:v49-gaussian-saddle}
As \(t\downarrow0\), the law of \(X_t\) converges weakly, with every
fixed absolute moment, to the standard normal law.  In particular,
\[
 \boxed{
 \kappa_{3,t}\longrightarrow0,
 \qquad
 \kappa_{4,t}\longrightarrow3,
 \qquad
 \tau_t^2\longrightarrow2.}
 \tag{V49.7}
\]
\end{proposition}

\begin{proof}
The PNT partial-summation reduction in \textup{(V28.7)} applies after
inserting any fixed continuous function of
\(\sqrt{2t}(u-1/(4t))\) with polynomial growth.  Completing the square as
in \textup{(V28.8)} gives the standard normal density after normalization.
The progression-removal error is exponentially smaller than every fixed
polynomial moment.  Equations \textup{(V28.5)}--\textup{(V28.6)} permit
replacement of the approximate center and scale by the exact
\(\mu_\chi(t)\) and \(\sigma_\chi(t)\).  This proves weak convergence and
uniform integrability of every fixed moment, hence \textup{(V49.7)}.
\end{proof}

\subsection{The exact Bessel ellipsoid}

Set
\[
 Q_{2,t}(X):=X^2-1-\kappa_{3,t}X.
 \tag{V49.8}
\]
The functions
\[
 1,qquad X_t,qquad \frac{Q_{2,t}(X_t)}{\tau_t}
 \tag{V49.9}
\]
are orthonormal in real \(L^2(\nu_{\chi,t})\) whenever \(\tau_t>0\).

\begin{theorem}[Sharp three-coordinate Bessel body]
\label{thm:v49-bessel-body}
For every \(t>0\) with \(\tau_t>0\), every real \(x\), and more generally
every measurable phasor \(\zeta\) with \(|\zeta|=1\),
\[
 \boxed{
 |z|^2+|c|^2+
 \frac{|d-z-\kappa_{3,t}c|^2}{\tau_t^2}
 \le1.}
 \tag{V49.10}
\]
Here \(z=\mathbb E_t\zeta\),
\(c=\mathbb E_t(X_t\zeta)\), and
\(d=\mathbb E_t(X_t^2\zeta)\).  Equivalently, the residual after
orthogonal projection of \(\zeta\) onto the three functions in
\textup{(V49.9)} has squared norm
\[
 \boxed{
 1-|z|^2-|c|^2-
 \frac{|d-z-\kappa_{3,t}c|^2}{\tau_t^2}.}
 \tag{V49.11}
\]
\end{theorem}

\begin{proof}
Since \(\mathbb E_tX_t=0\) and \(\mathbb E_tX_t^2=1\),
\[
 \mathbb E_tQ_{2,t}(X_t)=0,
 \qquad
 \mathbb E_t[X_tQ_{2,t}(X_t)]=0.
 \tag{V49.12}
\]
Direct expansion gives
\[
 \mathbb E_tQ_{2,t}(X_t)^2
 =\kappa_{4,t}-1-\kappa_{3,t}^2=\tau_t^2.
 \tag{V49.13}
\]
The three projection coefficients of \(\zeta\) are
\(z\), \(c\), and
\((d-z-\kappa_{3,t}c)/\tau_t\).  Bessel's inequality and
\(\|\zeta\|_2^2=1\) prove \textup{(V49.10)}.  Pythagoras gives
\textup{(V49.11)}.
\end{proof}

By Proposition~\ref{prop:v49-gaussian-saddle}, the limiting ellipsoid is
\[
 \boxed{
 |z|^2+|c|^2+\frac12|d-z|^2\le1.}
 \tag{V49.14}
\]
This is the strongest quadratic constraint furnished by the first three
orthogonal saddle coordinates.

\subsection{Critical contact as an affine slice of the body}

At a nondegenerate critical contact \((t,x)\), define the normalized
archimedean targets
\[
 r_t(x):=\frac{a_{\chi,t}(x)}{2Z_\chi(t)},
 \qquad
 s_t(x):=-\frac{\partial_xa_{\chi,t}(x)}{2Z_\chi(t)},
 \qquad
 h_t(x):=-\frac{\partial_x^2a_{\chi,t}(x)}{2Z_\chi(t)}.
 \tag{V49.15}
\]

\begin{theorem}[Exact three-moment contact slice]
\label{thm:v49-contact-slice}
The value equation, zero-slope equation, and positive-opening inequality
are respectively
\[
 \boxed{\Re z_t(x)=r_t(x),}
 \tag{V49.16}
\]
\[
 \boxed{
 \Im\bigl(\mu_\chi(t)z_t(x)+\sigma_\chi(t)c_t(x)\bigr)
 =s_t(x),}
 \tag{V49.17}
\]
and
\[
 \boxed{
 \Re\bigl(
 \mu_\chi(t)^2z_t(x)
 +2\mu_\chi(t)\sigma_\chi(t)c_t(x)
 +\sigma_\chi(t)^2d_t(x)
 \bigr)>h_t(x).}
 \tag{V49.18}
\]
On every fixed compact \(x\)-interval, all three targets in
\textup{(V49.15)} satisfy
\[
 |r_t(x)|+|s_t(x)|+|h_t(x)|
 \le\exp\left(-\frac1{16t}+O_{\chi,I}(\log(1/t))\right)
 \tag{V49.19}
\]
as \(t\downarrow0\).
\end{theorem}

\begin{proof}
Equations \textup{(V49.16)}--\textup{(V49.18)} follow by inserting
\textup{(V49.3)}--\textup{(V49.5)} into
\[
 m=a-2\Re P^{[0]},
 \qquad
 m_x=a_x+2\Im P^{[1]},
 \qquad
 m_{xx}=a_{xx}+2\Re P^{[2]},
 \tag{V49.20}
\]
and using \(m=m_x=0\), \(m_{xx}>0\).  On a compact interval, every
fixed derivative of \(a_{\chi,t}\) has at most polynomial growth in
\(t^{-1}\), whereas \(Z_\chi(t)\) has the asymptotic
\textup{(V28.4)}.  This proves \textup{(V49.19)}.
\end{proof}

The zero-slope equation removes the unrestricted imaginary carrier from
V28, but it constrains an imaginary coordinate.  The opening inequality
acts on an independent real coordinate, leaving room inside the Bessel
body.

\subsection{A robust Gaussian witness}

Let \(G\) be standard normal, let \(\Phi\) and \(\varphi\) be its
distribution function and density, and define
\[
 \epsilon_a(G):=
 \begin{cases}
 1,&G>a,\\
 -1,&G\le a.
 \end{cases}
 \qquad
 \zeta_{a,\phi}(G):=e^{i\phi}\epsilon_a(G).
 \tag{V49.21}
\]
Its first three correlations are
\[
 z_0(a)=1-2\Phi(a),
 \qquad
 c_0(a)=2\varphi(a),
 \qquad
 d_0(a)=z_0(a)+2a\varphi(a),
 \tag{V49.22}
\]
multiplied by \(e^{i\phi}\).

\begin{theorem}[Interior feasibility of value, slope, and opening]
\label{thm:v49-gaussian-witness}
Fix \(\mu>0\) and \(\sigma>0\).  For every sufficiently small pair of
real targets \((r,s)\), there are real \(a,\phi\), depending smoothly on
\((r,s)\), such that the unit phasor
\(\zeta_{a,\phi}\) satisfies
\[
 \boxed{
 \Re\mathbb E\zeta_{a,\phi}=r,}
 \tag{V49.23}
\]
\[
 \boxed{
 \Im\mathbb E[(\mu+\sigma G)\zeta_{a,\phi}]=s.}
 \tag{V49.24}
\]
At \((r,s)=(0,0)\), one may take \((a,\phi)=(0,0)\), and then
\[
 \boxed{
 z=0,
 \qquad
 c=\sqrt{\frac2\pi},
 \qquad
 d=0,}
 \tag{V49.25}
\]
\[
 \boxed{
 \Re\mathbb E[(\mu+\sigma G)^2\zeta_{0,0}]
 =2\mu\sigma\sqrt{\frac2\pi}>0.}
 \tag{V49.26}
\]
The point \textup{(V49.25)} lies strictly inside the limiting Bessel
ellipsoid, with slack
\[
 \boxed{1-\frac2\pi>0.}
 \tag{V49.27}
\]
Consequently the positive opening in \textup{(V49.26)} persists for all
sufficiently small \((r,s)\) and every target \(h=o(\mu\sigma)\).
\end{theorem}

\begin{proof}
Equation \textup{(V49.22)} follows from
\[
 \mathbb E[G\mathbf1_{G>a}]=\varphi(a),
 \qquad
 \mathbb E[G^2\mathbf1_{G>a}]
 =a\varphi(a)+1-\Phi(a).
 \tag{V49.28}
\]
Consider the two real maps
\[
 F_1(a,\phi):=z_0(a)\cos\phi,
 \tag{V49.29}
\]
\[
 F_2(a,\phi):=
 \bigl(\mu z_0(a)+\sigma c_0(a)\bigr)\sin\phi.
 \tag{V49.30}
\]
At \((0,0)\), their Jacobian is diagonal with entries
\[
 \partial_aF_1(0,0)=-2\varphi(0)\ne0,
 \qquad
 \partial_\phi F_2(0,0)=2\sigma\varphi(0)>0.
 \tag{V49.31}
\]
The implicit-function theorem proves \textup{(V49.23)}--
\textup{(V49.24)} for small \((r,s)\).  At the origin,
\(z_0(0)=d_0(0)=0\) and
\(c_0(0)=2\varphi(0)=\sqrt{2/\pi}\), proving
\textup{(V49.25)}--\textup{(V49.26)}.  Substitution in
\textup{(V49.14)} gives \textup{(V49.27)}.  Continuity gives the final
assertion.
\end{proof}

For the saddle parameters,
\[
 \mu_\chi(t)\sigma_\chi(t)
 \asymp t^{-3/2},
 \tag{V49.32}
\]
while \textup{(V49.19)} is exponentially small.  Hence the limiting
three-moment contact slice passes through a region with a polynomially
large positive opening and fixed ellipsoid slack.  There is no asymptotic
moment-body gap to exploit.

\subsection{What remains arithmetic}

The actual prime phases form the one-parameter curve
\[
 \Gamma_{\chi,t}(x):=\bigl(z_t(x),c_t(x),d_t(x)\bigr)
 \tag{V49.33}
\]
inside the body \textup{(V49.10)}.  It satisfies
\[
 z_t'(x)=i\bigl(\mu_\chi z_t+\sigma_\chi c_t\bigr),
 \tag{V49.34}
\]
\[
 z_t''(x)=-\bigl(
 \mu_\chi^2z_t+2\mu_\chi\sigma_\chi c_t
 +\sigma_\chi^2d_t\bigr).
 \tag{V49.35}
\]
Thus the contact slice asks whether the real projection of this specific
arithmetic curve meets a prescribed value with prescribed tangent and
positive curvature.  The Gaussian witness is not claimed to lie on
\(\Gamma_{\chi,t}\).

\begin{theorem}[V49 three-moment verdict]
\label{thm:v49-verdict}
The zeroth--first--second reader problem has the following exact status.
\begin{enumerate}[label=\textup{(\roman*)},leftmargin=3em]
\item Unit modulus and the first four absolute saddle moments place
      \((z,c,d)\) in the Bessel ellipsoid \textup{(V49.10)}.
\item Critical value, slope, and opening cut this body by the affine
      conditions \textup{(V49.16)}--\textup{(V49.18)}.
\item The Gaussian saddle limit contains a unit-phasor witness satisfying
      zero value and slope with strictly positive opening and fixed
      interior slack.
\item Small archimedean targets perturb this witness without closing the
      gap.
\item Any remaining exclusion must distinguish the actual arithmetic
      curve \(\Gamma_{\chi,t}\) from generic unit phasors; fixed-conductor
      geometry alone already reduced that task to progression errors in
      V29.
\end{enumerate}
Consequently the sharp three-moment body supplies no unconditional
contact exclusion and no reduction of the required arithmetic phase
cancellation.
\end{theorem}

\begin{proof}
Combine Theorems~\ref{thm:v49-bessel-body},
\ref{thm:v49-contact-slice}, and
\ref{thm:v49-gaussian-witness} with the arithmetic-curve identities
\textup{(V49.34)}--\textup{(V49.35)} and the V29 conductor projection.
\end{proof}

\begin{firewall}[Moment feasibility is not an arithmetic phase model]
The Gaussian threshold phasor proves that no strict gap follows from the
retained moment data.  It is not the Dirichlet phasor
\(\chi(n)e^{ix\log n}\), and it does not prove that the actual arithmetic
curve reaches every feasible point.  Conversely, showing that the actual
curve avoids the contact slice would be new arithmetic cancellation of
GRH strength.  No off-line zero is excluded and GRH remains open.
\end{firewall}

\begin{assessment}[V49 closes the generic three-moment contact body]
The zero-slope condition removes the imaginary carrier left free in V28,
but the positive opening uses an independent real covariance direction.
The sign-threshold witness occupies that direction with fixed slack, so
adding the second reader does not create generic coercivity.  Higher fixed
moments would repeat the orthogonal-projection hierarchy and retain the
same full Gaussian type.

The next continuation is the scheduled V50 corpus sweep required by the
ten-version protocol.  It should inventory every current parallel note,
search additions since the V40 sweep for a genuinely source-specific
inequality acting on the actual curve \(\Gamma_{\chi,t}\), and compare
that input against the closed branches V41--V49 before choosing a new
upstream direction.  No further generic moment extension should be made
before that audit.
\end{assessment}

\section*{Version V49 append-only record}
\addcontentsline{toc}{section}{Version V49 append-only record}
The complete V48 source was retained before this continuation.  It had
\(742219\) bytes, \(22434\) lines, and SHA--256
\[
 \texttt{4F570737A826453181DF3C02C5FFCAB50EDFBB9EE0B37AC8498CCDB3E89AB1BE}.
\]
V49 derived the exact three-coordinate Bessel body, translated critical
value--slope--opening data into its affine slice, constructed a robust
Gaussian unit-phasor witness with positive opening, and isolated the
remaining difficulty to the actual multiplicative phase curve.

% =============================================================================
% END APPENDED V49 MATERIAL
% =============================================================================
% =============================================================================
% BEGIN APPENDED V50 MATERIAL
% =============================================================================

\section{V50: scheduled corpus sweep and stationary-contact phase collapse}
\label{sec:grh-v50}

V49 closed the generic three-moment route: the sharp Bessel body contains
an open family of value--slope--opening witnesses, so any further exclusion
must act on the actual Dirichlet phase curve.  This is the scheduled
ten-version corpus sweep.  Every non-GRH \texttt{.tex} file present at the
cutoff below was hashed; the two rolling successors were screened through
their newest append-only sections; and every candidate import was tested
against the closed V41--V49 branches.

The sweep supplies no estimate of the signed prime phase.  It does sharpen
one obstruction exactly.  The SM phase-reweighting compiler requires a
uniform lower bound for an average phase.  At a small-time stationary
zero-level contact, the prime average phase is instead forced to collapse:
its modulus is at most \(\sqrt{8t}(1+o(1))\).  Consequently the normalized
complex reweighting has total variation at least
\((8t)^{-1/2}(1+o(1))\), while the centered unit-phasor variance tends to
one.  Thus the phase-floor route does not merely lack a proof in the GRH
problem; its key hypothesis is incompatible with the contact geometry that
an off-line zero produces.

\subsection{Timestamped V50 corpus ledger}

The preceding full sweep, V40, froze \(30\) non-GRH \texttt{.tex} files at
\[
 \texttt{2026-09-11T00:42:04.3520709Z}
\]
with canonical manifest SHA--256
\[
 \texttt{4F42F003B6739E2AABC063B0AD8E818FC189DD06E02FA86F58DCE8B0D151583E}.
 \tag{V50.1}
\]
The V50 cutoff was taken at
\[
 \boxed{\texttt{2026-09-11T02:03:31.7050564Z}.}
 \tag{V50.2}
\]
There were again \(30\) non-GRH \texttt{.tex} files.  For the canonical
UTF--8 manifest with rows
\[
 \texttt{filename|byte-count|SHA--256}\backslash\texttt{n}
\]
in ordinal filename order, the hash was
\[
 \boxed{\texttt{B5E9BC96AF2DE47F4367ED1D2D960944995F1FB327753AFFD899D817454F2D3B}.}
 \tag{V50.3}
\]

The \(27\) archive rows hash-gated at V30 and the SM f14 archive row
separately recorded at V40 remain byte-identical.  Thus exactly \(28\) of
the \(30\) rows need no rereading.  The two rolling V40 endpoints and their
successors at the V50 cutoff are:
\begin{center}
\tiny
\begin{tabular}{p{0.50\textwidth}|r|p{0.31\textwidth}}
\textbf{endpoint}&\textbf{bytes}&\textbf{SHA--256}\\ \hline
\path{Renewal_SM_Einstein_Continuum_Research_Notes_V61.tex}
 &595097&\texttt{CD34919A1695F4E2D4EDEE26E0EEECE}\newline
 \texttt{27A52357415C89E3511964BBF15575CA0}\\
\path{Renewal_SM_Einstein_Continuum_Research_Notes_V90.tex}
 &890431&\texttt{D5C4E3CED26C2D0A3A841DD5BF1BCC}\newline
 \texttt{5539A56EA0F62E15C2B2A54822F2AE1B99}\\
\path{Renewal_Yang_Mills_Mass_Gap_Research_Notes_V39.tex}
 &329898&\texttt{9B44B0FBCA6037F3319E86036753B0F3}\newline
 \texttt{B76BDD06F6045F00747E05DC01E23F9D}\\
\path{Renewal_Yang_Mills_Mass_Gap_Research_Notes_V41.tex}
 &356260&\texttt{0C111A921ABD8DC8970DB8E19412172D}\newline
 \texttt{72469898CC94670F92D12CC9BB8DD498}
\end{tabular}
\end{center}
The SM successor contains every append-only section V62--V90 and the
YM successor contains V40--V41, so the rolling replacement loses no
intervening section.  The cutoff and manifest deliberately describe one
instant of a concurrently moving directory.

\subsection{Typed transfer screen}

The SM delta divides into four potentially relevant families.
V62--V78 develops Schwinger--Dyson integration by parts, uniform
integrability, score-square/Fisher identities, logarithmic Sobolev
concentration, quotient Hessians, and local-to-global entropy compilers.
V79--V85 develops complex phase reweighting, determinant-line and anomaly
ledgers, and reflection-real sewing.  V86--V88 develops common carriers,
nuclear smoothers, uniform heat traces, and trace-norm convergence.  V89
lifts those positive heat factors to free fermionic and conditionally gapped
bosonic Fock spaces; V90 audits YM V41 and NS V26 and proves precise
local-symbol and restricted-input non-import criteria.
The YM delta V40--V41 proves exact rational positivity first on a widened
third-coordinate prism and then on one selected transverse quarter cell;
its own firewall stops before all-octant, full-torus, global covariant, or
continuum conclusions.

\begin{assessment}[V50 typed transfer verdict]
\label{ass:v50-transfer}
None of these results supplies a GRH estimate.
\begin{enumerate}[label=\textup{(\roman*)},leftmargin=3em]
\item The Schwinger--Dyson and logarithmic-Sobolev rows require a positive
      differentiable Gibbs density, an integration-by-parts direction, and
      a Dirichlet-form or Hessian coercivity constant.  The prime saddle is
      a discrete measure and the factor
      \(\chi(n)e^{ix\log n}\) is the oscillatory observable.  No companion
      theorem constructs a derivative on that support which estimates the
      actual curve \(\Gamma_{\chi,t}\).
\item SM V79 proves that normalized complex reweighting is stable only under
      a uniform average-phase floor.  Theorem~\ref{thm:v50-phase-collapse}
      below proves the opposite at stationary zero contacts.  Semicircle,
      small-lift-variance, and phase-floor hypotheses therefore cannot be
      imported without assuming away the contact.
\item SM V83 obtains positivity only after a reflection-real factorization
      into a norm square.  Applied to the character reader, this produces
      the two-prime correlations already isolated in V40--V43; their ratio
      diagonalization reduces to the progression errors of V29.  It gives no
      linear signed constraint on the contact phase.
\item The common-carrier, heat-trace, and Fock nuclearity results concern a
      fixed positive heat time and positive compact operators.  The GRH
      obstruction is a moving \(t\downarrow0\) contact with unsigned saddle
      mass of type \(e^{1/(16t)}\).  Trace ideals and determinant continuity
      are phase-blind and provide neither a common finite-rank tail at that
      scale nor cancellation of the character phase.
\item The YM certificates are finite-dimensional and pointwise in momentum.
      Their exact local atlas does not transport between all orientations
      even inside that programme, and it cannot replace cancellation in an
      infinite signed prime sum.
\end{enumerate}
Accordingly V50 imports no companion conclusion.  It retains the V79
average-phase criterion as a diagnostic and proves that its hypothesis
fails quantitatively on the relevant contact set.
\end{assessment}

\subsection{Stationarity collapses the complex average phase}

Retain the V49 standardized coordinates
\[
 z_t=\mathbb E_t\zeta_{n,x},\qquad
 c_t=\mathbb E_t(X_t\zeta_{n,x}),
 \tag{V50.4}
\]
and write
\[
 z_t=r_t+iy_t,\qquad
 \mu_t:=\mu_\chi(t),\qquad
 \sigma_t:=\sigma_\chi(t).
 \tag{V50.5}
\]
At a stationary zero-level contact, \(m_{\chi,t}(x)=0\) and
\(\partial_xm_{\chi,t}(x)=0\), equations \textup{(V49.16)}--
\textup{(V49.17)} give
\[
 r_t=\frac{a_{\chi,t}(x)}{2Z_\chi(t)},
 \qquad
 \Im(\mu_tz_t+\sigma_tc_t)
 =s_t:=-\frac{\partial_xa_{\chi,t}(x)}{2Z_\chi(t)}.
 \tag{V50.6}
\]

\begin{theorem}[Stationary-contact phase collapse]
\label{thm:v50-phase-collapse}
For every sufficiently small \(t>0\) and every stationary zero-level
contact,
\[
 \boxed{
 |z_t|^2\le r_t^2+
 \left(
  \frac{|s_t|+\sigma_t\sqrt{1-r_t^2}}{\mu_t}
 \right)^2.}
 \tag{V50.7}
\]
Consequently, along any sequence of such contacts with \(t\downarrow0\)
and \(x\) in a fixed compact interval,
\[
 \boxed{|z_t|^2\le8t\bigl(1+o_{\chi,I}(1)\bigr),
 \qquad
 |z_t|\le2\sqrt{2t}\bigl(1+o_{\chi,I}(1)\bigr).}
 \tag{V50.8}
\]
Moreover the exact centered phase variance satisfies
\[
 \boxed{
 \mathbb E_t|\zeta_{n,x}-z_t|^2
 =1-|z_t|^2
 \ge1-8t\bigl(1+o_{\chi,I}(1)\bigr).}
 \tag{V50.9}
\]
Thus the prime phasor is asymptotically maximally dispersed at a stationary
contact even though its mean is asymptotically zero.
\end{theorem}

\begin{proof}
The first two coordinates of the V49 Bessel body give
\[
 |c_t|^2\le1-|z_t|^2=1-r_t^2-y_t^2.
 \tag{V50.10}
\]
Taking imaginary parts in \textup{(V50.6)} and using
\textup{(V50.10)} yields
\[
 \mu_t|y_t|
 \le |s_t|+\sigma_t|\Im c_t|
 \le |s_t|+\sigma_t\sqrt{1-r_t^2-y_t^2}
 \le |s_t|+\sigma_t\sqrt{1-r_t^2}.
 \tag{V50.11}
\]
The saddle mean is positive for small \(t\), so squaring this inequality and
adding \(r_t^2\) proves \textup{(V50.7)}.

On a fixed compact interval, \textup{(V49.19)} makes \(r_t\) and \(s_t\)
exponentially small.  Proposition~\ref{prop:v28-saddle-moments} gives
\[
 \mu_t=\frac1{4t}+o_\chi(t^{-1/2}),\qquad
 \sigma_t^2=\frac1{2t}(1+o_\chi(1)),
 \qquad
 \frac{\sigma_t^2}{\mu_t^2}=8t(1+o_\chi(1)).
 \tag{V50.12}
\]
Substitution proves \textup{(V50.8)}.  Finally
\(\mathbb E_t|\zeta-z_t|^2=1-|z_t|^2\) is the exact identity
\textup{(V28.16)}, and \textup{(V50.9)} follows.
\end{proof}

The estimate is stronger than the zero-level value equation alone.  That
equation controls only \(\Re z_t\); stationarity couples the formerly free
imaginary carrier to the centered first moment, and the Bessel bound costs
exactly the saddle ratio \(\sigma_t/\mu_t\asymp\sqrt t\).  This consequence
was not used in V28's optimal unrestricted-slope envelope.

\subsection{The exact reweighting obstruction}

The complex phase-reweighting formalism normalizes the signed measure by its
mean phase.  Here, whenever \(z_t(x)\ne0\), define
\[
 d\rho_{t,x}:=\frac{\zeta_{n,x}}{z_t(x)}\,d\nu_{\chi,t}.
 \tag{V50.13}
\]
Then \(\rho_{t,x}(\Omega)=1\), but it is not a positive probability measure.

\begin{corollary}[Polynomial blow-up of normalized phase reweighting]
\label{cor:v50-reweighting}
At every stationary zero-level contact with \(z_t(x)\ne0\),
\[
 \boxed{
 \|\rho_{t,x}\|_{\mathrm{TV}}
 =\frac1{|z_t(x)|}.}
 \tag{V50.14}
\]
Along compact small-time contact sequences,
\[
 \boxed{
 \|\rho_{t,x}\|_{\mathrm{TV}}
 \ge\frac1{\sqrt{8t}}\bigl(1+o_{\chi,I}(1)\bigr).}
 \tag{V50.15}
\]
If \(z_t(x)=0\), the normalization \textup{(V50.13)} is undefined.  Hence
no uniform positive average-phase floor and no uniformly bounded normalized
complex reweighting can hold on the stationary contact set.
\end{corollary}

\begin{proof}
Because \(|\zeta_{n,x}|=1\), the total variation of
\textup{(V50.13)} is
\[
 \int\left|\frac{\zeta_{n,x}}{z_t(x)}\right|d\nu_{\chi,t}
 =\frac1{|z_t(x)|}.
 \tag{V50.16}
\]
Theorem~\ref{thm:v50-phase-collapse} gives \textup{(V50.15)} and the stated
dichotomy.
\end{proof}

This does not contradict the V49 feasibility theorem.  Its central Gaussian
witness has \(z=0\), exactly the maximally ill-conditioned endpoint of phase
reweighting, while retaining a strictly positive opening.  Nor is the
polynomial loss in \textup{(V50.15)} an exponential GRH contradiction: on
the normalized heat scale it has rate zero.  The new result closes a proposed
transfer mechanism, not the arithmetic phase problem.

\begin{theorem}[V50 sweep and phase-floor verdict]
\label{thm:v50-verdict}
The additions to the parallel corpus since V40 provide no theorem acting on
the signed contact curve \(\Gamma_{\chi,t}\).  Among their reusable
diagnostics, the strongest candidate is the SM average-phase floor.  On the
actual prime probability space, however, every compact small-time stationary
zero-contact sequence obeys
\[
 |\mathbb E_t\zeta_{n,x}|=O_{\chi,I}(\sqrt t),
 \qquad
 \mathbb E_t|\zeta_{n,x}-\mathbb E_t\zeta_{n,x}|^2=1-O_{\chi,I}(t),
 \tag{V50.17}
\]
and normalized phase reweighting is either undefined or has diverging total
variation.  Reflection squaring returns the already closed two-prime branch;
positive trace ideals, Fock determinants, and finite Wilson atlases remain
phase-blind.  Therefore no companion result excludes a contact or improves
the zero rate of the remaining arithmetic cancellation problem.
\end{theorem}

\begin{proof}
The corpus assertion is the hash-gated transfer screen above.  The
quantitative assertions are Theorem~\ref{thm:v50-phase-collapse} and
Corollary~\ref{cor:v50-reweighting}.  The branch identifications are those of
V29 and V40--V49.
\end{proof}

\begin{firewall}[Phase collapse is not phase cancellation]
V50 performs the full scheduled corpus sweep and proves that stationarity
forces the normalized complex prime phasor to collapse at least on the
\(\sqrt t\) scale.  This rules out a uniform average-phase floor and a
uniformly bounded normalized complex reweighting at contacts.  It does not
bound the unnormalized signed prime reader, prevent the V49 positive-opening
witness, exclude an off-line zero, or prove GRH.
\end{firewall}

\begin{assessment}[V50 acquisition and next upstream route]
The V41--V49 branches and the present sweep now close three broad detours:
quadratic reflection positivity returns unresolved two-prime correlations;
smooth or nuclear heat compilers control positive fixed-time operators but
not the moving signed source; and normalized complex reweighting becomes
singular precisely on stationary contacts.

The next continuation should therefore stay unnormalized and
source-specific.  V51 should study the actual multiplicative curve under the
simultaneous constraints
\[
 \Re z_t=r_t,qquad
 \Im(\mu_tz_t+\sigma_tc_t)=s_t,qquad
 |z_t|=O(\sqrt t),
 \tag{V50.18}
\]
and ask whether the character progression structure forces any deficit from
the V49 threshold witness in the \emph{real} first covariance coordinate
that carries the positive opening.  A useful result must distinguish
\(\chi(n)e^{ix\log n}\) from an arbitrary unit phasor before taking a norm.
If progression projection again leaves only the known error term, the branch
should be recorded as circular and the search moved further upstream rather
than extended to higher generic moments.
\end{assessment}

\section*{Version V50 append-only record}
\addcontentsline{toc}{section}{Version V50 append-only record}
The complete V49 source was retained before this continuation.  It had
\(755438\) bytes, \(22877\) lines, and SHA--256
\[
 \texttt{0D50E3C411030745D3E64AE885AF979F7AB093C84EED1FC5E546FD3FB6FCE90D}.
\]
V50 performed the scheduled full corpus sweep, proved stationary-contact
collapse of the complex mean prime phasor, identified asymptotically maximal
phase dispersion, and quantified the resulting total-variation blow-up of
normalized complex phase reweighting.

% =============================================================================
% END APPENDED V50 MATERIAL
% =============================================================================
% =============================================================================
% BEGIN APPENDED V51 MATERIAL
% =============================================================================

\section{V51: progression-scale collapse and the shrinking-witness no-go}
\label{sec:grh-v51}

V50 asked whether the actual progression structure forces a deficit from
V49's constant Gaussian threshold witness in the real first covariance
coordinate.  It does: on every compact ordinate interval, the entire actual
standardized triple \((z_t,c_t,d_t)\) tends to zero faster than every power
of \(t\).  This follows by centering the exact V29 residue-error projection,
not from the generic Bessel body.

The gain is rigorous but not yet coercive.  Its available quantitative form
is only \(t^{-M}e^{-b/\sqrt t}\), which has zero resolution on the decisive
\(1/t\) scale.  Two exact Gaussian constructions show why the shrinking
magnitude alone cannot determine the opening sign.  A rotated plane wave has
zero first-reader imaginary part, all three moments at an arbitrarily small
Gaussian Fourier scale, and strictly positive second-reader opening.  A
diluted V49 threshold witness additionally meets the exact exponentially
smaller archimedean value and slope targets while remaining inside any fixed
\(e^{-B/\sqrt t}\) envelope.  Therefore the proposed progression route
closes at a precise boundary: a useful next estimate must control the signed
direction of the character-projected error, not merely prove that its moment
vector collapses.

\subsection{The centered progression-error projection}

For \(a\in G_q\), retain the V29 progression errors
\(E_{a,t}^{[r]}(x)\), and write
\[
 \mathbf E_t^{[r]}(x)
 :=\bigl(E_{a,t}^{[r]}(x)\bigr)_{a\in G_q},
 \qquad
 \boldsymbol\chi:=\bigl(\chi(a)\bigr)_{a\in G_q}.
 \tag{V51.1}
\]
Define their first two saddle-centered combinations by
\[
 \widehat{\mathbf E}_t^{[1]}
 :=\mathbf E_t^{[1]}-\mu_t\mathbf E_t^{[0]},
 \tag{V51.2}
\]
\[
 \widehat{\mathbf E}_t^{[2]}
 :=\mathbf E_t^{[2]}-2\mu_t\mathbf E_t^{[1]}
   +\mu_t^2\mathbf E_t^{[0]}.
 \tag{V51.3}
\]

\begin{proposition}[Exact centered residue projection]
\label{prop:v51-centered-projection}
For every \(t>0\) and real \(x\),
\[
 \boxed{
 z_t(x)=\frac{\boldsymbol\chi\cdot\mathbf E_t^{[0]}(x)}{Z_\chi(t)},}
 \tag{V51.4}
\]
\[
 \boxed{
 c_t(x)=
 \frac{\boldsymbol\chi\cdot\widehat{\mathbf E}_t^{[1]}(x)}
      {\sigma_tZ_\chi(t)},}
 \tag{V51.5}
\]
\[
 \boxed{
 d_t(x)=
 \frac{\boldsymbol\chi\cdot\widehat{\mathbf E}_t^{[2]}(x)}
      {\sigma_t^2Z_\chi(t)}.}
 \tag{V51.6}
\]
In particular, the real covariance coordinate proposed in V50 is exactly
\[
 \boxed{
 \Re c_t(x)=\frac1{\sigma_tZ_\chi(t)}
 \Re\!\left[
  \boldsymbol\chi\cdot
  \bigl(\mathbf E_t^{[1]}(x)-\mu_t\mathbf E_t^{[0]}(x)\bigr)
 \right].}
 \tag{V51.7}
\]
It is a centered version of the same single nonprincipal character mode as
V29, not a coupling to a new residue mode.
\end{proposition}

\begin{proof}
Theorem~\ref{thm:v29-residue-projection} gives
\[
 P_{\chi,t}^{[r]}(x)
 =\boldsymbol\chi\cdot\mathbf E_t^{[r]}(x)
 \qquad(r=0,1,2).
 \tag{V51.8}
\]
The definitions \textup{(V49.1)}--\textup{(V49.5)} give
\[
 z_t=\frac{P_{\chi,t}^{[0]}}{Z_\chi},\qquad
 c_t=\frac{P_{\chi,t}^{[1]}-\mu_tP_{\chi,t}^{[0]}}
            {\sigma_tZ_\chi},
 \tag{V51.9}
\]
\[
 d_t=\frac{P_{\chi,t}^{[2]}-2\mu_tP_{\chi,t}^{[1]}
                    +\mu_t^2P_{\chi,t}^{[0]}}
            {\sigma_t^2Z_\chi}.
 \tag{V51.10}
\]
Substitution proves all four identities.
\end{proof}

The apparent covariance contribution to the opening also recombines without
creating a new arithmetic quantity:
\[
 \boxed{
 2\mu_t\sigma_tc_t+\sigma_t^2d_t
 =\frac{P_{\chi,t}^{[2]}}{Z_\chi(t)}-\mu_t^2z_t.}
 \tag{V51.11}
\]
Thus the complete expression in \textup{(V49.18)} is simply the raw second
reader \(P_{\chi,t}^{[2]}/Z_\chi\), as it must be.  Separating \(c_t\) from
\(d_t\) exposes moment geometry, but does not manufacture a second
character projection.

\subsection{The actual arithmetic curve collapses superpolynomially}

\begin{theorem}[Uniform progression-scale collapse]
\label{thm:v51-actual-collapse}
For every compact interval \(I\subset\mathbb R\), there are constants
\(C,M,b,t_0>0\), depending on \(\chi\) and \(I\), such that
\[
 \boxed{
 \sup_{x\in I}
 \bigl(|z_t(x)|+|c_t(x)|+|d_t(x)|\bigr)
 \le Ct^{-M}\exp\!\left(-\frac b{\sqrt t}\right)}
 \tag{V51.12}
\]
for \(0<t<t_0\).  Equivalently, for every \(A>0\),
\[
 \boxed{
 \sup_{x\in I}
 \bigl(|z_t(x)|+|c_t(x)|+|d_t(x)|\bigr)
 =O_{\chi,I,A}(t^A).}
 \tag{V51.13}
\]
Consequently the actual arithmetic curve \(\Gamma_{\chi,t}(I)\) is uniformly
separated from V49's constant threshold point:
\[
 \boxed{
 \inf_{x\in I}
 \left|c_t(x)-\sqrt{\frac2\pi}\right|
 \longrightarrow\sqrt{\frac2\pi}.}
 \tag{V51.14}
\]
\end{theorem}

\begin{proof}
Apply Proposition~\ref{prop:v29-ap-gaussian} with \(r=0,1,2\), reduce its
three positive saving constants to one \(b_0>0\), and enlarge one polynomial
power \(M_0\).  The finite character projection then gives
\[
 \sup_{x\in I}|P_{\chi,t}^{[r]}(x)|
 \le C_0t^{-M_0}
 \exp\!\left(\frac1{16t}-\frac{b_0}{\sqrt t}\right)
 \qquad(r=0,1,2).
 \tag{V51.15}
\]
By Proposition~\ref{prop:v28-saddle-moments},
\[
 Z_\chi(t)\asymp_\chi t^{-1/2}e^{1/(16t)},
 \qquad
 \mu_t\asymp t^{-1},
 \qquad
 \sigma_t\asymp t^{-1/2}.
 \tag{V51.16}
\]
Insert \textup{(V51.15)}--\textup{(V51.16)} into
\textup{(V51.9)}--\textup{(V51.10)}.  Every centering and standardization
factor costs only a fixed power of \(t^{-1}\).  After enlarging \(M\) and,
if necessary, decreasing \(b_0\) to \(b\), this proves
\textup{(V51.12)}.  The exponential \(e^{-b/\sqrt t}\) dominates every
fixed power of \(t\), proving \textup{(V51.13)}.  Finally \(c_t\to0\)
uniformly, and the reverse triangle inequality gives \textup{(V51.14)}.
\end{proof}

This is the source-specific deficit sought in V50.  It uses the actual
Dirichlet progression error and rules out the particular constant-size
Gaussian witness.  It still has normalized logarithmic rate zero:
\[
 \lim_{t\downarrow0}4t\log
 \left(1+t^{-M}e^{-b/\sqrt t}\right)=0.
 \tag{V51.17}
\]
The issue is whether a shrinking vector can approach the contact slice from
the positive-opening side.

\subsection{A deterministic shrinking plane-wave witness}

Let \(G\) be standard normal.  Fix \(\mu,\sigma>0\), \(\lambda\ge0\), and
put
\[
 R_\lambda:=\sqrt{\mu^2+\sigma^2\lambda^2},
 \qquad
 \phi_\lambda:=-\arctan\frac{\sigma\lambda}{\mu},
 \qquad
 \zeta_\lambda(G):=e^{i(\lambda G+\phi_\lambda)}.
 \tag{V51.18}
\]

\begin{proposition}[Rotated Gaussian plane wave]
\label{prop:v51-plane-wave}
The first three standardized moments of \(\zeta_\lambda\) are
\[
 \boxed{
 z_\lambda=e^{-\lambda^2/2}e^{i\phi_\lambda},\qquad
 c_\lambda=i\lambda z_\lambda,\qquad
 d_\lambda=(1-\lambda^2)z_\lambda.}
 \tag{V51.19}
\]
They satisfy the exact stationary-slope identity
\[
 \boxed{
 \Im(\mu z_\lambda+\sigma c_\lambda)=0,}
 \tag{V51.20}
\]
the positive value identity
\[
 \boxed{
 r_\lambda:=\Re z_\lambda
 =e^{-\lambda^2/2}\frac\mu{R_\lambda}>0,}
 \tag{V51.21}
\]
and the strictly positive opening
\[
 \boxed{
 \Re\bigl(\mu^2z_\lambda+2\mu\sigma c_\lambda
                  +\sigma^2d_\lambda\bigr)
 =r_\lambda\bigl(\mu^2+\sigma^2(1+\lambda^2)\bigr)>0.}
 \tag{V51.22}
\]
For saddle parameters \(\mu=\mu_t\), \(\sigma=\sigma_t\) and
\(\lambda_t=\kappa t^{-1/4}\), \(\kappa>0\),
\[
 |z_{\lambda_t}|+|c_{\lambda_t}|+|d_{\lambda_t}|
 \le C_\kappa t^{-1/2}
 \exp\!\left(-\frac{\kappa^2}{2\sqrt t}\right),
 \tag{V51.23}
\]
while \textup{(V51.20)} and the positive sign in
\textup{(V51.22)} remain exact.
\end{proposition}

\begin{proof}
The standard Gaussian characteristic function and its first two derivatives
give
\[
 \mathbb E e^{i\lambda G}=e^{-\lambda^2/2},\quad
 \mathbb E(Ge^{i\lambda G})=i\lambda e^{-\lambda^2/2},\quad
 \mathbb E(G^2e^{i\lambda G})=(1-\lambda^2)e^{-\lambda^2/2}.
 \tag{V51.24}
\]
This proves \textup{(V51.19)}.  Since
\[
 e^{i\phi_\lambda}(\mu+i\sigma\lambda)=R_\lambda,
 \tag{V51.25}
\]
the first reader is positive real, proving \textup{(V51.20)}--
\textup{(V51.21)}.  Multiplication by
\(e^{i\phi_\lambda}=(\mu-i\sigma\lambda)/R_\lambda\) gives
\[
 \begin{aligned}
 &\Re\left[e^{i\phi_\lambda}
  \bigl(\mu^2+\sigma^2(1-\lambda^2)
                     +2i\mu\sigma\lambda\bigr)\right]\\
 &\hspace{3em}=
 \frac\mu{R_\lambda}
 \bigl(\mu^2+\sigma^2(1+\lambda^2)\bigr),
 \end{aligned}
 \tag{V51.26}
\]
which proves \textup{(V51.22)} after multiplication by
\(e^{-\lambda^2/2}\).  Finally \textup{(V51.19)} and
\(1+\lambda_t+\lambda_t^2\ll_\kappa t^{-1/2}\) prove
\textup{(V51.23)}.
\end{proof}

The witness even has the same one-frequency form as a Gaussian
characteristic function.  It is not the nonprincipal Dirichlet phasor: the
periodic character twist is exactly what removes the continuous main term
and leaves the progression error.  Its role is narrower but decisive:
superpolynomial moment collapse, one-parameter Fourier kinematics, and
stationarity do not imply a nonpositive opening.

\subsection{Exact archimedean targets survive inside a shrinking body}

The preceding plane wave has a positive value target.  The following
construction handles both signs and the exact small value and slope targets.
It makes explicit why adding the bound \textup{(V51.12)} to the generic
moment data still supplies no contact exclusion.

\begin{theorem}[Diluted shrinking contact witness]
\label{thm:v51-diluted-witness}
Let \(\mu_t,\sigma_t\) have the saddle asymptotics
\textup{(V50.12)}, and let real targets \(r_t,s_t,h_t\) satisfy
\[
 |r_t|+|s_t|+|h_t|
 \le\exp\!\left(-\frac1{16t}+O(\log(1/t))\right).
 \tag{V51.27}
\]
For every fixed \(B>0\), put \(\epsilon_t=e^{-B/\sqrt t}\).  For all
sufficiently small \(t\), there is a probability space carrying a standard
normal variable \(G\) and a unit phasor \(\zeta_t\) such that, with
\[
 z_t^*=\mathbb E\zeta_t,qquad
 c_t^*=\mathbb E(G\zeta_t),qquad
 d_t^*=\mathbb E(G^2\zeta_t),
 \tag{V51.28}
\]
one has
\[
 \boxed{\Re z_t^*=r_t,}
 \tag{V51.29}
\]
\[
 \boxed{\Im(\mu_tz_t^*+\sigma_tc_t^*)=s_t,}
 \tag{V51.30}
\]
\[
 \boxed{
 \Re(\mu_t^2z_t^*+2\mu_t\sigma_tc_t^*+\sigma_t^2d_t^*)>h_t,}
 \tag{V51.31}
\]
and
\[
 \boxed{|z_t^*|+|c_t^*|+|d_t^*|\le C\epsilon_t}
 \tag{V51.32}
\]
for an absolute constant \(C\).  Thus the exact contact slice with positive
opening meets a moment body shrinking faster than every power of \(t\).
\end{theorem}

\begin{proof}
Divide the first two targets by \(\epsilon_t\).  Condition
\textup{(V51.27)} implies
\[
 \frac{|r_t|+|s_t|}{\epsilon_t}=o(t^A)
 \qquad\hbox{for every fixed }A>0.
 \tag{V51.33}
\]
Use the two-parameter threshold phasor \(\eta_t(G)\) from
Theorem~\ref{thm:v49-gaussian-witness} with targets
\((r_t/\epsilon_t,s_t/\epsilon_t)\).  The quantitative implicit-function
argument remains valid for the varying saddle parameters: its Jacobian at
the origin has diagonal entries
\(-2\varphi(0)\) and \(2\sigma_t\varphi(0)\), while all parameter growth is
polynomial in \(t^{-1}\); the targets in \textup{(V51.33)} are smaller than
every such polynomial scale.  Hence
\[
 \Re\mathbb E\eta_t=\frac{r_t}{\epsilon_t},qquad
 \Im\mathbb E[(\mu_t+\sigma_tG)\eta_t]
 =\frac{s_t}{\epsilon_t},
 \tag{V51.34}
\]
and, by continuity from \textup{(V49.26)},
\[
 \Re\mathbb E[(\mu_t+\sigma_tG)^2\eta_t]
 =2\mu_t\sigma_t\sqrt{\frac2\pi},(1+o(1)).
 \tag{V51.35}
\]

On an independent factor, let \(J_t\) be Bernoulli with mean
\(\epsilon_t\), and let \(U\) be uniform on the unit circle, independent of
\((G,J_t)\).  Set
\[
 \zeta_t:=
 \begin{cases}
  \eta_t(G),&J_t=1,\\
  U,&J_t=0.
 \end{cases}
 \tag{V51.36}
\]
Then \(|\zeta_t|=1\), and independence with \(\mathbb EU=0\) gives
\[
 \mathbb E(G^j\zeta_t)
 =\epsilon_t\mathbb E(G^j\eta_t)
 \qquad(j=0,1,2).
 \tag{V51.37}
\]
Equations \textup{(V51.29)}--\textup{(V51.30)} follow from
\textup{(V51.34)}.  Since the first three absolute Gaussian moments are
fixed, \textup{(V51.37)} gives \textup{(V51.32)} after enlarging \(C\).
Finally the left side of \textup{(V51.31)} equals
\[
 2\epsilon_t\mu_t\sigma_t\sqrt{\frac2\pi},(1+o(1)).
 \tag{V51.38}
\]
This dominates \(h_t\), because \(\mu_t\sigma_t\asymp t^{-3/2}\) and
\(e^{-1/(16t)+O(\log(1/t))}=o(e^{-B/\sqrt t}t^{-3/2})\).
\end{proof}

The auxiliary Bernoulli and circle variables are a capacity device, not a
model of primes.  They preserve the exact Gaussian law of \(G\), unit
modulus, all three contact equations, and an arbitrarily prescribed
superpolynomial collapse scale.  Hence any proposed contradiction using only
those data would also reject this explicit witness and is false.

\begin{theorem}[V51 progression-deficit verdict]
\label{thm:v51-verdict}
The fixed-conductor progression structure yields the following exact status.
\begin{enumerate}[label=\textup{(\roman*)},leftmargin=3em]
\item The actual arithmetic moment curve obeys the uniform
      superpolynomial collapse \textup{(V51.12)} and therefore avoids V49's
      constant-size threshold witness.
\item Its real first covariance is exactly the centered target-character
      projection \textup{(V51.7)}; no principal or auxiliary residue mode
      enters.
\item The classical progression saving has no nonzero \(1/t\)-rate and
      supplies no sign for that projection.
\item Plane-wave and diluted threshold phasors retain stationary slope and
      positive opening inside moment bodies with the same or faster
      superpolynomial collapse.
\end{enumerate}
Therefore the magnitude deficit furnished by PNT in progressions cannot
exclude a stationary contact.  Any successful continuation must estimate a
signed projective direction of the actual character error, which is
information beyond fixed-conductor orthogonality and componentwise PNT.
\end{theorem}

\begin{proof}
Items (i)--(ii) are Proposition~\ref{prop:v51-centered-projection} and
Theorem~\ref{thm:v51-actual-collapse}.  The rate statement is
\textup{(V51.17)}.  Proposition~\ref{prop:v51-plane-wave} and
Theorem~\ref{thm:v51-diluted-witness} prove item (iv).  Together they show
that no universal sign conclusion follows from the listed data.
\end{proof}

\begin{firewall}[A shrinking arithmetic curve need not have the right sign]
V51 proves a genuinely source-specific superpolynomial collapse of the
actual standardized character moments and rules out the constant V49
threshold point as an approximation to that curve.  It also proves that
collapse, stationarity, and unit modulus remain compatible with positive
opening at the exact archimedean target scale.  It does not determine the
sign of the actual progression-error direction, exclude an off-line zero,
or prove GRH.
\end{firewall}

\begin{assessment}[V51 acquisition and the projective upstream move]
The V50 question is now answered without ambiguity.  Progression structure
does force a deficit in \(\Re c_t\), but only as part of the simultaneous
collapse of the same V29 target mode.  Centering, standardization, and
separation into \((z,c,d)\) do not create an independent arithmetic row, and
the collapse scale admits exact positive-opening capacity witnesses.

V52 should stop estimating the norm of \(\Gamma_{\chi,t}\) and instead
projectivize it along the centered exposed-zero sequences already constructed
in V23--V27.  The next exact calculation is to insert the dominant Gaussian
zero packet before taking absolute values, compute the limiting ratios among
the normalized zeroth, centered-first, and centered-second readers, and test
whether stationary contact selects a signed ray incompatible with positive
opening.  This is upstream of PNT magnitude bounds and genuinely
source-specific.  If the ray is again freely phase-rotatable or reconstructs
the V24 max-plus zero signal, that loop should be closed explicitly rather
than hidden inside another covariance inequality.
\end{assessment}

\section*{Version V51 append-only record}
\addcontentsline{toc}{section}{Version V51 append-only record}
The complete V50 source was retained before this continuation.  It had
\(770054\) bytes, \(23233\) lines, and SHA--256
\[
 \texttt{17DD4E4447F15F2C9A4ADBB3D5D92572AA947FF5F804D43780D313624655D17D}.
\]
V51 derived the exact centered progression-error projection, proved uniform
superpolynomial collapse of the actual three-moment curve, constructed exact
shrinking positive-opening witnesses, and localized the remaining problem to
the signed projective direction of the target character error.

% =============================================================================
% END APPENDED V51 MATERIAL
% =============================================================================
% =============================================================================
% BEGIN APPENDED V52 MATERIAL
% =============================================================================

\section{V52: analytic extraction of the exposed-zero projective ray}
\label{sec:grh-v52}

V51 proved that the actual standardized moment curve collapses
superpolynomially, and proposed projectivizing it before taking another norm.
The scalar zero expansion of V22 appears at first to provide only alternating
real projections of the complex prime readers.  There is, however, a
canonical way to recover the missing quadratures: the prime reader contains
only positive logarithmic frequencies.  Applying the positive-frequency
projection to one reflected-pair Gaussian extracts one complex Gaussian wave;
the discarded Fourier leakage is only polynomial, while a center-exposed
off-line wave is exponentially large before normalization.

This gives an exact asymptotic projective signature at the V37 simple
contacts.  If the exposed displacement is \(0<\delta_*<1/2\), then
\[
 \left(1,\sqrt t\,\frac{c_t}{z_t},
              t\,\frac{d_t}{z_t}\right)
 \longrightarrow
 \left(1,
       \frac{2\delta_*-1}{2\sqrt2},
       \frac{(2\delta_*-1)^2}{8}\right),
\]
and the endpoint phases of \(z_t\) tend to \(-i\) and \(+i\).  Thus the
complex moment ray reconstructs the horizontal zero displacement.  It does
not exclude it: the ray is generated by the assumed zero, approaches the
interior origin of the V49 Bessel body, and occurs at simple boundary
contacts rather than at the stationary positive-opening contacts of V49--V51.

There is also a typographical correction to the V50 handoff display.  In
\textup{(V50.18)}, the literal word \texttt{qquad} is the spacing command
\(\backslash\texttt{qquad}\); the three mathematical conditions themselves
are unchanged.

\subsection{Positive-frequency recovery of the complex reader}

Let \(\Pi_+\) denote the Fourier projection which retains frequencies
\(u>0\) in the convention \(e^{ixu}\), and put
\[
 D:=-i\partial_x.
 \tag{V52.1}
\]
The absolutely convergent prime readers satisfy
\[
 P_{\chi,t}^{[r]}=D^rP_{\chi,t}^{[0]}.
 \tag{V52.2}
\]
Because \(a_{\chi,t}-m_{\chi,t}=2\Re P_{\chi,t}^{[0]}\) and every prime
frequency is \(\log n>0\), one has the exact distributional identity
\[
 \boxed{
 P_{\chi,t}^{[0]}=\Pi_+(a_{\chi,t}-m_{\chi,t}).}
 \tag{V52.3}
\]

For a reflected pair \(p=(\gamma,\delta,m_p)\), \(\delta>0\), retain its
V25 contribution
\[
 K_{p,t}(x)=
 \frac{4\pi m_p}{\sqrt{4\pi t}}
 \exp\!\left(\frac{\delta^2-(x-\gamma)^2}{4t}\right)
 \cos\!\left(\frac{\delta(x-\gamma)}{2t}\right).
 \tag{V52.4}
\]
Define its outgoing complex half by
\[
 W_{p,t}(x):=
 -\frac{2\pi m_p}{\sqrt{4\pi t}}
 \exp\!\left(\frac{\delta^2-(x-\gamma)^2}{4t}\right)
 \exp\!\left(i\frac{\delta(x-\gamma)}{2t}\right).
 \tag{V52.5}
\]
The minus sign is present because the pair occurs in \(m\), whereas
\(P^{[0]}=\Pi_+(a-m)\).

\begin{lemma}[One-pair analytic extraction]
\label{lem:v52-analytic-half}
For every fixed integer \(r\ge0\),
\[
 \boxed{
 D^r\Pi_+(-K_{p,t})=D^rW_{p,t}+O_{p,r}(t^{-M_r})}
 \tag{V52.6}
\]
uniformly for real \(x\), for some \(M_r<\infty\).  In particular, on
\(|x-\gamma|=O(t)\),
\[
 D^r\Pi_+(-K_{p,t})
 =D^rW_{p,t}
  +O_{p,r}\!\left(
    t^{-M_r}e^{-\delta^2/(4t)}|W_{p,t}(x)|
  \right).
 \tag{V52.7}
\]
\end{lemma}

\begin{proof}
Put \(v=x-\gamma\) and \(\omega=\delta/(2t)\).  Splitting the cosine gives
\[
 H_t(v)\cos(\omega v)
 =\frac12H_t(v)e^{i\omega v}
  +\frac12H_t(v)e^{-i\omega v}.
 \tag{V52.8}
\]
The Fourier densities of the two terms are Gaussians centered at
\(+\omega\) and \(-\omega\), respectively.  The difference between the
positive-frequency projection and the first full Gaussian wave consists of
the negative tail of the \(+\omega\) Gaussian and the positive tail of the
\(-\omega\) Gaussian.  In both cases the center is
\[
 \omega\sqrt t=\frac\delta{2\sqrt t}
 \tag{V52.9}
\]
standard Fourier widths from the origin.  Gaussian tail bounds, with a
factor \(|u|^r\) inserted after applying \(D^r\), give
\[
 \int_{\text{wrong half-line}}|u|^r
       e^{-t(u\mp\omega)^2},du
 \le C_{p,r}t^{-M_r}e^{-\delta^2/(4t)}.
 \tag{V52.10}
\]
The prefactor \(e^{\delta^2/(4t)}\) in \(K_{p,t}\) cancels this exponential,
proving \textup{(V52.6)}.  On \(|v|=O(t)\),
\(|W_{p,t}|\asymp_p t^{-1/2}e^{\delta^2/(4t)}\); division by this amplitude
gives \textup{(V52.7)} after changing \(M_r\).
\end{proof}

The archimedean positive-frequency part is harmless here.  The Binet
integral for the digamma function in \textup{(V19.9)}, followed by the heat
multiplier \(e^{-tu^2}\), gives for every compact \(I\) and fixed \(r\)
\[
 \sup_{x\in I}|D^r\Pi_+a_{\chi,t}(x)|
 \le C_{\chi,I,r}t^{-M_r}.
 \tag{V52.11}
\]
Only polynomial control is needed, because an exposed off-line half-wave has
amplitude \(e^{\delta_*^2/(4t)}\) up to powers of \(t\).

\subsection{Complex dominance of a center-exposed pair}

Let \(p_*=(\gamma_*,\delta_*,m_*)\) be the center-exposed pair of
Theorem~\ref{thm:v37-centered-exposure}, and let
\(t_k\downarrow0\), \(x_{k,-}<x_{k,+}\) be the adjacent simple contacts of
Corollary~\ref{cor:v37-captured-contacts}.  Thus
\[
 x_{k,\pm}=\gamma_*+\frac{(2\pm1)\pi t_k}{\delta_*}
 +O\!\left(t_ke^{-c/t_k}\right).
 \tag{V52.12}
\]

\begin{theorem}[Exposed complex prime wave]
\label{thm:v52-complex-wave}
There are \(c>0\) and \(M<\infty\) such that, for
\(r=0,1,2\),
\[
 \boxed{
 P_{\chi,t_k}^{[r]}(x_{k,\pm})
 =D^rW_{p_*,t_k}(x_{k,\pm})
  +O\!\left(t_k^{-M}e^{-c/t_k}
                 |W_{p_*,t_k}(x_{k,\pm})|\right).}
 \tag{V52.13}
\]
Consequently the scalar dominance of V27 upgrades to dominance of the full
complex prime reader and its first two logarithmic moments.
\end{theorem}

\begin{proof}
Apply \(D^r\Pi_+\) to the zero-side expansion
\(m_{\chi,t}=\sum_pK_{p,t}\), interpreted with the critical-line Gaussian
as the \(\delta_p=0\) case.  Lemma~\ref{lem:v52-analytic-half} extracts the
main complex wave of every off-line pair.  At
\(x=\gamma_*+O(t)\), the center-exposure gap
\textup{(V37.5)} makes every other main half-wave smaller than the
\(p_*\)-wave by \(O(t^{-M}e^{-\eta/(8t)})\).  The zero-counting shell bound
used in Lemma~\ref{lem:v23-remainder} absorbs all differentiated distant
waves into the same estimate.  Critical-line waves, wrong-half-line leakage
from \textup{(V52.10)}, and the archimedean term
\textup{(V52.11)} are only polynomial before division by the
\(e^{\delta_*^2/(4t)}\) main amplitude.  Taking \(c\) smaller than the
resulting finite collection of positive gaps proves \textup{(V52.13)}.
\end{proof}

This theorem is not an additional explicit formula.  It is the analytic
quadrature missing from the real paired formula: positivity of the prime
frequencies chooses one of the two complex Gaussian halves canonically.

\subsection{The projective moment ray}

For brevity put
\[
 q_{k,\pm}:=
 \frac{\delta_*+i(x_{k,\pm}-\gamma_*)}{2t_k}.
 \tag{V52.14}
\]
Direct differentiation of \textup{(V52.5)} gives
\[
 DW_{p_*,t_k}=q_{k,\pm}W_{p_*,t_k},
 \qquad
 D^2W_{p_*,t_k}
 =\left(q_{k,\pm}^2+\frac1{2t_k}\right)W_{p_*,t_k}
 \tag{V52.15}
\]
at \(x=x_{k,\pm}\).

\begin{theorem}[Anisotropic exposed-zero ray]
\label{thm:v52-projective-ray}
Along both endpoint sequences, \(z_{t_k}(x_{k,\pm})\ne0\) for all large
\(k\), and
\[
 \boxed{
 \frac{P_{\chi,t_k}^{[1]}}{P_{\chi,t_k}^{[0]}}
 =q_{k,\pm}+O(t_k^{-M}e^{-c/t_k}),}
 \tag{V52.16}
\]
\[
 \boxed{
 \frac{P_{\chi,t_k}^{[2]}}{P_{\chi,t_k}^{[0]}}
 =q_{k,\pm}^2+\frac1{2t_k}
  +O(t_k^{-M}e^{-c/t_k}).}
 \tag{V52.17}
\]
The normalized prime phase and its absolute rate satisfy
\[
 \boxed{
 \frac{z_{t_k}(x_{k,\pm})}{|z_{t_k}(x_{k,\pm})|}
 \longrightarrow\pm i,}
 \tag{V52.18}
\]
\[
 \boxed{
 \lim_{k\to\infty}4t_k\log|z_{t_k}(x_{k,\pm})|
 =\delta_*^2-\frac14<0.}
 \tag{V52.19}
\]
Most importantly,
\[
 \boxed{
 \sqrt{t_k}\,\frac{c_{t_k}(x_{k,\pm})}
                         {z_{t_k}(x_{k,\pm})}
 \longrightarrow
 \alpha_*:=\frac{2\delta_*-1}{2\sqrt2},}
 \tag{V52.20}
\]
\[
 \boxed{
 t_k\,\frac{d_{t_k}(x_{k,\pm})}
                   {z_{t_k}(x_{k,\pm})}
 \longrightarrow
 \alpha_*^2=\frac{(2\delta_*-1)^2}{8}.}
 \tag{V52.21}
\]
Thus \(-1/(2\sqrt2)<\alpha_*<0\), and the projective ray itself recovers
\[
 \boxed{\delta_*=\frac12+\sqrt2\,\alpha_*.}
 \tag{V52.22}
\]
The three quantities \(|z_t|,|c_t|,|d_t|\) all have the same logarithmic
rate \(\delta_*^2-1/4\) after multiplication of their logarithms by \(4t\).
\end{theorem}

\begin{proof}
Equations \textup{(V52.16)}--\textup{(V52.17)} follow from
Theorem~\ref{thm:v52-complex-wave} and \textup{(V52.15)}.  At the nominal
minus and plus endpoints, the phases
\(\delta_*(x-\gamma_*)/(2t_k)\) are respectively \(\pi/2\) and
\(3\pi/2\).  The minus sign in \textup{(V52.5)} and the exponentially small
root displacement therefore give the phases \(-i\) and \(+i\), proving
\textup{(V52.18)}.

The modulus of \(W_{p_*,t_k}\) has exponential part
\(e^{\delta_*^2/(4t_k)}\), because
\(x_{k,\pm}-\gamma_*=O(t_k)\).  Divide by
\(Z_\chi(t_k)\asymp t_k^{-1/2}e^{1/(16t_k)}\) from
\textup{(V28.4)}.  Constants and powers of \(t_k\) disappear after
multiplication by \(4t_k\), proving \textup{(V52.19)}.

The exact standardized identities \textup{(V51.9)}--\textup{(V51.10)} give
\[
 \frac{c_t}{z_t}
 =\frac{P^{[1]}/P^{[0]}-\mu_t}{\sigma_t},
 \tag{V52.23}
\]
\[
 \frac{d_t}{z_t}
 =\frac{P^{[2]}/P^{[0]}-2\mu_tP^{[1]}/P^{[0]}+\mu_t^2}
        {\sigma_t^2}.
 \tag{V52.24}
\]
By \textup{(V52.12)},
\[
 q_{k,\pm}
 =\frac{\delta_*}{2t_k}+O(1)i,
 \tag{V52.25}
\]
while \(\mu_t=1/(4t)+o(t^{-1/2})\) and
\(\sigma_t=(2t)^{-1/2}(1+o(1))\).  Hence
\[
 \sqrt t\,\frac{q_{k,\pm}-\mu_t}{\sigma_t}
 \longrightarrow\frac{2\delta_*-1}{2\sqrt2}.
 \tag{V52.26}
\]
Using \textup{(V52.17)} in \textup{(V52.24)} yields
\[
 \frac{d_t}{z_t}
 =\frac{(q_{k,\pm}-\mu_t)^2+1/(2t)}{\sigma_t^2}
  +O(t^{-M}e^{-c/t}).
 \tag{V52.27}
\]
After multiplication by \(t\), the \(1/(2t)\) term and the bounded
imaginary part in \textup{(V52.25)} vanish, while the square of the real
term gives \(\alpha_*^2\).  This proves \textup{(V52.20)}--
\textup{(V52.22)}.  Polynomial projective factors do not change the rate in
\textup{(V52.19)}, proving the last assertion.
\end{proof}

The leading projective parabola
\[
 \left(\alpha,\alpha^2\right),
 \qquad -\frac1{2\sqrt2}<\alpha<0,
 \tag{V52.28}
\]
has a simple large-deviation meaning.  The exposed zero wave has logarithmic
frequency \(\delta_*/(2t)\), whereas the unsigned prime saddle is centered at
\(1/(4t)\).  Their separation is
\((2\delta_*-1)/(4t)\), or \(\alpha_*/\sqrt t\) standardized units.  The
off-line signal therefore approaches the origin of the V49 moment body from
a tail direction which escapes every fixed Gaussian standardized window.

\subsection{What the ray detects and what it cannot do}

\begin{corollary}[Projective simple-contact criterion]
\label{cor:v52-projective-criterion}
Paired GRH is equivalent to the nonexistence of compact simple-contact
sequences \(t_k\downarrow0\), \(x_{k,\pm}\), for which
\[
 \frac{z_{t_k}(x_{k,\pm})}{|z_{t_k}(x_{k,\pm})|}\to\pm i,
 \qquad
 \left(\sqrt{t_k}\frac{c_{t_k}}{z_{t_k}},
       t_k\frac{d_{t_k}}{z_{t_k}}\right)
 \to(\alpha,\alpha^2)
 \tag{V52.29}
\]
for some \(-1/(2\sqrt2)<\alpha<0\).
\end{corollary}

\begin{proof}
If GRH holds, every center in \textup{(V22.12)} is real and every summand is
a strictly positive Gaussian; hence \(m_{\chi,t}(x)>0\) for all real \(x\)
and \(t>0\), so no zero-level contact exists.  If GRH fails,
Theorem~\ref{thm:v37-centered-exposure} and
Theorem~\ref{thm:v52-projective-ray} produce the required sequences.
\end{proof}

The criterion is reconstructive rather than simpler than GRH.  In fact its
parameter \(\delta_*\) is already recoverable from endpoint spacing in
\textup{(V27.26)}.  The new content is that the full complex reader has a
canonical phase and projective ray, not merely the alternating real jets of
the scalar symbol.

There is also a contact-type obstruction.  The two points
\(x_{k,-},x_{k,+}\) are simple endpoints with exponentially large nonzero
slopes.  The stationary point between them is a negative local minimum, not
a zero-level contact.  Therefore the ray cannot be inserted into the V49
stationary value--slope--opening slice.  A stationary zero contact occurs at
a different heat-scale bifurcation and is not asymptotically isolated by the
small-time exposed-pair construction.

\begin{theorem}[V52 projective verdict]
\label{thm:v52-verdict}
Positive-frequency extraction upgrades the center-exposed scalar zero packet
to a full complex prime wave and proves the anisotropic projective limits
\textup{(V52.18)}--\textup{(V52.21)}.  These limits identify the off-line
displacement and sharpen V27 tomography, but they yield no contradiction:
\begin{enumerate}[label=\textup{(\roman*)},leftmargin=3em]
\item all unscaled moments still tend to the interior origin of the V49
      Bessel body with fixed slack tending to one;
\item the projective ray is generated by, and reconstructs, the assumed
      off-line zero;
\item the available contacts are simple endpoints, whereas positive opening
      in V49 concerns a stationary zero contact.
\end{enumerate}
Thus the V51 projective route is exact but reconstructive.  It does not
provide the missing signed exclusion.
\end{theorem}

\begin{firewall}[A zero signature is not a zero exclusion]
V52 recovers the dominant complex Gaussian half-wave, its endpoint phases,
its common normalized heat rate, and its complete leading standardized
projective ray.  The resulting criterion is equivalent to GRH because an
off-line zero produces the ray.  It does not bound that ray away, turn a
simple endpoint into a stationary contact, exclude an off-line zero, or
prove GRH.
\end{firewall}

\begin{assessment}[V52 acquisition and the next upstream distinction]
The projective calculation closes the immediate V51 proposal.  It also
identifies why the V49 Gaussian body looked too large: the actual off-line
signal is not a fixed standardized perturbation but a moderate-to-large
deviation whose projective coordinates diverge before anisotropic scaling.
The divergence itself is fully determined by \(\delta_*\) and therefore
contains no extra coercivity.

V53 should compare contact types rather than add moments.  The exact next
question is whether a global heat-time continuation from the small-time
simple endpoint pair to its first stationary zero-level bifurcation preserves
any orientation, winding, or index supplied by the analytic prime reader.
Such an invariant would connect the exposed-zero ray to the critical contact
without demanding pointwise phase cancellation at every intermediate time.
The first task is topological and exact: derive the signed creation/annihilation
index for real zeros of the caloric symbol under generic bifurcations, then
test whether the positive-frequency analytic completion fixes that index or
whether V48's caloric realization theorem permits both signs.  If both signs
are realizable, this continuation route should be closed before returning to
arithmetic estimates.
\end{assessment}

\section*{Version V52 append-only record}
\addcontentsline{toc}{section}{Version V52 append-only record}
The complete V51 source was retained before this continuation.  It had
\(786678\) bytes, \(23713\) lines, and SHA--256
\[
 \texttt{D9C86FE4A7AF6F9C45C6BF96EF8F8A49D1B36A31ED6530D7937D95620F53DBCE}.
\]
V52 extracted the canonical positive-frequency half of an exposed reflected
zero pair, proved complex dominance of that wave, computed the full
anisotropically scaled \((z,c,d)\) ray, and showed that the resulting GRH
criterion is reconstructive rather than excluding.

% =============================================================================
% END APPENDED V52 MATERIAL
% =============================================================================
% =============================================================================
% BEGIN APPENDED V53 MATERIAL
% =============================================================================

\section{V53: zero-set index, Hardy nodes, and the topological no-go}
\label{sec:grh-v53}

V52 attached opposite limiting complex phases to the two simple endpoints of
an exposed negative lobe and asked whether those phases survive continuation
to a stationary zero-level bifurcation.  The real heat geometry has a fixed
local answer: a generic threshold minimum annihilates one negatively oriented
and one positively oriented simple zero, so its signed spatial degree is
zero.  There is therefore no parity or Brouwer-degree obstruction.

The positive-frequency completion nevertheless detects something exact.
After adding the half zero-frequency convention, the completed analytic
signal
\[
 {\cal F}_{\chi,t}:={\cal C}_+m_{\chi,t}
                   ={\cal C}_+a_{\chi,t}-P_{\chi,t}^{[0]}
\]
has real part \(m_{\chi,t}/2\).  It is consequently purely imaginary on the
entire real zero set.  At the V52 exposed endpoints its phases tend to
\(+i\) and \(-i\).  Hence every connected zero-set arc joining those
endpoints contains a point where the \emph{whole} analytic signal vanishes:
\[
 {\cal F}_{\chi,t}(x)=0,
 \qquad\text{equivalently}\qquad
 P_{\chi,t}^{[0]}(x)={\cal C}_+a_{\chi,t}(x).
\]
This is a genuine complex strengthening of the scalar contact equation.

It is not a contradiction.  An explicit one-parameter family of periodic
heat flows has a genuine positivity threshold and realizes the forced Hardy
node on the left branch, on the right branch, or at the terminal fold.
Thus neither the signed zero index nor analytic completion fixes the required
orientation.  For the arithmetic symbol, the unresolved issue is global:
one must first prove that the two exposed branches belong to a zero-set
network component whose boundary signs force such a node.

\subsection{The universal fold and its trivial signed degree}

Let \(u_t=u_{xx}\) be real analytic near \((t_*,x_*)\), and suppose
\[
 u(t_*,x_*)=u_x(t_*,x_*)=0,
 \qquad
 u_{xx}(t_*,x_*)=\kappa>0.
 \tag{V53.1}
\]
Put \(\tau=t-t_*\) and \(y=x-x_*\).

\begin{theorem}[Caloric fold normal form and signed zero index]
\label{thm:v53-fold-index}
Under \textup{(V53.1)},
\[
 \boxed{
 u(t_*+\tau,x_*+y)
 =\kappa\left(\tau+\frac{y^2}{2}\right)
  +O\!\left(|y|^3+|\tau y|+\tau^2\right).}
 \tag{V53.2}
\]
For every sufficiently small \(\tau<0\), the zero set consists locally of
two simple branches
\[
 y_\pm(\tau)=\pm\sqrt{-2\tau}+O(\tau),
 \tag{V53.3}
\]
with
\[
 \operatorname{sgn}u_x(t_*+\tau,x_*+y_-)
 =-1,
 \qquad
 \operatorname{sgn}u_x(t_*+\tau,x_*+y_+)
 =+1.
 \tag{V53.4}
\]
They annihilate for increasing time at \(\tau=0\), and their total signed
spatial degree is
\[
 \boxed{(-1)+(+1)=0.}
 \tag{V53.5}
\]
\end{theorem}

\begin{proof}
The heat equation gives \(u_t(t_*,x_*)=u_{xx}(t_*,x_*)=\kappa\).
Taylor expansion proves \textup{(V53.2)}.  Divide by \(\kappa\), set
\(y=\sqrt{-2\tau}\,v\), and apply the implicit-function theorem near
\(v=\pm1\) to obtain \textup{(V53.3)}.  Differentiating
\textup{(V53.2)} in \(y\) gives
\[
 u_x=\kappa y+O(y^2+|\tau|),
 \tag{V53.6}
\]
which proves \textup{(V53.4)}.  The degree identity
\textup{(V53.5)} follows.
\end{proof}

Thus the zero number drops by two at a generic threshold minimum, while the
oriented zero count does not change.  This is the local form of the
one-dimensional variation-diminishing principle for the heat equation.
It supplies annihilation, not a nonzero conserved charge.

\subsection{The completed positive-frequency signal}

For a real tempered distribution \(f\), let \({\cal C}_+f\) denote the
Fourier multiplier equal to \(1\) on \(u>0\), \(1/2\) at \(u=0\), and
\(0\) on \(u<0\).  Equivalently, with the Hilbert-transform convention
\[
 \widehat{{\cal H}f}(u)=-i\,\operatorname{sgn}(u)\widehat f(u),
 \tag{V53.7}
\]
one has
\[
 {\cal C}_+f=\frac12(f+i{\cal H}f),
 \qquad
 2\Re{\cal C}_+f=f.
 \tag{V53.8}
\]
The harmless polynomial normalization of the archimedean distribution is
fixed so that the same half-zero-frequency convention is used in both
terms.

Define
\[
 A_{\chi,t}:={\cal C}_+a_{\chi,t},
 \qquad
 {\cal F}_{\chi,t}:=A_{\chi,t}-P_{\chi,t}^{[0]}.
 \tag{V53.9}
\]
Since \(a_{\chi,t}-m_{\chi,t}=2\Re P_{\chi,t}^{[0]}\) has only the
nonzero frequencies \(\pm\log n\), the positive-frequency identity
\textup{(V52.3)} gives
\[
 \boxed{
 {\cal F}_{\chi,t}={\cal C}_+m_{\chi,t},
 \qquad
 2\Re{\cal F}_{\chi,t}=m_{\chi,t}.}
 \tag{V53.10}
\]
Both \(A_{\chi,t}\) and \(P_{\chi,t}^{[0]}\) satisfy the heat equation, so
\[
 \partial_t{\cal F}_{\chi,t}
 =\partial_x^2{\cal F}_{\chi,t}.
 \tag{V53.11}
\]

\begin{proposition}[Zero level is an imaginary-axis constraint]
\label{prop:v53-imaginary-axis}
At every real zero-level point,
\[
 m_{\chi,t}(x)=0
 \quad\Longrightarrow\quad
 {\cal F}_{\chi,t}(x)\in i\mathbb R.
 \tag{V53.12}
\]
Moreover,
\[
 {\cal F}_{\chi,t}(x)=0
 \quad\Longleftrightarrow\quad
 \begin{cases}
 m_{\chi,t}(x)=0,\\
 {\cal H}m_{\chi,t}(x)=0,
 \end{cases}
 \quad\Longleftrightarrow\quad
 P_{\chi,t}^{[0]}(x)=A_{\chi,t}(x).
 \tag{V53.13}
\]
\end{proposition}

\begin{proof}
The first statement is the real-part identity in \textup{(V53.10)}.
Equation \textup{(V53.8)} then gives the first equivalence in
\textup{(V53.13)}, and the definition \textup{(V53.9)} gives the second.
\end{proof}

The equality in \textup{(V53.13)} is strictly stronger than the ordinary
contact equation
\(\Re P_{\chi,t}^{[0]}=a_{\chi,t}/2\): it also fixes the missing Hilbert
quadrature.

\subsection{Opposite exposed phases force a Hardy node on a joining arc}

At the V52 endpoint sequences, the archimedean completion is polynomial in
\(t^{-1}\) on compact \(x\)-sets, whereas the exposed wave has size
\(e^{\delta_*^2/(4t)}\) up to powers of \(t\).  Consequently
\[
 {\cal F}_{\chi,t_k}(x_{k,\pm})
 =-P_{\chi,t_k}^{[0]}(x_{k,\pm})
   \left(1+O(t_k^{-M}e^{-c/t_k})\right).
 \tag{V53.14}
\]
Together with \textup{(V52.18)}, this gives
\[
 \boxed{
 \frac{{\cal F}_{\chi,t_k}(x_{k,\pm})}
      {|{\cal F}_{\chi,t_k}(x_{k,\pm})|}
 \longrightarrow\mp i.}
 \tag{V53.15}
\]
In particular, for all large \(k\),
\[
 \Im{\cal F}_{\chi,t_k}(x_{k,-})>0,
 \qquad
 \Im{\cal F}_{\chi,t_k}(x_{k,+})<0.
 \tag{V53.16}
\]

\begin{theorem}[Conditional continuation forces a full Hardy node]
\label{thm:v53-hardy-node}
Fix a sufficiently large \(k\).  Suppose there is a connected path
\[
 \Gamma_k\subset
 \{(t,x):t\ge t_k,\ m_{\chi,t}(x)=0\}
 \tag{V53.17}
\]
containing both \((t_k,x_{k,-})\) and \((t_k,x_{k,+})\).  Then there is
\((t_k^\circ,x_k^\circ)\in\Gamma_k\) such that
\[
 \boxed{
 {\cal F}_{\chi,t_k^\circ}(x_k^\circ)=0,
 \qquad
 P_{\chi,t_k^\circ}^{[0]}(x_k^\circ)
 =A_{\chi,t_k^\circ}(x_k^\circ).}
 \tag{V53.18}
\]
Equivalently, the real zero-set path contains a simultaneous zero of
\(m_{\chi,t}\) and its Hilbert transform.
\end{theorem}

\begin{proof}
By Proposition~\ref{prop:v53-imaginary-axis}, the continuous function
\({\cal F}_{\chi,t}(x)\) takes values in \(i\mathbb R\) on \(\Gamma_k\).
Its imaginary part has opposite signs at the two distinguished points by
\textup{(V53.16)}.  The intermediate value theorem on the connected path
gives a point where the imaginary part vanishes.  The real part already
vanishes on \(\Gamma_k\), proving \textup{(V53.18)}.
\end{proof}

This conclusion does not require the joining point to be a stationary
contact.  If the two simple branches persist until they meet at one generic
fold, their union is such a path and the theorem applies.  If they reconnect
with other zero branches at earlier multiple zeros, the hypothesis becomes
a question about the global zero-set network rather than local heat
kinematics.

\subsection{The fractional Hopf sign at a threshold node}

There is a fixed derivative orientation when the forced node happens exactly
at a nonconstant global threshold minimum.  Write
\[
 |D|:={\cal H}\partial_x,
 \tag{V53.19}
\]
whose Fourier multiplier is \(|u|\).

\begin{proposition}[Fractional Hopf orientation]
\label{prop:v53-fractional-hopf}
Suppose \(m_{\chi,t_*}\ge0\) on \(\mathbb R\),
\(m_{\chi,t_*}(x_*)=0\), the standard singular-integral representation of
\(|D|m_{\chi,t_*}(x_*)\) converges, and \(m_{\chi,t_*}\) is not identically
zero.  Then
\[
 \boxed{
 |D|m_{\chi,t_*}(x_*)<0.}
 \tag{V53.20}
\]
If in addition \({\cal F}_{\chi,t_*}(x_*)=0\), then
\[
 \boxed{
 \Im\partial_x{\cal F}_{\chi,t_*}(x_*)
 =\frac12|D|m_{\chi,t_*}(x_*)<0.}
 \tag{V53.21}
\]
\end{proposition}

\begin{proof}
With a positive normalization constant \(c\), the one-dimensional
half-Laplacian has the pointwise formula
\[
 |D|f(x_*)
 =c\,\operatorname{pv}\!\int_{\mathbb R}
 \frac{f(x_*)-f(y)}{(x_*-y)^2}\,dy.
 \tag{V53.22}
\]
Every integrand is nonpositive under the stated minimum hypothesis, and it
is strictly negative on a set of positive measure because the function is
not constant.  This proves \textup{(V53.20)}.  Differentiating
\({\cal F}=(m+i{\cal H}m)/2\), and using
\({\cal H}\partial_x=|D|\), proves \textup{(V53.21)}.
\end{proof}

The sign \textup{(V53.21)} is compatible with a Hardy node; it does not
forbid one.  The next model realizes it exactly.

\subsection{A threshold model with a branch-selectable Hardy node}

Work on the circle and fix \(|\varepsilon|<1\).  Set
\[
 f_\varepsilon(x)
 :=(1-\cos x)(1+\varepsilon\sin x)\ge0
 \tag{V53.23}
\]
and define the entire periodic heat flow
\[
 \begin{aligned}
 u_\varepsilon(t,x)
 &:=e^{t\partial_x^2}f_\varepsilon(x)\\
 &=1-e^{-t}\cos x
   +\varepsilon e^{-t}\sin x
   -\frac{\varepsilon}{2}e^{-4t}\sin2x.
 \end{aligned}
 \tag{V53.24}
\]

\begin{theorem}[Both Hardy-node orientations occur at genuine thresholds]
\label{thm:v53-node-realization}
For every \(|\varepsilon|<1\), the flow \textup{(V53.24)} satisfies
\[
 u_\varepsilon(t,x)\ge0
 \quad(t\ge0,\ x\in\mathbb T),
 \tag{V53.25}
\]
while
\[
 u_\varepsilon(t,0)=1-e^{-t}<0
 \quad(t<0).
 \tag{V53.26}
\]
Thus \(t_*=0\) is a genuine positivity threshold, with
\[
 u_\varepsilon(0,0)=\partial_xu_\varepsilon(0,0)=0,
 \qquad
 \partial_x^2u_\varepsilon(0,0)=1.
 \tag{V53.27}
\]

Its completed positive-frequency signal is
\[
 \boxed{
 F_\varepsilon(t,x)
 =\frac12
  -\left(\frac12+\frac{i\varepsilon}{2}\right)e^{-t}e^{ix}
  +\frac{i\varepsilon}{4}e^{-4t}e^{2ix}.}
 \tag{V53.28}
\]
For all sufficiently small \(\varepsilon\), it has a unique nearby real
Hardy node
\[
 F_\varepsilon(t_\varepsilon,x_\varepsilon)=0
 \tag{V53.29}
\]
with
\[
 \boxed{
 x_\varepsilon=-\frac{\varepsilon}{2}+O(\varepsilon^2),
 \qquad
 t_\varepsilon=-\frac{\varepsilon^2}{8}
                 +O(\varepsilon^3).}
 \tag{V53.30}
\]
Hence the node lies on the left zero branch when \(\varepsilon>0\), on the
right zero branch when \(\varepsilon<0\), and at the fold when
\(\varepsilon=0\).
\end{theorem}

\begin{proof}
Equation \textup{(V53.25)} follows from positivity preservation of the
forward heat semigroup applied to the nonnegative function
\textup{(V53.23)}.  Equation \textup{(V53.26)} is immediate from
\textup{(V53.24)}, and direct differentiation gives
\textup{(V53.27)}.

Taking the half constant coefficient and all positive Fourier modes of
\textup{(V53.24)} gives \textup{(V53.28)}.  At
\((t,x,\varepsilon)=(0,0,0)\),
\[
 F_0=0,
 \qquad
 \partial_tF_0=\frac12,
 \qquad
 \partial_xF_0=-\frac i2,
 \qquad
 \partial_\varepsilon F_0=-\frac i4.
 \tag{V53.31}
\]
The real Jacobian of \((t,x)\mapsto F_0(t,x)\) is invertible.  The
implicit-function theorem therefore gives a unique real node
\((t_\varepsilon,x_\varepsilon)\).  Differentiating
\textup{(V53.29)} at \(\varepsilon=0\) yields
\[
 t_\varepsilon'=0,
 \qquad
 x_\varepsilon'=-\frac12.
 \tag{V53.32}
\]
Because \(u_\varepsilon=2\Re F_\varepsilon\), the node also lies on the
real zero set.  Substitution into the fold expansion
\[
 u_\varepsilon(t,x)
 =t+\frac{x^2}{2}
  +O\!\left(|t|^2+|t x|+|x|^3
             +|\varepsilon|\,|x|^3\right)
 \tag{V53.33}
\]
then gives
\(t_\varepsilon=-x_\varepsilon^2/2+O(\varepsilon^3)\), proving
\textup{(V53.30)}.  The local zero branches are
\(x=\pm\sqrt{-2t}+O(t)\), so the sign of \(x_\varepsilon\) proves the
last assertion.
\end{proof}

At the threshold itself,
\[
 F_\varepsilon(0,0)=-\frac{i\varepsilon}{4}.
 \tag{V53.34}
\]
Thus a nonsymmetric threshold need not be a Hardy node.  The node moves
smoothly onto one of the incoming simple branches, with either choice
realized by the sign of \(\varepsilon\).  In the symmetric case
\(\varepsilon=0\), it sits at the fold and
\(\Im\partial_xF_0(0,0)=-1/2\), exactly as required by
Proposition~\ref{prop:v53-fractional-hopf}.

\subsection{Exact status of the continuation route}

\begin{theorem}[V53 topological-continuation verdict]
\label{thm:v53-verdict}
The contact-continuation proposal has the following exact status.
\begin{enumerate}[label=\textup{(\roman*)},leftmargin=3em]
\item A generic positive-threshold minimum annihilates a pair of simple
      zeros with opposite spatial orientations and zero total signed degree.
\item The completed positive-frequency signal is purely imaginary on the
      real zero set.
\item The V52 exposed endpoints have opposite signs on that imaginary axis.
      Therefore any zero-set path joining them contains a full Hardy node,
      equivalently the complex prime--archimedean equality
      \(P_{\chi,t}^{[0]}=A_{\chi,t}\).
\item If that node occurs at a global threshold minimum, its transverse
      analytic derivative has the fixed fractional Hopf sign
      \textup{(V53.21)}.
\item Periodic caloric threshold models realize the node on either incoming
      branch and at the fold, so none of these topological or analytic signs
      excludes the bifurcation.
\end{enumerate}
Consequently analytic completion upgrades a connected continuation to a
complex node condition, but topology supplies no contradiction and no
proof that the two distinguished arithmetic branches actually join in the
required zero-set component.
\end{theorem}

\begin{proof}
Combine Theorem~\ref{thm:v53-fold-index}, Proposition
\ref{prop:v53-imaginary-axis}, Theorem~\ref{thm:v53-hardy-node},
Proposition~\ref{prop:v53-fractional-hopf}, and Theorem
\ref{thm:v53-node-realization}.
\end{proof}

\begin{firewall}[A forced Hardy node is not an exclusion]
Theorem~\ref{thm:v53-hardy-node} is conditional on zero-set connectivity.
Even when it applies, the resulting equality
\(P_{\chi,t}^{[0]}=A_{\chi,t}\) is compatible with ordinary heat flow, as
the explicit periodic family proves.  The model does not reproduce the
Dirichlet coefficients, and the node equality may still be useful as an
arithmetic target.  No off-line zero is excluded and GRH remains open.
\end{firewall}

\begin{assessment}[V53 closes the generic topological-index route]
The signed real-zero index is trivial, but positive-frequency completion
does retain one nonlocal trace: opposite exposed endpoint phases force a
simultaneous zero of \(m\) and \({\cal H}m\) whenever the endpoint branches
join.  This converts the proposed winding invariant into the concrete
complex equation \(P=A\).  The branch-selectable threshold family shows
that the location and orientation of this node are not constrained by
caloricity, positivity above threshold, or the Hardy projection alone.

The next upstream question is global rather than local.  V54 should build
the finite zero-set continuation graph on a strip
\([t_0,T]\times\mathbb R\), using the archimedean growth at spatial infinity
and the one-dimensional zero-number principle.  It should decide whether
the two V37 exposed endpoints must lie in a component forcing a Hardy node,
or whether reconnection with other zero branches can separate them.  If a
node is unavoidable, the remaining arithmetic target is the full complex
equality \(P_{\chi,t}^{[0]}=A_{\chi,t}\) together with its fractional
derivative orientation; if reconnection defeats the implication, the
continuation route should be closed.
\end{assessment}

\section*{Version V53 append-only record}
\addcontentsline{toc}{section}{Version V53 append-only record}
The complete V52 source was retained before this continuation.  It had
\(802565\) bytes, \(24147\) lines, and SHA--256
\[
 \texttt{231874E300A8FFE25FA4FAB0A300DD4FC2032DD96884D3BBC083567C52DC0CDC}.
\]
V53 computed the generic caloric fold index, introduced the completed Hardy
signal of the Weil symbol, proved that opposite exposed endpoint phases force
a full analytic-signal node on every joining zero-set path, derived the
fractional Hopf orientation at a threshold node, and constructed genuine
threshold flows placing that node on either branch or at the fold.

% =============================================================================
% END APPENDED V53 MATERIAL
% =============================================================================
% =============================================================================
% BEGIN APPENDED V54 MATERIAL
% =============================================================================

\section{V54: the compact zero graph and the partner-exchange dichotomy}
\label{sec:grh-v54}

V53 showed that a zero-set path joining the two exposed endpoints must
contain a full Hardy node, but left the joining hypothesis unresolved.
This section constructs the global continuation graph on every positive-time
strip and determines exactly what one-dimensional zero ordering does and
does not imply.

The symbol tends to \(+\infty\) uniformly at spatial infinity on such a
strip, so its zero set is compact.  The classical one-dimensional
zero-number theorem makes the number of real zeros nonincreasing and forces a
strict drop at every multiple zero.  After resolving higher-order collision
vertices, the graph therefore gives a noncrossing matching of the simple
zeros on the initial time slice.  Every matched pair has opposite spatial
slope orientation.

This does \emph{not} imply that the endpoints of an initially negative
interval are matched to one another.  A positive interval between two
negative intervals may disappear first at a zero maximum; the two negative
components then merge, and the original endpoint partners are exchanged.
An explicit two-mode heat flow realizes exactly this event before its final
positivity threshold.

There is nevertheless a precise alternative.  At a simple zero define its
Hardy--slope sign by the product
\[
 \Im{\cal F}_{\chi,t}(x)\,\partial_xm_{\chi,t}(x).
\]
Both V52 exposed endpoints have negative product.  If a continuation arc
contains no Hardy node, the sign of \(\Im{\cal F}\) is constant on it, while
the matched endpoint has the opposite spatial slope.  Thus each exposed
endpoint must either lead to a Hardy node or be matched to a contact with
positive Hardy--slope product.  Partner exchange does not erase the
analytic obstruction; it moves it to a sign-defect contact elsewhere on the
same small-time slice.

\subsection{Uniform strip compactness and zero-number monotonicity}

Fix
\[
 0<t_0<T<\infty.
 \tag{V54.1}
\]
The uniform archimedean asymptotic \textup{(V21.4)} and the absolutely
convergent prime majorant at time \(t_0\) give
\[
 \boxed{
 m_{\chi,t}(x)=\log(2+|x|)+O_{\chi,t_0,T}(1)}
 \tag{V54.2}
\]
uniformly for \(t\in[t_0,T]\) and \(x\in\mathbb R\).

\begin{lemma}[Compact strip zero set]
\label{lem:v54-strip-compact}
There is \(R=R(\chi,t_0,T)<\infty\) such that
\[
 m_{\chi,t}(x)>0
 \qquad
 (t\in[t_0,T],\ |x|\ge R).
 \tag{V54.3}
\]
Consequently
\[
 {\cal Z}[t_0,T]
 :=\{(t,x)\in[t_0,T]\times\mathbb R:m_{\chi,t}(x)=0\}
 \tag{V54.4}
\]
is compact, and every time slice contains only finitely many zeros unless
the slice is identically zero.
\end{lemma}

\begin{proof}
Choose \(R\) in \textup{(V54.2)} so that the logarithmic term dominates
the uniform error.  This proves \textup{(V54.3)} and bounds the zero set.
Joint continuity makes it closed.  For \(t>0\), the Gaussian convolution
and the damped prime series make \(m_{\chi,t}\) real analytic in \(x\).
It is not identically zero by \textup{(V54.2)}, so its zeros in
\([-R,R]\) are finite.
\end{proof}

Let \(N_\chi(t)\) denote the number of distinct real zeros of
\(m_{\chi,t}\), counted without multiplicity.

\begin{lemma}[Sturm--Angenent zero-number rule]
\label{lem:v54-zero-number}
On a strip \textup{(V54.1)}, \(N_\chi(t)\) is finite and nonincreasing
apart from the harmless convention at a multiple-zero time.  More
precisely, it is locally constant while all zeros are simple and drops
strictly as increasing time passes a multiple spatial zero.
\end{lemma}

\begin{proof}
By Lemma~\ref{lem:v54-strip-compact}, choose \(R\) with strict positive
boundary values at \(x=\pm R\) throughout the strip.  Apply the classical
one-dimensional parabolic zero-number theorem to
\[
 \partial_tm_{\chi,t}=\partial_x^2m_{\chi,t}
 \tag{V54.5}
\]
on \([-R,R]\).  The theorem states that the number of zeros of a nontrivial
solution is finite, nonincreasing, and strictly decreases through a
multiple zero.  The positive boundary values prevent boundary entry or
exit, so the interval count is the full real-line count.\footnote{This is
the special heat-equation case of S.~Angenent, \emph{The zero set of a
solution of a parabolic equation}, J. Reine Angew. Math. \textbf{390}
(1988), 79--96.}
\end{proof}

If paired GRH fails, then \(\vartheta_\chi>0\) by
Theorem~\ref{thm:v20-threshold}.  For every \(T>\vartheta_\chi\),
\[
 m_{\chi,T}>0
 \tag{V54.6}
\]
strictly: at the threshold the symbol is nonnegative and nonzero, and
convolution with the strictly positive heat kernel for time
\(T-\vartheta_\chi\) is everywhere positive.  Hence the strip graph from a
regular \(t_0<\vartheta_\chi\) has no zero endpoints on its top or vertical
sides.

\subsection{The noncrossing matching induced by the zero graph}

Choose \(t_0\) so that all zeros on the bottom slice are simple, and write
\[
 x_1<x_2<\cdots<x_{2N}.
 \tag{V54.7}
\]
The number is even because the symbol is positive at both spatial ends.
Its slopes alternate:
\[
 \boxed{
 \operatorname{sgn}\partial_xm_{\chi,t_0}(x_j)=(-1)^j.}
 \tag{V54.8}
\]

\begin{theorem}[Zero-graph noncrossing matching]
\label{thm:v54-noncrossing-matching}
Let \(T>\vartheta_\chi\), so that \textup{(V54.6)} holds.  The compact
zero graph on \([t_0,T]\) induces, after resolving every higher-order
collision vertex into its order-preserving local branches, a noncrossing
perfect matching
\[
 \pi:\{1,\ldots,2N\}\longrightarrow\{1,\ldots,2N\}.
 \tag{V54.9}
\]
For every matched pair \(j<\pi(j)\),
\[
 \boxed{
 \pi(j)-j\ \text{is odd},
 \qquad
 \operatorname{sgn}m_x(t_0,x_{\pi(j)})
 =-\operatorname{sgn}m_x(t_0,x_j).}
 \tag{V54.10}
\]
\end{theorem}

\begin{proof}
At a simple spacetime point, the implicit-function theorem gives a unique
zero arc.  A generic double spatial zero is the fold of
Theorem~\ref{thm:v53-fold-index}; higher-order vertices have the
order-preserving resolution supplied by the one-dimensional zero-number
rule.  No component can be born in the interior, because that would
increase the zero number.  No component can be closed entirely inside the
strip: its lowest-time point would again be a multiple zero creating
zeros forward in time.  Compactness, the absence of top and side zeros,
and the strict zero-number drop therefore pair all bottom endpoints.

Distinct zero arcs cannot cross at a simple point.  Resolving a collision
according to spatial order consequently produces a noncrossing matching.
If \(j<\ell=\pi(j)\), every bottom endpoint strictly between \(j\) and
\(\ell\) must be matched to another endpoint in the same interval;
otherwise its arc would cross the \(j\)-to-\(\ell\) arc.  Thus
\(\ell-j-1\) is even and \(\ell-j\) is odd.  Equation
\textup{(V54.8)} then proves the slope identity.
\end{proof}

The theorem determines the parity of every partner but not the partner
itself.  In particular, an adjacent pair \((j,j+1)\) bounding one negative
interval need not satisfy \(\pi(j)=j+1\).

\subsection{Hardy sign transport and the exact dichotomy}

At a simple zero with \({\cal F}_{\chi,t}(x)\ne0\), define
\[
 \eta_{\chi,t}(x)
 :=\operatorname{sgn}\Im{\cal F}_{\chi,t}(x),
 \tag{V54.11}
\]
and call the zero \emph{Hardy coherent} when
\[
 \eta_{\chi,t}(x)\,
 \operatorname{sgn}\partial_xm_{\chi,t}(x)=-1,
 \tag{V54.12}
\]
and a \emph{Hardy sign defect} when the product is \(+1\).

By \textup{(V53.16)} and the fact that the exposed lobe is negative between
its endpoints,
\[
 \boxed{
 \eta_{\chi,t_k}(x_{k,\pm})\,
 \operatorname{sgn}\partial_xm_{\chi,t_k}(x_{k,\pm})=-1.}
 \tag{V54.13}
\]
Thus both exposed endpoints are coherent.

\begin{theorem}[Hardy-node or sign-defect alternative]
\label{thm:v54-node-defect}
Let \(x_j<x_{j+1}\) be the two exposed endpoints on a regular sufficiently
small time slice \(t_0\), and let \(\pi\) be the zero-graph matching of
Theorem~\ref{thm:v54-noncrossing-matching}.  For each
\(r\in\{j,j+1\}\), at least one of the following occurs:
\begin{enumerate}[label=\textup{(\alph*)},leftmargin=3em]
\item the continuation component issuing from \(x_r\) contains a Hardy
      node
      \[
      {\cal F}_{\chi,t}(x)=0;
      \tag{V54.14}
      \]
\item its matched bottom endpoint \(x_{\pi(r)}\) is a Hardy sign defect:
      \[
      \boxed{
      \Im{\cal F}_{\chi,t_0}(x_{\pi(r)})
      \,\partial_xm_{\chi,t_0}(x_{\pi(r)})>0.}
      \tag{V54.15}
      \]
\end{enumerate}
If the exposed endpoints are matched to each other, alternative
\textup{(a)} necessarily occurs on their common component.  If neither
continuation component contains a node, the exposed endpoints are matched
to two external contacts, both satisfying \textup{(V54.15)}.
\end{theorem}

\begin{proof}
On the zero graph, Proposition~\ref{prop:v53-imaginary-axis} gives
\({\cal F}\in i\mathbb R\).  If a component contains no Hardy node, the
continuous nonzero function \(\Im{\cal F}\) has constant sign along that
component.  Hence
\[
 \eta_{\chi,t_0}(x_{\pi(r)})
 =\eta_{\chi,t_0}(x_r).
 \tag{V54.16}
\]
The matched slopes have opposite signs by \textup{(V54.10)}, while the
exposed endpoint is coherent by \textup{(V54.13)}.  Multiplying the two
sign reversals proves \textup{(V54.15)}.

The two exposed endpoints have opposite values of \(\eta\) by
\textup{(V53.16)}.  If they are matched, a nonvanishing imaginary signal
could not retain both signs on one connected component, so a Hardy node
must occur.  If neither component has a node, they cannot be matched
together and \textup{(V54.15)} applies to each distinct external partner.
\end{proof}

At a Hardy node, \textup{(V53.13)} gives the full complex equality
\[
 P_{\chi,t}^{[0]}=A_{\chi,t}.
 \tag{V54.17}
\]
At a sign-defect contact, one instead has
\[
 m_{\chi,t}=0,
 \qquad
 \Im(A_{\chi,t}-P_{\chi,t}^{[0]})\,
 \partial_xm_{\chi,t}>0.
 \tag{V54.18}
\]
Theorem~\ref{thm:v54-node-defect} is therefore an unconditional graph
dichotomy once a regular exposed bottom slice is chosen, but neither branch
is excluded by heat geometry.

\subsection{An explicit heat flow exchanges the original partners}

The failure of adjacent-pair preservation already occurs in a two-mode
periodic heat flow.  Put
\[
 U(t,x)
 :=1-\frac52e^{-t}\cos x+\frac85e^{-4t}\cos2x.
 \tag{V54.19}
\]
It satisfies \(U_t=U_{xx}\).

\begin{theorem}[Caloric partner exchange before positivity]
\label{thm:v54-partner-exchange}
At \(t=0\), the function \(U\) has exactly four simple zeros on the circle
and two negative intervals separated by positive intervals around
\(x=0\) and \(x=\pi\).  As time increases, the two inner zeros first
annihilate at a zero local maximum at \(x=0\), merging the negative
intervals.  The two outer zeros later annihilate at a zero local minimum at
\(x=0\), after which \(U\) is strictly positive.

Consequently each negative interval present at \(t=0\) has one endpoint in
each of the two zero-graph matching components.  Its original adjacent
endpoints do not continue to the same fold.
\end{theorem}

\begin{proof}
At \(t=0\), write \(y=\cos x\).  Then
\[
 U(0,x)=\frac{16}{5}y^2-\frac52y-\frac35,
 \tag{V54.20}
\]
whose roots are
\[
 y_\pm=\frac{25\pm\sqrt{1393}}{64}.
 \tag{V54.21}
\]
Both lie strictly between \(-1\) and \(1\), so each has two circle
preimages and all four zeros are simple.  The quadratic is negative between
its roots and positive at \(y=\pm1\), giving the asserted two negative
intervals.

At the central critical point,
\[
 G(t):=U(t,0)
 =1-\frac52e^{-t}+\frac85e^{-4t}.
 \tag{V54.22}
\]
One has \(G(0)=1/10>0\), \(G(t)\to1\) as \(t\to\infty\), and
\[
 G'(t)=\frac52e^{-t}-\frac{32}{5}e^{-4t}.
 \tag{V54.23}
\]
The unique critical time is
\[
 t_m=\frac13\log\frac{64}{25}.
 \tag{V54.24}
\]
Writing \(r=e^{-t_m}=(25/64)^{1/3}\), the critical-point relation gives
\[
 G(t_m)=1-\frac{15}{8}r<0.
 \tag{V54.25}
\]
Thus \(G\) has two positive zeros
\[
 0<t_a<t_m<t_b.
 \tag{V54.26}
\]

Since \(U_t(t,0)=U_{xx}(t,0)=G'(t)\), the first contact at \(t_a\) has
negative second derivative and is a local maximum crossing downward.  It
annihilates the inner pair and merges the two negative intervals.  At
\(t_b\), the second derivative is positive and the remaining outer pair
annihilates at a minimum crossing upward.

For completeness, at time \(t\) the expression in \(y\) is the convex
quadratic
\[
 \frac{16}{5}e^{-4t}y^2-\frac52e^{-t}y
 +1-\frac85e^{-4t}.
 \tag{V54.27}
\]
For \(t>t_m\), its vertex lies to the right of \(y=1\), so its minimum on
\([-1,1]\) is \(G(t)\).  Hence \(U(t_b,\cdot)\ge0\), and strict positivity
for \(t>t_b\) follows either directly or from the positive heat kernel.
This proves the full event ordering.
\end{proof}

The example is even, so its completed Hardy signal vanishes at both
central folds.  It is used only to disprove preservation of the original
adjacent pairing.  The asymmetric family of
Theorem~\ref{thm:v53-node-realization} supplies the complementary local
fact: after its branch node but before its fold, a regular time slice has
two zeros with the same sign of \(\Im F_\varepsilon\), so one is a Hardy
sign defect.  Therefore both outcomes in
Theorem~\ref{thm:v54-node-defect} are compatible with genuine positivity
thresholds.

\subsection{Consequences for the exposed arithmetic lobe}

\begin{corollary}[Global continuation alternative under GRH failure]
\label{cor:v54-grh-alternative}
Assume paired GRH fails and choose a regular sufficiently small exposed
time \(t_k\).  Before the zero graph reaches any
\(T>\vartheta_\chi\), at least one of the following must occur:
\[
 \boxed{
 \begin{aligned}
 &P_{\chi,t}^{[0]}(x)=A_{\chi,t}(x)
 &&\text{at a zero-level point with }t\in[t_k,T],\\
 &\Im(A_{\chi,t_k}-P_{\chi,t_k}^{[0]})(y)\,
   \partial_xm_{\chi,t_k}(y)>0
 &&\text{at another simple zero }y.
 \end{aligned}}
 \tag{V54.28}
\]
If no Hardy node occurs on either exposed continuation component, there are
two distinct contacts of the second type.
\end{corollary}

\begin{proof}
Failure gives \(\vartheta_\chi>0\), while V37 and V52 give the exposed
simple endpoints and their coherent signs.  Lemma
\ref{lem:v54-strip-compact}, strict positivity above the threshold, and
Theorem~\ref{thm:v54-node-defect} give the stated alternative.
\end{proof}

The sign-defect partners need not remain in the V37 exposure neighbourhood
as \(t_k\downarrow0\).  They may escape to growing ordinates or lie where
several zero packets have comparable exponential weight.  Thus
\textup{(V54.28)} replaces the false same-pair continuation claim by a
correct global dichotomy, but does not yet supply a compact sign-defect
sequence.

\begin{theorem}[V54 zero-graph verdict]
\label{thm:v54-verdict}
On every positive-time strip, the Weil heat symbol has a compact finite
zero graph and obeys the one-dimensional zero-number rule.  That graph
induces a noncrossing matching of a regular initial zero slice, with
opposite slope orientations at matched endpoints.  However:
\begin{enumerate}[label=\textup{(\roman*)},leftmargin=3em]
\item adjacent endpoints of one negative interval need not be matched;
\item an explicit two-mode heat flow exchanges their partners through a
      zero-maximum merger before final positivity;
\item the V52 opposite phases survive globally only as the alternative
      Hardy node or Hardy sign defect in
      Theorem~\ref{thm:v54-node-defect};
\item both alternatives occur in ordinary caloric threshold models.
\end{enumerate}
Therefore the zero-number theorem does not close the V53 connectivity gap.
It converts that gap into the exact arithmetic alternative
\textup{(V54.28)}.
\end{theorem}

\begin{firewall}[Noncrossing zero graphs do not preserve negative-lobe
partners]
Zero-number monotonicity prevents creation and crossing of real zero
branches, but it permits the loss of a positive interval at a zero maximum.
That event exchanges which initial zeros are paired.  The resulting Hardy
node/sign-defect dichotomy is necessary under a hypothetical off-line zero,
not contradictory.  The periodic examples establish only the absence of a
generic heat-topological obstruction; they do not model the Dirichlet prime
phases.  GRH remains open.
\end{firewall}

\begin{assessment}[V54 isolates the next arithmetic localization problem]
The global continuation graph is now exact enough to remove the unsupported
same-lobe joining assumption.  The exposed endpoints cannot disappear
without leaving one of two traces: a full complex prime--archimedean match
at an intermediate Hardy node, or wrong Hardy--slope orientation at external
contacts on the same small-time slice.  This is stronger than an arbitrary
partner-exchange statement because the analytic sign is transported along
each zero component.

V55 should return to the zero-side Gaussian expansion and classify where a
small-time Hardy sign defect can occur.  Unique center exposure should force
coherence by the same complex half-wave calculation as V52.  Therefore a
defect must either leave every fixed ordinate window or occur in a tie set
where at least two reflected zero packets have equal leading exponent.
The useful target is a compactness-or-tie theorem for the partners in
\textup{(V54.28)}.  If escape and equal-exponent interference both remain
freely realizable, the global continuation branch should be closed and the
remaining Hardy-node equality treated directly as GRH-strength
cancellation.
\end{assessment}

\section*{Version V54 append-only record}
\addcontentsline{toc}{section}{Version V54 append-only record}
The complete V53 source was retained before this continuation.  It had
\(819433\) bytes, \(24651\) lines, and SHA--256
\[
 \texttt{FCACD350AE512C08F088C1B2AFE92D3D50C9A8B9663C54F146C3450759704BD3}.
\]
V54 proved strip compactness and zero-number monotonicity, constructed the
noncrossing matching of initial zeros, derived the Hardy-node/sign-defect
dichotomy for exposed endpoints, and exhibited an explicit heat flow that
exchanges the original negative-lobe partners before reaching positivity.

% =============================================================================
% END APPENDED V54 MATERIAL
% =============================================================================
% =============================================================================
% BEGIN APPENDED V55 MATERIAL
% =============================================================================

\section{V55: Hardy coherence under unique exposure and the tie-set trichotomy}
\label{sec:grh-v55}

V54 reduced global partner exchange to a concrete alternative: either an
intermediate Hardy node satisfies the full complex equality \(P=A\), or
external simple contacts on the same small-time slice have positive
Hardy--slope product.  This section locates the latter contacts on the
zero-side Gaussian envelope.

A uniformly exposed off-line packet with positive action cannot produce a
Hardy sign defect.  Its completed analytic half is a single rotating
Gaussian wave.  At a real zero, the imaginary quadrature is proportional to
\(\sin\theta\), while the scalar slope is proportional to
\(-\sin\theta\); their product is strictly negative.  All other packets and
the wrong Fourier half are exponentially smaller.

Consequently a compact sequence of sign defects with action bounded away
from zero must approach a point where at least two reflected zero packets
tie for the leading score.  Applied to the V54 partners, failure of GRH
therefore forces one of four traces: a Hardy node, ordinate escape, collapse
to the zero-action boundary, or positive-action packet interference.  A
two-frequency caloric analytic signal realizes the reversed sign at an
interference point, so the tie alternative is not excluded by Hardy
kinematics alone.

\subsection{The pointwise score envelope and its tie set}

For each grouped reflected pair \(p=(\gamma_p,\delta_p,m_p)\), retain
\[
 S_p(x):=\delta_p^2-(x-\gamma_p)^2
 \tag{V55.1}
\]
from \textup{(V37.2)}, including \(\delta_p=0\) for critical-line packets.
Define the upper score envelope
\[
 D_\chi(x):=\sup_p S_p(x)
 \tag{V55.2}
\]
and the positive-action tie set
\[
 {\cal T}_\chi^+
 :=\left\{x:
 D_\chi(x)>0,\ 
 S_p(x)=S_q(x)=D_\chi(x)
 \text{ for distinct }p,q\right\}.
 \tag{V55.3}
\]
Same-ordinate packets with smaller displacement are discarded as in V37.

\begin{lemma}[Local finiteness of positive-action competitors]
\label{lem:v55-local-finite}
For every compact interval \(I\) and every \(\varepsilon>0\), only finitely
many pairs satisfy
\[
 \sup_{x\in I}S_p(x)\ge\varepsilon.
 \tag{V55.4}
\]
Hence, on
\(\{x\in I:D_\chi(x)\ge\varepsilon\}\), the envelope is the maximum of
finitely many parabolas.  A unique maximizer has a positive gap on a
neighbourhood, and the tie set there is finite.
\end{lemma}

\begin{proof}
If \(I\subset[-R,R]\) and \textup{(V55.4)} holds, then
\[
 \operatorname{dist}(\gamma_p,I)^2
 \le\delta_p^2-\varepsilon<\frac14.
 \tag{V55.5}
\]
Thus \(\gamma_p\) lies in a fixed bounded interval.  Local finiteness of
the zero set leaves only finitely many pairs.  The remaining assertions
follow from continuity.  For distinct ordinates, equality of two scores is
the single affine condition
\[
 x=\frac{\gamma_p^2-\gamma_q^2+\delta_q^2-\delta_p^2}
          {2(\gamma_p-\gamma_q)}.
 \tag{V55.6}
\]
\end{proof}

\subsection{A single exposed analytic half forces coherence}

Write the scalar packet and its positive analytic half as
\[
 K_{p,t}(x)=2B_{p,t}(x)\cos\theta_{p,t}(x),
 \tag{V55.7}
\]
\[
 {\cal C}_+K_{p,t}(x)
 =B_{p,t}(x)e^{i\theta_{p,t}(x)}
   +\text{wrong-half leakage},
 \tag{V55.8}
\]
where
\[
 B_{p,t}(x)
 :=\frac{2\pi m_p}{\sqrt{4\pi t}}
   \exp\!\left(\frac{S_p(x)}{4t}\right),
 \qquad
 \theta_{p,t}(x)
 :=\frac{\delta_p(x-\gamma_p)}{2t}.
 \tag{V55.9}
\]
Equation \textup{(V55.8)} is the sign-reversed form of
Lemma~\ref{lem:v52-analytic-half}, because
\({\cal F}={\cal C}_+m\), whereas V52 defined \(W\) as the contribution
to \(P={\cal C}_+(a-m)\).

\begin{theorem}[Uniform unique exposure forces Hardy coherence]
\label{thm:v55-exposed-coherence}
Let \(t_k\downarrow0\) and let \(y_k\) remain in a compact interval.
Suppose there are an off-line pair \(p_*\), constants \(E,\eta>0\), and
all sufficiently large \(k\) such that
\[
 S_{p_*}(y_k)\ge E,
 \qquad
 S_q(y_k)\le S_{p_*}(y_k)-\eta
 \quad(q\ne p_*).
 \tag{V55.10}
\]
If
\[
 m_{\chi,t_k}(y_k)=0,
 \qquad
 \partial_xm_{\chi,t_k}(y_k)\ne0,
 \tag{V55.11}
\]
then
\[
 \boxed{
 \Im{\cal F}_{\chi,t_k}(y_k)\,
 \partial_xm_{\chi,t_k}(y_k)<0}
 \tag{V55.12}
\]
for all sufficiently large \(k\).
\end{theorem}

\begin{proof}
The exposure gap, the zero-counting shell estimate, and differentiated
Gaussian summation give
\[
 m_{\chi,t_k}(y_k)
 =2B_k\cos\theta_k
  +O(t_k^{-M}e^{-\eta/(8t_k)}B_k),
 \tag{V55.13}
\]
\[
 \partial_xm_{\chi,t_k}(y_k)
 =-\frac{\delta_*}{t_k}B_k\sin\theta_k
  +O(t_k^{-M}e^{-\eta/(8t_k)}B_k),
 \tag{V55.14}
\]
where \(B_k=B_{p_*,t_k}(y_k)\) and
\(\theta_k=\theta_{p_*,t_k}(y_k)\).  In
\textup{(V55.14)}, the derivative of the Gaussian amplitude is multiplied
by \(\cos\theta_k\) and is absorbed using \textup{(V55.13)}.

Apply the positive-frequency projection packetwise.  The one-pair estimate
\textup{(V52.6)}, the same exposure gap, and
\(B_k\ge c\,t_k^{-1/2}e^{E/(4t_k)}\) give
\[
 {\cal F}_{\chi,t_k}(y_k)
 =B_ke^{i\theta_k}
  +O(t_k^{-M}e^{-c/t_k}B_k)
 \tag{V55.15}
\]
for some \(c>0\).  The contact equation
\textup{(V55.13)} implies
\[
 \cos\theta_k=O(t_k^{-M}e^{-c/t_k}),
 \qquad
 |\sin\theta_k|=1+O(t_k^{-M}e^{-c/t_k}).
 \tag{V55.16}
\]
Multiplying \textup{(V55.14)} and the imaginary part of
\textup{(V55.15)} yields
\[
 \Im{\cal F}_{\chi,t_k}(y_k)\,
 \partial_xm_{\chi,t_k}(y_k)
 =-\frac{\delta_*}{t_k}B_k^2
  \left(1+O(t_k^{-M}e^{-c/t_k})\right)<0.
 \tag{V55.17}
\]
\end{proof}

The theorem is not restricted to the centered pair of V37.  It applies at
any compact contact sequence on which one positive-score packet has a
uniform action gap.

\subsection{Compact sign defects accumulate only at ties or zero action}

\begin{theorem}[Compact Hardy-defect localization]
\label{thm:v55-defect-localization}
Let \(t_k\downarrow0\) and \(y_k\in I\), where \(I\) is compact.  Suppose
the \(y_k\) are simple zero-level contacts and
\[
 \Im{\cal F}_{\chi,t_k}(y_k)\,
 \partial_xm_{\chi,t_k}(y_k)>0.
 \tag{V55.18}
\]
After passing to a subsequence, \(y_k\to y_\infty\), and at least one of
the following holds:
\[
 \boxed{
 D_\chi(y_\infty)\le0,
 \qquad\text{or}\qquad
 y_\infty\in{\cal T}_\chi^+.}
 \tag{V55.19}
\]
Equivalently, a compact sign-defect sequence cannot converge to a point
with a unique positive-action maximizer.
\end{theorem}

\begin{proof}
Compactness gives a convergent subsequence.  Suppose
\(D_\chi(y_\infty)>0\) and the maximizing pair \(p_*\) is unique.
Lemma~\ref{lem:v55-local-finite} gives \(E,\eta>0\) and a neighbourhood
of \(y_\infty\) on which
\[
 S_{p_*}\ge E,
 \qquad
 S_q\le S_{p_*}-\eta
 \quad(q\ne p_*).
 \tag{V55.20}
\]
Theorem~\ref{thm:v55-exposed-coherence} then makes the product in
\textup{(V55.18)} negative for all large \(k\), a contradiction.
\end{proof}

The alternative \(D_\chi(y_\infty)\le0\) is a genuine boundary of this
argument.  There the desired analytic half-wave is not exponentially larger
than its wrong-half leakage, so V52's projective extraction supplies no
relative asymptotic.

\subsection{Two positive frequencies can reverse the Hardy sign}

Tie localization alone is not an exclusion.  Consider the positive-frequency
caloric polynomial
\[
 F(t,x):=
 2i\,e^{-t}e^{ix}-i\,e^{-9t}e^{3ix},
 \qquad
 u(t,x):=2\Re F(t,x).
 \tag{V55.21}
\]
Both satisfy the heat equation.

\begin{proposition}[Two-wave Hardy sign-defect witness]
\label{prop:v55-two-wave-defect}
At \((t,x)=(0,0)\),
\[
 u=0,
 \qquad
 F=i,
 \qquad
 \partial_xF=1,
 \qquad
 \partial_xu=2.
 \tag{V55.22}
\]
Hence
\[
 \boxed{\Im F(0,0)\,\partial_xu(0,0)=2>0.}
 \tag{V55.23}
\]
Thus interference of two positive-frequency heat waves can realize a
Hardy sign defect.
\end{proposition}

\begin{proof}
At the indicated point, the two wave values are \(2i\) and \(-i\).
Their sum is \(i\), while
\[
 \partial_xF(0,0)
 =i\bigl(2i+3(-i)\bigr)=1.
 \tag{V55.24}
\]
Taking twice the real part gives the remaining identities.
\end{proof}

The witness does not reproduce Gaussian zero packets or Dirichlet
coefficients.  It proves the narrower and necessary point: positivity of
all Fourier frequencies does not determine the Hardy--slope sign once two
leading complex waves interfere.

\subsection{The global small-time trichotomy}

\begin{corollary}[Node, escape, or packet-degeneracy trichotomy]
\label{cor:v55-global-trichotomy}
Assume paired GRH fails and take the exposed times \(t_k\downarrow0\) from
V37, adjusted within their asymptotic windows so the full bottom zero slice
is regular.  Then, after passing to a subsequence, at least one of the
following occurs:
\begin{enumerate}[label=\textup{(\roman*)},leftmargin=3em]
\item an exposed continuation component contains a Hardy node
      \[
      P_{\chi,t}^{[0]}(x)=A_{\chi,t}(x);
      \tag{V55.25}
      \]
\item a sign-defect partner \(y_k\) satisfies
      \[
      |y_k|\longrightarrow\infty;
      \tag{V55.26}
      \]
\item the partners remain bounded and converge to \(y_\infty\) with
      \[
      D_\chi(y_\infty)\le0;
      \tag{V55.27}
      \]
\item the partners remain bounded and converge to a positive-action tie
      \[
      y_\infty\in{\cal T}_\chi^+.
      \tag{V55.28}
      \]
\end{enumerate}
\end{corollary}

\begin{proof}
Apply Theorem~\ref{thm:v54-node-defect}.  If its Hardy-node alternative
does not occur, select one of the resulting sign-defect partners.  Either
that sequence is unbounded, giving \textup{(V55.26)}, or it has a bounded
subsequence.  Theorem~\ref{thm:v55-defect-localization} gives
\textup{(V55.27)} or \textup{(V55.28)}.
\end{proof}

\begin{theorem}[V55 exposure-coherence verdict]
\label{thm:v55-verdict}
The positive-frequency completion supplies a strict local orientation at
every uniformly uniquely exposed positive-action zero:
\[
 \Im{\cal F}_{\chi,t}\,\partial_xm_{\chi,t}<0.
\]
Therefore the sign-defect branch of the V54 graph alternative is forced
out of all such regions.  Its only remaining asymptotic locations are
spatial infinity, the zero-action boundary, and the positive-action
tropical tie set.  The last location admits sign reversal for generic
two-wave analytic interference, so unique exposure alone does not exclude
the global partner exchange.
\end{theorem}

\begin{firewall}[Tropical localization is not arithmetic cancellation]
The tie set records equality of zero-side exponential scores.  It does not
control the multiplicities, phases, or lower-order contributions of the
tied packets, and Proposition~\ref{prop:v55-two-wave-defect} shows that
positive frequencies alone allow the wrong sign.  The escape and
zero-action alternatives are also untreated.  The trichotomy is necessary
under GRH failure, not contradictory, and no off-line zero is excluded.
\end{firewall}

\begin{assessment}[V55 closes single-packet sign defects]
V55 proves the strongest orientation statement available from the V52
analytic half-wave: a single uniformly exposed positive-action packet is
always Hardy coherent.  Combined with the global zero graph, this forces
any no-node continuation into escape or a degenerate score geometry.  The
new obstruction is therefore finite-dimensional on compact positive-action
sets: interference at explicit pairwise bisectors \textup{(V55.6)}.

V56 should analyze the tie geometry itself.  At a positive-action bisector,
retain the two leading Gaussian halves, impose the scalar zero equation,
and compute the exact sign-defect region in their amplitude ratio and phase
difference.  The key arithmetic question is whether equal zero scores also
fix those lower-order data through multiplicities and ordinates, or whether
the two-packet model contains an open sign-defect set.  In parallel, the
zero-action and escape branches should be recognized as separate regimes
rather than folded into the positive-action calculation.
\end{assessment}

\section*{Version V55 append-only record}
\addcontentsline{toc}{section}{Version V55 append-only record}
The complete V54 source was retained before this continuation.  It had
\(837680\) bytes, \(25127\) lines, and SHA--256
\[
 \texttt{7B3227FC3A3A1CD5716E2CFC8B92AB91E8EB9EE230446C6DADBC6329E0401A8B}.
\]
V55 proved Hardy coherence under unique positive-action exposure, localized
compact sign defects to the zero-action boundary or positive-action score
ties, exhibited a two-frequency sign-defect witness, and derived the global
node--escape--degeneracy trichotomy.

% =============================================================================
% END APPENDED V55 MATERIAL
% =============================================================================
% =============================================================================
% BEGIN APPENDED V56 MATERIAL
% =============================================================================

\section{V56: exact two-packet tie geometry and the open defect chamber}
\label{sec:grh-v56}

V55 proved that a compact positive-action Hardy sign defect must approach a
score tie.  This section resolves the leading two-packet geometry at such a
tie.  The microscopic coordinate is \(x=x_0+ty\): score differences are
linear in \(x\), so this is the unique scale on which two tied Gaussian
packets retain comparable amplitude.

After division by one packet amplitude, the completed analytic signal is
\[
 r e^{i\phi_1}+e^{i\phi_2},
\]
where \(r>0\) is an explicit exponential function of \(y\).  The scalar
contact equation removes its real part.  On that contact slice, the
Hardy--slope sign is an explicit quadratic expression in \(r\), the two
phases, displacements, and ordinate gap.  Its positive region is open and
nonempty for every pair of distinct off-line packets.  Thus equal leading
score does not inherit the single-packet coherence of V55.

What remains is not another local inequality.  The actual phases are tied
to the two variables \(1/t\) and \(y\).  Their continuous phase map has full
rank, but keeping \(y\) bounded while \(1/t\to\infty\) is a torus-recurrence
problem depending on the two zero displacements.  V56 therefore closes a
universal sign argument at ties and isolates the exact phase-access question.

\subsection{Microscopic normal form at a positive-action bisector}

Let \(p_j=(\gamma_j,\delta_j,m_j)\), \(j=1,2\), be distinct grouped
off-line pairs with \(\gamma_1\ne\gamma_2\).  Suppose that at \(x_0\)
\[
 S_{p_1}(x_0)=S_{p_2}(x_0)=E>0
 \tag{V56.1}
\]
and that there is \(\eta>0\) such that
\[
 S_q(x_0)\le E-\eta
 \qquad(q\notin\{p_1,p_2\}).
 \tag{V56.2}
\]
The equality point is the bisector \textup{(V55.6)}.  Put
\[
 \Delta\gamma:=\gamma_1-\gamma_2.
 \tag{V56.3}
\]
For \(x=x_0+ty\), the score difference is exactly
\[
 S_{p_1}(x)-S_{p_2}(x)=2t\,\Delta\gamma\,y.
 \tag{V56.4}
\]
Hence the ratio of the positive analytic amplitudes in
\textup{(V55.9)} is
\[
 \boxed{
 r(y):=\frac{B_{p_1,t}(x_0+ty)}
             {B_{p_2,t}(x_0+ty)}
 =\frac{m_1}{m_2}
  \exp\!\left(\frac{\Delta\gamma\,y}{2}\right).}
 \tag{V56.5}
\]
It is independent of \(t\) and ranges over all positive values as
\(y\) ranges over \(\mathbb R\).

Define
\[
 \theta_j(t,y)
 :=\frac{\delta_j(x_0-\gamma_j)}{2t}
   +\frac{\delta_jy}{2},
 \qquad j=1,2.
 \tag{V56.6}
\]

\begin{theorem}[Two-packet analytic normal form]
\label{thm:v56-two-packet-normal}
Uniformly for \(y\) in a fixed compact interval,
\[
 \boxed{
 \frac{{\cal F}_{\chi,t}(x_0+ty)}
      {B_{p_2,t}(x_0+ty)}
 =r(y)e^{i\theta_1(t,y)}
  +e^{i\theta_2(t,y)}
  +O(t^{-M}e^{-c/t})}
 \tag{V56.7}
\]
and
\[
 \boxed{
 \frac{t\,\partial_xm_{\chi,t}(x_0+ty)}
      {B_{p_2,t}(x_0+ty)}
 =
 \Re\!\left[
  a_1r(y)e^{i\theta_1(t,y)}
  +a_2e^{i\theta_2(t,y)}
 \right]
 +O(t^{-M}e^{-c/t}),}
 \tag{V56.8}
\]
where
\[
 a_j:=\gamma_j-x_0+i\delta_j.
 \tag{V56.9}
\]
The constants \(c>0\) and \(M<\infty\) may depend on the tie data and the
fixed \(y\)-interval.
\end{theorem}

\begin{proof}
Apply the packetwise analytic extraction
\textup{(V55.8)} to the two tied pairs.  On \(x=x_0+ty\), their amplitude
ratio is \textup{(V56.5)}.  The fixed score gap
\textup{(V56.2)} persists with gap at least \(\eta/2\) for bounded \(y\)
and small \(t\), so the zero-counting shell bound and wrong-half estimate
absorb every other contribution into the error in \textup{(V56.7)}.

For one analytic packet \(B_je^{i\theta_j}\),
\[
 \partial_x(B_je^{i\theta_j})
 =\frac{\gamma_j-x+i\delta_j}{2t}
  B_je^{i\theta_j}.
 \tag{V56.10}
\]
Since \(m=2\Re{\cal F}\), multiply the sum of
\textup{(V56.10)} by \(2t\), divide by \(B_{p_2,t}\), and replace
\(x=x_0+ty\) by \(x_0\) in the bounded coefficient.  The resulting
\(O(t)\) term is harmless and gives \textup{(V56.8)}.
\end{proof}

\subsection{The exact contact and sign-defect body}

Consider a sequence \(t_k\downarrow0\), \(y_k\to y\), along which
\[
 \theta_j(t_k,y_k)\longrightarrow\phi_j
 \pmod{2\pi},
 \qquad j=1,2.
 \tag{V56.11}
\]
Put \(r=r(y)\), and write \(c_j=\cos\phi_j\),
\(s_j=\sin\phi_j\).

\begin{theorem}[Exact leading sign on the two-packet contact slice]
\label{thm:v56-contact-sign}
Suppose \(x_k=x_0+t_ky_k\) are simple zero-level contacts.  Then their
limiting phases satisfy
\[
 \boxed{r c_1+c_2=0.}
 \tag{V56.12}
\]
Define
\[
 Y:=r s_1+s_2
 \tag{V56.13}
\]
and
\[
 \boxed{
 R:=r\Delta\gamma\,c_1-r\delta_1s_1-\delta_2s_2.}
 \tag{V56.14}
\]
If \(YR\ne0\), then for all large \(k\),
\[
 \boxed{
 \operatorname{sgn}\!\left(
 \Im{\cal F}_{\chi,t_k}(x_k)
 \,\partial_xm_{\chi,t_k}(x_k)\right)
 =\operatorname{sgn}(YR).}
 \tag{V56.15}
\]
In particular, the limiting contact is a Hardy sign defect precisely when
\[
 \boxed{YR>0.}
 \tag{V56.16}
\]
\end{theorem}

\begin{proof}
Take the real part of \textup{(V56.7)} and use
\(m=2\Re{\cal F}=0\) to obtain \textup{(V56.12)}.  The normalized
imaginary part tends to \(Y\).  Taking the real part in
\textup{(V56.8)} gives
\[
 \begin{aligned}
 &r(\gamma_1-x_0)c_1-r\delta_1s_1\\
 &\qquad+(\gamma_2-x_0)c_2-\delta_2s_2.
 \end{aligned}
 \tag{V56.17}
\]
Substitute \(c_2=-rc_1\) from \textup{(V56.12)}.  The ordinate terms
collapse to \(r(\gamma_1-\gamma_2)c_1\), yielding
\textup{(V56.14)}.  The positive normalizations
\(B_{p_2,t}\) and \(B_{p_2,t}/t\) do not affect the sign, which proves
\textup{(V56.15)}--\textup{(V56.16)}.
\end{proof}

For one packet, the \(R\)-factor is just
\(-\delta\sin\phi\), forcing \(YR<0\).  The new term
\(r\Delta\gamma c_1\) and the unequal displacement weights are precisely
what permit reversal at a tie.

\subsection{Every genuine two-packet tie has an open defect chamber}

\begin{theorem}[Nonempty open two-packet defect region]
\label{thm:v56-open-defect}
For every \(\delta_1,\delta_2>0\) and
\(\Delta\gamma\ne0\), the parameter set
\[
 \left\{(r,\phi_1,\phi_2):
 r>0,\ rc_1+c_2=0,\ YR>0\right\}
 \tag{V56.18}
\]
is nonempty and relatively open in the nondegenerate contact slice.
\end{theorem}

\begin{proof}
First suppose \(\delta_1\ne\delta_2\).  Choose opposite quadratures
\[
 \phi_1=\frac\pi2,
 \qquad
 \phi_2=-\frac\pi2.
 \tag{V56.19}
\]
The contact equation holds for every \(r>0\), and
\[
 YR=(r-1)(\delta_2-r\delta_1).
 \tag{V56.20}
\]
This is positive whenever \(r\) lies strictly between
\(1\) and \(\delta_2/\delta_1\).

Now suppose \(\delta_1=\delta_2=\delta\).  Take \(r=1\) and
\[
 \phi_2=\pi-\phi_1.
 \tag{V56.21}
\]
Then the contact equation holds and
\[
 YR
 =2\sin\phi_1
  \left(\Delta\gamma\cos\phi_1
        -2\delta\sin\phi_1\right).
 \tag{V56.22}
\]
Choose a sufficiently small nonzero \(\phi_1\) with the same sign as
\(\Delta\gamma\).  The expression is positive.  In both cases the
inequalities are strict, so continuity gives a relatively open defect
neighbourhood.
\end{proof}

Because \(r(y)\) in \textup{(V56.5)} ranges over \((0,\infty)\), the
amplitude coordinate itself never excludes these chambers.  Any exclusion
must come from the actual joint phase orbit.

\subsection{The remaining phase-access problem}

The two phases are affine functions of \(s=1/t\) and \(y\):
\[
 \begin{pmatrix}\theta_1\\\theta_2\end{pmatrix}
 =
 \frac12
 \begin{pmatrix}
 \delta_1(x_0-\gamma_1)&\delta_1\\
 \delta_2(x_0-\gamma_2)&\delta_2
 \end{pmatrix}
 \begin{pmatrix}s\\y\end{pmatrix}.
 \tag{V56.23}
\]
Their Jacobian determinant is
\[
 \boxed{
 \frac{\delta_1\delta_2(\gamma_2-\gamma_1)}4\ne0.}
 \tag{V56.24}
\]
Thus there is no local continuous rank loss: heat scale and microscopic
tie offset independently move the two phases.

The asymptotic problem is subtler.  One needs \(s\to+\infty\) while \(y\)
remains bounded so that \(r(y)\) stays in a fixed nondegenerate defect
interval.  Modulo \(2\pi\), this asks whether the affine phase strip
\[
 \left\{
 (\theta_1(s,y),\theta_2(s,y))\bmod2\pi:
 s\ge S,\ |y|\le C
 \right\}
 \tag{V56.25}
\]
meets the open contact-defect chamber for arbitrarily large \(S\).  Full
rank over \(\mathbb R^2\) does not by itself prove this bounded-\(y\)
recurrence.

\begin{theorem}[V56 tie-geometry verdict]
\label{thm:v56-verdict}
At an isolated positive-action tie of two off-line packets:
\begin{enumerate}[label=\textup{(\roman*)},leftmargin=3em]
\item \(x=x_0+ty\) is the exact comparability scale, with amplitude ratio
      \textup{(V56.5)};
\item the scalar contact and Hardy sign reduce to the explicit finite
      system \textup{(V56.12)}--\textup{(V56.16)};
\item the sign-defect region is open and nonempty for every genuine pair;
\item the phase map in \((1/t,y)\) has full rank;
\item realizing a defect along \(t\downarrow0\) with bounded \(y\) remains
      a torus-recurrence problem rather than a sign inequality.
\end{enumerate}
Hence two-packet tie geometry supplies no universal Hardy coherence and no
contradiction to the V55 tie alternative.
\end{theorem}

\begin{firewall}[An open phase chamber is not an arithmetic orbit hit]
Theorem~\ref{thm:v56-open-defect} varies the limiting amplitude and phases
as independent finite-dimensional parameters.  The actual zero packets
follow the affine orbit \textup{(V56.23)}, and the contact itself must be
solved with exponentially small remainders.  Full Jacobian rank does not
prove arbitrarily late intersection with the open defect chamber at bounded
microscopic offset.  No actual tie defect has been forced, no off-line zero
has been excluded, and GRH remains open.
\end{firewall}

\begin{assessment}[V56 closes universal coherence at score ties]
The V55 localization is now sharp at the two-packet level.  Unique exposure
forces coherence, but the first tie already contains a robust defect chamber.
The exact obstruction has moved from analysis to recurrence: whether the
one large heat phase and one bounded microscopic offset can repeatedly enter
that chamber while satisfying the scalar contact equation.

V57 should solve this bounded-offset phase-access problem without repeating
the finite-prime Kronecker argument of V8.  Eliminating \(y\) from
\textup{(V56.23)} leaves a one-parameter torus translation whose rotation
number is an explicit ratio of
\(\delta_j(x_0-\gamma_j)\) and \(\delta_j\).  The next step is to classify
the rational and irrational cases, determine whether every open defect
chamber is recurrent, and retain the scalar-contact transversality needed
to perturb phase hits to exact zeros.  A failure of recurrence would itself
be an explicit algebraic relation among the two off-line zero coordinates.
\end{assessment}

\section*{Version V56 append-only record}
\addcontentsline{toc}{section}{Version V56 append-only record}
The complete V55 source was retained before this continuation.  It had
\(850321\) bytes, \(25514\) lines, and SHA--256
\[
 \texttt{D670FFC8C0D074247DBEAEA465BF657D367003D77C8A128D29FC319D45B90F3C}.
\]
V56 derived the exact two-packet microscopic normal form, computed the
contact-slice Hardy sign, proved that every genuine positive-action tie has
an open sign-defect chamber, and isolated bounded-offset phase recurrence as
the remaining obstruction.

% =============================================================================
% END APPENDED V56 MATERIAL
% =============================================================================
% =============================================================================
% BEGIN APPENDED V57 MATERIAL
% =============================================================================

\section{V57: bounded-offset phase recurrence at a two-packet tie}
\label{sec:grh-v57}

V56 found an open Hardy sign-defect chamber in the leading two-packet
contact body but did not prove that the actual heat phases enter it with
bounded microscopic offset.  The missing transversality is already encoded
in the sign formula: the factor \(R\) in \textup{(V56.14)} is exactly twice
the derivative of the scalar contact equation with respect to the tie
coordinate \(y\).  Every strict defect point is therefore a transverse
contact.

This turns phase recurrence into an exact dichotomy.  If the two
\(1/t\)-carrier coefficients have irrational ratio, their continuous flow is
dense on the two-torus at every fixed \(y\).  Choose \(y\) so that the
amplitude ratio lies in one of the V56 defect chambers, approximate its
phases by arbitrarily late heat times, and use the nonzero \(y\)-derivative
to perturb each hit to an exact zero of the full symbol.  The defect sign
survives the exponentially smaller packets.

If the carrier ratio is rational, the phase orbit is a closed circle.
Every intersection of that circle with the defect chamber repeats
periodically at arbitrarily small heat times, but the circle can in
principle avoid the chamber.  Thus the only obstruction left at an isolated
two-packet tie is an explicit rational locking relation plus a compact
circle-avoidance condition.

\subsection{The contact derivative is the Hardy sign factor}

Retain the isolated positive-action tie data of
\textup{(V56.1)}--\textup{(V56.6)}.  Define the leading normalized analytic
signal
\[
 G(s,y):=
 r(y)e^{i\theta_1(s,y)}+e^{i\theta_2(s,y)},
 \qquad s:=\frac1t,
 \tag{V57.1}
\]
where
\[
 \theta_j(s,y)
 =A_js+\frac{\delta_jy}{2},
 \qquad
 A_j:=\frac{\delta_j(x_0-\gamma_j)}2.
 \tag{V57.2}
\]

\begin{lemma}[Defect transversality identity]
\label{lem:v57-transversality}
At a limiting contact
\[
 \Re G(s,y)=0,
 \tag{V57.3}
\]
with phases \(\phi_j=\theta_j(s,y)\), one has
\[
 \boxed{
 2\partial_y\Re G
 =r\Delta\gamma\cos\phi_1
  -r\delta_1\sin\phi_1
  -\delta_2\sin\phi_2
 =R.}
 \tag{V57.4}
\]
Consequently every strict Hardy defect \(YR>0\) is transverse in the
microscopic offset \(y\).
\end{lemma}

\begin{proof}
Differentiate \(r(y)\) from \textup{(V56.5)}:
\[
 r'(y)=\frac{\Delta\gamma}{2}r(y).
 \tag{V57.5}
\]
Together with \(\partial_y\theta_j=\delta_j/2\), this gives
\[
 2\partial_y\Re G
 =r\Delta\gamma\cos\phi_1
  -r\delta_1\sin\phi_1
  -\delta_2\sin\phi_2,
 \tag{V57.6}
\]
which is \textup{(V56.14)}.
\end{proof}

Thus the contact equation and sign condition are not competing
codimension-one constraints.  Positivity of \(YR\) automatically supplies
the derivative needed to correct a nearby phase hit to an exact contact.

\subsection{Irrational carrier slope forces recurrent defects}

Call the tie \emph{irrational} when
\[
 \boxed{\frac{A_1}{A_2}\notin\mathbb Q,}
 \tag{V57.7}
\]
with the usual convention that a zero denominator belongs to the rationally
locked case.  Equivalently, the continuous torus flow
\[
 s\longmapsto(A_1s,A_2s)\pmod{2\pi}
 \tag{V57.8}
\]
is dense in \(\mathbb T^2\).

\begin{theorem}[Recurrent exact sign defects at an irrational tie]
\label{thm:v57-irrational-defects}
Assume the isolated positive-action tie
\textup{(V56.1)}--\textup{(V56.2)} and the irrationality condition
\textup{(V57.7)}.  Then there are
\[
 t_k\downarrow0,
 \qquad
 y_k\longrightarrow y_0,
 \qquad
 x_k=x_0+t_ky_k,
 \tag{V57.9}
\]
such that
\[
 m_{\chi,t_k}(x_k)=0,
 \qquad
 \partial_xm_{\chi,t_k}(x_k)\ne0,
 \tag{V57.10}
\]
and
\[
 \boxed{
 \Im{\cal F}_{\chi,t_k}(x_k)\,
 \partial_xm_{\chi,t_k}(x_k)>0.}
 \tag{V57.11}
\]
Thus every irrational isolated two-packet tie actually generates
arbitrarily small-time Hardy sign defects.
\end{theorem}

\begin{proof}
By Theorem~\ref{thm:v56-open-defect}, choose \(y_0\) and phases
\((\phi_1,\phi_2)\) satisfying
\[
 r(y_0)\cos\phi_1+\cos\phi_2=0,
 \qquad
 YR>0.
 \tag{V57.12}
\]
For unequal displacements, choose \(r(y_0)\) strictly between
\(1\) and \(\delta_2/\delta_1\) and use the opposite quadratures
\textup{(V56.19)}.  For equal displacements, choose \(r(y_0)=1\) and the
small-angle phases in \textup{(V56.21)}--\textup{(V56.22)}.

Density of \textup{(V57.8)} gives \(s_k\to+\infty\) such that
\[
 \theta_j(s_k,y_0)\longrightarrow\phi_j
 \pmod{2\pi},
 \qquad j=1,2.
 \tag{V57.13}
\]
At the target point, Lemma~\ref{lem:v57-transversality} gives
\[
 \partial_y\Re G\ne0.
 \tag{V57.14}
\]
The implicit-function theorem therefore provides
\(\widetilde y_k\to y_0\) with
\[
 \Re G(s_k,\widetilde y_k)=0
 \tag{V57.15}
\]
and with the strict sign \(YR>0\) unchanged.

Set \(\widetilde t_k=s_k^{-1}\).  The two-packet normal form
\textup{(V56.7)}--\textup{(V56.8)} differs from \(G\), together with its
\(y\)-derivative, by \(O(t^{-M}e^{-c/t})\).  A second
implicit-function correction gives \(y_k-\widetilde y_k=
O(t_k^{-M}e^{-c/t_k})\) such that the \emph{full} scalar symbol vanishes.
Equation \textup{(V56.15)} preserves the positive sign and the nonzero
slope.  This proves \textup{(V57.9)}--\textup{(V57.11)}.
\end{proof}

The theorem is a feasibility result, not an exclusion: the irrational tie
realizes precisely the sign defects allowed by V54.

\subsection{Rational locking is a closed-orbit test}

Suppose now that \(A_1/A_2\in\mathbb Q\), including the cases in which one
carrier coefficient vanishes.  There is a least positive period \(L\), up
to an integer multiple, such that
\[
 (A_1L,A_2L)\in(2\pi\mathbb Z)^2.
 \tag{V57.16}
\]
For fixed \(y\), define the closed carrier orbit
\[
 {\cal O}_y
 :=\left\{
 (\theta_1(s,y),\theta_2(s,y))\pmod{2\pi}:
 0\le s<L
 \right\}.
 \tag{V57.17}
\]
Let \({\cal D}_{r(y)}\) be the strict contact-defect subset of the phase
torus defined by
\[
 r(y)\cos\phi_1+\cos\phi_2=0,
 \qquad
 YR>0.
 \tag{V57.18}
\]

\begin{theorem}[Exact rational recurrence criterion]
\label{thm:v57-rational-criterion}
In the rationally locked case, an isolated two-packet tie has a recurrent
bounded-offset strict defect sequence with a nondegenerate limiting sign
\(YR>0\) if and only if there is a finite \(y_0\) such that
\[
 \boxed{{\cal O}_{y_0}\cap{\cal D}_{r(y_0)}\ne\varnothing.}
 \tag{V57.19}
\]
Whenever \textup{(V57.19)} holds, the resulting leading contact repeats at
\[
 s_k=s_0+kL,
 \qquad
 t_k=s_k^{-1}\downarrow0,
 \tag{V57.20}
\]
and perturbs to exact full-symbol sign defects as in
Theorem~\ref{thm:v57-irrational-defects}.
\end{theorem}

\begin{proof}
If \textup{(V57.19)} holds, periodicity gives the identical leading phase
point at every \(s_0+kL\).  Strict defect implies \(R\ne0\) by
Lemma~\ref{lem:v57-transversality}, so the exponentially small full-symbol
remainder is removed by a \(y\)-correction, preserving the sign.

Conversely, any recurrent bounded-offset sequence has a convergent
subsequence \(y_k\to y_0\) and \(s_k\bmod L\to s_0\).  Passing to the limit
in the contact and strict-sign conditions puts the phase point in
\({\cal O}_{y_0}\cap\overline{{\cal D}_{r(y_0)}}\).  If the sequence has a
uniform strict sign margin after normalization, the limit lies in
\({\cal D}_{r(y_0)}\), proving \textup{(V57.19)}.  Degeneration to
\(YR=0\) is a boundary case rather than a recurrent strict chamber.
\end{proof}

Thus rationality does not itself prevent defects.  It replaces density by
the explicit compact intersection problem \textup{(V57.19)}.

\subsection{The arithmetic meaning of a locked tie}

Using \textup{(V57.2)}, rational locking is the relation
\[
 \boxed{
 \frac{\delta_1(x_0-\gamma_1)}
      {\delta_2(x_0-\gamma_2)}
 \in\mathbb Q.}
 \tag{V57.21}
\]
The tie equation simultaneously gives
\[
 \delta_1^2-(x_0-\gamma_1)^2
 =\delta_2^2-(x_0-\gamma_2)^2>0.
 \tag{V57.22}
\]
These are two explicit algebraic constraints on the reflected zero
coordinates and their bisector.  No known zero theorem forbids them.

\begin{corollary}[Refined V55 tie alternative]
\label{cor:v57-refined-tie}
In the positive-action tie branch of
Corollary~\ref{cor:v55-global-trichotomy}, an isolated two-packet limit has
the following exact status:
\begin{enumerate}[label=\textup{(\roman*)},leftmargin=3em]
\item if \textup{(V57.21)} is irrational, arbitrarily small-time sign
      defects occur near the tie;
\item if it is rational and \textup{(V57.19)} holds, periodic recurrent
      sign defects occur;
\item the only two-packet phase obstruction is rational locking together
      with avoidance of every strict defect chamber by the closed orbits
      \({\cal O}_y\).
\end{enumerate}
\end{corollary}

\begin{theorem}[V57 phase-recurrence verdict]
\label{thm:v57-verdict}
The bounded-offset recurrence problem at an isolated two-packet
positive-action tie is completely reduced to torus dynamics:
\begin{enumerate}[label=\textup{(\roman*)},leftmargin=3em]
\item strict Hardy defect automatically gives contact transversality;
\item irrational carrier slope forces exact recurrent full-symbol defects;
\item rational carrier slope is periodic and is decided by the explicit
      circle--chamber intersection \textup{(V57.19)};
\item neither case yields a universal contradiction.
\end{enumerate}
Accordingly the tie branch is generically feasible.  Any exclusion must
prove the exceptional rational relation and simultaneous chamber avoidance
for the actual zero pair, or introduce arithmetic information beyond the
two leading packets.
\end{theorem}

\begin{firewall}[Phase recurrence produces defects, not GRH]
Theorem~\ref{thm:v57-irrational-defects} works under the hypothetical
existence of an isolated positive-action tie of off-line packets.  Producing
the sign defects required by the global graph is consistent with GRH
failure and cannot exclude the assumed zeros.  The rational criterion is
also descriptive, not prohibitive.  No off-line zero is excluded and GRH
remains open.
\end{firewall}

\begin{assessment}[V57 closes the generic isolated-tie branch]
The open V56 defect chamber is dynamically accessible at every irrational
tie, and at every rational tie whose periodic phase circle meets it.
Because the sign factor is the contact derivative, no separate zero-solving
loss remains.  Thus the generic positive-action tie branch of the V55
trichotomy supports, rather than contradicts, the partner-exchange geometry.

V58 should stop extending the generic tie model.  The remaining branches
are structurally different: Hardy nodes \(P=A\), zero-action contacts, and
partners escaping to large ordinates.  The most upstream of these is the
Hardy node, because it is an exact complex equality at positive heat time.
The next calculation should differentiate \(P=A\) along the zero graph and
combine the fractional Hopf sign with the logarithmic reader moment
\(P^{[1]}\), testing whether the node has an arithmetic orientation not
present in the caloric countermodels.  If that identity again reduces to a
free Hilbert-transform jet, the Hardy continuation route should be closed.
\end{assessment}

\section*{Version V57 append-only record}
\addcontentsline{toc}{section}{Version V57 append-only record}
The complete V56 source was retained before this continuation.  It had
\(861718\) bytes, \(25866\) lines, and SHA--256
\[
 \texttt{1765C36A926F7F11D0C4D03185F343894F7F1A091E25C21639D4D6D8D80DEF03}.
\]
V57 identified the sign factor with microscopic contact transversality,
proved recurrent exact defects at every irrational two-packet tie, reduced
rational ties to an explicit periodic circle--chamber intersection, and
closed the generic isolated-tie branch as a source of exclusion.

% =============================================================================
% END APPENDED V57 MATERIAL
% =============================================================================
% =============================================================================
% BEGIN APPENDED V58 MATERIAL
% =============================================================================

\section{V58: Hardy-node reader jets and the fractional orientation no-go}
\label{sec:grh-v58}

V57 closed the generic isolated-tie branch and returned to the other trace
in the global zero graph: a Hardy node
\[
 {\cal F}_{\chi,t}(x)=A_{\chi,t}(x)-P_{\chi,t}^{[0]}(x)=0.
\]
This section computes its complete first spacetime orientation.

At every Hardy node the zeroth prime reader equals the full positive-frequency
archimedean completion, not merely its real part.  Spatial and heat
derivatives then identify the first and second reader jets.  If the node
occurs at a nondegenerate global threshold minimum, the fractional
half-Laplacian minimum principle forces a strict first-reader inequality and
a negative spacetime Jacobian.

These signs are exact but not new coercivity.  The first-reader inequality
is precisely the identity
\(-|D|m/2>0\) written on the prime side.  Periodic threshold flows attain it
for every Fourier carrier.  Away from the threshold minimum,
positive-frequency caloric polynomials realize both spacetime node degrees
at simple real zeros.  Thus the analytic-node orientation supplies no
additional obstruction beyond global nonnegativity at the threshold.

\subsection{Exact prime-reader identities at a Hardy node}

Recall
\[
 {\cal F}_{\chi,t}=A_{\chi,t}-P_{\chi,t}^{[0]},
 \qquad
 A_{\chi,t}:={\cal C}_+a_{\chi,t}.
 \tag{V58.1}
\]
The positive-frequency convention gives
\[
 A_{\chi,t}
 =\frac12\left(a_{\chi,t}+i{\cal H}a_{\chi,t}\right),
 \tag{V58.2}
\]
up to the fixed harmless distributional constant already normalized in
V53.  Since
\[
 \partial_xP_{\chi,t}^{[0]}=iP_{\chi,t}^{[1]},
 \qquad
 \partial_tP_{\chi,t}^{[0]}=-P_{\chi,t}^{[2]},
 \tag{V58.3}
\]
one has
\[
 \boxed{
 \partial_x{\cal F}_{\chi,t}
 =\partial_xA_{\chi,t}-iP_{\chi,t}^{[1]},}
 \tag{V58.4}
\]
\[
 \boxed{
 \partial_t{\cal F}_{\chi,t}
 =\partial_tA_{\chi,t}+P_{\chi,t}^{[2]}.}
 \tag{V58.5}
\]

\begin{theorem}[Complete zeroth--second reader system at a threshold node]
\label{thm:v58-reader-node}
Suppose \((t_*,x_*)\) is a nondegenerate global threshold minimum and a
Hardy node:
\[
 m_{\chi,t_*}\ge0,
 \quad
 m_{\chi,t_*}(x_*)=\partial_xm_{\chi,t_*}(x_*)=0,
 \quad
 \kappa:=\partial_x^2m_{\chi,t_*}(x_*)>0,
 \tag{V58.6}
\]
\[
 {\cal F}_{\chi,t_*}(x_*)=0.
 \tag{V58.7}
\]
Then, at \((t_*,x_*)\),
\[
 \boxed{P_{\chi,t_*}^{[0]}=A_{\chi,t_*},}
 \tag{V58.8}
\]
\[
 \boxed{
 \Im P_{\chi,t_*}^{[1]}
 =-\frac12\partial_xa_{\chi,t_*},}
 \tag{V58.9}
\]
\[
 \boxed{
 \Re P_{\chi,t_*}^{[1]}
 >\frac12|D|a_{\chi,t_*},}
 \tag{V58.10}
\]
and
\[
 \boxed{
 2\Re\left(
 P_{\chi,t_*}^{[2]}+\partial_tA_{\chi,t_*}
 \right)=\kappa>0.}
 \tag{V58.11}
\]
\end{theorem}

\begin{proof}
The node definition gives \textup{(V58.8)}.  Since
\(2\Re{\cal F}=m\), the stationary condition gives
\(\Re\partial_x{\cal F}=0\).  From \textup{(V58.2)} and
\textup{(V58.4)},
\[
 \Re\partial_x{\cal F}
 =\frac12\partial_xa_{\chi,t_*}
  +\Im P_{\chi,t_*}^{[1]},
 \tag{V58.12}
\]
which proves \textup{(V58.9)}.

The imaginary part of the same identity is
\[
 \Im\partial_x{\cal F}
 =\frac12|D|a_{\chi,t_*}
  -\Re P_{\chi,t_*}^{[1]}
 =\frac12|D|m_{\chi,t_*}.
 \tag{V58.13}
\]
Proposition~\ref{prop:v53-fractional-hopf} makes the last member strictly
negative, proving \textup{(V58.10)}.  Finally,
\(m_t=m_{xx}=\kappa\) at the contact and
\(m_t=2\Re\partial_t{\cal F}\).  Equation \textup{(V58.5)} proves
\textup{(V58.11)}.
\end{proof}

Equation \textup{(V58.10)} looks like a conditional first-moment bound, but
\textup{(V58.13)} shows that it is exactly the fractional minimum principle
for the already completed symbol.  It is not an independent estimate on the
Dirichlet phases.

\subsection{The spacetime degree of a Hardy node}

For any isolated node define
\[
 {\mathfrak J}_{\chi}(t,x)
 :=\det D_{(t,x)}
 \bigl(\Re{\cal F}_{\chi,t}(x),
       \Im{\cal F}_{\chi,t}(x)\bigr).
 \tag{V58.14}
\]
Because
\[
 {\cal F}=\frac12(m+i{\cal H}m),
 \tag{V58.15}
\]
direct differentiation gives
\[
 \boxed{
 {\mathfrak J}_{\chi}
 =\frac14\left(
 m_t\,|D|m-m_x\,{\cal H}m_t
 \right).}
 \tag{V58.16}
\]

\begin{theorem}[Negative degree at a nondegenerate threshold node]
\label{thm:v58-threshold-degree}
Under \textup{(V58.6)}--\textup{(V58.7)},
\[
 \boxed{
 {\mathfrak J}_{\chi}(t_*,x_*)
 =\frac{\kappa}{4}|D|m_{\chi,t_*}(x_*)<0.}
 \tag{V58.17}
\]
In particular, the complex Hardy node is nondegenerate and has local
topological degree \(-1\).
\end{theorem}

\begin{proof}
At the threshold minimum,
\[
 m_x=0,
 \qquad
 m_t=m_{xx}=\kappa>0.
 \tag{V58.18}
\]
Substitute these identities into \textup{(V58.16)} and apply the strict
fractional Hopf inequality \textup{(V53.20)}.
\end{proof}

Thus a threshold Hardy node has a fixed spacetime orientation even though
the two incoming real zero branches have total signed spatial degree zero.

\subsection{Threshold carriers realize the negative degree exactly}

For each integer \(N\ge1\), consider on the circle
\[
 u_N(t,x):=1-e^{-N^2t}\cos(Nx)
 \tag{V58.19}
\]
and its positive-frequency completion
\[
 F_N(t,x):=
 \frac12\left(1-e^{-N^2t}e^{iNx}\right).
 \tag{V58.20}
\]

\begin{proposition}[Exact threshold-node carrier family]
\label{prop:v58-threshold-family}
The function \(u_N\) solves the heat equation and satisfies
\[
 u_N(t,x)\ge0
 \qquad(t\ge0),
 \tag{V58.21}
\]
with a nondegenerate threshold Hardy node at \((0,0)\).  There
\[
 \partial_x^2u_N=N^2,
 \qquad
 |D|u_N=-N,
 \tag{V58.22}
\]
and
\[
 \boxed{
 \det D_{(t,x)}(\Re F_N,\Im F_N)
 =-\frac{N^3}{4}.}
 \tag{V58.23}
\]
If the constant archimedean model is \(a=1\), then
\[
 P_N=\frac12e^{-N^2t}e^{iNx},
 \qquad
 P_N^{[1]}(0,0)=\frac N2,
 \tag{V58.24}
\]
which realizes \textup{(V58.10)} for every \(N\).
\end{proposition}

\begin{proof}
All statements follow by direct differentiation.  Positivity in
\textup{(V58.21)} uses \(e^{-N^2t}\le1\).  At the origin,
\[
 \partial_tF_N=\frac{N^2}{2},
 \qquad
 \partial_xF_N=-\frac{iN}{2},
 \tag{V58.25}
\]
whose real Jacobian is \(-N^3/4\).
\end{proof}

The family proves that arbitrarily large first-reader scale and the negative
node degree are fully compatible with an ordinary positive threshold.

\subsection{Simple branch nodes have no universal degree}

The threshold sign uses global nonnegativity.  Remove that hypothesis and
even the spacetime degree becomes free within finite positive-frequency
heat sums.

\begin{theorem}[Both simple-node degrees occur with positive frequencies]
\label{thm:v58-both-degrees}
Define
\[
 F_-(t,x)
 :=i e^{-t}e^{ix}-i e^{-4t}e^{2ix},
 \tag{V58.26}
\]
\[
 F_+(t,x)
 :=i\left(
 3e^{-t}e^{ix}
 -5e^{-4t}e^{2ix}
 +2e^{-9t}e^{3ix}
 \right).
 \tag{V58.27}
\]
Both are positive-frequency caloric polynomials.  At \((0,0)\),
\[
 F_-=F_+=0,
 \qquad
 \partial_xF_-=\partial_xF_+=1.
 \tag{V58.28}
\]
Thus \(u_\pm:=2\Re F_\pm\) have simple real zero-level crossings there.
Their spacetime degrees have opposite signs:
\[
 \boxed{
 \det D(\Re F_-,\Im F_-)=-3,
 \qquad
 \det D(\Re F_+,\Im F_+)=+1.}
 \tag{V58.29}
\]
\end{theorem}

\begin{proof}
For \(F_-\), direct differentiation gives
\[
 \partial_tF_-(0,0)=3i,
 \qquad
 \partial_xF_-(0,0)=1,
 \tag{V58.30}
\]
so the determinant is \(-3\).  For \(F_+\), the coefficient identities
\[
 3-5+2=0,
 \quad
 3-10+6=-1,
 \quad
 3-20+18=1
 \tag{V58.31}
\]
give
\[
 F_+(0,0)=0,
 \qquad
 \partial_xF_+(0,0)=1,
 \qquad
 \partial_tF_+(0,0)=-i.
 \tag{V58.32}
\]
Its determinant is \(+1\).
\end{proof}

Thus the positive-frequency support that recovered the V52 quadrature does
not orient a Hardy node on a generic simple zero branch.

\subsection{Exact status of the Hardy-node route}

\begin{theorem}[V58 Hardy-node verdict]
\label{thm:v58-verdict}
At a Hardy node of the Weil heat symbol:
\begin{enumerate}[label=\textup{(\roman*)},leftmargin=3em]
\item the zeroth reader obeys the full complex equality \(P^{[0]}=A\);
\item at a nondegenerate global threshold node, the first reader satisfies
      the strict inequality \textup{(V58.10)} and the spacetime node has
      degree \(-1\);
\item both statements are exact rewritings of the half-Laplacian minimum
      principle for \(m\);
\item periodic threshold carriers realize the full sign and arbitrary
      carrier scale;
\item at simple branch nodes, positive-frequency caloric sums realize both
      spacetime degrees.
\end{enumerate}
Therefore differentiating the Hardy-node equality creates no independent
arithmetic orientation capable of excluding a positive heat threshold.
\end{theorem}

\begin{firewall}[Fractional Hopf orientation is not a prime cancellation
estimate]
The strict reader inequality \textup{(V58.10)} follows only after inserting
the completed symbol into the nonlocal minimum principle; using it backward
would assume the threshold positivity one is trying to exploit.  The
caloric models do not reproduce Dirichlet coefficients, but they prove that
the Hardy projection, heat equation, node equality, and threshold geometry
are mutually compatible.  No off-line zero is excluded and GRH remains
open.
\end{firewall}

\begin{assessment}[V58 closes the local Hardy-node orientation route]
The full node jet is now explicit through the second reader.  Its only
forced sign at a threshold is exactly \(|D|m(x_*)<0\); away from a global
minimum, even the complex spacetime degree is unrestricted.  More reader
derivatives would again be differentiated Hilbert jets and would not add an
independent equation.

V59 should address the two regimes deliberately left outside the
positive-action analysis: compact contacts with \(D_\chi(x)\le0\) and
sign-defect partners escaping to \(|x|\to\infty\) as \(t\downarrow0\).
The first question is whether the zero-side Gaussian sum permits any
small-time real zero where every score is nonpositive.  The second is to
identify the correct joint scaling of \(x\) and \(t\) at spatial escape.
These regimes must be separated before the mandatory V60 sweep of all
parallel notes.
\end{assessment}

\section*{Version V58 append-only record}
\addcontentsline{toc}{section}{Version V58 append-only record}
The complete V57 source was retained before this continuation.  It had
\(873698\) bytes, \(26209\) lines, and SHA--256
\[
 \texttt{7CC7C0C8FC2717B0DB340B68A9A674DA5C9A5B8CA893BA9FF149ACD8D2C86B84}.
\]
V58 derived the complete first Hardy-node reader jet, proved the negative
fractional degree at a threshold node, realized it by arbitrary periodic
carriers, and showed by explicit positive-frequency heat polynomials that
simple branch nodes admit both spacetime orientations.

% =============================================================================
% END APPENDED V58 MATERIAL
% =============================================================================
% =============================================================================
% BEGIN APPENDED V59 MATERIAL
% =============================================================================

\section{V59: the zero-action Stokes tail and the escape ceiling}
\label{sec:grh-v59}

V58 closed the local Hardy-node orientation route.  Two branches of the V55
trichotomy remained deliberately separate from positive-action exposure:
compact contacts with nonpositive score and contacts whose ordinates escape
as \(t\downarrow0\).  They have different mechanisms.

Outside a single off-line packet's influence cone, its scalar Gaussian is
exponentially small.  Its positive-frequency completion is not: the wrong
Fourier half has an algebraic endpoint tail.  At the packet's alternating
cosine zeros, this nonlocal Hardy tail has a fixed side sign while the local
scalar slope alternates.  Both Hardy coherence and Hardy sign defect
therefore occur.  This explicitly explains why V55 required positive action.

For a full-symbol zero whose ordinate escapes, the archimedean logarithm is
bounded by the absolute prime saddle.  This yields
\[
 \log\log(3+|x|)
 \le\frac1{16t}+O_\chi(\log(1/t)).
\]
The ceiling is double exponential in \(1/t\), far too weak to compactify the
partners from V54.  Hence neither residual regime supplies a universal
exclusion.

\subsection{The analytic Stokes tail of one reflected pair}

Retain one off-line packet \(p=(\gamma,\delta,m_p)\), and put
\[
 d:=x-\gamma.
 \tag{V59.1}
\]
Its scalar contribution is
\[
 K_{p,t}(x)=
 \frac{4\pi m_p}{\sqrt{4\pi t}}
 e^{(\delta^2-d^2)/(4t)}
 \cos\!\left(\frac{\delta d}{2t}\right).
 \tag{V59.2}
\]
Let
\[
 F_{p,t}:={\cal C}_+K_{p,t}.
 \tag{V59.3}
\]

\begin{lemma}[Algebraic Hardy tail outside the influence cone]
\label{lem:v59-stokes-tail}
Fix \(d\) with \(|d|>\delta\).  Then
\[
 \boxed{
 F_{p,t}(\gamma+d)
 =i\,\frac{2m_p d}{d^2+\delta^2}+O_{p,d}(t)
  +O_{p,d}\!\left(
    t^{-1/2}e^{-(d^2-\delta^2)/(4t)}
  \right).}
 \tag{V59.4}
\]
In particular,
\[
 \boxed{
 \Im F_{p,t}(\gamma+d)
 \longrightarrow\frac{2m_p d}{d^2+\delta^2}.}
 \tag{V59.5}
\]
\end{lemma}

\begin{proof}
Split the cosine into its two complex Gaussian halves and write the
positive-frequency projection as the sum of half-line Fourier integrals
\[
 \int_0^\infty
 e^{-tu^2+(\delta+id)u}\,du,
 \qquad
 \int_0^\infty
 e^{-tu^2+(-\delta+id)u}\,du,
 \tag{V59.6}
\]
with the common normalization inherited from \textup{(V59.2)}.
Because \(|d|>\delta\), steepest-descent deformation places the Gaussian
saddle below the real exponential scale.  Repeated endpoint integration by
parts gives
\[
 \frac12\left(
 -\frac1{\delta+id}+\frac1{\delta-id}
 \right)
 =\frac{id}{d^2+\delta^2}
 \tag{V59.7}
\]
as the leading half-line contribution.  Restoring the packet normalization
gives the first term in \textup{(V59.4)}.  The next endpoint term is
\(O(t)\), while the displaced complex saddle has the displayed
negative-action exponential size.

Equivalently, the same constant follows from the Cauchy kernels at the two
complex centers:
\[
 \frac12{\cal H}K_{p,t}(\gamma+d)
 \longrightarrow
 m_p\left(\frac1{d-i\delta}+\frac1{d+i\delta}\right)
 =\frac{2m_pd}{d^2+\delta^2}.
 \tag{V59.8}
\]
\end{proof}

The scalar packet itself tends exponentially to zero at such a point, while
its analytic quadrature tends to a nonzero algebraic value.  This is the
wrong-half leakage that was negligible only under the positive-action
hypothesis of V55.

\subsection{Alternating zero-action Hardy defects}

Fix \(d>\delta\), and choose
\[
 t_k:=\frac{\delta d}{(2k+1)\pi},
 \qquad k\ge0.
 \tag{V59.9}
\]
Then
\[
 \frac{\delta d}{2t_k}
 =\left(k+\frac12\right)\pi,
 \tag{V59.10}
\]
so \(K_{p,t_k}(\gamma+d)=0\).

\begin{theorem}[One packet realizes both Hardy orientations at negative
action]
\label{thm:v59-alternating-defects}
At the exact zero-level points \((t_k,\gamma+d)\),
\[
 \operatorname{sgn}\Im F_{p,t_k}(\gamma+d)=+1
 \tag{V59.11}
\]
for all sufficiently large \(k\), while
\[
 \operatorname{sgn}\partial_xK_{p,t_k}(\gamma+d)
 =-(-1)^k.
 \tag{V59.12}
\]
Consequently
\[
 \boxed{
 \operatorname{sgn}\left(
 \Im F_{p,t_k}\,\partial_xK_{p,t_k}
 \right)=-(-1)^k.}
 \tag{V59.13}
\]
Infinitely many of these contacts are Hardy coherent and infinitely many
are Hardy sign defects.
\end{theorem}

\begin{proof}
The side sign in \textup{(V59.11)} follows from
Lemma~\ref{lem:v59-stokes-tail}.  At a cosine zero, differentiation of
\textup{(V59.2)} gives
\[
 \partial_xK_{p,t_k}(\gamma+d)
 =-\frac{\delta}{2t_k}
 \frac{4\pi m_p}{\sqrt{4\pi t_k}}
 e^{(\delta^2-d^2)/(4t_k)}
 \sin\!\left(\left(k+\frac12\right)\pi\right),
 \tag{V59.14}
\]
which proves \textup{(V59.12)}--\textup{(V59.13)}.
\end{proof}

This exact one-packet model is not a full-symbol zero: other zero packets can
dominate at the same point.  It proves that nonpositive action destroys the
universal Hardy orientation even before packet interference is introduced.

\subsection{A double-exponential ceiling for escaping contacts}

Recall the absolute prime majorant
\[
 R(t):=
 2\sum_{n\ge2}\frac{\Lambda(n)}{\sqrt n}
 e^{-t(\log n)^2}.
 \tag{V59.15}
\]
The prime number theorem and Laplace's method at
\(\log n=1/(4t)\) give
\[
 \boxed{
 \log R(t)=\frac1{16t}+O(\log(1/t)).}
 \tag{V59.16}
\]

\begin{theorem}[Universal spatial ceiling for zero-level contacts]
\label{thm:v59-escape-ceiling}
If
\[
 m_{\chi,t}(x)=0,
 \qquad 0<t\le1,
 \tag{V59.17}
\]
then
\[
 \boxed{
 \log(2+|x|)
 \le R(t)+C_\chi,}
 \tag{V59.18}
\]
and therefore
\[
 \boxed{
 \log\log(3+|x|)
 \le\frac1{16t}+O_\chi(\log(1/t)).}
 \tag{V59.19}
\]
Equivalently,
\[
 \limsup_{t\downarrow0}
 16t\,\log\log(3+|x(t)|)\le1
 \tag{V59.20}
\]
along every zero-level escape curve.
\end{theorem}

\begin{proof}
The contact equation gives
\[
 a_{\chi,t}(x)
 =2\Re P_{\chi,t}^{[0]}(x)
 \le R(t).
 \tag{V59.21}
\]
The uniform archimedean asymptotic \textup{(V21.4)} gives
\[
 a_{\chi,t}(x)\ge\log(2+|x|)-C_\chi
 \tag{V59.22}
\]
for \(0<t\le1\).  Combining these inequalities proves
\textup{(V59.18)}.  Take one more logarithm and apply
\textup{(V59.16)} to obtain \textup{(V59.19)}--\textup{(V59.20)}.
\end{proof}

The estimate allows
\[
 |x(t)|
 \le\exp\!\left(
 \exp\!\left(\frac1{16t}+O_\chi(\log(1/t))\right)
 \right),
 \tag{V59.23}
\]
so it gives no compactness on the ordinate scale relevant to the zero
graph.  Improving it would require cancellation in the actual character
sum at a contact, not its absolute saddle.

\subsection{Closure of the residual continuation regimes}

\begin{theorem}[V59 residual-regime verdict]
\label{thm:v59-verdict}
The two regimes outside V55's positive-action analysis have the following
status.
\begin{enumerate}[label=\textup{(\roman*)},leftmargin=3em]
\item At fixed negative action, the scalar packet is exponentially small
      but its Hardy completion has an algebraic Stokes tail.
\item Exact cosine zeros of one off-line packet alternate between coherent
      and defective Hardy--slope signs.
\item Escaping full-symbol contacts satisfy the universal ceiling
      \textup{(V59.19)}, but that ceiling is double exponential.
\item Excluding either regime requires signed arithmetic cancellation not
      present in the heat equation, zero graph, or one-packet geometry.
\end{enumerate}
Thus zero action and spatial escape do not repair the continuation attack.
\end{theorem}

\begin{firewall}[A one-packet Stokes defect is not a full arithmetic
contact]
Theorem~\ref{thm:v59-alternating-defects} isolates the analytic mechanism
that defeats Hardy coherence outside the influence cone.  It does not claim
that the full zero sum vanishes at those same points.  Conversely, the
escape ceiling uses the absolute prime majorant and is far too weak to
control actual partner ordinates.  No off-line zero is excluded and GRH
remains open.
\end{firewall}

\begin{assessment}[V59 completes the pre-sweep continuation audit]
The continuation programme has now been tested in every asymptotic regime:
unique positive action is coherent; positive-action ties generically realize
defects; Hardy nodes have no independent local orientation; nonpositive
action has an alternating algebraic Stokes tail; and spatial escape is
bounded only at double-exponential scale.  No generic heat--Hardy mechanism
forces the contradiction needed for GRH.

V60 is the mandatory parallel-note sweep.  It should inventory all notes
created or advanced since the V50 sweep, test their newest lemmas against
the node, tie, Stokes, and escape alternatives above, and select a genuinely
upstream arithmetic direction.  The sweep should not continue the now
closed generic contact-topology branch unless another programme supplies
new coefficient-sensitive input.
\end{assessment}

\section*{Version V59 append-only record}
\addcontentsline{toc}{section}{Version V59 append-only record}
The complete V58 source was retained before this continuation.  It had
\(884560\) bytes, \(26605\) lines, and SHA--256
\[
 \texttt{559E69F18A5AE8ABBC9AD3D8CF127F66D2773D94F3F0D5AEC495002F156694D5}.
\]
V59 computed the algebraic Hardy Stokes tail outside a packet's influence
cone, proved alternating single-packet sign defects at negative action, and
derived the double-exponential spatial ceiling for escaping full-symbol
contacts.

% =============================================================================
% END APPENDED V59 MATERIAL
% =============================================================================
% =============================================================================
% BEGIN APPENDED V60 MATERIAL
% =============================================================================

\section{V60: mandatory corpus sweep and the coefficient-sensitive reset}
\label{sec:grh-v60}

V51--V59 completed the projective, zero-graph, Hardy-node, packet-tie,
negative-action, and escape analyses proposed at the V50 cutoff.  The user
requested a fresh sweep of all parallel notes every ten versions.  This
section performs that sweep before choosing another direction.

The corpus changed concurrently during the first inventory: the rolling SM
file advanced from sixteenth-cycle V12 to V13 while hashes were being read.
That mixed snapshot was discarded.  The recorded cutoff below was accepted
only after filenames, byte counts, modification ticks, and hashes remained
unchanged across the complete pass.

\subsection{Stable cutoff and manifest}

The V60 cutoff is
\[
 \boxed{\texttt{2026-09-11T03:16:15.8841028Z}.}
 \tag{V60.1}
\]
At that instant, \texttt{latex\_notes} contained \(31\) non-GRH
\texttt{.tex} files.  For the canonical UTF--8 manifest with rows
\[
 \texttt{filename|byte-count|SHA--256}\backslash\texttt{n}
 \tag{V60.2}
\]
in ordinal filename order, the hash was
\[
 \boxed{\texttt{D4CF7A659386ECDD8362C956915299269FA4E927F1052B010BF3C198E4FE80EF}.}
 \tag{V60.3}
\]

To verify the nonduplication boundary, remove the three current successor
rows below and restore the V50 rolling rows SM V90 and YM V41 with their
recorded byte counts and hashes.  The reconstructed \(30\)-row manifest has
SHA--256
\[
 \boxed{\texttt{B5E9BC96AF2DE47F4367ED1D2D960944995F1FB327753AFFD899D817454F2D3B},}
 \tag{V60.4}
\]
exactly equal to \textup{(V50.3)}.  Thus the \(28\) archive rows retained at
V50 are still byte-identical, and only the following three files require
reading:
\begin{center}
\tiny
\begin{tabular}{p{0.48\textwidth}|r|p{0.33\textwidth}}
\textbf{successor}&\textbf{bytes}&\textbf{SHA--256}\\ \hline
\path{Renewal_SM_Einstein_Continuum_Research_Notes_V100_f15.tex}
 &988883&\texttt{9B4C6C8CE2E53B5270A95E99AD8CE7}\newline
 \texttt{E622CAD1645B007B37AE885DD0879CEDCE}\\
\path{Renewal_SM_Einstein_Continuum_Research_Notes_V13.tex}
 &132567&\texttt{ED2799DF0A0442230D71D514E2A0CFE}\newline
 \texttt{A156A0131C89D12513B76A35D98B6896F}\\
\path{Renewal_Yang_Mills_Mass_Gap_Research_Notes_V46.tex}
 &396051&\texttt{B1E5580A117A5BED92CDB4AC94DDEB6}\newline
 \texttt{E5992F93E18D5DB46F0A38FAC6D3E5469}
\end{tabular}
\end{center}

The first successor contains SM fifteenth-cycle V91--V100, the second
contains the complete new sixteenth-cycle V1--V13, and the third contains
YM V42--V46 beyond the V50 endpoint.  No intervening append-only section is
missing.

\subsection{Typed transfer screen}

The SM fifteenth-cycle delta has five potentially transferable themes.
V91--V94 develops conditional interacting Gibbs compilers, stable polynomial
references, and the rule that a positive covariant square must not be split
into separately estimated signed pieces.  V96--V97 proves local quotient
coercivity and determinant control on constant-rank strata, while showing
that rank jumps create singular interface data.  V98--V100 proves nuclear
propagation on compact regular quotients and assembles a conditional
multi-sector transfer theorem, with all noncompact, singular, and
interaction rows left explicit.

The SM sixteenth-cycle delta develops compact-volume TT heat nuclearity,
thermodynamic local covariance, Gaussian Newton-source determinants,
nonlinear changes of measure, moving pseudoinverses, weighted
Hilbert--Schmidt sandwiches, and strong-truncation transfer.  Its newest
V12--V13 result is especially precise: a Gaussian exponential moment has a
strict spectral reserve, while comparison with an interacting density
\(h\in L^p\) consumes the reserve by the exact factor
\[
 q=\frac{p}{p-1}.
 \tag{V60.5}
\]
The comparison is finite exactly when
\[
 2qs\|C^{1/2}GC^{1/2}\|_{\rm op}<1.
 \tag{V60.6}
\]
Normalization or relative entropy alone supplies no such \(p>1\) reserve.

The YM V42--V46 delta independently builds the oppositely oriented rational
chart atlas.  It certifies both response signs on paired regions reaching
two transverse quarter faces, but explicitly refuses to infer an
orientation by symmetry or transport.  The third coordinate and the full
torus remain open.

\begin{theorem}[V60 typed non-import verdict]
\label{thm:v60-nonimport}
None of the post-V50 companion results supplies a new GRH estimate.
\begin{enumerate}[label=\textup{(\roman*)},leftmargin=3em]
\item The interacting Gibbs and Newton compilers require positive densities,
      common carriers, and strict cutoff-uniform spectral margins.  The
      character factor
      \(\chi(n)e^{ix\log n}\) is a complex observable on the unsigned prime
      saddle, not a positive comparison density.
\item Applying the V13 \(L^p\) theorem after taking absolute values removes
      the character phase.  Its exact factor \(q\) then spends rather than
      creates the exponential margin missing at a GRH contact.
\item Strong--Hilbert--Schmidt truncation needs a fixed compact factor and a
      uniform higher moment.  The GRH obstruction moves with
      \(t\downarrow0\), and V51 showed that its normalized moment curve
      collapses while its projective direction lives at a diverging
      standardized scale.
\item The command not to split a positive covariant square has already been
      honored by retaining the complete Weil symbol \(m\).  Splitting
      \({\cal F}=A-P\) in V53--V59 was an exact analytic identity, not a
      positivity estimate; recombining it produces no new inequality.
\item Constant-rank quotient determinants are stable only away from rank
      jumps.  Positive-action packet ties are precisely interface points
      where the leading rank changes, and V56--V57 already computed their
      independent local data.
\item The independently constructed YM opposite atlas reinforces rather
      than repairs the Hardy-sign obstruction: response orientation cannot
      be transported by symmetry, and the two-packet tie body contains an
      open region of each relevant sign.
\end{enumerate}
\end{theorem}

\begin{proof}
Items \textup{(i)}--\textup{(iii)} compare the exact hypotheses of SM
V91, sixteenth-cycle V11, and V13 with the prime-saddle identities
\textup{(V49.1)}--\textup{(V51.10)}.  Items \textup{(iv)}--\textup{(vi)}
compare the SM V94 and V97 firewalls and the YM V42--V46 independent atlas
with Theorems~\ref{thm:v55-exposed-coherence},
\ref{thm:v56-open-defect}, and
\ref{thm:v57-irrational-defects}.  In every case a required hypothesis is
exactly the missing coefficient-sensitive cancellation, not a consequence
of the companion theorem.
\end{proof}

\subsection{What the completed contact programme established}

\begin{theorem}[V51--V59 branch ledger]
\label{thm:v60-branch-ledger}
Under a hypothetical failure of paired GRH, the contact-continuation
programme now has the following exhaustive proved ledger.
\begin{enumerate}[label=\textup{(\roman*)},leftmargin=3em]
\item A center-exposed off-line pair produces complex projective endpoint
      rays that reconstruct its displacement but do not exclude it.
\item The global real zero graph yields a Hardy-node or Hardy-sign-defect
      alternative; adjacent negative-lobe endpoints need not remain paired.
\item Unique positive-action exposure forces the coherent Hardy sign.
\item Positive-action packet ties contain open defect chambers, and every
      irrational isolated two-packet tie recurrently realizes exact defects.
\item A threshold Hardy node has negative fractional degree, but that sign
      is exactly the half-Laplacian minimum principle; simple nodes admit
      either degree.
\item At nonpositive action, the algebraic analytic Stokes tail makes
      single-packet zero orientations alternate.
\item Escaping contacts are bounded only by the double-exponential ceiling
      \textup{(V59.19)}.
\end{enumerate}
Thus every generic heat, topology, Hardy-projection, and finite-packet sign
route is either reconstructive or explicitly realizable under the forbidden
geometry.
\end{theorem}

\begin{proof}
Combine Theorems~\ref{thm:v52-projective-ray},
\ref{thm:v54-node-defect}, \ref{thm:v55-exposed-coherence},
\ref{thm:v56-open-defect}, \ref{thm:v57-irrational-defects},
\ref{thm:v58-threshold-degree}, \ref{thm:v58-both-degrees},
\ref{thm:v59-alternating-defects}, and
\ref{thm:v59-escape-ceiling}.
\end{proof}

\subsection{The upstream coefficient-sensitive reset}

The positive-frequency object used throughout the last nine versions is the
heat evolution of one fixed spectral distribution:
\[
 {\cal F}_{\chi,t}
 ={\cal C}_+m_{\chi,t}
 =e^{t\partial_x^2}{\cal F}_{\chi,0}.
 \tag{V60.7}
\]
On positive frequencies its spectral data are exactly
\[
 \widehat{\cal F}_{\chi,t}(u)
 =e^{-tu^2}\widehat{\cal F}_{\chi,0}(u),
 \qquad u>0,
 \tag{V60.8}
\]
with
\[
 {\cal F}_{\chi,t}=A_{\chi,t}-P_{\chi,t}^{[0]}.
 \tag{V60.9}
\]
All previous no-go models varied an arbitrary positive-frequency
distribution.  The only remaining place for rigidity is the actual
coefficient relation inside \(\widehat{\cal F}_{\chi,0}\): the completed
gamma distribution minus the von Mangoldt character atoms.

\begin{criterion}[Coefficient-sensitive Hardy target]
\label{crit:v60-spectral-target}
A genuinely new continuation argument must prove a property of the fixed
distribution \(\widehat{\cal F}_{\chi,0}\) that is not shared by arbitrary
positive-frequency heat data and that excludes at least one necessary trace
in Theorem~\ref{thm:v60-branch-ledger}.  Examples of admissible targets are:
\begin{enumerate}[label=\textup{(\alph*)},leftmargin=3em]
\item an outer or zero-free factorization preventing
      \({\cal F}_{\chi,t}(x)=0\) on the real boundary;
\item a coefficient-sensitive sign law for
      \(\Im{\cal F}_{\chi,t}\,m_x\) that survives packet ties;
\item a uniform cancellation estimate compactifying the escape partners
      below the ceiling \textup{(V59.19)}.
\end{enumerate}
Any proof using only heat evolution, positive-frequency support, absolute
Gaussian moments, or zero-set topology is ruled out by V51--V59.
\end{criterion}

\begin{theorem}[V60 sweep verdict]
\label{thm:v60-verdict}
The mandatory corpus sweep imports no companion theorem into the GRH proof.
It does, however, sharpen the strategic boundary:
\begin{enumerate}[label=\textup{(\roman*)},leftmargin=3em]
\item positive comparison and nuclear-transfer methods require a strict
      reserve absent at the contact;
\item quotient-interface and opposite-atlas results confirm that tie
      orientation must be computed independently;
\item the completed contact programme has exhausted all generic
      heat--Hardy mechanisms;
\item the next attack must act directly on the gamma-minus-von-Mangoldt
      positive-frequency distribution in \textup{(V60.8)}.
\end{enumerate}
\end{theorem}

\begin{firewall}[The spectral reset is a target, not a factorization theorem]
No outer factorization, zero-free boundary theorem, global Hardy sign law,
or escape compactification has been proved for
\({\cal F}_{\chi,0}\).  The companion notes provide disciplined sufficient
conditions for positive transfer problems, but none verifies those
conditions for the complex character distribution.  GRH remains open.
\end{firewall}

\begin{assessment}[V60 acquisition and V61 direction]
V60 fulfills the scheduled full-corpus audit at a hash-stable cutoff and
proves that the 28 unchanged rows exactly reconstruct the V50 manifest.
The three live successors contribute useful non-import criteria but no
arithmetic cancellation.  In particular, the exact \(L^p\) comparison price
confirms that replacing the character phase by a positive density cannot
create a missing exponential reserve.

V61 should study the fixed distribution in \textup{(V60.8)} through
Hardy/outer factorization rather than another contact model.  The first
task is exact and falsifiable: determine the canonical half-plane analytic
extension of \({\cal F}_{\chi,t}\), factor its inner and outer parts at one
positive heat time, and express every real Hardy node as a boundary zero or
singular inner factor.  The Euler-product region should then be used only
where absolutely convergent to test whether the von Mangoldt atoms impose a
minimum-phase property.  If arbitrary inner factors remain compatible with
the completed coefficient distribution, that route should be closed before
returning to direct character cancellation.
\end{assessment}

\section*{Version V60 append-only record}
\addcontentsline{toc}{section}{Version V60 append-only record}
The complete V59 source was retained before this continuation.  It had
\(894012\) bytes, \(26918\) lines, and SHA--256
\[
 \texttt{7E2C09D92FB7A0D506E231528BEEEC27B83BD88593886F22A4CAB67CFD507733}.
\]
V60 performed the mandatory stable corpus sweep, hash-gated all unchanged
rows against V50, screened the three live successor deltas, consolidated the
V51--V59 contact-branch ledger, and reset the next attack to the fixed
gamma-minus-von-Mangoldt positive-frequency distribution.

% =============================================================================
% END APPENDED V60 MATERIAL
% =============================================================================

% =============================================================================
% BEGIN APPENDED V61 MATERIAL
% =============================================================================

\section{V61: Hardy factorization, the Euler anchor, and the boundary-zero no-go}
\label{sec:grh-v61}

V60 proposed canonical factorization as the first coefficient-sensitive
test of the fixed positive-frequency distribution.  The test can be carried
out, and it corrects an important premise of that proposal.  After a harmless
nonvanishing outer weight, the heat-regularized signal belongs to the upper
half-plane Hardy class and therefore has a canonical inner--outer
factorization.  Nevertheless, outerness does not prevent a zero on the real
boundary.  Such a zero may be encoded entirely by the outer boundary modulus,
without a Blaschke zero or singular inner factor.  An exact one-frequency heat
solution shows that upper-half-plane zeros can reach the boundary at an outer
function and then leave the half-plane.

Throughout this section let
\[
 s(z):=\frac12-iz,
 \qquad z=x+iy,
 \qquad \Re s(z)=\frac12+y,
 \tag{V61.1}
\]
and retain the V19 normalization
\[
 A_{\chi,0}(z)
 =\frac12\log\frac{f}{\pi}
  +\frac12\psi\!\left(\frac{\frac12+\kappa-iz}{2}\right),
 \qquad
 P_{\chi,0}^{[0]}(z)
 =\sum_{n\ge2}\frac{\Lambda(n)\chi(n)}{\sqrt n}
    e^{iz\log n}.
 \tag{V61.2}
\]
For a principal channel the standard pole term is to be subtracted; this
does not affect the factorization obstruction isolated below.

\subsection{The exact Hardy object}

\begin{theorem}[Weighted Hardy-class extension]
\label{thm:v61-hardy-class}
Fix \(t>0\).  The positive-frequency heat signal
\[
 {\cal F}_{\chi,t}(z)
 =A_{\chi,t}(z)-P_{\chi,t}^{[0]}(z)
 =e^{t\partial_z^2}{\cal F}_{\chi,0}(z)
 \tag{V61.3}
\]
extends holomorphically to \(\mathbb C_+=\{\Im z>0\}\) and continuously to
the real boundary.  Moreover
\[
 G_{\chi,t}(z):=\frac{{\cal F}_{\chi,t}(z)}{(z+i)^2}
 \tag{V61.4}
\]
belongs to \(H^2(\mathbb C_+)\).  The denominator is zero-free on
\(\overline{\mathbb C_+}\), so \(G_{\chi,t}\) and
\({\cal F}_{\chi,t}\) have exactly the same real zeros, with the same
multiplicities.
\end{theorem}

\begin{proof}
On every positive frequency \(u\), the analytic extension multiplies the
spectral datum by \(e^{izu}=e^{ixu-yu}\).  The heat factor \(e^{-tu^2}\)
dominates every exponential-linear factor on horizontal strips.  The
von Mangoldt atomic series and the completed gamma distribution therefore
define a holomorphic function in \(\mathbb C_+\), with continuous boundary
values.  Stirling's formula for the digamma contribution and absolute
Gaussian convergence of the prime series give, uniformly for \(y\ge0\),
\[
 |{\cal F}_{\chi,t}(x+iy)|\ll_{\chi,t}1+\log(2+|x|).
 \tag{V61.5}
\]
Consequently
\[
 \sup_{y>0}\int_{\mathbb R}|G_{\chi,t}(x+iy)|^2\,dx
 \ll_{\chi,t}
 \int_{\mathbb R}\frac{(1+\log(2+|x|))^2}{(1+x^2)^2}\,dx<\infty.
 \tag{V61.6}
\]
This is the defining half-plane \(H^2\) bound.  Finally \(z+i\ne0\) for
\(\Im z\ge0\), proving the zero statement.
\end{proof}

By canonical factorization, a nonzero \(G_{\chi,t}\in H^2(\mathbb C_+)\)
has the form
\[
 G_{\chi,t}=B_{\chi,t}S_{\chi,t}O_{\chi,t},
 \tag{V61.7}
\]
where \(B\) is the Blaschke product of its upper-half-plane zeros, \(S\)
is singular inner, and \(O\) is outer.  The decisive point is that this
factorization describes interior zeros and almost-everywhere boundary
modulus.  It does not turn every isolated boundary zero into an inner
factor.

\subsection{The completed logarithmic-derivative anchor}

\begin{theorem}[Exact Euler-region identity]
\label{thm:v61-euler-anchor}
For \(\Im z>\frac12\), equivalently \(\Re s(z)>1\), absolute convergence of
the Euler product gives
\[
 P_{\chi,0}^{[0]}(z)
 =-\frac{L'}{L}(s(z),\chi).
 \tag{V61.8}
\]
Hence, with
\[
 \Lambda(s,\chi)
 =\left(\frac f\pi\right)^{(s+\kappa)/2}
   \Gamma\!\left(\frac{s+\kappa}{2}\right)L(s,\chi),
 \tag{V61.9}
\]
one has the exact identity
\[
 {\cal F}_{\chi,0}(z)
 =A_{\chi,0}(z)-P_{\chi,0}^{[0]}(z)
 =\frac{\Lambda'}{\Lambda}(s(z),\chi)
 \qquad (\Im z>\tfrac12),
 \tag{V61.10}
\]
with the usual pole subtraction in a principal channel.
\end{theorem}

\begin{proof}
Because \(e^{iz\log n}=n^{iz}=n^{-s(z)+1/2}\), the summand in
\textup{(V61.2)} equals \(\Lambda(n)\chi(n)n^{-s(z)}\).  The absolutely
convergent logarithmic derivative of the Euler product proves
\textup{(V61.8)}.  Differentiating \textup{(V61.9)} gives
\[
 \frac{\Lambda'}{\Lambda}(s,\chi)
 =\frac12\log\frac f\pi
  +\frac12\psi\!\left(\frac{s+\kappa}{2}\right)
  +\frac{L'}{L}(s,\chi),
\]
which is \textup{(V61.10)} after inserting \(s=s(z)\).
\end{proof}

Theorem~\ref{thm:v61-euler-anchor} is the desired coefficient-sensitive
anchor: it identifies the unheated analytic signal with a completed
logarithmic derivative in a zero-free Euler half-plane.  It does not by
itself determine the inner factor of the heat-regularized boundary signal.
The heat multiplier acts on Fourier frequency, not by composition or
multiplication in \(z\), and therefore need not preserve inner--outer type.

\subsection{An exact outer boundary crossing}

\begin{theorem}[Outer functions can carry Hardy nodes]
\label{thm:v61-outer-node}
Fix an integer \(N\ge1\) and a time \(t_*>0\).  The positive-frequency
caloric polynomial
\[
 H^{(N)}(t,z)
 :=1-e^{N^2(t_*-t)}e^{iNz}
 \tag{V61.11}
\]
satisfies \(\partial_tH^{(N)}=\partial_z^2H^{(N)}\).  Its zeros are
\[
 z_k(t)=\frac{2\pi k}{N}+iN(t_*-t),
 \qquad k\in\mathbb Z.
 \tag{V61.12}
\]
Thus the zeros lie in \(\mathbb C_+\) for \(t<t_*\), meet the real boundary
at \(t=t_*\), and leave \(\overline{\mathbb C_+}\) for \(t>t_*\).
At the crossing time,
\[
 H^{(N)}(t_*,z)=1-e^{iNz}
 \tag{V61.13}
\]
is outer in \(H^\infty(\mathbb C_+)\), despite its infinitely many simple
real zeros.  After multiplication by \((z+i)^{-2}\), it is outer in
\(H^2(\mathbb C_+)\) and retains those zeros.
\end{theorem}

\begin{proof}
Direct differentiation proves the heat equation.  Solving
\(e^{N^2(t_*-t)}e^{iNz}=1\) gives \textup{(V61.12)}.  At \(t=t_*\), put
\(w=e^{iz}\); then \(|w|<1\) on \(\mathbb C_+\), and
\(1-e^{iNz}=1-w^N\).  The disk polynomial \(1-w^N\) has no zeros in the
open disk, has no singular inner factor, and its boundary logarithm is
integrable.  Its canonical factorization is therefore outer.  The standard
half-plane weight \((z+i)^{-2}\) is also outer and nonvanishing on the
closed upper half-plane, so the final assertion follows.
\end{proof}

This example has exactly the structural data available in
Theorem~\ref{thm:v61-hardy-class}: positive frequencies, heat evolution,
Hardy membership after the same weight, and even outerness at the instant
of contact.  Therefore a proof that \(G_{\chi,t}\) is outer would not
exclude \({\cal F}_{\chi,t}(x)=0\).

\subsection{What Poisson--Jensen can and cannot charge}

\begin{proposition}[The boundary-node budget]
\label{prop:v61-poisson-budget}
Let \(G\in H^2(\mathbb C_+)\) be nonzero, with factorization \(G=BSO\).
At an interior anchor \(z_0=x_0+iy_0\), the modulus identity has the form
\[
 \log|G(z_0)|
 =\log|B(z_0)|+\log|S(z_0)|
  +\frac1\pi\int_{\mathbb R}
    \frac{y_0\log|G(x)|}{(x-x_0)^2+y_0^2}\,dx.
 \tag{V61.14}
\]
An isolated boundary zero of finite order contributes locally
\(m\log|x-x_*|+O(1)\) to \(\log|G(x)|\), which is integrable.  Hence a
finite value at an Euler-region anchor does not prohibit such a boundary
zero.  Exclusion requires additional quantitative control of the negative
boundary logarithm, or an analytic reciprocal near every possible contact.
\end{proposition}

\begin{proof}
Equation \textup{(V61.14)} is the canonical factorization modulus formula;
the inner factors contribute through their values at \(z_0\), while the
outer factor is the Poisson integral of the boundary logarithm.  Since
\(\int_{-\varepsilon}^{\varepsilon}|\log|u||\,du<\infty\), a finite-order
boundary zero has finite Poisson cost at every fixed interior point.
Therefore neither an interior value nor outerness alone supplies a positive
lower boundary bound.  If instead \(1/G\) extends analytically and remains
locally bounded across a boundary interval, then continuity gives a
strictly positive lower bound there, which does exclude a node.
\end{proof}

The proposition identifies the missing arithmetic estimate exactly.  One
would need, uniformly up to the putative first-contact location, a bound
such as
\[
 \int_I \log^-|G_{\chi,t}(x)|\,dx<C_I
 \quad\hbox{with a scale-sensitive reserve excluding concentration,}
 \tag{V61.15}
\]
or local Smirnov/Hardy control of \(1/G_{\chi,t}\).  Mere finiteness of the
Szegő integral is insufficient, because the logarithmic singularity of a
node is integrable.

\begin{theorem}[V61 factorization verdict]
\label{thm:v61-verdict}
The first fixed-distribution factorization attack has the following exact
outcome.
\begin{enumerate}[label=\textup{(\roman*)},leftmargin=3em]
\item The positive-time arithmetic signal becomes an \(H^2(\mathbb C_+)\)
      function after the nonvanishing outer weight \((z+i)^{-2}\).
\item In the absolutely convergent Euler region, the unheated signal is
      exactly \(\Lambda'/\Lambda(\frac12-iz,\chi)\).
\item Canonical outerness cannot exclude real Hardy nodes: boundary zeros
      of finite order may belong entirely to the outer modulus.
\item Heat flow can move Blaschke zeros to an outer boundary node and then
      out of the upper half-plane, as \textup{(V61.11)} proves exactly.
\item The Euler anchor can exclude a boundary node only after a new
      coefficient-sensitive estimate controls the negative boundary
      logarithm or the analytic reciprocal.
\end{enumerate}
Thus the minimum-phase version of criterion
\ref{crit:v60-spectral-target}\textup{(a)} is false as stated.  Its valid
replacement is a quantitative reciprocal or negative-logarithm criterion.
\end{theorem}

\begin{firewall}[Factorization does not prove GRH]
No estimate in this section controls
\(\log^-|G_{\chi,t}|\) uniformly near a possible real contact, and no
analytic continuation theorem for \(1/G_{\chi,t}\) has been proved.  The
Euler identity \textup{(V61.10)} holds at \(t=0\) and \(\Im z>1/2\); it
cannot be transported to the real boundary through the heat semigroup by
formal inner--outer arguments.  GRH remains open.
\end{firewall}

\begin{assessment}[V61 closure and V62 direction]
V61 completes the first falsifiable task set by V60.  It constructs the
canonical Hardy object and the exact Euler anchor, then closes the naive
outer-factor route with an explicit caloric counterexample.  This is a
genuine narrowing: future factorization arguments must pay for a boundary
node quantitatively rather than merely classify the function as outer.

V62 should formulate the sharp Poisson--Jensen negative-log budget for
\(G_{\chi,t}\).  The first test is whether the Euler anchor together with
coefficient norms can control concentration of \(\log^-|G_{\chi,t}|\) on
shrinking boundary intervals.  If the exact heat-crossing family can match
all such norm data while concentrating one integrable logarithmic defect,
that coefficient-norm route should be closed; otherwise the surviving
scale-sensitive reserve is the next arithmetic target.
\end{assessment}

\section*{Version V61 append-only record}
\addcontentsline{toc}{section}{Version V61 append-only record}
The complete V60 source was retained before this continuation.  It had
\(907404\) bytes, \(27207\) lines, and SHA--256
\[
 \texttt{9911776AC8295800DC171D99ABDF55410268767068951904456733D41DA79CD6}.
\]
V61 weighted the positive-time analytic signal into
\(H^2(\mathbb C_+)\), identified its unheated Euler-region value with the
completed logarithmic derivative, and proved by an exact caloric polynomial
that outerness does not exclude real Hardy nodes.  The surviving target is
a scale-sensitive negative-logarithm or reciprocal estimate.

% =============================================================================
% END APPENDED V61 MATERIAL
% =============================================================================

% =============================================================================
% BEGIN APPENDED V62 MATERIAL
% =============================================================================

\section{V62: the corrected Hardy bound and the vanishing-cost outer node}
\label{sec:grh-v62}

V61 identified the right Hardy factorization obstruction, but its displayed
growth estimate \textup{(V61.5)} suppressed the dependence on the imaginary
coordinate.  The Hardy conclusion is correct; the uniform norm proof needs
the radial logarithmic bound recorded below.  This version supplies that
repair and then proves a sharper no-go: even a boundary node whose negative
logarithm has arbitrarily small local mass can be invisible to upper Hardy
norms and finitely many interior anchors.

\begin{lemma}[Correction to the V61 Hardy estimate]
\label{lem:v62-hardy-repair}
For fixed \(t>0\), the estimate used in Theorem
\ref{thm:v61-hardy-class} is
\[
 |{\cal F}_{\chi,t}(x+iy)|
 \ll_{\chi,t}1+\log\!\left(2+\sqrt{x^2+y^2}\right),
 \qquad y\ge0,
 \tag{V62.1}
\]
rather than the \(x\)-only right side printed in \textup{(V61.5)}.  This
still implies
\[
 \sup_{y>0}\int_{\mathbb R}
 \left|\frac{{\cal F}_{\chi,t}(x+iy)}{(x+i(y+1))^2}\right|^2dx<\infty.
 \tag{V62.2}
\]
Consequently the membership
\(G_{\chi,t}\in H^2(\mathbb C_+)\) and all factorization conclusions of V61
remain valid.
\end{lemma}

\begin{proof}
The Gaussian-smoothed prime atoms are absolutely controlled on horizontal
strips, and Stirling's formula gives the radial logarithm in
\textup{(V62.1)}.  Put \(a=y+1\ge1\) and substitute \(x=au\).  Since
\[
 \log(2+\sqrt{x^2+y^2})
 \ll \log(2+a)+\log(2+|u|),
\]
the integral in \textup{(V62.2)} is bounded by a constant times
\[
 a^{-3}\int_{\mathbb R}
 \frac{\{1+\log(2+a)+\log(2+|u|)\}^2}{(1+u^2)^2}\,du.
 \tag{V62.3}
\]
The integral in \(u\) is finite, and
\(a^{-3}\{1+\log(2+a)\}^2\) is bounded on \([1,\infty)\).  This proves the
uniform Hardy norm.
\end{proof}

\subsection{A localized boundary node with vanishing logarithmic charge}

\begin{theorem}[Finite anchor and norm data cannot see a boundary node]
\label{thm:v62-vanishing-node}
Fix \(x_*\in\mathbb R\).  For \(a>0\), define
\[
 E_{a,x_*}(z):=\frac{z-x_*}{z-x_*+ia}.
 \tag{V62.4}
\]
Then:
\begin{enumerate}[label=\textup{(\roman*)},leftmargin=3em]
\item \(E_{a,x_*}\) is outer in \(H^\infty(\mathbb C_+)\), satisfies
      \(\|E_{a,x_*}\|_\infty\le1\), and has the real boundary zero
      \(E_{a,x_*}(x_*)=0\).
\item For every compact \(K\subset\overline{\mathbb C_+}\setminus\{x_*\}\),
      \(E_{a,x_*}\to1\) uniformly on \(K\) as \(a\downarrow0\).
      In particular, every finite list of interior anchor values converges
      to the corresponding values of the zero-free constant function.
\item For every fixed \(R>0\), the local negative logarithmic charge is
\[
 \int_{x_*-R}^{x_*+R}\log^-|E_{a,x_*}(x)|\,dx
 =a\int_{-R/a}^{R/a}\frac12\log(1+u^{-2})\,du
 =O(a).
 \tag{V62.5}
\]
\item The weighted functions
      \(Q_{a,x_*}(z)=E_{a,x_*}(z)(z+i)^{-2}\) are outer in
      \(H^2(\mathbb C_+)\), have uniformly bounded \(H^2\) norms, retain
      the node at \(x_*\), and converge at every interior anchor to
      \((z+i)^{-2}\).
\end{enumerate}
Thus no criterion based only on upper Hardy norms, finitely many interior
values, or mere finiteness of a negative-log integral can exclude a real
Hardy node.
\end{theorem}

\begin{proof}
The only pole of \(E_{a,x_*}\) is \(x_*-ia\), below the boundary, and its
only zero is the boundary point \(x_*\).  Moreover
\[
 |E_{a,x_*}(x+iy)|^2
 =\frac{(x-x_*)^2+y^2}{(x-x_*)^2+(y+a)^2}\le1.
 \tag{V62.6}
\]
It is a zero-free bounded analytic rational function in the open upper
half-plane, with no singular inner factor and limit one at infinity; hence
it is outer.  Uniform convergence off \(x_*\) follows directly from
\[
 E_{a,x_*}(z)-1=\frac{-ia}{z-x_*+ia}.
 \tag{V62.7}
\]
On the boundary,
\[
 \log^-|E_{a,x_*}(x)|
 =\frac12\log\!\left(1+\frac{a^2}{(x-x_*)^2}\right).
 \tag{V62.8}
\]
The substitution \(x-x_*=au\) proves \textup{(V62.5)}, because
\(\log(1+u^{-2})\) is integrable on \(\mathbb R\).  Finally
\(|Q_{a,x_*}(x+iy)|\le|x+i(y+1)|^{-2}\), giving a uniform \(H^2\) bound;
products of outer factors are outer here.
\end{proof}

\begin{corollary}[The scale-sensitive reserve is indispensable]
\label{cor:v62-scale-reserve}
An Euler-region anchor value, an \(H^2\) upper bound, and a finite or even
vanishing local negative-log mass do not yield a pointwise lower bound at
the real boundary.  Any successful use of
Theorem~\ref{thm:v61-euler-anchor} must control the concentration scale,
for example through a lower bound on a harmonic-measure average at every
shrinking interval coupled to a coefficient-sensitive derivative bound, or
through locally uniform analytic control of \(1/G_{\chi,t}\).
\end{corollary}

\begin{proof}
The family in Theorem~\ref{thm:v62-vanishing-node} satisfies all three
unscaled data with constants converging to the zero-free model while
retaining an exact node for every \(a>0\).  Therefore an excluding estimate
must distinguish the shrinking transition layer of width \(a\).
\end{proof}

\begin{firewall}[The rational model is a Hardy no-go, not arithmetic data]
The functions \(E_{a,x_*}\) are positive-frequency boundary functions, but
they have not been represented by the fixed gamma-minus-von-Mangoldt
distribution or its heat orbit.  They disprove only arguments using the
listed generic Hardy data.  A coefficient-sensitive scale estimate may
still separate the arithmetic signal from this model.  GRH remains open.
\end{firewall}

\begin{assessment}[V62 closure and V63 direction]
V62 repairs the sole imprecise bound in V61 and sharpens its no-go from an
isolated example to a concentrating outer family.  The next admissible
question is now narrower: can the actual smoothed von Mangoldt coefficients
bound the spatial scale on which \({\cal F}_{\chi,t}\) approaches zero?
V63 should derive the strongest coefficient-level Bernstein or logarithmic
derivative estimate available at fixed \(t\), compare its cost as
\(t\downarrow0\) with the V59 escape ceiling, and test whether it supplies
the missing scale-sensitive reserve.
\end{assessment}

\section*{Version V62 append-only record}
\addcontentsline{toc}{section}{Version V62 append-only record}
The complete V61 source was retained before this continuation.  It had
\(919482\) bytes, \(27508\) lines, and SHA--256
\[
 \texttt{7913FD27ABCFF1C6AF8EF6C250BB3B319EA760660AAA04719A8A9AD29A217C9C}.
\]
V62 corrected the radial growth estimate underlying the Hardy norm and
proved that localized outer boundary nodes can have uniformly bounded
Hardy norms, asymptotically unchanged interior anchors, and vanishing local
negative-logarithm charge.  The surviving target is necessarily
coefficient-sensitive and scale-sensitive.

% =============================================================================
% END APPENDED V62 MATERIAL
% =============================================================================

% =============================================================================
% BEGIN APPENDED V63 MATERIAL
% =============================================================================

\section{V63: the coefficient Bernstein scale and its exact no-go}
\label{sec:grh-v63}

V62 showed that generic Hardy data cannot detect a boundary zero whose
transition layer concentrates.  The next possible input is the actual
Gaussian von Mangoldt coefficient envelope.  This version computes its
sharp derivative scale.  The result is exact but nonexcluding: successive
absolute derivative budgets cost \((4t)^{-1}\), so they resolve spatial
layers of order \(t\).  That is precisely the width of the exposed-pair
lobes already constructed in V27.  At the unnormalized level the first
derivative ceiling permits even exponentially thinner outer-node layers.

For every integer \(r\ge0\), define the absolute coefficient moment
\[
 R_r(t):=2\sum_{n\ge2}\frac{\Lambda(n)(\log n)^r}{\sqrt n}
 e^{-t(\log n)^2}.
 \tag{V63.1}
\]
Thus \(R_0=R\) in \textup{(V59.15)}.

\subsection{Sharp absolute moment asymptotics}

\begin{theorem}[The von Mangoldt Bernstein saddle]
\label{thm:v63-moment-saddle}
For every fixed integer \(r\ge0\), as \(t\downarrow0\),
\[
 \boxed{
 R_r(t)\sim
 2\sqrt\pi\,4^{-r}t^{-r-1/2}e^{1/(16t)}.}
 \tag{V63.2}
\]
In particular,
\[
 \log R_r(t)=\frac1{16t}+O_r(\log(1/t)),
 \qquad
 \boxed{\frac{R_{r+1}(t)}{R_r(t)}
 \sim\frac1{4t}.}
 \tag{V63.3}
\]
For all real \(x\),
\[
 \boxed{
 |\partial_x^rP_{\chi,t}^{[0]}(x)|
 =|P_{\chi,t}^{[r]}(x)|\le\frac12R_r(t).}
 \tag{V63.4}
\]
\end{theorem}

\begin{proof}
The prime number theorem and Stieltjes partial summation reduce
\(R_r(t)/2\), up to a relative \(o(1)\), to
\[
 \int_1^\infty y^{-1/2}(\log y)^r
 e^{-t(\log y)^2}\,dy
 =\int_0^\infty u^re^{-tu^2+u/2}\,du.
 \tag{V63.5}
\]
Complete the square:
\[
 -tu^2+\frac u2
 =\frac1{16t}-t\left(u-\frac1{4t}\right)^2.
 \tag{V63.6}
\]
The saddle \(u_t=(4t)^{-1}\) is a distance \(1/(4\sqrt t)\) from the
lower endpoint in Gaussian units.  Substitution
\(v=\sqrt t(u-u_t)\) therefore gives
\[
 \int_0^\infty u^re^{-tu^2+u/2}\,du
 \sim e^{1/(16t)}(4t)^{-r}t^{-1/2}
 \int_{-\infty}^{\infty}e^{-v^2}\,dv,
 \tag{V63.7}
\]
which is \textup{(V63.2)}.  Division proves the ratio in
\textup{(V63.3)}.  Finally termwise differentiation gives
\(\partial_x^rP^{[0]}=i^rP^{[r]}\), and the triangle inequality proves
\textup{(V63.4)}.
\end{proof}

The ratio in \textup{(V63.3)} is the sharp coefficient-level Bernstein
frequency.  Normalizing the prime reader by its absolute mass yields only
\[
 \frac{\|\partial_xP_{\chi,t}^{[0]}\|_\infty}
      {R_0(t)}\le\frac{R_1(t)}{2R_0(t)}
 \sim\frac1{8t}.
 \tag{V63.8}
\]
No character cancellation appears in this inequality.

\subsection{Comparison with the actual forbidden geometry}

\begin{proposition}[The Bernstein scale matches exposed-pair lobes]
\label{prop:v63-lobe-match}
Under the hypothetical exposed off-line pair of V27, the two adjacent
zero-level endpoints satisfy
\[
 x_{k,+}-x_{k,-}
 =\frac{2\pi t_k}{\delta_*}+o(t_k).
 \tag{V63.9}
\]
Hence their spatial scale is \(\Theta(t_k)\), exactly the reciprocal of the
absolute effective frequency in \textup{(V63.3)}.  Finite-order Bernstein
control obtained from the moments \(R_r\) is therefore compatible with the
forbidden off-line-zero geometry.
\end{proposition}

\begin{proof}
Equation \textup{(V63.9)} is Theorem
\ref{thm:v27-endpoint-asymptotics}.  The reciprocal consecutive-moment
ratio is \(R_r/R_{r+1}\sim4t\).  Both scales differ only by a fixed factor
depending on \(\delta_*\), so no asymptotic separation is available.
\end{proof}

This comparison also explains why taking more absolute derivatives does not
help.  For every fixed \(r\), all moments retain the same exponential type
\(e^{1/(16t)}\); differentiation changes only powers of \(t^{-1}\).  This
recovers the nonduplication boundary of V27--V31 from the sharper asymptotic
ratio rather than another covariance argument.

\subsection{A derivative-budget capacity witness}

Let
\[
 M(t):=1+R_1(t),
 \qquad a(t):=M(t)^{-1}.
 \tag{V63.10}
\]
For an arbitrary real center \(X(t)\), use the V62 outer factor
\[
 E_t(z):=\frac{z-X(t)}{z-X(t)+ia(t)}.
 \tag{V63.11}
\]

\begin{theorem}[The absolute derivative ceiling permits exact nodes]
\label{thm:v63-derivative-capacity}
The family \(E_t\) is outer in \(H^\infty(\mathbb C_+)\), has an exact
boundary zero at \(X(t)\), and satisfies
\[
 \|E_t\|_\infty\le1,
 \qquad
 \sup_{x\in\mathbb R}|E_t'(x)|=a(t)^{-1}=M(t).
 \tag{V63.12}
\]
Its transition width obeys
\[
 \boxed{
 \log a(t)=-\frac1{16t}+O(\log(1/t)).}
 \tag{V63.13}
\]
For every interior anchor whose Euclidean distance from \(X(t)\) is bounded
below,
\[
 E_t(z)=1+O(a(t)).
 \tag{V63.14}
\]
Thus an upper derivative budget of the exact absolute coefficient size does
not exclude a boundary node; it merely prevents its outer transition layer
from being thinner than an exponentially small scale.
\end{theorem}

\begin{proof}
The outer and node assertions follow from Theorem
\ref{thm:v62-vanishing-node}.  Differentiation gives
\[
 E_t'(z)=\frac{ia(t)}{(z-X(t)+ia(t))^2},
 \tag{V63.15}
\]
whose real-boundary maximum is \(a(t)^{-1}\) at \(x=X(t)\).
Theorem~\ref{thm:v63-moment-saddle} proves \textup{(V63.13)}, and
\textup{(V63.14)} follows from \textup{(V62.7)}.
\end{proof}

The witness is deliberately normalized in amplitude.  Multiplication by a
fixed bounded outer background preserves the node and the derivative
exponent, while its values at every fixed interior anchor converge to the
background values.  What is missing from an absolute ceiling is a
lower-modulus or phase relation tied to the actual character coefficients.

\subsection{No interaction with the escape ceiling}

\begin{corollary}[Translation defeats compactness]
\label{cor:v63-translation-escape}
The estimates \textup{(V63.12)}--\textup{(V63.14)} are invariant under the
choice of center \(X(t)\).  In particular one may choose \(X(t)\) anywhere
under the V59 allowance
\[
 \log\log(3+|X(t)|)
 \le\frac1{16t}+O(\log(1/t)).
 \tag{V63.16}
\]
Therefore coefficient Bernstein bounds do not improve the
double-exponential spatial ceiling or compactify escaping contacts.
\end{corollary}

\begin{proof}
Translation changes neither the supremum norms nor the transition width of
\(E_t\).  The displayed ceiling constrains only the center, while the
Bernstein estimate constrains only local variation.  There is no inequality
between them that bounds \(X(t)\) independently of \(t\).
\end{proof}

\begin{theorem}[V63 coefficient-scale verdict]
\label{thm:v63-verdict}
The full absolute von Mangoldt hierarchy supplies the sharp local frequency
scale \((4t)^{-1}\), but no GRH exclusion:
\begin{enumerate}[label=\textup{(\roman*)},leftmargin=3em]
\item every fixed derivative moment has exponential type \(1/16\);
\item its normalized spatial resolution is only order \(t\), matching the
      lobe width created by an exposed off-line zero;
\item its unnormalized first-derivative ceiling permits outer nodes on
      layers of width \(e^{-1/(16t)}\) up to powers of \(t\);
\item translation-invariant local bounds do not improve V59's escape
      ceiling.
\end{enumerate}
Any successful scale-sensitive reserve must therefore be conditional on the
complex contact equation and use signed character cancellation.
\end{theorem}

\begin{firewall}[The capacity witness is not the arithmetic signal]
The rational outer family \(E_t\) is not asserted to equal
\({\cal F}_{\chi,t}\) or to arise from its fixed heat orbit.  It proves that
the sharp absolute coefficient derivative data alone are insufficient.
The needed conditional signed estimate has not been proved, and GRH remains
open.
\end{firewall}

\begin{assessment}[V63 closure and V64 direction]
V63 completes the coefficient-level Bernstein test proposed in V62.  The
sharp saddle ratio yields exactly the spatial scale already realized by a
hypothetical off-line pair, so the route closes at the absolute level.

V64 should impose the full complex Hardy-node constraint
\(P_{\chi,t}^{[0]}(x)=A_{\chi,t}(x)\) before estimating
\(P_{\chi,t}^{[1]}(x)\).  The precise target is a conditional signed
value--derivative inequality whose exponential rate is strictly below
\(1/16\).  Its relation to the V27 real-contact regression must be made
explicit so that the stronger complex constraint either yields a genuine
saving or is proved equivalent to an already closed formulation.
\end{assessment}

\section*{Version V63 append-only record}
\addcontentsline{toc}{section}{Version V63 append-only record}
The complete V62 source was retained before this continuation.  It had
\(926516\) bytes, \(27681\) lines, and SHA--256
\[
 \texttt{39E9D19049CDA4823EF1C17D43FD15902CA4A460F5D16693265C3D362817BB86}.
\]
V63 computed the sharp absolute von Mangoldt moment asymptotics and
successive derivative ratio, matched that scale to exposed-pair lobes, and
proved that the resulting derivative ceiling neither excludes localized
outer nodes nor improves the double-exponential escape bound.

% =============================================================================
% END APPENDED V63 MATERIAL
% =============================================================================

% =============================================================================
% BEGIN APPENDED V64 MATERIAL
% =============================================================================

\section{V64: the complex-contact regression and its logical insufficiency}
\label{sec:grh-v64}

V63 closed every value-free coefficient Bernstein estimate and proposed
conditioning the first reader on the full Hardy-node equation
\[
 P_{\chi,t}^{[0]}(x)=A_{\chi,t}(x).
 \tag{V64.1}
\]
This is genuinely stronger than the real zero-level constraint, but two
distinct questions must not be conflated.  First, does the complex equality
itself force a derivative saving?  Second, would such a saving exclude the
global trace of an off-line zero?  The answer to both questions is negative
without additional arithmetic input.  A two-frequency positive-frequency
heat solution realizes arbitrary value and first reader at one point.
Moreover, the V54 zero graph allows the forbidden continuation to terminate
in Hardy sign defects without producing a Hardy node at all.

\subsection{The exact normalized complex regression}

Retain the absolute saddle quantities of V28,
\[
 Z_\chi(t)=\sum_{n\ge2}\frac{\Lambda(n)}{\sqrt n}
 e^{-t(\log n)^2},
 \quad
 \mu_t=\mathbb E_tU,
 \quad
 \sigma_t^2=\mathbb E_t(U-\mu_t)^2,
 \tag{V64.2}
\]
and define
\[
 z_t(x):=\frac{P_{\chi,t}^{[0]}(x)}{Z_\chi(t)},
 \qquad
 c_t(x):=
 \frac{P_{\chi,t}^{[1]}(x)-\mu_tP_{\chi,t}^{[0]}(x)}
      {\sigma_tZ_\chi(t)}.
 \tag{V64.3}
\]

\begin{proposition}[Exact Hardy-node regression ledger]
\label{prop:v64-complex-regression}
At every Hardy node \((t,x)\),
\[
 \boxed{
 z_t(x)=\frac{A_{\chi,t}(x)}{Z_\chi(t)},}
 \tag{V64.4}
\]
and
\[
 \boxed{
 P_{\chi,t}^{[1]}(x)
 =\mu_tA_{\chi,t}(x)+\sigma_tZ_\chi(t)c_t(x).}
 \tag{V64.5}
\]
Consequently
\[
 \boxed{
 \partial_x{\cal F}_{\chi,t}(x)
 =\partial_xA_{\chi,t}(x)
  -i\mu_tA_{\chi,t}(x)
  -i\sigma_tZ_\chi(t)c_t(x).}
 \tag{V64.6}
\]
The full complex value constraint fixes \(z_t\), but it gives no algebraic
condition on the centered first-reader coordinate \(c_t\).
\end{proposition}

\begin{proof}
Equation \textup{(V64.4)} is \textup{(V64.1)} divided by \(Z_\chi(t)\).
Solving the definition of \(c_t\) for \(P^{[1]}\) gives
\textup{(V64.5)}.  Insert that identity into
\(\partial_x{\cal F}=\partial_xA-iP^{[1]}\), which was proved in
\textup{(V58.4)}, to obtain \textup{(V64.6)}.
\end{proof}

On a fixed compact \(x\)-interval, V51 gives only
\[
 |c_t(x)|\le Ct^{-M}e^{-b/\sqrt t},
 \tag{V64.7}
\]
whereas
\[
 \sigma_tZ_\chi(t)
 =t^{-O(1)}e^{1/(16t)}.
 \tag{V64.8}
\]
Thus the known progression estimate inserted into \textup{(V64.6)} still
has exponential type \(1/16\).  To make the centered term subexponential
in \(1/t\), one would need the conditional full-exponential saving
\[
 \boxed{
 |c_t(x)|
 \le \exp\!\left(-\frac1{16t}+o(1/t)\right)
 \quad\hbox{whenever }z_t(x)=A_{\chi,t}(x)/Z_\chi(t).}
 \tag{V64.9}
\]
Neither the fixed-character PNT nor the absolute moment hierarchy approaches
this scale.

\subsection{Point interpolation leaves the first reader free}

\begin{theorem}[Two frequencies realize arbitrary complex value and
derivative]
\label{thm:v64-two-frequency-freedom}
Fix \(t_0>0\), \(x_0\in\mathbb R\), two distinct frequencies
\(0<\lambda_1<\lambda_2\), and arbitrary targets \(A_0,D\in\mathbb C\).
There is a positive-frequency heat solution
\[
 P(t,z)=\sum_{j=1}^2c_j
 e^{-\lambda_j^2(t-t_0)}e^{i\lambda_j(z-x_0)}
 \tag{V64.10}
\]
such that
\[
 \boxed{P(t_0,x_0)=A_0,
 \qquad P^{[1]}(t_0,x_0)=D.}
 \tag{V64.11}
\]
Explicitly,
\[
 c_1=\frac{\lambda_2A_0-D}{\lambda_2-\lambda_1},
 \qquad
 c_2=\frac{D-\lambda_1A_0}{\lambda_2-\lambda_1}.
 \tag{V64.12}
\]
Taking the constant archimedean background \(A(t,z)=A_0\), the signal
\({\cal F}=A-P\) has a Hardy node at \((t_0,x_0)\), while
\[
 \partial_x{\cal F}(t_0,x_0)=-iD
 \tag{V64.13}
\]
is arbitrary.
\end{theorem}

\begin{proof}
Each summand in \textup{(V64.10)} satisfies
\(P_t=P_{zz}\) and has positive frequency.  At \((t_0,x_0)\), the two
conditions are the Vandermonde system
\[
 c_1+c_2=A_0,
 \qquad
 \lambda_1c_1+\lambda_2c_2=D.
 \tag{V64.14}
\]
Its determinant is \(\lambda_2-\lambda_1\ne0\), and inversion gives
\textup{(V64.12)}.  The final assertions follow by subtraction and spatial
differentiation.
\end{proof}

The theorem is stronger than a sign counterexample: at a simple Hardy node
both components of the first reader remain freely prescribable.  If the
real signal \(m=2\Re{\cal F}\) is also stationary, then
\(\Im D=0\), fixing one real component but leaving \(\Re D\) arbitrary.
At a global threshold minimum, V58 adds the fractional one-sided inequality,
yet its carrier family realizes arbitrarily large allowed first readers.
Thus ordinary caloricity and positive-frequency support do not prove
\textup{(V64.9)}.

\subsection{Why a node-only estimate is not a GRH criterion}

For a compact interval \(I\), define the Hardy-node first-reader envelope
\[
 {\cal N}_{\chi,I}(t)
 :=\sup\left\{
 |P_{\chi,t}^{[1]}(x)|:
 x\in I,\ P_{\chi,t}^{[0]}(x)=A_{\chi,t}(x)
 \right\},
 \tag{V64.15}
\]
with value zero when the set is empty.

\begin{theorem}[Strict logical weakness of the node envelope]
\label{thm:v64-node-not-criterion}
Every estimate valid at all V27 real contacts also applies at Hardy nodes,
because every Hardy node is a zero-level contact.  The converse does not
follow.  More precisely, the global continuation alternative under a hypothetical
failure of paired GRH permits:
\begin{enumerate}[label=\textup{(\roman*)},leftmargin=3em]
\item a Hardy node satisfying \textup{(V64.1)}; or
\item no Hardy node on either exposed continuation component, but two
      external Hardy sign-defect contacts.
\end{enumerate}
Therefore even a zero-rate bound for \({\cal N}_{\chi,I}(t)\), or its
identical vanishing, is not by itself equivalent to GRH.  It leaves the
second branch of the proved V54 alternative untouched.
\end{theorem}

\begin{proof}
The inclusion of node sets follows from
\(2\Re{\cal F}=m\).  The two alternatives are exactly
Theorem~\ref{thm:v54-node-defect} and
Corollary~\ref{cor:v54-grh-alternative}.  Those results do not assert that
the first branch must occur.  Hence a hypothesis quantified only over nodes
can be vacuous along a forbidden continuation realized through the second
branch.
\end{proof}

This logical point prevents an illegitimate strengthening of V27.  Replacing
its real contact constraint by the smaller complex node set makes an
estimate easier to state but destroys the established equivalence unless
the sign-defect branch is independently excluded.

\begin{theorem}[V64 complex-contact verdict]
\label{thm:v64-verdict}
The full complex contact produces an exact regression identity but no
generic saving and no standalone GRH criterion:
\begin{enumerate}[label=\textup{(\roman*)},leftmargin=3em]
\item it fixes the normalized zeroth reader exactly as in
      \textup{(V64.4)};
\item the centered first reader remains an independent coordinate;
\item two positive frequencies realize arbitrary value and first-reader
      data at a Hardy node;
\item the known arithmetic bound for that coordinate misses the required
      scale \textup{(V64.9)} by a full \(e^{1/(16t)}\) factor at logarithmic
      scale;
\item a node-only estimate leaves the Hardy sign-defect continuation branch
      untouched.
\end{enumerate}
Thus the V63 proposal closes as a generic value--derivative route.  New
progress must either prove a fixed-coefficient conditional correlation far
beyond PNT or attack the node and sign-defect traces simultaneously.
\end{theorem}

\begin{firewall}[Interpolation capacity is not a Dirichlet model]
The coefficients \(c_1,c_2\) in Theorem
\ref{thm:v64-two-frequency-freedom} are freely chosen and are not von
Mangoldt character coefficients.  The theorem rules out deductions from
the complex equality, heat equation, and positive-frequency support alone;
it does not disprove a source-specific arithmetic correlation.  No such
correlation is proved here, and GRH remains open.
\end{firewall}

\begin{assessment}[V64 closure and V65 direction]
V64 separates algebraic strength from logical sufficiency.  The complex
node equation supplies one additional real value constraint, but no
derivative relation; moreover, nodes are only one branch of the established
global alternative.

V65 should seek one coefficient-sensitive quantity defined on the entire
real zero set that distinguishes both traces at once.  A natural candidate
is the signed phase flux
\[
 \Im{\cal F}_{\chi,t}(x)\,\partial_xm_{\chi,t}(x),
 \tag{V64.16}
\]
with nodes included as its zero boundary.  The task is to express its sum
over all zero-level contacts by an exact contour or argument-principle
identity for the fixed Hardy function.  If the identity reduces to the
already closed zero-graph degree or admits the V56--V59 defect models, that
combined route must also be closed.
\end{assessment}

\section*{Version V64 append-only record}
\addcontentsline{toc}{section}{Version V64 append-only record}
The complete V63 source was retained before this continuation.  It had
\(935779\) bytes, \(27940\) lines, and SHA--256
\[
 \texttt{501A4169521C69DDEB6E000E8E23AE45D42B140EB3E97B4B8629F98FC6D4FAD5}.
\]
V64 derived the exact normalized complex-contact regression, proved
two-frequency freedom of the first reader under a full Hardy-node value
constraint, and showed that a node-only derivative estimate is logically
weaker than the established node/sign-defect alternative.

% =============================================================================
% END APPENDED V64 MATERIAL
% =============================================================================

% =============================================================================
% BEGIN APPENDED V65 MATERIAL
% =============================================================================

\section{V65: phase-flux crossings and the contour no-go}
\label{sec:grh-v65}

V64 proposed a quantity defined on the whole real zero set,
\[
 \Im{\cal F}_{\chi,t}(x)\,\partial_xm_{\chi,t}(x),
 \tag{V65.1}
\]
in the hope that a contour identity might treat Hardy nodes and Hardy sign
defects simultaneously.  There is an exact intersection identity, but it
reveals two obstructions.  The argument principle sees only crossing
orientation, not the magnitude in \textup{(V65.1)}.  Moreover, the product
vanishes at every Hardy node, so it cannot record the other branch of the
V54 alternative without adding boundary-zero multiplicities separately.

Fix \(t>0\) and write on the real boundary
\[
 F(x):={\cal F}_{\chi,t}(x)=u(x)+iv(x),
 \qquad u(x)=\frac12m_{\chi,t}(x).
 \tag{V65.2}
\]
The positive-time functions are real analytic.  We first work on a compact
interval \([a,b]\) whose endpoints lie off both coordinate axes and assume
that all crossings are simple and that \(F\ne0\) on the interval.

\subsection{Local angular meaning of the phase flux}

\begin{lemma}[A zero-level crossing is an angular crossing]
\label{lem:v65-angular-crossing}
If \(u(x_j)=0\), \(u'(x_j)\ne0\), and \(v(x_j)\ne0\), then
\[
 \left.\frac{d}{dx}\arg F(x)\right|_{x=x_j}
 =-\frac{u'(x_j)}{v(x_j)}
 =-\frac{\partial_xm_{\chi,t}(x_j)}
         {2\Im{\cal F}_{\chi,t}(x_j)}.
 \tag{V65.3}
\]
Consequently
\[
 \boxed{
 \operatorname{sgn}\!\left(
 \Im{\cal F}_{\chi,t}(x_j)\,
 \partial_xm_{\chi,t}(x_j)
 \right)
 =-\operatorname{sgn}\left.\frac{d}{dx}\arg F(x)\right|_{x_j}.}
 \tag{V65.4}
\]
\end{lemma}

\begin{proof}
At the crossing, \(F(x_j)=iv(x_j)\).  Since
\[
 \frac{d}{dx}\arg F=\Im\frac{F'}F,
 \tag{V65.5}
\]
division of \(u'+iv'\) by \(iv\) gives \(-u'/v\) as its imaginary part.
Use \(m'=2u'\) to obtain \textup{(V65.3)}--\textup{(V65.4)}.
\end{proof}

Thus the sign-defect convention is precisely a local orientation of the
boundary curve \(F(\mathbb R)\) as it crosses the imaginary axis.  A local
orientation, however, is not yet a winding number.

\subsection{The exact quadrant-balance identity}

Let
\[
 Z_u:=\{x\in(a,b):u(x)=0\},
 \qquad
 Z_v:=\{x\in(a,b):v(x)=0\}.
 \tag{V65.6}
\]
For the assumed simple nonintersecting crossings, define
\[
 \epsilon(x):=\operatorname{sgn}(v(x)u'(x))
 \quad(x\in Z_u),
 \tag{V65.7}
\]
\[
 \eta(x):=\operatorname{sgn}(u(x)v'(x))
 \quad(x\in Z_v).
 \tag{V65.8}
\]

\begin{theorem}[Quadrant balance for the Hardy boundary curve]
\label{thm:v65-quadrant-balance}
Under the hypotheses above,
\[
 \boxed{
 \sum_{x\in Z_u}\epsilon(x)
 +\sum_{x\in Z_v}\eta(x)
 =\frac12\left[
 \operatorname{sgn}(u(b))\operatorname{sgn}(v(b))
 -\operatorname{sgn}(u(a))\operatorname{sgn}(v(a))
 \right].}
 \tag{V65.9}
\]
Equivalently, the algebraic sum of phase-flux orientations at zero-level
contacts is determined only after adding the oriented crossings of the
orthogonal equation \(\Im F=0\).
\end{theorem}

\begin{proof}
Away from the two crossing sets, put
\[
 q(x):=\operatorname{sgn}(u(x))\operatorname{sgn}(v(x)).
 \tag{V65.10}
\]
At a simple zero of \(u\), its sign jumps by \(2\operatorname{sgn}u'\),
so \(q\) jumps by \(2\epsilon\).  At a simple zero of \(v\), it jumps by
\(2\eta\).  Summing all jumps of the piecewise constant function \(q\)
from \(a\) to \(b\) proves \textup{(V65.9)}.
\end{proof}

The theorem is the discrete Stokes formula for the four quadrants.  It
contains no positivity: an undesired imaginary-axis crossing can be balanced
by a horizontal-axis crossing or by endpoint quadrant data.  The V56 open
defect chambers and V59 alternating orientations are therefore compatible
with the identity.

\subsection{The amplitude identity has an uncontrolled bulk term}

Let \(H\) denote the Heaviside function.  Distributionally,
\(\frac d{dx}H(m)=\delta(m)m'\).  Hence simple zero-level contacts obey
the exact coarea formula
\[
 \boxed{
 \sum_{m(x_j)=0}v(x_j)\operatorname{sgn}m'(x_j)
 =\int_a^b\delta(m(x))v(x)m'(x)\,dx}
 \tag{V65.11}
\]
and integration by parts gives
\[
 \boxed{
 \sum_{m(x_j)=0}v(x_j)\operatorname{sgn}m'(x_j)
 =[vH(m)]_a^b-\int_a^bH(m(x))v'(x)\,dx.}
 \tag{V65.12}
\]

\begin{proposition}[No contour closure for the weighted phase flux]
\label{prop:v65-weighted-no-go}
The weighted sum proposed in V64 has the coarea representation
\[
 \boxed{
 \sum_{m(x_j)=0}v(x_j)m'(x_j)
 =\int_a^b\delta(m(x))|m'(x)|v(x)m'(x)\,dx.}
 \tag{V65.13}
\]
It is not determined by the winding number, boundary-zero count, crossing
orientations, or endpoint values of \(\arg F\).  In particular, replacing
\(F\) by \(cF\) with \(c>0\) preserves all those topological data while
multiplying \textup{(V65.13)} by \(c^2\).
\end{proposition}

\begin{proof}
The delta-function change of variables at every simple root gives
\textup{(V65.13)}.  Under \(F\mapsto cF\), one has
\(v\mapsto cv\) and \(m'\mapsto cm'\), while zeros, arguments, and
orientations are unchanged.  Thus no identity depending only on those
contour data can determine the weighted sum.
\end{proof}

Even the oriented amplitude in \textup{(V65.12)} is not a boundary
invariant: it contains the bulk integral of \(v'\) over the positivity set
of \(m\).  For the arithmetic signal that term is another first-reader
integral, not a new cancellation mechanism.

\subsection{Hardy nodes are invisible to the proposed flux}

At a Hardy node \(x_*\),
\[
 u(x_*)=v(x_*)=0,
 \tag{V65.14}
\]
so the product in \textup{(V65.1)} is exactly zero regardless of
\(m'(x_*)\).  It has neither a sign-defect orientation nor a nonzero weight.

\begin{theorem}[Nodes require separate boundary multiplicity]
\label{thm:v65-node-multiplicity}
Suppose \(F\) extends holomorphically across a neighbourhood of a real node
\(x_*\) and has order \(q\ge1\):
\[
 F(z)=c(z-x_*)^q+O((z-x_*)^{q+1}),
 \qquad c\ne0.
 \tag{V65.15}
\]
Then a small semicircular indentation changes the boundary argument by
magnitude \(q\pi\), and the node contributes half its multiplicity to the
corresponding argument-principle count.  This information is absent from
\textup{(V65.1)}, which vanishes at the node.
\end{theorem}

\begin{proof}
On a semicircle \(z-x_*=\varepsilon e^{i\theta}\), the leading factor in
\textup{(V65.15)} changes argument by \(q\) times the semicircle angle.
The error is uniform little-oh of one as \(\varepsilon\downarrow0\).
Thus the local change has magnitude \(q\pi\), half the \(2\pi q\) of a
full circle.  The zero factor makes \textup{(V65.1)} vanish directly.
\end{proof}

Adding these half-multiplicities recovers the ordinary boundary-corrected
argument principle.  It does not exclude the nodes: V61--V62 proved that
finite-order real zeros can reside entirely in the outer boundary modulus.
Thus the proposed combined statistic separates back into the two already
known pieces---axis-crossing orientation for sign defects and canonical
boundary-zero multiplicity for Hardy nodes.

\begin{theorem}[V65 contour verdict]
\label{thm:v65-verdict}
No new simultaneous node/sign-defect obstruction follows from the phase
flux:
\begin{enumerate}[label=\textup{(\roman*)},leftmargin=3em]
\item its sign at a nonnode contact is the negative angular orientation of
      an imaginary-axis crossing;
\item the algebraic orientation sum is balanced by horizontal-axis
      crossings and endpoint quadrants through \textup{(V65.9)};
\item its magnitude is not topological and has the uncontrolled bulk form
      \textup{(V65.12)}--\textup{(V65.13)};
\item it vanishes at Hardy nodes, whose missing contribution is the usual
      boundary-zero multiplicity;
\item adjoining that multiplicity reduces the contour argument to the
      canonical factorization framework already closed in V61--V62.
\end{enumerate}
Therefore an argument-principle identity cannot force a common favorable
sign on both branches of the V54 alternative using only boundary topology.
\end{theorem}

\begin{firewall}[The crossing identity has no arithmetic sign]
The quadrant balance is exact for the arithmetic Hardy curve, but it is a
kinematic identity.  It supplies no restriction on the horizontal crossings
or endpoint quadrants, and the coefficient-dependent bulk integral remains
unestimated.  No Hardy node or sign defect has been excluded, and GRH
remains open.
\end{firewall}

\begin{assessment}[V65 closure and V66 direction]
V65 carries out the contour test proposed in V64 and closes it at the
generic level.  The only exact contour invariant is crossing orientation;
the desired amplitude is noninvariant, while nodes require a separate term
already known to be nonexcluding.

The next upstream coefficient test should ask whether the fixed completed
logarithmic derivative belongs to a Pick or Herglotz class in the Euler
half-plane and whether Gaussian heat evolution preserves any such cone.
V66 should compute the relevant imaginary-part kernel directly from the
Euler coefficients.  If complex character phases destroy the sign already
in the absolutely convergent region, or if the heat multiplier fails to
preserve the class, that analytic-monotonicity route should be closed before
returning to character-specific cancellation.
\end{assessment}

\section*{Version V65 append-only record}
\addcontentsline{toc}{section}{Version V65 append-only record}
The complete V64 source was retained before this continuation.  It had
\(945485\) bytes, \(28213\) lines, and SHA--256
\[
 \texttt{7190A40FC287A04BCDD845D3FA878909B9563FA5DC8C8A01F6E2629E19FEC5A5}.
\]
V65 derived the local angular identity, the global quadrant-balance formula,
and the coarea ledger for phase flux; it proved that weighted flux is not a
contour invariant and that Hardy nodes require separate boundary
multiplicity already covered by the outer-factor no-go.

% =============================================================================
% END APPENDED V65 MATERIAL
% =============================================================================

% =============================================================================
% BEGIN APPENDED V66 MATERIAL
% =============================================================================

\section{V66: the Pick reformulation and the heat-cone tautology}
\label{sec:grh-v66}

V65 closed the generic contour-flux route and proposed testing whether the
fixed completed logarithmic derivative belongs to a Pick or Herglotz class.
The test has an exact answer.  The unrotated function has both signs of its
imaginary part even in the absolutely convergent Euler region.  After the
correct quarter-turn, however, the Pick property is equivalent to paired
GRH: critical-line zeros become positive boundary atoms, whereas a zero to
the right of the line becomes an interior pole.  Gaussian heat evolution
does preserve this cone, but the resulting statement is exactly the scalar
Weil positivity criterion already isolated in V22.

For each \(\psi\in\{\chi,\overline\chi\}\), use the principal-channel pole
subtraction when necessary and put
\[
 F_{\psi,0}(z)
 :=\frac{\Lambda'}{\Lambda}
   \!\left(\frac12-iz,\psi\right).
 \tag{V66.1}
\]
By V61 this agrees with the unheated positive-frequency signal in the Euler
region \(\Im z>1/2\) and supplies its meromorphic continuation.

\subsection{The unrotated Euler function is not Pick}

Write \(z=x+iy\), \(\sigma=1/2+y>1\).  Absolute convergence gives
\[
 \begin{aligned}
 F_{\chi,0}(x+iy)
 &=\frac12\log\frac f\pi
  +\frac12\psi\!\left(\frac{\sigma+\kappa-ix}{2}\right)\\
 &\quad-
 \sum_{n\ge2}\frac{\Lambda(n)\chi(n)}{n^\sigma}
 e^{ix\log n}.
 \end{aligned}
 \tag{V66.2}
\]

\begin{proposition}[Both imaginary signs occur in the Euler region]
\label{prop:v66-euler-both-signs}
For every fixed primitive character \(\chi\), there are arbitrarily high
points in the Euler region at which
\(\Im F_{\chi,0}\) is positive and arbitrarily high points at which it is
negative.  More precisely, as \(\sigma\to\infty\),
\[
 \Im F_{\chi,0}(\sigma+i(\sigma-1/2))
 \longrightarrow-\frac\pi8,
 \tag{V66.3}
\]
\[
 \Im F_{\chi,0}(-\sigma+i(\sigma-1/2))
 \longrightarrow+\frac\pi8.
 \tag{V66.4}
\]
Thus \(F_{\chi,0}\) itself is neither a Pick function nor an anti-Pick
function.
\end{proposition}

\begin{proof}
At \(x=\pm\sigma\), the Euler series in \textup{(V66.2)} is
\(O_\chi(2^{-\sigma})\).  Stirling's formula gives
\[
 \psi\!\left(\frac{\sigma+\kappa\mp i\sigma}{2}\right)
 =\log\!\left(\frac{\sigma+\kappa\mp i\sigma}{2}\right)+o(1).
 \tag{V66.5}
\]
The two arguments tend to \(\mp\pi/4\).  Multiplication by \(1/2\) proves
\textup{(V66.3)}--\textup{(V66.4)}.
\end{proof}

The complex Euler coefficients are not the main reason for this sign
change; the gamma factor already forces it.  The relevant cone is positive
\emph{real} part, equivalently the Pick property after multiplication by
\(i\).

\subsection{The correct Pick statement is GRH-complete}

Recall that a Caratheodory function on \(\mathbb C_+\) is holomorphic and
has nonnegative real part.  Equivalently, \(iF\) is Pick/Herglotz because
\(\Im(iF)=\Re F\).

\begin{theorem}[Pick equivalence for the conjugate pair]
\label{thm:v66-pick-equivalence}
The following are equivalent.
\begin{enumerate}[label=\textup{(\roman*)},leftmargin=3em]
\item Paired GRH holds for \(\chi,\overline\chi\).
\item For both \(\psi=\chi\) and \(\psi=\overline\chi\), the function
      \(F_{\psi,0}\) is holomorphic on \(\mathbb C_+\) and
\[
 \Re F_{\psi,0}(z)\ge0
 \qquad(z\in\mathbb C_+).
 \tag{V66.6}
\]
\item For both channels, \(iF_{\psi,0}\) belongs to the Pick class on
      \(\mathbb C_+\).
\end{enumerate}
Under these conditions the boundary real-part measure is
\[
 \boxed{
 \Re F_{\psi,0}(x+i0)
 =\pi\sum_{\rho=1/2+i\gamma}
 m_\rho\,\delta_{-\gamma}(x),}
 \tag{V66.7}
\]
with the paired normalization of V22.
\end{theorem}

\begin{proof}
Assume paired GRH.  A zero \(\rho=1/2+i\gamma\) of the \(\psi\)-channel
maps under \(s=1/2-iz\) to the real point \(x_\rho=-\gamma\).  Its local
logarithmic-derivative term is
\[
 \frac{m_\rho}{s-\rho}=\frac{im_\rho}{z-x_\rho},
 \tag{V66.8}
\]
whose real part is the positive Poisson atom
\(m_\rho y/\{(x-x_\rho)^2+y^2\}\) and whose boundary real part is
\(\pi m_\rho\delta_{x_\rho}\).  The regularized Hadamard sum, with the
normalization already fixed in V22 and V53, gives \textup{(V66.7)}; its
remaining constant is purely imaginary.  Hence its Poisson extension is
nonnegative, \textup{(V66.6)} holds, and \(iF_{\psi,0}\) is Pick.

Conversely, a nontrivial zero \(\rho=\beta+i\gamma\) of
\(\Lambda(s,\psi)\) produces a pole of positive integer residue in its
logarithmic derivative.  Under \(s=1/2-iz\), its location is
\[
 z_\rho=-\gamma+i(\beta-1/2).
 \tag{V66.9}
\]
If \(\beta>1/2\), this is an interior pole of \(F_{\psi,0}\), contradicting
holomorphy in \textup{(ii)} or \textup{(iii)}.  Thus neither channel has a
zero to the right of the critical line.  The functional equation maps every
left-side zero of one channel to a right-side zero of the conjugate channel.
Requiring the condition for both channels therefore places every
nontrivial zero on \(\Re s=1/2\).  The equivalence of \textup{(ii)} and
\textup{(iii)} is \(\Im(iF)=\Re F\).
\end{proof}

Holomorphy alone in both channels already excludes off-line zeros.  The
positive-real assertion records the stronger boundary orientation of their
critical-line residues, but it is not an independently accessible Euler
estimate.

\subsection{Heat evolution preserves the cone and recovers V22}

\begin{theorem}[Gaussian preservation of the Caratheodory cone]
\label{thm:v66-heat-preservation}
Suppose \(F_0\) is a positive-real analytic signal on \(\mathbb C_+\) with
the tempered growth occurring here.  Then for every \(t>0\),
\[
 F_t(z):=e^{t\partial_z^2}F_0(z)
 =\int_{\mathbb R}H_t(r)F_0(z-r)\,dr
 \tag{V66.10}
\]
has nonnegative real part on \(\mathbb C_+\).  If its boundary real-part
measure is nonzero and positive, then
\[
 \Re F_t(x)>0
 \qquad(x\in\mathbb R).
 \tag{V66.11}
\]
For the arithmetic signal,
\[
 \boxed{2\Re F_{\chi,t}(x)=m_{\chi,t}(x),}
 \tag{V66.12}
\]
so the heat-preserved Pick condition is exactly positive-time scalar Weil
positivity.
\end{theorem}

\begin{proof}
The Gaussian kernel is positive and has total mass one.  Taking real parts
in \textup{(V66.10)} expresses \(\Re F_t(z)\) as a positive average of
translates of \(\Re F_0\), proving nonnegativity.  On the boundary it is
the convolution of a nonzero positive measure with the strictly positive
Gaussian, proving \textup{(V66.11)}.  Equation \textup{(V66.12)} is the
positive-frequency identity \textup{(V53.10)}.
\end{proof}

\begin{corollary}[All-time preservation is GRH; one-time preservation is
not]
\label{cor:v66-all-time}
For the paired arithmetic channel, positivity of
\(\Re F_{\chi,t}\) and \(\Re F_{\overline\chi,t}\) for every \(t>0\) is
equivalent to paired GRH.  Positivity at one fixed heat time is not: under a
hypothetical GRH failure, V22 permits the upper interval
\([\vartheta_\chi,\infty)\) on which the heat symbol is nonnegative.
\end{corollary}

\begin{proof}
The forward implication under GRH follows from
Theorems~\ref{thm:v66-pick-equivalence} and
\ref{thm:v66-heat-preservation}.  Conversely, positivity at all positive
times and the approximate-identity limit make the boundary real-part
distribution positive; equivalently, V23's exposed-pair theorem shows that
any off-line zero creates negative values along a sequence
\(t_k\downarrow0\).  The last assertion is the positivity-threshold
structure \textup{(V22.18)}.
\end{proof}

\subsection{Why the Euler product does not prove the Pick condition}

In the absolutely convergent region, the desired real-part inequality is
the coefficient-sensitive statement
\[
 \Re\left[
 \frac12\log\frac f\pi
 +\frac12\psi\!\left(\frac{\sigma+\kappa-ix}{2}\right)
 -\sum_{n\ge2}\frac{\Lambda(n)\chi(n)}{n^\sigma}
 e^{ix\log n}
 \right]\ge0.
 \tag{V66.13}
\]
The prime summands have no common sign for a nonprincipal character.  More
importantly, even proving \textup{(V66.13)} only for \(\sigma>1\) does not
exclude a pole in the unsampled strip \(1/2<\Re s\le1\).  Transporting the
inequality through that strip requires precisely the zero-free continuation
that Theorem~\ref{thm:v66-pick-equivalence} identifies with GRH.

\begin{theorem}[V66 Pick-route verdict]
\label{thm:v66-verdict}
The Pick/Herglotz route has the following exact status.
\begin{enumerate}[label=\textup{(\roman*)},leftmargin=3em]
\item The unrotated completed logarithmic derivative has both imaginary
      signs already in the Euler region.
\item The correctly rotated function \(iF_{\chi,0}\) is Pick in both
      conjugate channels if and only if paired GRH holds.
\item Gaussian heat flow preserves this cone by positive convolution.
\item For the arithmetic signal, that preservation is exactly the known
      identity \(m_{\chi,t}=2\Re F_{\chi,t}\) and the scalar Weil positivity
      criterion.
\item Euler-region positivity alone cannot cross the critical strip without
      assuming the missing zero-free continuation.
\end{enumerate}
Thus Pick theory gives a clean equivalent formulation but no new estimate
and no shortcut around the off-line-zero obstruction.
\end{theorem}

\begin{firewall}[A class equivalence is not a proof]
No Pick property has been established from the Euler coefficients in the
strip \(0<\Im z\le1/2\).  Theorem
\ref{thm:v66-pick-equivalence} shows that doing so for both channels would
already prove GRH.  Heat preservation starts only after positivity of the
initial boundary measure is known.  GRH remains open.
\end{firewall}

\begin{assessment}[V66 closure and V67 direction]
V66 resolves the analytic-monotonicity test completely.  It identifies the
correct quarter-turn, proves its GRH equivalence, and shows that the heat
cone is simply another form of the V22 scalar criterion.

V67 should test the associated Pick kernel
\[
 K_F(z,w):=\frac{F(z)+\overline{F(w)}}{-i(z-\overline w)}.
 \tag{V66.14}
\]
Finite kernel matrices in the Euler half-plane are directly computable from
the absolutely convergent coefficients.  The precise question is whether
their positivity plus a quantitative continuation principle can control
the lower strip, or whether finite Pick data admit rational positive-real
interpolants with arbitrary hidden poles below the sampling line.  The
latter would close finite interpolation certificates while identifying the
necessary infinite-scale input.
\end{assessment}

\section*{Version V66 append-only record}
\addcontentsline{toc}{section}{Version V66 append-only record}
The complete V65 source was retained before this continuation.  It had
\(955568\) bytes, \(28484\) lines, and SHA--256
\[
 \texttt{6AE1417BBCB253FD953CAF7323F28AC87D622C8C50F600E9A7A7E21F919015EE}.
\]
V66 proved that the unrotated Euler function is not Pick, that the correct
quarter-turned completed logarithmic derivative is Pick exactly under
paired GRH, and that Gaussian preservation of this cone reproduces the
existing scalar Weil positivity criterion.

% =============================================================================
% END APPENDED V66 MATERIAL
% =============================================================================

% =============================================================================
% BEGIN APPENDED V67 MATERIAL
% =============================================================================

\section{V67: finite Pick certificates and hidden lower-strip poles}
\label{sec:grh-v67}

V66 identified paired GRH with the global Pick property of the quarter-turned
completed logarithmic derivatives.  It proposed testing finite Pick kernels
at points where the Euler series converges absolutely.  This version carries
out that test and proves its exact limitation.  A positive finite Pick matrix
certifies that the sampled data admit \emph{some} positive-real interpolant;
it does not certify that the meromorphic continuation of the arithmetic
function is that interpolant.  Finite values and finite jets can be preserved
exactly while inserting a pole anywhere in the unsampled strip
\(0<\Im z<1/2\).

For a function \(F\) and points \(z,w\in\mathbb C_+\), put
\[
 K_F(z,w):=
 \frac{F(z)+\overline{F(w)}}{-i(z-\overline w)}.
 \tag{V67.1}
\]

\subsection{What a finite Pick matrix actually proves}

\begin{theorem}[Finite positive-real interpolation certificate]
\label{thm:v67-pick-certificate}
Let \(z_1,\ldots,z_N\in\mathbb C_+\) be distinct and let
\(w_1,\ldots,w_N\in\mathbb C\).  If a holomorphic function \(F\) satisfies
\(\Re F\ge0\) on \(\mathbb C_+\) and \(F(z_j)=w_j\), then
\[
 \boxed{
 \mathsf K_N:=
 \left[
 \frac{w_j+\overline{w_k}}{-i(z_j-\overline{z_k})}
 \right]_{j,k=1}^N\succeq0.}
 \tag{V67.2}
\]
Conversely, \(\mathsf K_N\succeq0\) is sufficient for the existence of at
least one positive-real holomorphic interpolant \(G\) with
\(G(z_j)=w_j\).  It neither identifies that interpolant uniquely nor
controls a different meromorphic function having the same sampled values.
\end{theorem}

\begin{proof}
The Cayley transform
\[
 \phi(z):=\frac{F(z)-1}{F(z)+1}
 \tag{V67.3}
\]
maps the positive-real class to the Schur class.  The usual half-plane Pick
matrix for \(\phi\) is positive semidefinite.  The identity
\[
 1-\phi(z_j)\overline{\phi(z_k)}
 =\frac{2\{w_j+\overline{w_k}\}}
        {(w_j+1)(\overline{w_k}+1)}
 \tag{V67.4}
\]
shows by diagonal congruence that its positivity is equivalent to
\textup{(V67.2)}.  The Nevanlinna--Pick interpolation theorem gives the
converse.  Nonuniqueness is immediate whenever the Pick problem is not
extremal; the hidden-pole construction below proves the stronger final
assertion without relying on nonuniqueness inside the positive class.
\end{proof}

For the arithmetic samples
\(w_j=F_{\chi,0}(z_j)\), a negative principal minor would disprove the
global Pick condition and hence paired GRH.  Positive minors are only
necessary evidence: they say that the finite table is compatible with GRH,
not that GRH holds.

\subsection{Exact finite-jet invisibility of a hidden pole}

\begin{theorem}[Rational hidden-pole interpolation]
\label{thm:v67-hidden-pole}
Fix distinct sample points \(z_1,\ldots,z_N\) in the Euler region
\(\Im z>1/2\), nonnegative jet orders \(r_1,\ldots,r_N\), a hidden point
\[
 p\in\mathbb C_+,
 \qquad 0<\Im p<\frac12,
 \tag{V67.5}
\]
and an integer \(L\ge0\).  Define
\[
 Q(z):=\prod_{j=1}^N(z-z_j)^{r_j+1}
 \tag{V67.6}
\]
and choose an integer
\[
 M>\deg Q+L+1.
 \tag{V67.7}
\]
Then
\[
 R_p(z):=\frac{Q(z)}{(z-p)(z+i)^M}
 \tag{V67.8}
\]
has the following properties:
\begin{enumerate}[label=\textup{(\roman*)},leftmargin=3em]
\item it has a simple pole at \(p\) and no other pole in
      \(\mathbb C_+\);
\item for every \(j\) and \(0\le k\le r_j\),
\[
 R_p^{(k)}(z_j)=0;
 \tag{V67.9}
\]
\item as \(z\to\infty\),
\[
 R_p(z)=O(|z|^{-L-2}).
 \tag{V67.10}
\]
Consequently, for any meromorphic base function \(B\) and any
\(\varepsilon\ne0\),
\[
 \widetilde B(z):=B(z)+\varepsilon R_p(z)
 \tag{V67.11}
\]
has exactly the same prescribed finite jets and the same asymptotic
expansion through order \(L\) as \(B\), while possessing a pole at \(p\)
unless a pre-existing residue is tuned to cancel it.
\end{enumerate}
\end{theorem}

\begin{proof}
The factor \((z-z_j)^{r_j+1}\) proves \textup{(V67.9)}.  The pole
\(z=-i\) lies below the closed upper half-plane, while \(Q(p)\ne0\) because
all sample points lie above the sampling line and \(p\) lies below it.
Thus \(p\) is the only upper-half-plane pole and it is simple.  Degree
counting with \textup{(V67.7)} gives \textup{(V67.10)}.  Addition proves
the remaining assertions.
\end{proof}

The residue at \(p\) can be prescribed arbitrarily by choosing
\[
 \varepsilon=\mathfrak r\,\frac{(p+i)^M}{Q(p)}.
 \tag{V67.12}
\]
In particular one may impose the logarithmic-derivative residue
\(\mathfrak r=i m\), \(m\in\mathbb N\), corresponding under
\(s=1/2-iz\) to a zero of multiplicity \(m\).  Thus the obstruction is not
an artifact of allowing a nonarithmetical residue.

\begin{corollary}[Confluent finite Pick tests cannot prove GRH]
\label{cor:v67-confluent-no-go}
No certificate using finitely many Euler-region values and finitely many
derivatives of \(F_{\chi,0}\), even together with finitely many terms of its
asymptotic expansion, can exclude a pole in
\(0<\Im z<1/2\).  All corresponding ordinary or confluent Pick matrices are
unchanged by \textup{(V67.11)}.
\end{corollary}

\begin{proof}
Every matrix entry is a function only of the retained finite jets.  Those
jets and the stated asymptotic data are identical for \(B\) and
\(\widetilde B\), while Theorem~\ref{thm:v67-hidden-pole} inserts the hidden
pole.
\end{proof}

\subsection{Even full Euler-region sign does not cross the gap}

Finite invisibility is not the only issue.  Pointwise positive real part on
the entire Euler half-plane is itself compatible with a pole below that
half-plane.

\begin{proposition}[A hidden pole positive above its height]
\label{prop:v67-pole-positive-above}
Let \(p=a+ib\) with \(0<b<1/2\), and put
\[
 H_p(z):=\frac{i}{z-p}.
 \tag{V67.13}
\]
Then \(H_p\) has a pole in \(\mathbb C_+\), but
\[
 \boxed{
 \Re H_p(x+iy)
 =\frac{y-b}{(x-a)^2+(y-b)^2}>0
 \qquad(y>1/2).}
 \tag{V67.14}
\]
Thus positivity on the complete absolutely convergent Euler region does not
imply holomorphy or positive-real continuation on all of \(\mathbb C_+\).
\end{proposition}

\begin{proof}
Multiply numerator and denominator in \textup{(V67.13)} by the conjugate of
\((x-a)+i(y-b)\).  The real part is exactly
\textup{(V67.14)}, and \(y>1/2>b\) makes it positive.
\end{proof}

The example is the local logarithmic-derivative term of an off-line zero.
It explains geometrically why Euler-region positivity cannot see that zero:
above the pole, the pole contributes with the favorable Poisson sign.  The
failure occurs only when continuation descends through its height.

\subsection{The Pick-kernel defect of an interior pole}

For completeness, the standard upper-half-plane kernel of
\textup{(V67.13)} is
\[
 \boxed{
 K_{H_p}(z,w)
 =\frac1{(z-p)(\overline w-\overline p)}
 -\frac{2\Im p}
 {(z-p)(\overline w-\overline p)\{-i(z-\overline w)\}}.}
 \tag{V67.15}
\]
When \(p\) lies on the real boundary, the second term vanishes and the
kernel is the positive rank-one boundary-atom kernel.  Moving the pole into
the upper half-plane creates the negative defect proportional to
\(2\Im p\).  Detecting that defect requires sample geometry rich enough to
distinguish the actual function from its finite positive interpolants.

\begin{proof}[Verification of \textup{(V67.15)}]
Writing \(p-\overline p=2i\Im p\), one has
\[
 \frac{i}{z-p}-\frac{i}{\overline w-\overline p}
 =\frac{-i(z-\overline w)-2\Im p}
 {(z-p)(\overline w-\overline p)}.
 \tag{V67.16}
\]
Division by \(-i(z-\overline w)\) proves the formula.
\end{proof}

\begin{theorem}[V67 finite-certificate verdict]
\label{thm:v67-verdict}
Euler-region Pick sampling has the following exact scope.
\begin{enumerate}[label=\textup{(\roman*)},leftmargin=3em]
\item A negative finite Pick minor would contradict paired GRH.
\item A positive finite Pick matrix certifies only the existence of a
      positive-real interpolant for that finite table.
\item Arbitrarily many but finitely many values, jets, and asymptotic
      coefficients can be preserved while inserting a correctly oriented
      simple logarithmic-derivative pole below the sampling line.
\item Even pointwise positive real part throughout \(\Im z>1/2\) does not
      exclude such a pole.
\item The pole appears in the global Pick kernel through the negative defect
      in \textup{(V67.15)}, but finite Euler data need not resolve it.
\end{enumerate}
Therefore finite Pick certificates cannot transport the Euler product
through the lower half of the critical strip.
\end{theorem}

\begin{firewall}[Interpolation compatibility is not continuation]
The rational perturbation in Theorem~\ref{thm:v67-hidden-pole} does not have
the fixed Euler coefficients or the complete functional equation of the
Dirichlet channel.  It proves that finite Pick data and finite asymptotics do
not encode those global constraints.  Excluding the hidden pole still
requires a coefficient-sensitive quantitative continuation theorem.  GRH
remains open.
\end{firewall}

\begin{assessment}[V67 closure and V68 direction]
V67 closes finite Euler-region Pick certificates, including confluent jet
tests.  The exact kernel formula nevertheless exposes a quantitative signal:
an interior pole produces a negative term proportional to its height
\(b=\Im p=\Re\rho-1/2\).

V68 should compute the smallest sampling configuration that detects the
model kernel \textup{(V67.15)} and quantify how its negative eigenvalue
decays when all samples remain above \(\Im z=1/2\).  Comparing that
detectability scale with the Euler truncation error will decide whether a
growing, rather than finite, certificate has any noncircular arithmetic
content.
\end{assessment}

\section*{Version V67 append-only record}
\addcontentsline{toc}{section}{Version V67 append-only record}
The complete V66 source was retained before this continuation.  It had
\(966794\) bytes, \(28770\) lines, and SHA--256
\[
 \texttt{4C6551CDBC5A7680E0F601A274197DDF05EC8DFD5B3F85B994EE4ACFE4CC005D}.
\]
V67 proved the finite Pick interpolation criterion, constructed exact
finite-jet hidden-pole perturbations with prescribed logarithmic residue,
and showed that even full Euler-region real-part positivity does not exclude
a pole in the lower strip.

% =============================================================================
% END APPENDED V67 MATERIAL
% =============================================================================

% =============================================================================
% BEGIN APPENDED V68 MATERIAL
% =============================================================================

\section{V68: exhaustive Pick detection and the nonuniform pole margin}
\label{sec:grh-v68}

V67 proved that no prescribed finite Euler-region sample can exclude a
hidden pole.  That statement does not imply that a fixed hidden pole evades
\emph{every} finite sample.  The distinction matters.  If all finite Pick
matrices on a dense Euler set were positive, infinite Pick compactness would
produce a global positive-real interpolant.  Agreement on the Euler region
would then force that interpolant to equal the meromorphic completed
logarithmic derivative, excluding every hidden pole by uniqueness.

Thus an off-line zero has some finite negative Pick witness consisting
entirely of absolutely convergent Euler-region values.  The witness is
nonuniform: its size, locations, and negative spectral margin can depend on
the entire analytic background, not only on the offending zero.  This
version proves the exhaustive detection theorem, computes the exact
two-sample margin for one isolated pole, and compares that margin with a
rigorous Euler-tail error.

\subsection{A dense Euler sampling set}

Let
\[
 {\cal D}:=\{q+ir:q,r\in\mathbb Q,\ r>1/2\}
 =\{z_1,z_2,\ldots\}
 \tag{V68.1}
\]
be any enumeration without repetition.  For a meromorphic function \(F\)
which is holomorphic in the Euler half-plane, define its \(N\)-sample Pick
matrix
\[
 {\mathsf K}_N(F):=
 \left[
 \frac{F(z_j)+\overline{F(z_k)}}
 {-i(z_j-\overline{z_k})}
 \right]_{j,k=1}^N.
 \tag{V68.2}
\]

\begin{theorem}[Exhaustive Euler sampling detects every hidden pole]
\label{thm:v68-exhaustive-detection}
Let \(F\) be holomorphic on \(\{\Im z>1/2\}\) and admit a meromorphic
continuation to \(\mathbb C_+\).  If
\[
 {\mathsf K}_N(F)\succeq0
 \qquad\hbox{for every }N\ge1,
 \tag{V68.3}
\]
then there is a positive-real holomorphic function \(G\) on
\(\mathbb C_+\) such that
\[
 G(z)=F(z)
 \qquad(\Im z>1/2).
 \tag{V68.4}
\]
Consequently the meromorphic continuation of \(F\) has no pole in
\(\mathbb C_+\).

Equivalently, if \(F\) has any pole in
\(0<\Im z<1/2\), then
\[
 \boxed{\lambda_{\min}{\mathsf K}_N(F)<0}
 \tag{V68.5}
\]
for at least one finite \(N\).
\end{theorem}

\begin{proof}
For every \(N\), Theorem~\ref{thm:v67-pick-certificate} gives a
positive-real function \(G_N\) on \(\mathbb C_+\) interpolating
\(F(z_1),\ldots,F(z_N)\).  Fix the first value.  The Cayley transforms of
the \(G_N\) form a normal family of Schur functions; equivalently, the
positive-real family is locally normal after this fixed normalization.
Choose a diagonal subsequence converging locally uniformly to a
positive-real holomorphic function \(G\).  For each fixed \(j\), all
sufficiently late members satisfy \(G_N(z_j)=F(z_j)\), hence
\[
 G(z_j)=F(z_j).
 \tag{V68.6}
\]
The dense set \({\cal D}\) has accumulation points inside the connected
Euler half-plane.  The identity theorem gives \textup{(V68.4)}.

Both \(G\) and the meromorphic continuation of \(F\) agree on a nonempty
open set, so uniqueness of meromorphic continuation makes them identical
wherever the latter is defined.  A holomorphic \(G\) cannot equal a
function with an interior pole.  This proves the first statement; its
contrapositive gives \textup{(V68.5)}.
\end{proof}

\begin{corollary}[A semidecision procedure for failure of paired GRH]
\label{cor:v68-semidecision}
Apply Theorem~\ref{thm:v68-exhaustive-detection} to
\[
 F_{\psi,0}(z)=
 \frac{\Lambda'}{\Lambda}\left(\frac12-iz,\psi\right),
 \qquad \psi\in\{\chi,\overline\chi\}.
 \tag{V68.7}
\]
If paired GRH fails, at least one channel has a finite Euler-region Pick
matrix with a strictly negative eigenvalue.  Since all entries are
absolutely computable from the gamma factor and Euler series, sufficiently
accurate interval evaluation eventually certifies that negative eigenvalue.

Thus exhaustive growing Pick tests semidecide failure of paired GRH.
Passing every finite stage is not a proof of GRH.
\end{corollary}

\begin{proof}
An off-line zero to the right of the critical line is an upper-half-plane
pole by \textup{(V66.9)}; the functional equation places such a pole in one
of the two conjugate channels whenever paired GRH fails.  Apply
Theorem~\ref{thm:v68-exhaustive-detection}.  A strictly negative eigenvalue
is stable under sufficiently small entrywise perturbations, so rigorous
interval approximations eventually verify its sign.
\end{proof}

This result does not contradict V67.  Given a prescribed finite sample,
V67 constructs a different meromorphic function with a hidden pole and
the same table.  Theorem~\ref{thm:v68-exhaustive-detection} instead fixes
one meromorphic function and lets the sample grow without bound.

\subsection{The exact two-sample defect of one pole}

Let
\[
 p=a+ib,\qquad b>0,\qquad H_p(z)=\frac{i}{z-p}.
 \tag{V68.8}
\]
For two distinct samples \(z_j=x_j+iy_j\), put
\[
 A:=y_1+y_2,\qquad
 R:=A^2+(x_1-x_2)^2.
 \tag{V68.9}
\]

\begin{theorem}[Two samples detect an isolated interior pole]
\label{thm:v68-two-sample}
If \(y_1,y_2>b\), then each one-point Pick test for \(H_p\) is positive,
but the two-point determinant is strictly negative:
\[
 K_{H_p}(z_j,z_j)
 =\frac{y_j-b}{y_j|z_j-p|^2}>0,
 \tag{V68.10}
\]
\[
 \boxed{
 \det\left[K_{H_p}(z_j,z_k)\right]_{j,k=1}^2
 =
 -\frac{
 b(A-b)|z_1-z_2|^2}
 {y_1y_2R\,|z_1-p|^2|z_2-p|^2}<0.}
 \tag{V68.11}
\]
Hence two is the smallest sample size capable of detecting the bare pole.
\end{theorem}

\begin{proof}
Formula \textup{(V67.15)} shows that diagonal congruence by
\(\operatorname{diag}((z_j-p)^{-1})\) reduces the kernel matrix to
\[
 L_{jk}=1-\frac{2b}{-i(z_j-\overline{z_k})}.
 \tag{V68.12}
\]
Its diagonal entries are \(1-b/y_j\), proving \textup{(V68.10)} after
undoing the congruence.  Direct expansion gives
\[
 \begin{aligned}
 \det L
 &=(1-b/y_1)(1-b/y_2)
   -\left|1-\frac{2b}{A-i(x_1-x_2)}\right|^2\\
 &=-\frac{b(A-b)|z_1-z_2|^2}{y_1y_2R}.
 \end{aligned}
 \tag{V68.13}
\]
The determinant acquires the additional factor
\((|z_1-p|^2|z_2-p|^2)^{-1}\) under inverse congruence, proving
\textup{(V68.11)}.
\end{proof}

For a symmetric pair at height \(Y>b\),
\[
 z_\pm=a\pm Y+iY,
 \qquad
 D_Y:=Y^2+(Y-b)^2,
 \tag{V68.14}
\]
the two diagonal entries coincide:
\[
 d_Y=\frac{Y-b}{YD_Y},
 \tag{V68.15}
\]
and
\[
 \det{\mathsf K}_2(H_p)
 =-\frac{b(2Y-b)}{2Y^2D_Y^2}.
 \tag{V68.16}
\]
The negative eigenvalue is therefore
\[
 \boxed{
 \lambda_-(Y)
 =d_Y-\sqrt{
 d_Y^2+\frac{b(2Y-b)}{2Y^2D_Y^2}}<0.}
 \tag{V68.17}
\]
As \(Y\to\infty\),
\[
 \boxed{\lambda_-(Y)=-\frac{b}{4Y^3}
 +O_b(Y^{-4}).}
 \tag{V68.18}
\]

\begin{proof}
Equations \textup{(V68.15)}--\textup{(V68.16)} follow from
Theorem~\ref{thm:v68-two-sample}.  A Hermitian \(2\times2\) matrix with
equal diagonal \(d_Y\) and negative determinant has eigenvalues
\(d_Y\pm\sqrt{d_Y^2+|\det|}\).  Since
\[
 d_Y=\frac1{2Y^2}+O_b(Y^{-3}),
 \qquad
 |\det|=\frac{b}{4Y^5}+O_b(Y^{-6}),
 \tag{V68.19}
\]
rationalization of the square root proves \textup{(V68.18)}.
\end{proof}

Thus high sampling does not erase the bare-pole defect, but its most
negative eigenvalue decays cubically with the sampling height and linearly
with the horizontal displacement \(b=\Re\rho-1/2\).

\subsection{Euler truncation versus the model margin}

At height \(Y>1/2\), put \(\sigma=Y+1/2\).  If the Euler logarithmic
derivative is truncated at \(n\le X\), the absolute tail satisfies
\[
 \begin{aligned}
 {\cal E}_X(Y)
 &:=\sup_x\left|
 \sum_{n>X}\frac{\Lambda(n)\chi(n)}{n^\sigma}
 e^{ix\log n}\right|\\
 &\le C
 X^{1-\sigma}\left(
 \frac{\log X}{\sigma-1}
 +\frac1{(\sigma-1)^2}\right).
 \end{aligned}
 \tag{V68.20}
\]

\begin{proposition}[Precision needed for the symmetric pole test]
\label{prop:v68-truncation-margin}
Suppose both sampled values at \textup{(V68.14)} are known with absolute
error at most \(\varepsilon\).  The induced \(2\times2\) Pick-matrix
operator error is at most \(2\varepsilon/Y\).  Hence the negative eigenvalue
is rigorously detected whenever
\[
 \boxed{\frac{2\varepsilon}{Y}<|\lambda_-(Y)|.}
 \tag{V68.21}
\]
For the isolated-pole model at large \(Y\), it is sufficient that
\[
 \varepsilon=o\left(\frac{b}{Y^2}\right).
 \tag{V68.22}
\]
Using the Euler truncation alone, a sufficient fixed-height condition is
\[
 X^{-(Y-1/2)}\log X=o(b),
 \tag{V68.23}
\]
with constants depending on \(Y\).  At \(Y=1\), for example,
\[
 X\ge b^{-2}\{\log(1/b)\}^{2+\delta}
 \tag{V68.24}
\]
for any fixed \(\delta>0\) is asymptotically sufficient.
\end{proposition}

\begin{proof}
Each kernel numerator uses two sampled values, while
\[
 |-i(z_j-\overline{z_k})|\ge2Y
\tag{V68.25}
\]
for the symmetric pair.  Every entry error is at most
\(\varepsilon/Y\), so the \(2\times2\) operator norm is at most
\(2\varepsilon/Y\).  Weyl's eigenvalue inequality proves
\textup{(V68.21)}.  Insert \textup{(V68.18)} to obtain
\textup{(V68.22)}, and insert \textup{(V68.20)} for
\textup{(V68.23)}.  Formula \textup{(V68.24)} makes
\(X^{-1/2}\log X=o(b)\).
\end{proof}

The estimate is fully effective for a known isolated pole.  It is not a
uniform GRH test.  The pole ordinate \(a\), displacement \(b\), and the
negative margin are unknown, while the remaining gamma and zero background
can mask the bare two-point defect.  V67's exact interpolation construction
shows that no lower bound depending only on \(b\), \(Y\), and the number of
samples can hold over arbitrary analytic backgrounds.

\begin{theorem}[V68 growing-certificate verdict]
\label{thm:v68-verdict}
Growing Euler Pick tests have a rigorous but one-sided role.
\begin{enumerate}[label=\textup{(\roman*)},leftmargin=3em]
\item Every off-line zero forces a strictly negative finite Pick matrix on
      a dense Euler sample set.
\item Rigorous Euler truncation eventually certifies that matrix once its
      locations and negative margin are reached.
\item A bare pole is detected by two distinct samples, with the exact
      determinant \textup{(V68.11)} and high-sample eigenvalue
      \(-b/(4Y^3)+O_b(Y^{-4})\).
\item An arbitrary analytic background destroys every uniform bound on the
      required sample count or spectral margin.
\item Consequently the procedure semidecides GRH failure but cannot certify
      GRH after any finite amount of successful testing.
\end{enumerate}
\end{theorem}

\begin{firewall}[Eventual detection is not a zero-free proof]
The exhaustive theorem is a compactness argument with no effective stopping
time under GRH.  Failure to find a negative minor at any finite stage is
compatible with a hidden pole whose detecting configuration lies later.
The isolated-pole margin does not survive an uncontrolled analytic
background.  No off-line zero has been excluded, and GRH remains open.
\end{firewall}

\begin{assessment}[V68 acquisition and V69 direction]
V68 converts the finite Pick no-go into an exact semidecision theorem and
quantifies the simplest negative certificate.  This is genuine algorithmic
information, but its nonuniformity prevents a proof of GRH.

V69 should ask whether the fixed completed logarithmic derivative supplies
an a priori bound on the analytic background in a Pick-kernel norm.  The
natural test is a Schur-complement estimate separating one candidate pole
kernel from the remaining zero measure.  If classical zero counting permits
the background to mask an arbitrarily small displacement \(b\), then no
uniform stopping rule exists without a zero-free estimate already of
GRH strength.
\end{assessment}

\section*{Version V68 append-only record}
\addcontentsline{toc}{section}{Version V68 append-only record}
The complete V67 source was retained before this continuation.  It had
\(977294\) bytes, \(29052\) lines, and SHA--256
\[
 \texttt{4F6C47E92D9FFDF6B0C03E326C6714BFA87F3F41AF49130D3EDE379E7B8097FF}.
\]
V68 proved exhaustive finite detection of every hidden pole on a dense
Euler sample set, derived the exact two-sample pole determinant and its
high-sampling negative eigenvalue, and compared the margin with rigorous
Euler truncation error.

% =============================================================================
% END APPENDED V68 MATERIAL
% =============================================================================

% =============================================================================
% BEGIN APPENDED V69 MATERIAL
% =============================================================================

\section{V69: Schur masking and the insufficiency of zero counting}
\label{sec:grh-v69}

V68 proved that every fixed off-line pole has some finite negative Euler
Pick witness, but found no uniform bound for the witness in the presence of
the remaining analytic background.  This version tests whether classical
zero counting can supply such a bound.  The answer is negative.

The isolated pole kernel splits into a positive rank-one term and a negative
Szego-kernel defect of size \(2b\), where
\(b=\Re\rho-1/2\).  A positive-real background can dominate that defect by
an exact Schur-complement inequality.  Classical zero counting bounds the
background on a two-point sample only by
\(O(\log(qT)/Y)\), while the optimized isolated-pole negative eigenvalue is
only \(b/(4Y^3)\).  The available upper bound is therefore larger than the
signal by the factor \(Y^2\log(qT)/b\), and the gap worsens as an off-line
zero approaches the critical line.

\subsection{The finite pole/background kernel ledger}

Fix distinct nodes
\[
 Z=(z_1,\ldots,z_N)\subset\mathbb C_+,
 \tag{V69.1}
\]
and put
\[
 {\mathsf S}_Z:=
 \left[\frac1{-i(z_j-\overline{z_k})}\right]_{j,k=1}^N.
 \tag{V69.2}
\]
This is the positive-definite Szego matrix of the upper half-plane.  For
\(p=a+ib\in\mathbb C_+\), define
\[
 {\mathsf D}_p:=
 \operatorname{diag}\left(
 \frac1{z_1-p},\ldots,\frac1{z_N-p}\right),
 \qquad
 {\bf d}_p:=
 \left(\frac1{z_j-p}\right)_{j=1}^N.
 \tag{V69.3}
\]

\begin{proposition}[Exact pole-kernel decomposition]
\label{prop:v69-pole-decomposition}
For \(H_p(z)=i/(z-p)\), its Pick matrix on \(Z\) is
\[
 \boxed{
 {\mathsf K}_Z(H_p)
 ={\bf d}_p{\bf d}_p^*
 -2b\,{\mathsf D}_p{\mathsf S}_Z{\mathsf D}_p^*.}
 \tag{V69.4}
\]
The first term is the positive boundary-atom kernel that would occur at
\(b=0\); the second is the entire off-line defect.
\end{proposition}

\begin{proof}
Equation \textup{(V67.15)} gives entrywise
\[
 K_{H_p}(z_j,z_k)
 =\frac1{(z_j-p)(\overline{z_k}-\overline p)}
 -\frac{2b}
 {(z_j-p)(\overline{z_k}-\overline p)
 \{-i(z_j-\overline{z_k})\}}.
 \tag{V69.5}
\]
This is exactly the matrix identity \textup{(V69.4)}.
\end{proof}

Let \(B\) be any function whose Pick matrix
\({\mathsf K}_Z(B)\) is positive semidefinite.  Then
\[
 {\mathsf K}_Z(B+H_p)
 ={\mathsf K}_Z(B)+{\bf d}_p{\bf d}_p^*
 -2b{\mathsf D}_p{\mathsf S}_Z{\mathsf D}_p^*.
 \tag{V69.6}
\]
Detection requires a vector \({\bf c}\) for which the negative final term
exceeds both positive terms.  Zero counting supplies no such direction.

\subsection{An exact Schur masking threshold}

\begin{theorem}[A positive constant masks any fixed finite pole test]
\label{thm:v69-constant-mask}
For the positive-real constant function \(B_C(z)=C\), \(C>0\),
\[
 {\mathsf K}_Z(B_C)=2C{\mathsf S}_Z.
 \tag{V69.7}
\]
Define the finite generalized defect norm
\[
 \Gamma_Z(p):=
 \left\|
 {\mathsf S}_Z^{-1/2}
 {\mathsf D}_p{\mathsf S}_Z{\mathsf D}_p^*
 {\mathsf S}_Z^{-1/2}
 \right\|_{\rm op}.
 \tag{V69.8}
\]
If
\[
 \boxed{C\ge b\,\Gamma_Z(p),}
 \tag{V69.9}
\]
then
\[
 \boxed{{\mathsf K}_Z(B_C+H_p)\succeq0.}
 \tag{V69.10}
\]
If the inequality in \textup{(V69.9)} is strict, the masking is stable
under sufficiently small perturbations of every sampled value.
\end{theorem}

\begin{proof}
By congruence with \({\mathsf S}_Z^{-1/2}\), condition
\textup{(V69.9)} is exactly
\[
 2C{\mathsf S}_Z
 -2b{\mathsf D}_p{\mathsf S}_Z{\mathsf D}_p^*
 \succeq0.
 \tag{V69.11}
\]
Adding the positive matrix
\({\bf d}_p{\bf d}_p^*\) and applying
\textup{(V69.4)} proves \textup{(V69.10)}.  Strict positivity has a
positive smallest eigenvalue and is stable under small Hermitian
perturbations.
\end{proof}

This theorem does not contradict the exhaustive result of V68.  The
constant needed to mask a \emph{fixed} node set changes when the set grows.
No one finite \(C\) is asserted to make every Pick matrix positive while
retaining the pole.  The result proves precisely that a finite pole margin
cannot be bounded below without controlling the background.

\subsection{Positive boundary atoms can approximate the mask}

The constant kernel is not alien to the Herglotz cone.  A positive-real
constant has a Herglotz representation with a positive absolutely
continuous boundary measure, and its kernel is a positive superposition of
the rank-one boundary-atom kernels
\[
 {\bf v}_\gamma{\bf v}_\gamma^*,
 \qquad
 {\bf v}_\gamma:=
 \left(\frac1{z_j-\gamma}\right)_{j=1}^N,
 \qquad \gamma\in\mathbb R.
 \tag{V69.12}
\]

\begin{proposition}[Discrete critical-line masking capacity]
\label{prop:v69-discrete-mask}
For every fixed finite sample \(Z\), every \(C>0\), and every
\(\varepsilon>0\), there are finitely many real points
\(\gamma_\ell\) and positive weights \(w_\ell\) such that
\[
 \left\|
 2C{\mathsf S}_Z
 -\sum_\ell w_\ell
 {\bf v}_{\gamma_\ell}{\bf v}_{\gamma_\ell}^*
 \right\|_{\rm op}<\varepsilon.
 \tag{V69.13}
\]
Consequently every strict constant mask in
Theorem~\ref{thm:v69-constant-mask} can be replaced by a finite positive
boundary-atom background while preserving positivity of the sampled Pick
matrix.
\end{proposition}

\begin{proof}
The Herglotz representation expresses \(2C{\mathsf S}_Z\) as the integral
of the continuous matrix-valued function
\({\bf v}_\gamma{\bf v}_\gamma^*\) against a positive measure.  First
truncate the integral, whose tails are uniformly integrable on the finite
node set, and then approximate the compact integral by a positive Riemann
sum.  This proves \textup{(V69.13)}.  If
\textup{(V69.9)} is strict, choose \(\varepsilon\) below the positive
spectral margin and apply Weyl's inequality.
\end{proof}

The points \(\gamma_\ell\) model critical-line zeros.  The following
translation observation makes the compatibility with ordinary counting
size precise.

\begin{corollary}[Integer-atom masks within a zero-counting envelope]
\label{cor:v69-integer-mask}
For a strict finite mask one may enlarge \(C\) arbitrarily.  For all
sufficiently large \(C\), the approximation in
\textup{(V69.13)} can be made with unit integer weights by a positive
quadrature of mesh \(O(C^{-1})\).  Translating the pole, sample nodes, and
all boundary atoms by the same real number \(U\) preserves every kernel
matrix.  If \(\log(2+|U|)\gg C\), the resulting local atom density is
compatible with the classical envelope
\(N(U+1)-N(U)=O(\log(2+|U|))\).
\end{corollary}

\begin{proof}
Use the positive boundary density representing \(2C{\mathsf S}_Z\) and
partition a large compact truncation into intervals having measure one.
Choosing one quadrature point in each interval gives unit weights; as
\(C\to\infty\), the mesh tends to zero and the matrix Riemann sums converge.
The strict masking margin absorbs the truncation and quadrature errors.
Real translation leaves every difference \(z_j-\gamma\),
\(z_j-p\), and \(z_j-\overline{z_k}\) unchanged.  The unit-interval atom
count is \(O(C)\), which is below the stated classical allowance after the
indicated translation.
\end{proof}

This remains a capacity theorem, not a claim that the actual Dirichlet zero
multiset occupies the selected locations.  It proves that positivity,
integer multiplicity, and ordinary zero-counting size are jointly compatible
with finite masking.

\subsection{The classical zero-counting scale}

Let \({\mathfrak Z}_0\) denote the critical-line part of a paired zero
multiset and consider its positive kernel
\[
 {\mathsf B}_0
 :=\sum_{\gamma\in{\mathfrak Z}_0}m_\gamma
 {\bf v}_\gamma{\bf v}_\gamma^*.
 \tag{V69.14}
\]
The standard local consequence of
\(N_\chi(T)=O_\chi(T\log(2+T))\) is
\[
 N_\chi(U+1)-N_\chi(U)
 =O_\chi(\log(2+|U|)).
 \tag{V69.15}
\]

\begin{lemma}[Zero-counting background bound]
\label{lem:v69-background-bound}
For \(z=x+iY\) with \(Y\ge1\),
\[
 \boxed{
 \sum_{\gamma\in{\mathfrak Z}_0}
 \frac{m_\gamma}{(x-\gamma)^2+Y^2}
 \ll_\chi
 \frac{\log(2+|x|+Y)}{Y}.}
 \tag{V69.16}
\]
For a fixed-size sample configuration at common height \(Y\), centered at
ordinates of size at most \(T\),
\[
 \|{\mathsf B}_0\|_{\rm op}
 \ll_{\chi,N}
 \frac{\log(2+T+Y)}{Y}.
 \tag{V69.17}
\]
\end{lemma}

\begin{proof}
Partition the real line into unit intervals according to distance from
\(x\), apply \textup{(V69.15)} on each interval, and sum the resulting
majorant
\[
 \sum_{k\in\mathbb Z}
 \frac{\log(2+|x|+Y+|k|)}
 {k^2+Y^2}.
 \tag{V69.18}
\]
Comparison with the corresponding integral gives
\textup{(V69.16)}.  Every diagonal entry of
\({\mathsf B}_0\) has this form.  Positivity gives
\(\|{\mathsf B}_0\|_{\rm op}\le\operatorname{tr}{\mathsf B}_0\), proving
\textup{(V69.17)}.
\end{proof}

At the optimized symmetric two-point sample of V68, a bare pole of height
\(b\) has
\[
 |\lambda_-(Y)|=\frac{b}{4Y^3}+O_b(Y^{-4}).
 \tag{V69.19}
\]
The zero-counting background bound instead has size
\[
 O_\chi\left(\frac{\log(2+T+Y)}Y\right).
 \tag{V69.20}
\]
Their ratio is bounded only by the nonvanishing scale
\[
 \boxed{
 O_\chi\left(
 \frac{Y^2\log(2+T+Y)}b
 \right).}
 \tag{V69.21}
\]
It grows with \(Y\), with the ordinate range, and as \(b\downarrow0\).

\begin{theorem}[Zero counting cannot isolate the Pick defect]
\label{thm:v69-counting-no-go}
An argument which controls the remaining zero background only through
\textup{(V69.15)} or the Poisson bound \textup{(V69.16)} cannot guarantee a
negative Schur complement for an off-line pole.
\begin{enumerate}[label=\textup{(\roman*)},leftmargin=3em]
\item At high symmetric samples, the available background norm exceeds the
      isolated-pole negative margin by the factor
      \textup{(V69.21)}.
\item At bounded Euler height, the pole margin is \(O(b)\), while the
      counting bound does not tend to zero with \(b\).
\item Positive integer boundary-atom configurations within ordinary counting
      envelopes can mask every prescribed finite pole test by
      Corollary~\ref{cor:v69-integer-mask}.
\end{enumerate}
Therefore zero counting supplies neither a uniform negative eigenvalue nor
a finite stopping rule for the exhaustive procedure of V68.
\end{theorem}

\begin{proof}
Items \textup{(i)} and \textup{(ii)} are the direct comparison of
\textup{(V69.16)} with the exact V68 pole margin.  They show that the
sufficient perturbative detection condition
\[
 \|{\mathsf K}_Z(B)\|_{\rm op}<|\lambda_{\min}
 {\mathsf K}_Z(H_p)|
 \tag{V69.22}
\]
cannot follow from the classical bound, uniformly over
\(0<b<1/2\).  Item \textup{(iii)} shows this mismatch is genuine capacity,
not merely a loose use of Weyl's inequality.
\end{proof}

\subsection{What a successful background estimate would require}

The exact detection condition is not an unsigned count.  From
\textup{(V69.6)}, one needs a vector \({\bf c}\) such that
\[
 {\bf c}^*{\mathsf K}_Z(B){\bf c}
 +|{\bf d}_p^*{\bf c}|^2
 <
 2b\,{\bf c}^*
 {\mathsf D}_p{\mathsf S}_Z{\mathsf D}_p^*{\bf c}.
 \tag{V69.23}
\]
This is a directional cancellation statement for the complete background
kernel.  Trace bounds, local zero counts, and absolute Poisson envelopes
discard precisely that direction.

\begin{theorem}[V69 Schur-complement verdict]
\label{thm:v69-verdict}
The V68 nonuniformity cannot be removed with classical zero counting.
\begin{enumerate}[label=\textup{(\roman*)},leftmargin=3em]
\item The off-line pole contributes the exact negative Szego defect
      \(-2b{\mathsf D}_p{\mathsf S}_Z{\mathsf D}_p^*\).
\item A positive-real background masks every fixed finite test above the
      explicit threshold \(C=b\Gamma_Z(p)\).
\item Finite positive critical-line atom sums approximate such masks.
\item Classical zero counting bounds the background far above, not below,
      the isolated-pole margin, especially as \(b\downarrow0\).
\item A useful estimate must control the directional quadratic form in
      \textup{(V69.23)}, not merely the number or absolute mass of zeros.
\end{enumerate}
\end{theorem}

\begin{firewall}[Masking capacity is not the actual zero configuration]
The constant and discretized backgrounds are comparison models, not claims
about the zeros of a particular Dirichlet \(L\)-function.  They prove that
the information retained by classical zero counting is insufficient to
force detection at a prescribed finite stage.  The exhaustive semidecision
of V68 remains valid, but it has no uniform stopping time.  GRH remains
open.
\end{firewall}

\begin{assessment}[V69 closure and the mandatory V70 sweep]
V69 carries out the Schur-complement test proposed in V68.  It identifies
the exact negative pole direction and proves that classical zero-density
input cannot separate it from the positive background.

V70 is the next scheduled ten-version corpus sweep.  It should hash-gate all
parallel notes against the V60 cutoff, inspect every successor created since
then, and test whether any new companion result supplies the directional
kernel control \textup{(V69.23)} or a genuinely coefficient-sensitive
continuation estimate.  If none does, the post-sweep attack should move
upstream from finite Pick matrices to the fixed Euler coefficient relation
for the full kernel.
\end{assessment}

\section*{Version V69 append-only record}
\addcontentsline{toc}{section}{Version V69 append-only record}
The complete V68 source was retained before this continuation.  It had
\(989593\) bytes, \(29411\) lines, and SHA--256
\[
 \texttt{640A2A6CD8F8C48B51046C9C588A36B6D9B1CA53A325A4197E0245EA76839DC8}.
\]
V69 decomposed the pole Pick kernel into positive and negative parts,
derived the exact constant-background masking threshold, approximated masks
by positive boundary atoms, and proved that classical zero counting cannot
separate the pole defect from the analytic background.

% =============================================================================
% END APPENDED V69 MATERIAL
% =============================================================================
% =============================================================================
% BEGIN APPENDED V70 MATERIAL
% =============================================================================

\section{V70: mandatory corpus sweep and the directional-kernel reset}
\label{sec:grh-v70}

V69 ended at a genuine information boundary.  A single off-line pole has an
explicit negative Pick direction, but a positive critical-line background can
mask every prescribed finite test, and classical zero counting does not
control the required Schur complement.  This version first performs the
scheduled ten-version audit.  It then tests the only new companion machinery
against the precise V69 bottleneck and records an exact quadratic-form
identity which fixes the post-audit direction.

\subsection{Hash-gated V60--V70 corpus audit}

The V60 audit used the UTC cutoff
\[
 \texttt{2026-09-11T03:16:15.8841028Z}
 \tag{V70.1}
\]
and a sorted 31-file non-GRH manifest with SHA--256
\[
 \texttt{D4CF7A659386ECDD8362C956915299269FA4E927F1052B010BF3C198E4FE80EF}.
 \tag{V70.2}
\]
For the present audit, each manifest row was again
\[
 \texttt{filename|byte-count|SHA--256}\backslash\texttt{n}.
 \tag{V70.3}
\]
The rows were sorted by filename, joined with LF terminators, and the
resulting UTF--8 byte string was hashed.  Two complete passes ended at
\texttt{04:13:23.7527975Z} and \texttt{04:13:37.7953693Z} on
\texttt{2026-09-11}, respectively.  They
both found 31 non-GRH \texttt{.tex} files and produced
\[
 \boxed{
 \texttt{56B7195509E847D2F16E08DF3580056CD5EB5C88D73A28D2D7DCA7B885E84AC6}.}
 \tag{V70.4}
\]
Several earlier passes were discarded because a companion writer advanced
its rolling file between enumeration and hashing.  No conclusion below uses
one of those mixed snapshots.

The live manifest differs from the V60 manifest only at the following two
rolling endpoints:
\begin{center}
\tiny
\begin{tabular}{p{0.48\textwidth}|r|p{0.33\textwidth}}
\toprule
file & bytes & SHA--256 \\
\midrule
\path{Renewal_SM_Einstein_Continuum_Research_Notes_V32.tex}
 & 300957 & \texttt{13A7201F618C149E6CA94105006528E74}\newline
 \texttt{7CE1FA3C20A553652B6885F4C4B5514} \\
\path{Renewal_Yang_Mills_Mass_Gap_Research_Notes_V48.tex}
 & 419467 & \texttt{740675B97D48DF77BB6F24D55D1C2E13}\newline
 \texttt{E33404DB639113BBFA4DD6326D7801FE} \\
\bottomrule
\end{tabular}
\end{center}
The unchanged long SM endpoint relevant to the previous sweep is
\path{Renewal_SM_Einstein_Continuum_Research_Notes_V100_f15.tex},
with 988883 bytes and SHA--256
\[
 \texttt{9B4C6C8CE2E53B5270A95E99AD8CE7E622CAD1645B007B37AE885DD0879CEDCE}.
 \tag{V70.5}
\]
As an independent continuity check, removing the two live endpoints from
the V70 row list and inserting the two archived rows below
\begin{center}
\tiny
\begin{tabular}{p{0.48\textwidth}|r|p{0.33\textwidth}}
\path{Renewal_SM_Einstein_Continuum_Research_Notes_V13.tex}
 &132567&\texttt{ED2799DF0A0442230D71D514E2A0CFE}\newline
 \texttt{A156A0131C89D12513B76A35D98B6896F}\\
\path{Renewal_Yang_Mills_Mass_Gap_Research_Notes_V46.tex}
 &396051&\texttt{B1E5580A117A5BED92CDB4AC94DDEB6}\newline
 \texttt{E5992F93E18D5DB46F0A38FAC6D3E5469}
\end{tabular}
\end{center}
reproduces exactly the 31-row digest in \textup{(V70.2)}.  Thus the actual
successor ranges are SM V14--V32 and Yang--Mills V47--V48; there is neither a
gap nor a silent change in any other row.

\subsection{Typed transfer from the SM successor}

The SM successor develops a coherent conditional chain: product
R\'enyi-volume growth and local tilts; positive block Schur complements;
Newton collar/far splittings; projected Gaussian and \(L^p\) densities;
soft-mode normalization; variational Schur identities; Mosco convergence;
and trace-norm convergence of retained influences.  Its most portable
operator statements have the following shape.  For a positive block
operator
\[
 \mathcal A_n=
 \begin{pmatrix}A_n&B_n\\B_n^*&C_n\end{pmatrix}\succeq0,
 \qquad
 J_n:=B_n^*A_n^+B_n,
 \tag{V70.6}
\]
one obtains \(0\le J_n\le C_n\).  Under fixed-rank soft-mode transport,
Mosco convergence, equicoercivity after retaining the soft coordinate, and
a trace-class domination such as \(C_n\to C\) in trace norm, weak convergence
of \(J_n\) can be upgraded to trace-norm convergence.  The final Gaussian
determinant and projected \(L^p\)-moment conclusions additionally spend a
strict spectral reserve.

These results are mathematically substantial, but their hypotheses do not
match the GRH obstruction.

\begin{theorem}[Positive Schur machinery does not transfer to the V69 defect]
\label{thm:v70-schur-nontransfer}
The SM V14--V32 operator chain does not supply the directional estimate
\textup{(V69.23)}.  More precisely:
\begin{enumerate}[label=\textup{(\roman*)},leftmargin=3em]
\item its order bound \(0\le J_n\le C_n\) is derived from positivity of the
      full block \(\mathcal A_n\), whereas positivity of the arithmetic Pick
      block in the lower strip is equivalent to the missing GRH conclusion;
\item the pole term
      \(-2b\mathsf D_p\mathsf S_Z\mathsf D_p^*\) is indefinite after the
      rank-one positive part is restored, and the ordinate background is not
      provided with the required global trace-class domination;
\item the strict spectral reserve used for projected-density stability is
      precisely absent at a first Pick-cone contact;
\item retaining an off-line soft coordinate records the pole displacement
      and its coupling, but gives no sign forcing that displacement to vanish.
\end{enumerate}
Consequently importing the positive Schur or Mosco conclusions would either
assume the desired positivity or replace the arithmetic continuation by a
different positive comparison problem.
\end{theorem}

\begin{proof}
For a positive block, testing \(\mathcal A_n\) on
\((-A_n^+B_n y,y)\) gives the Schur inequality
\(0\le J_n\le C_n\).  In the GRH problem, the analogous full Pick matrix is
the object whose possible negative direction is under investigation, so the
same test cannot be invoked before its premise is proved.  V69 gives the
exact signed pole decomposition, not a positive influence operator.  Its
critical-line background is controlled only through local Poisson mass,
which V69 showed to be much larger than the pole margin and which is not a
trace-norm convergence hypothesis on the unbounded ordinate axis.  Finally,
the determinant and \(L^p\) arguments require a uniform gap; at contact that
gap is zero.  These are type mismatches in the hypotheses, not missing
technical estimates.
\end{proof}

The earlier companion warning therefore survives: positive comparison
density controls absolute size, while the missing estimate is signed and
directional and must use the fixed Dirichlet coefficients.

\subsection{Typed transfer from the Yang--Mills successor}

Yang--Mills V47 repairs two false isotropic charts by anisotropic ones and
reaches paired two-step prisms without obtaining an orientation-transport
principle.  V48 closes one opposite transverse quarter-cell pair with 16
explicit power-three charts and rational floors.  The resulting direct and
inverse Wilson response pencils are positive on that paired inversion cell;
the other three inversion pairs and the full torus remain open.

This supplies a useful methodological warning: an asserted symmetry
transport must be checked chart by chart, and anisotropic repair can be
essential.  It supplies no GRH transfer.  A finite Wilson-pencil positivity
statement has no Dirichlet-coefficient identity, no full-domain passage, and
no estimate on the V69 negative direction.  Treating one paired cell as a
global symmetry certificate would repeat exactly the finite-to-global error
closed in V67--V69.

\begin{theorem}[V70 companion-corpus verdict]
\label{thm:v70-corpus-verdict}
No successor created after the V60 cutoff proves or strengthens a hypothesis
which closes the V69 Schur-complement inequality.  The transferable content
is diagnostic:
\begin{enumerate}[label=\textup{(\roman*)},leftmargin=3em]
\item positive block convergence cannot be used before arithmetic Pick
      positivity is known;
\item strict-reserve stability cannot decide a zero-gap contact;
\item retaining a soft coordinate describes an obstruction but does not
      exclude it;
\item finite paired-chart positivity does not propagate without a proved
      transport law.
\end{enumerate}
Hence the post-sweep attack must remain inside the arithmetic coefficient
relation rather than import a generic positive-operator compactness theorem.
\end{theorem}

\subsection{The post-V60 closure ledger}

For clarity, the nine intervening versions closed the following routes.
\begin{center}
\small
\begin{tabular}{c p{0.78\linewidth}}
\toprule
version & rigorous outcome \\
\midrule
V61 & weighted \(H^2\) factorization anchors the outer factor in the Euler
region, but outer boundary zeros do not forbid an off-line inner zero; \\
V62 & the corrected radial Hardy estimate allows a localized outer node to
carry vanishing negative-log charge; \\
V63 & the sharp absolute von Mangoldt moment ratio has Bernstein scale
\(R_{r+1}/R_r\sim(4t)^{-1}\), exactly the scale of the forbidden lobes; \\
V64 & two complex frequencies give independent value/derivative freedom, so
a node-only regression estimate is logically weaker than the required
global node-or-sign-defect alternative; \\
V65 & phase-flux quadrant balance is not contour invariant, and nodal
multiplicity remains separate data; \\
V66 & the quarter-turned completed logarithmic derivative is Pick in both
paired channels exactly under paired GRH; heat preservation is the same Weil
positivity condition in different language; \\
V67 & every finite value or jet table can be preserved by a rational
perturbation with a hidden lower-strip pole; full Euler-region positivity
does not continue across the unsampled strip; \\
V68 & dense rational Pick sampling semidecides failure and the bare pole has
an explicit two-sample negative eigenvalue, but no uniform stopping time was
obtained; \\
V69 & exact Schur masking and positive boundary-atom approximation show that
classical zero counting is too weak by the scale
\(O_\chi(Y^2\log(2+T+Y)/b)\). \\
\bottomrule
\end{tabular}
\end{center}
The common failure is now sharp.  Absolute envelopes, compactness, finite
interpolation, and generic positivity all erase the arithmetic direction in
which an off-line pole is negative.

\subsection{An exact directional form of the Pick test}

There is a simple identity which packages that direction without taking an
operator norm.  It is elementary, but it specifies exactly where the Euler
coefficients must enter next.

\begin{lemma}[Hardy integral for a finite Pick quadratic form]
\label{lem:v70-hardy-pick-form}
Let \(F\) be defined at \(Z=(z_1,\ldots,z_N)\subset\mathbb C_+\), let
\({\bf c}=(c_1,\ldots,c_N)^T\), and put
\[
 A_{Z,{\bf c}}(u):=\sum_{j=1}^N c_j e^{iz_j u},
 \qquad
 B_{F,Z,{\bf c}}(u):=\sum_{j=1}^N c_jF(z_j)e^{iz_j u}.
 \tag{V70.7}
\]
Then
\[
 \boxed{
 {\bf c}^*\mathsf K_Z(F){\bf c}
 =2\Re\int_0^\infty
 B_{F,Z,{\bf c}}(u)\overline{A_{Z,{\bf c}}(u)}\,du.}
 \tag{V70.8}
\]
If, at the sample points, an absolutely convergent expansion
\(F=G+\sum_\nu a_\nu\phi_\nu\) is available, then the right side decomposes
termwise as
\[
 {\bf c}^*\mathsf K_Z(G){\bf c}
 +2\Re\sum_\nu a_\nu
 \int_0^\infty
 \left(\sum_jc_j\phi_\nu(z_j)e^{iz_ju}\right)
 \overline{A_{Z,{\bf c}}(u)}\,du.
 \tag{V70.9}
\]
\end{lemma}

\begin{proof}
Since \(z_j,z_k\in\mathbb C_+\),
\[
 \frac1{-i(z_j-\overline{z_k})}
 =\int_0^\infty e^{iz_ju}\overline{e^{iz_ku}}\,du.
 \tag{V70.10}
\]
Insert this identity into the definition of \(\mathsf K_Z(F)\).  The terms
containing \(F(z_j)\) give
\(\int B\overline A\), while those containing
\(\overline{F(z_k)}\) give its complex conjugate.  Absolute convergence at
the finite sample set justifies the termwise expansion in
\textup{(V70.9)}.
\end{proof}

Unlike a trace or zero count, \textup{(V70.8)} preserves the test vector.
Unlike a generic interpolation theorem, \textup{(V70.9)} preserves the
fixed arithmetic coefficients.  It does not prove positivity: in the Euler
half-plane it merely rewrites a computable quantity, and continuing the
coefficient expansion into the lower strip is forbidden.  Its value is to
identify the legitimate target.  One must derive, from functional equation
and Euler data together, a bound for the particular correlation
\[
 \Re\int_0^\infty B_{F,Z,{\bf c}}(u)
                 \overline{A_{Z,{\bf c}}(u)}\,du
 \tag{V70.11}
\]
in the V69 pole-negative direction \({\bf c}\), without replacing the
coefficients by their absolute values and without first assuming Pick
positivity.

\begin{theorem}[Directional-kernel reset]
\label{thm:v70-directional-reset}
After the V60--V70 audit, the next nonduplicative branch is the following.
Start with the exact completed Euler formula at samples strictly inside the
absolute-convergence half-plane, insert it into \textup{(V70.9)}, and retain
the prime and prime-power phases in the directional quadratic form.  Then
test whether the functional equation supplies a coefficient-sensitive
continuation or comparison across the line \(\Im z=1/2\).

Any proposed estimate is admissible only if it satisfies all three checks:
\begin{enumerate}[label=\textup{(\roman*)},leftmargin=3em]
\item it is directional in \({\bf c}\), rather than a trace, counting, or
      absolute-moment bound;
\item it uses the fixed Euler coefficients, rather than an arbitrary
      positive-real interpolant or boundary-atom model;
\item it does not assume positivity, a zero-free continuation, or a strict
      spectral reserve in the lower strip.
\end{enumerate}
Failure of these checks routes the argument back to one of the closed
V61--V69 branches.
\end{theorem}

\begin{firewall}[Audit and reset are not GRH progress by implication]
The corpus audit proves provenance and prevents duplicate or circular
imports.  Lemma~\ref{lem:v70-hardy-pick-form} is an exact re-expression of a
finite Pick form, not a sign theorem and not a continuation principle.  No
new lower-strip positivity has been proved.  GRH remains open.
\end{firewall}

\begin{assessment}[V70 closure and V71 direction]
V70 completes the mandatory sweep on a stable, twice-hashed snapshot and
verifies continuity back to the V60 manifest.  The new SM Schur/Mosco chain
and Yang--Mills anisotropic chart repair are non-importable for precise typed
reasons.  Together with the V61--V69 ledger, this rules out another cycle of
generic positivity or absolute-background estimates.

V71 should expand \textup{(V70.9)} for the actual completed logarithmic
derivative in the Euler half-plane, separating gamma/conductor, prime, and
prime-power terms.  Its first task is algebraic: derive the exact
coefficient-weighted matrix or integral for a general finite test vector and
identify which phase correlations survive the paired-character functional
equation.  Only after that identity is fixed should one seek an inequality.
\end{assessment}

\section*{Version V70 append-only record}
\addcontentsline{toc}{section}{Version V70 append-only record}
The complete V69 source was retained before this continuation.  It had
1003582 bytes, 29810 lines, and SHA--256
\[
 \texttt{6C30381E9E063E875AEACA7751487C8F86CD34EFBF727F748629FB257EC7895F}.
\]
V70 hash-gated the mandatory companion sweep, proved the typed non-transfer
of the new positive Schur/Mosco and finite-chart results, consolidated the
V61--V69 closure ledger, and reset the attack to the exact
coefficient-sensitive directional Pick form.

% =============================================================================
% END APPENDED V70 MATERIAL
% =============================================================================
% =============================================================================
% BEGIN APPENDED V71 MATERIAL
% =============================================================================

\section{V71: the exact Euler shift form and the prime-layer obstruction}
\label{sec:grh-v71}

V70 reset the attack to the coefficient-sensitive quadratic form of the
completed logarithmic derivative.  This version performs that algebra.  A
conjugation slip in the displayed V70 test vector is corrected first.  The
correct formula then separates the conductor, gamma, prime, and higher
prime-power contributions exactly.  The separation gives a genuine
narrowing: the gamma term and every \(p^m\)-layer with \(m\ge2\) are
holomorphic throughout the target half-plane \(\Im z>0\); the only Euler
series that meets the line \(\Im z=1/2\) is the first-prime layer.

Throughout, \(\chi\) is a primitive nonprincipal Dirichlet character of
conductor \(f\) and parity \(\kappa\in\{0,1\}\).  The principal zeta channel
requires the standard separate pole terms already firewalled in V11 and V22.
Put
\[
 s(z):=\frac12-iz,
 \qquad
 F_{\chi,0}(z):=\frac{\Lambda'}{\Lambda}(s(z),\chi).
 \tag{V71.1}
\]

\subsection{Correction of the V70 coefficient convention}

Let \(Z=(z_1,\ldots,z_N)\subset\mathbb C_+\), let
\({\bf c}=(c_1,\ldots,c_N)^T\), and retain
\[
 [\mathsf K_Z(F)]_{jk}
 =\frac{F(z_j)+\overline{F(z_k)}}{-i(z_j-\overline{z_k})}.
 \tag{V71.2}
\]

\begin{remark}[V70 conjugation correction]
\label{rem:v71-v70-correction}
For the conventional quadratic form \({\bf c}^*\mathsf K_Z(F){\bf c}\),
the functions in \textup{(V70.7)} must be
\[
 A_{Z,{\bf c}}(u):=\sum_{j=1}^N\overline{c_j}e^{iz_ju},
 \qquad
 B_{F,Z,{\bf c}}(u):=
 \sum_{j=1}^N\overline{c_j}F(z_j)e^{iz_ju}.
 \tag{V71.3}
\]
With these definitions,
\[
 \boxed{
 {\bf c}^*\mathsf K_Z(F){\bf c}
 =2\Re\int_0^\infty
 B_{F,Z,{\bf c}}(u)\overline{A_{Z,{\bf c}}(u)}\,du.}
 \tag{V71.4}
\]
The uncorrected V70 display computes the same formula with the test vector
\(\overline{\bf c}\).  Since coordinatewise conjugation is a bijection, none
of the V70 positivity or non-transfer conclusions changes, but the direction
attached to a named vector must use \textup{(V71.3)}.
\end{remark}

Indeed,
\[
 \frac1{-i(z_j-\overline{z_k})}
 =\int_0^\infty e^{iz_ju}e^{-i\overline{z_k}u}\,du,
 \tag{V71.5}
\]
and direct expansion of \({\bf c}^*\mathsf K_Z(F){\bf c}\) proves
\textup{(V71.4)}.

Define the energy and one-sided autocorrelation
\[
 E_{Z,{\bf c}}:=\int_0^\infty|A_{Z,{\bf c}}(t)|^2dt,
 \qquad
 C_{Z,{\bf c}}(u):=
 \int_0^\infty A_{Z,{\bf c}}(t+u)
             \overline{A_{Z,{\bf c}}(t)}\,dt.
 \tag{V71.6}
\]
Then \(E_{Z,{\bf c}}=C_{Z,{\bf c}}(0)\ge0\), and a single exponential
\(e^{i\lambda z}\) contributes exactly
\[
 2\Re C_{Z,{\bf c}}(\lambda)
 \tag{V71.7}
\]
to the quadratic form when its coefficient is one.  More generally a
coefficient \(a\in\mathbb C\) contributes
\(2\Re\{aC_{Z,{\bf c}}(\lambda)\}\).

\subsection{The gamma term as an autocorrelation integral}

Write
\[
 \alpha_\kappa:=\frac12+\kappa,
 \qquad
 d_f:=\frac12\log\frac f\pi-\frac{\gamma}{2},
 \tag{V71.8}
\]
where \(\gamma\) is Euler's constant.  The standard integral for the
digamma function gives, for every \(z\in\mathbb C_+\),
\[
 \begin{aligned}
 A_{\chi,0}(z)
 &:=\frac12\log\frac f\pi
   +\frac12\psi\!\left(
      \frac{\frac12+\kappa-iz}{2}\right)\\
 &=d_f+\int_0^\infty
 \frac{e^{-2u}-e^{-\alpha_\kappa u}e^{izu}}
      {1-e^{-2u}}\,du.
 \end{aligned}
 \tag{V71.9}
\]
The two numerator terms in this integral must remain paired at \(u=0\).

\begin{lemma}[Exact archimedean quadratic form]
\label{lem:v71-archimedean-form}
For every finite \(Z\subset\mathbb C_+\) and every \({\bf c}\), define
\[
 \begin{aligned}
 \mathfrak A_{f,\kappa}[C]
 &:=(\log(f/\pi)-\gamma)E\\
 &\quad+2\int_0^\infty
 \frac{e^{-2u}E-e^{-\alpha_\kappa u}\Re C(u)}
      {1-e^{-2u}}\,du,
 \end{aligned}
 \tag{V71.10}
\]
where \(E=C(0)\) and \(C=C_{Z,{\bf c}}\).  Then the integral converges and
\[
 {\bf c}^*\mathsf K_Z(A_{\chi,0}){\bf c}
 =\mathfrak A_{f,\kappa}[C].
 \tag{V71.11}
\]
\end{lemma}

\begin{proof}
Insert \textup{(V71.9)} into \textup{(V71.4)}.  A real constant \(r\)
contributes \(2rE\), while
\(-q e^{iuz}\), with \(q\) real, contributes
\(-2q\Re C(u)\).  This gives \textup{(V71.10)}.
Near zero, differentiability of the finite exponential sum gives
\(C(u)=E+uC'(0)+O(u^2)\); hence the numerator in
\textup{(V71.10)} is \(O(u)\), cancelling
\(1-e^{-2u}\sim2u\).  At infinity both terms decay exponentially.
\end{proof}

Thus the archimedean term is already a completely explicit functional of
the same directional autocorrelation that samples the primes.  It has no
singularity anywhere in \(\mathbb C_+\).

\subsection{Exact prime and prime-power separation}

For \(\Im z>1/2\), absolute convergence gives
\[
 P_{\chi,0}^{[0]}(z)
 =\sum_p\sum_{m\ge1}
   \frac{(\log p)\chi(p)^m}{p^{m/2}}
   e^{izm\log p}
 =-\frac{L'}{L}(s(z),\chi).
 \tag{V71.12}
\]
Define
\[
 P_\chi^{(1)}(z)
 :=\sum_p\frac{(\log p)\chi(p)}{\sqrt p}e^{iz\log p},
 \tag{V71.13}
\]
and
\[
 P_\chi^{(\ge2)}(z)
 :=\sum_p\sum_{m\ge2}
 \frac{(\log p)\chi(p)^m}{p^{m/2}}e^{izm\log p}.
 \tag{V71.14}
\]

\begin{theorem}[Exact Euler shift formula]
\label{thm:v71-euler-shift}
Suppose \(y_*:=\min_j\Im z_j>1/2\).  With \(C=C_{Z,{\bf c}}\),
\[
 \boxed{
 \begin{aligned}
 {\bf c}^*\mathsf K_Z(F_{\chi,0}){\bf c}
 &=\mathfrak A_{f,\kappa}[C]\\
 &\quad-2\Re\sum_p\sum_{m\ge1}
 \frac{(\log p)\chi(p)^m}{p^{m/2}}
 C(m\log p).
 \end{aligned}}
 \tag{V71.15}
\]
Equivalently, the arithmetic part is the sum of the first-prime form
\[
 \mathfrak P_\chi^{(1)}[C]
 :=-2\Re\sum_p
 \frac{(\log p)\chi(p)}{\sqrt p}C(\log p)
 \tag{V71.16}
\]
and the higher-power form
\[
 \mathfrak P_\chi^{(\ge2)}[C]
 :=-2\Re\sum_p\sum_{m\ge2}
 \frac{(\log p)\chi(p)^m}{p^{m/2}}C(m\log p).
 \tag{V71.17}
\]
All sums and exchanges in \textup{(V71.15)} are absolute.
\end{theorem}

\begin{proof}
For \(y_*>0\), the finite exponential sum in \textup{(V71.3)} satisfies
\[
 |C(u)|\le
 \frac{(\sum_j|c_j|)^2}{2y_*}e^{-y_*u}.
 \tag{V71.18}
\]
Consequently the absolute majorant for \textup{(V71.15)} is a constant
times
\[
 \sum_p\sum_{m\ge1}(\log p)
 p^{-m(1/2+y_*)},
 \tag{V71.19}
\]
which converges when \(y_*>1/2\).  Applying
\textup{(V71.7)} term by term to
\(F_{\chi,0}=A_{\chi,0}-P_{\chi,0}^{[0]}\) proves the formula.
\end{proof}

The formula has a useful scalar check.  For \(N=1\), \(z=iY\), and
\(c=1\),
\[
 A(t)=e^{-Yt},\qquad E=\frac1{2Y},\qquad
 C(u)=\frac{e^{-Yu}}{2Y},
 \tag{V71.20}
\]
and \textup{(V71.15)} becomes exactly
\[
 {\bf c}^*\mathsf K_{\{iY\}}(F_{\chi,0}){\bf c}
 =\frac{\Re F_{\chi,0}(iY)}{Y}.
 \tag{V71.21}
\]

\subsection{Only the first-prime layer reaches the barrier}

The split at \(m=1\) is not cosmetic.

\begin{theorem}[Prime-layer localization of the continuation obstruction]
\label{thm:v71-prime-layer}
The following statements hold.
\begin{enumerate}[label=\textup{(\roman*)},leftmargin=3em]
\item The archimedean function \(A_{\chi,0}\) is holomorphic on
      \(\mathbb C_+\).
\item The series \(P_\chi^{(\ge2)}\) converges normally and defines a
      holomorphic function on all of \(\mathbb C_+\).
\item The first-prime series \(P_\chi^{(1)}\) is absolutely convergent on
      \(\Im z>1/2\), but its absolute majorant fails at and below that line.
      For the one-node test \textup{(V71.20)},
\[
 \sum_{p\nmid f}
 \left|
 \frac{\log p}{\sqrt p}C(\log p)
 \right|
 =\frac1{2Y}
 \sum_{p\nmid f}\frac{\log p}{p^{1/2+Y}}
 =\infty
 \quad(0<Y\le1/2).
 \tag{V71.22}
\]
\item On \(\Im z>1/2\),
\[
 P_\chi^{(1)}(z)
 =-\frac{L'}{L}(s(z),\chi)-P_\chi^{(\ge2)}(z).
 \tag{V71.23}
\]
Hence every pole of its meromorphic continuation in \(\mathbb C_+\) is a
nontrivial-zero pole; the gamma and higher-power layers create none.
\end{enumerate}
Thus the coefficient-side GRH obstruction is localized entirely in the
continued \(m=1\) prime layer.
\end{theorem}

\begin{proof}
Item \textup{(i)} follows because the digamma argument has positive real
part on \(\mathbb C_+\).  On a compact subset with
\(\Im z\ge\varepsilon>0\), the absolute majorant for
\textup{(V71.14)} is
\[
 \sum_p\sum_{m\ge2}(\log p)p^{-m(1/2+\varepsilon)}<\infty,
 \tag{V71.24}
\]
so the Weierstrass test proves \textup{(ii)}.  The corresponding \(m=1\)
majorant converges exactly when \(1/2+\Im z>1\).  Formula
\textup{(V71.22)} uses \textup{(V71.20)} and the divergence of
\(\sum_p(\log p)p^{-\sigma}\) for \(\sigma\le1\); omitting the finitely
many primes dividing \(f\) changes nothing.  Finally
\textup{(V71.23)} is \textup{(V71.12)} with the normally convergent tail
removed.  Since the other two terms are holomorphic on \(\mathbb C_+\),
only zeros of \(L\) can produce the stated poles.
\end{proof}

This result improves the V63 absolute-moment diagnosis.  It does not merely
say that an absolute bound loses the needed scale; it identifies the unique
Euler layer at which absolute control stops.  Any future coefficient
argument may treat the gamma and \(m\ge2\) contributions directly and reserve
all continuation work for \(P_\chi^{(1)}\).

\subsection{What pairing and the functional equation preserve}

The conjugate character has the same conductor and parity.  Applying
\textup{(V71.15)} to both channels at the same \((Z,{\bf c})\) gives an
exact phase ledger.

\begin{proposition}[Paired-character phase ledger]
\label{prop:v71-paired-ledger}
Let \(Q_\chi(Z,{\bf c})={\bf c}^*\mathsf K_Z(F_{\chi,0}){\bf c}\) in the
Euler region.  Then
\[
 \begin{aligned}
 Q_\chi+Q_{\overline\chi}
 &=2\mathfrak A_{f,\kappa}[C]\\
 &\quad-4\sum_p\sum_{m\ge1}
 \frac{\log p}{p^{m/2}}
 \Re(\chi(p)^m)\Re C(m\log p),
 \end{aligned}
 \tag{V71.25}
\]
whereas
\[
 Q_\chi-Q_{\overline\chi}
 =4\sum_p\sum_{m\ge1}
 \frac{\log p}{p^{m/2}}
 \Im(\chi(p)^m)\Im C(m\log p).
 \tag{V71.26}
\]
Thus pairing does not cancel the prime layer.  It separates its real and
imaginary character phases into the real and imaginary parts of the
directional autocorrelation.
\end{proposition}

\begin{proof}
Use
\[
 \Re(aC)+\Re(\overline aC)=2\Re a\Re C,
 \qquad
 \Re(aC)-\Re(\overline aC)=-2\Im a\Im C
 \tag{V71.27}
\]
in \textup{(V71.15)}.  The archimedean forms agree and therefore cancel in
the difference.
\end{proof}

There are also the exact meromorphic identities
\[
 F_{\overline\chi,0}(-\overline z)
 =\overline{F_{\chi,0}(z)},
 \qquad
 F_{\chi,0}(z)=-F_{\overline\chi,0}(-z).
 \tag{V71.28}
\]
The first is coefficient conjugation.  It maps an upper-half-plane node set
to its reflected upper-half-plane set and only duplicates the kernel test:
the two matrices are transposes and have the same spectrum.  The second is
the completed functional equation, but it maps \(\mathbb C_+\) to
\(\mathbb C_-\).  It therefore supplies no same-half-plane continuation of
the first-prime form across \(\Im z=1/2\).

\begin{theorem}[V71 coefficient-expansion verdict]
\label{thm:v71-verdict}
The V70 directional reset has the following exact outcome.
\begin{enumerate}[label=\textup{(\roman*)},leftmargin=3em]
\item Every finite Pick direction in the Euler region is the archimedean
      autocorrelation form \textup{(V71.10)} minus the character-weighted
      shift correlations \textup{(V71.15)}.
\item Gamma and higher prime powers are already holomorphic and absolutely
      controlled on the full target half-plane.
\item The first-prime layer alone loses absolute convergence at
      \(\Im z=1/2\), and its continued poles are exactly the nontrivial-zero
      obstruction.
\item Conjugate-channel pairing retains, rather than removes, both real and
      imaginary prime phases.
\item The functional equation reflects the continued first-prime layer into
      the opposite half-plane and supplies no positivity inequality on its
      own.
\end{enumerate}
Consequently the next admissible estimate must act on
\(P_\chi^{(1)}\) itself.  Recombining it immediately with
\(-L'/L\), taking absolute values, or invoking the paired functional equation
alone returns respectively to the GRH-equivalent pole statement, the V63
barrier, or the V66 equivalence.
\end{theorem}

\begin{firewall}[Prime-layer localization is not pole exclusion]
Theorem~\ref{thm:v71-prime-layer} identifies which normally convergent part
of the Euler expansion can and cannot cross the target half-plane.  It does
not continue the first-prime series holomorphically and does not control the
sign of \textup{(V71.16)}.  Equation \textup{(V71.23)} shows why: such a
continuation must already confront the zero poles of \(L'/L\).  GRH remains
open.
\end{firewall}

\begin{assessment}[V71 closure and V72 direction]
V71 corrects the finite-direction convention, derives the exact
gamma--prime shift formula, and localizes the entire coefficient-side
continuation obstruction to the \(m=1\) prime layer.  This removes the gamma
factor and all repeated prime powers from the difficult part of the next
argument.

V72 should test the first-prime layer through its finite character Fourier
structure rather than through \(-L'/L\).  The first calculation is to split
primes by reduced residue class modulo \(f\), diagonalize the character
phases, and determine whether a directional combination of the paired forms
can cancel the class-average prime density while retaining a coercive
nonprincipal remainder.  The test must be made at the correlation level
\(C(\log p)\), where V71 shows the phases actually survive; a bound on the
unweighted prime count would be irrelevant.
\end{assessment}

\section*{Version V71 append-only record}
\addcontentsline{toc}{section}{Version V71 append-only record}
The complete V70 source was retained before this continuation.  It had
1019338 bytes, 30165 lines, and SHA--256
\[
 \texttt{A6B488C772E2950FFFFA99DFCD521D7B7DC8A1E7484448300C1CE4BDE454FA0B}.
\]
V71 corrected the V70 vector convention, derived the exact directional
gamma--prime shift form, proved normal convergence of every higher
prime-power layer throughout \(\mathbb C_+\), isolated the first-prime layer
as the sole Euler continuation obstruction, and computed the phases retained
by conjugate-character pairing.

% =============================================================================
% END APPENDED V71 MATERIAL
% =============================================================================
% =============================================================================
% BEGIN APPENDED V72 MATERIAL
% =============================================================================

\section{V72: residue-class boundary landing and the Gevrey-two barrier}
\label{sec:grh-v72}

V71 isolated the first-prime layer
\[
 \mathcal P_\chi(s):=
 \sum_p\frac{\chi(p)\log p}{p^s},
 \qquad \Re s>1,
 \tag{V72.1}
\]
as the only part of the Euler expansion whose absolute convergence stops at
the line \(\Im z=1/2\), where \(s=1/2-iz\).  It proposed residue-class
Fourier decomposition at the directional correlation level.  That
calculation has one unconditional gain: cancellation between reduced residue
classes gives conditional convergence, with all tangential derivatives, on
\(\Re s=1\).  It does not open a strip.  The derivative bounds have Gevrey
order two, which is nonquasianalytic, and the power saving that would open the
whole target half-plane is itself GRH-complete.

Throughout this section \(\chi\) is primitive and nonprincipal, of conductor
\(f\), and
\[
 G_f:=(\mathbb Z/f\mathbb Z)^\times,
 \qquad d_f:=|G_f|=\varphi(f).
 \tag{V72.2}
\]

\subsection{Exact class cancellation at a Pick direction}

For \(a\in G_f\), a cutoff \(X\ge2\), and a V71 autocorrelation \(C\), put
\[
 \Pi_a[C;X]
 :=\sum_{\substack{p\le X\\p\equiv a\;({\rm mod}\ f)}}
 \frac{\log p}{\sqrt p}C(\log p),
 \qquad
 \overline\Pi[C;X]:=\frac1{d_f}\sum_{a\in G_f}\Pi_a[C;X].
 \tag{V72.3}
\]

\begin{proposition}[Directional residue-error projection]
\label{prop:v72-directional-residue}
For every finite cutoff,
\[
 \boxed{
 \sum_{p\le X}\frac{\chi(p)\log p}{\sqrt p}C(\log p)
 =\sum_{a\in G_f}\chi(a)
  \bigl(\Pi_a[C;X]-\overline\Pi[C;X]\bigr).}
 \tag{V72.4}
\]
More generally, the empirical average may be replaced by any scalar common
to all residue classes.
\end{proposition}

\begin{proof}
Group the unramified primes by their reduced residue class and use
\(\sum_{a\in G_f}\chi(a)=0\).  Primes dividing \(f\) have
\(\chi(p)=0\) and contribute nothing.
\end{proof}

This is the exact V29 projection with the Gaussian reader replaced by the
V71 correlation weight.  In particular the common prime-density profile
does cancel.  Finite Fourier geometry supplies no coercive remainder after
that cancellation.

\begin{proposition}[No target-mode coercivity from the class transform]
\label{prop:v72-no-class-coercivity}
Let \(E=(E_a)_{a\in G_f}\in\mathbb C^{d_f}\).  Character orthogonality gives
\[
 \sum_{\eta\in\widehat G_f}
 \left|\sum_{a\in G_f}\eta(a)E_a\right|^2
 =d_f\sum_{a\in G_f}|E_a|^2.
 \tag{V72.5}
\]
It gives no upper bound for the prescribed \(\chi\)-coordinate in terms of
the principal coordinate.  Indeed, for any \(w\in\mathbb C\),
\[
 E_a:=\frac{\overline{\chi(a)}}{d_f}w
 \tag{V72.6}
\]
has \(\chi\)-coordinate \(w\) and principal coordinate zero.
Consequently residue-class grouping, orthogonality, and Parseval alone cannot
bound the right side of \textup{(V72.4)}.
\end{proposition}

\begin{proof}
Equation \textup{(V72.5)} is finite character Parseval.  Substitution of
\textup{(V72.6)} uses
\(\sum_a\overline{\chi(a)}=0\) and
\(|\chi(a)|=1\) on \(G_f\).
\end{proof}

The vector in \textup{(V72.6)} is an information-capacity witness, not a
claim about the actual prime vector.  It proves that any useful bound must
use correlations as \(X\) varies, not merely the fixed conductor transform.
This agrees with the V29 and V43 no-go results and prevents the V71 proposal
from duplicating those arguments.

\subsection{Prime PNT error and conditional boundary convergence}

Define the prime-only Chebyshev sums
\[
 \vartheta_a(X):=
 \sum_{\substack{p\le X\\p\equiv a\;({\rm mod}\ f)}}\log p,
 \qquad
 \vartheta_\chi(X):=sum_{p\le X}\chi(p)\log p.
 \tag{V72.7}
\]
The fixed-modulus prime number theorem in progressions gives constants
\(A_\chi,c_\chi>0\) such that
\[
 \boxed{
 |\vartheta_\chi(X)|
 \le A_\chi X e^{-c_\chi\sqrt{\log X}}
 \qquad(X\ge2).}
 \tag{V72.8}
\]
This follows equally by projecting the progression errors: the common term
\(X/d_f\) vanishes in the nonprincipal character direction.

\begin{theorem}[Conditional landing on the Euler boundary]
\label{thm:v72-boundary-landing}
The first-prime series \textup{(V72.1)} converges conditionally at every
point \(s=1+i\tau\), uniformly for \(\tau\) in compact intervals.  Hence it
has a continuous boundary value on \(\Re s=1\).

Let \(Z=(z_1,\ldots,z_N)\) satisfy \(\Im z_j\ge1/2\), and let
\(C=C_{Z,{\bf c}}\) be the corrected V71 autocorrelation.  Then
\[
 \lim_{X\to\infty}
 \sum_{p\le X}\frac{\chi(p)\log p}{\sqrt p}C(\log p)
 \tag{V72.9}
\]
exists.  Thus the exact Euler shift form \textup{(V71.15)} has a
coefficient-sensitive, conditionally convergent limit on the closed Euler
region \(\min_j\Im z_j\ge1/2\).
\end{theorem}

\begin{proof}
Stieltjes summation gives, initially for \(\Re s>1\),
\[
 \mathcal P_\chi(s)
 =\int_{2^-}^{\infty}x^{-s}\,d\vartheta_\chi(x).
 \tag{V72.10}
\]
At \(s=1+i\tau\), integration by parts produces an integrand bounded on a
compact \(\tau\)-interval by a constant times
\[
 \frac{|\vartheta_\chi(x)|}{x^2}
 \ll_\chi\frac{e^{-c_\chi\sqrt{\log x}}}{x},
 \tag{V72.11}
\]
which is integrable after \(v=\log x\).  The boundary term tends to zero,
and the estimates are locally uniform in \(\tau\).

From \textup{(V71.3)} and \textup{(V71.6)},
\[
 C(u)=\sum_{j,k=1}^N
 \overline{c_j}c_k
 \frac{e^{iz_ju}}{-i(z_j-\overline{z_k})}.
 \tag{V72.12}
\]
Consequently the truncated expression in \textup{(V72.9)} is a finite
linear combination of
\(\mathcal P_{\chi,X}(s(z_j))\).  Every
\(\Re s(z_j)\ge1\), and the preceding boundary convergence proves the
claim.  The V71 gamma and higher-power terms already converge normally
there.
\end{proof}

This is a strict improvement over the absolute majorant in
\textup{(V71.22)}: arithmetic cancellation reaches the boundary where that
majorant diverges.  But \textup{(V72.8)} has effective exponent one in
\(X\), so the same Stieltjes argument supplies no convergence at any fixed
\(\Re s<1\).

\subsection{The complete boundary jet is only Gevrey two}

The error in \textup{(V72.8)} is strong enough to differentiate the boundary
value arbitrarily many times.

\begin{theorem}[Gevrey-two boundary jet]
\label{thm:v72-gevrey-two}
For every \(T>0\), the boundary function
\(\tau\mapsto\mathcal P_\chi(1+i\tau)\) belongs to
\(C^\infty([-T,T])\).  There are constants \(K_{\chi,T},B_\chi>0\) such
that, for every integer \(r\ge0\),
\[
 \boxed{
 \sup_{|\tau|\le T}
 \left|
 \frac{d^r}{d\tau^r}\mathcal P_\chi(1+i\tau)
 \right|
 \le K_{\chi,T}B_\chi^{,r+1}(2r+2)!.}
 \tag{V72.13}
\]
Thus the unconditional boundary control has Gevrey order two.  Its
Denjoy--Carleman sum converges:
\[
 \sum_{r\ge1}
 \left(B_\chi^{,r+1}(2r+2)!\right)^{-1/r}<\infty.
 \tag{V72.14}
\]
Accordingly this derivative bound is nonquasianalytic and cannot, by itself,
force a holomorphic continuation across \(\Re s=1\).
\end{theorem}

\begin{proof}
Formally the \(r\)-th tangential derivative inserts
\((-i\log p)^r\).  To justify it, apply Stieltjes summation to
\(x^{-1-i\tau}(\log x)^r\).  Apart from a fixed compact initial interval,
the absolute integral is bounded by a constant times
\[
 \int_0^\infty e^{-c_\chi\sqrt v}
 \bigl((1+T)v^r+rv^{r-1}\bigr)\,dv.
 \tag{V72.15}
\]
With \(w=\sqrt v\),
\[
 \int_0^\infty e^{-c_\chi\sqrt v}v^r,dv
 =\frac{2\Gamma(2r+2)}{c_\chi^{,2r+2}}.
 \tag{V72.16}
\]
The same bound for the adjacent moment proves uniform convergence of every
differentiated integral and yields \textup{(V72.13)} after enlarging the
constants.  Stirling's formula gives
\(((2r+2)!)^{1/r}\asymp r^2\), so the series in
\textup{(V72.14)} is comparable to \(\sum r^{-2}\).  The
Denjoy--Carleman criterion therefore does not give quasianalyticity; indeed
the corresponding Gevrey-two class contains nonzero flat functions.
\end{proof}

More generally, an error
\[
 \vartheta_\chi(X)ll X
 e^{-c(\log X)^\alpha},
 \qquad 0<\alpha<1,
 \tag{V72.17}
\]
produces moments of size
\[
 \int_0^\infty e^{-cv^\alpha}v^r,dv
 =\frac1\alpha c^{-(r+1)/\alpha}
   \Gamma\!\left(\frac{r+1}{\alpha}\right),
 \tag{V72.18}
\]
and hence Gevrey order \(1/\alpha>1\), still nonquasianalytic.  Analytic
order requires \(\alpha=1\), which is a power saving in the prime counting
error and therefore opens an actual left neighborhood by Stieltjes
summation.  This explains quantitatively why an arbitrarily smooth Euler
boundary is not an analytic-continuation theorem.

\subsection{The continuation-depth dictionary is GRH-complete}

The exact amount of prime cancellation needed to enter the strip can be
stated without Pick terminology.

\begin{proposition}[Prime-counting exponent versus continuation depth]
\label{prop:v72-depth-dictionary}
If, for some \(\beta<1\) and \(A\ge0\),
\[
 \vartheta_\chi(X)=O\bigl(X^\beta(\log(2X))^A\bigr),
 \tag{V72.19}
\]
then \(\mathcal P_\chi(s)\) converges locally uniformly and is holomorphic in
\(\Re s>\beta\).  In the \(z\)-coordinate this covers
\[
 \Im z>\beta-\frac12.
 \tag{V72.20}
\]
To cover every \(\Im z>0\) by this mechanism one needs
\(\beta=1/2+\varepsilon\) for every \(\varepsilon>0\).
\end{proposition}

\begin{proof}
Integrate \textup{(V72.10)} by parts.  The tail is majorized by
\(\int x^{\beta-\sigma-1}(\log x)^A dx\), which converges locally uniformly
for \(\sigma>\beta\).  Since \(\Re s(z)=1/2+\Im z\), this is
\textup{(V72.20)}.
\end{proof}

\begin{theorem}[The prime-only square-root estimate is paired-GRH complete]
\label{thm:v72-prime-grh-complete}
For the conjugate pair \(\chi,\overline\chi\), the estimates
\[
 \vartheta_\chi(X),\ \vartheta_{\overline\chi}(X)
 =O_{\chi,\varepsilon}(X^{1/2+\varepsilon})
 \quad\hbox{for every }\varepsilon>0
 \tag{V72.21}
\]
are equivalent to paired GRH.
\end{theorem}

\begin{proof}
Let
\(\Psi_\chi(X)=\sum_{n\le X}\Lambda(n)\chi(n)\).  The higher prime powers
satisfy the elementary bound
\[
 \Psi_\chi(X)-\vartheta_\chi(X)
 =O(\sqrt X\log^2(2X)).
 \tag{V72.22}
\]
Therefore \textup{(V72.21)} is equivalent to the paired square-root
Chebyshev estimates of V29.  Theorem
\ref{thm:v29-projected-error-grh} proves that those estimates are equivalent
to paired GRH.
\end{proof}

\begin{theorem}[V72 residue-class verdict]
\label{thm:v72-verdict}
The V71 residue-class proposal is now resolved.
\begin{enumerate}[label=\textup{(\roman*)},leftmargin=3em]
\item The common residue-class density cancels exactly at every directional
      correlation \(C(\log p)\).
\item Fixed-conductor Fourier geometry does not control the remaining target
      coordinate, exactly as in V29 and V43.
\item Classical progression cancellation nevertheless extends the
      first-prime form conditionally to the Euler boundary.
\item Its complete boundary jet has only Gevrey-two growth and is
      nonquasianalytic, so boundary smoothness does not cross the line.
\item A prime-counting power saving deep enough to cover all
      \(\mathbb C_+\) is equivalent to the desired paired GRH conclusion.
\end{enumerate}
Thus residue classes produce a rigorous boundary landing but neither a
coercive nonprincipal remainder nor a continuation principle.
\end{theorem}

\begin{firewall}[Smooth landing is not strip continuation]
Conditional convergence and all tangential derivatives on \(\Re s=1\) do
not imply holomorphic extension to \(\Re s<1\).  The factorial scale in
\textup{(V72.13)} lies in a nonquasianalytic class, and
\textup{(V72.21)} shows that the missing power saving is already
GRH-complete.  No off-line zero has been excluded.  GRH remains open.
\end{firewall}

\begin{assessment}[V72 closure and V73 direction]
V72 transports the V29 residue projection to the exact V71 prime
autocorrelation, obtains the unconditional boundary value, and identifies
the Gevrey-two obstruction to crossing the Euler line.  This prevents both a
duplicate finite-Fourier argument and an invalid appeal to infinite boundary
smoothness.

V73 should test whether the Euler product couples the boundary jets through
the prime-zeta identity
\[
 \sum_p\frac{\chi(p)}{p^s}
 =\sum_{m\ge1}\frac{\mu(m)}m
   \log L(ms,\chi^m),
 \tag{V72.23}
\]
initially for \(\Re s>1\).  This is exponent-space M\"obius inversion, not
the integer-factor convolution studied in V44.  The required decision is
whether finite character order makes the \(m\ge2\) power-character terms
feed back a quasianalytic constraint on the \(m=1\) jet, or whether they are
holomorphic spectators and the identity collapses exactly to the original
\(L'/L\) pole obstruction.
\end{assessment}

\section*{Version V72 append-only record}
\addcontentsline{toc}{section}{Version V72 append-only record}
The complete V71 source was retained before this continuation.  It had
1033635 bytes, 30596 lines, and SHA--256
\[
 \texttt{276120C88CD4363E14B4FFB844F05620242846AC946018C6C6005E9A6CC5AD6B}.
\]
V72 projected the first-prime directional form into residue classes, proved
conditional convergence and a complete boundary jet at \(\Re s=1\),
computed its Gevrey-two nonquasianalytic growth, and proved that the
prime-only square-root estimate needed to traverse the target half-plane is
paired-GRH complete.

% =============================================================================
% END APPENDED V72 MATERIAL
% =============================================================================
% =============================================================================
% BEGIN APPENDED V73 MATERIAL
% =============================================================================

\section{V73: exponent-space M\"obius inversion and the boundary-moment radius criterion}
\label{sec:grh-v73}

V72 reached the Euler boundary by residue-class cancellation and obtained a
Gevrey-two estimate for the boundary prime moments.  It then proposed testing
the exponent-space M\"obius inversion across the character powers.  The
inversion can be completed exactly.  It is triangular, not recursive: every
power-character term with exponent \(m\ge2\) is holomorphic on
\(\Re s>1/2\), while the \(m=1\) term retains every zero pole of
\(L(s,\chi)\).  Finite order of \(\chi\) does not feed information back to
the undilated argument.

This calculation also sharpens the boundary target.  The condition needed
for paired GRH is exactly an analytic-order bound on the conditionally
convergent prime moments at \(s=1+it\): their Taylor-root rate must be at most
two at every ordinate.  An off-line zero of real part \(\beta>1/2\) forces
rate at least \((1-\beta)^{-1}>2\) at its ordinate.

As in V71--V72, \(\chi\) is primitive and nonprincipal.  For a positive
integer \(m\), \(\chi^m\) denotes the possibly imprimitive Dirichlet
character modulo the same conductor support; its Euler factors are the ones
obtained by raising \(\chi(p)\) to the \(m\)-th power.

\subsection{Prime-zeta inversion with character powers}

Put
\[
 \mathfrak P_\chi(s):=\sum_p\frac{\chi(p)}{p^s},
 \qquad
 \mathcal P_\chi(s):=-\mathfrak P_\chi'(s)
 =\sum_p\frac{\chi(p)\log p}{p^s},
 \tag{V73.1}
\]
initially for \(\Re s>1\).

\begin{theorem}[Exponent-space M\"obius inversion]
\label{thm:v73-mobius-inversion}
For \(\Re s>1\), with Euler-product logarithms,
\[
 \boxed{
 \mathfrak P_\chi(s)
 =\sum_{m\ge1}\frac{\mu(m)}m
   \log L(ms,\chi^m).}
 \tag{V73.2}
\]
Differentiation gives
\[
 \boxed{
 \mathcal P_\chi(s)
 =-\sum_{m\ge1}\mu(m)
   \frac{L'}{L}(ms,\chi^m).}
 \tag{V73.3}
\]
Both identities are locally normally convergent in their initial domain.
\end{theorem}

\begin{proof}
Absolute convergence permits the expansion
\[
 \log L(ms,\chi^m)
 =\sum_p\sum_{k\ge1}
 \frac{\chi(p)^{mk}}{k p^{mks}}.
 \tag{V73.4}
\]
For a fixed exponent \(n=mk\), the coefficient on the right of
\textup{(V73.2)} is
\[
 \sum_{m\mid n}\frac{\mu(m)}{m(n/m)}
 =\frac1n\sum_{m\mid n}\mu(m)
 =\mathbf1_{n=1}.
 \tag{V73.5}
\]
Only the first-prime term survives.  Termwise differentiation is justified
on compact subsets of \(\Re s>1\) and proves \textup{(V73.3)}.
\end{proof}

This is not the V44 identity \(\Lambda=\mu*\log\).  V44 factorized the
integer variable inside a heat reader.  Here M\"obius inversion acts on the
Euler exponent and replaces repeated powers of one prime by dilated
logarithms of the power characters.

\subsection{The power-character correction is holomorphic}

Define
\[
 \mathfrak G_\chi(s)
 :=\sum_{m\ge2}\frac{\mu(m)}m
   \log L(ms,\chi^m)
 \tag{V73.6}
\]
using the zero-free Euler logarithm, and
\[
 \mathcal H_\chi(s):=-\mathfrak G_\chi'(s)
 =-\sum_{m\ge2}\mu(m)
   \frac{L'}{L}(ms,\chi^m).
 \tag{V73.7}
\]

\begin{theorem}[Triangular holomorphy of the power correction]
\label{thm:v73-triangular}
The series \textup{(V73.6)} and \textup{(V73.7)} converge locally normally
and define holomorphic functions on
\[
 \boxed{\Re s>\frac12.}
 \tag{V73.8}
\]
On that half-plane the meromorphic continuation of the first-prime layer is
\[
 \boxed{
 \mathcal P_\chi(s)
 =-\frac{L'}{L}(s,\chi)+\mathcal H_\chi(s).}
 \tag{V73.9}
\]
Its poles in \(\Re s>1/2\) are exactly the zeros of \(L(s,\chi)\), with
residue minus their multiplicity.
\end{theorem}

\begin{proof}
On a compact subset of \(\Re s>1/2\), choose \(\delta>0\) with
\(\Re s\ge1/2+\delta\).  For every \(m\ge2\),
\(\Re(ms)\ge1+2\delta\), and hence the Euler logarithm and logarithmic
derivative of \(L(ms,\chi^m)\) converge absolutely.  They are majorized by
the corresponding zeta series.  The large-\(m\) terms decay geometrically
like a polynomial in \(m\) times \(2^{-m(1/2+\delta)}\), uniformly with all
derivatives on a smaller compact set.  The Weierstrass theorem proves local
normal convergence and holomorphy.

Equation \textup{(V73.9)} is \textup{(V73.3)} with its \(m=1\) term
separated.  A primitive nonprincipal \(L(s,\chi)\) is entire, and its trivial
zeros lie outside \(\Re s>1/2\).  Since \(\mathcal H_\chi\) is holomorphic
there, a nontrivial zero \(\rho\) of multiplicity \(m_\rho\) contributes
\[
 -\frac{m_\rho}{s-\rho}
 \tag{V73.10}
\]
to \(\mathcal P_\chi\), with no possible cancellation from the power
correction.
\end{proof}

If \(\chi\) has finite order \(h\), the sequence \(\chi^m\) is periodic,
but the arguments \(ms\) are not.  Even when \(\chi^m\) is principal, one
has \(\Re(ms)>1\) in \textup{(V73.8)}; its pole at \(ms=1\) lies on or to
the left of the boundary.  Periodicity in character space therefore creates
no equation returning a dilated term to \(L(s,\chi)\) at the same point.

There is an equivalent monodromy statement.  On any simply connected
zero-free piece of \(\Re s>1/2\),
\[
 \mathfrak P_\chi(s)=\log L(s,\chi)+\mathfrak G_\chi(s).
 \tag{V73.11}
\]
A zero \(\rho\) of multiplicity \(m_\rho\) gives logarithmic monodromy
\(2\pi i m_\rho\) in \(\mathfrak P_\chi\).  The single-valued holomorphic
power correction cannot remove it.

\subsection{What the V72 Gevrey bound does and does not say}

\begin{remark}[The V72 estimate is epistemic, not an exact regularity order]
\label{rem:v73-v72-clarification}
The Gevrey-two bound \textup{(V72.13)} is the derivative growth obtained
from the classical progression error \textup{(V72.8)}.  It does not assert
that the actual boundary function fails to be analytic.  Equation
\textup{(V73.9)}, together with the classical zero-free line
\(\Re s=1\), shows that \(\mathcal P_\chi\) is holomorphic in some
neighborhood of each point \(1+it\).  What is unknown is whether the radius
of every such neighborhood reaches the full distance \(1/2\) to the
critical line.  The nonquasianalytic conclusion of V72 remains correct: the
PNT-derived bound alone cannot prove that radius.
\end{remark}

For \(t\in\mathbb R\), let \(R_\chi(t)\) be the radius of convergence of the
Taylor series of \(\mathcal P_\chi\) at
\[
 s_t:=1+it.
 \tag{V73.12}
\]
The line \(\Re s=1\) is zero-free, so \(R_\chi(t)>0\).

\begin{lemma}[The power correction already has the target radius]
\label{lem:v73-correction-radius}
For every \(t\), the Taylor series of \(\mathcal H_\chi\) at \(s_t\) has
radius at least \(1/2\).  Equivalently,
\[
 \limsup_{r\to\infty}
 \left(
 \frac{|\mathcal H_\chi^{(r)}(s_t)|}{r!}
 \right)^{1/r}\le2.
 \tag{V73.13}
\]
\end{lemma}

\begin{proof}
Every open disk \(|s-s_t|<1/2\) lies in \(\Re s>1/2\), where
Theorem~\ref{thm:v73-triangular} gives holomorphy.  The radius assertion and
\textup{(V73.13)} follow from the Taylor theorem and Cauchy--Hadamard.
\end{proof}

Thus the dilated power characters already satisfy exactly the analytic scale
needed for the GRH half-plane.  Only the undilated logarithmic derivative can
violate it.

\subsection{A coefficient-level Taylor-radius criterion for paired GRH}

V72 justifies the conditionally convergent boundary moments
\[
 \mathcal M_{\chi,r}(t)
 :=\sum_p
 \frac{\chi(p)(\log p)^{r+1}}{p^{1+it}},
 \qquad r=0,1,2,\ldots,
 \tag{V73.14}
\]
and termwise Stieltjes differentiation gives
\[
 \mathcal P_\chi^{(r)}(1+it)
 =(-1)^r\mathcal M_{\chi,r}(t).
 \tag{V73.15}
\]

\begin{theorem}[Boundary prime-moment criterion]
\label{thm:v73-moment-criterion}
The following are equivalent for the channel \(\chi\).
\begin{enumerate}[label=\textup{(\roman*)},leftmargin=3em]
\item \(L(s,\chi)\) has no zero with \(\Re s>1/2\).
\item \(R_\chi(t)\ge1/2\) for every \(t\in\mathbb R\).
\item For every \(t\in\mathbb R\),
\[
 \boxed{
 \limsup_{r\to\infty}
 \left(
 \frac{|\mathcal M_{\chi,r}(t)|}{r!}
 \right)^{1/r}\le2.}
 \tag{V73.16}
\]
\end{enumerate}
Consequently, imposing \textup{(V73.16)} for both
\(\chi\) and \(\overline\chi\) is equivalent to paired GRH.
\end{theorem}

\begin{proof}
If \(L(s,\chi)\) has no zero in \(\Re s>1/2\), then
\textup{(V73.9)} is holomorphic there.  The disk
\(|s-s_t|<1/2\) lies in that half-plane, proving \textup{(ii)}.

Conversely, suppose \(\rho=\beta+i\gamma\) is a zero with
\(\beta>1/2\).  At \(s_\gamma=1+i\gamma\), the pole
\textup{(V73.10)} is at distance
\[
 |s_\gamma-\rho|=1-\beta<\frac12.
 \tag{V73.17}
\]
It cannot be cancelled by \(\mathcal H_\chi\), so
\(R_\chi(\gamma)\le1-\beta<1/2\), contradicting \textup{(ii)}.
The equivalence of \textup{(ii)} and \textup{(iii)} is
Cauchy--Hadamard applied to \textup{(V73.15)}.  Finally, applying the channel
statement to both conjugate characters and using the functional equation is
the paired-GRH equivalence of V66.
\end{proof}

The criterion is directional and coefficient-sensitive.  It uses only the
prime values \(\chi(p)\) in \textup{(V73.14)}, not the zero set in its
statement.  It is nevertheless an equivalence, not yet an estimate.

\begin{corollary}[Exact high-moment signature of an off-line zero]
\label{cor:v73-zero-signature}
If \(L(s,\chi)\) has a zero \(\rho=\beta+i\gamma\) with
\(\beta>1/2\), then
\[
 \boxed{
 \limsup_{r\to\infty}
 \left(
 \frac{|\mathcal M_{\chi,r}(\gamma)|}{r!}
 \right)^{1/r}
 \ge\frac1{1-\beta}>2.}
 \tag{V73.18}
\]
For a zero of multiplicity \(m_\rho\), its local contribution is
\[
 -m_\rho\frac{(-1)^r r!}
 {(1-\beta)^{r+1}}
 \tag{V73.19}
\]
to \(\mathcal P_\chi^{(r)}(1+i\gamma)\), before addition of the other
singularities and the holomorphic correction.
\end{corollary}

\begin{proof}
The pole at distance \(1-\beta\) gives
\(R_\chi(\gamma)\le1-\beta\).  Cauchy--Hadamard and
\textup{(V73.15)} prove \textup{(V73.18)}.  Differentiating
\textup{(V73.10)} gives \textup{(V73.19)}.
\end{proof}

This avoids the finite-sample masking problem of V69 in one precise sense:
an off-line zero forces an all-order boundary-moment growth rate.  It does
not avoid the arithmetic difficulty, because proving the upper rate two is
another exact form of the required cancellation.

\begin{theorem}[V73 power-character verdict]
\label{thm:v73-verdict}
Exponent-space M\"obius inversion has the following complete outcome.
\begin{enumerate}[label=\textup{(\roman*)},leftmargin=3em]
\item It removes all repeated prime powers by an exact, normally convergent
      family of dilated power-character logarithms.
\item Every \(m\ge2\) correction is holomorphic on \(\Re s>1/2\) and already
      has boundary Taylor-root rate at most two.
\item Finite character order does not create feedback, because character
      repetition is accompanied by the dilation \(s\mapsto ms\).
\item Every off-line zero and every logarithmic monodromy remains entirely
      in the undilated \(m=1\) term.
\item Paired GRH is equivalent to the prime-moment rate
      \textup{(V73.16)} in both channels.
\end{enumerate}
Thus the power-character family sharpens the coefficient target but supplies
no bootstrap estimate for it.
\end{theorem}

\begin{firewall}[A Taylor-radius equivalence is not a moment bound]
Equation \textup{(V73.16)} is a new prime-only boundary formulation of the
zero-free half-plane.  No argument here proves its upper bound.  The
Gevrey-two estimate from V72 is far too large in the analytic root scale,
while the holomorphic power correction cannot control the singular
undilated term.  GRH remains open.
\end{firewall}

\begin{assessment}[V73 closure and V74 direction]
V73 proves that exponent-space M\"obius inversion is triangular and closes
the proposed finite-order feedback mechanism.  Its durable output is the
exact prime-moment radius threshold two and the forced rate
\((1-\beta)^{-1}\) of an off-line zero.

V74 should derive the zero-side power-sum formula for the normalized moments
\(\mathcal M_{\chi,r}(t)/r!\) and apply a quantitative Tur\'an power-sum
lemma on blocks of consecutive \(r\).  The purpose is not to restate
Cauchy--Hadamard, but to test whether a nonnegative trigonometric polynomial
in the prime phases, possibly involving the power characters, can bound
those blocks at root rate two.  If the unavoidable principal-character pole
or absolute prime mass only recovers the classical zero-free region, that
positivity route should be closed explicitly.
\end{assessment}

\section*{Version V73 append-only record}
\addcontentsline{toc}{section}{Version V73 append-only record}
The complete V72 source was retained before this continuation.  It had
1046971 bytes, 30973 lines, and SHA--256
\[
 \texttt{B4988D65BFFDA9BC505AB526AF7DC2FA98A59A7463198DB8113106E83012B734}.
\]
V73 completed exponent-space M\"obius inversion, proved the holomorphic
triangularity of every power-character correction, clarified the scope of
the V72 Gevrey bound, and established the prime-only boundary-moment Taylor
rate two as an exact paired-GRH criterion.

% =============================================================================
% END APPENDED V73 MATERIAL
% =============================================================================
% =============================================================================
% BEGIN APPENDED V74 MATERIAL
% =============================================================================

\section{V74: zero power sums, block visibility, and the principal-pole positivity barrier}
\label{sec:grh-v74}

V73 converted paired GRH into the boundary prime-moment root bound two and
proposed a Tur\'an power-sum test.  The zero-side power sum can be written
exactly.  It yields more than the Cauchy--Hadamard limsup: if an off-line zero
lies closer than the critical-line distance \(1/2\), then every sufficiently
high block of a fixed finite number of consecutive moment orders contains an
exponentially large witness.

The natural prime-side positivity test nevertheless fails for a precise
reason.  Every nonzero nonnegative trigonometric polynomial has a positive
constant Fourier coefficient.  In the Euler logarithmic derivative that
coefficient is the principal-character channel, whose pole at \(s=1\) is
closer to any evaluation point \(\sigma>1\) than a target zero
\(\beta+i\gamma\) with \(\beta<1\).  Its high-order power sum therefore
dominates the target zero exponentially.  This closes the fixed-polynomial
positivity test while leaving a sharply defined pole-compensation question.

\subsection{The exact zero-side moment expansion}

Retain the first-prime layer \(\mathcal P_\chi\), its boundary moments
\(\mathcal M_{\chi,r}(t)\), and the holomorphic power correction
\(\mathcal H_\chi\) from V73.  Put
\[
 A_\chi(s):=\frac12\log\frac f\pi
 +\frac12\psi\!\left(\frac{s+\kappa}{2}\right),
 \qquad
 J_\chi(s):=A_\chi(s)+\mathcal H_\chi(s).
 \tag{V74.1}
\]
Since
\(\Lambda'/\Lambda=A_\chi+L'/L\), equation
\textup{(V73.9)} becomes
\[
 \mathcal P_\chi(s)
 =-\frac{\Lambda'}{\Lambda}(s,\chi)+J_\chi(s).
 \tag{V74.2}
\]

\begin{lemma}[Normalized zero power sum]
\label{lem:v74-zero-power-sum}
Let \(s_t=1+it\).  For every integer \(r\ge1\),
\[
 \boxed{
 \frac{\mathcal M_{\chi,r}(t)}{r!}
 =-\sum_{\rho}
   \frac{m_\rho}{(s_t-\rho)^{r+1}}
 +\mathcal E_{\chi,r}(t),}
 \tag{V74.3}
\]
where the sum is over the nontrivial zeros of the completed
\(L\)-function, with multiplicity, and
\[
 \mathcal E_{\chi,r}(t)
 :=\frac{(-1)^r}{r!}J_\chi^{(r)}(s_t).
 \tag{V74.4}
\]
The zero sum is absolutely convergent, and
\[
 \boxed{
 \limsup_{r\to\infty}
 |\mathcal E_{\chi,r}(t)|^{1/r}\le2.}
 \tag{V74.5}
\]
\end{lemma}

\begin{proof}
The order-one Hadamard product gives
\[
 \frac{\Lambda'}{\Lambda}(s,\chi)
 =B_\chi+\sum_\rho m_\rho
 \left(\frac1{s-\rho}+\frac1\rho\right).
 \tag{V74.6}
\]
After one or more derivatives the zero sum is absolute and
\[
 \left(\frac{\Lambda'}{\Lambda}\right)^{(r)}(s,\chi)
 =(-1)^r r!\sum_\rho
  \frac{m_\rho}{(s-\rho)^{r+1}}.
 \tag{V74.7}
\]
Apply \((-1)^r/r!\) to \textup{(V74.2)} and use
\textup{(V73.15)}.  The gamma factor is holomorphic on a neighborhood of
the closed disk \(|s-s_t|\le R\) for every \(R<1/2\), and
Theorem~\ref{thm:v73-triangular} gives the same statement for
\(\mathcal H_\chi\).  Cauchy's estimate followed by \(R\uparrow1/2\)
proves \textup{(V74.5)}.
\end{proof}

Thus every root rate strictly larger than two comes from the zero power sum,
not from gamma or repeated-prime corrections.

\subsection{A Vandermonde block form of Tur\'an visibility}

Suppose that, for a fixed \(t\), the nearest zero distance
\[
 \delta_t:=\min_\rho|s_t-\rho|
 \tag{V74.8}
\]
satisfies \(\delta_t<1/2\).  The minimum exists because the zero set is
discrete and has the standard order-one counting bound.  Let
\[
 \Omega_t:=\{\rho:|s_t-\rho|=\delta_t\}
 \tag{V74.9}
\]
be the finite set of distinct closest zeros, and put \(q_t:=|\Omega_t|\).

\begin{theorem}[Every high block sees a closest off-line zero]
\label{thm:v74-block-visibility}
There are constants \(c_t>0\) and \(n_t\) such that, for every integer
\(n\ge n_t\),
\[
 \boxed{
 \max_{0\le\ell<q_t}
 \left\{
 \delta_t^{\,n+\ell}
 \left|
 \frac{\mathcal M_{\chi,n+\ell-1}(t)}
      {(n+\ell-1)!}
 \right|
 \right\}
 \ge c_t.}
 \tag{V74.10}
\]
Hence an off-line zero closer than \(1/2\) cannot be hidden for
\(q_t\) consecutive sufficiently high orders.  In particular, at least one
order in every such block has root scale asymptotic to at least
\(\delta_t^{-1}>2\).
\end{theorem}

\begin{proof}
For \(\rho\in\Omega_t\), set
\[
 u_\rho:=\frac{\delta_t}{s_t-\rho},
 \qquad |u_\rho|=1,
 \tag{V74.11}
\]
and define
\(T_k:=\sum_{\rho\in\Omega_t}m_\rho u_\rho^k\).
The \(q_t\times q_t\) Vandermonde matrix
\[
 V_t:=\bigl[u_\rho^\ell\bigr]_{
 0\le\ell<q_t,\ \rho\in\Omega_t}
 \tag{V74.12}
\]
is invertible.  For every \(n\),
\[
 (T_n,\ldots,T_{n+q_t-1})^T
 =V_t\,{\rm diag}(u_\rho^n)_\rho\,(m_\rho)_\rho.
 \tag{V74.13}
\]
The diagonal matrix is unitary, so the Euclidean norm of this block is
bounded below by
\(\sigma_{\min}(V_t)\|(m_\rho)_\rho\|_2>0\), uniformly in \(n\).

There is a next zero distance \(d_t>\delta_t\).  Absolute convergence of
\(\sum_\rho|s_t-\rho|^{-2}\) gives, uniformly for bounded \(\ell\),
\[
 \delta_t^{\,n+\ell}
 \sum_{\rho\notin\Omega_t}
 \frac{m_\rho}{(s_t-\rho)^{n+\ell}}
 \longrightarrow0.
 \tag{V74.14}
\]
The normalized holomorphic remainder tends to zero on the same scale by
\textup{(V74.5)}, because \(2\delta_t<1\).  Insert these facts into
\textup{(V74.3)} with \(r+1=n+\ell\).  At least one component of each
Vandermonde block remains bounded below, proving \textup{(V74.10)}.
\end{proof}

For a zero \(\rho=\beta+i\gamma\) and \(t=\gamma\), one has
\(\delta_t\le1-\beta<1/2\).  Thus Theorem
\ref{thm:v74-block-visibility} strengthens V73's limsup signature at exactly
the ordinate of every forbidden zero.

\subsection{The Euler-positive trigonometric family}

Let
\[
 P(\theta)=a_0+2\Re\sum_{k=1}^K a_ke^{ik\theta}
 \tag{V74.15}
\]
be a real trigonometric polynomial with \(P(\theta)\ge0\).  Its Fourier
coefficients satisfy
\[
 a_0=\frac1{2\pi}\int_0^{2\pi}P(\theta)d\theta,
 \qquad |a_k|\le a_0.
 \tag{V74.16}
\]
If \(P\not\equiv0\), then \(a_0>0\).

For a character \(\eta\), \(r\ge0\), and \(\Re s>1\), define the normalized
logarithmic-derivative moment
\[
 \mathcal Q_{\eta,r}(s)
 :=\frac{(-1)^r}{r!}
 \frac{d^r}{ds^r}
 \left(-\frac{L'}{L}(s,\eta)\right).
 \tag{V74.17}
\]
For the principal mode \(\mathbf1_f\), use the Euler product with the primes
dividing \(f\) removed.

\begin{proposition}[Prime-phase positivity identity]
\label{prop:v74-prime-positivity}
For every \(\sigma>1\), \(\tau\in\mathbb R\), and \(r\ge0\),
\[
 \boxed{
 \begin{aligned}
 &a_0\mathcal Q_{\mathbf1_f,r}(\sigma)
 +2\Re\sum_{k=1}^K
 a_k\mathcal Q_{\chi^k,r}(\sigma+ik\tau)\\
 &\quad=\frac1{r!}\sum_p\sum_{m\ge1}
 \frac{(\log p)(m\log p)^r}{p^{m\sigma}}
 P\!\left(m(\arg\chi(p)-\tau\log p)\right)
 \ge0.
 \end{aligned}}
 \tag{V74.18}
\]
The identity is absolutely convergent.
\end{proposition}

\begin{proof}
The Euler expansion of \textup{(V74.17)} is
\[
 \mathcal Q_{\eta,r}(s)
 =\frac1{r!}\sum_p\sum_{m\ge1}
 (\log p)(m\log p)^r\eta(p)^m p^{-ms}.
 \tag{V74.19}
\]
Insert \(s=\sigma+ik\tau\) and \(\eta=\chi^k\), combine conjugate Fourier
coefficients, and use \textup{(V74.15)}.  Every remaining scalar weight is
nonnegative.
\end{proof}

This is the full family of fixed finite prime-phase positivity certificates
suggested in V73.  It includes the classical de la Vall\'ee Poussin type
polynomials, now at every derivative order.

\subsection{Why the principal pole dominates every target zero}

The principal factor has a simple pole at \(s=1\), so for fixed
\(\sigma>1\),
\[
 \mathcal Q_{\mathbf1_f,r}(\sigma)
 =\frac1{(\sigma-1)^{r+1}}
 +\mathcal R_{f,r}(\sigma),
 \tag{V74.20}
\]
where the normalized remainder has a strictly smaller root rate on every
fixed zero-free disk about \(\sigma\).

If \(L(s,\chi)\) has a zero \(\rho=\beta+i\gamma\), then at
\(s=\sigma+i\gamma\) its local contribution to
\(\mathcal Q_{\chi,r}\) is
\[
 -\frac{m_\rho}{(\sigma-\beta)^{r+1}}.
 \tag{V74.21}
\]
But
\[
 0<\sigma-1<\sigma-\beta
 \qquad(\beta<1),
 \tag{V74.22}
\]
and hence
\[
 \frac{(\sigma-\beta)^{-(r+1)}}
      {(\sigma-1)^{-(r+1)}}
 =\left(\frac{\sigma-1}{\sigma-\beta}\right)^{r+1}
 \longrightarrow0.
 \tag{V74.23}
\]

\begin{theorem}[Fixed-polynomial principal-pole no-go]
\label{thm:v74-principal-pole-no-go}
Let \(P\not\equiv0\) be any fixed nonnegative trigonometric polynomial whose
first harmonic is used to test the \(\chi\)-channel.  In the positivity
identity \textup{(V74.18)}, the forced principal contribution has coefficient
\(a_0>0\), with \(|a_1|\le a_0\).  For every fixed
\(\sigma>1\) and every zero \(\beta+i\gamma\) with \(\beta<1\), the
principal-pole term eventually exceeds the magnitude of that target-zero
term by an exponential factor:
\[
 \boxed{
 \frac{a_0(\sigma-1)^{-(r+1)}}
 {|a_1|m_\rho(\sigma-\beta)^{-(r+1)}}
 \longrightarrow\infty.}
 \tag{V74.24}
\]
Therefore Theorem~\ref{thm:v74-block-visibility} cannot be converted into a
contradiction by a fixed Euler-positive trigonometric polynomial and
asymptotic dominance alone.  At \(\sigma=1\) the Euler-positive principal
series diverges, so taking a boundary limit does not remove the obstruction.
\end{theorem}

\begin{proof}
Nonnegativity and nontriviality force \(a_0>0\), while Fourier averaging
gives \(|a_1|\le a_0\).  Equations \textup{(V74.20)}--
\textup{(V74.23)} prove the ratio.  The principal Euler series contains
\(\sum_p(\log p)p^{-1}\) at the boundary and therefore diverges.
\end{proof}

The theorem is deliberately scoped.  It does not rule out an identity that
uses arithmetic relations among several zero channels, or an order-dependent
signed compensation of the principal pole.  It proves that nonnegative
prime-phase algebra by itself spends a larger high-moment budget than the
off-line zero it is meant to detect.

\begin{theorem}[V74 Tur\'an-positivity verdict]
\label{thm:v74-verdict}
The first boundary-moment power-sum attack has the following exact outcome.
\begin{enumerate}[label=\textup{(\roman*)},leftmargin=3em]
\item The normalized prime moments equal a nontrivial-zero power sum plus a
      holomorphic remainder of root rate at most two.
\item Closest off-line zeros survive all cancellation in every sufficiently
      high block of \(q_t\) consecutive orders.
\item Nonnegative trigonometric polynomials give exact Euler-positive
      combinations of all character-power logarithmic derivatives.
\item Every such nonzero polynomial necessarily imports the principal mode.
\item The principal pole is closer than every target zero with \(\beta<1\),
      and therefore dominates its high-order contribution exponentially.
\end{enumerate}
Thus zero-side block visibility is rigorous, but fixed prime-phase positivity
cannot supply the required upper root rate two.
\end{theorem}

\begin{firewall}[Power-sum visibility is not exclusion]
Theorem~\ref{thm:v74-block-visibility} says that a forbidden zero cannot hide
from all high boundary moments.  It does not bound those moments from the
prime side.  Proposition~\ref{prop:v74-prime-positivity} supplies positivity,
but Theorem~\ref{thm:v74-principal-pole-no-go} proves that its unavoidable
principal budget is asymptotically larger than the target signal.  GRH
remains open.
\end{firewall}

\begin{assessment}[V74 closure and V75 direction]
V74 upgrades the V73 limsup criterion to finite-block Tur\'an visibility and
then closes the most direct positive trigonometric-polynomial proof.  The
failure is not phase cancellation among the target zeros; it is the closer
principal pole forced by every nonnegative prime-phase certificate.

V75 should test principal-pole compensation itself.  The first exact problem
is to combine logarithmic-derivative moments at two radial points
\(\sigma\) and \(\sigma+a\) so that the pole terms cancel, and then inspect
the induced Euler weight
\(1-\lambda p^{-ma}\).  The decision is whether any pole-cancelling choice
can preserve nonnegativity for all prime powers, or whether cancellation
necessarily introduces a negative low-prime sector whose capacity matches
the removed pole.  No further fixed-polynomial optimization should be
attempted before this sign question is settled.
\end{assessment}

\section*{Version V74 append-only record}
\addcontentsline{toc}{section}{Version V74 append-only record}
The complete V73 source was retained before this continuation.  It had
1060122 bytes, 31332 lines, and SHA--256
\[
 \texttt{50B8F361FA370C5B4F4AA1FE2D4713B3A3B616AAF27C671557B0C79853E28578}.
\]
V74 derived the exact zero power sum for the boundary prime moments, proved
Vandermonde block visibility of every closest off-line zero, constructed the
full Euler-positive trigonometric family at all derivative orders, and
proved that its unavoidable principal pole dominates the target-zero signal.

% =============================================================================
% END APPENDED V74 MATERIAL
% =============================================================================
% =============================================================================
% BEGIN APPENDED V75 MATERIAL
% =============================================================================

\section{V75: two-radius pole compensation and the unavoidable Euler sign split}
\label{sec:grh-v75}

V74 proved that every fixed nonnegative prime-phase certificate imports a
principal pole whose high-order moment dominates the target zero.  It
proposed cancelling that pole by comparing two radial points.  This version
solves the sign problem exactly.

At derivative order \(r\), exact Laurent-pole cancellation fixes a unique
coefficient
\[
 \lambda_r=
 \left(\frac{\sigma+a-1}{\sigma-1}\right)^{r+1}.
 \tag{V75.1}
\]
On the Euler side it creates the weight
\(1-\lambda_r n^{-a}\).  For every fixed separation \(a>0\), this weight
splits the principal mass into a negative low-frequency saddle and a
positive high-frequency saddle; each has the full uncancelled principal
root rate.  If the separation is allowed to grow so fast that every
arithmetic weight remains nonnegative, then the compensating radial moment
is superexponentially negligible and the original principal budget remains.
Moreover, in the orientation with positive high-prime weights, an off-line
zero itself acquires the favorable, not contradictory, sign.

\subsection{The unique two-radius compensator}

Retain the normalized logarithmic-derivative moments
\[
 \mathcal Q_{\eta,r}(s)
 :=\frac{(-1)^r}{r!}\frac{d^r}{ds^r}
 \left(-\frac{L'}{L}(s,\eta)\right)
 \tag{V75.2}
\]
from V74.  Fix \(\sigma>1\), set
\[
 d:=\sigma-1>0,
 \tag{V75.3}
\]
and let \(a>0\).  The principal Laurent terms at \(\sigma\) and
\(\sigma+a\) are \(d^{-(r+1)}\) and
\((d+a)^{-(r+1)}\), respectively.

\begin{lemma}[Uniqueness of normalized pole cancellation]
\label{lem:v75-unique-compensator}
Up to an overall scalar, the unique two-radius combination which cancels
the principal Laurent pole at order \(r\) is
\[
 \boxed{
 \mathcal D_{\eta,r}^{\sigma,a}(s)
 :=\mathcal Q_{\eta,r}(s)
 -\lambda_r\mathcal Q_{\eta,r}(s+a),
 \qquad
 \lambda_r=\left(\frac{d+a}{d}\right)^{r+1}.}
 \tag{V75.4}
\]
For \(\Re s=\sigma\), its Euler expansion is
\[
 \boxed{
 \mathcal D_{\eta,r}^{\sigma,a}(s)
 =\frac1{r!}\sum_p\sum_{m\ge1}
 \frac{(\log p)(m\log p)^r\eta(p)^m}{p^{ms}}
 \left(1-\lambda_r p^{-ma}\right).}
 \tag{V75.5}
\]
\end{lemma}

\begin{proof}
If \(\alpha d^{-(r+1)}+\beta(d+a)^{-(r+1)}=0\), then
\(\beta/\alpha=-((d+a)/d)^{r+1}\).  Normalize \(\alpha=1\).
Equation \textup{(V75.5)} follows by subtracting the two absolutely
convergent Euler expansions.
\end{proof}

Put
\[
 U_r:=\frac{\log\lambda_r}{a}
 =\frac{r+1}{a}\log\left(1+\frac ad\right).
 \tag{V75.6}
\]
The Euler multiplier in \textup{(V75.5)} is negative for
\(m\log p<U_r\), zero at equality, and positive for
\(m\log p>U_r\).  Since
\[
 \frac1{d+a}
 <\frac1a\log\left(1+\frac ad\right)
 <\frac1d,
 \tag{V75.7}
\]
the sign-change point lies strictly between the two radial Laplace saddles
\((r+1)/(d+a)\) and \((r+1)/d\).

\subsection{The principal moment as a tilted probability law}

Let \(\mathbf1_f\) denote the principal Euler product with ramified primes
removed, and write
\[
 Q_r(\sigma):=\mathcal Q_{\mathbf1_f,r}(\sigma)
 =\frac1{r!}\sum_{\substack{n\ge2\\ \gcd(n,f)=1}}
 \frac{\Lambda(n)(\log n)^r}{n^\sigma}.
 \tag{V75.8}
\]
The pole at one is the unique nearest singularity to the real point
\(\sigma>1\), so
\[
 \boxed{
 Q_r(\sigma)=d^{-(r+1)}(1+o(1))}
 \qquad(r\to\infty).
 \tag{V75.9}
\]

Normalize the summands in \textup{(V75.8)} to a probability measure
\(\nu_{r,d}\) on \(u=\log n\).  For every real \(\theta<d\),
\[
 \int e^{\theta u}\,d\nu_{r,d}(u)
 =\frac{Q_r(\sigma-\theta)}{Q_r(\sigma)}
 =\left(\frac d{d-\theta}\right)^{r+1}(1+o(1)).
 \tag{V75.10}
\]
Chernoff's inequality therefore gives exponential concentration
\[
 \frac{u}{r+1}\longrightarrow\frac1d
 \quad\hbox{in }\nu_{r,d}\hbox{-probability}.
 \tag{V75.11}
\]
The analogous law at \(\sigma+a\) concentrates at \(1/(d+a)\).

\begin{lemma}[The two sign sectors retain the full pole scale]
\label{lem:v75-two-saddles}
Fix \(a>0\).  Split the absolute principal compensator at \(U_r\):
\[
 \begin{aligned}
 N_r&:=\frac1{r!}
 \sum_{\substack{(n,f)=1\\\log n<U_r}}
 \frac{\Lambda(n)(\log n)^r}{n^\sigma}
 \left(\lambda_r n^{-a}-1\right),\\
 P_r&:=\frac1{r!}
 \sum_{\substack{(n,f)=1\\\log n>U_r}}
 \frac{\Lambda(n)(\log n)^r}{n^\sigma}
 \left(1-\lambda_r n^{-a}\right).
 \end{aligned}
 \tag{V75.12}
\]
Then
\[
 \boxed{
 N_r\sim Q_r(\sigma),\qquad
 P_r\sim Q_r(\sigma),\qquad
 N_r+P_r\sim\frac{2}{d^{r+1}}.}
 \tag{V75.13}
\]
Thus cancellation of the signed Laurent pole does not reduce the total
variation root rate.
\end{lemma}

\begin{proof}
By \textup{(V75.7)}, \(U_r/(r+1)\) is below the
\(\nu_{r,d}\) concentration point and above the
\(\nu_{r,d+a}\) concentration point.  Hence
\[
 \nu_{r,d}([0,U_r])\to0,
 \qquad
 \nu_{r,d+a}([U_r,\infty))\to0
 \tag{V75.14}
\]
exponentially.  Also, by \textup{(V75.1)} and
\textup{(V75.9)},
\[
 \lambda_r Q_r(\sigma+a)\sim Q_r(\sigma).
 \tag{V75.15}
\]
In the positive sector, the first radial measure contributes
\((1+o(1))Q_r(\sigma)\), while the scaled second radial measure contributes
\(o(Q_r(\sigma))\).  In the negative sector their roles reverse.  This proves
\textup{(V75.13)}.
\end{proof}

For fixed \(a\), \(U_r\to\infty\).  The negative sector therefore contains
an increasing set of actual prime powers; it is not an artifact below the
arithmetic support.

\subsection{The variable-separation dichotomy}

Let
\[
 \ell_f:=\min\{p\text{ prime}:p\nmid f\},
 \qquad L_f:=\log\ell_f.
 \tag{V75.16}
\]
All supported prime powers satisfy \(m\log p\ge L_f\).

\begin{theorem}[Positivity or compensation, but not both]
\label{thm:v75-positivity-dichotomy}
Allow \(a=a_r>0\) to depend on \(r\), and retain the exact pole-cancelling
\(\lambda_r\) from \textup{(V75.1)}.
\begin{enumerate}[label=\textup{(\roman*)},leftmargin=3em]
\item If \(U_r>L_f\), then the Euler multiplier
      \(1-\lambda_r p^{-ma_r}\) is negative on a nonempty low-prime sector;
      termwise Euler positivity is lost.
\item If \(U_r\le L_f\), then every supported Euler multiplier is
      nonnegative, but
\[
 \boxed{
 \frac{\lambda_r Q_r(\sigma+a_r)}{Q_r(\sigma)}
 \longrightarrow0,\qquad
 \frac{Q_r(\sigma)-\lambda_r Q_r(\sigma+a_r)}
      {Q_r(\sigma)}
 \longrightarrow1.}
 \tag{V75.17}
\]
\end{enumerate}
Thus a separation large enough to preserve all Euler signs moves the
compensating radius outside the high-order saddle and leaves the full
principal moment behind.
\end{theorem}

\begin{proof}
The multiplier is increasing in \(m\log p\), so its sign on the arithmetic
support is determined by whether \(U_r\) exceeds \(L_f\).  This proves
\textup{(i)}.

Assume \(U_r\le L_f\).  Equivalently,
\[
 \lambda_r e^{-a_rL_f}\le1.
 \tag{V75.18}
\]
The same inequality forces \(a_r>r/L_f\) for all sufficiently large \(r\).
Indeed, \(a\mapsto a^{-1}\log(1+a/d)\) is decreasing.  If
\(a_r\le r/L_f\), then
\[
 \frac{r+1}{a_r}\log(1+a_r/d)
 \ge\frac{(r+1)L_f}{r}
 \log\left(1+\frac{r}{dL_f}\right)\longrightarrow\infty,
\]
contradicting the bound by \(L_f\) obtained from \textup{(V75.18)}.
Consequently \(u^re^{-a_ru}\) is decreasing for \(u\ge L_f\), and
\[
 \begin{aligned}
 Q_r(\sigma+a_r)
 &\le
 \frac{L_f^re^{-a_rL_f}}{r!}
 \sum_{\substack{n\ge2\\ \gcd(n,f)=1}}\frac{\Lambda(n)}{n^\sigma}\\
 &\ll_{\sigma,f}\frac{L_f^re^{-a_rL_f}}{r!}.
 \end{aligned}
 \tag{V75.19}
\]
Multiply by \(\lambda_r\), use \textup{(V75.18)}, and compare with
\textup{(V75.9)}:
\[
 \frac{\lambda_rQ_r(\sigma+a_r)}{Q_r(\sigma)}
 \ll_{\sigma,f}
 \frac{d^{r+1}L_f^r}{r!}\longrightarrow0.
 \tag{V75.20}
\]
This proves \textup{(V75.17)}.
\end{proof}

The apparent Laurent cancellation in case \textup{(ii)} occurs below the
smallest Euler frequency.  The analytic remainder at the remote radial point
then cancels its own formal pole term, while the actual Dirichlet moment is
too small to affect \(Q_r(\sigma)\).

\subsection{The compensated sign of a target zero}

There is a second obstruction independent of total variation.  Suppose
\(\rho=\beta+i\gamma\) is a zero of \(L(s,\chi)\), and evaluate the
two-radius combination at \(s=\sigma+i\gamma\).  Put
\[
 d_\beta:=\sigma-\beta>d.
 \tag{V75.21}
\]
The target-zero contribution to \(\mathcal D_{\chi,r}^{\sigma,a}\) is
\[
 m_\rho\left\{
 \frac{\lambda_r}{(d_\beta+a)^{r+1}}
 -\frac1{d_\beta^{r+1}}
 \right\}.
 \tag{V75.22}
\]
Since
\[
 \frac{d+a}{d}\frac{d_\beta}{d_\beta+a}>1,
 \tag{V75.23}
\]
the expression in braces is strictly positive.

\begin{theorem}[Pole compensation reverses the target sign]
\label{thm:v75-target-sign}
In the orientation
\[
 \mathcal Q_{\eta,r}(s)
 -\lambda_r\mathcal Q_{\eta,r}(s+a),
 \tag{V75.24}
\]
whose Euler multiplier is positive at high prime frequencies, exact
principal-pole cancellation makes every aligned zero with \(\beta<1\)
contribute positively.  Reversing the orientation makes the zero contribution
negative but also makes every sufficiently high Euler multiplier negative.
Thus the two-radius compensator cannot simultaneously cancel the principal
pole, preserve the high-prime sign, and retain a negative target-zero
witness.
\end{theorem}

\begin{proof}
Equation \textup{(V75.22)} follows from the local zero term
\(-m_\rho(s-\rho)^{-(r+1)}\) in V74.  Inequality
\textup{(V75.23)} is equivalent to
\(a(d_\beta-d)>0\).  Reversing the combination reverses both the zero
contribution and the limiting Euler weight.
\end{proof}

\begin{theorem}[V75 radial-compensation verdict]
\label{thm:v75-verdict}
Every two-radius cancellation of the principal Laurent pole obeys the
following exhaustive ledger.
\begin{enumerate}[label=\textup{(\roman*)},leftmargin=3em]
\item Its coefficient is uniquely fixed by \textup{(V75.1)}.
\item For fixed radial separation, its Euler weight changes sign between two
      distinct full-scale Laplace saddles.
\item The absolute masses of the negative and positive sectors are each
      asymptotic to the uncancelled principal moment.
\item If the separation varies so as to preserve every arithmetic sign, the
      compensating moment is negligible and the principal budget survives.
\item In the high-frequency-positive orientation, the target-zero
      contribution has the wrong sign for a contradiction.
\end{enumerate}
Therefore two-radius compensation neither preserves the V74 positivity cone
nor reduces its principal root-rate obstruction.
\end{theorem}

\begin{firewall}[Laurent cancellation is not Euler cancellation]
The pole terms in the meromorphic expansion cancel exactly, but the Euler
series records that cancellation through two large signed frequency sectors.
Discarding the negative sector is invalid; taking total variation restores
the full principal rate.  Moving the second radius far enough to avoid the
negative sector makes it arithmetically ineffective.  GRH remains open.
\end{firewall}

\begin{assessment}[V75 closure and V76 direction]
V75 settles the sign question posed in V74 and closes all two-radius
principal-pole compensators.  The obstruction is stronger than a poor choice
of separation: exact pole cancellation, Euler nonnegativity, and a negative
target-zero sign are mutually incompatible.

V76 should test whether a positive mixture of three or more radial points can
do better.  The correct formulation is a signed exponential polynomial
\[
 W_r(u)=\sum_{j=0}^J b_{j,r}e^{-a_{j,r}u}
 \tag{V75.25}
\]
whose coefficients annihilate the principal pole moments
\(\sum_jb_{j,r}(d+a_{j,r})^{-(r+1)}\), while \(W_r(m\log p)\ge0\) on every
supported prime power and the aligned target-zero functional remains
negative.  The first task is a convex-duality/separation theorem deciding
whether these three constraints are compatible.  If the positive
principal moment measure separates them, radial compensation is closed at
all finite orders at once.
\end{assessment}

\section*{Version V75 append-only record}
\addcontentsline{toc}{section}{Version V75 append-only record}
The complete V74 source was retained before this continuation.  It had
1073040 bytes, 31705 lines, and SHA--256
\[
 \texttt{5DB5B7BC666B6997E9C22679F1643E8F9D9FBC117B687735710D35BFC133E96E}.
\]
V75 derived the unique two-radius pole compensator, proved the full-scale
positive/negative saddle split, established the variable-separation
positivity dichotomy, and showed that the high-frequency-positive
orientation reverses the target-zero sign.

% =============================================================================
% END APPENDED V75 MATERIAL
% =============================================================================
% =============================================================================
% BEGIN APPENDED V76 MATERIAL
% =============================================================================

\section{V76: the multi-radius moment cone and the sub-prime feasibility gap}
\label{sec:grh-v76}

V75 proved that no two-radius functional can simultaneously cancel the
principal Laurent pole, preserve the Euler sign, and retain a negative
target-zero contribution.  It proposed the corresponding finite
multi-radius problem.  The convex geometry has a different answer: once
three or more radii are allowed, the three sign constraints are compatible.

The reason is a support mismatch.  A principal Laurent moment is an integral
over every \(u>0\), whereas the Euler logarithmic derivative samples only
\(u=m\log p\), beginning at the first unramified prime frequency.  A finite
exponential polynomial can place a positive--negative--positive oscillation
entirely below that first frequency.  Its negative continuous mass cancels
the principal pole, while every actual Euler multiplier remains positive and
the target-zero moment has the desired sign.  Thus the hoped-for convex
separation theorem is false without an additional coefficient or
holomorphic-remainder budget.

\subsection{The continuous moment functionals}

Let
\[
 k:=r+1\ge1,\qquad 0<d<x,
 \tag{V76.1}
\]
where \(d=\sigma-1\) is the principal-pole distance and
\(x=\sigma-\beta\) is the aligned target-zero distance.  For a finite real
exponential polynomial
\[
 W(u)=\sum_{j=0}^J b_je^{-a_ju},
 \qquad 0\le a_0<a_1<\cdots<a_J,
 \tag{V76.2}
\]
define
\[
 \mathcal L_y^{(k)}[W]
 :=\frac1{(k-1)!}\int_0^\infty
 u^{k-1}e^{-yu}W(u)\,du
 =\sum_{j=0}^J\frac{b_j}{(y+a_j)^k}.
 \tag{V76.3}
\]
The principal-pole cancellation and negative target-zero conditions are
\[
 \mathcal L_d^{(k)}[W]=0,
 \qquad
 \mathcal L_x^{(k)}[W]>0,
 \tag{V76.4}
\]
because a target zero contributes
\(-m_\rho\mathcal L_x^{(k)}[W]\).

Let
\[
 \ell_f:=\min\{p\text{ prime}:p\nmid f\},
 \qquad L_f:=\log\ell_f.
 \tag{V76.5}
\]
Every Euler frequency \(m\log p\) lies in \([L_f,\infty)\).

\subsection{What continuous positivity would imply}

\begin{proposition}[The full continuous cone is separated]
\label{prop:v76-continuous-separation}
If \(W(u)\ge0\) for every \(u\ge0\) and
\(\mathcal L_d^{(k)}[W]=0\), then \(W\equiv0\).  In particular,
\(\mathcal L_x^{(k)}[W]=0\).
\end{proposition}

\begin{proof}
The density \(u^{k-1}e^{-du}/(k-1)!\) is strictly positive for \(u>0\).
A nonnegative continuous function having zero integral against this density
vanishes on \((0,\infty)\), and hence identically by continuity.
\end{proof}

Thus a separating positive measure exists if Euler positivity is available
at every frequency.  It is not available below \(L_f\).

\begin{proposition}[At least two hidden sign changes are necessary]
\label{prop:v76-two-sign-changes}
Suppose
\[
 W(u)\ge0\quad(u\ge L_f),\qquad
 \mathcal L_d^{(k)}[W]=0,\qquad
 \mathcal L_x^{(k)}[W]>0.
 \tag{V76.6}
\]
Then \(W\) has at least two sign changes in \((0,L_f)\).  Consequently a
nontrivial exponential polynomial satisfying \textup{(V76.6)} has at least
three exponential terms.
\end{proposition}

\begin{proof}
The first two conditions force \(W<0\) somewhere below \(L_f\).  If \(W\)
had at most one sign change, its eventual nonnegative sign would give a
point \(c<L_f\) such that \(W\le0\) on \((0,c)\) and \(W\ge0\) on
\((c,\infty)\), up to zero intervals.  Since \(x>d\), the ratio of the two
moment densities is the strictly decreasing function
\[
 e^{-(x-d)u}.
 \tag{V76.7}
\]
It is at least \(e^{-(x-d)c}\) on the negative part and at most that value
on the positive part.  Therefore
\[
 \mathcal L_x^{(k)}[W]
 \le e^{-(x-d)c}\mathcal L_d^{(k)}[W]=0,
 \tag{V76.8}
\]
contrary to \textup{(V76.6)}.

The functions \(e^{-a_ju}\) with distinct real exponents form a Chebyshev
system: a nonzero linear combination of \(J+1\) such functions has at most
\(J\) real zeros counted with multiplicity.  Two sign changes therefore
require \(J+1\ge3\).
\end{proof}

This recovers the V75 obstruction conceptually: two radii permit only one
sign change, which necessarily gives the target zero the wrong sign.

\subsection{Finite multi-radius feasibility}

The converse is the decisive point.

\begin{theorem}[Sub-prime feasibility of the three constraints]
\label{thm:v76-feasibility}
For every \(k\ge1\), \(0<d<x\), and \(L>0\), there is a finite real
exponential polynomial with integer-spaced radii,
\[
 W(u)=\sum_{j=0}^Jb_je^{-ju},
 \tag{V76.9}
\]
such that
\[
 \boxed{
 W(u)>0\quad(u\ge L),\qquad
 \mathcal L_d^{(k)}[W]=0,\qquad
 \mathcal L_x^{(k)}[W]>0.}
 \tag{V76.10}
\]
In particular, with \(L=L_f\), every actual prime-power multiplier
\(W(m\log p)\) is strictly positive while the principal Laurent moment
vanishes and an aligned target zero contributes negatively.
\end{theorem}

\begin{proof}
Choose
\[
 0<u_1<u_2<L
 \tag{V76.11}
\]
and nonnegative continuous bump functions \(\phi_1,\phi_2\) with disjoint
small supports about \(u_1,u_2\), respectively.  Normalize them so that
\[
 \mathcal L_d^{(k)}[\phi_1]
 =\mathcal L_d^{(k)}[\phi_2]=1.
 \tag{V76.12}
\]
By making the supports sufficiently small and using the strict decrease of
\(e^{-(x-d)u}\), one has
\[
 \mathcal L_x^{(k)}[\phi_1]
 >\mathcal L_x^{(k)}[\phi_2].
 \tag{V76.13}
\]

For a sufficiently small \(\varepsilon>0\), put
\[
 w(u):=\phi_1(u)
 -\bigl(1+\varepsilon d^{-k}\bigr)\phi_2(u)
 +\varepsilon.
 \tag{V76.14}
\]
Because \(\mathcal L_d^{(k)}[1]=d^{-k}\), one has
\[
 \mathcal L_d^{(k)}[w]=0.
 \tag{V76.15}
\]
At \(\varepsilon=0\), its \(x\)-moment is strictly positive by
\textup{(V76.13)}, so it remains positive for all sufficiently small
\(\varepsilon\).  Outside the two bump supports, and in particular for
\(u\ge L\), one has \(w(u)=\varepsilon\).

Set \(q=e^{-u}\) and define the continuous function
\[
 g(q):=
 \begin{cases}
 w(-\log q),&0<q\le1,\\
 \varepsilon,&q=0.
 \end{cases}
 \tag{V76.16}
\]
The Weierstrass theorem supplies a real polynomial \(P(q)\) with
\[
 \sup_{0\le q\le1}|P(q)-g(q)|<\delta.
 \tag{V76.17}
\]
Choose \(\delta>0\) so small that the strict \(x\)-moment margin survives
and \(2\delta<\varepsilon\).  Uniform approximation gives
\[
 |\mathcal L_y^{(k)}[P(e^{-u})-w(u)]|
 \le\delta y^{-k}
 \qquad(y>0).
 \tag{V76.18}
\]
Define
\[
 c:=d^k\mathcal L_d^{(k)}[P(e^{-u})],
 \qquad
 W(u):=P(e^{-u})-c.
 \tag{V76.19}
\]
Since \(\mathcal L_d^{(k)}[w]=0\), equation
\textup{(V76.18)} gives \(|c|\le\delta\), while
\(\mathcal L_d^{(k)}[W]=0\) exactly.  For \(u\ge L\),
\[
 W(u)\ge\varepsilon-2\delta>0.
 \tag{V76.20}
\]
The total change in the \(x\)-moment is at most \(2\delta x^{-k}\), so
\(\mathcal L_x^{(k)}[W]>0\) after the initial choice of \(\delta\).
Finally, expanding \(P\) in powers of \(q=e^{-u}\) gives
\textup{(V76.9)}.
\end{proof}

The theorem is stronger than positivity merely at the discrete prime-power
set: it gives positivity on the whole continuous arithmetic half-line.
Therefore no refinement based only on denser sampling of large prime
frequencies can restore the missing separation.

\subsection{Where the cancellation is hidden}

\begin{corollary}[The pole debit lies entirely below the first prime]
\label{cor:v76-subprime-debit}
Every \(W\) produced by Theorem~\ref{thm:v76-feasibility} satisfies
\[
 \boxed{
 \frac1{(k-1)!}\int_0^{L_f}
 u^{k-1}e^{-du}W(u)\,du
 =
 -\frac1{(k-1)!}\int_{L_f}^{\infty}
 u^{k-1}e^{-du}W(u)\,du<0.}
 \tag{V76.21}
\]
Thus all negative mass financing the principal-pole cancellation occurs in
an interval containing no Euler logarithmic-derivative frequency.
\end{corollary}

\begin{proof}
Split \(\mathcal L_d^{(k)}[W]=0\) at \(L_f\).  The second integral is
strictly positive by \textup{(V76.10)}.
\end{proof}

At least one positive component also occurs below \(L_f\), by
Proposition~\ref{prop:v76-two-sign-changes}.  The target density
\(e^{-xu}u^{k-1}\) weights that earlier positive component more strongly
than the later negative one, producing the sign in
\(\mathcal L_x^{(k)}[W]\).

\subsection{The exact multi-radius Euler functional}

For a polynomial \(W\) as in \textup{(V76.2)}, define
\[
 \mathcal D_{\eta,r}^{W}(s)
 :=\sum_{j=0}^Jb_j\mathcal Q_{\eta,r}(s+a_j).
 \tag{V76.22}
\]
Absolute Euler expansion gives
\[
 \boxed{
 \mathcal D_{\eta,r}^{W}(s)
 =\frac1{r!}\sum_p\sum_{m\ge1}
 \frac{(\log p)(m\log p)^r\eta(p)^m}{p^{ms}}
 W(m\log p).}
 \tag{V76.23}
\]
For the principal character this is nonnegative whenever
\(W\ge0\) on \([L_f,\infty)\).  More generally, multiplying the V74
prime-phase identity by \(W(m\log p)\) preserves its termwise positivity.
On the zero side, the principal Laurent contribution is
\(\mathcal L_d^{(k)}[W]\), while an aligned target zero contributes
\[
 -m_\rho\mathcal L_x^{(k)}[W].
 \tag{V76.24}
\]
Theorem~\ref{thm:v76-feasibility} therefore satisfies exactly the three
formal constraints posed at the end of V75.

This is not a GRH contradiction.  The same linear combination also acts on
every other zero in every power-character channel and on the holomorphic
gamma and Euler corrections.  The approximation proof places no useful
bound on the degree, coefficient norm, or shifted analytic remainder of
\(W\).  Those uncontrolled terms can mask
\textup{(V76.24)}, just as the unrestricted positive background masked the
finite Pick defect in V69.

\begin{theorem}[Failure of the proposed convex separator]
\label{thm:v76-no-separator}
Let \(\mathcal C_{L_f}\) be the cone of finite real exponential polynomials
nonnegative on \([L_f,\infty)\).  On the hyperplane
\(\mathcal L_d^{(k)}[W]=0\), the functional
\(\mathcal L_x^{(k)}\) takes positive values.  Consequently no dual
separation argument using only
\[
 W\in\mathcal C_{L_f},
 \qquad \mathcal L_d^{(k)}[W]=0
 \tag{V76.25}
\]
can prove \(\mathcal L_x^{(k)}[W]\le0\).  A valid obstruction must add a
quantitative constraint controlling the coefficients or all shifted
holomorphic remainders.
\end{theorem}

\begin{proof}
Theorem~\ref{thm:v76-feasibility} gives a strict counterexample to the
proposed separating inequality.
\end{proof}

\begin{theorem}[V76 multi-radius verdict]
\label{thm:v76-verdict}
The finite multi-radius feasibility problem has the following complete
qualitative answer.
\begin{enumerate}[label=\textup{(\roman*)},leftmargin=3em]
\item Positivity on the full frequency half-line \([0,\infty)\) forces a
      pole-annihilating exponential polynomial to vanish.
\item Euler positivity begins only at \(L_f>0\), leaving a genuine unsampled
      interval.
\item The desired pole and target moment signs require at least two hidden
      sign changes and hence at least three radial terms.
\item With sufficiently many finite radii, all three formal constraints are
      compatible even under positivity on the entire interval
      \([L_f,\infty)\).
\item The compensation is financed by negative continuous moment mass below
      the first prime and carries no coefficient or remainder control.
\end{enumerate}
Thus the two-radius no-go does not extend by convexity alone; the missing
invariant is quantitative stability of the radial combination.
\end{theorem}

\begin{firewall}[Formal feasibility is not zero exclusion]
The polynomial in Theorem~\ref{thm:v76-feasibility} is an existence
certificate in the radial moment cone.  It is not shown to have bounded
degree, bounded coefficients, or a controlled action on the remaining zero
and gamma terms as \(r\to\infty\).  Its negative pole debit lies where the
Euler series has no samples.  No off-line zero has been excluded, and GRH
remains open.
\end{firewall}

\begin{assessment}[V76 closure and V77 direction]
V76 answers the V75 convex-separation question: the constraints are feasible,
and the obstruction to a proof is not qualitative sign but quantitative
conditioning.  The exact sub-prime support gap explains why the continuous
principal measure cannot separate the discrete Euler cone.

V77 should introduce a stability cost adapted to the holomorphic remainder,
for example
\[
 \mathfrak C_k(W;R)
 :=\sum_{j=0}^J|b_j|(R+a_j)^{-k},
 \tag{V76.26}
\]
and normalize \(\mathcal L_x^{(k)}[W]=1\).  The next decision is whether every
feasible pole-annihilating \(W\ge0\) on \([L_f,\infty)\) must have
\(\mathfrak C_k(W;R)\) at least as large as the target signal for some
\(R\ge1/2\).  Such a lower bound would show that the gamma/power correction
necessarily has enough capacity to mask the constructed zero witness; a
counterexample with subcritical cost would supply the first genuinely usable
multi-radius certificate.
\end{assessment}

\section*{Version V76 append-only record}
\addcontentsline{toc}{section}{Version V76 append-only record}
The complete V75 source was retained before this continuation.  It had
1085693 bytes, 32083 lines, and SHA--256
\[
 \texttt{D1D08656C0DF5C1F885C421A3568BCB31FC963C995DEC09BDAB3A3279E222EA1}.
\]
V76 solved the qualitative multi-radius moment problem: it proved the
necessity of two hidden sub-prime sign changes, constructed finite
exponential-polynomial certificates satisfying all three formal constraints,
and showed that convex separation fails without a coefficient/remainder
stability budget.

% =============================================================================
% END APPENDED V76 MATERIAL
% =============================================================================
% =============================================================================
% BEGIN APPENDED V77 MATERIAL
% =============================================================================

\section{V77: the quantitative stability price of sub-prime compensation}
\label{sec:grh-v77}

V76 proved that the qualitative multi-radius constraints are feasible:
a finite exponential polynomial can be positive at every Euler frequency,
annihilate the principal Laurent moment, and give an aligned off-line zero
the negative sign.  The construction left its conditioning uncontrolled.
This version computes a universal lower bound for exactly the stability cost
proposed there.

The result closes the apparent loophole.  After the target moment is
normalized to one, the coefficient-weighted Cauchy cost cannot decay
exponentially with the moment order.  At the natural distance from
\(\Re s=1\) to the critical line, it is bounded below by the fixed constant
\(\ell_f^{-1/2}\).  In addition, the ordinary coefficient norm is forced to
grow at least like \(k!/(\log\ell_f)^k\).  Thus the sub-prime construction is
feasible only through a quantitatively ill-conditioned radial combination.

\subsection{The scale-invariant feasible cone}

Retain the V76 notation
\[
 k\ge1,\qquad 0<d<x,\qquad L>0,
 \tag{V77.1}
\]
and let
\[
 W(u)=\sum_{j=0}^Jb_je^{-a_ju},
 \qquad a_j\ge0,
 \tag{V77.2}
\]
be a finite real exponential polynomial satisfying
\[
 W(u)\ge0\quad(u\ge L),\qquad
 \mathcal L_d^{(k)}[W]=0,\qquad
 T:=\mathcal L_x^{(k)}[W]>0.
 \tag{V77.3}
\]
For \(R>d\), define the absolute radial Cauchy cost
\[
 \mathfrak C_k(W;R)
 :=\sum_{j=0}^J|b_j|(R+a_j)^{-k}.
 \tag{V77.4}
\]
Both \(T\) and \(\mathfrak C_k\) scale linearly with \(W\), so their ratio is
the intrinsic quantity.

\subsection{A dual inequality from the support gap}

Put \(\Delta:=x-d>0\).  Since
\(\mathcal L_d^{(k)}[W]=0\),
\[
 \begin{aligned}
 T
 &=\frac1{(k-1)!}\int_0^\infty
 u^{k-1}e^{-du}
 \bigl(e^{-\Delta u}-e^{-\Delta L}\bigr)W(u)\,du.
 \end{aligned}
 \tag{V77.5}
\]
On \(u\ge L\), the parenthesis is nonpositive while \(W(u)\ge0\).
Consequently
\[
 T\le\frac1{(k-1)!}\int_0^L
 u^{k-1}e^{-du}|W(u)|\,du.
 \tag{V77.6}
\]
This inequality quantifies the statement that every positive target moment
must be financed inside the sub-prime interval.

\begin{lemma}[Truncated gamma evaluation bound]
\label{lem:v77-truncated-gamma}
For \(a\ge0\) and \(R>d\),
\[
 \boxed{
 \frac1{(k-1)!}\int_0^L
 u^{k-1}e^{-(d+a)u}\,du
 \le e^{(R-d)L}(R+a)^{-k}.}
 \tag{V77.7}
\]
\end{lemma}

\begin{proof}
Set \(v=(R+a)u\) and \(c=R-d>0\).  The left side equals
\[
 \frac{(R+a)^{-k}}{(k-1)!}
 \int_0^{(R+a)L}
 v^{k-1}e^{-v}
 e^{cv/(R+a)}\,dv.
 \tag{V77.8}
\]
On the integration interval,
\(e^{cv/(R+a)}\le e^{cL}\).  Extending the remaining gamma integral to
\([0,\infty)\) proves \textup{(V77.7)}.
\end{proof}

\begin{theorem}[Universal multi-radius stability lower bound]
\label{thm:v77-stability-lower}
Every feasible \(W\) in \textup{(V77.3)} satisfies
\[
 \boxed{
 T\le e^{(R-d)L}\mathfrak C_k(W;R),\qquad
 \mathfrak C_k(W;R)\ge e^{-(R-d)L}T.}
 \tag{V77.9}
\]
The constants are independent of the number of radii, their locations, the
coefficients, and the moment order \(k\).  In particular, for every feasible
sequence normalized by \(T=1\),
\[
 \boxed{
 \liminf_{k\to\infty}
 \mathfrak C_k(W_k;R)^{1/k}\ge1.}
 \tag{V77.10}
\]
\end{theorem}

\begin{proof}
Expand \(W\) in \textup{(V77.6)}, take absolute values term by term, and
apply Lemma~\ref{lem:v77-truncated-gamma}:
\[
 \begin{aligned}
 T
 &\le\sum_j|b_j|
 \frac1{(k-1)!}\int_0^L
 u^{k-1}e^{-(d+a_j)u}\,du\\
 &\le e^{(R-d)L}
 \sum_j|b_j|(R+a_j)^{-k}.
 \end{aligned}
 \tag{V77.11}
\]
This is \textup{(V77.9)}.  Taking \(k\)-th roots with \(T=1\) gives
\textup{(V77.10)}.
\end{proof}

For the arithmetic support \(L=L_f=\log\ell_f\), the natural half-plane
distance at the point \(s=\sigma+it\) is
\[
 R_0:=\sigma-\frac12=d+\frac12.
 \tag{V77.12}
\]
Theorem~\ref{thm:v77-stability-lower} then gives
\[
 \boxed{
 \mathfrak C_k(W;R_0)
 \ge e^{-L_f/2}T
 =\ell_f^{-1/2}T.}
 \tag{V77.13}
\]
For Cauchy estimates one may first use
\(R_0-\varepsilon\) and then let \(\varepsilon\downarrow0\); the same
root-rate conclusion holds.

\subsection{Factorial coefficient amplification}

The weighted cost bound has an elementary unweighted companion.

\begin{theorem}[Sub-prime amplitude and coefficient lower bounds]
\label{thm:v77-factorial-cost}
Every feasible \(W\) satisfies
\[
 \boxed{
 \sup_{0\le u\le L}|W(u)|
 \ge\frac{k!}{L^k}T,\qquad
 \sum_{j=0}^J|b_j|
 \ge\frac{k!}{L^k}T.}
 \tag{V77.14}
\]
If all radii obey \(a_j\le A_k\), then
\[
 \boxed{
 \mathfrak C_k(W;R)
 \ge
 \frac{k!}{L^k(R+A_k)^k}T.}
 \tag{V77.15}
\]
Consequently, a family with \(A_k=o(k)\) has superexponentially large
\(\mathfrak C_k(W_k;R)\) after normalization \(T=1\).
\end{theorem}

\begin{proof}
From \textup{(V77.6)}, after discarding \(e^{-du}\le1\),
\[
 T\le
 \frac{\sup_{[0,L]}|W|}{(k-1)!}
 \int_0^L u^{k-1}du
 =\frac{L^k}{k!}\sup_{[0,L]}|W|.
 \tag{V77.16}
\]
Since \(a_j,u\ge0\),
\[
 |W(u)|\le\sum_j|b_j|e^{-a_ju}
 \le\sum_j|b_j|.
 \tag{V77.17}
\]
This proves \textup{(V77.14)}.  If \(a_j\le A_k\), then
\[
 \mathfrak C_k(W;R)
 \ge(R+A_k)^{-k}\sum_j|b_j|,
 \tag{V77.18}
\]
which gives \textup{(V77.15)}.  Stirling's formula finishes the last
assertion.
\end{proof}

Thus any bounded-root-cost construction must send at least some radial
locations to order \(k\) or beyond.  The V76 Weierstrass approximation cannot
remain simultaneously low degree, bounded radius, and stably normalized.

\subsection{Interpretation as a holomorphic-remainder barrier}

Suppose \(G\) is holomorphic at all shifted centers and one controls it only
through normalized Cauchy bounds of the form
\[
 \frac{|G^{(k-1)}(s+a_j)|}{(k-1)!}
 \le M(R+a_j)^{-k}.
 \tag{V77.19}
\]
Then
\[
 \left|
 \sum_jb_j\frac{G^{(k-1)}(s+a_j)}{(k-1)!}
 \right|
 \le M\mathfrak C_k(W;R).
 \tag{V77.20}
\]
After normalizing the target-zero moment to \(T=1\), Theorem
\ref{thm:v77-stability-lower} shows that the right-side coefficient budget
cannot have root rate below one.  Therefore absolute Cauchy control alone
cannot yield an exponential reserve separating the holomorphic remainder
from the target signal.

This conclusion is informational.  A lower bound on the available absolute
budget does not assert that the actual gamma and power-character remainders
attain it.  It proves that triangle inequalities and individual Cauchy
estimates cannot certify the required cancellation.

\begin{corollary}[No subcritical certificate in the V76 cost]
\label{cor:v77-no-subcritical}
Let \(W_k\) be any sequence satisfying
\[
 W_k\ge0\text{ on }[L_f,\infty),\qquad
 \mathcal L_d^{(k)}[W_k]=0,\qquad
 \mathcal L_x^{(k)}[W_k]=1.
 \tag{V77.21}
\]
Then for every fixed \(R>d\),
\[
 \liminf_{k\to\infty}
 \mathfrak C_k(W_k;R)^{1/k}\ge1.
 \tag{V77.22}
\]
In particular, the counterexample option proposed in the V76 assessment
with subcritical cost does not exist.
\end{corollary}

\subsection{A dual-cone reading}

Inequality \textup{(V77.9)} is the missing quantitative separator, but it
separates in the adverse direction.  On the feasible slice
\(\mathcal L_d^{(k)}[W]=0\), positivity above \(L\) makes evaluation at
\(x>d\) a bounded functional of the coefficient variation norm, with norm at
most \(e^{(R-d)L}\).  Equivalently, its dual certificate is the truncated
gamma kernel on the unsampled interval \([0,L]\).  The same interval that
made qualitative feasibility possible forces the stability debit.

\begin{theorem}[V77 quantitative-conditioning verdict]
\label{thm:v77-verdict}
The V76 stability problem has the following exact resolution.
\begin{enumerate}[label=\textup{(\roman*)},leftmargin=3em]
\item Every positive target moment is controlled entirely by the sub-prime
      interval through \textup{(V77.6)}.
\item The target-to-Cauchy-cost ratio is uniformly bounded by
      \(e^{(R-d)L_f}\), independently of moment order and radial complexity.
\item At the critical-line holomorphy distance, the normalized cost is at
      least \(\ell_f^{-1/2}\).
\item The ordinary coefficient norm grows at least factorially unless the
      target normalization is allowed to vanish.
\item Keeping all radial shifts \(o(k)\) makes the weighted cost
      superexponential; sending shifts farther does not evade the universal
      root-rate lower bound one.
\end{enumerate}
Thus finite multi-radius feasibility supplies no exponentially stable
certificate against the holomorphic background.
\end{theorem}

\begin{firewall}[A stability no-go is not an actual masking theorem]
Theorem~\ref{thm:v77-stability-lower} proves that an absolute
coefficient-by-coefficient remainder estimate cannot be smaller than the
target on the exponential scale.  It does not prove that the actual
arithmetic remainder has the adverse phase or saturates this budget.
Exploiting cancellation inside that fixed remainder remains logically open.
GRH remains open.
\end{firewall}

\begin{assessment}[V77 closure and V78 direction]
V77 converts V76's qualitative support-gap construction into a sharp
quantitative obstruction: every normalized feasible certificate has
Cauchy-cost root rate at least one and factorial unweighted conditioning.
This closes all multi-radius arguments that estimate gamma and
power-character corrections separately in absolute value.

V78 should now insert the radial weight \(W\) into the exact polygamma
expansion rather than apply Cauchy bounds.  For \(r=k-1\ge1\),
\[
 \frac{(-1)^r}{r!}A_\chi^{(r)}(s)
 =-\sum_{n\ge0}(s+\kappa+2n)^{-k}.
 \tag{V77.23}
\]
The radial combination is therefore
\(-\sum_{n\ge0}\mathcal L_{s+\kappa+2n}^{(k)}[W]\).
The next decision is whether the sign-change geometry forced below \(L_f\)
determines the sign of this exact archimedean response, or whether it can
still be tuned independently of the target moment.  Only a signed result can
improve on the absolute-cost barrier proved here.
\end{assessment}

\section*{Version V77 append-only record}
\addcontentsline{toc}{section}{Version V77 append-only record}
The complete V76 source was retained before this continuation.  It had
1099072 bytes, 32473 lines, and SHA--256
\[
 \texttt{FEC73A42C1E5D4832B3A8D0F6CCD07E307C813C4EA79238DA21BBB667CD9B52B}.
\]
V77 proved a universal target-to-Cauchy-cost inequality for every feasible
multi-radius weight, obtained the exact lower bound
\(\ell_f^{-1/2}\) at the critical-line distance, and established factorial
sub-prime amplitude and coefficient amplification.

% =============================================================================
% END APPENDED V77 MATERIAL
% =============================================================================
% =============================================================================
% BEGIN APPENDED V78 MATERIAL
% =============================================================================

\section{V78: exact polygamma response and the loss of sign beyond minimal radial complexity}
\label{sec:grh-v78}

V77 proved that every feasible multi-radius certificate has critical
Cauchy-cost root rate at least one.  It left open the possibility that the
fixed archimedean remainder might have a favorable signed response, invisible
to absolute estimates.  This version evaluates that response exactly.

For a minimal feasible weight with three radial terms, total positivity does
give a sign at a real ordinate: every archimedean moment lies beyond the
target moment and is positive before the logarithmic-derivative minus sign.
Thus the polygamma contribution reinforces a real off-line-zero witness.
That sign is not stable under general finite radial refinement.  Additional
sub-prime oscillations can tune the archimedean functional to either sign
while preserving principal cancellation, target normalization, and strict
Euler positivity.  At a nonzero ordinate the exact kernel is oscillatory
already before such tuning.

\subsection{The exact radial polygamma formula}

Let
\[
 A_\chi(s):=\frac12\log\frac f\pi
 +\frac12\psi\!\left(\frac{s+\kappa}{2}\right).
 \tag{V78.1}
\]
Fix \(r\ge1\), put \(k=r+1\ge2\), and let
\[
 W(u)=\sum_{j=0}^Jb_je^{-a_ju}
 \tag{V78.2}
\]
be a finite real radial weight, where equal radii have been combined and
\(0\le a_0<\cdots<a_J\), with every \(b_j\ne0\).  Define its normalized archimedean response
at \(s\) by
\[
 \mathfrak A_{\chi,k}^{W}(s)
 :=\sum_{j=0}^Jb_j
 \frac{(-1)^r}{r!}A_\chi^{(r)}(s+a_j).
 \tag{V78.3}
\]

\begin{theorem}[Exact polygamma mixture]
\label{thm:v78-polygamma-mixture}
For \(\Re s>0\),
\[
 \boxed{
 \mathfrak A_{\chi,k}^{W}(s)
 =-\sum_{n\ge0}
 \mathcal L_{s+\kappa+2n}^{(k)}[W]}
 \tag{V78.4}
\]
and, equivalently,
\[
 \boxed{
 \mathfrak A_{\chi,k}^{W}(s)
 =-\frac1{(k-1)!}\int_0^\infty
 \frac{u^{k-1}e^{-(s+\kappa)u}}
      {1-e^{-2u}}W(u)\,du.}
 \tag{V78.5}
\]
Both expressions converge absolutely.
\end{theorem}

\begin{proof}
For \(r\ge1\), the polygamma series is
\[
 \psi^{(r)}(z)
 =(-1)^{r+1}r!\sum_{n\ge0}(n+z)^{-(r+1)}.
 \tag{V78.6}
\]
The chain-rule factor \(2^{-(r+1)}\) in
\(A_\chi^{(r)}\) cancels the factor \(2^{r+1}\) obtained from
\(n+(s+\kappa+a_j)/2\).  Therefore
\[
 \frac{(-1)^r}{r!}A_\chi^{(r)}(s+a_j)
 =-\sum_{n\ge0}(s+\kappa+a_j+2n)^{-k}.
 \tag{V78.7}
\]
Summing in \(j\) gives \textup{(V78.4)}.  Insert the gamma integral for
each denominator and sum
\(\sum_{n\ge0}e^{-2nu}=(1-e^{-2u})^{-1}\) to obtain
\textup{(V78.5)}.  Near zero the integrand is \(O(u^{k-2})\), and at
infinity it decays exponentially.
\end{proof}

The conductor constant disappears after one derivative.  The exact
archimedean problem is therefore a continuous Laplace functional of the same
radial weight \(W\) whose negative mass is hidden below the first Euler
frequency.

\subsection{A sign theorem at minimal radial complexity}

Assume now that \(s=\sigma>1\) is real and that a hypothetical real
off-line zero has \(1/2<\beta<1\).  Put
\[
 d:=\sigma-1,\qquad x:=\sigma-\beta,
 \tag{V78.8}
\]
so \(0<d<x<d+1/2\).  Suppose \(W\) is feasible in the V76 sense:
\[
 W\ge0\text{ on }[L_f,\infty),\qquad
 \mathcal L_d^{(k)}[W]=0,\qquad
 \mathcal L_x^{(k)}[W]>0.
 \tag{V78.9}
\]

\begin{lemma}[Variation-diminishing continuation of the target sign]
\label{lem:v78-minimal-transform-sign}
If \(W\) has exactly three nonzero exponential terms, then
\[
 \boxed{
 \mathcal L_y^{(k)}[W]>0
 \qquad(y\ge x).}
 \tag{V78.10}
\]
\end{lemma}

\begin{proof}
The three exponentials form a strict Chebyshev system, so \(W\) has at most
two real zeros counted with multiplicity.  Proposition
\ref{prop:v76-two-sign-changes} forces two genuine sign changes below
\(L_f\).  Since \(W\) is eventually positive, its sign pattern is
positive--negative--positive; in particular it is positive immediately to
the right of zero.

The Stieltjes kernels \((y+a_j)^{-k}\), as functions of \(y>0\) with
distinct \(a_j\), also form a strict Chebyshev system.  Hence
\[
 F(y):=\mathcal L_y^{(k)}[W]
 =\sum_jb_j(y+a_j)^{-k}
 \tag{V78.11}
\]
has at most two positive zeros counted with multiplicity.  It has a zero at
\(d\), is positive at \(x\), and is positive for all sufficiently large
\(y\) by the Laplace localization at \(u=0\) and the initial positive sign of
\(W\).  A zero at or beyond \(x\) would either be an even zero or would
require a second later crossing to recover the positive limiting sign.
Together with the zero at \(d\), this would give at least three zeros counted
with multiplicity.  Therefore no such zero exists.
\end{proof}

\begin{theorem}[Minimal three-radius archimedean sign]
\label{thm:v78-minimal-sign}
Under \textup{(V78.8)}--\textup{(V78.9)}, if \(W\) has exactly three
nonzero radial terms, then
\[
 \boxed{
 \mathfrak A_{\chi,k}^{W}(\sigma)<0.}
 \tag{V78.12}
\]
Thus at a real off-line zero the exact archimedean contribution has the same
negative sign as the target-zero term
\(-m_\rho\mathcal L_x^{(k)}[W]\).
\end{theorem}

\begin{proof}
Every parameter in \textup{(V78.4)} satisfies
\[
 \sigma+\kappa+2n>x,
 \tag{V78.13}
\]
because \(\sigma-x=\beta>0\).  Lemma
\ref{lem:v78-minimal-transform-sign} makes every term
\(\mathcal L_{\sigma+\kappa+2n}^{(k)}[W]\) strictly positive.  The minus sign
in \textup{(V78.4)} proves the claim.
\end{proof}

This is a genuine signed improvement over the V77 Cauchy budget, but it is
confined to real ordinates and minimal radial complexity.

\subsection{Why the sign is not stable under radial refinement}

For real \(\sigma\), remove the harmless positive common factor
\(u^{k-1}/(k-1)!\) and consider on \((0,L_f)\) the three kernels
\[
 e^{-du},\qquad e^{-xu},\qquad
 \frac{e^{-(\sigma+\kappa)u}}{1-e^{-2u}}.
 \tag{V78.14}
\]
They are linearly independent.  Indeed, the third has an infinite
exponential expansion and a \(u^{-1}\) singularity at zero, whereas a linear
combination of the first two does not.

\begin{lemma}[A sub-prime archimedean null-direction]
\label{lem:v78-arch-null-direction}
There exists a real continuous function
\(h\), compactly supported in \((0,L_f)\), such that
\[
 \mathcal L_d^{(k)}[h]=0,\qquad
 \mathcal L_x^{(k)}[h]=0,
 \tag{V78.15}
\]
but
\[
 \frac1{(k-1)!}\int_0^\infty
 \frac{u^{k-1}e^{-(\sigma+\kappa)u}}
      {1-e^{-2u}}h(u)\,du\ne0.
 \tag{V78.16}
\]
\end{lemma}

\begin{proof}
The three continuous linear functionals induced by the kernels in
\textup{(V78.14)} are linearly independent on
\(C_c(0,L_f)\).  If the third vanished on the intersection of the kernels
of the first two, elementary linear algebra would place it in their span,
contrary to that independence.  Choose \(h\) in the intersection on which
the third is nonzero.
\end{proof}

\begin{theorem}[Archimedean sign tunability at unrestricted finite complexity]
\label{thm:v78-sign-tunability}
Fix \(k,d,x,L_f,\sigma,\kappa\) as above.  There exist finite real
exponential polynomials \(W_+\) and \(W_-\) satisfying
\[
 W_\pm(u)>0\quad(u\ge L_f),\qquad
 \mathcal L_d^{(k)}[W_\pm]=0,\qquad
 \mathcal L_x^{(k)}[W_\pm]>0,
 \tag{V78.17}
\]
for which
\[
 \boxed{
 \mathfrak A_{\chi,k}^{W_+}(\sigma)>0,\qquad
 \mathfrak A_{\chi,k}^{W_-}(\sigma)<0.}
 \tag{V78.18}
\]
Thus no sign of the exact polygamma response follows from the unrestricted
V76 feasibility constraints.
\end{theorem}

\begin{proof}
Start with the strict continuous feasible function \(w\) in the proof of
Theorem~\ref{thm:v76-feasibility}.  Add \(th\), where \(h\) is supplied by
Lemma~\ref{lem:v78-arch-null-direction}.  For every real \(t\), the principal
and target moments remain fixed and the tail on \([L_f,\infty)\) is
unchanged, while the archimedean functional is a nonconstant affine function
of \(t\).  Choose two finite values of \(t\) giving opposite strict signs.

For each chosen function, pass to \(q=e^{-u}\) and approximate uniformly on
\([0,1]\) by a real polynomial as in V76.  Correct the two small moment
errors by adding
\[
 \alpha+\beta e^{-u}.
 \tag{V78.19}
\]
The correction is unique and small because
\[
 \det
 \begin{pmatrix}
 d^{-k}&(d+1)^{-k}\\
 x^{-k}&(x+1)^{-k}
 \end{pmatrix}\ne0.
 \tag{V78.20}
\]
Sufficiently accurate approximation preserves strict tail positivity,
strict target positivity, and the selected archimedean sign, while making
the principal moment exactly zero.  Expanding the approximating polynomial
in \(e^{-u}\) produces the finite exponential polynomials \(W_\pm\).
\end{proof}

The additional sign freedom necessarily uses more than the two hidden sign
changes of the minimal case.  It is another form of the V76 sub-prime
capacity, now detected by the exact gamma distribution rather than an
absolute norm.

\subsection{Nonzero ordinates are oscillatory}

At \(s=\sigma+i\gamma\), equation \textup{(V78.5)} gives
\[
 \begin{aligned}
 \Re\mathfrak A_{\chi,k}^{W}(\sigma+i\gamma)
 &=-\frac1{(k-1)!}\int_0^\infty
 \frac{u^{k-1}e^{-(\sigma+\kappa)u}\cos(\gamma u)}
      {1-e^{-2u}}W(u)\,du,\\
 \Im\mathfrak A_{\chi,k}^{W}(\sigma+i\gamma)
 &=\frac1{(k-1)!}\int_0^\infty
 \frac{u^{k-1}e^{-(\sigma+\kappa)u}\sin(\gamma u)}
      {1-e^{-2u}}W(u)\,du.
 \end{aligned}
 \tag{V78.21}
\]
Even when \(W\) has minimal sign complexity, the factors
\(\cos(\gamma u)\) and \(\sin(\gamma u)\) are not ordered Laplace kernels.
The real-axis variation-diminishing argument therefore gives no sign at the
ordinate of a general complex zero.

\begin{theorem}[V78 signed-polygamma verdict]
\label{thm:v78-verdict}
The exact archimedean screen has the following outcome.
\begin{enumerate}[label=\textup{(\roman*)},leftmargin=3em]
\item The complete radial polygamma response is the signed Laplace mixture
      \textup{(V78.4)}--\textup{(V78.5)}.
\item For a minimal three-radius feasible weight and a real target zero,
      strict total positivity forces the archimedean response to reinforce
      the negative zero contribution.
\item Additional sub-prime oscillations create null-directions for the
      principal and target moments along which the archimedean sign varies
      freely.
\item At a nonzero ordinate the exact kernel is oscillatory and has no
      order sign even before radial refinement.
\item Hence the V77 absolute-cost barrier can be beaten in sign only in the
      minimal real-ordinate subcase, not uniformly for the GRH problem.
\end{enumerate}
\end{theorem}

\begin{firewall}[A favorable real archimedean sign is not a zero-free theorem]
Theorem~\ref{thm:v78-minimal-sign} controls one fixed remainder for a real
zero and a three-term weight.  It does not control the other zero channels,
and its order argument disappears at nonzero ordinate or higher radial
complexity.  The tuning theorem proves that the unrestricted feasibility
class retains both archimedean signs.  GRH remains open.
\end{firewall}

\begin{assessment}[V78 closure and V79 direction]
V78 replaces the Cauchy estimate by the exact polygamma distribution.  It
finds a favorable sign in the smallest real case but proves that no general
archimedean sign survives the full multi-radius cone.  This closes the
polygamma-only repair of V77.

V79 should perform the remaining exact signed calculation for the
power-character correction
\(\mathcal H_\chi=-P_\chi^{(\ge2)}\).  After radial combination its
normalized derivative is a higher-prime-power Euler sum weighted by
\(-W(m\log p)\).  The first test is whether the V74 nonnegative
prime-phase combination makes this entire correction nonpositive, thereby
removing it from the masking background, or whether character-power mixing
reintroduces an uncontrolled sign.  V80 will then perform the next mandatory
corpus sweep before any further branch is chosen.
\end{assessment}

\section*{Version V78 append-only record}
\addcontentsline{toc}{section}{Version V78 append-only record}
The complete V77 source was retained before this continuation.  It had
1109795 bytes, 32808 lines, and SHA--256
\[
 \texttt{6F8D4944AD8A5C75AAE59A9B325887B39F27F3CD63043812825E12FF3DDED8C1}.
\]
V78 derived the exact radial polygamma mixture, proved its favorable sign
for minimal three-radius real-zero certificates, and established both
sub-prime sign tunability at unrestricted complexity and the oscillatory
loss of order at nonzero ordinates.

% =============================================================================
% END APPENDED V78 MATERIAL
% =============================================================================
% =============================================================================
% BEGIN APPENDED V79 MATERIAL
% =============================================================================

\section{V79: the power correction is an exact negative prime-power tail}
\label{sec:grh-v79}

V78 evaluated the gamma contribution exactly and found that its favorable
sign is confined to minimal real-ordinate radial certificates.  The other
holomorphic term in the first-prime explicit formula is the power-character
correction \(\mathcal H_\eta\) introduced in V73.  Its sign can be settled
without an estimate.

The M\"obius definition of \(\mathcal H_\eta\) collapses to minus the full
repeated-prime tail.  Consequently, after any V76-admissible radial weighting,
the V74 phase-positive character combination makes the complete correction
nonpositive term by term.  Character powers do not reintroduce an adverse
sign.  In fact the correction subtracts the \(m\ge2\) part of the full
logarithmic derivative and leaves exactly the nonnegative \(m=1\) prime
layer.  Thus the power correction is removed from the masking background;
the remaining uncontrolled terms are archimedean and zero-spectral.

\subsection{M\"obius collapse of the correction}

Let \(\eta\) be a Dirichlet character modulo \(f\), with the convention that
Euler factors at primes dividing \(f\) are omitted.  Recall the V73
definition
\[
 \mathcal H_\eta(s)
 :=-\sum_{n\ge2}\mu(n)\frac{L'}{L}(ns,\eta^n),
 \qquad \Re s>\frac12.
 \tag{V79.1}
\]

\begin{theorem}[Exact repeated-prime formula]
\label{thm:v79-mobius-collapse}
On \(\Re s>1/2\), locally normally with all derivatives,
\[
 \boxed{
 \mathcal H_\eta(s)
 =-\sum_{p\nmid f}\sum_{q\ge2}
   (\log p)\eta(p)^q p^{-qs}.}
 \tag{V79.2}
\]
Equivalently, \(\mathcal H_\eta\) is minus the higher-prime-power part of
\(-L'/L(s,\eta)\).
\end{theorem}

\begin{proof}
On every compact subset of \(\Re s>1/2\), the logarithmic derivatives in
\textup{(V79.1)} have absolutely convergent Euler series, uniformly after
the outer sum is grouped by the total exponent.  Hence
\[
 \begin{aligned}
 \mathcal H_\eta(s)
 &=\sum_{n\ge2}\mu(n)
   \sum_{p\nmid f}\sum_{m\ge1}
   (\log p)\eta(p)^{nm}p^{-nms}\\
 &=\sum_{p\nmid f}\sum_{q\ge2}
   (\log p)\eta(p)^q p^{-qs}
   \sum_{\substack{n\mid q\\n\ge2}}\mu(n).
 \end{aligned}
 \tag{V79.3}
\]
For \(q\ge2\),
\[
 \sum_{\substack{n\mid q\\n\ge2}}\mu(n)
 =-\mu(1)=-1.
 \tag{V79.4}
\]
This proves \textup{(V79.2)}.  On \(\Re s\ge1/2+\delta\), the terms with
\(q\ge2\) are dominated, after any fixed number of derivatives, by a
convergent series of the form
\(\sum_p (\log p)^A p^{-1-2\delta}\) plus a geometrically decreasing
\(q\)-tail.  This also proves the stated local normal convergence.
\end{proof}

This identity reconciles the two descriptions used earlier: the triangular
power-character sum in V73 and the higher Euler powers in V71 are exactly
the same holomorphic correction.

\subsection{Radial differentiation preserves the negative tail}

Fix \(r\ge1\), put \(k=r+1\), and let
\[
 W(u)=\sum_{j=0}^Jb_je^{-a_ju},
 \qquad 0\le a_0<\cdots<a_J,
 \tag{V79.5}
\]
be a finite real exponential polynomial.  Define
\[
 \mathfrak H_{\eta,k}^{W}(s)
 :=\sum_{j=0}^Jb_j\frac{(-1)^r}{r!}
 \mathcal H_\eta^{(r)}(s+a_j).
 \tag{V79.6}
\]

\begin{corollary}[Exact radially weighted correction]
\label{cor:v79-radial-correction}
For \(\Re s>1/2\),
\[
 \boxed{
 \mathfrak H_{\eta,k}^{W}(s)
 =-\frac1{r!}\sum_{p\nmid f}\sum_{q\ge2}
  \frac{(\log p)(q\log p)^r\eta(p)^q}{p^{qs}}
  W(q\log p).}
 \tag{V79.7}
\]
The series is locally normally convergent.
\end{corollary}

\begin{proof}
Differentiate \textup{(V79.2)} termwise.  The two factors
\((-1)^r\), one from the normalization and one from differentiating
\(p^{-qs}\), cancel, while the initial minus sign remains.  Summing over
the radii gives
\(\sum_jb_jp^{-qa_j}=W(q\log p)\).  Local normal convergence follows from
the majorant in the proof of Theorem~\ref{thm:v79-mobius-collapse}, since
the finite radial shifts have nonnegative real parts.
\end{proof}

\subsection{Exact sign under the V74 phase cone}

Let
\[
 T(\theta)=c_0+2\Re\sum_{\ell=1}^Kc_\ell e^{i\ell\theta}
 \tag{V79.8}
\]
be a real trigonometric polynomial satisfying \(T(\theta)\ge0\) for every
real \(\theta\).  Let \(\chi\) be a character modulo \(f\), fix
\(\sigma>1\) and \(\tau\in\mathbb R\), and put
\[
 \theta_p:=\arg\chi(p)-\tau\log p
 \qquad(p\nmid f).
 \tag{V79.9}
\]
For the zeroth character power use the principal character
\(\mathbf1_f\).

\begin{theorem}[Phase-positive power-correction sign]
\label{thm:v79-phase-sign}
Suppose
\[
 W(q\log p)\ge0
 \qquad(p\nmid f,\ q\ge2).
 \tag{V79.10}
\]
Then
\[
 \boxed{
 \begin{aligned}
 \mathfrak S^{H}_{T,W}
 &:=
 c_0\mathfrak H_{\mathbf1_f,k}^{W}(\sigma)
 +2\Re\sum_{\ell=1}^Kc_\ell
   \mathfrak H_{\chi^\ell,k}^{W}(\sigma+i\ell\tau)\\
 &=-\frac1{r!}\sum_{p\nmid f}\sum_{q\ge2}
  \frac{(\log p)(q\log p)^r}{p^{q\sigma}}
  W(q\log p)T(q\theta_p)
 \le0.
 \end{aligned}}
 \tag{V79.11}
\]
In particular, every V76-admissible weight, which is positive on
\([L_f,\infty)\), satisfies the hypothesis.
\end{theorem}

\begin{proof}
For \(p\nmid f\),
\[
 (\chi(p)^\ell)^q p^{-iq\ell\tau}
 =e^{i\ell q\theta_p}.
 \tag{V79.12}
\]
Insert \textup{(V79.7)} in the character combination and combine conjugate
Fourier coefficients.  The bracket is precisely \(T(q\theta_p)\), giving
the equality in \textup{(V79.11)}.  Every factor after the displayed minus
sign is nonnegative, which proves the inequality.  Finally,
\(q\log p\ge2L_f>L_f\) for every supported term, so V76 tail positivity is
more than sufficient.
\end{proof}

There is no hidden character-power phase defect: the same integer \(q\)
simultaneously dilates the character phase and the vertical shift, producing
exactly the phase \(q\theta_p\) already controlled by \(T\ge0\).

\subsection{The signed sieve sandwich}

For \(\Re s>1\), define the radially weighted first-prime response
\[
 \mathfrak U_{\eta,k}^{W}(s)
 :=\sum_{j=0}^Jb_j\frac{(-1)^r}{r!}
 \mathcal P_\eta^{(r)}(s+a_j),
 \qquad
 \mathcal P_\eta(s):=\sum_{p\nmid f}(\log p)\eta(p)p^{-s}.
 \tag{V79.13}
\]
Absolute expansion gives
\[
 \mathfrak U_{\eta,k}^{W}(s)
 =\frac1{r!}\sum_{p\nmid f}
  \frac{(\log p)^{r+1}\eta(p)}{p^s}W(\log p).
 \tag{V79.14}
\]
Retain the full logarithmic-derivative radial response
\(\mathcal D_{\eta,r}^{W}\) from V76.

\begin{theorem}[Exact first-prime sieve and sandwich]
\label{thm:v79-sieve-sandwich}
If \(W(q\log p)\ge0\) for every \(q\ge1\) and \(p\nmid f\), then
\[
 \begin{aligned}
 \mathfrak S^{(1)}_{T,W}
 &:=
 c_0\mathfrak U_{\mathbf1_f,k}^{W}(\sigma)
 +2\Re\sum_{\ell=1}^Kc_\ell
   \mathfrak U_{\chi^\ell,k}^{W}(\sigma+i\ell\tau)\\
 &=\frac1{r!}\sum_{p\nmid f}
  \frac{(\log p)^{r+1}}{p^\sigma}W(\log p)T(\theta_p)
 \ge0.
 \end{aligned}
 \tag{V79.15}
\]
Moreover, if \(\mathfrak S^{\mathrm{all}}_{T,W}\) denotes the corresponding V76
phase-combination of the full responses, then
\[
 \boxed{
 \mathfrak S^{(1)}_{T,W}=\mathfrak S^{\mathrm{all}}_{T,W}+\mathfrak S^{H}_{T,W},
 \qquad
 0\le\mathfrak S^{(1)}_{T,W}\le\mathfrak S^{\mathrm{all}}_{T,W}.}
 \tag{V79.16}
\]
Thus \(\mathcal H\) removes exactly the nonnegative \(q\ge2\) portion of
the full Euler certificate and leaves its nonnegative \(q=1\) layer.
\end{theorem}

\begin{proof}
Equation \textup{(V79.15)} follows from \textup{(V79.14)} exactly as in
the proof of Theorem~\ref{thm:v79-phase-sign}.  The Euler expansion
\textup{(V76.23)} splits into its \(q=1\) and \(q\ge2\) parts.  Corollary
\ref{cor:v79-radial-correction} is the negative of the latter, proving the
identity in \textup{(V79.16)}.  The two inequalities follow term by term
from \(T\ge0\) and the assumed arithmetic positivity of \(W\).
\end{proof}

For a V76 feasible weight, all actual prime-power samples are strictly
positive, while its principal Laurent moment vanishes.  On the zero side of
the first-prime explicit formula
\[
 \mathcal P_\eta(s)
 =-\frac{\Lambda'}{\Lambda}(s,\eta)
  +A_\eta(s)+\mathcal H_\eta(s),
 \tag{V79.17}
\]
an aligned target zero contributes
\(-m_\rho\mathcal L_x^{(k)}[W]<0\).  Theorem
\ref{thm:v79-phase-sign} shows that the combined power correction is also
nonpositive.  It therefore cannot be the positive term that masks the target
inside the nonnegative prime-side quantity \(\mathfrak S^{(1)}_{T,W}\).

This does not yield a contradiction.  The phase combination of
\textup{(V79.17)} still contains the gamma responses and every non-target
zero in the principal and power-character channels.  V78 proves that the
gamma combination has no uniform sign at nonzero ordinates, and no theorem
in the present ledger controls the aggregate sign of the remaining zero
channels after radial compensation.

\begin{theorem}[V79 power-correction verdict]
\label{thm:v79-verdict}
The exact signed calculation has the following consequences.
\begin{enumerate}[label=\textup{(\roman*)},leftmargin=3em]
\item The triangular V73 correction is identically minus the repeated-prime
      Euler tail.
\item Radial differentiation inserts the same multiplier
      \(W(q\log p)\) as in the full Euler formula and preserves the initial
      minus sign.
\item Every V74 nonnegative phase polynomial makes the complete correction
      nonpositive for a V76-admissible weight; character-power mixing causes
      no loss of sign.
\item Adding this correction to the full logarithmic derivative performs an
      exact signed sieve: it deletes all \(q\ge2\) terms and leaves a
      nonnegative first-prime certificate.
\item Hence the power correction cannot positively mask a negative target
      zero.  Any compensating positive contribution must instead come from
      the archimedean or non-target-zero terms.
\end{enumerate}
\end{theorem}

\begin{firewall}[A favorable correction sign is not zero separation]
Theorem~\ref{thm:v79-phase-sign} disposes of one holomorphic remainder, but
it does not assign a sign to the full zero sum or to the oscillatory gamma
response.  The nonnegative first-prime identity can still be balanced by
those terms.  No zero has been excluded, and GRH remains open.
\end{firewall}

\begin{assessment}[V79 closure and the mandatory V80 sweep]
V79 closes the exact signed power-correction question in the favorable
direction.  Unlike the polygamma term, its sign survives arbitrary finite
radial complexity because it is supported only at sampled prime-power
frequencies, where V76 positivity applies.

V80 is the next mandatory ten-version corpus sweep.  It should inventory all
notes added or changed since the V70 manifest and compare their strongest
new mechanisms with the now-reduced obstruction.  The sweep should search
especially for a rigorous way to separate one closest zero block across the
principal and power-character channels, or for an identity coupling that
zero block to the oscillatory gamma response.  If no such mechanism exists,
the next attack should move upstream from radial certificate design to the
joint zero-channel geometry rather than optimize \(W\) further.
\end{assessment}

\section*{Version V79 append-only record}
\addcontentsline{toc}{section}{Version V79 append-only record}
The complete V78 source was retained before this continuation.  It had
1122362 bytes, 33156 lines, and SHA--256
\[
 \texttt{6042850440EDF38F8BAFDA896BD693459AF79A504F67B442973207D691C620A1}.
\]
V79 identified the power correction with the exact negative repeated-prime
tail, proved its nonpositive phase-combined radial sign, and established the
signed sieve sandwich leaving precisely the nonnegative first-prime layer.

% =============================================================================
% END APPENDED V79 MATERIAL
% =============================================================================
% =============================================================================
% BEGIN APPENDED V80 MATERIAL
% =============================================================================

\section{V80: mandatory corpus sweep and the two-channel compensation reset}
\label{sec:grh-v80}

V79 proved that the power-character correction is not an uncontrolled
positive remainder: under the V74 phase cone it is the exact negative
repeated-prime tail.  This version performs the scheduled ten-version corpus
audit before choosing a successor branch.  The audit finds no transferable
zero-separation theorem, but one new companion warning is directly relevant:
convergence of a positive Gram operator does not determine the orientation of
the underlying response amplitude.

After the audit, the phase polynomial can be reduced to
\(1+\cos\theta\).  This keeps Euler positivity while using only the
principal and target-character channels.  Combining V76 and V79 then gives
an exact compensation law: if an off-line zero exists, the remaining
archimedean and two-channel spectral terms must contribute positively by at
least the full radial target moment.  This isolates the next obstruction
without importing higher character powers.

\subsection{Hash-gated V70--V80 corpus audit}

The stable V70 manifest contained \(31\) non-GRH \texttt{.tex} files and had
SHA--256
\[
 \texttt{56B7195509E847D2F16E08DF3580056CD5EB5C88D73A28D2D7DCA7B885E84AC6}.
 \tag{V80.1}
\]
As in the earlier audits, the present canonical rows are
\[
 \texttt{filename|byte-count|SHA--256}\backslash\texttt{n},
 \tag{V80.2}
\]
sorted by filename, joined with LF terminators, and hashed as UTF--8.

The two complete passes had the following UTC endpoints:
\begin{center}
\small
pass 1 start: \texttt{2026-09-11T05:47:12.3822408Z}\\
pass 1 end: \texttt{2026-09-11T05:47:12.4701987Z}\\
pass 2 start: \texttt{2026-09-11T05:47:12.4739579Z}\\
pass 2 end: \texttt{2026-09-11T05:47:12.5194727Z}.
\end{center}
Each pass rechecked byte counts and modification
ticks after hashing, found \(31\) stable non-GRH files, and produced
\[
 \boxed{
 \texttt{600CBF5481F3BC63CF12C87FEF596C2BFDD57F92F6DD1E820E7D6CEA8D544BCF}.}
 \tag{V80.3}
\]
Two earlier pairs of internally stable passes were discarded when the SM
writer advanced its endpoint before the content audit was complete.  No
conclusion below uses those superseded snapshots.  The SM writer advanced again
after the cutoff in \\textup{(V80.3)}; that later successor is outside this
frozen audit and is reserved for the next scheduled sweep.

Relative to V70, the final stable manifest changes only the following two
rolling endpoints:
\begin{center}
\tiny
\begin{tabular}{p{0.50\textwidth}|r|p{0.32\textwidth}}
\toprule
file & bytes & SHA--256\\
\midrule
\path{Renewal_SM_Einstein_Continuum_Research_Notes_V53.tex}
&544386&\texttt{84160EA8FAA35C64BA824BB5C641FABD}\newline
\texttt{F703AD2CFE41BE40069B794FC499F337}\\
\path{Renewal_Yang_Mills_Mass_Gap_Research_Notes_V54.tex}
&467684&\texttt{20F0DA83E09A4609C73A76F9F198569}\newline
\texttt{21AD4802715506933C3ADE489D9FC9569}\\
\bottomrule
\end{tabular}
\end{center}
Replacing these rows by the documented V70 endpoints
\[
\begin{array}{l|r}
\texttt{Renewal\_SM\_Einstein\_Continuum\_Research\_Notes\_V32.tex}
 &300957\\
\texttt{Renewal\_Yang\_Mills\_Mass\_Gap\_Research\_Notes\_V48.tex}
 &419467
\end{array}
\tag{V80.4}
\]
and their recorded hashes reproduces the digest in \textup{(V80.1)}
exactly.  Thus all other \(29\) rows are byte-identical, and the only new
material requiring review is SM V33--V53 and Yang--Mills V49--V54.

\subsection{Typed audit of SM V33--V53}

The SM successor is a coherent conditional construction, but its objects and
hypotheses must not be confused with the GRH zero-side form.
\begin{enumerate}[label=\textup{(\roman*)},leftmargin=3em]
\item V33--V36 prove projective local-state compactness, quantitative
      marginal compatibility, and finite or countable assembly of positive
      density owners, conditional on asymptotic compatibility and summable
      state-distance debits.
\item V37--V44 replace false support cancellation by conditional Gibbs
      cancellation, derive buffered Dobrushin transmission, split soft and
      hard carriers, and compile connected shell response from a strict
      contraction reserve, exponential moments, and reconstruction tails.
\item V45--V50 audit the physical owner ledger and derive metric
      counterterm, collar-trace, Sobolev reconstruction, and finite-rank
      trace-norm loaders.  Their remaining inputs are physical coefficient
      tails and aligned carrier localization.
\item V51 identifies the augmented response Gram and proves that
      trace-norm convergence of \(G_n^*G_n\) does not orient \(G_n\).
      Recovering a trace-class reconstructed response requires aligned
      response columns and a second Hilbert--Schmidt Sobolev factor.
\item V52 derives an exact entropy cost for positive Gibbs insertions and
      loads it from a uniform interpolated Poincar\'e inequality.  That
      inequality, response orientation, and the Sobolev half-factor remain
      unproved even in the source programme.
\item V53 proves a two-scale Poincar\'e theorem and derives the soft marginal
      constant from an effective Hessian only under the strict Schur margin
      \(a-C_hb^2>0\).  Its source explicitly leaves positivity of the full
      law, the Hessian floor, and that margin unproved.
\end{enumerate}

The V51 orientation counterexample is the most relevant item.  A positive
Gram or covariance object may retain norms while discarding a moving partial
isometry.  In the GRH problem the missing datum is likewise the signed
direction of a specific zero response, not its trace or absolute size.
Consequently the SM convergence theorems cannot supply that direction.
The V52--V53 entropy and two-scale identities live on positive probability
laws and spend Poincar\'e and strict Schur reserves; the meromorphic zero sum
is neither such a law nor equipped with either reserve.

\subsection{Typed audit of Yang--Mills V49--V54}

The Yang--Mills successor constructs a mixed-sign seed cube and enlarges it
through exact rational charts to two quarter-face directions.  V54 performs
an exact intersection ledger, reuses four old charts, and gives rational
negative-vector witnesses against a selected \(90\times90\) sufficient
factorization at four new centres.  The witnesses remain after one
overlap-compatible anisotropic radius repair.

The source correctly does not infer negativity of the underlying
\(45\times45\) Wilson pencil from failure of that sufficient factorization.
It moves upstream to the splitting or subdivision.  This is a useful
methodological match to the present problem: failure of one radial or
phase-positive certificate is not a theorem about the underlying
zero-channel form.  The finite Wilson matrices, however, contain no
Dirichlet coefficient identity, completed functional equation, or
zero-spectrum comparison, so none of their chart estimates transfers.

\begin{theorem}[V80 companion-corpus verdict]
\label{thm:v80-corpus-verdict}
No result added to the companion corpus after the V70 snapshot supplies a
hypothesis-closing estimate for the V79 obstruction.  In particular:
\begin{enumerate}[label=\textup{(\roman*)},leftmargin=3em]
\item positive-state compactness and owner assembly do not assign a sign to
      a meromorphic zero sum;
\item Dobrushin, Poincar\'e, and strict-reserve conclusions cannot be used
      without proving their analogues for the fixed arithmetic response;
\item trace or Gram convergence loses precisely the response orientation
      needed for a target-zero inequality;
\item finite rational chart positivity has no transport law to the
      unbounded zero ordinates;
\item failure of a sufficient factorization does not prove failure of the
      underlying arithmetic form.
\end{enumerate}
The transferable content is therefore diagnostic: retain the signed
coefficient direction and move upstream when a chosen certificate
factorization stalls.
\end{theorem}

\begin{proof}
Each item follows by matching the explicit hypotheses in the successor
theorems to the V79 identity.  The SM results assume positive densities,
conditional kernels, contraction or Poincar\'e reserves, and aligned
operator tails; none is established for the completed Dirichlet zero
response.  V51 gives an explicit rank-one family with fixed Gram and moving
output orientation, so a Gram conclusion alone cannot imply the required
signed response.  The Yang--Mills results concern finitely many specified
rational matrices and explicitly disclaim a conclusion about even their
underlying Wilson pencil when the sufficient block fails.  They therefore
cannot imply a global sign for the distinct arithmetic kernel.
\end{proof}

\subsection{The V71--V79 nonduplication ledger}

The nine versions since the preceding audit have the following exact
outcomes.
\begin{center}
\small
\begin{tabular}{c p{0.79\linewidth}}
\toprule
version & rigorous outcome\\
\midrule
V71 & corrected the conjugation convention and derived the exact
directional gamma/prime shift correlation; only the first-prime layer reaches
the continuation barrier;\\
V72 & proved boundary convergence by residue class and a PNT-derived
Gevrey-two estimate, which is insufficient for the GRH Taylor radius;\\
V73 & performed exponent-space M\"obius inversion, proved correction
holomorphy on \(\Re s>1/2\), and made the prime-moment root bound two
equivalent to paired GRH;\\
V74 & derived the exact zero power sums and finite-block visibility, while
showing that every fixed nonnegative phase polynomial imports a closer
principal pole;\\
V75 & proved the exhaustive two-radius pole-compensation sign split and the
wrong target orientation in the high-frequency-positive choice;\\
V76 & proved qualitative feasibility with at least three radial terms by
hiding two sign changes below the first prime frequency;\\
V77 & proved the critical Cauchy-cost and factorial coefficient lower bounds
for every such feasible radial certificate;\\
V78 & evaluated the polygamma response exactly, obtaining a favorable sign
only for minimal real-ordinate weights and sign tunability in the general
finite cone;\\
V79 & identified the power correction with the exact negative repeated-prime
tail and proved the signed first-prime sieve sandwich.\\
\bottomrule
\end{tabular}
\end{center}
This ledger rules out another optimization of an absolute remainder bound or
another attempt to infer the zero-side sign solely from Euler positivity.

\subsection{A degree-one two-channel reduction}

The general V74 polynomial is no longer needed to settle the power
correction.  Take
\[
 T_\star(\theta):=1+\cos\theta,
 \qquad c_0=1,\quad c_1=\frac12,\quad c_\ell=0\ (\ell\ge2).
 \tag{V80.5}
\]
This polynomial is nonnegative and uses only the principal character and
\(\chi\).

Fix a zero
\(\rho=\beta+i\gamma\) of \(L(s,\chi)\), with
\(1/2<\beta<1\), and choose \(\sigma>1\).  Put
\[
 d:=\sigma-1,\qquad x:=\sigma-\beta.
 \tag{V80.6}
\]
Let \(r\ge1\), \(k=r+1\), and let \(W\) be a finite real exponential
polynomial satisfying
\[
 W(q\log p)\ge0\quad(p\nmid f,\ q\ge1),\qquad
 \mathcal L_d^{(k)}[W]=0,\qquad
 \mathcal L_x^{(k)}[W]>0.
 \tag{V80.7}
\]
Such weights exist by V76.

Set
\[
 \mathfrak S_{\star}^{(1)}
 :=\mathfrak U_{\mathbf1_f,k}^{W}(\sigma)
   +\Re\mathfrak U_{\chi,k}^{W}(\sigma+i\gamma),
 \tag{V80.8}
\]
and likewise
\[
 \mathfrak H_\star
 :=\mathfrak H_{\mathbf1_f,k}^{W}(\sigma)
   +\Re\mathfrak H_{\chi,k}^{W}(\sigma+i\gamma).
 \tag{V80.9}
\]
V79 gives
\[
 \mathfrak S_{\star}^{(1)}\ge0,
 \qquad
 \mathfrak H_\star\le0.
 \tag{V80.10}
\]

In the completed explicit formula for \textup{(V80.8)}, extract the
principal Laurent term and the contribution of the selected zero \(\rho\).
The first vanishes by \textup{(V80.7)}, while the second is
\[
 -m_\rho\mathcal L_x^{(k)}[W].
 \tag{V80.11}
\]
Let \(\mathfrak R_{\star,k}^{W}(\rho)\) denote the real sum of every
remaining completed contribution: both archimedean responses, all other
zeros in the principal and \(\chi\) channels, and any harmless completed
principal endpoint term under the convention of V74.  Thus this is a
defined, fixed arithmetic remainder, not an absolute majorant.

\begin{theorem}[Two-channel forced-compensation law]
\label{thm:v80-forced-compensation}
Under \textup{(V80.5)}--\textup{(V80.7)}, the exact first-prime formula is
\[
 \boxed{
 \mathfrak S_{\star}^{(1)}
 =-m_\rho\mathcal L_x^{(k)}[W]
  +\mathfrak R_{\star,k}^{W}(\rho)+\mathfrak H_\star.}
 \tag{V80.12}
\]
Consequently every hypothetical off-line zero forces
\[
 \boxed{
 \mathfrak R_{\star,k}^{W}(\rho)
 \ge m_\rho\mathcal L_x^{(k)}[W]-\mathfrak H_\star
 \ge m_\rho\mathcal L_x^{(k)}[W]>0.}
 \tag{V80.13}
\]
No higher character-power zero channel occurs in this inequality.
\end{theorem}

\begin{proof}
Apply the V79 phase identity with \(T=T_\star\) and
\(\tau=\gamma\).  Its character combination is exactly
\textup{(V80.8)}, and its power-correction combination is
\textup{(V80.9)}.  Apply the V74 completed zero expansion at
\(\sigma\) in the principal channel and at \(\sigma+i\gamma\) in the
\(\chi\)-channel, then apply the same radial coefficients.  Principal-pole
cancellation is \(\mathcal L_d^{(k)}[W]=0\).  Since
\((\sigma+i\gamma)-\rho=x\), the selected-zero term is
\textup{(V80.11)}; all other terms are precisely the definition of
\(\mathfrak R_{\star,k}^{W}(\rho)\).  This proves
\textup{(V80.12)}.  Now use
\(\mathfrak S_{\star}^{(1)}\ge0\) and
\(\mathfrak H_\star\le0\) from \textup{(V80.10)} and rearrange.
\end{proof}

The theorem turns the residual difficulty into a falsifiable arithmetic
claim.  To exclude \(\rho\), it is enough to construct one admissible
\((k,W)\) for which the fixed two-channel remainder is strictly smaller than
the first quantity on the right of \textup{(V80.13)}.  Conversely, if the
zero exists, every admissible radial certificate must be met by a positive
response of that size from the remaining principal/character zeros and
gamma terms.

\begin{theorem}[V80 audit-and-reset verdict]
\label{thm:v80-verdict}
The mandatory sweep and the V71--V79 ledger imply:
\begin{enumerate}[label=\textup{(\roman*)},leftmargin=3em]
\item no new companion theorem supplies the missing signed arithmetic
      orientation;
\item the higher prime powers have a favorable exact sign and require no
      further character-power optimization;
\item the polynomial \(1+\cos\theta\) reduces the live explicit formula to
      the principal and target-character channels;
\item an off-line zero forces the exact positive compensation
      \textup{(V80.13)} from the remaining two-channel spectrum and gamma
      response;
\item the next nonduplicative task is to test this fixed compensation
      directly, not to optimize generic positivity, traces, or absolute
      costs.
\end{enumerate}
\end{theorem}

\begin{firewall}[Forced compensation is not a contradiction]
Equation \textup{(V80.13)} is a necessary consequence of a hypothetical
off-line zero and Euler positivity.  The present work does not prove an
upper bound contradicting it.  In particular, V78 permits the oscillatory
gamma response to have either sign, and V74 controls the magnitude but not
the favorable alignment of a closest zero block.  GRH remains open.
\end{firewall}

\begin{assessment}[V80 closure and V81 direction]
V80 completes a stable, continuity-verified ten-version sweep.  The new
companion material reinforces the need for orientation-sensitive estimates
but supplies none with arithmetic hypotheses.  The exact degree-one reset
removes all higher character-power zero channels and shows that the
remaining obstruction is a positive compensation law involving only the
principal spectrum, the rest of the \(\chi\)-spectrum, and gamma.

V81 should decompose \(\mathfrak R_{\star,k}^{W}(\rho)\) into its nearest
two-channel zero block and its strictly more distant tail.  It must combine
the V74 Vandermonde block theorem with the radial transform rather than
repeat the one-order Cauchy bound of V77.  The decisive question is whether
one can choose an admissible sequence \(W_k\) and a bounded block of orders
so that the selected off-line block has a uniformly negative orientation
while the principal/gamma and distant-zero terms are subcritical.  If the
nearest block can realize both orientations under all such choices, the
programme should move further upstream to a genuine cross-order
polarization identity.
\end{assessment}

\section*{Version V80 append-only record}
\addcontentsline{toc}{section}{Version V80 append-only record}
The complete V79 source was retained before this continuation.  It had
1134129 bytes, 33481 lines, and SHA--256
\[
 \texttt{4CADA0DF38C4844D6AAA5ADE1FCD2285DB31EFCA67C5FA65DDC83E634967EC44}.
\]
V80 performed the mandatory stable corpus sweep, proved typed nontransfer of
the new SM and Yang--Mills mechanisms, recorded the V71--V79 closure ledger,
and derived the degree-one two-channel forced-compensation law.

% =============================================================================
% END APPENDED V80 MATERIAL
% =============================================================================
% =============================================================================
% BEGIN APPENDED V81 MATERIAL
% =============================================================================

\section{V81: finite spectral interpolation and escape beyond every zero block}
\label{sec:grh-v81}

V80 reduced the phase-positive explicit formula to the principal and
target-character channels and proved that a hypothetical off-line zero
forces a positive compensation from the remaining zeros and gamma terms.
The proposed next step was to isolate the nearest two-channel zero block and
combine its V74 Vandermonde visibility with radial weighting.

There is an exact obstruction to that plan.  The sub-prime freedom found in
V76 is not limited to two scalar moments.  Finite real exponential
polynomials which are strictly positive at every arithmetic frequency can
interpolate an arbitrary finite conjugation-compatible table of complex
Laplace moments.  The same weight can do this for any fixed finite block of
consecutive derivative orders.  It can therefore cancel every undesired zero
in a prescribed nearest block, cancel finitely many gamma nodes, cancel the
principal pole throughout the order block, and retain a normalized target
zero throughout that block.

Thus the V74 Vandermonde theorem remains correct at one radial point, but it
does not survive unrestricted finite radial preprocessing.  Under the V80
forced-compensation law, the required positive response simply escapes
beyond every prescribed finite spectral window.  The obstruction is now
quantitative and infinite-window: the interpolating coefficients become
ill-conditioned, and the uncancelled distant tail can absorb that cost.

\subsection{A finite Laplace-moment interpolation theorem}

Write
\[
 \mathbb H_R:=\{y\in\mathbb C:\Re y>0\}
 \tag{V81.1}
\]
and, for \(n\ge0\), define
\[
 \Phi_{y,n}[w]
 :=\frac1{n!}\int_0^\infty u^ne^{-yu}w(u)\,du.
 \tag{V81.2}
\]
Thus \(\Phi_{y,n}=\mathcal L_y^{(n+1)}\) in the notation of V76.

Let \(\mathcal I\subset\mathbb H_R\times\mathbb N_0\) be finite and closed
under \((y,n)\mapsto(\overline y,n)\).  A table
\((c_{y,n})_{(y,n)\in\mathcal I}\) is called compatible when
\[
 c_{\overline y,n}=\overline{c_{y,n}},
 \qquad
 c_{y,n}\in\mathbb R\quad(y\in\mathbb R).
 \tag{V81.3}
\]

\begin{lemma}[Independence of finite exponential monomials]
\label{lem:v81-exponential-independence}
For distinct pairs \((y,n)\), the functions
\[
 u^ne^{-yu}
 \tag{V81.4}
\]
are linearly independent on every nonempty open interval.  Consequently the
real moment map
\[
 C_c(0,L;\mathbb R)\longrightarrow
 \{(c_{y,n})_{\mathcal I}:\textup{\(c\) satisfies \textup{(V81.3)}}\},
 \qquad
 h\longmapsto(\Phi_{y,n}[h])_{\mathcal I},
 \tag{V81.5}
\]
is onto for every \(L>0\).
\end{lemma}

\begin{proof}
A finite linear combination of \textup{(V81.4)} is an exponential
polynomial.  Group terms with the same exponent \(y\).  Applying, for each
other exponent \(z\), a sufficiently high power of
\(D+z\), where \(D=d/du\), leaves one polynomial times \(e^{-yu}\).
It cannot vanish on an interval unless that polynomial is zero.  Iterating
over the exponents proves complex linear independence.  Pairing conjugate
terms gives the corresponding real linear independence under
\textup{(V81.3)}.

If the map in \textup{(V81.5)} were not onto, a nonzero real covector on its
finite-dimensional target would annihilate its range.  The associated real
linear combination of the kernels in \textup{(V81.4)} would then integrate
to zero against every real \(h\in C_c(0,L)\), so it would vanish on
\((0,L)\).  This contradicts the independence just proved.
\end{proof}

\begin{theorem}[Tail-positive finite spectral interpolation]
\label{thm:v81-tail-positive-interpolation}
Fix \(L>0\), a finite conjugation-closed set \(\mathcal I\) as above, and
arbitrary compatible data \(c=(c_{y,n})_{\mathcal I}\).  There is a real
polynomial \(P\) such that
\[
 W(u):=P(e^{-u})
 \tag{V81.6}
\]
satisfies
\[
 \boxed{
 W(u)>0\quad(u\ge L),\qquad
 \Phi_{y,n}[W]=c_{y,n}
 \quad((y,n)\in\mathcal I).}
 \tag{V81.7}
\]
In particular \(W\) is a finite real exponential polynomial with
integer-spaced nonnegative radii.
\end{theorem}

\begin{proof}
Choose \(\varepsilon>0\).  By Lemma
\ref{lem:v81-exponential-independence}, there is
\(h\in C_c(0,L;\mathbb R)\) such that
\[
 \Phi_{y,n}[h]
 =c_{y,n}-\Phi_{y,n}[\varepsilon]
 \quad((y,n)\in\mathcal I).
 \tag{V81.8}
\]
Then \(w:=\varepsilon+h\) has the required moment table and equals
\(\varepsilon\) on \([L,\infty)\).

Compactify \([0,\infty)\) by \(q=e^{-u}\).  The function
\[
 g(q):=w(-\log q),\qquad g(0):=\varepsilon,
 \tag{V81.9}
\]
is real and continuous on \([0,1]\).  Real polynomials are uniformly dense
there.  Each functional \(\Phi_{y,n}\) is continuous in the uniform norm,
because
\[
 |\Phi_{y,n}[v]|
 \le \|v\|_\infty(\Re y)^{-(n+1)}.
 \tag{V81.10}
\]

It remains to make the moments exact after approximation.  Let
\(\mathcal E_{\mathcal I}\) be the finite-dimensional real compatible-data
space in \textup{(V81.5)}.  The images under \(\Phi\) of the functions
\(Q(e^{-u})\), with \(Q\) a real polynomial, span
\(\mathcal E_{\mathcal I}\).  Indeed, their uniform closure contains every
continuous function on the compactified half-line; if their images lay in a
proper subspace, a nonzero covector would annihilate all continuous
functions, contradicting Lemma
\ref{lem:v81-exponential-independence}.  Select finitely many polynomials
\(Q_1,\ldots,Q_M\) whose moment vectors form a basis of
\(\mathcal E_{\mathcal I}\).

Choose polynomial approximants \(P_\nu(e^{-u})\to w(u)\) uniformly.  Their
moment errors
\[
 e_\nu:=c-\Phi[P_\nu(e^{-u})]
 \tag{V81.11}
\]
tend to zero.  Expand \(e_\nu\) in the fixed basis
\(\Phi[Q_j(e^{-u})]\), and let \(R_\nu(e^{-u})\) be the corresponding
linear combination of the \(Q_j(e^{-u})\).  Finite-dimensional norm
equivalence gives
\[
 \|R_\nu(e^{-u})\|_\infty\longrightarrow0,
 \qquad
 \Phi[R_\nu(e^{-u})]=e_\nu.
 \tag{V81.12}
\]
For all sufficiently large \(\nu\),
\[
 W:=P_\nu(e^{-u})+R_\nu(e^{-u})
 \tag{V81.13}
\]
has the exact table \(c\) and is uniformly within \(\varepsilon/2\) of
\(w\).  Hence \(W(u)>\varepsilon/2\) for \(u\ge L\).  Combining equal powers
of \(e^{-u}\) gives \textup{(V81.6)}.
\end{proof}

The theorem is stronger than V76's three-constraint feasibility result.
It is also distinct from V67: V67 perturbed an arbitrary analytic function
while preserving a finite jet table, whereas \textup{(V81.7)} constructs an
actual radial Euler multiplier, preserves positivity on the whole arithmetic
half-line, and prescribes its Laplace responses.

\subsection{Simultaneous cancellation on a block of orders}

\begin{corollary}[Bounded-order spectral block interpolation]
\label{cor:v81-order-block}
Let \(0<d<x\), let
\(\mathcal Y\subset\mathbb H_R\setminus\{d,x\}\) be finite and closed under
conjugation, and fix integers \(k\ge1\) and \(B\ge1\).  For every \(L>0\)
there is a finite real exponential polynomial \(W\), strictly positive on
\([L,\infty)\), such that for every \(0\le\ell<B\),
\[
 \boxed{
 \mathcal L_d^{(k+\ell)}[W]=0,\qquad
 \mathcal L_x^{(k+\ell)}[W]=1,\qquad
 \mathcal L_y^{(k+\ell)}[W]=0\quad(y\in\mathcal Y).}
 \tag{V81.14}
\]
\end{corollary}

\begin{proof}
Apply Theorem~\ref{thm:v81-tail-positive-interpolation} to the pairs
\[
 (d,k+\ell-1),\quad(x,k+\ell-1),\quad
 (y,k+\ell-1)
 \tag{V81.15}
\]
for \(y\in\mathcal Y\) and \(0\le\ell<B\), prescribing respectively the
compatible values \(0,1,0\).
\end{proof}

This is precisely the interaction missed by applying V74 before V76.
The Vandermonde matrix in V74 prevents a finite closest-zero power sum from
vanishing throughout a block of consecutive raw orders.  Corollary
\ref{cor:v81-order-block} adds radial shifts and chooses their coefficients
to interpolate the transformed powers simultaneously; it can annihilate
the unwanted block while leaving the target transform equal to one.

\subsection{Application to the two-channel zero window}

Fix the setting of V80 and a selected off-line zero
\(\rho=\beta+i\gamma\) of \(L(s,\chi)\).  Let
\(\mathcal Z_0\) be any finite set of nontrivial zeros in the principal
channel, and let \(\mathcal Z_1\) be any finite set of zeros of
\(L(s,\chi)\) not containing \(\rho\).  Define
\[
 \begin{aligned}
 \mathcal Y_{\rm zero}
 :=\;&\{\sigma-\rho_0:\rho_0\in\mathcal Z_0\}\\
 &{}\cup
 \{\sigma+i\gamma-\rho_1:\rho_1\in\mathcal Z_1\},
 \end{aligned}
 \tag{V81.16}
\]
and adjoin the complex conjugate of every member.  Repeated entries are
combined.  Every member has positive real part, and none equals
\(x=\sigma-\beta\) after any accidental collision with the target response
is removed from the \emph{undesired} list.

For a finite gamma cutoff \(M\), also adjoin
\[
 \begin{aligned}
 \mathcal Y_{\rm gam}(M)
 :=\;&\{\sigma+\kappa_0+2m:0\le m<M\}\\
 &{}\cup
 \{\sigma+i\gamma+\kappa_\chi+2m:0\le m<M\},
 \end{aligned}
 \tag{V81.17}
\]
and their conjugates.  These are the nodes in the exact V78 polygamma
mixtures for the principal and target-character channels.

\begin{theorem}[Finite-window escape of the V80 compensation]
\label{thm:v81-finite-window-escape}
Fix \(B\ge1\), finite zero sets
\(\mathcal Z_0,\mathcal Z_1\), and a finite gamma cutoff \(M\).  There is one
finite real exponential polynomial \(W\), strictly positive at every
prime-power frequency, such that throughout a block of \(B\) consecutive
orders:
\begin{enumerate}[label=\textup{(\roman*)},leftmargin=3em]
\item the principal Laurent moment vanishes;
\item the selected target-zero moment equals one;
\item every zero response from \(\mathcal Z_0\cup\mathcal Z_1\) vanishes;
\item the first \(M\) nodes of both exact gamma mixtures vanish.
\end{enumerate}
For this \(W\), the V80 forced-compensation law implies at every order in
the block that the uncancelled spectral/gamma tail satisfies
\[
 \boxed{
 \mathfrak R_{\star,\mathrm{tail}}
 \ge m_\rho-\mathfrak H_\star
 \ge m_\rho.}
 \tag{V81.18}
\]
Thus no fixed finite nearest-zero and gamma window can be the source of a
contradiction in the unrestricted radial cone.
\end{theorem}

\begin{proof}
Take \(L=L_f\) in Corollary~\ref{cor:v81-order-block} and use
\[
 \mathcal Y
 =\mathcal Y_{\rm zero}\cup\mathcal Y_{\rm gam}(M).
 \tag{V81.19}
\]
The resulting \(W\) is positive on \([L_f,\infty)\), hence at every actual
prime-power frequency.  Equations \textup{(V81.14)},
\textup{(V81.16)}, and \textup{(V81.17)} give items
\textup{(i)}--\textup{(iv)}.  Apply Theorem
\ref{thm:v80-forced-compensation} at each order.  The target transform is
one, the named finite components of
\(\mathfrak R_{\star,k}^{W}(\rho)\) are zero, and V79 gives
\(\mathfrak H_\star\le0\).  What remains is exactly the uncancelled tail,
which proves \textup{(V81.18)}.
\end{proof}

The set \(\mathcal Z_0\cup\mathcal Z_1\) may in particular contain every
zero in the closest two-channel distance block and every zero in any fixed
larger compact window.  Hence increasing the V74 block length by a fixed
amount cannot restore visibility after radial interpolation.

\subsection{The price is conditioning, not finite geometry}

The construction is qualitative.  It does not control the degree or
coefficients of \(W\).  In fact V77 applies separately to every order in
\textup{(V81.14)}.  With target normalization one, it gives
\[
 \sum_j|b_j|
 \ge\frac{(k+\ell)!}{L_f^{\,k+\ell}}
 \qquad(0\le\ell<B),
 \tag{V81.20}
\]
together with the corresponding critical Cauchy-cost lower bound.  The finite-window cancellations can therefore
transfer a factorial coefficient budget into the first unprescribed zeros
or gamma nodes.  Equation \textup{(V81.18)} proves that, under the
hypothetical zero, such a transfer must occur.

\begin{theorem}[V81 radial-block verdict]
\label{thm:v81-verdict}
The nearest-block attack has the following exact outcome.
\begin{enumerate}[label=\textup{(\roman*)},leftmargin=3em]
\item finite exponential monomials give an onto real moment map on every
      sub-prime interval;
\item polynomial approximation plus finite-dimensional correction realizes
      every compatible finite table by an integer-radius exponential
      polynomial positive on the full arithmetic half-line;
\item one radial weight can impose these conditions on any fixed finite
      block of consecutive derivative orders;
\item consequently every prescribed finite two-channel zero block and
      finite gamma prefix can be cancelled while the pole is annihilated and
      the off-line target remains normalized;
\item the V80 compensation is forced to escape to the infinite unprescribed
      tail, paid for by the coefficient conditioning already detected in
      V77.
\end{enumerate}
Thus finite Vandermonde visibility alone cannot close the two-channel
compensation law.
\end{theorem}

\begin{firewall}[Finite interpolation is not spectral control]
Theorem~\ref{thm:v81-finite-window-escape} uses the hypothetical zero
locations to construct a qualitative radial interpolant and gives no
uniform coefficient bound as the order or spectral window grows.  It neither
controls the infinite zero tail nor constructs a predeclared arithmetic
certificate.  The forced tail inequality is conditional on an off-line zero,
not a contradiction.  GRH remains open.
\end{firewall}

\begin{assessment}[V81 closure and V82 direction]
V81 proves that the remaining V80 obstruction is genuinely
infinite-dimensional.  Any finite closest-zero block, any bounded block of
orders, and any finite gamma prefix can be neutralized inside the
Euler-positive radial cone.  This closes a direct transplantation of the V74
Vandermonde block theorem after radial compensation.

V82 should quantify the interpolation cost as the cancelled spectral window
grows.  The correct object is the quotient norm of the target Laplace kernel
modulo the pole, undesired-zero, and gamma kernels on the sub-prime interval,
equivalently the smallest singular value of the associated generalized
Cauchy--Vandermonde moment map with the polynomial-synthesis norm retained.
One must compare that cost with the first omitted zero/gamma tail using
actual zero counting and the exact V78 kernels.  A subcritical uniform
quotient bound would make \textup{(V81.18)} impossible; a matching lower
bound would close the entire finite-window radial route and force a genuinely
global polarization identity.
\end{assessment}

\section*{Version V81 append-only record}
\addcontentsline{toc}{section}{Version V81 append-only record}
The complete V80 source was retained before this continuation.  It had
1151410 bytes, 33871 lines, and SHA--256
\[
 \texttt{35DCD55C8E4024AB220849949B9C47D8CD7F78DDB026B92D163839FD00B25E07}.
\]
V81 proved tail-positive finite spectral interpolation, simultaneous
bounded-order block cancellation, and escape of the V80 forced compensation
beyond every prescribed finite two-channel zero and gamma window.

% =============================================================================
% END APPENDED V81 MATERIAL
% =============================================================================
% =============================================================================
% BEGIN APPENDED V82 MATERIAL
% =============================================================================

\section{V82: generalized Cauchy quotients and the gamma-window cost}
\label{sec:grh-v82}

V81 proved that every fixed finite zero and gamma window can be cancelled by
a tail-positive radial exponential polynomial.  Its construction was
qualitative: neither the number of radial channels nor the polynomial
synthesis norm was controlled.  This continuation quantifies both costs for
the exact gamma ladder.

The first possible estimate, based only on approximation of kernels on the
sub-prime interval, is not valid for the full moment constraints: subtracting
an oscillatory gamma kernel destroys the one-sided sign on the arithmetic
tail.  The correct global object lives instead in coefficient space.  It is
the uniform quotient distance of one Cauchy column from the first \(M\)
gamma columns over all admissible radial shifts.  A partial-fraction identity
gives an explicit dual approximant.  It proves polynomial collapse of this
quotient, a polynomial lower bound for the coefficient norm, and a linear
lower bound for radial channel count.  Classical zero counting then gives
the exact coefficientwise tail envelope, but leaves a two-parameter escape
corridor.  Thus the quantitative audit narrows the obstruction without
turning it into a GRH proof.

\subsection{The rational response and its zero capacity}

Retain the V76 notation.  For \(k\ge1\) and a nonzero finite radial weight
\[
 W(u)=\sum_{j=0}^{J}b_j e^{-a_j u},
 \qquad 0\le a_0<\cdots<a_J,
 \tag{V82.1}
\]
put
\[
 F_{W,k}(y):=\mathcal L_y^{(k)}[W]
 =\sum_{j=0}^{J}\frac{b_j}{(y+a_j)^k}.
 \tag{V82.2}
\]

\begin{lemma}[Radial channel count forced by a cancelled ladder]
\label{lem:v82-rational-zero-capacity}
The rational function \(F_{W,k}\) is nonzero and has at most \(kJ\) finite
zeros, counted with multiplicity.  Consequently, if
\[
 F_{W,k}(z+hm)=0\quad(0\le m<M),
 \qquad F_{W,k}(x)\ne0,
 \tag{V82.3}
\]
where the displayed nodes are distinct and avoid the poles, then
\[
 \boxed{J\ge \left\lceil\frac{M}{k}\right\rceil.}
 \tag{V82.4}
\]
In particular, no fixed finite radial weight can cancel an infinite gamma
ladder while retaining a target response.
\end{lemma}

\begin{proof}
With
\[
 D(y):=\prod_{j=0}^{J}(y+a_j)^k,
\]
the numerator \(D(y)F_{W,k}(y)\) has degree at most \(kJ\).  It is not the
zero polynomial: at the pole \(y=-a_j\), its leading Laurent coefficient is
\(b_j\ne0\).  Hence it has at most \(kJ\) finite zeros counted with
multiplicity.  Conditions \textup{(V82.3)} provide \(M\) distinct zeros and
also rule out the identically zero case, proving \textup{(V82.4)}.  Infinitely
many distinct ladder zeros are therefore impossible for fixed \(J\).
\end{proof}

This is a complexity statement, not a conditioning statement.  The latter
comes from the dual generalized Cauchy quotient.

\subsection{An exact coefficient-space quotient}

Fix
\[
 0<x<\Re z,\qquad h>0,qquad M\ge1,qquad k\ge1.
 \tag{V82.5}
\]
For a finite radius set \(A\subset[0,\infty)\), define
\[
 \delta_{A,M,k}(x,z;h)
 :=\inf_{c\in\mathbb C^M}\max_{a\in A}
 \left|
 \frac1{(a+x)^k}
 -\sum_{m=0}^{M-1}\frac{c_m}{(a+z+hm)^k}
 \right|.
 \tag{V82.6}
\]
The continuous-radius version is
\[
 \delta_{M,k}(x,z;h)
 :=\inf_{c\in\mathbb C^M}\sup_{a\ge0}
 \left|
 \frac1{(a+x)^k}
 -\sum_{m=0}^{M-1}\frac{c_m}{(a+z+hm)^k}
 \right|.
 \tag{V82.7}
\]

\begin{proposition}[Finite-radius quotient duality]
\label{prop:v82-quotient-duality}
Let \(A=\{a_0,\ldots,a_J\}\).  Among complex coefficient vectors
\(b=(b_j)\) satisfying
\[
 \sum_j\frac{b_j}{(x+a_j)^k}=1,
 \qquad
 \sum_j\frac{b_j}{(z+hm+a_j)^k}=0
 \quad(0\le m<M),
 \tag{V82.8}
\]
the infimum of \(\sum_j|b_j|\) is
\[
 \boxed{\delta_{A,M,k}(x,z;h)^{-1},}
 \tag{V82.9}
\]
with the convention \(0^{-1}=+\infty\).  Every real coefficient vector
satisfying the same complex moment equations therefore obeys
\[
 \sum_j|b_j|\ge \delta_{M,k}(x,z;h)^{-1}.
 \tag{V82.10}
\]
\end{proposition}

\begin{proof}
In the finite-dimensional space \(\ell^\infty(A)\), let \(v_x\) be the
column \(((x+a)^{-k})_{a\in A}\) and let \(G\) be the span of the \(M\)
columns \(v_{z+hm}\).  The distance from \(v_x\) to \(G\) is exactly
\(\delta_{A,M,k}\).  A coefficient vector in \(\ell^1(A)\) satisfying
\textup{(V82.8)} is a functional which annihilates \(G\) and takes value one
on \(v_x\).  Finite-dimensional Hahn--Banach duality says that the smallest
norm of such a functional is the reciprocal distance in
\textup{(V82.9)}.  Since the supremum in \textup{(V82.7)} is no smaller than
the maximum on \(A\), one has
\(\delta_{A,M,k}\le\delta_{M,k}\), which proves \textup{(V82.10)}.
\end{proof}

Unlike an interval-only Laplace approximation, this quotient tests every
radial shift \(a\ge0\), so no uncontrolled arithmetic-tail integral has been
discarded.

\subsection{A partial-fraction dual certificate}

Set
\[
 A_M:=\prod_{m=0}^{M-1}(z+hm-x),
 \qquad
 R_M(t):=
 \frac{A_M}{(t+x)\prod_{m=0}^{M-1}(t+z+hm)}.
 \tag{V82.11}
\]

\begin{theorem}[Global gamma-ladder quotient bound]
\label{thm:v82-gamma-quotient}
Under \textup{(V82.5)},
\[
 \boxed{
 \delta_{M,k}(x,z;h)
 \le
 \frac1x
 \prod_{m=0}^{M-1}\frac{|z+hm-x|}{|z+hm|}
 \left(
 \frac1x+\sum_{m=0}^{M-1}\frac1{|z+hm|}
 \right)^{k-1}.}
 \tag{V82.12}
\]
Consequently every \(W\) satisfying the normalized target and gamma-prefix
conditions \textup{(V82.8)} has coefficient variation
\[
 \boxed{
 \sum_j|b_j|
 \ge
 x\prod_{m=0}^{M-1}\frac{|z+hm|}{|z+hm-x|}
 \left(
 \frac1x+\sum_{m=0}^{M-1}\frac1{|z+hm|}
 \right)^{-(k-1)}.}
 \tag{V82.13}
\]
For fixed \(x,z,h,k\), there is \(C_{x,z,h,k}>0\) such that
\[
 \boxed{
 \sum_j|b_j|
 \ge C_{x,z,h,k}
 \frac{M^{x/h}}{(1+\log M)^{k-1}}.}
 \tag{V82.14}
\]
\end{theorem}

\begin{proof}
The residue of \(R_M\) at \(t=-x\) is one.  Its partial-fraction expansion
therefore has the form
\[
 R_M(t)=\frac1{t+x}-\sum_{m=0}^{M-1}\frac{c_m}{t+z+hm}
 \tag{V82.15}
\]
for suitable complex \(c_m\).  Differentiating \(k-1\) times gives
\[
 \frac1{(t+x)^k}
 -\sum_{m=0}^{M-1}\frac{c_m}{(t+z+hm)^k}
 =\frac{(-1)^{k-1}}{(k-1)!}R_M^{(k-1)}(t).
 \tag{V82.16}
\]

For numbers \(q_0,\ldots,q_M\), let \(H_r(q_0,\ldots,q_M)\) be the complete
homogeneous polynomial of degree \(r\).  Repeated differentiation of the
product in \textup{(V82.11)} yields
\[
 \frac{|R_M^{(k-1)}(t)|}{(k-1)!}
 \le |R_M(t)|
 H_{k-1}\left(
 \frac1{|t+x|},\frac1{|t+z|},\ldots,
 \frac1{|t+z+h(M-1)|}
 \right).
 \tag{V82.17}
\]
All real parts in the denominators are positive.  Hence, for \(t\ge0\),
every displayed reciprocal and \(|R_M(t)|\) is at most its value at \(t=0\).
Moreover \(H_r(q)\le(\sum q_j)^r\) for nonnegative \(q_j\).  Equations
\textup{(V82.11)}, \textup{(V82.16)}, and \textup{(V82.17)} now give
\textup{(V82.12)}.  Proposition~\ref{prop:v82-quotient-duality} gives
\textup{(V82.13)}.

Finally,
\[
 \prod_{m=0}^{M-1}\frac{z+hm}{z+hm-x}
 =
 \frac{\Gamma(M+z/h)\Gamma((z-x)/h)}
      {\Gamma(z/h)\Gamma(M+(z-x)/h)}.
 \tag{V82.18}
\]
The fixed-parameter gamma-ratio asymptotic has modulus
\[
 \left|\frac{\Gamma((z-x)/h)}{\Gamma(z/h)}\right|
 M^{x/h}(1+O_{x,z,h}(M^{-1})).
 \tag{V82.19}
\]
Also
\[
 \frac1x+\sum_{m=0}^{M-1}\frac1{|z+hm|}
 \le C_{x,z,h}+h^{-1}\log(1+M).
 \tag{V82.20}
\]
Substitution in \textup{(V82.13)} proves \textup{(V82.14)}.
\end{proof}

For \(k=1\), the right side of \textup{(V82.12)} is simply
\(|R_M(0)|\), and the dual quotient collapses at least as
\(M^{-x/h}\).  Higher derivative order pays the additional logarithmic
factor visible in \textup{(V82.12)}.  The theorem is a lower bound for the
physical synthesis norm, not merely for a change-of-basis matrix.

\subsection{Insertion into the V81 two-channel window}

For the target-character gamma ladder in V81, take
\[
 z=\sigma+i\gamma+\kappa_\chi,
 \qquad h=2,
 \qquad x=\sigma-\beta.
 \tag{V82.21}
\]
Then
\[
 \Re z-x=\beta+\kappa_\chi>0.
 \tag{V82.22}
\]

\begin{corollary}[Quantitative cost of the V81 gamma prefix]
\label{cor:v82-v81-cost}
Let \(W\) be a V81 interpolant which, at order \(k\), normalizes the selected
target zero to one and cancels the first \(M\) target-character gamma nodes.
If \(B(W):=\sum_j|b_j|\), then
\[
 \boxed{
 J\ge\left\lceil\frac{M}{k}\right\rceil,
 \qquad
 B(W)\ge
 C_{\rho,\sigma,k}
 \frac{M^{(\sigma-\beta)/2}}{(1+\log M)^{k-1}}.}
 \tag{V82.23}
\]
If the same weight is tail-positive, cancels the principal Laurent moment,
and has \(L=L_f\), then V77 independently gives
\[
 \boxed{B(W)\ge\frac{k!}{L_f^k}.}
 \tag{V82.24}
\]
Thus the valid combined lower bound is the maximum of
\textup{(V82.23)} and \textup{(V82.24)}, not their product.
\end{corollary}

\begin{proof}
Equations \textup{(V82.21)}--\textup{(V82.22)} permit Lemma
\ref{lem:v82-rational-zero-capacity} and Theorem
\ref{thm:v82-gamma-quotient} with \(h=2\).  The V77 factorial estimate is
Theorem~\ref{thm:v77-factorial-cost} with target normalization one.  Since
the two estimates arise from different dual tests, only their maximum is
automatic.
\end{proof}

This is the promised quantitative replacement for the qualitative
finite-dimensional norm equivalence in V81.  It also proves that the degree
of the polynomial \(P(e^{-u})\) cannot remain fixed as the cancelled gamma
window grows.

\subsection{The first omitted gamma and zero tails}

The exact V78 gamma tail admits the following coefficientwise envelope.  Put
\[
 G_{k,M}(A;h):=\sum_{m\ge M}(A+hm)^{-k},
 \qquad A>0.
 \tag{V82.25}
\]
For \(k\ge2\), monotone integral comparison gives
\[
 \frac{(A+hM)^{1-k}}{h(k-1)}
 \le G_{k,M}(A;h)
 \le (A+hM)^{-k}
 +\frac{(A+hM)^{1-k}}{h(k-1)}.
 \tag{V82.26}
\]
Therefore
\[
 \sum_{m\ge M}
 \left|\mathcal L_{z+hm}^{(k)}[W]\right|
 \le B(W)G_{k,M}(\Re z;h).
 \tag{V82.27}
\]

The zero tail has the analogous classical estimate.

\begin{lemma}[Power tail from local Dirichlet zero counting]
\label{lem:v82-zero-power-tail}
Let \(\gamma_\rho=\Im\rho\), with zeros counted with multiplicity.  For
\(k\ge2\), \(T\ge2\), and \(t\in\mathbb R\),
\[
 \boxed{
 \sum_{|\gamma_\rho-t|\ge T}
 \frac{m_\rho}{|\gamma_\rho-t|^k}
 \ll_\chi
 T^{1-k}\log(2+|t|+T).}
 \tag{V82.28}
\]
Consequently, for \(s=\sigma+it\) with \(\sigma>1\),
\[
 \sum_{|\gamma_\rho-t|\ge T}m_\rho
 \left|
 \sum_j\frac{b_j}{(s+a_j-\rho)^k}
 \right|
 \ll_\chi
 B(W)T^{1-k}\log(2+|t|+T).
 \tag{V82.29}
\]
\end{lemma}

\begin{proof}
Partition the ordinates into dyadic shells
\[
 2^qT\le|\gamma_\rho-t|<2^{q+1}T,
 \qquad q\ge0.
\]
The standard local estimate
\[
 N_\chi(U+1)-N_\chi(U)\ll_\chi\log(2+|U|)
 \tag{V82.30}
\]
shows that the \(q\)-th shell contains at most
\[
 \ll_\chi 2^qT\log(2+|t|+2^qT)
\]
zeros with multiplicity.  Its contribution is therefore at most
\[
 \ll_\chi (2^qT)^{1-k}
 \log(2+|t|+2^qT).
\]
For \(k\ge2\), summing the resulting geometric series, including the harmless
linear factor in \(q\) from the logarithm, proves \textup{(V82.28)}.  Finally
\[
 |s+a_j-\rho|\ge|t-\gamma_\rho|
\]
because \(a_j\ge0\) and \(0<\Re\rho<1<\sigma\).  Summing first in \(j\)
proves \textup{(V82.29)}.
\end{proof}

Lemma~\ref{lem:v82-zero-power-tail} is a tail estimate for the actual
Dirichlet zero multiset, not an arbitrary masking configuration.  It reuses
the local count already invoked in V23 and V69, but applies it to the present
high-derivative Cauchy kernel.

\subsection{The remaining two-parameter escape corridor}

Define the coefficientwise gamma and zero budgets
\[
 \mathfrak B_{\rm gam}:=B(W)G_{k,M}(A;h),
 \qquad
 \mathfrak B_{\rm zero}:=
 B(W)T^{1-k}\log(2+|t|+T).
 \tag{V82.31}
\]
They are the right sides available to a termwise absolute-value proof.  The
V77 lower bound \textup{(V82.24)} and the first term of the gamma tail give
\[
 \mathfrak B_{\rm gam}
 \ge \frac{k!}{L_f^k(A+hM)^k}.
 \tag{V82.32}
\]
If \(M/k\to\lambda\in(0,\infty)\), Stirling's formula yields
\[
 \boxed{
 \liminf_{k\to\infty}
 \mathfrak B_{\rm gam}^{1/k}
 \ge\frac1{eL_fh\lambda}.}
 \tag{V82.33}
\]
Likewise, if \(T/k\to\tau\in(0,\infty)\), then
\[
 \boxed{
 \liminf_{k\to\infty}
 \mathfrak B_{\rm zero}^{1/k}
 \ge\frac1{eL_f\tau}.}
 \tag{V82.34}
\]
The logarithm has \(k\)-th root tending to one.

Equations \textup{(V82.33)}--\textup{(V82.34)} exclude a vanishing
coefficientwise envelope when
\[
 \lambda<\frac1{eL_fh}
 \quad\hbox{or}\quad
 \tau<\frac1{eL_f},
 \tag{V82.35}
\]
respectively.  They do not exclude the simultaneous region
\[
 \boxed{
 \lambda>\frac1{eL_fh},
 \qquad
 \tau>\frac1{eL_f}.}
 \tag{V82.36}
\]
The quotient bound \textup{(V82.23)} does not close this region: for fixed
\(k>1+x/h\), its polynomial growth is weaker than the polynomial decay
\(G_{k,M}(A;h)\asymp M^{1-k}\), while for \(M\asymp k\) its logarithmic
factor gives no additional exponential lower rate.

This conclusion must be read in the correct direction.  The lower bounds do
not construct a family whose actual tails vanish in \textup{(V82.36)}.  They
show that the presently proved conditioning obstructions fail to rule such a
family out.  Conversely, a large triangle budget does not prove that the
fixed signed spectral tail attains it.  The only rigorous conclusion is that
coefficientwise absolute estimates close the small-window regimes in
\textup{(V82.35)} and leave the large-window corridor
\textup{(V82.36)} undecided.

\begin{theorem}[V82 quantitative-window verdict]
\label{thm:v82-verdict}
The V81 finite-window obstruction has the following quantitative resolution.
\begin{enumerate}[label=\textup{(\roman*)},leftmargin=3em]
\item Cancelling \(M\) distinct gamma nodes at derivative order \(k\) needs
      at least \(1+\lceil M/k\rceil\) radial terms.
\item The exact generalized Cauchy quotient over all radial shifts has the
      partial-fraction bound \textup{(V82.12)}; hence normalized synthesis
      variation grows at least as
      \(M^{x/h}(1+\log M)^{-(k-1)}\) at fixed order.
\item Tail positivity and principal cancellation independently force the V77
      factorial synthesis cost \(k!/L_f^k\).
\item The exact gamma tail decays as \(M^{1-k}\), and the actual Dirichlet
      zero tail has the classical envelope
      \(T^{1-k}\log(2+|t|+T)\).
\item Combining these estimates excludes undersized gamma or zero windows,
      but leaves the simultaneous large-window corridor
      \textup{(V82.36)}.  Classical zero counting controls cardinality and
      the omitted tail, not the least quotient distance generated by the
      cancelled actual zero nodes.
\end{enumerate}
Thus a contradiction from the V80 compensation law still requires either a
new lower bound for the actual zero-node Cauchy quotient or a signed global
identity which avoids coefficientwise tail estimates.
\end{theorem}

\begin{firewall}[A conditioning theorem is not zero exclusion]
The partial-fraction certificate is a dual lower bound on radial synthesis
cost.  It neither supplies a matching upper bound nor determines the phase of
the uncancelled spectral tail.  The corridor \textup{(V82.36)} is an
inconclusive regime, not a construction and not evidence for an off-line
zero.  No zero has been excluded; GRH remains open.
\end{firewall}

\begin{assessment}[V82 closure and V83 direction]
V82 replaces V81's qualitative finite-dimensional norm equivalence by an
explicit global generalized Cauchy quotient.  It proves unavoidable growth
of both channel count and synthesis variation, and it inserts the actual
Dirichlet zero-counting tail into the comparison.  The resulting bounds do
not close the cofinal large-window regime.

The next attack should not repeat another finite-window inversion.  Lemma
\ref{lem:v82-rational-zero-capacity} shows the upstream rigidity: a fixed
finite rational response cannot vanish on the full gamma ladder.  V83 should
test an all-order or generating-function polarization which preserves that
infinite-ladder uniqueness before truncation.  The target is a predeclared
positive combination of derivative orders whose zero and gamma responses
remain analytic as a whole.  If its tail can be controlled without replacing
it by coefficientwise absolute values, it may close the corridor; if not, a
counterexample in the corresponding Hardy or Stieltjes class would close the
entire radial programme and force a return to the global completed--Euler
metric.
\end{assessment}

\section*{Version V82 append-only record}
\addcontentsline{toc}{section}{Version V82 append-only record}
The complete V81 source was retained before this continuation.  It had
1166500 bytes, 34254 lines, and SHA--256
\[
 \texttt{B47B9A2010B305AE5BC2D41FB1F0568FEAAAF5F33F414A9C262D805D3E4EAF4C}.
\]
V82 proved the rational radial-channel lower bound, an explicit global
gamma-ladder Cauchy-quotient estimate, polynomial synthesis amplification,
the high-derivative zero-tail bound from actual zero counting, and the exact
large-window corridor left open by all current coefficientwise estimates.

% =============================================================================
% END APPENDED V82 MATERIAL
% =============================================================================
% =============================================================================
% BEGIN APPENDED V83 MATERIAL
% =============================================================================

\section{V83: all-order spectral shifts and the radial Borel barrier}
\label{sec:grh-v83}

V82 suggested preserving the rigidity of the infinite gamma ladder by
summing derivative orders before truncation.  There are two inequivalent
ways to do this.  For one fixed radial weight, the ordinary generating
function is exactly a translation of the spectral parameter.  It reaches
the principal pole before it can reach the target zero.  The exponential
generating function retains prime-side positivity through a modified Bessel
kernel, but the same principal term has strictly larger positive-direction
exponential type than the target.

One can evade this fixed-weight obstruction only by changing the radial
weight with the derivative order.  That destroys the translation identity.
It also imports the V77 factorial synthesis cost into the coefficient
family: the ordinary measure-valued series has radius zero, and even its
natural Borel transform has total-variation radius at most the first-prime
gap \(L_f=\log\ell_f\).  Thus an analytic all-order polarization cannot be
obtained from the V81 interpolants by formal summation.

These results concern the radial Euler-positive attack.  They do not repeat
V73's prime-moment Taylor-radius equivalence or V74's closest-zero
Vandermonde theorem.  Here the variables being summed are the normalized
logarithmic-derivative orders after radial preprocessing, and the obstruction
is the physical radial synthesis norm.

\subsection{Ordinary summation is only a spectral shift}

Recall the normalized logarithmic-derivative moment
\[
 \mathcal Q_{\eta,r}(s)
 :=\frac{(-1)^r}{r!}\frac{d^r}{ds^r}
 \left(-\frac{L'}{L}(s,\eta)\right)
 \tag{V83.1}
\]
and, for a fixed finite radial weight
\[
 W(u)=\sum_{j=0}^{J}b_j e^{-a_j u},
 \qquad 0\le a_0<\cdots<a_J,\qquad b_j\ne0,
 \tag{V83.2}
\]
put
\[
 \mathcal D_{\eta,r}^{W}(s)
 :=\sum_{j=0}^{J}b_j\mathcal Q_{\eta,r}(s+a_j).
 \tag{V83.3}
\]

\begin{theorem}[Ordinary all-order shift identity]
\label{thm:v83-ordinary-shift}
If
\[
 |q|<\min_j\{\Re(s+a_j)-1\},
 \tag{V83.4}
\]
then the following series converge absolutely and
\[
 \boxed{
 \sum_{r\ge0}q^r\mathcal D_{\eta,r}^{W}(s)
 =
 \sum_{j=0}^{J}b_j
 \left(-\frac{L'}{L}\right)(s+a_j-q).}
 \tag{V83.5}
\]
On the Euler side this is
\[
 \boxed{
 \sum_{r\ge0}q^r\mathcal D_{\eta,r}^{W}(s)
 =
 \sum_p\sum_{m\ge1}
 \frac{(\log p)\eta(p)^m}{p^{ms}}
 e^{q m\log p}W(m\log p).}
 \tag{V83.6}
\]
For every \(y\in\mathbb H_R\) with
\(|q|<\min_j|y+a_j|\), the corresponding singular response is
\[
 \boxed{
 \sum_{r\ge0}q^r\mathcal L_y^{(r+1)}[W]
 =\sum_{j=0}^{J}\frac{b_j}{y+a_j-q}.}
 \tag{V83.7}
\]
\end{theorem}

\begin{proof}
The Euler expansion gives
\[
 \mathcal Q_{\eta,r}(s+a_j)
 =\frac1{r!}\sum_p\sum_{m\ge1}
 (\log p)(m\log p)^r
 \eta(p)^m p^{-m(s+a_j)}.
 \tag{V83.8}
\]
Under \textup{(V83.4)}, absolute convergence permits summation of
\[
 \sum_{r\ge0}\frac{(q\,m\log p)^r}{r!}
 =e^{q m\log p}.
 \tag{V83.9}
\]
This proves \textup{(V83.5)}--\textup{(V83.6)}.  Likewise
\[
 \sum_{r\ge0}\frac{q^r}{(y+a_j)^{r+1}}
 =\frac1{y+a_j-q},
 \tag{V83.10}
\]
which proves \textup{(V83.7)}.
\end{proof}

For real \(q\ge0\), equation \textup{(V83.6)} preserves every V74
prime-phase positivity identity whenever \(W\) is nonnegative at the Euler
frequencies and all shifted Euler products remain absolutely convergent.
The identity supplies no new spectral variable: it is literally the
translation \(s\mapsto s-q\).

\subsection{The principal singularity is encountered first}

Put, as before,
\[
 d:=\sigma-1,\qquad x:=\sigma-\beta,
 \qquad 0<d<x,
 \tag{V83.11}
\]
and define the two rational responses
\[
 P_d(q):=\sum_j\frac{b_j}{d+a_j-q},
 \qquad
 P_x(q):=\sum_j\frac{b_j}{x+a_j-q}.
 \tag{V83.12}
\]

\begin{theorem}[Fixed-radial all-order principal barrier]
\label{thm:v83-fixed-principal-barrier}
For every nonzero fixed weight \textup{(V83.2)}, the nearest positive
singularity of \(P_d\) is
\[
 q=d+a_0,
 \tag{V83.13}
\]
whereas the nearest positive target singularity is
\[
 q=x+a_0>d+a_0.
 \tag{V83.14}
\]
More precisely,
\[
 P_d(q)=\frac{b_0}{d+a_0-q}+O_W(1)
 \quad(q\uparrow d+a_0),
 \tag{V83.15}
\]
while \(P_x\) is holomorphic there.  No nonzero fixed radial weight can
cancel the principal response as an analytic germ: if
\[
 \mathcal L_d^{(r+1)}[W]=0
 \qquad(r=0,1,2,\ldots),
 \tag{V83.16}
\]
then \(W=0\), and hence every target response also vanishes.
\end{theorem}

\begin{proof}
Distinct radial shifts give distinct simple poles in
\textup{(V83.12)}.  The pole at \(d+a_0\) has nonzero residue
\(-b_0\), so it cannot be cancelled by terms whose poles occur elsewhere.
Since \(d<x\), the target poles are all shifted to the right by \(x-d\).
This proves \textup{(V83.13)}--\textup{(V83.15)}.

If \textup{(V83.16)} holds, the Taylor series at \(q=0\) in
\textup{(V83.7)} with \(y=d\) is identically zero.  Thus \(P_d\) vanishes
on a disk and hence as a rational function.  Its distinct pole residues then
give \(b_j=0\) for every \(j\), contrary to nontriviality.
\end{proof}

The finite-block form gives a second complexity bound complementary to V82.

\begin{corollary}[Every cancelled order consumes a radial channel]
\label{cor:v83-order-channel-count}
Suppose a nonzero weight with \(J+1\) radial terms satisfies
\[
 \mathcal L_d^{(k+\ell)}[W]=0
 \qquad(0\le\ell<B).
 \tag{V83.17}
\]
Then
\[
 \boxed{J\ge B.}
 \tag{V83.18}
\]
Consequently a V81 interpolant cancelling \(M\) gamma nodes throughout a
block of \(B\) orders satisfies
\[
 \boxed{
 J\ge
 \max\left\{B,\left\lceil\frac{M}{k}\right\rceil\right\}.}
 \tag{V83.19}
\]
\end{corollary}

\begin{proof}
Set
\[
 \lambda_j:=(d+a_j)^{-1},
 \qquad c_j:=b_j\lambda_j^k.
\]
The first \(\min\{B,J+1\}\) equations in \textup{(V83.17)} are
\[
 \sum_{j=0}^{J}c_j\lambda_j^\ell=0.
 \tag{V83.20}
\]
If \(B\ge J+1\), the first \(J+1\) rows form a square Vandermonde matrix
with distinct positive nodes \(\lambda_j\), forcing every \(c_j\), and hence
every \(b_j\), to vanish.  Therefore a nonzero weight requires \(B\le J\).
Combine this with Lemma~\ref{lem:v82-rational-zero-capacity} at the first
order \(k\) to obtain \textup{(V83.19)}.
\end{proof}

\subsection{Borel summation retains positivity but not target dominance}

The extra factorial in an exponential generating function changes the
prime-side kernel but preserves its sign.

\begin{theorem}[Fixed-weight Borel identity and exponential-type barrier]
\label{thm:v83-fixed-borel}
For every fixed finite \(W\), every \(y\in\mathbb H_R\), and every
\(q\in\mathbb C\),
\[
 \boxed{
 \sum_{r\ge0}\frac{q^r}{r!}
 \mathcal L_y^{(r+1)}[W]
 =
 \sum_{j=0}^{J}
 \frac{b_j}{y+a_j}
 \exp\left(\frac{q}{y+a_j}\right).}
 \tag{V83.21}
\]
In an absolutely convergent Euler half-plane,
\[
 \boxed{
 \sum_{r\ge0}\frac{q^r}{r!}
 \mathcal D_{\eta,r}^{W}(s)
 =
 \sum_p\sum_{m\ge1}
 \frac{(\log p)\eta(p)^m}{p^{ms}}
 I_0\!\left(2\sqrt{q\,m\log p}\right)
 W(m\log p),}
 \tag{V83.22}
\]
where
\[
 I_0(2\sqrt z)=\sum_{r\ge0}\frac{z^r}{(r!)^2}.
 \tag{V83.23}
\]
Thus the prime side remains termwise nonnegative for \(q\ge0\) in every
phase-positive channel combination.

Nevertheless, for \(0<d<x\),
\[
 \boxed{
 \frac{
 \left|\sum_j\frac{b_j}{d+a_j}
 e^{q/(d+a_j)}\right|}
 {
 \left|\sum_j\frac{b_j}{x+a_j}
 e^{q/(x+a_j)}\right|}
 \sim
 \frac{x+a_0}{d+a_0}
 \exp\left[
 q\left(\frac1{d+a_0}-\frac1{x+a_0}\right)
 \right]
 \longrightarrow\infty}
 \tag{V83.24}
\]
as \(q\to+\infty\).  A fixed-weight positive Borel polarization therefore
amplifies the principal response faster than the target response.
\end{theorem}

\begin{proof}
For each \(j\),
\[
 \sum_{r\ge0}
 \frac{q^r}{r!(y+a_j)^{r+1}}
 =\frac1{y+a_j}\exp\left(\frac q{y+a_j}\right),
\]
proving \textup{(V83.21)}.  Insert \textup{(V83.8)}, multiply by
\(q^r/r!\), and use \textup{(V83.23)} to prove
\textup{(V83.22)}.  Positivity follows from the nonnegative Taylor
coefficients of \(I_0(2\sqrt z)\) for \(z\ge0\).

Because \(a_0<a_j\) for \(j>0\), the \(j=0\) summand is the unique term of
largest positive exponential type in both the numerator and denominator of
\textup{(V83.24)}.  Divide each sum by its \(j=0\) term.  All other terms
tend exponentially to zero, and the displayed asymptotic follows after
cancelling \(|b_0|\).
\end{proof}

This argument is stronger than comparing one derivative order: it shows
that a positive entire summation of all fixed-weight orders still points in
the wrong direction.

\subsection{Order-dependent interpolation has a finite Borel radius}

It remains possible to select a different radial weight at every order.
Let
\[
 W_r(u)=\sum_{j=0}^{J_r}b_{j,r}e^{-a_{j,r}u}
 \tag{V83.25}
\]
satisfy the normalized V76 conditions
\[
 W_r(u)\ge0\quad(u\ge L_f),\qquad
 \mathcal L_d^{(r+1)}[W_r]=0,\qquad
 \mathcal L_x^{(r+1)}[W_r]=1.
 \tag{V83.26}
\]
Encode its radial coefficients as the finite signed measure
\[
 \mu_r:=\sum_{j=0}^{J_r}b_{j,r}\delta_{a_{j,r}}
 \in\mathcal M([0,\infty)),
 \qquad
 \|\mu_r\|_{\rm TV}=\sum_j|b_{j,r}|.
 \tag{V83.27}
\]

\begin{theorem}[Zero ordinary radius and the first-prime Borel wall]
\label{thm:v83-borel-wall}
Every family satisfying \textup{(V83.26)} obeys
\[
 \boxed{
 \|\mu_r\|_{\rm TV}
 \ge\frac{(r+1)!}{L_f^{\,r+1}}.}
 \tag{V83.28}
\]
Consequently the measure-valued ordinary series
\[
 \sum_{r\ge0}q^r\mu_r
 \tag{V83.29}
\]
has radius of convergence zero in total variation.  Its natural Borel
transform
\[
 \mathfrak W(q):=
 \sum_{r\ge0}\frac{q^r}{r!}\mu_r
 \tag{V83.30}
\]
has total-variation radius at most \(L_f\), and cannot converge at any
\(|q|\ge L_f\).  Quantitatively, the absolute synthesis majorant satisfies,
for \(0\le q<L_f\),
\[
 \boxed{
 \sum_{r\ge0}\frac{q^r}{r!}\|\mu_r\|_{\rm TV}
 \ge
 \frac{L_f}{(L_f-q)^2}.}
 \tag{V83.31}
\]
\end{theorem}

\begin{proof}
Equation \textup{(V83.28)} is Theorem
\ref{thm:v77-factorial-cost} with \(k=r+1\), \(L=L_f\), and target
normalization one.  Hence
\[
 \liminf_{r\to\infty}\|\mu_r\|_{\rm TV}^{1/r}
 =+\infty,
\]
so \textup{(V83.29)} has radius zero.  Also
\[
 \left\|\frac{\mu_r}{r!}\right\|_{\rm TV}
 \ge\frac{r+1}{L_f^{\,r+1}}.
 \tag{V83.32}
\]
The Cauchy--Hadamard formula for Banach-valued power series gives a radius
at most \(L_f\).  At \(|q|\ge L_f\), the norm of the \(r\)-th term in
\textup{(V83.30)} is at least
\[
 \frac{r+1}{L_f}\left(\frac{|q|}{L_f}\right)^r,
\]
which does not tend to zero.  Finally,
\[
 \sum_{r\ge0}\frac{q^r}{r!}\|\mu_r\|_{\rm TV}
 \ge\frac1{L_f}\sum_{r\ge0}(r+1)
 \left(\frac q{L_f}\right)^r
 =\frac{L_f}{(L_f-q)^2},
\]
proving \textup{(V83.31)}.
\end{proof}

If the series \textup{(V83.30)} converged, the normalized target responses
would sum to \(e^q\), while the principal response would vanish
coefficientwise.  Theorem~\ref{thm:v83-borel-wall} shows that the physical
coefficient family cannot be continued in total variation beyond the
finite point \(q=L_f\).  In particular it cannot be sent to
\(q\to+\infty\), where Theorem~\ref{thm:v83-fixed-borel} detects
exponential type.

\begin{remark}[Higher-order Borel weights are not yet controlled]
\label{rem:v83-higher-borel}
Dividing the \(r\)-th coefficient by \((r!)^\alpha\) with \(\alpha>1\)
formally weakens the lower bound \textup{(V83.28)} enough that it no longer
forces a finite radius.  This is not a construction: V81 and V82 provide no
matching upper bound on \(\|\mu_r\|_{\rm TV}\), no common radial support, and
no control of the signed zero and gamma tails after such summation.
Moreover the ordinary spectral-shift identity is then lost.  A
higher-order Borel transform is therefore a genuinely new problem, not an
escape already supplied by the finite interpolation theorem.
\end{remark}

\begin{theorem}[V83 all-order radial verdict]
\label{thm:v83-verdict}
The proposed all-order radial repair has the following exact status.
\begin{enumerate}[label=\textup{(\roman*)},leftmargin=3em]
\item For one fixed radial weight, ordinary summation of the normalized
      derivative orders is exactly the spectral shift \(s\mapsto s-q\).
\item The principal singularity is encountered before the target singularity,
      and annihilating its complete Taylor germ forces the radial weight to
      vanish.
\item Cancelling a block of \(B\) pole orders requires at least \(B+1\)
      radial channels; together with V82, a \(B\)-by-\(M\) V81 window needs
      at least \(1+\max\{B,\lceil M/k\rceil\}\) channels.
\item Fixed-weight Borel summation preserves Euler positivity through
      \(I_0\), but its principal response has strictly larger positive
      exponential type than the target response.
\item Allowing the weight to depend on the order avoids fixed rational
      rigidity, but the ordinary coefficient series has radius zero and the
      natural total-variation Borel transform has radius at most \(L_f\).
\end{enumerate}
Thus neither fixed-weight all-order summation nor the natural Borel sum of
the V81 order-dependent interpolants closes the V82 large-window corridor.
\end{theorem}

\begin{firewall}[A summability obstruction is not GRH]
Theorems~\ref{thm:v83-fixed-principal-barrier} and
\ref{thm:v83-borel-wall} exclude two precise implementations of the
all-order radial idea.  They do not exclude a higher-order Borel family, a
non-radial completed-formula polarization, or cancellation using the actual
signed infinite spectral measure.  No upper bound of the V73 prime-moment
root rate has been proved, and GRH remains open.
\end{firewall}

\begin{assessment}[V83 closure and V84 direction]
V83 shows that the most direct upstream repair proposed after V82 is
circular.  Fixed radial data turn the order variable into a mere spectral
translation, while order-dependent pole compensation destroys ordinary
analytic summability and hits the first-prime wall under the natural Borel
normalization.

V84 should test the only summation class not closed here: a
Gevrey-two-compatible positive kernel with more than one factorial in the
order weight.  The required calculation is exact.  One must derive its
principal, target-zero, gamma, and Euler responses before estimating any of
them, then decide whether orderwise pole cancellation and the signed V79
correction survive in a common Banach space.  If the physical synthesis
series still fails to converge, the radial programme is closed.  If it
converges, its exponential-type ordering must be compared with the V72
Gevrey-two background rather than with coefficientwise Cauchy bounds.
\end{assessment}

\section*{Version V83 append-only record}
\addcontentsline{toc}{section}{Version V83 append-only record}
The complete V82 source was retained before this continuation.  It had
1184111 bytes, 34779 lines, and SHA--256
\[
 \texttt{9294B933909F76627A35B3D061170452B645B357C0210E5E13E7F18FDE1DB22D}.
\]
V83 proved the ordinary spectral-shift identity, the fixed-radial
principal-first barrier, the order-block radial-channel lower bound, the
positive Bessel-kernel Borel identity, and the first-prime total-variation
Borel wall for order-dependent pole-compensating interpolants.

% =============================================================================
% END APPENDED V83 MATERIAL
% =============================================================================
% =============================================================================
% BEGIN APPENDED V84 MATERIAL
% =============================================================================

\section{V84: quadratic Borel kernels and the positive-order no-go}
\label{sec:grh-v84}

V83 closed ordinary and first-Borel all-order summation for the direct radial
programme.  It left open a higher-order transform compatible with the
Gevrey-two scale visible in V72.  This continuation evaluates the first such
transform exactly and then proves a summation theorem which applies to every
nonnegative choice of order weights, not only to a particular Borel kernel.

The outcome is again adverse.  For a fixed radial weight, quadratic Borel
summation preserves Euler positivity but the principal pole has strictly
larger square-root exponential type than the target zero.  If the radial
weight is allowed to depend on the order, V77's critical Cauchy-cost lower
bound survives every nonnegative summation term by term.  The quadratic
Borel absolute synthesis majorant grows at least like \(e^{q/L_f}\), whereas
the normalized target grows only like \(e^{2\sqrt q}\).  Finally, the exact
polygamma contribution remains freely tunable at any one positively weighted
order.  Thus higher Borel smoothing does not create either an absolute
reserve or a signed archimedean inequality.

\subsection{The quadratic Borel kernel}

Define the positive entire function
\[
 \mathfrak F_2(z):=
 \sum_{r\ge0}\frac{z^r}{(r!)^3}.
 \tag{V84.1}
\]
The subscript records the two additional factorials in the order weight;
the normalized logarithmic derivative already contains the third factorial
on the Euler side.

\begin{theorem}[Exact quadratic Borel responses]
\label{thm:v84-quadratic-borel}
Let \(W\) be the fixed finite radial weight \textup{(V83.2)}.  For every
\(y\in\mathbb H_R\) and \(q\in\mathbb C\),
\[
 \boxed{
 \sum_{r\ge0}\frac{q^r}{(r!)^2}
 \mathcal L_y^{(r+1)}[W]
 =
 \sum_{j=0}^{J}\frac{b_j}{y+a_j}
 I_0\!\left(2\sqrt{\frac q{y+a_j}}\right).}
 \tag{V84.2}
\]
In every absolutely convergent Euler half-plane,
\[
 \boxed{
 \sum_{r\ge0}\frac{q^r}{(r!)^2}
 \mathcal D_{\eta,r}^{W}(s)
 =
 \sum_p\sum_{m\ge1}
 \frac{(\log p)\eta(p)^m}{p^{ms}}
 \mathfrak F_2(q\,m\log p)W(m\log p).}
 \tag{V84.3}
\]
For \(q\ge0\), the kernel \(\mathfrak F_2(q\,u)\) is positive.  Hence
quadratic Borel summation preserves every V74 phase-positive Euler
combination whenever \(W\) is nonnegative at the Euler frequencies.
\end{theorem}

\begin{proof}
The defining series for \(I_0\) gives
\[
 \sum_{r\ge0}\frac{q^r}{(r!)^2(y+a_j)^{r+1}}
 =\frac1{y+a_j}
 I_0\!\left(2\sqrt{\frac q{y+a_j}}\right),
\]
which proves \textup{(V84.2)} after the finite sum in \(j\).
Equation \textup{(V83.8)} and absolute convergence permit interchange of
the order and Euler sums.  The resulting kernel is
\[
 \sum_{r\ge0}\frac{(q\,m\log p)^r}{(r!)^3}
 =\mathfrak F_2(q\,m\log p),
\]
proving \textup{(V84.3)}.  Its Taylor coefficients are nonnegative.
\end{proof}

The response in \textup{(V84.2)} has an exact type ordering.

\begin{lemma}[The required \(I_0\) asymptotic]
\label{lem:v84-i0-asymptotic}
As \(t\to+\infty\),
\[
 \boxed{
 I_0(t)=\frac{e^t}{\sqrt{2\pi t}}
 \left(1+O(t^{-1})\right).}
 \tag{V84.4}
\]
\end{lemma}

\begin{proof}
Use the integral representation
\[
 I_0(t)=\frac1\pi\int_0^\pi e^{t\cos\theta}\,d\theta.
\]
Split at a fixed small \(\theta_0>0\).  On
\([\theta_0,\pi]\), one has \(\cos\theta\le1-c_{\theta_0}\), so that part is
\(O(e^{(1-c_{\theta_0})t})\).  On \([0,\theta_0]\), put
\(\theta=v/\sqrt t\) and use
\[
 t(\cos(v/\sqrt t)-1)
 =-\frac{v^2}{2}+O(v^4/t)
\]
on \(v\le t^{1/8}\); the complementary portion is bounded by a Gaussian
tail.  Dominated convergence with one further Taylor term gives
\[
 \int_0^{\theta_0}e^{t\cos\theta}\,d\theta
 =\frac{e^t}{\sqrt t}
 \left(\int_0^\infty e^{-v^2/2}\,dv+O(t^{-1})\right).
\]
Since the Gaussian integral is \(\sqrt{\pi/2}\), equation
\textup{(V84.4)} follows.
\end{proof}

\begin{theorem}[Quadratic-Borel principal-type barrier]
\label{thm:v84-principal-type}
Let \(0<d<x\) and let \(W\) be any nonzero fixed finite radial weight
\textup{(V83.2)}.  As \(q\to+\infty\),
\[
 \boxed{
 \frac{
 \left|\displaystyle\sum_j\frac{b_j}{d+a_j}
 I_0\!\left(2\sqrt{\frac q{d+a_j}}\right)\right|}
 {
 \left|\displaystyle\sum_j\frac{b_j}{x+a_j}
 I_0\!\left(2\sqrt{\frac q{x+a_j}}\right)\right|}
 \sim
 \left(\frac{x+a_0}{d+a_0}\right)^{3/4}
 \exp\!\left[
 2\sqrt q\left(
 \frac1{\sqrt{d+a_0}}-\frac1{\sqrt{x+a_0}}
 \right)\right]
 \longrightarrow\infty.}
 \tag{V84.5}
\]
Thus a fixed radial weight cannot reverse principal-first dominance by
quadratic Borel summation.
\end{theorem}

\begin{proof}
Lemma~\ref{lem:v84-i0-asymptotic} gives, for fixed \(c>0\),
\[
 \frac1c I_0\!\left(2\sqrt{\frac q c}\right)
 =
 \frac{c^{-3/4}q^{-1/4}}{\sqrt{4\pi}}
 e^{2\sqrt{q/c}}\left(1+O_c(q^{-1/2})\right).
 \tag{V84.6}
\]
Because \(a_0<a_j\) for \(j>0\), the \(j=0\) term is the unique term of
largest square-root exponential type in both sums in
\textup{(V84.5)}.  Divide by those leading terms and apply
\textup{(V84.6)}.  The factors \(b_0\) and \(q^{-1/4}\) cancel, giving the
displayed asymptotic.
\end{proof}

\subsection{Every positive order sum inherits the V77 debit}

The preceding fixed-weight theorem does not apply when the radial weight
depends on \(r\).  The V77 dual estimate does.

For \(r\ge0\), let \(W_r\) be any finite radial weight satisfying
\[
 W_r(u)\ge0\quad(u\ge L),\qquad
 \mathcal L_d^{(r+1)}[W_r]=0,\qquad
 T_r:=\mathcal L_x^{(r+1)}[W_r]>0.
 \tag{V84.7}
\]
For \(R>d\), put
\[
 C_r(R):=
 \mathfrak C_{r+1}(W_r;R)
 =\sum_j|b_{j,r}|(R+a_{j,r})^{-(r+1)}.
 \tag{V84.8}
\]

\begin{theorem}[Positive-order summation cannot improve the Cauchy ratio]
\label{thm:v84-positive-order-no-go}
Let \(\omega_r\ge0\) be arbitrary.  For every finite \(N\),
\[
 \boxed{
 \sum_{r=0}^{N}\omega_r C_r(R)
 \ge e^{-(R-d)L}
 \sum_{r=0}^{N}\omega_r T_r.}
 \tag{V84.9}
\]
The same inequality holds for infinite sums in \([0,\infty]\).  At the
critical-line distance \(R_0=d+1/2\), \(L=L_f\),
\[
 \boxed{
 \sum_{r\ge0}\omega_r C_r(R_0)
 \ge\ell_f^{-1/2}
 \sum_{r\ge0}\omega_r T_r.}
 \tag{V84.10}
\]
Consequently no nonnegative order weighting---ordinary, Borel, quadratic
Borel, or otherwise---can turn the coefficientwise holomorphic-remainder
budget into \(o(1)\) of the normalized target aggregate.
\end{theorem}

\begin{proof}
Theorem~\ref{thm:v77-stability-lower}, applied at order \(r+1\), gives
\[
 C_r(R)\ge e^{-(R-d)L}T_r.
\]
Multiply by \(\omega_r\ge0\) and sum.  Finite partial sums prove
\textup{(V84.9)}, and monotone convergence proves the extended-valued
infinite version.  Substituting \(R_0-d=1/2\) and
\(e^{-L_f/2}=\ell_f^{-1/2}\) gives \textup{(V84.10)}.
\end{proof}

This is a summation-stable statement about an absolute budget.  It does not
say that the fixed gamma or zero remainder attains that budget; it says that
positive summation cannot manufacture the missing asymptotic reserve from
the V77 estimate.

\subsection{Quadratic-Borel synthesis is larger than the target scale}

Now impose the normalized order-dependent conditions
\[
 T_r=1,\qquad L=L_f,
 \tag{V84.11}
\]
and retain the coefficient measures \(\mu_r\) from
\textup{(V83.27)}.  Define the quadratic-Borel absolute synthesis majorant
\[
 \mathfrak S_2(q):=
 \sum_{r\ge0}\frac{q^r}{(r!)^2}\|\mu_r\|_{\rm TV},
 \qquad q\ge0,
 \tag{V84.12}
\]
allowing the value \(+\infty\).

\begin{theorem}[Exponential synthesis excess]
\label{thm:v84-synthesis-excess}
For every \(q\ge0\),
\[
 \boxed{
 \mathfrak S_2(q)
 \ge
 \frac1{L_f}
 \left(1+\frac q{L_f}\right)e^{q/L_f}.}
 \tag{V84.13}
\]
The normalized target aggregate is
\[
 \boxed{
 \sum_{r\ge0}\frac{q^r}{(r!)^2}T_r
 =I_0(2\sqrt q).}
 \tag{V84.14}
\]
Consequently
\[
 \boxed{
 \frac{\mathfrak S_2(q)}{I_0(2\sqrt q)}
 \longrightarrow+\infty}
 \tag{V84.15}
\]
at exponential rate \(e^{q/L_f-2\sqrt q}\), up to polynomial factors.
\end{theorem}

\begin{proof}
The V77 factorial bound gives
\[
 \|\mu_r\|_{\rm TV}
 \ge\frac{(r+1)!}{L_f^{\,r+1}}.
\]
Therefore, with \(z=q/L_f\),
\[
 \begin{aligned}
 \mathfrak S_2(q)
 &\ge\frac1{L_f}\sum_{r\ge0}
 \frac{r+1}{r!}z^r\\
 &=\frac1{L_f}(1+z)e^z,
 \end{aligned}
\]
which is \textup{(V84.13)}.  Equation \textup{(V84.14)} is the defining
series of \(I_0\).  Lemma~\ref{lem:v84-i0-asymptotic} then proves
\textup{(V84.15)}.
\end{proof}

The comparison is deliberately made with the absolute synthesis majorant,
not with the norm of a conditionally combined signed measure.  Any
coefficient-by-coefficient tail estimate pays \(\mathfrak S_2\); exploiting
cancellation across orders would require a new common signed identity.

\subsection{Positive summation does not restore an archimedean sign}

V79's higher-prime-power correction is nonpositive for each tail-positive
weight, so nonnegative order weights preserve that favorable sign.  The
polygamma term behaves differently.

\begin{lemma}[Unbounded real archimedean null direction]
\label{lem:v84-unbounded-arch}
In the real-ordinate setting of V78, fix one order \(k\ge2\).  Subject to
\[
 W(u)>0\quad(u\ge L_f),\qquad
 \mathcal L_d^{(k)}[W]=0,\qquad
 \mathcal L_x^{(k)}[W]=1,
 \tag{V84.16}
\]
the exact real polygamma response
\(\mathfrak A_{\chi,k}^{W}(\sigma)\) takes every real value.
\end{lemma}

\begin{proof}
The compactly supported function \(h\) in Lemma
\ref{lem:v78-arch-null-direction} annihilates the principal and target
moments but has nonzero archimedean response.  Begin with the strict
continuous feasible precursor from V76, rescaled to target moment one.
The family \(w+t h\), \(t\in\mathbb R\), preserves both moment constraints
and is unchanged on \([L_f,\infty)\), while its archimedean response is a
nonconstant affine function of \(t\).

For two parameters \(t_-,t_+\) whose responses straddle a prescribed real
number, use the polynomial approximation and exact two-moment correction
from the proof of Theorem~\ref{thm:v78-sign-tunability}.  Approximation can
be chosen sufficiently accurate to preserve strict tail positivity and the
strict inequalities around the prescribed value.  The convex segment
between the two resulting exponential polynomials preserves
\textup{(V84.16)} and tail positivity.  Continuity of the archimedean
functional gives a point on that segment with exactly the prescribed
response.
\end{proof}

\begin{corollary}[No archimedean sign from a positive order sum]
\label{cor:v84-summed-arch}
Let \(\omega_r\ge0\) be any order weights with
\(\omega_{r_0}>0\) for some \(r_0\ge1\).  Allow the feasible radial weight
\(W_r\) to depend on \(r\).  Holding all \(W_r\), \(r\ne r_0\), fixed, the
summed real polygamma response
\[
 \sum_r\omega_r
 \mathfrak A_{\chi,r+1}^{W_r}(\sigma)
 \tag{V84.17}
\]
can be assigned either sign, or zero, by changing only \(W_{r_0}\).
\end{corollary}

\begin{proof}
Apply Lemma~\ref{lem:v84-unbounded-arch} at order \(r_0+1\).  Since
\(\omega_{r_0}>0\), its freely prescribed real response can offset the
fixed contribution from all other orders.  For an infinite sum this
statement applies whenever the fixed complementary sum is defined.
\end{proof}

Thus the exact favorable V79 sign survives positive order summation, but the
V78 archimedean ambiguity survives as well.  At nonzero ordinate, the
oscillatory kernel in \textup{(V78.21)} provides no stronger sign.

\begin{theorem}[V84 higher-Borel verdict]
\label{thm:v84-verdict}
The Gevrey-compatible positive-summation route has the following exact
outcome.
\begin{enumerate}[label=\textup{(\roman*)},leftmargin=3em]
\item Quadratic Borel summation has the exact positive Euler kernel
      \(\mathfrak F_2(q\,u)\) and the exact \(I_0\) singular kernel
      \textup{(V84.2)}.
\item For every fixed radial weight, the principal response has larger
      square-root exponential type than the target response.
\item For order-dependent weights, every nonnegative summation of orders
      inherits the V77 target-to-Cauchy-cost ratio without improvement.
\item Under quadratic Borel weights, the absolute radial synthesis majorant
      grows at least like \(e^{q/L_f}\), while the normalized target grows
      like \(e^{2\sqrt q}\).
\item Positive order summation preserves the favorable signed V79
      higher-prime-power correction, but it cannot impose a sign on the
      exact V78 polygamma response.
\end{enumerate}
Therefore no scalar nonnegative summation of derivative orders closes the
V82 large-window corridor by fixed-weight dominance, absolute Cauchy
control, or an archimedean sign.
\end{theorem}

\begin{firewall}[A positive-summation no-go is not GRH]
Theorem~\ref{thm:v84-positive-order-no-go} concerns the available absolute
remainder norm, and Corollary~\ref{cor:v84-summed-arch} concerns freedom
inside the unrestricted order-dependent radial cone.  Neither determines
the signed contribution of the actual infinite zero multiset.  A coupled
identity using cancellations across orders or channels remains logically
possible.  GRH remains open.
\end{firewall}

\begin{assessment}[V84 closure and V85 direction]
V84 closes the remaining scalar positive Borel variant proposed after V83.
Adding factorials changes the special function but not the principal-first
ordering for a fixed weight, and V77's dual debit is invariant under every
positive summation when the weights vary with the order.

The next noncircular possibility must permit signed order coefficients while
retaining nonnegativity of the \emph{aggregate} prime kernel on the
arithmetic half-line.  V85 should formulate the cone of entire order kernels
\[
 K(u)=\sum_{r\ge0}\omega_r\frac{u^r}{r!},
 \qquad K(u)\ge0\quad(u\ge L_f),
\]
without requiring \(\omega_r\ge0\) termwise.  It should then determine
whether the same sub-prime support gap recreates the V76 interpolation
freedom, or whether analyticity and controlled exponential type yield a
genuine dual separator.  Any successful separator must be applied to the
full signed zero and gamma distributions before absolute values are taken.
\end{assessment}

\section*{Version V84 append-only record}
\addcontentsline{toc}{section}{Version V84 append-only record}
The complete V83 source was retained before this continuation.  It had
1199468 bytes, 35253 lines, and SHA--256
\[
 \texttt{BFBE558655FD5DEF6269F6D4398871BABD390116BB32C8663F184F0C1839F02E}.
\]
V84 evaluated the quadratic Borel Euler and singular kernels, proved the
square-root principal-type barrier, established the positive-order summation
no-go for the V77 Cauchy ratio, quantified exponential synthesis excess, and
proved persistence of exact archimedean sign tunability under positive
order summation.

% =============================================================================
% END APPENDED V84 MATERIAL
% =============================================================================
% =============================================================================
% BEGIN APPENDED V85 MATERIAL
% =============================================================================

\section{V85: bounded-band signed order kernels and return of finite-window escape}
\label{sec:grh-v85}

V84 closed every scalar order sum whose coefficients are nonnegative term by
term.  It proposed allowing signed order coefficients while requiring only
the aggregate prime kernel to be nonnegative on the arithmetic half-line.
That larger cone can be identified exactly.

An exponential polynomial
\[
 K(u)=\sum_j b_j e^{-a_j u}
\]
is an entire order kernel with signed Taylor coefficients
\[
 \omega_r=K^{(r)}(0)=\sum_j b_j(-a_j)^r.
\]
When all \(a_j\) lie strictly below the Euler convergence margin
\(\sigma-1\), the signed sum of normalized logarithmic-derivative orders is
absolutely convergent and equals the corresponding finite radial shift
combination.  The apparent new order operation is therefore the old radial
operation in a different functional calculus.

There is a stronger fact.  For every bandwidth \(A>0\), finite exponential
polynomials using only exponents \(0\le a<A\) are uniformly dense on the
one-point compactification of the positive half-line.  The V81
tail-positive interpolation theorem can consequently be carried out inside
an arbitrarily narrow exponent band.  One may simultaneously impose
\(K(0)=0\), which removes the order-zero coefficient and regularizes both
the undifferentiated Hadamard zero sum and the digamma ladder.  Taking
\(A<\sigma-1\) therefore recreates finite-window escape inside the strict
absolute-convergence domain of the signed order sum.

\subsection{Signed order coefficients are bounded radial shifts}

Fix a finite real exponential polynomial
\[
 K(u)=\sum_{j=0}^{J}b_j e^{-a_j u},
 \qquad 0\le a_j<A,
 \tag{V85.1}
\]
with distinct exponents after equal terms are combined, and define
\[
 \omega_r:=\sum_{j=0}^{J}b_j(-a_j)^r
 \qquad(r\ge0).
 \tag{V85.2}
\]
Then
\[
 K(u)=\sum_{r\ge0}\omega_r\frac{u^r}{r!}
 \tag{V85.3}
\]
with locally uniform, indeed entire, convergence.

\begin{theorem}[Absolute signed-order/radial equivalence]
\label{thm:v85-order-radial-equivalence}
Let \(s\in\mathbb C\) satisfy
\[
 \Re s-A>1.
 \tag{V85.4}
\]
For every Dirichlet character \(\eta\),
\[
 \boxed{
 \sum_{r\ge0}\omega_r\mathcal Q_{\eta,r}(s)
 =
 \sum_{j=0}^{J}b_j
 \left(-\frac{L'}{L}\right)(s+a_j,\eta)}
 \tag{V85.5}
\]
and both sides have the absolutely convergent Euler representation
\[
 \boxed{
 \sum_p\sum_{m\ge1}
 \frac{(\log p)\eta(p)^m}{p^{ms}}K(m\log p).}
 \tag{V85.6}
\]
Moreover, if \(|y|>A\), then the singular response is absolutely summable:
\[
 \boxed{
 \sum_{r\ge0}\frac{\omega_r}{y^{r+1}}
 =\sum_{j=0}^{J}\frac{b_j}{y+a_j}
 =\int_0^\infty e^{-yu}K(u)\,du.}
 \tag{V85.7}
\]
\end{theorem}

\begin{proof}
For \(u\ge0\),
\[
 \sum_{r\ge0}|\omega_r|\frac{u^r}{r!}
 \le\sum_j|b_j|e^{a_j u}
 \le B(K)e^{Au},
 \qquad B(K):=\sum_j|b_j|.
 \tag{V85.8}
\]
Insert the Euler expansion \textup{(V83.8)}.  Condition
\textup{(V85.4)} majorizes the absolute triple sum by the logarithmic
derivative Euler series at \(\Re s-A>1\).  Summing in \(r\) gives
\[
 \sum_{r\ge0}\omega_r\frac{(m\log p)^r}{r!}
 =K(m\log p),
\]
which proves \textup{(V85.5)}--\textup{(V85.6)}.  Equivalently, the two
minus signs in \(\omega_r=(-a_j)^r\) and in the definition of
\(\mathcal Q_{\eta,r}\) turn the order series into the Taylor expansion at
\(s+a_j\).

For \(|y|>A\),
\[
 \sum_{r\ge0}\frac{|\omega_r|}{|y|^{r+1}}
 \le\sum_j\frac{|b_j|}{|y|-a_j}<\infty.
 \tag{V85.9}
\]
The geometric sum is
\[
 \sum_{r\ge0}\frac{(-a_j)^r}{y^{r+1}}
 =\frac1{y+a_j},
\]
and the Laplace integral gives the same value.
\end{proof}

If \(K(0)=0\), then \(\omega_0=\sum_jb_j=0\), so the left side of
\textup{(V85.5)} begins at \(r=1\).  Thus no undifferentiated
logarithmic-derivative row has been silently added.

\subsection{Narrow exponent bands are uniformly dense}

Let \(X=[0,\infty]\) be the one-point compactification of the positive
half-line.  For \(A>0\), write
\[
 \mathcal E_A:=
 \operatorname{span}_{\mathbb R}
 \{u\mapsto e^{-au}:0\le a<A\}
 \subset C(X;\mathbb R).
 \tag{V85.10}
\]

\begin{lemma}[Density of every positive exponent band]
\label{lem:v85-band-density}
For every \(A>0\),
\[
 \boxed{\overline{\mathcal E_A}^{\,\|\cdot\|_\infty}
 =C(X;\mathbb R).}
 \tag{V85.11}
\]
\end{lemma}

\begin{proof}
Suppose a finite signed Borel measure \(\nu\) on \(X\) annihilates
\(\mathcal E_A\).  Its Laplace transform
\[
 F(z):=\int_{[0,\infty)}e^{-zu}\,d\nu(u)
\]
is holomorphic on \(\Re z>0\) and vanishes for every real
\(z\in(0,A)\).  The identity theorem gives \(F(z)=0\) throughout that
half-plane.  Fix \(t>0\) and push \(\nu|_{[0,\infty)}\) forward by
\(q=e^{-tu}\).  Its moments of every positive integer order are
\(F(nt)=0\).  The constant function \(a=0\) is also in
\(\mathcal E_A\), so the zeroth moment, including the mass at infinity, is
zero.  Polynomial density on \([0,1]\) now gives \(\nu=0\).
Hahn--Banach and the Riesz representation theorem prove
\textup{(V85.11)}.
\end{proof}

The endpoint constraint needed below does not destroy moment surjectivity.
Retain the V81 functionals
\[
 \Phi_{y,n}[K]
 =\frac1{n!}\int_0^\infty u^n e^{-yu}K(u)\,du.
 \tag{V85.12}
\]

\begin{lemma}[Band-limited moments plus one endpoint value]
\label{lem:v85-augmented-surjectivity}
Let \(\mathcal I\subset\mathbb H_R\times\mathbb N_0\) be finite and
conjugation closed.  The augmented real map
\[
 \mathcal E_A\longrightarrow
 \mathcal E_{\mathcal I}\times\mathbb R,
 \qquad
 K\longmapsto
 \left((\Phi_{y,n}[K])_{(y,n)\in\mathcal I},K(0)\right)
 \tag{V85.13}
\]
is onto.
\end{lemma}

\begin{proof}
If a real covector annihilated the range, then for every \(a\in[0,A)\)
one would have
\[
 c+\sum_{(y,n)\in\mathcal I}
 c_{y,n}(y+a)^{-(n+1)}=0,
 \tag{V85.14}
\]
with conjugate coefficients paired so that the expression is real on the
real \(a\)-interval.  The left side is meromorphic in \(a\) and vanishes on
an interval, hence identically away from its poles.  Letting
\(a\to+\infty\) gives \(c=0\).  Uniqueness of the principal parts at the
distinct poles \(-y\), including their different orders, gives
\(c_{y,n}=0\) for every pair.  Thus no nonzero covector annihilates the
range, proving surjectivity.
\end{proof}

\subsection{Tail-positive interpolation below any convergence margin}

\begin{theorem}[Bounded-band tail-positive spectral interpolation]
\label{thm:v85-bounded-band-interpolation}
Fix \(A,L>0\), a finite conjugation-closed index set
\(\mathcal I\subset\mathbb H_R\times\mathbb N_0\), and arbitrary compatible
moment data \(c=(c_{y,n})_{\mathcal I}\).  There is a real finite
exponential polynomial
\[
 K(u)=\sum_{j=0}^{J}b_j e^{-a_j u},
 \qquad 0\le a_j<A,
 \tag{V85.15}
\]
such that
\[
 \boxed{
 K(0)=0,\qquad
 K(u)>0\quad(u\ge L),\qquad
 \Phi_{y,n}[K]=c_{y,n}
 \quad((y,n)\in\mathcal I).}
 \tag{V85.16}
\]
\end{theorem}

\begin{proof}
Choose \(a_\star\in(0,A)\) and \(\varepsilon>0\).  The continuous baseline
\[
 g_0(u):=\varepsilon(1-e^{-a_\star u})
 \tag{V85.17}
\]
vanishes at zero and has the positive tail floor
\[
 g_0(u)\ge\varepsilon(1-e^{-a_\star L})
 \qquad(u\ge L).
 \tag{V85.18}
\]
By Lemma~\ref{lem:v81-exponential-independence}, choose
\(h\in C_c(0,L;\mathbb R)\) whose moments equal
\[
 \Phi_{y,n}[h]=c_{y,n}-\Phi_{y,n}[g_0].
 \tag{V85.19}
\]
Then \(w:=g_0+h\) has the desired moment table, satisfies \(w(0)=0\), and
equals \(g_0\) on \([L,\infty)\).

Lemma~\ref{lem:v85-band-density} supplies
\(P_\nu\in\mathcal E_A\) converging uniformly to \(w\) on \(X\).
Each moment functional is uniformly continuous by
\textup{(V81.10)}, and endpoint evaluation is uniformly continuous.
Lemma~\ref{lem:v85-augmented-surjectivity} supplies finitely many fixed
functions \(Q_1,\ldots,Q_N\in\mathcal E_A\) whose augmented data vectors
form a basis.  Correct the small augmented-data error of \(P_\nu\) by a
linear combination \(R_\nu\) of the \(Q_j\).  Finite-dimensional norm
equivalence gives
\[
 \|R_\nu\|_\infty\longrightarrow0.
 \tag{V85.20}
\]
For sufficiently large \(\nu\), \(K:=P_\nu+R_\nu\) has the exact moments,
has \(K(0)=0\), and remains strictly positive on \([L,\infty)\) by
\textup{(V85.18)}.  Combining equal exponents gives
\textup{(V85.15)}.
\end{proof}

This strengthens V81 in a direction needed for summability: the exponents
can be confined to an arbitrarily short interval adjacent to zero.  The
price is uncontrolled coefficient variation and clustering of the selected
exponents.

\subsection{Absolute regularization of the order-zero completed formula}

Take \(\sigma>1\), put \(d=\sigma-1\), and choose
\[
 0<A<d.
 \tag{V85.21}
\]
Every nontrivial-zero response node occurring at \(\sigma\) or
\(\sigma+i\gamma\) has real part strictly larger than \(d\); the same is true
of every gamma node.  Hence \(|y|>A\) for all such nodes.

\begin{lemma}[Endpoint cancellation regularizes zeros and gamma]
\label{lem:v85-order-zero-regularization}
Let \(K\) satisfy \textup{(V85.15)} and \(K(0)=\sum_jb_j=0\).
Then:
\begin{enumerate}[label=\textup{(\roman*)},leftmargin=3em]
\item for fixed \(s\) away from zeros,
\[
 \sum_\rho m_\rho
 \left|\sum_j\frac{b_j}{s+a_j-\rho}\right|<\infty;
 \tag{V85.22}
\]
\item the conductor constant cancels in
\(\sum_jb_jA_\chi(s+a_j)\);
\item the exact digamma response is
\[
 \boxed{
 \sum_jb_jA_\chi(s+a_j)
 =
 -\int_0^\infty
 \frac{e^{-(s+\kappa_\chi)u}}{1-e^{-2u}}K(u)\,du
 =
 -\sum_{n\ge0}
 \mathcal L_{s+\kappa_\chi+2n}^{(1)}[K],}
 \tag{V85.23}
\]
and both the integral and the final series converge absolutely.
\end{enumerate}
\end{lemma}

\begin{proof}
For large \(|\rho|\), uniformly in the finite shift set,
\[
 \sum_j\frac{b_j}{s+a_j-\rho}
 =
 \frac{\sum_jb_j}{s-\rho}
 +O_K(|\rho|^{-2})
 =O_K(|\rho|^{-2}).
 \tag{V85.24}
\]
The standard order-one zero count makes the last majorant summable, proving
\textup{(V85.22)}.

The constant \(\frac12\log(f/\pi)\) in \(A_\chi\) is multiplied by
\(\sum_jb_j=0\).  Apply the classical integral formula for the digamma
function to the remaining shifted combination.  Its shift-independent term
also cancels, and the change of variable \(t=2u\) gives the integral in
\textup{(V85.23)}.  Near zero, \(K(u)=O_K(u)\), so
\[
 \frac{K(u)}{1-e^{-2u}}=O_K(1).
 \]
At infinity the integrand decays exponentially.  Expanding
\((1-e^{-2u})^{-1}=\sum_{n\ge0}e^{-2nu}\) gives the last expression.
Because \(K(0)=0\), its Laplace transform is \(O_K(y^{-2})\) as
\(y\to+\infty\), so the \(n\)-series is absolutely convergent.
\end{proof}

Thus the restriction \(r\ge1\) used in V78 is replaced here by an exact
endpoint cancellation.  No conditionally rearranged zero or gamma sum is
being used.

\subsection{Signed-order finite-window escape}

Fix the two-channel setting of V80 and a hypothetical zero
\(\rho=\beta+i\gamma\), \(1/2<\beta<1\).  Let
\[
 x:=\sigma-\beta>d.
 \tag{V85.25}
\]
Choose any finite principal zero set, any finite target-character zero set
not containing \(\rho\), and any finite prefixes of the two gamma ladders.
Form the corresponding finite conjugation-closed response-node set
\(\mathcal Y\), exactly as in \textup{(V81.16)}--\textup{(V81.17)}.

\begin{theorem}[Finite-window escape inside the signed-order convergence disk]
\label{thm:v85-signed-order-escape}
For every \(0<A<d\), there is a real finite exponential order kernel \(K\)
with exponents in \([0,A)\) such that
\[
 \boxed{
 K(0)=0,\qquad K(u)>0\ (u\ge L_f),\qquad
 \mathcal L_d^{(1)}[K]=0,\qquad
 \mathcal L_x^{(1)}[K]=1,}
 \tag{V85.26}
\]
and
\[
 \boxed{
 \mathcal L_y^{(1)}[K]=0
 \qquad(y\in\mathcal Y).}
 \tag{V85.27}
\]
Its signed order coefficients \(\omega_r=K^{(r)}(0)\) satisfy
\(\omega_0=0\), and the aggregate order formula is absolutely convergent on
the Euler side and at every named zero and gamma node.

In the degree-one phase channel \(1+\cos\theta\), the aggregate prime side
is nonnegative, the principal Laurent response vanishes, the selected-zero
response is \(-m_\rho\), the named finite zero and gamma responses vanish,
and the signed higher-prime-power correction is nonpositive.  Hence the
uncancelled two-channel zero/gamma tail must satisfy
\[
 \boxed{
 \mathfrak R_{\star,\mathrm{tail}}^{K}
 \ge m_\rho-\mathfrak H_\star^{K}
 \ge m_\rho.}
 \tag{V85.28}
\]
Thus aggregate Euler positivity does not prevent the V80 compensation from
escaping beyond every prescribed finite spectral window, even when the
signed order sum has a strict absolute-convergence margin.
\end{theorem}

\begin{proof}
Apply Theorem~\ref{thm:v85-bounded-band-interpolation} with \(L=L_f\),
the pairs \((d,0)\), \((x,0)\), and \((y,0)\) for
\(y\in\mathcal Y\), prescribing respectively \(0,1,0\).  This gives
\textup{(V85.26)}--\textup{(V85.27)}.

Since \(A<d=\sigma-1\), Theorem
\ref{thm:v85-order-radial-equivalence} gives absolute convergence of the
aggregate Euler formula.  Every named spectral node has real part larger
than \(d\), so its order series is also absolute.  Lemma
\ref{lem:v85-order-zero-regularization} justifies the complete zero and
gamma decomposition.

The prime-phase identity multiplied by \(K(m\log p)>0\) is termwise
nonnegative.  The principal and named finite responses have the prescribed
values, while the selected zero contributes \(-m_\rho\).  The V79
higher-prime-power calculation is an exact negative Euler tail multiplied
by the same positive kernel \(K\), so
\(\mathfrak H_\star^K\le0\).  Rearranging the completed two-channel identity
gives \textup{(V85.28)}.
\end{proof}

The order coefficients are genuinely signed.  If they were all
nonnegative, V84 would apply.  Their signs cooperate only after the complete
entire kernel \(K\) has been summed, but Theorem
\ref{thm:v85-signed-order-escape} shows that this added freedom is exactly
large enough to restore the sub-prime interpolation mechanism.

\begin{theorem}[V85 signed-order verdict]
\label{thm:v85-verdict}
The aggregate-positive signed-order programme has the following exact
outcome.
\begin{enumerate}[label=\textup{(\roman*)},leftmargin=3em]
\item A finite exponential aggregate kernel has signed order coefficients
      and is exactly equivalent to a finite radial shift combination.
\item The equivalence is absolutely convergent whenever its exponent band
      lies below the Euler margin \(\sigma-1\).
\item Every positive exponent band, however narrow, is uniformly dense on
      the compactified half-line and supports arbitrary finite compatible
      Laplace interpolation.
\item The interpolation can impose \(K(0)=0\), which removes order zero and
      makes the undifferentiated zero and gamma differences absolute.
\item Consequently every prescribed finite two-channel zero/gamma window
      can be cancelled inside the aggregate-positive cone while retaining a
      normalized off-line target and a strict convergence margin.
\end{enumerate}
Thus signed scalar order coefficients do not supply the separator missing
from V84; they reproduce the V81 finite-window escape in a summable
bounded-band form.
\end{theorem}

\begin{firewall}[Finite-window escape is not a zero construction]
Theorem~\ref{thm:v85-signed-order-escape} chooses its interpolation table
using a hypothetical finite zero window.  It proves capacity of the
aggregate-positive test cone, not existence of an off-line zero.  The
uncontrolled coefficient variation and the actual infinite tail remain.
Equation \textup{(V85.28)} is a necessary compensation law, not a
contradiction.  GRH remains open.
\end{firewall}

\begin{assessment}[V85 closure and V86 direction]
V85 closes the signed scalar-order alternative proposed after V84.  The
closure is stronger than a formal identification: bounded-band interpolation
shows that the obstruction persists with absolute convergence and exact
order-zero regularization.

The remaining quantitative object is now a bounded-frequency
superresolution constant.  V86 should estimate the least coefficient
variation required to interpolate a growing actual zero/gamma window using
exponents \(0\le a<A<d\), with the endpoint constraint
\(\sum_jb_j=0\) retained.  This is a compact-band generalized Cauchy
quotient, different from V82's unrestricted-radius quotient.  A lower bound
strong enough to dominate the zero-counting tail would close the scalar
programme; a matching construction would prove that every scalar
Euler-positive route is intrinsically maskable and force a genuinely
matrix-valued completed-spectrum polarization.
\end{assessment}

\section*{Version V85 append-only record}
\addcontentsline{toc}{section}{Version V85 append-only record}
The complete V84 source was retained before this continuation.  It had
1214330 bytes, 35676 lines, and SHA--256
\[
 \texttt{8542BD9F7BABB2EA0A305FE17F2F9F861CED417D2AFEDE3865C6B425EA046617}.
\]
V85 proved bounded-band density and tail-positive interpolation with endpoint
cancellation, identified signed order kernels with absolutely convergent
radial shifts, regularized the order-zero zero and gamma sums, and proved
finite-window compensation escape inside every strict Euler convergence
margin.

% =============================================================================
% END APPENDED V85 MATERIAL
% =============================================================================
% =============================================================================
% BEGIN APPENDED V86 MATERIAL
% =============================================================================

\section{V86: Chebyshev superresolution and the growing endpoint-jet corridor}
\label{sec:grh-v86}

V85 proved qualitative interpolation inside every strict exponent band
\(0\le a<A<\sigma-1\).  This continuation computes a quantitative dual
bound for that compact-band problem.  The endpoint cancellation
\(\sum_jb_j=0\) is decisive: it allows the target Cauchy column to be
approximated not only by the cancelled gamma columns but also by a constant.
Chebyshev interpolation on the radius interval then gives an exponentially
small quotient distance and hence an exponentially large coefficient norm.

The endpoint cancellation can be strengthened.  For every finite
\(\nu\), exponential polynomials whose first \(\nu+1\) derivatives vanish
at zero remain uniformly dense among continuous functions vanishing at
zero.  Arbitrary finite spectral interpolation is still possible in that
subspace.  Each additional endpoint moment improves the distant zero and
gamma decay by one power, but the Chebyshev coefficient cost acquires
another factor \(4x/A\).  The exact balance shows that a coefficientwise
gamma-tail argument cannot even enter an undecided regime unless the
endpoint jet grows at least on the scale \(M/\log M\).

\subsection{Endpoint-jet exponent bands}

For \(A>0\) and an integer \(\nu\ge0\), define
\[
 \mathcal E_A^{(\nu)}
 :=
 \left\{
 K(u)=\sum_{j=0}^{J}b_j e^{-a_j u}:
 0\le a_j<A,\quad
 \sum_jb_j a_j^r=0\ (0\le r\le\nu)
 \right\}.
 \tag{V86.1}
\]
Equivalently,
\[
 K^{(r)}(0)=0\qquad(0\le r\le\nu).
 \tag{V86.2}
\]

\begin{lemma}[Every finite endpoint jet is uniformly invisible]
\label{lem:v86-jet-density}
Let \(X=[0,\infty]\) be the one-point compactification.  For every
\(A>0\) and finite \(\nu\),
\[
 \boxed{
 \overline{\mathcal E_A^{(\nu)}}^{\,\|\cdot\|_\infty}
 =
 \{f\in C(X;\mathbb R):f(0)=0\}.}
 \tag{V86.3}
\]
\end{lemma}

\begin{proof}
Every member of \(\mathcal E_A^{(\nu)}\) vanishes at zero, so only density
in the displayed codimension-one space needs proof.  Let a finite signed
measure \(\mu\) on \(X\) annihilate \(\mathcal E_A^{(\nu)}\), and put
\[
 F(a):=\int_X e^{-au}\,d\mu(u),
 \qquad 0<a<A,
 \tag{V86.4}
\]
with \(e^{-a\infty}=0\).

Put \(\phi_0\equiv1\) on \(X\), and put
\(\phi_a(u)=e^{-au}\) for \(a>0\), with \(\phi_a(\infty)=0\).  Thus
\(F(a)=\int_X\phi_a\,d\mu\), including at \(a=0\).  Choose distinct
anchors \(a_0,\ldots,a_\nu\in[0,A)\).  For any other
\(a\in[0,A)\), subtract from the point mass at \(a\) its Lagrange
interpolant on the anchors.  The resulting finitely supported coefficient
vector has zero moments through degree \(\nu\), so annihilation gives
\[
 F(a)=P_\nu(a)
 \tag{V86.5}
\]
for the polynomial of degree at most \(\nu\) interpolating the anchor
values.  On \(a>0\), the transform \(F\) is holomorphic on \(\Re a>0\);
therefore \textup{(V86.5)} holds throughout that half-plane.  Since \(F\)
is bounded on the positive real axis, \(P_\nu\) is constant, say \(C\).
Dominated convergence as \(a\downarrow0\) gives
\(C=\mu(X\setminus\{\infty\})\), whereas the already established identity
at \(a=0\) gives \(C=F(0)=\mu(X)\).  Thus \(\mu(\{\infty\})=0\).
Subtracting \(C\) times the atom at \(u=0\) now leaves a measure on
\([0,\infty)\) with identically zero Laplace transform.  Laplace uniqueness,
as in Lemma \ref{lem:v85-band-density}, makes that measure zero.  Hence the annihilator
is exactly the span of evaluation at \(u=0\).  Hahn--Banach proves
\textup{(V86.3)}.
\end{proof}

The same Lagrange argument gives spectral surjectivity inside the jet
subspace.

\begin{lemma}[Finite spectral moments survive every finite endpoint jet]
\label{lem:v86-jet-moment-surjectivity}
For every finite conjugation-closed
\(\mathcal I\subset\mathbb H_R\times\mathbb N_0\), the real map
\[
 \mathcal E_A^{(\nu)}
 \longrightarrow\mathcal E_{\mathcal I},
 \qquad
 K\longmapsto(\Phi_{y,n}[K])_{(y,n)\in\mathcal I}
 \tag{V86.6}
\]
is onto.
\end{lemma}

\begin{proof}
If a nonzero compatible covector annihilated the range, the meromorphic
function
\[
 H(a):=\sum_{(y,n)\in\mathcal I}
 c_{y,n}(y+a)^{-(n+1)}
 \tag{V86.7}
\]
would annihilate every finitely supported coefficient vector on
\([0,A)\) whose moments through degree \(\nu\) vanish.  Lagrange
interpolation then forces \(H\) to agree on \((0,A)\) with a polynomial of
degree at most \(\nu\).  Analytic continuation gives the same identity
away from the poles.  Since \(H(a)\to0\) as \(a\to+\infty\), that polynomial
is zero.  Uniqueness of the principal parts at the distinct poles \(-y\)
then forces every \(c_{y,n}=0\), a contradiction.
\end{proof}

\begin{theorem}[Jet-regularized bounded-band interpolation]
\label{thm:v86-jet-interpolation}
Fix \(A,L>0\), a finite \(\nu\), and arbitrary compatible data on a finite
index set \(\mathcal I\).  There is \(K\in\mathcal E_A^{(\nu)}\) such that
\[
 \boxed{
 K(u)>0\quad(u\ge L),\qquad
 \Phi_{y,n}[K]=c_{y,n}
 \quad((y,n)\in\mathcal I).}
 \tag{V86.8}
\]
\end{theorem}

\begin{proof}
Use the continuous precursor \(w=g_0+h\) from the proof of Theorem
\ref{thm:v85-bounded-band-interpolation}.  It vanishes at zero, has the
prescribed spectral moments, and has a fixed positive floor on
\([L,\infty)\).  Lemma~\ref{lem:v86-jet-density} approximates \(w\)
uniformly by members of \(\mathcal E_A^{(\nu)}\).  By Lemma
\ref{lem:v86-jet-moment-surjectivity}, choose finitely many fixed members of
that same subspace whose moment vectors form a basis.  Correct the moment
error exactly as in V81 and V85.  The correction tends to zero uniformly,
preserves the endpoint jet, and therefore preserves the positive tail for
a sufficiently accurate approximation.
\end{proof}

Thus arbitrarily high but finite endpoint regularization does not remove
the qualitative V85 masking capacity.

\subsection{The endpoint-constrained Cauchy quotient}

Fix
\[
 0<A<x<\Re z,\qquad h>0,\qquad M\ge1,\qquad \nu\ge0,
 \tag{V86.9}
\]
and put \(y_m=z+hm\).  Define
\[
 \begin{aligned}
 \delta_{A,M}^{(\nu)}(x,z;h)
 :=\inf_{\substack{P\in\mathbb C[a]\\ \deg P\le\nu\\
                   c_0,\ldots,c_{M-1}\in\mathbb C}}
 \sup_{0\le a\le A}
 \left|
 \frac1{a+x}-P(a)
 -\sum_{m=0}^{M-1}\frac{c_m}{a+y_m}
 \right|.
 \end{aligned}
 \tag{V86.10}
\]
The polynomial component is dual to the endpoint-jet conditions in
\textup{(V86.1)}.

\begin{proposition}[Endpoint-jet quotient duality]
\label{prop:v86-jet-quotient-duality}
Suppose
\[
 K(u)=\sum_{j=0}^{J}b_j e^{-a_j u},
 \qquad 0\le a_j<A,
 \tag{V86.11}
\]
satisfies
\[
 \sum_jb_j a_j^r=0\ (0\le r\le\nu),\qquad
 \sum_j\frac{b_j}{x+a_j}=1,\qquad
 \sum_j\frac{b_j}{y_m+a_j}=0\ (0\le m<M).
 \tag{V86.12}
\]
Then
\[
 \boxed{
 B(K):=\sum_j|b_j|
 \ge\delta_{A,M}^{(\nu)}(x,z;h)^{-1}.}
 \tag{V86.13}
\]
\end{proposition}

\begin{proof}
For every approximant in \textup{(V86.10)}, the conditions
\textup{(V86.12)} give
\[
 1=\sum_jb_j
 \left[
 \frac1{x+a_j}-P(a_j)
 -\sum_{m=0}^{M-1}\frac{c_m}{y_m+a_j}
 \right].
\]
Taking absolute values and then the infimum proves
\textup{(V86.13)}.
\end{proof}

\subsection{The Chebyshev dual certificate}

Let
\[
 n:=M+\nu+1
 \tag{V86.14}
\]
and let \(Q_n\) be the monic Chebyshev polynomial on \([0,A]\).  Thus its
zeros lie strictly inside that interval and
\[
 \|Q_n\|_{L^\infty[0,A]}
 =\frac{A^n}{2^{2n-1}}
 =2\left(\frac A4\right)^n.
 \tag{V86.15}
\]

\begin{theorem}[Exponential bounded-band superresolution cost]
\label{thm:v86-chebyshev-cost}
Under \textup{(V86.9)},
\[
 \boxed{
 \delta_{A,M}^{(\nu)}(x,z;h)
 \le
 \frac2x\left(\frac A{4x}\right)^{M+\nu+1}
 \prod_{m=0}^{M-1}
 \frac{|z+hm-x|}{|z+hm|}.}
 \tag{V86.16}
\]
Consequently every \(K\) satisfying \textup{(V86.12)} obeys
\[
 \boxed{
 B(K)\ge
 \frac x2\left(\frac{4x}{A}\right)^{M+\nu+1}
 \prod_{m=0}^{M-1}
 \frac{|z+hm|}{|z+hm-x|}.}
 \tag{V86.17}
\]
For fixed \(x,z,h\),
\[
 \boxed{
 B(K)\ge
 C_{x,z,h}
 \left(\frac{4x}{A}\right)^{M+\nu+1}
 M^{x/h}.}
 \tag{V86.18}
\]
\end{theorem}

\begin{proof}
Set
\[
 D_M(a):=(a+x)\prod_{m=0}^{M-1}(a+y_m)
\]
and
\[
 E_n(a):=
 \frac{\displaystyle
       \prod_{m=0}^{M-1}(y_m-x)}
      {Q_n(-x)}
 \frac{Q_n(a)}{D_M(a)}.
 \tag{V86.19}
\]
Its residue at \(a=-x\) is one.  Since the numerator degree exceeds the
denominator degree by \(\nu\), partial fractions and polynomial division
give
\[
 E_n(a)=\frac1{a+x}-P_\nu(a)
 -\sum_{m=0}^{M-1}\frac{c_m}{a+y_m}
 \tag{V86.20}
\]
for some polynomial \(P_\nu\) of degree at most \(\nu\) and some
coefficients \(c_m\).

If \(t_1,\ldots,t_n\) are the zeros of \(Q_n\), then
\[
 |Q_n(-x)|=\prod_{i=1}^{n}(x+t_i)\ge x^n.
 \tag{V86.21}
\]
For \(0\le a\le A\), positivity of all denominator real parts gives
\[
 a+x\ge x,\qquad |a+y_m|\ge|y_m|.
 \tag{V86.22}
\]
Insert \textup{(V86.15)}, \textup{(V86.21)}, and
\textup{(V86.22)} into \textup{(V86.19)}.  This gives
\textup{(V86.16)}, and Proposition
\ref{prop:v86-jet-quotient-duality} gives \textup{(V86.17)}.
Finally the gamma-ratio calculation \textup{(V82.18)}--\textup{(V82.19)}
gives
\[
 \prod_{m=0}^{M-1}
 \frac{|z+hm|}{|z+hm-x|}
 \asymp_{x,z,h}M^{x/h},
\]
proving \textup{(V86.18)}.
\end{proof}

The base \(4x/A\) is uniformly larger than four in the V85 convergence
regime \(A<d<x\).  Thus compacting the exponent band changes V82's
polynomial quotient loss into an exponential one.

\subsection{Channel count with endpoint regularization}

\begin{lemma}[One channel per zero and endpoint moment]
\label{lem:v86-channel-count}
After consolidating repeated exponents and deleting zero coefficients, let
\(N=J+1\) be the number of distinct radial exponents in \textup{(V86.11)}.
Under \textup{(V86.12)}, it satisfies
\[
 \boxed{N\ge M+\nu+2.}
 \tag{V86.23}
\]
\end{lemma}

\begin{proof}
The rational function
\[
 F_K(y):=\sum_{j=0}^{J}\frac{b_j}{y+a_j}
 \tag{V86.24}
\]
has denominator of degree \(N\).  Ordinarily its numerator has degree at
most \(N-1\).  The moment cancellations
\(\sum_jb_ja_j^r=0\), \(0\le r\le\nu\), cancel the first
\(\nu+1\) terms in its expansion at infinity, so the numerator degree is at
most \(N-\nu-2\).  It has the \(M\) distinct zeros \(y_m\), but it is not
identically zero because \(F_K(x)=1\).  Hence
\[
 M\le N-\nu-2,
\]
which is \textup{(V86.23)}.
\end{proof}

The lower bounds \textup{(V86.18)} and \textup{(V86.23)} are independent:
one controls synthesis variation, the other algebraic complexity.

\subsection{Exact tail gain from the endpoint jet}

\begin{lemma}[Endpoint-jet decay of the Laplace response]
\label{lem:v86-jet-tail}
If \(K\) satisfies the first conditions in \textup{(V86.12)}, then for
\(|y|>A\),
\[
 \boxed{
 F_K(y)=
 \sum_j
 \frac{b_j(-a_j)^{\nu+1}}
 {y^{\nu+1}(y+a_j)},\qquad
 |F_K(y)|
 \le
 \frac{B(K)A^{\nu+1}}
 {|y|^{\nu+1}(|y|-A)}.}
 \tag{V86.25}
\]
\end{lemma}

\begin{proof}
The finite geometric identity gives
\[
 \frac1{y+a}
 =
 \sum_{r=0}^{\nu}\frac{(-a)^r}{y^{r+1}}
 +\frac{(-a)^{\nu+1}}{y^{\nu+1}(y+a)}.
 \tag{V86.26}
\]
Sum in \(j\).  Every term in the finite sum vanishes by the endpoint
moments, leaving \textup{(V86.25)}.
\end{proof}

\begin{corollary}[Regularized gamma and zero tail envelopes]
\label{cor:v86-regularized-tails}
For \(M\) large enough that \(\Re z+hM>A\),
\[
 \boxed{
 \sum_{m\ge M}|F_K(z+hm)|
 \ll_h
 B(K)A^{\nu+1}
 \left\{
 \frac1{(\Re z+hM-A)^{\nu+2}}
 +\frac1{(\nu+1)(\Re z+hM-A)^{\nu+1}}
 \right\}.}
 \tag{V86.27}
\]
For a Dirichlet zero channel centered at ordinate \(t\), and
\(T\ge2A+2\),
\[
 \boxed{
 \sum_{|\gamma_\rho-t|\ge T}
 m_\rho|F_K(s-\rho)|
 \ll_\chi
 B(K)A^{\nu+1}
 T^{-\nu-1}\log(2+|t|+T).}
 \tag{V86.28}
\]
\end{corollary}

\begin{proof}
Apply \textup{(V86.25)} to the gamma arithmetic progression.  Isolate its
first term and compare the remainder of the resulting power \(\nu+2\) sum
with its integral.  This proves
\textup{(V86.27)}.  For \(|\gamma_\rho-t|\ge T\), one has
\(|s-\rho|\ge|\gamma_\rho-t|\) and
\(|s-\rho|-A\ge|\gamma_\rho-t|/2\).  Sum the bound
\[
 |F_K(s-\rho)|
 \le 2B(K)A^{\nu+1}|\gamma_\rho-t|^{-\nu-2}
\]
over dyadic ordinate shells and use the local count
\textup{(V82.30)}.  The same calculation as Lemma
\ref{lem:v82-zero-power-tail} gives \textup{(V86.28)}.
\end{proof}

\subsection{The growing-jet phase boundary}

The right side of \textup{(V86.27)} is the standard coefficientwise gamma
tail budget.  Combining it with \textup{(V86.18)} shows how large the
endpoint jet must be before that budget can possibly become small.

\begin{theorem}[A fixed or slowly growing endpoint jet cannot close]
\label{thm:v86-jet-phase-boundary}
Fix \(0<A<x<\Re z\) and \(h>0\).  Let \(M\to\infty\) and let
\(\nu=\nu_M\ge0\).  For every normalized kernel satisfying
\textup{(V86.12)}, define the numerical gamma-tail envelope
\[
 \mathfrak B_{\rm gam}^{(\nu)}(K;M)
 :=B(K)A^{\nu+1}
 \left\{
 \frac1{(\Re z+hM-A)^{\nu+2}}
 +\frac1{(\nu+1)(\Re z+hM-A)^{\nu+1}}
 \right\}.
 \tag{V86.29}
\]
If
\[
 \nu_M=o\!\left(\frac{M}{\log M}\right),
 \tag{V86.30}
\]
then there is \(c_{A,x}>0\) such that
\[
 \boxed{
 \mathfrak B_{\rm gam}^{(\nu_M)}(K;M)
 \ge e^{c_{A,x}M}}
 \tag{V86.31}
\]
for all sufficiently large \(M\), uniformly over all such kernels.
More precisely,
\[
 \log\mathfrak B_{\rm gam}^{(\nu)}(K;M)
 \ge
 M\log\frac{4x}{A}
 -\nu\log M
 -O_{A,x,z,h}(\nu+\log M).
 \tag{V86.32}
\]
Therefore a necessary condition for the coefficientwise gamma envelope not
to be forced exponentially large is
\[
 \boxed{
 \nu_M\ge
 \left(\log\frac{4x}{A}+o(1)\right)
 \frac{M}{\log M}.}
 \tag{V86.33}
\]
\end{theorem}

\begin{proof}
Insert \textup{(V86.18)} into the second, positive summand of
\textup{(V86.29)}.  Since \(\Re z+hM-A\asymp_h M\),
\[
 \begin{aligned}
 \log\mathfrak B_{\rm gam}^{(\nu)}(K;M)
 \ge{}&
 (M+\nu+1)\log\frac{4x}{A}
 +\frac{x}{h}\log M\\
 &+(\nu+1)\log A
 -(\nu+1)\log(hM+O_{z,A}(1))
 -\log(\nu+1)+O(1).
 \end{aligned}
\]
Collecting the terms gives \textup{(V86.32)}.  Under
\textup{(V86.30)}, its negative term is \(o(M)\), while
\(\log(4x/A)>0\), proving \textup{(V86.31)}.  Rearranging the leading terms
gives the necessary condition \textup{(V86.33)}.
\end{proof}

The theorem concerns the absolute envelope in
\textup{(V86.27)}, not the signed actual gamma tail.  A jet satisfying
\textup{(V86.33)} is only an escape from this lower obstruction; it is not
a construction with a small tail.  Theorem
\ref{thm:v86-jet-interpolation} proves qualitative feasibility for every
finite pair \((M,\nu)\), but gives no coefficient upper bound in the
growing regime.

\begin{theorem}[V86 bounded-band verdict]
\label{thm:v86-verdict}
The compact-band superresolution problem has the following exact outcome.
\begin{enumerate}[label=\textup{(\roman*)},leftmargin=3em]
\item Every finite endpoint jet is uniformly invisible and is compatible
      with arbitrary finite tail-positive spectral interpolation.
\item Cancelling \(M\) gamma nodes while killing \(\nu+1\) endpoint moments
      requires at least \(M+\nu+2\) radial channels.
\item The endpoint-constrained generalized Cauchy quotient has the explicit
      Chebyshev bound \textup{(V86.16)}.
\item Normalized coefficient variation is at least
      \((4x/A)^{M+\nu+1}M^{x/h}\), up to a fixed constant.
\item The endpoint jet improves distant gamma and zero tails by
      \(\nu+1\) powers, but a fixed or
      \(o(M/\log M)\) jet leaves the coefficientwise gamma budget
      exponentially large.
\end{enumerate}
Thus V85's strict convergence margin does not produce a stable scalar
certificate.  The only remaining coefficientwise escape uses a growing
endpoint jet at least on the scale \textup{(V86.33)}.
\end{theorem}

\begin{firewall}[Superresolution cost is not spectral orientation]
The Chebyshev theorem is a lower bound on coefficient variation and on the
corresponding absolute tail envelope.  It does not prove that the actual
signed gamma or zero tail is large, nor does it preclude cancellations
available with a growing endpoint jet.  The compensation law remains
consistent, and GRH remains open.
\end{firewall}

\begin{assessment}[V86 closure and V87 direction]
V86 resolves the bounded-band quotient requested after V85 and replaces the
old qualitative conditioning warning by an explicit exponential constant.
It also identifies the sole scalar corridor not closed by the present
absolute estimates: simultaneous growth of the cancelled gamma prefix and
the endpoint vanishing jet.

V87 should optimize that corridor.  Set
\(\nu\asymp M/\log M\), retain the exact Chebyshev quotient rather than its
coarse asymptotic, and compare the regularized gamma tail with the two
actual zero tails at a height \(T=T(M)\).  The decisive question is whether
one can prove a matching upper synthesis construction in
\(\mathcal E_A^{(\nu)}\).  If every construction still pays a larger
exponent than the tail gain, the scalar programme closes.  If a matching
construction exists, its signed finite-window response should be inserted
directly into the V80 compensation law to determine where the actual
spectral mass must escape.
\end{assessment}

\section*{Version V86 append-only record}
\addcontentsline{toc}{section}{Version V86 append-only record}
The complete V85 source was retained before this continuation.  It had
1231631 bytes, 36170 lines, and SHA--256
\[
 \texttt{DA422827CACACABA9909EB28A157CBD8C3EB81E4E513C6B25A9AF7FA8B604200}.
\]
V86 proved finite endpoint-jet density and interpolation, the exact
Chebyshev bounded-band quotient bound, exponential coefficient and channel
costs, regularized gamma and zero tail estimates, and the
\(M/\log M\) endpoint-jet phase boundary.

% =============================================================================
% END APPENDED V86 MATERIAL
% =============================================================================
% =============================================================================
% BEGIN APPENDED V87 MATERIAL
% =============================================================================

\section{V87: exact rational sign repair and superexponential signed-tail escape}
\label{sec:grh-v87}

V86 measured a compact-band interpolation by coefficient variation.  That
quantity is necessarily exponential, but it enters the spectral remainder
only after a triangle inequality.  This continuation computes the signed
rational response itself.  The result sharply separates three effects which
the coefficient envelope conflates.

First, endpoint moments and finitely many response zeros determine the
numerator of the Cauchy transform.  The principal-pole zero \(d<x\) then
forces the slowest residue of a minimally supported response to have the
wrong sign.  One additional radial channel is both necessary and sufficient
to repair that sign: its one free numerator zero may be placed in
\((0,x)\).  Second, after this repair, the actual uncancelled gamma-ladder
tail is superexponentially small, even though the coefficient variation is
at least exponentially ill-conditioned.  Third, the repaired inverse Laplace
kernel has at least one sign change for every cancelled positive gamma node.
Euler admissibility therefore requires an increasing oscillation block to be
compressed below the fixed first-prime threshold.  This last geometric
compression, together with conditioning at the actual zero nodes, is the
remaining obstruction.

\subsection{Rational factorization and the wrong minimal sign}

Fix
\[
 0<d<x,\qquad A>0,\qquad \nu\ge0.
 \tag{V87.1}
\]
Let
\[
 0\le a_0<a_1<\cdots<a_{N-1}<A,\qquad
 K(u)=\sum_{j=0}^{N-1}b_j e^{-a_j u},
 \tag{V87.2}
\]
where every \(b_j\ne0\), and put
\[
 F_K(y):=\int_0^\infty e^{-yu}K(u)\,du
 =\sum_{j=0}^{N-1}\frac{b_j}{y+a_j}.
 \tag{V87.3}
\]
Let \(\mathcal Y=\{y_1,\ldots,y_Q\}\) be a finite set, closed under
conjugation, contained in \(\Re y>0\), disjoint from \(d,x\), and such that
every real member of \(\mathcal Y\) is larger than \(x\).

\begin{lemma}[The principal zero gives the minimal response the wrong tail]
\label{lem:v87-wrong-minimal-sign}
Suppose the coefficients in \textup{(V87.2)} are real and
\[
 \sum_jb_ja_j^r=0\quad(0\le r\le\nu),\qquad
 F_K(d)=0,\qquad F_K(x)=1,\qquad
 F_K(y_q)=0\quad(1\le q\le Q).
 \tag{V87.4}
\]
Then
\[
 \boxed{N\ge Q+\nu+3.}
 \tag{V87.5}
\]
If equality holds, \(b_0<0\), and consequently \(K(u)<0\) for all
sufficiently large \(u\).  Hence every eventually positive solution of
\textup{(V87.4)} satisfies the sharper bound
\[
 \boxed{N\ge Q+\nu+4.}
 \tag{V87.6}
\]
\end{lemma}

\begin{proof}
Write
\[
 D(y):=\prod_{j=0}^{N-1}(y+a_j),\qquad
 F_K(y)=\frac{R(y)}{D(y)}.
\]
Expansion at infinity gives
\[
 F_K(y)
 =\sum_{r\ge0}\frac{(-1)^r\sum_jb_ja_j^r}{y^{r+1}},
 \qquad |y|>A.
 \tag{V87.7}
\]
The endpoint moments therefore imply
\(\deg R\le N-\nu-2\).  The \(Q+1\) distinct forced zeros
\(d,y_1,\ldots,y_Q\) divide \(R\), while \(R(x)\ne0\).  This proves
\textup{(V87.5)}.

If equality holds, normalization forces
\[
 R(y)=
 \frac{D(x)}{(x-d)\prod_q(x-y_q)}
 (y-d)\prod_q(y-y_q).
 \tag{V87.8}
\]
The slowest residue is
\[
 b_0=\frac{R(-a_0)}
 {\prod_{j=1}^{N-1}(a_j-a_0)}.
 \tag{V87.9}
\]
The denominator is positive.  In the quotient \(R(-a_0)/R(x)\), the
factor from \(d\) is
\[
 \frac{-a_0-d}{x-d}<0.
\]
Every nonreal conjugate pair from \(\mathcal Y\) contributes a positive
modulus square, and every real \(y_q>x\) contributes
\(( -a_0-y_q)/(x-y_q)>0\).  Thus \(b_0<0\).  Since \(a_0\) is the unique
smallest exponent,
\[
 e^{a_0u}K(u)\longrightarrow b_0,
\]
so \(K\) is eventually negative.  An eventually positive solution cannot
have equality in \textup{(V87.5)}, proving \textup{(V87.6)}.
\end{proof}

The extra channel in \textup{(V87.6)} is not an artifact of the proof.

\begin{theorem}[One free numerator zero repairs the tail sign]
\label{thm:v87-one-zero-repair}
Choose any \(\xi\in(0,x)\setminus\{d\}\), any
\[
 N=Q+\nu+4
\]
distinct exponents satisfying \textup{(V87.2)}, and define
\[
 \begin{aligned}
 P(y)&:=(y-d)(y-\xi)\prod_{q=1}^{Q}(y-y_q),\\
 D(y)&:=\prod_{j=0}^{N-1}(y+a_j),\\
 F(y)&:=\frac{D(x)}{P(x)}\frac{P(y)}{D(y)}.
 \end{aligned}
 \tag{V87.10}
\]
Let
\[
 F(y)=\sum_{j=0}^{N-1}\frac{b_j}{y+a_j},
 \qquad
 K(u):=\sum_{j=0}^{N-1}b_j e^{-a_j u}
 \tag{V87.11}
\]
be its partial-fraction and inverse-Laplace expansions.  Then every \(b_j\)
is real and nonzero,
\[
 \boxed{
 \sum_jb_ja_j^r=0\ (0\le r\le\nu),\quad
 F(d)=F(\xi)=0,\quad F(x)=1,\quad
 F(y_q)=0,}
 \tag{V87.12}
\]
and \(b_0>0\).  In particular \(K(u)>0\) for all sufficiently large
\(u\).  Thus the channel bound \textup{(V87.6)} is sharp for eventual
positivity.
\end{theorem}

\begin{proof}
The set of numerator zeros is conjugation closed, so \(P,D,F\) have real
coefficients.  All poles are simple and none is a numerator zero, hence the
residues in \textup{(V87.11)} are real and nonzero.  Since
\[
 \deg D-\deg P=\nu+2,
\]
comparison with \textup{(V87.7)} gives the endpoint moments.  The remaining
identities in \textup{(V87.12)} are immediate from
\textup{(V87.10)}.

For the slowest residue, the factor from \(d\) has negative quotient sign
between \(-a_0\) and \(x\), and the factor from \(\xi\) has the same negative
sign.  Their product is positive.  Every factor from \(\mathcal Y\) has
positive quotient sign exactly as in Lemma
\ref{lem:v87-wrong-minimal-sign}.  Thus \(b_0>0\), and
\(e^{a_0u}K(u)\to b_0\).
\end{proof}

\begin{corollary}[A free response zero is necessary for sign repair]
\label{cor:v87-free-zero}
Under the hypotheses of Lemma \ref{lem:v87-wrong-minimal-sign}, write
\[
 R(y)=C(y-d)\prod_q(y-y_q)S(y).
 \tag{V87.13}
\]
If \(K\) is eventually positive, then
\[
 \deg S\ge1,\qquad S(-a_0)S(x)<0.
 \tag{V87.14}
\]
Consequently \(S\) has a real zero in \((-a_0,x)\).
\end{corollary}

\begin{proof}
The sign computation in Lemma \ref{lem:v87-wrong-minimal-sign} shows that
the forced factors alone contribute a negative quotient between
\(-a_0\) and \(x\).  Eventual positivity is equivalent to \(b_0>0\), so
the remaining factor must contribute a negative quotient.  This is
\textup{(V87.14)}, and the intermediate value theorem supplies the zero.
\end{proof}

\subsection{A coefficient-free bound for the actual gamma tail}

The repaired response can now be evaluated without expanding its ill
conditioned residues.  Fix \(R\ge1\), \(h>0\), and starts
\[
 z_1,\ldots,z_R,\qquad \Re z_\ell>x,
 \tag{V87.15}
\]
whose set is closed under conjugation and whose arithmetic progressions are
pairwise disjoint.  Real starts are therefore larger than \(x\).  Put
\[
 \mathcal Y_M
 :=\{z_\ell+hm:1\le\ell\le R,\ 0\le m<M\}.
 \tag{V87.16}
\]
For any \(\nu=\nu_M\ge0\), choose arbitrary distinct
\[
 0\le a_0<\cdots<a_{N-1}<A,\qquad
 N=RM+\nu+4,
 \tag{V87.17}
\]
and let \(F_{M,\nu}\) and \(K_{M,\nu}\) be the construction of Theorem
\ref{thm:v87-one-zero-repair} with
\(\mathcal Y=\mathcal Y_M\).

\begin{theorem}[The signed omitted gamma tail is superexponentially small]
\label{thm:v87-signed-gamma-tail}
There are constants \(C_0,C_1,B_0>0\), depending only on the fixed data in
\textup{(V87.1)}, \textup{(V87.15)}, and on \(\xi\), such that for every
\(M\ge2\), every \(\nu\ge0\), every exponent set
\textup{(V87.17)}, and every \(1\le\ell\le R\),
\[
 \boxed{
 \sum_{k\ge0}
 \left|F_{M,\nu}\bigl(z_\ell+h(M+k)\bigr)\right|
 \le
 C_0 C_1^{RM+\nu}
 M^{B_0-RM-\nu-1}.}
 \tag{V87.18}
\]
In particular,
\[
 \boxed{
 \log\sum_{k\ge0}
 \left|F_{M,\nu}\bigl(z_\ell+h(M+k)\bigr)\right|
 \le
 -(RM+\nu)\log M+O(M+\nu+\log M).}
 \tag{V87.19}
\]
Thus the actual signed response beyond the cancelled gamma prefixes tends
to zero superexponentially for every choice of \(\nu_M\ge0\).  This estimate
contains no coefficient variation \(B(K)\).
\end{theorem}

\begin{proof}
Write \(y_{\ell,k}=z_\ell+h(M+k)\).  From
\textup{(V87.10)},
\[
 \begin{aligned}
 |F_{M,\nu}(y_{\ell,k})|
 ={}&
 \left|\frac{(y_{\ell,k}-d)(y_{\ell,k}-\xi)}
 {(x-d)(x-\xi)}\right|
 \prod_{r=1}^{R}\prod_{m=0}^{M-1}
 \frac{|y_{\ell,k}-z_r-hm|}{|x-z_r-hm|}\\
 &{}\times
 \prod_{j=0}^{N-1}
 \frac{|x+a_j|}{|y_{\ell,k}+a_j|}.
 \end{aligned}
 \tag{V87.20}
\]
Let
\[
 c:=\min_r\frac{\Re z_r-x}{h}>0.
\]
Because \(M+k-m\ge1\), the fixed differences \(z_\ell-z_r\) may be
absorbed into a constant, giving
\[
 \prod_{m=0}^{M-1}
 \frac{|y_{\ell,k}-z_r-hm|}{|x-z_r-hm|}
 \le
 C^M\frac{(k+1)_M}{(c)_M}.
 \tag{V87.21}
\]
The gamma-ratio identity implies, for a fixed \(B\),
\[
 \frac{(k+1)_M}{(c)_M}
 \le C M^B {M+k\choose M}.
 \tag{V87.22}
\]
Also,
\[
 \left|\frac{(y_{\ell,k}-d)(y_{\ell,k}-\xi)}
 {(x-d)(x-\xi)}\right|
 \le C(M+k+1)^2
 \tag{V87.23}
\]
and the bounded exponent band gives
\[
 \prod_{j=0}^{N-1}
 \frac{|x+a_j|}{|y_{\ell,k}+a_j|}
 \le
 \left(\frac{C}{M+k+1}\right)^{RM+\nu+4}.
 \tag{V87.24}
\]
Finally,
\[
 {M+k\choose M}
 \le\left(\frac{e(M+k+1)}{M}\right)^M.
 \tag{V87.25}
\]
Insert \textup{(V87.21)}--\textup{(V87.25)} into
\textup{(V87.20)}.  After enlarging fixed constants,
\[
 |F_{M,\nu}(y_{\ell,k})|
 \le
 C_0C_1^{RM+\nu}
 M^{B_0-RM}(M+k+1)^{-\nu-2}.
 \tag{V87.26}
\]
The elementary integral comparison
\[
 \sum_{k\ge0}(M+k+1)^{-\nu-2}
 \le C M^{-\nu-1}
 \tag{V87.27}
\]
proves \textup{(V87.18)}.  Taking logarithms gives
\textup{(V87.19)}.
\end{proof}

The point is not that the residues are small.  V86 proves that their total
variation is exponentially large.  Equation \textup{(V87.18)} proves that
the fixed rational product cancels those large residues before the omitted
gamma values are estimated.

\subsection{The actual distant-zero tail}

The same product representation improves the coefficientwise zero estimate
outside the scale of the cancelled gamma block.

\begin{corollary}[Superexponential distant-zero response]
\label{cor:v87-distant-zero-tail}
Fix \(s=\sigma+it\) with \(\sigma>1\).  There are fixed constants
\(C_2,C_3>0\) such that, whenever \(T\ge C_2M\),
\[
 \boxed{
 \sum_{|\gamma_\rho-t|\ge T}
 m_\rho\,|F_{M,\nu}(s-\rho)|
 \ll_\chi
 \frac{C_3^{RM+\nu}}{((M-1)!)^R}
 T^{-\nu-1}\log(2+|t|+T).}
 \tag{V87.28}
\]
Consequently, for \(T\ge C_2M\),
\[
 \log\left(
 \sum_{|\gamma_\rho-t|\ge T}
 m_\rho\,|F_{M,\nu}(s-\rho)|
 \right)
 \le
 -RM\log M-(\nu+1)\log T
 O(M+\nu+\log\log(3+T)).
 \tag{V87.29}
\]
\end{corollary}

\begin{proof}
At the normalization point,
\[
 \prod_{r=1}^{R}\prod_{m=0}^{M-1}|x-z_r-hm|
 \ge C^{-RM}((M-1)!)^R.
 \tag{V87.30}
\]
The numerator \(D(x)\) is at most \((x+A)^{RM+\nu+4}\); the two remaining
normalizing factors are fixed and nonzero.  Hence
\[
 \left|\frac{D(x)}{P(x)}\right|
 \le
 \frac{C^{RM+\nu}}{((M-1)!)^R}.
 \tag{V87.31}
\]
All zeros of \(P\) have modulus \(O(M)\), while all poles lie in
\([-A,0]\).  If \(T\ge C_2M\) and
\(|\gamma_\rho-t|\ge T\), fixed-factor comparison in the product gives
\[
 \left|\frac{P(s-\rho)}{D(s-\rho)}\right|
 \le C^{RM+\nu}|s-\rho|^{-\nu-2}.
 \tag{V87.32}
\]
Combine \textup{(V87.31)}--\textup{(V87.32)}, enlarge \(C\), and use
the local zero-counting tail estimate from Lemma
\ref{lem:v82-zero-power-tail}, with power \(\nu+2\).  This proves
\textup{(V87.28)}.  Stirling's formula proves
\textup{(V87.29)}.
\end{proof}

\subsection{Euler admissibility compresses an increasing oscillation block}

The repaired response is eventually positive, but Theorem
\ref{thm:v87-one-zero-repair} does not place the beginning of that positive
tail.  The real principal gamma ladder shows why this distinction is
substantial.

\begin{lemma}[Laplace zeros require time-side sign changes]
\label{lem:v87-laplace-variation}
Let \(K\) be a nonzero real exponential polynomial and suppose
\[
 F(y)=\int_0^\infty e^{-yu}K(u)\,du
\]
has \(q\) distinct zeros on the positive real axis.  Then \(K\) has at
least \(q\) sign changes on \((0,\infty)\).
\end{lemma}

\begin{proof}
Suppose instead that \(K\) has \(s<q\) sign changes, at
\(0<u_1<\cdots<u_s\).  Choose \(s+1\) distinct positive zeros
\(\lambda_0,\ldots,\lambda_s\) of \(F\).  The strict Chebyshev system
\(\{e^{-\lambda_j u}\}_{j=0}^{s}\) contains a nonzero linear combination
\[
 G(u)=\sum_{j=0}^{s}c_je^{-\lambda_j u}
\]
vanishing at the \(s\) prescribed points.  A member of this system has at
most \(s\) real zeros counted with multiplicity; hence these zeros are
simple, there are no others, and the sign of \(G\) alternates exactly where
the sign of \(K\) alternates.  After multiplying \(G\) by \(-1\) if
necessary, \(K(u)G(u)\ge0\) and is positive on a nonempty interval.
Therefore
\[
 0<\int_0^\infty K(u)G(u)\,du
 =\sum_{j=0}^{s}c_jF(\lambda_j)=0,
\]
a contradiction.
\end{proof}

\begin{theorem}[Sub-prime oscillation compression is necessary]
\label{thm:v87-oscillation-compression}
Assume one of the starts in \textup{(V87.15)}, say \(z_1\), is real.  Then
the repaired kernel \(K_{M,\nu}\) has at least \(M+2\) sign changes, because
its transform vanishes at
\[
 d,\quad \xi,\quad z_1,\quad z_1+h,\quad\ldots,\quad z_1+h(M-1).
 \tag{V87.33}
\]
If it is Euler admissible in the V85 sense,
\[
 K_{M,\nu}(u)>0\qquad(u\ge L_f),
 \tag{V87.34}
\]
all \(M+2\) sign changes must lie in the fixed interval \((0,L_f)\).
\end{theorem}

\begin{proof}
The displayed values are distinct positive zeros of \(F_{M,\nu}\).
Lemma~\ref{lem:v87-laplace-variation} gives the sign-change count.
Condition \textup{(V87.34)} leaves no sign change at or beyond \(L_f\).
\end{proof}

This theorem turns the uncontrolled phrase \emph{eventually positive} into
an exact missing estimate.  The V85 density theorem proves that such
compression is qualitatively possible for every fixed \(M\), but gives no
uniform coefficient, exponent-separation, or tail-response control as
\(M\to\infty\).

\subsection{Where the actual zero block re-enters}

Let \(\mathcal Z=\{w_1,\ldots,w_S\}\) be an additional finite,
conjugation-closed response-node set, none equal to \(x\).  The construction
extends verbatim by multiplying \(P\) by
\(\prod_q(y-w_q)\) and adding \(S\) poles.  Its normalization contains the
exact factor
\[
 \boxed{
 \Delta_x(\mathcal Z)^{-1},
 \qquad
 \Delta_x(\mathcal Z):=\prod_{q=1}^{S}|x-w_q|.}
 \tag{V87.35}
\]
Thus cancellation of an actual zero block preserves the rational
factorization, but every tail estimate inherits the reciprocal target-node
separation product.  In the target-character channel,
\[
 x-w_q=\rho_q-\rho_\star.
 \tag{V87.36}
\]
Classical zero counting bounds the number of factors in
\textup{(V87.35)} but gives no lower bound for their product.  Multiplicity
at \(\rho_\star\) belongs to the selected target contribution; arbitrarily
near distinct zeros remain a genuine conditioning problem.

\begin{theorem}[V87 exact-response verdict]
\label{thm:v87-verdict}
For the bounded-band signed-order programme:
\begin{enumerate}[label=\textup{(\roman*)},leftmargin=3em]
\item endpoint moments and prescribed Cauchy zeros give an exact rational
      numerator factorization;
\item principal-pole cancellation makes the minimally supported inverse
      response eventually negative;
\item one additional channel and one auxiliary zero in \((0,x)\) repair
      that sign, and this support count is sharp for eventual positivity;
\item the repaired actual gamma tail and the zero tail beyond \(T\gg M\)
      are superexponentially small without using coefficientwise absolute
      values;
\item cancelling the real principal gamma prefix forces at least \(M+2\)
      time-side sign changes, all of which must lie below \(L_f\) for Euler
      admissibility;
\item cancellation of the actual compact zero block introduces the
      uncontrolled separation product \textup{(V87.35)}.
\end{enumerate}
Therefore the V86 coefficient envelope is not an obstruction to the signed
gamma or distant-zero tails.  The live scalar problem is now the joint
control of fixed-threshold oscillation compression and the target-zero
separation determinant.
\end{theorem}

\begin{firewall}[Eventual positivity is not Euler positivity]
The explicit repaired kernel is positive after a construction-dependent
threshold.  The prime-side argument requires positivity already at the fixed
frequency \(L_f\).  Theorem \ref{thm:v87-oscillation-compression} proves that
this demands \(M+2\) sign changes below that threshold, but does not prove
that the demand is impossible.  Nor does zero counting control
\(\Delta_x(\mathcal Z)\).  Hence the superexponential signed-tail estimate
does not contradict the V80 compensation law, and GRH remains open.
\end{firewall}

\begin{assessment}[V87 closure and V88 direction]
V87 closes the coefficientwise interpretation of the growing-jet question
posed after V86 on the exact gamma-response side: coefficient growth is
canceled by the rational product, so increasing the endpoint jet is not
needed to make this explicit signed omitted gamma response small.  The same
product controls zeros beyond a linearly growing ordinate window.  The
fixed-threshold Euler-positivity requirement remains open.

V88 should attack the newly isolated fixed-interval problem.  A useful
quantitative theorem would bound the cost of placing \(M+2\) sign changes of
an exponent-band polynomial inside \((0,L_f)\) while retaining
\(F(d)=0\), \(F(x)=1\), and the endpoint jet.  The natural variables are the
clustered pole gaps and the divided-difference representation of the inverse
Laplace transform.  If that compression forces a loss larger than the exact
product saving in \textup{(V87.18)}, the scalar route closes.  If a matching
compressed construction exists, it should be inserted into the V80 identity;
then the only remaining escape is the actual local zero-separation factor
\textup{(V87.35)}, which must be treated by an orientation-sensitive block
identity rather than zero counting.
\end{assessment}

\section*{Version V87 append-only record}
\addcontentsline{toc}{section}{Version V87 append-only record}
The complete V86 source was retained before this continuation.  It had
1249414 bytes, 36745 lines, and SHA--256
\[
 \texttt{8142369D68B181B4077723A032A2F3CA30741CC7E4C4A472E4BA79F20727FADA}.
\]
V87 proved the sharp extra-channel sign repair, constructed the exact
eventually positive rational response, bounded its signed gamma and distant
zero tails superexponentially, and isolated fixed-threshold sign-change
compression and the target-zero separation product as the surviving
obstructions.

% =============================================================================
% END APPENDED V87 MATERIAL
% =============================================================================
% =============================================================================
% BEGIN APPENDED V88 MATERIAL
% =============================================================================

\section{V88: moving-line compression and the exact gamma phase transition}
\label{sec:grh-v88}

V87 showed that a single free numerator zero repairs eventual positivity and
that, on a fixed evaluation line, the resulting signed gamma tail is
superexponentially small.  Eventual positivity is not enough for the Euler
side: all sign changes must occur before the fixed first-prime frequency
\(L_f\).  The evaluation line \(\sigma\) is itself free, however.  Moving it
right enlarges the admissible exponent band \(A<\sigma-1\) and compresses
the inverse-Laplace oscillations.

This continuation makes that idea quantitative for an exact, equally spaced
pole construction.  The residues alternate, and their adjacent ratios give
an explicit Euler-positivity threshold.  The same movement of \(\sigma\)
weakens the rational gamma cancellation.  At the linear scale
\(\sigma-1\sim cM\), the response at the second omitted gamma block has an
exact large-deviation rate with one critical constant.  The elementary
positivity certificate uses the larger scale \(M\log M\), where that gamma
response grows exponentially.  Thus merely moving far to the right does not
close the programme.  The remaining question becomes the smallest positive
zero of one explicit generalized-Eulerian polynomial at the linear scale.

\subsection{The moving two-channel parameters}

Fix a hypothetical selected zero
\[
 \rho_\star=\beta+i\gamma,\qquad \frac12<\beta<1,
 \tag{V88.1}
\]
a number \(\eta\) with
\[
 \beta<\eta<1,
 \tag{V88.2}
\]
and a band fraction \(0<\theta<1\).  For \(d>0\), put
\[
 \sigma=d+1,\qquad
 x_d=d+1-\beta,\qquad
 \xi_d=d+1-\eta.
 \tag{V88.3}
\]
Then
\[
 0<d<\xi_d<x_d.
 \tag{V88.4}
\]

To include the principal and target-character gamma ladders uniformly, fix
a conjugation-closed set of offsets
\[
 \zeta_1,\ldots,\zeta_R,\qquad
 \Re\zeta_\ell>-\beta,
 \tag{V88.5}
\]
containing at least one real member.  Assume the progressions below are
pairwise disjoint, and set
\[
 z_{\ell,d}:=d+1+\zeta_\ell,\qquad
 \mathcal Y_{M,d}:=
 \{z_{\ell,d}+hm:1\le\ell\le R,\ 0\le m<M\}.
 \tag{V88.6}
\]
For the V80 application one takes the distinct members generated by
\(\zeta=\kappa_0\) and
\(\zeta=\kappa_\chi\pm i\gamma\), with \(h=2\).

Set
\[
 N:=RM+4,\qquad n:=N-1=RM+3,\qquad
 \varepsilon:=\frac{\theta d}{N},
 \qquad a_j:=j\varepsilon\quad(0\le j\le n).
 \tag{V88.7}
\]
Thus
\[
 0=a_0<\cdots<a_n
 =\frac{n}{N}\theta d<\theta d<d.
 \tag{V88.8}
\]
Define
\[
 \begin{aligned}
 P_{M,d}(y)
 &:=(y-d)(y-\xi_d)
   \prod_{\ell=1}^{R}\prod_{m=0}^{M-1}
   (y-z_{\ell,d}-hm),\\
 D_{M,d}(y)&:=\prod_{j=0}^{n}(y+a_j),\\
 F_{M,d}(y)&:=
 \frac{D_{M,d}(x_d)}{P_{M,d}(x_d)}
 \frac{P_{M,d}(y)}{D_{M,d}(y)}.
 \end{aligned}
 \tag{V88.9}
\]
Let
\[
 F_{M,d}(y)=\sum_{j=0}^{n}\frac{b_j}{y+a_j},
 \qquad
 K_{M,d}(u):=\sum_{j=0}^{n}b_je^{-a_ju}.
 \tag{V88.10}
\]
Theorem \ref{thm:v87-one-zero-repair} applies: the coefficients are real,
\(b_0>0\), \(K_{M,d}(0)=0\), the pole response at \(d\) vanishes, the
target response at \(x_d\) is one, and the first \(M\) nodes of every
gamma ladder vanish.

\subsection{Alternating residues and an explicit positivity onset}

All roots of \(P_{M,d}\) are either positive real numbers or nonreal
conjugate pairs in the right half-plane.  Consequently
\[
 \frac{P_{M,d}(-a)}{P_{M,d}(0)}>0
 \qquad(a\ge0).
 \tag{V88.11}
\]

\begin{lemma}[Exact alternating-residue polynomial]
\label{lem:v88-alternating-polynomial}
Put
\[
 c_j:={n\choose j}
 \frac{P_{M,d}(-j\varepsilon)}{P_{M,d}(0)}>0.
 \tag{V88.12}
\]
Then
\[
 \boxed{
 \frac{b_j}{b_0}=(-1)^jc_j,\qquad
 \frac{K_{M,d}(u)}{b_0}
 =H_{M,d}(e^{-\varepsilon u}),}
 \tag{V88.13}
\]
where
\[
 H_{M,d}(q):=\sum_{j=0}^{n}(-1)^jc_jq^j.
 \tag{V88.14}
\]
If \(\Lambda_{M,d}\) is the multiset of roots of \(P_{M,d}\) and
\(\mathcal D_q=q\,d/dq\), then the same polynomial has the exact operator
representation
\[
 \boxed{
 H_{M,d}(q)
 =
 \prod_{\lambda\in\Lambda_{M,d}}
 \left(1+\frac{\varepsilon}{\lambda}\mathcal D_q\right)
 (1-q)^n.}
 \tag{V88.15}
\]
Conjugate operator factors are paired, so the right side is real.
\end{lemma}

\begin{proof}
Since
\[
 D_{M,d}'(-j\varepsilon)
 =\varepsilon^n(-1)^j j!(n-j)!,
\]
the residue formula gives
\[
 \frac{b_j}{b_0}
 =
 \frac{P_{M,d}(-j\varepsilon)}{P_{M,d}(0)}
 \frac{D_{M,d}'(0)}{D_{M,d}'(-j\varepsilon)}
 =
 (-1)^j{n\choose j}
 \frac{P_{M,d}(-j\varepsilon)}{P_{M,d}(0)}.
\]
Substitution into \textup{(V88.10)} proves
\textup{(V88.13)}--\textup{(V88.14)}.

Moreover,
\[
 \frac{P_{M,d}(-j\varepsilon)}{P_{M,d}(0)}
 =
 \prod_{\lambda\in\Lambda_{M,d}}
 \left(1+\frac{j\varepsilon}{\lambda}\right).
 \tag{V88.16}
\]
The operator \(\mathcal D_q\) multiplies the coefficient of \(q^j\) by
\(j\).  Apply the product of operators in \textup{(V88.15)} to
\[
 (1-q)^n=\sum_{j=0}^{n}(-1)^j{n\choose j}q^j
\]
and use \textup{(V88.16)}.
\end{proof}

The adjacent coefficients provide a completely explicit, though not sharp,
positive-tail threshold.

\begin{theorem}[A constructive \(M\log M\) Euler-positivity bound]
\label{thm:v88-explicit-onset}
For every \(0\le j<n\),
\[
 \boxed{
 \frac{c_{j+1}}{c_j}
 \le e^\theta\frac{n-j}{j+1}
 \le e^\theta n.}
 \tag{V88.17}
\]
Consequently,
\[
 \boxed{
 K_{M,d}(u)>0
 \quad\text{whenever}\quad
 u>\mathcal L_{M,d}:=
 \frac{\log n+\theta}{\varepsilon}
 =
 \frac{N(\log n+\theta)}{\theta d}.}
 \tag{V88.18}
\]
In particular, the exact rational interpolant is Euler admissible,
\[
 K_{M,d}(u)>0\qquad(u\ge L_f),
 \tag{V88.19}
\]
whenever
\[
 \boxed{
 d>
 \frac{N(\log n+\theta)}{\theta L_f}.}
 \tag{V88.20}
\]
\end{theorem}

\begin{proof}
Every root \(\lambda\) of \(P_{M,d}\) has
\(\Re\lambda\ge d\).  Pairing conjugate roots when necessary,
\[
 \frac{P_{M,d}(-(j+1)\varepsilon)}
      {P_{M,d}(-j\varepsilon)}
 =
 \prod_{\lambda\in\Lambda_{M,d}}
 \left|
 1+\frac{\varepsilon}{\lambda+j\varepsilon}
 \right|
 \le
 \left(1+\frac{\varepsilon}{d}\right)^{RM+2}.
 \tag{V88.21}
\]
Here the displayed product of moduli is understood pairwise; real factors
are positive.  Since
\[
 (RM+2)\frac{\varepsilon}{d}
 =\frac{(RM+2)\theta}{RM+4}<\theta,
\]
\textup{(V88.21)} is at most \(e^\theta\).  Combining this with the
adjacent binomial ratio proves \textup{(V88.17)}.

Put \(q=e^{-\varepsilon u}\).  If
\[
 0<q<\frac{e^{-\theta}}{n},
 \tag{V88.22}
\]
then the positive term magnitudes \(c_jq^j\) in
\textup{(V88.14)} are strictly decreasing.  The finite alternating-series
test gives \(H_{M,d}(q)>0\).  Since \(b_0>0\), this proves
\textup{(V88.18)}.  Condition \textup{(V88.20)} makes
\(\mathcal L_{M,d}<L_f\), proving \textup{(V88.19)}.
\end{proof}

This is a quantitative construction, not merely a density argument.  It
simultaneously has a strict exponent margin, exact target and pole moments,
exact gamma-prefix zeros, and prime-frequency positivity.  Its line movement
is \(d\gg M\log M\).  The next calculation shows why that scale is too far
right for absolute gamma-tail control.

\subsection{The exact linear-scale gamma phase}

For a fixed ladder index \(\ell\), test the first point of the block after
one further block of length \(M\):
\[
 Y_{\ell,M,d}:=z_{\ell,d}+2hM.
 \tag{V88.23}
\]
This is an omitted gamma node, since only indices \(0,\ldots,M-1\) were
cancelled.

\begin{theorem}[Moving-line gamma phase transition]
\label{thm:v88-gamma-phase}
Suppose
\[
 \frac{d_M}{M}\longrightarrow c\in(0,\infty).
 \tag{V88.24}
\]
Then, uniformly for each fixed \(\ell\),
\[
 \boxed{
 \lim_{M\to\infty}\frac1M
 \log|F_{M,d_M}(Y_{\ell,M,d_M})|
 =
 \Psi_{\theta,h,R}(c),}
 \tag{V88.25}
\]
where
\[
 \boxed{
 \Psi_{\theta,h,R}(c)
 :=
 R\left[
 \log4+
 \int_0^1
 \log\frac{c(1+\theta t)}
 {c(1+\theta t)+2h}\,dt
 \right].}
 \tag{V88.26}
\]
The function in brackets is strictly increasing from \(-\infty\) to
\(\log4\).  Hence there is a unique
\[
 c_\star=c_\star(\theta,h)>0
 \tag{V88.27}
\]
at which \(\Psi_{\theta,h,R}(c_\star)=0\).  In particular, the single
omitted response in \textup{(V88.23)} decays exponentially if
\(c<c_\star\) and grows exponentially if \(c>c_\star\).
\end{theorem}

\begin{proof}
Use the exact product
\[
 \begin{aligned}
 |F_{M,d}(Y_{\ell,M,d})|
 ={}&
 \left|
 \frac{(Y_{\ell,M,d}-d)(Y_{\ell,M,d}-\xi_d)}
 {(x_d-d)(x_d-\xi_d)}
 \right|\\
 &{}\times
 \prod_{r=1}^{R}\prod_{m=0}^{M-1}
 \frac{|Y_{\ell,M,d}-z_{r,d}-hm|}
 {|x_d-z_{r,d}-hm|}\\
 &{}\times
 \prod_{j=0}^{n}
 \frac{|x_d+a_j|}{|Y_{\ell,M,d}+a_j|}.
 \end{aligned}
 \tag{V88.28}
\]
The first line has logarithm \(O(\log M)\).  For each \(r\), cancellation
of the common \(d+1\) shift gives
\[
 \prod_{m=0}^{M-1}
 \frac{|\zeta_\ell-\zeta_r+h(2M-m)|}
 {|\beta+\zeta_r+hm|}.
 \tag{V88.29}
\]
The gamma-ratio formula and Stirling's theorem give
\[
 \frac1M\log\textup{(V88.29)}\longrightarrow\log4.
 \tag{V88.30}
\]

For the last line of \textup{(V88.28)}, condition
\textup{(V88.24)} and the equally spaced poles give the Riemann sum
\[
 \begin{aligned}
 \frac1M\sum_{j=0}^{n}
 \log\frac{|x_{d_M}+a_j|}
 {|Y_{\ell,M,d_M}+a_j|}
 \longrightarrow
 R\int_0^1
 \log\frac{c(1+\theta t)}
 {c(1+\theta t)+2h}\,dt.
 \end{aligned}
 \tag{V88.31}
\]
Adding \textup{(V88.30)} over \(r=1,\ldots,R\) proves
\textup{(V88.25)}.

If
\[
 \phi(c):=\log4+
 \int_0^1
 \log\frac{c(1+\theta t)}
 {c(1+\theta t)+2h}\,dt,
\]
then
\[
 \phi'(c)=
 \int_0^1
 \frac{2h}
 {c\{c(1+\theta t)+2h\}}\,dt>0.
 \tag{V88.32}
\]
Also \(\phi(c)\to-\infty\) as \(c\downarrow0\), whereas
\(\phi(c)\to\log4\) as \(c\to\infty\).  This proves existence and
uniqueness of \(c_\star\) and the final assertions.
\end{proof}

The phase transition remains meaningful without assuming a finite linear
limit.

\begin{corollary}[Overshifting destroys absolute gamma-tail control]
\label{cor:v88-overshift}
If \(d_M/M\to\infty\), then
\[
 \boxed{
 \frac1M
 \log|F_{M,d_M}(Y_{\ell,M,d_M})|
 \longrightarrow R\log4.}
 \tag{V88.33}
\]
In particular, every sequence satisfying the explicit Euler-positivity
condition \textup{(V88.20)} has an omitted gamma response of size
\[
 \exp\{(R\log4+o(1))M\}.
 \tag{V88.34}
\]
\end{corollary}

\begin{proof}
In \textup{(V88.28)}, the gamma products still contribute
\(R\log4+o(1)\) after division by \(M\).  For the pole product, uniformly
in \(j\),
\[
 \log\frac{|x_d+a_j|}{|Y_{\ell,M,d}+a_j|}
 =O\!\left(\frac{M}{d}\right).
\]
There are \(N=O(M)\) terms, so its logarithm divided by \(M\) is
\(O(M/d)=o(1)\).  The two fixed repair factors again contribute
\(o(M)\).  This proves \textup{(V88.33)}.  Condition
\textup{(V88.20)} has \(d_M/M\gg\log M\), so
\textup{(V88.34)} follows.
\end{proof}

\subsection{The exact remaining root-location problem}

Theorem \ref{thm:v88-explicit-onset} used only adjacent absolute
coefficients and therefore paid the factor \(\log n\).  It did not use the
cancellations internal to \(H_{M,d}\).  The exact equivalence is
\[
 \boxed{
 K_{M,d}(u)>0\ (u\ge L_f)
 \quad\Longleftrightarrow\quad
 H_{M,d}(q)>0
 \quad(0<q\le e^{-\varepsilon L_f}).}
 \tag{V88.35}
\]
At the linear scale \(d\sim cM\),
\[
 \varepsilon\longrightarrow\frac{\theta c}{R},
 \qquad
 e^{-\varepsilon L_f}
 \longrightarrow
 \exp\!\left(-\frac{\theta cL_f}{R}\right).
 \tag{V88.36}
\]
Thus a proof that the smallest positive zero of
\textup{(V88.14)} stays uniformly above the right side of
\textup{(V88.36)} for some \(c<c_\star\) would produce an
Euler-admissible rational filter with a decaying second-block gamma
response.  Conversely, a root forced below that threshold for every
\(c\le c_\star\) would close this explicit equally spaced construction.
The operator formula \textup{(V88.15)} is the exact finite object on which
that decision rests.

\begin{theorem}[V88 moving-line verdict]
\label{thm:v88-verdict}
The moving-line version of the V87 filter has the following rigorous phase
diagram.
\begin{enumerate}[label=\textup{(\roman*)},leftmargin=3em]
\item Equally spaced poles give an exact alternating-residue polynomial and
      the differential-operator representation \textup{(V88.15)}.
\item Adjacent coefficient control proves Euler positivity once
      \(d\gg M\log M\).
\item At \(d\sim cM\), the second omitted gamma block has the exact rate
      \(\Psi_{\theta,h,R}(c)\), with one critical constant \(c_\star\).
\item Every \(M\log M\) overshift certified by the elementary positivity
      argument makes that omitted gamma response grow like \(4^{RM}\) on
      the exponential scale.
\item The undecided regime is therefore the linear line \(d\sim cM\) with
      \(c<c_\star\), where Euler positivity is exactly the root-location
      condition \textup{(V88.35)}.
\end{enumerate}
Thus line movement converts the qualitative fixed-threshold obstruction of
V87 into one explicit polynomial root problem and one explicit gamma phase;
it does not by itself prove GRH.
\end{theorem}

\begin{firewall}[A sufficient overshift is not a necessary one]
The logarithmic onset in \textup{(V88.18)} is obtained by forcing the
alternating terms to decrease individually.  Internal cancellation may make
\(H_{M,d}\) positive at a much larger \(q\).  Therefore
\textup{(V88.34)} rules out only this elementary \(M\log M\) certificate,
not every Euler-admissible placement of the same poles on a linear line.
The gamma phase is exact, but the needed root inequality is not yet proved.
GRH remains open.
\end{firewall}

\begin{assessment}[V88 closure and V89 direction]
V88 supplies the first quantitative construction that simultaneously
enforces the endpoint cancellation, the principal-pole zero, target
normalization, complete finite gamma-prefix cancellation, a strict Euler
convergence margin, and positivity from an explicit time onward.  It also
shows exactly why the currently proved onset is too late: the required line
movement crosses the gamma large-deviation phase transition.

V89 should analyze \textup{(V88.15)} at \(d=cM\).  The product contains
linear multiplier operators for real gamma roots and conjugate quadratic
pairs for the target gamma ladder.  A root-preserving or saddle-point
estimate uniform on
\[
 0<q\le \exp(-\theta cL_f/R)
\]
would decide whether the smallest positive root remains above the Euler
threshold for any \(c<c_\star\).  If it does, insert the resulting filter
into the V80 compensation law and return to the local zero-separation
determinant.  If it does not, the equally spaced scalar escape is closed and
the next route must vary the pole distribution or move to a matrix-valued
orientation.
\end{assessment}

\section*{Version V88 append-only record}
\addcontentsline{toc}{section}{Version V88 append-only record}
The complete V87 source was retained before this continuation.  It had
1267938 bytes, 37306 lines, and SHA--256
\[
 \texttt{D8E233B97C95D9D7FBE16D3C9DD6D457467877774AC3197CE21F7CD56A9C2954}.
\]
V88 constructed the equally spaced moving-line rational filter, proved its
alternating-residue positivity onset, derived the exact linear-scale gamma
large-deviation phase and its unique transition, ruled out the elementary
\(M\log M\) overshift, and reduced the surviving linear regime to the
smallest positive zero of an explicit generalized-Eulerian polynomial.

% =============================================================================
% END APPENDED V88 MATERIAL
% =============================================================================

% =============================================================================
% BEGIN APPENDED V89 MATERIAL
% =============================================================================

\section{Finite multiplier sequences close the linear equispaced escape}
\label{sec:grh-v89}

V88 left one exact question.  For the equally spaced poles, does the
generalized Eulerian polynomial
\[
 H_{M,d}(q)=\sum_{j=0}^{n}(-1)^j{n\choose j}
 \frac{P_{M,d}(-j\varepsilon)}{P_{M,d}(0)}q^j
 \tag{V89.1}
\]
remain positive up to the Euler threshold when \(d\sim cM\) and the
second omitted gamma block still decays?  The answer is no after harmless
real-ladder padding.  The reason is not an asymptotic saddle calculation:
it is a finite P\'olya--Schur theorem.  Each real numerator zero supplies a
linear finite multiplier sequence, and each conjugate pair supplies a
quadratic one whose Jensen discriminant is positive on every moving line
with \(d^2/M\to\infty\).  Thus every zero of \(H_{M,d}\) is positive.
Its exact root product then forces the last time-domain zero beyond the
first-prime frequency on the whole decaying side of the V88 gamma phase.

The finite multiplier criterion used below is the algebraic
characterization of \(n\)-sequences due to Craven and Csordas.\footnote{T.
Craven and G. Csordas, ``Location of Zeros. Part I: Real Polynomials and
Entire Functions,'' \emph{Illinois J. Math.} \textbf{27} (1983), Theorem
3.1.}  We state the precise form needed here so that no conjectural
real-rootedness assertion is imported.

\subsection{The one-polynomial finite multiplier test}

For a real vector
\(
 \boldsymbol\mu=(\mu_0,\ldots,\mu_n)
\)
write
\[
 T_{\boldsymbol\mu}\!\left[\sum_{j=0}^{n}a_jq^j\right]
 :=\sum_{j=0}^{n}\mu_ja_jq^j.
 \tag{V89.2}
\]

\begin{theorem}[Finite multiplier-sequence criterion]
\label{thm:v89-finite-multiplier}
The transform \(T_{\boldsymbol\mu}\) carries every real-rooted real
polynomial of degree at most \(n\) to a real-rooted polynomial if and only
if
\[
 J_{\boldsymbol\mu,n}(x)
 :=T_{\boldsymbol\mu}[(1+x)^n]
 =\sum_{j=0}^{n}{n\choose j}\mu_jx^j
 \tag{V89.3}
\]
has only real zeros of one sign.  In particular, if
\(
 \mu_j\ge0
\)
and \(J_{\boldsymbol\mu,n}\) has only real negative zeros, then
\(T_{\boldsymbol\mu}\) preserves real-rootedness through degree \(n\).
Pointwise products of such vectors again have this property.
\end{theorem}

\begin{proof}
The equivalence is the Craven--Csordas algebraic characterization of
finite \(n\)-sequences, itself a finite Grace--Szeg\H{o} composition
theorem.  For the last statement observe that
\[
 T_{\boldsymbol\mu\boldsymbol\nu}
 =T_{\boldsymbol\mu}\circ T_{\boldsymbol\nu},
 \qquad
 (\boldsymbol\mu\boldsymbol\nu)_j=\mu_j\nu_j.
 \tag{V89.4}
\]
The composition of two real-root preservers is again a real-root
preserver.
\end{proof}

The two elementary test vectors needed for \(P_{M,d}\) can be decided
exactly.

\begin{lemma}[Linear and conjugate-quadratic factors]
\label{lem:v89-factor-tests}
Fix \(u>0\), \(v\in\mathbb R\), and \(n\ge2\).
\begin{enumerate}[label=\textup{(\roman*)},leftmargin=3em]
\item The vector
      \(
       \mu_j=j+u
      \)
      is a finite multiplier sequence through degree \(n\).
\item The vector
      \[
       \nu_j=(j+u)^2+v^2
       \tag{V89.5}
      \]
      is a finite multiplier sequence through degree \(n\) whenever
      \[
       \boxed{v^2\le u+\frac14.}
       \tag{V89.6}
      \]
\end{enumerate}
\end{lemma}

\begin{proof}
The linear Jensen test is
\[
 \sum_{j=0}^{n}{n\choose j}(j+u)x^j
 =(1+x)^{n-1}\{u+(u+n)x\}.
 \tag{V89.7}
\]
Both zeros are negative, counting the repeated zero at \(-1\).

For the quadratic vector, the binomial identities for \(j\) and \(j^2\)
give
\[
 \sum_{j=0}^{n}{n\choose j}\{(j+u)^2+v^2\}x^j
 =(1+x)^{n-2}Q_{n,u,v}(x),
 \tag{V89.8}
\]
where, with \(A=u^2+v^2\),
\[
 \begin{aligned}
 Q_{n,u,v}(x)
 ={}&A+\{n(1+2u)+2A\}x
       +\{(n+u)^2+v^2\}x^2.
 \end{aligned}
 \tag{V89.9}
\]
Its discriminant is exactly
\[
 \boxed{
 \Delta_{n,u,v}
 =n\left[n(1+4u-4v^2)+4(u^2+v^2)\right].}
 \tag{V89.10}
\]
Condition \textup{(V89.6)} makes this positive.  All three coefficients
of \(Q_{n,u,v}\) are positive, so its two real zeros are negative.  Theorem
\ref{thm:v89-finite-multiplier} proves both claims.
\end{proof}

\subsection{Real-rootedness of the V88 residue polynomial}

Let
\[
 B:=\max_{1\le\ell\le R}|\Im\zeta_\ell|.
 \tag{V89.11}
\]
The two repair roots \(d,\xi_d\) are positive real, and a nonreal root of
\(P_{M,d}\) has the form
\[
 \lambda=d+1+\zeta_\ell+hm=a+ib,
 \qquad a>d,\quad |b|\le B.
 \tag{V89.12}
\]

\begin{theorem}[Positive real roots on the upper moving scales]
\label{thm:v89-positive-roots}
If
\[
 \boxed{
 \frac{B^2N^2}{\theta^2d^2}
 \le \frac{N}{\theta}+\frac14,}
 \tag{V89.13}
\]
then \(H_{M,d}\) has exactly \(n\) positive real zeros, counted with
multiplicity.  In particular this holds, for all sufficiently large \(M\),
along every sequence for which
\[
 \frac{d^2}{M}\longrightarrow\infty.
 \tag{V89.14}
\]
If all offsets are real, no lower bound on \(d\) beyond the V88 construction
is needed.
\end{theorem}

\begin{proof}
For a positive real root \(\lambda\) of \(P_{M,d}\), its contribution to
\textup{(V88.16)} is, up to a positive constant,
\[
 1+\frac{j\varepsilon}{\lambda}
 =\frac{\varepsilon}{\lambda}
 \left(j+\frac{\lambda}{\varepsilon}\right).
 \tag{V89.15}
\]
Lemma \ref{lem:v89-factor-tests}(i) applies.  For a conjugate pair
\(\lambda=a+ib,\overline\lambda\), the paired contribution is, again up to
a positive constant,
\[
 \left(1+\frac{j\varepsilon}{\lambda}\right)
 \left(1+\frac{j\varepsilon}{\overline\lambda}\right)
 =\frac{\varepsilon^2}{|\lambda|^2}
 \left\{\left(j+\frac a\varepsilon\right)^2
       +\left(\frac b\varepsilon\right)^2\right\}.
 \tag{V89.16}
\]
Here
\[
 u=\frac a\varepsilon>\frac d\varepsilon=\frac N\theta,
 \qquad
 v^2=\frac{b^2}{\varepsilon^2}
 \le\frac{B^2N^2}{\theta^2d^2}.
 \tag{V89.17}
\]
Thus \textup{(V89.13)} implies \(v^2\le u+1/4\), and Lemma
\ref{lem:v89-factor-tests}(ii) applies.  Theorem
\ref{thm:v89-finite-multiplier} and composition now show that the complete
coefficient vector
\[
 \mu_j:=\frac{P_{M,d}(-j\varepsilon)}{P_{M,d}(0)}
 \tag{V89.18}
\]
preserves real-rootedness through degree \(n\).  Applying it to
\((1-q)^n\) proves that \(H_{M,d}\) is real-rooted.

Every coefficient magnitude in \textup{(V89.1)} is positive.  Hence
\[
 H_{M,d}(-q)=\sum_{j=0}^{n}{n\choose j}\mu_jq^j>0
 \qquad(q\ge0).
 \tag{V89.19}
\]
There are no nonpositive roots.  The constant and leading coefficients are
nonzero, so all \(n\) roots are positive.

Since \(N=RM+4\), condition \textup{(V89.14)} makes the left side of
\textup{(V89.13)} \(o(M)\), while its right side is asymptotic to
\(RM/\theta\).  This proves the asymptotic assertion.  If \(B=0\), only
the linear test is present.
\end{proof}

This theorem settles the root-location question posed in V88 on every
linear line.  More importantly, it gives access to the extreme root without
having to locate all the other roots.

\subsection{The root product and the last time-domain zero}

Write the positive zeros of \(H_{M,d}\) as
\[
 0<r_{1,M,d}\le\cdots\le r_{n,M,d},
 \qquad
 U_{M,d}:=-\frac1\varepsilon\log r_{1,M,d}.
 \tag{V89.20}
\]
Thus \(U_{M,d}\) is the largest real time obtained from a zero by the map
\(q=e^{-\varepsilon u}\).

\begin{lemma}[Exact product forcing]
\label{lem:v89-root-product}
Under \textup{(V89.13)},
\[
 \prod_{k=1}^{n}r_{k,M,d}
 =\frac1{c_n},
 \qquad
 c_n=\frac{P_{M,d}(-n\varepsilon)}{P_{M,d}(0)},
 \tag{V89.21}
\]
and consequently
\[
 \boxed{
 U_{M,d}\ge
 \frac{\log c_n}{n\varepsilon}.}
 \tag{V89.22}
\]
If \(d_M/M\to c\in(0,\infty)\), then
\[
 \boxed{
 \liminf_{M\to\infty}U_{M,d_M}
 \ge \Lambda_{\theta,h,R}(c),}
 \tag{V89.23}
\]
where
\[
 \boxed{
 \Lambda_{\theta,h,R}(c)
 :=\frac{R}{\theta c}
 \int_0^1
 \log\left(1+\frac{\theta c}{c+ht}\right)dt.}
 \tag{V89.24}
\]
\end{lemma}

\begin{proof}
The constant coefficient of \(H_{M,d}\) is one and its leading coefficient
is \((-1)^nc_n\), which proves \textup{(V89.21)}.  The smallest positive
root is no larger than the geometric mean, so
\[
 r_{1,M,d}\le c_n^{-1/n}.
 \tag{V89.25}
\]
Taking logarithms gives \textup{(V89.22)}.

The two repair roots contribute only \(O(1)\) to \(\log c_n\).  The fixed
offsets in the gamma ladders disappear uniformly after scaling by \(M\),
and \(n\varepsilon=(n/N)\theta d\).  Hence the \(RM\) gamma roots give the
Riemann-sum limit
\[
 \frac{\log c_n}{n\varepsilon}
 \longrightarrow
 \frac{R}{\theta c}
 \int_0^1
 \log\left(1+\frac{\theta c}{c+ht}\right)dt.
 \tag{V89.26}
\]
For conjugate pairs this is obtained by pairing the two logarithms before
taking the limit.  Equations \textup{(V89.22)} and \textup{(V89.26)} prove
the claim.
\end{proof}

\subsection{Comparison with the gamma phase}

The forcing length \(\Lambda\) is largest on the side where the V88 gamma
response is smallest.

\begin{lemma}[Monotonicity and a uniform value at the phase boundary]
\label{lem:v89-lambda-bound}
The function \(c\mapsto\Lambda_{\theta,h,R}(c)\) is strictly decreasing on
\((0,\infty)\).  If \(c_\star=c_\star(\theta,h)\) is the V88 phase boundary,
then
\[
 c_\star\le\frac{2h}{3}
 \tag{V89.27}
\]
and
\[
 \boxed{
 \Lambda_{\theta,h,R}(c_\star)
 >\frac Rh\log\frac74.}
 \tag{V89.28}
\]
\end{lemma}

\begin{proof}
For fixed \(t>0\), put
\[
 g_t(c):=\log\left(1+\frac{\theta c}{c+ht}\right).
 \tag{V89.29}
\]
This function is increasing and strictly concave, with \(g_t(0)=0\).
Therefore \(g_t(c)/c\) is strictly decreasing.  The \(t=0\) integrand
contributes the constant \(\log(1+\theta)/c\), also strictly decreasing.
Integration proves the first assertion.

The defining phase equation \textup{(V88.26)} is
\[
 \log4
 =\int_0^1
 \log\left(1+\frac{2h}{c_\star(1+\theta t)}\right)dt
 \le\log\left(1+\frac{2h}{c_\star}\right).
 \tag{V89.30}
\]
This implies \textup{(V89.27)}.  Next use
\(
 \log(1+s)\ge s/(1+s)
\)
to obtain
\[
 \begin{aligned}
 \int_0^1\log\left(1+\frac{\theta c}{c+ht}\right)dt
 &\ge
 \int_0^1\frac{\theta c}{c(1+\theta)+ht}\,dt\\
 &=\frac{\theta c}{h}
 \log\left(1+\frac{h}{c(1+\theta)}\right).
 \end{aligned}
 \tag{V89.31}
\]
At \(c=c_\star\), condition \(0<\theta<1\) and
\textup{(V89.27)} give
\[
 \frac{h}{c_\star(1+\theta)}>\frac34.
 \tag{V89.32}
\]
Insert this into \textup{(V89.24)} and \textup{(V89.31)} to prove
\textup{(V89.28)}.
\end{proof}

The number of ladders is not intrinsic to the cancellation problem.  One
may add finitely many real offsets, in distinct residue classes modulo
\(h\), and cancel the corresponding arithmetic progressions as well.  These
extra zeros of \(F_{M,d}\) are harmless: they add no response that must be
recovered and preserve all normalization and endpoint conditions.

\begin{lemma}[Real-ladder padding]
\label{lem:v89-padding}
Starting from any V88 conjugation-closed offset set, one may enlarge it by
finitely many real offsets so that
\[
 \boxed{
 \frac Rh\log\frac74>L_f.}
 \tag{V89.33}
\]
The enlarged progressions can be chosen pairwise disjoint, and the critical
constant \(c_\star(\theta,h)\) is unchanged.
\end{lemma}

\begin{proof}
Choose distinct real offsets greater than \(-\beta\), avoiding the finitely
many congruence classes modulo \(h\) already used.  This gives any prescribed
finite enlargement of \(R\).  The extra roots are positive real and hence
do not alter Theorem \ref{thm:v89-positive-roots}.  Finally, the V88 rate is
\(
 \Psi_{\theta,h,R}(c)=R\phi_{\theta,h}(c)
\);
its zero is the zero of \(\phi_{\theta,h}\), independent of \(R\).
\end{proof}

\begin{theorem}[No linear overlap for the padded equispaced family]
\label{thm:v89-no-linear-overlap}
Fix a padded family satisfying \textup{(V89.33)}, and suppose
\[
 \frac{d_M}{M}\longrightarrow c\in(0,\infty).
 \tag{V89.34}
\]
Then the following two requirements cannot both hold for all sufficiently
large \(M\):
\begin{enumerate}[label=\textup{(\roman*)},leftmargin=3em]
\item \(K_{M,d_M}(u)>0\) for every \(u\ge L_f\);
\item the first response in the second omitted gamma block tends to zero in
      absolute value.
\end{enumerate}
More precisely, if \(c\le c_\star\), then \(K_{M,d_M}\) has a zero at some
\(u>L_f\); if \(c>c_\star\), then every fixed-ladder response
\(F_{M,d_M}(Y_{\ell,M,d_M})\) grows exponentially.
\end{theorem}

\begin{proof}
Every linear line satisfies \(d_M^2/M\to\infty\), so Theorem
\ref{thm:v89-positive-roots} applies.  If \(c\le c_\star\), Lemmas
\ref{lem:v89-root-product} and \ref{lem:v89-lambda-bound}, monotonicity of
\(\Lambda\), and \textup{(V89.33)} give
\[
 \liminf_{M\to\infty}U_{M,d_M}
 \ge\Lambda_{\theta,h,R}(c)
 \ge\Lambda_{\theta,h,R}(c_\star)
 >L_f.
 \tag{V89.35}
\]
Thus \(r_{1,M,d_M}=e^{-\varepsilon U_{M,d_M}}\) is a zero in the Euler
interval from \textup{(V88.35)}, and strict positivity fails.

If \(c>c_\star\), Theorem \ref{thm:v88-gamma-phase} gives
\[
 |F_{M,d_M}(Y_{\ell,M,d_M})|
 =\exp\{(\Psi_{\theta,h,R}(c)+o(1))M\},
 \qquad \Psi_{\theta,h,R}(c)>0.
 \tag{V89.36}
\]
Requirement (ii) fails.  The boundary \(c=c_\star\) is already covered by
the strict root-product obstruction.
\end{proof}

This is the desired decision of the V88 linear root problem.  Padding is
used only to make the universal lower bound exceed the externally fixed
frequency \(L_f\); it does not manufacture the phase transition.

\subsection{The upper mesoscopic regime}

The same mechanism also reaches below the linear scale.

\begin{corollary}[No Euler positivity when
\(\sqrt M\ll d\ll M\)]
\label{cor:v89-mesoscopic}
Suppose
\[
 \frac{d_M}{M}\longrightarrow0,
 \qquad
 \frac{d_M^2}{M}\longrightarrow\infty.
 \tag{V89.37}
\]
Then, even without padding,
\[
 U_{M,d_M}\longrightarrow\infty.
 \tag{V89.38}
\]
Consequently \(K_{M,d_M}\) has a zero after \(L_f\) for all sufficiently
large \(M\).
\end{corollary}

\begin{proof}
The second condition in \textup{(V89.37)} invokes Theorem
\ref{thm:v89-positive-roots}.  Put \(c_M=d_M/M\).  The uniform version of
the gamma-root Riemann sum in \textup{(V89.26)} gives
\[
 \frac{\log c_n}{n\varepsilon}
 =\Lambda_{\theta,h,R}(c_M)+o(1).
 \tag{V89.39}
\]
From \textup{(V89.31)},
\[
 \Lambda_{\theta,h,R}(c_M)
 \ge\frac Rh
 \log\left(1+\frac{h}{c_M(1+\theta)}\right)
 \longrightarrow\infty.
 \tag{V89.40}
\]
Now use \textup{(V89.22)}.
\end{proof}

\begin{theorem}[V89 equispaced-filter verdict]
\label{thm:v89-verdict}
For the padded equally spaced V88 rational filters:
\begin{enumerate}[label=\textup{(\roman*)},leftmargin=3em]
\item the finite multiplier criterion proves that all residue-polynomial
      zeros are positive whenever \(d^2/M\to\infty\);
\item their exact product forces a time-domain zero beyond \(L_f\) on every
      linear line \(d/M\to c\le c_\star\);
\item on every line \(d/M\to c>c_\star\), a single omitted gamma response
      grows exponentially;
\item the same late-zero obstruction holds throughout
      \(\sqrt M\ll d\ll M\);
\item V88 already closes \(d/M\to\infty\) by exponential gamma-response
      growth.
\end{enumerate}
Thus this scalar equispaced escape has no admissible linear or upper
mesoscopic scale.  The unresolved part of this family is
\[
 \boxed{d=O(\sqrt M),}
 \tag{V89.41}
\]
where the conjugate-quadratic discriminant in \textup{(V89.10)} can cease
to be positive.  This theorem neither treats arbitrary pole distributions
nor proves GRH.
\end{theorem}

\begin{firewall}[Finite real-rootedness is the decisive input]
The root-product identity alone does not locate a zero: without
real-rootedness, the geometric mean could be carried by complex roots.
Conversely, the finite multiplier theorem does not say that the smallest
root is beyond the Euler threshold.  The obstruction requires both exact
pieces.  It also uses padding by auxiliary real gamma ladders; therefore it
closes the padded equally spaced construction, not every rational filter
with the V80 moment data.  GRH remains open.
\end{firewall}

\begin{assessment}[V89 closure and mandatory V90 sweep]
V89 resolves the explicit polynomial problem left by V88 on all linear
moving lines and, in fact, down to the square-root boundary.  The outcome
is adverse to the proposed scalar escape: on the phase side where an
individual gamma response can decay, a rigorously positive real root moves
past the first-prime time; on the other side, that response grows.

V90 is the next ten-iteration checkpoint and must begin with a stable sweep
of all parallel notes created since the V80 sweep.  After incorporating any
new arguments, the live alternatives are: analyze the critical and lower
scales \(d=O(\sqrt M)\), replace the uniform pole law by a non-equispaced
distribution, or go upstream to a matrix-valued orientation/local
zero-separation determinant.  Any return to the scalar equispaced family
must confront the exact quadratic Jensen discriminant
\textup{(V89.10)} rather than repeat the already closed linear analysis.
\end{assessment}

\section*{Version V89 append-only record}
\addcontentsline{toc}{section}{Version V89 append-only record}
The complete V88 source was retained before this continuation.  It had
1283344 bytes, 37837 lines, and SHA--256
\[
 \texttt{CDF11514E834EAA745956137F526B35E8C18E4D7B79C557370C5A3B4A37B4DF6}.
\]
V89 invoked the finite \(n\)-sequence characterization, proved positive
real-rootedness of the V88 residue polynomial on every scale
\(d^2/M\to\infty\), combined the exact root product with the gamma phase,
and ruled out the padded equally spaced filter on all linear and upper
mesoscopic moving lines.  The remaining scale is \(d=O(\sqrt M)\), together
with nonuniform or matrix-valued constructions outside this family.

% =============================================================================
% END APPENDED V89 MATERIAL
% =============================================================================

% =============================================================================
% BEGIN APPENDED V90 MATERIAL
% =============================================================================

\section{V90: mandatory corpus sweep and the square-root Jensen boundary}
\label{sec:grh-v90}

V89 closed the padded equally spaced scalar filter on every linear moving
line and throughout \(\sqrt M\ll d\ll M\).  This version first performs the
scheduled ten-version audit of the companion notes.  The audit supplies no
arithmetic zero-separation theorem, but it reinforces two proof disciplines:
an oriented response cannot be replaced by its Gram size, and failure of one
sufficient factorization calls for a change of factorization rather than a
claim about the underlying form.

Returning to the exact V89 Jensen polynomial then improves the mathematical
boundary.  The sufficient inequality \(v^2\le u+1/4\) used there was not the
exact quadratic criterion.  Keeping the positive \(4u^2\) term in the
discriminant proves real-rootedness on a nontrivial part of the critical line
\(d\asymp\sqrt M\).  The factors which fail the criterion are also localized
exactly to an initial gamma prefix.  This closes the large-constant
square-root regime and identifies the finite-factor obstruction that remains
below it.

\subsection{Hash-gated V80--V90 corpus audit}

The stable V80 manifest contained \(31\) non-GRH \texttt{.tex} files and had
SHA--256
\[
 \texttt{600CBF5481F3BC63CF12C87FEF596C2BFDD57F92F6DD1E820E7D6CEA8D544BCF}.
 \tag{V90.1}
\]
As before, canonical rows are
\[
 \texttt{filename|byte-count|SHA--256}\backslash\texttt{n},
 \tag{V90.2}
\]
sorted by filename, joined with LF terminators, and hashed as UTF--8.

One preliminary snapshot was discarded when the SM writer advanced from
V86 during the content audit.  The final clean snapshot had one live SM
endpoint and one live Yang--Mills endpoint.  Two complete passes had the UTC
endpoints
\begin{center}
\small
pass 1 start: \texttt{2026-09-11T08:24:58.2324370Z}\\
pass 1 end: \texttt{2026-09-11T08:24:58.3532694Z}\\
pass 2 start: \texttt{2026-09-11T08:24:58.3560729Z}\\
pass 2 end: \texttt{2026-09-11T08:24:58.4012058Z}.
\end{center}
Each pass rechecked byte counts and modification ticks after hashing, found
\(31\) stable non-GRH files, and produced
\[
 \boxed{
 \texttt{F6BFD1D3E054643B8154F9E9758085110330023AF3D0473D6BB475CA6D1A439D}.}
 \tag{V90.3}
\]
The changed endpoint rows were
\begin{center}
\tiny
\begin{tabular}{p{0.54\textwidth}|r|p{0.30\textwidth}}
\toprule
file & bytes & SHA--256\\
\midrule
\path{Renewal_SM_Einstein_Continuum_Research_Notes_V89.tex}
&923213&\texttt{CA7FC92EC97B6D862FE75A06C92273F0}\newline
\texttt{09D0E207C581EC2377F725BC6473A600}\\
\path{Renewal_Yang_Mills_Mass_Gap_Research_Notes_V64.tex}
&544441&\texttt{69B9A7675E3F46414962EF6CD74AAD10}\newline
\texttt{E17E1EEA975C543B78419885D2B34808}\\
\bottomrule
\end{tabular}
\end{center}
Replacing these rows by the V80 endpoints
\[
\begin{array}{c|r|c}
\text{SM V53}
 &544386&
\texttt{84160EA8FAA35C64BA824BB5C641FABD}\cdots\\
\text{Yang--Mills V54}
 &467684&
\texttt{20F0DA83E09A4609C73A76F9F198569}\cdots
\end{array}
\tag{V90.4}
\]
and their recorded full hashes reproduces \textup{(V90.1)} exactly.
Therefore the other \(29\) rows are byte-identical to the V80 corpus, and
the unique new material is SM V54--V89 and Yang--Mills V55--V64.

\subsection{Typed audit of SM V54--V89}

The SM successor contains four coherent chains.
\begin{enumerate}[label=\textup{(\roman*)},leftmargin=3em]
\item V54--V65 derive reconstructed-action Hessian ledgers, nonconvex
      Poincar\'e compilers, exact constraint reactions, local unimodular
      gauge charts, and intensive determinant-response owners.  Every
      positive-law conclusion remains conditional on physical coercivity,
      determinant orientation, or trace-ideal hypotheses.
\item V66--V76 pass to block conditional gaps and exact scalar or matrix
      conductance projections.  V71's moving-factor identity separates
      right-factor response from unitary commutators; V72--V76 then identify
      neutrality and paired boundary flux as the low-frequency owners.
\item V77--V83 retain rigid rotations, remove a checkerboard hourglass with
      full-face moments, prove curved quadrature and Schur domination under
      declared reconstruction hypotheses, and assemble a single
      finite-cutoff positive-law certificate.  The physical entries of that
      certificate are not verified.
\item V84--V89 derive localized covariance from a Witten resolvent, state
      the physical clustering corridor, prove the conditional
      Osterwalder--Schrader spectral handoff, preserve reflection positivity
      under a convergent cutoff passage, and obtain negative-Sobolev
      tightness and moments from log-Sobolev inputs.  Existence,
      reflection positivity of the cutoff law, entropy tensorization, and
      the required uniform means remain open in that source.
\end{enumerate}

The response-orientation warning is relevant, but none of these theorems
contains a Dirichlet coefficient identity, a completed \(L\)-function
functional equation, or a comparison of the principal and character zero
spectra.  The conductance and moving-square-factor results assume positive
self-adjoint factors and controlled kernels; the GRH residue polynomial is
a signed scalar transform with no such identified factor.  The probability,
compactness, and OS theorems assume a positive measure or reflected Gram
form that the zero-side explicit formula does not provide.  Their
conclusions therefore do not close any premise of V89.

\subsection{Typed audit of Yang--Mills V55--V64}

Yang--Mills V55 localizes four failures of a selected \(90\times90\)
sufficient splitting.  V56 repairs those centres with exact anisotropic
rational bridge charts; V57 closes one mixed direct-plus quarter cell.
V58--V62 build the inverted fourth orientation through its first two
quarter coordinates.  V63 again isolates four failures of the original
full-radius certificate, and V64 crosses them by conjugate rational bridges,
proving both direct signs on paired third-coordinate slabs through depth
\(2t\).

The source correctly never turns failure of a sufficient splitting into
negativity of the underlying Wilson pencil.  Its useful procedural lesson is
to isolate the obstructing factors and change the local cover.  The Wilson
matrices have no arithmetic Euler coefficients, gamma ladders, or
zero-spectrum transform, so their positivity floors and bridge charts are
not importable as GRH estimates.

\begin{theorem}[V90 companion-corpus verdict]
\label{thm:v90-corpus-verdict}
No post-V80 companion theorem supplies a missing hypothesis for the V89
equispaced-filter problem or for the upstream V80 two-channel compensation
law.  The transferable content is diagnostic:
\begin{enumerate}[label=\textup{(\roman*)},leftmargin=3em]
\item preserve the signed response direction rather than only a Gram norm;
\item isolate the factors at which a sufficient certificate fails;
\item do not infer failure of the underlying arithmetic form from failure
      of that certificate;
\item compactness, clustering, and spectral handoff theorems cannot be used
      before an arithmetic positive form with their hypotheses is proved.
\end{enumerate}
\end{theorem}

\begin{proof}
The SM results require positive conditional laws, self-adjoint square
factors, coercive influence matrices, reflection-positive limits, or
log-Sobolev constants.  The Yang--Mills results concern exact finite Wilson
pencils and local phase covers.  None supplies the Euler positivity, gamma
tail, or zero-separation inequality appearing in V80--V89.  Matching the
stated hypotheses object by object proves the nontransfer claims.
\end{proof}

\subsection{The exact conjugate-quadratic criterion}

Return to Lemma~\ref{lem:v89-factor-tests}.  Its Jensen polynomial is
\[
 (1+x)^{n-2}Q_{n,u,v}(x),
 \tag{V90.5}
\]
and V89 computed
\[
 \Delta_{n,u,v}
 =n\{n(1+4u-4v^2)+4(u^2+v^2)\}.
 \tag{V90.6}
\]
Keeping all terms gives an equivalence.

\begin{lemma}[Exact quadratic finite-multiplier boundary]
\label{lem:v90-exact-quadratic}
For \(n\ge2\), \(u>0\), and \(v\in\mathbb R\), the vector
\[
 \nu_j=(j+u)^2+v^2,\qquad 0\le j\le n,
 \tag{V90.7}
\]
is a finite multiplier sequence through degree \(n\) if and only if
\[
 \boxed{
 v^2\le {\mathcal Q}_n(u):=
 \frac{n(1+4u)+4u^2}{4(n-1)}.}
 \tag{V90.8}
\]
\end{lemma}

\begin{proof}
The coefficients of the quadratic \(Q_{n,u,v}\) in
\textup{(V89.9)} are strictly positive.  Hence, if its discriminant is
nonnegative, both roots are real and negative; if the discriminant is
negative, they are nonreal.  Formula \textup{(V90.6)} is equivalent to
\[
 n(1+4u)+4u^2-4(n-1)v^2\ge0.
 \tag{V90.9}
\]
The finite \(n\)-sequence criterion
\ref{thm:v89-finite-multiplier} now proves the equivalence.
\end{proof}

The function \({\mathcal Q}_n(u)\) is strictly increasing for \(u>0\).
Consequently, within one gamma ladder the smallest real part is the only
one that can determine whether all its conjugate factors pass the test.

\subsection{A sharp factor threshold on \(d\asymp\sqrt M\)}

For a conjugate ladder offset write
\[
 \zeta_\ell=\alpha_\ell+ib_\ell,\qquad
 a_{\ell,m}:=d+1+\alpha_\ell+hm.
 \tag{V90.10}
\]
Its V89 factor parameters are
\[
 u_{\ell,m}=\frac{a_{\ell,m}}{\varepsilon},
 \qquad
 v_\ell=\frac{b_\ell}{\varepsilon}.
 \tag{V90.11}
\]
Let
\[
 B:=\max_\ell|b_\ell|,
 \qquad
 \kappa_{\rm J}:=
 B\sqrt{\frac{R}{1+\theta}}.
 \tag{V90.12}
\]

\begin{theorem}[Critical-scale positive real-rootedness]
\label{thm:v90-critical-real-rootedness}
Suppose
\[
 \frac{d_M}{\sqrt M}\longrightarrow\kappa\in(0,\infty)
 \quad\text{and}\quad
 \boxed{\kappa>\kappa_{\rm J}.}
 \tag{V90.13}
\]
Then, for all sufficiently large \(M\), every conjugate factor of
\(P_{M,d_M}\) is a finite multiplier sequence through degree \(n\).
Consequently \(H_{M,d_M}\) has exactly \(n\) positive real zeros.
\end{theorem}

\begin{proof}
For each fixed ladder, \({\mathcal Q}_n\) is increasing, so it is enough to
test \(m=0\).  Since \(N=RM+4\), \(n=RM+3\), and
\(\varepsilon=\theta d/N\),
\[
 \frac{u_{\ell,0}}{n}\longrightarrow\frac1\theta,
 \qquad
 \frac{v_\ell^2}{n}
 \longrightarrow
 \frac{b_\ell^2R}{\theta^2\kappa^2}.
 \tag{V90.14}
\]
On the other hand,
\[
 \frac{{\mathcal Q}_n(u_{\ell,0})}{n}
 \longrightarrow
 \frac1\theta+\frac1{\theta^2}
 =\frac{1+\theta}{\theta^2}.
 \tag{V90.15}
\]
Condition \textup{(V90.13)} makes the limit in the second part of
\textup{(V90.14)} strictly smaller than \textup{(V90.15)}, uniformly over
the finitely many ladders.  Lemma~\ref{lem:v90-exact-quadratic} applies to
all conjugate pairs.  The positive real roots \(d,\xi_d\) and any real
gamma offsets give the linear multiplier factors already treated in
Lemma~\ref{lem:v89-factor-tests}.  Composition and
\textup{(V89.19)} show that all roots of \(H_{M,d_M}\) are positive.
\end{proof}

This improves the V89 condition \(d^2/M\to\infty\) to an explicit open
part of the critical line.  No padding is needed.

\subsection{Root-product growth below the linear scale}

The V89 product bound becomes unbounded whenever \(d\to\infty\) but
\(d=o(M)\).

\begin{lemma}[Sublinear last-zero scale]
\label{lem:v90-sublinear-product}
Suppose \(d_M\to\infty\) and \(d_M/M\to0\).  For the V88 equally spaced
filter,
\[
 \boxed{
 \frac{\log c_n}{n\varepsilon}
 =\frac Rh\log\frac{M}{d_M}+O(1),}
 \tag{V90.16}
\]
where the constant may depend on the fixed offsets,
\(\theta,h,R\), but not on \(M\).
\end{lemma}

\begin{proof}
Put \(A=n\varepsilon=\theta d(n/N)\), so \(A\asymp d\).  A fixed complex
offset changes the paired logarithmic gamma factor by a summable amount:
\[
 \sum_{m=0}^{M-1}
 \left[
 \log\left|1+\frac{A}{d+1+\zeta_\ell+hm}\right|
 -\log\left(1+\frac{A}{d+hm}\right)
 \right]
 =O(1).
 \tag{V90.17}
\]
Indeed the termwise difference is bounded by a constant times
\(A/(d+hm)^2\), whose sum is \(O(A/d)=O(1)\).  For the real model summand,
monotone sum--integral comparison gives
\[
 \begin{aligned}
 \sum_{m=0}^{M-1}\log\left(1+\frac{A}{d+hm}\right)
 &=
 \frac1h\int_d^{d+hM}\log\left(1+\frac A x\right)\,dx+O(1)\\
 &=\frac Ah\log\frac Md+O(A).
 \end{aligned}
 \tag{V90.18}
\]
There are \(R\) ladders.  The two repair roots add \(O(1)\) to
\(\log c_n\).  Divide the resulting identity by \(A=n\varepsilon\) to
obtain \textup{(V90.16)}.
\end{proof}

\begin{corollary}[Large-constant square-root obstruction]
\label{cor:v90-critical-obstruction}
Under the hypotheses of
Theorem~\ref{thm:v90-critical-real-rootedness},
\[
 U_{M,d_M}\ge
 \frac{R}{2h}\log M+O(1)\longrightarrow\infty.
 \tag{V90.19}
\]
Therefore \(K_{M,d_M}\) has a zero at some \(u>L_f\) for all sufficiently
large \(M\), and the V88 Euler-positivity requirement fails.
\end{corollary}

\begin{proof}
The exact root-product inequality \textup{(V89.22)} applies because all
roots are positive.  Insert \(d_M=\kappa\sqrt M(1+o(1))\) into
\textup{(V90.16)}.
\end{proof}

\subsection{Localization of the unresolved factors}

The same discriminant gives an exact physical-space cutoff for every
conjugate gamma node.

\begin{proposition}[Exact bad-prefix endpoint]
\label{prop:v90-bad-prefix}
For a conjugate pair with imaginary part \(b\), define
\[
 a_{\rm J}(n,\varepsilon,b)
 :=
 \frac{
 -n\varepsilon+
 \sqrt{(n-1)(n\varepsilon^2+4b^2)}
 }{2}.
 \tag{V90.20}
\]
Its factor at real part \(a>0\) is a finite multiplier sequence through
degree \(n\) if and only if
\[
 \boxed{a\ge \max\{0,a_{\rm J}(n,\varepsilon,b)\}.}
 \tag{V90.21}
\]
Thus the noncertified nodes in the ladder are exactly the initial indices
\[
 0\le m<
 \frac{a_{\rm J}(n,\varepsilon,b)-d-1-\alpha_\ell}{h},
 \tag{V90.22}
\]
when the right side is positive.
\end{proposition}

\begin{proof}
Substitute \(u=a/\varepsilon\) and \(v=b/\varepsilon\) into
\textup{(V90.9)} and multiply by \(\varepsilon^2/4\).  The condition is
\[
 a^2+n\varepsilon a+\frac{n\varepsilon^2}{4}
 -(n-1)b^2\ge0.
 \tag{V90.23}
\]
The quadratic is increasing for \(a>0\), and its positive root is
\textup{(V90.20)}.  Inserting \(a=a_{\ell,m}\) proves
\textup{(V90.22)}.
\end{proof}

\begin{corollary}[Size of the dangerous prefix]
\label{cor:v90-prefix-size}
If \(d_M/\sqrt M\to\kappa\), then
\[
 \frac{a_{\rm J}(n,\varepsilon,b)}{\sqrt M}
 \longrightarrow
 \frac{
 -\theta\kappa+
 \sqrt{\theta^2\kappa^2+4Rb^2}
 }{2}.
 \tag{V90.24}
\]
If \(d=o(\sqrt M)\) and \(b\ne0\), the number of noncertified factors in
that ladder is
\[
 \frac{|b|\sqrt{RM}}{h}+o(\sqrt M).
 \tag{V90.25}
\]
In particular, below the threshold
\[
 \kappa<|b|\sqrt{\frac{R}{1+\theta}},
 \tag{V90.26}
\]
the first conjugate factor is not an \(n\)-sequence.
\end{corollary}

\begin{proof}
Use \(n\varepsilon/\sqrt M\to\theta\kappa\) and
\(n\varepsilon^2\to\theta^2\kappa^2/R\) in
\textup{(V90.20)}.  If \(d=o(\sqrt M)\), the first term is lower order and
\(a_{\rm J}=|b|\sqrt{RM}+o(\sqrt M)\).  Formula
\textup{(V90.22)} gives the count.  Finally, comparison of the limit in
\textup{(V90.24)} with \(d/\sqrt M=\kappa\) is equivalent to
\textup{(V90.26)}.
\end{proof}

\begin{firewall}[Failure of a factor test is not failure of the product]
When \textup{(V90.26)} holds, the Jensen polynomial for one conjugate pair
has two nonreal roots.  It does not follow that the complete
\(H_{M,d}\) has nonreal roots: composition with the other diagonal factors
could restore real-rootedness or could produce a positive interval without
global real-rootedness.  Conversely, the root-product bound cannot force a
late real zero until enough real-rootedness is recovered.  This is exactly
the sufficient-factorization distinction emphasized by the companion
audit.  GRH remains open.
\end{firewall}

\begin{theorem}[V90 square-root verdict]
\label{thm:v90-verdict}
For the V88 equally spaced family:
\begin{enumerate}[label=\textup{(\roman*)},leftmargin=3em]
\item V89 closes \(d/\sqrt M\to\infty\);
\item V90 closes every critical line
      \(d/\sqrt M\to\kappa>\kappa_{\rm J}\), without padding;
\item below \(\kappa_{\rm J}\), every failed quadratic certificate is
      confined to the explicit initial prefix
      \textup{(V90.22)};
\item for \(d=o(\sqrt M)\), that prefix has order \(\sqrt M\), whereas the
      remaining gamma factors are certified finite multiplier sequences.
\end{enumerate}
The unresolved equispaced range is therefore
\[
 \boxed{
 d=O(\sqrt M),\qquad
 \limsup\frac d{\sqrt M}\le
 B\sqrt{\frac{R}{1+\theta}},}
 \tag{V90.27}
\]
and within it the only obstruction to the V89 real-rootedness argument is
the localized conjugate prefix.  This does not resolve arbitrary pole laws
or GRH.
\end{theorem}

\begin{assessment}[V90 closure and V91 direction]
V90 completes the mandatory stable corpus sweep and verifies exact
continuity with the V80 manifest.  The newest companion mechanisms do not
transfer an arithmetic estimate, but their factorization warning prevents a
false conclusion below the Jensen threshold.

The exact discriminant nevertheless advances the attack: the large-constant
part of \(d\asymp\sqrt M\) is closed, the late zero grows at least
logarithmically, and every uncertified factor is now listed by the explicit
endpoint \textup{(V90.20)}.  V91 should factor
\[
 H_{M,d}
 =T_{\rm bad}\,T_{\rm good}[(1-q)^n],
\]
where \(T_{\rm good}\) is the certified tail and \(T_{\rm bad}\) contains
only \textup{(V90.22)}.  The next decision is whether a sector or
Hermite--Biehler estimate controls this \(O(\sqrt M)\) prefix.  If it does
not, the Yang--Mills audit suggests changing the pole mesh locally near the
complex gamma nodes rather than repeating the failed uniform
factorization.
\end{assessment}

\section*{Version V90 append-only record}
\addcontentsline{toc}{section}{Version V90 append-only record}
The complete V89 source was retained before this continuation.  It had
1301157 bytes, 38405 lines, and SHA--256
\[
 \texttt{9ACF8F389E075E37E9C4D7C74AB8F97C494D49DC7B0C6C7763F15040A2F11EEB}.
\]
V90 performed the mandatory stable companion-corpus sweep, proved the exact
quadratic finite-multiplier threshold, closed the large-constant
square-root regime of the equally spaced filter, and localized every
remaining noncertified conjugate factor to an explicit initial gamma
prefix.

% =============================================================================
% END APPENDED V90 MATERIAL
% =============================================================================
% =============================================================================
% BEGIN APPENDED V91 MATERIAL
% =============================================================================

\section{V91: a local zero-count certificate at the first-prime horizon}
\label{sec:grh-v91}

V90 localized the failure of the factorwise multiplier proof to an initial
conjugate prefix of order at most \(\sqrt M\).  Global real-rootedness is
stronger than the arithmetic application needs.  The V87 variation theorem
requires only a decision about how many sign-changing zeros can occur before
the fixed time \(L_f\).  This version recodes that local question exactly.

The alternating-residue polynomial has a reversed Newton expansion whose
coefficients are strictly alternating and explicitly positive in magnitude.
After the endpoint zero is removed, the Budan--Fourier theorem gives a finite
derivative-sign certificate which, by itself, can disprove Euler
admissibility.  A second calculation expresses the action of one failed
conjugate factor on the certified positive-rooted tail through its logarithmic
derivative.  These two formulas replace the open-ended request for a sector
estimate by two concrete inequalities.

\subsection{The reversed Newton polynomial}

Extend the V88 multiplier data from the integers to the polynomial
\[
 \mu(x):=\frac{P_{M,d}(-\varepsilon x)}{P_{M,d}(0)},
 \qquad \deg\mu=n-1,
 \tag{V91.1}
\]
and write
\[
 \delta_r:=\Delta^r\mu(0)
 =\sum_{k=0}^r(-1)^{r-k}{r\choose k}\mu(k),
 \qquad 0\le r\le n-1.
 \tag{V91.2}
\]

\begin{lemma}[Strict positivity of all Newton data]
\label{lem:v91-positive-newton}
For the conjugation-closed V88 family,
\[
 \boxed{\delta_r>0\qquad(0\le r\le n-1).}
 \tag{V91.3}
\]
\end{lemma}

\begin{proof}
Each positive real root \(\lambda\) of \(P_{M,d}\) contributes the factor
\[
 1+\frac{\varepsilon}{\lambda}x,
\]
whose coefficients are positive.  A conjugate pair \(a\pm ib\), with
\(a>0\), contributes
\[
 \frac{(a+\varepsilon x)^2+b^2}{a^2+b^2},
 \tag{V91.4}
\]
again with strictly positive constant, linear, and quadratic coefficients.
Consequently
\(
 \mu(x)=\sum_{m=0}^{n-1}p_mx^m
\)
has \(p_m>0\) at every degree.  The standard identity
\[
 \Delta^r x^m\big|_{x=0}=r!\,{m\brace r}
 \tag{V91.5}
\]
then gives
\[
 \delta_r=r!\sum_{m=r}^{n-1}p_m{m\brace r}>0,
\]
where \({m\brace r}\) is a Stirling number of the second kind.
\end{proof}

Put
\[
 y:=\frac{1-q}{q},\qquad
 A_s:={n\choose s}\delta_{n-s}>0\quad(1\le s\le n),
 \tag{V91.6}
\]
and define
\[
 \mathcal R_{M,d}(y)
 :=\sum_{s=1}^{n}(-1)^{s+1}A_sy^s,
 \qquad
 \mathcal S_{M,d}(y):=\frac{\mathcal R_{M,d}(y)}y.
 \tag{V91.7}
\]

\begin{theorem}[Exact first-prime recoding]
\label{thm:v91-newton-recoding}
For every \(q>0\),
\[
 \boxed{
 H_{M,d}(q)=(-1)^{n+1}q^n
 \mathcal R_{M,d}\!\left(\frac{1-q}{q}\right).}
 \tag{V91.8}
\]
In particular, \(q=1\) is a simple zero of \(H_{M,d}\).  Under
\(q=e^{-\varepsilon u}\), the nonendpoint zeros of \(K_{M,d}\) in
\((0,L_f)\) correspond, with multiplicity and sign change preserved, to
the positive zeros of \(\mathcal S_{M,d}\) in
\[
 \boxed{0<y<Y_f,
 \qquad Y_f:=e^{\varepsilon L_f}-1.}
 \tag{V91.9}
\]
\end{theorem}

\begin{proof}
Newton interpolation and \(\deg\mu=n-1\) give
\[
 \mu(j)=\sum_{r=0}^{n-1}{j\choose r}\delta_r.
 \tag{V91.10}
\]
Using
\(
 {n\choose j}{j\choose r}={n\choose r}{n-r\choose j-r}
\)
in \textup{(V88.14)}, one obtains
\[
 \begin{aligned}
 H_{M,d}(q)
 &=\sum_{r=0}^{n-1}{n\choose r}\delta_r
   \sum_{j=r}^{n}(-1)^j{n-r\choose j-r}q^j\\
 &=\sum_{r=0}^{n-1}(-1)^r{n\choose r}\delta_r
   q^r(1-q)^{n-r}.
 \end{aligned}
 \tag{V91.11}
\]
Set \(s=n-r\), factor \((-1)^{n+1}q^n\), and use
\textup{(V91.6)}--\textup{(V91.7)} to prove \textup{(V91.8)}.
The coefficient of \(y\) in \(\mathcal R\) is
\(A_1=n\delta_{n-1}>0\), so its zero at the origin is simple.
Finally,
\[
 \frac{1-e^{-\varepsilon u}}{e^{-\varepsilon u}}
 =e^{\varepsilon u}-1
 \tag{V91.12}
\]
is strictly increasing and has positive derivative.  The nonvanishing
prefactors in \textup{(V88.13)} and \textup{(V91.8)} preserve zero
multiplicity and sign changes, proving \textup{(V91.9)}.
\end{proof}

This identity is materially different from applying the finite multiplier
test to each root of \(P_{M,d}\).  It uses the complete product first and
then asks only for zeros in a shrinking interval.  In the remaining range
\(d=O(\sqrt M)\),
\[
 Y_f=\frac{\theta dL_f}{RM}+O\!\left(\frac{d^2}{M^2}\right)
 =O(M^{-1/2}),
 \tag{V91.13}
\]
and for fixed \(d\) the interval is of order \(M^{-1}\).

\subsection{A finite Budan--Fourier obstruction}

For a real polynomial \(G\) of degree \(m\), let
\[
 V_G(x):=\operatorname{var}
 \bigl(G(x),G'(x),\ldots,G^{(m)}(x)\bigr),
 \tag{V91.14}
\]
where zero entries are deleted before counting sign variations.  The
derivatives needed here are the explicit finite sums
\[
 \mathcal S_{M,d}^{(k)}(Y)
 =\sum_{t=k}^{n-1}(-1)^t\frac{t!}{(t-k)!}
 {n\choose t+1}\delta_{n-t-1}Y^{t-k}.
 \tag{V91.15}
\]

\begin{theorem}[Endpoint-variation certificate]
\label{thm:v91-budan-certificate}
Let \(Z_{M,d}(Y)\) denote the number of real zeros of
\(\mathcal S_{M,d}\) in \((0,Y]\), counted with multiplicity.  Then
\[
 \boxed{
 Z_{M,d}(Y)\le n-1-V_{\mathcal S_{M,d}}(Y).}
 \tag{V91.16}
\]
Since \(n=RM+3\), either of the conditions
\[
 \boxed{
 V_{\mathcal S_{M,d}}(Y_f)\ge (R-1)M+1}
 \tag{V91.17}
\]
or
\[
 \boxed{Z_{M,d}(Y_f)\le M+1}
 \tag{V91.18}
\]
is sufficient to prove that the V88 filter is not Euler admissible.
Condition \textup{(V91.18)} can be decided exactly by the Sturm chain of
\(\mathcal S_{M,d}\); condition \textup{(V91.17)} is the cheaper
derivative-sign certificate.
\end{theorem}

\begin{proof}
From \textup{(V91.7)},
\[
 \mathcal S_{M,d}^{(k)}(0)
 =k!(-1)^k A_{k+1}\ne0,
 \qquad 0\le k\le n-1.
 \tag{V91.19}
\]
Thus the derivative signs at zero alternate and
\[
 V_{\mathcal S_{M,d}}(0)=n-1.
 \tag{V91.20}
\]
The Budan--Fourier theorem bounds the number of zeros in \((0,Y]\),
counted with multiplicity, by
\(
 V_{\mathcal S}(0)-V_{\mathcal S}(Y)
\), proving \textup{(V91.16)}.  Under \textup{(V91.17)}, its right side is
at most
\[
 RM+2-\{(R-1)M+1\}=M+1.
 \tag{V91.21}
\]
On the other hand, Theorem~\ref{thm:v87-oscillation-compression} requires
at least \(M+2\) sign changes of \(K_{M,d}\) in \((0,L_f)\) if the filter
is Euler admissible.  Theorem~\ref{thm:v91-newton-recoding} sends them to at
least \(M+2\) real zeros of \(\mathcal S_{M,d}\) in \((0,Y_f)\), a
contradiction.  An exact Sturm count gives the same contradiction directly
under \textup{(V91.18)}.
\end{proof}

\begin{corollary}[The remaining equispaced problem is a sign ledger]
\label{cor:v91-sign-ledger}
To close the full unresolved range \textup{(V90.27)}, it is enough to prove
uniformly there, for all sufficiently large \(M\), either
\[
 V_{\mathcal S_{M,d}}(e^{\varepsilon L_f}-1)
 \ge (R-1)M+1
 \tag{V91.22}
\]
or the sharper Sturm bound \textup{(V91.18)}.  No global real-rootedness
statement for \(H_{M,d}\) is required.
\end{corollary}

\begin{proof}
This is Theorem~\ref{thm:v91-budan-certificate} applied to every parameter
sequence in \textup{(V90.27)}.
\end{proof}

The entries in \textup{(V91.15)} involve cancellation, but their inputs do
not: Lemma~\ref{lem:v91-positive-newton} expresses every \(\delta_r\) as a
positive Stirling transform.  Therefore interval bounds on the coefficients
of \(\mu\), rather than subtraction of nearly equal sampled values, are the
appropriate route to a rigorous asymptotic sign ledger.

\subsection{One failed pair as a logarithmic-derivative budget}

There is a complementary way to exploit the good tail from V90.  Let
\[
 G(q)=C\prod_{k=1}^{n}\left(1-\frac q{r_k}\right),
 \qquad C>0,\quad r_k>0,
 \tag{V91.23}
\]
and suppose \(0<q<r_1:=\min r_k\).  Put
\[
 L_1(q):=\sum_{k=1}^{n}\frac q{r_k-q},
 \qquad
 L_2(q):=\sum_{k=1}^{n}\frac{qr_k}{(r_k-q)^2}.
 \tag{V91.24}
\]

\begin{lemma}[Exact action of a conjugate quadratic]
\label{lem:v91-log-derivative}
For \(\mathcal D_q=q\,d/dq\), \(u>0\), and \(v\in\mathbb R\),
\[
 \boxed{
 \frac{\{(\mathcal D_q+u)^2+v^2\}G(q)}{G(q)}
 =(u-L_1(q))^2+v^2-L_2(q).}
 \tag{V91.25}
\]
Moreover,
\[
 L_2=L_1+\sum_{k=1}^{n}
 \left(\frac q{r_k-q}\right)^2
 \le L_1+L_1^2,
 \tag{V91.26}
\]
and hence
\[
 \boxed{
 \frac{\{(\mathcal D_q+u)^2+v^2\}G(q)}{G(q)}
 \ge u^2+v^2-(2u+1)L_1(q).}
 \tag{V91.27}
\]
In particular, the failed conjugate factor is nevertheless positive at
\(q\) whenever
\[
 L_1(q)<\frac{u^2+v^2}{2u+1}.
 \tag{V91.28}
\]
\end{lemma}

\begin{proof}
Logarithmic differentiation of \textup{(V91.23)} gives
\[
 \frac{\mathcal D_qG}{G}=-L_1,
 \qquad
 \mathcal D_q^2\log G=-L_2.
 \tag{V91.29}
\]
Thus
\(
 \mathcal D_q^2G/G=L_1^2-L_2
\), which proves \textup{(V91.25)} after expansion.  If
\(x_k=q/(r_k-q)>0\), then
\(
 qr_k/(r_k-q)^2=x_k+x_k^2
\).  Hence \textup{(V91.26)} follows from
\(
 \sum x_k^2\le(\sum x_k)^2
\).  Substitution in \textup{(V91.25)} proves
\textup{(V91.27)}--\textup{(V91.28)}.
\end{proof}

For a V90 conjugate root \(a\pm ib\), the paired operator in
\textup{(V88.15)} is the positive constant
\(\varepsilon^2/(a^2+b^2)\) times
\[
 (\mathcal D_q+a/\varepsilon)^2+(b/\varepsilon)^2.
 \tag{V91.30}
\]
Thus Lemma~\ref{lem:v91-log-derivative} applies without a sign change after
all certified factors have been composed.  It does not yet iterate through
the entire bad prefix: after one failed pair, the output need not retain
positive real roots.  Any iteration must preserve an explicit invariant,
not silently reuse \textup{(V91.23)}.

A recent generalized Eulerian result shows that certain complex-root
quadratics can still yield real-rooted Eulerian numerators.\footnote{D.
Karp and A. Vishnyakova, ``The series of generalized \(f\)-Eulerian
polynomials,'' arXiv:2509.13897 (2025),
\url{https://arxiv.org/abs/2509.13897}.}  Its recursive hypotheses concern
ordinary \(f\)-Eulerian transforms with real parameter restrictions, so it
does not imply real-rootedness of the finite binomial transform
\textup{(V88.14)}.  The comparison supports testing the complete product,
but supplies no missing theorem here.

\begin{firewall}[A numerical sign ledger is not an asymptotic proof]
High-precision experiments can evaluate \textup{(V91.15)} or a Sturm chain
for fixed \(M\), but floating-point signs near the endpoint are not a
uniform certificate.  The required result is an analytic lower bound for
the variation count, or validated interval/Sturm arithmetic together with
a proof that the certificate persists throughout a moving parameter range.
The exact recoding does not itself prove GRH.
\end{firewall}

\begin{theorem}[V91 local-zero verdict]
\label{thm:v91-verdict}
For the unresolved V88 equispaced family:
\begin{enumerate}[label=\textup{(\roman*)},leftmargin=3em]
\item the kernel polynomial has the exact reversed Newton representation
      \textup{(V91.8)}, with strictly positive Newton data;
\item Euler admissibility forces at least \(M+2\) zeros of the reduced
      polynomial \(\mathcal S_{M,d}\) inside the explicit shrinking
      interval \((0,e^{\varepsilon L_f}-1)\);
\item the derivative-variation inequality \textup{(V91.17)}, or an exact
      Sturm count at most \(M+1\), rules out admissibility without any
      assertion about the other zeros;
\item one failed conjugate factor acting on the certified tail is governed
      exactly by the logarithmic-derivative budget \textup{(V91.25)}.
\end{enumerate}
The unresolved task is now the uniform sign ledger
\textup{(V91.22)} or a replacement invariant that permits iteration of
\textup{(V91.28)}.  Arbitrary pole meshes, the target-zero separation
product, and GRH remain open.
\end{theorem}

\begin{assessment}[V91 closure and V92 direction]
V91 removes a genuine excess requirement from V90.  The complete
equispaced obstruction no longer depends on proving that all \(n\) roots
of \(H_{M,d}\) are positive.  It is enough to keep more than
\((R-1)M\) derivative sign variations at one explicit endpoint.  All
coefficients feeding that endpoint test are positive Stirling transforms,
which gives a cancellation-free starting point for asymptotics.

V92 should estimate \textup{(V91.15)} after the scale
\(Y=\varepsilon u+O(\varepsilon^2)\), first for fixed \(d\) and then
uniformly for \(d\le\kappa_{\rm J}\sqrt M\).  If the derivative ledger
does not remain above \((R-1)M\), the upstream alternative is to bound the
good-tail logarithmic derivative in \textup{(V91.28)} and identify an
invariant stable under the \(O(\sqrt M)\) bad-pair composition.  Repeating
the factorwise Jensen test would not address either criterion.
\end{assessment}

\section*{Version V91 append-only record}
\addcontentsline{toc}{section}{Version V91 append-only record}
The complete V90 source was retained before this continuation.  It had
1319609 bytes, 38915 lines, and SHA--256
\[
 \texttt{18EF6AE9F7C4128A04B51179FE88BB860D0FFE8ACC3A1AB7856D368E15BBE381}.
\]
V91 derived the exact reversed Newton representation of the filter,
converted first-prime oscillation compression into a finite
Budan--Fourier/Sturm certificate, and expressed one failed conjugate factor
as an explicit logarithmic-derivative budget on the certified tail.

% =============================================================================
% END APPENDED V91 MATERIAL
% =============================================================================
% =============================================================================
% BEGIN APPENDED V92 MATERIAL
% =============================================================================

\section{V92: the order ladder and its moving conjugate frontier}
\label{sec:grh-v92}

V91 converted the first-prime question into the signs of all derivatives of
one reversed Newton polynomial at
\(Y_f=e^{\varepsilon L_f}-1\).  The coefficients in those derivatives are
positive Stirling transforms, but their alternating sums still obscure the
required variation count.  This version identifies the entire derivative
ledger with a ladder of lower-order Jensen transforms.  That ladder obeys an
exact triangular recurrence and a discrete Riccati law: its sign changes are
precisely threshold crossings of one ratio against \(Y_f\).

The V90 conjugate-prefix calculation also persists at every order on this
ladder.  At order \(m\), only \(O(\sqrt m)\) conjugate factors fail the finite
multiplier test, and their exact entry frontier is computed below.  Hence the
V91 certificate has a triangular defect region, rather than an unrelated
degree-\(n\) polynomial at each derivative.

\subsection{Every derivative is a lower-order Jensen transform}

Keep the polynomial \(\mu\) from \textup{(V91.1)} and define, for
\(0\le m\le n\),
\[
 \mathcal J_m(q)
 :=\sum_{j=0}^{m}(-1)^j{m\choose j}\mu(j)q^j.
 \tag{V92.1}
\]
Thus \(\mathcal J_n=H_{M,d}\), whereas \(\mathcal J_0=1\).

\begin{theorem}[Derivative--order duality]
\label{thm:v92-derivative-order}
Let \(q=(1+Y)^{-1}\).  For \(0\le m\le n\),
\[
 \boxed{
 \mathcal R_{M,d}^{(n-m)}(Y)
 =(-1)^{n+1}(n-m)!{n\choose m}
  (1+Y)^m\mathcal J_m(q).}
 \tag{V92.2}
\]
Consequently
\[
 \boxed{
 V_{\mathcal R_{M,d}}(Y)
 =\operatorname{var}\bigl(
 \mathcal J_0(q),\mathcal J_1(q),\ldots,
 \mathcal J_n(q)\bigr),}
 \tag{V92.3}
\]
where zero entries are deleted.  In particular, the sufficient V91
obstruction has the equivalent order-ladder form
\[
 \boxed{
 \operatorname{var}\bigl(
 \mathcal J_0(e^{-\varepsilon L_f}),\ldots,
 \mathcal J_n(e^{-\varepsilon L_f})\bigr)
 \ge (R-1)M+1.}
 \tag{V92.4}
\]
\end{theorem}

\begin{proof}
Formula \textup{(V91.8)} is equivalent to
\[
 \mathcal R_{M,d}(Y)
 =(-1)^{n+1}
 \sum_{j=0}^{n}(-1)^j{n\choose j}\mu(j)(1+Y)^{n-j}.
 \tag{V92.5}
\]
Differentiate \(k=n-m\) times.  Terms with \(j>m\) vanish, and
\[
 {n\choose j}(n-j)_k
 =k!{n\choose k}{n-k\choose j}
 =(n-m)!{n\choose m}{m\choose j}.
 \tag{V92.6}
\]
Factoring \((1+Y)^m\) proves \textup{(V92.2)}.  All prefactors have the
same nonzero sign, and reversing a finite list does not change its number
of sign variations.  This proves \textup{(V92.3)}.

The polynomial \(\mathcal R=y\mathcal S\) has the same zeros as
\(\mathcal S\) away from zero, while
\(
 V_{\mathcal R}(0)=n-1
\)
after its initial zero is deleted.  Applying Budan--Fourier directly to
\(\mathcal R\) therefore gives the same upper bound as
Theorem~\ref{thm:v91-budan-certificate}.  At \(Y=Y_f\), one has
\(q=(1+Y_f)^{-1}=e^{-\varepsilon L_f}\), proving
\textup{(V92.4)}.
\end{proof}

The order ladder avoids forming high derivatives of a degree-\(n\)
polynomial.  It also exposes the common multiplier data \(\mu(0),\mu(1),
\ldots\) at every level.

\subsection{A triangular recurrence and exact threshold crossings}

For \(0\le r\le n\) and \(0\le m\le n\), set
\[
 \mathcal J_m^{[r]}(q)
 :=\sum_{j=0}^{m}(-1)^j{m\choose j}
 \Delta^r\mu(j)q^j.
 \tag{V92.7}
\]
with the convention \(\Delta^n\mu\equiv0\).  The superscript is a
forward-difference level, not a derivative.  In particular,
\[
 \mathcal J_m^{[0]}=\mathcal J_m,
 \qquad
 \mathcal J_0^{[r]}=\delta_r>0\quad(0\le r<n),
 \qquad \mathcal J_m^{[n]}=0.
 \tag{V92.8}
\]

\begin{lemma}[Pascal recurrence]
\label{lem:v92-pascal}
For \(0\le m<n\) and \(0\le r<n\),
\[
 \boxed{
 \mathcal J_{m+1}^{[r]}(q)
 =(1-q)\mathcal J_m^{[r]}(q)
 -q\mathcal J_m^{[r+1]}(q).}
 \tag{V92.9}
\]
Equivalently, with \(Y=(1-q)/q\),
\[
 \boxed{
 \mathcal J_{m+1}^{[r]}(q)
 =q\{Y\mathcal J_m^{[r]}(q)
      -\mathcal J_m^{[r+1]}(q)\}.}
 \tag{V92.10}
\]
\end{lemma}

\begin{proof}
Use
\(
 {m+1\choose j}={m\choose j}+{m\choose j-1}
\)
in \textup{(V92.7)}.  After shifting the second sum, use
\[
 \Delta^r\mu(j+1)
 =\Delta^r\mu(j)+\Delta^{r+1}\mu(j).
 \tag{V92.11}
\]
The result is \textup{(V92.9)}, and
\(1-q=qY\) gives \textup{(V92.10)}.
\end{proof}

Whenever the denominators are nonzero, define, for \(0\le r<n\),
\[
 X_m^{[r]}(Y)
 :=\frac{\mathcal J_m^{[r+1]}(q)}
          {\mathcal J_m^{[r]}(q)},
 \qquad q=(1+Y)^{-1}.
 \tag{V92.12}
\]

\begin{theorem}[Discrete Riccati ledger]
\label{thm:v92-riccati}
The ratio array starts from cancellation-free positive data,
\[
 \boxed{X_0^{[r]}=\frac{\delta_{r+1}}{\delta_r}>0
 \quad(0\le r<n-1),}
 \tag{V92.13}
\]
while \(X_0^{[n-1]}=0\), and it obeys, wherever the displayed ratios are
defined and \(0\le r<n-1\),
\[
 \boxed{
 X_{m+1}^{[r]}
 =X_m^{[r]}
  \frac{Y-X_m^{[r+1]}}{Y-X_m^{[r]}}.}
 \tag{V92.14}
\]
If none of
\(\mathcal J_0(q),\ldots,\mathcal J_n(q)\) vanishes, then
\[
 \boxed{
 \operatorname{var}(\mathcal J_0(q),\ldots,\mathcal J_n(q))
 =\#\{0\le m<n:X_m^{[0]}(Y)>Y\}.}
 \tag{V92.15}
\]
Therefore, in the nondegenerate case, it suffices for the GRH filter
obstruction to prove
\[
 \boxed{
 \#\{0\le m<n:X_m^{[0]}(Y_f)>Y_f\}
 \ge(R-1)M+1.}
 \tag{V92.16}
\]
At a degenerate parameter the certificate is evaluated by deleting zero
entries in the order ladder, exactly as in \textup{(V91.14)}.
\end{theorem}

\begin{proof}
Equation \textup{(V92.13)} follows from \textup{(V92.8)} and
Lemma~\ref{lem:v91-positive-newton}.  Apply \textup{(V92.10)} at levels
\(r\) and \(r+1\), then divide, to obtain \textup{(V92.14)}.
Also,
\[
 \frac{\mathcal J_{m+1}^{[0]}(q)}
      {\mathcal J_m^{[0]}(q)}
 =q\{Y-X_m^{[0]}(Y)\}.
 \tag{V92.17}
\]
Because \(q>0\), two adjacent nonzero order-ladder entries have opposite
signs exactly when \(X_m^{[0]}>Y\).  Summing these adjacent sign changes
proves \textup{(V92.15)}.  Combine it with
\textup{(V92.4)} for \textup{(V92.16)}.  The convention for zero entries
is precisely the convention in \textup{(V91.14)}.
\end{proof}

Thus the difficult cancellations have been moved into a deterministic
ratio dynamics whose entire initial edge consists of positive Stirling
transforms.  A proof of \textup{(V92.16)} may use invariant regions for
\textup{(V92.14)}; direct floating-point subtraction of the
\(\mathcal J_m\) is unnecessary.

The initial edge has additional root geometry which is not visible in the
Stirling formula.

\begin{lemma}[Forward differences preserve the Hurwitz half-plane]
\label{lem:v92-hurwitz-differences}
Let \(f\) be a nonconstant polynomial all of whose zeros satisfy
\(\Re\rho<0\).  Then
\[
 \Delta f(x):=f(x+1)-f(x)
 \tag{V92.17a}
\]
has degree \(\deg f-1\), and all of its zeros also satisfy
\(\Re\rho<0\).  Consequently, for \(0\le r<n\),
\[
 \boxed{
 f_r(x):=\Delta^r\mu(x)
 \text{ has degree }n-1-r
 \text{ and all zeros in }\Re x<0.}
 \tag{V92.17b}
\]
Writing \(\alpha_-(g)\) and \(\alpha_+(g)\) for the least and greatest
real parts of the zeros of a nonconstant polynomial \(g\), one has the
quantitative enclosure
\[
 \boxed{
 \alpha_-(\mu)-\frac r2
 \le \Re\rho\le
 \alpha_+(\mu)-\frac r2
 \quad\text{for every zero }\rho\text{ of }f_r.}
\]
If \(0\le r<n-1\) and \(\mathcal Z_r\) is the zero multiset of \(f_r\), then
\[
 \boxed{
 X_0^{[r]}
 =\frac{f_r(1)}{f_r(0)}-1
 =\prod_{\rho\in\mathcal Z_r}
   \left|1-\frac1\rho\right|-1>0,}
 \tag{V92.17c}
\]
where conjugate factors are paired and a negative real factor is read as
the corresponding positive real ratio.  If \(D=n-1\), then the root
centroids satisfy the exact drift law
\[
 \boxed{
 \frac1{D-r}\sum_{\rho\in\mathcal Z_r}\rho
 =\frac1D\sum_{\rho\in\mathcal Z_0}\rho-\frac r2.}
 \tag{V92.17c'}
\]
Moreover, at the first-prime endpoint,
\[
 \boxed{
 X_0^{[r]}>Y_f
 \quad\Longleftrightarrow\quad
 \sum_{\rho\in\mathcal Z_r}
 \log\left|1-\frac1\rho\right|>\varepsilon L_f.}
 \tag{V92.17c''}
\]
\end{lemma}

\begin{proof}
Write \(f(z)=c\prod_\rho(z-\rho)\).  If \(\Re z\ge0\), then
\(w=z-\rho\) has positive real part and
\[
 |w+1|^2-|w|^2=2\Re w+1>0.
 \tag{V92.17d}
\]
If \(\Re z>\alpha_+(f)-1/2\), then every factor in
\(f(z+1)/f(z)\) has modulus greater than one by
\textup{(V92.17d)}.  If \(\Re z<\alpha_-(f)-1/2\), every factor has
modulus less than one.  Thus \(\Delta f(z)\ne0\) outside the closed strip
\[
 \alpha_-(f)-\frac12\le\Re z\le\alpha_+(f)-\frac12.
\]
In particular, \(|f(z+1)|>|f(z)|\) on the closed right half-plane, so
\(f(z+1)-f(z)\ne0\) there.  Cancellation of the leading terms lowers the
degree by exactly one.  The roots of
\(\mu(x)=P_{M,d}(-\varepsilon x)/P_{M,d}(0)\) are
\(-\lambda/\varepsilon\), with \(\Re\lambda>0\); induction proves
\textup{(V92.17b)} and the displayed strip enclosure.

Finally, \(X_0^{[r]}=\delta_{r+1}/\delta_r
=\{f_r(1)-f_r(0)\}/f_r(0)\).  Factoring the quotient
\(f_r(1)/f_r(0)\) proves \textup{(V92.17c)}.  Every paired modulus ratio is
strictly greater than one by \textup{(V92.17d)}, as is every negative-real
root ratio.

If \(g(x)=a_Dx^D+a_{D-1}x^{D-1}+\cdots\), the two leading coefficients of
\(\Delta g\) are
\[
 Da_D,
 \qquad {D\choose2}a_D+(D-1)a_{D-1}.
\]
Vieta's formula therefore shows that the centroid of the zeros of
\(\Delta g\) is the centroid of the zeros of \(g\) minus \(1/2\).
Iteration proves \textup{(V92.17c')}.  Finally, take logarithms in
\textup{(V92.17c)} and use
\(1+Y_f=e^{\varepsilon L_f}\) to prove
\textup{(V92.17c'')}.
\end{proof}

\begin{firewall}[Imaginary spreading survives the real drift]
The strip enclosure controls real parts only.  There is no corresponding
convex-hull bound in the imaginary direction.  Indeed,
\[
 \Delta x^D=(x+1)^D-x^D
\]
has the \(D-1\) zeros
\[
 \boxed{
 \rho_k=\frac1{e^{2\pi i k/D}-1}
 =-\frac12-\frac i2\cot\frac{\pi k}{D},
 \qquad 1\le k<D.}
 \tag{V92.17e}
\]
Although the original zero set is \(\{0\}\), the largest imaginary part
has magnitude \(D/(2\pi)+O(D^{-1})\).  Thus the real-part drift in the
lemma cannot by itself turn the product \textup{(V92.17c)} into a useful
lower bound; additional structure of the gamma-root polynomial is required.
\end{firewall}
Thus estimating the positive starting ratios is equivalent to tracking how
successive forward differences move a Hurwitz zero cloud.  General finite
difference preservation theorems require careful alignment between the
step and the boundary line; the direct argument above verifies that
alignment for the present left half-plane and needs no external
hypothesis.\footnote{For the general line, strip, and half-plane theory see
O. Katkova, M. Tyaglov, and A. Vishnyakova, ``Hermite--Poulain theorems for
linear finite difference operators,'' arXiv:1901.06398 (2019),
\url{https://arxiv.org/abs/1901.06398}.}
\subsection{The conjugate defect frontier at every order}

For a conjugate gamma factor at physical real part \(a>0\) and imaginary
part \(b\), put as before
\[
 u=\frac a\varepsilon,
 \qquad v=\frac b\varepsilon.
 \tag{V92.18}
\]
The V90 criterion applies with the transform degree \(m\), not only with
the top degree \(n\).

\begin{proposition}[Order-adaptive factor threshold]
\label{prop:v92-order-threshold}
For \(m\ge2\), the conjugate factor is a finite multiplier sequence
through degree \(m\) if and only if
\[
 \boxed{
 a\ge\max\{0,a_{\rm J}(m,\varepsilon,b)\},}
 \tag{V92.19}
\]
where
\[
 a_{\rm J}(m,\varepsilon,b)
 :=\frac{-m\varepsilon+
 \sqrt{(m-1)(m\varepsilon^2+4b^2)}}2.
 \tag{V92.20}
\]
Equivalently, if
\(b^2>\varepsilon a+\varepsilon^2/4\), the factor is certified exactly
for
\[
 \boxed{
 m\le m_\star(a,b):=
 \frac{a^2+b^2}{b^2-\varepsilon a-\varepsilon^2/4};}
 \tag{V92.21}
\]
if the denominator is nonpositive, it is certified at every order.
Furthermore,
\[
 \boxed{a_{\rm J}(m,\varepsilon,b)
 \le |b|\sqrt{m-1}.}
 \tag{V92.22}
\]
\end{proposition}

\begin{proof}
Replace \(n\) by \(m\) in \textup{(V90.23)}.  The exact condition is
\[
 a^2+m\varepsilon a+\frac{m\varepsilon^2}{4}
 -(m-1)b^2\ge0.
 \tag{V92.23}
\]
Solving the increasing quadratic in \(a>0\) gives
\textup{(V92.19)}--\textup{(V92.20)}.  Rearranging it as
\[
 a^2+b^2
 \ge m\{b^2-\varepsilon a-\varepsilon^2/4\}
 \tag{V92.24}
\]
proves \textup{(V92.21)} and the denominator alternative.  Finally,
\[
 \begin{aligned}
 \sqrt{(m-1)(m\varepsilon^2+4b^2)}
 &\le\varepsilon\sqrt{m(m-1)}+2|b|\sqrt{m-1}\\
 &\le m\varepsilon+2|b|\sqrt{m-1},
 \end{aligned}
 \tag{V92.25}
\]
which proves \textup{(V92.22)}.
\end{proof}

For each conjugate-pair ladder \(\ell\), write
\[
 a_{\ell,r}=d+1+\alpha_\ell+hr,
 \qquad 0\le r<M,
 \tag{V92.26}
\]
and let \(B_m\) be the total number of conjugate factors which fail at
order \(m\).  The pair is counted once, not once per conjugate member.

\begin{theorem}[Triangular bad-prefix law]
\label{thm:v92-triangular-prefix}
The exact contribution of ladder \(\ell\) to \(B_m\) is
\[
 \boxed{
 B_{m,\ell}
 =\min\!\left\{M,
 \max\!\left\{0,
 \left\lceil
 \frac{a_{\rm J}(m,\varepsilon,b_\ell)
       -d-1-\alpha_\ell}{h}
 \right\rceil\right\}\right\}.}
 \tag{V92.27}
\]
In particular,
\[
 \boxed{
 B_m\le\frac{\sqrt{m-1}}h
 \sum_{\ell}|b_\ell|+O(1),}
 \tag{V92.28}
\]
uniformly in \(0\le m\le n\).

Suppose \(d/\sqrt M\to\kappa\ge0\) and \(m/M\to\tau\in(0,R]\).  Define
\[
 A_{\rm J}(\tau,\kappa,b)
 :=\frac{-\tau\theta\kappa/R+
 \sqrt{\tau^2\theta^2\kappa^2/R^2+4\tau b^2}}2.
 \tag{V92.29}
\]
Then
\[
 \boxed{
 \frac{B_m}{\sqrt M}\longrightarrow
 \frac1h\sum_{\ell}
 \bigl[A_{\rm J}(\tau,\kappa,b_\ell)-\kappa\bigr]_+.}
 \tag{V92.30}
\]
In the lower-scale case \(d=o(\sqrt M)\), this becomes
\[
 \boxed{
 \frac{B_m}{\sqrt M}\longrightarrow
 \frac{\sqrt\tau}{h}\sum_{\ell}|b_\ell|.}
 \tag{V92.31}
\]
\end{theorem}

\begin{proof}
A factor fails precisely when
\(
 a_{\ell,r}<a_{\rm J}(m,\varepsilon,b_\ell)
\).  Counting the integers \(0\le r<M\) satisfying this strict inequality
gives \textup{(V92.27)}.  Bound \textup{(V92.22)} then proves
\textup{(V92.28)}; fixed offsets contribute only \(O(1)\).

Since \(N/M\to R\),
\[
 \frac{m\varepsilon}{\sqrt M}
 \longrightarrow\frac{\tau\theta\kappa}{R},
 \qquad
 m\varepsilon^2
 \longrightarrow\frac{\tau\theta^2\kappa^2}{R^2}.
 \tag{V92.32}
\]
Divide \textup{(V92.20)} by \(\sqrt M\) to obtain
\(a_{\rm J}/\sqrt M\to A_{\rm J}\).  Also
\(a_{\ell,0}/\sqrt M\to\kappa\), and the upper truncation by \(M\)
is inactive because the count is \(O(\sqrt M)\).  This proves
\textup{(V92.30)}.  Setting \(\kappa=0\) gives
\(A_{\rm J}=|b|\sqrt\tau\), proving \textup{(V92.31)}.
\end{proof}

Let \(T_{{\rm good},m}\) be the product of the linear factors and the
conjugate factors outside \textup{(V92.27)}, and let
\(T_{{\rm bad},m}\) be the remaining product.  Proposition
\ref{prop:v92-order-threshold} and composition of finite multiplier
sequences give the exact factorization
\[
 \boxed{
 \mathcal J_m(q)
 =T_{{\rm bad},m}\,G_m(q),
 \qquad
 G_m:=T_{{\rm good},m}[(1-q)^m],}
 \tag{V92.33}
\]
where every zero of \(G_m\) is positive real.  Thus every order-ladder
entry has a certified core and only the triangular defect count
\textup{(V92.27)}.

One might try to control the logarithmic derivative in V91 by invoking the
fact that finite multiplier sequences do not decrease logarithmic root
mesh.\footnote{O. Katkova, B. Shapiro, and A. Vishnyakova,
``Multiplier sequences and logarithmic mesh,'' arXiv:1010.6052 (2010),
\url{https://arxiv.org/abs/1010.6052}.}  That theorem is applicable to the
certified composition, but it gives no strict bound here: the elementary
degree-\(m\) Jensen polynomials contain \((1+x)^{m-1}\) or
\((1+x)^{m-2}\), so for large \(m\) their logarithmic mesh is exactly one.
This route therefore does not by itself bound \(L_1\) in
\textup{(V91.28)}.

\begin{firewall}[A square-root defect is not automatically perturbative]
The estimate \(B_m=O(\sqrt m)\) is a factor count, not a norm bound.  A
single uncertified quadratic may create a nonreal pair or reverse a sign,
and the Riccati recursion can approach a pole when
\(X_m^{[r]}=Y_f\).  Therefore neither \textup{(V92.28)} nor the positive
core \textup{(V92.33)} proves the variation target
\textup{(V92.16)} without an invariant region or a quantitative
logarithmic-derivative estimate.  GRH remains open.
\end{firewall}

\begin{theorem}[V92 order-ladder verdict]
\label{thm:v92-verdict}
For the unresolved V88 equispaced family:
\begin{enumerate}[label=\textup{(\roman*)},leftmargin=3em]
\item the V91 endpoint derivative sequence is exactly the order ladder
      \(\mathcal J_0(q_f),\ldots,\mathcal J_n(q_f)\);
\item that ladder is generated from positive Newton data by the triangular
      recurrence \textup{(V92.10)};
\item away from zero denominators, the required variation count is exactly
      the number of Riccati threshold crossings in \textup{(V92.16)};
\item at every order \(m\), all but the explicit \(O(\sqrt m)\) conjugate
      prefix are certified finite multipliers;
\item on critical lines the complete moving frontier is the limit
      \textup{(V92.30)}.
\end{enumerate}
This does not decide the threshold-crossing count, arbitrary pole meshes,
or the target-zero separation product, and hence does not prove GRH.
\end{theorem}

\begin{assessment}[V92 closure and V93 direction]
V92 converts V91's collection of derivative signs into one coherent
two-dimensional dynamics.  The initial edge
\(\delta_{r+1}/\delta_r\) is positive and cancellation-free, every sign
flip has the exact owner \(X_m^{[0]}>Y_f\), and the only failure of the
factorwise real-rooted core lies below a moving square-root frontier.

V93 should seek invariant regions for \textup{(V92.14)} after scaling
\(Y_f\sim\theta dL_f/(RM)\).  The first concrete test is whether a block of
at least \((R-1)M+1\) indices can be forced above threshold using upper and
lower bounds for the positive initial ratios
\(\delta_{r+1}/\delta_r\).  If Riccati poles prevent that induction, the
same order-adaptive factorization should be combined with the V91
logarithmic-derivative identity; logarithmic mesh alone has now been ruled
out as an adequate estimate.
\end{assessment}

\section*{Version V92 append-only record}
\addcontentsline{toc}{section}{Version V92 append-only record}
The complete V91 source was retained before this continuation.  It had
1332974 bytes, 39324 lines, and SHA--256
\[
 \texttt{84D439F6BB54405CAF5A8D29AF85E81755E61DC51B7DBDA206D78E9811EC8F68}.
\]
V92 identified the endpoint derivative sequence with a lower-order Jensen
ladder, derived its Pascal and Riccati recurrences, and proved an exact
order-adaptive square-root law for every uncertified conjugate prefix.

% =============================================================================
% END APPENDED V92 MATERIAL
% =============================================================================
% =============================================================================
% BEGIN APPENDED V93 MATERIAL
% =============================================================================

\section{V93: Peano potentials and the guaranteed Riccati input block}
\label{sec:grh-v93}

V92 reduced the local zero-count obstruction to threshold crossings in a
two-dimensional Riccati array.  Its initial edge is
\(X_0^{[r]}=\delta_{r+1}/\delta_r\), where the \(\delta_r\) are positive
Newton coefficients.  The forward-difference zero cloud supplied an exact
product formula, but it can spread by order \(n\) in the imaginary
direction.  This version changes representation rather than trying to
repair that false bound.

The same Newton coefficients admit a positive Peano-cube integral involving
ordinary derivatives of \(\mu\).  Those derivative zeros remain in the
convex hull of the original gamma-root cloud by Gauss--Lucas.  This gives a
uniform lower bound for the complete initial Riccati edge and identifies an
explicit macroscopic block which lies above the first-prime threshold.

\subsection{A positive cube formula for Newton coefficients}

Put
\[
 D:=\deg\mu=n-1=RM+2,
 \qquad g_r(x):=\mu^{(r)}(x).
 \tag{V93.1}
\]

\begin{lemma}[Peano-cube representation]
\label{lem:v93-peano-cube}
For \(0\le r\le D\),
\[
 \boxed{
 \delta_r
 =\int_{[0,1]^r}
 g_r(t_1+\cdots+t_r)\,dt_1\cdots dt_r,}
 \tag{V93.2}
\]
where the right side is \(\mu(0)\) when \(r=0\).  For \(0\le r<D\),
define a probability measure on \([0,1]^r\) by
\[
 d\nu_r(\mathbf t)
 :=\frac{g_r(t_1+\cdots+t_r)}{\delta_r}\,d\mathbf t.
 \tag{V93.3}
\]
Then
\[
 \boxed{
 X_0^{[r]}
 =\int_{[0,1]^r}
 \left\{
 \frac{g_r(S+1)}{g_r(S)}-1
 \right\}d\nu_r(\mathbf t),
 \qquad S=t_1+\cdots+t_r.}
 \tag{V93.4}
\]
\end{lemma}

\begin{proof}
The fundamental theorem of calculus gives
\[
 \Delta f(x)=\int_0^1 f'(x+t)\,dt.
 \tag{V93.5}
\]
Iteration and Fubini's theorem give
\[
 \Delta^r f(0)
 =\int_{[0,1]^r}f^{(r)}(t_1+\cdots+t_r)\,d\mathbf t.
 \tag{V93.6}
\]
Apply this to \(f=\mu\) and use \(\delta_r=\Delta^r\mu(0)\) to prove
\textup{(V93.2)}.  The polynomial \(g_r\) has positive coefficients, so
the density in \textup{(V93.3)} is positive and has total mass one.

Applying \textup{(V93.6)} to \(\Delta\mu\), or integrating once more,
gives
\[
 \delta_{r+1}
 =\int_{[0,1]^r}\{g_r(S+1)-g_r(S)\}\,d\mathbf t.
 \tag{V93.7}
\]
Divide by \(\delta_r\) to obtain \textup{(V93.4)}.
\end{proof}

Unlike an alternating finite-difference sum, every quantity in
\textup{(V93.2)}--\textup{(V93.4)} is positive.

\subsection{A derivative-cloud lower bound}

Let \(\Lambda_{M,d}\) be the root multiset of \(P_{M,d}\), as in V88, and
set
\[
 A_-:=\frac1\varepsilon
       \min_{\lambda\in\Lambda_{M,d}}\Re\lambda,
 \qquad
 A_+:=\frac1\varepsilon
       \max_{\lambda\in\Lambda_{M,d}}\Re\lambda,
 \qquad
 B_*:=\frac1\varepsilon
       \max_{\lambda\in\Lambda_{M,d}}|\Im\lambda|.
 \tag{V93.8}
\]
All three quantities are finite, and \(A_->0\).

\begin{theorem}[Uniform initial-edge potential bound]
\label{thm:v93-initial-potential}
For every \(0\le r<D\),
\[
 \boxed{
 \log(1+X_0^{[r]})
 \ge
 \frac{D-r}
 {r+1+A_++B_*^2/A_-}.}
 \tag{V93.9}
\]
Consequently,
\[
 \boxed{
 D-r>\varepsilon L_f
 \left(r+1+A_++\frac{B_*^2}{A_-}\right)}
 \tag{V93.10}
\]
is a sufficient condition for
\(X_0^{[r]}>Y_f=e^{\varepsilon L_f}-1\).
\end{theorem}

\begin{proof}
The roots of \(\mu\) are \(-\lambda/\varepsilon\).  They lie in the
rectangle
\[
 -A_+\le\Re z\le-A_-<0,
 \qquad |\Im z|\le B_*.
 \tag{V93.11}
\]
Repeated application of the Gauss--Lucas theorem shows that every zero
\(\rho=-a+ib\) of \(g_r=\mu^{(r)}\) lies in the same rectangle.  Hence,
for \(0\le x\le r+1\),
\[
 \begin{aligned}
 \frac{g_r'(x)}{g_r(x)}
 &=\sum_{\rho}\Re\frac1{x-\rho}\\
 &=\sum_{\rho}\frac{x+a}{(x+a)^2+b^2}\\
 &=\sum_{\rho}
   \frac1{x+a+b^2/(x+a)}\\
 &\ge
 \frac{D-r}{r+1+A_++B_*^2/A_-}.
 \end{aligned}
 \tag{V93.12}
\]
The sum is real because nonreal zeros are paired.  For every
\(S\in[0,r]\), integration over \([S,S+1]\) gives
\[
 \frac{g_r(S+1)}{g_r(S)}
 \ge
 \exp\!\left\{
 \frac{D-r}{r+1+A_++B_*^2/A_-}
 \right\}.
 \tag{V93.13}
\]
Insert this pointwise bound into the probability average
\textup{(V93.4)} to prove \textup{(V93.9)}.  If
\textup{(V93.10)} holds, then
\(
 1+X_0^{[r]}>e^{\varepsilon L_f}=1+Y_f
\), proving the final assertion.
\end{proof}

The estimate uses the derivative cloud, whose imaginary part is controlled
by convexity.  It does not contradict the V92 example
\(\Delta x^D\): forward differences and ordinary derivatives have
different zero geometry even though they encode the same Newton data
through \textup{(V93.2)}.

\subsection{The macroscopic guaranteed block}

The exact V88 root list gives, uniformly for \(d=o(M)\) and
\(d\ge d_0>0\),
\[
 \begin{aligned}
 A_+&=\frac{hRM^2}{\theta d}\{1+o(1)\},\\
 A_-&\asymp M,
 \qquad
 B_*=O(M/d).
 \end{aligned}
 \tag{V93.14}
\]
The fixed offsets only change lower-order terms.

\begin{theorem}[Scaled initial-crossing frontier]
\label{thm:v93-scaled-frontier}
Suppose
\[
 d_0\le d_M=O(\sqrt M),
 \qquad
 \frac{r_M}{M}\longrightarrow\tau\in[0,R).
 \tag{V93.15}
\]
Then
\[
 \boxed{
 \liminf_{M\to\infty}
 \frac1\varepsilon\log(1+X_0^{[r_M]})
 \ge\frac{R-\tau}{h}.}
 \tag{V93.16}
\]
In particular, for every \(\sigma>0\), all sufficiently large \(M\)
satisfy
\[
 \boxed{
 X_0^{[r]}>Y_f
 \quad\text{whenever}\quad
 0\le r\le
 (R-hL_f-\sigma)M,}
 \tag{V93.17}
\]
provided the right endpoint is nonnegative.
\end{theorem}

\begin{proof}
The largest real part among the physical roots of \(P_{M,d}\) is
\(d+hM+O(1)\), while their imaginary parts are bounded by the fixed
constant \(B\).  Since
\(\varepsilon=\theta d/(RM+4)\), this proves the first and third estimates
in \textup{(V93.14)}.  Positivity of all physical real parts and the fixed
lower offset bound give \(A_-\asymp M\), uniformly for
\(d_0\le d=O(\sqrt M)\).  Hence, for \(r/M\to\tau\),
\[
 r+1+\frac{B_*^2}{A_-}=o(A_+).
 \tag{V93.18}
\]
Since \(D=RM+2\), Theorem~\ref{thm:v93-initial-potential} yields
\[
 \log(1+X_0^{[r]})
 \ge
 \frac{\theta d(R-\tau)}{hRM}\{1+o(1)\}.
 \tag{V93.19}
\]
Division by
\(
 \varepsilon=\theta d/(RM)\{1+o(1)\}
\)
proves \textup{(V93.16)}.  If
\(r/M\le R-hL_f-\sigma\), the right side of
\textup{(V93.16)} exceeds \(L_f\) by a fixed margin.  Thus
\(\log(1+X_0^{[r]})>\varepsilon L_f\), uniformly over the displayed
block, which is equivalent to \textup{(V93.17)}.
\end{proof}

For the V80 arithmetic spacing \(h=2\) and first-prime frequency
\(L_f=\log2\), the guaranteed input block has asymptotic length
\[
 \boxed{(R-2\log2)M.}
 \tag{V93.20}
\]
The V92 Budan obstruction requires \((R-1)M+1\) sign changes on the
\emph{order} edge.  Since \(2\log2>1\), the lower bound
\textup{(V93.17)} alone leaves a macroscopic gap
\[
 (2\log2-1)M+O(1).
 \tag{V93.21}
\]
This comparison does not assert that the remaining initial ratios are
below threshold.  It identifies exactly what the present universal
rectangle estimate fails to guarantee and what the Riccati transport must
supply.

\begin{firewall}[Initial-edge growth is not order-edge variation]
The V92 obstruction counts \(X_m^{[0]}>Y_f\) as \(m\) varies.  Theorem
\ref{thm:v93-scaled-frontier} controls \(X_0^{[r]}\) as \(r\) varies.
The Riccati recurrence couples these two edges but is not monotone across a
pole.  Therefore the block \textup{(V93.17)} cannot simply be relabelled as
the required sign variations.  A transport invariant is still necessary,
and GRH remains open.
\end{firewall}

\begin{theorem}[V93 Peano-potential verdict]
\label{thm:v93-verdict}
For the remaining equispaced regime:
\begin{enumerate}[label=\textup{(\roman*)},leftmargin=3em]
\item every positive Newton coefficient is a positive Peano-cube integral;
\item each initial Riccati ratio is an average of a derivative-polynomial
      unit-step growth factor;
\item Gauss--Lucas gives the explicit uniform lower bound
      \textup{(V93.9)};
\item the whole input block \(r/M<R-hL_f\) lies above the first-prime
      threshold asymptotically;
\item for \(h=2\), the universal bound alone is short of the order-edge
      target by \((2\log2-1)M+O(1)\).
\end{enumerate}
Thus the next missing estimate is a Riccati transport law or a sharper
gamma-specific derivative-cloud potential, not positivity of the Newton
data.  Arbitrary pole meshes, target-zero separation, and GRH remain open.
\end{theorem}

\begin{assessment}[V93 closure and V94 direction]
V93 turns the qualitative Hurwitz information of V92 into a quantitative,
uniform first-prime comparison.  The cancellation-free input block is
macroscopic and its limiting boundary is explicit.  The constant comparison
also shows why a rectangle-only zero-cloud estimate cannot finish the V92
variation count at arithmetic spacing \(h=2\).

V94 should analyze the Riccati map
\[
 X_{m+1}^{[r]}=X_m^{[r]}
 \frac{Y_f-X_m^{[r+1]}}{Y_f-X_m^{[r]}}
\]
on the block supplied by \textup{(V93.17)}.  The immediate question is
whether threshold crossings propagate southwest through the array unless a
denominator pole itself creates an additional order-edge sign change.  If
that dichotomy fails, the derivative potential must use the actual gamma
root density rather than only its enclosing rectangle.
\end{assessment}

\section*{Version V93 append-only record}
\addcontentsline{toc}{section}{Version V93 append-only record}
The complete V92 source was retained before this continuation.  It had
1351770 bytes, 39916 lines, and SHA--256
\[
 \texttt{DA7FE953C3A1E03757207997DA665EC70DC4E1E5A0E7C5A9CA69FECCAF8726ED}.
\]
V93 derived a positive Peano-cube representation for the Newton data,
bounded its derivative-zero logarithmic potential, and proved the explicit
macroscopic initial crossing frontier \(r/M<R-hL_f\).

% =============================================================================
% END APPENDED V93 MATERIAL
% =============================================================================
% =============================================================================
% BEGIN APPENDED V94 MATERIAL
% =============================================================================

\section{V94: shifted-difference transport and the finite monotonicity defect}
\label{sec:grh-v94}

V93 produced a macroscopic block of positive initial inequalities
\(\delta_{r+1}>Y_f\delta_r\), but the V92 certificate requires sign
changes along the perpendicular, order edge.  The Riccati quotients are
useful away from their poles; for sign transport, however, the underlying
linear array is cleaner.  This version conjugates the entire Pascal array
to powers of one shifted forward-difference operator.

That representation also exposes a tempting but invalid shortcut.  A
Hurwitz polynomial need not have log-concave forward differences at the
origin, so neither ratio monotonicity nor total positivity follows from
the half-plane location proved in V92.  For the actual gamma factors a
more precise statement survives: every possible increase of the
unit-step factor ratio belongs to an explicit, \(M\)-independent
conjugate prefix.  Thus the square-root multiplier defect of V92 and the
adjacent-ratio monotonicity defect are genuinely different objects.

\subsection{The Riccati array is a shifted-difference orbit}

Let \(E_x f(x)=f(x+1)\), let \(\Delta_x=E_x-I\), and put
\[
 q=(1+Y)^{-1}.
 \tag{V94.1}
\]
Define the checkerboard-conjugated array
\[
 \mathcal C_m^{[r]}(Y)
 :=(-1)^m q^{-m}\mathcal J_m^{[r]}(q).
 \tag{V94.2}
\]

\begin{theorem}[Exact shifted-difference representation]
\label{thm:v94-shifted-difference}
For every \(m,r\ge0\) for which the entries occur in the V92 triangle,
\[
 \boxed{
 \mathcal C_m^{[r]}(Y)
 =\Delta_x^r\bigl(E_x-(1+Y)I\bigr)^m\mu(0)
 =\Delta_x^r\bigl(\Delta_x-YI\bigr)^m\mu(0).}
 \tag{V94.3}
\]
Equivalently, the array obeys
\[
 \boxed{
 \mathcal C_{m+1}^{[r]}
 =\mathcal C_m^{[r+1]}-Y\mathcal C_m^{[r]},
 \qquad
 \mathcal C_0^{[r]}=\delta_r.}
 \tag{V94.4}
\]
If no order-edge entry vanishes, then
\[
 \boxed{
 \operatorname{var}\bigl(
 \mathcal J_0(q),\ldots,\mathcal J_n(q)\bigr)
 =
 \#\{0\le m<n:
 \mathcal C_m^{[0]}\mathcal C_{m+1}^{[0]}>0\}.}
 \tag{V94.5}
\]
In particular, a constant-sign run
\(\mathcal C_a^{[0]},\ldots,\mathcal C_b^{[0]}\)
contributes exactly \(b-a\) order-edge variations.
\end{theorem}

\begin{proof}
Because \(E_x\) and \(\Delta_x\) commute, the binomial theorem gives
\[
 \begin{aligned}
 \Delta_x^r\bigl(E_x-(1+Y)I\bigr)^m\mu(0)
 &=
 \sum_{j=0}^m {m\choose j}
       \bigl(-(1+Y)\bigr)^{m-j}\Delta^r\mu(j)\\
 &=(-1)^m(1+Y)^m
   \sum_{j=0}^m(-1)^j{m\choose j}
       \Delta^r\mu(j)(1+Y)^{-j}\\
 &=(-1)^m q^{-m}\mathcal J_m^{[r]}(q).
 \end{aligned}
 \tag{V94.6}
\]
This proves the first identity in \textup{(V94.3)}; the second follows
from \(E_x-(1+Y)I=\Delta_x-YI\).  Multiplication by that operator proves
\textup{(V94.4)}, which also follows by conjugating
\textup{(V92.10)}.

Finally,
\[
 \mathcal J_m(q)=(-1)^m q^m\mathcal C_m^{[0]}(Y).
 \tag{V94.7}
\]
Thus two adjacent nonzero \(\mathcal J\)-entries have opposite signs
exactly when the corresponding \(\mathcal C\)-entries have the same
sign.  Summing the adjacent changes proves \textup{(V94.5)}.
\end{proof}

At \(Y=Y_f\), one has \(1+Y_f=e^{\varepsilon L_f}\), so
\[
 \boxed{
 \mathcal C_m^{[r]}(Y_f)
 =\Delta_x^r\bigl(E_x-e^{\varepsilon L_f}I\bigr)^m\mu(0).}
 \tag{V94.8}
\]
Moreover V93 supplies not merely positive input data but a positive
two-row seed:
\[
 \mathcal C_0^{[r]}=\delta_r>0,\qquad
 \mathcal C_1^{[r]}
 =\delta_{r+1}-Y_f\delta_r>0
 \quad\text{when}\quad
 \frac rM<R-hL_f-\sigma.
 \tag{V94.9}
\]
The unresolved problem is now exact: determine how much of this positive
seed remains constant-sign after repeated application of
\(\Delta_x-Y_fI\).

\subsection{Hurwitz stability does not give the required total positivity}

The forward-difference polynomials \(\Delta^r\mu\) are Hurwitz and their
values at the nonnegative integers are positive.  A natural shortcut would be to
know that
\[
 \delta_r^2\ge\delta_{r-1}\delta_{r+1},
 \tag{V94.10}
\]
because then the initial ratios \(X_0^{[r]}=\delta_{r+1}/\delta_r\)
would be nonincreasing.  But even this first-order consequence is false
for general Hurwitz polynomials.

\begin{proposition}[Quadratic obstruction to the log-concavity shortcut]
\label{prop:v94-quadratic-obstruction}
There is a monic polynomial with positive coefficients and all zeros in
the open left half-plane whose positive forward differences at the
origin are not log-concave.  Explicitly, let
\[
 f(x)=(x+1)^2+9=x^2+2x+10.
 \tag{V94.11}
\]
Then the zeros are \(-1\pm3i\), while
\[
 \boxed{
 \Delta^0f(0)=10,\qquad
 \Delta f(0)=3,\qquad
 \Delta^2f(0)=2,}
 \tag{V94.12}
\]
and
\[
 3^2<10\cdot2.
 \tag{V94.13}
\]
Consequently Hurwitz stability and coefficient positivity alone imply
neither log-concavity of the Newton data nor a Pólya-frequency property
strong enough to transport \textup{(V94.9)}.
\end{proposition}

\begin{proof}
The root and coefficient assertions are immediate.  Direct subtraction
gives
\[
 f(1)-f(0)=13-10=3,\qquad
 f(2)-2f(1)+f(0)=18-26+10=2.
\]
This proves \textup{(V94.12)}--\textup{(V94.13)}.  Log-concavity is a
necessary \(2\times2\) minor condition for the usual
Pólya-frequency conclusions, so those conclusions cannot follow from
the stated hypotheses.
\end{proof}

\begin{firewall}[No monotone-ratio induction from Hurwitz location]
For the example in Proposition~\ref{prop:v94-quadratic-obstruction},
\[
 \frac{\Delta f(0)}{f(0)}=\frac3{10}
 \quad\text{but}\quad
 \frac{\Delta^2f(0)}{\Delta f(0)}=\frac23.
 \tag{V94.14}
\]
Thus the positive initial ratios can increase with the difference level.
The V92 half-plane preservation and V93 Peano positivity do not justify
the southwest threshold propagation suggested at the end of V93.  Any
such propagation must use the specific gamma-ladder factor geometry.
\end{firewall}

\subsection{Only a fixed conjugate prefix can increase a factor ratio}

The relevant gamma geometry can be tested before taking forward
differences.  A real physical factor with real part \(a>0\) contributes
\[
 R_{a,0}(x)
 :=\frac{a+\varepsilon(x+1)}{a+\varepsilon x}.
 \tag{V94.15}
\]
A conjugate pair at \(a\pm ib\) contributes
\[
 R_{a,b}(x)
 :=\frac{(a+\varepsilon(x+1))^2+b^2}
         {(a+\varepsilon x)^2+b^2}.
 \tag{V94.16}
\]

\begin{theorem}[Finite monotonicity-defect prefix]
\label{thm:v94-finite-defect}
For every \(x\ge0\),
\[
 \boxed{
 \frac{d}{dx}\log R_{a,0}(x)
 =-\frac{\varepsilon^2}
 {(a+\varepsilon x)(a+\varepsilon(x+1))}<0,}
 \tag{V94.17}
\]
and
\[
 \boxed{
 \frac{d}{dx}\log R_{a,b}(x)
 =
 \frac{2\varepsilon^2
 \{b^2-(a+\varepsilon x)(a+\varepsilon(x+1))\}}
 {\bigl((a+\varepsilon x)^2+b^2\bigr)
  \bigl((a+\varepsilon(x+1))^2+b^2\bigr)}.}
 \tag{V94.18}
\]
Hence a conjugate factor can increase only when
\[
 \boxed{
 a<
 T_b(x):=\sqrt{b^2+\frac{\varepsilon^2}{4}}
           -\varepsilon\left(x+\frac12\right).}
 \tag{V94.19}
\]
For the gamma ladder
\[
 a_{\ell,k}=d+1+\alpha_\ell+hk,\qquad 0\le k<M,
 \tag{V94.20}
\]
the exact number of factors which can increase at \(x\) is
\[
 \boxed{
 D_\ell(x)=
 \min\!\left\{M,
 \max\!\left\{0,
 \left\lceil
 \frac{T_{b_\ell}(x)-d-1-\alpha_\ell}{h}
 \right\rceil\right\}\right\}.}
 \tag{V94.21}
\]
It is nonincreasing in \(x\), and, for fixed gamma data in the V88
scaling regime, it is \(O(1)\) uniformly in \(M,d\), and \(x\ge0\).
\end{theorem}

\begin{proof}
Logarithmic differentiation of \textup{(V94.15)} gives
\textup{(V94.17)}.  Put \(A=a+\varepsilon x\).  For the conjugate pair,
\[
 \begin{aligned}
 \frac{d}{dx}\log R_{a,b}(x)
 &=
 \frac{2\varepsilon(A+\varepsilon)}
      {(A+\varepsilon)^2+b^2}
 -\frac{2\varepsilon A}{A^2+b^2}\\
 &=
 \frac{2\varepsilon^2\{b^2-A(A+\varepsilon)\}}
      {(A^2+b^2)((A+\varepsilon)^2+b^2)},
 \end{aligned}
\]
which proves \textup{(V94.18)}.  Since
\[
 A(A+\varepsilon)
 =\left(A+\frac{\varepsilon}{2}\right)^2
   -\frac{\varepsilon^2}{4},
\]
positivity is equivalent to \textup{(V94.19)}.  Counting the integers
\(0\le k<M\) satisfying the strict inequality
\(a_{\ell,k}<T_{b_\ell}(x)\) proves \textup{(V94.21)}.
The threshold decreases with \(x\).  Also
\(T_b(x)\le\sqrt{b^2+\varepsilon^2/4}\), so for fixed
\(b_\ell,\alpha_\ell,h\), the count is bounded independently of
\(M,d,x\).
\end{proof}

For the complete physical product, the derivative of
\(\log\{\mu(x+1)/\mu(x)\}\) is the sum of
\textup{(V94.17)}--\textup{(V94.18)} over its factors.  The theorem
therefore localizes every positive summand to
\[
 \sum_\ell D_\ell(x)=O(1)
 \tag{V94.22}
\]
conjugate pairs.  This is much smaller than the
\(O(\sqrt m)\) prefix in \textup{(V92.28)} because the two tests ask
different questions: V92 asks whether a factor is a finite multiplier
through order \(m\), whereas V94 asks whether its adjacent-value ratio
can increase at one real argument.

\begin{firewall}[A finite positive part need not have negative total sum]
Equation \textup{(V94.22)} does not by itself prove that
\(\mu(x+1)/\mu(x)\) decreases.  The finitely many positive terms in
\textup{(V94.18)} can dominate many small negative terms, especially
near a low-lying conjugate factor.  What has been proved is an ownership
statement: any failure of monotonicity is caused by a fixed physical
prefix.  A quantitative comparison of that prefix with the negative
gamma-ladder bulk is still required.
\end{firewall}

\subsection{The exact one-pair shifted-difference threshold}

The operator in \textup{(V94.3)} can also be computed on one factor.
This does not factorize the full transformed product, but it identifies
the natural ordinate scale seen by the first transport step.

\begin{proposition}[Single conjugate-pair discriminant]
\label{prop:v94-pair-discriminant}
Let
\[
 g(x)=(x+u)^2+v^2,\qquad Y>0.
\]
Then
\[
 \boxed{
 (\Delta_x-YI)g(x)
 =2(x+u)+1-Y\{(x+u)^2+v^2\}.}
 \tag{V94.23}
\]
Its two zeros are real if and only if
\[
 \boxed{Y^2v^2\le1+Y,}
 \tag{V94.24}
\]
and, in that case,
\[
 \boxed{
 x+u=\frac{1\pm\sqrt{1+Y-Y^2v^2}}{Y}.}
 \tag{V94.25}
\]
At the first-prime endpoint, with \(v=b/\varepsilon\),
\[
 Y_f^2v^2\longrightarrow b^2L_f^2.
 \tag{V94.26}
\]
Thus a fixed pair has real zeros after the first shifted-difference step
for all sufficiently small \(\varepsilon\) when \(|b|L_f<1\), and has a
nonreal pair when \(|b|L_f>1\).
\end{proposition}

\begin{proof}
Expanding \(g(x+1)-(1+Y)g(x)\) proves
\textup{(V94.23)}.  In the variable \(z=x+u\), its zero equation is
\[
 Yz^2-2z+(Yv^2-1)=0.
\]
The discriminant is \(4(1+Y-Y^2v^2)\), proving
\textup{(V94.24)}--\textup{(V94.25)}.  Finally,
\(Y_f/\varepsilon\to L_f\), which proves
\textup{(V94.26)} and the strict alternatives.
\end{proof}

\begin{firewall}[The one-factor test is not multiplicative]
The operator \(\Delta_x-YI\) is linear but does not satisfy a product
rule that preserves factorization.  Proposition
\ref{prop:v94-pair-discriminant} therefore gives a sharp scale test for
one conjugate pair, not a zero theorem for the full gamma product.
In particular, multiplying the individual discriminant conditions is
invalid.
\end{firewall}

\begin{theorem}[V94 shifted-transport verdict]
\label{thm:v94-verdict}
For the unresolved equispaced family:
\begin{enumerate}[label=\textup{(\roman*)},leftmargin=3em]
\item the full Riccati sign problem is exactly the linear
      shifted-difference orbit \textup{(V94.3)};
\item order-edge variations count adjacent equal signs on the conjugated
      edge, as in \textup{(V94.5)};
\item V93 supplies a macroscopic positive two-row seed
      \textup{(V94.9)};
\item Hurwitz stability alone cannot propagate that seed, because even a
      stable quadratic violates Newton log-concavity;
\item every possible increase of the original gamma-product unit-step
      ratio is owned by the explicit \(O(1)\) prefix
      \textup{(V94.21)};
\item the first shifted difference of one conjugate pair has the exact
      physical threshold \(|b|L_f=1\).
\end{enumerate}
The missing estimate is now narrower: compare the fixed positive prefix
in \textup{(V94.18)} with the negative gamma-ladder bulk, and determine
whether the resulting adjacent-ratio shape survives enough iterates of
\(\Delta_x-Y_fI\) to force \((R-1)M+1\) equal-sign adjacencies.
Arbitrary pole meshes, target-zero separation, and GRH remain open.
\end{theorem}

\begin{assessment}[V94 closure and V95 direction]
V94 removes a quotient singularity from the transport problem by
returning to its linear shifted-difference operator.  It also closes the
most immediate abstract route: neither Hurwitz stability nor positivity
supplies the needed log-concavity.  The gamma structure nevertheless
gives a stronger localization than the order-adaptive multiplier test:
only finitely many physical conjugate factors can oppose adjacent-ratio
decay.

V95 should estimate
\[
 \frac{d}{dx}\log\frac{\mu(x+1)}{\mu(x)}
\]
by summing the negative real and good-conjugate terms explicitly along
each arithmetic gamma ladder, while retaining the finite bad prefix
without absolute-value inflation.  A useful first target is a one-turning
point theorem for \(\mu(x+1)/\mu(x)\).  If such a theorem fails even for
the original product, the correct upstream object is the gamma-ladder
sum in \textup{(V94.18)}, not another general stability surrogate.
\end{assessment}

\section*{Version V94 append-only record}
\addcontentsline{toc}{section}{Version V94 append-only record}
The complete V93 source was retained before this continuation.  It had
1361735 bytes, 40245 lines, and SHA--256
\[
 \texttt{5C5E159415D2D76C0AB53B5B5E8242FB7EFE10D8A6041A14C0186D3CF0EF7F23}.
\]
V94 conjugated the Riccati array to a shifted-difference orbit, disproved
the general Newton-log-concavity shortcut, and localized every possible
increase of an original gamma-factor ratio to an explicit fixed prefix.

% =============================================================================
% END APPENDED V94 MATERIAL
% =============================================================================
% =============================================================================
% BEGIN APPENDED V95 MATERIAL
% =============================================================================

\section{V95: matched gamma triples and an improved square-root boundary}
\label{sec:grh-v95}

V94 localized every possible increase of an individual conjugate-pair
ratio to a fixed prefix, but it did not compare that prefix with the
negative factors.  In the actual two-channel construction there is a
canonical comparison: the principal character is even, so its real gamma
ladder begins no farther right than either member of the target
conjugate ladder.  Matching equal ladder indices shows that one principal
real factor dominates the entire positive logarithmic derivative of one
target pair.  The full original adjacent ratio is therefore strictly
decreasing, with no exceptional prefix.

More importantly, the same matching improves the finite-multiplier
argument.  Instead of testing the conjugate quadratic alone, combine it
with its matched real factor and test one cubic vector.  Its Jensen
polynomial has an explicit cubic residual and an exact discriminant.
On \(d\asymp\sqrt M\), the discriminant has a universal scaled limit.  For
an even target character this enlarges the V90 certified range by the
factor \(\sqrt3\); for an odd target it gives a new explicit
parity-dependent boundary.

\subsection{A real factor dominates a matched conjugate pair}

For \(A>0\), \(b\in\mathbb R\), and step
\(e=\varepsilon>0\), define
\[
 \begin{aligned}
 D_0(A)
 &:=-\frac{e^2}{A(A+e)},\\
 D_b(A)
 &:=
 \frac{2e^2\{b^2-A(A+e)\}}
 {(A^2+b^2)((A+e)^2+b^2)}.
 \end{aligned}
 \tag{V95.1}
\]
These are exactly the real-factor and conjugate-pair contributions in
\textup{(V94.17)}--\textup{(V94.18)} after putting
\(A=a+\varepsilon x\).

\begin{lemma}[Matched-factor domination]
\label{lem:v95-matched-domination}
For every \(A>0\) and \(b\in\mathbb R\),
\[
 \boxed{
 D_0(A)+D_b(A)
 =
 -\frac{e^2\{3p^2+e^2b^2+b^4\}}
 {p(A^2+b^2)((A+e)^2+b^2)}<0,
 \qquad p:=A(A+e).}
 \tag{V95.2}
\]
If \(0<A_0\le A\), then
\[
 \boxed{D_0(A_0)+D_b(A)<0.}
 \tag{V95.3}
\]
\end{lemma}

\begin{proof}
Since
\[
 (A^2+b^2)((A+e)^2+b^2)
 =p^2+b^2(2p+e^2)+b^4,
 \tag{V95.4}
\]
putting the two terms in \textup{(V95.1)} over a common denominator
gives numerator
\[
 e^2\left[
 2p(b^2-p)
 -\{p^2+b^2(2p+e^2)+b^4\}\right],
\]
which is the numerator in \textup{(V95.2)}.  Also
\(A\mapsto\{A(A+e)\}^{-1}\) is strictly decreasing, so
\(D_0(A_0)\le D_0(A)\) when \(A_0\le A\).  Combine this with
\textup{(V95.2)}.
\end{proof}

For the V80 application,
\[
 \kappa_0=0,\qquad \kappa_\chi\in\{0,1\},
 \tag{V95.5}
\]
and the matched physical real parts are
\[
 a_{0,k}=d+1+\kappa_0+2k,\qquad
 a_{\chi,k}=d+1+\kappa_\chi+2k.
 \tag{V95.6}
\]
Thus \(a_{0,k}\le a_{\chi,k}\) for every \(k\).

\begin{theorem}[Strict decay of the complete original adjacent ratio]
\label{thm:v95-complete-ratio-decay}
For the minimal two-channel V80 offset set, and also after adjoining any
number of real padding ladders, one has
\[
 \boxed{
 \frac{d}{dx}\log\frac{\mu(x+1)}{\mu(x)}<0
 \qquad(x\ge0).}
 \tag{V95.7}
\]
Consequently the sampled positive sequence \(\{\mu(j)\}_{j\ge0}\) is
strictly log-concave.
\end{theorem}

\begin{proof}
Match the principal real factor at \(a_{0,k}\) with the target
conjugate pair at \(a_{\chi,k}\).  At a fixed \(x\), put
\[
 A_0=a_{0,k}+\varepsilon x,\qquad
 A=a_{\chi,k}+\varepsilon x.
\]
Equations \textup{(V95.3)} and \textup{(V95.6)} show that each matched
triple contributes a strictly negative logarithmic derivative.  The two
repair factors \(d,\xi_d\), every additional real gamma ladder, and every
other positive real factor contribute terms of the form \(D_0<0\).
Summing proves \textup{(V95.7)}.  Evaluating the decreasing adjacent
ratio at consecutive integers is precisely strict discrete
log-concavity.
\end{proof}

\begin{firewall}[Log-concavity before inversion is still insufficient]
Theorem~\ref{thm:v95-complete-ratio-decay} concerns the values
\(\mu(j)\), not their inverse binomial transform
\(\delta_r=\Delta^r\mu(0)\).  Even the matched-triple geometry does not
make the latter log-concave.  For example,
\[
 f(x)=(x+1)\{(x+1)^2+100\}
 \tag{V95.8}
\]
has a strictly decreasing adjacent ratio by
Lemma~\ref{lem:v95-matched-domination}, but
\[
 \bigl(\Delta^0f(0),\Delta f(0),\Delta^2f(0),\Delta^3f(0)\bigr)
 =(101,107,12,6),
 \tag{V95.9}
\]
and \(12^2<107\cdot6\).  Hence \textup{(V95.7)} does not by itself
propagate the positive two-row seed in \textup{(V94.9)}.
\end{firewall}

\subsection{The exact cubic finite-multiplier test}

Let \(n\ge3\), \(u,w>0\), and \(v\in\mathbb R\).  The matched normalized
coefficient vector differs by a positive scalar from
\[
 \tau_j=(j+w)\{(j+u)^2+v^2\},\qquad 0\le j\le n.
 \tag{V95.10}
\]
Put
\[
 \begin{aligned}
 A_0&:=w(u^2+v^2),\\
 A_1&:=u^2+v^2+2uw+2u+w+1,\\
 A_2&:=2u+w+3.
 \end{aligned}
 \tag{V95.11}
\]

\begin{theorem}[Matched cubic Jensen polynomial]
\label{thm:v95-cubic-jensen}
The Jensen test for \(\{\tau_j\}_{j=0}^n\) factors as
\[
 \boxed{
 \sum_{j=0}^n{n\choose j}\tau_jq^j
 =(1+q)^{n-3}Q_{n,u,w,v}(q),}
 \tag{V95.12}
\]
where
\[
 \boxed{
 \begin{aligned}
 Q_{n,u,w,v}(q)
 ={}&
 A_0(1+q)^3
 +nA_1q(1+q)^2\\
 &+n(n-1)A_2q^2(1+q)
 +n(n-1)(n-2)q^3.
 \end{aligned}}
 \tag{V95.13}
\]
All four coefficients of \(Q_{n,u,w,v}\) are positive.  Therefore
\(\{\tau_j\}_{j=0}^n\) is a finite multiplier sequence through degree
\(n\) if and only if
\[
 \boxed{\operatorname{Disc}_q Q_{n,u,w,v}\ge0.}
 \tag{V95.14}
\]
\end{theorem}

\begin{proof}
Expanding the cubic in the falling-factorial basis gives
\[
 \tau_j=(j)_3+A_2(j)_2+A_1(j)_1+A_0.
 \tag{V95.15}
\]
Use
\[
 \sum_{j=0}^n{n\choose j}(j)_rq^j
 =(n)_rq^r(1+q)^{n-r}
 \tag{V95.16}
\]
for \(0\le r\le3\) to obtain
\textup{(V95.12)}--\textup{(V95.13)}.
Every displayed coefficient is positive.  A real cubic with positive
coefficients and nonzero constant term has three real zeros if and only
if its discriminant is nonnegative; in that case Descartes' rule excludes
positive zeros, so all three are negative.  Now apply the finite
multiplier criterion in Theorem~\ref{thm:v89-finite-multiplier}.
\end{proof}

\subsection{Critical scaling of the cubic discriminant}

The cancellation in the discriminant is best resolved at the triple
limiting root.  Suppose
\[
 \frac wn\longrightarrow s>0,\qquad
 \frac un\longrightarrow s,\qquad
 \frac{v^2}{n}\longrightarrow t\ge0,\qquad
 \frac{u-w}{\sqrt n}\longrightarrow\varrho\in\mathbb R,
 \tag{V95.17}
\]
and put
\[
 K:=s(s+1).
 \tag{V95.18}
\]

\begin{lemma}[Universal matched-cubic discriminant]
\label{lem:v95-cubic-discriminant}
Under \textup{(V95.17)},
\[
 \boxed{
 \frac{\operatorname{Disc}_q Q_{n,u,w,v}}{n^9}
 \longrightarrow4\Phi(K,t,\varrho),}
 \tag{V95.19}
\]
where
\[
 \boxed{
 \begin{aligned}
 \Phi(K,t,\varrho)
 :={}&(K-t)\varrho^4
 +(9K^2-6Kt-2t^2)\varrho^2\\
 &+(3K-t)^3.
 \end{aligned}}
 \tag{V95.20}
\]
Moreover,
\[
 \boxed{
 \frac{\partial\Phi}{\partial K}
 =\{\varrho^2+3(3K-t)\}^2\ge0.}
 \tag{V95.21}
\]
\end{lemma}

\begin{proof}
Set \(s_n=w/n\), \(S_n=1+s_n\), and use the local coordinate
\[
 q=-\frac{s_n}{S_n}
   +\frac{z}{S_n^2\sqrt n}.
 \tag{V95.22}
\]
Substitution into \textup{(V95.13)}, followed by division by
\(n^{3/2}/S_n^3\), gives coefficientwise convergence to
\[
 \boxed{
 \mathcal Q_{K,t,\varrho}(z)
 =z^3+2\varrho z^2+
   (t+\varrho^2-3K)z-2K\varrho.}
 \tag{V95.23}
\]
This is a direct expansion: the \(n^3\) terms form the cube associated
with the triple root \(-s/(1+s)\); the \(n^{5/2}\) and \(n^2\) terms
cancel after centering, and the remaining \(n^{3/2}\) coefficient is
\textup{(V95.23)}.

For a monic cubic \(z^3+az^2+bz+c\), the discriminant is
\[
 a^2b^2-4b^3-4a^3c-27c^2+18abc.
 \tag{V95.24}
\]
Insert
\[
 a=2\varrho,\quad
 b=t+\varrho^2-3K,\quad
 c=-2K\varrho.
\]
Expansion gives \(4\Phi(K,t,\varrho)\).  The affine change
\textup{(V95.22)} and the leading-coefficient scaling cancel exactly in
the discriminant normalization, yielding \textup{(V95.19)}.
Finally, direct differentiation of \textup{(V95.20)} gives
\[
 \varrho^4+6(3K-t)\varrho^2+9(3K-t)^2,
\]
which is the square in \textup{(V95.21)}.
\end{proof}

\subsection{The improved critical-line obstruction}

Return to the actual principal/target triple with \(h=2\), and suppose
\[
 \frac{d_M}{\sqrt M}\longrightarrow\kappa>0.
 \tag{V95.25}
\]
At the first matched ladder index, \textup{(V95.17)} has
\[
 \boxed{
 s_0=\frac1\theta,\qquad
 K_0=\frac{1+\theta}{\theta^2},\qquad
 t_\gamma=\frac{\gamma^2R}{\theta^2\kappa^2},\qquad
 \varrho_\chi=
 \frac{\kappa_\chi\sqrt R}{\theta\kappa}.}
 \tag{V95.26}
\]

\begin{theorem}[Matched critical-line closure]
\label{thm:v95-critical-closure}
For the minimal V80 two-channel offset set, with any additional real
padding ladders allowed, assume
\[
 \boxed{
 \Phi(K_0,t_\gamma,\varrho_\chi)>0.}
 \tag{V95.27}
\]
Then for all sufficiently large \(M\), every matched
principal/target triple is a finite multiplier sequence through degree
\(n\).  Consequently \(H_{M,d_M}\) has exactly \(n\) positive real
zeros, and its last time-domain zero satisfies
\[
 \boxed{
 U_{M,d_M}\ge\frac{R}{2h}\log M+O(1)
 \longrightarrow\infty.}
 \tag{V95.28}
\]
In particular, the V88 Euler-positivity requirement fails for all
sufficiently large \(M\).
\end{theorem}

\begin{proof}
Let \(k/\sqrt M\to\xi\ge0\).  The common scaled real part in
\textup{(V95.17)} becomes
\[
 s_\xi=\frac{\kappa+h\xi}{\theta\kappa},
 \qquad K_\xi=s_\xi(s_\xi+1)\ge K_0,
 \tag{V95.29}
\]
whereas \(t=t_\gamma\) and
\(\varrho=\varrho_\chi\) are independent of \(\xi\).
By \textup{(V95.21)} and the strict margin in
\textup{(V95.27)},
\[
 \Phi(K_\xi,t_\gamma,\varrho_\chi)
 \ge\Phi(K_0,t_\gamma,\varrho_\chi)>0.
 \tag{V95.30}
\]
The expansion in Lemma~\ref{lem:v95-cubic-discriminant} is uniform for
\(\xi\) in compact intervals.  Hence all matched triples with
\(k\le C\sqrt M\) pass the cubic criterion for large \(M\), for each
fixed \(C\).

Choose \(C\) larger than the V90 bad-prefix endpoint divided by
\(\sqrt M\).  For \(k>C\sqrt M\), the conjugate quadratic itself passes
the exact V90 test, while the principal real factor is a linear
multiplier.  Thus every ladder index is certified.  The two repair
factors and all real padding factors are also linear multipliers.
Composition, followed by \textup{(V89.19)}, proves that all zeros of
\(H_{M,d_M}\) are positive real.

Now apply the exact root-product inequality \textup{(V89.22)} and the
sublinear asymptotic \textup{(V90.16)} with
\(d_M\asymp\sqrt M\).  This is \textup{(V95.28)}, and a zero after
\(L_f\) contradicts Euler positivity.
\end{proof}

For an even target character, \(\varrho_\chi=0\), so
\textup{(V95.27)} reduces to
\[
 t_\gamma<3K_0.
\]
Equivalently,
\[
 \boxed{
 \kappa>
 |\gamma|\sqrt{\frac{R}{3(1+\theta)}}.}
 \tag{V95.31}
\]
The V90 separate-quadratic threshold was
\(
 |\gamma|\sqrt{R/(1+\theta)}
\);
matching therefore lowers the certified boundary by exactly
\(\sqrt3\).

For an odd target character, put
\[
 z_\kappa:=\frac{R}{\theta^2\kappa^2}.
 \tag{V95.32}
\]
Then \(\varrho_\chi^2=z_\kappa\),
\(t_\gamma=\gamma^2z_\kappa\), and the condition is the explicit cubic
inequality
\[
 \boxed{
 \begin{aligned}
 0<{}&
 27K_0^3
 +9K_0^2(1-3\gamma^2)z_\kappa\\
 &+K_0(1-3\gamma^2)^2z_\kappa^2
 -\gamma^2(1+\gamma^2)^2z_\kappa^3.
 \end{aligned}}
 \tag{V95.33}
\]
Whenever the separate V90 quadratic test holds, one has
\(t_\gamma<K_0\), and every term in \textup{(V95.20)} is positive.
Thus V95 contains the old range and, by continuity, certifies a strictly
larger open range below its boundary.

\begin{firewall}[A failed cubic test is still only a failed factorization]
At equality in \textup{(V95.27)}, lower-order terms decide the finite
discriminant.  When \(\Phi<0\), one matched cubic has a nonreal pair in
its Jensen test, but composition with the remaining factors could still
restore real-rootedness or could provide the local variation bound without
global real-rootedness.  Therefore the complement of
\textup{(V95.27)} is unresolved, not excluded.  The theorem closes an
expanded part of the critical line; it does not prove GRH.
\end{firewall}

\begin{theorem}[V95 matched-triple verdict]
\label{thm:v95-verdict}
For the actual two-channel equispaced family:
\begin{enumerate}[label=\textup{(\roman*)},leftmargin=3em]
\item the entire original unit-step ratio
      \(\mu(x+1)/\mu(x)\) is strictly decreasing;
\item this follows from exact indexwise domination of every target
      conjugate pair by its principal real factor;
\item original-value log-concavity does not imply Newton-data
      log-concavity, as the matched cubic \textup{(V95.8)} shows;
\item grouping the same factors gives the exact cubic Jensen test
      \textup{(V95.13)};
\item its critical discriminant is the universal polynomial
      \(\Phi\) in \textup{(V95.20)};
\item every square-root line satisfying \textup{(V95.27)} is closed by a
      late positive zero of the residue polynomial;
\item for even target characters, the new boundary improves the V90
      threshold by the factor \(\sqrt3\).
\end{enumerate}
The remaining equispaced range consists of the lower critical region
\(\Phi\le0\), its boundary layers, and \(d=o(\sqrt M)\).  Arbitrary pole
meshes, the upstream target-zero separation problem, and GRH remain open.
\end{theorem}

\begin{assessment}[V95 closure and V96 direction]
V95 completes the one-turning-point task proposed in V94 with the stronger
answer of strict decay.  More importantly, it changes the sufficient
factorization in the way suggested by the V90 audit: the formerly bad
conjugate quadratic is grouped with a canonical real gamma factor.  The
resulting cubic closes a provably larger portion of the square-root
boundary.

V96 should replace \(n\) by the order variable \(m\) in
\textup{(V95.13)} and derive the order-adaptive matched-cubic frontier.
That frontier can then be compared directly with the V92 variation target.
If a macroscopic order interval remains uncertified, the next useful
quantity is the local zero count of the cubic-composed positive tail, not
another unpaired quadratic bound.
\end{assessment}

\section*{Version V95 append-only record}
\addcontentsline{toc}{section}{Version V95 append-only record}
The complete V94 source was retained before this continuation.  It had
1375776 bytes, 40675 lines, and SHA--256
\[
 \texttt{507EDD7B7A340A9CEC64734046879685A25D0A1D6D79D84F8BB0C67FD8205654}.
\]
V95 proved strict adjacent-ratio decay by matching the principal gamma
ladder with the target conjugate ladder, derived the exact matched-cubic
finite-multiplier test and its universal critical discriminant, and
improved the closed even-character square-root range by a factor
\(\sqrt3\).

% =============================================================================
% END APPENDED V95 MATERIAL
% =============================================================================
% =============================================================================
% BEGIN APPENDED V96 MATERIAL
% =============================================================================

\section{V96: the order-adaptive matched frontier and the tilted sequence}
\label{sec:grh-v96}

V95 grouped the principal real gamma factor with the target conjugate
pair and thereby enlarged the top-degree finite-multiplier region.
The V92 obstruction, however, is an order-edge statement: it needs signs
for \(\mathcal J_m(q_f)\) across a macroscopic interval of \(m\).
This version replaces the top degree \(n\) in the matched-cubic
calculation by \(m\), derives the complete critical-scale frontier, and
counts its remaining bad prefix.

For an even target character the frontier becomes elementary.  It gives
an exact order fraction below which every factor is certified, even on
some critical lines where the top-degree test fails.  A second exact
conjugation identifies the still-missing signs with forward differences
of the exponentially tilted sequence
\(e^{-\varepsilon L_fx}\mu(x)\).  V95 proves this sequence is strictly
log-concave and unimodal, but not absolutely monotone; thus factor
certification and sign transport are now cleanly separated.

\subsection{Critical parameters at a moving order}

Assume
\[
 \frac{d_M}{\sqrt M}\longrightarrow\kappa>0,\qquad
 \frac{m_M}{M}\longrightarrow\tau\in(0,R].
 \tag{V96.1}
\]
For the matched triple at ladder index \(k\), use the V95 parameters
\[
 u_k=\frac{d+1+\kappa_\chi+hk}{\varepsilon},\qquad
 w_k=\frac{d+1+\kappa_0+hk}{\varepsilon},\qquad
 v=\frac{\gamma}{\varepsilon},
 \tag{V96.2}
\]
where \(\kappa_0=0\), \(\kappa_\chi\in\{0,1\}\), and \(h=2\) in the
two-channel application.

If \(k/\sqrt M\to\xi\ge0\), define
\[
 \begin{aligned}
 s_{\tau,\xi}
 &:=
 \frac{R(\kappa+h\xi)}{\theta\kappa\tau},&
 K_{\tau,\xi}
 &:=s_{\tau,\xi}(s_{\tau,\xi}+1),\\
 t_\tau
 &:=
 \frac{\gamma^2R^2}{\theta^2\kappa^2\tau},&
 \varrho_{\chi,\tau}
 &:=
 \frac{\kappa_\chi R}{\theta\kappa\sqrt\tau}.
 \end{aligned}
 \tag{V96.3}
\]

\begin{theorem}[Order-adaptive matched discriminant]
\label{thm:v96-order-discriminant}
Under \textup{(V96.1)}, if \(k/\sqrt M\to\xi\), then
\[
 \boxed{
 \frac{
 \operatorname{Disc}_q
 Q_{m,u_k,w_k,v}}{m^9}
 \longrightarrow
 4\Phi(K_{\tau,\xi},t_\tau,\varrho_{\chi,\tau}),}
 \tag{V96.4}
\]
where \(\Phi\) is the universal polynomial in
\textup{(V95.20)}.  The right side is nondecreasing in \(\xi\), since
\[
 \boxed{
 \frac{\partial\Phi}{\partial K}
 =\{\varrho_{\chi,\tau}^2+
 3(3K-t_\tau)\}^2\ge0,
 \qquad
 \frac{\partial K_{\tau,\xi}}{\partial\xi}>0.}
 \tag{V96.5}
\]
If it is positive at \(\xi=0\), every matched triple is a finite
multiplier sequence through degree \(m\), for all sufficiently large
\(M\).
\end{theorem}

\begin{proof}
From \(N/M\to R\), \(n/N\to1\), and
\(\varepsilon=\theta d/N\), one obtains
\[
 \frac{u_k}{m},\frac{w_k}{m}\longrightarrow s_{\tau,\xi},
 \qquad
 \frac{v^2}{m}\longrightarrow t_\tau,
 \qquad
 \frac{u_k-w_k}{\sqrt m}
 \longrightarrow\varrho_{\chi,\tau}.
 \tag{V96.6}
\]
Apply Lemma~\ref{lem:v95-cubic-discriminant} with \(n\) there replaced
by \(m\).  This proves \textup{(V96.4)}.  Equation
\textup{(V96.5)} is \textup{(V95.21)} together with the strict increase
of \(s(s+1)\) for \(s>0\).

If the value at zero is positive, compact-uniformity of the coefficient
expansion certifies \(k\le C\sqrt M\) for every fixed \(C\).  Choose
\(C\) beyond the V92 quadratic bad-prefix endpoint.  Above it the
conjugate quadratic and the principal linear factor are separately
finite multiplier sequences through degree \(m\).  This covers every
\(k\).
\end{proof}

\subsection{The unique matched bad-prefix boundary}

For \(t>0\) and fixed \(\varrho\), regard
\(\Phi(K,t,\varrho)\) as a function of \(K\ge0\).  At the origin,
\[
 \Phi(0,t,\varrho)
 =-t(t+\varrho^2)^2<0,
 \tag{V96.7}
\]
whereas \(\Phi(K,t,\varrho)\to+\infty\) as \(K\to\infty\).

\begin{lemma}[Unique cubic threshold]
\label{lem:v96-unique-threshold}
For \(t>0\) and \(\varrho\in\mathbb R\), there is a unique
\[
 K_\star(t,\varrho)>0
 \tag{V96.8}
\]
such that
\[
 \boxed{
 \Phi(K,t,\varrho)<0\ (K<K_\star),\qquad
 \Phi(K,t,\varrho)=0\ (K=K_\star),\qquad
 \Phi(K,t,\varrho)>0\ (K>K_\star).}
 \tag{V96.9}
\]
If
\[
 s_\star(t,\varrho)
 :=\frac{\sqrt{1+4K_\star(t,\varrho)}-1}{2},
 \tag{V96.10}
\]
then \(s(s+1)>K_\star\) exactly when \(s>s_\star\).
\end{lemma}

\begin{proof}
Equation \textup{(V96.5)} shows that \(\Phi\) is nondecreasing in
\(K\).  Its derivative is the square of an affine function and hence
cannot vanish on a nonempty interval.  Therefore \(\Phi\) is strictly
increasing.  Combine this with \textup{(V96.7)} and its positive cubic
growth at infinity.  Solving \(s(s+1)=K_\star\) for \(s>0\) proves the
last assertion.
\end{proof}

Let \(B_{m,\chi}^{(3)}\) count the matched triple indices whose cubic
Jensen discriminant is negative at order \(m\).  Boundary-zero indices
may be placed on either side; they affect the count by only \(o(\sqrt M)\)
under the strict scaling below.

\begin{theorem}[Matched bad-prefix law]
\label{thm:v96-bad-prefix}
Under \textup{(V96.1)},
\[
 \boxed{
 \frac{B_{m,\chi}^{(3)}}{\sqrt M}
 \longrightarrow
 \frac{\theta\kappa\tau}{Rh}
 \left[
 s_\star(t_\tau,\varrho_{\chi,\tau})
 -\frac{R}{\theta\tau}
 \right]_+.}
 \tag{V96.11}
\]
Thus the uncertified matched factors still form an initial
\(O(\sqrt M)\) prefix, but its endpoint is governed by the cubic
threshold rather than the V92 unpaired quadratic threshold.
\end{theorem}

\begin{proof}
By \textup{(V96.3)},
\[
 s_{\tau,\xi}
 =\frac{R}{\theta\tau}
 +\frac{Rh}{\theta\kappa\tau}\xi.
 \tag{V96.12}
\]
Lemmas~\ref{lem:v95-cubic-discriminant} and
\ref{lem:v96-unique-threshold} show, uniformly away from the unique
boundary, that a triple fails precisely when
\(s_{\tau,\xi}<s_\star(t_\tau,\varrho_{\chi,\tau})\).
Solving for \(\xi\) gives the right side of \textup{(V96.11)}.
The number of indices in a boundary window \(o(\sqrt M)\) is
\(o(\sqrt M)\), which proves the limit.
\end{proof}

\subsection{Closed formulas for an even target character}

When \(\kappa_\chi=0\), one has
\(\varrho_{\chi,\tau}=0\) and
\[
 \Phi(K,t,0)=(3K-t)^3.
 \tag{V96.13}
\]
Hence
\[
 K_\star(t,0)=\frac t3,\qquad
 s_\star(t,0)=\frac{\sqrt{1+4t/3}-1}{2}.
 \tag{V96.14}
\]

\begin{corollary}[Even-character certified order interval]
\label{cor:v96-even-orders}
For an even target character, every matched triple is certified through
order \(m\), asymptotically with a strict margin, exactly when
\[
 \boxed{
 \kappa^2>
 \frac{\gamma^2R\tau}{3(R+\theta\tau)}.}
 \tag{V96.15}
\]
Equivalently, if
\[
 \tau_3(\kappa):=
 \begin{cases}
 R,
 &\gamma^2R\le3\kappa^2\theta,\\[2mm]
 \displaystyle
 \min\!\left\{R,
 \frac{3\kappa^2R}{\gamma^2R-3\kappa^2\theta}
 \right\},
 &\gamma^2R>3\kappa^2\theta,
 \end{cases}
 \tag{V96.16}
\]
then every compact order interval
\[
 0<\frac mM<\tau_3(\kappa)
 \tag{V96.17}
\]
is fully certified for all sufficiently large \(M\).
The bad-prefix constant becomes
\[
 \boxed{
 \frac{B_{m,\chi}^{(3)}}{\sqrt M}
 \longrightarrow
 \frac{\theta\kappa\tau}{Rh}
 \left[
 \frac{\sqrt{1+4t_\tau/3}-1}{2}
 -\frac{R}{\theta\tau}
 \right]_+.}
 \tag{V96.18}
\]
\end{corollary}

\begin{proof}
At \(\xi=0\),
\[
 K_{\tau,0}
 =\frac{R}{\theta\tau}
  \left(\frac{R}{\theta\tau}+1\right).
\]
The inequality \(3K_{\tau,0}>t_\tau\) simplifies to
\textup{(V96.15)}.  Solving it for \(\tau\) gives
\textup{(V96.16)}--\textup{(V96.17)}.  Finally insert
\textup{(V96.14)} into \textup{(V96.11)}.
\end{proof}

The Budan--Fourier target in V92 has length \((R-1)M+1\).  The
order-adaptive cubic has enough \emph{factor-certification capacity} to
cover that many order levels whenever
\[
 \boxed{
 \kappa^2>
 \frac{\gamma^2R(R-1)}
 {3\{R+\theta(R-1)\}}.}
 \tag{V96.19}
\]
Indeed, this is \textup{(V96.15)} at \(\tau=R-1\).
It is weaker than the top-degree condition
\textup{(V95.31)}.

\begin{firewall}[Certification capacity is not the variation count]
Condition \textup{(V96.19)} says that, for at least the required number
of lower order levels, the complete diagonal factor vector preserves
real-rootedness.  It does not determine
\(\operatorname{sgn}\mathcal J_m(q_f)\).  A positive-rooted polynomial
can have either sign at \(q_f\), depending on how many of its roots lie
below \(q_f\).  Therefore the numerical coincidence between
\(\tau_3(\kappa)M\) and the V92 target is not yet a
Budan--Fourier certificate.
\end{firewall}

\subsection{Exponential tilting removes the shifted operator}

The order signs admit a second exact normalization.  Put
\[
 \lambda_f:=\varepsilon L_f,\qquad
 g_f(x):=e^{-\lambda_fx}\mu(x).
 \tag{V96.20}
\]
Let \(M_{\lambda_f}\) denote multiplication by \(e^{\lambda_fx}\).
Then
\[
 M_{\lambda_f}^{-1}
 (E_x-e^{\lambda_f}I)M_{\lambda_f}
 =e^{\lambda_f}\Delta_x.
 \tag{V96.21}
\]

\begin{theorem}[Tilted forward-difference identity]
\label{thm:v96-tilted-difference}
For every \(m\ge0\),
\[
 \boxed{
 \mathcal C_m^{[0]}(Y_f)
 =e^{m\lambda_f}\Delta_x^m g_f(0),}
 \tag{V96.22}
\]
and therefore
\[
 \boxed{
 \mathcal J_m(q_f)
 =(-1)^m\Delta_x^m g_f(0),\qquad
 q_f=e^{-\lambda_f}.}
 \tag{V96.23}
\]
The V92 sufficient obstruction is consequently equivalent, away from
zero entries, to
\[
 \boxed{
 \#\{0\le m<n:
 \Delta^m g_f(0)\Delta^{m+1}g_f(0)>0\}
 \ge(R-1)M+1.}
 \tag{V96.24}
\]
\end{theorem}

\begin{proof}
Conjugating and iterating \textup{(V96.21)} gives
\[
 (E_x-e^{\lambda_f}I)^m\mu(x)
 =e^{\lambda_f(x+m)}\Delta_x^mg_f(x).
 \tag{V96.25}
\]
Set \(x=0\) and use \textup{(V94.8)} to prove
\textup{(V96.22)}.  Since
\(\mathcal J_m=(-1)^mq_f^m\mathcal C_m^{[0]}\) and
\(q_f^me^{m\lambda_f}=1\), equation
\textup{(V96.23)} follows.  Substitute it into
\textup{(V94.5)} to obtain \textup{(V96.24)}.
\end{proof}

The tilt has a rigid one-peak shape.

\begin{corollary}[Strict unimodality of the tilted samples]
\label{cor:v96-tilted-unimodal}
The ratio
\[
 \frac{g_f(x+1)}{g_f(x)}
 =e^{-\lambda_f}\frac{\mu(x+1)}{\mu(x)}
 \tag{V96.26}
\]
is strictly decreasing for \(x\ge0\).  It tends to
\(e^{-\lambda_f}<1\) as \(x\to\infty\).  Hence
\(\{g_f(j)\}_{j\ge0}\) is strictly log-concave and has at most one
change from increasing to decreasing.
\end{corollary}

\begin{proof}
Theorem~\ref{thm:v95-complete-ratio-decay} gives strict decrease, and
the polynomial asymptotic
\(\mu(x+1)/\mu(x)\to1\) gives the limit.  The ratio criterion is
equivalent to strict log-concavity and permits at most one crossing of
one.
\end{proof}

\begin{firewall}[Unimodality does not control all forward differences]
Even a positive strictly log-concave sequence need not have
\(\Delta^mg_f(0)\) of one sign for a macroscopic range of \(m\).
The counterexamples in V94--V95 already rule out the corresponding
Newton-log-concavity inference before tilting.  Thus
\textup{(V96.22)} is an exact reduction, not an absolute-monotonicity
claim.  The missing input is now a root-location or saddle estimate that
turns the gamma-product shape into signs of the high forward differences.
\end{firewall}

\begin{theorem}[V96 order-frontier verdict]
\label{thm:v96-verdict}
For the critical two-channel equispaced family:
\begin{enumerate}[label=\textup{(\roman*)},leftmargin=3em]
\item every order fraction \(\tau=m/M\) has the explicit matched-cubic
      discriminant \textup{(V96.4)};
\item the uncertified matched triples form the exact asymptotic prefix
      \textup{(V96.11)};
\item for even target characters, the fully certified order interval is
      given by \textup{(V96.15)}--\textup{(V96.17)};
\item condition \textup{(V96.19)} gives enough certified order levels
      to match the length of the V92 variation target, but not their
      signs;
\item those signs are exactly the forward-difference signs of one
      strictly log-concave, unimodal tilted sequence, as in
      \textup{(V96.23)}--\textup{(V96.24)}.
\end{enumerate}
The remaining problem is no longer factorwise: within the certified
order block, locate the positive roots of
\(\mathcal J_m\) relative to \(q_f\), equivalently determine the signs
of \(\Delta^mg_f(0)\).  The lower critical region, boundary cases,
arbitrary pole meshes, target-zero separation, and GRH remain open.
\end{theorem}

\begin{assessment}[V96 closure and V97 direction]
V96 converts the V95 top-degree improvement into a complete moving
frontier and shows precisely where factor certification has enough
macroscopic capacity for the Budan target.  The exponential tilt then
removes the shifted operator entirely and identifies the residual
question as high-order forward differences of a one-peak gamma-product
sequence.

V97 should analyze \(\Delta^mg_f(0)\) by its Cauchy or binomial integral
at \(m/M\to\tau\), using the paired gamma-product logarithmic phase rather
than only real-rootedness.  A uniform saddle separated from the nearest
positive root would decide the sign on the certified interval.  If two
competing saddles occur, their exchange curve is the next exact boundary
to compute.
\end{assessment}

\section*{Version V96 append-only record}
\addcontentsline{toc}{section}{Version V96 append-only record}
The complete V95 source was retained before this continuation.  It had
1390730 bytes, 41166 lines, and SHA--256
\[
 \texttt{42EB712360C94D28CBF67EA523C63926E8693B01E1A42DD876CCA25CA26F0713}.
\]
V96 derived the order-adaptive matched-cubic discriminant and bad-prefix
law, obtained a closed even-character order frontier, and converted the
remaining sign ledger to ordinary forward differences of one strictly
log-concave exponentially tilted sequence.

% =============================================================================
% END APPENDED V96 MATERIAL
% =============================================================================
% =============================================================================
% BEGIN APPENDED V97 MATERIAL
% =============================================================================

\section{V97: the Hermite heat envelope and parity-neutral padding}
\label{sec:grh-v97}

V96 reduced the order signs to forward differences of a tilted sequence,
but its matched-cubic calculation also points upstream: a fixed group of
nearby gamma factors has a universal critical Jensen limit.  This version
derives that limit for an arbitrary fixed block.  The limit is not another
ad hoc discriminant; it is the backward heat transform of the polynomial
formed by the scaled physical offsets.

For one real factor and one conjugate pair, the heat polynomial recovers
the V95 cubic and proves that exact real-part alignment is optimal among
all placements of that one real factor.  For \(p\) coalescing real
factors, the heat polynomial is a two-term Hermite combination with an
exact real-rootedness threshold.  Real padding ladders may be chosen
arbitrarily close while remaining distinct, so a strict Hermite margin
survives.  This removes the odd-character parity loss and closes every
critical square-root line above one parity-independent constant.

\subsection{A universal fixed-block Jensen limit}

Fix an integer \(r\ge1\).  For each transform degree \(m\), let
\[
 \pi_m(j):=\prod_{\nu=1}^{r}(j+u_{\nu,m}),
 \qquad
 u_{\nu,m}=sm+\alpha_\nu\sqrt m+o(\sqrt m),
 \tag{V97.1}
\]
where \(s>0\), the multiset
\(\{\alpha_1,\ldots,\alpha_r\}\) is closed under conjugation, and the
errors are uniform over the fixed block.  Write
\[
 \sum_{j=0}^{m}{m\choose j}\pi_m(j)q^j
 =(1+q)^{m-r}Q_m(q).
 \tag{V97.2}
\]
The residual \(Q_m\) has degree \(r\).  Put
\[
 S:=1+s,\qquad K:=s(1+s),
 \qquad
 \Pi_{\boldsymbol\alpha}(z):=
 \prod_{\nu=1}^{r}(z+\alpha_\nu).
 \tag{V97.3}
\]

\begin{theorem}[Backward-heat cluster limit]
\label{thm:v97-heat-limit}
Uniformly for \(z\) in compact subsets of \(\mathbb C\),
\[
 \boxed{
 \frac{S^r}{m^{r/2}}\,
 Q_m\!\left(
 -\frac{s}{S}+\frac{z}{S^2\sqrt m}
 \right)
 \longrightarrow
 \mathcal H_{\boldsymbol\alpha,K}(z)
 :=
 \exp\!\left(-\frac K2D_z^2\right)
 \Pi_{\boldsymbol\alpha}(z).}
 \tag{V97.4}
\]
The exponential on the right terminates after
\(\lfloor r/2\rfloor\) terms.

If \(\mathcal H_{\boldsymbol\alpha,K}\) has \(r\) simple real zeros,
then \(Q_m\) has \(r\) negative real zeros for all sufficiently large
\(m\), and \(\{\pi_m(j)\}_{j=0}^m\) is a finite multiplier sequence
through degree \(m\).  If the heat polynomial has a nonreal simple zero,
then the finite multiplier property fails for all sufficiently large
\(m\).
\end{theorem}

\begin{proof}
The operator identity
\[
 \sum_{j=0}^{m}{m\choose j}\pi_m(j)q^j
 =
 \prod_{\nu=1}^{r}(\mathcal D_q+u_{\nu,m})(1+q)^m,
 \qquad
 \mathcal D_q=q\frac d{dq},
 \tag{V97.5}
\]
shows at once that the left side of \textup{(V97.2)} has the asserted
factor.  Center at the common zero of the leading linear tests,
\[
 q_0=-\frac{s}{1+s},
 \qquad q=q_0+\frac{z}{S^2\sqrt m}.
 \tag{V97.6}
\]
Apply the factors in \textup{(V97.5)} successively, removing one power
of \(1+q\) at each step.  After multiplication by \(S/\sqrt m\), the
first-order part of a factor tends to multiplication by
\(z+\alpha_\nu\).  Whenever two differentiations strike powers already
removed from \((1+q)^m\), their surviving contraction is
\(-K D_z^2/2\).  A contraction of \(2a\) factors has coefficient
\[
 \frac1{a!}\left(-\frac K2\right)^aD_z^{2a}.
 \tag{V97.7}
\]
Summing over the finitely many partial pairings gives
\[
 \sum_{a=0}^{\lfloor r/2\rfloor}
 \frac1{a!}\left(-\frac K2\right)^a
 D_z^{2a}\Pi_{\boldsymbol\alpha}(z),
\]
which is \textup{(V97.4)}.  This argument is a finite induction in
\(r\); the \(o(\sqrt m)\) errors contribute \(o(1)\) after the displayed
normalization.

Coefficientwise convergence of fixed-degree polynomials is equivalent
to locally uniform convergence.  Simple roots therefore persist with
their conjugation type.  A real simple limiting root has one real nearby
root, while a nonreal simple root has a nonreal nearby root.  All local
real roots correspond through \textup{(V97.6)} to negative \(q\) for
large \(m\).  The finite multiplier criterion
\ref{thm:v89-finite-multiplier} proves the final assertions.
\end{proof}

\begin{lemma}[Backward heat preserves real-rootedness]
\label{lem:v97-heat-preserver}
If \(P\) is a real polynomial with only real zeros and \(a\ge0\), then
\[
 \exp(-aD^2)P
 \tag{V97.7a}
\]
also has only real zeros.  Consequently, if
\(\mathcal H_{\boldsymbol\alpha,K_0}\) is real-rooted, then
\(\mathcal H_{\boldsymbol\alpha,K}\) is real-rooted for every
\(K\ge K_0\).
\end{lemma}

\begin{proof}
For \(c\ge0\),
\[
 I-cD^2=(I-\sqrt c\,D)(I+\sqrt c\,D).
\]
Each operator \(I+aD\) preserves real-rootedness: the zeros of \(P'\)
interlace those of \(P\), and the equation \(P'/P=-1/a\) has the required
real solutions between consecutive zeros and in the appropriate exterior
interval.  Therefore
\[
 \left(I-\frac aND^2\right)^NP
\]
is real-rooted for every positive integer \(N\).  Since \(P\) has fixed
degree, these polynomials converge coefficientwise to
\(\exp(-aD^2)P\), and the real-rooted cone is closed.  The semigroup
identity
\[
 \mathcal H_{\boldsymbol\alpha,K}
 =
 \exp\!\left(-\frac{K-K_0}{2}D^2\right)
 \mathcal H_{\boldsymbol\alpha,K_0}
 \tag{V97.7b}
\]
proves the last assertion.
\end{proof}

For \(r=3\), take
\[
 \Pi(z)=(z+p)\{(z+q)^2+t\}.
 \tag{V97.8}
\]
Then
\[
 e^{-(K/2)D^2}\Pi(z)
 =\Pi(z)-\frac K2\Pi''(z).
 \tag{V97.9}
\]
After translating \(z\) so that \(p=0\) and writing
\(\varrho=q-p\), this is exactly
\[
 z^3+2\varrho z^2+
 (t+\varrho^2-3K)z-2K\varrho,
 \tag{V97.10}
\]
the V95 limiting cubic.

\subsection{Alignment is optimal with one real ladder}

\begin{theorem}[Sharp one-real-factor envelope]
\label{thm:v97-one-real-optimal}
Let \(K,t>0\).  Among all real relative placements
\(\varrho\) of one real factor and a conjugate pair, the heat polynomial
\textup{(V97.10)} has three simple real zeros for some
\(\varrho\) if and only if
\[
 \boxed{t<3K.}
 \tag{V97.11}
\]
The threshold is attained by exact alignment \(\varrho=0\).
At \(t=3K,\varrho=0\), the polynomial has a multiple zero; if
\(t>3K\), no real placement restores real-rootedness.
\end{theorem}

\begin{proof}
The discriminant is \(4\Phi(K,t,\varrho)\).  As a polynomial in
\(y=\varrho^2\ge0\),
\[
 \Phi
 =(K-t)y^2+
 (9K^2-6Kt-2t^2)y+
 (3K-t)^3.
 \tag{V97.12}
\]
If \(t>3K\), all three coefficients are strictly negative.  At
\(t=3K\), the last coefficient is zero and the first two are negative,
so only \(y=0\) can give a nonnegative discriminant.  If \(t<3K\),
setting \(y=0\) gives \((3K-t)^3>0\).  A real cubic has three simple
real roots exactly when its discriminant is positive, which proves the
claim.
\end{proof}

Thus the even-character improvement in V95 is not an accident of one
chosen pairing: no horizontal shift of the one available real ladder can
beat its factor \(3\).

\subsection{Coalesced real padding and a Hermite threshold}

Define the monic variance-\(K\) Hermite polynomials by
\[
 H_j^{(K)}(z):=
 \exp\!\left(-\frac K2D_z^2\right)z^j.
 \tag{V97.13}
\]
They obey
\[
 H_{j+1}^{(K)}(z)
 =zH_j^{(K)}(z)-jK H_{j-1}^{(K)}(z),
 \qquad
 (H_j^{(K)})'=jH_{j-1}^{(K)}.
 \tag{V97.14}
\]
For \(p\) coalesced real factors and one centered conjugate pair,
\[
 \Pi_{p,t}(z)=z^p(z^2+t),
 \tag{V97.15}
\]
so Theorem~\ref{thm:v97-heat-limit} gives
\[
 \boxed{
 \mathcal H_{p,t,K}(z)
 =H_{p+2}^{(K)}(z)+tH_p^{(K)}(z).}
 \tag{V97.16}
\]
Put
\[
 \nu_p:=
 \begin{cases}
 p+1,&p\ \hbox{even},\\
 p+2,&p\ \hbox{odd}.
 \end{cases}
 \tag{V97.17}
\]

\begin{theorem}[Exact coalesced-padding threshold]
\label{thm:v97-hermite-threshold}
For \(p\ge0\), \(K>0\), and \(t\ge0\), the polynomial
\(\mathcal H_{p,t,K}\) has only real zeros if and only if
\[
 \boxed{t\le\nu_pK.}
 \tag{V97.18}
\]
The zeros are simple when the inequality is strict.  At equality, the
two central zeros collide at the origin; for \(t>\nu_pK\), precisely
that pair is nonreal and the other \(p\) zeros remain real.
\end{theorem}

\begin{proof}
The Hermite polynomials form a Sturm chain under
\textup{(V97.14)}.  The positive zeros of
\(H_p^{(K)}\) strictly interlace the positive zeros of
\(H_{p+2}^{(K)}\), apart from the one central pair added by the degree
increase.  At every nonzero zero of \(H_p^{(K)}\), the value of
\(\mathcal H_{p,t,K}\) equals that of \(H_{p+2}^{(K)}\), independently
of \(t\).  The alternating signs therefore force \(p\) real outer zeros
for every \(t\ge0\).  Only the central pair can leave the real axis.

If \(p=2a\), the polynomial is even and
\[
 \begin{aligned}
 \mathcal H_{p,t,K}(0)
 &=
 H_{2a+2}^{(K)}(0)+tH_{2a}^{(K)}(0)\\
 &=(-1)^a(2a-1)!!K^a
   \{t-(2a+1)K\}.
 \end{aligned}
 \tag{V97.19}
\]
The central pair is real on the side
\(t\le(2a+1)K=(p+1)K\), and it collides at zero at equality.

If \(p=2a+1\), the polynomial is odd and already has a zero at zero.
The additional central pair is decided by
\[
 \begin{aligned}
 \mathcal H_{p,t,K}'(0)
 &=(p+2)H_{p+1}^{(K)}(0)
   +tpH_{p-1}^{(K)}(0)\\
 &=
 (-1)^a(2a+1)(2a-1)!!K^a
 \{t-(2a+3)K\}.
 \end{aligned}
 \tag{V97.20}
\]
It is real for
\(t\le(2a+3)K=(p+2)K\) and forms a triple zero at equality.
The forced outer zeros account for all remaining roots, completing the
proof.
\end{proof}

Although distinct real progressions cannot literally coincide with one
another, their offsets may be chosen arbitrarily close while remaining
pairwise disjoint.  Under a strict inequality in
\textup{(V97.18)}, the Hermite zeros are simple.  Root continuity
therefore gives the following usable version.

\begin{corollary}[Distinct near-aligned padding]
\label{cor:v97-near-padding}
If \(t<\nu_pK\), there is a neighborhood of the coalesced real-offset
configuration such that every conjugation-closed block of \(p\) distinct
real offsets and the fixed conjugate pair has a simple real-rooted heat
polynomial.  The corresponding finite Jensen blocks are multiplier
sequences for all sufficiently large transform degrees in the critical
scaling.
\end{corollary}

\begin{proof}
The discriminant of a fixed-degree polynomial is continuous in its
coefficients, and simple real roots persist under sufficiently small real
coefficient perturbations.  Apply Theorem
\ref{thm:v97-heat-limit}.
\end{proof}

\subsection{A parity-neutral critical-line obstruction}

Suppose first that the target character is even.  The principal real
ladder is already aligned with the target pair, so \(p=1\) and the
minimal family has \(R=3=\nu_1\).  At top degree,
\[
 K=\frac{1+\theta}{\theta^2},
 \qquad
 t=\frac{\gamma^2R}{\theta^2\kappa^2}.
 \tag{V97.21}
\]
Condition \(t<\nu_1K\) is exactly
\[
 \kappa^2>\frac{\gamma^2}{1+\theta}.
 \tag{V97.22}
\]

For an odd target, retain the separate principal ladder and adjoin
\(p\) distinct real padding ladders with offsets arbitrarily close to
\(\kappa_\chi=1\).  Group those \(p\) ladders with the target pair.
Then the total ladder count is
\[
 R=p+3.
 \tag{V97.23}
\]
Choose \(p\) odd, so \(\nu_p=p+2=R-1\).  The strict heat condition is
\[
 \kappa^2>
 \frac{\gamma^2(p+3)}
 {(p+2)(1+\theta)}.
 \tag{V97.24}
\]
The fraction on the right decreases to
\(\gamma^2/(1+\theta)\).

\begin{theorem}[Parity-neutral padded square-root closure]
\label{thm:v97-parity-neutral}
Fix a hypothetical target zero ordinate \(\gamma\ne0\) and
\(0<\theta<1\).  Suppose
\[
 \boxed{
 \frac{d_M}{\sqrt M}\longrightarrow\kappa,
 \qquad
 \kappa>\frac{|\gamma|}{\sqrt{1+\theta}}.}
 \tag{V97.25}
\]
For either parity of the target character, there is a fixed finite
conjugation-closed V88 offset set, containing the required principal and
target ladders and only real padding beyond them, such that for all
sufficiently large \(M\):
\begin{enumerate}[label=\textup{(\roman*)},leftmargin=3em]
\item every grouped gamma block is a finite multiplier sequence through
      degree \(n\);
\item \(H_{M,d_M}\) has exactly \(n\) positive real zeros;
\item its last time-domain zero tends to infinity;
\item Euler positivity fails beyond \(L_f\).
\end{enumerate}
\end{theorem}

\begin{proof}
For an even character, use the minimal aligned triple and
\textup{(V97.22)}.  For an odd character, the strict inequality in
\textup{(V97.25)} permits an odd \(p\) large enough that
\textup{(V97.24)} holds.  Choose the \(p\) real offsets sufficiently
close to \(1\) that Corollary~\ref{cor:v97-near-padding} applies; keep
them distinct modulo the step \(h=2\).  The principal ladder and the two
repair factors are separate linear multiplier sequences.

Choose \(C\) beyond the separate quadratic bad-prefix endpoint.
For ladder indices \(k\le C\sqrt M\), the heat-limit argument is uniform
and the common heat parameter only increases with \(k\).
Lemma~\ref{lem:v97-heat-preserver} and compactness of this \(K\)-interval
allow the distinct offsets to be chosen sufficiently close that the
strict certificate persists throughout it.  For \(k>C\sqrt M\), V90
certifies the target quadratic and every real factor separately.
Composition proves
real-rootedness of \(H_{M,d_M}\), and coefficient positivity places every
root on the positive axis.

Finally, \textup{(V90.16)} and the exact root-product inequality
\textup{(V89.22)} give
\[
 U_{M,d_M}\ge\frac{R}{2h}\log M+O(1)\to\infty.
 \tag{V97.26}
\]
Thus a positive root corresponds to a kernel zero after \(L_f\), which
contradicts Euler positivity.
\end{proof}

For the simplest odd-character improvement, take \(p=1\): add one real
ladder exactly aligned in real part with the target pair.  Then \(R=4\)
and
\[
 \boxed{
 \kappa>
 \frac{2|\gamma|}{\sqrt{3(1+\theta)}}}
 \tag{V97.27}
\]
already suffices.  Increasing odd \(p\) approaches the
parity-neutral threshold \textup{(V97.25)} from above.

\begin{firewall}[The aligned-padding threshold is not a universal no-go]
Theorem~\ref{thm:v97-parity-neutral} is a closure result for every line
above \(|\gamma|/\sqrt{1+\theta}\).  It does not prove that padding below
that line is impossible.  Noncoalesced real offsets produce the more
general heat polynomial \textup{(V97.4)} and may have a different
real-rootedness region.  Nor does failure of any fixed block test exclude
local endpoint variation without global real-rootedness.  The region
\(\kappa\le|\gamma|/\sqrt{1+\theta}\), lower scales, and GRH remain open.
\end{firewall}

\begin{theorem}[V97 heat-envelope verdict]
\label{thm:v97-verdict}
For fixed gamma-factor groups in the critical Jensen scaling:
\begin{enumerate}[label=\textup{(\roman*)},leftmargin=3em]
\item the residual Jensen polynomial converges to the backward-heat
      transform \textup{(V97.4)} of its scaled offset polynomial;
\item the V95 cubic is the three-factor instance of this universal limit;
\item exact real-part alignment is optimal when only one real factor is
      available;
\item \(p\) coalesced real factors give the exact Hermite threshold
      \(t\le\nu_pK\);
\item distinct sufficiently near-aligned padding retains every strict
      Hermite certificate;
\item finite real padding removes the parity penalty and closes all
      critical lines satisfying \textup{(V97.25)}.
\end{enumerate}
The remaining sign problem lies below the parity-neutral square-root
threshold or within local variation questions not requiring global
real-rootedness.  Arbitrary pole meshes, target-zero separation, and GRH
remain open.
\end{theorem}

\begin{assessment}[V97 closure and V98 direction]
V97 changes the factorization upstream rather than forcing the V96
forward-difference saddle prematurely.  The heat envelope both explains
the cubic discriminant and supplies a systematic design space for real
padding.  It yields a parity-independent closed region and proves the
one-real matching was already optimally placed.

V98 should optimize the general real shift vector
\(\boldsymbol\alpha\) in \textup{(V97.4)} below
\(|\gamma|/\sqrt{1+\theta}\).  The first rigorous question is whether a
second-moment or discriminant invariant bounds the largest admissible
imaginary pair uniformly over all real padding shifts.  If no such bound
exists, construct an explicit noncoalesced heat polynomial with all real
zeros; if it does, return to the V96 tilted-difference saddle inside the
remaining certified order block.
\end{assessment}

\section*{Version V97 append-only record}
\addcontentsline{toc}{section}{Version V97 append-only record}
The complete V96 source was retained before this continuation.  It had
1404486 bytes, 41610 lines, and SHA--256
\[
 \texttt{A8E82C5D9E76F24C360D822F4BFED893D8071262547C325E1A8E21147B597439}.
\]
V97 derived the universal backward-heat cluster limit for fixed grouped
gamma factors, proved exact one-real and coalesced-Hermite thresholds,
and used distinct near-aligned real padding to close every critical
square-root line above the parity-neutral boundary
\(|\gamma|/\sqrt{1+\theta}\).

% =============================================================================
% END APPENDED V97 MATERIAL
% =============================================================================
% =============================================================================
% BEGIN APPENDED V98 MATERIAL
% =============================================================================

\section{V98: noncoalesced padding and the backward-heat speed limit}
\label{sec:grh-v98}

V97 derived the fixed-block heat envelope and left open whether separated
real offsets can beat the coalesced Hermite threshold.  They can: an
explicit four-factor polynomial with two real padding factors has a
simple real-rooted heat image at \(t=7K/2\), whereas the coalesced
threshold for two real factors is only \(3K\).

There is nevertheless a universal restriction.  Under backward heat,
zeros obey an exact repulsive flow.  If the initial offset polynomial has
one conjugate pair at height \(\sqrt t\) and \(p\) real zeros, the squared
height can decrease at speed at most \(2p+1\).  Hence no choice of the
real offsets can become real-rooted by heat time \(K\) unless
\(t\le(2p+1)K\).  Translated back to the gamma construction, this gives
a parity-independent lower barrier for every finite fixed-block padding
certificate.

\subsection{A separated quartic beats the coalesced threshold}

Take \(K>0\) and put
\[
 \Pi_K(z):=
 \left(z+\frac{\sqrt K}{2}\right)^2
 \left(z^2+\frac{7K}{2}\right).
 \tag{V98.1}
\]
This has two real zeros at \(-\sqrt K/2\) and one conjugate pair at
\(\pm i\sqrt{7K/2}\).

\begin{proposition}[Exact noncoalesced quartic certificate]
\label{prop:v98-quartic}
The backward-heat image
\[
 \mathcal Q_K(z):=
 \exp\!\left(-\frac K2D_z^2\right)\Pi_K(z)
 \tag{V98.2}
\]
has four distinct real zeros.  Consequently, after an arbitrarily small
real perturbation separating the repeated real input zero, the input has
two distinct real zeros and one conjugate pair, while its heat image
remains simple real-rooted.
\end{proposition}

\begin{proof}
By the scaling \(z=\sqrt K\,x\), it is enough to take \(K=1\).  Direct
differentiation gives
\[
 \mathcal Q_1(x)
 =x^4+x^3-\frac94x^2+\frac12x+\frac18.
 \tag{V98.3}
\]
It factors as
\[
 \boxed{
 \mathcal Q_1(x)
 =
 \left(x-\frac12\right)
 \left(x^3+\frac32x^2-\frac32x-\frac14\right).}
 \tag{V98.4}
\]
The cubic factor has discriminant
\[
 \frac{243}{8}>0
 \tag{V98.5}
\]
and hence three distinct real zeros.  Its value at \(x=1/2\) is
\(-1/2\), so the displayed linear zero is different from all three.
This proves simple real-rootedness.  Polynomial roots depend
continuously on the coefficients, and the simple real-rooted cone is
open among real fixed-degree polynomials.  A sufficiently small
perturbation of the two real input shifts therefore preserves the
conclusion.
\end{proof}

For \(p=2\), Theorem~\ref{thm:v97-hermite-threshold} gives only
\(t\le3K\) when both real inputs are centered on the conjugate pair.
Proposition~\ref{prop:v98-quartic} reaches \(t=7K/2\).  Thus the
coalesced Hermite family is not globally optimal once two real shifts are
available.

\subsection{The exact zero flow}

Let
\[
 P_\tau(z):=\exp\!\left(-\frac{\tau}{2}D_z^2\right)P_0(z),
 \qquad \tau\ge0.
 \tag{V98.6}
\]
Suppose initially
\[
 P_0(z)=
 \{(z-x_0)^2+y_0^2\}
 \prod_{j=1}^{p}(z-r_{j,0}),
 \qquad y_0>0,
 \tag{V98.7}
\]
with distinct real \(r_{j,0}\).  Multiple real zeros are obtained by a
limit.

\begin{lemma}[Backward-heat zero dynamics]
\label{lem:v98-zero-flow}
As long as all zeros are simple, every zero \(z_i(\tau)\) of \(P_\tau\)
obeys
\[
 \boxed{
 z_i'(\tau)
 =\sum_{j\ne i}\frac1{z_i(\tau)-z_j(\tau)}.}
 \tag{V98.8}
\]
The real zeros remain real and cannot collide with one another before
the conjugate pair reaches the real axis.  Writing that pair as
\[
 x(\tau)\pm iy(\tau),\qquad y(\tau)>0,
 \]
one has
\[
 \boxed{
 \frac d{d\tau}y(\tau)^2
 =
 -1-
 2\sum_{j=1}^{p}
 \frac{y(\tau)^2}
 {(x(\tau)-r_j(\tau))^2+y(\tau)^2}.}
 \tag{V98.9}
\]
\end{lemma}

\begin{proof}
The heat equation is
\[
 \partial_\tau P_\tau=-\frac12P_\tau''.
 \]
Differentiating \(P_\tau(z_i(\tau))=0\) gives
\[
 z_i'=\frac{P_\tau''(z_i)}{2P_\tau'(z_i)}.
 \]
For a polynomial with simple zeros,
\[
 \frac{P_\tau''(z_i)}{2P_\tau'(z_i)}
 =\sum_{j\ne i}\frac1{z_i-z_j},
 \]
which proves \textup{(V98.8)}.

At a real zero, the contributions of the conjugate pair add to a real
number, and all other contributions are real.  Hence each simple real
zero stays on the real axis.  The singular mutual terms between two
nearby real zeros point in opposite outward directions, so their squared
separation cannot first vanish; equivalently, uniqueness of the
Calogero system prevents a collision before another zero joins them.

For \(z=x+iy\), the contribution of its conjugate is
\[
 \frac1{2iy}=-\frac{i}{2y}.
 \]
The contribution of a real zero \(r_j\) has imaginary part
\[
 -\frac{y}{(x-r_j)^2+y^2}.
 \]
Taking imaginary parts in \textup{(V98.8)} and multiplying by \(2y\)
proves \textup{(V98.9)}.
\end{proof}

\subsection{A universal finite-block speed budget}

Let \(T\) be the first time at which the conjugate pair in
\textup{(V98.7)} reaches the real axis.  Equation
\textup{(V98.9)} gives
\[
 1\le
 -\frac d{d\tau}y(\tau)^2
 \le2p+1.
 \tag{V98.10}
\]

\begin{theorem}[One-pair heat-time obstruction]
\label{thm:v98-speed-limit}
If
\[
 \exp\!\left(-\frac K2D^2\right)P_0
 \tag{V98.11}
\]
is real-rooted, then necessarily
\[
 \boxed{
 y_0^2\le(2p+1)K.}
 \tag{V98.12}
\]
\end{theorem}

\begin{proof}
Until time \(T\), Lemma~\ref{lem:v98-zero-flow} gives exactly one
conjugate pair and \(p\) real zeros.  Integrating
\textup{(V98.10)} yields
\[
 y_0^2
 =
 \int_0^T
 \left(-\frac d{d\tau}y(\tau)^2\right)d\tau
 \le(2p+1)T.
 \tag{V98.13}
\]
If the polynomial is real-rooted at time \(K\), then \(T\le K\):
after the first all-real time, Lemma~\ref{lem:v97-heat-preserver} keeps
the polynomial real-rooted.  This proves \textup{(V98.12)}.
\end{proof}

For \(p=1\), the upper speed is \(3\), and the aligned cubic in V95
attains the threshold \(t=3K\).  Hence the V97 one-real optimality
theorem is also sharp from the dynamical viewpoint.  For \(p=2\), the
quartic construction reaches \(7K/2\), while the universal ceiling is
\(5K\); a genuine optimization interval remains.

\subsection{Translation to critical gamma padding}

Consider a fixed heat block containing the target conjugate pair and
\(p\) real gamma padding factors.  At top degree
\(n/M\to R\), its critical parameters are
\[
 K=\frac{1+\theta}{\theta^2},
 \qquad
 t=\frac{\gamma^2R}{\theta^2\kappa^2}.
 \tag{V98.14}
\]
The total number \(R\) of gamma ladders satisfies
\[
 R\ge p+2,
 \tag{V98.15}
\]
with a further ungrouped principal ladder included when parity or the
chosen grouping requires it.

\begin{theorem}[Universal fixed-block padding barrier]
\label{thm:v98-padding-barrier}
If a fixed block with one target conjugate pair and \(p\) real gamma
factors is to have a real-rooted critical heat envelope, then
\[
 \boxed{
 \kappa^2\ge
 \frac{\gamma^2R}{(2p+1)(1+\theta)}.}
 \tag{V98.16}
\]
Consequently no choice of finitely many real padding offsets, grouped
with the target pair in this fixed-block heat scheme, can certify
top-degree real-rootedness when
\[
 \boxed{
 \kappa\le
 \frac{|\gamma|}{\sqrt{2(1+\theta)}}.}
 \tag{V98.17}
\]
\end{theorem}

\begin{proof}
Insert \textup{(V98.14)} into the necessary heat-time inequality
\(t\le(2p+1)K\).  This is \textup{(V98.16)}.  Since
\[
 \frac{R}{2p+1}\ge
 \frac{p+2}{2p+1}>\frac12
 \tag{V98.18}
\]
for every finite \(p\), condition \textup{(V98.17)} contradicts
\textup{(V98.16)}.
\end{proof}

The V97 construction closes
\[
 \kappa>\frac{|\gamma|}{\sqrt{1+\theta}},
 \tag{V98.19}
\]
whereas the speed limit rules out every fixed one-pair block certificate
at or below the smaller boundary in \textup{(V98.17)}.  Thus the exact
optimization window for separated padding is now
\[
 \boxed{
 \frac{|\gamma|}{\sqrt{2(1+\theta)}}
 <
 \kappa
 \le
 \frac{|\gamma|}{\sqrt{1+\theta}}.}
 \tag{V98.20}
\]

\begin{firewall}[A block-speed obstruction is not a filter obstruction]
Theorem~\ref{thm:v98-padding-barrier} applies to a fixed factor group
containing one conjugate pair and finitely many real factors.  It does
not exclude groups whose size grows with \(M\), joint groups containing
several target-pair ladder indices, recovery of only the local
Budan--Fourier variation count, or a nonuniform pole mesh.  Therefore
\textup{(V98.17)} closes a finite-multiplier proof architecture, not the
underlying Euler-admissible filter and not GRH.
\end{firewall}

\begin{theorem}[V98 heat-flow verdict]
\label{thm:v98-verdict}
For critical fixed-block gamma padding:
\begin{enumerate}[label=\textup{(\roman*)},leftmargin=3em]
\item separated real offsets can strictly improve the coalesced Hermite
      threshold;
\item the explicit quartic \textup{(V98.1)} reaches \(t=7K/2\) with a
      simple real-rooted heat image;
\item the conjugate height obeys the exact zero-flow identity
      \textup{(V98.9)};
\item \(p\) real factors can remove squared height at speed at most
      \(2p+1\);
\item every finite one-pair block is therefore subject to the universal
      critical barrier \textup{(V98.16)};
\item the remaining padding-optimization window is exactly
      \textup{(V98.20)} at the level of currently proved sufficient and
      necessary bounds.
\end{enumerate}
The lower half of the critical region requires growing factor blocks,
local variation without global real-rootedness, or a different pole
geometry.  Target-zero separation and GRH remain open.
\end{theorem}

\begin{assessment}[V98 closure and V99 direction]
V98 answers the first V97 question in both directions.  Coalescence is
not optimal once two real shifts are available, but arbitrary fixed
padding has a rigorous dynamical speed limit.  This prevents an
unbounded search over offsets from being mistaken for a route through
the entire lower critical region.

V99 should optimize the first nontrivial quartic block exactly, or
construct a sequence of \(p\)-factor blocks whose admissible height per
factor approaches the speed ceiling.  Either result will decide whether
the window \textup{(V98.20)} is a real construction zone or merely a
gap between two estimates.  V100 must then perform the scheduled stable
sweep of all parallel notes before the direction changes again.
\end{assessment}

\section*{Version V98 append-only record}
\addcontentsline{toc}{section}{Version V98 append-only record}
The complete V97 source was retained before this continuation.  It had
1421584 bytes, 42113 lines, and SHA--256
\[
 \texttt{9F0428C6FC757E3F60656F85C47602356AFB8D285A8BE586D65F2A723B4F586F}.
\]
V98 exhibited an exact separated-padding quartic beyond the coalesced
Hermite threshold, derived the conjugate-pair backward-heat zero flow,
and proved the universal speed-limit barrier for every finite
one-pair padding block.

% =============================================================================
% END APPENDED V98 MATERIAL
% =============================================================================

% =============================================================================
% BEGIN APPENDED V99 MATERIAL
% =============================================================================

\section{V99: exact optimization of the translated double-padding quartic}
\label{sec:grh-v99}

V98 proved that separated real padding can beat the centered Hermite
threshold and exhibited a quartic at normalized height \(7/2\).  This
continuation optimizes the first rigid separated family exactly: the two
real padding zeros are still coincident with each other, but their common
location is allowed to move away from the real part of the conjugate
pair.  The optimal height is an algebraic constant
\[
 c_\star=3.814511867\ldots,
 \qquad
 4c_\star^3-21c_\star^2+24c_\star-8=0.
 \tag{V99.1}
\]
The extremal heat image has a triple real zero.  This gives an exact
quartic benchmark for the remaining asymmetric optimization, but it
does not by itself enlarge the V97 critical closure region.

\subsection{The translated double-padding family}

After translating the conjugate pair to have real part zero and scaling
the heat time to one, put
\[
 P_{a,t}(z):=(z+a)^2(z^2+t),
 \qquad a\in\mathbb R,\quad t\ge0,
 \tag{V99.2}
\]
and
\[
 H_{a,t}(z):=
 \exp\!\left(-\frac12D_z^2\right)P_{a,t}(z).
 \tag{V99.3}
\]
Since \(P_{a,t}\) is quartic, the exponential terminates after the
second derivative term.  Direct differentiation gives
\[
 \boxed{
 H_{a,t}(z)=z^4+2az^3+(a^2+t-6)z^2
       +2a(t-3)z+a^2t-a^2-t+3.}
 \tag{V99.4}
\]
Thus \(t\) is the squared height of the sole conjugate pair, while the
two real input zeros coincide at \(-a\).

Define
\[
 \mathcal T_{\rm eq}:=
 \sup\{t\ge0:H_{a,t}\text{ is real-rooted for some }a\in\mathbb R\}.
 \tag{V99.5}
\]
Only \(A=a^2\) enters the discriminant.  Dividing that discriminant by
the harmless positive factor \(16\), write
\[
 \frac1{16}\operatorname{Disc}_z H_{a,t}=F(A,t),
 \tag{V99.6}
\]
where the exact expansion is
\begin{align}
 F(A,t)={}&(1-t)A^4+(-4t^2-14t+20)A^3 \\
 &+(-6t^3-12t^2-120t+180)A^2 \\
 &+(-4t^4+22t^3+132t^2-720t+864)A \\
 &-t^5+19t^4-160t^3+720t^2-1728t+1728.
 \tag{V99.7}
\end{align}

\begin{lemma}[Stationary discriminant values]
\label{lem:v99-resultant}
If \(A>0\) and
\[
 F(A,t)=\partial_A F(A,t)=0,
 \tag{V99.8}
\]
then
\[
 t^2(t-1)(4t-3)
 \bigl(4t^3-21t^2+24t-8\bigr)^3=0.
 \tag{V99.9}
\]
The cubic in \textup{(V99.9)} has exactly one real zero, namely
\(c_\star\in(3,4)\).
\end{lemma}

\begin{proof}
Take the Sylvester determinant in \(A\) of the degree-four polynomial
\(F\) and its derivative.  Exact integer row elimination gives
\[
 \operatorname{Res}_A(F,F_A)
 =C\,t^2(t-1)(4t-3)
 \bigl(4t^3-21t^2+24t-8\bigr)^3,
 \tag{V99.10}
\]
where \(C\ne0\) is an integer.  This proves the first assertion.  For
completeness, if
\[
 h(t):=4t^3-21t^2+24t-8,
\]
then
\[
 h'(t)=6(2t^2-7t+4).
\]
At the two critical points \((7\pm\sqrt{17})/4\), direct substitution
gives a negative value.  Since \(h(3)=-17\), \(h(4)=8\), and the leading
coefficient is positive, \(h\) has one real zero and that zero lies in
\((3,4)\).

Equation \textup{(V99.10)} is also readily checked without a
computer-algebra convention: insert the five coefficient polynomials
from \textup{(V99.7)} into the \(7\)-by-\(7\) Sylvester matrix and use
fraction-free elimination.  No numerical root classification enters
the identity.
\end{proof}

\begin{theorem}[Sharp equal-shift quartic height]
\label{thm:v99-equal-quartic}
The translated double-padding family satisfies
\[
 \boxed{\mathcal T_{\rm eq}=c_\star.}
 \tag{V99.11}
\]
Equivalently, for heat time \(K>0\),
\[
 \exp\!\left(-\frac K2D_z^2\right)
 (z+a)^2(z^2+t)
\]
can be real-rooted only if \(t\le c_\star K\), and equality is attained
after scaling the extremizer below.
\end{theorem}

\begin{proof}
V98 already supplies a member of this family at \(t=7/2\), namely
\(a=1/2\), whose heat image has four simple real zeros.  Hence
\(\mathcal T_{\rm eq}\ge7/2\).

We first show that a maximizer exists above \(7/2\).  For \(t>1\), the
leading \(A\)-coefficient in \textup{(V99.7)} is negative, so
\[
 F(A,t)\longrightarrow-\infty
 \qquad(A\longrightarrow+\infty).
 \tag{V99.12}
\]
Every real-rooted quartic has nonnegative discriminant, while
Theorem~\ref{thm:v98-speed-limit} gives the uniform bound \(t\le5\).
On the compact interval \(7/2\le t\le5\), the coefficient \(1-t\) is
bounded above by \(-5/2\), and all lower \(A\)-coefficients in
\textup{(V99.7)} are bounded.  Hence the real-rooted parameter set in
this strip is uniformly bounded in \(A\).  It is closed by continuity
of polynomial zeros, so its largest \(t\)-coordinate is attained.

The boundary \(A=0\) cannot contain that maximizer.  Indeed,
\[
 H_{0,t}(z)=z^4+(t-6)z^2+3-t.
 \tag{V99.13}
\]
If all four zeros are real, the two zeros of
\(u^2+(t-6)u+3-t\) must be nonnegative.  Their product is \(3-t\), so
\(t\le3\).

At a maximizing point with \(A>0\), the quartic cannot have four simple
real zeros, since simple real-rootedness is open and \(t\) could then be
increased.  Hence \(F=0\).  The local boundary of the hyperbolicity
region is a discriminant arc.  If \(F_A\ne0\), the implicit-function
theorem writes this arc as \(A=A(t)\), and the side carrying four real
zeros continues to parameters with slightly larger \(t\), again a
contradiction.  At a higher-order collision both first discriminant
derivatives vanish automatically.  Thus in every case \(F_A=0\) at the
maximizer.

Lemma~\ref{lem:v99-resultant} now lists every possible maximizing
height.  The values \(0,3/4,1\) are below \(7/2\), and the only real zero
of the cubic is \(c_\star\).  Therefore
\[
 \mathcal T_{\rm eq}\le c_\star.
 \tag{V99.14}
\]

It remains to attain this value.  Let \(x\) be the unique positive zero
of
\[
 x^3+3x^2-8=0,
 \tag{V99.15}
\]
and choose \(r>0\), \(a\), and \(t\) by
\[
 r^2=\frac{3x}{(x+1)^2(x+4)},
 \qquad a=xr,
 \qquad
 t=6-r^2(x^2+6x+6).
 \tag{V99.16}
\]
Using \textup{(V99.15)} to simplify gives
\[
 t=\frac{3x^2+12x+16}{x^2+3x+4}.
 \tag{V99.17}
\]
Substitution into \textup{(V99.4)} yields the exact factorization
\[
 \boxed{
 H_{a,t}(z)=(z-r)^3\bigl(z+(2x+3)r\bigr).}
 \tag{V99.18}
\]
Thus the heat image is real-rooted.  Eliminating \(x\) between
\textup{(V99.15)} and \textup{(V99.17)} gives
\[
 (t-3)^2\bigl(4t^3-21t^2+24t-8\bigr)=0.
 \tag{V99.19}
\]
Since the value in \textup{(V99.17)} lies in \((3,4)\), it is
\(c_\star\).  This proves equality in \textup{(V99.11)}.  Finally the
change \(z=\sqrt K\,w\) proves the general-\(K\) statement.
\end{proof}

\subsection{Why the extremizer is a triple collision}

The algebra in the preceding proof has a useful dynamical
interpretation.  Suppose \(H_{a,t}\) has a triple zero at \(r\ne0\), and
write \(a=xr\).  The equations \(H''(r)=H'(r)=H(r)=0\) first give
\[
 t=6-r^2(x^2+6x+6),
 \qquad
 r^2=\frac{3x}{(x+1)^2(x+4)},
 \tag{V99.20}
\]
and the remaining equation reduces exactly to
\[
 (x+1)^3(x^3+3x^2-8)=0.
 \tag{V99.21}
\]
The branch \(x=-1\) is excluded by the denominator in
\textup{(V99.20)}.  Hence the positive solution in
\textup{(V99.15)} is forced.  The stationary discriminant calculation
therefore says geometrically that the top of the equal-shift
hyperbolicity component is the point where the incoming conjugate pair
and one real heat-flow zero meet simultaneously.

At equality the image is not simple.  Nevertheless every ratio
\(t/K<c_\star\) sufficiently close to \(c_\star\) has a simple
real-rooted certificate: start from the real-rooted polynomial
\textup{(V99.18)}, continue the backward heat flow for a slightly
longer time \(K>1\), and then rescale to unit heat time.  Real-rootedness
is preserved by Lemma~\ref{lem:v97-heat-preserver}, while multiple
collisions split for all but isolated continuation times.

\subsection{Translation to the critical gamma block}

For a target conjugate pair and two equal-shift real padding factors,
V98 gives
\[
 K=\frac{1+\theta}{\theta^2},
 \qquad
 t=\frac{\gamma^2R}{\theta^2\kappa^2}.
 \tag{V99.22}
\]
Theorem~\ref{thm:v99-equal-quartic} therefore yields the exact
family-specific condition
\[
 \boxed{
 \kappa^2\ge
 \frac{\gamma^2R}{c_\star(1+\theta)}.}
 \tag{V99.23}
\]
There are at least four gamma ladders in this block, so \(R\ge4\).
Since
\[
 c_\star<4,
 \tag{V99.24}
\]
even the optimized equal-shift quartic requires
\[
 \kappa>
 \frac{|\gamma|}{\sqrt{1+\theta}}
 \tag{V99.25}
\]
at minimal ladder count.  It approaches the V97 boundary more closely
than the V98 value \(7/2\), but cannot cross it.

\begin{firewall}[Equal-shift optimality is not quartic optimality]
Theorem~\ref{thm:v99-equal-quartic} assumes the two real input zeros
coincide with each other.  It does not prove that an asymmetric pair of
real shifts cannot exceed \(c_\star\).  Numerical discrimination is not
a substitute for that missing inequality.  The unrestricted quartic
still has the V98 ceiling \(t\le5K\), and crossing the critical gamma
boundary at \(R=4\) would require an asymmetric construction with
\(t/K>4\).
\end{firewall}

\begin{theorem}[V99 quartic verdict]
\label{thm:v99-verdict}
For the heat block
\[
 (z+a)^2(z^2+t)
\]
with two coincident translated real padding zeros:
\begin{enumerate}[label=\textup{(\roman*)},leftmargin=3em]
\item the sharp normalized height is the algebraic constant
      \(c_\star=3.814511867\ldots\);
\item \(c_\star\) is the unique real zero of
      \(4c^3-21c^2+24c-8\);
\item the extremal heat image has the explicit triple-root
      factorization \textup{(V99.18)};
\item this improves the V98 witness \(7/2\) but remains strictly below
      the four-ladder target \(4\);
\item the unresolved quartic question is now sharply localized to
      unequal real shifts.
\end{enumerate}
Target-zero separation and GRH remain open.
\end{theorem}

\begin{assessment}[V99 closure and V100 direction]
V99 converts the first quartic search from a numerical experiment into
an exact algebraic theorem.  Translating a repeated real padding zero
does improve the centered Hermite threshold, but its sharp constant is
still below the value needed to cross the already proved V97 critical
line.  Any genuinely new quartic region must therefore break equality
of the two real shifts.

V100 must first perform the scheduled stable sweep of all notes in
\texttt{latex\_notes}, including parallel notes created since the V90
sweep.  It should then attack the unrestricted quartic with
\(s=a+b\), \(p=ab\), and the exact constraint \(s^2-4p\ge0\).  The
decisive target is a rigorous proof or counterexample for \(t/K>4\),
not merely a numerical maximizer of the quartic discriminant.
\end{assessment}

\section*{Version V99 append-only record}
\addcontentsline{toc}{section}{Version V99 append-only record}
The complete V98 source was retained before this continuation.  It had
1432521 bytes, 42463 lines, and SHA--256
\[
 \texttt{F040D80DCDA7D417BB70E172E6F1C4F54188CD4D40F6EB4EB4AD76C2EFCFF7D0}.
\]
V99 exactly optimized the translated double-padding quartic, identified
the sharp algebraic height \(c_\star\), exhibited its triple-root heat
image, and proved that this symmetric quartic family still cannot cross
the four-ladder critical boundary.

% =============================================================================
% END APPENDED V99 MATERIAL
% =============================================================================

% =============================================================================
% BEGIN APPENDED V100 MATERIAL
% =============================================================================

\section{V100: mandatory corpus sweep and the sharp unrestricted quartic}
\label{sec:grh-v100}

V99 found the exact optimum \(c_\star=3.814511867\ldots\) when the
two real padding zeros coincide.  It left open whether unequal real
shifts could exceed that value, and in particular whether they could
cross the four-ladder target \(t/K=4\).

This version first performs the scheduled stable corpus sweep.  The
only transferable new method is a finite semialgebraic atlas which
separates regular strata, singular crossings, and boundary blow-ups.
Applying that discipline at the first backward-heat collision avoids
the unwieldy full quartic discriminant.  The input real-root
discriminant becomes a one-variable polynomial \(G(w,t)\).  At
\(t=c_\star\) it has an exact positive double-factor decomposition,
and it increases strictly with \(t\).  Consequently \(c_\star\) is
the sharp height for the unrestricted quartic as well: asymmetric
two-real-factor padding gives no improvement.

\subsection{Hash-gated V90--V100 companion audit}

The V90 snapshot contained \(31\) non-GRH \texttt{.tex} files and had
SHA--256
\[
 \texttt{F6BFD1D3E054643B8154F9E9758085110330023AF3D0473D6BB475CA6D1A439D}.
 \tag{V100.1}
\]
For V100, canonical manifest rows again have the form
\[
 \texttt{filename|byte-count|SHA--256}\backslash\texttt{n},
 \tag{V100.2}
\]
sorted by filename and joined with LF terminators.  Two complete
passes over all non-GRH notes had UTC bounds
\begin{center}
\small
pass 1:
\texttt{2026-09-11T11:05:24.961674Z}--%
\texttt{2026-09-11T11:05:24.991168Z},\\
pass 2:
\texttt{2026-09-11T11:05:25.091571Z}--%
\texttt{2026-09-11T11:05:25.119586Z}.
\end{center}
Both passes found \(32\) files and produced
\[
 \boxed{
 \texttt{300B05CE0D47278FC7234579B2861EA06AAF3245AAB213012B14AE04FDF014E6}.}
 \tag{V100.3}
\]
The live endpoints included
\begin{center}
\tiny
\begin{tabular}{p{0.58\textwidth}|r|p{0.28\textwidth}}
\toprule
file&bytes&SHA--256\\
\midrule
\path{Renewal_SM_Einstein_Continuum_Research_Notes_V27.tex}
&273056&
\texttt{5AC1CE6FCF09649E3FFBB7371ECA0DD9}\newline
\texttt{3E6966A19BDC7ED3CEB1CEEC761667DE}\\
\path{Renewal_Yang_Mills_Mass_Gap_Research_Notes_V78.tex}
&639688&
\texttt{BD4FDF9EA3D739EA1E5FDCE20FC7F932}\newline
\texttt{9E34AAFCE6E61363735AFA809676EDEF}\\
\path{Renewal_YM_Parallel_Research_Notes_V5.tex}
&54174&
\texttt{306F1BB84CFA8A8D47EA569EC17E9545}\newline
\texttt{F9867C8D32B548EC9D9C444B4F3C0512}\\
\bottomrule
\end{tabular}
\end{center}
The companion writers advanced during the preliminary read.  Their
successors were reread through SM V27 and Yang--Mills V78 before the
stable manifest \textup{(V100.3)} was frozen.

\subsection{Typed content audit}

The post-V90 SM material contains two relevant chains.  Its
sixteenth-cycle endpoint consolidates conditional finite-cutoff
positivity, mixing, compactness, renormalization, reflection
positivity, and dense-channel spectral-gap gates, but leaves the
programme-specific hypotheses unpopulated.  The current seventeenth
cycle studies polynomial shell stability.  It proves finite atlases
for compact angular zero sets, quadratic distance floors on regular
sum-of-squares strata, ideal-power vanishing tests, and separate
blow-up treatment of singular strata.  Its later V25--V27 results
concern exact conformal chord moments and finite-color graph blocking.

The Yang--Mills successor constructs further exact rational charts for
constrained Wilson-response pencils.  Through V78 it proves both signs
on paired first quarter faces and paired
\(\pi/4\times2t\times t\) slabs, while expressly leaving full-torus
positivity, cofinal control, the continuum measure, reflection
positivity, clustering, and the Hamiltonian gap open.

The parallel Yang--Mills V5 file has two completed decisions.  Its
first programme partitions connected owner words exactly and proves
that nonlinear remainder terms preserve a supplied symbolic margin on
an explicit weak-coupling domain.  Its second programme proves that a
volume-uniform tail for an extensive global integer charge is
incompatible with tensorization unless nonzero sectors have zero
probability.  Neither statement contains an arithmetic Euler
coefficient or an \(L\)-function zero estimate.  The Navier--Stokes
endpoint is unchanged.

\begin{theorem}[V100 companion-corpus verdict]
\label{thm:v100-corpus-verdict}
No post-V90 companion theorem supplies target-zero separation, an
Euler-admissible GRH filter, or a positivity theorem for the signed
Dirichlet coefficient transform.  No physical numerical constant is
importable.

One proof method transfers: on a compact semialgebraic parameter set,
separate the interior regular stratum, the constraint boundary, and
higher-order collision strata; certify each with its own exact
identity.  For the quartic heat block this means using the first
multiple-zero equations together with the real-padding discriminant,
instead of maximizing the full quartic discriminant numerically.
\end{theorem}

\begin{proof}
The SM theorems assume positive Gibbs laws, coercive polynomial
shells, conditional moment windows, or programme-specific geometric
actions.  The Yang--Mills theorems concern finite Hermitian Wilson
pencils, rational phase charts, or tensorized topological sectors.
None has Dirichlet characters, gamma ladders, or the target/principal
zero spectra as objects.  Their constants therefore fail the typed
input test for the GRH transform.

The finite-atlas and stratum-splitting arguments use only compactness,
polynomial identities, and separately verified boundary charts.
Those hypotheses can be checked directly for the quartic parameter
cone below, so the stated method transfers without transferring a
physical conclusion.
\end{proof}

\subsection{The first-collision normal form}

Let
\[
 P_0(z)=\bigl((z-x)^2+y^2\bigr)(z-a)(z-b),
 \qquad a,b,x\in\mathbb R,\quad y>0,
 \tag{V100.4}
\]
and
\[
 P_\tau(z)=
 \exp\!\left(-\frac{\tau}{2}D_z^2\right)P_0(z).
 \tag{V100.5}
\]
Suppose the conjugate pair first reaches the real axis at time
\(T>0\).  Translate its collision point to zero and scale
\(z=\sqrt T\,\zeta\).  The normalized input then has the form
\[
 \Pi(\zeta)
 =
 \bigl((\zeta-c)^2+t\bigr)
 \bigl(\zeta^2+\alpha\zeta+\beta\bigr),
 \qquad
 \alpha^2-4\beta\ge0,
 \tag{V100.6}
\]
where
\[
 t=\frac{y^2}{T}.
 \tag{V100.7}
\]
Its unit-time image
\[
 Q(\zeta):=
 \exp\!\left(-\frac12D_\zeta^2\right)\Pi(\zeta)
 \tag{V100.8}
\]
has a multiple zero at the origin:
\[
 Q(0)=Q'(0)=0.
 \tag{V100.9}
\]

Put
\[
 w=c^2,\qquad q=c^2+t=w+t.
 \tag{V100.10}
\]
Because \(\Pi\) is monic quartic,
\[
 Q=\Pi-\frac12\Pi''+3.
 \tag{V100.11}
\]
The coefficient expansion of \(\Pi\) gives the two collision
equations
\[
 \begin{split}
 2c\alpha+(q-1)\beta&=q-3,\\
 (q-3)\alpha-2c\beta&=-6c.
 \end{split}
 \tag{V100.12}
\]

\begin{lemma}[Collision discriminant identity]
\label{lem:v100-collision-discriminant}
Assume \(t>3\), and define
\[
 D:=4w+(q-1)(q-3)>0.
 \tag{V100.13}
\]
Then \textup{(V100.12)} forces
\[
 \alpha=-\frac{4cq}{D},
 \qquad
 \beta=\frac{12w+(q-3)^2}{D}.
 \tag{V100.14}
\]
Moreover,
\[
 \boxed{
 \alpha^2-4\beta=-\frac{4G(w,t)}{D^2},}
 \tag{V100.15}
\]
where
\begin{align}
 G(w,t):={}&w^4+(4t+2)w^3+6(t^2-t+2)w^2 \notag\\
 &+(4t^3-18t^2+18)w+(t-3)^3(t-1).
 \tag{V100.16}
\end{align}
\end{lemma}

\begin{proof}
For \(t>3\), one has \(q>3\), so both terms in
\textup{(V100.13)} are nonnegative and the second is positive.
The determinant of the linear system \textup{(V100.12)} is
\(-D\).  Cramer's rule gives \textup{(V100.14)}.

Consequently
\[
 \alpha^2-4\beta
 =
 \frac4{D^2}
 \left\{
 4wq^2-\bigl(12w+(q-3)^2\bigr)D
 \right\}.
 \tag{V100.17}
\]
Substitute \(q=w+t\) and expand.  The expression in braces is
\(-G(w,t)\), which proves \textup{(V100.15)}.
\end{proof}

\subsection{The exact positive factor at the V99 constant}

Retain the V99 definitions.  Let \(x_\star>0\) be the unique root of
\[
 x^3+3x^2-8=0,
 \tag{V100.18}
\]
and put
\[
 c_\star=
 \frac{3x_\star^2+12x_\star+16}
      {x_\star^2+3x_\star+4},
 \qquad
 w_\star=
 \frac{x_\star}{x_\star^2+3x_\star+4}.
 \tag{V100.19}
\]
Thus \(c_\star\) is the unique real root of
\[
 4c^3-21c^2+24c-8=0
 \tag{V100.20}
\]
and \(3<c_\star<4\).

\begin{lemma}[Sharp collision factorization]
\label{lem:v100-sharp-factor}
Define
\[
 A_\star:=4c_\star+2+2w_\star,
 \tag{V100.21}
\]
\[
 B_\star:=
 6(c_\star^2-c_\star+2)
 +w_\star(8c_\star+4+3w_\star).
 \tag{V100.22}
\]
Then \(A_\star>0\), \(B_\star>0\), and
\[
 \boxed{
 G(w,c_\star)
 =
 (w-w_\star)^2
 (w^2+A_\star w+B_\star).}
 \tag{V100.23}
\]
Furthermore, for \(t>3\) and \(w\ge0\),
\[
 \boxed{
 \partial_tG(w,t)
 =
 4w^3+(12t-6)w^2+12t(t-3)w
 +2(t-3)^2(2t-3)>0.}
 \tag{V100.24}
\]
\end{lemma}

\begin{proof}
Substitute \textup{(V100.19)} into \(G\) and its \(w\)-derivative,
clear the positive denominator
\((x_\star^2+3x_\star+4)^4\), and reduce powers of \(x_\star\)
using \textup{(V100.18)}.  Both remainders are zero:
\[
 G(w_\star,c_\star)=0,
 \qquad
 \partial_wG(w_\star,c_\star)=0.
 \tag{V100.25}
\]
Since \(G\) is monic of degree four in \(w\), division by
\((w-w_\star)^2\) gives a monic quadratic.  Matching the
\(w^3\) and \(w^2\) coefficients gives exactly
\textup{(V100.21)}--\textup{(V100.22)}, proving
\textup{(V100.23)}.  Every term in \textup{(V100.22)} is positive,
so the quadratic factor is positive for \(w\ge0\).

Differentiating \textup{(V100.16)} gives
\textup{(V100.24)}.  All four displayed terms are nonnegative for
\(t>3,w\ge0\), and the last is strictly positive.
\end{proof}

\subsection{Sharp unrestricted two-padding theorem}

\begin{theorem}[Sharp quartic one-pair heat-time obstruction]
\label{thm:v100-sharp-quartic}
Let \(P_0\) be the quartic in \textup{(V100.4)}.  If
\[
 \exp\!\left(-\frac K2D_z^2\right)P_0(z)
 \tag{V100.26}
\]
is real-rooted for some \(K>0\), then necessarily
\[
 \boxed{y^2\le c_\star K.}
 \tag{V100.27}
\]
The constant is sharp.  Equality is possible only, up to real
translation, reflection, and heat scaling, when the two real input
zeros coincide and the unit-time image has the V99 triple collision.
Thus unequal real shifts never improve the V99 constant.
\end{theorem}

\begin{proof}
By Lemma~\ref{lem:v98-zero-flow}, the two real zeros stay real until
the conjugate pair reaches the axis.  If the image at time \(K\) is
real-rooted, there is a first collision time \(0<T\le K\).  Translate
and scale as in \textup{(V100.6)}.  If \(t=y^2/T\le3\), then
\(t<c_\star\) and there is nothing to prove.

Suppose \(t>3\).  The two real input zeros give
\[
 \alpha^2-4\beta\ge0.
 \tag{V100.28}
\]
Lemma~\ref{lem:v100-collision-discriminant} therefore implies
\[
 G(w,t)\le0.
 \tag{V100.29}
\]
But Lemma~\ref{lem:v100-sharp-factor} gives
\[
 G(w,c_\star)\ge0
 \qquad(w\ge0),
 \tag{V100.30}
\]
with equality only at \(w=w_\star\), and
\(\partial_tG>0\) for \(t>3\).  Hence \textup{(V100.29)} is
impossible when \(t>c_\star\).  Thus
\[
 \frac{y^2}{T}=t\le c_\star.
 \]
Since \(T\le K\), this proves \textup{(V100.27)}.

The V99 factorization \textup{(V99.18)} attains equality, proving
sharpness.  If equality holds in \textup{(V100.27)}, then
\(T=K\), \(t=c_\star\), and \textup{(V100.30)} forces
\(w=w_\star\).  Equation \textup{(V100.15)} then gives
\(\alpha^2-4\beta=0\), so the two real input zeros coincide.
The collision is precisely the V99 triple-root configuration.
\end{proof}

\begin{corollary}[Exact unrestricted quartic optimum]
\label{cor:v100-quartic-optimum}
Among all monic quartics having one conjugate pair and two real
zeros, the supremum of the squared conjugate height removable by
unit backward heat is
\[
 \boxed{c_\star=3.814511867\ldots.}
 \tag{V100.31}
\]
This supremum is attained only on the repeated-real-zero boundary;
with two distinct real padding zeros it is approached but not
attained.
\end{corollary}

\begin{proof}
The upper bound is Theorem~\ref{thm:v100-sharp-quartic}.  Attainment
and perturbative approximation by distinct real zeros follow from
Theorem~\ref{thm:v99-equal-quartic} and
Proposition~\ref{prop:v98-quartic}.
\end{proof}

\subsection{Critical gamma consequence}

For a fixed block containing one target conjugate pair and exactly
two real gamma padding factors, the critical parameters remain
\[
 K=\frac{1+\theta}{\theta^2},
 \qquad
 t=\frac{\gamma^2R}{\theta^2\kappa^2},
 \qquad R\ge4.
 \tag{V100.32}
\]

\begin{theorem}[Two-real-padding architecture is closed]
\label{thm:v100-two-padding-closed}
Every such quartic gamma block must satisfy
\[
 \boxed{
 \kappa^2\ge
 \frac{\gamma^2R}{c_\star(1+\theta)}.}
 \tag{V100.33}
\]
Because \(R\ge4>c_\star\), it follows necessarily that
\[
 \boxed{
 \kappa>
 \frac{|\gamma|}{\sqrt{1+\theta}}.}
 \tag{V100.34}
\]
Thus neither equal nor unequal placement of two real gamma factors
can enlarge the V97 critical region.
\end{theorem}

\begin{proof}
Insert \textup{(V100.32)} into the sharp heat-time bound
\(t\le c_\star K\).  This gives \textup{(V100.33)}.  The strict
inequality \textup{(V100.34)} follows from \(R/c_\star\ge4/c_\star>1\).
\end{proof}

\begin{firewall}[Quartic closure is not finite-padding closure]
Theorem~\ref{thm:v100-two-padding-closed} settles every quartic
one-pair block, not blocks with three or more real padding factors.
It also does not treat factor groups growing with \(M\), several
target-pair ladder indices in one group, local variation recovery
without global real-rootedness, or nonuniform pole meshes.  The
universal V98 speed barrier and the remaining GRH target-separation
problem are unchanged.
\end{firewall}

\begin{theorem}[V100 heat-block verdict]
\label{thm:v100-verdict}
\begin{enumerate}[label=\textup{(\roman*)},leftmargin=3em]
\item the scheduled companion sweep is complete on the stable
      manifest \textup{(V100.3)};
\item no companion theorem supplies an arithmetic GRH hypothesis;
\item the first-collision equations reduce unrestricted quartic
      feasibility to the real-padding discriminant
      \textup{(V100.15)};
\item the exact positive factorization \textup{(V100.23)} proves that
      \(c_\star\) is sharp for arbitrary, not merely equal, real shifts;
\item equality forces the two real input zeros to coincide and the
      heat image to have the V99 triple collision;
\item every two-real-padding gamma architecture is therefore closed
      at the already known V97 boundary.
\end{enumerate}
Target-zero separation and GRH remain open.
\end{theorem}

\begin{assessment}[V100 closure and rollover direction]
V100 resolves the first nontrivial separated-padding optimization
completely.  The full quartic discriminant was unnecessary: moving
upstream to the first collision exposes a \(2\times2\) linear system,
and the only input constraint needed is the discriminant of the two
real padding roots.  This converts the numerical observation that
asymmetry does not help into the sharp theorem
\textup{(V100.27)}.

The next cycle should apply the same collision-first philosophy to
three real padding factors.  The decisive question is whether the
sharp removable height for \(p=3\) exceeds the minimal five-ladder
count \(R=5\).  A negative answer would close the next finite block;
a positive construction would enter the still-open band between the
V97 sufficient boundary and the V98 universal speed ceiling.
\end{assessment}

\section*{Version V100 append-only record}
\addcontentsline{toc}{section}{Version V100 append-only record}
The complete V99 source was retained before this continuation.  It
had \(1443866\) bytes, \(42807\) lines, and SHA--256
\[
 \texttt{99C8B0A7AF7CFE8C70678292F80A942581ADCB27FA367FDC744D018EB472AE16}.
\]
V100 performed the stable ten-version companion sweep and proved the
sharp unrestricted quartic heat-time obstruction, closing every
two-real-padding critical gamma block at the V97 boundary.

% =============================================================================
% END APPENDED V100 MATERIAL
% =============================================================================
\end{document}
